% Auto-generated from data/segments.json + layout.json (+ transforms) — do not edit by hand.
% volume: 17An03
% mode: reading
% segments: 3662
% transform_rules: 0
% toc_marks: 556

% Volume layout defaults (layout.json ``layout``).
\csromanlayoutapply{1}{1.5}{21.6pt}{5pt}{6.3pt}{65pt}{2.5em}%

% Reading physical-page layout (layout.json ``page_layout_reading_mode``).
\csromanreadingpagelayoutvolume{1}{1.5}{21.6pt}{5pt}{6.3pt}{65pt}{2.5em}%
\csromanreadingpagelayoutenable

\csromanheader{2}{\textbf{องฺคุตฺตรนิกาย}}
\gambhira{\textbf{อฏฺฐกนิปาตปาฬิ}}
\namakkaram{นโม ตสฺส ภควโต อรหโต สมฺมาสมฺพุทฺธสฺส.}
\csromanmatikaanchor{mtk.0}
\csromanmatikaanchor{mtk.21}
\csromantocmarknopage{chapter}{นวกนิปาตปาฬิ}{mtk.0}
\bookmark[dest={mtk.0},level=0]{นวกนิปาตปาฬิ}
\chapterhead{1. ปฐมปณฺณาสก}
\csromanheader{2}{1. เมตฺตาสุตฺต}
\proseitem{1}{\csromanfolio{1}เอวํ เม สุตํ\csromandash{}เอกํ สมยํ ภควา สาวตฺถิยํ วิหรติ เชตวเน อนาถปิณฺฑิกสฺส อาราเม.\csromanspacer{} ตตฺร โข ภควา ภิกฺขู อามนฺเตสิ “ภิกฺขโว”ติ. “ภทนฺเต”ติ เต ภิกฺขู ภควโต ปจฺจสฺโสสุํ.\csromanspacer{} ภควา เอตทโวจ\csromandash{}}

\prose{\csromansymbolfootnote{*}{อุปริ 542 ปิฏฺเฐปิ.}เมตฺตาย ภิกฺขเว เจโตวิมุตฺติยา อาเสวิตาย ภาวิตาย พหุลีกตาย ยานีกตาย วตฺถุกตาย อนุฏฺฐิตาย ปริจิตาย สุสมารทฺธาย อฏฺฐานิสํสา ปาฏิกงฺขา.\csromanspacer{} กตเม อฏฺฐ? สุขํ สุปติ, สุขํ ปฏิพุชฺฌติ, น ปาปกํ สุปินํ ปสฺสติ, มนุสฺสานํ ปิโย โหติ, อมนุสฺสานํ ปิโย โหติ, เทวตา รกฺขนฺติ, นาสฺส อคฺคิ วา วิสํ วา สตฺถํ วา กมติ, อุตฺตริํ อปฺปฏิวิชฺฌนฺโต พฺรหฺมโลกูปโค โหติ.\csromanspacer{} เมตฺตาย ภิกฺขเว เจโตวิมุตฺติยา อาเสวิตาย ภาวิตาย พหุลีกตาย ยานีกตาย วตฺถุกตาย อนุฏฺฐิตาย ปริจิตาย สุสมารทฺธาย อิเม อฏฺฐานิสํสา ปาฏิกงฺขาติ.}

\csromangathameasure{โย จ เมตฺตํ ภาวยติ, อปฺปมาณํ ปฏิสฺสโต. \\ ตนู สํโยชนา โหนฺติ, ปสฺสโต อุปธิกฺขยํ. \\ เอกมฺปิ เจ ปาณมทุฏฺฐจิตฺโต, \\ เมตฺตายติ กุสลี เตน โหติ. \\ สพฺเพ จ ปาเณ มนสานุกมฺปี, \\ ปหูตมริโย ปกโรติ ปุญฺญํ. \\ เย สตฺตสณฺฑํ ปถวิํ วิเชตฺวา, \\ ราชิสโย ยชมานา อนุปริยคา. \\ อสฺสเมธํ ปุริสเมธํ, \\ สมฺมาปาสํ วาชเปยฺยํ นิรคฺคฬํ. \\ เมตฺตสฺส จิตฺตสฺส สุภาวิตสฺส, \\ กลมฺปิ เต นานุภวนฺติ โสฬสิํ. \\ จนฺทปฺปภา ตารคณาว สพฺเพ, \\ ยถา น อคฺฆนฺติ กลมฺปิ โสฬสิํ. \\ โย น หนฺติ น ฆาเตติ, น ชินาติ น ชาปเย. \\ เมตฺตํโส สพฺพภูตานํ, เวรํ ตสฺส น เกนจีติ.ปฐมํ.}
\csromangathagroup{\csromangathastanza{\csromanfolio{2}โย จ เมตฺตํ ภาวยติ, อปฺปมาณํ ปฏิสฺสโต\footnote{ปติสฺสโต (สี)}. \\ ตนู สํโยชนา โหนฺติ, ปสฺสโต อุปธิกฺขยํ.}\csromangathastanza{เอกมฺปิ เจ ปาณมทุฏฺฐจิตฺโต, \\ เมตฺตายติ กุสลี เตน โหติ. \\ สพฺเพ จ ปาเณ มนสานุกมฺปี, \\ ปหูตมริโย ปกโรติ ปุญฺญํ.}\csromangathastanza{เย สตฺตสณฺฑํ ปถวิํ วิเชตฺวา, \\ ราชิสโย ยชมานา อนุปริยคา. \\ อสฺสเมธํ ปุริสเมธํ, \\ สมฺมาปาสํ วาชเปยฺยํ นิรคฺคฬํ.}\csromangathastanza{เมตฺตสฺส จิตฺตสฺส สุภาวิตสฺส, \\ กลมฺปิ เต นานุภวนฺติ โสฬสิํ. \\ จนฺทปฺปภา ตารคณาว สพฺเพ, \\ ยถา น อคฺฆนฺติ กลมฺปิ โสฬสิํ\footnote{อยํ ปาโท พหูสุ น ทิสฺสติ.}. \\ โย น หนฺติ น ฆาเตติ, น ชินาติ น ชาปเย. \\ เมตฺตํโส สพฺพภูตานํ, เวรํ ตสฺส น เกนจีติ.\csromanspacer{}\csromanspacer{}ปฐมํ.}}

\csromanheader{2}{2. ปญฺญาสุตฺต}
\proseitem{2}{อฏฺฐิเม ภิกฺขเว เหตู อฏฺฐ ปจฺจยา อาทิพฺรหฺมจริยิกาย ปญฺญาย อปฺปฏิลทฺธาย ปฏิลาภาย, ปฏิลทฺธาย ภิยฺโยภาวาย เวปุลฺลาย ภาวนาย ปาริปูริยา สํวตฺตนฺติ.\csromanspacer{} กตเม อฏฺฐ? อิธ ภิกฺขเว ภิกฺขุ สตฺถารํ อุปนิสฺสาย วิหรติ อญฺญตรํ วา ครุฏฺฐานิยํ สพฺรหฺมจาริํ, ยตฺถสฺส ติพฺพํ หิโรตฺตปฺปํ ปจฺจุปฏฺฐิตํ โหติ เปมญฺจ คารโว จ.\csromanspacer{} อยํ ภิกฺขเว ปฐโม เหตุ ปฐโม ปจฺจโย อาทิพฺรหฺมจริยิกาย ปญฺญาย อปฺปฏิลทฺธาย ปฏิลาภาย ปฏิลทฺธาย ภิยฺโยภาวาย เวปุลฺลาย}

\prose{โส ตํ สตฺถารํ อุปนิสฺสาย วิหรนฺโต อญฺญตรํ วา ครุฏฺฐานิยํ สพฺรหฺมจาริํ, ยตฺถสฺส ติพฺพํ หิโรตฺตปฺปํ ปจฺจุปฏฺฐิตํ โหติ เปมญฺจ \csromanfolio{3}คารโว จ, เต กาเลน กาลํ อุปสงฺกมิตฺวา ปริปุจฺฉติ ปริปญฺหติ “อิทํ ภนฺเต กถํ, อิมสฺส โก อตฺโถ”ติ.\csromanspacer{} ตสฺส เต อายสฺมนฺโต อวิวฏญฺเจว วิวรนฺติ, อนุตฺตานีกตญฺจ อุตฺตานี กโรนฺติ, อเนกวิหิเตสุ จ กงฺขาฐานิเยสุ ธมฺเมสุ กงฺขํ ปฏิวิโนเทนฺติ.\csromanspacer{} อยํ ภิกฺขเว ทุติโย เหตุ ทุติโย ปจฺจโย อาทิพฺรหฺมจริยิกาย ปญฺญาย อปฺปฏิลทฺธาย ปฏิลาภาย ปฏิลทฺธาย ภิยฺโยภาวาย เวปุลฺลาย ภาวนาย ปาริปูริยา สํวตฺตติ.}

\prose{โส ตํ ธมฺมํ สุตฺวา ทฺวเยน วูปกาเสน สมฺปาเทติ กายวูปกาเสน จ จิตฺตวูปกาเสน จ.\csromanspacer{} อยํ ภิกฺขเว ตติโย เหตุ ตติโย ปจฺจโย อาทิพฺรหฺมจริยิกาย ปญฺญาย อปฺปฏิลทฺธาย ปฏิลาภาย ปฏิลทฺธาย ภิยฺโยภาวาย เวปุลฺลาย ภาวนาย ปาริปูริยา สํวตฺตติ.}

\prose{สีลวา โหติ, ปาติโมกฺขสํวรสํวุโต วิหรติ อาจารโคจรสมฺปนฺโน, อณุมตฺเตสุ วชฺเชสุ ภยทสฺสาวี สมาทาย สิกฺขติ สิกฺขาปเทสุ.\csromanspacer{} อยํ ภิกฺขเว จตุตฺโถ เหตุ จตุตฺโถ ปจฺจโย อาทิพฺรหฺมจริยิกาย ปญฺญาย อปฺปฏิลทฺธาย ปฏิลาภาย ปฏิลทฺธาย ภิยฺโยภาวาย เวปุลฺลาย}

\prose{พหุสฺสุโต โหติ สุตธโร สุตสนฺนิจโย, เย เต ธมฺมา อาทิกลฺยาณา มชฺเฌกลฺยาณา ปริโยสานกลฺยาณา สาตฺถํ สพฺยญฺชนํ\footnote{สตฺถา สพฺยญฺชนา (ก, สี)} เกวลปริปุณฺณํ ปริสุทฺธํ พฺรหฺมจริยํ อภิวทนฺติ, ตตฺถารูปาสฺส ธมฺมา พหุสฺสุตา โหนฺติ ธาตา\footnote{ธตา (สี, สฺยา, กํ, อิ)} วจสา ปริจิตา มนสานุเปกฺขิตา ทิฏฺฐิยา สุปฺปฏิวิทฺธา.\csromanspacer{} อยํ ภิกฺขเว ปญฺจโม เหตุ ปญฺจโม ปจฺจโย อาทิพฺรหฺมจริยิกาย ปญฺญาย อปฺปฏิลทฺธาย ปฏิลาภาย ปฏิลทฺธาย ภิยฺโยภาวาย เวปุลฺลาย}

\prose{อารทฺธวีริโย วิหรติ อกุสลานํ ธมฺมานํ ปหานาย กุสลานํ ธมฺมานํ อุปสมฺปทาย, ถามวา ทฬฺหปรกฺกโม อนิกฺขิตฺตธุโร กุสเลสุ ธมฺเมสุ.\csromanspacer{} อยํ ภิกฺขเว ฉฏฺโฐ เหตุ ฉฏฺโฐ ปจฺจโย อาทิพฺรหฺมจริยิกาย ปญฺญาย อปฺปฏิลทฺธาย ปฏิลาภาย, ปฏิลทฺธาย ภิยฺโยภาวาย เวปุลฺลาย ภาวนาย ปาริปูริยา สํวตฺตติ.}

\prose{\csromanfolio{4}สํฆคโต โข ปน อนานากถิโก โหติ อติรจฺฉานกถิโก, สามํ วา ธมฺมํ ภาสติ, ปรํ วา อชฺเชสติ, อริยํ วา ตุณฺหีภาวํ นาติมญฺญติ.\csromanspacer{} อยํ ภิกฺขเว สตฺตโม เหตุ สตฺตโม ปจฺจโย อาทิพฺรหฺมจริยิกาย ปญฺญาย อปฺปฏิลทฺธาย ปฏิลาภาย ปฏลทฺธาย ภิยฺโยภาวาย เวปุลฺลาย}

\prose{ปญฺจสุ โข ปน อุปาทานกฺขนฺเธสุ อุทยพฺพยานุปสฺสี วิหรติ “อิติ รูปํ อิติ รูปสฺส สมุทโย อิติ รูปสฺส อตฺถงฺคโม, อิติ เวทนา อิติ เวทนาย สมุทโย อิติ เวทนาย อตฺถงฺคโม, อิติ สญฺญา ฯเปฯ อิติ สงฺขารา.\csromanspacer{} อิติ วิญฺญาณํ อิติ วิญฺญาณสฺส สมุทโย อิติ วิญฺญาณสฺส อตฺถงฺคโม”ติ.\csromanspacer{} อยํ ภิกฺขเว อฏฺฐโม เหตุ อฏฺฐโม ปจฺจโย อาทิพฺรหฺมจริยิกาย ปญฺญาย อปฺปฏิลทฺธาย ปฏิลาภาย ปฏิลทฺธาย ภิยฺโยภาวาย เวปุลฺลาย ภาวนาย ปาริปูริยา สํวตฺตติ.}

\prose{ตเมนํ สพฺรหฺมจารี เอวํ สมฺภาเวนฺติ “อยํ โข อายสฺมา สตฺถารํ อุปนิสฺสาย วิหรติ, อญฺญตรํ วา ครุฏฺฐานิยํ สพฺรหฺมจาริํ, ยตฺถสฺส ติพฺพํ หิโรตฺตปฺปํ ปจฺจุปฏฺฐิตํ โหติ เปมญฺจ คารโว จ, อทฺธา อยมายสฺมา ชานํ ชานาติ, ปสฺสํ ปสฺสตี”ติ.\csromanspacer{} อยมฺปิ ธมฺโม ปิยตฺตาย ครุตฺตาย\footnote{ปิยตาย ครุตาย (สฺยา)} ภาวนาย สามญฺญาย เอกีภาวาย สํวตฺตติ. (1)}

\prose{ตํ โข ปนายมายสฺมา สตฺถารํ อุปนิสฺสาย วิหรนฺโต อญฺญตรํ วา ครุฏฺฐานิยํ สพฺรหฺมจาริํ, ยตฺถสฺส ติพฺพํ หิโรตฺตปฺปํ ปจฺจุปฏฺฐิตํ โหติ เปมญฺจ คารโว จ, เต กาเลน กาลํ อุปสงฺกมิตฺวา ปริปุจฺฉติ ปริปญฺหติ “อิทํ ภนฺเต กถํ, อิมสฺส โก อตฺโถ”ติ.\csromanspacer{} ตสฺส เต อายสฺมนฺโต อวิวฏญฺเจว วิวรนฺติ, อนุตฺตานีกตญฺจ อุตฺตานี กโรนฺติ, อเนกวิหิเตสุ จ กงฺขาฐานิเยสุ ธมฺเมสุ กงฺขํ ปฏิวิโนเทนฺติ.\csromanspacer{} อทฺธา อยมายสฺมา ชานํ ชานาติ, ปสฺสํ ปสฺสตีติ.\csromanspacer{} อยมฺปิ ธมฺโม ปิยตฺตาย ครุตฺตาย ภาวนาย สามญฺญาย เอกีภาวาย สํวตฺตติ. (2)}

\prose{ตํ โข ปนายมายสฺมา ธมฺมํ สุตฺวา ทฺวเยน วูปกาเสน สมฺปาเทติ กายวูปกาเสน จ จิตฺตวูปกาเสน จ.\csromanspacer{} อทฺธา อยมายสฺมา ชานํ ชานาติ, ปสฺสํ ปสฺสตีติ.\csromanspacer{} อยมฺปิ ธมฺโม ปิยตฺตาย ครุตฺตาย ภาวนาย สามญฺญาย เอกีภาวาย สํวตฺตติ. (3)}

\prose{\csromanfolio{5}สีลวา โข ปนายมายสฺมา ปาติโมกฺขสํวรสํวุโต วิหรติ อาจารโคจรสมฺปนฺโน, อณุมตฺเตสุ วชฺเชสุ ภยทสฺสาวี สมาทาย สิกฺขติ สิกฺขาปเทสุ.\csromanspacer{} อทฺธา อยมายสฺมา ชานํ ชานาติ, ปสฺสํ ปสฺสตีติ.\csromanspacer{} อยมฺปิ ธมฺโม ปิยตฺตาย ครุตฺตาย ภาวนาย สามญฺญาย เอกีภาวาย สํวตฺตติ. (4)}

\prose{พหุสฺสุโต โข ปนายมายสฺมา สุตธโร สุตสนฺนิจโย, เย เต ธมฺมา อาทิกลฺยาณา มชฺเฌกลฺยาณา ปริโยสานกลฺยาณา, สาตฺถํ สพฺยญฺชนํ เกวลปริปุณฺณํ ปริสุทฺธํ พฺรหฺมจริยํ อภิวทนฺติ, ตถารูปาสฺส ธมฺมา พหุสฺสุตา โหนฺติ ธาตา วจสา ปริจิตา มนสานุเปกฺขิตา ทิฏฺฐิยา สุปฺปฏิวิทฺธา.\csromanspacer{} อทฺธา อยมายสฺมา ชานํ ชานาติ, ปสฺสํ ปสฺสตีติ.\csromanspacer{} อยมฺปิ ธมฺโม ปิยตฺตาย ครุตฺตาย ภาวนาย สามญฺญาย เอกีภาวาย สํวตฺตติ. (5)}

\prose{อารทฺธวีริโย โข ปนายมายสฺมา วิหรติ อกุสลานํ ธมฺมานํ ปหานาย กุสลานํ ธมฺมานํ อุปสมฺปทาย ถามวา ทฬฺหปรกฺกโม อนิกฺขิตฺตธุโร กุสเลสุ ธมฺเมสุ.\csromanspacer{} อทฺธา อยมายสฺมา ชานํ ชานาติ, ปสฺสํ ปสฺสตีติ.\csromanspacer{} อยมฺปิ ธมฺโม ปิยตฺตาย ครุตฺตาย ภาวนาย สามญฺญาย เอกีภาวาย สํวตฺตติ. (6)}

\prose{สํฆคโต โข ปนายมายสฺมา อนานากถิโก โหติ อติรจฺฉานกถิโก, สามํ วา ธมฺมํ ภาสติ, ปรํ วา อชฺเฌสติ, อริยํ วา ตุณฺหีภาวํ นาติมญฺญติ.\csromanspacer{} อทฺธา อยมายสฺมา ชานํ ชานาติ, ปสฺสํ ปสฺสตีติ.\csromanspacer{} อยมฺปิ ธมฺโม ปิยตฺตาย ครุตฺตาย ภาวนาย สามญฺญาย เอกีภาวาย สํวตฺตติ. (7)}

\prose{ปญฺจสุ โข ปนายมายสฺมา อุปาทานกฺขนฺเธสุ อุทยพฺพยานุปสฺสี วิหรติ “อิติ รูปํ อิติ รูปสฺส สมุทโย อิติ รูปสฺส อตฺถงฺคโม, อิติ เวทนา ฯเปฯ อิติ สญฺญา.\csromanspacer{} อิติ สงฺขารา.\csromanspacer{} อิติ วิญฺญาณํ อิติ วิญฺญาณสฺส สมุทโย อิติ วิญฺญาณสฺส อตฺถงฺคโม”ติ.\csromanspacer{} อทฺธา อยมายสฺมา ชานํ ชานาติ, ปสฺสํ ปสฺสตีติ.\csromanspacer{} อยมฺปิ ธมฺโม ปิยตฺตาย ครุตฺตาย ภาวนาย สามญฺญาย เอกีภาวาย สํวตฺตติ. (8)}

\prose{อิเม โข ภิกฺขเว อฏฺฐ เหตู อฏฺฐ ปจฺจยา อาทิพฺรหฺมจริยิกาย ปญฺญาย อปฺปฏิลทฺธาย ปฏิลาภาย, ปฏิลทฺธาย ภิยฺโยภาวาย เวปุลฺลาย ภาวนาย ปาริปูริยา สํวตฺตนฺตีติ.\csromanspacer{}\csromanspacer{}ทุติยํ.}

\csromanheader{2}{3. ปฐม-อปฺปิยสุตฺต}
\proseitem{3}{\csromanfolio{6}อฏฺฐหิ ภิกฺขเว ธมฺเมหิ สมนฺนาคโต ภิกฺขุ สพฺรหฺมจารีนํ อปฺปิโย จ โหติ อมนาโป จ อครุ จ อภาวนีโย จ.\csromanspacer{} กตเมหิ อฏฺฐหิ? อิธ ภิกฺขเว ภิกฺขุ อปฺปิยปสํสี จ โหติ ปิยครหี จ ลาภกาโม จ สกฺการกาโม จ อหิริโก จ อโนตฺตปฺปี จ ปาปิจฺโฉ จ มิจฺฉาทิฏฺฐิ จ.\csromanspacer{} อิเมหิ โข ภิกฺขเว อฏฺฐหิ ธมฺเมหิ สมนฺนาคโต ภิกฺขุ สพฺรหฺมจารีนํ อปฺปิโย จ โหติ อมนาโป จ อครุ จ อภาวนีโย จ.}

\prose{อฏฺฐหิ ภิกฺขเว ธมฺเมหิ สมนฺนาคโต ภิกฺขุ สพฺรหฺมจารีนํ ปิโย จ โหติ มนาโป จ ครุ จ ภาวนีโย จ.\csromanspacer{} กตเมหิ อฏฺฐหิ? อิธ ภิกฺขเว ภิกฺขุ น อปฺปิยปสํสี จ โหติ น ปิยครหี จ น ลาภกาโม จ น สกฺการกาโม จ หิรีมา จ โหติ โอตฺตปฺปี จ อปฺปิจฺโฉ จ สมฺมาทิฏฺฐิ จ.\csromanspacer{} อิเมหิ โข ภิกฺขเว อฏฺฐหิ ธมฺเมหิ สมนฺนาคโต ภิกฺขุ สพฺรหฺมจารีนํ ปิโย จ โหติ มนาโป จ ครุ จ ภาวนีโย จาติ.\csromanspacer{}\csromanspacer{}ตติยํ.}

\csromanheader{2}{4. ทุติย-อปฺปิยสุตฺต}
\proseitem{4}{อฏฺฐหิ ภิกฺขเว ธมฺเมหิ สมนฺนาคโต ภิกฺขุ สพฺรหฺมจารีนํ อปฺปิโย จ โหติ อมนาโป จ อครุ จ อภาวนีโย จ.\csromanspacer{} กตเมหิ อฏฺฐหิ? อิธ ภิกฺขเว ภิกฺขุ ลาภกาโม จ โหติ สกฺการกาโม จ อนวญฺญตฺติกาโม จ อกาลญฺญู จ อมตฺตญฺญู จ อสุจิ จ พหุภาณี จ อกฺโกสกปริภาสโก จ สพฺรหฺมจารีนํ.\csromanspacer{} อิเมหิ โข ภิกฺขเว อฏฺฐหิ ธมฺเมหิ สมนฺนาคโต ภิกฺขุ สพฺรหฺมจารีนํ อปฺปิโย จ โหติ อมนาโป จ อครุ จ อภาวนีโย จ.}

\prose{อฏฺฐหิ ภิกฺขเว ธมฺเมหิ สมนฺนาคโต ภิกฺขุ สพฺรหฺมจารีนํ ปิโย จ โหติ มนาโป จ ครุ จ ภาวนีโย จ.\csromanspacer{} กตเมหิ อฏฺฐหิ? อิธ ภิกฺขเว ภิกฺขุ น ลาภกาโม จ โหติ น สกฺการกาโม จ น อนวญฺญตฺติกาโม จ กาลญฺญู จ มตฺตญฺญู จ สุจิ จ น พหุภาณี จ อนกฺโกสกปริภาสโก จ สพฺรหฺมจารีนํ.\csromanspacer{} อิเมหิ โข ภิกฺขเว อฏฺฐหิ ธมฺเมหิ สมนฺนาคโต ภิกฺขุ สพฺรหฺมจารีนํ ปิโย จ โหติ มนาโป จ ครุ จ ภาวนีโย จาติ.\csromanspacer{}\csromanspacer{}จตุตฺถํ.}

\csromanheader{2}{5. ปฐมโลกธมฺมสุตฺต}
\proseitem{5}{\csromanfolio{7}อฏฺฐิเม ภิกฺขเว โลกธมฺมา โลกํ อนุปริวตฺตนฺติ, โลโก จ อฏฺฐ โลกธมฺเม อนุปริวตฺตติ.\csromanspacer{} กตเม อฏฺฐ? ลาโภ จ อลาโภ จ ยโส จ อยโส จ นินฺทา จ ปสํสา จ สุขญฺจ ทุกฺขญฺจ.\csromanspacer{} อิเม โข ภิกฺขเว อฏฺฐ โลกธมฺมา โลกํ อนุปริวตฺตนฺติ, โลโก จ อิเม อฏฺฐ โลกธมฺเม อนุปริวตฺตตีติ.}

\csromangathameasure{ลาโภ อลาโภ จ ยสายโส จ, \\ นินฺทา ปสํสา จ สุขํ ทุขญฺจ. \\ เอเต อนิจฺจา มนุเชสุ ธมฺมา, \\ อสสฺสตา วิปริณามธมฺมา. \\ เอเต จ ญตฺวา สติมา สุเมโธ, \\ อเวกฺขติ วิปริณามธมฺเม. \\ อิฏฺฐสฺส ธมฺมา น มเถนฺติ จิตฺตํ, \\ อนิฏฺฐโต โน ปฏิฆาตเมติ. \\ ตสฺสานุโรธา อถ วา วิโรธา, \\ วิธูปิตา อตฺถงฺคตา น สนฺติ. \\ ปทญฺจ ญตฺวา วิรชํ อโสกํ, \\ สมฺมปฺปชานาติ ภวสฺส ปารคูติ.ปญฺจมํ.}
\csromangathagroup{\csromangathastanza{ลาโภ อลาโภ จ ยสายโส จ, \\ นินฺทา ปสํสา จ สุขํ ทุขญฺจ. \\ เอเต อนิจฺจา มนุเชสุ ธมฺมา, \\ อสสฺสตา วิปริณามธมฺมา.}\csromangathastanza{เอเต จ ญตฺวา สติมา สุเมโธ, \\ อเวกฺขติ วิปริณามธมฺเม. \\ อิฏฺฐสฺส ธมฺมา น มเถนฺติ จิตฺตํ, \\ อนิฏฺฐโต โน ปฏิฆาตเมติ.}\csromangathastanza{ตสฺสานุโรธา อถ วา วิโรธา, \\ วิธูปิตา อตฺถงฺคตา น สนฺติ. \\ ปทญฺจ ญตฺวา วิรชํ อโสกํ, \\ สมฺมปฺปชานาติ ภวสฺส ปารคูติ.\csromanspacer{}\csromanspacer{}ปญฺจมํ.}}

\csromanheader{2}{6. ทุติยโลกธมฺมสุตฺต}
\proseitem{6}{อฏฺฐิเม ภิกฺขเว โลกธมฺมา โลกํ อนุปริวตฺตนฺติ, โลโก จ อฏฺฐ โลกธมฺเม อนุปริวตฺตติ.\csromanspacer{} กตเม อฏฺฐ? ลาโภ จ อลาโภ จ ยโส จ อยโส จ นินฺทา จ ปสํสา จ สุขญฺจ ทุกฺขญฺจ.\csromanspacer{} อิเม โข ภิกฺขเว อฏฺฐ โลกธมฺมา โลกํ อนุปริวตฺตนฺติ, โลโก จ อิเม อฏฺฐ โลกธมฺเม อนุปริวตฺตติ.}

\prose{อสฺสุตวโต ภิกฺขเว ปุถุชฺชนสฺส อุปฺปชฺชติ ลาโภปิ อลาโภปิ ยโสปิ อยโสปิ นินฺทาปิ ปสํสาปิ สุขมฺปิ ทุกฺขมฺปิ.\csromanspacer{} สุตวโตปิ ภิกฺขเว อริยสาวกสฺส อุปฺปชฺชติ ลาโภปิ อลาโภปิ ยโสปิ อยโสปิ นินฺทาปิ ปสํสาปิ สุขมฺปิ ทุกฺขมฺปิ.\csromanspacer{} ตตฺร ภิกฺขเว โก วิเสโส \csromanfolio{8}โก อธิปฺปยาโส\footnote{อธิปฺปาโย, (สี) อธิปฺปายโส (สฺยา, กํ), อธิ + ป + ยสุ + ณ = อธิปฺปยาโส.} กิํ นานากรณํ สุตวโต อริยสาวกสฺส อสฺสุตวตา ปุถุชฺชเนนาติ.\csromanspacer{} ภควํมูลกา โน ภนฺเต ธมฺมา ภควํเนตฺติกา ภควํปฏิสรณา, สาธุ วต ภนฺเต ภควนฺตํเยว ปฏิภาตุ เอตสฺส ภาสิตสฺส อตฺโถ, ภควโต สุตฺวา ภิกฺขู ธาเรสฺสนฺตีติ.}

\prose{เตน หิ ภิกฺขเว สุณาถ สาธุกํ มนสิ กโรถ, ภาสิสฺสามีติ. “เอวํ ภนฺเต”ติ โข เต ภิกฺขู ภควโต ปจฺจสฺโสสุํ.\csromanspacer{} ภควา เอตทโวจ\csromandash{}อสฺสุตวโต ภิกฺขเว ปุถุชฺชนสฺส อุปฺปชฺชติ ลาโภ.\csromanspacer{} โส น อิติ ปฏิสญฺจิกฺขติ “อุปฺปนฺโน โข เม อยํ ลาโภ, โสจ โข อนิจฺโจ ทุกฺโข วิปริณามธมฺโม”ติ ยถาภูตํ นปฺปชานาติ.\csromanspacer{} อุปฺปชฺชติ อลาโภ ฯเปฯ อุปฺปชฺชติ ยโส.\csromanspacer{} อุปฺปชฺชติ อยโส.\csromanspacer{} อุปฺปชฺชติ นินฺทา.\csromanspacer{} อุปฺปชฺชติ ปสํสา.\csromanspacer{} อุปฺปชฺชติ สุขํ.\csromanspacer{} อุปฺปชฺชติ ทุกฺขํ.\csromanspacer{} โส น อิติ ปฏิสญฺจิกฺขติ “อุปฺปนฺนํ โข เม อิทํ ทุกฺขํ, ตญฺจ โข อนิจฺจํ ทุกฺขํ วิปริณามธมฺมนฺ”ติ ยถาภูตํ นปฺปชานาติ.}

\prose{ตสฺส ลาโภปิ จิตฺตํ ปริยาทาย ติฏฺฐติ, อลาโภปิ จิตฺตํ ปริยาทาย ติฏฺฐติ, ยโสปิ จิตฺตํ ปริยาทาย ติฏฺฐติ, อยโสปิ จิตฺตํ ปริยาทาย ติฏฺฐติ, นินฺทาปิ จิตฺตํ ปริยาทาย ติฏฺฐติ, ปสํสาปิ จิตฺตํ ปริยาทาย ติฏฺฐติ, สุขมฺปิ จิตฺตํ ปริยาทาย ติฏฺฐติ, ทุกฺขมฺปิ จิตฺตํ ปริยาทาย ติฏฺฐติ.\csromanspacer{} โส อุปฺปนฺนํ ลาภํ อนุรุชฺฌติ, อลาเภ ปฏิวิรุชฺฌติ.\csromanspacer{} อุปฺปนฺนํ ยสํ อนุรุชฺฌติ, อยเส ปฏิวิรุชฺฌติ.\csromanspacer{} อุปฺปนฺนํ ปสํสํ อนุรุชฺฌติ, นินฺทาย ปฏิวิรุชฺฌติ.\csromanspacer{} อุปฺปนฺนํ สุขํ อนุรุชฺฌติ, ทุกฺเข ปฏิวิรุชฺฌติ.\csromanspacer{} โส เอวํ อนุโรธวิโรธสมาปนฺโน น ปริมุจฺจติ ชาติยา ชราย มรเณน โสเกหิ ปริเทเวหิ ทุกฺเขหิ โทมนสฺเสหิ อุปายาเสหิ, น ปริมุจฺจติ ทุกฺขสฺมาติ วทามิ.}

\prose{สุตวโต จ โข ภิกฺขเว อริยสาวกสฺส อุปฺปชฺชติ ลาโภ.\csromanspacer{} โส อิติ ปฏิสญฺจิกฺขติ “อุปฺปนฺโน โข เม อยํ ลาโภ, โส จ โข อนิจฺโจ ทุกฺโข วิปริณามธมฺโม”ติ ยถาภูตํ ปชานาติ.\csromanspacer{} อุปฺปชฺชติ อลาโภ ฯเปฯ อุปฺปชฺชติ ยโส.\csromanspacer{} อุปฺปชฺชติ อยโส.\csromanspacer{} อุปฺปชฺชติ นินฺทา.\csromanspacer{} อุปฺปชฺชติ ปสํสา.\csromanspacer{} อุปฺปชฺชติ สุขํ.\csromanspacer{} อุปฺปชฺชติ ทุกฺขํ.\csromanspacer{} โส อิติ ปฏิสญฺจิกฺขติ “อุปฺปนฺนํ โข เม อิทํ ทุกฺขํ, ตญฺจ โข อนิจฺจํ ทุกฺขํ วิปริณามธมฺมนฺ”ติ ยถาภูตํ ปชานาติ.}

\prose{\csromanfolio{9}ตสฺส ลาโภปิ จิตฺตํ น ปริยาทาย ติฏฺฐติ, อลาโภปิ จิตฺตํ น ปริยาทาย ติฏฺฐติ, ยโสปิ จิตฺตํ น ปริยาทาย ติฏฺฐติ, อยโสปิ จิตฺตํ น ปริยาทาย ติฏฺฐติ, นินฺทาปิ จิตฺตํ น ปริยาทาย ติฏฺฐติ, ปสํสาปิ จิตฺตํ น ปริยาทาย ติฏฺฐติ, สุขมฺปิ จิตฺตํ น ปริยาทาย ติฏฺฐติ, ทุกฺขมฺปิ จิตฺตํ น ปริยาทาย ติฏฺฐติ.\csromanspacer{} โส อุปฺปนฺนํ ลาภํ นานุรุชฺฌติ, อลาเภ นปฺปฏิวิรุชฺฌติ.\csromanspacer{} อุปฺปนฺนํ ยสํ นานุรุชฺฌติ, อยเส นปฺปฏิวิรุชฺฌติ.\csromanspacer{} อุปฺปนฺนํ ปสํสํ นานุรุชฺฌติ, นินฺทาย นปฺปฏิวิรุชฺฌติ.\csromanspacer{} อุปฺปนฺนํ สุขํ นานุรุชฺฌติ, ทุกฺเข นปฺปฏิวิรุชฺฌติ.\csromanspacer{} โส เอวํ อนุโรธวิโรธวิปฺปหีโน ปริมุจฺจติ ชาติยา ชราย มรเณน โสเกหิ ปริเทเวหิ ทุกฺเขหิ โทมนสฺเสหิ อุปายาเสหิ, ปริมุจฺจติ ทุกฺขสฺมาติ วทามิ.\csromanspacer{} อยํ โข ภิกฺขเว วิเสโส อยํ อธิปฺปยาโส อิทํ นานากรณํ สุตวโต อริยสาวกสฺส อสฺสุตวตา ปุถุชฺชเนนาติ.}

\csromangathameasure{ลาโภ อลาโภ จ ยสายโส จ, \\ นินฺทา ปสํสา จ สุขํ ทุขญฺจ. \\ เอเต อนิจฺจา มนุเชสุ ธมฺมา, \\ อสสฺสตา วิปริณามธมฺมา. \\ เอเต จ ญตฺวา สติมา สุเมโธ, \\ อเวกฺขติ วิปริณามธมฺเม. \\ อิฏฺฐสฺส ธมฺมา น มเถนฺติ จิตฺตํ, \\ อนิฏฺฐโต โน ปฏิฆาตเมติ. \\ ตสฺสานุโรธา อถ วา วิโรธา, \\ วิธูปิตา อตฺถงฺคตา น สนฺติ. \\ ปทญฺจ ญตฺวา วิรชํ อโสกํ, \\ สมฺมปฺปชานาติ ภวสฺส ปารคูติ.ฉฏฺฐํ.}
\csromangathagroup{\csromangathastanza{ลาโภ อลาโภ จ ยสายโส จ, \\ นินฺทา ปสํสา จ สุขํ ทุขญฺจ. \\ เอเต อนิจฺจา มนุเชสุ ธมฺมา, \\ อสสฺสตา วิปริณามธมฺมา.}\csromangathastanza{เอเต จ ญตฺวา สติมา สุเมโธ, \\ อเวกฺขติ วิปริณามธมฺเม. \\ อิฏฺฐสฺส ธมฺมา น มเถนฺติ จิตฺตํ, \\ อนิฏฺฐโต โน ปฏิฆาตเมติ.}\csromangathastanza{ตสฺสานุโรธา อถ วา วิโรธา, \\ วิธูปิตา อตฺถงฺคตา น สนฺติ. \\ ปทญฺจ ญตฺวา วิรชํ อโสกํ, \\ สมฺมปฺปชานาติ ภวสฺส ปารคูติ.\csromanspacer{}\csromanspacer{}ฉฏฺฐํ.}}

\csromanheader{2}{7. เทวทตฺตวิปตฺติสุตฺต}
\proseitem{7}{เอกํ สมยํ ภควา ราชคเห วิหรติ คิชฺฌกูเฏ ปพฺพเต อจิรปกฺกนฺเต เทวทตฺเต.\csromanspacer{} ตตฺร ภควา เทวทตฺตํ อารพฺภ ภิกฺขู อามนฺเตสิ\csromandash{}สาธุ ภิกฺขเว ภิกฺขุ กาเลน กาลํ อตฺตวิปตฺติํ ปจฺจเวกฺขิตา โหติ, สาธุ ภิกฺขเว ภิกฺขุ กาเลน กาลํ ปรวิปตฺติํ ปจฺจเวกฺขิตา โหติ, สาธุ ภิกฺขเว ภิกฺขุ กาเลน กาลํ อตฺตสมฺปตฺติํ ปจฺจเวกฺขิตา โหติ, \csromanfolio{10}สาธุ ภิกฺขเว ภิกฺขุ กาเลน กาลํ ปรสมฺปตฺติํ ปจฺจเวกฺขิตา โหติ.\csromanspacer{} อฏฺฐหิ ภิกฺขเว อสทฺธมฺเมหิ อภิภูโต ปริยาทินฺนจิตฺโต เทวทตฺโต อาปายิโก เนรยิโก กปฺปฏฺโฐ อเตกิจฺโฉ.}

\prose{\csromansymbolfootnote{*}{วิ 4. 365 ปิฏฺเฐ.}กตเมหิ อฏฺฐหิ? ลาเภน หิ ภิกฺขเว อภิภูโต ปริยาทินฺนจิตฺโต เทวทตฺโต อาปายิโก เนรยิโก กปฺปฏฺโฐ อเตกิจฺโฉ.\csromanspacer{} อลาเภน ภิกฺขเว ฯเปฯ.\csromanspacer{} ยเสน ภิกฺขเว.\csromanspacer{} อยเสน ภิกฺขเว.\csromanspacer{} สกฺกาเรน ภิกฺขเว.\csromanspacer{} อสกฺกาเรน ภิกฺขเว.\csromanspacer{} ปาปิจฺฉตาย ภิกฺขเว.\csromanspacer{} ปาปมิตฺตตาย ภิกฺขเว อภิภูโต ปริยาทินฺนจิตฺโต เทวทตฺโต อาปายิโก เนรยิโก กปฺปฏฺโฐ อเตกิจฺโฉ.\csromanspacer{} อิเมหิ โข ภิกฺขเว อฏฺฐหิ อสทฺธมฺเมหิ อภิภูโต ปริยาทินฺนจิตฺโต เทวทตฺโต อาปายิโก เนรยิโก กปฺปฏฺโฐ อเตกิจฺโฉ.}

\prose{สาธุ ภิกฺขเว ภิกฺขุ อุปฺปนฺนํ ลาภํ อภิภุยฺย อภิภุยฺย วิหเรยฺย.\csromanspacer{} อุปฺปนฺนํ อลาภํ ฯเปฯ อุปฺปนฺนํ ยสํ.\csromanspacer{} อุปฺปนฺนํ อยสํ.\csromanspacer{} อุปฺปนฺนํ สกฺการํ.\csromanspacer{} อุปฺปนฺนํ อสกฺการํ.\csromanspacer{} อุปฺปนฺนํ ปาปิจฺฉตํ.\csromanspacer{} อุปฺปนฺนํ ปาปมิตฺตตํ อภิภุยฺย อภิภุยฺย วิหเรยฺย.}

\prose{กิญฺจ\footnote{กถญฺจ(ก)} ภิกฺขเว ภิกฺขุ อตฺถวสํ ปฏิจฺจ อุปฺปนฺนํ ลาภํ อภิภุยฺย อภิภุยฺย วิหเรยฺย.\csromanspacer{} อุปฺปนฺนํ อลาภํ ฯเปฯ อุปฺปนฺนํ ยสํ.\csromanspacer{} อุปฺปนฺนํ อยสํ.\csromanspacer{} อุปฺปนฺนํ สกฺการํ.\csromanspacer{} อุปฺปนฺนํ อสกฺการํ.\csromanspacer{} อุปฺปนฺนํ ปาปิจฺฉตํ.\csromanspacer{} อุปฺปนฺนํ ปาปมิตฺตตํ อภิภุยฺย อภิภุยฺย วิหเรยฺย.}

\prose{ยํ หิสฺส ภิกฺขเว อุปฺปนฺนํ ลาภํ อนภิภุยฺย\footnote{อนภิภุยฺย อนภิภุยฺย (ก)} วิหรโต อุปฺปชฺเชยฺยุํ อาสวา วิฆาตปริฬาหา, อุปฺปนฺนํ ลาภํ อภิภุยฺย\footnote{อภิภุยฺย อภิภุยฺย (ก)} วิหรโต เอวํส เต อาสวา วิฆาตปริฬาหา น โหนฺติ.\csromanspacer{} ยํ หิสฺส ภิกฺขเว อุปฺปนฺนํ อลาภํ ฯเปฯ อุปฺปนฺนํ ยสํ.\csromanspacer{} อุปฺปนฺนํ อยสํ.\csromanspacer{} อุปฺปนฺนํ สกฺการํ.\csromanspacer{} อุปฺปนฺนํ อสกฺการํ.\csromanspacer{} อุปฺปนฺนํ ปาปิจฺฉตํ.\csromanspacer{} อุปฺปนฺนํ ปาปมิตฺตตํ อนภิภุยฺย วิหรโต อุปฺปชฺเชยฺยุํ อาสวา วิฆาตปริฬาหา, อุปฺปนฺนํ ปาปมิตฺตตํ อภิภุยฺย วิหรโต เอวํส เต อาสวา วิฆาตปริฬาหา น โหนฺติ.\csromanspacer{} อิทํ โข ภิกฺขเว ภิกฺขุ อตฺถวสํ ปฏิจฺจ อุปฺปนฺนํ ลาภํ อภิภุยฺย อภิภุยฺย วิหเรยฺย.\csromanspacer{} อุปฺปนฺนํ อลาภํ ฯเปฯ อุปฺปนฺนํ ยสํ.\csromanspacer{} อุปฺปนฺนํ อยสํ.\csromanspacer{} อุปฺปนฺนํ สกฺการํ.\csromanspacer{} อุปฺปนฺนํ อสกฺการํ.\csromanspacer{} อุปฺปนฺนํ ปาปิจฺฉตํ.\csromanspacer{} อุปฺปนฺนํ ปาปมิตฺตตํ อภิภุยฺย อภิภุยฺย วิหเรยฺย.}

\prose{\csromanfolio{11}ตสฺมาติห ภิกฺขเว เอวํ สิกฺขิตพฺพํ “อุปฺปนฺนํ ลาภํ อภิภุยฺย อภิภุยฺย วิหริสฺสาม.\csromanspacer{} อุปฺปนฺนํ อลาภํ ฯเปฯ อุปฺปนํ ยสํ.\csromanspacer{} อุปฺปนฺนํ อยสํ.\csromanspacer{} อุปฺปนฺนํ สกฺการํ.\csromanspacer{} อุปฺปนฺนํ อสกฺการํ.\csromanspacer{} อุปฺปนฺนํ ปาปิจฺฉตํ.\csromanspacer{} อุปฺปนฺนํ ปาปมิตฺตตํ อภิภุยฺย อภิภุยฺย วิหริสฺสามา”ติ.\csromanspacer{} เอวํ หิ โว ภิกฺขเว สิกฺขิตพฺพนฺติ.\csromanspacer{}\csromanspacer{}สตฺตมํ.}

\csromanheader{2}{8. อุตฺตรวิปตฺติสุตฺต}
\proseitem{8}{เอกํ สมยํ อายสฺมา อุตฺตโร มหิสวตฺถุสฺมิํ วิหรติ สงฺเขยฺยเก ปพฺพเต วฏชาลิกายํ\footnote{ธวชาลิกายํ (สี), วฏฺฏชาลิกายํ (สฺยา)}.\csromanspacer{} ตตฺร โข อายสฺมา อุตฺตโร ภิกฺขู อามนฺเตสิ “สาธาวุโส ภิกฺขุ กาเลน กาลํ อตฺตวิปตฺติํ ปจฺจเวกฺขิตา โหติ, สาธาวุโส ภิกฺขุ กาเลน กาลํ ปรวิปตฺติํ ปจฺจเวกฺขิตา โหติ, สาธาวุโส ภิกฺขุ กาเลน กาลํ อตฺตสมฺปตฺติํ ปจฺจเวกฺขิตา โหติ, สาธาวุโส ภิกฺขุ กาเลน กาลํ ปรสมฺปตฺติํ ปจฺจเวกฺขิตา โหตี”ติ.}

\prose{เตน โข ปน สมเยน เวสฺสวโณ มหาราชา อุตฺตราย ทิสาย ทกฺขิณํ ทิสํ คจฺฉติ เกนจิเทว กรณีเยน.\csromanspacer{} อสฺโสสิ โข เวสฺสวโณ มหาราชา อายสฺมโต อุตฺตรสฺส มหิสวตฺถุสฺมิํ สงฺเขยฺยเก ปพฺพเต วฏชาลิกายํ ภิกฺขูนํ เอวํ ธมฺมํ เทเสนฺตสฺส “สาธาวุโส ภิกฺขุ กาเลน กาลํ อตฺตวิปตฺติํ ปจฺจเวกฺขิตา โหติ, สาธาวุโส ภิกฺขุ กาเลน กาลํ ปรวิปตฺติํ ปจฺจเวกฺขิตา โหติ, สาธาวุโส ภิกฺขุ กาเลน กาลํ อตฺตสมฺปตฺติํ ปจฺจเวกฺขิตา โหติ, สาธาวุโส ภิกฺขุ กาเลน กาลํ ปรสมฺปตฺติํ ปจฺจเวกฺขิตา โหตี”ติ.}

\prose{อถ โข เวสฺสวโณ มหาราชา เสยฺยถาปิ นาม พลวา ปุริโส สมิญฺชิตํ\footnote{สมฺมิญฺชิตํ (สี, สฺยา, กํ, อิ)} วา พาหํ ปสาเรยฺย, ปสาริตํ วา พาหํ สมิญฺเชยฺย\footnote{สมฺมิญฺเชยฺย (สี, สฺยา, กํ, อิ)}.\csromanspacer{} เอวเมวํ มหิสวตฺถุสฺมิํ สงฺเขยฺยเก ปพฺพเต วฏชาลิกายํ อนฺตรหิโต เทเวสุ ตาวติํเสสุ ปาตุรโหสิ.\csromanspacer{} อถ โข เวสฺสวโณ มหาราชา เยน สกฺโก เทวานมินฺโท เตนุปสงฺกมิ, อุปสงฺกมิตฺวา สกฺกํ เทวานมินฺทํ เอตทโวจ “ยคฺเฆ มาริส ชาเนยฺยาสิ, เอโส อายสฺมา อุตฺตโร มหิสวตฺถุสฺมิํ สงฺเขยฺยเก ปพฺพเต วฏชาลิกายํ ภิกฺขูนํ เอวํ ธมฺมํ เทเสติ “สาธาวุโส ภิกฺขุ กาเลน กาลํ อตฺตวิปตฺติํ ปจฺจเวกฺขิตา \csromanfolio{12}โหติ, สาธาวุโส ภิกฺขุ กาเลน กาลํ ปรวิปตฺติํ ฯเปฯ อตฺตสมฺปตฺติํ.\csromanspacer{} ปรสมฺปตฺติํ ปจฺจเวกฺขิตา โหตี’ติ”.}

\prose{อถ โข สกฺโก เทวานมินฺโท เสยฺยถาปิ นาม พลวา ปุริโส สมิญฺชิตํ วา พาหํ ปสาเรยฺย, ปสาริตํ วา พาหํ สมิญฺเชยฺย.\csromanspacer{} เอวเมวํ เทเวสุ ตาวติํเสสุ อนฺตรหิโต มหิสวตฺถุสฺมิํ สงฺเขยฺยเก ปพฺพเต วฏชาลิกายํ อายสฺมโต อุตฺตรสฺส สมฺมุเข ปาตุรโหสิ.\csromanspacer{} อถ โข สกฺโก เทวานมินฺโท เยนายสฺมา อุตฺตโร เตนุปสงฺกมิ, อุปสงฺกมิตฺวา อายสฺมนฺตํ อุตฺตรํ อภิวาเทตฺวา เอกมนฺตํ อฏฺฐาสิ, เอกมนฺตํ ฐิโต โข สกฺโก เทวานมินฺโท อายสฺมนฺตํ อุตฺตรํ เอตทโวจ\csromandash{}}

\prose{สจฺจํ กิร ภนฺเต อายสฺมา อุตฺตโร ภิกฺขูนํ เอวํ ธมฺมํ เทเสสิ “สาธาวุโส ภิกฺขุ กาเลน กาลํ อตฺตวิปตฺติํ ปจฺจเวกฺขิตา โหติ, สาธาวุโส ภิกฺขุ กาเลน กาลํ ปรวิปตฺติํ ฯเปฯ อตฺตสมฺปตฺติํ.\csromanspacer{} ปรสมฺปตฺติํ ปจฺจเวกฺขิตา โหตี”ติ.\csromanspacer{} เอวํ เทวานมินฺทาติ.\csromanspacer{} กิํ ปนิทํ\footnote{กิํ ปน (สฺยา)} ภนฺเต อายสฺมโต อุตฺตรสฺส สกํ ปฏิภานํ\footnote{สกปฏิภานํ อุปาทาย (ก)} อุทาหุ ตสฺส ภควโต วจนํ อรหโต สมฺมาสมฺพุทฺธสฺสาติ.\csromanspacer{} เตน หิ เทวานมินฺท อุปมํ เต กริสฺสามิ, อุปมาย มิเธกจฺเจ วิญฺญู ปุริสา ภาสิตสฺส อตฺถํ อาชานนฺติ.}

\prose{เสยฺยถาปิ เทวานมินฺท คามสฺส วา นิคมสฺส วา อวิทูเร มหาธญฺญราสิ, ตโต มหาชนกาโย ธญฺญํ อาหเรยฺย กาเชหิปิ ปิฏเกหิปิ อุจฺฉงฺเคหิปิ อญฺชลีหิปิ.\csromanspacer{} โย นุ โข เทวานมินฺท ตํ มหาชนกายํ อุปสงฺกมิตฺวา เอวํ ปุจฺเฉยฺย “กุโต อิมํ ธญฺญํ อาหรถา”ติ.\csromanspacer{} กถํ พฺยากรมาโน นุ โข เทวานมินฺท โส มหาชนกาโย สมฺมา พฺยากรมาโน พฺยากเรยฺยาติ. “อมุมฺหา มหาธญฺญราสิมฺหา อาหรามา”ติ โข ภนฺเต โส มหาชนกาโย สมฺมา พฺยากรมาโน พฺยากเรยฺยาติ.\csromanspacer{} เอวเมวํ โข เทวานมินฺท ยํ กิญฺจิ สุภาสิตํ, สพฺพํ ตํ ตสฺส ภควโต วจนํ อรหโต สมฺมาสมฺพุทฺธสฺส.\csromanspacer{} ตโต อุปาทายุปาทาย มยญฺจญฺเญ จ ภณามาติ.}

\prose{อจฺฉริยํ ภนฺเต, อพฺภุตํ ภนฺเต, ยาว สุภาสิตญฺจิทํ อายสฺมตา อุตฺตเรน ยํ กิญฺจิ สุภาสิตํ, สพฺพํ ตํ ตสฺส ภควโต วจนํ อรหโต \csromanfolio{13}สมฺมาสมฺพุทฺธสฺส.\csromanspacer{} ตโต อุปาทายุปาทาย มยญฺจญฺเญ จ ภณามาติ.\csromanspacer{} เอกมิทํ ภนฺเต อุตฺตร สมยํ ภควา ราชคเห วิหรติ คิชฺฌกูเฏ ปพฺพเต อจิรปกฺกนฺเต เทวทตฺเต.\csromanspacer{} ตตฺร โข ภควา เทวทตฺตํ อารพฺภ ภิกฺขู อามนฺเตสิ\csromandash{}}

\prose{สาธุ ภิกฺขเว ภิกฺขุ กาเลน กาลํ อตฺตวิปตฺติํ ปจฺจเวกฺขิตา โหติ, สาธุ ภิกฺขเว ภิกฺขุ กาเลน กาลํ ปรวิปตฺติํ ฯเปฯ อตฺตสมฺปตฺติํ.\csromanspacer{} ปรสมฺปตฺติํ ปจฺจเวกฺขิตา โหติ.\csromanspacer{} อฏฺฐหิ ภิกฺขเว อสทฺธมฺเมหิ อภิภูโต ปริยาทินฺนจิตฺโต เทวทตฺโต อาปายิโก เนรยิโก กปฺปฏฺโฐ อเตกิจฺโฉ.\csromanspacer{} กตเมหิ อฏฺฐหิ? ลาเภน หิ ภิกฺขเว อภิภูโต ปริยาทินฺนจิตฺโต เทวทตฺโต อาปายิโก เนรยิโก กปฺปฏฺโฐ อเตกิจฺโฉ.\csromanspacer{} อลาเภน ภิกฺขเว ฯเปฯ.\csromanspacer{} ยเสน ภิกฺขเว.\csromanspacer{} อยเสน ภิกฺขเว.\csromanspacer{} สกฺกาเรน ภิกฺขเว.\csromanspacer{} อสกฺกาเรน ภิกฺขเว.\csromanspacer{} ปาปิจฺฉตาย ภิกฺขเว.\csromanspacer{} ปาปมิตฺตตาย ภิกฺขเว อภิภูโต ปริยาทินฺนจิตฺโต เทวทตฺโต อาปายิโก เนรยิโก กปฺปฏฺโฐ อเตกิจฺโฉ.\csromanspacer{} อิเมหิ โข ภิกฺขเว อฏฺฐหิ อสทฺธมฺเมหิ อภิภูโต ปริยาทินฺนจิตฺโต เทวทตฺโต อาปายิโก เนรยิโก กปฺปฏฺโฐ อเตกิจฺโฉ.}

\prose{สาธุ ภิกฺขเว ภิกฺขุ อุปฺปนฺนํ ลาภํ อภิภุยฺย อภิภุยฺย วิหเรยฺย.\csromanspacer{} อุปฺปนฺนํ อลาภํ ฯเปฯ อุปฺปนฺนํ ยสํ.\csromanspacer{} อุปฺปนฺนํ อยสํ.\csromanspacer{} อุปฺปนฺนํ สกฺการํ.\csromanspacer{} อุปฺปนฺนํ อสกฺการํ.\csromanspacer{} อุปฺปนฺนํ ปาปิจฺฉตํ.\csromanspacer{} อุปฺปนฺนํ ปาปมิตฺตตํ อภิภุยฺย อภิภุยฺย วิหเรยฺย.}

\prose{กิญฺจ ภิกฺขเว ภิกฺขุ อตฺถวสํ ปฏิจฺจ อุปฺปนฺนํ ลาภํ อภิภุยฺย อภิภุยฺย วิหเรยฺย.\csromanspacer{} อุปฺปนฺนํ อลาภํ ฯเปฯ อุปฺปนฺนํ ยสํ.\csromanspacer{} อุปฺปนฺนํ อยสํ.\csromanspacer{} อุปฺปนฺนํ สกฺการํ.\csromanspacer{} อุปฺปนฺนํ อสกฺการํ.\csromanspacer{} อุปฺปนฺนํ ปาปิจฺฉตํ.\csromanspacer{} อุปฺปนฺนํ ปาปมิตฺตตํ อภิภุยฺย อภิภุยฺย วิหเรยฺย.}

\prose{ยํ หิสฺส ภิกฺขเว อุปฺปนฺนํ ลาภํ อนภิภุยฺย วิหรโต อุปฺปชฺเชยฺยุํ อาสวา วิฆาตปริฬาหา, อุปฺปนฺนํ ลาภํ อภิภุยฺย วิหรโต เอวํส เต อาสวา วิฆาตปริฬาหา น โหนฺติ.\csromanspacer{} ยํ หิสฺส ภิกฺขเว อุปฺปนฺนํ อลาภํ ฯเปฯ อุปฺปนฺนํ ยสํ.\csromanspacer{} อุปฺปนฺนํ อยสํ.\csromanspacer{} อุปฺปนฺนํ สกฺการํ.\csromanspacer{} อุปฺปนฺนํ อสกฺการํ.\csromanspacer{} อุปฺปนฺนํ ปาปิจฺฉตํ.\csromanspacer{} อุปฺปนฺนํ ปาปมิตฺตตํ อนภิภุยฺย วิหรโต อุปฺปชฺเชยฺยุํ อาสวา วิฆาตปริฬาหา, อุปฺปนฺนํ ปาปมิตฺตตํ อภิภุยฺย วิหรโต \csromanfolio{14}เอวํส เต อาสวา วิฆาตปริฬาหา น โหนฺติ.\csromanspacer{} อิทํ โข ภิกฺขเว ภิกฺขุ อตฺถวสํ ปฏิจฺจ อุปฺปนฺนํ ลาภํ อภิภุยฺย อภิภุยฺย วิหเรยฺย.\csromanspacer{} อุปฺปนฺนํ อลาภํ ฯเปฯ อุปฺปนฺนํ ยสํ.\csromanspacer{} อุปฺปนฺนํ อยสํ.\csromanspacer{} อุปฺปนฺนํ สกฺการํ.\csromanspacer{} อุปฺปนฺนํ อสกฺการํ.\csromanspacer{} อุปฺปนฺนํ ปาปิจฺฉตํ.\csromanspacer{} อุปฺปนฺนํ ปาปมิตฺตตํ อภิภุยฺย อภิภุยฺย วิหเรยฺย.}

\prose{ตสฺมาติห ภิกฺขเว เอวํ สิกฺขิตพฺพํ “อุปฺปนฺนํ ลาภํ อภิภุยฺย อภิภุยฺย วิหริสฺสาม.\csromanspacer{} อุปฺปนฺนํ อลาภํ ฯเปฯ อุปฺปนฺนํ ยสํ.\csromanspacer{} อุปฺปนฺนํ อยสํ.\csromanspacer{} อุปฺปนฺนํ สกฺการํ.\csromanspacer{} อุปฺปนฺนํ อสกฺการํ.\csromanspacer{} อุปฺปนฺนํ ปาปิจฺฉตํ.\csromanspacer{} อุปฺปนฺนํ ปาปมิตฺตตํ อภิภุยฺย อภิภุยฺย วิหริสฺสามา”ติ.\csromanspacer{} เอวํ หิ โว ภิกฺขเว สิกฺขิตพฺพนฺติ.}

\prose{เอตฺตาวตา ภนฺเต อุตฺตร มนุสฺเสสุ จตสฺโส ปริสา ภิกฺขู ภิกฺขุนิโย อุปาสกา อุปาสิกาโย.\csromanspacer{} นายํ ธมฺมปริยาโย กิสฺมิญฺจิ อุปฏฺฐิโต\footnote{ปติฏฺฐิโต (สี, สฺยา)}.\csromanspacer{} อุคฺคณฺหตุ ภนฺเต อายสฺมา อุตฺตโร อิมํ ธมฺมปริยายํ, ปริยาปุณาตุ ภนฺเต อายสฺมา อุตฺตโร อิมํ ธมฺมปริยายํ, ธาเรตุ ภนฺเต อายสฺมา อุตฺตโร อิมํ ธมฺมปริยายํ.\csromanspacer{} อตฺถสํหิโต อยํ ภนฺเต ธมฺมปริยาโย อาทิพฺรหฺมจริยโกติ\footnote{อาทิพฺรหฺมจริยิโก (สี, ก)}.\csromanspacer{}\csromanspacer{}อฏฺฐมํ.}

\csromanheader{2}{9. นนฺทสุตฺต}
\proseitem{9}{“กุลปุตฺโต”ติ ภิกฺขเว นนฺทํ สมฺมา วทมาโน วเทยฺย, “พลวา”ติ ภิกฺขเว นนฺทํ สมฺมา วทมาโน วเทยฺย, “ปาสาทิโก”ติ ภิกฺขเว นนฺทํ สมฺมา วทมาโน วเทยฺย, “ติพฺพราโค”ติ ภิกฺขเว นนฺทํ สมฺมา วทมาโน วเทยฺย.\csromanspacer{} กิมญฺญตฺร ภิกฺขเว นนฺโท อินฺทฺริเยสุ คุตฺตทฺวาโร โภชเน มตฺตญฺญู ชาคริยํ อนุยุตฺโต สติสมฺปชญฺเญน สมนฺนาคโต, เยหิ\footnote{เยน (ก)} นนฺโท สกฺโกติ ปริปุณฺณํ ปริสุทฺธํ พฺรหฺมจริยํ จริตุํ.\csromanspacer{} ตตฺริทํ ภิกฺขเว นนฺทสฺส อินฺทฺริเยสุ คุตฺตทฺวารตาย โหติ.\csromanspacer{} สเจ ภิกฺขเว นนฺทสฺส ปุรตฺถิมา ทิสา อาโลเกตพฺพา โหติ, สพฺพํ เจตสา สมนฺนาหริตฺวา นนฺโท ปุรตฺถิมํ ทิสํ อาโลเกติ “เอวํ เม ปุรตฺถิมํ ทิสํ อาโลกยโต นาภิชฺฌาโทมนสฺสา ปาปกา อกุสลา ธมฺมา อนฺวาสฺสวิสฺสนฺตี”ติ.\csromanspacer{} อิติห ตตฺถ สมฺปชาโน โหติ.}

\prose{\csromanfolio{15}สเจ ภิกฺขเว นนฺทสฺส ปจฺฉิมา ทิสา อาโลเกตพฺพา โหติ ฯเปฯ อุตฺตรา ทิสา อาโลเกตพฺพา โหติ.\csromanspacer{} ทกฺขิณา ทิสา อาโลเกตพฺพา โหติ.\csromanspacer{} อุทฺธํ อุลฺโลเกตพฺพา โหติ.\csromanspacer{} อโธ โอโลเกตพฺพา โหติ.\csromanspacer{} อนุทิสา อนุวิโลเกตพฺพา โหติ, สพฺพํ เจตสา สมนฺนาหริตฺวา นนฺโท อนุทิสํ อนุวิโลเกติ “เอวํ เม อนุทิสํ อนุวิโลกยโต นาภิชฺฌาโทมนสฺสา ปาปกา อกุสลา ธมฺมา อนฺวาสฺสวิสฺสนฺตี”ติ.\csromanspacer{} อิติห ตตฺถ สมฺปชาโน โหติ.\csromanspacer{} อิทํ โข ภิกฺขเว นนฺทสฺส อินฺทฺริเยสุ คุตฺตทฺวารตาย โหติ.}

\prose{ตตฺริทํ ภิกฺขเว นนฺทสฺส โภชเน มตฺตญฺญุตาย โหติ, อิธ ภิกฺขเว นนฺโท ปฏิสงฺขา โยนิโส อาหารํ อาหาเรติ, เนว ทวาย น มทาย น มณฺฑนาย น วิภูสนาย, ยาวเทว อิมสฺส กายสฺส ฐิติยา ยาปนาย วิหิํสูปรติยา พฺรหฺมจริยานุคฺคหาย, อิติ ปุราณญฺจ เวทนํ ปฏิหงฺขามิ, นวญฺจ เวทนํ น อุปฺปาเทสฺสามิ, ยาตฺรา จ เม ภวิสฺสติ อนวชฺชตา จ ผาสุวิหาโร จาติ.\csromanspacer{} อิทํ โข ภิกฺขเว นนฺทสฺส โภชเน มตฺตญฺญุตาย โหติ.}

\prose{ตตฺริทํ ภิกฺขเว นนฺทสฺส ชาคริยานุโยคสฺมิํ โหติ, อิธ ภิกฺขเว นนฺโท ทิวสํ จงฺกเมน นิสชฺชาย อาวรณีเยหิ ธมฺเมหิ จิตฺตํ ปริโสเธติ, รตฺติยา ปฐมํ ยามํ จงฺกเมน นิสชฺชาย อาวรณีเยหิ ธมฺเมหิ จิตฺตํ ปริโสเธติ, รตฺติยา มชฺฌิมํ ยามํ ทกฺขิเณน ปสฺเสน สีหเสยฺยํ กปฺเปติ ปาเท ปาทํ อจฺจาธาย สโต สมฺปชาโน อุฏฺฐานสญฺญํ มนสิ กริตฺวา, รตฺติยา ปจฺฉิมํ ยามํ ปจฺจุฏฺฐาย จงฺกเมน นิสชฺชาย อาวรณีเยหิ ธมฺเมหิ จิตฺตํ ปริโสเธติ.\csromanspacer{} อิทํ โข ภิกฺขเว นนฺทสฺส ชาคริยานุโยคสฺมิํ โหติ.}

\prose{ตตฺริทํ ภิกฺขเว นนฺทสฺส สติสมฺปชญฺญสฺมิํ โหติ, อิธ ภิกฺขเว นนฺทสฺส วิทิตา เวทนา อุปฺปชฺชนฺติ, วิทิตา อุปฏฺฐหนฺติ, วิทิตา อพฺภตฺถํ คจฺฉนฺติ.\csromanspacer{} วิทิตา สญฺญา ฯเปฯ วิทิตา วิตกฺกา ฯเปฯ อพฺภตฺถํ คจฺฉนฺติ.\csromanspacer{} อิทํ โข ภิกฺขเว นนฺทสฺส สติสมฺปชญฺญสฺมิํ โหติ.}

\prose{กิมญฺญตฺร ภิกฺขเว นนฺโท อินฺทฺริเยสุ คุตฺตทฺวาโร โภชเน มตฺตญฺญู ชาคริยํ อนุยุตฺโต สติสมฺปชญฺเญน สมนฺนาคโต, เยหิ นนฺโท สกฺโกติ ปริปุณฺณํ ปริสุทฺธํ พฺรหฺมจริยํ จริตุนฺติ.\csromanspacer{}\csromanspacer{}นวมํ.}

\csromanheader{2}{10. การณฺฑวสุตฺต}
\proseitem{10}{\csromanfolio{16}เอกํ สมยํ ภควา จมฺปายํ วิหรติ คคฺคราย โปกฺขรณิยา ตีเร.\csromanspacer{} เตน โข ปน สมเยน ภิกฺขู ภิกฺขุํ อาปตฺติยา โจเทนฺติ.\csromanspacer{} โส ภิกฺขุ ภิกฺขูหิ อาปตฺติยา โจทิยมาโน อญฺเญนาญฺญํ ปฏิจรติ, พหิทฺธา กถํ อปนาเมติ, โกปญฺจ โทสญฺจ อปฺปจฺจยญฺจ ปาตุกโรติ.}

\prose{อถ โข ภควา ภิกฺขู อามนฺเตสิ “นิทฺธมเถตํ ภิกฺขเว ปุคฺคลํ, นิทฺธมเถตํ ภิกฺขเว ปุคฺคลํ.\csromanspacer{} อปเนยฺเยโส\footnote{อปเนยฺโย โส (สี), อปเนยฺโย (สฺยา)} ภิกฺขเว ปุคฺคโล, กิํ โว เตน ปรปุตฺเตน วิโสธิเตน\footnote{กิํ โว ปรปุตฺโต วิเหฐิยติ (สี), กิํ ปรปุตฺโต วิเหเฐติ (สฺยา), กิํ โว ปรปุตฺตา วิเหเฐติ (อิ), กิํ โส ปรปุตฺโต วิโสเธติ (ก)}.\csromanspacer{} อิธ ภิกฺขเว เอกจฺจสฺส ปุคฺคลสฺส ตาทิสํเยว โหติ อภิกฺกนฺตํ ปฏิกฺกนฺตํ อาโลกิตํ วิโลกิตํ สมิญฺชิตํ ปสาริตํ สํฆาฏิปตฺตจีวรธารณํ, เสยฺยถาปิ อญฺเญสํ ภทฺทกานํ ภิกฺขูนํ, ยาวสฺส ภิกฺขู อาปตฺติํ น ปสฺสนฺติ.\csromanspacer{} ยโต จ ขฺวสฺส ภิกฺขู อาปตฺติํ ปสฺสนฺติ, ตเมนํ เอวํ ชานนฺติ “สมณทูสีวายํ\footnote{สมณรูปี จายํ (ก)} สมณปลาโป สมณการณฺฑโว”ติ\footnote{สมณกรณฺฑโวติ (ก)}.\csromanspacer{} ตเมนํ อิติ วิทิตฺวา พหิทฺธา นาเสนฺติ.\csromanspacer{} ตํ กิสฺส เหตุ? “มา อญฺเญ ภทฺทเก ภิกฺขู ทูเสสี”ติ.}

\prose{เสยฺยถาปิ ภิกฺขเว สมฺปนฺเน ยวกรเณ ยวทูสี\footnote{ยวรูปี (ก)} ชาเยถ ยวปลาโป ยวการณฺฑโวติ.\csromanspacer{} ตสฺส ตาทิสํเยว มูลํ โหติ เสยฺยถาปิ อญฺเญสํ ภทฺทกานํ ยวานํ, ตาทิสํเยว นาฬํ โหติ เสยฺยถาปิ อญฺเญสํ ภทฺทกานํ ยวานํ, ตาทิสํเยว ปตฺตํ โหติ เสยฺยถาปิ อญฺเญสํ ภทฺทกานํ ยวานํ, ยาวสฺส สีสํ น นิพฺพตฺตติ.\csromanspacer{} ยโต จ ขฺวสฺส สีสํ นิพฺพตฺตติ, ตเมนํ เอวํ “ชานนฺติ ยวทูสีวายํ ยวปลาโป ยวการณฺฑโว”ติ.\csromanspacer{} ตเมนํ อิติ วิทิตฺวา สมูลํ อุปฺปาเฏตฺวา พหิทฺธา ยวกรณสฺส ฉฑฺเฑนฺติ.\csromanspacer{} ตํ กิสฺส เหตุ? “มา อญฺเญ ภทฺทเก ยเว ทูเสสี”ติ.}

\prose{เอวเมวํ โข ภิกฺขเว อิเธกจฺจสฺส ปุคฺคลสฺส ตาทิสํเยว โหติ อภิกฺกนฺตํ ปฏิกฺกนฺตํ อาโลกิตํ วิโลกิตํ สมิญฺชิตํ ปสาริตํ สํฆาฏิปตฺตจีวรธารณํ, เสยฺยถาปิ อญฺเญสํ ภทฺทกานํ ภิกฺขูนํ, \csromanfolio{17}ยาวสฺส ภิกฺขู อาปตฺติํ น ปสฺสนฺติ.\csromanspacer{} ยโต จ ขฺวสฺส ภิกฺขู อาปตฺติํ ปสฺสนฺติ, ตเมนํ เอวํ ชานนฺติ “สมณทูสีวายํ สมณปลาโป สมณการณฺฑโว”ติ.\csromanspacer{} ตเมนํ อิติ วิทิตฺวา พหิทฺธา นาเสนฺติ.\csromanspacer{} ตํ กิสฺส เหตุ? “มา อญฺเญ ภทฺทเก ภิกฺขู ทูเสสี”ติ.}

\prose{เสยฺยถาปิ ภิกฺขเว มหโต ธญฺญราสิสฺส ผุณมานสฺส\footnote{วุยฺหมานสฺส (สี, อิ), ผุสยมานสฺส (สฺยา), ปุนมานสฺส (?)} ตตฺถ ยานิ ตานิ ธญฺญานิ ทฬฺหานิ สารวนฺตานิ, ตานิ เอกมนฺตํ ปุญฺชํ โหติ.\csromanspacer{} ยานิ ปน ตานิ ธญฺญานิ ทุพฺพลานิ ปลาปานิ, ตานิ วาโต เอกมนฺตํ อปวหติ\footnote{อปกสฺสติ (สี)}.\csromanspacer{} ตเมนํ สามิกา สมฺมชฺชนิํ คเหตฺวา ภิยฺโยโสมตฺตาย อปสมฺมชฺชนฺติ.\csromanspacer{} ตํ กิสฺส เหตุ? “มา อญฺเญ ภทฺทเก ธญฺเญ ทูเสสี”ติ.\csromanspacer{} เอวเมวํ โข ภิกฺขเว อิเธกจฺจสฺส ปุคฺคลสฺส ตาทิสํเยว โหติ อภิกฺกนฺตํ ปฏิกฺกนฺตํ อาโลกิตํ วิโลกิตํ สมิญฺชิตํ ปสาริตํ สํฆาฏิปตฺตจีวรธารณํ, เสยฺยถาปิ อญฺเญสํ ภทฺทกานํ ภิกฺขูนํ, ยาวสฺส ภิกฺขู อาปตฺติํ น ปสฺสนฺติ.\csromanspacer{} ยโต จ ขฺวสฺส ภิกฺขู อาปตฺติํ ปสฺสนฺติ, ตเมนํ เอวํ ชานนฺติ “สมณทูสีวายํ สมณปลาโป สมณการณฺฑโว”ติ.\csromanspacer{} ตเมนํ อิติ วิทิตฺวา พหิทฺธา นาเสนฺติ.\csromanspacer{} ตํ กิสฺส เหตุ? “มา อญฺเญ ภทฺทเก ภิกฺขู ทูเสสี”ติ.}

\prose{เสยฺยถาปิ ภิกฺขเว ปุริโส อุทปานปนาฬิยตฺถิโก ติณฺหํ กุฐาริํ\footnote{กุธาริํ (สฺยา, กํ, ก)} อาทาย วนํ ปวิเสยฺย, โส ยํ ยเทว รุกฺขํ กุฐาริปาเสน อาโกเฏยฺย.\csromanspacer{} ตตฺถ ยานิ ตานิ รุกฺขานิ ทฬฺหานิ สารวนฺตานิ, ตานิ กุฐาริปาเสน อาโกฏิตานิ กกฺขฬํ ปฏินทนฺติ.\csromanspacer{} ยานิ ปน ตานิ รุกฺขานิ อนฺโตปูตีนิ อวสฺสุตานิ กสมฺพุชาตานิ, ตานิ กุฐาริปาเสน อาโกฏิตานิ ททฺทรํ ปฏินทนฺติ.\csromanspacer{} ตเมนํ มูเล ฉินฺทติ, มูเล ฉินฺทิตฺวา อคฺเค ฉินฺทติ, อคฺเค ฉินฺทิตฺวา อนฺโต สุวิโสธิตํ วิโสเธติ, อนฺโต สุวิโสธิตํ วิโสเธตฺวา อุทปานปนาฬิํ โยเชติ.\csromanspacer{} เอวเมวํ โข ภิกฺขเว อิเธกจฺจสฺส ปุคฺคลสฺส ตาทิสํเยว โหติ อภิกฺกนฺตํ ปฏิกฺกนฺตํ อาโลกิตํ วิโลกิตํ สมิญฺชิตํ ปสาริตํ สํฆาฏิปตฺตจีวรธารณํ, เสยฺยถาปิ อญฺเญสํ ภทฺทกานํ ภิกฺขูนํ, ยาวสฺส ภิกฺขู อาปตฺติํ น ปสฺสนฺติ.\csromanspacer{} ยโต จ ขฺวสฺส ภิกฺขู อาปตฺติํ ปสฺสนฺติ, ตเมนํ \csromanfolio{18}เอวํ ชานนฺติ “สมณทูสีวายํ สมณปลาโป สมณการณฺฑโว”ติ.\csromanspacer{} ตเมนํ อิติ วิทิตฺวา พหิทฺธา นาเสนฺติ.\csromanspacer{} ตํ กิสฺส เหตุ? “มา อญฺเญ ภทฺทเก ภิกฺขู ทูเสสี”ติ.}

\csromangathameasure{สํวาสา’ยํ วิชานาถ, ปาปิจฺโฉ โกธโน อิติ. \\ มกฺขี ถมฺภี ปฬาสี จ, อิสฺสุกี มจฺฉรี สโฐ. \\ สนฺตวาโจ ชนวติ, สมโณ วิย ภาสติ. \\ รโห กโรติ กรณํ, ปาปทิฏฺฐิ อนาทโร. \\ สํสปฺปี จ มุสาวาที, ตํ วิทิตฺวา ยถาตถํ. \\ สพฺเพ สมคฺคา หุตฺวาน, อภินิพฺพชฺชยาถ นํ. \\ การณฺฑวํ นิทฺธมถ, กสมฺพุํ อปกสฺสถ. \\ ตโต ปลาเป วาเหถ, อสฺสมเณ สมณมานิเน. \\ นิทฺธมิตฺวาน ปาปิจฺเฉ, ปาป-อาจารโคจเร. \\ สุทฺธาสุทฺเธหิ สํวาสํ, กปฺปยโวฺห ปติสฺสตา. \\ ตโต สมคฺคา นิปกา, ทุกฺขสฺสนฺตํ กริสฺสถาติ.ทสมํ. เมตฺตาวคฺโค ปฐโม.}
\csromangathagroup{\csromangathastanza{สํวาสา’ยํ วิชานาถ, ปาปิจฺโฉ โกธโน อิติ. \\ มกฺขี ถมฺภี ปฬาสี จ, อิสฺสุกี มจฺฉรี สโฐ.}\csromangathastanza{สนฺตวาโจ ชนวติ, สมโณ วิย ภาสติ. \\ รโห กโรติ กรณํ, ปาปทิฏฺฐิ อนาทโร.}\csromangathastanza{สํสปฺปี จ มุสาวาที, ตํ วิทิตฺวา ยถาตถํ. \\ สพฺเพ สมคฺคา หุตฺวาน, อภินิพฺพชฺชยาถ\footnote{อภินิพฺพิชฺชเยถ (ก)} นํ.}\csromangathastanza{การณฺฑวํ\footnote{กรณฺฑวํ (ก) ขุ 1. 322 ปิฏฺเฐปิ ปสฺสิตพฺพํ.} นิทฺธมถ, กสมฺพุํ อปกสฺสถ\footnote{อวกสฺสถ (ก)}. \\ ตโต ปลาเป วาเหถ, อสฺสมเณ สมณมานิเน.}\csromangathastanza{นิทฺธมิตฺวาน ปาปิจฺเฉ, ปาป-อาจารโคจเร. \\ สุทฺธาสุทฺเธหิ สํวาสํ, กปฺปยโวฺห ปติสฺสตา. \\ ตโต สมคฺคา นิปกา, ทุกฺขสฺสนฺตํ กริสฺสถาติ.\csromanspacer{}\csromanspacer{}ทสมํ.\csromanspacer{} เมตฺตาวคฺโค ปฐโม.}}

\titlehead{\textbf{ตสฺสุทฺทานํ}}
\csromangathameasure{เมตฺตํ ปญฺญา จ เทฺว ปิยา, เทฺว โลกา เทฺว วิปตฺติโย. \\ เทวทตฺโต จ อุตฺตโร, นนฺโท การณฺฑเวน จาติ.}
\csromangathagroup{\csromangathastanza{เมตฺตํ ปญฺญา จ เทฺว ปิยา, เทฺว โลกา เทฺว วิปตฺติโย. \\ เทวทตฺโต จ อุตฺตโร, นนฺโท การณฺฑเวน จาติ.}}

\chapterhead{2. มหาวคฺค}
\csromanheader{2}{1. เวรญฺชสุตฺต}
\proseitem{11}{เอวํ เม สุตํ\csromandash{} * เอกํ สมยํ ภควา เวรญฺชายํ วิหรติ นเฬรุปุจิมนฺทมูเล.\csromanspacer{} อถ โข เวรญฺโช พฺราหฺมโณ เยน ภควา เตนุปสงฺกมิ, อุปสงฺกมิตฺวา ภควตา สทฺธิํ สมฺโมทิ, สมฺโมทนียํ กถํ สารณียํ\footnote{สาราณียํ (สี, สฺยา, กํ, อิ)}}

\vspace{\tipitakaparskip}
\csromancenter{มหาวคฺค วีติสาเรตฺวา เอกมนฺตํ นิสีทิ, เอกมนฺตํ นิสินฺโน โข เวรญฺโช พฺราหฺมโณ ภควนฺตํ เอตทโวจ\csromandash{}}
\prose{\csromanfolio{19}สุตํ เมตํ โภ โคตม “น สมโณ โคตโม พฺราหฺมเณ ชิณฺเณ วุฑฺเฒ มหลฺลเก อทฺธคเต วโย-อนุปฺปตฺเต อภิวาเทติ วา ปจฺจุฏฺเฐติ วา อาสเนน วา นิมนฺเตตี”ติ.\csromanspacer{} ตยิทํ โภ โคตม ตเถว, น หิ ภวํ โคตโม พฺราหฺมเณ ชิณฺเณ วุฑฺเฒ มหลฺลเก อทฺธคเต วโย-อนุปฺปตฺเต อภิวาเทติ วา ปจฺจุฏฺเฐติ วา อาสเนน วา นิมนฺเตติ, ตยิทํ โภ โคตม น สมฺปนฺนเมวาติ.\csromanspacer{} นาหํ ตํ พฺราหฺมณ ปสฺสามิ สเทวเก โลเก สมารเก สพฺรหฺมเก สสฺสมณพฺราหฺมณิยา ปชาย สเทวมนุสฺสาย, ยมหํ อภิวาเทยฺยํ วา ปจฺจุฏฺเฐยฺยํ วา อาสเนน วา นิมนฺเตยฺยํ.\csromanspacer{} ยํ หิ พฺราหฺมณ ตถาคโต อภิวาเทยฺย วา ปจฺจุฏฺเฐยฺย วา อาสเนน วา นิมนฺเตยฺย, มุทฺธาปิ ตสฺส วิปเตยฺยาติ.}

\prose{อรสรูโป ภวํ โคตโมติ.\csromanspacer{} อตฺถิ เขฺวส พฺราหฺมณ ปริยาโย, เยน มํ ปริยาเยน สมฺมา วทมาโน วเทยฺย “อรสรูโป สมโณ โคตโม”ติ.\csromanspacer{} เย เต พฺราหฺมณ รูปรสา สทฺทรสา คนฺธรสา รสรสา โผฏฺฐพฺพรสา, เต ตถาคตสฺส ปหีนา อุจฺฉินฺนมูลา ตาลาวตฺถุกตา อนภาวํกตา\footnote{อนภาวกตา (สี, อิ)} อายติํ อนุปฺปาทธมฺมา.\csromanspacer{} อยํ โข พฺราหฺมณ ปริยาโย, เยน มํ ปริยาเยน สมฺมา วทมาโน วเทยฺย “อรสรูโป สมโณ โคตโม”ติ.\csromanspacer{} โน จ โข ยํ ตฺวํ สนฺธาย วเทสีติ\footnote{วเทสิ (สี, ก)}. (1)}

\prose{นิพฺโภโค ภวํ โคตโมติ.\csromanspacer{} อตฺถิ เขฺวส พฺราหฺมณ ปริยาโย, เยน มํ ปริยาเยน สมฺมา วทมาโน วเทยฺย “นิพฺโภโค สมโณ โคตโม”ติ.\csromanspacer{} เย เต พฺราหฺมณ รูปโภคา สทฺทโภคา คนฺธโภคา รสโภคา โผฏฺฐพฺพโภคา, เต ตถาคตสฺส ปหีนา อุจฺฉินฺนมูลา ตาลาวตฺถุกตา อนภาวํกตา อายติํ อนุปฺปาทธมฺมา.\csromanspacer{} อยํ โข พฺราหฺมณ ปริยาโย, เยน มํ ปริยาเยน สมฺมา วทมาโน วเทยฺย “นิพฺโภโค สมโณ โคตโม”ติ.\csromanspacer{} โน จ โข ยํ ตฺวํ สนฺธาย วเทสีติ. (2)}

\prose{\csromanfolio{20}อกิริยวาโท ภวํ โคตโมติ.\csromanspacer{} อตฺถิ เขฺวส พฺราหฺมณ ปริยาโย, เยน มํ ปริยาเยน สมฺมา วทมาโน วเทยฺย “อกิริยวาโท สมโณ โคตโม”ติ.\csromanspacer{} อหํ หิ พฺราหฺมณ อกิริยํ วทามิ กายทุจฺจริตสฺส วจีทุจฺจริตสฺส มโนทุจฺจริตสฺส, อเนกวิหิตานํ ปาปกานํ อกุสลานํ ธมฺมานํ อกิริยํ วทามิ.\csromanspacer{} อยํ โข พฺราหฺมณ ปริยาโย, เยน มํ ปริยาเยน สมฺมา วทมาโน วเทยฺย “อกิริยวาโท สมโณ โคตโม”ติ.\csromanspacer{} โน จ โข ยํ ตฺวํ สนฺธาย วเทสีติ. (3)}

\prose{อุจฺเฉทวาโท ภวํ โคตโมติ.\csromanspacer{} อตฺถิ เขฺวส พฺราหฺมณ ปริยาโย, เยน มํ ปริยาเยน สมฺมา วทมาโน วเทยฺย “อุจฺเฉทวาโท สมโณ โคตโม”ติ.\csromanspacer{} อหํ หิ พฺราหฺมณ อุจฺเฉทํ วทามิ ราคสฺส โทสสฺส โมหสฺส, อเนกวิหิตานํ ปาปกานํ อกุสลานํ ธมฺมานํ อุจฺเฉทํ วทามิ.\csromanspacer{} อยํ โข พฺราหฺมณ ปริยาโย, เยน มํ ปริยาเยน สมฺมา วทมาโน วเทยฺย “อุจฺเฉทวาโท สมโณ โคตโม”ติ.\csromanspacer{} โน จ โข ยํ ตฺวํ สนฺธาย วเทสีติ. (4)}

\prose{เชคุจฺฉี ภวํ โคตโมติ.\csromanspacer{} อตฺถิ เขฺวส พฺราหฺมณ ปริยาโย, เยน มํ ปริยาเยน สมฺมา วทมาโน วเทยฺย “เชคุจฺฉี สมโณ โคตโม”ติ.\csromanspacer{} อหํ หิ พฺราหฺมณ ชิคุจฺฉามิ กายทุจฺจริเตน วจีทุจฺจริเตน มโนทุจฺจริเตน, ชิคุจฺฉามิ อเนกวิหิตานํ ปาปกานํ อกุสลานํ ธมฺมานํ สมาปตฺติยา.\csromanspacer{} อยํ โข พฺราหฺมณ ปริยาโย, เยน มํ ปริยาเยน สมฺมา วทมาโน วเทยฺย “เชคุจฺฉี สมโณ โคตโม”ติ.\csromanspacer{} โน จ โข ยํ ตฺวํ สนฺธาย วเทสีติ. (5)}

\prose{เวนยิโก ภวํ โคตโมติ.\csromanspacer{} อตฺถิ เขฺวส พฺราหฺมณ ปริยาโย, เยน มํ ปริยาเยน สมฺมา วทมาโน วเทยฺย “เวนยิโก สมโณ โคตโม”ติ.\csromanspacer{} อหํ หิ พฺราหฺมณ วินยาย ธมฺมํ เทเสมิ ราคสฺส โทสสฺส โมหสฺส, อเนกวิหิตานํ ปาปกานํ อกุสลานํ ธมฺมานํ วินยาย ธมฺมํ เทเสมิ.\csromanspacer{} อยํ โข พฺราหฺมณ ปริยาโย, เยน มํ ปริยาเยน สมฺมา วทมาโน วเทยฺย “เวนยิโก สมโณ โคตโม”ติ.\csromanspacer{} โน จ โข ยํ ตฺวํ สนฺธาย วเทสีติ. (6)}

\prose{ตปสฺสี ภวํ โคตโมติ.\csromanspacer{} อตฺถิ เขฺวส พฺราหฺมณ ปริยาโย, เยน มํ ปริยาเยน สมฺมา วทมาโน วเทยฺย “ตปสฺสี สมโณ โคตโม”ติ.}

\vspace{\tipitakaparskip}
\csromancenter{มหาวคฺค ตปนียาหํ พฺราหฺมณ ปาปเก อกุสเล ธมฺเม วทามิ กายทุจฺจริตํ วจีทุจฺจริตํ มโนทุจฺจริตํ, ยสฺส โข พฺราหฺมณ ตปนียา ปาปกา อกุสลา ธมฺมา ปหีนา อุจฺฉินฺนมูลา ตาลาวตฺถุกตา อนภาวํกตา อายติํ อนุปฺปาทธมฺมา, ตมหํ ‘ตปสฺสี’ติ วทามิ.\csromanspacer{} ตถาคตสฺส โข พฺราหฺมณ ตปนียา ปาปกา อกุสลา ธมฺมา ปหีนา อุจฺฉินฺนมูลา ตาลาวตฺถุกตา อนภาวํกตา อายติํ อนุปฺปาทธมฺมา.\csromanspacer{} อยํ โข พฺราหฺมณ ปริยาโย, เยน มํ ปริยาเยน สมฺมา วทมาโน วเทยฺย “ตปสฺสี สมโณ โคตโม”ติ.\csromanspacer{} โน จ โข ยํ ตฺวํ สนฺธาย วเทสีติ. (7)}
\prose{\csromanfolio{21}อปคพฺโภ ภวํ โคตโมติ.\csromanspacer{} อตฺถิ เขฺวส พฺราหฺมณ ปริยาโย, เยน มํ ปริยาเยน สมฺมา วทมาโน วเทยฺย “อปคพฺโภ สมโณ โคตโม”ติ.\csromanspacer{} ยสฺส โข พฺราหฺมณ อายติํ คพฺภเสยฺยา ปุนพฺภวาภินิพฺพตฺติ ปหีนา อุจฺฉินฺนมูลา ตาลาวตฺถุกตา อนภาวํกตา อายติํ อนุปฺปาทธมฺมา, ตมหํ ‘อปคพฺโภ’ติ วทามิ.\csromanspacer{} ตถาคตสฺส โข พฺราหฺมณ อายติํ คพฺภเสยฺยา ปุนพฺภวาภินิพฺพตฺติ ปหีนา อุจฺฉินฺนมูลา ตาลาวตฺถุกตา อนภาวํกตา อายติํ อนุปฺปาทธมฺมา.\csromanspacer{} อยํ โข พฺราหฺมณ ปริยาโย, เยน มํ ปริยาเยน สมฺมา วทมาโน วเทยฺย “อปคพฺโภ สมโณ โคตโม”ติ.\csromanspacer{} โน จ โข ยํ ตฺวํ สนฺธาย วเทสิ. (8)}

\prose{เสยฺยถาปิ พฺราหฺมณ กุกฺกุฏิยา อณฺฑานิ อฏฺฐ วา ทส วา ทฺวาทส วา, ตานาสฺสุ กุกฺกุฏิยา สมฺมา อธิสยิตานิ สมฺมา ปริเสทิตานิ สมฺมา ปริภาวิตานิ.\csromanspacer{} โย นุ โข เตสํ กุกฺกุฏจฺฉาปกานํ ปฐมตรํ ปาทนขสิขาย วา มุขตุณฺฑเกน วา อณฺฑโกสํ ปทาเลตฺวา โสตฺถินา อภินิพฺภิชฺเชยฺย, กินฺติ สฺวาสฺส วจนีโย “เชฏฺโฐ วา กนิฏฺโฐ วา”ติ. “เชฏฺโฐ”ติสฺส โภ โคตม วจนีโย.\csromanspacer{} โส หิ เนสํ โภ โคตม เชฏฺโฐ โหตีติ.}

\prose{เอวเมวํ โข อหํ พฺราหฺมณ อวิชฺชาคตาย ปชาย อณฺฑภูตาย ปริโยนทฺธาย อวิชฺชณฺฑโกสํ ปทาเลตฺวา เอโกว โลเก อนุตฺตรํ สมฺมาสมฺโพธิํ อภิสมฺพุทฺโธ, อหํ หิ พฺราหฺมณ เชฏฺโฐ เสฏฺโฐ โลกสฺส, อารทฺธํ โข ปน เม พฺราหฺมณ วีริยํ อโหสิ อสลฺลีนํ, อุปฏฺฐิตา สติ อสมฺมุฏฺฐา, ปสฺสทฺโธ กาโย อสารทฺโธ, สมาหิตํ จิตฺตํ เอกคฺคํ.}

\prose{\csromanfolio{22}โส โข อหํ พฺราหฺมณ วิวิจฺเจว กาเมหิ วิวิจฺจ อกุสเลหิ ธมฺเมหิ สวิตกฺกํ สวิจารํ วิเวกชํ ปีติสุขํ ปฐมํ ฌานํ อุปสมฺปชฺช วิหรามิ.\csromanspacer{} วิตกฺกวิจารานํ วูปสมา อชฺฌตฺตํ สมฺปสาทนํ เจตโส เอโกทิภาวํ อวิตกฺกํ อวิจารํ สมาธิชํ ปีติสุขํ ทุติยํ ฌานํ อุปสมฺปชฺช วิหรามิ.\csromanspacer{} ปีติยา จ วิราคา อุเปกฺขโก จ วิหรามิ, สโต จ สมฺปชาโน สุขญฺจ กาเยน ปฏิสํเวเทมิ, ยํ ตํ อริยา อาจิกฺขนฺติ “อุเปกฺขโก สติมา สุขวิหารี”ติ, ตติยํ ฌานํ อุปสมฺปชฺช วิหรามิ.\csromanspacer{} สุขสฺส จ ปหานา ทุกฺขสฺส จ ปหานา ปุพฺเพว โสมนสฺสโทมนสฺสานํ อตฺถงฺคมา อทุกฺขมสุขํ อุเปกฺขาสติปาริสุทฺธิํ จตุตฺถํ ฌานํ อุปสมฺปชฺช วิหรามิ.}

\prose{โส เอวํ สมาหิเต จิตฺเต ปริสุทฺเธ ปริโยทาเต อนงฺคเณ วิคตูปกฺกิเลเส มุทุภูเต กมฺมนิเย ฐิเต อาเนญฺชปฺปตฺเต ปุพฺเพนิวาสานุสฺสติญาณาย จิตฺตํ อภินินฺนาเมสิํ.\csromanspacer{} โส อเนกวิหิตํ ปุพฺเพนิวาสํ อนุสฺสรามิ.\csromanspacer{} เสยฺยถิทํ, เอกมฺปิ ชาติํ เทฺวปิ ชาติโย ติสฺโสปิ ชาติโย จตสฺโสปิ ชาติโย ปญฺจปิ ชาติโย ทสปิ ชาติโย วีสมฺปิ ชาติโย ติํสมฺปิ ชาติโย จตฺตาลีสมฺปิ ชาติโย ปญฺญาสมฺปิ ชาติโย ชาติสตมฺปิ ชาติสหสฺสมฺปิ ชาติสตสหสฺสมฺปิ อเนเกปิ สํวฏฺฏกปฺเป อเนเกปิ วิวฏฺฏกปฺเป อเนเกปิ สํวฏฺฏวิวฏฺฏกปฺเป, อมุตฺราสิํ เอวํนาโม เอวํโคตฺโต เอวํวณฺโณ เอวมาหาโร เอวํสุขทุกฺขปฺปฏิสํเวที เอวมายุปริยนฺโต.\csromanspacer{} โส ตโต จุโต อมุตฺร อุทปาทิํ, ตตฺราปาสิํ “เอวํนาโม เอวํโคตฺโต เอวํวณฺโณ เอวมาหาโร เอวํสุขทุกฺขปฺปฏิสํเวที เอวมายุปริยนฺโต.\csromanspacer{} โส ตโต จุโต อิธูปปนฺโน”ติ.\csromanspacer{} อิติ สาการํ ส-อุทฺเทสํ อเนกวิหิตํ ปุพฺเพนิวาสํ อนุสฺสรามิ.}

\prose{อยํ โข เม พฺราหฺมณ รตฺติยา ปฐเม ยาเม ปฐมา วิชฺชา อธิคตา, อวิชฺชา วิหตา, วิชฺชา อุปฺปนฺนา, ตโม วิหโต, อาโลโก อุปฺปนฺโน, ยถา ตํ อปฺปมตฺตสฺส อาตาปิโน ปหิตตฺตสฺส วิหรโต.\csromanspacer{} อยํ โข เม พฺราหฺมณ ปฐมา อภินิพฺภิทา อโหสิ กุกฺกุฏจฺฉาปกสฺเสว อณฺฑโกสมฺหา. (1)}

\prose{โส เอวํ สมาหิเต จิตฺเต ปริสุทฺเธ ปริโยทาเต อนงฺคเณ วิคตูปกฺกิเลเส มุทุภูเต กมฺมนิเย ฐิเต อาเนญฺชปฺปตฺเต สตฺตานํ จุตูปปาตญาณาย จิตฺตํ อภินินฺนาเมสิํ.\csromanspacer{} โส ทิพฺเพน จกฺขุนา วิสุทฺเธน อติกฺกนฺตมานุสเกน สตฺเต ปสฺสามิ จโมเน อุปปชฺชมาเน หีเน}

\vspace{\tipitakaparskip}
\csromancenter{มหาวคฺค ปณีเต สุวณฺเณ ทุพฺพณฺเณ สุคเต ทุคฺคเต, ยถากมฺมูปเค สตฺเต ปชานามิ “อิเม วต โภนฺโต สตฺตา กายทุจฺจริเตน สมนฺนาคตา วจีทุจฺจริเตน สมนฺนาคตา มโนทุจฺจริเตน สมนฺนาคตา, อริยานํ อุปวาทกา มิจฺฉาทิฏฺฐิกา มิจฺฉาทิฏฺฐิกมฺมสมาทานา, เต กายสฺส เภทา ปรํ มรณา อปายํ ทุคฺคติํ วินิปาตํ นิรยํ อุปปนฺนา”ติ. “อิเม วา ปน โภนฺโต สตฺตา กายสุจริเตน สมนฺนาคตา วจีสุจริเตน สมนฺนาคตา มโนสุจริเตน สมนฺนาคตา, อริยานํ อนุปวาทกา.\csromanspacer{} สมฺมาทิฏฺฐิกา สมฺมาทิฏฺฐิกมฺมสมาทานา.\csromanspacer{} เต กายสฺส เภทา ปรํ มรณา สุคติํ สคฺคํ โลกํ อุปปนฺนา”ติ.\csromanspacer{} อิติ ทิพฺเพน จกฺขุนา วิสุทฺเธน อติกฺกนฺตมานุสเกน สตฺเต ปสฺสามิ จวมาเน อุปปชฺชมาเน หีเน ปณีเต สุวณฺเณ ทุพฺพณฺเณ สุคเต ทุคฺคเต, ยถากมฺมูปเค สตฺเต ปชานามิ.}
\prose{\csromanfolio{23}อยํ โข เม พฺราหฺมณ รตฺติยา มชฺฌิเม ยาเม ทุติยา วิชฺชา อธิคตา, อวิชฺชา วิหตา, วิชฺชา อุปฺปนฺนา, ตโม วิหโต, อาโลโก อุปฺปนฺโน, ยถา ตํ อปฺปมตฺตสฺส อาตาปิโน ปหิตตฺตสฺส วิหรโต.\csromanspacer{} อยํ โข เม พฺราหฺมณ ทุติยา อภินิพฺภิทา อโหสิ กุกฺกุฏจฺฉาปกสฺเสว อณฺฑโกสมฺหา. (2)}

\prose{โส เอวํ สมาหิเต จิตฺเต ปริสุทฺเธ ปริโยทาเต อนงฺคเณ วิคตูปกฺกิเลเส มุทุภูเต กมฺมนิเย ฐิเต อาเนญฺชปฺปตฺเต อาสวานํ ขยญาณาย จิตฺตํ อภินินฺนาเมสิํ.\csromanspacer{} โส “อิทํ ทุกฺขนฺ”ติ ยถาภูตํ อพฺภญฺญาสิํ, “อยํ ทุกฺขสมุทโย”ติ ยถาภูตํ อพฺภญฺญาสิํ, “อยํ ทุกฺขนิโรโธ”ติ ยถาภูตํ อพฺภญฺญาสิํ, “อยํ ทุกฺขนิโรธคามินี ปฏิปทา”ติ ยถาภูตํ อพฺภญฺญาสิํ, “อิเม อาสวา”ติ ยถาภูตํ อพฺภญฺญาสิํ, “อยํ อาสวสมุทโย”ติ ยถาภูตํ อพฺภญฺญาสิํ, “อยํ อาสวนิโรโธ”ติ ยถาภูตํ อพฺภญฺญาสิํ, “อยํ อาสวนิโรธคามินี ปฏิปทา”ติ ยถาภูตํ อพฺภญฺญาสิํ.\csromanspacer{} ตสฺส เม เอวํ ชานโต เอวํ ปสฺสโต กามาสวาปิ จิตฺตํ วิมุจฺจิตฺถ, ภวาสวาปิ จิตฺตํ วิมุจฺจิตฺถ, อวิชฺชาสวาปิ จิตฺตํ วิมุจฺจิตฺถ, วิมุตฺตสฺมิํ “วิมุตฺตมฺ”อิติ ญาณํ อโหสิ, “ขีณา ชาติ, วุสิตํ พฺรหฺมจริยํ, กตํ กรณียํ, นาปรํ อิตฺถตฺตายา”ติ อพฺภญฺญาสิํ.}

\prose{อยํ โข เม พฺราหฺมณ รตฺติยา ปจฺฉิเม ยาเม ตติยา วิชฺชา อธิคตา, อวิชฺชา วิหตา, วิชฺชา อุปฺปนฺนา, ตโม วิหโต, อาโลโก \csromanfolio{24}อุปฺปนฺโน, ยถา ตํ อปฺปมตฺตสฺส อาตาปิโน ปหิตตฺตสฺส วิหรโต.\csromanspacer{} อยํ โข เม พฺราหฺมณ ตติยา อภินิพฺภิทา อโหสิ กุกฺกุฏจฺฉาปกสฺเสว อณฺฑโกสมฺหาติ. (3)}

\prose{เอวํ วุตฺเต เวรญฺโช พฺราหฺมโณ ภควนฺตํ เอตทโวจ “เชฏฺโฐ ภวํ โคตโม, เสฏฺโฐ ภวํ โคตโม.\csromanspacer{} อภิกฺกนฺตํ โภ โคตม, อภิกฺกนฺตํ โภ โคตม, เสยฺยถาปิ โภ โคตม นิกฺกุชฺชิตํ\footnote{นิกุชฺชิตํ (ก)} วา อุกฺกุชฺเชยฺย, ปฏิจฺฉนฺนํ วา วิวเรยฺย, มูฬฺหสฺส วา มคฺคํ อาจิกฺเขยฺย, อนฺธกาเร วา เตลปชฺโชตํ ธาเรยฺย ‘จกฺขุมนฺโต รูปานิ ทกฺขนฺตี’ติ.\csromanspacer{} เอวเมวํ โภตา โคตเมน อเนกปริยาเยน ธมฺโม ปกาสิโต, เอสาหํ ภวนฺตํ โคตมํ สรณํ คจฺฉามิ ธมฺมญฺจ ภิกฺขุสํฆญฺจ, อุปาสกํ มํ ภวํ โคตโม ธาเรตุ อชฺชตคฺเค ปาณุเปตํ สรณํ คตนฺ”ติ.\csromanspacer{}\csromanspacer{}ปฐมํ.}

\csromanheader{2}{2. สีหสุตฺต}
\proseitem{12}{เอกํ สมยํ ภควา เวสาลิยํ วิหรติ มหาวเน กูฏาคารสาลายํ.\csromanspacer{} เตน โข ปน สมเยน สมฺพหุลา อภิญฺญาตา อภิญฺญาตา ลิจฺฉวี สนฺถาคาเร\footnote{สนฺธาคาเร (ก)} สนฺนิสินฺนา สนฺนิปติตา อเนกปริยาเยน พุทฺธสฺส วณฺณํ ภาสนฺติ, ธมฺมสฺส วณฺณํ ภาสนฺติ, สํฆสฺส วณฺณํ ภาสนฺติ.}

\prose{เตน โข ปน สมเยน สีโห เสนาปติ นิคณฺฐสาวโก ตสฺสํ ปริสายํ นิสินฺโน โหติ, อถ โข สีหสฺส เสนาปติสฺส เอตทโหสิ “นิสฺสํสยํ โข โส ภควา อรหํ สมฺมาสมฺพุทฺโธ ภวิสฺสติ, ตถา หิเม สมฺพหุลา อภิญฺญาตา อภิญฺญาตา ลิจฺฉวี สนฺถาคาเร สนฺนิสินฺนา สนฺนิปติตา อเนกปริยาเยน พุทฺธสฺส วณฺณํ ภาสนฺติ, ธมฺมสฺส วณฺณํ ภาสนฺติ, สํฆสฺส วณฺณํ ภาสนฺติ, ยํนูนาหํ ตํ ภควนฺตํ ทสฺสนาย อุปสงฺกเมยฺยํ อรหนฺตํ สมฺมาสมฺพุทฺธนฺ”ติ.\csromanspacer{} อถ โข สีโห เสนาปติ เยน นิคณฺโฐ นาฏปุตฺโต\footnote{นาถปุตฺโต (ก-สี), นาตปุตฺโต (ก-สี)} เตนุปสงฺกมิ, อุปสงฺกมิตฺวา นิคณฺฐํ นาฏปุตฺตํ เอตทโวจ “อิจฺฉามหํ ภนฺเต สมณํ โคตมํ ทสฺสนาย อุปสงฺกมิตุนฺ”ติ.}

\chapterheadpageread{2. มหาวคฺค}
\prose{\csromanfolio{25}กิํ ปน ตฺวํ สีห กิริยวาโท สมาโน อกิริยวาทํ สมณํ โคตมํ ทสฺสนาย อุปสงฺกมิสฺสสิ, สมโณ หิ สีห โคตโม อกิริยวาโท, อกิริยาย ธมฺมํ เทเสติ, เตน จ สาวเก วิเนตีติ.\csromanspacer{} อถ โข สีหสฺส เสนาปติสฺส โย อโหสิ คมิยาภิสงฺขาโร\footnote{คมิกาภิสงฺขาโร (ก-สี) วิ 3. 328-9 ปิฏฺเฐปิ.} ภควนฺตํ ทสฺสนาย, โส ปฏิปฺปสฺสมฺภิ.}

\prose{ทุติยมฺปิ โข สมฺพหุลา อภิญฺญาตา อภิญฺญาตา ลิจฺฉวี สนฺถาคาเร สนฺนิสินฺนา สนฺนิปติตา อเนกปริยาเยน พุทฺธสฺส ฯเปฯ ธมฺมสฺส ฯเปฯ สํฆสฺส วณฺณํ ภาสนฺติ.\csromanspacer{} ทุติยมฺปิ โข สีหสฺส เสนาปติสฺส เอตทโหสิ “นิสฺสํสยํ โข โส ภควา อรหํ สมฺมาสมฺพุทฺโธ ภวิสฺสติ, ตถา หิเม สมฺพหุลา อภิญฺญาตา อภิญฺญาตา ลิจฺฉวี สนฺถาคาเร สนฺนิสินฺนา สนฺนิปติตา อเนกปริยาเยน พุทฺธสฺส วณฺณํ ภาสนฺติ, ธมฺมสฺส ฯเปฯ สํฆสฺส วณฺณํ ภาสนฺติ, ยํนูนาหํ ตํ ภควนฺตํ ทสฺสนาย อุปสงฺกเมยฺยํ อรหนฺตํ สมฺมาสมฺพุทฺธนฺ”ติ.\csromanspacer{} อถ โข สีโห เสนาปติ เยน นิคณฺโฐ นาฏปุตฺโต เตนุปสงฺกมิ, อุปสงฺกมิตฺวา นิคณฺฐํ นาฏปุตฺตํ เอตทโวจ “อิจฺฉามหํ ภนฺเต สมณํ โคตมํ ทสฺสนาย อุปสงฺกมิตุนฺ”ติ.}

\prose{กิํ ปน ตฺวํ สีห กิริยวาโท สมาโน อกิริยวาทํ สมณํ โคตมํ ทสฺสนาย อุปสงฺกมิสฺสสิ, สมโณ หิ สีห โคตโม อกิริยวาโท, อกิริยาย ธมฺมํ เทเสติ, เตน จ สาวเก วิเนตี”ติ.\csromanspacer{} ทุติยมฺปิ โข สีหสฺส เสนาปติสฺส โย อโหสิ คมิยาภิสงฺขาโร ภควนฺตํ ทสฺสนาย, โส ปฏิปฺปสฺสมฺภิ.}

\prose{ตติยมฺปิ โข สมฺพหุลา อภิญฺญาตา อภิญฺญาตา ลิจฺฉวี สนฺถาคาเร สนฺนิสินฺนา สนฺนิปติตา อเนกปริยาเยน พุทฺธสฺส ฯเปฯ ธมฺมสฺส ฯเปฯ สํฆสฺส วณฺณํ ภาสนฺติ.\csromanspacer{} ตติยมฺปิ โข สีหสฺส เสนาปติสฺส เอตทโหสิ “นิสฺสํสยํ โข โส ภควา อรหํ สมฺมาสมฺพุทฺโธ ภวิสฺสติ, ตถา หิเม สมฺพหุลา อภิญฺญาตา อภิญฺญาตา ลิจฺฉวี สนฺถาคาเร สนฺนิสินฺนา สนฺนิปติตา อเนกปริยาเยน พุทฺธสฺส วณฺณํ ภาสนฺติ, ธมฺมสฺส วณฺณํ ภาสนฺติ, สํฆสฺส วณฺณํ ภาสนฺติ.\csromanspacer{} กิํ หิ เม กริสฺสนฺติ}

\prose{\csromanfolio{26}นิคณฺฐา อปโลกิตา วา อนปโลกิตา วา, ยํนูนาหํ อนปโลเกตฺวาว นิคณฺเฐ\footnote{นิคณฺฐํ (สฺยา, ก) วิ 3. 329 ปิฏฺเฐ ปสฺสิตพฺพํ.} ตํ ภควนฺตํ ทสฺสนาย อุปสงฺกเมยฺยํ อรหนฺตํ สมฺมาสมฺพุทฺธนฺ”ติ.}

\prose{อถ โข สีโห เสนาปติ ปญฺจมตฺเตหิ รถสเตหิ ทิวา ทิวสฺส เวสาลิยา นิยฺยาสิ ภควนฺตํ ทสฺสนาย.\csromanspacer{} ยาวติกา ยานสฺส ภูมิ, ยาเนน คนฺตฺวา ยานา ปจฺโจโรหิตฺวา ปตฺติโกว อคมาสิ.\csromanspacer{} อถ โข สีโห เสนาปติ เยน ภควา เตนุปสงฺกมิ, อุปสงฺกมิตฺวา ภควนฺตํ อภิวาเทตฺวา เอกมนฺตํ นิสีทิ, เอกมนฺตํ นิสินฺโน โข สีโห เสนาปติ ภควนฺตํ เอตทโวจ\csromandash{}}

\prose{สุตํ เมตํ ภนฺเต “อกิริยวาโท สมโณ โคตโม, อกิริยาย ธมฺมํ เทเสติ, เตน จ สาวเก วิเนตี”ติ.\csromanspacer{} เย เต ภนฺเต เอวมาหํสุ “อกิริยวาโท สมโณ โคตโม, อกิริยาย ธมฺมํ เทเสติ, เตน จ สาวเก วิเนตี”ติ.\csromanspacer{} กจฺจิ เต ภนฺเต ภควโต วุตฺตวาทิโน, น จ ภควนฺตํ อภูเตน อพฺภาจิกฺขนฺติ, ธมฺมสฺส จานุธมฺมํ พฺยากโรนฺติ, น จ โกจิ สหธมฺมิโก วาทานุวาโท\footnote{วาทานุปาโต (ก, สี, สฺยา) อํ 1. 161; อํ 2. 3 ปิฏฺเฐสุปิ.} คารยฺหํ ฐานํ อาคจฺฉติ.\csromanspacer{} อนพฺภกฺขาตุกามา หิ มยํ ภนฺเต ภควนฺตนฺติ.}

\prose{อตฺถิ สีห ปริยาโย, เยน มํ ปริยาเยน สมฺมา วทมาโน วเทยฺย “อกิริยวาโท สมโณ โคตโม, อกิริยาย ธมฺมํ เทเสติ, เตน จ สาวเก วิเนตี”ติ. (1)}

\prose{อตฺถิ สีห ปริยาโย, เยน มํ ปริยาเยน สมฺมา วทมาโน วเทยฺย “กิริยวาโท สมโณ โคตโม, กิริยาย ธมฺมํ เทเสติ, เตน จ สาวเก วิเนตี”ติ. (2)}

\prose{อตฺถิ สีห ปริยาโย, เยน มํ ปริยาเยน สมฺมา วทมาโน วเทยฺย “อุจฺเฉทวาโท สมโณ โคตโม, อุจฺเฉทาย ธมฺมํ เทเสติ, เตน จ สาวเก วิเนตี”ติ. (3)}

\prose{อตฺถิ สีห ปริยาโย, เยน มํ ปริยาเยน สมฺมา วทมาโน วเทยฺย “เชคุจฺฉี สมโณ โคตโม, เชคุจฺฉิตาย ธมฺมํ เทเสติ, เตน จ สาวเก วิเนตี”ติ. (4)}

\chapterheadpageread{2. มหาวคฺค}
\prose{\csromanfolio{27}อตฺถิ สีห ปริยาโย, เยน มํ ปริยาเยน สมฺมา วทมาโน วเทยฺย “เวนยิโก สมโณ โคตโม,วินยาย ธมฺมํ เทเสติ, เตน จ สาวเก วิเนตี”ติ. (5)}

\prose{อตฺถิ สีห ปริยาโย, เยน มํ ปริยาเยน สมฺมา วทมาโน วเทยฺย “ตปสฺสี สมโณ โคตโม, ตปสฺสิตาย ธมฺมํ เทเสติ, เตน จ สาวเก วิเนตี”ติ. (6)}

\prose{อตฺถิ สีห ปริยาโย, เยน มํ ปริยาเยน สมฺมา วทมาโน วเทยฺย “อปคพฺโภ สมโณ, อปคพฺภตาย ธมฺมํ เทเสติ, เตน จ สาวเก วิเนตี”ติ. (7)}

\prose{อตฺถิ สีห ปริยาโย, เยน มํ ปริยาเยน สมฺมา วทมาโน วเทยฺย “อสฺสาสโก สมโณ โคตโม, อสฺสาสาย ธมฺมํ เทเสติ, เตน จ สาวเก วิเนตี”ติ. (8)}

\prose{กตโม จ สีห ปริยาโย, เยน มํ ปริยาเยน สมฺมา วทมาโน วเทยฺย “อกิริยวาโท สมโณ โคตโม, อกิริยาย ธมฺมํ เทเสติ, เตน จ สาวเก วิเนตี”ติ.\csromanspacer{} อหํ หิ สีห อกิริยํ วทามิ กายทุจฺจริตสฺส วจีทุจฺจริตสฺส มโนทุจฺจริตสฺส, อเนกวิหิตานํ ปาปกานํ อกุสลานํ ธมฺมานํ อกิริยํ วทามิ.\csromanspacer{} อยํ โข สีห ปริยาโย, เยน มํ ปริยาเยน สมฺมา วทมาโน วเทยฺย “อกิริยวาโท สมโณ โคตโม, อกิริยาย ธมฺมํ เทเสติ, เตน จ สาวเก วิเนตี”ติ. (1)}

\prose{กตโม จ สีห ปริยาโย, เยน มํ ปริยาเยน สมฺมา วทมาโน วเทยฺย “กิริยวาโท สมโณ โคตโม, กิริยาย ธมฺมํ เทเสติ, เตน จ สาวเก วิเนตี”ติ.\csromanspacer{} อหํ หิ สีห กิริยํ วทามิ กายสุจริตสฺส วจีสุจริตสฺส มโนสุจริตสฺส, อเนกวิหิตานํ กุสลานํ ธมฺมานํ กิริยํ วทามิ.\csromanspacer{} อยํ โข สีห ปริยาโย, เยน มํ ปริยาเยน สมฺมา วทมาโน วเทยฺย “กิริยวาโท สมโณ โคตโม, กิริยาย ธมฺมํ เทเสติ, เตน จ สาวเก วิเนตี”ติ. (2)}

\prose{กตโม จ สีห ปริยาโย, เยน มํ ปริยาเยน สมฺมา วทมาโน วเทยฺย “อุจฺเฉทวาโท สมโณ โคตโม, อุจฺเฉทาย ธมฺมํ เทเสติ, เตน จ สาวเก วิเนตี”ติ.\csromanspacer{} อหํ หิ สีห อุจฺเฉทํ วทามิ ราคสฺส}

\prose{\csromanfolio{28}โทสสฺส โมหสฺส, อเนกวิหิตานํ ปาปกานํ อกุสลานํ ธมฺมานํ อุจฺเฉทํ วทามิ.\csromanspacer{} อยํ โข สีห ปริยาโย, เยน มํ ปริยาเยน สมฺมา วทมาโน วเทยฺย “อุจฺเฉทวาโท สมโณ โคตโม, อุจฺเฉทาย ธมฺมํ เทเสติ, เตน จ สาวเก วิเนตี”ติ. (3)}

\prose{กตโม จ สีห ปริยาโย, เยน มํ ปริยาเยน สมฺมา วทมาโน วเทยฺย “เชคุจฺฉี สมโณ โคตโม, เชคุจฺฉิตาย ธมฺมํ เทเสติ, เตน จ สาวเก วิเนตี”ติ.\csromanspacer{} อหํ หิ สีห ชิคุจฺฉามิ กายทุจฺจริเตน วจีทุจฺจริเตน มโนทุจฺจริเตน, ชิคุจฺฉามิ อเนกวิหิตานํ ปาปกานํ อกุสลานํ ธมฺมานํ สมาปตฺติยา.\csromanspacer{} อยํ โข สีห ปริยาโย, เยน มํ ปริยาเยน สมฺมา วทมาโน วเทยฺย “เชคุจฺฉี สมโณ โคตโม, เชคุจฺฉิตาย ธมฺมํ เทเสติ, เตน จ สาวเก วิเนตี”ติ. (4)}

\prose{กตโม จ สีห ปริยาโย, เยน มํ ปริยาเยน สมฺมา วทมาโน วเทยฺย “เวนยิโก สมโณ โคตโม, วินยาย ธมฺมํ เทเสติ, เตน จ สาวเก วิเนตี”ติ.\csromanspacer{} อหํ หิ สีห วินยาย ธมฺมํ เทเสมิ ราคสฺส โทสสฺส โมหสฺส, อเนกวิหิตานํ ปาปกานํ อกุสลานํ ธมฺมานํ วินยาย ธมฺมํ เทเสมิ.\csromanspacer{} อยํ โข สีห ปริยาโย, เยน มํ ปริยาเยน สมฺมา วทมาโน วเทยฺย “เวนยิโก สมโณ โคตโม, วินยาย ธมฺมํ เทเสติ, เตน จ สาวเก วิเนตี”ติ. (5)}

\prose{กตโม จ สีห ปริยาโย, เยน มํ ปริยาเยน สมฺมา วทมาโน วเทยฺย “ตปสฺสี สมโณ โคตโม, ตปสฺสิตาย ธมฺมํ เทเสติ, เตน จ สาวเก วิเนตี”ติ.\csromanspacer{} ตปนียาหํ สีห ปาปเก อกุสเล ธมฺเม วทามิ กายทุจฺจริตํ วจีทุจฺจริตํ มโนทุจฺจริตํ, ยสฺส โข สีห ตปนียา ปาปกา อกุสลา ธมฺมา ปหีนา อุจฺฉินฺนมูลา ตาลาวตฺถุกตา อนภาวํกตา อายติํ อนุปฺปาทธมฺมา, ตมหํ ‘ตปสฺสี’ติ วทามิ.\csromanspacer{} ตถาคตสฺส โข สีห ตปนียา ปาปกา อกุสลา ธมฺมา ปหีนา อุจฺฉินฺนมูลา ตาลาวตฺถุกตา อนภาวํกตา อายติํ อนุปฺปาทธมฺมา.\csromanspacer{} อยํ โข สีห ปริยาโย, เยน มํ ปริยาเยน สมฺมา วทมาโน วเทยฺย “ตปสฺสี สมโณ โคตโม, ตปสฺสิตาย ธมฺมํ เทเสติ, เตน จ สาวเก วิเนตี”ติ. (6)}

\chapterheadpageread{2. มหาวคฺค}
\prose{\csromanfolio{29}กตโม จ สีห ปริยาโย, เยน มํ ปริยาเยน สมฺมา วทมาโน วเทยฺย “อปคพฺโภ สมโณ โคตโม, อปคพฺภตาย ธมฺมํ เทเสติ, เตน จ สาวเก วิเนตี”ติ.\csromanspacer{} ยสฺส โข สีห อายติํ คพฺภเสยฺยา ปุนพฺภวาภินิพฺพตฺติ ปหีนา อุจฺฉินฺนมูลา ตาลาวตฺถุกตา อนภาวํกตา อายติํ อนุปฺปาทธมฺมา, ตมหํ ‘อปคพฺโภ’ติ วทามิ.\csromanspacer{} ตถาคตสฺส โข สีห อายติํ คพฺภเสยฺยา ปุนพฺภวาภินิพฺพตฺติ ปหีนา อุจฺฉินฺนมูลา ตาลาวตฺถุกตา อนภาวํกตา อายติํ อนุปฺปาทธมฺมา.\csromanspacer{} อยํ โข สีห ปริยาโย, เยน มํ ปริยาเยน สมฺมา วทมาโน วเทยฺย “อปคพฺโภ สมโณ โคตโม, อปคพฺภตาย ธมฺมํ เทเสติ, เตน จ สาวเก วิเนตี”ติ. (7)}

\prose{กตโม จ สีห ปริยาโย, เยน มํ ปริยาเยน สมฺมา วทมาโน วเทยฺย “อสฺสาสโก สมโณ โคตโม, อสฺสาสาย ธมฺมํ เทเสติ, เตน จ สาวเก วิเนตี”ติ.\csromanspacer{} อหํ หิ สีห อสฺสาสโก ปรเมน อสฺสาเสน อสฺสาสาย ธมฺมํ เทเสมิ, เตน จ สาวเก วิเนมิ.\csromanspacer{} อยํ โข สีห ปริยาโย, เยน มํ ปริยาเยน สมฺมา วทมาโน วเทยฺย “อสฺสาสโก สมโณ โคตโม, อสฺสาสาย ธมฺมํ เทเสติ, เตน จ สาวเก วิเนตี”ติ. (8)}

\prose{เอวํ วุตฺเต สีโห เสนาปติ ภควนฺตํ เอตทโวจ “อภิกฺกนฺตํ ภนฺเต, อภิกฺกนฺตํ ภนฺเต ฯเปฯ อุปาสกํ มํ ภนฺเต ภควา ธาเรตุ อชฺชตคฺเค ปาณุเปตํ สรณํ คตนฺ”ติ.}

\prose{อนุวิจฺจการํ โข สีห กโรหิ, อนุวิจฺจกาโร ตุมฺหาทิสานํ ญาตมนุสฺสานํ สาธุ โหตีติ.\csromanspacer{} อิมินาปาหํ ภนฺเต ภควโต ภิยฺโยโสมตฺตาย อตฺตมโน อภิรทฺโธ.\csromanspacer{} ยํ มํ ภควา เอวมาห “อนุวิจฺจการํ โข สีห กโรหิ, อนุวิจฺจกาโร ตุมฺหาทิสานํ ญาตมนุสฺสานํ สาธุ โหตี”ติ.\csromanspacer{} มํ หิ ภนฺเต อญฺญติตฺถิยา สาวกํ ลภิตฺวา เกวลกปฺปํ เวสาลิํ ปฏากํ ปริหเรยฺยุํ “สีโห อมฺหากํ เสนาปติ สาวกตฺตํ อุปคโต”ติ.\csromanspacer{} อถ จ ปน ภควา เอวมาห “อนุวิจฺจการํ สีห กโรหิ, อนุวิจฺจกาโร ตุมฺหาทิสานํ ญาตมนุสฺสานํ สาธุ โหตี”ติ.\csromanspacer{} เอสาหํ ภนฺเต ทุติยมฺปิ ภควนฺตํ สรณํ คจฺฉามิ ธมฺมญฺจ ภิกฺขุสํฆญฺจ, อุปาสกํ มํ ภควา ธาเรตุ, อชฺชตคฺเค ปาณุเปตํ สรณํ คตนฺติ.}

\prose{\csromanfolio{30}ทีฆรตฺตํ โข เต สีห นิคณฺฐานํ โอปานภูตํ กุลํ, เยน เนสํ อุปคตานํ ปิณฺฑกํ ทาตพฺพํ มญฺเญยฺยาสีติ.\csromanspacer{} อิมินาปาหํ ภนฺเต ภควโต ภิยฺโยโสมตฺตาย อตฺตมโน อภิรทฺโธ, ยํ มํ ภควา เอวมาห “ทีฆรตฺตํ โข เต สีห นิคณฺฐานํ โอปานภูตํ กุลํ, เยน เนสํ อุปคตานํ ปิณฺฑกํ ทาตพฺพํ มญฺเญยฺยาสี”ติ.\csromanspacer{} สุตํ เมตํ ภนฺเต สมโณ โคตโม เอวมาห “มยฺหเมว ทานํ ทาตพฺพํ, มยฺหเมว สาวกานํ ทาตพฺพํ, มยฺหเมว ทินฺนํ มหปฺผลํ, น อญฺเญสํ ทินฺนํ มหปฺผลํ, มยฺหเมว สาวกานํ ทินฺนํ มหปฺผลํ, น อญฺเญสํ สาวกานํ ทินฺนํ มหปฺผลนฺ”ติ.\csromanspacer{} อถ จ ปน มํ ภควา นิคณฺเฐสุปิ ทาเน สมาทเปติ\footnote{สมาทาเปติ (?)}, อปิ จ ภนฺเต มยเมตฺถ กาลํ ชานิสฺสาม, เอสาหํ ภนฺเต ตติยมฺปิ ภควนฺตํ สรณํ คจฺฉามิ ธมฺมญฺจ ภิกฺขุสํฆญฺจ, อุปาสกํ มํ ภนฺเต ภควา ธาเรตุ อชฺชตคฺเค ปาณุเปตํ สรณํ คตนฺติ.}

\prose{อถ โข ภควา สีหสฺส เสนาปติสฺส อนุปุพฺพิํ กถํ\footnote{อนุปุพฺพิกถํ (สพฺพตฺถ)} กเถสิ.\csromanspacer{} เสยฺยถิทํ, ทานกถํ สีลกถํ สคฺคกถํ กามานํ อาทีนวํ โอการํ สํกิเลสํ เนกฺขมฺเม อานิสํสํ ปกาเสสิ, ยทา ภควา อญฺญาสิ สีหํ เสนาปติํ กลฺลจิตฺตํ มุทุจิตฺตํ วินีวรณจิตฺตํ อุทคฺคจิตฺตํ ปสนฺนจิตฺตํ.\csromanspacer{} อถ ยา พุทฺธานํ สามุกฺกํสิกา ธมฺมเทสนา, ตํ ปกาเสสิ\csromandash{}ทุกฺขํ สมุทยํ นิโรธํ มคฺคํ.\csromanspacer{} เสยฺยถาปิ นาม สุทฺธํ วตฺถํ อปคตกาฬกํ สมฺมเทว รชนํ ปฏิคฺคเณฺหยฺย.\csromanspacer{} เอวเมวํ สีหสฺส เสนาปติสฺส ตสฺมิํเยว อาสเน วิรชํ วีตมลํ ธมฺมจกฺขุํ อุทปาทิ “ยํ กิญฺจิ สมุทยธมฺมํ, สพฺพํ ตํ นิโรธธมฺมนฺ”ติ.}

\prose{อถ โข สีโห เสนาปติ ทิฏฺฐธมฺโม ปตฺตธมฺโม วิทิตธมฺโม ปริโยคาฬฺหธมฺโม ติณฺณวิจิกิจฺโฉ วิคตกถํกโถ เวสารชฺชปฺปตฺโต อปรปฺปจฺจโย สตฺถุสาสเน ภควนฺตํ เอตทโวจ “อธิวาเสตุ เม ภนฺเต ภควา สฺวาตนาย ภตฺตํ สทฺธิํ ภิกฺขุสํเฆนา”ติ.\csromanspacer{} อธิวาเสสิ ภควา ตุณฺหีภาเวน.}

\prose{อถ โข สีโห เสนาปติ ภควโต อธิวาสนํ วิทิตฺวา อุฏฺฐายาสนา ภควนฺตํ อภิวาเทตฺวา ปทกฺขิณํ กตฺวา ปกฺกามิ, อถ โข สีโห เสนาปติ อญฺญตรํ ปุริสํ อามนฺเตสิ “คจฺฉ ตฺวํ อมฺโภ}

\vspace{\tipitakaparskip}
\csromancenter{มหาวคฺค ปุริส ปวตฺตมํสํ ชานาหี”ติ.\csromanspacer{} อถ โข สีโห เสนาปติ ตสฺสา รตฺติยา อจฺจเยน สเก นิเวสเน ปณีตํ ขาทนียํ โภชนียํ ปฏิยาทาเปตฺวา ภควโต กาลํ อาโรจาเปสิ “กาโล ภนฺเต นิฏฺฐิตํ ภตฺตนฺ”ติ.}
\prose{\csromanfolio{31}อถ โข ภควา ปุพฺพณฺหสมยํ นิวาเสตฺวา ปตฺตจีวรมาทาย เยน สีหสฺส เสนาปติสฺส นิเวสนํ เตนุปสงฺกมิ, อุปสงฺกมิตฺวา ปญฺญตฺเต อาสเน นิสีทิ สทฺธิํ ภิกฺขุสํเฆน.\csromanspacer{} เตน โข ปน สมเยน สมฺพหุลา นิคณฺฐา เวสาลิยํ รถิกาย รถิกํ\footnote{รถิยาย รถิยํ (พหูสุ)} สิงฺฆาฏเกน สิงฺฆาฏกํ พาหา ปคฺคยฺห กนฺทนฺติ “อชฺช สีเหน เสนาปตินา ถูลํ ปสุํ วธิตฺวา สมณสฺส โคตมสฺส ภตฺตํ กตํ, ตํ สมโณ โคตโม ชานํ อุทฺทิสฺสกตํ มํสํ ปริภุญฺชติ, ปฏิจฺจกมฺมนฺ”ติ.}

\prose{อถ โข อญฺญตโร ปุริโส เยน สีโห เสนาปติ เตนุปสงฺกมิ, อุปสงฺกมิตฺวา สีหสฺส เสนาปติสฺส อุปกณฺณเก อาโรเจสิ “ยคฺเฆ ภนฺเต ชาเนยฺยาสิ เอเต สมฺพหุลา นิคณฺฐา เวสาลิยํ รถิกาย รถิกํ สิงฺฆาฏเกน สิงฺฆาฏกํ พาหา ปคฺคยฺห กนฺทนฺติ ‘อชฺช สีเหน เสนาปตินา ถูลํ ปสุํ วธิตฺวา สมณสฺส โคตมสฺส ภตฺตํ กตํ, ตํ สมโณ โคตโม ชานํ อุทฺทิสฺสกตํ มํสํ ปริภุญฺชติ, ปฏิจฺจกมฺมนฺ’ติ.\csromanspacer{} อลํ อยฺโย ทีฆรตฺตํ หิ เต อายสฺมนฺโต อวณฺณกามา พุทฺธสฺส, อวณฺณกามา ธมฺมสฺส, อวณฺณกามา สํฆสฺส, น จ ปเนเต อายสฺมนฺโต ชิริทนฺติ ตํ ภควนฺตํ อสตา ตุจฺฉา มุสา อภูเตน อพฺภาจิกฺขิตุํ, น จ มยํ ชีวิตเหตุปิ สญฺจิจฺจ ปาณํ ชีวิตา โวโรเปยฺยามา”ติ.}

\prose{อถ โข สีโห เสนาปติ พุทฺธปฺปมุขํ ภิกฺขุสํฆํ ปณีเตน ขาทนีเยน โภชนีเยน สหตฺถา สนฺตปฺเปสิ สมฺปวาเรสิ.\csromanspacer{} อถ โข สีโห เสนาปติ ภควนฺตํ ภุตฺตาวิํ โอนีตปตฺตปาณิํ เอกมนฺตํ นิสีทิ, เอกมนฺตํ นิสินฺนํ โข สีหํ เสนาปติํ ภควา ธมฺมิยา กถาย สนฺทสฺเสตฺวา สมาทเปตฺวา สมุตฺเตเชตฺวา สมฺปหํเสตฺวา อุฏฺฐายาสนา ปกฺกามีติ.\csromanspacer{}\csromanspacer{}ทุติยํ.}

\csromanheader{2}{3. อสฺสาชานียสุตฺต}
\proseitem{13}{\csromanfolio{32}อฏฺฐหิ ภิกฺขเว องฺเคหิ สมนฺนาคโต รญฺโญ ภทฺโท อสฺสาชานีโย ราชารโห โหติ ราชโภคฺโค, รญฺโญ องฺคนฺเตว สงฺขํ คจฺฉติ.\csromanspacer{} กตเมหิ อฏฺฐหิ? อิธ ภิกฺขเว รญฺโญ ภทฺโท อสฺสาชานีโย อุภโต สุชาโต โหติ มาติโต จ ปิติโต จ.\csromanspacer{} ยสฺสํ ทิสายํ อญฺเญปิ ภทฺทา อสฺสาชานียา ชายนฺติ, ตสฺสํ ทิสายํ ชาโต โหติ.\csromanspacer{} ยํ โข ปนสฺส โภชนํ เทนฺติ อลฺลํ วา สุกฺขํ วา, ตํ สกฺกจฺจํเยว ปริภุญฺชติ อวิกิรนฺโต.\csromanspacer{} เชคุจฺฉี โหติ อุจฺฉารํ วา ปสฺสาวํ วา อภินิสีทิตุํ วา อภินิปชฺชิตุํ วา.\csromanspacer{} โสรโต โหติ สุขสํวาโส, น จ อญฺเญ อสฺเส อุพฺเพเชตา.\csromanspacer{} ยานิ โข ปนสฺส โหนฺติ\footnote{ยานิ โข ปนสฺส ตานิ (สฺยา)} สาเฐยฺยานิ กูเฏยฺยานิ ชิเมฺหยฺยานิ วงฺเกยฺยานิ, ตานิ ยถาภูตํ สารถิสฺส อาวิกตฺตา โหติ, เตสมสฺส สารถิ อภินิมฺมทนาย วายมติ.\csromanspacer{} วาหี โข ปน โหติ, “กามญฺเญ อสฺสา วหนฺตุ วา มา วา, อหเมตฺถ วหิสฺสามี”ติ จิตฺตํ อุปฺปาเทติ.\csromanspacer{} คจฺฉนฺโต โข ปน อุชุมคฺเคเนว คจฺฉติ, ถามวา โหติ ยาว ชีวิตมรณปริยาทานา, ถามํ อุปทํเสตา.\csromanspacer{} อิเมหิ โข ภิกฺขเว อฏฺฐหิ องฺเคหิ สมนฺนาคโต รญฺโญ ภทฺโท อสฺสาชานีโย ราชารโห โหติ ราชโภคฺโค, รญฺโญ องฺคนฺเตว สงฺขํ คจฺฉติ.}

\prose{เอวเมวํ โข ภิกฺขเว อฏฺฐหิ ธมฺเมหิ สมนฺนาคโต ภิกฺขุ อาหุเนยฺโย โหติ ฯเปฯ อนุตฺตรํ ปุญฺญกฺเขตฺตํ โลกสฺส.\csromanspacer{} กตเมหิ อฏฺฐหิ? อิธ ภิกฺขเว ภิกฺขุ สีลวา โหติ, ปาติโมกฺขสํวรสํวุโต วิหรติ อาจารโคจรสมฺปนฺโน, อณุมตฺเตสุ วชฺเชสุ ภยทสฺสาวี สมาทาย สิกฺขติ สิกฺขาปเทสุ.\csromanspacer{} ยํ โข ปนสฺส โภชนํ เทนฺติ ลูขํ วา ปณีตํ วา, ตํ สกฺกจฺจํเยว ปริภุญฺชติ อวิหญฺญมาโน.\csromanspacer{} เชคุจฺฉี โหติ กายทุจฺจริเตน วจีทุจฺจริเตน มโนทุจฺจริเตน, เชคุจฺฉี โหติ อเนกวิหิตานํ ปาปกานํ อกุสลานํ ธมฺมานํ สมาปตฺติยา.\csromanspacer{} โสรโต โหติ สุขสํวาโส, น อญฺเญ ภิกฺขู อุพฺเพเชตา.\csromanspacer{} ยานิ โข ปนสฺส โหนฺติ สาเฐยฺยานิ กูเฏยฺยานิ ชิเมฺหยฺยานิ วงฺเกยฺยานิ, ตานิ ยถาภูตํ อาวิกตฺตา โหติ สตฺถริ วา วิญฺญูสุ วา สพฺรหฺมจารีสุ, เตสมสฺส สตฺถา วา วิญฺญู วา สพฺรหฺมจารี อภินิมฺมทนาย วายมติ.\csromanspacer{} สิกฺขิตา โข ปน โหติ,}

\vspace{\tipitakaparskip}
\csromancenter{มหาวคฺค “กามญฺเญ ภิกฺขู สิกฺขนฺตุ วา มา วา, อหเมตฺถ สิกฺขิสฺสามี”ติ จิตฺตํ อุปฺปาเทติ.\csromanspacer{} คจฺฉนฺโต โข ปน อุชุมคฺเคเนว คจฺฉติ, ตตฺรายํ อุชุมคฺโค, เสยฺยถิทํ, สมฺมาทิฏฺฐิ ฯเปฯ สมฺมาสมาธิ.\csromanspacer{} อารทฺธวีริโย วิหรติ “กามํ ตโจ จ นฺหารุ\footnote{นหารุ (สี, สฺยา, กํ, อิ)} จ อฏฺฐิ จ อวสิสฺสตุ สรีเร อุปสุสฺสตุ มํสโลหิตํ, ยํ ตํ ปุริสถาเมน ปุริสวีริเยน ปุริสปรกฺกเมน ปตฺตพฺพํ, น ตํ อปาปุณิตฺวา วีริยสฺส สณฺฐานํ ภวิสฺสตี”ติ.\csromanspacer{} อิเมหิ โข ภิกฺขเว อฏฺฐหิ ธมฺเมหิ สมนฺนาคโต ภิกฺขุ อาหุเนยฺโย โหติ ฯเปฯ อนุตฺตรํ ปุญฺญกฺเขตฺตํ โลกสฺสาติ.\csromanspacer{}\csromanspacer{}ตติยํ.}
\csromanheader{2}{4. อสฺสขฬุงฺกสุตฺต}
\proseitem{14}{\csromanfolio{33}อฏฺฐ จ\footnote{อฏฺฐ (สฺยา)} ภิกฺขเว อสฺสขฬุงฺเก\footnote{อสฺสขลุงฺเก (สี)} เทเสสฺสามิ อฏฺฐ จ อสฺสโทเส, อฏฺฐ จ ปุริสขฬุงฺเก อฏฺฐ จ ปุริสโทเส, ตํ สุณาถ สาธุกํ มนสิ กโรถ, ภาสิสฺสามีติ. “เอวํ ภนฺเต”ติ โข เต ภิกฺขู ภควโต ปจฺจสฺโสสุํ.\csromanspacer{} ภควา เอตทโวจ\csromandash{}}

\prose{กตเม จ ภิกฺขเว อฏฺฐ อสฺสขฬุงฺกา? อฏฺฐ จ อสฺสโทสา, อิธ ภิกฺขเว เอกจฺโจ อสฺสขฬุงฺโก “เปหี”ติ วุตฺโต วิทฺโธ สมาโน โจทิโต สารถินา ปจฺฉโต ปฏิกฺกมติ, ปิฏฺฐิโต รถํ ปวตฺเตติ, เอวรูโปปิ ภิกฺขเว อิเธกจฺโจ อสฺสขฬุงฺโก โหติ.\csromanspacer{} อยํ ภิกฺขเว ปฐโม อสฺสโทโส.}

\prose{ปุน จปรํ ภิกฺขเว อิเธกจฺโจ อสฺสขฬุงฺโก “เปหี”ติ วุตฺโต วิทฺโธ สมาโน โจทิโต สารถินา ปจฺฉา ลงฺฆติ, กุพฺพรํ หนติ, ติทณฺฑํ ภญฺชติ, เอวรูโปปิ ภิกฺขเว อิเธกจฺโจ อสฺสขฬุงฺโก โหติ.\csromanspacer{} อยํ ภิกฺขเว ทุติโย อสฺสโทโส.}

\prose{ปุน จปรํ ภิกฺขเว อิเธกจฺโจ อสฺสขฬุงฺโก “เปหี”ติ วุตฺโต วิทฺโธ สมาโน โจทิโต สารถินา รถีสาย สตฺถิํ อุสฺสชฺชิตฺวา รถีสํเยว อชฺโฌมทฺทติ, เอวรูโปปิ ภิกฺขเว อิเธกจฺโจ อสฺสขฬุงฺโก โหติ.\csromanspacer{} อยํ ภิกฺขเว ตติโย อสฺสโทโส.}

\prose{ปุน จปรํ ภิกฺขเว อิเธกจฺโจ อสฺสขฬุงฺโก “เปหี”ติ วุตฺโต วิทฺโธ สมาโน โจทิโต สารถินา อุมฺมคฺคํ คณฺหติ, อุพฺพฏุมํ รถํ กโรติ, \csromanfolio{34}เอวรูโปปิ ภิกฺขเว อิเธกจฺโจ อสฺสขฬุงฺโก โหติ.\csromanspacer{} อยํ ภิกฺขเว จตุตฺโถ อสฺสโทโส.}

\prose{ปุน จปรํ ภิกฺขเว อิเธกจฺโจ อสฺสขฬุงฺโก “เปหี”ติ วุตฺโต วิทฺโธ สมาโน โจทิโต สารถินา ลงฺฆติ ปุริมกายํ, ปคฺคณฺหติ ปุริเม ปาเท, เอวรูโปปิ ภิกฺขเว อิเธกจฺโจ อสฺสขฬุงฺโก โหติ.\csromanspacer{} อยํ ภิกฺขเว ปญฺจโม อสฺสโทโส.}

\prose{ปุน จปรํ ภิกฺขเว อิเธกจฺโจ อสฺสขฬุงฺโก “เปหี”ติ วุตฺโต วิทฺโธ สมาโน โจทิโต สารถินา อนาทิยิตฺวา สารถิํ อนาทิยิตฺวา ปโตทลฏฺฐิํ\footnote{ปโตทํ (สี, อิ), ปโตทยฏฺฐิํ (สฺยา, กํ)} ทนฺเตหิ มุขาธานํ\footnote{มุขาฐานํ (ก)} วิธํสิตฺวา เยน กามํ ปกฺกมติ, เอวรูโปปิ ภิกฺขเว อิเธกจฺโจ อสฺสขฬุงฺโก โหติ.\csromanspacer{} อยํ ภิกฺขเว ฉฏฺโฐ อสฺสโทโส.}

\prose{ปุน จปรํ ภิกฺขเว อิเธกจฺโจ อสฺสขฬุงฺโก “เปหี”ติ วุตฺโต วิทฺโธ สมาโน โจทิโต สารถินา เนว อภิกฺกมติ, โน ปฏิกฺกมติ, ตตฺเถว ขีลฏฺฐายี ฐิโต โหติ, เอวรูโปปิ ภิกฺขเว อิเธกจฺโจ อสฺสขฬุงฺโก โหติ.\csromanspacer{} อยํ ภิกฺขเว สตฺตโม อสฺสโทโส.}

\prose{ปุน จปรํ ภิกฺขเว อิเธกจฺโจ อสฺสขฬุงฺโก “เปหี”ติ วุตฺโต วิทฺโธ สมาโน โจทิโต สารถินา ปุริเม จ ปาเท สํหริตฺวา ปจฺฉิเม จ ปาเท สํหริตฺวา\footnote{สงฺขริตฺวา (ก)} ตตฺเถว จตฺตาโร ปาเท อภินิสีทติ, เอวรูโปปิ ภิกฺขเว อิเธกจฺโจ อสฺสขฬุงฺโก โหติ.\csromanspacer{} อยํ ภิกฺขเว อฏฺฐโม อสฺสโทโส.\csromanspacer{} อิเม โข ภิกฺขเว อฏฺฐ อสฺสขฬุงฺกา, อฏฺฐ จ อสฺสโทสา.}

\prose{\csromansymbolfootnote{*}{อภิ 2. 402 ปิฏฺเฐปิ.}กตเม จ ภิกฺขเว อฏฺฐ ปุริสขฬุงฺกา, อฏฺฐ จ ปุริสโทสา, อิธ ภิกฺขเว ภิกฺขู ภิกฺขุํ อาปตฺติยา โจเทนฺติ, โส ภิกฺขุ ภิกฺขูหิ อาปตฺติยา โจทิยมาโน “น สรามี”ติ อสติยา นิพฺเพเฐติ.\csromanspacer{} เสยฺยถาปิ โส ภิกฺขเว อสฺสขฬุงฺโก “เปหี”ติ วุตฺโต วิทฺโธ สมาโน โจทิโต สารถินา ปจฺฉโต ปฏิกฺกมติ, ปิฏฺฐิโต รถํ วตฺเตติ.\csromanspacer{} ตถูปมาหํ ภิกฺขเว อิมํ ปุคฺคลํ วทามิ.\csromanspacer{} เอวรูโปปิ ภิกฺขเว อิเธกจฺโจ ปุริสขฬุงฺโก โหติ.\csromanspacer{} อยํ ภิกฺขเว ปฐโม ปุริสโทโส.}

\chapterheadpageread{2. มหาวคฺค}
\prose{\csromanfolio{35}ปุน จปรํ ภิกฺขเว ภิกฺขู ภิกฺขุํ อาปตฺติยา โจเทนฺติ, โส ภิกฺขุ ภิกฺขูหิ อาปตฺติยา โจทิยมาโน โจทกํเยว ปฏิปฺผรติ “กิํ นุ โข ตุยฺหํ พาลสฺส อพฺยตฺตสฺส ภณิเตน, ตฺวมฺปิ นาม ภณิตพฺพํ มญฺญสี”ติ.\csromanspacer{} เสยฺยถาปิ โส ภิกฺขเว อสฺสขฬุงฺโก “เปหี”ติ วุตฺโต วิทฺโธ สมาโน โจทิโต สารถินา ปจฺฉา ลงฺฆติ, กุพฺพรํ หนติ, ติทณฺฑํ ภญฺชติ, ตถูปมาหํ ภิกฺขเว อิมํ ปุคฺคลํ วทามิ, เอวรูโปปิ ภิกฺขเว อิเธกจฺโจ ปุริสขฬุงฺโก โหติ.\csromanspacer{} อยํ ภิกฺขเว ทุติโย ปุริสโทโส.}

\prose{ปุน จปรํ ภิกฺขเว ภิกฺขู ภิกฺขุํ อาปตฺติยา โจเทนฺติ, โส ภิกฺขุ ภิกฺขูหิ อาปตฺติยา โจทิยมาโน โจทกสฺเสว ปจฺจาโรเปติ “ตฺวํ โขสิ อิตฺถนฺนามํ อาปตฺติํ อาปนฺโน, ตฺวํ ตาว ปฐมํ ปฏิกโรหี”ติ.\csromanspacer{} เสยฺยถาปิ โส ภิกฺขเว อสฺสขฬุงฺโก “เปหี”ติ วุตฺโต วิทฺโธ สมาโน โจทิโต สารถินา รถีสาย สตฺถิํ อุสฺสชฺชิตฺวา รถีสํเยว อชฺโฌมทฺทติ, ตถูปมาหํ ภิกฺขเว อิมํ ปุคฺคลํ วทามิ, เอวรูโปปิ ภิกฺขเว อิเธกจฺโจ ปุริสขฬุงฺโก โหติ.\csromanspacer{} อยํ ภิกฺขเว ตติโย ปุริสโทโส.}

\prose{ปุน จปรํ ภิกฺขเว ภิกฺขู ภิกฺขุํ อาปตฺติยา โจเทนฺติ, โส ภิกฺขุ ภิกฺขูหิ อาปตฺติยา โจทิยมาโน อญฺเญนาญฺญํ ปฏิจรติ, พหิทฺธา กถํ อปนาเมติ, โกปญฺจ โทสญฺจ อปฺปจฺจยญฺจ ปาตุกโรติ.\csromanspacer{} เสยฺยถาปิ โส ภิกฺขเว อสฺสขฬุงฺโก “เปหี”ติ วุตฺโต วิทฺโธ สมาโน โจทิโต สารถินา อุมฺมคฺคํ คณฺหติ, อุพฺพฏุมํ รถํ กโรติ, ตถูปมาหํ ภิกฺขเว อิมํ ปุคฺคลํ วทามิ, เอวรูโปปิ ภิกฺขเว อิเธกจฺโจ ปุริสขฬุงฺโก โหติ.\csromanspacer{} อยํ ภิกฺขเว จตุตฺโถ ปุริสโทโส.}

\prose{ปุน จปรํ ภิกฺขเว ภิกฺขู ภิกฺขุํ อาปตฺติยา โจเทนฺติ, โส ภิกฺขุ ภิกฺขูหิ อาปตฺติยา โจทิยมาโน สํฆมชฺเฌ พาหุวิกฺเขปํ กโรติ.\csromanspacer{} เสยฺยถาปิ โส ภิกฺขเว อสฺสขฬุงฺโก “เปหี”ติ วุตฺโต วิทฺโธ สมาโน โจทิโต สารถินา ลงฺฆติ ปุริมกายํ, ปคฺคณฺหติ ปุริเม ปาเท, ตถูปมาหํ ภิกฺขเว อิมํ ปุคฺคลํ วทามิ, เอวรูโปปิ ภิกฺขเว อิเธกจฺโจ ปุริสขฬุงฺโก โหติ.\csromanspacer{} อยํ ภิกฺขเว ปญฺจโม ปุริสโทโส.}

\prose{\csromanfolio{36}ปุน จปรํ ภิกฺขเว ภิกฺขู ภิกฺขุํ อาปตฺติยา โจเทนฺติ, โส ภิกฺขุ ภิกฺขูหิ อาปตฺติยา โจทิยมาโน อนาทิยิตฺวา สํฆํ อนาทิยิตฺวา โจทกํ สาปตฺติโกว เยน กามํ ปกฺกมติ.\csromanspacer{} เสยฺยถาปิ โส ภิกฺขเว อสฺสขฬุงฺโก “เปหี”ติ วุตฺโต วิทฺโธ สมาโน โจทิโต สารถินา อนาทิยิตฺวา สารถิํ อนาทิยิตฺวา ปโตทลฏฺฐิํ ทนฺเตหิ มุขาธานํ วิธํสิตฺวา เยน กามํ ปกฺกมติ, ตถูปมาหํ ภิกฺขเว อิมํ ปุคฺคลํ วทามิ, เอวรูโปปิ ภิกฺขเว อิเธกจฺโจ ปุริสขฬุงฺโก โหติ.\csromanspacer{} อยํ ภิกฺขเว ฉฏฺโฐ ปุริสโทโส.}

\prose{ปุน จปรํ ภิกฺขเว ภิกฺขู ภิกฺขุํ อาปตฺติยา โจเทนฺติ, โส ภิกฺขุ ภิกฺขูหิ อาปตฺติยา โจทิยมาโน “เนวาหํ อาปนฺโนมฺหิ, น ปนาหํ อาปนฺโนมฺหี”ติ, โส ตุณฺหีภาเวน สํฆํ วิเหเฐติ\footnote{วิเหเสติ (อิ, ก)}.\csromanspacer{} เสยฺยถาปิ โส ภิกฺขเว อสฺสขฬุงฺโก “เปหี”ติ วุตฺโต วิทฺโธ สมาโน โจทิโต สารถินา เนว อภิกฺกมติ, โน ปฏิกฺกมติ, ตตฺเถว ขีลฏฺฐายี ฐิโต โหติ, ตถูปมาหํ ภิกฺขเว อิมํ ปุคฺคลํ วทามิ, เอวรูโปปิ ภิกฺขเว อิเธกจฺโจ ปุริสขฬุงฺโก โหติ.\csromanspacer{} อยํ ภิกฺขเว สตฺตโม ปุริสโทโส.}

\prose{ปุน จปรํ ภิกฺขเว ภิกฺขู ภิกฺขุํ อาปตฺติยา โจเทนฺติ, โส ภิกฺขุ ภิกฺขูหิ อาปตฺติยา โจทิยมาโน เอวมาห “กิํ นุ โข ตุเมฺห อายสฺมนฺโต อติพาฬฺหํ มยิ พฺยาวฏา, ยาว\footnote{อิทํ ปทํ สีหฬโปตฺถเก นตฺถิ.} อิทานาหํ สิกฺขํ ปจฺจกฺขาย หีนายาวตฺติสฺสามี”ติ.\csromanspacer{} โส สิกฺขํ ปจฺจกฺขาย หีนายาวตฺติตฺวา เอวมาห “อิทานิ โข ตุเมฺห อายสฺมนฺโต อตฺตมนา โหถา”ติ.\csromanspacer{} เสยฺยถาปิ โส ภิกฺขเว อสฺสขฬุงฺโก “เปหี”ติ วุตฺโต วิทฺโธ สมาโน โจทิโต สารถินา ปุริเม จ ปาเท สํหริตฺวา ปจฺฉิเม จ ปาเท สํหริตฺวา ตตฺเถว จตฺตาโร ปาเท อภินิสีทติ, ตถูปมาหํ ภิกฺขเว อิมํ ปุคฺคลํ วทามิ, เอวรูโปปิ ภิกฺขเว อิเธกจฺโจ ปุริสขฬุงฺโก โหติ.\csromanspacer{} อยํ ภิกฺขเว อฏฺฐโม ปุริสโทโส, อิเม โข ภิกฺขเว อฏฺฐ ปุริสขฬุงฺกา, อฏฺฐ จ ปุริสโทสาติ.\csromanspacer{}\csromanspacer{}จตุตฺถํ.}

\csromanheader{1}{2. มหาวคฺค \textbf{5. มลสุตฺต}}
\proseitem{15}{\csromanfolio{37}อฏฺฐิมานิ ภิกฺขเว มลานิ.\csromanspacer{} กตมานิ อฏฺฐ? อสชฺฌายมลา ภิกฺขเว มนฺตา, อนุฏฺฐานมลา ภิกฺขเว ฆรา.\csromanspacer{} มลํ ภิกฺขเว วณฺณสฺส โกสชฺชํ, ปมาโท ภิกฺขเว รกฺขโต มลํ.\csromanspacer{} มลํ ภิกฺขเว อิตฺถิยา ทุจฺจริตํ, มจฺเฉรํ ภิกฺขเว ททโต มลํ, มลา ภิกฺขเว ปาปกา อกุสลา ธมฺมา อสฺมิํ โลเก ปรมฺหิ จ, ตโต\footnote{ตโต จ (สฺยา, อิ)} ภิกฺขเว มลา มลตรํ อวิชฺชา ปรมํ มลํ, อิมานิ โข ภิกฺขเว อฏฺฐ มลานีติ.}

\csromangathameasure{อสชฺฌายมลา มนฺตา, อนุฏฺฐานมลา ฆรา. \\ มลํ วณฺณสฺส โกสชฺชํ, ปมาโท รกฺขโต มลํ. \\ มลิตฺถิยา ทุจฺจริตํ, มจฺเฉรํ ททโต มลํ. \\ มลา เว ปาปกา ธมฺมา, อสฺมิํ โลเก ปรมฺหิ จ. \\ ตโต มลา มลตรํ, อวิชฺชาปรมํ มลนฺติ.ปญฺจมํ.}
\csromangathagroup{\csromangathastanza{อสชฺฌายมลา มนฺตา, อนุฏฺฐานมลา ฆรา. \\ มลํ วณฺณสฺส โกสชฺชํ, ปมาโท รกฺขโต มลํ.}\csromangathastanza{มลิตฺถิยา ทุจฺจริตํ, มจฺเฉรํ ททโต มลํ. \\ มลา เว ปาปกา ธมฺมา, อสฺมิํ โลเก ปรมฺหิ จ. \\ ตโต มลา มลตรํ, อวิชฺชาปรมํ มลนฺติ.\csromanspacer{}\csromanspacer{}ปญฺจมํ.}}

\csromanheader{2}{6. ทูเตยฺยสุตฺต}
\proseitem{16}{\csromansymbolfootnote{*}{วิ 4. 365 ปิฏฺเฐปิ.}อฏฺฐหิ ภิกฺขเว ธมฺเมหิ สมนฺนาคโต ภิกฺขุ ทูเตยฺยํ คนฺตุมรหติ.\csromanspacer{} กตเมหิ อฏฺฐหิ? อิธ ภิกฺขเว ภิกฺขุ โสตา จ โหติ สาเวตา จ อุคฺคเหตา จ ธาเรตา จ วิญฺญาตา จ วิญฺญาเปตา จ, กุสโล จ สหิตาสหิตสฺส, โน จ กลหการโก.\csromanspacer{} อิเมหิ โข ภิกฺขเว อฏฺฐหิ ธมฺเมหิ สมนฺนาคโต ภิกฺขุ ทูเตยฺยํ คนฺตุมรหติ.\csromanspacer{} อฏฺฐหิ ภิกฺขเว ธมฺเมหิ สมนฺนาคโต สาริปุตฺโต ทูเตยฺยํ คนฺตุมรหติ.\csromanspacer{} กตเมหิ อฏฺฐหิ? อิธ ภิกฺขเว สาริปุตฺโต โสตา จ โหติ สาเวตา จ อุคฺคเหตา จ ธาเรตา จ วิญฺญาตา จ วิญฺญาเปตา จ, กุสโล จ สหิตาสหิตสฺส, โน จ กลหการโก.\csromanspacer{} อิเมหิ โข ภิกฺขเว อฏฺฐหิ ธมฺเมหิ สมนฺนาคโต สาริปุตฺโต ทูเตยฺยํ คนฺตุมรหตีติ.}

\csromangathameasure{โย เว น พฺยถติ ปตฺวา, ปริสํ อุคฺควาทินิํ3. \\ น จ หาเปติ วจนํ, น จ ฉาเทติ สาสนํ. \\ อสนฺทิทฺธญฺจ ภณติ, ปุจฺฉิโต น จ กุปฺปติ. \\ ส เว ตาทิสโก ภิกฺขุ, ทูเตยฺยํ คนฺตุมรหตีติ. ฉฏฺฐํ.}
\csromangathagroup{\csromangathastanza{โย เว น พฺยถติ\footnote{น เวธติ (สี), น พฺยาธติ (สฺยา, อิ) 3. อุคฺควาทินํ (สี), อุคฺคหวาทินํ (สฺยา, อิ), อุคฺคตวาทินิํ (ก)} ปตฺวา, ปริสํ อุคฺควาทินิํ3. \\ น จ หาเปติ วจนํ, น จ ฉาเทติ สาสนํ.}\csromangathastanza{\csromanfolio{38}อสนฺทิทฺธญฺจ ภณติ\footnote{อสนฺทิทฺโธ จ อกฺขาติ(วิ 4. 365 ปิฏฺเฐ)}, ปุจฺฉิโต น จ กุปฺปติ. \\ ส เว ตาทิสโก ภิกฺขุ, ทูเตยฺยํ คนฺตุมรหตีติ.\csromanspacer{} ฉฏฺฐํ.}}

\csromanheader{2}{7. ปฐมพนฺธนสุตฺต}
\csromanmatikaanchor{mtk.1}
\csromanmatikaanchor{mtk.46}
\csromantocmarknopage{section}{(7) 2. สติปฏฺฐานวคฺค}{mtk.1}
\bookmark[dest={mtk.1},level=1]{(7) 2. สติปฏฺฐานวคฺค}
\proseitem{17}{อฏฺฐหิ ภิกฺขเว อากาเรหิ อิตฺถี ปุริสํ พนฺธติ.\csromanspacer{} กตเมหิ อฏฺฐหิ? รุณฺเณน ภิกฺขเว อิตฺถี ปุริสํ พนฺธติ, หสิเตน ภิกฺขเว อิตฺถี ปุริสํ พนฺธติ, ภณิเตน ภิกฺขเว อิตฺถี ปุริสํ พนฺธติ, อากปฺเปน ภิกฺขเว อิตฺถี ปุริสํ พนฺธติ, วนภงฺเคน ภิกฺขเว อิตฺถี ปุริสํ พนฺธติ, คนฺเธน ภิกฺขเว อิตฺถี ปุริสํ พนฺธติ, รเสน ภิกฺขเว อิตฺถี ปุริสํ พนฺธติ, ผสฺเสน ภิกฺขเว อิตฺถี ปุริสํ พนฺธติ.\csromanspacer{} อิเมหิ โข ภิกฺขเว อฏฺฐหากาเรหิ อิตฺถี ปุริสํ พนฺธติ.\csromanspacer{} เต ภิกฺขเว สตฺตา สุพทฺธา\footnote{สุพนฺธา (สี, สฺยา, ก)}, เย\footnote{เยว (สฺยา, อิ, ก)} ผสฺเสน พทฺธาติ\footnote{พนฺธาติ (สี, สฺยา, ก)}.\csromanspacer{}\csromanspacer{}สตฺตมํ.}

\csromanheader{2}{8. ทุติยพนฺธนสุตฺต}
\proseitem{18}{อฏฺฐหิ ภิกฺขเว อากาเรหิ ปุริโส อิตฺถิํ พนฺธติ.\csromanspacer{} กตเมหิ อฏฺฐหิ? รุณฺเณน ภิกฺขเว ปุริโส อิตฺถิํ พนฺธติ, หสิเตน ภิกฺขเว ปุริโส อิตฺถิํ พนฺธติ, ภณิเตน ภิกฺขเว ปุริโส อิตฺถิํ พนฺธติ อากปฺเปน ภิกฺขเว ปุริโส อิตฺถิํ พนฺธติ, วนภงฺเคน ภิกฺขเว ปุริโส อิตฺถิํ พนฺธติ, คนฺเธน ภิกฺขเว ปุริโส อิตฺถิํ พนฺธติ, รเสน ภิกฺขเว ปุริโส อิตฺถิํ พนฺธติ, ผสฺเสน ภิกฺขเว ปุริโส อิตฺถิํ พนฺธติ.\csromanspacer{} อิเมหิ โข ภิกฺขเว อฏฺฐหากาเรหิ ปุริโส อิตฺถิํ พนฺธติ.\csromanspacer{} เต ภิกฺขเว สตฺตา สุพทฺธา, เย ผสฺเสน พทฺธาติ.\csromanspacer{}\csromanspacer{}อฏฺฐมํ.}

\csromanheader{2}{9. ปหาราทสุตฺต}
\proseitem{19}{เอกํ สมยํ ภควา เวรญฺชายํ วิหรติ นเฬรุปุจิมนฺทมูเล.\csromanspacer{} อถ โข ปหาราโท อสุรินฺโท เยน ภควา เตนุปสงฺกมิ, อุปสงฺกมิตฺวา ภควนฺตํ อภิวาเทตฺวา เอกมนฺตํ อฏฺฐาสิ, เอกมนฺตํ ฐิตํ โข ปหาราทํ อสุรินฺทํ ภควา เอตทโวจ\csromandash{}}

\chapterheadpageread{2. มหาวคฺค}
\prose{\csromanfolio{39}อปิ\footnote{กิํ (ก)} ปน ปหาราท อสุรา มหาสมุทฺเท อภิรมนฺตีติ.\csromanspacer{} อภิรมนฺติ ภนฺเต อสุรา มหาสมุทฺเทติ.\csromanspacer{} กติ ปน ปหาราท มหาสมุทฺเท อจฺฉริยา อพฺภุตา ธมฺมา\footnote{อพฺภุตธมฺมา (สฺยา, ก) วิ 4. 419 ปิฏฺเฐ ปสฺสิตพฺพํ.}, เย ทิสฺวา ทิสฺวา อสุรา มหาสมุทฺเท อภิรมนฺตีติ.\csromanspacer{} อฏฺฐ ภนฺเต มหาสมุทฺเท อจฺฉริยา อพฺภุตา ธมฺมา, เย ทิสฺวา ทิสฺวา อสุรา มหาสมุทฺเท อภิรมนฺติ.\csromanspacer{} กตเม อฏฺฐ? มหาสมุทฺโท ภนฺเต อนุปุพฺพนินฺโน อนุปุพฺพโปโณ อนุปุพฺพปพฺภาโร น อายตเกเนว ปปาโต.\csromanspacer{} ยมฺปิ ภนฺเต มหาสมุทฺโท อนุปุพฺพนินฺโน อนุปุพฺพโปโณ อนุปุพฺพปพฺภาโร น อายตเกเนว ปปาโต.\csromanspacer{} อยํ ภนฺเต มหาสมุทฺเท ปฐโม อจฺฉริโย อพฺภุโต ธมฺโม, ยํ ทิสฺวา ทิสฺวา อสุรา มหาสมุทฺเท อภิรมนฺติ.}

\prose{ปุน จปรํ ภนฺเต มหาสมุทฺโท ฐิตธมฺโม เวลํ นาติวตฺตติ, ยมฺปิ ภนฺเต มหาสมุทฺโท ฐิตธมฺโม เวลํ นาติวตฺตติ.\csromanspacer{} อยํ\footnote{อยมฺปิ (ก)} ภนฺเต มหาสมุทฺเท ทุติโย อจฺฉริโย อพฺภุโต ธมฺโม, ยํ ทิสฺวา ทิสฺวา อสุรา มหาสมุทฺเท อภิรมนฺติ.}

\prose{ปุน จปรํ ภนฺเต มหาสมุทฺโท น มเตน กุณเปน สํวสติ\footnote{สํวตฺตติ (สฺยา)}.\csromanspacer{} ยํ โหติ มหาสมุทฺเท มตํ กุณปํ, ตํ ขิปฺปเมว\footnote{ขิปฺปํเยว (สี), ขิปฺปํเอว (อิ), ขิปฺปญฺเญว (วิ 4. 420 ปิฏฺเฐ)} ตีรํ วาเหติ ถลํ อุสฺสาเรติ.\csromanspacer{} ยมฺปิ ภนฺเต มหาสมุทฺโท น มเตน กุณเปน สํวสติ.\csromanspacer{} ยํ โหติ มหาสมุทฺเท มตํ กุณปํ, ตํ ขิปฺปเมว ตีรํ วาเหติ ถลํ อุสฺสาเรติ.\csromanspacer{} อยํ ภนฺเต มหาสมุทฺเท ตติโย อจฺฉริโย อพฺภุโต ธมฺโม, ยํ ทิสฺวา ทิสฺวา อสุรา มหาสมุทฺเท อภิรมนฺติ.}

\prose{ปุน จปรํ ภนฺเต ยา กาจิ มหานทิโย.\csromanspacer{} เสยฺยถิทํ, คงฺคา ยมุนา อจิรวตี สรภู มหี.\csromanspacer{} ตา มหาสมุทฺทํ ปตฺวา\footnote{ปตฺตา (ก, วิ 4. 420 ปิฏฺเฐ)} ชหนฺติ ปุริมานิ นามโคตฺตานิ “มหาสมุทฺโท”เตฺวว สงฺขํ คจฺฉนฺติ.\csromanspacer{} ยมฺปิ ภนฺเต ยา กาจิ มหานทิโย.\csromanspacer{} เสยฺยถิทํ, คงฺคา ยมุนา อจิรวตี สรภู มหี.\csromanspacer{} ตา มหาสมุทฺทํ ปตฺวา ชหนฺติ ปุริมานิ นามโคตฺตานิ “มหาสมุทฺโท”เตฺวว สงฺขํ คจฺฉนฺติ.\csromanspacer{} อยํ ภนฺเต มหาสมุทฺเท จตุตฺโถ อจฺฉริโย อพฺภุโต ธมฺโม, ยํ ทิสฺวา ทิสฺวา อสุรา มหาสมุทฺเท อภิรมนฺติ.}

\prose{\csromanfolio{40}ปุน จปรํ ภนฺเต ยา จ\footnote{ยา กาจิ (สฺยา, อิ, ก)} โลเก สวนฺติโย มหาสมุทฺทํ อปฺเปนฺติ, ยา จ อนฺตลิกฺขา ธารา ปปตนฺติ, น เตน มหาสมุทฺทสฺส อูนตฺตํ คา ปูรตฺตํ วา ปญฺญายติ.\csromanspacer{} ยมฺปิ ภนฺเต ยา จ โลเก สวนฺติโย มหาสมุทฺทํ อปฺเปนฺติ, ยา จ อนฺตลิกฺขา ธารา ปปตนฺติ.\csromanspacer{} น เตน มหาสมุทฺทสฺส อูนตฺตํ วา ปูรตฺตํ วา ปญฺญายติ.\csromanspacer{} อยํ ภนฺเต มหาสมุทฺเท ปญฺจโม อจฺฉริโย อพฺภุโต ธมฺโม, ยํ ทิสฺวา ทิสฺวา อสุรา มหาสมุทฺเท อภิรมนฺติ.}

\prose{ปุน จปรํ ภนฺเต มหาสมุทฺโท เอกรโส โลณรโส.\csromanspacer{} ยมฺปิ ภนฺเต มหาสมุทฺโท เอกรโส โลณรโส.\csromanspacer{} อยํ ภนฺเต มหาสมุทฺเท ฉฏฺโฐ อจฺฉริโย อพฺภุโต ธมฺโม, ยํ ทิสฺวา ทิสฺวา อสุรา มหาสมุทฺเท อภิรมนฺติ.}

\prose{ปุน จปรํ ภนฺเต มหาสมุทฺโท พหุรตโน\footnote{ปหูตรตโน (ก)} อเนกรตโน.\csromanspacer{} ตตฺริมานิ รตนานิ, เสยฺยถิทํ, มุตฺตา มณิ เวฬุริโย สงฺโข สิลา ปวาฬํ รชตํ ชาตรูปํ โลหิตโก มสารคลฺลํ.\csromanspacer{} ยมฺปิ ภนฺเต มหาสมุทฺโท พหุรตโน อเนกรตโน.\csromanspacer{} ตตฺริมานิ รตนานิ, เสยฺยถิทํ, มุตฺตา มณิ เวฬุริโย สงฺโข สิลา ปวาฬํ รชตํ ชาตรูปํ โลหิตโก มสารคลฺลํ.\csromanspacer{} อยํ ภนฺเต มหาสมุทฺเท สตฺตโม อจฺฉริโย อพฺภุโต ธมฺโม, ยํ ทิสฺวา ทิสฺวา อสุรา มหาสมุทฺเท อภิรมนฺติ.}

\prose{ปุน จปรํ ภนฺเต มหาสมุทฺโท มหตํ ภูตานํ อาวาโส.\csromanspacer{} ตตฺริเม ภูตา, ติมิ ติมิงฺคโล ติมิรปิงฺคโล\footnote{ติมิติมิงฺคลา ติมิรปิงฺคลา (สี), ติมิติมิงฺคลา ติมิรมิงฺคลา (สฺยา, อิ)} อสุรา นาคา คนฺธพฺพา, สนฺติ มหาสมุทฺเท โยชนสติกาปิ อตฺตภาวา, ทฺวิโยชนสติกาปิ อตฺตภาวา, ติโยชนสติกาปิ อตฺตภาวา, จตุโยชนสติกาปิ อตฺตภาวา, ปญฺจโยชนสติกาปิ อตฺตภาวา.\csromanspacer{} ยมฺปิ ภนฺเต มหาสมุทฺโท มหตํ ภูตานํ อาวาโส.\csromanspacer{} ตตฺริเม ภูตา, ติมิ ติมิงฺคโล ติมิรปิงฺคโล อสุรา นาคา คนฺธพฺพา, สนฺติ มหาสมุทฺเท โยชนสติกาปิ อตฺตภาวา ฯเปฯ ติโยชน.\csromanspacer{} จตุโยชน.\csromanspacer{} ปญฺจโยชนสติกาปิ อตฺตภาวา.\csromanspacer{} อยํ ภนฺเต มหาสมุทฺเท อฏฺฐโม อจฺฉริโย อพฺภุโต ธมฺโม, ยํ ทิสฺวา ทิสฺวา อสุรา มหาสมุทฺเท อภิรมนฺติ.\csromanspacer{} อิเม โข}

\vspace{\tipitakaparskip}
\csromancenter{มหาวคฺค ภนฺเต มหาสมุทฺเท อฏฺฐ อจฺฉริยา อพฺภุตา ธมฺมา, เย ทิสฺวา ทิสฺวา อสุรา มหาสมุทฺเท อภิรมนฺตีติ.}
\prose{\csromanfolio{41}อปิ ปน ภนฺเต ภิกฺขู อิมสฺมิํ ธมฺมวินเย อภิรมนฺตีติ.\csromanspacer{} อภิรมนฺติ ปหาราท ภิกฺขู อิมสฺมิํ ธมฺมวินเยติ.\csromanspacer{} กติ ปน ภนฺเต อิมสฺมิํ ธมฺมวินเย อจฺฉริยา อพฺภุตา ธมฺมา, เย ทิสฺวา ทิสฺวา ภิกฺขู อิมสฺมิํ ธมฺมวินเย อภิรมนฺตีติ.\csromanspacer{} อฏฺฐ ปหาราท อิมสฺมิํ ธมฺมวินเย อจฺฉริยา อพฺภุตา ธมฺมา, เย ทิสฺวา ทิสฺวา ภิกฺขู อิมสฺมิํ ธมฺมวินเย อภิรมนฺติ.\csromanspacer{} กตเม อฏฺฐ? เสยฺยถาปิ ปหาราท มหาสมุทฺโท อนุปุพฺพนินฺโน อนุปุพฺพโปโณ อนุปุพฺพปพฺภาโร น อายตเกเนว ปปาโต.\csromanspacer{} เอวเมวํ โข ปหาราท อิมสฺมิํ ธมฺมวินเย อนุปุพฺพสิกฺขา อนุปุพฺพกิริยา อนุปุพฺพปฏิปทา น อายตเกเนว อญฺญาปฏิเวโธ.\csromanspacer{} ยมฺปิ ปหาราท อิมสฺมิํ ธมฺมวินเย อนุปุพฺพสิกฺขา อนุปุพฺพกิริยา อนุปุพฺพปฏิปทา น อายตเกเนว อญฺญาปฏิเวโธ.\csromanspacer{} อยํ ปหาราท อิมสฺมิํ ธมฺมวินเย ปฐโม อจฺฉริโย อพฺภุโต ธมฺโม, ยํ ทิสฺวา ทิสฺวา ภิกฺขู อิมสฺมิํ ธมฺมวินเย อภิรมนฺติ.}

\prose{เสยฺยถาปิ ปหาราท มหาสมุทฺโท ฐิตธมฺโม เวลํ นาติวตฺตติ.\csromanspacer{} เอวเมวํ โข ปหาราท ยํ มยา สาวกานํ สิกฺขาปทํ ปญฺญตฺตํ, ตํ มม สาวกา ชีวิตเหตุปิ นาติกฺกมนฺติ.\csromanspacer{} ยมฺปิ ปหาราท มยา สาวกานํ สิกฺขาปทํ ปญฺญตฺตํ, ตํ มม สาวกา ชีวิตเหตุปิ นาติกฺกมนฺติ.\csromanspacer{} อยํ ปหาราท อิมสฺมิํ ธมฺมวินเย ทุติโย อจฺฉริโย อพฺภุโต ธมฺโม, ยํ ทิสฺวา ทิสฺวา ภิกฺขู อิมสฺมิํ ธมฺมวินเย อภิรมนฺติ.}

\prose{เสยฺยถาปิ ปหาราท มหาสมุทฺโท น มเตน กุณเปน สํวสติ.\csromanspacer{} ยํ โหติ มหาสมุทฺเท มตํ กุณปํ, ตํ ขิปฺปเมว ตีรํ วาเหติ, ถลํ อุสฺสาเรติ.\csromanspacer{} เอวเมวํ โข ปหาราท โย โส ปุคฺคโล ทุสฺสีโล ปาปธมฺโม อสุจิสงฺกสฺสรสมาจาโร ปฏิจฺฉนฺนกมฺมนฺโต อสฺสมโณ สมณปฏิญฺโญ อพฺรหฺมจารี พฺรหฺมจาริปฏิญฺโญ อนฺโตปูติ อวสฺสุโต กสมฺพุชาโต, น เตน สํโฆ สํวสติ, ขิปฺปเมว นํ สนฺนิปติตฺวา อุกฺขิปติ.}

\prose{กิญฺจาปิ โส โหติ มชฺเฌ ภิกฺขุสํฆสฺส สนฺนิสินฺโน, อถ โข โส อารกาว สํฆมฺหา สํโฆ จ เตน.\csromanspacer{} ยมฺปิ ปหาราท โย โส \csromanfolio{42}ปุคฺคโล ทุสฺสีโล ปาปธมฺโม อสุจิสงฺกสฺสรสมาจาโร ปฏิจฺฉนฺนกมฺมนฺโต อสฺสมโณ สมณปฏิญฺโญ อพฺรหฺมจารี พฺรหฺมจาริปฏิญฺโญ อนฺโตปูติ อวสฺสุโต กสมฺพุชาโต, น เตน สํโฆ สํวสติ, ขิปฺปเมว นํ สนฺนิปติตฺวา อุกฺขิปติ.\csromanspacer{} กิญฺจาปิ โส โหติ มชฺเฌ ภิกฺขุสํฆสฺส สนฺนิสินฺโน, อถ โข โส อารกาว สํฆมฺหา สํโฆ จ เตน.\csromanspacer{} อยํ ปหาราท อิมสฺมิํ ธมฺมวินเย ตติโย อจฺฉริโย อพฺภุโต ธมฺโม, ยํ ทิสฺวา ทิสฺวา ภิกฺขู อิมสฺมิํ ธมฺมวินเย อภิรมนฺติ.}

\prose{เสยฺยถาปิ ปหาราท ยา กาจิ มหานทิโย.\csromanspacer{} เสยฺยถิทํ, คงฺคา ยมุนา อจิรวตี สรภู มหี.\csromanspacer{} ตา มหาสมุทฺทํ ปตฺวา ชหนฺติ ปุริมานิ นาม โคตฺตานิ, “มหาสมุทฺโท”เตฺวว สงฺขํ คจฺฉนฺติ.\csromanspacer{} เอวเมวํ โข ปหาราท จตฺตาโรเม วณฺณา ขตฺติยา พฺราหฺมณา เวสฺสา สุทฺทา, เต ตถาคตปฺปเวทิเต ธมฺมวินเย อคารสฺมา อนคาริยํ ปพฺพชิตฺวา ชหนฺติ ปุริมานิ นามโคตฺตานิ, “สมณา สกฺยปุตฺติยา”เตฺวว\footnote{สมโณ สกฺยปุตฺติโย เตฺวว (สฺยา, ก)} สงฺขํ คจฺฉนฺติ.\csromanspacer{} ยมฺปิ ปหาราท จตฺตาโรเม วณฺณา ขตฺติยา พฺราหฺมณา เวสฺสา สุทฺทา, เต ตถาคตปฺปเวทิเต ธมฺมวินเย อคารสฺมา อนคาริยํ ปพฺพชิตฺวา ชหนฺติ ปุริมานิ นามโคตฺตานิ, “สมณา สกฺยปุตฺติยา”เตฺวว สงฺขํ คจฺฉนฺติ.\csromanspacer{} อยํ ปหาราท อิมสฺมิํ ธมฺมวินเย จตุตฺโถ อจฺฉริโย อพฺภุโต ธมฺโม, ยํ ทิสฺวา ทิสฺวา ภิกฺขู อิมสฺมิํ ธมฺมวินเย อภิรมนฺติ.}

\prose{เสยฺยถาปิ ปหาราท ยา จ โลเก สวนฺติโย มหาสมุทฺทํ อปฺเปนฺติ, ยา จ อนฺตลิกฺขา ธารา ปปตนฺติ, น เตน มหาสมุทฺทสฺส อูนตฺตํ วา ปูรตฺตํ วา ปญฺญายติ.\csromanspacer{} เอวเมวํ โข ปหาราท พหู เจปิ ภิกฺขู อนุปาทิเสสาย นิพฺพานธาตุยา ปรินิพฺพายนฺติ, น เตน นิพฺพานธาตุยา อูนตฺตํ วา ปูรตฺตํ วา ปญฺญายติ.\csromanspacer{} ยมฺปิ ปหาราท พหู เจปิ ภิกฺขู อนุปาทิเสสาย นิพฺพานธาตุยา ปรินิพฺพายนฺติ, น เตน นิพฺพานธาตุยา อูนตฺตํ วา ปูรตฺตํ วา ปญฺญายติ.\csromanspacer{} อยํ ปหาราท อิมสฺมิํ ธมฺมวินเย ปญฺจโม อจฺฉริโย อพฺภุโต ธมฺโม, ยํ ทิสฺวา ทิสฺวา ภิกฺขู อิมสฺมิํ ธมฺมวินเย อภิรมนฺติ.}

\prose{เสยฺยถาปิ ปหาราท มหาสมุทฺโท เอกรโส โลณรโส, เอวเมวํ โข ปหาราท อยํ ธมฺมวินโย เอกรโส วิมุตฺติรโส.\csromanspacer{} ยมฺปิ}

\vspace{\tipitakaparskip}
\csromancenter{มหาวคฺค ปหาราท อยํ ธมฺมวินโย เอกรโส วิมุตฺติรโส.\csromanspacer{} อยํ ปหาราท อิมสฺมิํ ธมฺมวินเย ฉฏฺโฐ อจฺฉริโย อพฺภุโต ธมฺโม, ยํ ทิสฺวา ทิสฺวา ภิกฺขู อิมสฺมิํ ธมฺมวินเย อภิรมนฺติ.}
\prose{\csromanfolio{43}เสยฺยถาปิ ปหาราท มหาสมุทฺโท พหุรตโน อเนกรตโน, ตตฺริมานิ, รตนานิ, เสยฺยถิทํ, มุตฺตา มณิ เวฬุริโย สงฺโข สิลา ปวาฬํ รชตํ ชาตรูปํ โลหิตโก มสารคลฺลํ.\csromanspacer{} เอวเมวํ โข ปหาราท อยํ ธมฺมวินโย พหุรตโน อเนกรตโน, ตตฺริมานิ รตนานิ, เสยฺยถิทํ, จตฺตาโร สติปฏฺฐานา, จตฺตาโร สมฺมปฺปธานา, จตฺตาโร อิทฺธิปาทา, ปญฺจินฺทฺริยานิ, ปญฺจ พลานิ, สตฺต โพชฺฌงฺคา, อริโย อฏฺฐงฺคิโก มคฺโค.\csromanspacer{} ยมฺปิ ปหาราท อยํ ธมฺมวินโย พหุรตโน อเนกรตโน.\csromanspacer{} ตตฺริมานิ รตนานิ, เสยฺยถิทํ, จตฺตาโร สติปฏฺฐานา, จตฺตาโร สมฺมปฺปธานา, จตฺตาโร อิทฺธิปาทา, ปญฺจินฺทฺริยานิ, ปญฺจ พลานิ, สตฺต โพชฺฌงฺคา, อริโย อฏฺฐงฺคิโก มคฺโค.\csromanspacer{} อยํ ปหาราท อิมสฺมิํ ธมฺมวินเย สตฺตโม อจฺฉริโย อพฺภุโต ธมฺโม, ยํ ทิสฺวา ทิสฺวา ภิกฺขู อิมสฺมิํ ธมฺมวินเย อภิรมนฺติ.}

\prose{เสยฺยถาปิ ปหาราท มหาสมุทฺโท มหตํ ภูตานํ อาวาโส.\csromanspacer{} ตตฺริเม ภูตา, ติมิ ติมิงฺคโล ติมิรปิงฺคโล อสุรา นาคา คนฺธพฺพา.\csromanspacer{} สนฺติ มหาสมุทฺเท โยชนสติกาปิ อตฺตภาวา, ทฺวิโยชนสติกาปิ อตฺตภาวา, ติโยชนสติกาปิ อตฺตภาวา, จตุโยชนสติกาปิ อตฺตภาวา, ปญฺจโยชนสติกาปิ อตฺตภาวา, เอวเมวํ โข ปหาราท อยํ ธมฺมวินโย มหตํ ภูตานํ อาวาโส.\csromanspacer{} ตตฺริเม ภูตา, โสตาปนฺโน, โสตาปตฺติผลสจฺฉิกิริยาย ปฏิปนฺโน, สกทาคามี, สกทาคามิผล- สจฺฉิกิริยาย ปฏิปนฺโน, อนาคามี, อนาคามิผลสจฺฉิกิริยาย ปฏิปนฺโน, อรหา, อรหตฺตาย ปฏิปนฺโน.\csromanspacer{} ยมฺปิ ปหาราท อยํ ธมฺมวินโย มหตํ ภูตานํ อาวาโส.\csromanspacer{} ตตฺริเม ภูตา, โสตาปนฺโน, โสตาปตฺติผล- สจฺฉิกิริยาย ปฏิปนฺโน, สกทาคามี, สกทาคามิผลสจฺฉิกิริยาย ปฏิปนฺโน, อนาคามี, อนาคามิผลสจฺฉิกิริยาย ปฏิปนฺโน, อรหา, อรหตฺตาย ปฏิปนฺโน.\csromanspacer{} อยํ ปหาราท อิมสฺมิํ ธมฺมวินเย อฏฺฐโม อจฺฉริโย อพฺภุโต ธมฺโม, ยํ ทิสฺวา ทิสฺวา ภิกฺขู อิมสฺมิํ ธมฺมวินเย อภิรมนฺติ.\csromanspacer{} อิเม โข ปหาราท อิมสฺมิํ ธมฺมวินเย \csromanfolio{44}อฏฺฐ อจฺฉริยา อพฺภุตา ธมฺมา, เย ทิสฺวา ทิสฺวา ภิกฺขู อิมสฺมิํ ธมฺมวินเย อภิรมนฺตีติ.\csromanspacer{}\csromanspacer{}นวมํ.}

\csromanheader{2}{10. อุโปสถสุตฺต}
\proseitem{20}{\csromansymbolfootnote{*}{วิ 4. 418; ขุ 1. 138; อภิ 4. 168 ปิฏฺเฐสุ.}เอกํ สมยํ ภควา สาวตฺถิยํ วิหรติ ปุพฺพาราเม มิคารมาตุปาสาเท.\csromanspacer{} เตน โข ปน สมเยน ภควา ตทหุโปสเถ ภิกฺขุสํฆปริวุโต นิสินฺโน โหติ.\csromanspacer{} อถ โข อายสฺมา อานนฺโท อภิกฺกนฺตาย รตฺติยา นิกฺขนฺเต ปฐเม ยาเม อุฏฺฐายาสนา เอกํสํ อุตฺตราสงฺคํ กริตฺวา เยน ภควา เตนญฺชลิํ ปณาเมตฺวา ภควนฺตํ เอตทโวจ “อภิกฺกนฺตา ภนฺเต รตฺติ, นิกฺขนฺโต ปฐโม ยาโม, จิรนิสินฺโน ภิกฺขุสํโฆ, อุทฺทิสตุ ภนฺเต ภควา ภิกฺขูนํ ปาติโมกฺขนฺ”ติ.}

\prose{เอวํ วุตฺเต ภควา ตุณฺหี อโหสิ.\csromanspacer{} ทุติยมฺปิ โข อายสฺมา อานนฺโท อภิกฺกนฺตาย รตฺติยา นิกฺขนฺเต มชฺฌิเม ยาเม อุฏฺฐายาสนา เอกํสํ อุตฺตราสงฺคํ กริตฺวา เยน ภควา เตนญฺชลิํ ปณาเมตฺวา ภควนฺตํ เอตทโวจ “อภิกฺกนฺตา ภนฺเต รตฺติ, นิกฺขนฺโต มชฺฌิโม ยาโม, จิรนิสินฺโน ภิกฺขุสํโฆ, อุทฺทิสตุ ภนฺเต ภควา ภิกฺขูนํ ปาติโมกฺขนฺ”ติ.\csromanspacer{} ทุติยมฺปิ โข ภควา ตุณฺหี อโหสิ.\csromanspacer{} ตติยมฺปิ โข อายสฺมา อานนฺโท อภิกฺกนฺตาย รตฺติยา นิกฺขนฺเต ปจฺฉิเม ยาเม อุทฺธสฺเต อรุเณ นนฺทิมุขิยา รตฺติยา อุฏฺฐายาสนา เอกํสํ อุตฺตราสงฺคํ กริตฺวา เยน ภควา เตนญฺชลิํ ปณาเมตฺวา ภควนฺตํ เอตทโวจ “อภิกฺกนฺตา ภนฺเต รตฺติ, นิกฺขนฺโต ปจฺฉิโม ยาเม, อุทฺธสฺตํ อรุณํ, นนฺทิมุขี รตฺติ, จิรนิสินฺโน ภิกฺขุสํโฆ, อุทฺทิสตุ ภนฺเต ภควา ภิกฺขูนํ ปาติโมกฺขนฺ”ติ.\csromanspacer{} อปริสุทฺธา อานนฺท ปริสาติ.}

\prose{อถ โข อายสฺมโต มหาโมคฺคลฺลานสฺส เอตทโหสิ “กํ นุ โข ภควา ปุคฺคลํ สนฺธาย เอวมาห ‘อปริสุทฺธา อานนฺท ปริสา’ติ”.\csromanspacer{} อถ โข อายสฺมา มหาโมคฺคลฺลาโน สพฺพาวนฺตํ ภิกฺขุสํฆํ เจตสา เจโต ปริจฺจ มนสากาสิ, อทฺทสา โข อายสฺมา มหาโมคฺคลฺลาโน ตํ ปุคฺคลํ ทุสฺสีลํ ปาปธมฺมํ อสุจิํ สงฺกสฺสรสมาจารํ ปฏิจฺฉนฺนกมฺมนฺตํ อสฺสมณํ สมณปฏิญฺญํ อพฺรหฺมจาริํ พฺรหฺมจาริปฏิญฺญํ อนฺโตปูติํ อวสฺสุตํ กสมฺพุชาตํ มชฺเฌ ภิกฺขุสํฆสฺส นิสินฺนํ, ทิสฺวาน อุฏฺฐายาสนา เยน โส ปุคฺคโล เตนุปสงฺกมิ, อุปสงฺกมิตฺวา ตํ}

\vspace{\tipitakaparskip}
\csromancenter{มหาวคฺค ปุคฺคลํ เอตทโวจ “อุฏฺเฐหาวุโส, ทิฏฺโฐสิ ภควตา, นตฺถิ เต ภิกฺขูหิ สทฺธิํ สํวาโส”ติ.}
\prose{\csromanfolio{45}เอวํ วุตฺเต โส ปุคฺคโล ตุณฺหี อโหสิ, ทุติยมฺปิ โข อายสฺมา มหาโมคฺคลฺลาโน ตํ ปุคฺคลํ เอตทโวจ “อุฏฺเฐหาวุโส, ทิฏฺโฐสิ ภควตา, นตฺถิ เต ภิกฺขูหิ สทฺธิํ สํวาโส”ติ.\csromanspacer{} ทุติยมฺปิ โข โส ปุคฺคโล ตุณฺหี อโหสิ.\csromanspacer{} ตติยมฺปิ โข อายสฺมา มหาโมคฺคลฺลาโน ตํ ปุคฺคลํ เอตทโวจ “อุฏฺเฐหาวุโส, ทิฏฺโฐสิ ภควตา, นตฺถิ เต ภิกฺขูหิ สทฺธิํ สํวาโส”ติ.\csromanspacer{} ตติยมฺปิ โข โส ปุคฺคโล ตุณฺหี อโหสิ.}

\prose{อถ โข อายสฺมา มหาโมคฺคลฺลาโน ตํ ปุคฺคลํ พาหายํ คเหตฺวา พหิทฺวารโกฏฺฐกา นิกฺขาเมตฺวา สูจิฆฏิกํ ทตฺวา เยน ภควา เตนุปสงฺกมิ, อุปสงฺกมิตฺวา ภควนฺตํ เอตทโวจ “นิกฺขามิโต โส ภนฺเต ปุคฺคโล มยา, ปริสุทฺธา ปริสา, อุทฺทิสตุ ภนฺเต ภควา ภิกฺขูนํ ปาติโมกฺขนฺ”ติ.\csromanspacer{} อจฺฉริยํ โมคฺคลฺลาน, อพฺภุตํ โมคฺคลฺลาน, ยาว พาหา คหณาปิ นาม โส โมฆปุริโส อาคมิสฺสตีติ.}

\prose{อถ โข ภควา ภิกฺขู อามนฺเตสิ\csromandash{}ตุเมฺหว ทานิ ภิกฺขเว อุโปสถํ กเรยฺยาถ, ปาติโมกฺขํ อุทฺทิเสยฺยาถ, น ทานาหํ ภิกฺขเว อชฺชตคฺเค อุโปสถํ กริสฺสามิ, ปาติโมกฺขํ อุทฺทิสิสฺสามิ.\csromanspacer{} อฏฺฐานเมตํ ภิกฺขเว อนวกาโส ยํ ตถาคโต อปริสุทฺธาย ปริสาย ปาติโมกฺขํ อุทฺทิเสยฺย.}

\prose{อฏฺฐิเม ภิกฺขเว มหาสมุทฺเท อจฺฉริยา อพฺภุตา ธมฺมา, เย ทิสฺวา ทิสฺวา อสุรา มหาสมุทฺเท อภิรมนฺติ.\csromanspacer{} กตเม อฏฺฐ? มหาสมุทฺโท ภิกฺขเว อนุปุพฺพนินฺโน อนุปุพฺพโปโณ อนุปุพฺพปพฺภาโร น อายตเกเนว ปปาโต.\csromanspacer{} ยมฺปิ ภิกฺขเว มหาสมุทฺโท อนุปุพฺพนินฺโน อนุปุพฺพโปโณ อนุปุพฺพปพฺภาโร น อายตเกเนว ปปาโต.\csromanspacer{} อยํ ภิกฺขเว มหาสมุทฺเท ปฐโม อจฺฉริโย อพฺภุโต ธมฺโม, ยํ ทิสฺวา ทิสฺวา อสุรา มหาสมุทฺเท อภิรมนฺติ. (ยถา ปูริเม ตถา วิตฺถาโร) ฯเปฯ.}

\prose{ปุน จปรํ ภิกฺขเว มหาสมุทฺโท มหตํ ภูตานํ อาวาโส.\csromanspacer{} ตตฺริเม ภูตา, ติมิ ติมิงฺคโล ติมิรปิงฺคโล อสุรา นาคา คนฺธพฺพา, วสนฺติ มหาสมุทฺเท โยชนสติกาปิ อตฺตภาวา ฯเปฯ.\csromanspacer{} ปญฺจโยชนสติกาปิ \csromanfolio{46}อตฺตภาวา.\csromanspacer{} ยมฺปิ ภิกฺขเว มหาสมุทฺโท มหตํ ภูตานํ อาวาโส.\csromanspacer{} ตตฺริเม ภูตา, ติมิ ติมิงฺคโล ติมิรปิงฺคโล อสุรา นาคา คนฺธพฺพา, วสนฺติ มหาสมุทฺเท โยชนสติกาปิ อตฺตภาวา ฯเปฯ ปญฺจโยชนสติกาปิ อตฺตภาวา.\csromanspacer{} อยํ ภิกฺขเว มหาสมุทฺเท อฏฺฐโม อจฺฉริโย อพฺภุโต ธมฺโม, ยํ ทิสฺวา ทิสฺวา อสุรา มหาสมุทฺเท อภิรมนฺติ.\csromanspacer{} อิเม โข ภิกฺขเว มหาสมุทฺเท อฏฺฐ อจฺฉริยา อพฺภุตา ธมฺมา, ยํ ทิสฺวา ทิสฺวา อสุรา มหาสมุทฺเท อภิรมนฺติ.}

\prose{เอวเมวํ โข ภิกฺขเว อฏฺฐ อิมสฺมิํ ธมฺมวินเย อจฺฉริยา อพฺภุตา ธมฺมา, เย ทิสฺวา ทิสฺวา ภิกฺขู อิมสฺมิํ ธมฺมวินเย อภิรมนฺติ.\csromanspacer{} กตเม อฏฺฐ? เสยฺยถาปิ ภิกฺขเว มหาสมุทฺโท อนุปุพฺพนินฺโน อนุปุพฺพโปโณ อนุปุพฺพปพฺภาโร น อายตเกเนว ปปาโต.\csromanspacer{} เอวเมวํ โข ภิกฺขเว อิมสฺมิํ ธมฺมวินเย อนุปุพฺพสิกฺขา อนุปุพฺพกิริยา อนุปุพฺพปฏิปทา น อายตเกเนว อญฺญาปฏิเวโธ.\csromanspacer{} ยมฺปิ ภิกฺขเว อิมสฺมิํ ธมฺมวินเย อนุปุพฺพสิกฺขา อนุปุพฺพกิริยา อนุปุพฺพปฏิปทา น อายตเกเนว อญฺญาปฏิเวโธ.\csromanspacer{} อยํ ภิกฺขเว อิมสฺมิํ ธมฺมวินเย ปฐโม อจฺฉริโย อพฺภุโต ธมฺโม, ยํ ทิสฺวา ทิสฺวา ภิกฺขู อิมสฺมิํ ธมฺมวินเย อภิรมนฺติ ฯเปฯ.\csromanspacer{} เสยฺยถาปิ ภิกฺขเว มหาสมุทฺโท มหตํ ภูตานํ อาวาโส.\csromanspacer{} ตตฺริเม ภูตา, ติมิ ติมิงฺคโล ติมิรปิงฺคโล อสุรา นาคา คนฺธพฺพา, วสนฺติ มหาสมุทฺเท โยชนสติกาปิ อตฺตภาวา ฯเปฯ ปญฺจโยชนสติกาปิ อตฺตภาวา.\csromanspacer{} เอวเมวํ โข ภิกฺขเว อยํ ธมฺมวินโย มหตํ ภูตานํ อาวาโส.\csromanspacer{} ตตฺริเม ภูตา, โสตาปนฺโน, โสตาปตฺติผลสจฺฉิกิริยาย ปฏิปนฺโน ฯเปฯ อรหา, อรหตฺตาย ปฏิปนฺโน.\csromanspacer{} ยมฺปิ ภิกฺขเว อยํ ธมฺมวินโย มหตํ ภูตานํ อาวาโส.\csromanspacer{} ตตฺริเม ภูตา, โสตาปนฺโน, โสตปตฺติผลสจฺฉิกิริยาย ปฏิปนฺโน ฯเปฯ อรหา, อรหตฺตาย ปฏิปนฺโน.\csromanspacer{} อยํ ภิกฺขเว อิมสฺมิํ ธมฺมวินเย อฏฺฐโม อจฺฉริโย อพฺภุโต ธมฺโม, ยํ ทิสฺวา ทิสฺวา ภิกฺขู อิมสฺมิํ ธมฺมวินเย อภิรมนฺติ.\csromanspacer{} อิเม โข ภิกฺขเว อิมสฺมิํ ธมฺมวินเย อฏฺฐ อจฺฉริยา อพฺภุตา ธมฺมา, เยทิสฺวา ทิสฺวา ภิกฺขู อิมสฺมิํ ธมฺมวินเย อภิรมนฺตีติ.\csromanspacer{}\csromanspacer{}ทสมํ.}

\vspace{\tipitakaparskip}
\csromancenter{มหาวคฺโค ทุติโย.}
\chapterheadpageread{3. คหปติวคฺค \textbf{ตสฺสุทฺทานํ}}
\csromangathameasure{เวรญฺโช สีโห อาชญฺญํ, ขฬุงฺเกน มลานิ จ. \\ ทูเตยฺยํ เทฺว จ พนฺธนา, ปหาราโท อุโปสโถติ.}
\csromangathagroup{\csromangathastanza{\csromanfolio{47}เวรญฺโช สีโห อาชญฺญํ, ขฬุงฺเกน มลานิ จ.}\csromangathastanza{ทูเตยฺยํ เทฺว จ พนฺธนา, ปหาราโท อุโปสโถติ.}}

\chapterhead{3. คหปติวคฺค}
\csromanheader{2}{1. ปฐม-อุคฺคสุตฺต}
\proseitem{21}{เอกํ สมยํ ภควา เวสาลิยํ วิหรติ มหาวเน กูฏาคารสาลายํ.\csromanspacer{} ตตฺร โข ภควา ภิกฺขู อามนฺเตสิ “อฏฺฐหิ ภิกฺขเว อจฺฉริเยหิ อพฺภุเตหิ ธมฺเมหิ\footnote{อพฺภุตธมฺเมหิ (สฺยา, ก)} สมนฺนาคตํ อุคฺคํ คหปติํ เวสาลิกํ ธาเรถา”ติ.\csromanspacer{} อิทมโวจ ภควา, อิทํ วตฺวาน สุคโต อุฏฺฐายาสนา วิหารํ ปาวิสิ.}

\prose{อถ โข อญฺญตโร ภิกฺขุ ปุพฺพณฺหสมยํ นิวาเสตฺวา ปตฺตจีวรมาทาย เยน อุคฺคสฺส คหปติโน เวสาลิกสฺส นิเวสนํ เตนุปสงฺกมิ, อุปสงฺกมิตฺวา ปญฺญตฺเต อาสเน นิสีทิ.\csromanspacer{} อถ โข อุคฺโค คหปติ เวสาลิโก เยน โส ภิกฺขุ เตนุปสงฺกมิ, อุปสงฺกมิตฺวา ตํ ภิกฺขุํ อภิวาเทตฺวา เอกมนฺตํ นิสีทิ, เอกมนฺตํ นิสินฺนํ โข อุคฺคํ คหปติํ เวสาลิกํ โส ภิกฺขุ เอตทโวจ\csromandash{}}

\prose{อฏฺฐหิ โข ตฺวํ คหปติ อจฺฉริเยหิ อพฺภุเตหิ ธมฺเมหิ สมนฺนาคโต ภควตา พฺยากโต.\csromanspacer{} กตเม เต คหปติ อฏฺฐ อจฺฉริยา อพฺภุตา ธมฺมา, เยหิ ตฺวํ สมนฺนาคโต ภควตา พฺยากโตติ.\csromanspacer{} น โข อหํ ภนฺเต ชานามิ “กตเมหิ อฏฺฐหิ อจฺฉริเยหิ อพฺภุเตหิ ธมฺเมหิ สมนฺนาคโต ภควตา พฺยากโต”ติ.\csromanspacer{} อปิ จ ภนฺเต เย เม อฏฺฐ อจฺฉริยา อพฺภุตา ธมฺมา สํวิชฺชนฺติ, ตํ สุโณหิ สาธุกํ มนสิ กโรหิ, ภาสิสฺสามีติ. “เอวํ คหปตี”ติ โข โส ภิกฺขุ อุคฺคสฺส คหปติโน เวสาลิกสฺส ปจฺจสฺโสสิ.\csromanspacer{} อุคฺโค คหปติ เวสาลิโก เอตทโวจ “ยทาหํ ภนฺเต ภควนฺตํ ปฐมํ ทูรโตว อทฺทสํ, สห ทสฺสเนเนว เม \csromanfolio{48}ภนฺเต ภควโต จิตฺตํ ปสีทิ.\csromanspacer{} อยํ โข เม ภนฺเต ปฐโม อจฺฉริโย อพฺภุโต ธมฺโม สํวิชฺชติ.}

\prose{โส โข อหํ ภนฺเต ปสนฺนจิตฺโต ภควนฺตํ ปยิรุปาสิํ.\csromanspacer{} ตสฺส เม ภควา อนุปุพฺพิํ กถํ กเถสิ.\csromanspacer{} เสยฺยถิทํ, ทานกถํ สีลกถํ สคฺคกถํ กามานํ อาทีนวํ โอการํ สํกิเลสํ เนกฺขมฺเม อานิสํสํ ปกาเสสิ.\csromanspacer{} ยทา มํ ภควา อญฺญาสิ กลฺลจิตฺตํ มุทุจิตฺตํ วินีวรณจิตฺตํ อุทคฺคจิตฺตํ ปสนฺนจิตฺตํ, อถ ยา พุทฺธานํ สามุกฺกํสิกา ธมฺมเทสนา, ตํ ปกาเสสิ\csromandash{}ทุกฺขํ สมุทยํ นิโรธํ มคฺคํ.\csromanspacer{} เสยฺยถาปิ นาม สุทฺธํ วตฺถํ อปคตกาฬกํ สมฺมเทว รชนํ ปฏิคฺคเณฺหยฺย.\csromanspacer{} เอวเมวํ โข เม ตสฺมิํเยว อาสเน วิรชํ วีตมลํ ธมฺมจกฺขุํ อุทปาทิ “ยํ กิญฺจิ สมุทยธมฺมํ, สพฺพํ ตํ นิโรธธมฺมนฺ”ติ.\csromanspacer{} โส โข อหํ ภนฺเต ทิฏฺฐธมฺโม ปตฺตธมฺโม วิทิตธมฺโม ปริโยคาฬฺหธมฺโม ติณฺณวิจิกิจฺโฉ วิคตกถํกโถ เวสารชฺชปฺปตฺโต อปรปฺปจฺจโย สตฺถุสาสเน ตตฺเถว พุทฺธญฺจ ธมฺมญฺจ สํฆญฺจ สรณํ อคมาสิํ, พฺรหฺมจริยปญฺจมานิ จ สิกฺขาปทานิ สมาทิยิํ.\csromanspacer{} อยํ โข เม ภนฺเต ทุติโย อจฺฉริโย อพฺภุโต ธมฺโม สํวิชฺชติ.}

\prose{ตสฺส มยฺหํ ภนฺเต จตสฺโส โกมาริโย ปชาปติโย อเหสุํ.\csromanspacer{} อถ ขฺวาหํ ภนฺเต เยน ตา ปชาปติโย เตนุปสงฺกมิํ, อุปสงฺกมิตฺวา ตา ปชาปติโย เอตทวจํ “มยา โข ภคินิโย พฺรหฺมจริยปญฺจมานิ สิกฺขาปทานิ สมาทินฺนานิ\footnote{สมาทิณฺณานิ (สี, ก)}.\csromanspacer{} ยา อิจฺฉติ, สา อิเธว โภเค จ ภุญฺชตุ, ปุญฺญานิ จ กโรตุ, สกานิ วา ญาติกุลานิ คจฺฉตุ, โหติ วา ปน ปุริสาธิปฺปาโย ‘กสฺส โว ทมฺมี’ติ”.\csromanspacer{} เอวํ วุตฺเต สา ภนฺเต เชฏฺฐา ปชาปติ มํ เอตทโวจ “อิตฺถนฺนามสฺส มํ อยฺยปุตฺต ปุริสสฺส เทหี”ติ.\csromanspacer{} อถ โข อหํ ภนฺเต ตํ ปุริสํ ปกฺโกสาเปตฺวา วาเมน หตฺเถน ปชาปติํ คเหตฺวา ทกฺขิเณน หตฺเถน ภิงฺคารํ คเหตฺวา ตสฺส ปุริสสฺส โอโณเชสิํ.\csromanspacer{} โกมาริํ โข ปนาหํ ภนฺเต ทารํ ปริจฺจชนฺโต นาภิชานามิ จิตฺตสฺส อญฺญถตฺตํ.\csromanspacer{} อยํ โข เม ภนฺเต ตติโย อจฺฉริโย อพฺภุโต ธมฺโม สํวิชฺชติ.}

\prose{สํวิชฺชนฺติ โข ปน เม ภนฺเต กุเล โภคา, เต จ โข อปฺปฏิวิภตฺตา สีลวนฺเตหิ กลฺยาณธมฺเมหิ.\csromanspacer{} อยํ โข เม ภนฺเต จตุตฺโถ อจฺฉริโย อพฺภุโต ธมฺโม สํวิชฺชติ.}

\chapterheadpageread{3. คหปติวคฺค}
\prose{\csromanfolio{49}ยํ โข ปนาหํ ภนฺเต ภิกฺขุํ ปยิรุปาสามิ, สกฺกจฺจํเยว ปยิรุปาสามิ โน อสกฺกจฺจํ.\csromanspacer{} อยํ โข เม ภนฺเต ปญฺจโม อจฺฉริโย อพฺภุโต ธมฺโม สํวิชฺชติ.}

\prose{โส เจ ภนฺเต เม อายสฺมา ธมฺมํ เทเสติ, สกฺกจฺจํเยว สุโณมิ โน อสกฺกจฺจํ.\csromanspacer{} โน เจ เม โส อายสฺมา ธมฺมํ เทเสติ, อหมสฺส ธมฺมํ เทเสมิ.\csromanspacer{} อยํ โข เม ภนฺเต ฉฏฺโฐ อจฺฉริโย อพฺภุโต ธมฺโม สํวิชฺชติ.}

\prose{อนจฺฉริยํ โข ปน มํ ภนฺเต เทวตา อุปสงฺกมิตฺวา อาโรเจนฺติ “สฺวากฺขาโต คหปติ ภควตา ธมฺโม”ติ.\csromanspacer{} เอวํ วุตฺเต อหํ ภนฺเต ตา เทวตา เอวํ วทามิ “วเทยฺยาถ วา เอวํ โข ตุเมฺห เทวตา, โน วา วเทยฺยาถ, อถ โข สฺวากฺขาโต ภควตา ธมฺโม”ติ.\csromanspacer{} น โข ปนาหํ ภนฺเต อภิชานามิ “ตโตนิทานํ จิตฺตสฺส อุนฺนติํ\footnote{อุณฺณติํ (ก) อภิ1. 225; อภิ 2. 363 ปิฏฺเฐสุปิ ปสฺสิตพฺพํ.} ‘มํ วา เทวตา อุปสงฺกมนฺติ, อหํ วา เทวตาหิ สทฺธิํ สลฺลปามี’ติ”.\csromanspacer{} อยํ โข เม ภนฺเต สตฺตโม อจฺฉริโย อพฺภุโต ธมฺโม สํวิชฺชติ.}

\prose{ยานิมานิ ภนฺเต ภควตา เทสิตานิ ปญฺโจรมฺภาคิยานิ สํโยชนานิ, นาหํ เตสํ กิญฺจิ อตฺตนิ อปฺปหีนํ สมนุปสฺสามิ.\csromanspacer{} อยํ โข เม ภนฺเต อฏฺฐโม อจฺฉริโย อพฺภุโต ธมฺโม สํวิชฺชติ.\csromanspacer{} อิเม โข เม ภนฺเต อฏฺฐ อจฺฉริยา อพฺภุตา ธมฺมา สํวิชฺชนฺติ.\csromanspacer{} น จ โข อหํ ชานามิ “กตเมหิ จาหํ\footnote{กตเมหิปหํ (สี), กตเมหิปาหํ (อิ, ก)} อฏฺฐหิ อจฺฉริเยหิ อพฺภุเตหิ ธมฺเมหิ สมนฺนาคโต ภควตา พฺยากโต”ติ.}

\prose{อถ โข โส ภิกฺขุ อุคฺคสฺส คหปติโน เวสาลิกสฺส นิเวสเน ปิณฺฑปาตํ คเหตฺวา อุฏฺฐายาสนา ปกฺกามิ.\csromanspacer{} อถ โข โส ภิกฺขุ ปจฺฉาภตฺตํ ปิณฺฑปาตปฏิกฺกนฺโต เยน ภควา เตนุปสงฺกมิ, อุปสงฺกมิตฺวา ภควนฺตํ อภิวาเทตฺวา เอกมนฺตํ นิสีทิ, เอกมนฺตํ นิสินฺโน โข โส ภิกฺขุ ยาวตโก อโหสิ อุคฺเคน คหปตินา เวสาลิเกน สทฺธิํ กถาสลฺลาโป, ตํ สพฺพํ ภควโต อาโรเจสิ.}

\prose{สาธุ สาธุ ภิกฺขุ, ยถา ตํ อุคฺโค คหปติ เวสาลิโก สมฺมา พฺยากรมาโน พฺยากเรยฺย.\csromanspacer{} อิเมเหว โข ภิกฺขุ อฏฺฐหิ อจฺฉริเยหิ}

\prose{\csromanfolio{50}อพฺภุเตหิ ธมฺเมหิ สมนฺนาคโต อุคฺโค คหปติ เวสาลิโก มยา พฺยากโต.\csromanspacer{} อิเมหิ จ ปน ภิกฺขุ อฏฺฐหิ อจฺฉริเยหิ อพฺภุเตหิ ธมฺเมหิ สมนฺนาคตํ อุคฺคํ คหปติํ เวสาลิกํ ธาเรหีติ.\csromanspacer{}\csromanspacer{}ปฐมํ.}

\csromanheader{2}{2. ทุติย-อุคฺคสุตฺต}
\proseitem{22}{เอกํ สมยํ ภควา วชฺชีสุ วิหรติ หตฺถิคาเม.\csromanspacer{} ตตฺร โข ภควา ภิกฺขู อามนฺเตสิ “อฏฺฐหิ ภิกฺขเว อจฺฉริเยหิ อพฺภุเตหิ ธมฺเมหิ สมนฺนาคตํ อุคฺคํ คหปติํ หตฺถิคามกํ ธาเรถา”ติ.\csromanspacer{} อิทมโวจ ภควา, อิทํ วตฺวาน สุคโต อุฏฺฐายาสนา วิหารํ ปาวิสิ.}

\prose{อถ โข อญฺญตโร ภิกฺขุ ปุพฺพณฺหสมยํ นิวาเสตฺวา ปตฺตจีวรมาทาย เยน อุคฺคสฺส คหปติโน หตฺถิคามกสฺส นิเวสนํ เตนุปสงฺกมิ, อุปสงฺกมิตฺวา ปญฺญตฺเต อาสเน นิสีทิ.\csromanspacer{} อถ โข อุคฺโค คหปติ หตฺถิคามโก เยน โส ภิกฺขุ เตนุปสงฺกมิ, อุปสงฺกมิตฺวา ตํ ภิกฺขุํ อภิวาเทตฺวา เอกมนฺตํ นิสีทิ, เอกมนฺตํ นิสินฺนํ โข อุคฺคํ คหปติํ หตฺถิคามกํ โส ภิกฺขุ เอตทโวจ “อฏฺฐหิ โข ตฺวํ คหปติ อจฺฉริเยหิ อพฺภุเตหิ ธมฺเมหิ สมนฺนาคโต ภควตา พฺยากโต.\csromanspacer{} กตเม เต คหปติ อฏฺฐ อจฺฉริยา อพฺภุตา ธมฺมา, เยหิ ตฺวํ สมนฺนาคโต ภควตา พฺยากโต”ติ.}

\prose{น โข อหํ ภนฺเต ชานามิ “กตเมหิ อฏฺฐหิ อจฺฉริเยหิ อพฺภุเตหิ ธมฺเมหิ สมนฺนาคโต ภควตา พฺยากโต”ติ.\csromanspacer{} อปิ จ ภนฺเต เย เม อฏฺฐ อจฺฉริยา อพฺภุตา ธมฺมา สํวิชฺชนฺติ.\csromanspacer{} ตํ สุณาหิ สาธุกํ มนสิ กโรหิ, ภาสิสฺสามี”ติ. “เอวํ คหปตี”ติ โข โส ภิกฺขุ อุคฺคสฺส คหปติโน หตฺถิคามกสฺส ปจฺจสฺโสสิ.\csromanspacer{} อุคฺโค คหปติ หตฺถิคามโก เอตทโวจ “ยทาหํ ภนฺเต นาควเน ปริจรนฺโต ภควนฺตํ ปฐมํ ทูรโตว อทฺทสํ, สห ทสฺสเนเนว เม ภนฺเต ภควโต จิตฺตํ ปสีทิ, สุรามโท จ ปหียิ.\csromanspacer{} อยํ โข เม ภนฺเต ปฐโม อจฺฉริโย อพฺภุโต ธมฺโม สํวิชฺชติ.}

\prose{โส โข อหํ ภนฺเต ปสนฺนจิตฺโต ภควนฺตํ ปยิรุปาสิํ, ตสฺส เม ภควา อนุปุพฺพิํ กถํ กเถสิ.\csromanspacer{} เสยฺยถิทํ, ทานกถํ สีลกถํ กามานํ อาทีนวํ โอการํ สํกิเลสํ เนกฺขมฺเม อานิสํสํ ปกาเสสิ.\csromanspacer{} ยทา มํ ภควา อญฺญาสิ กลฺลจิตฺตํ มุทุจิตฺตํ วินีวรณจิตฺตํ อุทคฺคจิตฺตํ ปสนฺนจิตฺตํ.\csromanspacer{} อถ}

\vspace{\tipitakaparskip}
\csromancenter{คหปติวคฺค ยา พุทฺธานํ สามุกฺกํสิกา ธมฺมเทสนา, ตํ ปกาเสสิ\csromandash{}ทุกฺขํ สมุทยํ นิโรธํ มคฺคํ.\csromanspacer{} เสยฺยถาปิ นาม สุทฺธํ วตฺถํ อปคตกาฬกํ สมฺมเทว รชนํ ปฏิคฺคเณฺหยฺย.\csromanspacer{} เอวเมวํ โข เม ตสฺมิํ เยว อาสเน วิรชํ วีตมลํ ธมฺมจกฺขุํ อุทปาทิ “ยํ กิญฺจิ สมุทยธมฺมํ, สพฺพํ ตํ นิโรธธมฺมนฺ”ติ.\csromanspacer{} โส โข อหํ ภนฺเต ทิฏฺฐธมฺโม ปตฺตธมฺโม วิทิตธมฺโม ปริโยคาฬฺหธมฺโม ติณฺณวิจิกิจฺโฉ วิคตกถํกโถ เวสารชฺชปฺปตฺโต อปรปฺปจฺจโย สตฺถุสาสเน ตตฺเถว พุทฺธญฺจ ธมฺมญฺจ สํฆญฺจ สรณํ อคมาสิํ.\csromanspacer{} พฺรหฺมจริยปญฺจมานิ จ สิกฺขาปทานิ สมาทิยิํ.\csromanspacer{} อยํ โข เม ภนฺเต ทุติโย อจฺฉริโย อพฺภุโต ธมฺโม สํวิชฺชติ.}
\prose{\csromanfolio{51}ตสฺส มยฺหํ ภนฺเต จตสฺโส โกมาริโย ปชาปติโย อเหสุํ.\csromanspacer{} อถ ขฺวาหํ ภนฺเต เยน ตา ปชาปติโย เตนุปสงฺกมิํ, อุปสงฺกมิตฺวา ตา ปชาปติโย เอตทวจํ “มยา โข ภคินิโย พฺรหฺมจริยปญฺจมานิ สิกฺขาปทานิ สมาทินฺนานิ.\csromanspacer{} ยา อิจฺฉติ, สา อิเธว โภเค จ ภุญฺจตุ, ปุญฺญานิ จ กโรตุ, สกานิ วา ญาติกุลานิ คจฺฉตุ, โหติ วาปน ปุริสาธิปฺปาโย, ‘กสฺส โว ทมฺมี’ติ”.\csromanspacer{} เอวํ วุตฺเต สา ภนฺเต เชฏฺฐา ปชาปติ มํ เอตทโวจ “อิตฺถนฺนามสฺส มํ อยฺยปุตฺต ปุริสสฺส เทหี”ติ.\csromanspacer{} อถ โข อหํ ภนฺเต ตํ ปุริสํ ปกฺโกสาเปตฺวา วาเมน หตฺเถน ปชาปติํ คเหตฺวา ทกฺขิเณน หตฺเถน ภิงฺคารํ คเหตฺวา ตสฺส ปุริสสฺส โอโณเชสิํ.\csromanspacer{} โกมาริํ โข ปนาหํ ภนฺเต ทารํ ปริจฺจชนฺโต นาภิชานามิ จิตฺตสฺส อญฺญถตฺถํ.\csromanspacer{} อยํ โข เม ภนฺเต ตติโย อจฺฉริโย อพฺภุโต ธมฺโม.}

\prose{สํวิชฺชนฺติ โข ปน เม ภนฺเต กุเล โภคา, เต จ โข อปฺปฏิวิภตฺตา สีลวนฺเตหิ กลฺยาณธมฺเมหิ.\csromanspacer{} อยํ โข เม ภนฺเต จตุตฺโถ อจฺฉริโย อพฺภุโต ธมฺโม สํวิชฺชติ.}

\prose{ยํ โข ปนาหํ ภนฺเต ภิกฺขุ ปยิรุปาสามิ, สกฺกจฺจํเยว ปยิรุปาสามิ โน อสกฺกจฺจํ.\csromanspacer{} โส เจ เม อายสฺมา ธมฺมํ เทเสติ, สกฺกจฺจํเยว สุโณมิ โน อสกฺกจฺจํ.\csromanspacer{} โน เจ เม โส อายสฺมา ธมฺมํ เทเสติ, อหมสฺส ธมฺมํ เทเสมิ.\csromanspacer{} อยํ โข เม ภนฺเต ปญฺจโม อจฺฉริโย อพฺภุโต ธมฺโม สํวิชฺชติ.}

\prose{\csromanfolio{52}อนจฺฉริยํ โข ปน ภนฺเต สํเฆ นิมนฺติเต เทวตา อุปสงฺกมิตฺวา อาโรเจนฺติ “อสุโก คหปติ ภิกฺขุ อุภโตภาควิมุตฺโต, อสุโก ปญฺญาวิมุตฺโต, อสุโก กายสกฺขี, อสุโก ทิฏฺฐิปฺปตฺโต\footnote{ทิฏฺฐปฺปตฺโต (ก)}, อสุโก สทฺธาวิมุตฺโต, อสุโก ธมฺมานุสารี, อสุโก สทฺธานุสารี, อสุโก สีลวา กลฺยาณธมฺโม, อสุโก ทุสฺสีโล ปาปธมฺโม”ติ.\csromanspacer{} สํฆํ โข ปนาหํ ภนฺเต ปริวิสนฺโต นาภิชานามิ เอวํ จิตฺตํ อุปฺปาเทนฺโต “อิมสฺส วา โถกํ เทมิ, อิมสฺส วา พหุกนฺ”ติ.\csromanspacer{} อถ ขฺวาหํ ภนฺเต สมจิตฺโตว เทมิ.\csromanspacer{} อยํ โข เม ภนฺเต ฉฏฺโฐ อจฺฉริโย อพฺภุโต ธมฺโม สํวิชฺชติ.}

\prose{อนจฺฉริยํ โข ปน มํ ภนฺเต เทวตา อุปสงฺกมิตฺวา อาโรเจนฺติ “สฺวากฺขาโต คหปติ ภควตา ธมฺโม”ติ.\csromanspacer{} เอวํ วุตฺเต อหํ ภนฺเต ตา เทวตา เอวํ วเทมิ “วเทยฺยาถ วา เอวํ โข ตุเมฺห เทวตา, โน วา วเทยฺยาถ.\csromanspacer{} อถ โข สฺวากฺขาโต ภควตา ธมฺโม”ติ.\csromanspacer{} น โข ปนาหํ ภนฺเต อภิชานามิ ตโตนิทานํ จิตฺตสฺส อุนฺนติํ “มํ ตา เทวตา อุปสงฺกมนฺติ, อหํ วา เทวตาหิ สทฺธิํ สลฺลปามี”ติ.\csromanspacer{} อยํ โข เม ภนฺเต สตฺตโม อจฺฉริโย อพฺภุโต ธมฺโม สํวิชฺชติ.}

\prose{สเจ โข ปนาหํ ภนฺเต ภควโต ปฐมตรํ กาลํ กเรยฺยํ, อนจฺฉริยํ โข ปเนตํ.\csromanspacer{} ยํ มํ ภควา เอวํ พฺยากเรยฺย “นตฺถิ ตํ สํโยชนํ เยน สํยุตฺโต อุคฺโค คหปติ หตฺถิคามโก ปุน อิมํ โลกํ อาคจฺเฉยฺยา”ติ.\csromanspacer{} อยํ โข เม ภนฺเต อฏฺฐโม อจฺฉริโย อพฺภุโต ธมฺโม สํวิชฺชติ.\csromanspacer{} อิเม โข เม ภนฺเต อฏฺฐ อจฺฉริยา อพฺภุตา ธมฺมา สํวิชฺชนฺติ.\csromanspacer{} น จ โข อหํ ชานามิ “กตเมหิ จาหํ อฏฺฐหิ อจฺฉริเยหิ อพฺภุเตหิ ธมฺเมหิ สมนฺนาคโต ภควตา พฺยากโต”ติ.}

\prose{อถ โข โส ภิกฺขุ อุคฺคสฺส คหปติโน หตฺถิคามกสฺส นิเวสเน ปิณฺฑปาตํ คเหตฺวา อุฏฺฐายาสนา ปกฺกามิ.\csromanspacer{} อถ โข โส ภิกฺขุ ปจฺฉาภตฺตํ ปิณฺฑปาตปฏิกฺกนฺโต เยน ภควา เตนุปสงฺกมิ, อุปสงฺกมิตฺวา ภควนฺตํ อภิวาเทตฺวา เอกมนฺตํ นิสีทิ, เอกมนฺตํ นิสินฺโน โข โส ภิกฺขุ ยาวตโก อโหสิ อุคฺเคน คหปตินา หตฺถิคามเกน สทฺธิํ กถาสลฺลาโป, ตํ สพฺพํ ภควโต อาโรเจสิ.}

\chapterheadpageread{3. คหปติวคฺค}
\prose{\csromanfolio{53}สาธุ สาธุ ภิกฺขุ, ยถา ตํ อุคฺโค คหปติ หตฺถิคามโก สมฺมา พฺยากรมาโน พฺยากเรยฺย, อิเมเหว โข ภิกฺขุ อฏฺฐหิ อจฺฉริเยหิ อพฺภุเตหิ ธมฺเมหิ สมนฺนาคโต อุคฺโค คหปติ หตฺถิคามโก มยา พฺยากโต, อิเมหิ จ ปน ภิกฺขุ อฏฺฐหิ อจฺฉริเยหิ อพฺภุเตหิ ธมฺเมหิ สมนฺนาคตํ อุคฺคํ คหปติํ หตฺถิคามกํ ธาเรหีติ.\csromanspacer{}\csromanspacer{}ทุติยํ.}

\csromanheader{2}{3. ปฐมหตฺถกสุตฺต}
\proseitem{23}{เอกํ สมยํ ภควา อาฬวิยํ วิหรติ อคฺคาฬเว เจติเย.\csromanspacer{} ตตฺร โข ภควา ภิกฺขู อามนฺเตสิ\csromandash{}สตฺตหิ ภิกฺขเว อจฺฉริเยหิ อพฺภุเตหิ ธมฺเมหิ สมนฺนาคตํ หตฺถกํ อาฬวกํ ธาเรถ.\csromanspacer{} กตเมหิ สตฺตหิ? สทฺโธ หิ ภิกฺขเว หตฺถโก อาฬวโก, สีลวา ภิกฺขเว หตฺถโก อาฬวโก, หิรีมา ภิกฺขเว หตฺถโก อาฬวโก, โอตฺตปฺปี ภิกฺขเว หตฺถโก อาฬวโก, พหุสฺสุโต ภิกฺขเว หตฺถโก อาฬวโก, จาควา ภิกฺขเว หตฺถโก อาฬวโก, ปญฺญวา ภิกฺขเว หตฺถโก อาฬวโก.\csromanspacer{} อิเมหิ โข ภิกฺขเว สตฺตหิ อจฺฉริเยหิ อพฺภุเตหิ ธมฺเมหิ สมนฺนาคตํ หตฺถกํ อาฬวกํ ธาเรถาติ.\csromanspacer{} อิทมโวจ ภควา, อิทํ วตฺวาน สุคโต อุฏฺฐายาสนา วิหารํ ปาวิสิ.}

\prose{อถ โข อญฺญตโร ภิกฺขุ ปุพฺพณฺหสมยํ นิวาเสตฺวา ปตฺตจีวรมาทาย เยน หตฺถกสฺส อาฬวกสฺส นิเวสนํ เตนุปสงฺกมิ, อุปสงฺกมิตฺวา ปญฺญตฺเต อาสเน นิสีทิ.\csromanspacer{} อถ โข หตฺถโก อาฬวโก เยน โส ภิกฺขุ เตนุปสงฺกมิ, อุปสงฺกมิตฺวา ตํ ภิกฺขุํ อภิวาเทตฺวา เอกมนฺตํ นิสีทิ, เอกมนฺตํ นิสินฺนํ โข หตฺถกํ อาฬวกํ โส ภิกฺขุ เอตทโวจ\csromandash{}}

\prose{สตฺตหิ โข ตฺวํ อาวุโส อจฺฉริเยหิ อพฺภุเตหิ ธมฺเมหิ สมนฺนาคโต ภควตา พฺยากโต.\csromanspacer{} กตเมหิ สตฺตหิ? สทฺโธ ภิกฺขเว หตฺถโก อาฬวโก, สีลวา ฯเปฯ.\csromanspacer{} หิรีมา.\csromanspacer{} โอตฺตปฺปี.\csromanspacer{} พหุสฺสุโต.\csromanspacer{} จาควา, ปญฺญาวา ภิกฺขเว หตฺถโก อาฬวโกติ.\csromanspacer{} อิเมหิ โข ตฺวํ อาวุโส สตฺตหิ อจฺฉริเยหิ อพฺภุเตหิ ธมฺเมหิ สมนฺนาคโต ภควตา พฺยากโตติ.\csromanspacer{} กจฺจิตฺถ ภนฺเต น โกจิ คิหี อโหสิ โอทาตวสโนติ.\csromanspacer{} น เหตฺถ อาวุโส โกจิ คิหี อโหสิ โอทาตวสโนติ.\csromanspacer{} สาธุ ภนฺเต ยเทตฺถ น โกจิ คิหี อโหสิ โอทาตวสโนติ.}

\prose{\csromanfolio{54}อถ โข โส ภิกฺขุ หตฺถกสฺส อาฬวกสฺส นิเวสเน ปิณฺฑปาตํ คเหตฺวา อุฏฺฐายาสนา ปกฺกามิ.\csromanspacer{} อถ โข โส ภิกฺขุ ปจฺฉาภตฺตํ ปิณฺฑปาตปฏิกฺกนฺโต เยน ภควา เตนุปสงฺกมิ, อุปสงฺกมิตฺวา ภควนฺตํ อภิวาเทตฺวา เอกมนฺตํ นิสีทิ, เอกมนฺตํ นิสินฺโน โข โส ภิกฺขุ ภควนฺตํ เอตทโวจ\csromandash{}}

\prose{อิธาหํ ภนฺเต ปุพฺพณฺหสมยํ นิวาเสตฺวา ปตฺตจีวรมาทาย เยน หตฺถกสฺส อาฬวกสฺส นิเวสนํ เตนุปสงฺกมิํ, อุปสงฺกมิตฺวา ปญฺญตฺเต อาสเน นิสีทิํ.\csromanspacer{} อถ โข ภนฺเต หตฺถโก อาฬวโก เยนาหํ เตนุปสงฺกมิ, อุปสงฺกมิตฺวา มํ อภิวาเทตฺวา เอกมนฺตํ นิสีทิ.\csromanspacer{} เอกมนฺตํ นิสินฺนํ โข อหํ ภนฺเต หตฺถกํ อาฬวกํ เอตทวจํ “สตฺตหิ โข ตฺวํ อาวุโส อจฺฉริเยหิ อพฺภุเตหิ ธมฺเมหิ สมนฺนาคโต ภควตา พฺยากโต.\csromanspacer{} กตเมหิ สตฺตหิ สทฺโธ ภิกฺขเว หตฺถโก อาฬวโก, สีลวา ฯเปฯ.\csromanspacer{} หิรีมา.\csromanspacer{} โอตฺตปฺปี.\csromanspacer{} พหุสฺสุโต.\csromanspacer{} จาควา.\csromanspacer{} ปญฺญวา ภิกฺขเว หตฺถโก อาฬวโกติ.\csromanspacer{} อิเมหิ โข ตฺวํ อาวุโส สตฺตหิ อจฺฉริเยหิ อพฺภุเตหิ ธมฺเมหิ สมนฺนาคโต ภควตา พฺยากโต”ติ.}

\prose{เอวํ วุตฺเต ภนฺเต หตฺถโก มํ เอตทโวจ “กจฺจิตฺถ ภนฺเต น โกจิ คิหี อโหสิ โอทาตวสโน”ติ.\csromanspacer{} น เหตฺถ อาวุโส โกจิ คิหี อโหสิ โอทาตวสโนติ.\csromanspacer{} สาธุ ภนฺเต ยเทตฺถ น โกจิ คิหี อโหสิ โอทาตวสโนติ.}

\prose{สาธุ สาธุ ภิกฺขุ, อปฺปิจฺโฉ โส ภิกฺขุ กุลปุตฺโต สนฺเตเยว อตฺตนิ กุสลธมฺเม น อิจฺฉติ ปเรหิ ญายมาเน\footnote{ปญฺญาปยมาเน (ก)}.\csromanspacer{} เตน หิ ตฺวํ ภิกฺขุ อิมินาปิ อฏฺฐเมน อจฺฉริเยน อพฺภุเตน ธมฺเมน สมนฺนาคตํ หตฺถกํ อาฬวกํ ธาเรหิ, ยทิทํ อปฺปิจฺฉตายาติ.\csromanspacer{}\csromanspacer{}ตติยํ.}

\csromanheader{2}{4. ทุติยหตฺถกสุตฺต}
\proseitem{24}{เอกํ สมยํ ภควา อาฬวิยํ วิหรติ อคฺคาฬเว เจติเย.\csromanspacer{} อถ โข หตฺถโก อาฬวโก ปญฺจมตฺเตหิ อุปาสกสเตหิ ปริวุโต เยน ภควา เตนุปสงฺกมิ, อุปสงฺกมิตฺวา ภควนฺตํ อภิวาเทตฺวา เอกมนฺตํ นิสีทิ.\csromanspacer{} เอกมนฺตํ นิสินฺนํ โข หตฺถกํ อาฬวกํ}

\vspace{\tipitakaparskip}
\csromancenter{คหปติวคฺค ภควา เอตทโวจ “มหตี โข ตฺยายํ หตฺถก ปริสา, กถํ ปน ตฺวํ หตฺถก อิมํ มหติํ ปริสํ สงฺคณฺหาสี”ติ.\csromanspacer{} ยานิมานิ ภนฺเต ภควตา เทสิตานิ * จตฺตาริ สงฺคหวตฺถูนิ, เตหาหํ\footnote{เตนาหํ (สี)} อิมํ มหติํ ปริสํ สงฺคณฺหามิ, อหํ ภนฺเต ยํ ชานามิ “อยํ ทาเนน สงฺคเหตพฺโพ”ติ, ตํ ทาเนน สงฺคณฺหามิ, ยํ ชานามิ “อยํ เปยฺยวชฺเชน สงฺคเหตพฺโพ”ติ, ตํ เปยฺยวชฺเชน สงฺคณฺหามิ, ยํ ชานามิ “อยํ อตฺถจริยาย สงฺคเหตพฺโพ”ติ, ตํ อตฺถจริยาย สงฺคณฺหามิ, ยํ ชานามิ “อยํ สมานตฺตตาย สงฺคเหตพฺโพ”ติ, ตํ สมานตฺตตาย สงฺคณฺหามิ.\csromanspacer{} สํวิชฺชนฺติ โข ปน เม ภนฺเต กุเล โภคา, ทลิทฺทสฺส โข โน ตถา โสตพฺพํ มญฺญนฺตีติ.\csromanspacer{} สาธุ สาธุ หตฺถก, โยนิ โข ตฺยายํ หตฺถก มหติํ ปริสํ สงฺคเหตุํ.\csromanspacer{} เย หิ เกจิ หตฺถก อตีตมทฺธานํ มหติํ ปริสํ สงฺฆเหสุํ, สพฺเพเต อิเมเหว จตูหิ องฺคหวตฺถูหิ มหติํ ปริสํ สงฺคเหสุํ.\csromanspacer{} เยปิ หิ เกจิ หตฺถก อนาคตมทฺธานํ มหติํ ปริสํ สงฺคณฺหิสฺสนฺติ, สพฺเพเต อิเมเหว จตูหิ สงฺคหวตฺถูหิ มหติํ ปริสํ สงฺคณฺหิสฺสนฺติ.\csromanspacer{} เยปิ หิ เกจิ หตฺถก เอตรหิ มหติํ ปริสํ สงฺคณฺหนฺติ, สพฺเพเต อิเมเหว จตูหิ สงฺคหวตฺถูหิ มหติํ ปริสํ สงฺคณฺหนฺตีติ.}
\prose{\csromanfolio{55}อถ โข หตฺถโก อาฬวโก ภควตา ธมฺมิยา กถาย สนฺทสฺสิโต สมาทปิโต สมุตฺเตชิโต สมฺปหํสิโต อุฏฺฐายาสนา ภควนฺตํ อภิวาเทตฺวา ปทกฺขิณํ กตฺวา ปกฺกามิ.\csromanspacer{} อถ โข ภควา อจิรปกฺกนฺเต หตฺถเก อาฬวเก ภิกฺขู อามนฺเตสิ “อฏฺฐหิ ภิกฺขเว อจฺฉริเยหิ อพฺภุเตหิ ธมฺเมหิ สมนฺนาคตํ หตฺถกํ อาฬวกํ ธาเรถ.\csromanspacer{} กตเมหิ อฏฺฐหิ? สทฺโธ ภิกฺขเว หตฺถโก อาฬวโก, สีลวา ภิกฺขเว ฯเปฯ.\csromanspacer{} หิรีมา.\csromanspacer{} โอตฺตปฺปี.\csromanspacer{} พหุสฺสุโต.\csromanspacer{} จาควา.\csromanspacer{} ปญฺญวา ภิกฺขเว หตฺถโก อาฬวโก, อปฺปิจฺโฉ ภิกฺขเว หตฺถโก อาฬวโก.\csromanspacer{} อิเมหิ โข ภิกฺขเว อฏฺฐหิ อจฺฉริเยหิ อพฺภุเตหิ ธมฺเมหิ สมนฺนาคตํ หตฺถกํ อาฬวกํ ธาเรถาติ.\csromanspacer{}\csromanspacer{}จตุตฺถํ.}

\csromanheader{2}{5. มหานามสุตฺต}
\proseitem{25}{เอกํ สมยํ ภควา สกฺเกสุ วิหรติ กปิลวตฺถุสฺมิํ นิโคฺรธาราเม.\csromanspacer{} อถ โข มหานาโม สกฺโก เยน ภควา เตนุปสงฺกมิ, \csromanfolio{56}อุปสงฺกมิตฺวา ภควนฺตํ อภิวาเทตฺวา เอกมนฺตํ นิสีทิ, เอกมนฺตํ นิสินฺโน โข มหานาโม สกฺโก ภควนฺตํ เอตทโวจ “กิตฺตาวตา นุ โข ภนฺเต อุปาสโก โหตี”ติ.\csromanspacer{} ยโต โข มหานาม พุทฺธํ สรณํ คโต โหติ, ธมฺมํ สรณํ คโต โหติ, สํฆํ สรณํ คโต โหติ.\csromanspacer{} เอตฺตาวตา โข มหานาม อุปาสโก โหตีติ.}

\prose{“กิตฺตาวตา ปน ภนฺเต อุปาสโก สีลวา โหตี”ติ.\csromanspacer{} ยโต โข มหานาม อุปาสโก ปาณาติปาตา ปฏิวิรโต โหติ, อทินฺนาทานา ปฏิวิรโต โหติ, กาเมสุมิจฺฉาจารา ปฏิวิรโต โหติ, มุสาวาทา ปฏิวิรโต โหติ, สุราเมรยมชฺชปมาทฏฺฐานา ปฏิวิรโต โหติ.\csromanspacer{} เอตฺตาวตา โข มหานาม อุปาสโก สีลวา โหตีติ.}

\prose{“กิตฺตาวตา ปน ภนฺเต อุปาสโก อตฺตหิตาย ปฏิปนฺโน โหติ, โน ปรหิตายา”ติ.\csromanspacer{} ยโต โข มหานาม อุปาสโก อตฺตนาว สทฺธาสมฺปนฺโน โหติ, โน ปรํ สทฺธาสมฺปทาย สมาทเปติ\footnote{สมาทาเปติ (?)}.\csromanspacer{} อตฺตนาว สีลสมฺปนฺโน โหติ, โน ปรํ สีลสมฺปทาย สมาทเปติ.\csromanspacer{} อตฺตนาว จาคสมฺปนฺโน โหติ, โน ปรํ จาคสมฺปทาย สมาทเปติ.\csromanspacer{} อตฺตนาว ภิกฺขูนํ ทสฺสนกาโม โหติ, โน ปรํ ภิกฺขูนํ ทสฺสเน สมาทเปติ.\csromanspacer{} อตฺตนาว สทฺธมฺมํ โสตุกาโม โหติ, โน ปรํ สทฺธมฺมสฺสวเน สมาทเปติ.\csromanspacer{} อตฺตนาว สุตานํ ธมฺมานํ ธารณชาติโก โหติ, โน ปรํ ธมฺมธารณาย สมาทเปติ.\csromanspacer{} อตฺตนาว สุตานํ ธมฺมานํ อตฺถูปปริกฺขิตา โหติ, โน ปรํ อตฺถูปปริกฺขาย สมาทเปติ.\csromanspacer{} อตฺตนาว อตฺถมญฺญาย ธมฺมมญฺญาย ธมฺมานุธมฺมปฺปฏิปนฺโน โหติ, โน ปรํ ธมฺมานุธมฺมปฺปฏิ- ปตฺติยา สมาทเปติ.\csromanspacer{} เอตฺตาวตา โข มหานาม อุปาสโก อตฺตหิตาย ปฏิปนฺโน โหติ, โน ปรหิตายาติ.}

\prose{“กิตฺตาวตา ปน ภนฺเต อุปาสโก อตฺตหิตาย จ ปฏิปนฺโน โหติ ปรหิตาย จา”ติ.\csromanspacer{} ยโต โข มหานาม อุปาสโก อตฺตนา จ สทฺธาสมฺปนฺโน โหติ, ปรญฺจ สทฺธาสมฺปทาย สมาทเปติ.\csromanspacer{} อตฺตนา จ สีลสมฺปนฺโน โหติ, ปรญฺจ สีลสมฺปทาย สมาทเปติ.\csromanspacer{} อตฺตนา จ จาคสมฺปนฺโน โหติ, ปรญฺจ จาคสมฺปทาย สมาทเปติ.\csromanspacer{} อตฺตนา จ ภิกฺขูนํ ทสฺสนกาโม โหติ, ปรญฺจ ภิกฺขูนํ ทสฺสเน สมาทเปติ.\csromanspacer{} อตฺตนา จ}

\vspace{\tipitakaparskip}
\csromancenter{คหปติวคฺค สทฺธมฺมํ โสตุกาโม โหติ, ปรญฺจ สทฺธมฺมสฺสวเน สมาทเปติ.\csromanspacer{} อตฺตนา จ สุตานํ ธมฺมานํ ธารณชาติโก โหติ, ปรญฺจ ธมฺมธารณาย สมาทเปติ.\csromanspacer{} อตฺตนา จ สุตานํ ธมฺมานํ อตฺถูปปริกฺขิตา โหติ, ปรญฺจ อตฺถูปปริกฺขาย สมาทเปติ.\csromanspacer{} อตฺตนา จ อตฺถมญฺญาย ธมฺมมญฺญาย ธมฺมานุธมฺมปฺปฏิปนฺโน โหติ, ปรญฺจ ธมฺมานุธมฺมปฺปฏิปตฺติยา สมาทเปติ.\csromanspacer{} เอตฺตาวตา โข มหานาม อุปาสโก อตฺตหิตาย จ ปฏิปนฺโน โหติ ปรหิตาย จาติ.\csromanspacer{}\csromanspacer{}ปญฺจมํ.}
\csromanheader{2}{6. ชีวกสุตฺต}
\proseitem{26}{\csromanfolio{57}เอกํ สมยํ ภควา ราชคเห วิหรติ ชีวกมฺพวเน.\csromanspacer{} อถ โข ชีวโก โกมารภจฺโจ เยน ภควา เตนุปสงฺกมิ, อุปสงฺกมิตฺวา ภควนฺตํ อภิวาเทตฺวา เอกมนฺตํ นิสีทิ, เอกมนฺตํ นิสินฺโน โข ชีวโก โกมารภจฺโจ ภควนฺตํ เอตทโวจ “กิตฺตาวตา นุ โข ภนฺเต อุปาสโก โหตี”ติ.\csromanspacer{} ยโต โข ชีวก พุทฺธํ สรณํ คโต โหติ, ธมฺมํ สรณํ คโต โหติ, สํฆํ สรณํ คโต โหติ.\csromanspacer{} เอตฺตาวตา โข ชีวก อุปาสโก โหตีติ.}

\prose{“กิตฺตาวตา ปน ภนฺเต อุปาสโก สีลวา โหตี”ติ.\csromanspacer{} ยโต โข ชีวก อุปาสโก ปาณาติปาตา ปฏิวิรโต โหติ ฯเปฯ สุราเมรยมชฺชปมาทฏฺฐานา ปฏิวิรโต โหติ.\csromanspacer{} เอตฺตาวตา โข ชีวก อุปาสโก สีลวา โหตีติ.}

\prose{“กิตฺตาวตา ปน ภนฺเต อุปาสโก อตฺตหิตาย ปฏิปนฺโน โหติ, โน ปรหิตายา”ติ.\csromanspacer{} ยโต โข ชีวก อุปาสโก อตฺตนาว สทฺธาสมฺปนฺโน โหติ, โน ปรํ สทฺธาสมฺปทาย สมาทเปติ ฯเปฯ.\csromanspacer{} อตฺตนาว อตฺถมญฺญาย ธมฺมมญฺญาย ธมฺมานุธมฺมปฺปฏิปนฺโน โหติ, โน ปรํ ธมฺมานุธมฺมปฺปฏิปตฺติยา สมาทเปติ.\csromanspacer{} เอตฺตาวตา โข ชีวก อุปาสโก อตฺตหิตาย ปฏิปนฺโน โหติ, โน ปรหิตายาติ.}

\prose{“กิตฺตาวตา ปน ภนฺเต อุปาสโก อตฺตหิตาย จ ปฏิปนฺโน โหติ ปรหิตาย จา”ติ.\csromanspacer{} ยโต โข ชีวก อุปาสโก อตฺตนา จ สทฺธาสมฺปนฺโน โหติ, ปรญฺจ สทฺธาสมฺปทาย สมาทเปติ.\csromanspacer{} อตฺตนา จ สีลสมฺปนฺโน โหติ, ปรญฺจ สีลสมฺปทาย สมาทเปติ.\csromanspacer{} อตฺตนา จ \csromanfolio{58}จาคสมฺปนฺโน โหติ, ปรญฺจ จาคสมฺปทาย สมาทเปติ.\csromanspacer{} อตฺตนา จ ภิกฺขูนํ ทสฺสนกาโม โหติ, ปรญฺจ ภิกฺขูนํ ทสฺสเน สมาทเปติ.\csromanspacer{} อตฺตนา จ สทฺธมฺมํ โสตุกาโม โหติ, ปรญฺจ สทฺธมฺมสฺสวเน สมาทเปติ.\csromanspacer{} อตฺตนา จ สุตานํ ธมฺมานํ ธารณชาติโก โหติ, ปรญฺจ ธมฺมธารณาย สมาทเปติ.\csromanspacer{} อตฺตนา จ สุตานํ ธมฺมานํ อตฺถูปปริกฺขิตา โหติ, ปรญฺจ อตฺถูปปริกฺขาย สมาทเปติ.\csromanspacer{} อตฺตนา จ อตฺถมญฺญาย ธมฺมมญฺญาย ธมฺมานุธมฺมปฺปฏิปนฺโน โหติ, ปรญฺจ ธมฺมานุธมฺมปฺปฏิปตฺติยา สมาทเปติ.\csromanspacer{} เอตฺตาวตา โข ชีวก อุปาสโก อตฺตหิตาย จ ปฏิปนฺโน โหติ ปรหิตาย จาติ.\csromanspacer{}\csromanspacer{}ฉฏฺฐํ.}

\csromanheader{2}{7. ปฐมพลสุตฺต}
\proseitem{27}{อฏฺฐิมานิ ภิกฺขเว พลานิ.\csromanspacer{} กตมานิ อฏฺฐ? รุณฺณพลา ภิกฺขเว ทารกา, โกธพลา มาตุคามา, อาวุธพลา โจรา, อิสฺสริยพลา ราชาโน, อุชฺฌตฺติพลา พาลา, นิชฺฌตฺติพลา ปณฺฑิตา, ปฏิสงฺขานพลา พหุสฺสุตา, ขนฺติพลา สมณพฺราหฺมณา.\csromanspacer{} อิมานิ โข ภิกฺขเว อฏฺฐ พลานีติ.}

\prose{.สตฺตมํ.}

\csromanheader{2}{8. ทุติยพลสุตฺต}
\proseitem{28}{อถ โข อายสฺมา สาริปุตฺโต เยน ภควา เตนุปสงฺกมิ, อุปสงฺกมิตฺวา ภควนฺตํ อภิวาเทตฺวา เอกมนฺตํ นิสีทิ, เอกมนฺตํ นิสินฺนํ โข อายสฺมนฺตํ สาริปุตฺตํ ภควา เอตทโวจ “กติ นุ โข สาริปุตฺต ขีณาสวสฺส ภิกฺขุโน พลานิ, เยหิ พเลหิ สมนฺนาคโต ขีณาสโว ภิกฺขุ อาสวานํ ขยํ ปฏิชานาติ ‘ขีณา เม อาสวา’ติ”.\csromanspacer{} อฏฺฐ ภนฺเต ขีณาสวสฺส ภิกฺขุโน พลานิ, เยหิ พเลหิ สมนฺนาคโต ขีณาสโว ภิกฺขุ อาสวานํ ขยํ ปฏิชานาติ “ขีณา เม อาสวา”ติ.}

\prose{กตมานิ อฏฺฐ? * อิธ ภนฺเต ขีณาสวสฺส ภิกฺขุโน อนิจฺจโต สพฺเพ สงฺขารา ยถาภูตํ สมฺมปฺปญฺญาย สุทิฏฺฐา โหนฺติ.\csromanspacer{} ยมฺปิ ภนฺเต ขีณาสวสฺส ภิกฺขุโน อนิจฺจโต สพฺเพ สงฺขารา ยถาภูตํ สมฺมปฺปญฺญาย สุทิฏฺฐา โหนฺติ.\csromanspacer{} อิทมฺปิ ภนฺเต ขีณาสวสฺส ภิกฺขุโน พลํ โหติ, ยํ พลํ อาคมฺม ขีณาสโว ภิกฺขุ อาสวานํ ขยํ ปฏิชานาติ “ขีณา เม อาสวา”ติ. (1)}

\chapterheadpageread{3. คหปติวคฺค}
\prose{\csromanfolio{59}ปุน จปรํ ภนฺเต ขีณาสวสฺส ภิกฺขุโน องฺคารกาสูปมา กามา ยถาภูตํ สมฺมปฺปญฺญาย สุทิฏฺฐา โหนฺติ.\csromanspacer{} ยมฺปิ ภนฺเต ขีณาสวสฺส ภิกฺขุโน องฺคารกาสูปมา กามา ยถาภูตํ สมฺมปฺปญฺญาย สุทิฏฺฐา โหนฺติ.\csromanspacer{} อิทมฺปิ ภนฺเต ขีณาสวสฺส ภิกฺขุโน พลํ โหติ, ยํ พลํ อาคมฺม ขีณาสโว ภิกฺขุ อาสวานํ ขยํ ปฏิชานาติ “ขีณา เม อาสวา”ติ. (2)}

\prose{ปุน จปรํ ภนฺเต ขีณาสวสฺส ภิกฺขุโน วิเวกนินฺนํ จิตฺตํ โหติ วิเวกโปณํ วิเวกปพฺภารํ วิเวกฏฺฐํ เนกฺขมฺมาภิรตํ พฺยนฺติภูตํ สพฺพโส อาสวฏฺฐานิเยหิ ธมฺเมหิ.\csromanspacer{} ยมฺปิ ภนฺเต ขีณาสวสฺส ภิกฺขุโน วิเวกนินฺนํ จิตฺตํ โหติ วิเวกโปณํ วิเวกปพฺภารํ วิเวกฏฺฐํ เนกฺขมฺมาภิรตํ พฺยนฺติภูตํ สพฺพโส อาสวฏฺฐานิเยหิ ธมฺเมหิ.\csromanspacer{} อิทมฺปิ ภนฺเต ขีณาสวสฺส ภิกฺขุโน พลํ โหติ, ยํ พลํ อาคมฺม ขีณาสโว ภิกฺขุ อาสวานํ ขยํ ปฏิชานาติ “ขีณา เม อาสวา”ติ. (3)}

\prose{ปุน จปรํ ภนฺเต ขีณาสวสฺส ภิกฺขุโน จตฺตาโร สติปฏฺฐานา ภาวิตา โหนฺติ สุภาวิตา.\csromanspacer{} ยมฺปิ ภนฺเต ขีณาสวสฺส ภิกฺขุโน จตฺตาโร สติปฏฺฐานา ภาวิตา โหนฺติ สุภาวิตา.\csromanspacer{} อิทมฺปิ ภนฺเต ขีณาสวสฺส ภิกฺขุโน พลํ โหติ.\csromanspacer{} ยํ พลํ อาคมฺม ขีณาสโว ภิกฺขุ อาสวานํ ขยํ ปฏิชานาติ “ขีณา เม อาสวา”ติ. (4)}

\prose{ปุน จปรํ ภนฺเต ขีณาสวสฺส ภิกฺขุโน จตฺตาโร อิทฺธิปาทา ภาวิตา โหนฺติ สุภาวิตา ฯเปฯ ปญฺจินฺทฺริยานิ ภาวิตานิ โหนฺติ สุภาวิตานิ ฯเปฯ สตฺต โพชฺฌงฺคา ภาวิตา โหนฺติ สุภาวิตา ฯเปฯ อริโย อฏฺฐงฺคิโก มคฺโค ภาวิโต โหติ สุภาวิโต.\csromanspacer{} ยมฺปิ ภนฺเต ขีณาสวสฺส ภิกฺขุโน อริโย อฏฺฐงฺคิโก มคฺโค ภาวิโต โหติ สุภาวิโต.\csromanspacer{} อิทมฺปิ ภนฺเต ขีณาสวสฺส ภิกฺขุโน พลํ โหติ, ยํ พลํ อาคมฺม ขีณาสโว ภิกฺขุ อาสวานํ ขยํ ปฏิชานาติ “ขีณา เม อาสวา”ติ. (5-8)}

\prose{อิมานิ โข ภนฺเต อฏฺฐ ขีณาสวสฺส ภิกฺขุโน พลานิ, เยหิ พเลหิ สมนฺนาคโต ขีณาสโว ภิกฺขุ อาสวานํ ขยํ ปฏิชานาติ “ขีณา เม อาสวา”ติ.\csromanspacer{}\csromanspacer{}อฏฺฐมํ.}

\csromanheader{2}{9. อกฺขณสุตฺต}
\proseitem{29}{\csromanfolio{60}“ขณกิจฺโจ โลโก ขณกิจฺโจ โลโก”ติ ภิกฺขเว อสฺสุตวา ปุถุชฺชโน ภาสติ, โน จ โข โส ชานาติ ขณํ วา อกฺขณํ วา.\csromanspacer{} อฏฺฐิเม ภิกฺขเว อกฺขณา อสมยา พฺรหฺมจริยวาสาย.\csromanspacer{} กตเม อฏฺฐ? อิธ ภิกฺขเว ตถาคโต จ โลเก อุปฺปนฺโน โหติ อรหํ สมฺมาสมฺพุทฺโธ วิชฺชาจรณสมฺปนฺโน สุคโต โลกวิทู อนุตฺตโร ปุริสทมฺมสารถิ สตฺถา เทวมนุสฺสานํ พุทฺโธ ภควา.\csromanspacer{} ธมฺโม จ เทสิยติ โอปสมิโก ปรินิพฺพานิโก สมฺโพธคามี สุคตปฺปเวทิโต, อยญฺจ ปุคฺคโล นิรยํ อุปปนฺโน โหติ.\csromanspacer{} อยํ ภิกฺขเว ปฐโม อกฺขโณ อสมโย พฺรหฺมจริยวาสาย. (1)}

\prose{ปุน จปรํ ภิกฺขเว ตถาคโต จ โลเก อุปฺปนฺโน โหติ ฯเปฯ สตฺถา เทวมนุสฺสานํ พุทฺโธ ภควา.\csromanspacer{} ธมฺโม จ เทสิยติ โอปสมิโก ปรินิพฺพานิโก สมฺโพธคามี สุคตปฺปเวทิโต, อยญฺจ ปุคฺคโล ติรจฺฉานโยนิํ อุปปนฺโน โหติ ฯเปฯ. (2)}

\prose{ปุน จปรํ ภิกฺขเว ฯเปฯ อยญฺจ ปุคฺคโล เปตฺติวิสยํ อุปปนฺโนโหติ ฯเปฯ. (3)}

\prose{ปุน จปรํ ภิกฺขเว ฯเปฯ อยญฺจ ปุคฺคโล อญฺญตรํ ทีฆายุกํ เทวนิกายํ อุปปนฺโน โหติ ฯเปฯ. (4)}

\prose{ปุน จปรํ ภิกฺขเว ฯเปฯ อยญฺจ ปุคฺคโล ปจฺจนฺติเมสุ ชนปเทสุ ปจฺจาชาโต โหติ, โส จ โหติ อวิญฺญาตาเรสุ มิลกฺเขสุ\footnote{มิลกฺขูสุ (สฺยา, ก) ที 3. 249 ปิฏฺเฐปิ.}.\csromanspacer{} ยตฺถ นตฺถิ คติ ภิกฺขูนํ ภิกฺขุนีนํ อุปาสกานํ อุปาสิกานํ ฯเปฯ. (5)}

\prose{ปุน จปรํ ภิกฺขเว ฯเปฯ อยญฺจ ปุคฺคโล มชฺฌิเมสุ ชนปเทสุ ปจฺจาชาโต โหติ, โส จ โหติ มิจฺฉาทิฏฺฐิโก วิปรีตทสฺสโน “นตฺถิ ทินฺนํ, นตฺถิ ยิฏฺฐํ, นตฺถิ หุตํ, นตฺถิ สุกตทุกฺกฏานํ กมฺมานํ ผลํ วิปาโก, นตฺถิ อยํ โลโก, นตฺถิ ปโร โลโก, นตฺถิ มาตา, นตฺถิ ปิตา, นตฺถิ สตฺตา โอปปาติกา, นตฺถิ โลเก สมณพฺราหฺมณา สมฺมคฺคตา สมฺมา ปฏิปนฺนา, เย อิมญฺจ โลกํ ปรญฺจ โลกํ สยํ อภิญฺญา สจฺฉิกตฺวา ปเวเทนฺตีติ ฯเปฯ. (6)}

\chapterheadpageread{3. คหปติวคฺค}
\prose{\csromanfolio{61}ปุน จปรํ ภิกฺขเว ฯเปฯ อยญฺจ ปุคฺคโล มชฺฌิเมสุ ชนปเทสุ ปจฺจาชาโต โหติ, โส จ โหติ ทุปฺปญฺโญ ชโฬ เอฬมูโค อปฺปฏิพโล สุภาสิตทุพฺภาสิตสฺส อตฺถมญฺญาตุํ.\csromanspacer{} อยํ ภิกฺขเว สตฺตโม อกฺขโณ อสมโย พฺรหฺมจริยวาสาย. (7)}

\prose{ปุน จปรํ ภิกฺขเว ตถาคโต จ โลเก อนุปฺปนฺโน โหติ อรหํ สมฺมาสมฺพุทฺโธ ฯเปฯ สตฺถา เทวมนุสฺสานํ พุทฺโธ ภควา.\csromanspacer{} ธมฺโม จ น เทสิยติ โอปสมิโก ปรินิพฺพานิโก สมฺโพธคามี สุคตปฺปเวทิโต, อยญฺจ ปุคฺคโล มชฺฌิเมสุ ชนปเทสุ ปจฺจาชาโต โหติ, โส จ โหติ ปญฺญวา อชโฬ อเนฬมูโค ปฏิพโล สุภาสิตทุพฺภาสิตสฺส อตฺถมญฺญาตุํ.\csromanspacer{} อยํ ภิกฺขเว อฏฺฐโม อกฺขโณ อสมโย พฺรหฺมจริยวาสาย.\csromanspacer{} อิเม โข ภิกฺขเว อฏฺฐ อกฺขณา อสมยา พฺรหฺมจริยวาสาย. (8)}

\prose{เอโกว โข ภิกฺขเว ขโณ จ สมโย จ พฺรหฺมจริยวาสาย.\csromanspacer{} กตโม เอโก? อิธ ภิกฺขเว ตถาคโต จ โลเก อุปฺปนฺโน โหติ อรหํ สมฺมาสมฺพุทฺโธ วิชฺชาจรณสมฺปนฺโน สุคโต โลกวิทู อนุตฺตโร ปุริสทมฺมสารถิ สตฺถา เทวมนุสฺสานํ พุทฺโธ ภควา.\csromanspacer{} ธมฺโม จ เทสิยติ โอปสมิโก ปรินิพฺพานิโก สมฺโพธคามี สุคตปฺปเวทิโต, อยญฺจ ปุคฺคโล มชฺฌิเมสุ ชนปเทสุ ปจฺจาชาโต โหติ, โส จ โหติ ปญฺญวา อชโฬ อเนฬมูโค ปฏิพโล สุภาสิตทุพฺภาสิตสฺส อตฺถมญฺญาตุํ.\csromanspacer{} อยํ ภิกฺขเว เอโกว ขโณ จ สมโย จ พฺรหฺมจริยวาสายาติ. (1)}

\csromangathameasure{มนุสฺสลาภํ ลทฺธาน, สทฺธมฺเม สุปฺปเวทิเต. \\ เย ขณํ นาธิคจฺฉนฺติ, อตินาเมนฺติ เต ขณํ. \\ พหู หิ อกฺขณา วุตฺตา, มคฺคสฺส อนฺตรายิกา. \\ กทาจิ กรหจิ โลเก, อุปฺปชฺชนฺติ ตถาคตา. \\ ตยิทํ สมฺมุขีภูตํ, ยํ โลกสฺมิํ สุทุลฺลภํ. \\ มนุสฺสปฏิลาโภ จ, สทฺธมฺมสฺส จ เทสนา.}
\csromangathagroup{\csromangathastanza{มนุสฺสลาภํ\footnote{มนุสฺสโลกํ (สฺยา)} ลทฺธาน, สทฺธมฺเม สุปฺปเวทิเต. \\ เย ขณํ นาธิคจฺฉนฺติ, อตินาเมนฺติ เต ขณํ.}\csromangathastanza{พหู หิ อกฺขณา วุตฺตา, มคฺคสฺส อนฺตรายิกา. \\ กทาจิ กรหจิ โลเก, อุปฺปชฺชนฺติ ตถาคตา.}\csromangathastanza{ตยิทํ\footnote{ตสฺสิทํ (ก)} สมฺมุขีภูตํ, ยํ โลกสฺมิํ สุทุลฺลภํ. \\ มนุสฺสปฏิลาโภ จ, สทฺธมฺมสฺส จ เทสนา.}}

\prose{อลํ วายมิตุํ ตตฺถ, อตฺตกาเมน\footnote{อตฺถกาเมน (สี, สฺยา, ก)} ชนฺตุนา,}

\csromangathameasure{กถํ วิชญฺญา สทฺธมฺมํ, ขโณ เว มา อุปจฺจคา. \\ ขณาตีตา หิ โสจนฺติ, นิรยมฺหิ สมปฺปิตา. \\ อิธ เจ นํ วิราเธติ, สทฺธมฺมสฺส นิยามตํ. \\ วาณิโชว อตีตตฺโถ, จิรตฺตํ อนุตปิสฺสติ. \\ อวิชฺชานิวุโต โปโส, สทฺธมฺมํ อปราธิโก. \\ ชาติมรณสํสารํ, จิรํ ปจฺจนุโภสฺสติ. \\ เย จ ลทฺธา มนุสฺสตฺตํ, สทฺธมฺเม สุปฺปเวทิเต. \\ อกํสุ สตฺถุ วจนํ, กริสฺสนฺติ กโรนฺติ วา. \\ ขณํ ปจฺจวิทุํ โลเก, พฺรหฺมจริยํ อนุตฺตรํ. \\ เย มคฺคํ ปฏิปชฺชิํสุ, ตถาคตปฺปเวทิตํ. \\ เย สํวรา จกฺขุมตา, เทสิตาทิจฺจพนฺธุนา. \\ เตสุ คุตฺโต สทา สโต, วิหเร อนวสฺสุโต. \\ สพฺเพ อนุสเย เฉตฺวา, มารเธยฺยปรานุเค. \\ เต เว ปารงฺคตา โลเก, เย ปตฺตา อาสวกฺขยนฺติ. นวมํ.}
\csromangathagroup{\csromangathastanza{\csromanfolio{62}กถํ วิชญฺญา สทฺธมฺมํ, ขโณ เว\footnote{โว (สฺยา)} มา อุปจฺจคา. \\ ขณาตีตา หิ โสจนฺติ, นิรยมฺหิ สมปฺปิตา.}\csromangathastanza{อิธ เจ นํ วิราเธติ, สทฺธมฺมสฺส นิยามตํ\footnote{นิยามิตํ (สฺยา)}. \\ วาณิโชว อตีตตฺโถ, จิรตฺตํ\footnote{นิยามิตํ (สฺยา)} อนุตปิสฺสติ.}\csromangathastanza{อวิชฺชานิวุโต โปโส, สทฺธมฺมํ อปราธิโก. \\ ชาติมรณสํสารํ, จิรํ ปจฺจนุโภสฺสติ.}\csromangathastanza{เย จ ลทฺธา มนุสฺสตฺตํ, สทฺธมฺเม สุปฺปเวทิเต. \\ อกํสุ สตฺถุ วจนํ, กริสฺสนฺติ กโรนฺติ วา.}\csromangathastanza{ขณํ ปจฺจวิทุํ โลเก, พฺรหฺมจริยํ อนุตฺตรํ. \\ เย มคฺคํ ปฏิปชฺชิํสุ, ตถาคตปฺปเวทิตํ.}\csromangathastanza{เย สํวรา จกฺขุมตา, เทสิตาทิจฺจพนฺธุนา. \\ เตสุ\footnote{เตสํ(ก)} คุตฺโต สทา สโต, วิหเร อนวสฺสุโต.}\csromangathastanza{สพฺเพ อนุสเย เฉตฺวา, มารเธยฺยปรานุเค. \\ เต เว ปารงฺคตา\footnote{ปารคตา (สี, สฺยา, อิ)} โลเก, เย ปตฺตา อาสวกฺขยนฺติ.\csromanspacer{} นวมํ.}}

\csromanheader{2}{10. อนุรุทฺธมหาวิตกฺกสุตฺต}
\proseitem{30}{เอกํ สมยํ ภควา ภคฺเคสุ วิหรติ สุสุมารคิเร เภสกฬาวเน มิคทาเย.\csromanspacer{} เตน โข ปน สมเยน อายสฺมา อนุรุทฺโธ เจตีสุ วิหรติ ปาจีนวํสทาเย.\csromanspacer{} อถ โข อายสฺมโต อนุรุทฺธสฺส รโหคตสฺส ปฏิสลฺลีนสฺส เอวํ เจตโส ปริวิตกฺโก อุทปาทิ “อปฺปิจฺฉสฺสายํ ธมฺโม, นายํ ธมฺโม มหิจฺฉสฺส.\csromanspacer{} สนฺตุฏฺฐสฺสายํ ธมฺโม, นายํ ธมฺโม อสนฺตุฏฺฐสฺส.\csromanspacer{} ปวิวิตฺตสฺสายํ ธมฺโม, นายํ ธมฺโม สงฺคณิการามสฺส.\csromanspacer{} อารทฺธวีริยสฺสายํ ธมฺโม, นายํ ธมฺโม กุสีตสฺส.\csromanspacer{} อุปฏฺฐิตสฺสติสฺสายํ\footnote{อุปฏฺฐิตสติสฺสายํ (สี, สฺยา, อิ)} ธมฺโม, นายํ ธมฺโม มุฏฺฐสฺสติสฺส\footnote{มุฏฺฐสติสฺส (สี, สฺยา, อิ)}.\csromanspacer{} สมาหิตสฺสายํ ธมฺโม, นายํ ธมฺโม อสมาหิตสฺส.\csromanspacer{} ปญฺญวโต อยํ ธมฺโม, นายํ ธมฺโม ทุปฺปญฺญสฺสา”ติ.}

\chapterheadpageread{3. คหปติวคฺค}
\prose{\csromanfolio{63}อถ โข ภควา อายสฺมโต อนุรุทฺธสฺส เจตสา เจโตปริวิตกฺกมญฺญาย เสยฺยถาปิ นาม พลวา ปุริโส สมิญฺชิตํ วา พาหํ ปสาเรยฺย, ปสาริตํ วา พาหํ สมิญฺเชยฺย.\csromanspacer{} เอวเมวํ ภคฺเคสุ สุสุมารคิเร เภสกฬาวเน มิคทาเย อนฺตรหิโต เจตีสุ ปาจีนวํสทาเย อายสฺมโต อนุรุทฺธสฺส สมฺมุเข ปาตุรโหสิ.\csromanspacer{} นิสีทิ ภควา ปญฺญตฺเต อาสเน.\csromanspacer{} อายสฺมาปิ โข อนุรุทฺโธ ภควนฺตํ อภิวาเทตฺวา เอกมนฺตํ นิสีทิ.\csromanspacer{} เอกมนฺตํ นิสินฺนํ โข อายสฺมนฺตํ อนุรุทฺธํ ภควา เอตทโวจ\csromandash{}}

\prose{สาธุ สาธุ อนุรุทฺธ, สาธุ โข ตฺวํ อนุรุทฺธ (ยํ ตํ มหาปุริสวิตกฺกํ)\footnote{สตฺต มหาปุริสวิตกฺเก (สี, อิ) ที 3. 249 ปิฏฺเฐปิ.} วิตกฺเกสิ “อปฺปิจฺฉสฺสายํ ธมฺโม, นายํ ธมฺโม มหิจฺฉสฺส.\csromanspacer{} สนฺตุฏฺฐสฺสายํ ธมฺโม, นายํ ธมฺโม อสนฺตุฏฺฐสฺส.\csromanspacer{} ปวิวิตฺตสฺสายํ ธมฺโม, นายํ ธมฺโม สงฺคณิการามสฺส.\csromanspacer{} อารทฺธวีริยสฺสายํ ธมฺโม, นายํ ธมฺโม กุสีตสฺส.\csromanspacer{} อุปฏฺฐิตสฺสติสฺสายํ ธมฺโม, นายํ ธมฺโม มุฏฺฐสฺสติสฺส.\csromanspacer{} สมาหิตสฺสายํ ธมฺโม, นายํ ธมฺโม อสมาหิตสฺส.\csromanspacer{} ปญฺญวโต อยํ ธมฺโม, นายํ ธมฺโม ทุปฺปญฺญสฺสา”ติ.\csromanspacer{} เตน หิ ตฺวํ อนุรุทฺธ อิมมฺปิ อฏฺฐมํ มหาปุริสวิตกฺกํ วิตกฺเกหิ “นิปฺปปญฺจารามสฺสายํ ธมฺโม นิปฺปปญฺจรติโน, นายํ ธมฺโม ปปญฺจารามสฺส ปปญฺจรติโน”ติ.}

\prose{ยโต โข ตฺวํ อนุรุทฺธ อิเม อฏฺฐ มหาปุริสวิตกฺเก วิตกฺเกสฺสสิ, ตโต ตฺวํ อนุรุทฺธ ยาวเทว\footnote{ยาวเท (สํ 1. 412 ปิฏฺเฐ)} อากงฺขิสฺสสิ, วิวิจฺเจว กาเมหิ วิวิจฺจ อกุสเลหิ ธมฺเมหิ สวิตกฺกํ สวิจารํ วิเวกชํ ปีติสุขํ ปฐมํ ฌานํ อุปสมฺปชฺช วิหริสฺสสิ.}

\prose{ยโต โข ตฺวํ อนุรุทฺธ อิเม อฏฺฐ มหาปุริสวิตกฺเก วิตกฺเกสฺสสิ, ตโต ตฺวํ อนุรุทฺธ ยาวเทว อากงฺขิสฺสสิ, วิตกฺกวิจารานํ วูปสมา อชฺฌตฺตํ สมฺปสาทนํ เจตโส เอโกทิภาวํ อวิตกฺกํ อวิจารํ สมาธิชํ ปีติสุขํ ทุติยํ ฌานํ อุปสมฺปชฺช วิหริสฺสสิ.}

\prose{ยโต โข ตฺวํ อนุรุทฺธ อิเม อฏฺฐ มหาปุริสวิตกฺเก วิตกฺเกสฺสสิ, ตโต ตฺวํ อนุรุทฺธ ยาวเทว อากงฺขิสฺสสิ, ปีติยา จ วิราคา อุเปกฺขโก จ วิหริสฺสสิ, สโต จ สมฺปชาโน สุขญฺจ กาเยน}

\prose{\csromanfolio{64}ปฏิสํเวทิสฺสสิ, ยํ ตํ อริยา อาจิกฺขนฺติ “อุเปกฺขโก สติมา สุขวิหารี”ติ ตติยํ ฌานํ อุปสมฺปชฺช วิหริสฺสสิ.}

\prose{ยโต โข ตฺวํ อนุรุทฺธ อิเม อฏฺฐ มหาปุริสวิตกฺเก วิตกฺเกสฺสสิ, ตโต ตฺวํ อนุรุทฺธ ยาวเทว อากงฺขิสฺสสิ, สุขสฺส จ ปหานา ทุกฺขสฺส จ ปหานา ปุพฺเพว โสมนสฺสโทมนสฺสานํ อตฺถงฺคมา อทุกฺขมสุขํ อุเปกฺขาสติปาริสุทฺธิํ จตุตฺถํ ฌานํ อุปสมฺปชฺช วิหริสฺสสิ.}

\prose{ยโต โข ตฺวํ อนุรุทฺธ อิเม จ อฏฺฐ มหาปุริสวิตกฺเก วิตกฺเกสฺสสิ, อิเมสญฺจ จตุนฺนํ ฌานานํ อาภิเจตสิกานํ ทิฏฺฐธมฺมสุขวิหารานํ นิกามลาภี ภวิสฺสสิ อกิจฺฉลาภี อกสิรลาภี, ตโต ตุยฺหํ อนุรุทฺธ เสยฺยถาปิ นาม คหปติสฺส วา คหปติปุตฺตสฺส วา นานารตฺตานํ ทุสฺสานํ ทุสฺสกรณฺฑโก ปูโร.\csromanspacer{} เอวเมวํ เต ปํสุกูลจีวรํ ขายิสฺสติ สนฺตุฏฺฐสฺส วิหรโต รติยา อปริตสฺสาย ผาสุวิหาราย โอกฺกมนาย นิพฺพานสฺส.}

\prose{ยโต โข ตฺวํ อนุรุทฺธ อิเม จ อฏฺฐ มหาปุริสวิตกฺเก วิตกฺเกสฺสสิ, อิเมสญฺจ จตุนฺนํ ฌานานํ อาภิเจตสิกานํ ทิฏฺฐธมฺมสุขวิหารานํ นิกามลาภี ภวิสฺสสิ อกิจฺฉลาภี อกสิรลาภี, ตโต ตุยฺหํ อนุรุทฺธ เสยฺยถาปิ นาม คหปติสฺส วา คหปติปุตฺตสฺส วา สาลีนํ โอทโน วิจิตกาฬโก อเนกสูโป อเนกพฺยญฺชโน.\csromanspacer{} เอวเมวํ เต ปิณฺฑิยาโลปโภชนํ ขายิสฺสติ สนฺตุฏฺฐสฺส วิหรโต รติยา อปริตสฺสาย ผาสุวิหาราย โอกฺกมนาย นิพฺพานสฺส.}

\prose{ยโต โข ตฺวํ อนุรุทฺธ อิเม จ อฏฺฐ มหาปุริสวิตกฺเก วิตกฺเกสฺสสิ, อิเมสญฺจ จตุนฺนํ ฌานานํ อาภิเจตสิกานํ ทิฏฺฐธมฺมสุขวิหารานํ นิกามลาภี ภวิสฺสสิ อกิจฺฉลาภี อกสิรลาภี, ตโต ตุยฺหํ อนุรุทฺธ เสยฺยถาปิ นาม คหปติสฺส วา คหปติปุตฺตสฺส วา กูฏาคารํ อุลฺลิตฺตาวลิตฺตํ นิวาตํ ผุสิตคฺคฬํ ปิหิตวาตปานํ.\csromanspacer{} เอวเมวํ เต รุกฺขมูลเสนาสนํ ขายิสฺสติ สนฺตุฏฺฐสฺส วิหรโต รติยา อปิริตสฺสาย ผาสุวิหาราย โอกฺกมนาย นิพฺพานสฺส.}

\prose{ยโต โข ตฺวํ อนุรุทฺธ อิเม จ อฏฺฐ มหาปุริสวิตกฺเก วิตกฺเกสฺสสิ, อิเมสญฺจ จตุนฺนํ ฌานานํ อาภิเจตสิกานํ ทิฏฺฐธมฺมสุข- วิหารานํ นิกามลาภี ภวิสฺสสิ อกิจฺฉลาภี อกสิรลาภี, ตโต ตุยฺหํ}

\vspace{\tipitakaparskip}
\csromancenter{คหปติวคฺค อนุรุทฺธ เสยฺยถาปิ นาม คหปติสฺส วา คหปติปุตฺตสฺส วา ปลฺลงฺโก โคนกตฺถโต ปฏิกตฺถโต ปฏลิกตฺถโต กทลิมิคปวรปจฺจตฺถรโณ\footnote{กาทลิ...ปจฺจตฺถรโณ (สี)} ส-อุตฺตรจฺฉโท อุภโตโลหิตกูปธาโน.\csromanspacer{} เอวเมวํ เต ติณสนฺถารกสยนาสนํ ขายิสฺสติ สนฺตุฏฺฐสฺส วิหรโต รติยา อปริตสฺสาย ผาสุวิหาราย โอกฺกมนาย นิพฺพานสฺส.}
\prose{\csromanfolio{65}ยโต โข ตฺวํ อนุรุทฺธ อิเม จ อฏฺฐ มหาปุริสวิตกฺเก วิตกฺเกสฺสสิ, อิเมสญฺจ จตุนฺนํ ฌานานํ อาภิเจตสิกานํ ทิฏฺฐธมฺมสุขวิหารานํ นิกามลาภี ภวิสฺสสิ อกิจฺฉลาภี อกสิรลาภี, ตโต ตุยฺหํ อนุรุทฺธ เสยฺยถาปิ นาม คหปติสฺส วา คหปติปุตฺตสฺส วา นานาเภสชฺชานิ.\csromanspacer{} เสยฺยถิทํ, สปฺปิ นวนีตํ เตลํ มธุ ผาณิตํ.\csromanspacer{} เอวเมวํ เต ปูติมุตฺตเภสชฺชํ ขายิสฺสติ สนฺตุฏฺฐสฺส วิหรโต รติยา อปริตสฺสาย ผาสุวิหาราย โอกฺกมนาย นิพฺพานสฺส.\csromanspacer{} เตน หิ ตฺวํ อนุรุทฺธ อายติกมฺปิ วสฺสาวาสํ อิเธว เจตีสุ ปาจีนวํสทาเย วิหเรยฺยาสีติ. “เอวํ ภนฺเต”ติ โข อายสฺมา อนุรุทฺโธ ภควโต ปจฺจสฺโสสิ.}

\prose{อถ โข ภควา อายสฺมนฺตํ อนุรุทฺธํ อิมินา โอวาเทน โอวทิตฺวา เสยฺยถาปิ นาม พลวา ปุริโส สมิญฺชิตํ วา พาหํ ปสาเรยฺย, ปสาริตํ วา พาหํ สมิญฺเชยฺย.\csromanspacer{} เอวเมวํ เจตีสุ ปาจีนวํสทาเย อนฺตรหิโต ภคฺเคสุ สุสุมารคิเร เภสกฬาวเน มิคทาเย ปาตุรโหสีติ.\csromanspacer{} นิสีทิ ภควา ปญฺญตฺเต อาสเน, นิสชฺช โข ภควา ภิกฺขู อามนฺเตสิ อฏฺฐ โข ภิกฺขเว มหาปุริสวิตกฺเก เทเสสฺสามิ, ตํ สุณาถ ฯเปฯ.\csromanspacer{} กตเม จ ภิกฺขเว อฏฺฐ มหาปุริสวิตกฺกา? อปฺปิจฺฉสฺสายํ ภิกฺขเว ธมฺโม, นายํ ธมฺโม มหิจฺฉสฺส.\csromanspacer{} สนฺตุฏฺฐสฺสายํ ภิกฺขเว ธมฺโม, นายํ ธมฺโม อสนฺตุฏฺฐสฺส.\csromanspacer{} ปวิวิตฺตสฺสายํ ภิกฺขเว ธมฺโม, นายํ ธมฺโม สงฺคณิการามสฺส.\csromanspacer{} อารทฺธวีริยสฺสายํ ภิกฺขเว ธมฺโม, นายํ ธมฺโม กุสีตสฺส.\csromanspacer{} อุปฏฺฐิตสฺสติสฺสายํ ภิกฺขเว ธมฺโม, นายํ ธมฺโม มุฏฺฐสฺสติสฺส.\csromanspacer{} สมาหิตสฺสายํ ภิกฺขเว ธมฺโม, นายํ ธมฺโม อสมาหิตสฺส.\csromanspacer{} ปญฺญวโต อยํ ภิกฺขเว ธมฺโม, นายํ ธมฺโม ทุปฺปญฺญสฺส.\csromanspacer{} นิปฺปปญฺจารามสฺสายํ ภิกฺขเว ธมฺโม นิปฺปปญฺจรติโน, นายํ ธมฺโม ปปญฺจารามสฺส ปปญฺจรติโน.}

\prose{\csromanfolio{66}“อปฺปิจฺฉสฺสายํ ภิกฺขเว ธมฺโม, นายํ ธมฺโม มหิจฺฉสฺสา”ติ อิติ โข ปเนตํ วุตฺตํ, กิญฺเจตํ ปฏิจฺจ วุตฺตํ.\csromanspacer{} อิธ ภิกฺขเว ภิกฺขุ อปฺปิจฺโฉ สมาโน “อปฺปิจฺโฉติ มํ ชาเนยฺยุนฺ”ติ น อิจฺฉติ, สนฺตุฏฺโฐ สมาโน “สนฺตุฏฺโฐติ มํ ชาเนยฺยุนฺ”ติ น อิจฺฉติ, ปวิวิตฺโต สมาโน “ปวิวิตฺโตติ มํ ชาเนยฺยุนฺ”ติ น อิจฺฉติ, อารทฺธวีริโย สมาโน “อารทฺธวีริโยติ มํ ชาเนยฺยุนฺ”ติ น อิจฺฉติ, อุปฏฺฐิตสฺสติ สมาโน “อุปฏฺฐิตสฺสตีติ มํ ชาเนยฺยุนฺ”ติ น อิจฺฉติ, สมาหิโต สมาโน “สมาหิโตติ มํ ชาเนยฺยุนฺ”ติ น อิจฺฉติ, ปญฺญวา สมาโน “ปญฺญวาติ มํ ชาเนยฺยุนฺ”ติ น อิจฺฉติ, นิปฺปปญฺจาราโม สมาโน “นิปฺปปญฺจาราโมติ มํ ชาเนยฺยุนฺ”ติ น อิจฺฉติ, “อปฺปิจฺฉสฺสายํ ภิกฺขเว ธมฺโม, นายํ ธมฺโม มหิจฺฉสฺสา”ติ อิติ ยํ ตํ วุตฺตํ, อิทเมตํ ปฏิจฺจ วุตฺตํ. (1)}

\prose{“สนฺตุฏฺฐสฺสายํ ภิกฺขเว ธมฺโม, นายํ ธมฺโม อสนฺตุฏฺฐสฺสา”ติ อิติ โข ปเนตํ วุตฺตํ, กิญฺเจตํ ปฏิจฺจ วุตฺตํ.\csromanspacer{} อิธ ภิกฺขเว ภิกฺขุ สนฺตุฏฺโฐ โหติ อิตรีตร จีวร ปิณฺฑปาต เสนาสน คิลานปฺปจฺจย เภสชฺชปริกฺขาเรน. “สนฺตุฏฺฐสฺสายํ ภิกฺขเว ธมฺโม, นายํ ธมฺโม อสนฺตุฏฺฐสฺสา”ติ อิติ ยํ ตํ วุตฺตํ, อิทเมตํ ปฏิจฺจ วุตฺตํ. (2)}

\prose{“ปวิวิตฺตสฺสายํ ภิกฺขเว ธมฺโม, นายํ ธมฺโม สงฺคณิการามสฺสา”ติ อิติ โข ปเนตํ วุตฺตํ, กิญฺเจตํ ปฏิจฺจ วุตฺตํ.\csromanspacer{} อิธ ภิกฺขเว ภิกฺขุโน ปวิวิตฺตสฺส วิหรโต ภวนฺติ อุปสงฺกมิตาโร ภิกฺขู ภิกฺขุนิโย อุปาสกา อุปาสิกาโย ราชาโน ราชมหามตฺตา ติตฺถิยา ติตฺถิยสาวกา.\csromanspacer{} ตตฺร ภิกฺขุ วิเวกนินฺเนน จิตฺเตน วิเวกโปเณน วิเวกปพฺภาเรน วิเวกฏฺเฐน เนกฺขมฺมาภิรเตน อญฺญทตฺถุ อุยฺโยชนิกปฏิสํยุตฺตํเยว กถํ กตฺตา\footnote{ปวตฺตา (ก)} โหติ. “ปวิวิตฺตสฺสายํ ภิกฺขเว ธมฺโม, นายํ ธมฺโม สงฺคณิการามสฺสา”ติ อิติ ยํ ตํ วุตฺตํ, อิทเมตํ ปฏิจฺจ วุตฺตํ. (3)}

\prose{“อารทฺธวีริยสฺสายํ ภิกฺขเว ธมฺโม, นายํ ธมฺโม กุสีตสฺสา”ติ อิติ โข ปเนตํ วุตฺตํ, กิญฺเจตํ ปฏิจฺจ วุตฺตํ.\csromanspacer{} อิธ ภิกฺขเว ภิกฺขุ อารทฺธวีริโย วิหรติ อกุสลานํ ธมฺมานํ ปหานาย กุสลานํ ธมฺมานํ อุปสมฺปทาย, ถามวา ทฬฺหปรกฺกโม อนิกฺขิตฺตธุโร กุสเลสุ ธมฺเมสุ.}

\vspace{\tipitakaparskip}
\csromancenter{คหปติวคฺค “อารทฺธวีริยสฺสายํ ภิกฺขเว ธมฺโม, นายํ ธมฺโม กุสีตสฺสา”ติ อิติ ยํ ตํ วุตฺตํ, อิทเมตํ ปฏิจฺจ วุตฺตํ. (4)}
\prose{\csromanfolio{67}“อุปฏฺฐิตสฺสติสฺสายํ ภิกฺขเว ธมฺโม, นายํ ธมฺโม มุฏฺฐสฺสติสฺสา”ติ อิติ โข ปเนตํ วุตฺตํ, กิญฺเจตํ ปฏิจฺจ วุตฺตํ.\csromanspacer{} อิธ ภิกฺขเว ภิกฺขุ สติมา โหติ ปรเมน สติเนปกฺเกน สมนฺนาคโต จิรกตมฺปิ จิรภาสิตมฺปิ สริตา อนุสฺสริตา. “อุปฏฺฐิตสฺสติสฺสายํ ภิกฺขเว ธมฺโม, นายํ ธมฺโม มุฏฺฐสฺสติสฺสา”ติ อิติ ยํ ตํ วุตฺตํ, อิทเมตํ ปฏิจฺจ วุตฺตํ. (5)}

\prose{“สมาหิตสฺสายํ ภิกฺขเว ธมฺโม, นายํ ธมฺโม อสมาหิตสฺสา”ติ อิติ โข ปเนตํ วุตฺตํ, กิญฺเจตํ ปฏิจฺจ วุตฺตํ.\csromanspacer{} อิธ ภิกฺขเว ภิกฺขุ วิวิจฺเจว กาเมหิ ฯเปฯ จตุตฺถํ ฌานํ อุปสมฺปชฺช วิหรติ. “สมาหิตสฺสายํ ภิกฺขเว ธมฺโม, นายํ ธมฺโม อสมาหิตสฺสา”ติ อิติ ยํ ตํ วุตฺตํ, อิทเมตํ ปฏิจฺจ วุตฺตํ. (6)}

\prose{“ปญฺญวโต อยํ ภิกฺขเว ธมฺโม, นายํ ธมฺโม ทุปฺปญฺญสฺสา”ติ อิติ โข ปเนตํ วุตฺตํ, กิญฺเจตํ ปฏิจฺจ วุตฺตํ.\csromanspacer{} อิธ ภิกฺขเว ภิกฺขุ ปญฺญวา โหติ อุทยตฺถคามินิยา ปญฺญาย สมนฺนาคโต อริยาย นิพฺเพธิกาย สมฺมา ทุกฺขกฺขยคามินิยา. “ปญฺญวโต อยํ ภิกฺขเว ธมฺโม, นายํ ธมฺโม ทุปฺปญฺญสฺสา”ติ อิติ ยํ ตํ วุตฺตํ, อิทเมตํ ปฏิจฺจ วุตฺตํ. (7)}

\prose{“นิปฺปปญฺจารามสฺสายํ ภิกฺขเว ธมฺโม นิปฺปปญฺจรติโน, นายํ ธมฺโม ปปญฺจารามสฺส ปปญฺจรติโน”ติ อิติ โข ปเนตํ วุตฺตํ, กิญฺเจตํ ปฏิจฺจ วุตฺตํ.\csromanspacer{} อิธ ภิกฺขเว ภิกฺขุโน ปปญฺจนิโรเธ จิตฺตํ ปกฺขนฺทติ ปสีทติ สนฺติฏฺฐติ วิมุจฺจติ. “นิปฺปปญฺจารามสฺสายํ ภิกฺขเว ธมฺโม นิปฺปปญฺจรติโน, นายํ ธมฺโม ปปญฺจารามสฺส ปปญฺจรติโน”ติ อิติ ยํ ตํ วุตฺตํ, อิทเมตํ ปฏิจฺจ วุตฺตนฺติ. (8)}

\prose{อถ โข อายสฺมา อนุรุทฺโธ อายติกมฺปิ วสฺสาวาสํ ตตฺเถว เจตีสุ ปาจีนวํสทาเย วิหาสิ.\csromanspacer{} อถ โข อายสฺมา อนุรุทฺโธ เอโก วูปกฏฺโฐ อปฺปมตฺโต อาตาปี ปหิตตฺโต วิหรนฺโต นจิรสฺเสว ยสฺสตฺถาย กุลปุตฺตา สมฺมเทว อคารสฺมา อนคาริยํ ปพฺพชนฺติ, ตทนุตฺตรํ พฺรหฺมจริยปริโยสานํ ทิฏฺเฐว ธมฺเม สยํ อภิญฺญา สจฺฉิกตฺวา อุปสมฺปชฺช วิหาสิ, “ขีณา ชาติ, วุสิตํ พฺรหฺมจริยํ, กตํ กรณียํ, นาปรํ อิตฺถตฺตายา”ติ อพฺภญฺญาสิ.\csromanspacer{} อญฺญตโร จ ปนายสฺมา อนุรุทฺโธ \csromanfolio{68}อรหตํ อโหสีติ.\csromanspacer{} อถ โข อายสฺมา อนุรุทฺโธ อรหตฺตปฺปตฺโต ตายํ เวลายํ อิมา คาถาโย อภาสิ\csromandash{}}

\prose{\csromansymbolfootnote{*}{ขุ 2. 336 ปิฏฺเฐปิ.}“มม สงฺกปฺปมญฺญาย, สตฺถา โลเก อนุตฺตโร.}

\csromangathameasure{มโนมเยน กาเยน, อิทฺธิยา อุปสงฺกมิ. \\ ยถา เม อหุ สงฺกปฺโป, ตโต อุตฺตริ เทสยิ. \\ นิปฺปปญฺจรโต พุทฺโธ, นิปฺปปญฺจํ อเทสยิ. \\ ตสฺสาหํ ธมฺมมญฺญาย, วิหาสิํ สาสเน รโต. \\ ติสฺโส วิชฺชา อนุปฺปตฺตา, กตํ พุทฺธสฺส สาสนนฺ”ติ.ทสมํ. คหปติวคฺโค ตติโย.}
\csromangathagroup{\csromangathastanza{มโนมเยน กาเยน, อิทฺธิยา อุปสงฺกมิ. \\ ยถา เม อหุ สงฺกปฺโป, ตโต อุตฺตริ เทสยิ.}\csromangathastanza{นิปฺปปญฺจรโต พุทฺโธ, นิปฺปปญฺจํ อเทสยิ. \\ ตสฺสาหํ ธมฺมมญฺญาย, วิหาสิํ สาสเน รโต. \\ ติสฺโส วิชฺชา อนุปฺปตฺตา, กตํ พุทฺธสฺส สาสนนฺ”ติ.\csromanspacer{}\csromanspacer{}ทสมํ.\csromanspacer{} คหปติวคฺโค ตติโย.}}

\titlehead{\textbf{ตสฺสุทฺทานํ}}
\csromangathameasure{เทฺว อุคฺคา เทฺว จ หตฺถกา, มหานาเมน ชีวโก. \\ เทฺว พลา อกฺขณา วุตฺตา, อนุรุทฺเธน เต ทสาติ.}
\csromangathagroup{\csromangathastanza{เทฺว อุคฺคา เทฺว จ หตฺถกา, มหานาเมน ชีวโก. \\ เทฺว พลา อกฺขณา วุตฺตา, อนุรุทฺเธน เต ทสาติ.}}

\chapterhead{4. ทานวคฺค}
\csromanheader{2}{1. ปฐมทานสุตฺต}
\proseitem{31}{\csromansymbolfootnote{+}{ที 3. 214 ปิฏฺเฐปิ. ++ อภิ 4. 255 ปิฏฺเฐปิ.}อฏฺฐิมานิ ภิกฺขเว ทานานิ.\csromanspacer{} กตมานิ อฏฺฐ? อาสชฺช ทานํ เทติ, ภยา ทานํ เทติ, “อทาสิ เม”ติ ทานํ เทติ, “ทสฺสติ เม”ติ ทานํ เทติ, “สาหุ ทานนฺ”ติ ทานํ เทติ, “อหํ ปจามิ, อิเม น ปจนฺติ, นารหามิ ปจนฺโต อปจนฺตานํ ทานํ อทาตุนฺ”ติ ทานํ เทติ, “อิมํ เม ทานํ ททโต กลฺยาโณ กิตฺติสทฺโท อพฺภุคฺคจฺฉตี”ติ ทานํ เทติ จิตฺตาลงฺการจิตฺตปริกฺขารตฺถํ ทานํ เทติ.\csromanspacer{} อิมานิ โข ภิกฺขเว อฏฺฐ ทานานีติ.\csromanspacer{}\csromanspacer{}ปฐมํ.}

\csromanheader{2}{2. ทุติยทานสุตฺต}
\proseitem{32}{\csromansymbolmark{+}+ สทฺธา หิริยํ กุสลญฺจ ทานํ,}

\prose{ธมฺมา เอเต สปฺปุริสานุยาตา.}

\csromangathameasure{เอตํ หิ มคฺคํ ทิวิยํ วทนฺติ, \\ เอเตน หิ คจฺฉติ เทวโลกนฺติ.ทุติยํ.}
\csromangathagroup{\csromangathastanza{เอตํ หิ มคฺคํ ทิวิยํ วทนฺติ, \\ เอเตน หิ คจฺฉติ เทวโลกนฺติ.\csromanspacer{}\csromanspacer{}ทุติยํ.}}

\chapterheadpageread{4. ทานวคฺค \textbf{3. ทานวตฺถุสุตฺต}}
\proseitem{33}{\csromanfolio{69}อฏฺฐิมานิ ภิกฺขเว ทานวตฺถูนิ.\csromanspacer{} กตมานิ อฏฺฐ? ฉนฺทา ทานํ เทติ, โทสา ทานํ เทติ, โมหา ทานํ เทติ, ภยา ทานํ เทติ, “ทินฺนปุพฺพํ กตปุพฺพํ ปิตุปิตามเหหิ, นารหามิ โปราณํ กุลวํสํ หาเปตุนฺ”ติ ทานํ เทติ, “อิมาหํ ทานํ ทตฺวา กายสฺส เภทา ปรํ มรณา สุคติํ สคฺคํ โลกํ อุปปชฺชิสฺสามี”ติ ทานํ เทติ, “อิมํ เม ทานํ ททโต จิตฺตํ ปสีทติ, อตฺตมนตา โสมนสฺสํ อุปชายตี”ติ ทานํ เทติ, จิตฺตาลงฺการจิตฺตปริกฺขารตฺถํ ทานํ เทติ.\csromanspacer{} อิมานิ โข ภิกฺขเว อฏฺฐ ทานวตฺถูนีติ.\csromanspacer{}\csromanspacer{}ตติยํ.}

\csromanheader{2}{4. เขตฺตสุตฺต}
\proseitem{34}{อฏฺฐงฺคสมนฺนาคเต ภิกฺขเว เขตฺเต พีชํ วุตฺตํ น มหปฺผลํ โหติ น มหสฺสาทํ น ผาติเสยฺยํ\footnote{น ผาติเสยฺยนฺติ (สี, สฺยา, ก), น ผาติเสยฺยา (กตฺถจิ)}.\csromanspacer{} กถํ อฏฺฐงฺคสมนฺนาคเต? อิธ ภิกฺขเว เขตฺตํ อุนฺนามนินฺนามิ จ โหติ, ปาสาณสกฺขริกญฺจ โหติ, อูสรญฺจ โหติ, น จ คมฺภีรสิตํ โหติ, น อายสมฺปนฺนํ โหติ, น อปายสมฺปนฺนํ โหติ, น มาติกาสมฺปนฺนํ โหติ, น มริยาทสมฺปนฺนํ โหติ.\csromanspacer{} เอวํ อฏฺฐงฺคสมนฺนาคเต ภิกฺขเว เขตฺเต พีชํ วุตฺตํ น มหปฺผลํ โหติ น มหสฺสาทํ น ผาติเสยฺยํ.}

\prose{เอวเมวํ โข ภิกฺขเว อฏฺฐงฺคสมนฺนาคเตสุ สมณพฺราหฺมเณสุ ทานํ ทินฺนํ น มหปฺผลํ โหติ น มหานิสํสํ น มหาชุติกํ น มหาวิปฺผารํ.\csromanspacer{} กถํ อฏฺฐงฺคสมนฺนาคเตสุ? อิธ ภิกฺขเว สมณพฺราหฺมณา มิจฺฉาทิฏฺฐิกา โหนฺติ มิจฺฉาสงฺกปฺปา มิจฺฉาวาจา มิจฺฉากมฺมนฺตา มิจฺฉา-อาชีวา มิจฺฉาวายามา มิจฺฉาสติโน มิจฺฉาสมาธิโน.\csromanspacer{} เอวํ อฏฺฐงฺคสมนฺนาคเตสุ ภิกฺขเว สมณพฺราหฺมเณสุ ทานํ ทินฺนํ น มหปฺผลํ โหติ น มหานิสํสํ น มหาชุติกํ น มหาวิปฺผารํ.}

\prose{อฏฺฐงฺคสมนฺนาคเต ภิกฺขเว เขตฺเต พีชํ วุตฺตํ มหปฺผลํ โหติ มหสฺสาทํ ผาติเสยฺยํ.\csromanspacer{} กถํ อฏฺฐงฺคสมนฺนาคเต? อิธ ภิกฺขเว เขตฺตํ อนุนฺนามานินฺนามิ จ โหติ, อปาสาณสกฺขริกญฺจ โหติ, อนูสรญฺจ โหติ, คมฺภีรสิตํ โหติ, อายสมฺปนฺนํ โหติ, อปายสมฺปนฺนํ โหติ, \csromanfolio{70}มาติกาสมฺปนฺนํ โหติ, มริยาทสมฺปนฺนํ โหติ, เอวํ อฏฺฐงฺคสมนฺนาคเต ภิกฺขเว เขตฺเต พีชํ วุตฺตํ มหปฺผลํ โหติ มหสฺสาทํ ผาติเสยฺยํ.}

\prose{เอวเมวํ โข ภิกฺขเว อฏฺฐงฺคสมนฺนาคเตสุ สมณพฺราหฺมเณสุ ทานํ ทินฺนํ มหปฺผลํ โหติ มหานิสํสํ มหาชุติกํ มหาวิปฺผารํ.\csromanspacer{} กถํ อฏฺฐงฺคสมนฺนาคเตสุ? อิธ ภิกฺขเว สมณพฺราหฺมณา สมฺมาทิฏฺฐิกา โหนฺติ สมฺมาสงฺกปฺปา สมฺมาวาจา สมฺมากมฺมนฺตา สมฺมา-อาชีวา สมฺมาวายามา สมฺมาสติโน สมฺมาสมาธิโน.\csromanspacer{} เอวํ อฏฺฐงฺคสมนฺนาคเตสุ ภิกฺขเว สมณพฺราหฺมเณสุ ทานํ ทินฺนํ มหปฺผลํ โหติ มหานิสํสํ มหาชุติกํ มหาวิปฺผารนฺติ.}

\csromangathameasure{ยถาปิ เขตฺเต สมฺปนฺเน, ปวุตฺตา พีชสมฺปทา. \\ เทเว สมฺปาทยนฺตมฺหิ, โหติ ธญฺญสฺส สมฺปทา. \\ อนีติสมฺปทา โหติ, วิรูฬฺหี ภวติ สมฺปทา. \\ เวปุลฺลสมฺปทา โหติ, ผลํ เว โหติ สมฺปทา. \\ เอวํ สมฺปนฺนสีเลสุ, ทินฺนา โภชนสมฺปทา. \\ สมฺปทานํ อุปเนติ, สมฺปนฺนํ หิสฺส ตํ กตํ. \\ ตสฺมา สมฺปทมากงฺขี, สมฺปนฺนตฺถูธ ปุคฺคโล. \\ สมฺปนฺนปญฺเญ เสเวถ, เอวํ อิชฺฌนฺติ สมฺปทา. \\ วิชฺชาจรณสมฺปนฺเน, ลทฺธา จิตฺตสฺส สมฺปทํ. \\ กโรติ กมฺมสมฺปทํ, ลภติ จตฺถสมฺปทํ. \\ โลกํ ญตฺวา ยถาภูตํ, ปปฺปุยฺย ทิฏฺฐิสมฺปทํ. \\ มคฺคสมฺปทมาคมฺม, ยาติ สมฺปนฺนมานโส. \\ โอธุนิตฺวา มลํ สพฺพํ, ปตฺวา นิพฺพานสมฺปทํ. \\ มุจฺจติ สพฺพทุกฺเขหิ, สา โหติ สพฺพสมฺปทาติ.จตุตฺถํ.}
\csromangathagroup{\csromangathastanza{ยถาปิ เขตฺเต สมฺปนฺเน, ปวุตฺตา พีชสมฺปทา. \\ เทเว สมฺปาทยนฺตมฺหิ\footnote{สญฺชายนฺตมฺหิ (ก)}, โหติ ธญฺญสฺส สมฺปทา.}\csromangathastanza{อนีติสมฺปทา โหติ, วิรูฬฺหี ภวติ สมฺปทา. \\ เวปุลฺลสมฺปทา โหติ, ผลํ เว โหติ สมฺปทา.}\csromangathastanza{เอวํ สมฺปนฺนสีเลสุ, ทินฺนา โภชนสมฺปทา. \\ สมฺปทานํ อุปเนติ, สมฺปนฺนํ หิสฺส ตํ กตํ.}\csromangathastanza{ตสฺมา สมฺปทมากงฺขี, สมฺปนฺนตฺถูธ ปุคฺคโล. \\ สมฺปนฺนปญฺเญ เสเวถ, เอวํ อิชฺฌนฺติ สมฺปทา.}\csromangathastanza{วิชฺชาจรณสมฺปนฺเน, ลทฺธา จิตฺตสฺส สมฺปทํ. \\ กโรติ กมฺมสมฺปทํ, ลภติ จตฺถสมฺปทํ.}\csromangathastanza{โลกํ ญตฺวา ยถาภูตํ, ปปฺปุยฺย ทิฏฺฐิสมฺปทํ. \\ มคฺคสมฺปทมาคมฺม, ยาติ สมฺปนฺนมานโส.}\csromangathastanza{โอธุนิตฺวา มลํ สพฺพํ, ปตฺวา นิพฺพานสมฺปทํ. \\ มุจฺจติ สพฺพทุกฺเขหิ, สา โหติ สพฺพสมฺปทาติ.\csromanspacer{}\csromanspacer{}จตุตฺถํ.}}

\csromanheader{2}{5. ทานูปปตฺติสุตฺต}
\proseitem{35}{\csromansymbolfootnote{*}{ที 3. 214 ปิฏฺเฐปิ.}อฏฺฐิมา ภิกฺขเว ทานูปปตฺติโย.\csromanspacer{} กตมา อฏฺฐ? อิธ ภิกฺขเว เอกจฺโจ ทานํ เทติ สมณสฺส วา พฺราหฺมณสฺส วา อนฺนํ ปานํ วตฺถํ}

\vspace{\tipitakaparskip}
\csromancenter{ทานวคฺค ยานํ มาลาคนฺธวิเลปนํ เสยฺยาวสถปทีเปยฺยํ.\csromanspacer{} โส ยํ เทติ, ตํ ปจฺจาสีสติ\footnote{ปจฺจาสิํสติ (สี, สฺยา, กํ, อิ)}.\csromanspacer{} โส ปสฺสติ ขตฺติยมหาสาเล วา พฺราหฺมณมหาสาเล วา คหปติมหาสาเล วา ปญฺจหิ กามคุเณหิ สมปฺปิเต สมงฺคีภูเต ปริจารยมาเน, ตสฺส เอวํ โหติ “อโห วตาหํ กายสฺส เภทา ปรํ มรณา ขตฺติยมหาสาลานํ วา พฺราหฺมณมหาสาลานํ วา คหปติมหาสาลานํ วา สหพฺยตํ อุปปชฺเชยฺยนฺ”ติ.\csromanspacer{} โส ตํ จิตฺตํ ทหติ, ตํ จิตฺตํ อธิฏฺฐาติ, ตํ จิตฺตํ ภาเวติ.\csromanspacer{} ตสฺส ตํ จิตฺตํ หีเน วิมุตฺตํ\footnote{หีเนธิมุตฺตํ (สฺยา, อิ) วิมุตฺตนฺติ อธิมุตฺตํ, วิมุตฺตนฺติ วา วิสฺสฏฺฐํ (ฏีกาสํวณฺณนา)} อุตฺตริ อภาวิตํ ตตฺรูปปตฺติยา สํวตฺตติ, กายสฺส เภทา ปรํ มรณา ขตฺติยมหาสาลานํ วา พฺราหฺมณมหาสาลานํ วา คหปติมหาสาลานํ วา สหพฺยตํ อุปปชฺชติ.\csromanspacer{} ตญฺจ โข สีลวโต วทามิ, โน ทุสฺสีลสฺส.\csromanspacer{} อิชฺฌติ ภิกฺขเว สีลวโต เจโตปณิธิ วิสุทฺธตฺตา. (1)}
\prose{\csromanfolio{71}อิธ ปน ภิกฺขเว เอกจฺโจ ทานํ เทติ สมณสฺส วา พฺราหฺมณสฺส วา อนฺนํ ปานํ วตฺถํ ยานํ มาลาคนฺธวิเลปนํ เสยฺยาวสถปทีเปยฺยํ.\csromanspacer{} โส ยํ เทติ, ตํ ปจฺจาสีสติ.\csromanspacer{} ตสฺส สุตํ โหติ “จาตุมหาราชิกา\footnote{จาตุมฺมหาราชิกา (สี, สฺยา, กํ, อิ)} เทวา ทีฆายุกา วณฺณวนฺโต สุขพหุลา”ติ.\csromanspacer{} ตสฺส เอวํ โหติ “อโห วตาหํ กายสฺส เภทา ปรํ มรณา จาตุมหาราชิกานํ เทวานํ สหพฺยตํ อุปปชฺเชยฺยนฺ”ติ.\csromanspacer{} โส ตํ จิตฺตํ ทหติ, ตํ จิตฺตํ อธิฏฺฐาติ, ตํ จิตฺตํ ภาเวติ.\csromanspacer{} ตสฺส ตํ จิตฺตํ หีเน วิมุตฺตํ อุตฺตริ อภาวิตํ ตตฺรูปปตฺติยา สํวตฺตติ, กายสฺส เภทา ปรํ มรณา จาตุมหาราชิกานํ เทวานํ สหพฺยตํ อุปปชฺชติ.\csromanspacer{} ตญฺจ โข สีลวโต วทามิ, โน ทุสฺสีลสฺส.\csromanspacer{} อิชฺฌติ ภิกฺขเว สีลวโต เจโตปณิธิ วิสุทฺธตฺตา. (2)}

\prose{อิธ ปน ภิกฺขเว เอกจฺโจ ทานํ เทติ สมณสฺส วา พฺราหฺมณสฺส วา อนฺนํ ปานํ วตฺถํ ยานํ มาลาคนฺธวิเลปนํ เสยฺยาวสถปทีเปยฺยํ.\csromanspacer{} โส ยํ เทติ, ตํ ปจฺจาสีสติ.\csromanspacer{} ตสฺส สุตํ โหติ “ตาวติํสา เทวา ฯเปฯ ยามา เทวา.\csromanspacer{} ตุสิตา เทวา.\csromanspacer{} นิมฺมานรตี เทวา.\csromanspacer{} ปรนิมฺมิตวสวตฺตี เทวา ทีฆายุกา วณฺณวนฺโต สุขพหุลา”ติ.\csromanspacer{} ตสฺส เอวํ โหติ “อโห วตาหํ กายสฺส เภทา ปรํ มรณา ปรนิมฺมิตวสวตฺตีนํ เทวานํ สหพฺยตํ อุปปชฺเชยฺยนฺ”ติ.\csromanspacer{} โส ตํ จิตฺตํ ทหติ, ตํ จิตฺตํ \csromanfolio{72}อธิฏฺฐาติ, ตํ จิตฺตํ ภาเวติ.\csromanspacer{} ตสฺส ตํ จิตฺตํ หีเน วิมุตฺตํ อุตฺตริ อภาวิตํ ตตฺรูปปตฺติยา สํวตฺตติ, กายสฺส เภทา ปรํ มรณา ปรนิมฺมิตวสวตฺตีนํ เทวานํ สหพฺยตํ อุปปชฺชติ.\csromanspacer{} ตญฺจ โข สีลวโต วทามิ, โน ทุสฺสีลสฺส.\csromanspacer{} อิชฺฌติ ภิกฺขเว สีลวโต เจโตปณิธิ วิสุทฺธตฺตา. (3-7)}

\prose{อิธ ปน ภิกฺขเว เอกจฺโจ ทานํ เทติ สมณสฺส วา พฺราหฺมณสฺส วา อนฺนํ ปานํ วตฺถํ ยานํ มาลาคนฺธวิเลปนํ เสยฺยาวสถปทีเปยฺยํ.\csromanspacer{} โส ยํ เทติ, ตํ ปจฺจาสีสติ.\csromanspacer{} ตสฺส สุตํ โหติ “พฺรหฺมกายิกา เทวา ทีฆายุกา วณฺณวนฺโต สุขพหุลา”ติ.\csromanspacer{} ตสฺส เอวํ โหติ “อโห วตาหํ กายสฺส เภทา ปรํ มรณา พฺรหฺมกายิกานํ เทวานํ สหพฺยตํ อุปปชฺเชยฺยนฺ”ติ.\csromanspacer{} โส ตํ จิตฺตํ ทหติ, ตํ จิตฺตํ อธิฏฺฐาติ, ตํ จิตฺตํ ภาเวติ.\csromanspacer{} ตสฺส ตํ จิตฺตํ หีเน วิมุตฺตํ อุตฺตริ อภาวิตํ ตตฺรูปปตฺติยา สํวตฺตติ, กายสฺส เภทา ปรํ มรณา พฺรหฺมกายิกานํ เทวานํ สหพฺยตํ อุปปชฺชติ.\csromanspacer{} ตญฺจ โข สีลวโต วทามิ, โน ทุสฺสีลสฺส.\csromanspacer{} วีตราคสฺส, โน สราคสฺส.\csromanspacer{} อิชฺฌติ ภิกฺขเว สีลวโต เจโตปณิธิ วีตราคตฺตา.\csromanspacer{} อิมา โข ภิกฺขเว อฏฺฐ ทานูปปตฺติโยติ. (8).ปญฺจมํ.}

\csromanheader{2}{6. ปุญฺญกิริยวตฺถุสุตฺต}
\proseitem{36}{ตีณิมานิ ภิกฺขเว ปุญฺญกิริยวตฺถูนิ.\csromanspacer{} กตมานิ ตีณิ? ทานมยํ ปุญฺญกิริยวตฺถุ\footnote{ปุญฺญกิริยวตฺถุํ (สี, อิ) เอวมุปริปิ.}, สีลมยํ ปุญฺญกิริยวตฺถุ, ภาวนามยํ ปุญฺญกิริยวตฺถุ.\csromanspacer{} อิธ ภิกฺขเว เอกจฺจสฺส ทานมยํ ปุญฺญกิริยวตฺถุ ปริตฺตํ กตํ โหติ, สีลมยํ ปุญฺญกิริยวตฺถุ ปริตฺตํ กตํ โหติ, ภาวนามยํ ปุญฺญกิริยวตฺถุํ\footnote{ปุญฺญกิริยวตฺถุ (สฺยา)} นาภิสมฺโภติ.\csromanspacer{} โส กายสฺส เภทา ปรํ มรณา มนุสฺสโทภคฺยํ อุปปชฺชติ.}

\prose{อิธ ปน ภิกฺขเว เอกจฺจสฺส ทานมยํ ปุญฺญกิริยวตฺถุ มตฺตโส กตํ โหติ, สีลมยํ ปุญฺญกิริยวตฺถุ มตฺตโส กตํ โหติ, ภาวนามยํ ปุญฺญกิริยวตฺถุํ นาภิสมฺโภติ.\csromanspacer{} โส กายสฺส เภทา ปรํ มรณา มนุสฺสโสภคฺยํ อุปปชฺชติ.}

\chapterheadpageread{4. ทานวคฺค}
\prose{\csromanfolio{73}อิธ ปน ภิกฺขเว เอกจฺจสฺส ทานมยํ ปุญฺญกิริยวตฺถุ อธิมตฺตํ กตํ โหติ, สีลมยํ ปุญฺญกิริยวตฺถุ อธิมตฺตํ กตํ โหติ, ภาวนามยํ ปุญฺญกิริยวตฺถุํ นาภิสมฺโภติ.\csromanspacer{} โส กายสฺส เภทา ปรํ มรณา จาตุมหาราชิกานํ เทวานํ สหพฺยตํ อุปปชฺชติ.\csromanspacer{} ตตฺร ภิกฺขเว จตฺตาโร มหาราชาโน ทานมยํ ปุญฺญกิริยวตฺถุํ อติเรกํ กริตฺวา สีลมยํ ปุญฺญกิริยวตฺถุํ อติเรกํ กริตฺวา จาตุมหาราชิเก เทเว ทสหิ ฐาเนหิ อธิคณฺหนฺติ ทิพฺเพน อายุนา ทิพฺเพน วณฺเณน ทิพฺเพน สุเขน ทิพฺเพน ยเสน ทิพฺเพน อาธิปเตยฺเยน ทิพฺเพหิ รูเปหิ ทิพฺเพหิ สทฺเทหิ ทิพฺเพหิ คนฺเธหิ ทิพฺเพหิ รเสหิ ทิพฺเพหิ โผฏฺฐพฺเพหิ.}

\prose{อิธ ปน ภิกฺขเว เอกจฺจสฺส ทานมยํ ปุญฺญกิริยวตฺถุ อธิมตฺตํ กตํ โหติ, สีลมยํ ปุญฺญกิริยวตฺถุ อธิมตฺตํ กตํ โหติ, ภาวนามยํ ปุญฺญกิริยวตฺถุํ นาภิสมฺโภติ.\csromanspacer{} โส กายสฺส เภทา ปรํ มรณา ตาวติํสานํ เทวานํ สหพฺยตํ อุปปชฺชติ.\csromanspacer{} ตตฺร ภิกฺขเว สกฺโก เทวานมินฺโท ทานมยํ ปุญฺญกิริยวตฺถุํ อติเรกํ กริตฺวา สีลมยํ ปุญฺญกิริยวตฺถุํ อติเรกํ กริตฺวา ตาวติํเส เทเว ทสหิ ฐาเนหิ อธิคณฺหาติ ทิพฺเพน อายุนา ฯเปฯ ทิพฺเพหิ โผฏฺฐพฺเพหิ.}

\prose{อิธ ปน ภิกฺขเว เอกจฺจสฺส ทานมยํ ปุญฺญกิริยวตฺถุ อธิมตฺตํ กตํ โหติ, สีลมยํ ปุญฺญกิริยวตฺถุ อธิมตฺตํ กตํ โหติ, ภาวนามยํ ปุญฺญกิริยวตฺถุํ นาภิสมฺโภติ.\csromanspacer{} โส กายสฺส เภทา ปรํ มรณา ยามานํ เทวานํ สหพฺยตํ อุปปชฺชติ.\csromanspacer{} ตตฺร ภิกฺขเว สุยาโม เทวปุตฺโต ทานมยํ ปุญฺญกิริยวตฺถุํ อติเรกํ กริตฺวา สีลมยํ ปุญฺญกิริยวตฺถุํ อติเรกํ กริตฺวา ยาเม เทเว ทสหิ ฐาเนหิ อธิคณฺหาติ ทิพฺเพน อายุนา ฯเปฯ ทิพฺเพหิ โผฏฺฐพฺเพหิ.}

\prose{อิธ ปน ภิกฺขเว เอกจฺจสฺส ทานมยํ ปุญฺญกิริยวตฺถุ อธิมตฺตํ กตํ โหติ, สีลมยํ ปุญฺญกิริยวตฺถุ อธิมตฺตํ กตํ โหติ, ภาวนามยํ ปุญฺญกิริยวตฺถุํ นาภิสมฺโภติ.\csromanspacer{} โส กายสฺส เภทา ปรํ มรณา ตุสิตานํ เทวานํ สหพฺยตํ อุปปชฺชติ.\csromanspacer{} ตตฺร ภิกฺขเว สนฺตุสิโต เทวปุตฺโต ทานมยํ ปุญฺญกิริยวตฺถุํ อติเรกํ กริตฺวา สีลมยํ ปุญฺญกิริยวตฺถุํ อติเรกํ กริตฺวา ตุสิเต เทเว ทสหิ ฐาเนหิ อธิคณฺหาติ ทิพฺเพน อายุนา ฯเปฯ ทิพฺเพหิ โผฏฺฐพฺเพหิ.}

\prose{\csromanfolio{74}อิธ ปน ภิกฺขเว เอกจฺจสฺส ทานมยํ ปุญฺญกิริยวตฺถุ อธิมตฺตํ กตํ โหติ, สีลมยํ ปุญฺญกิริยวตฺถุ อธิมตฺตํ กตํ โหติ, ภาวนามยํ ปุญฺญกิริวตฺถุํ นาภิสมฺโภติ.\csromanspacer{} โส กายสฺส เภทา ปรํ มรณา นิมฺมานรตีนํ เทวานํ สหพฺยตํ อุปปชฺชติ.\csromanspacer{} ตตฺร ภิกฺขเว สุนิมฺมิโต เทวปุตฺโต ทานมยํ ปุญฺญกิริยวตฺถุํ อติเรกํ กริตฺวา สีลมยํ ปุญฺญกิริยวตฺถุํ อติเรกํ กริตฺวา นิมฺมานรตีเทเว ทสหิ ฐาเนหิ อธิคณฺหาติ ทิพฺเพน อายุนา ฯเปฯ ทิพฺเพหิ โผฏฺฐพฺเพหิ.}

\prose{อิธ ปน ภิกฺขเว เอกจฺจสฺส ทานมยํ ปุญฺญกิริยวตฺถุ อธิมตฺตํ กตํ โหติ, สีลมยํ ปุญฺญกิริยวตฺถุ อธิมตฺตํ กตํ โหติ, ภาวนามยํ ปุญฺญกิริยวตฺถุํ นาภิสมฺโภติ.\csromanspacer{} โส กายสฺส เภทา ปรํ มรณา ปรนิมฺมิตวสวตฺตีนํ เทวานํ สหพฺยตํ อุปปชฺชติ.\csromanspacer{} ตตฺร ภิกฺขเว วสวตฺตี เทวปุตฺโต ทานมยํ ปุญฺญกิริยวตฺถุํ อติเรกํ กริตฺวา สีลมยํ ปุญฺญกิริยวตฺถุํ อติเรกํ กริตฺวา ปรนิมฺมิตวสวตฺตีเทเว ทสหิ ฐาเนหิ อธิคณฺหาติ ทิพฺเพน อายุนา ทิพฺเพน วณฺเณน ทิพฺเพน สุเขน ทิพฺเพน ยเสน ทิพฺเพน อาธิปเตยฺเยน ทิพฺเพหิ รูเปหิ ทิพฺเพหิ สทฺเทหิ ทิพฺเพหิ คนฺเธหิ ทิพฺเพหิ รเสหิ ทิพฺเพหิ โผฏฺฐพฺเพหิ.\csromanspacer{} อิมานิ โข ภิกฺขเว ตีณิ ปุญฺญกิริยวตฺถูนีติ.\csromanspacer{}\csromanspacer{}ฉฏฺฐํ.}

\csromanheader{2}{7. สปฺปุริสทานสุตฺต}
\proseitem{37}{อฏฺฐิมานิ ภิกฺขเว สปฺปุริสทานานิ.\csromanspacer{} กตมานิ อฏฺฐ? สุจิํ เทติ, ปณีตํ เทติ, กาเลน เทติ, กปฺปิยํ เทติ, วิเจยฺย เทติ, อภิณฺหํ เทติ, ททํ จิตฺตํ ปสาเทติ, ทตฺวา อตฺตมโน โหติ.\csromanspacer{} อิมานิ โข ภิกฺขเว อฏฺฐ สปฺปุริสทานานีติ.}

\csromangathameasure{สุจิํ ปณีตํ กาเลน, กปฺปิยํ ปานโภชนํ. \\ อภิณฺหํ ททาติ ทานํ, สุเขตฺเตสุ พฺรหฺมจาริสุ. \\ เนว วิปฺปฏิสาริสฺส, จชิตฺวา อามิสํ พหุํ. \\ เอวํ ทินฺนานิ ทานานิ, วณฺณยนฺติ วิปสฺสิโน. \\ เอวํ ยชิตฺวา เมธาวี, สทฺโธ มุตฺเตน เจตสา. \\ อพฺยาพชฺฌํ สุขํ โลกํ, ปณฺฑิโต อุปปชฺชตีติ.สตฺตมํ.}
\csromangathagroup{\csromangathastanza{สุจิํ ปณีตํ กาเลน, กปฺปิยํ ปานโภชนํ. \\ อภิณฺหํ ททาติ ทานํ, สุเขตฺเตสุ\footnote{สุเขตฺเต (สี, อิ)} พฺรหฺมจาริสุ.}\csromangathastanza{เนว\footnote{น จ (สี, อิ)} วิปฺปฏิสาริสฺส, จชิตฺวา อามิสํ พหุํ. \\ เอวํ ทินฺนานิ ทานานิ, วณฺณยนฺติ วิปสฺสิโน.}\csromangathastanza{เอวํ ยชิตฺวา เมธาวี, สทฺโธ มุตฺเตน เจตสา. \\ อพฺยาพชฺฌํ\footnote{อพฺยาปชฺฌํ (ก) อํ 1. 353; อํ 2. 296 ปิฏฺเฐสุปิ.} สุขํ โลกํ, ปณฺฑิโต อุปปชฺชตีติ.\csromanspacer{}\csromanspacer{}สตฺตมํ.}}

\chapterheadpageread{4. ทานวคฺค \textbf{8. สปฺปุริสสุตฺต}}
\proseitem{38}{\csromanfolio{75}สปฺปุริโส ภิกฺขเว กุเล ชายมาโน พหุโน ชนสฺส อตฺถาย หิตาย สุขาย โหติ, มาตาปิตูนํ อตฺถาย หิตาย สุขาย โหติ, ปุตฺตทารสฺส อตฺถาย หิตาย สุขาย โหติ, ทาสกมฺมกรโปริสสฺส อตฺถาย หิตาย สุขาย โหติ, มิตฺตามจฺจานํ อตฺถาย หิตาย สุขาย โหติ, ปุพฺพเปตานํ อตฺถาย หิตาย สุขาย โหติ, รญฺโญ อตฺถาย หิตาย สุขาย โหติ, เทวตานํ อตฺถาย หิตาย สุขาย โหติ, สมณพฺราหฺมณานํ อตฺถาย หิตาย สุขาย โหติ.}

\prose{เสยฺยถาปิ ภิกฺขเว มหาเมโฆ สพฺพสสฺสานิ สมฺปาเทนฺโต พหุโน ชนสฺส อตฺถาย หิตาย สุขาย\footnote{หิตาย ฯเปฯ. (สฺยา, ก)} โหติ.\csromanspacer{} เอวเมวํ โข ภิกฺขเว สปฺปุริโส กุเล ชายมาโน พหุโน ชนสฺส อตฺถาย หิตาย สุขาย โหติ, มาตาปิตูนํ อตฺถาย หิตาย สุขาย โหติ, ปุตฺตทารสฺส อตฺถาย หิตาย สุขาย โหติ, ทาสกมฺมกรโปริสสฺส อตฺถาย หิตาย สุขาย โหติ, มิตฺตามจฺจานํ อตฺถาย หิตาย สุขาย โหติ, ปุพฺพเปตานํ อตฺถาย หิตาย สุขาย โหติ, รญฺโญ อตฺถาย หิตาย สุขาย โหติ, เทวตานํ อตฺถาย หิตาย สุขาย โหติ, สมณพฺราหฺมณานํ อตฺถาย หิตาย สุขาย โหตีติ.}

\csromangathameasure{พหูนํ วต อตฺถาย, สปฺปญฺโญ ฆรมาวสํ. \\ มาตรํ ปิตรํ ปุพฺเพ, รตฺตินฺทิวมตนฺทิโต. \\ ปูเชติ สหธมฺเมน, ปุพฺเพกตมนุสฺสรํ. \\ อนาคาเร ปพฺพชิเต, อปเจ พฺรหฺมจารโย. \\ นิวิฏฺฐสทฺโธ ปูเชติ, ญตฺวา ธมฺเม จ เปสโล. \\ รญฺโญ หิโต เทวหิโต, ญาตีนํ สขินํ หิโต. \\ สพฺเพสํ โส หิโต โหติ, สทฺธมฺเม สุปฺปติฏฺฐิโต. \\ วิเนยฺย มจฺเฉรมลํ, ส โลกํ ภชเต สิวนฺติ.อฏฺฐมํ.}
\csromangathagroup{\csromangathastanza{พหูนํ\footnote{พหุนฺนํ (สี, อิ)} วต อตฺถาย, สปฺปญฺโญ ฆรมาวสํ. \\ มาตรํ ปิตรํ ปุพฺเพ, รตฺตินฺทิวมตนฺทิโต.}\csromangathastanza{ปูเชติ สหธมฺเมน, ปุพฺเพกตมนุสฺสรํ. \\ อนาคาเร ปพฺพชิเต, อปเจ พฺรหฺมจารโย\footnote{พฺรหฺมจาริโน (สฺยา)}.}\csromangathastanza{นิวิฏฺฐสทฺโธ ปูเชติ, ญตฺวา ธมฺเม จ เปสโล\footnote{เปสเล (ก)}. \\ รญฺโญ หิโต เทวหิโต, ญาตีนํ สขินํ หิโต.}\csromangathastanza{สพฺเพสํ\footnote{สพฺเพสุ (ก)} โส\footnote{ส (สฺยา, อิ, ก)} หิโต โหติ, สทฺธมฺเม สุปฺปติฏฺฐิโต. \\ วิเนยฺย มจฺเฉรมลํ, ส โลกํ ภชเต สิวนฺติ.\csromanspacer{}\csromanspacer{}อฏฺฐมํ.}}

\csromanheader{2}{9. อภิสนฺทสุตฺต}
\proseitem{39}{\csromanfolio{76}อฏฺฐิเม ภิกฺขเว ปุญฺญาภิสนฺทา กุสลาภิสนฺทา สุขสฺสาหารา โสวคฺคิกา สุขวิปากา สคฺคสํวตฺตนิกา อิฏฺฐาย กนฺตาย มนาปาย หิตาย สุขาย สํวตฺตนฺติ.\csromanspacer{} กตเม อฏฺฐ? อิธ ภิกฺขเว อริยสาวโก พุทฺธํ สรณํ คโต โหติ, อยํ ภิกฺขเว ปฐโม ปุญฺญาภิสนฺโท กุสลาภิสนฺโท สุขสฺสาหาโร โสวคฺคิโก สุขวิปาโก สคฺคสํวตฺตนิโก อิฏฺฐาย กนฺตาย มนาปาย หิตาย สุขาย สํวตฺตติ.}

\prose{ปุน จปรํ ภิกฺขเว อริยสาวโก ธมฺมํ สรณํ คโต โหติ, อยํ ภิกฺขเว ทุติโย ปุญฺญาภิสนฺโท ฯเปฯ สํวตฺตติ.}

\prose{ปุน จปรํ ภิกฺขเว อริยสาวโก สํฆํ สรณํ คโต โหติ, อยํ ภิกฺขเว ตติโย ปุญฺญาภิสนฺโท กุสลาภิสนฺโท สุขสฺสาหาโร โสวคฺคิโก สุขวิปาโก สคฺคสํวตฺตนิโก อิฏฺฐาย กนฺตาย มนาปาย หิตาย สุขาย สํวตฺตติ.}

\prose{\csromansymbolfootnote{*}{อภิ 4. 256 ปิฏฺเฐ.}ปญฺจิมานิ ภิกฺขเว ทานานิ มหาทานานิ อคฺคญฺญานิ รตฺตญฺญานิ วํสญฺญานิ โปราณานิ อสํกิณฺณานิ อสํกิณฺณปุพฺพานิ น สํกิยนฺติ น สํกิยิสฺสนฺติ, อปฺปฏิกุฏฺฐานิ\footnote{อปฺปติกุฏฺฐานิ (สี)} สมเณหิ พฺราหฺมเณหิ วิญฺญูหิ.\csromanspacer{} กตมานิ ปญฺจ? อิธ ภิกฺขเว อริยสาวโก ปาณาติปาตํ ปหาย ปาณาติปาตา ปฏิวิรโต โหติ.\csromanspacer{} ปาณาติปาตา ปฏิวิรโต ภิกฺขเว อริยสาวโก อปริมาณานํ สตฺตานํ อภยํ เทติ, อเวรํ เทติ, อพฺยาพชฺฌํ\footnote{อพฺยาปชฺฌํ (ก) เอวมุปริปิ.} เทติ.\csromanspacer{} อปริมาณานํ สตฺตานํ อภยํ ทตฺวา อเวรํ ทตฺวา อพฺยาพชฺฌํ ทตฺวา อปริมาณสฺส อภยสฺส อเวรสฺส อพฺยาพชฺฌสฺส ภาคี โหติ.\csromanspacer{} อิทํ ภิกฺขเว ปฐมํ ทานํ มหาทานํ อคฺคญฺญํ รตฺตญฺญํ วํสญฺญํ โปราณํ อสํกิณฺณํ อสํกิณฺณปุพฺพํ น สํกิยติ น สํกิยิสฺสติ, อปฺปฏิกุฏฺฐํ สมเณหิ พฺราหฺมเณหิ วิญฺญูหิ.\csromanspacer{} อยํ ภิกฺขเว จตุตฺโถ ปุญฺญาภิสนฺโท กุสลาภิสนฺโท สุขสฺสาหาโร โสวคฺคิโก สุขวิปาโก สคฺคสํวตฺตนิโก อิฏฺฐาย กนฺตาย มนาปาย หิตาย สุขาย สํวตฺตติ.}

\prose{ปุน จปรํ ภิกฺขเว อริยสาวโก อทินฺนาทานํ ปหาย อทินฺนาทานา ปฏิวิรโต โหติ ฯเปฯ กาเมสุมิจฺฉาจารํ ปหาย กาเมสุมิจฺฉาจารา}

\vspace{\tipitakaparskip}
\csromancenter{ทานวคฺค ปฏิวิรโต โหติ ฯเปฯ มุสาวาทํ ปหาย มุสาวาทา ปฏิวิรโต โหติ ฯเปฯ สุราเมรยมชฺชปมาทฏฺฐานํ ปหาย สุราเมรยมชฺชปมาทฏฺฐานา ปฏิวิรโต โหติ.\csromanspacer{} สุราเมรยมชฺชปมาทฏฺฐานา ปฏิวิรโต ภิกฺขเว อริยสาวโก อปริมาณานํ สตฺตานํ อภยํ เทติ, อเวรํ เทติ, อพฺยาพชฺฌํ เทติ.\csromanspacer{} อปริมาณานํ สตฺตานํ อภยํ ทตฺวา อเวรํ ทตฺวา อพฺยาพชฺฌํ ทตฺวา อปริมาณสฺส อภยสฺส อเวรสฺส อพฺยาพชฺฌสฺส ภาคี โหติ.\csromanspacer{} อิทํ ภิกฺขเว ปญฺจมํ ทานํ มหาทานํ อคฺคญฺญํ รตฺตญฺญํ วํสญฺญํ โปราณํ อสํกิณฺณํ อสํกิณฺณปุพฺพํ น สํกิยติ น สํกิยิสฺสติ, อปฺปฏิกุฏฺฐํ สมเณหิ พฺราหฺมเณหิ วิญฺญูหิ.\csromanspacer{} อยํ โข ภิกฺขเว อฏฺฐโม ปุญฺญาภิสนฺโท กุสลาภิสนฺโท สุขสฺสาหาโร โสวคฺคิโก สุขวิปาโก สคฺคสํวตฺตนิโก อิฏฺฐาย กนฺตาย มนาปาย หิตาย สุขาย สํวตฺตติ.\csromanspacer{} อิเม โข ภิกฺขเว อฏฺฐ ปุญฺญาภิสนฺทา กุสลาภิสนฺทา สุขสฺสาหารา โสวคฺคิกา สุขวิปากา สคฺคสํวตฺตนิกา อิฏฺฐาย กนฺตาย มนาปาย หิตาย สุขาย สํวตฺตนฺตีติ.\csromanspacer{}\csromanspacer{}นวมํ.}
\csromanheader{2}{10. ทุจฺจริตวิปากสุตฺต}
\proseitem{40}{\csromanfolio{77}ปาณาติปาโต ภิกฺขเว อาเสวิโต ภาวิโต พหุลีกโต นิรยสํวตฺตนิโก, ติรจฺฉานโยนิสํวตฺตนิโก, เปตฺติวิสยสํวตฺตนิโก.\csromanspacer{} โย สพฺพลหุโส\footnote{สพฺพลหุโสติ สพฺพลหุโก (สฺยา-ฏฺฐ)} ปาณาติปาตสฺส วิปาโก มนุสฺสภูตสฺส อปฺปายุกสํวตฺตนิโก โหติ.}

\prose{อทินฺนาทานํ ภิกฺขเว อาเสวิตํ ภาวิตํ พหุลีกตํ นิรยสํวตฺตนิกํ, ติรจฺฉานโยนิสํวตฺตนิกํ, เปตฺติวิสยสํวตฺตนิกํ.\csromanspacer{} โย สพฺพลหุโส อทินฺนาทานสฺส วิปาโก มนุสฺสภูตสฺส โภคพฺยสนสํวตฺตนิโก โหติ.}

\prose{กาเมสุมิจฺฉาจาโร ภิกฺขเว อาเสวิโต ภาวิโต พหุลีกโต นิรยสํวตฺตนิโก, ติรจฺฉานโยนิสํวตฺตนิโก, เปตฺติวิสยสํวตฺตนิโก.\csromanspacer{} โย สพฺพลหุโส กาเมสุมิจฺฉาจารสฺส วิปาโก มนุสฺสภูตสฺส สปตฺตเวรสํวตฺตนิโก โหติ.}

\prose{\csromanfolio{78}มุสาวาโท ภิกฺขเว อาเสวิโต ภาวิโต พหุลีกโต นิรยสํวตฺตนิโก, ติรจฺฉานโยนิสํวตฺตนิโก, เปตฺติวิสยสํวตฺตนิโก.\csromanspacer{} โย สพฺพลหุโส มุสาวาทสฺส วิปาโก มนุสฺสภูตสฺส อภูตพฺภกฺขานสํวตฺตนิโก โหติ.}

\prose{ปิสุณา ภิกฺขเว วาจา อาเสวิตา ภาวิตา พหุลีกตา นิรยสํวตฺตนิกา, ติรจฺฉานโยนิสํวตฺตนิกา, เปตฺติวิสยสํวตฺตนิกา.\csromanspacer{} โย สพฺพลหุโส ปิสุณาย วาจาย วิปาโก มนุสฺสภูตสฺส มิตฺเตหิ เภทนสํวตฺตนิโก โหติ.}

\prose{ผรุสา ภิกฺขเว วาจา อาเสวิตา ภาวิตา พหุลีกตา นิรยสํวตฺตนิกา, ติรจฺฉานโยนิสํวตฺตนิกา, เปตฺติวิสยสํวตฺตนิกา.\csromanspacer{} โย สพฺพลหุโส ผรุสาย วาจาย วิปาโก มนุสฺสภูตสฺส อมนาปสทฺทสํวตฺตนิโก โหติ.}

\prose{สมฺผปฺปลาโป ภิกฺขเว อาเสวิโต ภาวิโต พหุลีกโต นิรยสํวตฺตนิโก, ติรจฺฉานโยนิสํวตฺตนิโก, เปตฺติวิสยสํวตฺตนิโก.\csromanspacer{} โย สพฺพลหุโส สมฺผปฺปลาปสฺส วิปาโก มนุสฺสภูตสฺส อนาเทยฺยวาจาสํวตฺตนิโก โหติ.}

\prose{สุราเมรยปานํ ภิกฺขเว อาเสวิตํ ภาวิตํ พหุลีกตํ นิรยสํวตฺตนิกํ, ติรจฺฉานโยนิสํวตฺตนิกํ, เปตฺติวิสยสํวตฺตนิกํ.\csromanspacer{} โย สพฺพลหุโส สุราเมรยปานสฺส วิปาโก มนุสฺสภูตสฺส อุมฺมตฺตกสํวตฺตนิโก โหตีติ.\csromanspacer{}\csromanspacer{}ทสมํ.\csromanspacer{} ทานวคฺโค จตุตฺโถ}

\titlehead{\textbf{ตสฺสุทฺทานํ}}
\csromangathameasure{เทฺว ทานานิ วตฺถุญฺจ, เขตฺตํ ทานูปปตฺติโย. \\ กิริยํ เทฺว สปฺปุริสา, อภิสนฺโท วิปาโก จาติ.}
\csromangathagroup{\csromangathastanza{เทฺว ทานานิ วตฺถุญฺจ, เขตฺตํ ทานูปปตฺติโย. \\ กิริยํ เทฺว สปฺปุริสา, อภิสนฺโท วิปาโก จาติ.}}

\chapterheadpageread{5. อุโปสถวคฺค \textbf{5. อุโปสถวคฺค}}
\csromanheader{2}{1. สํขิตฺตูโปสถสุตฺต}
\proseitem{41}{\csromanfolio{79}เอวํ เม สุตํ\csromandash{}เอกํ สมยํ ภควา สาวตฺถิยํ วิหรติ เชตวเน อนาถปิณฺฑิกสฺส อาราเม.\csromanspacer{} ตตฺร โข ภควา ภิกฺขู อามนฺเตสิ “ภิกฺขโว”ติ. “ภทนฺเต”ติ เต ภิกฺขู ภควโต ปจฺจสฺโสสุํ.\csromanspacer{} ภควา เอตทโวจ\csromandash{}}

\prose{อฏฺฐงฺคสมนฺนาคโต ภิกฺขเว อุโปสโถ อุปวุตฺโถ มหปฺผโล โหติ มหานิสํโส มหาชุติโก มหาวิปฺผาโร.\csromanspacer{} กถํ อุปวุตฺโถ จ ภิกฺขเว อฏฺฐงฺคสมนฺนาคโต อุโปสโถ มหปฺผโล โหติ มหานิสํโส มหาชุติโก มหาวิปฺโผโร? อิธ ภิกฺขเว อริยสาวโก อิติ ปฏิสญฺจิกฺขติ “ยาวชีวํ อรหนฺโต ปาณาติปาตํ ปหาย ปาณาติปาตา ปฏิวิรตา นิหิตทณฺฑา นิหิตสตฺถา ลชฺชี ทยาปนฺนา สพฺพปาณภูตหิตานุกมฺปิโน วิหรนฺติ, อหมฺปชฺช อิมญฺจ รตฺติํ อิมญฺจ ทิวสํ ปาณาติปาตํ ปหาย ปาณาติปาตา ปฏิวิรโต นิหิตทณฺโฑ นิหิตสตฺโถ ลชฺชี ทยาปนฺโน สพฺพปาณภูตหิตานุกมฺปี วิหรามิ, อิมินาปงฺเคน\footnote{อิมินาปิ องฺเคน (สี, อิ) อํ 1. 212 ปิฏฺเฐปิ.} อรหตํ อนุกโรมิ, อุโปสโถ จ เม อุปวุตฺโถ ภวิสฺสตี”ติ.\csromanspacer{} อิมินา ปฐเมน องฺเคน สมนฺนาคโต โหติ.}

\prose{“ยาวชีวํ อรหนฺโต อทินฺนาทานํ ปหาย อทินฺนาทานา ปฏิวิรตา ทินฺนาทายี ทินฺนปาฏิกงฺขี อเถเนน สุจิภูเตน อตฺตนา วิหรนฺติ, อหมฺปชฺช อิมญฺจ รตฺติํ อิมญฺจ ทิวสํ อทินฺนาทานํ ปหาย อทินฺนาทานา ปฏิวิรโต ทินฺนาทายี ทินฺนปาฏิกงฺขี อเถเนน สุจิภูเตน อตฺตนา วิหรามิ, อิมินาปงฺเคน อรหตํ อนุกโรมิ, อุโปสโถ จ เม อุปวุตฺโถ ภวิสฺสตี”ติ.\csromanspacer{} อิมินา ทุติเยน องฺเคน สมนฺนาคโต โหติ.}

\prose{“ยาวชีวํ อรหนฺโต อพฺรหฺมจริยํ ปหาย พฺรหฺมจาริโน อาราจาริโน วิรตา เมถุนา คามธมฺมา, อหมฺปชฺช อิมญฺจ รตฺติํ อิมญฺจ ทิวสํ อพฺรหฺมจริยํ ปหาย พฺรหฺมจารี อาราจารี\footnote{อนาจารี (ก)} วิรโต เมถุนา คามธมฺมา, อิมินาปงฺเคน อรหตํ อนุกโรมิ, อุโปสโถ จ เม อุปวุตฺโถ ภวิสฺสตี”ติ.\csromanspacer{} อิมินา ตติเยน องฺเคน สมนฺนาคโต โหติ.}

\prose{\csromanfolio{80}“ยาวชีวํ อรหนฺโต มุสาวาทํ ปหาย มุสาวาทา ปฏิวิรตา สจฺจวาทิโน สจฺจสนฺธา เถตา ปจฺจยิกา อวิสํวาทกา โลกสฺส, อหมฺปชฺช อิมญฺจ รตฺติํ อิมญฺจ ทิวสํ มุสาวาทํ ปหาย มุสาวาทา ปฏิวิรโต สจฺจวาที สจฺจสนฺโธ เถโต ปจฺจยิโก อวิสํวาทโก โลกสฺส, อิมินาปงฺเคน อรหตํ อนุกโรมิ, อุโปสโถ จ เม อุปวุตฺโถ ภวิสฺสตี”ติ.\csromanspacer{} อิมินา จตุตฺเถน องฺเคน สมนฺนาคโต โหติ.}

\prose{“ยาวชีวํ อรหนฺโต สุราเมรยมชฺชปมาทฏฺฐานํ ปหาย สุราเมรยมชฺชปมาทฏฺฐานา ปฏิวิรตา, อหมฺปชฺช อิมญฺจ รตฺติํ อิมญฺจ ทิวสํ สุรา เมรย มชฺช ปมาทฏฺฐานํ ปหาย สุรา เมรย มชฺช ปมาทฏฺฐานา ปฏิวิรโต, อิมินาปงฺเคน อรหตํ อนุกโรมิ, อุโปสโถ จ เม อุปวุตฺโถ ภวิสฺสตี”ติ.\csromanspacer{} อิมินา ปญฺจเมน องฺเคน สมนฺนาคโต โหติ.}

\prose{“ยาวชีวํ อรหนฺโต เอกภตฺติกา รตฺตูปรตา วิรตา วิกาลโภชนา, อหมฺปชฺช อิมญฺจ รตฺติํ อิมญฺจ ทิวสํ เอกภตฺติโก รตฺตูปรโต วิรโต วิกาลโภชนา, อิมินาปงฺเคน อรหตํ อนุกโรมิ, อุโปสโถ จ เม อุปวุตฺโถ ภวิสฺสตี”ติ.\csromanspacer{} อิมินา ฉฏฺเฐน องฺเคน สมนฺนาคโต โหติ.}

\prose{“ยาวชีวํ อรหนฺโต นจฺจคีตวาทิตวิสูกทสฺสนมาลาคนฺธ- วิเลปนธารณมณฺฑนวิภูสนฏฺฐานํ ปหาย นจฺจคีตวาทิตวิสูก- ทสฺสนมาลาคนฺธวิเลปนธารณมณฺฑนวิภูสนฏฺฐานา ปฏิวิรตา, อหมฺปชฺช อิมญฺจ รตฺติํ อิมญฺจ ทิวสํ นจฺจคีตวาทิตวิสูกทสฺสน- มาลาคนฺธวิเลปนธารณมณฺฑนวิภูสนฏฺฐานํ ปหาย นจฺจคีต- วาทิตวิสูกทสฺสนมาลาคนฺธวิเลปนธารณมณฺฑนวิภูสนฏฺฐานา ปฏิวิรโต, อิมินาปงฺเคน อรหตํ อนุกโรมิ, อุโปสโถ จ เม อุปวุตฺโถ ภวิสฺสตี”ติ.\csromanspacer{} อิมินา สตฺตเมน องฺเคน สมนฺนาคโต โหติ.}

\prose{“ยาวชีวํ อรหนฺโต อุจฺจาสยนมหาสยนํ ปหาย อุจฺจาสยนมหาสยนา ปฏิวิรตา นีจเสยฺยํ กปฺเปนฺติ มญฺจเก วา ติณสนฺถารเกวา, อหมฺปชฺช อิมญฺจ รตฺติํ อิมญฺจ ทิวสํ อุจฺจาสยนมหาสยนํ ปหาย อุจฺจาสยนมหาสยนา ปฏิวิรโต นีจเสยฺยํ กปฺเปมิ มญฺจเก วา ติณสนฺถารเก วา, อิมินาปงฺเคน อรหตํ อนุกโรมิ, อุโปสโถ จ เม อุปวุตฺโถ ภวิสฺสตี”ติ.\csromanspacer{} อิมินา อฏฺฐเมน องฺเคน สมนฺนาคโต}

\vspace{\tipitakaparskip}
\csromancenter{อุโปสถวคฺค โหติ.\csromanspacer{} เอวํ อุปวุตฺโถ โข ภิกฺขเว อฏฺฐงฺคสมนฺนาคโต อุโปสโถ มหปฺผโล โหติ มหานิสํโส มหาชุติโก มหาวิปฺผาโรติ.\csromanspacer{}\csromanspacer{}ปฏฺฐมํ.}
\csromanheader{2}{2. วิตฺถตูโปสถสุตฺต}
\proseitem{42}{\csromanfolio{81}อฏฺฐงฺคสมนฺนาคโต ภิกฺขเว อุโปสโถ อุปวุตฺโถ มหปฺผโล โหติ มหานิสํโส มหาชุติโก มหาวิปฺผาโร.\csromanspacer{} กถํ อุปวุตฺโถ จ ภิกฺขเว อฏฺฐงฺคสมนฺนาคโต อุโปสโถ มหปฺผโล โหติ มหานิสํโส มหาชุติโก มหาวิปฺผาโร? อิธ ภิกฺขเว อริยสาวโก อิติ ปฏิสญฺจิกฺขติ “ยาวชีวํ อรหนฺโต ปาณาติปาตํ ปหาย ปาณาติปาตา ปฏิวิรตา นิหิตทณฺฑา นิหิตสตฺถา ลชฺชี ทยาปนฺนา สพฺพปาณภูตหิตานุกมฺปิโน วิหรติ, อหมฺปชฺช อิมญฺจ รตฺติํ อิมญฺจ ทิวสํ ปาณาติปาตํ ปหาย ปาณาติปาตา ปฏิวิรโต นิหิตทณฺโฑ นิหิตสตฺโถ ลชฺชี ทยาปนฺโน สพฺพปาณภูต- หิตานุกมฺปี วิหรามิ, อิมินาปงฺเคน อรหตํ อนุกโรมิ, อุโปสโถ จ เม อุปวุตฺโถ ภวิสฺสตี”ติ.\csromanspacer{} อิมินา ปฐเมน องฺเคน สมนฺนาคโต โหติ ฯเปฯ.}

\prose{“ยาวชีวํ อรหนฺโต อุจฺจาสยนมหาสยนํ ปหาย อุจฺจาสยนมหาสยนา ปฏิวิรตา นีจเสยฺยํ กปฺเปนฺติ มญฺจเก วา ติณสนฺถารเก วา, อหมฺปชฺช อิมญฺจ รตฺติํ อิมญฺจ ทิวสํ อุจฺจาสยนมหาสยนํ ปหาย อุจฺจาสยนมหาสยนา ปฏิวิรโต นีจเสยฺยํ กปฺเปมิ มญฺจเก วา ติณสนฺถารเก วา, อิมินาปงฺเคน อรหตํ อนุกโรมิ, อุโปสโถ จ เม อุปวุตฺโถ ภวิสฺสตี”ติ.\csromanspacer{} อิมินา อฏฺฐเมน องฺเคน สมนฺนาคโต โหติ.\csromanspacer{} เอวํ อุปวุตฺโถ โข ภิกฺขเว อฏฺฐงฺคสมนฺนาคโต อุโปสโถ มหปฺผโล โหติ มหานิสํโส มหาชุติโก มหาวิปฺผาโร.}

\prose{กีวมหปฺผโล โหติ กีวมหานิสํโส กีวมหาชุติโก กีวมหาวิปฺผาโร? เสยฺยถาปิ ภิกฺขเว โย อิเมสํ โสฬสนฺนํ มหาชนปทานํ ปหูตรตฺตรตนานํ\footnote{ปหูตสตฺตรตนานํ (สี, สฺยา กํ, อิ) อํ 1. 213 ปิฏฺเฐ ปาฬิยา ฏีกายํ ทสฺสิตปาฬิเยว. ตทฏฺฐกถาปิ ปสฺสิตพฺพา.} อิสฺสริยาธิปจฺจํ รชฺชํ กาเรยฺย.\csromanspacer{} เสยฺยถิทํ, องฺคานํ มคธานํ กาสีนํ โกสลานํ วชฺชีนํ มลฺลานํ เจตีนํ วงฺคานํ กุรูนํ ปญฺจาลานํ มจฺฉานํ\footnote{มชฺชานํ (ก)} สูรเสนานํ อสฺสกานํ อวนฺตีนํ คนฺธารานํ กมฺโพชานํ, \csromanfolio{82}อฏฺฐงฺคสมนฺนาคตสฺส อุโปสถสฺส เอตํ\footnote{เอกํ (ก)} กลํ นาคฺฆติ โสฬสิํ.\csromanspacer{} ตํ กิสฺส เหตุ? กปณํ ภิกฺขเว มานุสกํ รชฺชํ ทิพฺพํ สุขํ อุปนิธาย.}

\prose{ยานิ ภิกฺขเว มานุสกานิ ปญฺญาส วสฺสานิ, จาตุมหาราชิกานํ เทวานํ เอโส เอโก รตฺตินฺทิโว\footnote{รตฺติทิโว (ก)}, ตาย รตฺติยา ติํสรตฺติโย มาโส, เตน มาเสน ทฺวาทสมาสิโย สํวจฺฉโร, เตน สํวจฺฉเรน ทิพฺพานิ ปญฺจ วสฺสสตานิ จาตุมหาราชิกานํ เทวานํ อายุปฺปมาณํ.\csromanspacer{} ฐานํ โข ปเนตํ ภิกฺขเว วิชฺชติ, ยํ อิเธกจฺโจ อิตฺถี วา ปุริโส วา อฏฺฐงฺคสมนฺนาคตํ อุโปสถํ อุปวสิตฺวา กายสฺส เภทา ปรํ มรณา จาตุมหาราชิกานํ เทวานํ สหพฺยตํ อุปปชฺเชยฺย.\csromanspacer{} อิทํ โข ปเนตํ ภิกฺขเว สนฺธาย ภาสิตํ กปณํ มานุสกํ รชฺชํ ทิพฺพํ สุขํ อุปนิธาย.}

\prose{ยานิ ภิกฺขเว มานุสกานิ วสฺสสตานิ, ตาวติํสานํ เทวานํ เอโส เอโก รตฺตินฺทิโว, ตาย รตฺติยา ติํสรตฺติโย มาโส, เตน มาเสน ทฺวาทสมาสิโย สํวจฺฉโร, เตน สํวจฺฉเรน ทิพฺพํ วสฺสสหสฺสํ ตาวติํสานํ เทวานํ อายุปฺปมาณํ.\csromanspacer{} ฐานํ โข ปเนตํ ภิกฺขเว วิชฺชติ, ยํ อิเธกจฺโจ อิตฺถี วา ปุริโส วา อฏฺฐงฺคสมนฺนาคตํ อุโปสถํ อุปวสิตฺวา กายสฺส เภทา ปรํ มรณา ตาวติํสานํ เทวานํ สหพฺยตํ อุปปชฺเชยฺย.\csromanspacer{} อิทํ โข ปเนตํ ภิกฺขเว สนฺธาย ภาสิตํ กปณํ มานุสกํ รชฺชํ ทิพฺพํ สุขํ อุปนิธาย.}

\prose{ยานิ ภิกฺขเว มานุสกานิ เทฺว วสฺสสตานิ, ยามานํ เทวานํ เอโส เอโก รตฺตินฺทิโว, ตาย รตฺติยา ติํสรตฺติโย มาโส, เตน มาเสน ทฺวาทสมาสิโย สํวจฺฉโร, เตน สํวจฺฉเรน ทิพฺพานิ เทฺว วสฺสสหสฺสานิ ยามานํ เทวานํ อายุปฺปมาณํ.\csromanspacer{} ฐานํ โข ปเนตํ ภิกฺขเว วิชฺชติ, ยํ อิเธกจฺโจ อิตฺถี วา ปุริโส วา อฏฺฐงฺคสมนฺนาคตํ อุโปสถํ อุปวสิตฺวา กายสฺส เภทา ปรํ มรณา ยามานํ เทวานํ สหพฺยตํ อุปปชฺเชยฺย.\csromanspacer{} อิทํ โข ปเนตํ ภิกฺขเว สนฺธาย ภาสิตํ กปณํ มานุสกํ รชฺชํ ทิพฺพํ สุขํ อุปนิธาย.}

\prose{ยานิ ภิกฺขเว มานุสกานิ จตฺตาริ วสฺสสตานิ, ตุสิตานํ เทวานํ เอโส เอโก รตฺตินฺทิโว, ตาย รตฺติยา ติํสรตฺติโย มาโส, เตน มาเสน ทฺวาทสมาสิโย สํวจฺฉโร, เตน สํวจฺฉเรน ทิพฺพานิ จตฺตาริ}

\vspace{\tipitakaparskip}
\csromancenter{อุโปสถวคฺค วสฺสสหสฺสานิ ตุสิตานํ เทวานํ อายุปฺปมาณํ.\csromanspacer{} ฐานํ โข ปเนตํ ภิกฺขเว วิชฺชติ, ยํ อิเธกจฺโจ อิตฺถี วา ปุริโส วา อฏฺฐงฺคสมนฺนาคตํ อุโปสถํ อุปวสิตฺวา กายสฺส เภทา ปรํ มรณา ตุสิตานํ เทวานํ สหพฺยตํ อุปปชฺเชยฺย.\csromanspacer{} อิทํ โข ปเนตํ ภิกฺขเว สนฺธาย ภาสิตํ กปณํ มานุสกํ รชฺชํ ทิพฺพํ สุขํ อุปนิธาย.}
\prose{\csromanfolio{83}ยานิ ภิกฺขเว มานุสกานิ อฏฺฐ วสฺสสตานิ, นิมฺมานรตีนํ เทวานํ เอโส เอโก รตฺตินฺทิโว, ตาย รตฺติยา ติํสรตฺติโย มาโส, เตน มาเสน ทฺวาทสมาสิโย สํวจฺฉโร, เตน สํวจฺฉเรน ทิพฺพานิ อฏฺฐ วสฺสสหสฺสานิ นิมฺมานรตีนํ เทวานํ อายุปฺปมาณํ.\csromanspacer{} ฐานํ โข ปเนตํ ภิกฺขเว วิชฺชติ, ยํ อิเธกจฺโจ อิตฺถี วา ปุริโส วา อฏฺฐงฺคสมนฺนาคตํ อุโปสถํ อุปวสิตฺวา กายสฺส เภทา ปรํ มรณา นิมฺมานรตีนํ เทวานํ สหพฺยตํ อุปปชฺเชยฺย.\csromanspacer{} อิทํ โข ปเนตํ ภิกฺขเว สนฺธาย ภาสิตํ กปณํ มานุสกํ รชฺชํ ทิพฺพํ สุขํ อุปนิธาย.}

\prose{ยานิ ภิกฺขเว มานุสกานิ โสฬส วสฺสสตานิ ปรนิมฺมิตวส- วตฺตีนํ เทวานํ เอโส เอโก รตฺตินฺทิโว, ตาย รตฺติยา ติํสรตฺติโย มาโส, เตน มาเสน ทฺวาทสมาสิโย สํวจฺฉโร, เตน สํวจฺฉเรน ทิพฺพานิ โสฬส วสฺสสหสฺสานิ ปรนิมฺมิตวสวตฺตีนํ เทวานํ อายุปฺปมาณํ.\csromanspacer{} ฐานํ โข ปเนตํ ภิกฺขเว วิชฺชติ, ยํ อิเธกจฺโจ อิตฺถี วา ปุริโส วา อฏฺฐงฺคสมนฺนาคตํ อุโปสถํ อุปวสิตฺวา กายสฺส เภทา ปรํ มรณา ปรนิมฺมิตวสวตฺตีนํ เทวานํ สหพฺยตํ อุปปชฺเชยฺย.\csromanspacer{} อิทํ โข ปเนตํ ภิกฺขเว สนฺธาย ภาสิตํ “กปณํ มานุสกํ รชฺชํ ทิพฺพํ สุขํ อุปนิธายา”ติ.}

\csromangathameasure{ปาณํ น หญฺเญ น จทินฺนมาทิเย, \\ มุสา น ภาเส น จ มชฺชโป สิยา. \\ อพฺรหฺมจริยา วิรเมยฺย เมถุนา, \\ รตฺติํ น ภุญฺเชยฺย วิกาลโภชนํ. \\ มาลํ น ธาเร น จ คนฺธมาจเร, \\ มญฺเจ ฉมายํ ว สเยถ สนฺถเต. \\ เอตํ หิ อฏฺฐงฺคิกมาหุโปสถํ, \\ พุทฺเธน ทุกฺขนฺตคุนา ปกาสิตํ. \\ จนฺโท จ สุริโย จ อุโภ สุทสฺสนา, \\ โอภาสยํ อนุปริยนฺติ ยาวตา. \\ ตโมนุทา เต ปน อนฺตลิกฺขคา, \\ นเภ ปภาสนฺติ ทิสาวิโรจนา. \\ เอตสฺมิํ ยํ วิชฺชติ อนฺตเร ธนํ, \\ มุตฺตา มณิ เวฬุริยญฺจ ภทฺทกํ. \\ สิงฺคีสุวณฺณํ อถ วาปิ กญฺจนํ, \\ ยํ ชาตรูปํ หฏกนฺติ วุจฺจติ. \\ อฏฺฐงฺคุเปตสฺส อุโปสถสฺส, \\ กลมฺปิ เต นานุภวนฺติ โสฬสิํ. \\ จนฺทปฺปภา ตารคณา จ สพฺเพ. \\ ตสฺมา หิ นารี จ นโร จ สีลวา, \\ อฏฺฐงฺคุเปตํ อุปวสฺสุโปสถํ. \\ ปุญฺญานิ กตฺวาน สุขุทฺรยานิ, \\ อนินฺทิตา สคฺคมุเปนฺติ ฐานนฺติ.ทุติยํ.}
\csromangathagroup{\csromangathastanza{ปาณํ น หญฺเญ\footnote{หาเน (สี), หเน (ก) อํ 1. 215 ปิฏฺเฐปิ.} น จทินฺนมาทิเย, \\ มุสา น ภาเส น จ มชฺชโป สิยา. \\ อพฺรหฺมจริยา วิรเมยฺย เมถุนา, \\ รตฺติํ น ภุญฺเชยฺย วิกาลโภชนํ.}\csromangathastanza{มาลํ น ธาเร น จ คนฺธมาจเร\footnote{คนฺธมาธเร (ก)}, \\ มญฺเจ ฉมายํ ว สเยถ สนฺถเต. \\ เอตํ หิ อฏฺฐงฺคิกมาหุโปสถํ, \\ พุทฺเธน ทุกฺขนฺตคุนา ปกาสิตํ.}\csromangathastanza{\csromanfolio{84}จนฺโท จ สุริโย จ อุโภ สุทสฺสนา, \\ โอภาสยํ อนุปริยนฺติ ยาวตา. \\ ตโมนุทา เต ปน อนฺตลิกฺขคา, \\ นเภ ปภาสนฺติ ทิสาวิโรจนา.}\csromangathastanza{เอตสฺมิํ ยํ วิชฺชติ อนฺตเร ธนํ, \\ มุตฺตา มณิ เวฬุริยญฺจ ภทฺทกํ. \\ สิงฺคีสุวณฺณํ อถ วาปิ กญฺจนํ, \\ ยํ ชาตรูปํ หฏกนฺติ วุจฺจติ.}\csromangathastanza{อฏฺฐงฺคุเปตสฺส อุโปสถสฺส, \\ กลมฺปิ เต นานุภวนฺติ โสฬสิํ. \\ จนฺทปฺปภา ตารคณา จ สพฺเพ. \\ ตสฺมา หิ นารี จ นโร จ สีลวา,}\csromangathastanza{อฏฺฐงฺคุเปตํ อุปวสฺสุโปสถํ. \\ ปุญฺญานิ กตฺวาน สุขุทฺรยานิ, \\ อนินฺทิตา สคฺคมุเปนฺติ ฐานนฺติ.\csromanspacer{}\csromanspacer{}ทุติยํ.}}

\csromanheader{2}{3. วิสาขาสุตฺต}
\proseitem{43}{\csromansymbolfootnote{*}{อํ 1. 206 ปิฏฺฐาทีสุปิ.}เอกํ สมยํ ภควา สาวตฺถิยํ วิหรติ ปุพฺพาราเม มิคารมาตุปาสาเท.\csromanspacer{} อถ โข วิสาขา มิคารมาตา เยน ภควา เตนุปสงฺกมิ, อุปสงฺกมิตฺวา ภควนฺตํ อภิวาเทตฺวา เอกมนฺตํ นิสีทิ, เอกมนฺตํ นิสินฺนํ โข วิสาขํ มิคารมาตรํ ภควา เอตทโวจ\csromandash{} อฏฺฐงฺคสมนฺนาคโต โข วิสาเข อุโปสโถ อุปวุตฺโถ มหปฺผโล โหติ มหานิสํโส มหาชุติโก มหาวิปฺผาโร.\csromanspacer{} กถํ อุปวุตฺโถ จ วิสาเข อฏฺฐงฺคสมนฺนาคโต อุโปสโถ มหปฺผโล โหติ มหานิสํโส มหาชุติโกมหาวิปฺผาโร? อิธ วิสาเข อริยสาวโก อิติ ปฏิสญฺจิกฺขติ “ยาวชีวํ อรหนฺโต ปาณาติปาตํ ปหาย ปาณาติปาตา ปฏิวิรตา นิหิตทณฺฑา นิหิตสตฺถา ลชฺชี ทยาปนฺนา สพฺพปาณภูตหิตานุกมฺปิโน วิหรนฺติ, อหมฺปชฺช อิมญฺจ รตฺติํ อิมญฺจ ทิวสํ ปาณาติปาตํ ปหาย ปาณาติปาตา ปฏิวิรโต นิหิตทณฺโฑ นิหิตสตฺโถ ลชฺชี ทยาปนฺโน สพฺพปาณภูตหิตา- นุกมฺปี วิหรามิ.\csromanspacer{} อิมินาปงฺเคน อรหตํ อนุกโรมิ,}

\vspace{\tipitakaparskip}
\csromancenter{อุโปสถวคฺค อุโปสโถ จ เม อุปวุตฺโถ ภวิสฺสตี”ติ.\csromanspacer{} อิมินา ปฐเมน องฺเคน สมนฺนาคโต โหติ ฯเปฯ.}
\prose{\csromanfolio{85}“ยาวชีวํ อรหนฺโต อุจฺจาสยนมหาสยนํ ปหาย อุจฺจาสยนมหาสยนา ปฏิวิรตา นีจเสยฺยํ กปฺเปนฺติ มญฺจเก วา ติณสนฺถารเก วา, อหมฺปชฺช อิมญฺจ รตฺติํ อิมญฺจ ทิวสํ อุจฺจาสยนมหาสยนํ ปหาย อุจฺจาสยนมหาสยนา ปฏิวิรโต นีจเสยฺยํ กปฺเปมิ มญฺจเก วา ติณสนฺถารเก วา.\csromanspacer{} อิมินาปงฺเคน อรหตํ อนุกโรมิ, อุโปสโถ จ เม อุปวุตฺโถ ภวิสฺสตี”ติ.\csromanspacer{} อิมินา อฏฺฐเมน องฺเคน สมนฺนาคโต โหติ.\csromanspacer{} เอวํ อุปวุตฺโถ โข วิสาเข อฏฺฐงฺคสมนฺนาคโต อุโปสโถ มหปฺผโล โหติ มหานิสํโส มหาชุติโก มหาวิปฺผาโร.}

\prose{กีวมหปฺผโล โหติ กีวมหานิสํโส กีวมหาชุติโก กีวมหาวิปฺผาโร? เสยฺยถาปิ วิสาเข โย อิเมสํ โสฬสนฺนํ มหาชนปทานํ ปหูตรตฺตรตนานํ อิสฺสริยาธิปจฺจํ รชฺชํ กาเรยฺย.\csromanspacer{} เสยฺยถิทํ, องฺคานํ มคธานํ กาสีนํ โกสลานํ วชฺชีนํ มลฺลานํ เจตีนํ วงฺคานํ กุรูนํ ปญฺจาลานํ มจฺฉานํ สูรเสนานํ อสฺสกานํ อวนฺตีนํ คนฺธารานํ กมฺโพชานํ, อฏฺฐงฺคสมนฺนาคตสฺส อุโปสถสฺส เอตํ กลํ นาคฺฆติ โสฬสิํ.\csromanspacer{} ตํ กิสฺส เหตุ? กปณํ วิสาเข มานุสกํ รชฺชํ ทิพฺพํ สุขํ อุปนิธาย.}

\prose{ยานิ วิสาเข มานุสกานิ ปญฺญาส วสฺสานิ จาตุมหาราชิกานํ เทวานํ เอโส เอโก รตฺตินฺทิโว, ตาย รตฺติยา ติํสรตฺติโย มาโส, เตน มาเสน ทฺวาทสมาสิโย สํวจฺฉโร, เตน สํวจฺฉเรน ทิพฺพานิ ปญฺจ วสฺสสตานิ จาตุมหาราชิกานํ เทวานํ อายุปฺปมาณํ.\csromanspacer{} ฐานํ โข ปเนตํ วิสาเข วิชฺชติ, ยํ อิเธกจฺโจ อิตฺถี วา ปุริโส วา อฏฺฐงฺคสมนฺนาคตํ อุโปสถํ อุปวสิตฺวา กายสฺส เภทา ปรํ มรณา จาตุมหาราชิกานํ เทวานํ สหพฺยตํ อุปปชฺเชยฺย.\csromanspacer{} อิทํ โข ปเนตํ วิสาเข สนฺธาย ภาสิตํ กปณํ มานุสกํ รชฺชํ ทิพฺพํ สุขํ อุปนิธาย.}

\prose{ยํ วิสาเข มานุสกํ วสฺสสตํ, ตาวติํสานํ เทวานํ เอโส เอโก รตฺตินฺทิโว, ตาย รตฺติยา ติํสรตฺติโย มาโส, เตน มาเสน ทฺวาทสมาสิโย สํวจฺฉโร, เตน สํวจฺฉเรน วสฺสสหสฺสํ ตาวติํสานํ เทวานํ อายุปฺปมาณํ.\csromanspacer{} ฐานํ โข ปเนตํ วิสาเข \csromanfolio{86}วิชฺชติ, ยํ อิเธกจฺโจ อิตฺถี วา ปุริโส วา อฏฺฐงฺคสมนฺนาคตํ อุโปสถํ อุปวสิตฺวา กายสฺส เภทา ปรํ มรณา ตาวติํสานํ เทวานํ สหพฺยตํ อุปปชฺเชยฺย.\csromanspacer{} อิทํ โข ปเน ปเนตํ วิสาเข สนฺธาย ภาสิตํ กปณํ มานุสกํ รชฺชํ ทิพฺพํ สุขํ อุปนิธาย.}

\prose{ยานิ วิสาเข มานุสกานิ เทฺว วสฺสสตานิ ฯเปฯ จตฺตาริ วสฺสสตานิ ฯเปฯ อฏฺฐ วสฺสสตานิ ฯเปฯ โสฬส วสฺสสตานิ ปรนิมฺมิตวสวตฺตีนํ เทวานํ เอโส เอโก รตฺตินฺทิโว, ตาย รตฺติยา ติํสรตฺติโย มาโส, เตน มาเสน ทฺวาทสมาสิโย สํวจฺฉโร, เตน สํวจฺฉเรน ทิพฺพานิ โสฬส วสฺสสหสฺสานิ ปรนิมฺมิตวสวตฺตีนํ เทวานํ อายุปฺปมาณํ.\csromanspacer{} ฐานํ โข ปเนตํ วิสาเข วิชฺชติ, ยํ อิเธกจฺโจ อิตฺถี วา ปุริโส วา อฏฺฐงฺคสมนฺนาคตํ อุโปสถํ อุปวสิตฺวา กายสฺส เภทา ปรํ มรณา ปรนิมฺมิตวสวตฺตีนํ เทวานํ สหพฺยตํ อุปปชฺเชยฺย.\csromanspacer{} อิทํ โข ปเนตํ วิสาเข สนฺธาย ภาสิตํ “กปณํ มานุสกํ รชฺชํ ทิพฺพํ สุขํ อุปนิธายา”ติ.}

\csromangathameasure{ปาณํ น หญฺเญ น จทินฺนมาทิเย, \\ มุสา น ภาเส น จ มชฺชโป สิยา. \\ อพฺรหฺมจริยา วิรเมยฺย เมถุนา, \\ รตฺติํ น ภุญฺเชยฺย วิกาลโภชนํ. \\ มาลํ น ธาเร น จ คนฺธมาจเร, \\ มญฺเจ ฉมายํ ว สเยถ สนฺถเต. \\ เอตํ หิ อฏฺฐงฺคิกมาหุโปสถํ, \\ พุทฺเธน ทุกฺขนฺตคุนา ปกาสิตํ. \\ จนฺโท จ สุริโย จ อุโภ สุทสฺสนา, \\ โอภาสยํ อนุปริยนฺติ ยาวตา. \\ ตโมนุทา เต ปน อนฺตลิกฺขคา, \\ นเภ ปภาสนฺติ ทิสาวิโรจนา. \\ เอตสฺมิํ ยํ วิชฺชติ อนฺตเร ธนํ, \\ มุตฺตา มณิ เวฬุริยญฺจ ภทฺทกํ. \\ สิงฺคีสุวณฺณํ อถ วาปิ กญฺจนํ, \\ ยํ ชาตรูปํ หฏกนฺติ วุจฺจติ.}
\csromangathagroup{\csromangathastanza{ปาณํ น หญฺเญ น จทินฺนมาทิเย, \\ มุสา น ภาเส น จ มชฺชโป สิยา. \\ อพฺรหฺมจริยา วิรเมยฺย เมถุนา, \\ รตฺติํ น ภุญฺเชยฺย วิกาลโภชนํ.}\csromangathastanza{มาลํ น ธาเร น จ คนฺธมาจเร, \\ มญฺเจ ฉมายํ ว สเยถ สนฺถเต. \\ เอตํ หิ อฏฺฐงฺคิกมาหุโปสถํ, \\ พุทฺเธน ทุกฺขนฺตคุนา ปกาสิตํ.}\csromangathastanza{จนฺโท จ สุริโย จ อุโภ สุทสฺสนา, \\ โอภาสยํ อนุปริยนฺติ ยาวตา. \\ ตโมนุทา เต ปน อนฺตลิกฺขคา, \\ นเภ ปภาสนฺติ ทิสาวิโรจนา.}\csromangathastanza{เอตสฺมิํ ยํ วิชฺชติ อนฺตเร ธนํ, \\ มุตฺตา มณิ เวฬุริยญฺจ ภทฺทกํ. \\ สิงฺคีสุวณฺณํ อถ วาปิ กญฺจนํ, \\ ยํ ชาตรูปํ หฏกนฺติ วุจฺจติ.}}

\chapterheadpageread{5. อุโปสถวคฺค}
\csromangathameasure{อฏฺฐงฺคุเปตสฺส อุโปสถสฺส, \\ กลมฺปิ เต นานุภวนฺติ โสฬสิํ. \\ จนฺทปฺปภา ตารคณา จ สพฺเพ. \\ ตสฺมา หิ นารี จ นโร จ สีลวา, \\ อฏฺฐงฺคุเปตํ อุปวสฺสุโปสถํ. \\ ปุญฺญานิ กตฺวาน สุขุทฺรยานิ, \\ อนินฺทิตา สคฺคมุเปนฺติ ฐานนฺติ.ตติยํ.}
\csromangathagroup{\csromangathastanza{\csromanfolio{87}อฏฺฐงฺคุเปตสฺส อุโปสถสฺส, \\ กลมฺปิ เต นานุภวนฺติ โสฬสิํ. \\ จนฺทปฺปภา ตารคณา จ สพฺเพ. \\ ตสฺมา หิ นารี จ นโร จ สีลวา,}\csromangathastanza{อฏฺฐงฺคุเปตํ อุปวสฺสุโปสถํ. \\ ปุญฺญานิ กตฺวาน สุขุทฺรยานิ, \\ อนินฺทิตา สคฺคมุเปนฺติ ฐานนฺติ.\csromanspacer{}\csromanspacer{}ตติยํ.}}

\csromanheader{2}{4. วาเสฏฺฐสุตฺต}
\proseitem{44}{เอกํ สมยํ ภควา เวสาลิยํ วิหรติ มหาวเน กูฏาคารสาลายํ.\csromanspacer{} อถ โข วาเสฏฺโฐ อุปาสโก เยน ภควา เตนุปสงฺกมิ, อุปสงฺกมิตฺวา ภควนฺตํ อภิวาเทตฺวา เอกมนฺตํ นีสีทิ, เอกมนฺตํ นิสินฺนํ โข วาเสฏฺฐํ อุปาสกํ ภควา เอตทโวจ “อฏฺฐงฺคสมนฺนาคโต วาเสฏฺฐ อุโปสโถ อุปวุตฺโถ มหปฺผโล โหติ ฯเปฯ อนินฺทิตา สคฺคมุเปนฺติ ฐานนฺ”ติ.}

\prose{เอวํ วุตฺเต วาเสฏฺโฐ อุปาสโก ภควนฺตํ เอตทโวจ “ปิยา เม ภนฺเต ญาติสาโลหิตา อฏฺฐงฺคสมนฺนาคตํ อุโปสถํ อุปวเสยฺยุํ, ปิยานมฺปิ เม อสฺส ญาติสาโลหิตานํ ทีฆรตฺตํ หิตาย สุขาย.\csromanspacer{} สพฺเพ เจปิ ภนฺเต ขตฺติยา อฏฺฐงฺคสมนฺนาคตํ อุโปสถํ อุปวเสยฺยุํ, สพฺเพสมฺปิสฺส ขตฺติยานํ ทีฆรตฺตํ หิตาย สุขาย.\csromanspacer{} สพฺเพ เจปิ ภนฺเต พฺราหฺมณา ฯเปฯ เวสฺสา.\csromanspacer{} สุทฺทา อฏฺฐงฺคสมนฺนาคตํ อุโปสถํ อุปวเสยฺยุํ, สพฺเพสมฺปิสฺส สุทฺทานํ ทีฆรตฺตํ หิตาย สุขายา”ติ.}

\prose{เอวเมตํ วาเสฏฺฐ, เอวเมตํ วาเสฏฺฐ, สพฺเพ เจปิ วาเสฏฺฐ ขตฺติยา อฏฺฐงฺคสมนฺนาคตํ อุโปสถํ อุปวเสยฺยุํ, สพฺเพสมฺปิสฺส ขตฺติยานํ ทีฆรตฺตํ หิตาย สุขาย.\csromanspacer{} สพฺเพ เจปิ วาเสฏฺฐ พฺราหฺมณา ฯเปฯ เวสฺสา.\csromanspacer{} สุทฺทา อฏฺฐงฺคสมนฺนาคตํ อุโปสถํ อุปวเสยฺยุํ, สพฺเพสมฺปิสฺส สุทฺทานํ ทีฆรตฺตํ หิตาย สุขาย.\csromanspacer{} สเทวโก เจปิ วาเสฏฺฐ โลโก สมารโก สพฺรหฺมโก สสฺสมณพฺราหฺมณี ปชา สเทวมนุสฺสา อฏฺฐงฺคสมนฺนาคตํ อุโปสถํ อุปวเสยฺยุํ\footnote{อุปวเสยฺย (?)}, สเทวกสฺสปิสฺส\footnote{สเทวกสฺส (สพฺพตฺถ) อํ 1. 515 ปิฏฺเฐ; ม 3. 46 ปิฏฺเฐ จ ปสฺสิตพฺพํ.} โลกสฺส \csromanfolio{88}สมารกสฺส สพฺรหฺมกสฺส สสฺสมณพฺราหฺมณิยา ปชาย สเทวมนุสฺสาย ทีฆรตฺตํ หิตาย สุขาย.\csromanspacer{} อิเม เจปิ วาเสฏฺฐ มหาสาลา อฏฺฐงฺคสมนฺนาคตํ อุโปสถํ อุปวเสยฺยุํ, อิเมสมฺปิสฺส มหาสาลานํ ทีฆรตฺตํ หิตาย สุขาย ( )\footnote{(สเจ เจเตยฺยุํ) กตฺถจิ อตฺถิ. อํ 1. 515 ปิฏฺเฐ ปสฺสิตพฺพํ.}, โก ปน วาโท มนุสฺสภูตสฺสาติ.\csromanspacer{}\csromanspacer{}จตุตฺถํ.}

\csromanheader{2}{5. โพชฺฌสุตฺต}
\proseitem{45}{เอกํ สมยํ ภควา สาวตฺถิยํ วิหรติ เชตวเน อนาถปิณฺฑิกสฺส อาราเม.\csromanspacer{} อถ โข โพชฺฌา อุปาสิกา เยน ภควา เตนุปสงฺกมิ, อุปสงฺกมิตฺวา ภควนฺตํ อภิวาเทตฺวา เอกมนฺตํ นิสีทิ, เอกมนฺตํ นิสินฺนํ โข โพชฺฌํ อุปาสิกํ ภควา เอตทโวจ\csromandash{}}

\prose{อฏฺฐงฺคสมนฺนาคโต โพชฺเฌ อุโปสโถ อุปวุตฺโถ มหปฺผโล โหติ มหานิสํโส มหาชุติโก มหาวิปฺผาโร.\csromanspacer{} กถํ อุปวุตฺโถ จ โพชฺเฌ อฏฺฐงฺคสมนฺนาคโต อุโปสโถ มหปฺผโล โหติ มหานิสํโส มหาชุติโก มหาวิปฺผาโร? อิธ โพชฺเฌ อริยสาวโก อิติ ปฏิสญฺจิกฺขติ “ยาวชีวํ อรหนฺโต ปาณาติปาตํ ปหาย ปาณาติปาตา ปฏิวิรตา นิหิตทณฺฑา นิหิตสตฺถา ลชฺชี ทยาปนฺนา สพฺพปาณภูตหิตานุกมฺปิโน วิหรนฺติ, อหมฺปชฺช อิมญฺจ รตฺติํ อิมญฺจ ทิวสํ ปาณาติปาตํ ปหาย ปาณาติปาตา ปฏิวิรโต นิหิตทณฺโฑ นิหิตสตฺโถ ลชฺชี ทยาปนฺโน สพฺพปาณภูตหิตานุกมฺปี วิหรามิ.\csromanspacer{} อิมินาปงฺเคน อรหตํ อนุกโรมิ, อุโปสโถ จ เม อุปวุตฺโถ ภวิสฺสตี”ติ.\csromanspacer{} อิมินา ปฐเมน องฺเคน สมนฺนาคโต โหติ ฯเปฯ.}

\prose{“ยาวชีวํ อรหนฺโต อุจฺจาสยนมหาสยนํ ปหาย อุจฺจาสยนมหาสยนา ปฏิวิรตา นีจเสยฺยํ กปฺเปนฺติ มญฺจเก วา ติณสนฺถารเก วา, อหมฺปชฺช อิมญฺจ รตฺติํ อิมญฺจ ทิวสํ อุจฺจาสยนมหาสยนํ ปหาย อุจฺจาสยนมหาสยนา ปฏิวิรโต นีจเสยฺยํ กปฺเปมิ มญฺจเก วา ติณสนฺถารเก วา.\csromanspacer{} อิมินาปงฺเคน อรหตํ อนุกโรมิ, อุโปสโถ จ เม อุปวุตฺโถ ภวิสฺสตี”ติ.\csromanspacer{} อิมินา อฏฺฐเมน องฺเคน สมนฺนาคโต โหติ.\csromanspacer{} เอวํ อุปวุตฺโถ โข โพชฺเฌ อฏฺฐงฺคสมนฺนาคโต อุโปสโถ มหปฺผโล โหติ มหานิสํโส มหาชุติโก มหาวิปฺผาโร.}

\chapterheadpageread{5. อุโปสถวคฺค}
\prose{\csromanfolio{89}กีวมหปฺผโล โหติ กีวมหานิสํโส กีวมหาชุติโก กีวมหาวิปฺผาโร? เสยฺยถาปิ โพชฺเฌ โย อิเมสํ โสฬสนฺนํ มหาชนปทานํ ปหูตรตฺตรตนานํ อิสฺสริยาธิปจฺจํ รชฺชํ กาเรยฺย.\csromanspacer{} เสยฺยถิทํ, องฺคานํ มคธานํ กาสีนํ โกสลานํ วชฺชีนํ มลฺลานํ เจตีนํ วงฺคานํ กุรูนํ ปญฺจาลานํ มจฺฉานํ สูรเสนานํ อสฺสกานํ อวนฺตีนํ คนฺธารานํ กมฺโพชานํ, อฏฺฐงฺคสมนฺนาคตสฺส อุโปสถสฺส เอตํ กลํ นาคฺฆติ โสฬสิํ.\csromanspacer{} ตํ กิสฺส เหตุ? กปณํ โพชฺเฌ มานุสกํ รชฺชํ ทิพฺพํ สุขํ อุปนิธาย.}

\prose{ยานิ โพชฺเฌ มานุสกานิ ปญฺญาส วสฺสานิ, จาตุมหาราชิกานํ เทวานํ เอโส เอโก รตฺตินฺทิโว, ตาย รตฺติยา ติํสรตฺติโย มาโส, เตน มาเสน ทฺวาทสมาสิโย สํวจฺฉโร, เตน สํวจฺฉเรน ทิพฺพานิ ปญฺจ วสฺสสตานิ จาตุมหาราชิกานํ เทวานํ อายุปฺปมาณํ.\csromanspacer{} ฐานํ โข ปเนตํ โพชฺเฌ วิชฺชติ, ยํ อิเธกจฺโจ อิตฺถี วา ปุริโส วา อฏฺฐงฺคสมนฺนาคตํ อุโปสถํ อุปวสิตฺวา กายสฺส เภทา ปรํ มรณา จาตุมหาราชิกานํ เทวานํ สหพฺยตํ อุปปชฺเชยฺย.\csromanspacer{} อิทํ โข ปเนตํ โพชฺเฌ สนฺธาย ภาสิตํ กปณํ มานุสกํ รชฺชํ ทิพฺพํ สุขํ อุปนิธาย.}

\prose{ยํ โพชฺเฌ มานุสกํ วสฺสสตํ ฯเปฯ ตานิ โพชฺเฌ มานุสกานิ เทฺว วสฺสสตานิ ฯเปฯ จตฺตาริ วสฺสสตานิ ฯเปฯ อฏฺฐ วสฺสสตานิ ฯเปฯ โสฬส วสฺสสตานิ, ปรนิมฺมิตวสวตฺตีนํ เทวานํ เอโส เอโก รตฺตินฺทิโว, ตาย รตฺติยา ติํสรตฺติโย มาโส, เตน มาเสน ทฺวาทสมาสิโย สํวจฺฉโร, เตน สํวจฺฉเรน ทิพฺพานิ โสฬส วสฺสสหสฺสานิ ปรนิมฺมิตวสวตฺตีนํ เทวานํ อายุปฺปมาณํ.\csromanspacer{} ฐานํ โข ปเนตํ โพชฺเฌ วิชฺชติ, ยํ อิเธกจฺโจ อิตฺถี วา ปุริโส วา อฏฺฐงฺคสมนฺนาคตํ อุโปสถํ อุปวสิตฺวา กายสฺส เภทา ปรํ มรณา ปรนิมฺมิตวสวตฺตีนํ เทวานํ สหพฺยตํ อุปปชฺเชยฺย.\csromanspacer{} อิทํ โข ปเนตํ โพชฺเฌ สนฺธาย ภาสิตํ “กปณํ มานุสกํ รชฺชํ ทิพฺพํ สุขํ อุปนิธายา”ติ.}

\csromangathameasure{ปาณํ น หญฺเญ น จทินฺนมาทิเย, \\ มุสา น ภาเส น จ มชฺชโป สิยา. \\ อพฺรหฺมจริยา วิรเมยฺย เมถุนา, \\ รตฺติํ น ภุญฺเชยฺย วิกาลโภชนํ. \\ มาลํ น ธาเร น จ คนฺธมาจเร, \\ มญฺเจ ฉมายํ ว สเยถ สนฺถเต. \\ เอตํ หิ อฏฺฐงฺคิกมาหุโปสถํ, \\ พุทฺเธน ทุกฺขนฺตคุนา ปกาสิตํ. \\ จนฺโท จ สุริโย จ อุโภ สุทสฺสนา, \\ โอภาสยํ อนุปริยนฺติ ยาวตา. \\ ตโมนุทา เต ปน อนฺตลิกฺขคา, \\ นเภ ปภาสนฺติ ทิสาวิโรจนา. \\ เอตสฺมิํ ยํ วิชฺชติ อนฺตเร ธนํ, \\ มุตฺตา มณิ เวฬุริยญฺจ ภทฺทกํ. \\ สิงฺคีสุวณฺณํ อถ วาปิ กญฺจนํ, \\ ยํ ชาตรูปํ หฏกนฺติ วุจฺจติ. \\ อฏฺฐงฺคุเปตสฺส อุโปสถสฺส, \\ กลมฺปิ เต นานุภวนฺติ โสฬสิํ. \\ จนฺทปฺปภา ตารคณา จ สพฺเพ. \\ ตสฺมา หิ นารี จ นโร จ สีลวา, \\ อฏฺฐงฺคุเปตํ อุปวสฺสุโปสถํ. \\ ปุญฺญานิ กตฺวาน สุขุทฺรยานิ, \\ อนินฺทิตา สคฺคมุเปนฺติ ฐานนฺติ.ปญฺจมํ.}
\csromangathagroup{\csromangathastanza{ปาณํ น หญฺเญ น จทินฺนมาทิเย, \\ มุสา น ภาเส น จ มชฺชโป สิยา. \\ อพฺรหฺมจริยา วิรเมยฺย เมถุนา, \\ รตฺติํ น ภุญฺเชยฺย วิกาลโภชนํ.}\csromangathastanza{\csromanfolio{90}มาลํ น ธาเร น จ คนฺธมาจเร, \\ มญฺเจ ฉมายํ ว สเยถ สนฺถเต. \\ เอตํ หิ อฏฺฐงฺคิกมาหุโปสถํ, \\ พุทฺเธน ทุกฺขนฺตคุนา ปกาสิตํ.}\csromangathastanza{จนฺโท จ สุริโย จ อุโภ สุทสฺสนา, \\ โอภาสยํ อนุปริยนฺติ ยาวตา. \\ ตโมนุทา เต ปน อนฺตลิกฺขคา, \\ นเภ ปภาสนฺติ ทิสาวิโรจนา.}\csromangathastanza{เอตสฺมิํ ยํ วิชฺชติ อนฺตเร ธนํ, \\ มุตฺตา มณิ เวฬุริยญฺจ ภทฺทกํ. \\ สิงฺคีสุวณฺณํ อถ วาปิ กญฺจนํ, \\ ยํ ชาตรูปํ หฏกนฺติ วุจฺจติ.}\csromangathastanza{อฏฺฐงฺคุเปตสฺส อุโปสถสฺส, \\ กลมฺปิ เต นานุภวนฺติ โสฬสิํ. \\ จนฺทปฺปภา ตารคณา จ สพฺเพ. \\ ตสฺมา หิ นารี จ นโร จ สีลวา,}\csromangathastanza{อฏฺฐงฺคุเปตํ อุปวสฺสุโปสถํ. \\ ปุญฺญานิ กตฺวาน สุขุทฺรยานิ, \\ อนินฺทิตา สคฺคมุเปนฺติ ฐานนฺติ.\csromanspacer{}\csromanspacer{}ปญฺจมํ.}}

\csromanheader{2}{6. อนุรุทฺธสุตฺต}
\proseitem{46}{เอกํ สมยํ ภควา โกสมฺพิยํ วิหรติ โฆสิตาราเม.\csromanspacer{} เตน โข ปน สมเยน อายสฺมา อนุรุทฺโธ ทิวาวิหารํ คโต โหติ ปฏิสลฺลีโน.\csromanspacer{} อถ โข สมฺพหุลา มนาปกายิกา เทวตา เยนายสฺมา อนุรุทฺโธ เตนุปสงฺกมิํสุ, อุปสงฺกมิตฺวา อายสฺมนฺตํ อนุรุทฺธํ อภิวาเทตฺวา เอกมนฺตํ อฏฺฐํสุ, เอกมนฺตํ ฐิตา โข ตา เทวตา อายสฺมนฺตํ อนุรุทฺธํ เอตทโวจุํ “มยํ ภนฺเต อนุรุทฺธ มนาปกายิกา นาม เทวตา ตีสุ ฐาเนสุ อิสฺสริยํ กาเรม วสํ วตฺเตม.\csromanspacer{} มยํ ภนฺเต อนุรุทฺธ ยาทิสกํ วณฺณํ อากงฺขาม, ตาทิสกํ วณฺณํ ฐานโส ปฏิลภาม.\csromanspacer{} ยาทิสกํ สรํ อากงฺขาม, ตาทิสกํ สรํ ฐานโส ปฏิลภาม.\csromanspacer{} ยาทิสกํ สุขํ}

\vspace{\tipitakaparskip}
\csromancenter{อุโปสถวคฺค อากงฺขาม, ตาทิสกํ สุขํ ฐานโส ปฏิลภาม.\csromanspacer{} มยํ ภนฺเต อนุรุทฺธ มนาปกายิกา นาม เทวตา อิเมสุ ตีสุ ฐาเนสุ อิสฺสริยํ กาเรม วสํ วตฺเตมา”ติ.}
\prose{\csromanfolio{91}อถ โข อายสฺมโต อนุรุทฺธสฺส เอตทโหสิ “อโห วติมา เทวตา สพฺพาว นีลา อสฺสุ นีลวณฺณา นีลวตฺถา นีลาลงฺการา”ติ.\csromanspacer{} อถ โข ตา เทวตา อายสฺมโต อนุรุทฺธสฺส จิตฺตมญฺญาย สพฺพาว นีลา อเหสุํ นีลวณฺณา นีลวตฺถา นีลาลงฺการา.}

\prose{อถ โข อายสฺมโต อนุรุทฺธสฺส เอตทโหสิ “อโห วติมา เทวตา สพฺพาว ปีตา อสฺสุ ฯเปฯ สพฺพาว โลหิตกา อสฺสุ.\csromanspacer{} สพฺพาว โอทาตา อสฺสุ โอทาตวณฺณา โอทาตวตฺถา โอทาตาลงฺการา”ติ.\csromanspacer{} อถ โข ตา เทวตา อายสฺมโต อนุรุทฺธสฺส จิตฺตมญฺญาย สพฺพาว โอทาตา อเหสุํ โอทาตวณฺณา โอทาตวตฺถา โอทาตาลงฺการา.}

\prose{อถ โข ตา เทวตา เอกา จ\footnote{เอกา (สี), เอกาว (สฺยา, อิ)} คายิ, เอกา จ\footnote{เอกา ปน (สี), เอกาว (สฺยา, อิ)} นจฺจิ, เอกา จ1 อจฺฉรํ วาเทสิ.\csromanspacer{} เสยฺยถาปิ นาม ปญฺจงฺคิกสฺส ตูริยสฺส\footnote{ตุริยสฺส (สี, สฺยา, อิ)} สุวินีตสฺส สุปฺปฏิปตาฬิตสฺส กุสเลหิ สุสมนฺนาหตสฺส สทฺโท โหติ วคฺคุ จ รชนีโย จ กมนีโย จ เปมนีโย จ มทนีโย จ.\csromanspacer{} เอวเมวํ ตาสํ เทวตานํ อลงฺการานํ สทฺโท โหติ วคฺคุ จ รชนีโย จ กมนีโย จ เปมนีโย จ มทนีโย จ.\csromanspacer{} อถ โข อายสฺมา อนุรุทฺโธ อินฺทฺริยานิ โอกฺขิปิ.}

\prose{อถ โข ตา เทวตา “น ขฺวยฺโย อนุรุทฺโธ สาทิยตี”ติ\footnote{สาทยตีติ (สทฺทนีติธาตุมาลา)} ตตฺเถวนฺตรธายิํสุ.\csromanspacer{} อถ โข อายสฺมา อนุรุทฺโธ สายนฺหสมยํ ปฏิสลฺลานา วุฏฺฐิโต เยน ภควา เตนุปสงฺกมิ, อุปสงฺกมิตฺวา ภควนฺตํ อภิวาเทตฺวา เอกมนฺตํ นิสีทิ, เอกมนฺตํ นิสินฺโน โข อายสฺมา อนุรุทฺโธ ภควนฺตํ เอตทโวจ\csromandash{}}

\prose{อิธาหํ ภนฺเต ทิวาวิหารํ คโต โหมิ ปฏิสลฺลีโน, อถ โข ภนฺเต สมฺพหุลา มนาปกายิกา เทวตา เยนาหํ เตนุปสงฺกมิํสุ, อุปสงฺกมิตฺวา มํ อภิวาเทตฺวา เอกมนฺตํ อฏฺฐํสุ, เอกมนฺตํ ฐิตา โข ภนฺเต \csromanfolio{92}ตา เทวตา มํ เอตทโวจุํ “มยํ ภนฺเต อนุรุทฺธ มนาปกายิกา นาม เทวตา ตีสุ ฐาเนสุ อิสฺสริยํ กาเรม วสํ วตฺเตม.\csromanspacer{} มยํ ภนฺเต อนุรุทฺธ ยาทิสกํ วณฺณํ อากงฺขาม, ตาทิสกํ วณฺณํ ฐานโส ปฏิลภาม.\csromanspacer{} ยาทิสกํ สรํ อากงฺขาม, ตาทิสกํ สรํ ฐานโส ปฏิลภาม.\csromanspacer{} ยาทิสกํ สุขํ อากงฺขาม, ตาทิสกํ สุขํ ฐานโส ปฏิลภาม.\csromanspacer{} มยํ ภนฺเต อนุรุทฺธ มนาปกายิกา นาม เทวตา อิเมสุ ตีสุ ฐาเนสุ อิสฺสริยํ กาเรม วสํ วตฺเตมา”ติ.\csromanspacer{} ตสฺส มยฺหํ ภนฺเต เอตทโหสิ “อโห วติมา เทวตา สพฺพาว นีลา อสฺสุ นีลวณฺณา นีลวตฺถา นีลาลงฺการา”ติ.\csromanspacer{} อถ โข ภนฺเต ตา เทวตา มม จิตฺตมญฺญาย สพฺพาว นีลา อเหสุํ นีลวณฺณา นีลวตฺถา นีลาลงฺการา.}

\prose{ตสฺส มยฺหํ ภนฺเต เอตทโหสิ “อโห วติมา เทวตา สพฺพาว ปีตา อสฺสุ ฯเปฯ สพฺพาว โลหิตกา อสฺสุ ฯเปฯ สพฺพาว โอทาตา อสฺสุ โอทาตวณฺณา โอทาตวตฺถา โอทาตาลงฺการา”ติ.\csromanspacer{} อถ โข ภนฺเต ตา เทวตา มม จิตฺตมญฺญาย สพฺพาว โอทาตา อเหสุํ โอทาตวณฺณา โอทาตวตฺถา โอทาตาลงฺการา.}

\prose{อถ โข ภนฺเต ตา เทวตา เอกา จ คายิ, เอกา จ นจฺจิ, เอกา จ อจฺฉรํ วาเทสิ.\csromanspacer{} เสยฺยถาปิ นาม ปญฺจงฺคิกสฺส ตูริยสฺส สุวินีตสฺส สุปฺปฏิปตาฬิตสฺส กุสเลหิ สุสมนฺนาหตสฺส สทฺโท โหติ วคฺคุ จ รชนีโย จ กมนีโย จ เปมนีโย จ มทนีโย จ.เอวเมวํ ตาสํ เทวตานํ อลงฺการานํ สทฺโท โหติ วคฺคุ จ รชนีโย จ กมนีโย จ เปมนีโย จ มทนีโย จ.\csromanspacer{} อถ ขฺวาหํ ภนฺเต อินฺทฺริยานิ โอกฺขิปิ.}

\prose{อถ โข ภนฺเต ตา เทวตา “น ขฺวยฺโย อนุรุทฺโธ สาทิยตี”ติ ตตฺเถวนฺตรธายิํสุ.\csromanspacer{} กติหิ นุ โข ภนฺเต ธมฺเมหิ สมนฺนาคโต มาตุคาโม กายสฺส เภทา ปรํ มรณา มนาปกายิกานํ เทวานํ สหพฺยตํ อุปปชฺชตีติ.}

\prose{อฏฺฐหิ โข อนุรุทฺธ ธมฺเมหิ สมนฺนาคโต มาตุคาโม กายสฺส เภทา ปรํ มรณา มนาปกายิกานํ เทวานํ สหพฺยตํ อุปปชฺชติ.\csromanspacer{} กตเมหิ อฏฺฐหิ? อิธ อนุรุทฺธ มาตุคาโม, ยสฺส มาตาปิตโร ภตฺตุโน เทนฺติ อตฺถกามา หิเตสิโน อนุกมฺปกา อนุกมฺปํ อุปาทาย,}

\vspace{\tipitakaparskip}
\csromancenter{อุโปสถวคฺค ตสฺส โหติ ปุพฺพุฏฺฐายินี ปจฺฉานิปาตินี กิํการปฏิสฺสาวินี มนาปจารินี ปิยวาทินี. (1)}
\prose{\csromanfolio{93}เย เต ภตฺตุ ครุโน\footnote{คุรุโน (ก)} โหนฺติ “มาตา”ติ วา “ปิตา”ติ วา “สมณพฺราหฺมณา”ติ วา, เต สกฺกโรติ ครุํ กโรติ\footnote{ครุกโรติ (สี, สฺยา, อิ) 3. ตตฺรูปายาย (สี), อํ 1. 344; อุปริ 539 ปิฏฺเฐปิ. 4. ปจฺจเยน (สฺยา), ปจฺจตฺตํเสน (ก) อํ 2. 32 ปิฏฺเฐปิ.} มาเนติ ปูเชติ, อพฺภาคเต จ อาสโนทเกน ปฏิปูเชติ. (2)}

\prose{เย เต ภตฺตุ อพฺภนฺตรา กมฺมนฺตา “อุณฺณา”ติ วา “กปฺปาสา”ติ วา, ตตฺถ ทกฺขา โหติ อนลสา, ตตฺรุปายาย3 วีมํสาย สมนฺนาคตา อลํ กาตุํ อลํ สํวิธาตุํ. (3)}

\prose{โย โส ภตฺตุ อพฺภนฺตโร อนฺโตชโน “ทาสา”ติ วา “เปสฺสา”ติ วา “กมฺมกรา”ติ วา, เตสํ กตญฺจ กตโต ชานาติ, อกตญฺจ อกตโต ชานาติ, คิลานกานญฺจ พลาพลํ ชานาติ, ขาทนียํ โภชนียญฺจสฺส ปจฺจํเสน4 สํวิภชติ. (4)}

\prose{ยํ ภตฺตุ อาหรติ ธนํ วา ธญฺญํ วา ชาตรูปํ วา, ตํ อารกฺเขน คุตฺติยา สมฺปาเทติ, ตตฺถ จ โหติ อธุตฺตี อเถนี อโสณฺฑี อวินาสิกา. (5)}

\prose{อุปาสิกา โข ปน โหติ พุทฺธํ สรณํ คตา, ธมฺมํ สรณํ คตา, สํฆํ สรณํ คตา. (6)}

\prose{สีลวตี โข ปน โหติ ปาณาติปาตา ปฏิวิรตา, อทินฺนาทานา ปฏิวิรตา, กาเมสุมิจฺฉาจารา ปฏิวิรตา, มุสาวาทา ปฏิวิรตา, สุราเมรยมชฺชปมาทฏฺฐานา ปฏิวิรตา. (7)}

\prose{จาควตี โข ปน โหติ วิคตมลมจฺเฉเรน เจตสา อคารํ อชฺฌาวสติ มุตฺตจาคา\footnote{มุตฺตจาคี (สฺยา) 6. ปยตปาณิ (สี), ปยตปาณี (สฺยา, อิ, ก)} ปยตปาณินี6 โวสคฺครตา ยาจโยคา ทานสํวิภาครตา. (8)}

\prose{อิเมหิ โข อนุรุทฺธ อฏฺฐหิ ธมฺเมหิ สมนฺนาคโต มาตุคาโม กายสฺส เภทา ปรํ มรณา มนาปกายิกานํ เทวานํ สหพฺยตํ อุปปชฺชตีติ.}

\csromangathameasure{โย นํ ภรติ สพฺพทา, นิจฺจํ อาตาปิ อุสฺสุโก. \\ ตํ สพฺพกามทํ โปสํ, ภตฺตารํ นาติมญฺญติ. \\ น จาปิ โสตฺถิ ภตฺตารํ, อิสฺสาวาเทน โรสเย. \\ ภตฺตุ จ ครุโน สพฺเพ, ปฏิปูเชติ ปณฺฑิตา. \\ อุฏฺฐาหิกา อนลสา, สงฺคหิตปริชฺชนา. \\ ภตฺตุ มนาปํ จรติ, สมฺภตํ อนุรกฺขติ. \\ ยา เอวํ วตฺตติ นารี, ภตฺตุ ฉนฺทวสานุคา. \\ มนาปา นาม เต เทวา, ยตฺถ สา อุปปชฺชตีติ.ฉฏฺฐํ.}
\csromangathagroup{\csromangathastanza{\csromanfolio{94}โย นํ ภรติ สพฺพทา, นิจฺจํ อาตาปิ อุสฺสุโก. \\ ตํ สพฺพกามทํ\footnote{ตํ สพฺพกามหรํ (สี, สฺยา, อิ) สพฺพกามหรํ (อํ 2. 32 ปิฏฺเฐ)} โปสํ, ภตฺตารํ นาติมญฺญติ.}\csromangathastanza{น จาปิ โสตฺถิ ภตฺตารํ, อิสฺสาวาเทน โรสเย. \\ ภตฺตุ จ ครุโน สพฺเพ, ปฏิปูเชติ ปณฺฑิตา.}\csromangathastanza{อุฏฺฐาหิกา\footnote{อุฏฺฐายิกา (ก)} อนลสา, สงฺคหิตปริชฺชนา. \\ ภตฺตุ มนาปํ จรติ, สมฺภตํ อนุรกฺขติ.}\csromangathastanza{ยา เอวํ วตฺตติ นารี, ภตฺตุ ฉนฺทวสานุคา. \\ มนาปา นาม เต\footnote{มนาปกายิกา (สี, ก)} เทวา, ยตฺถ สา อุปปชฺชตีติ.\csromanspacer{}\csromanspacer{}ฉฏฺฐํ.}}

\csromanheader{2}{7. ทุติยวิสาขาสุตฺต}
\proseitem{47}{เอกํ สมยํ ภควา สาวตฺถิยํ วิหรติ ปุพฺพาราเม มิคารมาตุปาสาเท.\csromanspacer{} อถ โข วิสาขา มิคารมาตา ฯเปฯ เอกมนฺตํ นิสินฺนํ โข วิสาขํ มิคารมาตรํ ภควา เอตทโวจ\csromandash{}}

\prose{อฏฺฐหิ โข วิสาเข ธมฺเมหิ สมนฺนาคโต มาตุคาโม กายสฺส เภทา ปรํ มรณา มนาปกายิกานํ เทวานํ สหพฺยตํ อุปปชฺชติ.\csromanspacer{} กตเมหิ อฏฺฐหิ? อิธ วิสาเข มาตุคาโม, ยสฺส มาตาปิตโร ภตฺตุโน เทนฺติ อตฺถกามา หิเตสิโน อนุกมฺปกา อนุกมฺปํ อุปาทาย, ตสฺส โหติ ปุพฺพุฏฺฐายินี ปจฺฉานิปาตินี กิํการปฏิสฺสาวินี มนาปจารินี ปิยวาทินี ฯเปฯ.}

\prose{จาควตี โข ปน โหติ, วิคตมลจฺเฉเรน เจตสา อคารํ อชฺฌาวสติ มุตฺตจาคา ปยตปาณินี โวสคฺครตา ยาจโยคา ทานสํวิภาครตา.\csromanspacer{} อิเมหิ โข วิสาเข อฏฺฐหิ ธมฺเมหิ สมนฺนาคโต มาตุคาโม กายสฺส เภทา ปรํ มรณา มนาปกายิกานํ เทวานํ สหพฺยตํ อุปปชฺชตีติ.}

\csromangathameasure{โย นํ ภรติ สพฺพทา, นิจฺจํ อาตาปิ อุสฺสุโก. \\ ตํ สพฺพกามทํ โปสํ, ภตฺตารํ นาติมญฺญติ.}
\csromangathagroup{\csromangathastanza{โย นํ ภรติ สพฺพทา, นิจฺจํ อาตาปิ อุสฺสุโก. \\ ตํ สพฺพกามทํ โปสํ, ภตฺตารํ นาติมญฺญติ.}}

\chapterheadpageread{5. อุโปสถวคฺค}
\csromangathameasure{น จาปิ โสตฺถิ ภตฺตารํ, อิสฺสาวาเทน โรสเย. \\ ภตฺตุ จ ครุโน สพฺเพ, ปฏิปูเชติ ปณฺฑิตา. \\ อุฏฺฐาหิกา อนลสา, สงฺคหิตปริชฺชนา. \\ ภตฺตุ มนาปํ จรติ, สมฺภตํ อนุรกฺขติ. \\ ยา เอวํ วตฺตติ นารี, ภตฺตุ ฉนฺทวสานุคา. \\ มนาปา นาม เต เทวา, ยตฺถ สา อุปปชฺชตีติ.สตฺตมํ.}
\csromangathagroup{\csromangathastanza{\csromanfolio{95}น จาปิ โสตฺถิ ภตฺตารํ, อิสฺสาวาเทน โรสเย. \\ ภตฺตุ จ ครุโน สพฺเพ, ปฏิปูเชติ ปณฺฑิตา.}\csromangathastanza{อุฏฺฐาหิกา อนลสา, สงฺคหิตปริชฺชนา. \\ ภตฺตุ มนาปํ จรติ, สมฺภตํ อนุรกฺขติ.}\csromangathastanza{ยา เอวํ วตฺตติ นารี, ภตฺตุ ฉนฺทวสานุคา. \\ มนาปา นาม เต\footnote{มนาปกายิกา (สี, ก)} เทวา, ยตฺถ สา อุปปชฺชตีติ.\csromanspacer{}\csromanspacer{}สตฺตมํ.}}

\csromanheader{2}{8. นกุลมาตาสุตฺต}
\proseitem{48}{เอกํ สมยํ ภควา ภคฺเคสุ วิหรติ สุสุมารคิเร\footnote{สุํสุมารคิเร (สี, สฺยา, อิ)} เภสกฬาวเน มิคทาเย.\csromanspacer{} อถ โข นกุลมาตา คหปตานี เยน ภควา เตนุปสงฺกมิ, อุปสงฺกมิตฺวา ฯเปฯ เอกมนฺตํ นิสินฺนํ โข นกุลมาตรํ คหปตานิํ ภควา เอตทโวจ\csromandash{}}

\prose{อฏฺฐหิ โข นกุลมาเต ธมฺเมหิ สมนฺนาคโต มาตุคาโม กายสฺส เภทา ปรํ มรณา มนาปกายิกานํ เทวานํ สหพฺยตํ อุปปชฺชติ.\csromanspacer{} กตเมหิ อฏฺฐหิ? อิธ นกุลมาเต มาตุคาโม, ยสฺส มาตาปิตโร ภตฺตุโน เทนฺติ อตฺถกามา หิเตสิโน อนุกมฺปกา อนุกมฺปํ อุปาทาย, ตสฺส โหติ ปุพฺพุฏฺฐายินี ปจฺฉานิปาตินี กิํการปฏิสฺสาวินี มนาปจารินี ปิยวาทินี. (1)}

\prose{เย เต ภตฺตุ ครุโน โหนฺติ “มาตา”ติ วา “ปิตา”ติ วา “สมณพฺราหฺมณา”ติ วา, เต สกฺกโรติ ครุํ กโรติ มาเนติ ปูเชติ, อพฺภาคเต จ อาสโนทเกน ปฏิปูเชติ. (2)}

\prose{เย เต ภตฺตุ อพฺภนฺตรา กมฺมนฺตา “อุณฺณา”ติ วา “กปฺปาสา”ติ วา, ตตฺถ ทกฺขา โหติ อนลสา, ตตฺรุปายาย วีมํสาย สมนฺนาคตา อลํ กาตุํ อลํ สํวิธาตุํ. (3)}

\prose{โย โส ภตฺตุ อพฺภนฺตโร อนฺโตชโน “ทาสา”ติ วา “เปสฺสา”ติ วา “กมฺมกรา”ติ วา, เตสํ กตญฺจ กตโต ชานาติ, อกตญฺจ อกตโต ชานาติ, คิลานกานญฺจ พลาพลํ ชานาติ, ขาทนียํ โภชนียญฺจสฺส ปจฺจํเสน สํวิภชติ. (4)}

\prose{\csromanfolio{96}ยํ ภตฺตา อาหรติ ธนํ วา ธญฺญํ วา รชตํ วา ชาตรูปํ วา, ตํ อารกฺเขน คุตฺติยา สมฺปาเทติ, ตตฺถ จ โหติ อธุตฺตี อเถนี อโสณฺฑี อวินาสิกา. (5)}

\prose{อุปาสิกา โข ปน โหติ พุทฺธํ สรณํ คตา, ธมฺมํ สรณํ คตา, สํฆํ สรณํ คตา. (6)}

\prose{สีลวตี โข ปน โหติ ปาณาติปาตา ปฏิวิรตา ฯเปฯ สุราเมรยมชฺช- ปมาทฏฺฐานา ปฏิวิรตา ฯเปฯ. (7)}

\prose{จาควตี โข ปน โหติ, วิคตมลมจฺเฉเรน เจตสา อคารํ อชฺฌาวสติ มุตฺตจาคา ปยตปาณินี โวสคฺครตา ยาจโยคา ทานสํวิภาครตา. (8)}

\prose{อิเมหิ โข นกุลมาเต อฏฺฐหิ ธมฺเมหิ สมนฺนาคโต มาตุคาโม กายสฺส เภทา ปรํ มรณา มนาปกายิกานํ เทวานํ สหพฺยตํ อุปปชฺชตีติ.}

\csromangathameasure{โย นํ ภรติ สพฺพทา, นิจฺจํ อาตาปิ อุสฺสุโก. \\ ตํ สพฺพกามทํ โปสํ, ภตฺตารํ นาติมญฺญติ. \\ น จาปิ โสตฺถิ ภตฺตารํ, อิสฺสาวาเทน โรสเย. \\ ภตฺตุ จ ครุโน สพฺเพ, ปฏิปูเชติ ปณฺฑิตา. \\ อุฏฺฐาหิกา อนลสา, สงฺคหิตปริชฺชนา. \\ ภตฺตุ มนาปํ จรติ, สมฺภตํ อนุรกฺขติ. \\ ยา เอวํ วตฺตติ นารี, ภตฺตุ ฉนฺทวสานุคา. \\ มนาปา นาม เต เทวา, ยตฺถ สา อุปปชฺชตีติ.อฏฺฐมํ.}
\csromangathagroup{\csromangathastanza{โย นํ ภรติ สพฺพทา, นิจฺจํ อาตาปิ อุสฺสุโก. \\ ตํ สพฺพกามทํ โปสํ, ภตฺตารํ นาติมญฺญติ.}\csromangathastanza{น จาปิ โสตฺถิ ภตฺตารํ, อิสฺสาวาเทน โรสเย. \\ ภตฺตุ จ ครุโน สพฺเพ, ปฏิปูเชติ ปณฺฑิตา.}\csromangathastanza{อุฏฺฐาหิกา อนลสา, สงฺคหิตปริชฺชนา. \\ ภตฺตุ มนาปํ จรติ, สมฺภตํ อนุรกฺขติ.}\csromangathastanza{ยา เอวํ วตฺตติ นารี, ภตฺตุ ฉนฺทวสานุคา. \\ มนาปา นาม เต\footnote{มนาปกายิกา (สี)} เทวา, ยตฺถ สา อุปปชฺชตีติ.\csromanspacer{}\csromanspacer{}อฏฺฐมํ.}}

\csromanheader{2}{9. ปฐม-อิธโลกิกสุตฺต}
\proseitem{49}{เอกํ สมยํ ภควา สาวตฺถิยํ วิหรติ ปุพฺพาราเม มิคารมาตุปาสาเท.\csromanspacer{} อถ โข วิสาขา มิคารมาตา เยน ภควา เตนุปสงฺกมิ ฯเปฯ เอกมนฺตํ นิสินฺนํ โข วิสาขํ มิคารมาตรํ ภควา เอตทโวจ\csromandash{}}

\chapterheadpageread{5. อุโปสถวคฺค}
\prose{\csromanfolio{97}จตูหิ โข วิสาเข ธมฺเมหิ สมนฺนาคโต มาตุคาโม อิธโลกวิชยาย ปฏิปนฺโน โหติ, อยํ’ส โลโก อารทฺโธ โหติ.\csromanspacer{} กตเมหิ จตูหิ? อิธ วิสาเข มาตุคาโม สุสํวิหิตกมฺมนฺโต โหติ, สงฺคหิตปริชโน, ภตฺตุ มนาปํ จรติ, สมฺภตํ อนุรกฺขติ.}

\prose{กถญฺจ วิสาเข มาตุคาโม สุสํวิหิตกมฺมนฺโต โหติ? อิธ วิสาเข มาตุคาโม, เย เต ภตฺตุ อพฺภนฺตรา กมฺมนฺตา “อุณฺณา”ติ วา “กปฺปาสา”ติ วา, ตตฺถ ทกฺขา โหติ อนลสา, ตตฺรุปายาย วีมํสาย สมนฺนาคตา อลํ กาตุํ อลํ สํวิธาตุํ.\csromanspacer{} เอวํ โข วิสาเข มาตุคาโม สุสํวิหิตกมฺมนฺโต โหติ.}

\prose{กถญฺจ วิสาเข มาตุคาโม สงฺคหิตปริชโน โหติ? อิธ วิสาเข มาตุคาโม โย โส ภตฺตุ อพฺภนฺตโร อนฺโตชโน “ทาสา”ติ วา “เปสฺสา”ติ วา “กมฺมกรา”ติ วา, เตสํ กตญฺจ กตโต ชานาติ, อกตญฺจ อกตโต ชานาติ, คิลานกานญฺจ พลาพลํ ชานาติ, ขาทนียโภชนียญฺจสฺส ปจฺจํเสน สํวิภชติ.\csromanspacer{} เอวํ โข วิสาเข มาตุคาโม สงฺคหิตปริชโน โหติ.}

\prose{กถญฺจ วิสาเข มาตุคาโม ภตฺตุ มนาปํ จรติ? อิธ วิสาเข มาตุคาโม, ยํ ภตฺตุ อมนาปสงฺขาตํ, ตํ ชีวิตเหตุปิ น อชฺฌาจรติ.\csromanspacer{} เอวํ โข วิสาเข มาตุคาโม ภตฺตุ มนาปํ จรติ.}

\prose{กถญฺจ วิสาเข มาตุคาโม สมฺภตํ อนุรกฺขติ? อิธ วิสาเข มาตุคาโม, ยํ ภตฺตา อาหรติ ธนํ วา ธญฺญํ วา รชตํ วา ชาตรูปํ วา, ตํ อารกฺเขน คุตฺติยา สมฺปาเทติ, ตตฺถ จ โหติ อธุตฺตี อเถนี อโสณฺฑี อวินาสิกา.\csromanspacer{} เอวํ โข วิสาเข มาตุคาโม สมฺภตํ อนุรกฺขติ.\csromanspacer{} อิเมหิ โข วิสาเข จตูหิ ธมฺเมหิ สมนฺนาคโต มาตุคาโม อิธโลกวิชยาย ปฏิปนฺโน โหติ, อยํ’ส โลโก อารทฺโธ โหติ.}

\prose{จตูหิ โข วิสาเข ธมฺเมหิ สมนฺนาคโต มาตุคาโม ปรโลกวิชยาย ปฏิปนฺโน โหติ, ปรโลโก อารทฺโธ โหติ.\csromanspacer{} กตเมหิ จตูหิ? อิธ วิสาเข มาตุคาโม สทฺธาสมฺปนฺโน โหติ, สีลสมฺปนฺโน โหติ, จาคสมฺปนฺโน โหติ, ปญฺญาสมฺปนฺโน โหติ.}

\prose{\csromanfolio{98}กถญฺจ วิสาเข มาตุคาโม สทฺธาสมฺปนฺโน โหติ? อิธ วิสาเข มาตุคาโม สทฺโธ โหติ, สทฺทหติ ตถาคตสฺส โพธิํ “อิติปิ โส ภควา อรหํ สมฺมาสมฺพุทฺโธ วิชฺชาจรณสมฺปนฺโน สุคโต โลกวิทู อนุตฺตโร ปุริสทมฺมสารถิ สตฺถา เทวมนุสฺสานํ พุทฺโธ ภควา”ติ.\csromanspacer{} เอวํ โข วิสาเข มาตุคาโม สทฺธาสมฺปนฺโน โหติ.}

\prose{กถญฺจ วิสาเข มาตุคาโม สีลสมฺปนฺโน โหติ? อิธ วิสาเข มาตุคาโม ปาณาติปาตา ปฏิวิรโต โหติ ฯเปฯ สุราเมรยมชฺชปมาทฏฺฐานา ปฏิวิรโต โหติ.\csromanspacer{} เอวํ โข วิสาเข มาตุคาโม สีลสมฺปนฺโน โหติ.}

\prose{กถญฺจ วิสาเข มาตุคาโม จาคสมฺปนฺโน โหติ? อิธ วิสาเข มาตุคาโม วิคตมลมจฺเฉเรน เจตสา อคารํ อชฺฌาวสติ มุตฺตจาคา ปยตปาณินี โวสคฺครตา ยาจโยคา ทานสํวิภาครตา.\csromanspacer{} เอวํ โข วิสาเข มาตุคาโม จาคสมฺปนฺโน โหติ.}

\prose{กถญฺจ วิสาเข มาตุคาโม ปญฺญาสมฺปนฺโน โหติ? อิธ วิสาเข มาตุคาโม ปญฺญวา โหติ ฯเปฯ.\csromanspacer{} เอวํ โข วิสาเข มาตุคาโม ปญฺญาสมฺปนฺโน โหติ.\csromanspacer{} อิเมหิ โข วิสาเข จตูหิ ธมฺเมหิ สมนฺนาคโต มาตุคาโม ปรโลกวิชยาย ปฏิปนฺโน โหติ, ปรโลโก อารทฺโธ โหตีติ.}

\csromangathameasure{สุสํวิหิตกมฺมนฺตา, สงฺคหิตปริชฺชนา. \\ ภตฺตุ มนาปํ จรติ, สมฺภตํ อนุรกฺขติ. \\ สทฺธา สีเลน สมฺปนฺนา, วทญฺญู วีตมจฺฉรา. \\ นิจฺจํ มคฺคํ วิโสเธติ, โสตฺถานํ สมฺปรายิกํ. \\ อิจฺเจเต อฏฺฐ ธมฺมา จ, ยสฺสา วิชฺชนฺติ นาริยา. \\ ตมฺปิ สีลวติํ อาหุ, ธมฺมฏฺฐํ สจฺจวาทินิํ. \\ โสฬสาการสมฺปนฺนา, อฏฺฐงฺคสุสมาคตา. \\ ตาทิสี สีลวตี อุปาสิกา, อุปปชฺชติ เทวโลกํ มนาปนฺติ. นวมํ.}
\csromangathagroup{\csromangathastanza{สุสํวิหิตกมฺมนฺตา, สงฺคหิตปริชฺชนา. \\ ภตฺตุ มนาปํ จรติ, สมฺภตํ อนุรกฺขติ.}\csromangathastanza{สทฺธา สีเลน สมฺปนฺนา, วทญฺญู วีตมจฺฉรา. \\ นิจฺจํ มคฺคํ วิโสเธติ, โสตฺถานํ สมฺปรายิกํ.}\csromangathastanza{อิจฺเจเต อฏฺฐ ธมฺมา จ, ยสฺสา วิชฺชนฺติ นาริยา. \\ ตมฺปิ สีลวติํ อาหุ, ธมฺมฏฺฐํ สจฺจวาทินิํ.}\csromangathastanza{โสฬสาการสมฺปนฺนา, อฏฺฐงฺคสุสมาคตา. \\ ตาทิสี สีลวตี อุปาสิกา, อุปปชฺชติ เทวโลกํ มนาปนฺติ.\csromanspacer{} นวมํ.}}

\chapterheadpageread{5. อุโปสถวคฺค \textbf{10. ทุติย-อิธโลกิกสุตฺต}}
\proseitem{50}{\csromanfolio{99}จตูหิ ภิกฺขเว ธมฺเมหิ สมนฺนาคโต มาตุคาโม อิธโลกวิชยาย ปฏิปนฺโน โหติ, อยํ’ส โลโก อารทฺโธ โหติ.\csromanspacer{} กตเมหิ จตูหิ? อิธ ภิกฺขเว มาตุคาโม สุสํวิหิตกมฺมนฺโต โหติ, สงฺคหิตปริชโน, ภตฺตุ มนาปํ จรติ, สมฺภตํ อนุรกฺขติ.}

\prose{กถญฺจ ภิกฺขเว มาตุคาโม สุสํวิหิตกมฺมนฺโต โหติ? อิธ ภิกฺขเว มาตุคาโม, เย เต ภตฺตุ อพฺภนฺตรา กมฺมนฺตา ฯเปฯ.\csromanspacer{} เอวํ โข ภิกฺขเว มาตุคาโม สุสํวิหิตกมฺมนฺโต โหติ.}

\prose{กถญฺจ ภิกฺขเว มาตุคาโม สงฺคหิตปริชโน โหติ? อิธ ภิกฺขเว มาตุคาโม, โย โส ภตฺตุ อพฺภนฺตโร อนฺโตชโน ฯเปฯ.\csromanspacer{} เอวํ โข ภิกฺขเว มาตุคาโม สงฺคหิตปริชโน โหติ.}

\prose{กถญฺจ ภิกฺขเว มาตุคาโม ภตฺตุ มนาปํ จรติ? อิธ ภิกฺขเว มาตุคาโม, ยํ ภตฺตุ อมนาปสงฺขาตํ, ตํ ชีวิตเหตุปิ น อชฺฌาจรติ.\csromanspacer{} เอวํ โข ภิกฺขเว มาตุคาโม ภตฺตุ มนาปํ จรติ.}

\prose{กถญฺจ ภิกฺขเว มาตุคาโม สมฺภตํ อนุรกฺขติ, อิธ ภิกฺขเว มาตุคาโม, ยํ ภตฺตา อาหรติ ฯเปฯ.\csromanspacer{} เอวํ โข ภิกฺขเว มาตุคาโม สมฺภตํ อนุรกฺขติ.\csromanspacer{} อิเมหิ โข ภิกฺขเว จตูหิ ธมฺเมหิ สมนฺนาคโต มาตุคาโม อิธโลกวิชยาย ปฏิปนฺโน โหติ, อยํ’ส โลโก อารทฺโธ โหติ.}

\prose{จตูหิ ภิกฺขเว ธมฺเมหิ สมนฺนาคโต มาตุคาโม ปรโลกวิชยาย ปฏิปนฺโน โหติ, ปรโลโก อารทฺโธ โหติ.\csromanspacer{} กตเมหิ จตูหิ, อิธ ภิกฺขเว มาตุคาโม สทฺธาสมฺปนฺโน โหติ, สีลสมฺปนฺโน โหติ, จาคสมฺปนฺโน โหติ, ปญฺญาสมฺปนฺโน โหติ.}

\prose{กถญฺจ ภิกฺขเว มาตุคาโม สทฺธาสมฺปนฺโน โหติ, อิธ ภิกฺขเว มาตุคาโม สทฺโธ โหติ ฯเปฯ.\csromanspacer{} เอวํ โข ภิกฺขเว มาตุคาโม สทฺธาสมฺปนฺโน โหติ.}

\prose{\csromanfolio{100}กถญฺจ ภิกฺขเว มาตุคาโม สีลสมฺปนฺโน โหติ, อิธ ภิกฺขเว มาตุคาโม ปาณาติปาตา ปฏิวิรโต โหติ ฯเปฯ สุราเมรยมชฺชปมาทฏฺฐานา ปฏิวิรโต โหติ.\csromanspacer{} เอวํ โข ภิกฺขเว มาตุคาโม สีลสมฺปนฺโน โหติ.}

\prose{กถญฺจ ภิกฺขเว มาตุคาโม จาคสมฺปนฺโน โหติ, อิธ ภิกฺขเว มาตุคาโม วิคตมลมจฺเฉเรน เจตสา อคารํ อชฺฌาวสติ ฯเปฯ.\csromanspacer{} เอวํ โข ภิกฺขเว มาตุคาโม จาคสมฺปนฺโน โหติ.}

\prose{กถญฺจ ภิกฺขเว มาตุคาโม ปญฺญาสมฺปนฺโน โหติ, อิธ ภิกฺขเว มาตุคาโม ปญฺญวา โหติ ฯเปฯ.\csromanspacer{} เอวํ โข ภิกฺขเว มาตุคาโม ปญฺญาสมฺปนฺโน โหติ.\csromanspacer{} อิเมหิ โข ภิกฺขเว จตูหิ ธมฺเมหิ สมนฺนาคโต มาตุคาโม ปรโลกวิชยาย ปฏิปนฺโน โหติ, ปรโลโก อารทฺโธ โหตีติ.}

\csromangathameasure{สุสํวิหิตกมฺมนฺตา, สงฺคหิตปริชฺชนา. \\ ภตฺตุ มนาปํ จรติ, สมฺภตํ อนุรกฺขติ. \\ สทฺธา สีเลน สมฺปนฺนา, วทญฺญู วีตมจฺฉรา. \\ นิจฺจํ มคฺคํ วิโสเธติ, โสตฺถานํ สมฺปรายิกํ. \\ อิจฺเจเต อฏฺฐ ธมฺมา จ, ยสฺสา วิชฺชนฺติ นาริยา. \\ ตมฺปิ สีลวติํ อาหุ, ธมฺมฏฺฐํ สจฺจวาทินิํ. \\ โสฬสาการสมฺปนฺนา, อฏฺฐงฺคสุสมาคตา. \\ ตาทิสี สีลวตี อุปาสิกา, อุปปชฺชติ เทวโลกํ มนาปนฺติ. ทสมํ. อุโปสถวคฺโค ปญฺจโม.}
\csromangathagroup{\csromangathastanza{สุสํวิหิตกมฺมนฺตา, สงฺคหิตปริชฺชนา. \\ ภตฺตุ มนาปํ จรติ, สมฺภตํ อนุรกฺขติ.}\csromangathastanza{สทฺธา สีเลน สมฺปนฺนา, วทญฺญู วีตมจฺฉรา. \\ นิจฺจํ มคฺคํ วิโสเธติ, โสตฺถานํ สมฺปรายิกํ.}\csromangathastanza{อิจฺเจเต อฏฺฐ ธมฺมา จ, ยสฺสา วิชฺชนฺติ นาริยา. \\ ตมฺปิ สีลวติํ อาหุ, ธมฺมฏฺฐํ สจฺจวาทินิํ.}\csromangathastanza{โสฬสาการสมฺปนฺนา, อฏฺฐงฺคสุสมาคตา. \\ ตาทิสี สีลวตี อุปาสิกา, อุปปชฺชติ เทวโลกํ มนาปนฺติ.\csromanspacer{} ทสมํ.\csromanspacer{} อุโปสถวคฺโค ปญฺจโม.}}

\titlehead{\textbf{ตสฺสุทฺทานํ}}
\csromangathameasure{สํขิตฺเต วิตฺถเต วิสาเข, วาเสฏฺโฐ โพชฺฌาย ปญฺจมํ. \\ อนุรุทฺธํ ปุน วิสาเข, นกุลา อิธโลกิกา เทฺวติ. \\ \textbf{ปฐโม ปณฺณาสโก สมตฺโต.}}
\csromangathagroup{\csromangathastanza{สํขิตฺเต วิตฺถเต วิสาเข, วาเสฏฺโฐ โพชฺฌาย ปญฺจมํ. \\ อนุรุทฺธํ ปุน วิสาเข, นกุลา อิธโลกิกา เทฺวติ. \\ \textbf{ปฐโม ปณฺณาสโก สมตฺโต.}}}

\chapterheadpageread{2. ทุติยปณฺณาสก}
\chapterhead{\textbf{(6) 1. โคตมีวคฺค}}
\csromanheader{2}{1. โคตมีสุตฺต}
\proseitem{51}{\csromanfolio{101}เอกํ สมยํ ภควา สกฺเกสุ วิหรติ กปิลวตฺถุสฺมิํ นิโคฺรธาราเม.\csromanspacer{} อถ โข มหาปชาปตี\footnote{มหาปชาปติ (สฺยา) วิ 4. 442 ปิฏฺเฐปิ.} โคตมี เยน ภควา เตนุปสงฺกมิ, อุปสงฺกมิตฺวา ภควนฺตํ อภิวาเทตฺวา เอกมนฺตํ อฏฺฐาสิ, เอกมนฺตํ ฐิตา โข มหาปชาปตี โคตมี ภควนฺตํ เอตทโวจ “สาธุ ภนฺเต ลเภยฺย มาตุคาโม ตถาคตปฺปเวทิเต ธมฺมวินเย อคารสฺมา อนคาริยํ ปพฺพชฺชนฺ”ติ.\csromanspacer{} อลํ โคตมิ, มา เต รุจฺจิ มาตุคามสฺส ตถาคตปฺปเวทิเต ธมฺมวินเย อคารสฺมา อนคาริยํ ปพฺพชฺชาติ.}

\prose{ทุติยมฺปิ โข มหาปชาปตี โคตมี ภควนฺตํ เอตทโวจ “สาธุ ภนฺเต ลเภยฺย มาตุคาโม ตถาคตปฺปเวทิเต ธมฺมวินเย อคารสฺมา อนคาริยํ ปพฺพชฺชนฺ”ติ.\csromanspacer{} อลํ โคตมิ, มา เต รุจฺจิ มาตุคามสฺส ตถาคตปฺปเวทิเต ธมฺมวินเย อคารสฺมา อนคาริยํ ปพฺพชฺชาติ.\csromanspacer{} ตติยมฺปิ โข มหาปชาปตี โคตมี ภควนฺตํ เอตทโวจ “สาธุ ภนฺเต ลเภยฺย มาตุคาโม ตถาคตปฺปเวทิเต ธมฺมวินเย อคารสฺมา อนคาริยํ ปพฺพชฺชนฺ”ติ.\csromanspacer{} อลํ โคตมิ, มา เต รุจฺจิ มาตุคามสฺส ตถาคตปฺปเวทิเต ธมฺมวินเย อคารสฺมา อนคาริยํ ปพฺพชฺชาติ.}

\prose{อถ โข มหาปชาปตี โคตมี “น ภควา อนุชานาติ มาตุคามสฺส ตถาคตปฺปเวทิเต ธมฺมวินเย อคารสฺมา อนคาริยํ ปพฺพชฺชนฺ”ติ ทุกฺขี ทุมฺมนา อสฺสุมุขี รุทมานา ภควนฺตํ อภิวาเทตฺวา ปทกฺขิณํ กตฺวา ปกฺกามิ.}

\prose{อถ โข ภควา กปิลวตฺถุสฺมิํ ยถาภิรนฺตํ วิหริตฺวา เยน เวสาลี เตน จาริกํ ปกฺกามิ, อนุปุพฺเพน จาริกํ จรมาโน เยน เวสาลี ตทวสริ.\csromanspacer{} ตตฺร สุทํ ภควา เวสาลิยํ วิหรติ มหาวเน \csromanfolio{102}กูฏาคารสาลายํ.\csromanspacer{} อถ โข มหาปชาปตี โคตมี เกเส เฉทาเปตฺวา กาสายานิ วตฺถานิ อจฺฉาเทตฺวา สมฺพหุลาหิ สากิยานีหิ สทฺธิํ เยน เวสาลี เตน ปกฺกามิ, อนุปุพฺเพน เยน เวสาลี มหาวนํ กูฏาคารสาลา เตนุปสงฺกมิ.\csromanspacer{} อถ โข มหาปชาปตี โคตมี สูเนหิ ปาเทหิ รโชกิณฺเณน คตฺเตน ทุกฺขี ทุมฺมนา อสฺสุมุขี รุทมานา พหิทฺวารโกฏฺฐเก อฏฺฐาสิ.}

\prose{อทฺทสา โข อายสฺมา อานนฺโท มหาปชาปติํ โคตมิํ สูเนหิ ปาเทหิ รโชกิณฺเณน คตฺเตน ทุกฺขิํ ทุมฺมนํ อสฺสุมุขิํ รุทมานํ พหิทฺวารโกฏฺฐเก ฐิตํ, ทิสฺวาน มหาปชาปติํ โคตมิํ เอตทโวจ “กิํ นุ ตฺวํ โคตมิ สูเนหิ ปาเทหิ รโชกิณฺเณน คตฺเตน ทุกฺขี ทุมฺมนา อสฺสุมุขี รุทมานา พหิทฺวารโกฏฺฐเก ฐิตา”ติ.\csromanspacer{} ตถา หิ ปน ภนฺเต อานนฺท น ภควา อนุชานาติ มาตุคามสฺส ตถาคตปฺปเวทิเต ธมฺมวินเย อคารสฺมา อนคาริยํ ปพฺพชฺชนฺติ.\csromanspacer{} เตน หิ ตฺวํ โคตมิ มุหุตฺตํ อิเธว ตาว โหหิ, ยาวาหํ ภควนฺตํ ยาจามิ มาตุคามสฺส ตถาคตปฺปเวทิเต ธมฺมวินเย อคารสฺมา อนคาริยํ ปพฺพชฺชนฺติ.}

\csromanmatikaanchor{mtk.2}
\csromanmatikaanchor{mtk.82}
\csromantocmarkref{subsection}{1. สิกฺขาทุพฺพลฺยสุตฺต}{mtk.2}
\bookmark[dest={mtk.2},level=2]{1. สิกฺขาทุพฺพลฺยสุตฺต}
\prose{อถ โข อายสฺมา อานนฺโท เยน ภควา เตนุปสงฺกมิ, อุปสงฺกมิตฺวา ภควนฺตํ อภิวาเทตฺวา เอกมนฺตํ นิสีทิ, เอกมนฺตํ นิสินฺโน โข อายสฺมา อานนฺโท ภควนฺตํ เอตทโวจ “เอสา ภนฺเต มหาปชาปตี โคตมี สูเนหิ ปาเทหิ รโชกิณฺเณน คตฺเตน ทุกฺขี ทุมฺมนา อสฺสุมุขี รุทมานา พหิทฺวารโกฏฺฐเก ฐิตา ‘น ภควา อนุชานาติ มฺตุคามสฺส ตถาคตปฺปเวทิเต ธมฺมวินเย อคารสฺมา อนคาริยํ ปพฺพชฺชนฺ’ติ, สาธุ ภนฺเต ลเภยฺย มาตุคาโม ตถาคตปฺปเวทิเต ธมฺมวินเย อคารสฺมา อนคาริยํ ปพฺพชฺชนฺ”ติ.\csromanspacer{} อลํ อานนฺท, มา เต รุจฺจิ มฺตุคามสฺส ตถาคตปฺปเวทิเต ธมฺมวินเย อคารสฺมา อนคาริยํ ปพฺพชฺชาติ.}

\prose{ทุติยมฺปิ โข ฯเปฯ.\csromanspacer{} ตติยมฺปิ โข อายสฺมา อานนฺโท ภควนฺตํ เอตทโวจ “สาธุ ภนฺเต ลเภยฺย มฺตุคาโม ตถาคตปฺปเวทิเต ธมฺมวินเย อคารสฺมา อนคาริยํ ปพฺพชฺชนฺ”ติ.\csromanspacer{} อลํ อานนฺท, มา เต รุจฺจิ มฺตุคามสฺส ตถาคตปฺปเวทิเต ธมฺมวินเย อคารสฺมา อนคาริยํ ปพฺพชฺชาติ.}

\chapterheadpageread{(6) 1. โคตมีวคฺค}
\prose{\csromanfolio{103}อถ โข อายสฺมโต อานนฺทสฺส เอตทโหสิ “น ภควา อนุชานาติ มาตุคามสฺส ตถาคตปฺปเวทิเต ธมฺมวินเย อคารสฺมา อนคาริยํ ปพฺพชฺชํ, ยํนูนาหํ อญฺเญนปิ ปริยาเยน ภควนฺตํ ยาเจยฺยํ มาตุคามสฺส ตถาคตปฺปเวทิเต ธมฺมวินเย อคารสฺมา อนคาริยํ ปพฺพชฺชนฺ”ติ.\csromanspacer{} อถ โข อายสฺมา อานนฺโท ภควนฺตํ เอตทโวจ “ภพฺโพ นุ โข ภนฺเต มาตุคาโม ตถาคตปฺปเวทิเต ธมฺมวินเย อคารสฺมา อนคาริยํ ปพฺพชิตฺวา โสตาปตฺติผลํ วา สกทาคามิผลํ วา อนาคามิ- ผลํ วา อรหตฺตผลํ วา สจฺฉิกาตุนฺ”ติ.\csromanspacer{} ภพฺโพ อานนฺท มาตุคาโม ตถาคตปฺปเวทิเต ธมฺมวินเย อคารสฺมา อนคาริยํ ปพฺพชิตฺวา โสตาปตฺติผลมฺปิ สกทาคามิผลมฺปิ อนาคามิผลมฺปิ อรหตฺตผลมฺปิ สจฺฉิกาตุนฺติ.\csromanspacer{} สเจ ภนฺเต ภพฺโพ มาตุคาโม ตถาคตปฺปเวทิเต ธมฺมวินเย อคารสฺมา อนคาริยํ ปพฺพชิตฺวา โสตาปตฺติผลมฺปิ ฯเปฯ อรหตฺตผลมฺปิ สจฺฉิกาตุํ.\csromanspacer{} พหุการา ภนฺเต มหาปชาปตี โคตมี ภควโต มาตุจฺฉา อาปาทิกา โปสิกา ขีรสฺส ทายิกา ภควนฺตํ ชเนตฺติยา กาลงฺกตาย ถญฺญํ ปาเยสิ, สาธุ ภนฺเต ลเภยฺย มาตุคาโม ตถาคตปฺปเวทิเต ธมฺมวินเย อคารสฺมา อนคาริยํ ปพฺพชฺชนฺติ.}

\prose{สเจ อานนฺท มหาปชาปตี โคตมี อฏฺฐ ครุธมฺเม ปฏิคฺคณฺหาติ, สาวสฺสา โหตุ อุปสมฺปทา.}

\prose{\csromansymbolfootnote{*}{วิ 2. 74; วิ 4. 444 ปิฏฺเฐสุปิ.}วสฺสสตูปสมฺปนฺนาย ภิกฺขุนิยา ตทหูปสมฺปนฺนสฺส ภิกฺขุโน อภิวาทนํ ปจฺจุฏฺฐานํ อญฺชลิกมฺมํ สามีจิกมฺมํ กตฺตพฺพํ, อยมฺปิ ธมฺโม สกฺกตฺวา ครุํ กตฺวา\footnote{ครุกตฺวา (สี, สฺยา, อิ)} มาเนตฺวา ปูเชตฺวา ยาวชีวํ อนติกฺกมนีโย. (1)}

\prose{น ภิกฺขุนิยา อภิกฺขุเก อาวาเส วสฺสํ อุปคนฺตพฺพํ, อยมฺปิ ธมฺโม สกฺกตฺวา ครุํ กตฺวา มาเนตฺวา ปูเชตฺวา ยาวชีวํ อนติกฺกมนีโย. (2)}

\prose{อนฺวฑฺฒมาสํ ภิกฺขุนิยา ภิกฺขุสํฆโต เทฺว ธมฺมา ปจฺจาสีสิตพฺพา\footnote{ปจฺจาสิํสิตพฺพา (สี, สฺยา, อิ)} อุโปสถปุจฺฉกญฺจ โอวาทูปสงฺกมนญฺจ, อยมฺปิ ธมฺโม สกฺกตฺวา ครุํ กตฺวา มาเนตฺวา ปูเชตฺวา ยาวชีวํ อนติกฺกมนีโย. (3)}

\prose{วสฺสํวุฏฺฐาย ภิกฺขุนิยา อุภโตสํเฆ ตีหิ ฐาเนหิ ปวาเรตพฺพํ ทิฏฺเฐน วา สุเตน วา ปริสงฺกาย วา, อยมฺปิ ธมฺโม สกฺกตฺวา ครุํ กตฺวา มาเนตฺวา ปูเชตฺวา ยาวชีวํ อนติกฺกมนีโย. (4)}

\prose{\csromanfolio{104}ครุธมฺมํ อชฺฌาปนฺนาย ภิกฺขุนิยา อุภโตสํเฆ ปกฺขมานตฺตํ จริตพฺพํ, อยมฺปิ ธมฺโม สกฺกตฺวา ครุํ กตฺวา มาเนตฺวา ปูเชตฺวา ยาวชีวํ อนติกฺกมนีโย. (5)}

\prose{เทฺว วสฺสานิ ฉสุ ธมฺเมสุ สิกฺขิตสิกฺขาย สิกฺขมานาย อุภโตสํเฆ อุปสมฺปทา ปริเยสิตพฺพา, อยมฺปิ ธมฺโม สกฺกตฺวา ครุํ กตฺวา มาเนตฺวา ปูเชตฺวา ยาวชีวํ อนติกฺกมนีโย. (6)}

\prose{น เกนจิ ปริยาเยน ภิกฺขุนิยา ภิกฺขุ อกฺโกสิตพฺโพ ปริภาสิตพฺโพ, อยมฺปิ ธมฺโม สกฺกตฺวา ครุํ กตฺวา มาเนตฺวา ปูเชตฺวา ยาวชีวํ อนติกฺกมนีโย. (7)}

\prose{อชฺชตคฺเค โอวโฏ ภิกฺขุนีนํ ภิกฺขูสุ วจนปโถ, อโนวโฏ ภิกฺขูนํ ภิกฺขุนีสุ วจนปโถ, อยมฺปิ ธมฺโม สกฺกตฺวา ครุํ กตฺวา มาเนตฺวา ปูเชตฺวา ยาวชีวํ อนติกฺกมนีโย. (8)}

\prose{สเจ อานนฺท มหาปชาปตี โคตมี อิเม อฏฺฐ ครุธมฺเม ปฏิคฺคณฺหาติ, สาว’สฺสา โหตุ อุปสมฺปทาติ.}

\prose{อถ โข อายสฺมา อานนฺโท ภควโต สนฺติเก อิเม อฏฺฐ ครุธมฺเม อุคฺคเหตฺวา เยน มหาปชาปตี โคตมี เตนุปสงฺกมิ, อุปสงฺกมิตฺวา มหาปชาปติํ โคตมิํ เอตทโวจ\csromandash{}}

\prose{สเจ โข ตฺวํ โคตมิ อฏฺฐ ครุธมฺเม ปฏิคฺคเณฺหยฺยาสิ, สาว เต ภวิสฺสติ อุปสมฺปทา.}

\prose{วสฺสสตูปสมฺปนฺนาย ภิกฺขุนิยา ตทหูปสมฺปนฺนสฺส ภิกฺขุโน อภิวาทนํ ปจฺจุฏฺฐานํ อญฺชลิกมฺมํ สามีจิกมฺมํ กตฺตพฺพํ, อยมฺปิ ธมฺโม สกฺกตฺวา ครุํ กตฺวา มาเนตฺวา ปูเชตฺวา ยาวชีวํ อนติกฺกมนีโย ฯเปฯ.}

\prose{อชฺชตคฺเค โอวโฏ ภิกฺขุนีนํ ภิกฺขูสุ วจนปโถ, อโนวโฏ ภิกฺขูนํ ภิกฺขุนีสุ วจนปโถ, อยมฺปิ ธมฺโม สกฺกตฺวา ครุํ กตฺวา มาเนตฺวา ปูเชตฺวา ยาวชีวํ อนติกฺกมนีโย.\csromanspacer{} สเจ โข ตฺวํ โคตมิ อิเม อฏฺฐ ครุธมฺเม ปฏิคฺคเณฺหยฺยาสิ, สาว เต ภวิสฺสติ อุปสมฺปทาติ.}

\chapterheadpageread{(6) 1. โคตมีวคฺค}
\prose{\csromanfolio{105}เสยฺยถาปิ ภนฺเต อานนฺท อิตฺถี วา ปุริโส วา ทหโร ยุวา มณฺฑนกชาติโก\footnote{มณฺฑนกชาติโย (สี, อิ)} สีสํนฺหาโต\footnote{สีสํนหาโต (สี, อิ), สีสนหาโต (สฺยา)} อุปฺปลมาลํ วา วสฺสิกมาลํ วา อธิมุตฺตกมาลํ\footnote{อติมุตฺตกมาลํ (สี)} วา ลภิตฺวา อุโภหิ หตฺเถหิ ปฏิคฺคเหตฺวา อุตฺตมงฺเค สิรสฺมิํ ปติฏฺฐาเปยฺย.\csromanspacer{} เอวเมวํ โข อหํ ภนฺเต อานนฺท อิเม อฏฺฐ ครุธมฺเม ปฏิคฺคณฺหามิ ยาวชีวํ อนติกฺกมนีเยติ.}

\prose{อถ โข อายสฺมา อานนฺโท เยน ภควา เตนุปสงฺกมิ, อุปสงฺกมิตฺวา ภควนฺตํ อภิวาเทตฺวา เอกมนฺตํ นิสีทิ, เอกมนฺตํ นิสินฺโน โข อายสฺมา อานนฺโท ภควนฺตํ เอตทโวจ “ปฏิคฺคหิตา ภนฺเต มหาปชาปติยา โคตมิยา อฏฺฐ ครุธมฺมา ยาวชีวํ อนติกฺกมนียา”ติ.}

\prose{สเจ อานนฺท นาลภิสฺส มาตุคาโม ตถาคตปฺปเวทิเต ธมฺมวินเย อคารสฺมา อนคาริยํ ปพฺพชฺชํ, จิรฏฺฐิติกํ อานนฺท พฺรหฺมจริยํ อภวิสฺส, วสฺสสหสฺสเมว สทฺธมฺโม ติฏฺเฐยฺย.\csromanspacer{} ยโต จ โข อานนฺท มาตุคาโม ตถาคตปฺปเวทิเต ธมฺมวินเย อคารสฺมา อนคาริยํ ปพฺพชิโต, น ทานิ อานนฺท พฺรหฺมจริยํ จิรฏฺฐิติกํ ภวิสฺสติ, ปญฺเจว ทานิ อานนฺท วสฺสสตานิ สทฺธมฺโม ฐสฺสติ.}

\prose{เสยฺยถาปิ อานนฺท ยานิ กานิจิ กุลานิ พหุตฺถิกานิ\footnote{พหุกิตฺถิกานิ (สี, อิ), พหุ-อิตฺถิกานิ (สฺยา)} อปฺปปุริสกานิ, ตานิ สุปฺปธํสิยานิ โหนฺติ โจเรหิ กุมฺภตฺเถนเกหิ.\csromanspacer{} เอวเมวํ โข อานนฺท ยสฺมิํ ธมฺมวินเย ลภติ มาตุคาโม อคารสฺมา อนคาริยํ ปพฺพชฺชํ, น ตํ พฺรหฺมจริยํ จิรฏฺฐิติกํ โหติ.}

\prose{เสยฺยถาปิ อานนฺท สมฺปนฺเน สาลิกฺเขตฺเต เสตฏฺฐิกา นาม โรคชาติ นิปตติ, เอวํ ตํ สาลิกฺเขตฺตํ น จิรฏฺฐิติกํ โหติ.\csromanspacer{} เอวเมวํ โข อานนฺท ยสฺมิํ ธมฺมวินเย ลภติ นาตุคาโม อคารสฺมา อนคาริยํ ปพฺพชฺชํ, น ตํ พฺรหฺมจริยํ จิรฏฺฐิติกํ โหติ.}

\prose{เสยฺยถาปิ อานนฺท สมฺปนฺเน อุจฺฉุกฺเขตฺเต มญฺชิฏฺฐิกา\footnote{มญฺเชฏฺฐิกา (สี, สฺยา)} นาม โรคชาติ นิปตติ, เอวํ ตํ อุจฺฉุกฺเขตฺตํ น จิรฏฺฐิติกํ โหติ.\csromanspacer{} เอวเมวํ โข อานนฺท ยสฺมิํ ธมฺมวินเย ลภติ มาตุคาโม อคารสฺมา อนคาริยํ ปพฺพชฺชํ, น ตํ พฺรหฺมจริยํ จิรฏฺฐิติกํ โหติ.}

\prose{\csromanfolio{106}เสยฺยถาปิ อานนฺท ปุริโส มหโต ตฬากสฺส ปฏิกจฺเจว\footnote{ปฏิคจฺเจว (สี, อิ)} อาฬิํ พนฺเธยฺย ยาวเทว อุทกสฺส อนติกฺกมนาย.\csromanspacer{} เอวเมวํ โข อานนฺท มยา ปฏิ กจฺเจว ภิกฺขุนีนํ อฏฺฐ ครุธมฺมา ปญฺญตฺตา ยาวชีวํ อนติกฺกมนียาติ.\csromanspacer{}\csromanspacer{}ปฐมํ.}

\csromanheader{2}{2. โอวาทสุตฺต}
\proseitem{52}{เอกํ สมยํ ภควา เวสาลิยํ วิหรติ มหาวเน กูฏาคารสาลายํ.\csromanspacer{} อถ โข อายสฺมา อานนฺโท เยน ภควา เตนุปสงฺกมิ, อุปสงฺกมิตฺวา ภควนฺตํ อภิวาเทตฺวา เอกมนฺตํ นิสีทิ, เอกมนฺตํ นิสินฺโน โข อายสฺมา อานนฺโท ภควนฺตํ เอตทโวจ “กติหิ นุ โข ภนฺเต ธมฺเมหิ สมนฺนาคโต ภิกฺขุ ภิกฺขุโนวาทโก สมฺมนฺนิตพฺโพ”ติ.}

\prose{\csromansymbolfootnote{+}{วิ 2. 73 ปิฏฺเฐปิ.}อฏฺฐหิ โข อานนฺท ธมฺเมหิ สมนฺนาคโต ภิกฺขุ ภิกฺขุโนวาทโก สมฺมนฺนิตพฺโพ.\csromanspacer{} กตเมหิ อฏฺฐหิ? อิธานนฺท ภิกฺขุ สีลวา โหติ ฯเปฯ สมาทาย สิกฺขติ สิกฺขาปเทสุ.\csromanspacer{} พหุสฺสุโต โหติ ฯเปฯ ทิฏฺฐิยา สุปฺปฏิวิทฺธา.\csromanspacer{} อุภยานิ โข ปนสฺส ปาติโมกฺขานิ วิตฺถาเรน สฺวาคตานิ โหนฺติ สุวิภตฺตานิ สุปฺปวตฺตีนิ สุวินิจฺฉิตานิ สุตฺตโส อนุพฺยญฺชนโส.\csromanspacer{} กลฺยาณวาโจ โหติ กลฺยาณวากฺกรโณ โปริยา วาจาย สมนฺนาคโต วิสฺสฏฺฐาย\footnote{วิสฏฺฐาย (ก)} อเนลคฬาย\footnote{อเนฬคฬาย (สี, ก)} อตฺถสฺส วิญฺญาปนิยา.\csromanspacer{} ปฏิพโล โหติ ภิกฺขุนิสํฆสฺส ธมฺมิยา กถาย สนฺทสฺเสตุํ สมาทเปตุํ สมุตฺเตเชตุํ สมฺปหํเสตุํ.\csromanspacer{} เยภุยฺเยน ภิกฺขุนีนํ ปิโย โหติ มนาโป.\csromanspacer{} น โข ปเนตํ ภควนฺตํ อุทฺทิสฺส ปพฺพชิตาย กาสายวตฺถวสนาย ครุธมฺมํ อชฺฌาปนฺนปุพฺโพ โหติ.\csromanspacer{} วีสติวสฺโส วา โหติ อติเรกวีสติวสฺโส วา.\csromanspacer{} อิเมหิ โข อานนฺท อฏฺฐหิ ธมฺเมหิ สมนฺนาคโต ภิกฺขุ ภิกฺขุโนวาทโก สมฺมนฺนิตพฺโพติ.\csromanspacer{}\csromanspacer{}ทุติยํ.}

\csromanheader{2}{3. สํขิตฺตสุตฺต}
\proseitem{53}{\csromansymbolfootnote{*}{วิ 4. 449 ปิฏฺเฐปิ.}เอกํ สมยํ ภควา เวสาลิยํ วิหรติ มหาวเน กูฏาคารสาลายํ.\csromanspacer{} อถ โข มหาปชาปตี โคตมี เยน ภควา เตนุปสงฺกมิ, อุปสงฺกมิตฺวา ภควนฺตํ อภิวาเทตฺวา เอกมนฺตํ อฏฺฐาสิ, เอกมนฺตํ ฐิตา โข สา มหาปชาปตี โคตมี ภควนฺตํ เอตทโวจ\csromandash{}}

\chapterheadpageread{(6) 1. โคตมีวคฺค}
\prose{\csromanfolio{107}สาธุ เม ภนฺเต ภควา สํขิตฺเตน ธมฺมํ เทเสตุ, ยมหํ ภควโต ธมฺมํ สุตฺวา เอกา วูปกฏฺฐา อปฺปมตฺตา อาตาปินี ปหิตตฺตา วิหเรยฺยนฺติ.\csromanspacer{} เย โข ตฺวํ โคตมิ ธมฺเม ชาเนยฺยาสิ “อิเม ธมฺมา สราคาย สํวตฺตนฺติ โน วิราคาย, สํโยคาย สํวตฺตนฺติ โน วิสํโยคาย, อาจยาย สํวตฺตนฺติ โน อปจยาย, มหิจฺฉตาย สํวตฺตนฺติ โน อปฺปิจฺฉตาย, อสนฺตุฏฺฐิยา สํวตฺตนฺติ โน สนฺตุฏฺฐิยา, สงฺคณิกาย สํวตฺตนฺติ โน ปวิเวกาย, โกสชฺชาย สํวตฺตนฺติ โน วีริยารมฺภาย, ทุพฺภรตาย สํวตฺตนฺติ โน สุภรตายา”ติ.\csromanspacer{} เอกํเสน โคตมิ ธาเรยฺยาสิ “เนโส ธมฺโม เนโส วินโย เนตํ สตฺถุสาสนนฺ”ติ.}

\prose{เย จ โข ตฺวํ โคตมิ ธมฺเม ชาเนยฺยาสิ “อิเม ธมฺมา วิราคาย สํวตฺตนฺติ โน สราคาย, วิสํโยคาย สํวตฺตนฺติ โน สํโยคาย, อปจยาย สํวตฺตนฺติ โน อาจยาย, อปฺปิจฺฉตาย สํวตฺตนฺติ โน มหิจฺฉตาย, สนฺตุฏฺฐิยา สํวตฺตนฺติ โน อสนฺตุฏฺฐิยา, ปวิเวกาย สํวตฺตนฺติ โน สงฺคณิกาย, วีริยารมฺภาย สํวตฺตนฺติ โน โกสชฺชาย, สุภรตาย สํวตฺตนฺติ โน ทุพฺภรตายา”ติ.\csromanspacer{} เอกํเสน โคตมิ ธาเรยฺยาสิ “เอโส ธมฺโม เอโส วินโย เอตํ สตฺถุสาสนนฺ”ติ.\csromanspacer{}\csromanspacer{}ตติยํ.}

\csromanheader{2}{4. ทีฆชาณุสุตฺต}
\proseitem{54}{เอกํ สมยํ ภควา โกลิเยสุ วิหรติ กกฺกรปตฺตํ นาม โกลิยานํ นิคโม.\csromanspacer{} อถ โข ทีฆชาณุ โกลิยปุตฺโต เยน ภควา เตนุปสงฺกมิ, อุปสงฺกมิตฺวา ภควนฺตํ อภิวาเทตฺวา เอกมนฺตํ นิสีทิ, เอกมนฺตํ นิสินฺโน โข ทีฆชาณุ โกลิยปุตฺโต ภควนฺตํ เอตทโวจ “มยํ ภนฺเต คิหี กามโภคิโน\footnote{กามโภคี (สี, สฺยา, อิ)} ปุตฺตสมฺพาธสยนํ อชฺฌาวสาม, กาสิกจนฺทนํ ปจฺจนุโภม, มาลาคนฺธวิเลปนํ ธารยาม, ชาตรูปรชตํ สาทยาม.\csromanspacer{} เตสํ โน ภนฺเต ภควา อมฺหากํ ตถา ธมฺมํ เทเสตุ, เย อมฺหากํ อสฺสุ ธมฺมา ทิฏฺฐธมฺมหิตาย ทิฏฺฐธมฺมสุขาย สมฺปรายหิตาย สมฺปรายสุขายา”ติ.}

\prose{จตฺตาโรเม พฺยคฺฆปชฺช ธมฺมา กุลปุตฺตสฺส ทิฏฺฐธมฺมหิตาย สํวตฺตนฺติ ทิฏฺฐธมฺมสุขาย.\csromanspacer{} กตเม จตฺตาโร? อุฏฺฐานสมฺปทา อารกฺขสมฺปทา กลฺยาณมิตฺตตา สมชีวิตา\footnote{สมชีวิกตา (สี), อุปริ 140 ปิฏฺเฐปิ.}.\csromanspacer{} กตมา จ พฺยคฺฆปชฺช อุฏฺฐานสมฺปทา? อิธ พฺยคฺฆปชฺช \csromanfolio{108}กุลปุตฺโต เยน กมฺมฏฺฐาเนน ชีวิกํ\footnote{ชีวิตํ (ก)} กปฺเปติ ยทิ กสิยา ยทิ วณิชฺชาย ยทิ โครกฺเขน ยทิ อิสฺสตฺเตน\footnote{อิสฺสตฺเถน (สี, สฺยา, อิ)} ยทิ ราชโปริเสน ยทิ สิปฺปญฺญตเรน, ตตฺถ ทกฺโข โหติ อนลโส, ตตฺรุปายาย วีมํสาย สมนฺนาคโต อลํ กาตุํ อลํ สํวิธาตุํ.\csromanspacer{} อยํ วุจฺจติ พฺยคฺฆปชฺช อุฏฺฐานสมฺปทา.}

\prose{กตมา จ พฺยคฺฆปชฺช อารกฺขสมฺปทา? อิธ พฺยคฺฆปชฺช กุลปุตฺตสฺส โภคา โหนฺติ อุฏฺฐานวีริยาธิคตา พาหาพลปริจิตา เสทาวกฺขิตฺตา ธมฺมิกา ธมฺมลทฺธา, เต อารกฺเขน คุตฺติยา สมฺปาเทติ “กินฺติ เม อิเม โภเค เนว ราชาโน หเรยฺยุํ, น โจรา หเรยฺยุํ, น อคฺคิ ฑเหยฺย, น อุทกํ วเหยฺย, น อปฺปิยา ทายาทา หเรยฺยุนฺ”ติ.\csromanspacer{} อยํ วุจฺจติ พฺยคฺฆปชฺช อารกฺขสมฺปทา.}

\prose{กตมา จ พฺยคฺฆปชฺช กลฺยาณมิตฺตตา? อิธ พฺยคฺฆปชฺช กุลปุตฺโต ยสฺมิํ คาเม วา นิคเม วา ปฏิวสติ, ตตฺถ เย เต โหนฺติ คหปตี วา คหปติปุตฺตา วา ทหรา วา วุทฺธสีลิโน วุทฺธา วา วุทฺธสีลิโน สทฺธาสมฺปนฺนา สีลสมฺปนฺนา จาคสมฺปนฺนา ปญฺญาสมฺปนฺนา, เตหิ สทฺธิํ สนฺติฏฺฐติ สลฺลปติ สากจฺฉํ สมาปชฺชติ, ยถารูปานํ สทฺธาสมฺปนฺนานํ สทฺธาสมฺปทํ อนุสิกฺขติ, ยถารูปานํ สีลสมฺปนฺนานํ สีลสมฺปทํ อนุสิกฺขติ, ยถารูปานํ จาคสมฺปนฺนานํ จาคสมฺปทํ อนุสิกฺขติ, ยถารูปานํ ปญฺญาสมฺปนฺนานํ ปญฺญาสมฺปทํ อนุสิกฺขติ.\csromanspacer{} อยํ วุจฺจติ พฺยคฺฆปชฺช กลฺยาณมิตฺตตา.}

\prose{กตมา จ พฺยคฺฆปชฺช สมชีวิตา? อิธ พฺยคฺฆปชฺช กุลปุตฺโต อายญฺจ โภคานํ วิทิตฺวา วยญฺจ โภคานํ วิทิตฺวา สมํ ชีวิกํ\footnote{สมชีวิกํ (สฺยา), สมชีวิตํ (ก) 4. โอณตํ (ก)} กปฺเปติ นาจฺโจคาฬฺหํ นาติหีนํ “เอวํ เม อาโย วยํ ปริยาทาย ฐสฺสติ, น จ เม วโย อายํ ปริยาทาย ฐสฺสตี”ติ.\csromanspacer{} เสยฺยถาปิ พฺยคฺฆปชฺช ตุลาธาโร วา ตุลาธารนฺเตวาสี วา ตุลํ ปคฺคเหตฺวา ชานาติ “เอตฺตเกน วา โอนตํ4 เอตฺตเกน วา อุนฺนตนฺ”ติ\footnote{อุณฺณตนฺติ (ก)}.\csromanspacer{} เอวเมวํ โข พฺยคฺฆปชฺช กุลปุตฺโต อายญฺจ โภคานํ วิทิตฺวา วยญฺจ โภคานํ วิทิตฺวา สมํ ชีวิกํ กปฺเปติ นาจฺโจคาฬฺหํ นาติหีนํ “เอวํ เม อาโย วยํ ปริยาทาย ฐสฺสติ, น จ เม วโย อายํ ปริยาทาย ฐสฺสตี”ติ.\csromanspacer{} สจายํ พฺยคฺฆปชฺช กุลปุตฺโต}

\vspace{\tipitakaparskip}
\csromancenter{(6) 1.\csromanspacer{} โคตมีวคฺค อปฺปาโย สมาโน อุฬารํ ชีวิกํ\footnote{ชีวิตํ (ก)} กปฺเปติ, ตสฺส ภวนฺติ วตฺตาโร “อุทุมฺพรขาทีวายํ\footnote{อุทุมฺพรขาทิกํ วายํ (สี, อิ), อุทุมฺพรขาทกํ จายํ (สฺยา)} กุลปุตฺโต โภเค ขาทตี”ติ.\csromanspacer{} สเจ ปนายํ พฺยคฺฆปชฺช กุลปุตฺโต มหาโย สมาโน กสิรํ ชีวิกํ1 กปฺเปติ, ตสฺส ภวนฺติ วตฺตาโร “อเชฏฺฐมรณํวายํ\footnote{อชทฺธุมาริกํ วายํ (สี, อิ), อทฺธมารกํ จายํ (สฺยา), เอตฺถ ชทฺธูติ อสนํ = ภตฺตภุญฺชนํ, ตสฺมา อชทฺธุมาริกนฺติ อนสนมรณนฺติ วุตฺตํ โหติ. ม 1. 312 ปิฏฺเฐปิ อโธลิปิยา “อชทฺธุกนฺ”ติ ปทํ ทสฺสิตํ.} กุลปุตฺโต มริสฺสตี”ติ.\csromanspacer{} ยโต จ โขยํ พฺยคฺฆปชฺช กุลปุตฺโต อายญฺจ โภคานํ วิทิตฺวา วยญฺจ โภคานํ วิทิตฺวา สมํ ชีวิกํ กปฺเปติ นาจฺโจคาฬฺหํ นาติหีนํ “เอวํ เม อาโย วยํ ปริยาทาย ฐสฺสติ, น จ เม วโย อายํ ปริยาทาย ฐสฺสตี”ติ.\csromanspacer{} อยํ วุจฺจติ พฺยคฺฆปชฺช สมชีวิตา.}
\prose{\csromanfolio{109}เอวํ สมุปฺปนฺนานํ พฺยคฺฆปชฺช โภคานํ จตฺตาริ อปายมุขานิ โหนฺติ อิตฺถิธุตฺโต สุราธุตฺโต อกฺขธุตฺโต ปาปมิตฺโต ปาปสหาโย ปาปสมฺปวงฺโก.\csromanspacer{} เสยฺยถาปิ พฺยคฺฆปชฺช มหโต ตฬากสฺส จตฺตาริ เจว อายมุขานิ จตฺตาริ จ อปายมุขานิ, ตสฺส ปุริโส ยานิ เจว อายมุขานิ, ตานิ ปิทเหยฺย.\csromanspacer{} ยานิ จ อปายมุขานิ, ตานิ วิวเรยฺย.\csromanspacer{} เทโว จ น สมฺมา ธารํ อนุปฺปเวจฺเฉยฺย.\csromanspacer{} เอวํ หิ ตสฺส พฺยคฺฆปชฺช มหโต ตฬากสฺส ปริหานิเยว ปาฏิกงฺขา โน วุทฺธิ.\csromanspacer{} เอวเมวํ พฺยคฺฆปชฺช เอวํ สมุปฺปนฺนานํ โภคานํ จตฺตาริ อปายมุขานิ โหนฺติ อิตฺถิธุตฺโต สุราธุตฺโต อกฺขธุตฺโต ปาปมิตฺโต ปาปสหาโย ปาปสมฺปวงฺโก.}

\prose{เอวํ สมุปฺปนฺนานํ พฺยคฺฆปชฺช โภคานํ จตฺตาริ อายมุขานิ โหนฺติ น อิตฺถิธุตฺโต น สุราธุตฺโต น อกฺขธุตฺโต กลฺยาณมิตฺโต กลฺยาณสหาโย กลฺยาณสมฺปวงฺโก.\csromanspacer{} เสยฺยถาปิ พฺยคฺฆปชฺช มหโต ตฬากสฺส จตฺตาริ เจว อายมุขานิ จตฺตาริ จ อปายมุขานิ, ตสฺส ปุริโส ยานิ เจว อายมุขานิ, ตานิ วิวเรยฺย.\csromanspacer{} ยานิ จ อปายมุขานิ, ตานิ ปิทเหยฺย.\csromanspacer{} เทโว จ สมฺมา ธารํ อนุปฺปเวจฺเฉยฺย.\csromanspacer{} เอวํ หิ ตสฺส พฺยคฺฆปชฺช มหโต ตฬากสฺส วุทฺธิเยว ปาฏิกงฺขา โน ปริหานิ.\csromanspacer{} เอวเมวํ โข พฺยคฺฆปชฺช เอวํ สมุปฺปนฺนานํ โภคานํ จตฺตาริ อายมุขานิ โหนฺติ น \csromanfolio{110}อิตฺถิธุตฺโต น สุราธุตฺโต น อกฺขธุตฺโต กลฺยาณมิตฺโต กลฺยาณสหาโย กลฺยาณสมฺปวงฺโก.\csromanspacer{} อิเม โข พฺยคฺฆปชฺช จตฺตาโร ธมฺมา กุลปุตฺตสฺส ทิฏฺฐธมฺมหิตาย สํวตฺตนฺติ ทิฏฺฐธมฺมสุขาย.}

\prose{จตฺตาโรเม พฺยคฺฆปชฺช ธมฺมา กุลปุตฺตสฺส สมฺปรายหิตาย สํวตฺตนฺติ สมฺปรายสุขาย.\csromanspacer{} กตเม จตฺตาโร? สทฺธาสมฺปทา สีลสมฺปทา จาคสมฺปทา ปญฺญาสมฺปทา.\csromanspacer{} กตมา จ พฺยคฺฆปชฺช สทฺธาสมฺปทา? อิธ พฺยคฺฆปชฺช กุลปุตฺโต สทฺโธ โหติ, สทฺทหติ ตถาคตสฺส โพธิํ “อิติปิ โส ภควา ฯเปฯ สตฺถา เทวมนุสฺสานํ พุทฺโธ ภควา”ติ.\csromanspacer{} อยํ วุจฺจติ พฺยคฺฆปชฺช สทฺธาสมฺปทา.}

\prose{กตมา จ พฺยคฺฆปชฺช สีลสมฺปทา? อิธ พฺยคฺฆปชฺช กุลปุตฺโต ปาณาติปาตา ปฏิวิรโต โหติ ฯเปฯ สุราเมรยมชฺชปมาทฏฺฐานา ปฏิวิรโต โหติ.\csromanspacer{} อยํ วุจฺจติ พฺยคฺฆปชฺช สีลสมฺปทา.}

\prose{กตมา จ พฺยคฺฆปชฺช จาคสมฺปทา? อิธ พฺยคฺฆปชฺช กุลปุตฺโต วิคตมลมจฺเฉเรน เจตสา อคารํ อชฺฌาวสติ มุตฺตจาโค ปยตปาณิ โวสคฺครโต ยาจโยโค ทานสํวิภาครโต.\csromanspacer{} อยํ วุจฺจติ พฺยคฺฆปชฺช จาคสมฺปทา.}

\prose{กตมา จ พฺยคฺฆปชฺช ปญฺญาสมฺปทา? อิธ พฺยคฺฆปชฺช กุลปุตฺโต ปญฺญวา โหติ, อุทยตฺถคามินิยา ปญฺญาย สมนฺนาคโต อริยาย นิพฺเพธิกาย สมฺมา ทุกฺขกฺขยคามินิยา.\csromanspacer{} อยํ วุจฺจติ พฺยคฺฆปชฺช ปญฺญาสมฺปทา.\csromanspacer{} อิเม โข พฺยคฺฆปชฺช จตฺตาโร ธมฺมา กุลปุตฺตสฺส สมฺปราย หิตาย สํวตฺตนฺติ สมฺปรายสุขายาติ.}

\csromangathameasure{อุฏฺฐาตา กมฺมเธยฺเยสุ, อปฺปมตฺโต วิธานวา. \\ สมํ กปฺเปติ ชีวิกํ, สมฺภตํ อนุรกฺขติ. \\ สทฺโธ สีเลน สมฺปนฺโน, วทญฺญู วีตมจฺฉโร. \\ นิจฺจํ มคฺคํ วิโสเธติ, โสตฺถานํ สมฺปรายิกํ. \\ อิจฺเจเต อฏฺฐ ธมฺมา จ, สทฺธสฺส ฆรเมสิโน. \\ อกฺขาตา สจฺจนาเมน, อุภยตฺถ สุขาวหา. \\ ทิฏฺฐธมฺมหิตตฺถาย, สมฺปรายสุขาย จ. \\ เอวเมตํ คหฏฺฐานํ, จาโค ปุญฺญํ ปวฑฺฒตีติ.จตุตฺถํ.}
\csromangathagroup{\csromangathastanza{อุฏฺฐาตา กมฺมเธยฺเยสุ, อปฺปมตฺโต วิธานวา. \\ สมํ กปฺเปติ ชีวิกํ\footnote{ชีวิตํ (ก)}, สมฺภตํ อนุรกฺขติ.}\csromangathastanza{สทฺโธ สีเลน สมฺปนฺโน, วทญฺญู วีตมจฺฉโร. \\ นิจฺจํ มคฺคํ วิโสเธติ, โสตฺถานํ สมฺปรายิกํ.}\csromangathastanza{อิจฺเจเต อฏฺฐ ธมฺมา จ, สทฺธสฺส ฆรเมสิโน. \\ อกฺขาตา สจฺจนาเมน, อุภยตฺถ สุขาวหา.}\csromangathastanza{ทิฏฺฐธมฺมหิตตฺถาย, สมฺปรายสุขาย จ. \\ เอวเมตํ คหฏฺฐานํ, จาโค ปุญฺญํ ปวฑฺฒตีติ.\csromanspacer{}\csromanspacer{}จตุตฺถํ.}}

\chapterheadpageread{(6) 1. โคตมีวคฺค \textbf{5. อุชฺชยสุตฺต}}
\proseitem{55}{\csromanfolio{111}อถ โข อุชฺชโย พฺราหฺมโณ เยน ภควา เตนุปสงฺกมิ, อุปสงฺกมิตฺวา ภควตา สทฺธิํ สมฺโมทิ, สมฺโมทนียํ กถํ สารณียํ วีติสาเรตฺวา เอกมนฺตํ นิสีทิ, เอกมนฺตํ นิสินฺโน โข อุชฺชโย พฺราหฺมโณ ภควนฺตํ เอตทโวจ “มยํ โภ โคตม ปวาสํ คนฺตุกามา, เตสํ โน ภวํ โคตโม อมฺหากํ ตถา ธมฺมํ เทเสตุ, เย อมฺหากํ อสฺสุ ธมฺมา ทิฏฺฐธมฺมหิตาย ทิฏฺฐธมฺมสุขาย สมฺปรายหิตาย สมฺปรายสุขายา”ติ.}

\prose{จตฺตาโรเม พฺราหฺมณ ธมฺมา กุลปุตฺตสฺส ทิฏฺฐธมฺมหิตาย สํวตฺตนฺติ ทิฏฺฐธมฺมสุขาย.\csromanspacer{} กตเม จตฺตาโร? อุฏฺฐานสมฺปทา อารกฺขสมฺปทา กลฺยาณมิตฺตตา สมชีวิตา.\csromanspacer{} กตมา จ พฺราหฺมณ อุฏฺฐานสมฺปทา? อิธ พฺราหฺมณ กุลปุตฺโต เยน กมฺมฏฺฐาเนน ชีวิกํ กปฺเปติ ยทิ กสิยา ยทิ วณิชฺชาย ยทิ โครกฺเขน ยทิ อิสฺสตฺเตน ยทิ ราชโปริเสน ยทิ สิปฺปญฺญตเรน.\csromanspacer{} ตตฺถ ทกฺโข โหติ อนลโส, ตตฺรุปายาย วีมํสาย สมนฺนาคโต อลํ กาตุํ อลํ สํวิธาตุํ.\csromanspacer{} อยํ วุจฺจติ พฺราหฺมณ อุฏฺฐานสมฺปทา.}

\prose{กตมา จ พฺราหฺมณ อารกฺขสมฺปทา? อิธ พฺราหฺมณ กุลปุตฺตสฺส โภคา โหนฺติ อุฏฺฐานวีริยาธิคตา พาหาพลปริจิตา เสทาวกฺขิตฺตา ธมฺมิกา ธมฺมลทฺธา.\csromanspacer{} เต อารกฺเขน คุตฺติยา สมฺปาเทติ “กินฺติ เม อิเม โภเค เนว ราชาโน หเรยฺยุํ, น โจรา หเรยฺยุํ, น อคฺคิ ฑเหยฺย, น อุทกํ วเหยฺย, น อปฺปิยา ทายาทา หเรยฺยุนฺ”ติ.\csromanspacer{} อยํ วุจฺจติ พฺราหฺมณ อารกฺขสมฺปทา.}

\prose{กตมา จ พฺราหฺมณ กลฺยาณมิตฺตตา? อิธ พฺราหฺมณ กุลปุตฺโต ยสฺมิํ คาเม วา นิคเม วา ปฏิวสติ, ตตฺร เย เต โหนฺติ คหปตี วา คหปติปุตฺตา วา ทหรา วา วุทฺธสีลิโน วุทฺธา วา วุทฺธสีลิโน สทฺธาสมฺปนฺนา สีลสมฺปนฺนา จาคสมฺปนฺนา ปญฺญาสมฺปนฺนา, เตหิ สทฺธิํ สนฺติฏฺฐติ สลฺลปติ สากจฺฉํ สมาปชฺชติ.\csromanspacer{} ยถารูปานํ สทฺธาสมฺปนฺนานํ สทฺธาสมฺปทํ อนุสิกฺขติ, ยถารูปานํ สีลสมฺปนฺนานํ สีลสมฺปทํ อนุสิกฺขติ, ยถารูปานํ จาคสมฺปนฺนานํ จาคสมฺปทํ อนุสิกฺขติ, ยถารูปานํ ปญฺญาสมฺปนฺนานํ ปญฺญาสมฺปทํ อนุสิกฺขติ.\csromanspacer{} อยํ วุจฺจติ พฺราหฺมณ กลฺยาณมิตฺตตา.}

\prose{\csromanfolio{112}กตมา จ พฺราหฺมณ สมชีวิตา? อิธ พฺราหฺมณ กุลปุตฺโต อายญฺจ โภคานํ วิทิตฺวา วยญฺจ โภคานํ วิทิตฺวา สมํ ชีวิกํ กปฺเปติ นาจฺโจคาฬฺหํ นาติหีนํ “เอวํ เม อาโย วยํ ปริยาทาย ฐสฺสติ, น จ เม วโย อายํ ปริยาทาย ฐสฺสตี”ติ.\csromanspacer{} เสยฺยถาปิ พฺราหฺมณ ตุลาธาโร วา ตุลาธารนฺเตวาสี วา ตุลํ ปคฺคเหตฺวา ชานาติ “เอตฺตเกน วา โอนตํ เอตฺตเกน วา อุนฺนตนฺ”ติ.\csromanspacer{} เอวเมวํ โข พฺราหฺมณ กุลปุตฺโต อายญฺจ โภคานํ วิทิตฺวา วยญฺจ โภคานํ วิทิตฺวา สมํ ชีวิกํ กปฺเปติ นาจฺโจคาฬฺหํ นาติหีนํ “เอวํ เม อาโย วยํ ปริยาทาย ฐสฺสติ, น จ เม วโย อายํ ปริยาทาย ฐสฺสตี”ติ.\csromanspacer{} สจายํ พฺราหฺมณ กุลปุตฺโต อปฺปาโย สมาโน อุฬารํ ชีวิกํ กปฺเปติ, ตสฺส ภวนฺติ วตฺตาโร “อุทุมฺพรขาทีวายํ กุลปุตฺโต โภเค ขาทตี”ติ.\csromanspacer{} สเจ ปนายํ พฺราหฺมณ กุลปุตฺโต มหาโย สมาโน กสิรํ ชีวิกํ กปฺเปติ, ตสฺส ภวนฺติ วตฺตาโร “อเชฏฺฐมรณํวายํ กุลปุตฺโต มริสฺสตี”ติ.\csromanspacer{} ยโต จ โขยํ พฺราหฺมณ กุลปุตฺโต อายญฺจ โภคานํ วิทิตฺวา วยญฺจ โภคานํ วิทิตฺวา สมํ ชีวิกํ กปฺเปติ นาจฺโจคาฬฺหํ นาติหีนํ “เอวํ เม อาโย วยํ ปริยาทาย ฐสฺสติ, น จ เม วโย อายํ ปริยาทาย ฐสฺสตี”ติ.\csromanspacer{} อยํ วุจฺจติ พฺราหฺมณ สมชีวิตา.}

\prose{เอวํ สมุปฺปนฺนานํ พฺราหฺมณ โภคานํ จตฺตาริ อปายมุขานิ โหนฺติ อิตฺถิธุตฺโต สุราธุตฺโต อกฺขธุตฺโต ปาปมิตฺโต ปาปสหาโย ปาปสมฺปวงฺโก.\csromanspacer{} เสยฺยถาปิ พฺราหฺมณ มหโต ตฬากสฺส จตฺตาริ เจว อายมุขานิ จตฺตาริ จ อปายมุขานิ, ตสฺส ปุริโส ยานิ เจว อายมุขานิ, ตานิ ปิทเหยฺย.\csromanspacer{} ยานิ จ อปายมุขานิ, ตานิ วิวเรยฺย.\csromanspacer{} เทโว จ น สมฺมา ธารํ อนุปฺปเวจฺเฉยฺย, เอวํ หิ ตสฺส พฺราหฺมณ มหโต ตฬากสฺส ปริหานิเยว ปาฏิกงฺขา โน วุทฺธิ.\csromanspacer{} เอวเมวํ โข พฺราหฺมณ เอวํ สมุปฺปนฺนานํ โภคานํ จตฺตาริ อปายมุขานิ โหนฺติ อิตฺถิธุตฺโต สุราธุตฺโต อกฺขธุตฺโต ปาปมิตฺโต ปาปสหาโย ปาปสมฺปวงฺโก.}

\prose{เอวํ สมุปฺปนฺนานํ พฺราหฺมณ โภคานํ จตฺตาริ อายมุขานิ โหนฺติ น อิตฺถิธุตฺโต น สุราธุตฺโต น อกฺขธุตฺโต กลฺยาณมิตฺโต กลฺยาณสหาโย กลฺยาณสมฺปวงฺโก.\csromanspacer{} เสยฺยถาปิ พฺราหฺมณ มหโต ตฬากสฺส จตฺตาริ เจว อายมุขานิ จตฺตาริ จ อปายมุขานิ, ตสฺส ปุริโส ยานิ เจว อายมุขานิ, ตานิ วิวเรยฺย.\csromanspacer{} ยานิ จ}

\vspace{\tipitakaparskip}
\csromancenter{(6) 1.\csromanspacer{} โคตมีวคฺค อปายมุขานิ, ตานิ ปิทเหยฺย.\csromanspacer{} เทโว จ สมฺมา ธารํ อนุปฺปเวจฺเฉยฺย.\csromanspacer{} เอวํ หิ ตสฺส พฺราหฺมณ มหโต ตฬากสฺส วุทฺธิเยว ปาฏิกงฺขา โน ปริหานิ.\csromanspacer{} เอวเมวํ โข พฺราหฺมณ เอวํ สมุปฺปนฺนานํ โภคานํ จตฺตาริ อายมุขานิ โหนฺติ น อิตฺถิธุตฺโต ฯเปฯ กลฺยาณสมฺปวงฺโก.\csromanspacer{} อิเม โข พฺราหฺมณ จตฺตาโร ธมฺมา กุลปุตฺตสฺส ทิฏฺฐธมฺมหิตาย สํวตฺตนฺติ ทิฏฺฐธมฺมสุขาย.}
\prose{\csromanfolio{113}จตฺตาโรเม พฺราหฺมณ กุลปุตฺตสฺส ธมฺมา สมฺปรายหิตาย สํวตฺตนฺติ สมฺปรายสุขาย.\csromanspacer{} กตเม จตฺตาโร? สทฺธาสมฺปทา สีลสมฺปทา จาคสมฺปทา ปญฺญาสมฺปทา.\csromanspacer{} กตมา จ พฺราหฺมณ สทฺธาสมฺปทา, อิธ พฺราหฺมณ กุลปุตฺโต สทฺโธ โหติ, สทฺทหติ ตถาคตสฺส โพธิํ “อิติปิ โส ภควา ฯเปฯ สตฺถา เทวมนุสฺสานํ พุทฺโธ ภควา”ติ.\csromanspacer{} อยํ วุจฺจติ พฺราหฺมณ สทฺธาสมฺปทา.}

\prose{กตมา จ พฺราหฺมณ สีลสมฺปทา? อิธ พฺราหฺมณ กุลปุตฺโต ปาณาติปาตา ปฏิวิรโต โหติ ฯเปฯ สุราเมรยมชฺชปมาทฏฺฐานา ปฏิวิรโต โหติ.\csromanspacer{} อยํ วุจฺจติ พฺราหฺมณ สีลสมฺปทา.}

\prose{กตมา จ พฺราหฺมณ จาคสมฺปทา? อิธ พฺราหฺมณ กุลปุตฺโต วิคตมลมจฺเฉเรน เจตสา อคารํ อชฺฌาวสติ มุตฺตจาโค ปยตปาณิ โวสคฺครโต ยาจโยโค ทานสํวิภาครโต.\csromanspacer{} อยํ วุจฺจติ พฺราหฺมณ จาคสมฺปทา.}

\prose{กตมา จ พฺราหฺมณ ปญฺญาสมฺปทา? อิธ พฺราหฺมณ กุลปุตฺโต ปญฺญวา โหติ ฯเปฯ สมฺมา ทุกฺขกฺขยคามินิยา.\csromanspacer{} อยํ วุจฺจติ พฺราหฺมณ ปญฺญาสมฺปทา.\csromanspacer{} อิเม โข พฺราหฺมณ จตฺตาโร ธมฺมา กุลปุตฺตสฺส สมฺปรายหิตาย สํวตฺตนฺติ สมฺปรายสุขายาติ.}

\prose{อุฏฺฐาตา กมฺมเธยฺเยสุ, อปฺปมตฺโต วิธานวา.}

\prose{สมํ กปฺเปติ ชีวิกํ, สมฺภตํ อนุรกฺขติ.}

\prose{สทฺโธ สีเลน สมฺปนฺโน, วทญฺญู วีตมจฺฉโร.}

\prose{นิจฺจํ มคฺคํ วิโสเธติ, โสตฺถานํ สมฺปรายิกํ.}

\prose{อิจฺเจเต อฏฺฐ ธมฺมา จ, สทฺธสฺส ฆรเมสิโน.}

\prose{อกฺขาตา สจฺจนาเมน, อุภยตฺถ สุขาวหา.}

\prose{ทิฏฺฐธมฺมหิตตฺถาย, สมฺปรายสุขาย จ.}

\prose{เอวเมตํ คหฏฺฐานํ, จาโค ปุญฺญํ ปวฑฺฒตีติ.\csromanspacer{}\csromanspacer{}ปญฺจมํ.}

\csromanheader{2}{6. ภยสุตฺต}
\proseitem{56}{\csromanfolio{114}“ภยนฺ”ติ * ภิกฺขเว กามานเมตํ อธิวจนํ. “ทุกฺขนฺ”ติ ภิกฺขเว กามนเมตํ อธิวจนํ. “โรโค”ติ ภิกฺขเว กามานเมตํ อธิวจนํ. “คณฺโฑ”ติ ภิกฺขเว กามานเมตํ อธิวจนํ. “สลฺลนฺ”ติ ภิกฺขเว กามานเมตํ อธิวจนํ. “สงฺโค”ติ ภิกฺขเว กามานเมตํ อธิวจนํ. “ปงฺโก”ติ ภิกฺขเว กามานเมตํ อธิวจนํ. “คพฺโภ”ติ ภิกฺขเว กามานเมตํ อธิวจนํ.\csromanspacer{} กสฺมา จ ภิกฺขเว “ภยนฺ”ติ กามานเมตํ อธิวจนํ.\csromanspacer{} ยสฺมา จ กามราครตฺตายํ ภิกฺขเว ฉนฺทราควินิพทฺโธ ทิฏฺฐธมฺมิกาปิ ภยา น ปริมุจฺจติ, สมฺปรายิกาปิ ภยา น ปริมุจฺจติ, ตสฺมา “ภยนฺ”ติ กามานเมตํ อธิวจนํ.\csromanspacer{} กสฺมา จ ภิกฺขเว “ทุกฺขนฺ”ติ ฯเปฯ “โรโค”ติ. “คณฺโฑ”ติ. “สลฺลนฺ”ติ. “สงฺโค”ติ. “ปงฺโก”ติ. “คพฺโภ”ติ กามานเมตํ อธิวจนํ.\csromanspacer{} ยสฺมา จ กามราครตฺตายํ ภิกฺขเว ฉนฺทราควินิพทฺโธ ทิฏฺฐธมฺมิกาปิ คพฺภา น ปริมุจฺจติ, สมฺปรายิกาปิ คพฺภา น ปริมุจฺจติ, ตสฺมา “คพฺโภ”ติ กามานเมตํ อธิวจนํ.}

\csromangathameasure{ภยํ ทุกฺขญฺจ โรโค จ, คณฺโฑ สลฺลญฺจ สงฺโค จ. \\ ปงฺโก คพฺโภ จ อุภยํ, เอเต กามา ปวุจฺจนฺติ. ยตฺถ สตฺโต ปุถุชฺชโน. \\ โอติณฺโณ สาตรูเปน, ปุน คพฺภาย คจฺฉติ. \\ ยโต จ ภิกฺขุ อาตาปี, สมฺปชญฺญํ น ริจฺจติ. \\ โส อิมํ ปลิปถํ ทุคฺคํ, อติกฺกมฺม ตถาวิโธ. \\ ปชํ ชาติชรูเปตํ, ผนฺทมานํ อเวกฺขตีติ.ฉฏฺฐํ.}
\csromangathagroup{\csromangathastanza{ภยํ ทุกฺขญฺจ โรโค จ, คณฺโฑ สลฺลญฺจ สงฺโค จ. \\ ปงฺโก คพฺโภ จ อุภยํ, เอเต กามา ปวุจฺจนฺติ.\csromanspacer{} ยตฺถ สตฺโต ปุถุชฺชโน.}\csromangathastanza{โอติณฺโณ สาตรูเปน, ปุน คพฺภาย คจฺฉติ. \\ ยโต จ ภิกฺขุ อาตาปี, สมฺปชญฺญํ\footnote{สมฺปชญฺเญ (สฺยา, ก) สํ 2. 408 ปิฏฺเฐ ปสฺสิตพฺพํ.} น ริจฺจติ.}\csromangathastanza{โส อิมํ ปลิปถํ ทุคฺคํ, อติกฺกมฺม ตถาวิโธ. \\ ปชํ ชาติชรูเปตํ, ผนฺทมานํ อเวกฺขตีติ.\csromanspacer{}\csromanspacer{}ฉฏฺฐํ.}}

\csromanheader{2}{7. ปฐม อาหุเนยฺยสุตฺต}
\proseitem{57}{อฏฺฐหิ ภิกฺขเว ธมฺเมหิ สมนฺนาคโต ภิกฺขุ อาหุเนยฺโย โหติ ปาหุเนยฺโย ทกฺขิเณยฺโย อญฺชลิกรณีโย อนุตฺตรํ ปุญฺญกฺเขตฺตํ โลกสฺส.\csromanspacer{} กตเมหิ อฏฺฐหิ? อิธ ภิกฺขเว ภิกฺขุ สีลวา โหติ ฯเปฯ สมาทาย สิกฺขติ สิกฺขาปเทสุ.\csromanspacer{} พหุสฺสุโต โหติ ฯเปฯ ทิฏฺฐิยา สุปฺปฏิวิทฺธา.\csromanspacer{} กลฺยาณมิตฺโต โหติ กลฺยาณสหาโย}

\vspace{\tipitakaparskip}
\csromancenter{(6) 1.\csromanspacer{} โคตมีวคฺค กลฺยาณสมฺปวงฺโก.\csromanspacer{} สมฺมาทิฏฺฐิโก โหติ สมฺมาทสฺสเนน สมนฺนาคโต.\csromanspacer{} จตุนฺนํ ฌานานํ อาภิเจตสิกานํ ทิฏฺฐธมฺมสุขวิหารานํ นิกามลาภี โหติ อกิจฺฉลาภี อกสิรลาภี.\csromanspacer{} อเนกวิหิตํ ปุพฺเพนิวาสํ อนุสฺสรติ.\csromanspacer{} เสยฺยถิทํ, เอกมฺปิ ชาติํ เทฺวปิ ชาติโย ฯเปฯ อิติ สาการํ ส-อุทฺเทสํ อเนกวิหิตํ ปุพฺเพนิวาสํ อนุสฺสรติ.\csromanspacer{} ทิพฺเพน จกฺขุนา วิสุทฺเธน อติกฺกนฺตมานุสเกน ฯเปฯ ยถากมฺมูปเค สตฺเต ปชานาติ.\csromanspacer{} อาสวานํ ขยา ฯเปฯ สจฺฉิกตฺวา อุปสมฺปชฺช วิหรติ.\csromanspacer{} อิเมหิ โข ภิกฺขเว อฏฺฐหิ ธมฺเมหิ สมนฺนาคโต ภิกฺขุ อาหุเนยฺโย โหติ ฯเปฯ อนุตฺตรํ ปุญฺญกฺเขตฺตํ โลกสฺสาติ.\csromanspacer{}\csromanspacer{}สตฺตมํ.}
\csromanheader{2}{8. ทุติย อาหุเนยฺยสุตฺต}
\proseitem{58}{\csromanfolio{115}อฏฺฐหิ ภิกฺขเว ธมฺเมหิ สมนฺนาคโต ภิกฺขุ อาหุเนยฺโย โหติ ฯเปฯ อนุตฺตรํ ปุญฺญกฺเขตฺตํ โลกสฺส.\csromanspacer{} กตเมหิ อฏฺฐหิ? อิธ ภิกฺขเว ภิกฺขุ สีลวา โหติ ฯเปฯ สมาทาย สิกฺขติ สิกฺขาปเทสุ.\csromanspacer{} พหุสฺสุโต โหติ ฯเปฯ ทิฏฺฐิยา สุปฺปฏิวิทฺธา.\csromanspacer{} อารทฺธวีริโย วิหรติ, ถามวา ทฬฺหปรกฺกโม อนิกฺขิตฺตธุโร กุสเลสุ ธมฺเมสุ.\csromanspacer{} อารญฺญิโก โหติ ปนฺตเสนาสโน.\csromanspacer{} อรติรติสโห โหติ, อุปฺปนฺนํ อรติํ อภิภุยฺย อภิภุยฺย วิหรติ.\csromanspacer{} ภยเภรวสโห โหติ, อุปฺปนฺนํ ภยเภรวํ อภิภุยฺย อภิภุยฺย วิหรติ.\csromanspacer{} จตุนฺนํ ฌานานํ อาภิเจตสิกานํ ทิฏฺฐธมฺมสุขวิหารานํ นิกามลาภี โหติ อกิจฺฉลาภี อกสิรลาภี.\csromanspacer{} อาสวานํ ขยา ฯเปฯ สจฺฉิกตฺวา อุปสมฺปชฺช วิหรติ.\csromanspacer{} อิเมหิ โข ภิกฺขเว อฏฺฐหิ ธมฺเมหิ สมนฺนาคโต ภิกฺขุ อาหุเนยฺโย ฯเปฯ อนุตฺตรํ ปุญฺญกฺเขตฺตํ โลกสฺสาติ.\csromanspacer{}\csromanspacer{}อฏฺฐมํ.}

\csromanheader{2}{9. ปฐมปุคฺคลสุตฺต}
\proseitem{59}{อฏฺฐิเม ภิกฺขเว ปุคฺคลา อาหุเนยฺยา ปาหุเนยฺยา ทกฺขิเณยฺยา อญฺชลิกรณียา อนุตฺตรํ ปุญฺญกฺเขตฺตํ โลกสฺส.\csromanspacer{} กตเม อฏฺฐ? โสตาปนฺโน โสตาปตฺติผลสจฺฉิกิริยาย ปฏิปนฺโน, สกทาคามี สกทาคามิผล- สจฺฉิกิริยาย ปฏิปนฺโน, อนาคามี อนาคามิผลสจฺฉิกิริยาย ปฏิปนฺโน, อรหา อรหตฺตาย ปฏิปนฺโน.\csromanspacer{} อิเม โข ภิกฺขเว อฏฺฐ ปุคฺคลา อาหุเนยฺยา ฯเปฯ อนุตฺตรํ ปุญฺญกฺเขตฺตํ โลกสฺสาติ.}

\csromangathameasure{จตฺตาโร จ ปฏิปนฺนา, จตฺตาโร จ ผเล ฐิตา. \\ เอส สํโฆ อุชุภูโต, ปญฺญาสีลสมาหิโต. \\ ยชมานานํ มนุสฺสานํ, \\ ปุญฺญเปกฺขาน ปาณินํ. \\ กโรตํ โอปธิกํ ปุญฺญํ,}
\csromangathagroup{\csromangathastanza{\csromanfolio{116}จตฺตาโร จ ปฏิปนฺนา, จตฺตาโร จ ผเล ฐิตา. \\ เอส สํโฆ อุชุภูโต, ปญฺญาสีลสมาหิโต.}\csromangathastanza{ยชมานานํ มนุสฺสานํ, \\ ปุญฺญเปกฺขาน ปาณินํ. \\ กโรตํ โอปธิกํ ปุญฺญํ,}}

\vspace{\tipitakaparskip}
\csromancenter{สํเฆ ทินฺนํ มหปฺผลนฺติ.\csromanspacer{}\csromanspacer{}นวมํ. \textbf{10.}\csromanspacer{}\textbf{ ทุติยปุคฺคลสุตฺต}}
\proseitem{60}{อฏฺฐิเม ภิกฺขเว ปุคฺคลา อาหุเนยฺยา ฯเปฯ อนุตฺตรํ ปุญฺญกฺเขตฺตํ โลกสฺส.\csromanspacer{} กตเม อฏฺฐ? โสตาปนฺโน โสตาปตฺติผลสจฺฉิ- กิริยาย ปฏิปนฺโน ฯเปฯ อรหา อรหตฺตาย ปฏิปนฺโน.\csromanspacer{} อิเม โข ภิกฺขเว อฏฺฐ ปุคฺคลา อาหุเนยฺยา ฯเปฯ อนุตฺตรํ ปุญฺญกฺเขตฺตํ โลกสฺสาติ.}

\csromangathameasure{จตฺตาโร จ ปฏิปนฺนา, จตฺตาโร จ ผเล ฐิตา. \\ เอส สํโฆ สมุกฺกฏฺโฐ, สตฺตานํ อฏฺฐ ปุคฺคลา. \\ ยชมานานํ มนุสฺสานํ, \\ ปุญฺญเปกฺขาน ปาณินํ. \\ กโรตํ โอปธิกํ ปุญฺญํ,}
\csromangathagroup{\csromangathastanza{จตฺตาโร จ ปฏิปนฺนา, จตฺตาโร จ ผเล ฐิตา. \\ เอส สํโฆ สมุกฺกฏฺโฐ, สตฺตานํ อฏฺฐ ปุคฺคลา.}\csromangathastanza{ยชมานานํ มนุสฺสานํ, \\ ปุญฺญเปกฺขาน ปาณินํ. \\ กโรตํ โอปธิกํ ปุญฺญํ,}}

\prose{เอตฺถ ทินฺนํ มหปฺผลนฺติ.\csromanspacer{}\csromanspacer{}ทสมํ.\csromanspacer{} โคตมีวคฺโค ปฐโม.\csromanspacer{} \textbf{ตสฺสุทฺทานํ}}

\csromangathameasure{โคตมี โอวาทํ สํขิตฺตํ, ทีฆชาณุ จ อุชฺชโย. \\ ภยา เทฺว อาหุเนยฺยา จ, เทฺว จ อฏฺฐ ปุคฺคลาติ.}
\csromangathagroup{\csromangathastanza{โคตมี โอวาทํ สํขิตฺตํ, ทีฆชาณุ จ อุชฺชโย.}\csromangathastanza{ภยา เทฺว อาหุเนยฺยา จ, เทฺว จ อฏฺฐ ปุคฺคลาติ.}}

\chapterhead{\textbf{(7) 2. ภูมิจาลวคฺค} \textbf{1. อิจฺฉาสุตฺต}}
\proseitem{61}{\csromansymbolfootnote{*}{อุปริ 143 ปิฏฺเฐปิ.}อฏฺฐิเม ภิกฺขเว ปุคฺคลา สนฺโต สํวิชฺชมานา โลกสฺมิํ.\csromanspacer{} กตเม อฏฺฐ? อิธ ภิกฺขเว ภิกฺขุโน ปวิวิตฺตสฺส วิหรโต นิรายตฺตวุตฺติโน อิจฺฉา อุปฺปชฺชติ ลาภาย, โส อุฏฺฐหติ ฆฏติ วายมติ ลาภาย, ตสฺส อุฏฺฐหโต ฆฏโต วายมโต ลาภาย ลาโภ นุปฺปชฺชติ.\csromanspacer{} โส เตน อลาเภน โสจติ กิลมติ ปริเทวติ อุรตฺตาฬิํ กนฺทติ สมฺโมหํ อาปชฺชติ.\csromanspacer{} อยํ วุจฺจติ ภิกฺขเว ภิกฺขุ อิจฺโฉ วิหรติ}

\vspace{\tipitakaparskip}
\csromancenter{(7) 2.\csromanspacer{} ภูมิจาลวคฺค ลาภาย, อุฏฺฐหติ ฆฏติ วายมติ ลาภาย, น จ ลาภี โสจี จ ปริเทวี จ จุโต จ สทฺธมฺมา.}
\prose{\csromanfolio{117}อิธ ปน ภิกฺขเว ภิกฺขุโน ปวิวิตฺตสฺส วิหรโต นิรายตฺตวุตฺติโน อิจฺฉา อุปฺปชฺชติ ลาภาย, โส อุฏฺฐหติ ฆฏติ วายมติ ลาภาย, ตสฺส อุฏฺฐหโต ฆฏโต วายมโต ลาภาย ลาโภ อุปฺปชฺชติ.\csromanspacer{} โส เตน ลาเภน มชฺชติ ปมชฺชติ ปมาทมาปชฺชติ.\csromanspacer{} อยํ วุจฺจติ ภิกฺขเว ภิกฺขุ อิจฺโฉ วิหรติ ลาภาย, อุฏฺฐหติ ฆฏติ วายมติ ลาภาย, ลาภี จ มที จ ปมาที จ จุโต จ สทฺธมฺมา.}

\prose{อิธ ปน ภิกฺขเว ภิกฺขุโน ปวิวิตฺตสฺส วิหรโต นิรายตฺตวุตฺติโน อิจฺฉา อุปฺปชฺชติ ลาภาย, โส น อุฏฺฐหติ น ฆฏติ น วายมติ ลาภาย, ตสฺส อนุฏฺฐหโต อฆฏโต อวายมโต ลาภาย ลาโภ นุปฺปชฺชติ.\csromanspacer{} โส เตน อลาเภน โสจติ กิลมติ ปริเทวติ อุรตฺตาฬิํ กนฺทติ สมฺโมหํ อาปชฺชติ.\csromanspacer{} อยํ วุจฺจติ ภิกฺขเว ภิกฺขุ อิจฺโฉ วิหรติ ลาภาย, น อุฏฺฐหติ น ฆฏติ น วายมติ ลาภาย, น จ ลาภี โสจี จ ปริเทวี จ จุโต จ สทฺธมฺมา.}

\prose{อิธ ปน ภิกฺขเว ภิกฺขุโน ปวิวิตฺตสฺส วิหรโต นิรายตฺตวุตฺติโน อิจฺฉา อุปฺปชฺชติ ลาภาย, โส น อุฏฺฐหติ น ฆฏติ น วายมติ ลาภาย, ตสฺส อนุฏฺฐหโต อฆฏโต อวายมโต ลาภาย ลาโภ อุปฺปชฺชติ.\csromanspacer{} โส เตน ลาเภน มชฺชติ ปมชฺชติ ปมาทมาปชฺชติ.\csromanspacer{} อยํ วุจฺจติ ภิกฺขเว ภิกฺขุ อิจฺโฉ วิหรติ ลาภาย, น อุฏฺฐหติ น ฆฏติ น วายมติ ลาภาย, ลาภี จ มที จ ปมาที จ จุโต จ สทฺธมฺมา.}

\prose{อิธ ปน ภิกฺขเว ภิกฺขุโน ปวิวิตฺตสฺส วิหรโต นิรายตฺตวุตฺติโน อิจฺฉา อุปฺปชฺชติ ลาภาย, โส อุฏฺฐหติ ฆฏติ วายมติ ลาภาย, ตสฺส อุฏฺฐหโต ฆฏโต วายมโต ลาภาย ลาโภ นุปฺปชฺชติ.\csromanspacer{} โส เตน อลาเภน น โสจติ น กิลมติ น ปริเทวติ น อุรตฺตาฬิํ กนฺทติ น สมฺโมหํ อาปชฺชติ.\csromanspacer{} อยํ วุจฺจติ ภิกฺขเว ภิกฺขุ อิจฺโฉ วิหรติ ลาภาย, อุฏฺฐหติ ฆฏติ วายมติ ลาภาย, น จ ลาภี น จ โสจี น จ ปริเทวี อจฺจุโต จ สทฺธมฺมา.}

\prose{อิธ ปน ภิกฺขเว ภิกฺขุโน ปวิวิตฺตสฺส วิหรโต นิรายตฺตวุตฺติโน อิจฺฉา อุปฺปชฺชติ ลาภาย, โส อุฏฺฐหติ ฆฏติ วายมติ ลาภาย, \csromanfolio{118}ตสฺส อุฏฺฐหโต ฆฏโต วายมโต ลาภาย ลาโภ อุปฺปชฺชติ.\csromanspacer{} โส เตน ลาเภน น มชฺชติ น ปมชฺชติ น ปมาทมาปชฺชติ.\csromanspacer{} อยํ วุจฺจติ ภิกฺขเว ภิกฺขุ อิจฺโฉ วิหรติ ลาภาย, อุฏฺฐหติ ฆฏติ วายมติ ลาภาย, ลาภี จ น จ มที น จ ปมาที อจฺจุโต จ สทฺธมฺมา.}

\prose{อิธ ปน ภิกฺขเว ภิกฺขุโน ปวิวิตฺตสฺส วิหรโต นิรายตฺตวุตฺติโน อิจฺฉา อุปฺปชฺชติ ลาภาย, โส น อุฏฺฐหติ น ฆฏติ น วายมติ ลาภาย, ตสฺส อนุฏฺฐหโต อฆฏโต อวายมโต ลภาย ลาโภ นุปฺปชฺชติ.\csromanspacer{} โส เตน อลาเภน น โสจติ น กิลมติ น ปริเทวติ น อุรตฺตาฬิํ กนฺทติ น สมฺโมหํ อาปชฺชติ.\csromanspacer{} อยํ วุจฺจติ ภิกฺขเว ภิกฺขุ อิจฺโฉ วิหรติ ลาภาย, น อุฏฺฐหติ น ฆฏติ น วายมติ ลาภาย, น จ ลาภี น จ โสจี น จ ปริเทวี อจฺจุโต จ สทฺธมฺมา.}

\prose{อิธ ปน ภิกฺขเว ภิกฺขุโน ปวิวิตฺตสฺส วิหรโต นิรายตฺตวุตฺติโน อิจฺฉา อุปฺปชฺชติ ลาภาย, โส น อุฏฺฐหติ น ฆฏติ น วายมติ ลาภาย, ตสฺส อนุฏฺฐหโต อฆฏโต อวายมโต ลาภาย ลาโภ อุปฺปชฺชติ.\csromanspacer{} โส เตน ลาเภน น มชฺชติ น ปมชฺชติ น ปมาทมาปชฺชติ.\csromanspacer{} อยํ วุจฺจติ ภิกฺขเว ภิกฺขุ อิจฺโฉ วิหรติ ลาภาย, น อุฏฺฐหติ น ฆฏติ น วายมติ ลาภาย, ลาภี จ น จ มที น จ ปมาที อจฺจุโต จ สทฺธมฺมา.\csromanspacer{} อิเม โข ภิกฺขเว อฏฺฐ ปุคฺคลา สนฺโต สํวิชฺชมานา โลกสฺมินฺติ.\csromanspacer{}\csromanspacer{}ปฏฺฐมํ.}

\csromanheader{2}{2. อลํสุตฺต}
\proseitem{62}{ฉหิ ภิกฺขเว ธมฺเมหิ สมนฺนาคโต ภิกฺขุ อลํ อตฺตโน อลํ ปเรสํ.\csromanspacer{} กตเมหิ ฉหิ? อิธ ภิกฺขเว ภิกฺขุ ขิปฺปนิสนฺติ จ โหติ กุสเลสุ ธมฺเมสุ.\csromanspacer{} สุตานญฺจ ธมฺมานํ ธารณชาติโก\footnote{ธารกชาติโก (สี, สฺยา, อิ) อุปริ145 ปิฏฺเฐปิ.} โหติ.\csromanspacer{} ธาตานญฺจ\footnote{ธตานญฺจ (สี, สฺยา, อิ)} ธมฺมานํ อตฺถูปปริกฺขิตา\footnote{อตฺถูปปริกฺขี (สี, สฺยา, อิ)} โหติ.\csromanspacer{} อตฺถมญฺญาย ธมฺมมญฺญาย ธมฺมานุธมฺมปฺปฏิปนฺโน จ โหติ.\csromanspacer{} กลฺยาณวาโจ จ โหติ กลฺยาณวากฺกรโณ โปริยา วาจาย สมนฺนาคโต วิสฺสฏฺฐาย อเนลคฬาย อตฺถสฺส วิญฺญาปนิยา.\csromanspacer{} สนฺทสฺสโก จ โหติ สมาทปโก\footnote{สมาทาปโก (?)}}

\vspace{\tipitakaparskip}
\csromancenter{(7) 2.\csromanspacer{} ภูมิจาลวคฺค สมุตฺเตชโก สมฺปหํสโก สพฺรหฺมจารีนํ.\csromanspacer{} อิเมหิ โข ภิกฺขเว ฉหิ ธมฺเมหิ สมนฺนาคโต ภิกฺขุ อลํ อตฺตโน อลํ ปเรสํ. (1)}
\prose{\csromanfolio{119}ปญฺจหิ ภิกฺขเว ธมฺเมหิ สมนฺนาคโต ภิกฺขุ อลํ อตฺตโน อลํ ปเรสํ.\csromanspacer{} กตเมหิ ปญฺจหิ? อิธ ภิกฺขเว ภิกฺขุ น เหว โข ขิปฺปนิสนฺติ จ โหติ กุสเลสุ ธมฺเมสุ.\csromanspacer{} สุตานญฺจ ธมฺมานํ ธารณชาติโก โหติ.\csromanspacer{} ธาตานญฺจ ธมฺมานํ อตฺถูปปริกฺขิตา โหติ.\csromanspacer{} อตฺถมญฺญาย ธมฺมมญฺญาย ธมฺมานุธมฺมปฺปฏิปนฺโน จ โหติ.\csromanspacer{} กลฺยาณวาโจ จ โหติ ฯเปฯ อตฺถสฺส วิญฺญาปนิยา.\csromanspacer{} สนฺทสฺสโก จ โหติ สมาทปโก สมุตฺเตชโก สมฺปหํสโก สพฺรหฺมจารีนํ.\csromanspacer{} อิเมหิ โข ภิกฺขเว ปญฺจหิ ธมฺเมหิ สมนฺนาคโต ภิกฺขุ อลํ อตฺตโน อลํ ปเรสํ. (2)}

\prose{จตูหิ ภิกฺขเว ธมฺเมหิ สมนฺนาคโต ภิกฺขุ อลํ อตฺตโน นาลํ ปเรสํ.\csromanspacer{} กตเมหิ จตูหิ? อิธ ภิกฺขเว ภิกฺขุ ขิปฺปนิสนฺติ จ โหติ กุสเลสุ ธมฺเมสุ.\csromanspacer{} สุตานญฺจ ธมฺมานํ ธารณชาติโก โหติ.\csromanspacer{} ธาตานญฺจ ธมฺมานํ อตฺถูปปริกฺขิตา โหติ.\csromanspacer{} อตฺถมญฺญาย ธมฺมมญฺญาย ธมฺมานุธมฺมปฺปฏิปนฺโน จ โหติ.\csromanspacer{} โน จ กลฺยาณวาโจ โหติ กลฺยาณวากฺกรโณ โปริยา วาจาย สมนฺนาคโต วิสฺสฏฺฐาย อเนลคฬาย อตฺถสฺส วิญฺญาปนิยา.\csromanspacer{} โน จ สนฺทสฺสโก โหติ สมาทปโก สมุตฺเตชโก สมฺปหํสโก สพฺรหฺมจารีนํ.\csromanspacer{} อิเมหิ โข ภิกฺขเว จตูหิ ธมฺเมหิ สมนฺนาคโต ภิกฺขุ อลํ อตฺตโน นาลํ ปเรสํ. (3)}

\prose{จตูหิ ภิกฺขเว ธมฺเมหิ สมนฺนาคโต ภิกฺขุ อลํ ปเรสํ นาลํ อตฺตโน.\csromanspacer{} กตเมหิ จตูหิ? อิธ ภิกฺขเว ภิกฺขุ ขิปฺปนิสนฺติ จ โหติ กุสเลสุ ธมฺเมสุ.\csromanspacer{} สุตานญฺจ ธมฺมานํ ธารณชาติโก โหติ.\csromanspacer{} โน จ ธาตานํ ธมฺมานํ อตฺถูปปริกฺขิตา โหติ.\csromanspacer{} น จ อตฺถมญฺญาย ธมฺมมญฺญาย ธมฺมานุธมฺมปฺปฏิปนฺโน โหติ.\csromanspacer{} กลฺยาณวาโจ จ โหติ กลฺยาณวากฺกรโณ ฯเปฯ อตฺถสฺส วิญฺญาปนิยา.\csromanspacer{} สนฺทสฺสโก จ โหติ ฯเปฯ สพฺรหฺมจารีนํ.\csromanspacer{} อิเมหิ โข ภิกฺขเว จตูหิ ธมฺเมหิ สมนฺนาคโต ภิกฺขุ อลํ ปเรสํ นาลํ อตฺตโน. (4)}

\prose{ตีหิ ภิกฺขเว ธมฺเมหิ สมนฺนาคโต ภิกฺขุ อลํ อตฺตโน นาลํ ปเรสํ.\csromanspacer{} กตเมหิ ตีหิ? อิธ ภิกฺขเว ภิกฺขุ น เหว โข ขิปฺปนิสนฺติ จ โหติ กุสเลสุ ธมฺเมสุ.\csromanspacer{} สุตานญฺจ ธมฺมานํ ธารณชาติโก โหติ. \csromanfolio{120}ธาตานญฺจ ธมฺมานํ อตฺถูปปริกฺขิตา โหติ.\csromanspacer{} อตฺถมญฺญาย ธมฺมมญฺญาย ธมฺมานุธมฺมปฺปฏิปนฺโน จ โหติ.\csromanspacer{} โน จ กลฺยาณวาโจ โหติ กลฺยาณวากฺกรโณ โปริยา วาจาย สมนฺนาคโต วิสฺสฏฺฐาย อเนลคฬาย อตฺถสฺส วิญฺญาปนิยา.\csromanspacer{} โน จ สนฺทสฺสโก โหติ สมาทปโก สมุตฺเตชโก สมฺปหํสโก สพฺรหฺมจารีนํ.\csromanspacer{} อิเมหิ โข ภิกฺขเว ตีหิ ธมฺเมหิ สมนฺนาคโต ภิกฺขุ อลํ อตฺตโน นาลํ ปเรสํ. (5)}

\prose{ตีหิ ภิกฺขเว ธมฺเมหิ สมนฺนาคโต ภิกฺขุ อลํ ปเรสํ นาลํ อตฺตโน.\csromanspacer{} กตเมหิ ตีหิ? อิธ ภิกฺขเว ภิกฺขุ น เหว โข ขิปฺปนิสนฺติ จ โหติ กุสเลสุ ธมฺเมสุ.\csromanspacer{} สุตานญฺจ ธมฺมานํ ธารณชาติโก โหติ.\csromanspacer{} โน จ ธาตานํ ธมฺมานํ อตฺถูปปริกฺขิตา โหติ.\csromanspacer{} โนจ อตฺถมญฺญาย ธมฺมมญฺญาย ธมฺมานุธมฺมปฺปฏิปนฺโน โหติ.\csromanspacer{} กลฺยาณวาโจ จ โหติ ฯเปฯ อตฺถสฺส วิญฺญาปนิยา.\csromanspacer{} สนฺทสฺสโก จ โหติ สมาทปโก สมุตฺเตชโก สมฺปหํสโก สพฺรหฺมจารีนํ.\csromanspacer{} อิเมหิ โข ภิกฺขเว ตีหิ ธมฺเมหิ สมนฺนาคโต ภิกฺขุ อลํ ปเรสํ นาลํ อตฺตโน. (6)}

\prose{ทฺวีหิ ภิกฺขเว ธมฺเมหิ สมนฺนาคโต ภิกฺขุ อลํ อตฺตโน นาลํ ปเรสํ.\csromanspacer{} กตเมหิ ทฺวีหิ? อิธ ภิกฺขเว ภิกฺขุ น เหว โข ขิปฺปนิสนฺติ จ โหติ กุสเลสุ ธมฺเมสุ.\csromanspacer{} โน จ สุตานํ ธมฺมานํ ธารณชาติโก โหติ.\csromanspacer{} ธาตานญฺจ ธมฺมานํ อตฺถูปปริกฺขิตา โหติ.\csromanspacer{} อตฺถมญฺญาย ธมฺมมญฺญาย ธมฺมานุธมฺมปฺปฏิปนฺโน จ โหติ.\csromanspacer{} โน จ กลฺยาณวาโจ โหติ ฯเปฯ อตฺถสฺส วิญฺญาปนิยา.\csromanspacer{} โน จ สนฺทสฺสโก โหติ ฯเปฯ สพฺรหฺมจารีนํ.\csromanspacer{} อิเมหิ โข ภิกฺขเว ทฺวีหิ ธมฺเมหิ สมนฺนาคโต ภิกฺขุ อลํ อตฺตโน นาลํ ปเรสํ. (7)}

\prose{ทฺวีหิ ภิกฺขเว ธมฺเมหิ สมนฺนาคโต ภิกฺขุ อลํ ปเรสํ นาลํ อตฺตโน.\csromanspacer{} กตเมหิ ทฺวีหิ? อิธ ภิกฺขเว ภิกฺขุ น เหว โข ขิปฺปนิสนฺติ จ โหติ กุสเลสุ ธมฺเมสุ.\csromanspacer{} โน จ สุตานํ ธมฺมานํ ธารณชาติโก โหติ.\csromanspacer{} โน จ ธาตานํ ธมฺมานํ อตฺถูปปริกฺขิตา โหติ.\csromanspacer{} โน จ อตฺถมญฺญาย ธมฺมมญฺญาย ธมฺมานุธมฺมปฺปฏิปนฺโน โหติ.\csromanspacer{} กลฺยาณวาโจ จ โหติ กลฺยาณวากฺกรโณ โปริยา วาจาย สมนฺนาคโต วิสฺสฏฺฐาย อเนลคฬาย อตฺถสฺส วิญฺญาปนิยา.\csromanspacer{} สนฺทสฺสโก จ โหติ สมาทปโก สมุตฺเตชโก สมฺปหํสโก สพฺรหฺมจารีนํ.\csromanspacer{} อิเมหิ โข}

\vspace{\tipitakaparskip}
\csromancenter{(7) 2.\csromanspacer{} ภูมิจาลวคฺค ภิกฺขเว ทฺวีหิ ธมฺเมหิ สมนฺนาคโต ภิกฺขุ อลํ ปเรสํ นาลํ อตฺตโนติ. (8).\csromanspacer{}\csromanspacer{}ทุติยํ.}
\csromanheader{2}{3. สํขิตฺตสุตฺต}
\proseitem{63}{\csromanfolio{121}อถ โข อญฺญตโร ภิกฺขุ เยน ภควา เตนุปสงฺกมิ ฯเปฯ เอกมนฺตํ นิสินฺโน โข โส ภิกฺขุ ภควนฺตํ เอตทโวจ “สาธุ เม ภนฺเต ภควา สํขิตฺเตน ธมฺมํ เทเสตุ, ยมหํ ภควโต ธมฺมํ สุตฺวา เอโก วูปกฏฺโฐ อปฺปมตฺโต อาตาปี ปหิตตฺโต วิหเรยฺยนฺ”ติ.\csromanspacer{} เอวเมว ปนิเธกจฺเจ โมฆปุริสา มมญฺเญว อชฺเฌสนฺติ, ธมฺเม จ ภาสิเต มมญฺเญว อนุพนฺธิตพฺพํ มญฺญนฺตีติ.\csromanspacer{} เทเสตุ เม ภนฺเต ภควา สํขิตฺเตน ธมฺมํ, เทเสตุ สุคโต สํขิตฺเตน ธมฺมํ, อปฺเปว นามาหํ ภควโต ภาสิตสฺส อตฺถํ อาชาเนยฺยํ, อปฺเปว นามาหํ ภควโต ภาสิตสฺส ทายาโท อสฺสนฺติ.\csromanspacer{} ตสฺมาติห เต ภิกฺขุ เอวํ สิกฺขิตพฺพํ “อชฺฌตฺตํ เม จิตฺตํ ฐิตํ ภวิสฺสติ สุสณฺฐิตํ, น จ อุปฺปนฺนา ปาปกา อกุสลา ธมฺมา จิตฺตํ ปริยาทาย ฐสฺสนฺตี”ติ.\csromanspacer{} เอวํ หิ เต ภิกฺขุ สิกฺขิตพฺพํ.}

\prose{ยโต โข เต ภิกฺขุ อชฺฌตฺตํ จิตฺตํ ฐิตํ โหติ สุสณฺฐิตํ, น จ อุปฺปนฺนา ปาปกา อกุสลา ธมฺมา จิตฺตํ ปริยาทาย ติฏฺฐนฺติ, ตโต เต ภิกฺขุ เอวํ สิกฺขิตพฺพํ “เมตฺตา เม เจโตวิมุตฺติ ภาวิตา ภวิสฺสติ พหุลีกตา ยานีกตา วตฺถุกตา อนุฏฺฐิตา ปริจิตา สุสมารทฺธา”ติ.\csromanspacer{} เอวํ หิ เต ภิกฺขุ สิกฺขิตพฺพํ.}

\prose{ยโต โข เต ภิกฺขุ อยํ สมาธิ เอวํ ภาวิโต โหติ พหุลีกโต, ตโต ตฺวํ ภิกฺขุ อิมํ สมาธิํ สวิตกฺกมฺปิ สวิจารํ\footnote{สวิตกฺกสวิจารมฺปิ (ก)} ภาเวยฺยาสิ, อวิตกฺกมฺปิ วิจารมตฺตํ\footnote{อวิตกฺกวิจารมตฺตมฺปิ (ก) วิสุทฺธิ 1. 315 ปิฏฺเฐ ปสฺสิตพฺพํ.} ภาเวยฺยาสิ, อวิตกฺกมฺปิ อวิจารํ\footnote{อวิตกฺก-อวิจารมฺปิ (ก)} ภาเวยฺยาสิ, สปฺปีติกมฺปิ ภาเวยฺยาสิ, นิปฺปีติกมฺปิ ภาเวยฺยาสิ, สาตสหคตมฺปิ ภาเวยฺยาสิ, อุเปกฺขาสหคตมฺปิ ภาเวยฺยาสิ.}

\prose{ยโต โข เต ภิกฺขุ อยํ สมาธิ เอวํ ภาวิโต โหติ สุภาวิโต, ตโต เต ภิกฺขุ เอวํ สิกฺขิตพฺพํ “กรุณา เม เจโตวิมุตฺติ.\csromanspacer{} มุทิตา เม เจโตวิมุตฺติ.\csromanspacer{} อุเปกฺขา เม เจโตวิมุตฺติ ภาวิตา \csromanfolio{122}ภวิสฺสติ พหุลีกตา ยานีกตา วตฺถุกตา อนุฏฺฐิตา ปริจิตา สุสมารทฺธา”ติ.\csromanspacer{} เอวํ หิ เต ภิกฺขุ สิกฺขิตพฺพํ.}

\prose{ยโต โข เต ภิกฺขุ อยํ สมาธิ เอวํ ภาวิโต โหติ สุภาวิโต, ตโต ตฺวํ ภิกฺขุ อิมํ สมาธิํ สวิตกฺกสวิจารมฺปิ ภาเวยฺยาสิ, อวิตกฺกวิจารมตฺตมฺปิ ภาเวยฺยาสิ, อวิตกฺก-อวิจารมฺปิ ภาเวยฺยาสิ, สปฺปีติกมฺปิ ภาเวยฺยาสิ, นิปฺปีติกมฺปิ ภาเวยฺยาสิ, สาตสหคตมฺปิ ภาเวยฺยาสิ, อุเปกฺขาสหคตมฺปิ ภาเวยฺยาสิ.}

\prose{ยโต โข เต ภิกฺขุ อยํ สมาธิ เอวํ ภาวิโต โหติ สุภาวิโต, ตโต เต ภิกฺขุ เอวํ สิกฺขิตพฺพํ “กาเย กายานุปสฺสี วิหริสฺสามิ อาตาปี สมฺปชาโน สติมา วิเนยฺย โลเก อภิชฺฌาโทมนสฺสนฺ”ติ.\csromanspacer{} เอวํ หิ เต ภิกฺขุ สิกฺขิตพฺพํ.}

\prose{ยโต โข เต ภิกฺขุ อยํ สมาธิ เอวํ ภาวิโต โหติ พหุลีกโต, ตโต ตฺวํ ภิกฺขุ อิมํ สมาธิํ สวิตกฺกสวิจารมฺปิ ภาเวยฺยาสิ, อวิตกฺกวิจารมตฺตมฺปิ ภาเวยฺยาสิ, อวิตกฺก-อวิจารมฺปิ ภาเวยฺยาสิ, สปฺปีติกมฺปิ ภาเวยฺยาสิ, นิปฺปีติกมฺปิ ภาเวยฺยาสิ, สาตสหคตมฺปิ ภาเวยฺยาสิ, อุเปกฺขาสหคตมฺปิ ภาเวยฺยาสิ.}

\prose{ยโต โข เต ภิกฺขุ อยํ สมาธิ เอวํ ภาวิโต โหติ สุภาวิโต, ตโต เต ภิกฺขุ เอวํ สิกฺขิตพฺพํ “เวทนาสุ.\csromanspacer{} จิตฺเต.\csromanspacer{} ธมฺเมสุ ธมฺมานุปสฺสี วิหริสฺสามิ อาตาปี สมฺปชาโน สติมา วิเนยฺย โลเก อภิชฺฌาโทมนสฺสนฺ”ติ.\csromanspacer{} เอวํ หิ เต ภิกฺขุ สิกฺขิตพฺพํ.}

\prose{ยโต โข เต ภิกฺขุ อยํ สมาธิ เอวํ ภาวิโต โหติ พหุลีกโต, ตโต ตฺวํ ภิกฺขุ อิมํ สมาธิํ สวิตกฺกสวิจารมฺปิ ภาเวยฺยาสิ, อวิตกฺกวิจารมตฺตมฺปิ ภาเวยฺยาสิ, อวิตกฺก-อวิจารมฺปิ ภาเวยฺยาสิ, สปฺปีติกมฺปิ ภาเวยฺยาสิ, นิปฺปีติกมฺปิ ภาเวยฺยาสิ, สาตสหคตมฺปิ ภาเวยฺยาสิ, อุเปกฺขาสหคตมฺปิ ภาเวยฺยาสิ.}

\prose{ยโต โข เต ภิกฺขุ อยํ สมาธิ เอวํ ภาวิโต โหติ สุภาวิโต, ตโต ตฺวํ ภิกฺขุ เยน เยเนว คคฺฆสิ, ผาสุํเยว คคฺฆสิ.\csromanspacer{} ยตฺถ ยตฺถ ฐสฺสสิ, ผาสุํเยว ฐสฺสสิ.\csromanspacer{} ยตฺถ ยตฺถ นิสีทิสฺสสิ, ผาสุํเยว นิสีทิสฺสสิ.\csromanspacer{} ยตฺถ ยตฺถ เสยฺยํ กปฺเปสฺสสิ, ผาสุํเยว เสยฺยํ กปฺเปสฺสสีติ.}

\chapterheadpageread{(7) 2. ภูมิจาลวคฺค}
\prose{\csromanfolio{123}อถ โข โส ภิกฺขุ ภควตา อิมินา โอวาเทน โอวทิโต อุฏฺฐายาสนา ภควนฺตํ อภิวาเทตฺวา ปทกฺขิณํ กตฺวา ปกฺกามิ.\csromanspacer{} อถ โข โส ภิกฺขุ เอโก วูปกฏฺโฐ อปฺปมตฺโต อาตาปี ปหิตตฺโต วิหรนฺโต นจิรสฺเสว ยสฺสตฺถาย กุลปุตฺตา สมฺมเทว อคารสฺมา อนคาริยํ ปพฺพชนฺติ, ตทนุตฺตรํ พฺรหฺมจริยปริโยสานํ ทิฏฺเฐว ธมฺเม สยํ อภิญฺญา สจฺฉิกตฺวา อุปสมฺปชฺช วิหาสิ, “ขีณา ชาติ, วุสิตํ พฺรหฺมจริยํ, กตํ กรณียํ, นาปรํ อิตฺถตฺตายา”ติ อพฺภญฺญาสิ.\csromanspacer{} อญฺญตโร จ ปน โส ภิกฺขุ อรหตํ อโหสีติ.\csromanspacer{}\csromanspacer{}ตติยํ.}

\csromanheader{2}{4. คยาสีสสุตฺต}
\proseitem{64}{เอกํ สมยํ ภควา คยายํ วิหรติ คยาสีเส.\csromanspacer{} ตตฺร โข ภควา ภิกฺขู อามนฺเตสิ ฯเปฯ.\csromanspacer{} ปุพฺพาหํ ภิกฺขเว สมฺโพธา อนภิสมฺพุทฺโธ โพธิสตฺโตว สมาโน โอภาสญฺเญว โข สญฺชานามิ, โน จ รูปานิ ปสฺสามิ. (1)}

\prose{ตสฺส มยฺหํ ภิกฺขเว เอตทโหสิ “สเจ โข อหํ โอภาสญฺเจว สญฺชาเนยฺยํ, รูปานิ จ ปสฺเสยฺยํ.\csromanspacer{} เอวํ เม อิทํ ญาณทสฺสนํ ปริสุทฺธตรํ อสฺสา”ติ.}

\prose{โส โข อหํ ภิกฺขเว อปเรน สมเยน อปฺปมตฺโต อาตาปี ปหิตตฺโต วิหรนฺโต โอภาสญฺเจว สญฺชานามิ, รูปานิ จ ปสฺสามิ.\csromanspacer{} โน จ โข ตาหิ เทวตาหิ สทฺธิํ สนฺติฏฺฐามิ สลฺลปามิ สากจฺฉํ สมาปชฺชามิ. (2)}

\prose{ตสฺส มยฺหํ ภิกฺขเว เอตทโหสิ “สเจ โข อหํ โอภาสญฺเจว สญฺชาเนยฺยํ, รูปานิ จ ปสฺเสยฺยํ.\csromanspacer{} ตาหิ จ เทวตาหิ สทฺธิํ สนฺติฏฺเฐยฺยํ สลฺลเปยฺยํ สากจฺฉํ สมาปชฺเชยฺยํ.\csromanspacer{} เอวํ เม อิทํ ญาณทสฺสนํ ปริสุทฺธตรํ อสฺสา”ติ.}

\prose{โส โข อหํ ภิกฺขเว อปเรน สมเยน อปฺปมตฺโต อาตาปี ปหิตตฺโต วิหรนฺโต โอภาสญฺเจว สญฺชานามิ, รูปานิ จ ปสฺสามิ.\csromanspacer{} ตาหิ จ เทวตาหิ สทฺธิํ สนฺติฏฺฐามิ สลฺลปามิ สากจฺฉํ สมาปชฺชามิ.\csromanspacer{} โน จ โข ตา เทวตา ชานามิ “อิมา เทวตา อมุกมฺหา วา อมุกมฺหา วา เทวนิกายา”ติ. (3)}

\prose{\csromanfolio{124}ตสฺส มยฺหํ ภิกฺขเว เอตทโหสิ “สเจ โข อหํ โอภาสญฺเจว สญฺชาเนยฺยํ, รูปานิ จ ปสฺเสยฺยํ.\csromanspacer{} ตาหิ จ เทวตาหิ สทฺธิํ สนฺติฏฺเฐยฺยํ สลฺลเปยฺยํ สากจฺฉํ สมาปชฺเชยฺยํ, ตา จ เทวตา ชาเนยฺยํ ‘อิมา เทวตา อมุกมฺหา วา อมุกมฺหา วา เทวนิกายา’ติ.\csromanspacer{} เอวํ เม อิทํ ญาณทสฺสนํ ปริสุทฺธตรํ อสฺสา”ติ.}

\prose{โส โข อหํ ภิกฺขเว อปเรน สมเยน อปฺปมตฺโต อาตาปี ปหิตตฺโต วิหรนฺโต โอภาสญฺเจว สญฺชานามิ, รูปานิ จ ปสฺสามิ.\csromanspacer{} ตาหิ จ เทวตาหิ สทฺธิํ สนฺติฏฺฐามิ สลฺลปามิ สากจฺฉํ สมาปชฺชามิ.\csromanspacer{} ตา จ เทวตา ชานามิ “อิมา เทวตา อมุกมฺหา วา อมุกมฺหา วา เทวนิกายา”ติ.\csromanspacer{} โน จ โข ตา เทวตา ชานามิ “อิมา เทวตา อิมสฺส กมฺมสฺส วิปาเกน อิโต จุตา ตตฺถ อุปปนฺนา”ติ ฯเปฯ.\csromanspacer{} ตา จ เทวตา ชานามิ “อิมา เทวตา อิมสฺส กมฺมสฺส วิปาเกน อิโต จุตา ตตฺถ อุปปนฺนา”ติ.\csromanspacer{} โน จ โข ตา เทวตา ชานามิ “อิมา เทวตา อิมสฺส กมฺมสฺส วิปาเกน เอวมาหารา เอวํสุขทุกฺขปฺปฏิสํเวทินิโย”ติ ฯเปฯ.\csromanspacer{} ตา จ เทวตา ชานามิ “อิมา เทวตา อิมสฺส กมฺมสฺส วิปาเกน เอวมาหารา เอวํสุขทุกฺขปฺปฏิสํเวทินิโย”ติ.\csromanspacer{} โน จ โข ตา เทวตา ชานามิ “อิมา เทวตา เอวํทีฆายุกา เอวํจิรฏฺฐิติกา”ติ ฯเปฯ.\csromanspacer{} ตา จ เทวตา ชานามิ “อิมา เทวตา เอวํทีฆายุกา เอวํจิรฏฺฐิติกา”ติ.\csromanspacer{} โน จ โข ตา เทวตา ชานามิ “ยทิ วา เม อิมาหิ เทวตาหิ สทฺธิํ สนฺนิวุตฺถปุพฺพํ, ยทิ วา น สนฺนิวุตฺถปุพฺพนฺ”ติ. (4-7)}

\prose{ตสฺส มยฺหํ ภิกฺขเว เอตทโหสิ “สเจ โข อหํ โอภาสญฺเจว สญฺชาเนยฺยํ, รูปานิ จ ปสฺเสยฺยํ.\csromanspacer{} ตาหิ จ เทวตาหิ สทฺธิํ สนฺติฏฺเฐยฺยํ สลฺลเปยฺยํ สากจฺฉํ สมาปชฺเชยฺยํ.\csromanspacer{} ตา จ เทวตา ชาเนยฺยํ ‘อิมา เทวตา อมุกมฺหา วา อมุกมฺหา วา เทวนิกายา’ติ.\csromanspacer{} ตา จ เทวตา ชาเนยฺยํ ‘อิมา เทวตา อิมสฺส กมฺมสฺส วิปาเกน อิโต จุตา ตตฺถ อุปปนฺนา’ติ.\csromanspacer{} ตา จ เทวตา ชาเนยฺยํ ‘อิมา เทวตา เอวมาหารา เอวํสุขทุกฺขปฺปฏิสํเวทินิโย’ติ.\csromanspacer{} ตา จ เทวตา ชาเนยฺยํ ‘อิมา เทวตา เอวํทีฆายุกา เอวํจิรฏฺฐิติกา’ติ.\csromanspacer{} ตา จ เทวตา ชาเนยฺยํ ‘ยทิ วา เม อิมาหิ เทวตาหิ สทฺธิํ สนฺนิวุตฺถปุพฺพํ, ยทิ วา น สนฺนิวุตฺถปุพฺพนฺติ.\csromanspacer{} เอวํ เม อิทํ ญาณทสฺสนํ ปริสุทฺธตรํ อสฺสา”ติ.}

\chapterheadpageread{(7) 2. ภูมิจาลวคฺค}
\prose{\csromanfolio{125}โส โข อหํ ภิกฺขเว อปเรน สมเยน อปฺปมตฺโต อาตาปี ปหิตตฺโต วิหรนฺโต โอภาสญฺเจว สญฺชานามิ, รูปานิ จ ปสฺสามิ.\csromanspacer{} ตาหิ จ เทวตาหิ สทฺธิํ สนฺติฏฺฐามิ สลฺลปามิ สากจฺฉํ สมาปชฺชามิ.\csromanspacer{} ตา จ เทวตา ชานามิ “อิมา เทวตา อมุกมฺหา วา อมุกมฺหา วา เทวนิกายา”ติ.\csromanspacer{} ตา จ เทวตา ชานามิ “อิมา เทวตา อิมสฺส กมฺมสฺส วิปาเกน อิโต จุตา ตตฺถ อุปปนฺนา”ติ.\csromanspacer{} ตา จ เทวตา ชานามิ “อิมา เทวตา เอวมาหารา เอวํสุขทุกฺขปฺ- ปฏิสํเวทินิโย”ติ.\csromanspacer{} ตา จ เทวตา ชานามิ “อิมา เทวตา เอวํทีฆายุกา เอวํจิรฏฺฐิติกา”ติ.\csromanspacer{} ตา จ เทวตา ชานามิ “ยทิ วา เม เทวตาหิ สทฺธิํ สนฺนิวุตฺถปุพฺพํ, ยทิ วา น สนฺนิวุตฺถปุพฺพนฺ”ติ. (8)}

\prose{ยาวกีวญฺจ เม ภิกฺขเว เอวํ อฏฺฐปริวฏฺฏํ อธิเทวญาณทสฺสนํ น สุวิสุทฺธํ อโหสิ, เนว ตาวาหํ ภิกฺขเว สเทวเก โลเก สมารเก สพฺรหฺมเก สสฺสมณพฺราหฺมณิยา ปชาย สเทวมนุสฺสาย อนุตฺตรํ สมฺมาสมฺโพธิํ อภิสมฺพุทฺโธติ\footnote{อภิสมฺพุทฺโธ (สี, สฺยา, อิ)} ปจฺจญฺญาสิํ.\csromanspacer{} ยโต จ โข เม ภิกฺขเว เอวํ อฏฺฐปริวฏฺฏํ อธิเทวญาณทสฺสนํ สุวิสุทฺธํ อโหสิ.\csromanspacer{} อถาหํ ภิกฺขเว สเทวเก โลเก สมารเก สพฺรหฺมเก สสฺสมณพฺราหฺมณิยา ปชาย สเทวมนุสฺสาย อนุตฺตรํ สมฺมาสมฺโพธิํ อภิสมฺพุทฺโธติ ปจฺจญฺญาสิํ.\csromanspacer{} ญาณญฺจ ปน เม ทสฺสนํ อุทปาทิ, อกุปฺปา เม เจโตวิมุตฺติ\footnote{วิมุตฺติ (ก-สี, ก)}, อยมนฺติมา ชาติ, นตฺถิ ทานิ ปุนพฺภโวติ.\csromanspacer{}\csromanspacer{}จตุตฺถํ.}

\csromanheader{2}{5. อภิภายตนสุตฺต}
\proseitem{65}{\csromansymbolfootnote{*}{ที 3. 215, 250 ปิฏฺเฐสุปิ; อุปริ 304 ปิฏฺเฐปิ.}อฏฺฐิมานิ ภิกฺขเว อภิภายตนานิ.\csromanspacer{} กตมานิ อฏฺฐ? อชฺฌตฺตํ รูปสญฺญี เอโก พหิทฺธา รูปานิ ปสฺสติ ปริตฺตานิ สุวณฺณทุพฺพณฺณานิ.\csromanspacer{} ตานิ อภิภุยฺย “ชานามิ ปสฺสามี”ติ เอวํสญฺญี โหติ, อิทํ ปฐมํ อภิภายตนํ.}

\prose{อชฺฌตฺตํ รูปสญฺญี เอโก พหิทฺธา รูปานิ ปสฺสติ อปฺปมาณานิ สุวณฺณทุพฺพณฺณานิ.\csromanspacer{} ตานิ อภิภุยฺย “ชานามิ ปสฺสามี”ติ เอวํสญฺญี โหติ, อิทํ ทุติยํ อภิภายตนํ.}

\prose{\csromanfolio{126}อชฺฌตฺตํ อรูปสญฺญี เอโก พหิทฺธา รูปานิ ปสฺสติ ปริตฺตานิ สุวณฺณทุพฺพณฺณานิ.\csromanspacer{} ตานิ อภิภุยฺย “ชานามิ ปสฺสามี”ติ เอวํสญฺญี โหติ, อิทํ ตติยํ อภิภายตนํ.}

\prose{อชฺฌตฺตํ อรูปสญฺญี เอโก พหิทฺธา รูปานิ ปสฺสติ อปฺปมาณานิ สุวณฺณทุพฺพณฺณานิ.\csromanspacer{} ตานิ อภิภุยฺย “ชานามิ ปสฺสามี”ติ เอวํสญฺญี โหติ, อิทํ จตุตฺถํ อภิภายตนํ.}

\prose{อชฺฌตฺตํ อรูปสญฺญี เอโก พหิทฺธา รูปานิ ปสฺสติ นีลานิ นีลวณฺณานิ นีลนิทสฺสนานิ นีลนิภาสานิ.\csromanspacer{} ตานิ อภิภุยฺย “ชานามิ ปสฺสามี”ติ เอวํสญฺญี โหติ, อิทํ ปญฺจมํ อภิภายตนํ.}

\prose{อชฺฌตฺตํ อรูปสญฺญี เอโก พหิทฺธา รูปานิ ปสฺสติ ปีตานิ ปีตวณฺณานิ ปีตนิทสฺสนานิ ปีตนิภาสานิ.\csromanspacer{} ตานิ อภิภุยฺย “ชานามิ ปสฺสามี”ติ เอวํสญฺญี โหติ, อิทํ ฉฏฺฐํ อภิภายตนํ.}

\prose{อชฺฌตฺตํ อรูปสญฺญี เอโก พหิทฺธา รูปานิ ปสฺสติ โลหิตกานิ โลหิตกวณฺณานิ โลหิตกนิทสฺสนานิ โลหิตกนิภาสานิ.\csromanspacer{} ตานิ อภิภุยฺย “ชานามิ ปสฺสามี”ติ เอวํสญฺญี โหติ, อิทํ สตฺตมํ อภิภายตนํ.}

\prose{อชฺฌตฺตํ อรูปสญฺญี เอโก พหิทฺธา รูปานิ ปสฺสติ โอทาตานิ โอทาตวณฺณานิ โอทาตนิทสฺสนานิ โอทาตนิภาสานิ.\csromanspacer{} ตานิ อภิภุยฺย “ชานามิ ปสฺสามี”ติ เอวํสญฺญี โหติ, อิทํ อฏฺฐมํ อภิภายตนํ.\csromanspacer{} อิมานิ โข ภิกฺขเว อฏฺฐ อภิภายตนานีติ.\csromanspacer{}\csromanspacer{}ปญฺจมํ.}

\csromanheader{2}{6. วิโมกฺขสุตฺต}
\proseitem{66}{อฏฺฐิเม ภิกฺขเว วิโมกฺขา.\csromanspacer{} กตเม อฏฺฐ? รูปี รูปานิ ปสฺสติ, อยํ ปฐโม วิโมกฺโข.}

\prose{อชฺฌตฺตํ อรูปสญฺญี พหิทฺธา\footnote{อรูปสญฺญี เอโก พหิทฺธา (สฺยา, อิ, ก) ที 2. 60, 94; ที 3. 217, 251; อุปริ 161 ปิฏฺเฐสุ; ม 2. 205 ปิฏฺเฐ จ ปสฺสิตพฺพํ.} รูปานิ ปสฺสติ, อยํ ทุติโย วิโมกฺโข.}

\prose{“สุภนฺ”เตว อธิมุตฺโต โหติ, อยํ ตติโย วิโมกฺโข.}

\chapterheadpageread{(7) 2. ภูมิจาลวคฺค}
\prose{\csromanfolio{127}สพฺพโส รูปสญฺญานํ สมติกฺกมา ปฏิฆสญฺญานํ อตฺถงฺคมา นานตฺตสญฺญานํ อมนสิการา “อนนฺโต อากาโส”ติ อากาสานญฺจายตนํ อุปสมฺปชฺช วิหรติ, อยํ จตุตฺโถ วิโมกฺโข.}

\prose{สพฺพโส อากาสานญฺจายตนํ สมติกฺกมฺม “อนนฺตํ วิญฺญาณนฺ”ติ วิญฺญาณญฺจายตนํ อุปสมฺปชฺช วิหรติ, อยํ ปญฺจโม วิโมกฺโข.}

\prose{สพฺพโส วิญฺญาณญฺจายตนํ สมติกฺกมฺม “นตฺถิ กิญฺจี”ติ อากิญฺจญฺญายตนํ อุปสมฺปชฺช วิหรติ, อยํ ฉฏฺโฐ วิโมกฺโข.}

\prose{สพฺพโส อากิญฺจญฺญายตนํ สมติกฺกมฺม เนวสญฺญานาสญฺญายตนํ อุปสมฺปชฺช วิหรติ, อยํ สตฺตโม วิโมกฺโข.}

\prose{สพฺพโส เนวสญฺญานาสญฺญายตนํ สมติกฺกมฺม สญฺญาเวทยิตนิโรธํ อุปสมฺปชฺช วิหรติ, อยํ อฏฺฐโม วิโมกฺโข.\csromanspacer{} อิเม โข ภิกฺขเว อฏฺฐ วิโมกฺขาติ.\csromanspacer{}\csromanspacer{}ฉฏฺฐํ.}

\csromanheader{2}{7. อนริยโวหารสุตฺต}
\proseitem{67}{อฏฺฐิเม ภิกฺขเว อนริยโวหารา.\csromanspacer{} กตเม อฏฺฐ? อทิฏฺเฐ ทิฏฺฐวาทิตา, อสุเต สุตวาทิตา, อมุเต มุตวาทิตา, อวิญฺญาเต วิญฺญาตวาทิตา, ทิฏฺเฐ อทิฏฺฐวาทิตา, สุเต อสุตวาทิตา, มุเต อมุตวาทิตา, วิญฺญาเต อวิญฺญาตวาทิตา.\csromanspacer{} อิเม โข ภิกฺขเว อฏฺฐ อนริยโวหาราติ.\csromanspacer{}\csromanspacer{}สตฺตมํ.}

\csromanheader{2}{8. อริยโวหารสุตฺต}
\proseitem{68}{อฏฺฐิเม ภิกฺขเว อริยโวหารา.\csromanspacer{} กตเม อฏฺฐ? อทิฏฺเฐ อทิฏฺฐวาทิตา, อสุเต อสุตวาทิตา, อมุเต อมุตวาทิตา, อวิญฺญาเต อวิญฺญาตวาทิตา, ทิฏฺเฐ ทิฏฺฐวาทิตา, สุเต สุตวาทิตา, มุเต มุตวาทิตา, วิญฺญาเต วิญฺญาตวาทิตา.\csromanspacer{} อิเม โข ภิกฺขเว อฏฺฐ อริยโวหาราติ.\csromanspacer{}\csromanspacer{}อฏฺฐมํ.}

\csromanheader{2}{9. ปริสาสุตฺต}
\proseitem{69}{อฏฺฐิมา ภิกฺขเว ปริสา.\csromanspacer{} กตมา อฏฺฐ? ขตฺติยปริสา พฺราหฺมณปริสา คหปติปริสา สมณปริสา จาตุมหาราชิกปริสา ตาวติํสปริสา มารปริสา พฺรหฺมปริสา.\csromanspacer{} อภิชานามิ โข ปนาหํ ภิกฺขเว \csromanfolio{128}อเนกสตํ ขตฺติยปริสํ อุปสงฺกมิตา, ตตฺรปิ มยา สนฺนิสินฺนปุพฺพญฺเจว สลฺลปิตปุพฺพญฺจ สากจฺฉา จ สมาปนฺนปุพฺพา.\csromanspacer{} ตตฺถ ยาทิสโก เตสํ วณฺโณ โหติ, ตาทิสโก มยฺหํ วณฺโณ โหติ.\csromanspacer{} ยาทิสโก เตสํ สโร โหติ, ตาทิสโก มยฺหํ สโร โหติ.\csromanspacer{} ธมฺมิยา จ กถาย สนฺทสฺเสมิ สมาทเปมิ สมุตฺเตเชมิ สมฺปหํเสมิ.\csromanspacer{} ภาสมานญฺจ มํ น ชานนฺติ “โก นุ โข อยํ ภาสติ เทโว วา มนุสฺโส วา”ติ.\csromanspacer{} ธมฺมิยา กถาย สนฺทสฺเสตฺวา สมาทเปตฺวา สมุตฺเตเชตฺวา สมฺปหํเสตฺวา อนฺตรธายามิ.\csromanspacer{} อนฺตรหิตญฺจ มํ น ชานนฺติ “โก นุ โข อยํ อนฺตรหิโต เทโว วา มนุสฺโส วา”ติ.}

\prose{อภิชานามิ โข ปนาหํ ภิกฺขเว อเนกสตํ พฺราหฺมณปริสํ ฯเปฯ คหปติปริสํ.\csromanspacer{} สมณปริสํ.\csromanspacer{} จาตุมหาราชิกปริสํ.\csromanspacer{} ตาวติํสปริสํ.\csromanspacer{} มารปริสํ.\csromanspacer{} พฺรหฺมปริสํ อุปสงฺกมิตา.\csromanspacer{} ตตฺรปิ มยา สนฺนิสินฺนปุพฺพญฺเจว สลฺลปิตปุพฺพญฺจ สากจฺฉา จ สมาปนฺนปุพฺพา.\csromanspacer{} ตตฺถ ยาทิสโก เตสํ วณฺโณ โหติ, ตาทิสโก มยฺหํ วณฺโณ โหติ.\csromanspacer{} ยาทิสโก เตสํ สโร โหติ, ตาทิสโก มยฺหํ สโร โหติ.\csromanspacer{} ธมฺมิยา จ กถาย สนฺทสฺเสมิ สมาทเปมิ สมุตฺเตเชมิ สมฺปหํเสมิ.\csromanspacer{} ภาสมานญฺจ มํ น ชานนฺติ “โก นุ โข อยํ ภาสติ เทโว วา มนุสฺโส วา”ติ.\csromanspacer{} ธมฺมิยา กถาย สนฺทสฺเสตฺวา สมาทเปตฺวา สมุตฺเตเชตฺวา สมฺปหํเสตฺวา อนฺตรธายามิ.\csromanspacer{} อนฺตรหิตญฺจ มํ น ชานนฺติ “โก นุ โข อยํ อนฺตรหิโต เทโว วา มนุสฺโส วา”ติ.\csromanspacer{} อิมา โข ภิกฺขเว อฏฺฐ ปริสาติ.\csromanspacer{}\csromanspacer{}นวมํ.}

\csromanheader{2}{10. ภูมิจาลสุตฺต}
\proseitem{70}{เอกํ สมยํ ภควา เวสาลิยํ วิหรติ มหาวเน กูฏาคารสาลายํ.\csromanspacer{} อถ โข ภควา ปุพฺพณฺหสมยํ นิวาเสตฺวา ปตฺตจีวรมาทาย เวสาลิํ ปิณฺฑาย ปาวิสิ, เวสาลิยํ ปิณฺฑาย จริตฺวา ปจฺฉาภตฺตํ ปิณฺฑปาตปฏิกฺกนฺโต อายสฺมนฺตํ อานนฺทํ อามนฺเตสิ “คณฺหาหิ อานนฺท นิสีทนํ, เยน จาปาลํ เจติยํ\footnote{ปาวาลเจติยํ (สฺยา), จาปาลเจติยํ (อิ, ก)} เตนุปสงฺกมิสฺสาม ทิวาวิหารายา”ติ. “เอวํ ภนฺเต”ติ โข อายสฺมา อานนฺโท ภควโต ปฏิสฺสุตฺวา นิสีทนํ อาทาย ภควนฺตํ ปิฏฺฐิโต ปิฏฺฐิโต อนุพนฺธิ.}

\chapterheadpageread{(7) 2. ภูมิจาลวคฺค}
\prose{\csromanfolio{129}อถ โข ภควา เยน จาปาลํ เจติยํ เตนุปสงฺกมิ, อุปสงฺกมิตฺวา ปญฺญตฺเต อาสเน นิสีทิ, นิสชฺช โข ภควา อายสฺมนฺตํ อานนฺทํ อามนฺเตสิ “รมณียา อานนฺท เวสาลี, รมณียํ อุเทนํ เจติยํ, รมณียํ โคตมกํ เจติยํ, รมณียํ สตฺตมฺพํ เจติยํ, รมณียํ พหุปุตฺตกํ เจติยํ, รมณียํ สารนฺททํ เจติยํ, รมณียํ จาปาลํ เจติยํ.\csromanspacer{} ยสฺส กสฺสจิ อานนฺท จตฺตาโร อิทฺธิปาทา ภาวิตา พหุลีกตา ยานีกตา วตฺถุกตา อนุฏฺฐิตา ปริจิตา สุสมารทฺธา, อากงฺขมาโน โส อานนฺท กปฺปํ วา ติฏฺเฐยฺย กปฺปาวเสสํ วา.\csromanspacer{} ตถาคตสฺส โข อานนฺท จตฺตาโร อิทฺธิปาทา ภาวิตา พหุลีกตา ยานีกตา วตฺถุกตา อนุฏฺฐิตา ปริจิตา สุสมารทฺธา, อากงฺขมาโน อานนฺท ตถาคโต กปฺปํ วา ติฏฺเฐยฺย กปฺปาวเสสํ วา”ติ.\csromanspacer{} เอวมฺปิ โข อายสฺมา อานนฺโท ภควตา โอฬาริเก นิมิตฺเต กยิรมาเน โอฬาริเก โอภาเส กยิรมาเน นาสกฺขิ ปฏิวิชฺฌิตุํ, น ภควนฺตํ ยาจิ “ติฏฺฐตุ ภนฺเต ภควา กปฺปํ, ติฏฺฐตุ สุคโต กปฺปํ พหุชนหิตาย พหุชนสุขาย โลกานุกมฺปาย อตฺถาย หิตาย สุขาย เทวมนุสฺสานนฺ”ติ.\csromanspacer{} ยถา ตํ มาเรน ปริยุฏฺฐิตจิตฺโต.}

\prose{ทุติยมฺปิ โข ภควา ฯเปฯ.\csromanspacer{} ตติยมฺปิ โข ภควา อายสฺมนฺตํ อานนฺทํ อามนฺเตสิ “รมณียา อานนฺท เวสาลี, รมณียํ อุเทนํ เจติยํ, รมณียํ โคตมกํ เจติยํ, รมณียํ สตฺตมฺพํ เจติยํ, รมณียํ พหุปุตฺตกํ เจติยํ, รมณียํ สารนฺททํ เจติยํ, รมณียํ จาปาลํ เจติยํ.\csromanspacer{} ยสฺส กสฺสจิ อานนฺท จตฺตาโร อิทฺธิปาทา ภาวิตา พหุลีกตา ยานีกตา วตฺถุกตา อนุฏฺฐิตา ปริจิตา สุสมารทฺธา, อากงฺขมาโน โส อานนฺท กปฺปํ วา ติฏฺเฐยฺย กปฺปาวเสสํ วา.\csromanspacer{} ตถาคตสฺส โข อานนฺท จตฺตาโร อิทฺธิปาทา ภาวิตา ฯเปฯ อากงฺขมาโน อานนฺท ตถาคโต กปฺปํ วา ติฏฺเฐยฺย, กปฺปาวเสสํ วา”ติ.\csromanspacer{} เอวมฺปิ โข อายสฺมา อานนฺโท ภควตา โอฬาริเก นิมิตฺเต กยิรมาเน โอฬาริเก โอภาเส กยิรมาเน นาสกฺขิ ปฏิวิชฺฌิตุํ, น ภควนฺตํ ยาจิ “ติฏฺฐตุ ภนฺเต ภควา กปฺปํ, ติฏฺฐตุ สุคโต กปฺปํ พหุชนหิตาย พหุชนสุขาย โลกานุกมฺปาย อตฺถาย หิตาย สุขาย เทวมนุสฺสานนฺ”ติ.\csromanspacer{} ยถา ตํ มาเรน ปริยุฏฺฐิตจิตฺโต.}

\prose{\csromanfolio{130}อถ โข ภควา อายสฺมนฺตํ อานนฺทํ อามนฺเตสิ “คจฺฉ ตฺวํ\footnote{คจฺฉ โข ตฺวํ (สํ 3. 228 ปิฏฺเฐ) ขุ 1. 152 ปิฏฺเฐปิ ปสฺสิตพฺพํ.} อานนฺท, ยสฺส ทานิ กาลํ มญฺญสี”ติ. “เอวํ ภนฺเต”ติ โข อายสฺมา อานนฺโท ภควโต ปฏิสฺสุตฺวา อุฏฺฐายาสนา ภควนฺตํ อภิวาเทตฺวา ปทกฺขิณํ กตฺวา ภควโต อวิทูเร อญฺญตรสฺมิํ รุกฺขมูเล นิสีทิ.\csromanspacer{} อถ โข มาโร ปาปิมา อจิรปกฺกนฺเต อายสฺมนฺเต อานนฺเท ภควนฺตํ เอตทโวจ\csromandash{}}

\prose{ปรินิพฺพาตุทานิ ภนฺเต ภควา, ปรินิพฺพาตุ สุคโต.\csromanspacer{} ปรินิพฺพานกาโล ทานิ ภนฺเต ภควโต.\csromanspacer{} ภาสิตา โข ปเนสา ภนฺเต ภควตา วาจา “น ตาวาหํ ปาปิม ปรินิพฺพายิสฺสามิ, ยาว เม ภิกฺขู น สาวกา ภวิสฺสนฺติ วิยตฺตา วินีตา วิสารทา ปตฺตโยคกฺเขมา\footnote{อิทํ ปทํ ที 2. 87 ปิฏฺเฐ จ สํ 3. 228 ปิฏฺเฐ จ น ทิสฺสติ.} พหุสฺสุตา ธมฺมธรา ธมฺมานุธมฺมปฺปฏิปนฺนา สามีจิปฺปฏิปนฺนา อนุธมฺมจาริโน สกํ อาจริยกํ อุคฺคเหตฺวา อาจิกฺขิสฺสนฺติ เทเสสฺสนฺติ ปญฺญเปสฺสนฺติ ปฏฺฐเปสฺสนฺติ วิวริสฺสนฺติ วิภชิสฺสนฺติ อุตฺตานีกริสฺสนฺติ, อุปฺปนฺนํ ปรปฺปวาทํ สหธมฺเมน สุนิคฺคหิตํ นิคฺคเหตฺวา สปฺปาฏิหาริยํ ธมฺมํ เทเสสฺสนฺตี”ติ.\csromanspacer{} เอตรหิ ภนฺเต ภิกฺขุ ภควโต สาวกา วิยตฺตา วินีตา วิสารทา ปตฺตโยคกฺเขมา พหุสฺสุตา ธมฺมธรา ธมฺมานุธมฺมปฺปฏิปนฺนา สามีจิปฺปฏิปนฺนา อนุธมฺมจาริโน สกํ อาจริยกํ อุคฺคเหตฺวา อาจิกฺขนฺติ เทเสนฺติ ปญฺญเปนฺติ ปฏฺฐเปนฺติ วิวรนฺติ วิภชนฺติ อุตฺตานีกโรนฺติ, อุปฺปนฺนํ ปรปฺปวาทํ สหธมฺเมน สุนิคฺคหิตํ นิคฺคเหตฺวา สปฺปาฏิหาริยํ ธมฺมํ เทเสนฺติ.}

\prose{ปรินิพฺพาตุ ทานิ ภนฺเต ภควา, ปรินิพฺพาตุ สุคโต.\csromanspacer{} ปรินิพฺพานกาโล ทานิ ภนฺเต ภควโต.\csromanspacer{} ภาสิตา โข ปเนสา ภนฺเต ภควตา วาจา “น ตาวาหํ ปาปิม ปรินิพฺพายิสฺสามิ, ยาว เม ภิกฺขุนิโย น สาวิกา ภวิสฺสนฺติ ฯเปฯ ยาว เม อุปาสกา น สาวกา ภวิสฺสนฺติ ฯเปฯ.\csromanspacer{} ยาว เม อุปาสิกา น สาวิกา ภวิสฺสนฺติ วิยตฺตา วินีตา วิสารทา ปตฺตโยคกฺเขมา พหุสฺสุตา ธมฺมธรา ธมฺมานุธมฺมปฺปฏิปนฺนา สามีจิปฺปฏิปนฺนา อนุธมฺมจารินิโย สกํ อาจริยกํ อุคฺคเหตฺวา อาจิกฺขิสฺสนฺติ เทเสสฺสนฺติ ปญฺญเปสฺสนฺติ ปฏฺฐเปสฺสนฺติ วิวริสฺสนฺติ วิภชิสฺสนฺติ อุตฺตานีกริสฺสนฺติ, อุปฺปนฺนํ ปรปฺปวาทํ สหธมฺเมน สุนิคฺคหิตํ นิคฺคเหตฺวา สปฺปาฏิหาริยํ ธมฺมํ เทเสสฺสนฺตี”ติ.\csromanspacer{} เอตรหิ ภนฺเต อุปาสิกา ภควโต สาวิกา วิยตฺตา วินีตา วิสารทา ปตฺตโยคกฺเขมา}

\vspace{\tipitakaparskip}
\csromancenter{(7) 2.\csromanspacer{} ภูมิจาลวคฺค พหุสฺสุตา ธมฺมธรา ธมฺมานุธมฺมปฺปฏิปนฺนา สามีจิปฺปฏิปนฺนา อนุธมฺมจารินิโย สกํ อาจริยกํ อุคฺคเหตฺวา อาจิกฺขนฺติ เทเสนฺติ ปญฺญเปนฺติ ปฏฺฐเปนฺติ วิวรนฺติ วิภชนฺติ อุตฺตานีกโรนฺติ, อุปฺปนฺนํ ปรปฺปวาทํ สหธมฺเมน สุนิคฺคหิตํ นิคฺคเหตฺวา สปฺปาฏิหาริยํ ธมฺมํ เทเสนฺติ.}
\prose{\csromanfolio{131}ปรินิพฺพาตุ ทานิ ภนฺเต ภควา, ปรินิพฺพาตุ สุคโต.\csromanspacer{} ปรินิพฺพานกาโล ทานิ ภนฺเต ภควโต.\csromanspacer{} ภาสิตา โข ปเนสา ภนฺเต ภควตา วาจา “น ตาวาหํ ปาปิม ปรินิพฺพายิสฺสามิ, ยาว เม อิทํ พฺรหฺมจริยํ น อิทฺธญฺเจว ภวิสฺสติ ผีตญฺจ วิตฺถาริกํ พาหุชญฺญํ ปุถุภูตํ ยาว เทวมนุสฺเสหิ สุปฺปกาสิตนฺ”ติ.\csromanspacer{} เอตรหิ ภนฺเต ภควโต พฺรหฺมจริยํ อิทฺธญฺเจว ผีตญฺจ วิตฺถาริกํ พาหุชญฺญํ ปุถุภูตํ ยาว เทวมนุสฺเสหิ สุปฺปกาสิตํ.}

\prose{ปรินิพฺพาตุ ทานิ ภนฺเต ภควา, ปรินิพฺพาตุ สุคโต.\csromanspacer{} ปรินิพฺพานกาโล ทานิ ภนฺเต ภควโตติ.\csromanspacer{} อปฺโปสฺสุกฺโก ตฺวํ ปาปิม โหหิ, นจิรํ ตถาคตสฺส ปรินิพฺพานํ ภวิสฺสติ, อิโต ติณฺณํ มาสานํ อจฺจเยน ตถาคโต ปรินิพฺพายิสฺสตีติ.}

\prose{อถ โข ภควา จาปาเล เจติเย สโต สมฺปชาโน อายุสงฺขารํ โอสฺสชิ, โอสฺสฏฺเฐ จ ภควตา อายุสงฺขาเร มหาภูมิจาโล อโหสิ ภิํสนโก สโลมหํโส, เทวทุนฺทุภิโย จ ผลิํสุ.\csromanspacer{} อถ โข ภควา เอตมตฺตํ วิทิตฺวา ตายํ เวลายํ อิมํ อุทานํ อุทาเนสิ\csromandash{}}

\csromangathameasure{“ตุลมตุลญฺจ สมฺภวํ, ภวสงฺขารมวสฺสชิ มุนิ. \\ อชฺฌตฺตรโต สมาหิโต, อภินฺทิ กวจมิวตฺตสมฺภวนฺ”ติ.}
\csromangathagroup{\csromangathastanza{“ตุลมตุลญฺจ สมฺภวํ, ภวสงฺขารมวสฺสชิ มุนิ. \\ อชฺฌตฺตรโต สมาหิโต, อภินฺทิ กวจมิวตฺตสมฺภวนฺ”ติ.}}

\prose{อถ โข อายสฺมโต อานนฺทสฺส เอตทโหสิ “มหา วตายํ ภูมิจาโล, สุมหา วตายํ ภูมิจาโล ภิํสนโก สโลมหํโส, เทวทุนฺทุภิโย จ ผลิํสุ.\csromanspacer{} โก นุ โข เหตุ โก ปจฺจโย มหโต ภูมิจาลสฺส ปาตุภาวายา”ติ.}

\prose{อถ โข อายสฺมา อานนฺโท เยน ภควา เตนุปสงฺกมิ, อุปสงฺกมิตฺวา ภควนฺตํ อภิวาเทตฺวา เอกมนฺตํ นิสีทิ, เอกมนฺตํ นิสินฺโน โข อายสฺมา อานนฺโท ภควนฺตํ เอตทโวจ “มหา วตายํ ภนฺเต}

\prose{\csromanfolio{132}ภูมิจาโล, สุมหา วตายํ ภนฺเต ภูมิจาโล ภิํสนโก สโลมหํโส, เทวทุนฺทุภิโย จ ผลิํสุ.\csromanspacer{} โก นุ โข ภนฺเต เหตุ โก ปจฺจโย มหโต ภูมิจาลสฺส ปาตุภาวายา”ติ.}

\prose{อฏฺฐิเม อานนฺท เหตู อฏฺฐ ปจฺจยา มหโต ภูมิจาลสฺส ปาตุภาวาย.\csromanspacer{} กตเม อฏฺฐ? อยํ อานนฺท มหาปถวี อุทเก ปติฏฺฐิตา, อุทกํ วาเต ปติฏฺฐิตํ, วาโต อากาสฏฺโฐ โหติ.\csromanspacer{} โส อานนฺท สมโย, ยํ มหาวาตา วายนฺติ, มหาวาตา วายนฺตา อุทกํ กมฺเปนฺติ, อุทกํ กมฺปิตํ ปถวิํ กมฺเปติ.\csromanspacer{} อยํ อานนฺท ปฐโม เหตุ ปฐโม ปจฺจโย มหโต ภูมิจาลสฺส ปาตุภาวาย.}

\prose{ปุน จปรํ อานนฺท สมโณ วา พฺราหฺมโณ วา อิทฺธิมา เจโตวสิปฺปตฺโต เทวตา วา มหิทฺธิกา มหานุภาวา, ตสฺส ปริตฺตา ปถวีสญฺญา ภาวิตา โหติ อปฺปมาณา อาโปสญฺญา.\csromanspacer{} โส อิมํ ปถวิํ กมฺเปติ สงฺกมฺเปติ สมฺปกมฺเปติ สมฺปเวเธติ.\csromanspacer{} อยํ อานนฺท ทุติโย เหตุ ทุติโย ปจฺจโย มหโต ภูมิจาลสฺส ปาตุภาวาย.}

\prose{ปุน จปรํ อานนฺท ยทา โพธิสตฺโต ตุสิตา กายา จวิตฺวา สโต สมฺปชาโน มาตุกุจฺฉิํ โอกฺกมติ, ตทายํ ปถวี กมฺปติ สงฺกมฺปติ สมฺปกมฺปติ สมฺปเวธติ.\csromanspacer{} อยํ อานนฺท ตติโย เหตุ ตติโย ปจฺจโย มหโต ภูมิจาลสฺส ปาตุภาวาย.}

\prose{ปุน จปรํ อานนฺท ยทา โพธิสตฺโต สโต สมฺปชาโน มาตุกุจฺฉิสฺมา นิกฺขมติ, ตทายํ ปถวี กมฺปติ สงฺกมฺปติ สมฺปกมฺปติ สมฺปเวธติ.\csromanspacer{} อยํ อานนฺท จตุตฺโถ เหตุ จตุตฺโถ ปจฺจโย มหโต ภูมิจาลสฺส ปาตุภาวาย.}

\prose{ปุน จปรํ อานนฺท ยทา ตถาคโต อนุตฺตรํ สมฺมาสมฺโพธิํ อภิสมฺพุชฺฌติ, ตทายํ ปถวี กมฺปติ สงฺกมฺปติ สมฺปกมฺปติ สมฺปเวธติ, อยํ อานนฺท ปญฺจโม เหตุ ปญฺจโม ปจฺจโย มหโต ภูมิจาลสฺส ปาตุภาวาย.}

\prose{ปุน จปรํ อานนฺท ยทา ตถาคโต อนุตฺตรํ ธมฺมจกฺกํ ปวตฺเตติ.\csromanspacer{} ตทายํ ปถวี กมฺปติ สงฺกมฺปติ สมฺปกมฺปติ สมฺปเวธติ.\csromanspacer{} อยํ อานนฺท ฉฏฺโฐ เหตุ ฉฏฺโฐ ปจฺจโย มหโต ภูมิจาลสฺส ปาตุภาวาย.}

\chapterheadpageread{\textbf{(8) 3. ยมกวคฺค}}
\prose{\csromanfolio{133}ปุน จปรํ อานนฺท ยทา ตถาคโต สโต สมฺปชาโน อายุสงฺขารํ โอสฺสชฺชติ, ตทายํ ปถวี กมฺปติ สงฺกมฺปติ สมฺปกมฺปติ สมฺปเวธติ.\csromanspacer{} อยํ อานนฺท สตฺตโม เหตุ สตฺตโม ปจฺจโย มหโต ภูมิจาลสฺส ปาตุภาวาย.}

\prose{ปุน จปรํ อานนฺท ยทา ตถาคโต อนุปาทิเสสาย นิพฺพานธาตุยา ปรินิพฺพายติ, ตทายํ ปถวี กมฺปติ สงฺกมฺปติ สมฺปกมฺปติ สมฺปเวธติ.\csromanspacer{} อยํ อานนฺท อฏฺฐโม เหตุ อฏฺฐโม ปจฺจโย มหโต ภูมิจาลสฺส ปาตุภาวาย.\csromanspacer{} อิเม โข อานนฺท อฏฺฐ เหตู อฏฺฐ ปจฺจยา มหโต ภูมิจาลสฺส ปาตุภาวายาติ.\csromanspacer{}\csromanspacer{}ทสมํ.\csromanspacer{} ภูมิจาลวคฺโค ทุติโย.}

\titlehead{\textbf{ตสฺสุทฺทานํ}}
\csromangathameasure{อิจฺฉา อลญฺจ สํขิตฺตํ, คยา อภิภุนา สห. \\ วิโมกฺโข เทฺว จ โวหารา, ปริสา ภูมิจาเลนาติ.}
\csromangathagroup{\csromangathastanza{อิจฺฉา อลญฺจ สํขิตฺตํ, คยา อภิภุนา สห. \\ วิโมกฺโข เทฺว จ โวหารา, ปริสา ภูมิจาเลนาติ.}}

\chapterhead{\textbf{(8) 3. ยมกวคฺค}}
\csromanheader{2}{1. ปฐมสทฺธาสุตฺต}
\proseitem{71}{สทฺโธ จ\footnote{สทฺโธ (สฺยา) เอตฺเถว. อุปริ 170, 264 ปิฏฺเฐสุปิ.} ภิกฺขเว ภิกฺขุ โหติ, โน จ\footnote{โน (สฺยา) เอวมุปริปิ “โน”เตฺวว ทิสฺสติ.} สีลวา, เอวํ โส เตนงฺเคน อปริปูโร โหติ.\csromanspacer{} เตน ตํ องฺคํ ปริปูเรตพฺพํ “กินฺตาหํ สทฺโธ จ อสฺสํ สีลวา จา”ติ.\csromanspacer{} ยโต จ โข ภิกฺขเว ภิกฺขุ สทฺโธ จ โหติ สีลวา จ, เอวํ โส เตนงฺเคน ปริปูโร โหติ.}

\prose{สทฺโธ จ ภิกฺขเว ภิกฺขุ โหติ สีลวา จ, โน จ พหุสฺสุโต, เอวํ โส เตนงฺเคน อปริปูโร โหติ.\csromanspacer{} เตน ตํ องฺคํ ปริปูเรตพฺพํ “กินฺตาหํ สทฺโธ จ อสฺสํ สีลวา จ พหุสฺสุโต จา”ติ.\csromanspacer{} ยโต จ โข ภิกฺขเว ภิกฺขุ สทฺโธ จ โหติ สีลวา จ พหุสฺสุโต จ, เอวํ โส เตนงฺเคน ปริปูโร โหติ.}

\prose{\csromanfolio{134}สทฺโธ จ ภิกฺขเว ภิกฺขุ โหติ สีลวา จ พหุสฺสุโต จ, โน จ ธมฺมกถิโก ฯเปฯ ธมฺมกถิโก จ, โน จ ปริสาวจโร ฯเปฯ ปริสาวจโร จ, โน จ วิสารโท ปริสาย ธมฺมํ เทเสติ ฯเปฯ วิสารโท จ ปริสาย ธมฺมํ เทเสติ, โน จ จตุนฺนํ ฌานานํ อาภิเจตสิกานํ ทิฏฺฐธมฺมสุขวิหารานํ นิกามลาภี โหติ อกิจฺฉลาภี อกสิรลาภี ฯเปฯ จตุนฺนํ ฌานานํ อาภิเจตสิกานํ ทิฏฺฐธมฺมสุขวิหารานํ นิกามลาภี โหติ อกิจฺฉลาภี อกสิรลาภี, โน จ อาสวานํ ขยา อนาสวํ เจโตวิมุตฺติํ ปญฺญาวิมุตฺติํ ทิฏฺเฐว ธมฺเม สยํ อภิญฺญา สจฺฉิกตฺวา อุปสมฺปชฺช วิหรติ, เอวํ โส เตนงฺเคน อปริปูโร โหติ.\csromanspacer{} เตน ตํ องฺคํ ปริปูเรตพฺพํ “กินฺตาหํ สทฺโธ จ อสฺสํ สีลวา จ พหุสฺสุโต จ ธมฺมกถิโก จ ปริสาวจโร จ, วิสารโท จ ปริสาย ธมฺมํ เทเสยฺยํ, จตุนฺนญฺจ ฌานานํ อาภิเจตสิกานํ ทิฏฺฐธมฺมสุขวิหารานํ นิกามลาภี อสฺสํ อกิจฺฉลาภี อกสิรลาภี, อาสวานญฺจ ขยา อนาสวํ เจโตวิมุตฺติํ ปญฺญาวิมุตฺติํ ทิฏฺเฐว ธมฺเม สยํ อภิญฺญา สจฺฉิกตฺวา อุปสมฺปชฺช วิหเรยฺยนฺ”ติ.}

\prose{ยโต จ โข ภิกฺขเว ภิกฺขุ สทฺโธ จ โหติ สีลวา จ พหุสฺสุโต จ ธมฺมกถิโก จ ปริสาวจโร จ, วิสารโท จ ปริสาย ธมฺมํ เทเสติ, จตุนฺนญฺจ ฌานานํ อาภิเจตสิกานํ ทิฏฺฐธมฺมสุขวิหารานํ นิกามลาภี โหติ อกิจฺฉลาภี อกสิรลาภี, อาสวานญฺจ ขยา อนาสวํ เจโตวิมุตฺติํ ปญฺญาวิมุตฺติํ ทิฏฺเฐว ธมฺเม สยํ อภิญฺญา สจฺฉิกตฺวา อุปสมฺปชฺช วิหรติ, เอวํ โส เตนงฺเคน ปริปูโร โหติ.\csromanspacer{} อิเมหิ โข ภิกฺขเว อฏฺฐหิ ธมฺเมหิ สมนฺนาคโต ภิกฺขุ สมนฺตปาสาทิโก จ โหติ สพฺพาการปริปูโร จาติ.\csromanspacer{}\csromanspacer{}ปฐมํ.}

\csromanheader{2}{2. ทุติยสทฺธาสุตฺต}
\proseitem{72}{สทฺโธ จ ภิกฺขเว ภิกฺขุ โหติ, โน จ สีลวา, เอวํ โส เตนงฺเคน อปริปูโร โหติ.\csromanspacer{} เตน ตํ องฺคํ ปริปูเรตพฺพํ “กินฺตาหํ สทฺโธ จ อสฺสํ สีลวา จา”ติ.\csromanspacer{} ยโต จ โข ภิกฺขเว ภิกฺขุ สทฺโธ จ โหติ สีลวา จ, เอวํ โส เตนงฺเคน ปริปูโร โหติ.}

\prose{สทฺโธ จ ภิกฺขเว ภิกฺขุ โหติ สีลวา จ, โน จ พหุสฺสุโต ฯเปฯ พหุสฺสุโต จ, โน จ ธมฺมกถิโก ฯเปฯ ธมฺมกถิโก จ, โน จ ปริสาวจโร ฯเปฯ ปริสาวจโร จ, โน จ วิสารโท ปริสาย ธมฺมํ}

\vspace{\tipitakaparskip}
\csromancenter{(8) 3.\csromanspacer{} ยมกวคฺค เทเสติ ฯเปฯ วิสารโท จ ปริสาย ธมฺมํ เทเสติ, โน จ เย เต สนฺตา วิโมกฺขา อติกฺกมฺม รูเป อารุปฺปา, เต กาเยน ผุสิตฺวา วิหรติ ฯเปฯ เย เต สนฺตา วิโมกฺขา อติกฺกมฺม รูเป อารุปฺปา, เต กาเยน ผุสิตฺวา วิหรติ, โน จ อาสวานํ ขยา อนาสวํ เจโตวิมุตฺติํ ปญฺญาวิมุตฺติํ ทิฏฺเฐว ธมฺเม สยํ อภิญฺญา สจฺฉิกตฺวา อุปสมฺปชฺช วิหรติ, เอวํ โส เตนงฺเคน อปริปูโร โหติ.\csromanspacer{} เตน ตํ องฺคํ ปริปูเรตพฺพํ “กินฺตาหํ สทฺโธ จ อสฺสํ สีลวา จ พหุสฺสุโต จ ธมฺมกถิโก จ ปริสาวจโร จ, วิสารโท จ ปริสาย ธมฺมํ เทเสยฺยํ, เย เต สนฺตา วิโมกฺขา อติกฺกมฺม รูเป อารุปฺปา, เต กาเยน ผุสิตฺวา วิหเรยฺยํ, อาสวานญฺจ ขยา อนาสวํ เจโตวิมุตฺติํ ปญฺญาวิมุตฺติํ ทิฏฺเฐว ธมฺเม สยํ อภิญฺญา สจฺฉิกตฺวา อุปสมฺปชฺช วิหเรยฺยนฺ”ติ.}
\prose{\csromanfolio{135}ยโต จ โข ภิกฺขเว ภิกฺขุ สทฺโธ จ โหติ สีลวา จ พหุสฺสุโต จ ธมฺมกถิโก จ ปริสาวจโร จ, วิสารโท จ ปริสาย ธมฺมํ เทเสติ, เย เต สนฺตา วิโมกฺขา อตฺติกฺกมฺม รูเป อารุปฺปา, เต จ กาเยน ผุสิตฺวา วิหรติ, อาสวานญฺจ ขยา ฯเปฯ สจฺฉิกตฺวา อุปสมฺปชฺช วิหรติ, เอวํ โส เตนงฺเคน ปริปูโร โหติ.\csromanspacer{} อิเมหิ โข ภิกฺขเว อฏฺฐหิ ธมฺเมหิ สมฺมาคโต ภิกฺขุ สมนฺตปาสาทิโก จ โหติ สพฺพาการปริปูโร จาติ.\csromanspacer{}\csromanspacer{}ทุติยํ.}

\csromanheader{2}{3. ปฐมมรณสฺสติสุตฺต}
\proseitem{73}{เอกํ สมยํ ภควา นาติเก\footnote{นาทิเก (สี, สฺยา), นาฏิเก (อิ) อํ 2. 268 ปิฏฺเฐปิ} วิหรติ คิญฺชกาวสเถ.\csromanspacer{} ตตฺร โข ภควา ภิกฺขู อามนฺเตสิ “ภิกฺขโว”ติ. “ภนฺเต”ติ เต ภิกฺขู ภควโต ปจฺจสฺโสสุํ.\csromanspacer{} ภควา เอตทโวจ “มรณสฺสติ ภิกฺขเว ภาวิตา พหุลีกตา มหปฺผลา โหติ มหานิสํสา อมโตคธา อมตปริโยสานา.\csromanspacer{} ภาเวถ โน ตุเมฺห ภิกฺขเว มรณสฺสตินฺ”ติ.}

\prose{เอวํ วุตฺเต อญฺญตโร ภิกฺขุ ภควนฺตํ เอตทโวจ “อหํ โข ภนฺเต ภาเวมิ มรณสฺสตินฺ”ติ.\csromanspacer{} ยถา กถํ ปน ตฺวํ ภิกฺขุ ภาเวสิ มรณสฺสตินฺติ.\csromanspacer{} อิธ มยฺหํ ภนฺเต เอวํ โหติ “อโห วตาหํ รตฺตินฺทิวํ \csromanfolio{136}ชีเวยฺยํ, ภควโต สาสนํ มนสิ กเรยฺยํ, พหุ\footnote{พหุํ (สี, อิ)} วต เม กตํ อสฺสา”ติ.\csromanspacer{} เอวํ โข อหํ ภนฺเต ภาเวมิ มรณสฺสตินฺติ. (1)}

\prose{อญฺญตโรปิ โข ภิกฺขุ ภควนฺตํ เอตทโวจ “อหมฺปิ โข ภนฺเต ภาเวมิ มรณสฺสตินฺ”ติ.\csromanspacer{} ยถา กถํ ปน ตฺวํ ภิกฺขุ ภาเวสิ มรณสฺสตินฺติ.\csromanspacer{} อิธ มยฺหํ ภนฺเต เอวํ โหติ “อโห วตาหํ ทิวสํ ชีเวยฺยํ, ภควโต สาสนํ มนสิ กเรยฺยํ, พหุ วต เม กตํ อสฺสา”ติ.\csromanspacer{} เอวํ โข อหํ ภนฺเต ภาเวมิ มรณสฺสตินฺติ. (2)}

\prose{อญฺญตโร โข ภิกฺขุ ภควนฺตํ เอตทโวจ “อหมฺปิ โข ภนฺเต ภาเวมิ มรณสฺสตินฺ”ติ.\csromanspacer{} ยถา กถํ ปน ตฺวํ ภิกฺขุ ภาเวสิ มรณสฺสตินฺติ.\csromanspacer{} อิธ มยฺหํ ภนฺเต เอวํ โหติ “อโห วตาหํ อุปฑฺฒทิวสํ ชีเวยฺยํ, ภควโต สาสนํ มนสิ กเรยฺยํ, พหุ วต เม กตํ อสฺสา”ติ.\csromanspacer{} เอวํ โข อหํ ภนฺเต ภาเวมิ มรณสฺสตินฺติ. (3)}

\prose{อญฺญตโรปิ โข ภิกฺขุ ภควนฺตํ เอตทโวจ “อหมฺปิ โข ภนฺเต ภาเวมิ มรณสฺสตินฺ”ติ.\csromanspacer{} ยถา กถํ ปน ตฺวํ ภิกฺขุ ภาเวสิ มรณสฺสตินฺติ.\csromanspacer{} อิธ มยฺหํ ภนฺเต เอวํ โหติ “อโห วตาหํ ตทนฺตรํ ชีเวยฺยํ, ยทนฺตรํ เอกปิณฺฑปาตํ ภุญฺชามิ, ภควโต สาสนํ มนสิ กเรยฺยํ, พหุ วต เม กตํ อสฺสา”ติ.\csromanspacer{} เอวํ โข อหํ ภนฺเต ภาเวมิ มรณสฺสตินฺติ. (4)}

\prose{อญฺญตโรปิ โข ภิกฺขุ ภควนฺตํ เอตทโวจ “อหมฺปิ โข ภนฺเต ภาเวมิ มรณสฺสตินฺ”ติ.\csromanspacer{} ยถา กถํ ปน ตฺวํ ภิกฺขุ ภาเวสิ มรณสฺสตินฺติ.\csromanspacer{} อิธ มยฺหํ ภนฺเต เอวํ โหติ “อโห วตาหํ ตทนฺตรํ ชีเวยฺยํ, ยทนฺตรํ อุปฑฺฒปิณฺฑปาตํ ภุญฺชามิ, ภควโต สาสนํ มนสิ กเรยฺยํ, พหุ วต เม กตํ อสฺสา”ติ.\csromanspacer{} เอวํ โข อหํ ภนฺเต ภาเวมิ มรณสฺสตินฺติ. (5)}

\prose{อญฺญตโรปิ โข ภิกฺขุ ภควนฺตํ เอตทโวจ “อหมฺปิ โข ภนฺเต ภาเวมิ มรณสฺสตินฺ”ติ.\csromanspacer{} ยถา กถํ ปน ตฺวํ ภิกฺขุ ภาเวสิ มรณสฺสตินฺติ.\csromanspacer{} อิธ มยฺหํ ภนฺเต เอวํ โหติ “อโห วตาหํ ตทนฺตรํ ชีเวยฺยํ, ยทนฺตรํ จตฺตาโร ปญฺจ อาโลเป สงฺขาทิตฺวา\footnote{สงฺขริตฺวา (ก)} อชฺโฌหรามิ,}

\vspace{\tipitakaparskip}
\csromancenter{(8) 3.\csromanspacer{} ยมกวคฺค ภควโต สาสนํ มนสิ กเรยฺยํ, พหุ วต เม กตํ อสฺสา”ติ.\csromanspacer{} เอวํ โข อหํ ภนฺเต ภาเวมิ มรณสฺสตินฺติ. (6)}
\prose{\csromanfolio{137}อญฺญตโรปิ โข ภิกฺขุ ภควนฺตํ เอตทโวจ “อหมฺปิ โข ภนฺเต ภาเวมิ มรณสฺสตินฺ”ติ.\csromanspacer{} ยถา กถํ ปน ตฺวํ ภิกฺขุ ภาเวสิ มรณสฺสตินฺติ.\csromanspacer{} อิธ มยฺหํ ภนฺเต เอวํ โหติ “อโห วตาหํ ตทนฺตรํ ชีเวยฺยํ, ยทนฺตรํ เอกํ อาโลปํ สงฺขาทิตฺวา\footnote{สงฺขริตฺวา (ก)} อชฺโฌหรามิ, ภควโต สาสนํ มนสิ กเรยฺยํ, พหุ วต เม กตํ อสฺสา”ติ.\csromanspacer{} เอวํ โข อหํ ภนฺเต ภาเวมิ มรณสฺสตินฺติ. (7)}

\prose{อญฺญตโรปิ โข ภิกฺขุ ภควนฺตํ เอตทโวจ “อหมฺปิ โข ภนฺเต ภาเวมิ มรณสฺสตินฺ”ติ.\csromanspacer{} ยถา กถํ ปน ตฺวํ ภิกฺขุ ภาเวสิ มรณสฺสตินฺติ.\csromanspacer{} อิธ มยฺหํ ภนฺเต เอวํ โหติ “อโห วตาหํ ตทนฺตรํ ชีเวยฺยํ, ยทนฺตรํ อสฺสสิตฺวา วา ปสฺสสามิ, ปสฺสสิตฺวา วา อสฺสสามิ, ภควโต สาสนํ มนสิ กเรยฺยํ, พหุ วต เม กตํ อสฺสา”ติ.\csromanspacer{} เอวํ โข อหํ ภนฺเต ภาเวมิ มรณสฺสตินฺติ. (8)}

\prose{เอวํ วุตฺเต ภควา เต ภิกฺขู เอตทโวจ\csromandash{}ยฺวายํ\footnote{โย จายํ (ก-สี)} ภิกฺขเว ภิกฺขุ เอวํ มรณสฺสติํ ภาเวติ “อโห วตาหํ รตฺตินฺทิวํ ชีเวยฺยํ, ภควโต สาสนํ มนสิ กเรยฺยํ, พหุ วต เม กตํ อสฺสา”ติ.\csromanspacer{} โย จายํ\footnote{โย ปายํ (ก) อํ 2. 269 ปิฏฺเฐ ปสฺสิตพฺพํ.} ภิกฺขเว ภิกฺขุ เอวํ มรณสฺสติํ ภาเวติ “อโห วตาหํ ทิวสํ ชีเวยฺยํ, ภควโต สาสนํ มนสิ กเรยฺยํ, พหุ วต เม กตํ อสฺสา”ติ.\csromanspacer{} โย จายํ3 ภิกฺขเว ภิกฺขุ เอวํ มรณสฺสติํ ภาเวติ “อโห วตาหํ อุปฑฺฒทิวสํ ชีเวยฺยํ, ภควโต สาสนํ มนสิ กเรยฺยํ, พหุ วต เม กตํ อสฺสา”ติ.\csromanspacer{} โย จายํ3 ภิกฺขเว ภิกฺขุ เอวํ มรณสฺสติํ ภาเวติ “อโห วตาหํ ตทนฺตรํ ชีเวยฺยํ, ยทนฺตรํ เอกปิณฺฑปาตํ ภุญฺชามิ, ภควโต สาสนํ มนสิ กเรยฺยํ, พหุ วต เม กตํ อสฺสา”ติ.\csromanspacer{} โย จายํ3 ภิกฺขเว ภิกฺขุ เอวํ มรณสฺสติํ ภาเวติ “อโห วตาหํ ตทนฺตรํ ชีเวยฺยํ, ยทนฺตรํ อุปฑฺฒปิณฺฑปาตํ ภุญฺชามิ, ภควโต สาสนํ มนสิ กเรยฺยํ, พหุ วต เม กตํ อสฺสา”ติ.\csromanspacer{} โย จายํ3 ภิกฺขเว ภิกฺขุ เอวํ มรณสฺสติํ ภาเวติ “อโห วตาหํ \csromanfolio{138}ตทนฺตรํ ชีเวยฺยํ, ยทนฺตรํ จตฺตาโร ปญฺจ อาโลเป สงฺขาทิตฺวา อชฺโฌหรามิ, ภควโต สาสนํ มนสิ กเรยฺยํ, พหุ วต เม กตํ อสฺสา”ติ.\csromanspacer{} อิเม วุจฺจนฺติ ภิกฺขเว ภิกฺขู ปมตฺตา วิหรนฺติ, ทนฺธํ มรณสฺสติํ ภาเวนฺติ อาสวานํ ขยาย.}

\prose{โย จ ขฺวายํ\footnote{โย จายํ (สฺยา), โย จ โข ยํ (ก)} ภิกฺขเว ภิกฺขุ เอวํ มรณสฺสติํ ภาเวติ “อโห วตาหํ ตทนฺตรํ ชีเวยฺยํ, ยทนฺตรํ เอกํ อาโลปํ สงฺขาทิตฺวา อชฺโฌหรามิ, ภควโต สาสนํ มนสิ กเรยฺยํ, พหุ วต เม กตํ อสฺสา”ติ.\csromanspacer{} โยจายํ\footnote{โย ปายํ (ก)} ภิกฺขเว ภิกฺขุ เอวํ มรณสฺสติํ ภาเวติ “อโห วตาหํ ตทนฺตรํ ชีเวยฺยํ, ยทนฺตรํ อสฺสสิตฺวา วา ปสฺสสามิ, ปสฺสสิตฺวา วา อสฺสสามิ, ภควโต สาสนํ มนสิ กเรยฺยํ, พหุ วต เม กตํ อสฺสา”ติ.\csromanspacer{} อิเม วุจฺจนฺติ ภิกฺขเว ภิกฺขู อปฺปมตฺตา วิหรนฺติ, ติกฺขํ มรณสฺสติํ ภาเวนฺติ อาสวานํ ขยาย.}

\prose{ตสฺมาติห ภิกฺขเว เอวํ สิกฺขิตพฺพํ “อปฺปมตฺตา วิหริสฺสาม, ติกฺขํ มรณสฺสติํ ภาวยิสฺสาม อาสวานํ ขยายา”ติ, เอวํ หิ โว ภิกฺขเว สิกฺขิตพฺพนฺติ.\csromanspacer{}\csromanspacer{}ตติยํ.}

\csromanheader{2}{4. ทุติยมรณสฺสติสุตฺต}
\proseitem{74}{เอกํ สมยํ ภควา นาติเก วิหรติ คิญฺชกาวสเถ.\csromanspacer{} ตตฺร โข ภควา ภิกฺขู อามนฺเตสิ ฯเปฯ.\csromanspacer{} มรณสฺสติ ภิกฺขเว ภาวิตา พหุลีกตา มหปฺผลา โหติ มหานิสํสา อมโตคธา อมตปริโยสานา.}

\prose{กถํ ภาวิตา จ ภิกฺขเว มรณสฺสติ กถํ พหุลีกตา มหปฺผลา โหติ มหานิสํสา อมโตคธา อมตปริโยสานา.\csromanspacer{} อิธ ภิกฺขเว ภิกฺขุ ทิวเส นิกฺขนฺเต รตฺติยา ปติหิตาย\footnote{ปฏิหิตาย (อิ), (อํ 2. 270 ปิฏฺเฐปิ ปสฺสิตพฺพํ) 4. มมสฺส (อํ 2. 271 ปิฏฺเฐ)} อิติ ปฏิสญฺจิกฺขติ “พหุกา โข เม ปจฺจยา มรณสฺส, อหิ วา มํ ฑํเสยฺย, วิจฺฉิโก วา มํ ฑํเสยฺย, สตปที วา มํ ฑํเสยฺย, เตน เม อสฺส กาลกิริยา, โส มม อสฺส4 อนฺตราโย.\csromanspacer{} อุปกฺขลิตฺวา วา ปปเตยฺยํ, ภตฺตํ วา เม ภุตฺตํ พฺยาปชฺเชยฺย, ปิตฺตํ วา เม กุปฺเปยฺย, เสมฺหํ วา เม กุปฺเปยฺย, สตฺถกา วา เม วาตา}

\vspace{\tipitakaparskip}
\csromancenter{(8) 3.\csromanspacer{} ยมกวคฺค กุปฺเปยฺยุํ, มนุสฺสา วา มํ อุปกฺกเมยฺยุํ, อมนุสฺสา วา มํ อุปกฺกเมยฺยุํ, เตน เม อสฺส กาลกิริยา, โส มม อสฺส อนฺตราโย”ติ.\csromanspacer{} เตน ภิกฺขเว ภิกฺขุนา อิติ ปฏิสญฺจิกฺขิตพฺพํ “อตฺถิ นุ โข เม ปาปกา อกุสลา ธมฺมา อปฺปหีนา, เย เม อสฺสุ รตฺติํ กาลํ กโรนฺตสฺส อนฺตรายายา”ติ.}
\prose{\csromanfolio{139}สเจ ภิกฺขเว ภิกฺขุ ปจฺจเวกฺขมาโน เอวํ ชานาติ “อตฺถิ เม ปาปกา อกุสลา ธมฺมา อปฺปหีนา, เย เม อสฺสุ รตฺติํ กาลํ กโรนฺตสฺส อนฺตรายายา”ติ.\csromanspacer{} เตน ภิกฺขเว ภิกฺขุนา เตสํเยว ปาปกานํ อกุสลานํ ธมฺมานํ ปหานาย อธิมตฺโต ฉนฺโท จ วายาโม จ อุสฺสาโห จ อุสฺโสฬฺหี จ อปฺปฏิวานี จ สติ จ สมฺปชญฺญญฺจ กรณียํ.}

\prose{เสยฺยถาปิ ภิกฺขเว อาทิตฺตเจโล วา อาทิตฺตสีโส วา ตสฺเสว เจลสฺส วา สีสสฺส วา นิพฺพาปนาย อธิมตฺตํ ฉนฺทญฺจ วายามญฺจ อุสฺสาหญฺจ อุสฺโสฬฺหิญฺจ อปฺปฏิวานิญฺจ สติญฺจ สมฺปชญฺญญฺจ กเรยฺย.\csromanspacer{} เอวเมวํ โข ภิกฺขเว เตน ภิกฺขุนา เตสํเยว ปาปกานํ อกุสลานํ ธมฺมานํ ปหานาย อธิมตฺโต ฉนฺโท จ วายาโม จ อุสฺสาโห จ อุสฺโสฬฺหี จ อปฺปฏิวานี จ สติ จ สมฺปชญฺญญฺจ กรณียํ.}

\prose{สเจ ปน ภิกฺขเว ภิกฺขุ ปจฺจเวกฺขมาโน เอวํ ชานาติ “นตฺถิ เม ปาปกา อกุสลา ธมฺมา อปฺปหีนา, เย เม อสฺสุ รตฺติํ กาลํ กโรนฺตสฺส อนฺตรายายา”ติ.\csromanspacer{} เตน ภิกฺขเว ภิกฺขุนา เตเนว ปีติปาโมชฺเชน วิหาตพฺพํ อโหรตฺตานุสิกฺขินา กุสเลสุ ธมฺเมสุ.}

\prose{อิธ ปน ภิกฺขเว ภิกฺขุ รตฺติยา นิกฺขนฺตาย ทิวเส ปติหิเต อิติ ปฏิสญฺจิกฺขติ “พหุกา โข เม ปจฺจยา มรณสฺส, อหิ วา มํ ฑํเสยฺย, วิจฺฉิโก วา มํ ฑํเสยฺย, สตปที วา มํ ฑํเสยฺย, เตน เม อสฺส กาลกิริยา, โส มม อสฺส อนฺตราโย.\csromanspacer{} อุปกฺขลิตฺวา วา ปปเตยฺยํ, ภตฺตํ วา เม ภุตฺตํ พฺยาปชฺเชยฺย, ปิตฺตํ วา เม กุปฺเปยฺย, เสมฺหํ วา เม กุปฺเปยฺย, สตฺถกา วา เม วาตา กุปฺเปยฺยุํ, มนุสฺสา วา มํ อุปกฺกเมยฺยุํ, อมนุสฺสา วา มํ อุปกฺกเมยฺยุํ, เตน เม อสฺส กาลกิริยา, โส มม อสฺส อนฺตราโย”ติ.\csromanspacer{} เตน ภิกฺขเว ภิกฺขุนา อิติ ปฏิสญฺจิกฺขิตพฺพํ “อตฺถิ นุ โข เม ปาปกา อกุสลา ธมฺมา อปฺปหีนา, เย เม อสฺสุ ทิวา กาลํ กโรนฺตสฺส อนฺตรายายา”ติ.}

\prose{\csromanfolio{140}สเจ ภิกฺขเว ภิกฺขุ ปจฺจเวกฺขมาโน เอวํ ชานาติ “อตฺถิ เม ปาปกา อกุสลา ธมฺมา อปฺปหีนา, เย เม อสฺสุ ทิวา กาลํ กโรนฺตสฺส อนฺตรายายา”ติ.\csromanspacer{} เตน ภิกฺขเว ภิกฺขุนา เตสํเยว ปาปกานํ อกุสลานํ ธมฺมานํ ปหานาย อธิมตฺโต ฉนฺโท จ วายาโม จ อุสฺสาโห จ อุสฺโสฬฺหี จ อปฺปฏิวานี จ สติ จ สมฺปชญฺญญฺจ กรณียํ.}

\prose{เสยฺยถาปิ ภิกฺขเว อาทิตฺตเจโล วา อาทิตฺตสีโส วา ตสฺเสว เจลสฺส วา สีสสฺส วา นิพฺพาปนาย อธิมตฺตํ ฉนฺทญฺจ วายามญฺจ อุสฺสาหญฺจ อุสฺโสฬฺหิญฺจ อปฺปฏิวานิญฺจ สติญฺจ สมฺปชญฺญญฺจ กเรยฺย.\csromanspacer{} เอวเมวํ โข ภิกฺขเว เตน ภิกฺขุนา เตสํเยว ปาปกานํ อกุสลานํ ธมฺมานํ ปหานาย อธิมตฺโต ฉนฺโท จ วายาโม จ อุสฺสาโห จ อุสฺโสฬฺหี จ อปฺปฏิวานี จ สติ จ สมฺปชญฺญญฺจ กรณียํ.}

\prose{สเจ ปน ภิกฺขเว ภิกฺขุ ปจฺจเวกฺขมาโน เอวํ ชานาติ “นตฺถิ เมปาปกา อกุสลา ธมฺมา อปฺปหีนา, เย เม อสฺสุ ทิวา กาลํ กโรนฺตสฺส อนฺตรายายา”ติ.\csromanspacer{} เตน ภิกฺขเว ภิกฺขุนา เตนว ปีติปาโมชฺเชน วิหาตพฺพํ อโหรตฺตานุสิกฺขินา กุสเลสุ ธมฺเมสุ.\csromanspacer{} เอวํ ภาวิตา โข ภิกฺขเว มรณสฺสติ เอวํ พหุลีกตา มหปฺผลา โหติ มหานิสํสา อมโตคธา อมตปริโยสานาติ.\csromanspacer{}\csromanspacer{}จตุตฺถํ.}

\csromanheader{2}{5. ปฐมสมฺปทาสุตฺต}
\proseitem{75}{อฏฺฐิมา ภิกฺขเว สมฺปทา.\csromanspacer{} กตมา อฏฺฐ? * อุฏฺฐานสมฺปทา อารกฺขสมฺปทา กลฺยาณมิตฺตตา สมชีวิตา สทฺธาสมฺปทา สีลสมฺปทา จาคสมฺปทา ปญฺญาสมฺปทา.\csromanspacer{} อิมา โข ภิกฺขเว อฏฺฐ สมฺปทาติ.}

\prose{อุฏฺฐาตา กมฺมเธยฺเยสุ, อปฺปมตฺโต วิธานวา.}

\prose{สมํ กปฺเปติ ชีวิกํ, สมฺภตํ อนุรกฺขติ.}

\prose{สทฺโธ สีเลน สมฺปนฺโน, วทญฺญู วีตมจฺฉโร.}

\prose{นิจฺจํ มคฺคํ วิโสเธติ, โสตฺถานํ สมฺปรายิกํ.}

\prose{อิจฺเฉเต อฏฺฐ ธมฺมา จ, สทฺธสฺส ฆรเมสิโน.}

\prose{อกฺขาตา สจฺจนาเมน, อุภยตฺถ สุขาวหา.}

\prose{ทิฏฺฐธมฺมหิตตฺถาย, สมฺปรายสุขาย จ.}

\prose{เอวเมตํ คหฏฺฐานํ, จาโค ปุญฺญํ ปวฑฺฒตีติ.\csromanspacer{}\csromanspacer{}ปญฺจมํ.}

\chapterheadpageread{(8) 3. ยมกวคฺค \textbf{6. ทุติยสมฺปทาสุตฺต}}
\proseitem{76}{\csromanfolio{141}อฏฺฐิมา ภิกฺขเว สมฺปทา.\csromanspacer{} กตมา อฏฺฐ? อุฏฺฐานสมฺปทา อารกฺขสมฺปทา กลฺยาณมิตฺตตา สมชีวิตา สทฺธาสมฺปทา สีลสมฺปทา จาคสมฺปทา ปญฺญาสมฺปทา.\csromanspacer{} กตมา จ ภิกฺขเว อุฏฺฐานสมฺปทา? อิธ ภิกฺขเว กุลปุตฺโต เยน กมฺมฏฺฐาเนน ชีวิตํ กปฺเปติ ยทิ กสิยา ยทิ วณิชฺชาย ยทิ โครกฺเขน ยทิ อิสฺสตฺเตน ยทิ ราชโปริเสน ยทิ สิปฺปญฺญตเรน.\csromanspacer{} ตตฺถ ทกฺโข โหติ อนลโส ตตฺรุปายาย วีมํสาย สมนฺนาคโต อลํ กาตุํ อลํ สํวิธาตุนฺติ.\csromanspacer{} อยํ วุจฺจติ ภิกฺขเว อุฏฺฐานสมฺปทา.}

\prose{กตมา จ ภิกฺขเว อารกฺขสมฺปทา? อิธ ภิกฺขเว กุลปุตฺตสฺส โภคา โหนฺติ อุฏฺฐานวีริยาธิคตา พาหาพลปริจิตา เสทาวกฺขิตฺตา ธมฺมิกา ธมฺมลทฺธา, เต อารกฺเขน คุตฺติยา สมฺปาเทติ “กินฺติเม โภเค เนว ราชาโน หเรยฺยุํ, น โจรา หเรยฺยุํ, น อคฺคิ ฑเหยฺย, น อุทกํ วเหยฺย, น อปฺปิยา ทายาทา หเรยฺยุนฺ”ติ.\csromanspacer{} อยํ วุจฺจติ ภิกฺขเว อารกฺขสมฺปทา.}

\prose{กตมา จ ภิกฺขเว กลฺยาณมิตฺตตา? อิธ ภิกฺขเว กุลปุตฺโต ยสฺมิํ คาเม วา นิคเม วา ปฏิวสติ, ตตฺถ เย เต โหนฺติ คหปตี วา คหปติปุตฺตา วา ทหรา วา วุทฺธสีลิโน วุทฺธา วา วุทฺธสีลิโน สทฺธาสมฺปนฺนา สีลสมฺปนฺนา จาคสมฺปนฺนา ปญฺญาสมฺปนฺนา, เตหิ สทฺธิํ สนฺติฏฺฐติ สลฺลปติ สากจฺฉํ สมาปชฺชติ, ยถารูปานํ สทฺธาสมฺปนฺนานํ สทฺธาสมฺปทํ อนุสิกฺขติ, ยถารูปานํ สีลสมฺปนฺนานํ สีลสมฺปทํ อนุสิกฺขติ, ยถารูปานํ จาคสมฺปนฺนานํ จาคสมฺปทํ อนุสิกฺขติ, ยถารูปานํ ปญฺญาสมฺปนฺนานํ ปญฺญาสมฺปทํ อนุสิกฺขติ.\csromanspacer{} อยํ วุจฺจติ ภิกฺขเว กลฺยาณมิตฺตตา.}

\prose{กตมา จ ภิกฺขเว สมชีวิตา? อิธ ภิกฺขเว กุลปุตฺโต อายญฺจ โภคานํ วิทิตฺวา วยญฺจ โภคานํ วิทิตฺวา สมํ ชีวิกํ กปฺเปติ นาจฺโจคาฬฺหํ นาติหีนํ “เอวํ เม อาโย วยํ ปริยาทาย ฐสฺสติ, น จ เม วโย อายํ ปริยาทาย ฐสฺสตี”ติ.\csromanspacer{} เสยฺยถาปิ ภิกฺขเว ตุลาธาโร วา ตุลาธารนฺเตวาสี วา ตุลํ ปคฺคเหตฺวา ชานาติ “เอตฺตเกน วา โอนตํ เอตฺตเกน วา อุนฺนตนฺ”ติ.\csromanspacer{} เอวเมวํ โข ภิกฺขเว กุลปุตฺโต อายญฺจ โภคานํ วิทิตฺวา วยญฺจ โภคานํ วิทิตฺวา สมํ ชีวิกํ กปฺเปติ นาจฺโจคาฬฺหํ นาติหีนํ “เอวํ เม อาโย วยํ ปริยาทาย ฐสฺสติ, น จ เม วโย อายํ ปริยาทาย ฐสฺสตี”ติ.\csromanspacer{} สจายํ ภิกฺขเว \csromanfolio{142}กุลปุตฺโต อปฺปาโย สมาโน อุฬารํ ชีวิกํ กปฺเปติ, ตสฺส ภวนฺติ วตฺตาโร “อุทุมฺพรขาทีวายํ กุลปุตฺโต โภเค ขาทตี”ติ.\csromanspacer{} สเจ ปนายํ ภิกฺขเว กุลปุตฺโต มหาโย สมาโน กสิรํ ชีวิกํ กปฺเปติ, ตสฺส ภวนฺติ วตฺตาโร “อเชฏฺฐมรณํวายํ กุลปุตฺโต มริสฺสตี”ติ.\csromanspacer{} ยโต จ โขยํ ภิกฺขเว กุลปุตฺโต อายญฺจ โภคานํ วิทิตฺวา วยญฺจ โภคานํ วิทิตฺวา สมํ ชีวิกํ กปฺเปติ นาจฺโจคาฬฺหํ นาติหีนํ “เอวํ เม อาโย วยํ ปริยาทาย ฐสฺสติ, น จ เม วโย อายํ ปริยาทาย ฐสฺสตี”ติ.\csromanspacer{} อยํ วุจฺจติ ภิกฺขเว สมชีวิตา.}

\prose{กตมา จ ภิกฺขเว สทฺธาสมฺปทา? อิธ ภิกฺขเว กุลปุตฺโต สทฺโธ โหติ, สทฺทหติ ตถาคตสฺส โพธิํ “อิติปิ โส ภควา ฯเปฯ สตฺถา เทวมนุสฺสานํ พุทฺโธ ภควา”ติ.\csromanspacer{} อยํ วุจฺจติ ภิกฺขเว สทฺธาสมฺปทา.}

\prose{กตมา จ ภิกฺขเว สีลสมฺปทา? อิธ ภิกฺขเว กุลปุตฺโต ปาณาติปาตา ปฏิวิรโต โหติ ฯเปฯ สุราเมรยมชฺชปมาทฏฺฐานา ปฏิวิรโต โหติ.\csromanspacer{} อยํ วุจฺจติ ภิกฺขเว สีลสมฺปทา.}

\prose{กตมา จ ภิกฺขเว จาคสมฺปทา? อิธ ภิกฺขเว กุลปุตฺโต วิคตมลมจฺเฉเรน เจตสา อคารํ อชฺฌาวสติ ฯเปฯ ยาจโยโค ทานสํวิภาครโต.\csromanspacer{} อยํ วุจฺจติ ภิกฺขเว จาคสมฺปทา.}

\prose{กตมา จ ภิกฺขเว ปญฺญาสมฺปทา? อิธ ภิกฺขเว กุลปุตฺโต ปญฺญวา โหติ ฯเปฯ สมฺมา ทุกฺขกฺขยคามินิยา.\csromanspacer{} อยํ วุจฺจติ ภิกฺขเว ปญฺญาสมฺปทา.\csromanspacer{} อิมา โข ภิกฺขเว อฏฺฐ สมฺปทาติ.}

\prose{อุฏฺฐาตา กมฺมเธยฺเยสุ, อปฺปมตฺโต วิธานวา.}

\prose{สมํ กปฺเปติ ชีวิกํ, สมฺภตํ อนุรกฺขติ.}

\prose{สทฺโธ สีเลน สมฺปนฺโน, วทญฺญู วีตมจฺฉโร.}

\prose{นิจฺจํ มคฺคํ วิโสเธติ, โสตฺถานํ สมฺปรายิกํ.}

\prose{อิจฺเจเต อฏฺฐ ธมฺมา จ, สทฺธสฺส ฆรเมสิโน.}

\prose{อกฺขาตา สจฺจนาเมน, อุภยตฺถ สุขาวหา.}

\prose{ทิฏฺฐธมฺมหิตตฺถาย, สมฺปรายสุขาย จ.}

\prose{เอวเมตํ คหฏฺฐานํ, จาโค ปุญฺญํ ปวฑฺฒตีติ.\csromanspacer{}\csromanspacer{}ฉฏฺฐํ.}

\chapterheadpageread{(8) 3. ยมกวคฺค \textbf{7. อิจฺฉาสุตฺต}}
\proseitem{77}{\csromanfolio{143}ตตฺร โข อายสฺมา สาริปุตฺโต ภิกฺขู อามนฺเตสิ "อาวุโส ภิกฺขโว”ติ. “อาวุโส”ติ โข เต ภิกฺขู อายสฺมโต สาริปุตฺตสฺส ปจฺจสฺโสสุํ.\csromanspacer{} อายสฺมา สาริปุตฺโต เอตทโวจ\csromandash{}}

\prose{\csromansymbolfootnote{*}{เหฏฺฐา 116 ปิฏฺเฐปิ.}อฏฺฐิเม อาวุโส ปุคฺคลา สนฺโต สํวิชฺชมานา โลกสฺมิํ.\csromanspacer{} กตเม อฏฺฐ? อิธาวุโส ภิกฺขุโน ปวิวิตฺตสฺส วิหรโต นิรายตฺตวุตฺติโน อิจฺฉา อุปฺปชฺชติ ลาภาย, โส อุฏฺฐหติ ฆฏติ วายมติ ลาภาย, ตสฺส อุฏฺฐหโต ฆฏโต วายมโต ลาภาย ลาโภ นุปฺปชฺชติ.\csromanspacer{} โส เตน อลาเภน โสจติ กิลมติ ปริเทวติ อุรตฺตาฬิํ กนฺทติ สมฺโมหํ อาปชฺชติ.\csromanspacer{} อยํ วุจฺจตาวุโส ภิกฺขุ อิจฺโฉ วิหรติ ลาภาย, อุฏฺฐหติ ฆฏติ วายมติ ลาภาย, น จ ลาภี โสจี จ ปริเทวี จ จุโต จ สทฺธมฺมา.}

\prose{อิธ ปนาวุโส ภิกฺขุโน ปวิวิตฺตสฺส วิหรโต นิรายตฺตวุตฺติโน อิจฺฉา อุปฺปชฺชติ ลาภาย, โส อุฏฺฐหติ ฆฏติ วายมติ ลาภาย, ตสฺส อุฏฺฐหโต ฆฏโต วายมโต ลาภาย ลาโภ อุปฺปชฺชติ.\csromanspacer{} โส เตน ลาเภน มชฺชติ ปมชฺชติ ปมาทมาปชฺชติ.\csromanspacer{} อยํ วุจฺจตาวุโส ภิกฺขุ อิจฺโฉ วิหรติ ลาภาย, อุฏฺฐหติ ฆฏติ วายมติ ลาภาย, ลาภี จ มที จ ปมาที จ จุโต จ สทฺธมฺมา.}

\prose{อิธ ปนาวุโส ภิกฺขุโน ปวิวิตฺตสฺส วิหรโต นิรายตฺตวุตฺติโน อิจฺฉา อุปฺปชฺชติ ลาภาย, โส น อุฏฺฐหติ น ฆฏติ น วายมติ ลาภาย, ตสฺส อนุฏฺฐหโต อฆฏโต อวายมโต ลาภาย ลาโภ นุปฺปชฺชติ.\csromanspacer{} โส เตน อลาเภน โสจติ กิลมติ ปริเทวติ อุรตฺตาฬิํ กนฺทติ สมฺโมหํ อาปชฺชติ.\csromanspacer{} อยํ วุจฺจตาวุโส ภิกฺขุ อิจฺโฉ วิหรติ ลาภาย, น อุฏฺฐหติ น ฆฏติ น วายมติ ลาภาย, น จ ลาภี โสจี จ ปริเทวี จ จุโต จ สทฺธมฺมา.}

\prose{อิธ ปนาวุโส ภิกฺขุโน ปวิวิตฺตสฺส วิหรโต นิรายตฺตวุตฺติโน อิจฺฉา อุปฺปชฺชติ ลาภาย, โส น อุฏฺฐหติ น ฆฏติ น วายมติ ลาภาย, ตสฺส อนุฏฺฐหโต อฆฏโต อวายมโต ลาภาย ลาโภ อุปฺปชฺชติ.\csromanspacer{} โส เตน ลาเภน มชฺชติ ปมชฺชติ ปมาทมาปชฺชติ.\csromanspacer{} อยํ วุจฺจตาวุโส ภิกฺขุ อิจฺโฉ วิหรติ ลาภาย, น อุฏฺฐหติ \csromanfolio{144}น ฆฏติ น วายมติ ลาภาย, ลาภี จ มที จ ปมาที จ จุโต จ สทฺธมฺมา.}

\prose{อิธ ปนาวุโส ภิกฺขุโน ปวิวิตฺตสฺส วิหรโต นิรายตฺตวุตฺติโน อิจฺฉา อุปฺปชฺชติ ลาภาย, โส อุฏฺฐหติ ฆฏติ วายมติ ลาภาย, ตสฺส อุฏฺฐหโต ฆฏโต วายมโต ลาภาย ลาโภ นุปฺปชฺชติ.\csromanspacer{} โส เตน อลาเภน น โสจติ น กิลมติ น ปริเทวติ น อุรตฺตาฬิํ กนฺทติ น สมฺโมหํ อาปชฺชติ.\csromanspacer{} อยํ วุจฺจตาวุโส ภิกฺขุ อิจฺโฉ วิหรติ ลาภาย, อุฏฺฐหติ ฆฏติ วายมติ ลาภาย, น จ ลาภี น จ โสจี น จ ปริเทวี อจฺจุโต จ สทฺธมฺมา.}

\prose{อิธ ปนาวุโส ภิกฺขุโน ปวิวิตฺตสฺส วิหรโต นิรายตฺตวุตฺติโน อิจฺฉา อุปฺปชฺชติ ลาภาย, โส อุฏฺฐหติ ฆฏติ วายมติ ลาภาย, ตสฺส อุฏฺฐหโต ฆฏโต วายมโต ลาภาย ลาโภ อุปฺปชฺชติ.\csromanspacer{} โส เตน ลาเภน น มชฺชติ น ปมชฺชติ น ปมาทมาปชฺชติ.\csromanspacer{} อยํ วุจฺจตาวุโส ภิกฺขุ อิจฺโฉ วิหรติ ลาภาย, อุฏฺฐหติ ฆฏติ วายมติ ลาภาย, ลาภี จ น จ มที น จ ปมาที อจฺจุโต จ สทฺธมฺมา.}

\prose{อิธ ปนาวุโส ภิกฺขุโน ปวิวิตฺตสฺส วิหรโต นิรายตฺตวุตฺติโน อิจฺฉา อุปฺปชฺชติ ลาภาย, โส น อุฏฺฐหติ น ฆฏติ น วายมติ ลาภาย, ตสฺส อนุฏฺฐหโต อฆฏโต อวายมโต ลาภาย ลาโภ นุปฺปชฺชติ.\csromanspacer{} โส เตน อลาเภน น โสจติ น กิลมติ น ปริเทวติ น อุรตฺตาฬิํ กนฺทติ น สมฺโมหํ อาปชฺชติ.\csromanspacer{} อยํ วุจฺจตาวุโส ภิกฺขุ อิจฺโฉ วิหรติ ลาภาย, น อุฏฺฐหติ น ฆฏติ น วายมติ ลาภาย, น จ ลาภี น จ โสจี น จ ปริเทวี อจฺจุโต จ สทฺธมฺมา.}

\prose{อิธ ปนาวุโส ภิกฺขุโน ปวิวิตฺตสฺส วิหรโต นิรายตฺตวุตฺติโน อิจฺฉา อุปฺปชฺชติ ลาภาย, โส น อุฏฺฐหติ น ฆฏติ น วายมติ ลาภาย, ตสฺส อนุฏฺฐหโต อฆฏโต อวายมโต ลาภาย ลาโภ อุปฺปชฺชติ.\csromanspacer{} โส เตน ลาเภน น มชฺชติ น ปมชฺชติ น ปมาทมาปชฺชติ.\csromanspacer{} อยํ วุจฺจตาวุโส ภิกฺขุ อิจฺโฉ วิหรติ ลาภาย, น อุฏฺฐหติ น ฆฏติ น วายมติ ลาภาย, ลาภี จ น จ มที น จ ปมาที อจฺจุโต จ สทฺธมฺมา.\csromanspacer{} อิเม โข อาวุโส อฏฺฐ ปุคฺคลา สนฺโต สํวิชฺชมานา โลกสฺมินฺติ.\csromanspacer{}\csromanspacer{}สตฺตมํ.}

\chapterheadpageread{(8) 3. ยมกวคฺค \textbf{8. อลํสุตฺต}}
\proseitem{78}{\csromanfolio{145}\csromansymbolfootnote{*}{เหฏฺฐา 118 ปิฏฺเฐปิ.}ตตฺร โข อายสฺมา สาริปุตฺโต ภิกฺขู อามนฺเตสิ ฯเปฯ.\csromanspacer{} ฉหาวุโส ธมฺเมหิ สมนฺนาคโต ภิกฺขุ อลํ อตฺตโน อลํ ปเรสํ.\csromanspacer{} กตเมหิ ฉหิ? อิธาวุโส ภิกฺขุ ขิปฺปนิสนฺติ จ โหติ กุสเลสุ ธมฺเมสุ.\csromanspacer{} สุตานญฺจ ธมฺมานํ ธารณชาติโก โหติ.\csromanspacer{} ธาตานญฺจ ธมฺมานํ อตฺถูปปริกฺขิตา โหติ.\csromanspacer{} อตฺถมญฺญาย ธมฺมมญฺญาย ธมฺมานุธมฺมปฺปฏิปนฺโน จ โหติ.\csromanspacer{} กลฺยาณวาโจ จ โหติ กลฺยาณวากฺกรโณ โปริยา วาจาย สมนฺนาคโต วิสฺสฏฺฐาย อเนลคฬาย อตฺถสฺส วิญฺญาปนิยา.\csromanspacer{} สนฺทสฺสโก จ โหติ สมาทปโก สมุตฺเตชโก สมฺปหํสโก สพฺรหฺมจารีนํ.\csromanspacer{} อิเมหิ โข อาวุโส ฉหิ ธมฺเมหิ สมนฺนาคโต ภิกฺขุ อลํ อตฺตโน อลํ ปเรสํ.}

\prose{ปญฺจหาวุโส ธมฺเมหิ สมนฺนาคโต ภิกฺขุ อลํ อตฺตโน อลํ ปเรสํ.\csromanspacer{} กตเมหิ ปญฺจหิ? อิธาวุโส ภิกฺขุ น เหว โข ขิปฺปนิสนฺติ จ โหติ กุสเลสุ ธมฺเมสุ.\csromanspacer{} สุตานญฺจ ธมฺมานํ ธารณชาติโก โหติ.\csromanspacer{} ธาตานญฺจ ธมฺมานํ อตฺถูปปริกฺขิตา โหติ.\csromanspacer{} อตฺถมญฺญาย ธมฺมมญฺญาย ธมฺมานุธมฺมปฺปฏิปนฺโน จ โหติ.\csromanspacer{} กลฺยาณวาโจ จ โหติ ฯเปฯ.\csromanspacer{} สนฺทสฺสโก จ โหติ ฯเปฯ สพฺรหฺมจารีนํ.\csromanspacer{} อิเมหิ โข อาวุโส ปญฺจหิ ธมฺเมหิ สมนฺนาคโต ภิกฺขุ อลํ อตฺตโน อลํ ปเรสํ.}

\prose{จตูหาวุโส ธมฺเมหิ สมนฺนาคโต ภิกฺขุ อลํ อตฺตโน นาลํ ปเรสํ.\csromanspacer{} กตเมหิ จตูหิ? อิธาวุโส ภิกฺขุ ขิปฺปนิสนฺติ จ โหติ กุสเลสุ ธมฺเมสุ.\csromanspacer{} สุตานญฺจ ธมฺมานํ ธารณชาติโก โหติ.\csromanspacer{} ธาตานญฺจ ธมฺมานํ อตฺถูปปริกฺขิตา โหติ.\csromanspacer{} อตฺถมญฺญาย ธมฺมมญฺญาย ธมฺมานุธมฺมปฺปฏิปนฺโน จ โหติ.\csromanspacer{} โน จ กลฺยาณวาโจ โหติ ฯเปฯ.\csromanspacer{} โน จ สนฺทสฺสโก โหติ ฯเปฯ สพฺรหฺมจารีนํ.\csromanspacer{} อิเมหิ โข อาวุโส จตูหิ ธมฺเมหิ สมนฺนาคโต ภิกฺขุ อลํ อตฺตโน นาลํ ปเรสํ.}

\prose{จตูหาวุโส ธมฺเมหิ สมนฺนาคโต ภิกฺขุ อลํ ปเรสํ นาลํ อตฺตโน.\csromanspacer{} กตเมหิ จตูหิ? อิธาวุโส ภิกฺขุ ขิปฺปนิสนฺติ จ โหติ กุสเลสุ ธมฺเมสุ.\csromanspacer{} สุตานญฺจ ธมฺมานํ ธารณชาติโก โหติ.\csromanspacer{} โน จ ธาตานํ ธมฺมานํ อตฺถูปปริกฺขิตา โหติ.\csromanspacer{} โน จ อตฺถมญฺญาย ธมฺมมญฺญาย ธมฺมานุธมฺมปฺปฏิปนฺโน โหติ.\csromanspacer{} กลฺยาณวาโจ จ โหติ ฯเปฯ \csromanfolio{146}สนฺทสฺสโก จ โหติ สพฺรหฺมจารีนํ.\csromanspacer{} อิเมหิ โข อาวุโส จตูหิ ธมฺเมหิ สมนฺนาคโต ภิกฺขุ อลํ ปเรสํ นาลํ อตฺตโน.}

\prose{ตีหาวุโส ธมฺเมหิ สมนฺนาคโต ภิกฺขุ อลํ อตฺตโน นาลํ ปเรสํ.\csromanspacer{} กตเมหิ ตีหิ? อิธาวุโส ภิกฺขุ น เหว โข ขิปฺปนิสนฺติ จ โหติ กุสเลสุ ธมฺเมสุ.\csromanspacer{} สุตานญฺจ ธมฺมานํ ธารณชาติโก โหติ.\csromanspacer{} ธาตานญฺจ ธมฺมานํ อตฺถูปปริกฺขิตา โหติ.\csromanspacer{} อตฺถมญฺญาย ธมฺมมญฺญาย ธมฺมานุธมฺมปฺปฏิปนฺโน จ โหติ.\csromanspacer{} โน จ กลฺยาณวาโจ โหติ ฯเปฯ.\csromanspacer{} โน จ สนฺทสฺสโก โหติ ฯเปฯ สพฺรหฺมจารีนํ.\csromanspacer{} อิเมหิ โข อาวุโส ตีหิ ธมฺเมหิ สมนฺนาคโต ภิกฺขุ อลํ อตฺตโน นาลํ ปเรสํ.}

\prose{ตีหาวุโส ธมฺเมหิ สมนฺนาคโต ภิกฺขุ อลํ ปเรสํ นาลํ อตฺตโน.\csromanspacer{} กตเมหิ ตีหิ? อิธาวุโส ภิกฺขุ น เหว โข ขิปฺปนิสนฺติ จ โหติ กุสเลสุ ธมฺเมสุ.\csromanspacer{} สุตานญฺจ ธมฺมานํ ธารณชาติโก โหติ.\csromanspacer{} โน จ ธาตานํ ธมฺมานํ อตฺถูปปริกฺขิตา โหติ.\csromanspacer{} โน จ อตฺถมญฺญาย ธมฺมมญฺญาย ธมฺมานุธมฺมปฺปฏิปนฺโน โหติ.\csromanspacer{} กลฺยาณวาโจ จ โหติ ฯเปฯ อตฺถสฺส วิญฺญาปนิยา.\csromanspacer{} สนฺทสฺสโก จ โหติ สพฺรหฺมจารีนํ.\csromanspacer{} อิเมหิ โข อาวุโส ตีหิ ธมฺเมหิ สมนฺนาคโต ภิกฺขุ อลํ ปเรสํ นาลํ อตฺตโน.}

\prose{ทฺวีหาวุโส ธมฺเมหิ สมนฺนาคโต ภิกฺขุ อลํ อตฺตโน นาลํ ปเรสํ.\csromanspacer{} กตเมหิ ทฺวีหิ? อิธาวุโส ภิกฺขุ น เหว โข ขิปฺปนิสนฺติ จ โหติ กุสเลสุ ธมฺเมสุ.\csromanspacer{} โน จ สุตานํ ธมฺมานํ ธารณชาติโก โหติ.\csromanspacer{} ธาตานญฺจ ธมฺมานํ อตฺถูปปริกฺขิตา โหติ.\csromanspacer{} อตฺถมญฺญาย ธมฺมมญฺญาย ธมฺมานุธมฺมปฺปฏิปนฺโน จ โหติ.\csromanspacer{} โน จ กลฺยาณวาโจ โหติ ฯเปฯ.\csromanspacer{} โน จ สนฺทสฺสโก โหติ ฯเปฯ สพฺรหฺมจารีนํ.\csromanspacer{} อิเมหิ โข อาวุโส ทฺวีหิ ธมฺเมหิ สมนฺนาคโต ภิกฺขุ อลํ อตฺตโน นาลํ ปเรสํ.}

\prose{ทฺวีหาวุโส ธมฺเมหิ สมนฺนาคโต ภิกฺขุ อลํ ปเรสํ นาลํ อตฺตโน.\csromanspacer{} กตเมหิ ทฺวีหิ? อิธาวุโส ภิกฺขุ น เหว โข ขิปฺปนิสนฺติ จ โหติ กุสเลสุ ธมฺเมสุ.\csromanspacer{} โน จ สุตานํ ธมฺมานํ ธารณชาติโก โหติ.\csromanspacer{} โน จ ธาตานํ ธมฺมานํ อตฺถูปปริกฺขิตา โหติ.\csromanspacer{} โน จ อตฺถมญฺญาย ธมฺมมญฺญาย ธมฺมานุธมฺมปฺปฏิปนฺโน โหติ.\csromanspacer{} กลฺยาณวาโจ จ โหติ}

\vspace{\tipitakaparskip}
\csromancenter{(8) 3.\csromanspacer{} ยมกวคฺค กลฺยาณวากฺกรโณ โปริยา วาจาย สมนฺนาคโต วิสฺสฏฺฐาย อเนลคฬาย อตฺถสฺส วิญฺญาปนิยา.\csromanspacer{} สนฺทสฺสโก จ โหติ สมาทปโก สมุตฺเตชโก สมฺปหํสโก สพฺรหฺมจารีนํ.\csromanspacer{} อิเมหิ โข อาวุโส ทฺวีหิ ธมฺเมหิ สมนฺนาคโต ภิกฺขุ อลํ ปเรสํ นาลํ อตฺตโนติ.\csromanspacer{}\csromanspacer{}อฏฺฐมํ.}
\csromanheader{2}{9. ปริหานสุตฺต}
\proseitem{79}{\csromanfolio{147}อฏฺฐิเม ภิกฺขเว ธมฺมา เสขสฺส ภิกฺขุโน ปริหานาย สํวตฺตนฺติ.\csromanspacer{} กตเม อฏฺฐ? กมฺมารามตา ภสฺสารามตา นิทฺทารามตา สงฺคณิการามตา อินฺทฺริเยสุ อคุตฺตทฺวารตา โภชเน อมตฺตญฺญุตา สํสคฺคารามตา ปปญฺจารามตา.\csromanspacer{} อิเม โข ภิกฺขเว อฏฺฐ ธมฺมา เสขสฺส ภิกฺขุโน ปริหานาย สํวตฺตนฺติ.}

\prose{อฏฺฐิเม ภิกฺขเว ธมฺมา เสขสฺส ภิกฺขุโน อปริหานาย สํวตฺตนฺติ.\csromanspacer{} กตเม อฏฺฐ? น กมฺมารามตา น ภสฺสารามตา น นิทฺทารามตา น สงฺคณิการามตา อินฺทฺริเยสุ คุตฺตทฺวารตา โภชเน มตฺตญฺญุตา อสํสคฺคารามตา นิปฺปปญฺจารามตา.\csromanspacer{} อิเม โข ภิกฺขเว อฏฺฐ ธมฺมา เสขสฺส ภิกฺขุโน อปริหานาย สํวตฺตนฺตีติ.\csromanspacer{}\csromanspacer{}นวมํ.}

\csromanheader{2}{10. กุสีตารมฺภวตฺถุสุตฺต}
\proseitem{80}{\csromansymbolfootnote{*}{ที 3. 211, 246 ปิฏฺเฐสุปิ.}อฏฺฐิมานิ ภิกฺขเว กุสีตวตฺถูนิ.\csromanspacer{} กตมานิ อฏฺฐ? อิธ ภิกฺขเว ภิกฺขุนา กมฺมํ กตฺตพฺพํ โหติ.\csromanspacer{} ตสฺส เอวํ โหติ “กมฺมํ โข เม กตฺตพฺพํ ภวิสฺสติ, กมฺมํ โข ปน เม กโรนฺตสฺส กาโย กิลมิสฺสติ, หนฺทาหํ นิปชฺชามี”ติ.\csromanspacer{} โส นิปชฺชติ, น วีริยํ อารภติ อปฺปตฺตสฺส ปตฺติยา อนธิคตสฺส อธิคมาย อสจฺฉิกตสฺส สจฺฉิกิริยาย.\csromanspacer{} อิทํ ภิกฺขเว ปฐมํ กุสีตวตฺถุ.}

\prose{ปุน จปรํ ภิกฺขเว ภิกฺขุนา กมฺมํ กตํ โหติ.\csromanspacer{} ตสฺส เอวํ โหติ “อหํ โข กมฺมํ อกาสิํ, กมฺมํ โข ปน เม กโรนฺตสฺส กาโย กิลนฺโต, หนฺทาหํ นิปชฺชามี”ติ.\csromanspacer{} โส นิปชฺชติ, น วีริยํ อารภติ อปฺปตฺตสฺส ปตฺติยา อนธิคตสฺส อธิคมาย อสจฺฉิกตสฺส สจฺฉิกิริยาย.\csromanspacer{} อิทํ ภิกฺขเว ทุติยํ กุสีตวตฺถุ.}

\prose{\csromanfolio{148}ปุน จปรํ ภิกฺขเว ภิกฺขุนา มคฺโค คนฺตพฺโพ โหติ.\csromanspacer{} ตสฺส เอวํ โหติ “มคฺโค เม คนฺตพฺโพ ภวิสฺสติ, มคฺคํ โข ปน เม คจฺฉนฺตสฺส กาโย กิลมิสฺสติ, หนฺทาหํ นิปชฺชามี”ติ.\csromanspacer{} โส นิปชฺชติ, น วีริยํ อารภติ อปฺปตฺตสฺส ปตฺติยา อนธิคตสฺส อธิคมาย อสจฺฉิกตสฺส สจฺฉิกิริยาย.\csromanspacer{} อิทํ ภิกฺขเว ตติยํ กุสีตวตฺถุ.}

\prose{ปุน จปรํ ภิกฺขเว ภิกฺขุนา มคฺโค คโต โหติ.\csromanspacer{} ตสฺส เอวํ โหติ “อหํ โข มคฺคํ อคมาสิํ, มคฺคํ โข ปน เม คจฺฉนฺตสฺส กาโย กิลนฺโต, หนฺทาหํ นิปชฺชามี”ติ.\csromanspacer{} โส นิปชฺชติ, น วีริยํ อารภติ อปฺปตฺตสฺส ปตฺติยา อนธิคตสฺส อธิคมาย อสจฺฉิกตสฺส สจฺฉิกิริยาย.\csromanspacer{} อิทํ ภิกฺขเว จตุตฺถํ กุสีตวตฺถุ.}

\prose{ปุน จปรํ ภิกฺขเว ภิกฺขุ คามํ วา นิคมํ วา ปิณฺฑาย จรนฺโต น ลภติ ลูขสฺส วา ปณีตสฺส วา โภชนสฺส ยาวทตฺถํ ปาริปูริํ.\csromanspacer{} ตสฺส เอวํ โหติ “อหํ โข คามํ วา นิคมํ วา ปิณฺฑาย จรนฺโต นาลตฺถํ ลูขสฺส วา ปณีตสฺส วา โภชนสฺส ยาวทตฺถํ ปาริปูริํ, ตสฺส เม กาโย กิลนฺโต อกมฺมญฺโญ, หนฺทาหํ นิปชฺชามี”ติ.\csromanspacer{} โส นิปชฺชติ, น วีริยํ อารภติ ฯเปฯ.\csromanspacer{} อิทํ ภิกฺขเว ปญฺจมํ กุสีตวตฺถุ.}

\prose{ปุน จปรํ ภิกฺขเว ภิกฺขุ คามํ วา นิคมํ วา ปิณฺฑาย จรนฺโต ลภติ ลูขสฺส วา ปณีตสฺส วา โภชนสฺส ยาวทตฺถํ ปาริปูริํ.\csromanspacer{} ตสฺส เอวํ โหติ “อหํ โข คามํ วา นิคมํ วา ปิณฺฑาย จรนฺโต อลตฺถํ ลูขสฺส วา ปณีตสฺส วา โภชนสฺส ยาวทตฺถํ ปาริปูริํ, ตสฺส เม กาโย ครุโก อกมฺมญฺโญ มาสาจิตํ มญฺเญ, หนฺทาหํ นิปชฺชามี”ติ.\csromanspacer{} โส นิปชฺชติ, น วีริยํ อารภติ ฯเปฯ.\csromanspacer{} อิทํ ภิกฺขเว ฉฏฺฐํ กุสีตวตฺถุ.}

\prose{ปุน จปรํ ภิกฺขเว ภิกฺขุโน อุปฺปนฺโน โหติ อปฺปมตฺตโก อาพาโธ.\csromanspacer{} ตสฺส เอวํ โหติ “อุปฺปนฺโน โข เม อยํ อปฺปมตฺตโก อาพาโธ, อตฺถิ กปฺโป นิปชฺชิตุํ, หนฺทาหํ นิปชฺชามี”ติ.\csromanspacer{} โส นิปชฺชติ, น วีริยํ อารภติ ฯเปฯ.\csromanspacer{} อิทํ ภิกฺขเว สตฺตมํ กุสีตวตฺถุ.}

\prose{ปุน จปรํ ภิกฺขเว ภิกฺขุ คิลานา วุฏฺฐิโต\footnote{อํ 2. 263 ปิฏฺเฐ สุตฺตวณฺณนา ฏีกา โอโลเกตพฺพา.} โหติ, อจิรวุฏฺฐิโต เคลญฺญา.\csromanspacer{} ตสฺส เอวํ โหติ “อหํ โข คิลานา วุฏฺฐิโต, อจิรวุฏฺฐิโต เคลญฺญา, ตสฺส เม กาโย ทุพฺพโล อกมฺมญฺโญ, หนฺทาหํ}

\vspace{\tipitakaparskip}
\csromancenter{(8) 3.\csromanspacer{} ยมกวคฺค นิปชฺชามี”ติ.\csromanspacer{} โส นิปชฺชติ, น วีริยํ อารภติ อปฺปตฺตสฺส ปตฺติยา อนธิคตสฺส อธิคมาย อสจฺฉิกตสฺส สจฺฉิกิริยาย.\csromanspacer{} อิทํ ภิกฺขเว อฏฺฐมํ กุสีตวตฺถุ.\csromanspacer{} อิมานิ โข ภิกฺขเว อฏฺฐ กุสีตวตฺถูนิ.}
\prose{\csromanfolio{149}\csromansymbolfootnote{*}{ที 3. 212, 247 ปิฏฺเฐสุปิ.}อฏฺฐิมานิ ภิกฺขเว อารมฺภวตฺถูนิ.\csromanspacer{} กตมานิ อฏฺฐ? อิธ ภิกฺขเว ภิกฺขุนา กมฺมํ กตฺตพฺพํ โหติ.\csromanspacer{} ตสฺส เอวํ โหติ “กมฺมํ โข เม กตฺตพฺพํ ภวิสฺสติ, กมฺมํ โข มยา กโรนฺเตน น สุกรํ พุทฺธานํ สาสนํ มนสิ กาตุํ, หนฺทาหํ ปฏิกจฺเจว วีริยํ อารภามิ อปฺปตฺตสฺส ปตฺติยา อนธิคตสฺส อธิคมาย อสจฺฉิกตสฺส สจฺฉิกิริยายา”ติ.\csromanspacer{} โส วีริยํ อารภติ อปฺปตฺตสฺส ปตฺติยา อนธิคตสฺส อธิคมาย อสจฺฉิกตสฺส สจฺฉิกิริยาย.\csromanspacer{} อิทํ ภิกฺขเว ปฐมํ อารมฺภวตฺถุ.}

\prose{ปุน จปรํ ภิกฺขเว ภิกฺขุนา กมฺมํ กตํ โหติ.\csromanspacer{} ตสฺส เอวํ โหติ “อหํ โข กมฺมํ อกาสิํ, กมฺมํ โข ปนาหํ กโรนฺโต นาสกฺขิํ พุทฺธานํ สาสนํ มนสิ กาตุํ, หนฺทาหํ วีริยํ อารภามิ อปฺปตฺตสฺส ปตฺติยา อนธิคตสฺส อธิคมาย อสจฺฉิกตสฺส สจฺฉิกิริยายา”ติ.\csromanspacer{} โส วีริยํ อารภติ.\csromanspacer{} อิทํ ภิกฺขเว ทุติยํ อารมฺภวตฺถุ.}

\prose{ปุน จปรํ ภิกฺขเว ภิกฺขุนา มคฺโค คนฺตพฺโพ โหติ.\csromanspacer{} ตสฺส เอวํ โหติ “มคฺโค โข เม คนฺตพฺโพ ภวิสฺสติ, มคฺคํ โข ปน เม คจฺฉนฺเตน น สุกรํ พุทฺธานํ สาสนํ มนสิ กาตุํ, หนฺทาหํ วีริยํ ฯเปฯ”.\csromanspacer{} อิทํ ภิกฺขเว ตติยํ อารมฺภวตฺถุ.}

\prose{ปุน จปรํ ภิกฺขเว ภิกฺขุนา มคฺโค คโต โหติ.\csromanspacer{} ตสฺส เอวํ โหติ “อหํ โข มคฺคํ อคมาสิํ, มคฺคํ โข ปนาหํ คจฺฉนฺโต นาสกฺขิํ พุทฺธานํ สาสนํ มนสิ กาตุํ, หนฺทาหํ วีริยํ อารภามิ ฯเปฯ”.\csromanspacer{} อิทํ ภิกฺขเว จตุตฺถํ อารมฺภวตฺถุ.}

\prose{ปุน จปรํ ภิกฺขเว ภิกฺขุ คามํ วา นิคมํ วา ปิณฺฑาย จรนฺโต น ลภติ ลูขสฺส วา ปณีตสฺส วา โภชนสฺส ยาวทตฺถํ ปาริปูริํ.\csromanspacer{} ตสฺส เอวํ โหติ “อหํ โข คามํ วา นิคมํ วา ปิณฺฑาย จรนฺโต นาลตฺถํ ลูขสฺส วา ปณีตสฺส วา โภชนสฺส ยาวทตฺถํ ปาริปูริํ, ตสฺส เม กาโย ลหุโก กมฺมญฺโญ, หนฺทาหํ วีริยํ อารภามิ ฯเปฯ”.\csromanspacer{} อิทํ ภิกฺขเว ปญฺจมํ อารมฺภวตฺถุ.}

\prose{\csromanfolio{150}ปุน จปรํ ภิกฺขเว ภิกฺขุ คามํ วา นิคมํ วา ปิณฺฑาย จรนฺโต ลภติ ลูขสฺส วา ปณีตสฺส วา โภชนสฺส ยาวทตฺถํ ปาริปูริํ.\csromanspacer{} ตสฺส เอวํ โหติ “อหํ โข คามํ วา นิคมํ วา ปิณฺฑาย จรนฺโต อลตฺถํ ลูขสฺส วา ปณีตสฺส วา โภชนสฺส ยาวทตฺถํ ปาริปูริํ, ตสฺส เม กาโย พลวา กมฺมญฺโญ, หนฺทาหํ วีริยํ อารภามิ ฯเปฯ”.\csromanspacer{} อิทํ ภิกฺขเว ฉฏฺฐํ อารมฺภวตฺถุ.}

\prose{ปุน จปรํ ภิกฺขเว ภิกฺขุโน อุปฺปนฺโน โหติ อปฺปมตฺตโก อาพาโธ.\csromanspacer{} ตสฺส เอวํ โหติ “อุปฺปนฺโน โข เม อยํ อปฺปมตฺตโก อาพาโธ.\csromanspacer{} ฐานํ โข ปเนตํ วิชฺชติ, ยํ เม อาพาโธ ปวฑฺเฒยฺย, หนฺทาหํ ปฏิกจฺเจว วีริยํ อารภามิ ฯเปฯ”.\csromanspacer{} อิทํ ภิกฺขเว สตฺตมํ อารมฺภวตฺถุ.}

\prose{ปุน จปรํ ภิกฺขเว ภิกฺขุ คิลานา วุฏฺฐิโต โหติ, อจิรวุฏฺฐิโต เคลญฺญา, ตสฺส เอวํ โหติ “อหํ โข คิลานา วุฏฺฐิโต, อจิรวุฏฺฐิโต เคลญฺญา.\csromanspacer{} ฐานํ โข ปเนตํ วิชฺชติ, ยํ เม อาพาโธ ปจฺจุทาวตฺเตยฺย, หนฺทาหํ ปฏิกจฺเจว วีริยํ อารภามิ อปฺปตฺตสฺส ปตฺติยา อนธิคตสฺส อธิคมาย อสจฺฉิกตสฺส สจฺฉิกิริยายา”ติ.\csromanspacer{} โส วีริยํ อารภติ อปฺปตฺตสฺส ปตฺติยา อนธิคตสฺส อธิคมาย อสจฺฉิกตสฺส สจฺฉิกิริยาย.\csromanspacer{} อิทํ ภิกฺขเว อฏฺฐมํ อารมฺภวตฺถุ.\csromanspacer{} อิมานิ โข ภิกฺขเว อฏฺฐ อารมฺภวตฺถูนีติ.\csromanspacer{}\csromanspacer{}ทสมํ.\csromanspacer{} ยมกวคฺโค ตติโย.}

\titlehead{\textbf{ตสฺสุทฺทานํ}}
\csromangathameasure{เทฺว สทฺธา เทฺว มรณสฺสตี, เทฺว สมฺปทา อถาปเร. \\ อิจฺฉา อลํ ปริหานํ, กุสีตารมฺภวตฺถูนีติ.}
\csromangathagroup{\csromangathastanza{เทฺว สทฺธา เทฺว มรณสฺสตี, เทฺว สมฺปทา อถาปเร. \\ อิจฺฉา อลํ ปริหานํ, กุสีตารมฺภวตฺถูนีติ.}}

\chapterhead{\textbf{(9) 4. สติวคฺค}}
\csromanheader{2}{1. สติสมฺปชญฺญสุตฺต}
\proseitem{81}{สติสมฺปชญฺเญ ภิกฺขเว อสติ สติสมฺปชญฺญวิปนฺนสฺส หตูปนิสํ โหติ หิโรตฺตปฺปํ, หิโรตฺตปฺเป อสติ หิโรตฺตปฺปวิปนฺนสฺส หตูปนิโส โหติ อินฺทฺริยสํวโร, อินฺทฺริยสํวเร อสติ อินฺทฺริยสํวรวิปนฺนสฺส หตูปนิสํ}

\vspace{\tipitakaparskip}
\csromancenter{(9) 4.\csromanspacer{} สติวคฺค โหติ สีลํ, สีเล อสติ สีลวิปนฺนสฺส ปตูปนิโส โหติ สมฺมาสมาธิ, สมฺมาสมาธิมฺหิ อสติ สมฺมาสมาธิวิปนฺนสฺส หตูปนิสํ โหติ ยถาภูตญาณทสฺสนํ, ยถาภูตญาณทสฺสเน อสติ ยถาภูตญาณทสฺสนวิปนฺนสฺส หตูปนิโส โหติ นิพฺพิทาวิโรโค, นิพฺพิทาวิราเค อสติ นิพฺพิทาวิราควิปนฺนสฺส หตูปนิสํ โหติ วิมุตฺติญาณทสฺสนํ.\csromanspacer{} เสยฺยถาปิ ภิกฺขเว รุกฺโข สาขาปลาสวิปนฺโน ตสฺส ปปฏิกาปิ น ปาริปูริํ คจฺฉติ, ตโจปิ เผคฺคุปิ สาโรปิ น ปาริปูริํ คจฺฉติ.\csromanspacer{} เอวเมวํ โข ภิกฺขเว สติสมฺปชญฺเญ อสติ สติสมฺปชญฺญวิปนฺนสฺส หตูปนิสํ โหติ หิโรตฺตปฺปํ, หิโรตฺตปฺเป อสติ หิโรตฺตปฺปวิปนฺนสฺส หตูปนิโส โหติ ฯเปฯ วิมุตฺติญาณทสฺสนํ.}
\prose{\csromanfolio{151}สติสมฺปชญฺเญ ภิกฺขเว สติ สติสมฺปชญฺญสมฺปนฺนสฺส อุปนิสสมฺปนฺนํ โหติ หิโรตฺตปฺปํ, หิโรตฺตปฺเป สติ หิโรตฺตปฺปสมฺปนฺนสฺส อุปนิสสมฺปนฺโน โหติ อินฺทฺริยสํวโร, อินฺทฺริยสํวเร สติ อินฺทฺริยสํวรสมฺปนฺนสฺส อุปนิสสมฺปนฺนํ โหติ สีลํ, สีเล สติ สีลสมฺปนฺนสฺส อุปนิสสมฺปนฺโน โหติ สมฺมาสมาธิ, สมฺมาสมาธิมฺหิ สติ สมฺมาสมาธิสมฺปนฺนสฺส อุปนิสสมฺปนฺนํ โหติ ยถาภูตญาณทสฺสนํ, ยถาภูตญาณทสฺสเน สติ ยถาภูตญาณทสฺสนสมฺปนฺนสฺส อุปนิสสมฺปนฺโน โหติ นิพฺพิทาวิราโค, นิพฺพิทาวิราเค สติ นิพฺพิทาวิราคสมฺปนฺนสฺส อุปนิสสมฺปนํ โหติ วิมุตฺติญาณทสฺสนํ.\csromanspacer{} เสยฺยถาปิ ภิกฺขเว รุกฺโข สาขาปลาสสมฺปนฺโน ตสฺส ปปฏิกาปิ ปาริปูริํ คจฺฉติ, ตโจปิ เผคฺคุปิ สาโรปิ ปาริปูริํ คจฺฉติ.\csromanspacer{} เอวเมวํ โข ภิกฺขเว สติสมฺปชญฺเญ สติ สติสมฺปชญฺญสมฺปนฺนสฺส อุปนิสสมฺปนฺนํ โหติ หิโรตฺตปฺปํ, หิโรตฺตปฺเป สติ หิโรตฺตปฺปสมฺปนฺนสฺส อุปนิสสมฺปนฺโน โหติ ฯเปฯ วิมุตฺติญาณทสฺสนนฺติ.\csromanspacer{}\csromanspacer{}ปฐมํ.}

\csromanheader{2}{2. ปุณฺณิยสุตฺต}
\proseitem{82}{อถ โข อายสฺมา ปุณฺณิโย เยน ภควา เตนุปสงฺกมิ, อุปสงฺกมิตฺวา ภควนฺตํ อภิวาเทตฺวา เอกมนฺตํ นิสีทิ, เอกมนฺตํ นิสินฺโน โข อายสฺมา ปุณฺณิโย ภควนฺตํ เอตทโวจ “โก นุ โข ภนฺเต เหตุ โก ปจฺจโย, เยน อปฺเปกทา ตถาคตํ ธมฺมเทสนา ปฏิภาติ, อปฺเปกทา น ปฏิภาตี”ติ.\csromanspacer{} สทฺโธ จ ปุณฺณิย ภิกฺขุ โหติ โน จุปสงฺกมิตา, เนว ตถาคตํ ธมฺมเทสนา ปฏิภาติ.\csromanspacer{} ยโต จ \csromanfolio{152}โข ปุณฺณิย ภิกฺขุ สทฺโธ จ โหติ อุปสงฺกมิตา จ, เอวํ ตถาคตํ ธมฺมเทสนา ปฏิภาติ.\csromanspacer{} สทฺโธ จ ปุณฺณิย ภิกฺขุ โหติ อุปสงฺกมิตา จ, โน จ ปยิรุปาสิตา ฯเปฯ ปยิรุปาสิตา จ, โน จ ปริปุจฺฉิตา.\csromanspacer{} ปริปุจฺฉิตา จ, โน จ โอหิตโสโต ธมฺมํ สุณาติ.\csromanspacer{} โอหิตโสโต จ ธมฺมํ สุณาติ, โน จ สุตฺวา ธมฺมํ ธาเรติ.\csromanspacer{} สุตฺวา จ ธมฺมํ ธาเรติ, โน จ ธาตานํ ธมฺมานํ อตฺถํ อุปปริกฺขติ.\csromanspacer{} ธาตานญฺจ ธมฺมานํ อตฺถํ อุปปริกฺขติ, โน จ อตฺถมญฺญาย ธมฺมมญฺญาย ธมฺมานุธมฺมปฺปฏิปนฺโน โหติ, เนว ตาว ตถาคตํ ธมฺมเทสนา ปฏิภาติ.}

\prose{ยโต จ โข ปุณฺณิย ภิกฺขุ สทฺโธ จ โหติ อุปสงฺกมิตา จ ปยิรุปาสิตา จ ปริปุจฺฉิตา จ, โอหิตโสโต จ ธมฺมํ สุณาติ สุตฺวา จ ธมฺมํ ธาเรติ, ธาตานญฺจ ธมฺมานํ อตฺถํ อุปปริกฺขติ, อตฺถมญฺญาย ธมฺมมญฺญาย ธมฺมานุธมฺมปฺปฏิปนฺโน จ โหติ, เอวํ ตถาคตํ ธมฺมเทสนา ปฏิภาติ.\csromanspacer{} อิเมหิ โข ปุณฺณิย อฏฺฐหิ ธมฺเมหิ สมนฺนาคตา\footnote{สมนฺนาคโต (สี, อิ), สมนฺนาคตํ (สฺยา, ก)} เอกนฺตปฏิภานา\footnote{เอกนฺตปฏิภานํ (สพฺพตฺถ) อุปริ 380 ปิฏฺเฐ ปน ปสฺสิตพฺพํ.} ตถาคตํ ธมฺมเทสนา โหตีติ.\csromanspacer{}\csromanspacer{}ทุติยํ.}

\csromanheader{2}{3. มูลกสุตฺต}
\proseitem{83}{\csromansymbolfootnote{*}{อุปริ 190, 340 ปิฏฺเฐสุปิ ปสฺสิตพฺพํ.}สเจ ภิกฺขเว อญฺญติตฺถิยา ปริพฺพาชกา เอวํ ปุจฺเฉยฺยุํ “กิํมูลกา อาวุโส สพฺเพ ธมฺมา, กิํสมฺภวา สพฺเพ ธมฺมา, กิํสมุทยา สพฺเพ ธมฺมา, กิํสโมสรณา สพฺเพ ธมฺมา, กิํปมุขา สพฺเพ ธมฺมา, กิํ อธิปเตยฺยา สพฺเพ ธมฺมา, กิํอุตฺตรา สพฺเพ ธมฺมา, กิํสารา สพฺเพ ธมฺมา”ติ.\csromanspacer{} เอวํ ปุฏฺฐา ตุเมฺห ภิกฺขเว เตสํ อญฺญติตฺถิยานํ ปริพฺพาชกานํ กินฺติ พฺยากเรยฺยาถาติ.\csromanspacer{} ภควํมูลกา โน ภนฺเต ธมฺมา, ภควํเนตฺติกา ภควํปฏิสรณา.\csromanspacer{} สาธุ ภนฺเต ภควนฺตํเยว ปฏิภาตุ เอตสฺส ภาสิตสฺส อตฺโถ, ภควโต สุตฺวา ภิกฺขู ธาเรสฺสนฺตีติ.\csromanspacer{} เตน หิ ภิกฺขเว เทเสสฺสามิ, ตํ สุณาถ สาธุกํ มนสิ กโรถ, ภาสิสฺสามีติ. “เอวํ ภนฺเต”ติ โข เต ภิกฺขู ภควโต ปจฺจสฺโสสุํ.\csromanspacer{} ภควา เอตทโวจ\csromandash{}สเจ ภิกฺขเว อญฺญติตฺถิยา ปริพฺพาชกา เอวํ ปุจฺเฉยฺยุํ “กิํ มูลกา อาวุโส สพฺเพ ธมฺมา, กิํสมฺภวา สพฺเพ ธมฺมา, กิํสมุทยา สพฺเพ ธมฺมา, กิํสโมสรณา สพฺเพ ธมฺมา.\csromanspacer{} กิํปมุขา สพฺเพ}

\vspace{\tipitakaparskip}
\csromancenter{(9) 4.\csromanspacer{} สติวคฺค ธมฺมา, กิํ อธิปเตยฺยา สพฺเพ ธมฺมา, กิํ อุตฺตรา สพฺเพ ธมฺมา, กิํสารา สพฺเพ ธมฺมา”ติ.\csromanspacer{} เอวํ ปุฏฺฐา ตุเมฺห ภิกฺขเว เตสํ อญฺญติตฺถิยานํ ปริพฺพาชกานํ เอวํ พฺยากเรยฺยาถ “ฉนฺทมูลกา อาวุโส สพฺเพ ธมฺมา, มนสิการสมฺภวา สพฺเพ ธมฺมา, ผสฺสสมุทยา สพฺเพ ธมฺมา, เวทนาสโมสรณา สพฺเพ ธมฺมา, สมาธิปฺปมุขา สพฺเพ ธมฺมา, สตาธิปเตยฺยา สพฺเพ ธมฺมา, ปญฺญุตฺตรา สพฺเพ ธมฺมา, วิมุตฺติสารา สพฺเพ ธมฺมา”ติ.\csromanspacer{} เอวํ ปุฏฺฐา ตุเมฺห ภิกฺขเว เตสํ อญฺญติตฺถิยานํ ปริพฺพาชกานํ เอวํ พฺยากเรยฺยาถาติ.\csromanspacer{}\csromanspacer{}ตติยํ.}
\csromanheader{2}{4. โจรสุตฺต}
\proseitem{84}{\csromanfolio{153}อฏฺฐหิ ภิกฺขเว องฺเคหิ สมนฺนาคโต มหาโจโร ขิปฺปํ ปริยาปชฺชติ, น จิรฏฺฐิติโก โหติ.\csromanspacer{} กตเมหิ อฏฺฐหิ? อปฺปหรนฺตสฺส ปหรติ, อนวเสสํ อาทิยติ, อิตฺถิํ หนติ, กุมาริํ ทูเสติ, ปพฺพชิตํ วิลุมฺปติ, ราชโกสํ วิลุมฺปติ, อจฺจาสนฺเน กมฺมํ กโรติ, น จ นิธานกุสโล โหติ.\csromanspacer{} อิเมหิ โข ภิกฺขเว อฏฺฐหงฺเคหิ สมนฺนาคโต มหาโจโร ขิปฺปํ ปริยาปชฺชติ, น จิรฏฺฐิติโก โหติ.}

\prose{อฏฺฐหิ ภิกฺขเว องฺเคหิ สมนฺนาคโต มหาโจโร น ขิปฺปํ ปริยาปชฺชติ, จิรฏฺฐิติโก โหติ.\csromanspacer{} กตเมหิ อฏฺฐหิ? น อปฺปหรนฺตสฺส ปหรติ, น อนวเสสํ อาทิยติ, น อิตฺถิํ หนติ, น กุมาริํ ทูเสติ, น ปพฺพชิตํ วิลุมฺปติ, น ราชโกสํ วิลุมฺปติ, น อจฺจาสนฺเน กมฺมํ กโรติ, นิธานกุสโล จ โหติ.\csromanspacer{} อิเมหิ โข ภิกฺขเว อฏฺฐหงฺเคหิ สมนฺนาคโต มหาโจโร น ขิปฺปํ ปริยาปชฺชติ, จิรฏฺฐิติโก โหตีติ.\csromanspacer{}\csromanspacer{}จตุตฺถํ.}

\csromanheader{2}{5. สมณสุตฺต}
\proseitem{85}{“สมโณ”ติ ภิกฺขเว ตถาคตสฺเสตํ อธิวจนํ อรหโต สมฺมาสมฺพุทฺธสฺส. “พฺราหฺมโณ”ติ ภิกฺขเว ตถาคตสฺเสตํ อธิวจนํ อรหโต สมฺมาสมฺพุทฺธสฺส. “เวทคู”ติ ภิกฺขเว ตถาคตสฺเสตํ อธิวจนํ อรหโต สมฺมาสมฺพุทฺธสฺส. “ภิสกฺโก”ติ ภิกฺขเว ตถาคตสฺเสตํ อธิวจนํ อรหโต สมฺมาสมฺพุทฺธสฺส. “นิมฺมโล”ติ ภิกฺขเว ตถาคตสฺเสตํ อธิวจนํ อรหโต สมฺมาสมฺพุทฺธสฺส. “วิมโล”ติ ภิกฺขเว ตถาคตสฺเสตํ อธิวจนํ อรหโต สมฺมาสมฺพุทฺธสฺส. “ญาณี”ติ ภิกฺขเว ตถาคตสฺเสตํ อธิวจนํ \csromanfolio{154}อรหโต สมฺมาสมฺพุทฺธสฺส. “วิมุตฺโต”ติ ภิกฺขเว ตถาคตสฺเสตํ อธิวจนํ อรหโต สมฺมาสมฺพุทฺธสฺสาติ.}

\csromangathameasure{ยํ สมเณน ปตฺตพฺพํ, พฺราหฺมเณน วุสีมตา. \\ ยํ เวทคุนา ปตฺตพฺพํ, ภิสกฺเกน อนุตฺตรํ. \\ ยํ นิมฺมเลน ปตฺตพฺพํ, วิมเลน สุจีมตา. \\ ยํ ญาณินา จ ปตฺตพฺพํ, วิมุตฺเตน อนุตฺตรํ. \\ โสหํ วิชิตสงฺคาโม, มุตฺโต โมเจมิ พนฺธนา. \\ นาโคมฺหิ ปรมทนฺโต, อเสโข ปรินิพฺพุโตติ.ปญฺจมํ.}
\csromangathagroup{\csromangathastanza{ยํ สมเณน ปตฺตพฺพํ, พฺราหฺมเณน วุสีมตา. \\ ยํ เวทคุนา ปตฺตพฺพํ, ภิสกฺเกน อนุตฺตรํ.}\csromangathastanza{ยํ นิมฺมเลน ปตฺตพฺพํ, วิมเลน สุจีมตา. \\ ยํ ญาณินา จ ปตฺตพฺพํ, วิมุตฺเตน อนุตฺตรํ.}\csromangathastanza{โสหํ วิชิตสงฺคาโม, มุตฺโต โมเจมิ พนฺธนา. \\ นาโคมฺหิ ปรมทนฺโต, อเสโข ปรินิพฺพุโตติ.\csromanspacer{}\csromanspacer{}ปญฺจมํ.}}

\csromanheader{2}{6. ยสสุตฺต}
\proseitem{86}{เอกํ สมยํ ภควา โกสเลสุ จาริกํ จรมาโน มหตา ภิกฺขุสํเฆน สทฺธิํ เยน อิจฺฉานงฺคลํ นาม โกสลานํ พฺราหฺมณคาโม ตทวสริ, ตตฺร สุทํ ภควา อิจฺฉานงฺคเล วิหรติ อิจฺฉานงฺคลวนสณฺเฑ.\csromanspacer{} อสฺโสสุํ โข อิจฺฉานงฺคลกา พฺราหฺมณคหปติกา “สมโณ ขลุ โภ โคตโม สกฺยปุตฺโต สกฺยกุลา ปพฺพชิโต อิจฺฉานงฺคลํ อนุปฺปตฺโต อิจฺฉานงฺคเล วิหรติ อิจฺฉานงฺคลวนสณฺเฑ.\csromanspacer{} ตํ โข ปน ภวนฺตํ โคตมํ เอวํ กลฺยาโณ กิตฺติสทฺโท อพฺภุคฺคโต อิติปิ โส ภควา อรหํ สมฺมาสมฺพุทฺโธ ฯเปฯ.\csromanspacer{} สาธุ โข ปน ตถารูปานํ อรหตํ ทสฺสนํ โหตี”ติ.}

\prose{อถ โข อิจฺฉานงฺคลกา พฺราหฺมณคหปติกา ตสฺสา รตฺติยา อจฺจเยน ปหุตํ ขาทนียํ โภชนียํ อาทาย เยน อิจฺฉานงฺคลวนสณฺโฑ เตนุปสงฺกมิํสุ, อุปสงฺกมิตฺวา พหิทฺวารโกฏฺฐเก อฏฺฐํสุ อุจฺจาสทฺทา มหาสทฺทา.\csromanspacer{} เตน โข ปน สมเยน อายสฺมา นาคิโต ภควโต อุปฏฺฐาโก โหติ.\csromanspacer{} อถ โข ภควา อายสฺมนฺตํ นาคิตํ อามนฺเตสิ “เก ปน เต นาคิต อุจฺจาสทฺทา มหาสทฺทา เกวฏฺฏา มญฺเญ มจฺฉวิโลเป”ติ.\csromanspacer{} เอเต ภนฺเต อิจฺฉานงฺคลกา พฺราหฺมณคหปติกา ปหุตํ ขาทนียํ โภชนียํ อาทาย พหิทฺวารโกฏฺฐเก ฐิตา ภควนฺตํเยว อุทฺทิสฺส ภิกฺขุ- สํฆญฺจาติ.\csromanspacer{} มาหํ นาคิต ยเสน สมาคมํ, มา จ มยา ยโส.\csromanspacer{} โย โข นาคิต นยิมสฺส เนกฺขมฺมสุขสฺส ปวิเวกสุขสฺส อุปสมสุขสฺส สมฺโพธสุขสฺส นิกามลาภี อสฺส อกิจฺฉลาภี อกสิรลาภี, ยสฺสาหํ เนกฺขมฺมสุขสฺส ปวิเวกสุขสฺส อุปสมสุขสฺส}

\vspace{\tipitakaparskip}
\csromancenter{(9) 4.\csromanspacer{} สติวคฺค สมฺโพธสุขสฺส นิกามลาภี\footnote{นิกามลาภี อสฺสํ (พหูสุ) อํ 2. 25 ปิฏฺเฐ ปสฺสิตพฺพํ. ตตฺถ หิ อยํ ปาฐเภโท นตฺถิ.} อกิจฺฉลาภี อกสิรลาภี, โส ตํ มีฬฺหสุขํ มิทฺธสุขํ ลาภสกฺการสิโลกสุขํ สาทิเยยฺยาติ.}
\prose{\csromanfolio{155}อธิวาเสตุ ทานิ ภนฺเต ภควา, อธิวาเสตุ สุคโต, อธิวาสนกาโล ทานิ ภนฺเต ภควโต.\csromanspacer{} เยน เยเนว ทานิ ภนฺเต ภควา คมิสฺสติ, ตนฺนินฺนาว ภวิสฺสนฺติ พฺราหฺมณคหปติกา เนคมา เจว ชานปทา จ.\csromanspacer{} เสยฺยถาปิ ภนฺเต ถุลฺลผุสิตเก เทเว วสฺสนฺเต ยถานินฺนํ อุทกานิ ปวตฺตนฺติ.\csromanspacer{} เอวเมวํ โข ภนฺเต เยน เยเนว ทานิ ภควา คมิสฺสติ, ตนฺนินฺนาว ภวิสฺสนฺติ พฺราหฺมณคหปติกา เนคมา เจว ชานปทา จ.\csromanspacer{} ตํ กิสฺส เหตุ? ตถา หิ ภนฺเต ภควโต สีลปญฺญาณนฺติ.}

\prose{มาหํ นาคิต ยเสน สมาคมํ, มา จ มยา ยโส.\csromanspacer{} โย โข นาคิต นยิมสฺส เนกฺขมฺมสุขสฺส ปวิเวกสุขสฺส อุปสมสุขสฺส สมฺโพธสุขสฺส นิกามลาภี อสฺส อกิจฺฉลาภี อกสิรลาภี, ยสฺสาหํ เนกฺขมฺมสุขสฺส ปวิเวกสุขสฺส อุปสมสุขสฺส สมฺโพธสุขสฺส นิกามลาภี อกิจฺฉลาภี อกสิรลาภี, โส ตํ มีฬฺหสุขํ มิทฺธสุขํ ลาภสกฺการสิโลกสุขํ สาทิเยยฺย.}

\prose{เทวตาปิ โข นาคิต เอกจฺจา นยิมสฺส\footnote{เอกจฺจา อิมสฺส (?)} เนกฺขมฺมสุขสฺส ปวิเวกสุขสฺส อุปสมสุขสฺส สมฺโพธสุขสฺส นิกามลาภินิโย อสฺสุ\footnote{อิทํ ปทํ กตฺถจิ นตฺถิ.} อกิจฺฉลาภินิโย\footnote{นิกามลาภินิโย อกิจฺฉลาภินิโย (?)} อกสิรลาภินิโย, ยสฺสาหํ เนกฺขมฺมสุขสฺส ปวิเวกสุขสฺส อุปสมสุขสฺส สมฺโพธสุขสฺส นิกามลาภี อกิจฺฉลาภี อกสิรลาภี.\csromanspacer{} ตุมฺหากมฺปิ\footnote{ตาสมฺปิ (?)} โข นาคิต สงฺคมฺม สมาคมฺม สงฺคณิกวิหารํ อนุยุตฺตานํ วิหรตํ\footnote{อนุยุตฺเต วิหรนฺเต ทิสฺวา (?)} เอวํ โหติ “น หิ นูนเม\footnote{น หนูนเม (สี, สฺยา, อิ)} อายสฺมนฺโต อิมสฺส เนกฺขมฺมสุขสฺส ปวิเวกสุขสฺส อุปสมสุขสฺส สมฺโพธสุขสฺส นิกามลาภิโน อสฺสุ3 อกิจฺฉลาภิโน อกสิรลาภิโน, ยสฺสาหํ เนกฺขมฺมสุขสฺส ปวิเวกสุขสฺส อุปสมสุขสฺส สมฺโพธสุขสฺส นิกามลาภี อกิจฺฉลาภี อกสิรลาภี.\csromanspacer{} ตถา \csromanfolio{156}หิ ปน’เม อายสฺมนฺโต สงฺคมฺม สมาคมฺม สงฺคณิกวิหารํ อนุยุตฺตา วิหรนฺติ”.}

\prose{อิธาหํ นาคิต ภิกฺขู ปสฺสามิ อญฺญมญฺญํ องฺคุลิปโตทเกน\footnote{องฺคุลิปโตทเกหิ (สี, อิ, ก)} สญฺชคฺฆนฺเต สํกีฬนฺเต.\csromanspacer{} ตสฺส มยฺหํ นาคิต เอวํ โหติ “น หิ นูนเม อายสฺมนฺโต อิมสฺส เนกฺขมฺมสุขสฺส ปวิเวกสุขสฺส อุปสมสุขสฺส สมฺโพธสุขสฺส นิกามลาภิโน อสฺสุ อกิจฺฉลาภิโน อกสิรลาภิโน, ยสฺสาหํ เนกฺขมฺมสุขสฺส ปวิเวกสุขสฺส อุปสมสุขสฺส สมฺโพธสุขสฺส นิกามลาภี อกิจฺฉลาภี อกสิรลาภี.\csromanspacer{} ตถา หิ ปน’เม อายสฺมนฺโต อญฺญมญฺญํ องฺคุลิปโตทเกน สญฺชคฺฆนฺติ สํกีฬนฺติ”. (1)}

\prose{อิธ ปนาหํ\footnote{อิธาหํ (สี, อิ, ก)} นาคิต ภิกฺขู ปสฺสามิ ยาวทตฺถํ อุทราวเทหกํ ภุญฺชิตฺวา เสยฺยสุขํ ปสฺสสุขํ มิทฺธสุขํ อนุยุตฺเต วิหรนฺเต.\csromanspacer{} ตสฺส มยฺหํ นาคิต เอวํ โหติ “น หิ นูนเม อายสฺมนฺโต อิมสฺส เนกฺขมฺมสุขสฺส ปวิเวกสุขสฺส อุปสมสุขสฺส สมฺโพธสุขสฺส นิกามลาภิโน อสฺสุ อกิจฺฉลาภิโน อกสิรลาภิโน, ยสฺสาหํ เนกฺขมฺมสุขสฺส ปวิเวกสุขสฺส อุปสมสุขสฺส สมฺโพธสุขสฺส นิกามลาภี อกิจฺฉลาภี อกสิรลาภี.\csromanspacer{} ตถา หิ ปน’เม อายสฺมนฺโต ยาวทตฺถํ อุทราวเทหกํ ภุญฺชิตฺวา เสยฺยสุขํ ปสฺสสุขํ มิทฺธสุขํ อนุยุตฺตา วิหรนฺติ”. (2)}

\prose{อิธาหํ\footnote{อิธ ปนาหํ(?)} นาคิต ภิกฺขุํ ปสฺสามิ คามนฺตวิหาริํ สมาหิตํ นิสินฺนํ.\csromanspacer{} ตสฺส มยฺหํ นาคิต เอวํ โหติ “อิทานิ อิมํ\footnote{อิทานิมํ (กตฺถจิ) อํ 2. 301 ปิฏฺเฐปิ.} อายสฺมนฺตํ อารามิโก วา อุปฏฺฐหิสฺสติ\footnote{ปจฺเจสฺสติ (สี, อิ), อุปฏฺฐหติ (ก)} สมณุทฺเทโส วา, ตํ ตมฺหา\footnote{โส ตมฺหา (ก-สี), โส ตํ ตมฺหา (?)} สมาธิมฺหา จาเวสฺสตี”ติ.\csromanspacer{} เตนาหํ นาคิต ตสฺส ภิกฺขุโน น อตฺถมโน โหมิ คามนฺตวิหาเรน. (3)}

\prose{อิธ ปนาหํ นาคิต ภิกฺขุํ ปสฺสามิ อารญฺญิกํ อรญฺเญ ปจลายมานํ นิสินฺนํ.\csromanspacer{} ตสฺส มยฺหํ นาคิต เอวํ โหติ “อิทานิ อยมายสฺมา อิมํ นิทฺทากิลมถํ ปฏิวิโนเทตฺวา อรญฺญสญฺญํเยว มนสิ กริสฺสติ เอกตฺตนฺ”ติ.}

\vspace{\tipitakaparskip}
\csromancenter{(9) 4.\csromanspacer{} สติวคฺค เตนาหํ นาคิต ตสฺส ภิกฺขุโน อตฺตมโน โหมิ อรญฺญวิหาเรน. (4)}
\prose{\csromanfolio{157}อิธ ปนาหํ นาคิต ภิกฺขุํ ปสฺสามิ อารญฺญิกํ อรญฺเญ อสมาหิตํ นิสินฺนํ.\csromanspacer{} ตสฺส มยฺหํ นาคิต เอวํ โหติ “อิทานิ อยมายสฺมา อสมาหิตํ วา จิตฺตํ สมาทหิสฺสติ\footnote{สมาทเหสฺสติ (กตฺถจิ) 2-2. [ ] เอตฺถนฺตเร ปาโฐ อํ 2. 301 ปิฏฺเฐ ฉกฺกนิปาเตเยว ทิสฺสติ, น เอตฺถ อฏฺฐกนิปาเต.}, สมาหิตํ วา จิตฺตํ อนุรกฺขิสฺสตี”ติ.\csromanspacer{} เตนาหํ นาคิต ตสฺส ภิกฺขุโน อตฺตมโน โหมิ อรญฺญวิหาเรน. (5)}

\prose{อิธ ปนาหํ นาคิต ภิกฺขุํ ปสฺสามิ อารญฺญิกํ อรญฺเญ สมาหิตํ นิสินฺนํ ตสฺส มยฺหํ นาคิต เอวํ โหติ “อิทานิ อยมายสฺมา อวิมุตฺตํ วา จิตฺตํ วิมุจฺจิสฺสติ, วิมุตฺตํ วา จิตฺตํ อนุรกฺขิสฺสตี”ติ.\csromanspacer{} เตนาหํ นาคิต ตสฺส ภิกฺขุโน อตฺตมโน โหมิ อรญฺญวิหาเรน. (6)}

\prose{2[อิธ ปนาหํ นาคิต ภิกฺขุํ ปสฺสามิ คามนฺตวิหาริํ ลาภิํ จีวรปิณฺฑปาต เสนาสน คิลานปจฺจย เภสชฺชปริกฺขารานํ.\csromanspacer{} โส ตํ ลาภสกฺการสิโลกํ นิกามยมาโน ริญฺจติ ปฏิสลฺลานํ, ริญฺจติ อรญฺญวนปตฺถานิ ปนฺตานิ เสนาสนานิ, คามนิคมราชธานิํ โอสริตฺวา วาสํ กปฺเปติ.\csromanspacer{} เตนาหํ นาคิต ตสฺส ภิกฺขุโน น อตฺตมโน โหมิ คามนฺตวิหาเรน. (7)}

\prose{อิธ ปนาหํ นาคิต ภิกฺขุํ ปสฺสามิ อารญฺญิกํ ลาภิํ จีวร ปิณฺฑปาตเสนาสน คิลานปจฺจย เภสชฺชปริกฺขารานํ, โส ตํ ลาภสกฺการสิโลกํ ปฏิปณาเมตฺวา น ริญฺจติ ปฏิสลฺลานํ, น ริญฺจติ อรญฺญวนปตฺถานิ ปนฺตานิ เสนาสนานิ.\csromanspacer{} เตนาหํ นาคิต ตสฺส ภิกฺขุโน อตฺตมโน โหมิ อรญฺญวิหาเรน. (8)]2}

\prose{ยสฺมาหํ\footnote{ยสฺมิํ หํ (กตฺถจิ)} นาคิต สมเย อทฺธานมคฺคปฺปฏิปนฺโน น กญฺจิ ปสฺสามิ ปุรโต วา ปจฺฉโต วา, ผาสุ เม นาคิต ตสฺมิํ สมเย โหติ อนฺตมโส อุจฺจารปสฺสาวกมฺมายาติ.\csromanspacer{}\csromanspacer{}ฉฏฺฐํ.}

\csromanheader{2}{7. ปตฺตนิกุชฺชนสุตฺต}
\proseitem{87}{\csromanfolio{158}\csromansymbolfootnote{*}{วิ 4. 263 ปิฏฺเฐปิ.}อฏฺฐหิ ภิกฺขเว องฺเคหิ สมนฺนาคตสฺส อุปาสกสฺส อากงฺขมาโน สํโฆ ปตฺตํ นิกฺกุชฺเชยฺย\footnote{นิกุชฺเชยฺย (ก)}.\csromanspacer{} กตเมหิ อฏฺฐหิ? ภิกฺขูนํ อลาภาย ปริสกฺกติ, ภิกฺขูนํ อนตฺถาย ปริสกฺกติ, ภิกฺขูนํ อวาสาย\footnote{อนาวาสาย (สี, สฺยา)} ปริสกฺกติ, ภิกฺขู อกฺโกสติ ปริภาสติ, ภิกฺขู ภิกฺขูหิ เภเทติ\footnote{วิเภเทติ (พหูสุ)}, พุทฺธสฺส อวณฺณํ ภาสติ, ธมฺมสฺส อวณฺณํ ภาสติ, สํฆสฺส อวณฺณํ ภาสติ.\csromanspacer{} อิเมหิ โข ภิกฺขเว อฏฺฐหงฺเคหิ สมนฺนาคตสฺส อุปาสกสฺส อากงฺขมาโน สํโฆ ปตฺตํ นิกุชฺเชยฺย.}

\prose{อฏฺฐหิ ภิกฺขเว องฺเคหิ สมนฺนาคตสฺส อุปาสกสฺส อากงฺขมาโน สํโฆ ปตฺตํ อุกฺกุชฺเชยฺย.\csromanspacer{} กตเมหิ อฏฺฐหิ? น ภิกฺขูนํ อลาภาย ปริสกฺกติ, น ภิกฺขูนํ อนตฺถาย ปริสกฺกติ, น ภิกฺขูนํ อวาสาย ปริสกฺกติ, น ภิกฺขู อกฺโกสติ ปริภาสสติ, น ภิกฺขู ภิกฺขูหิ เภเทติ, พุทฺธสฺส วณฺณํ ภาสติ, ธมฺมสฺส วณฺณํ ภาสติ, สํฆสฺส วณฺณํ ภาสติ.\csromanspacer{} อิเมหิ โข ภิกฺขเว อฏฺฐหงฺเคหิ สมนฺนาคตสฺส อุปาสกสฺส อากงฺขมาโน สํโฆ ปตฺตํ อุกฺกุชฺเชยฺยาติ.\csromanspacer{}\csromanspacer{}สตฺตมํ.}

\csromanheader{2}{8. อปฺปสาทปเวทนียสุตฺต}
\proseitem{88}{อฏฺฐหิ ภิกฺขเว ธมฺเมหิ สมนฺนาคตสฺส ภิกฺขุโน อากงฺขมาโน อุปาสกา อปฺปสาทํ ปเวเทยฺยุํ.\csromanspacer{} กตเมหิ อฏฺฐหิ? คิหีนํ อลาภาย ปริสกฺกติ, คิหีนํ อนตฺถาย ปริสกฺกติ, คิหี อกฺโกสติ ปริภาสติ, คิหี คิหีหิ เภเทติ, พุทฺธสฺส อวณฺณํ ภาสติ, ธมฺมสฺส อวณฺณํ ภาสติ, สํฆสฺส อวณฺณํ ภาสติ, อโคจเร จ นํ ปสฺสนฺติ.\csromanspacer{} อิเมหิ โข ภิกฺขเว อฏฺฐหิ ธมฺเมหิ สมนฺนาคตสฺส ภิกฺขุโน อากงฺขมานา อุปาสกา อปฺปสาทํ ปเวเทยฺยุํ.}

\prose{อฏฺฐหิ ภิกฺขเว ธมฺเมหิ สมนฺนาคตสฺส ภิกฺขุโน อากงฺขมานา อุปาสกา ปสาทํ ปเวเทยฺยุํ.\csromanspacer{} กตเมหิ อฏฺฐหิ? น คิหีนํ อลาภาย ปริสกฺกติ, น คิหีนํ อนตฺถาย ปริสกฺกติ, น คิหี อกฺโกสติ ปริภาสติ, น คิหี คิหีหิ เภเทติ, พุทฺธสฺส วณฺณํ ภาสติ, ธมฺมสฺส วณฺณํ ภาสติ, สํฆสฺส วณฺณํ ภาสติ, โคจเร จ นํ ปสฺสนฺติ.\csromanspacer{} อิเมหิ โข}

\vspace{\tipitakaparskip}
\csromancenter{(9) 4.\csromanspacer{} สติวคฺค ภิกฺขเว อฏฺฐหิ ธมฺเมหิ สมนฺนาคตสฺส ภิกฺขุโน อากงฺขมานา อุปาสกา ปสาทํ ปเวเทยฺยุนฺติ.\csromanspacer{}\csromanspacer{}อฏฺฐมํ.}
\csromanheader{2}{9. ปฏิสารณียสุตฺต}
\proseitem{89}{\csromanfolio{159}\csromansymbolfootnote{*}{วิ 4. 42 ปิฏฺเฐปิ โถกํ วิสทิสํ.}อฏฺฐหิ ภิกฺขเว ธมฺเมหิ สมนฺนาคตสฺส ภิกฺขุโน อากงฺขมาโน สํโฆ ปฏิสารณียกมฺมํ กเรยฺย.\csromanspacer{} กตเมหิ อฏฺฐหิ? คิหีนํ อลาภาย ปริสกฺกติ, คิหีนํ อนตฺถาย ปริสกฺกติ, คิหี อกฺโกสติ ปริภาสติ, คิหี คิหีหิ เภเทติ, พุทฺธสฺส อวณฺณํ ภาสติ, ธมฺมสฺส อวณฺณํ ภาสติ, สํฆสฺส อวณฺณํ ภาสติ, ธมฺมิกญฺจ คิหิปฏิสฺสวํ น สจฺจาเปติ.\csromanspacer{} อิเมหิ โข ภิกฺขเว อฏฺฐหิ ธมฺเมหิ สมนฺนาคตสฺส ภิกฺขุโน อากงฺขมาโน สํโฆ ปฏิสารณียํ กมฺมํ กเรยฺย.}

\prose{อฏฺฐหิ ภิกฺขเว ธมฺเมหิ สมนฺนาคตสฺส ภิกฺขุโน อากงฺขมาโน สํโฆ ปฏิสารณียกมฺมํ ปฏิปฺปสฺสมฺเภยฺย.\csromanspacer{} กตเมหิ อฏฺฐหิ? น คิหีนํ อลาภาย ปริสกฺกติ, น คิหีนํ อนตฺถาย ปริสกฺกติ, น คิหี อกฺโกสติ ปริภาสติ, น คิหี คิหีหิ เภเทติ, พุทฺธสฺส วณฺณํ ภาสติ, ธมฺมสฺส วณฺณํ ภาสติ, สํฆสฺส วณฺณํ ภาสติ, ธมฺมิกญฺจ คิหิปฏิสฺสวํ สจฺจาเปติ.\csromanspacer{} อิเมหิ โข ภิกฺขเว อฏฺฐหิ ธมฺเมหิ สมนฺนาคตสฺส ภิกฺขุโน อากงฺขมาโน สํโฆ ปฏิสารณียกมฺมํ ปฏิปฺปสฺสมฺเภยฺยาติ.\csromanspacer{}\csromanspacer{}นวมํ.}

\csromanheader{2}{10. สมฺมาวตฺตนสุตฺต}
\proseitem{90}{\csromansymbolfootnote{+}{วิ 4. 208 ปิฏฺเฐปิ.}ตสฺสปาปิยสิกกมฺมกเตน ภิกฺขเว ภิกฺขุนา อฏฺฐสุ ธมฺเมสุ สมฺมา วตฺติตพฺพํ.\csromanspacer{} น อุปสมฺปาเทตพฺโพ, น นิสฺสโย ทาตพฺโพ, น สามเณโร อุปฏฺฐาเปตพฺโพ, น ภิกฺขุโนวาทกสมฺมุติ สาทิตพฺพา, สมฺมเตนปิ ภิกฺขุนิโย น โอวทิตพฺพา, น กาจิ สํฆสมฺมุติ สาทิตพฺพา, น กิสฺมิญฺจิ ปจฺเจกฏฺฐาเน ฐเปตพฺโพ, น จ เตน มูเลน วุฏฺฐาเปตพฺโพ.\csromanspacer{} ตสฺสปาปิยสิกกมฺมกเตน ภิกฺขเว ภิกฺขุนา อิเมสุ อฏฺฐสุ ธมฺเมสุ สมฺมา วตฺติตพฺพนฺติ.\csromanspacer{}\csromanspacer{}ทสมํ.\csromanspacer{} สติวคฺโค จตุตฺโถ.}

\titlehead{\textbf{ตสฺสุทฺทานํ}}
\csromangathameasure{สติปุณฺณิยมูเลน, โจรสมเณน ปญฺจมํ. \\ ยโส ปตฺตปฺปสาเทน, ปฏิสารณียญฺจ วตฺตนนฺติ.}
\csromangathagroup{\csromangathastanza{\csromanfolio{160}สติปุณฺณิยมูเลน, โจรสมเณน ปญฺจมํ. \\ ยโส ปตฺตปฺปสาเทน, ปฏิสารณียญฺจ วตฺตนนฺติ.}}

\chapterhead{\textbf{(10) 5. สามญฺญวคฺค}}
\proseitem{91-116}{อถ โข1 โพชฺฌา\footnote{โพชฺฌงฺคา (ก-สี)} อุปาสิกา1, สิรีมา, ปทุมา, สุตนา\footnote{สุธนา (สี, อิ), สุธมฺมา (สฺยา)}, มนุชา, อุตฺตรา, มุตฺตา, เขมา, รุจี\footnote{รูปี (สี, อิ)}, จุนฺที, พิมฺพี, สุมนา, มลฺลิกา, ติสฺสา, ติสฺสมาตา\footnote{ติสฺสาย มาตา (สี, อิ)}, โสณา, โสณาย มาตา\footnote{โสณมาตา (สฺยา)}, กาณา, กาณมาตา\footnote{กาณาย มาตา (สี, อิ)}, อุตฺตรา นนฺทมาตา, วิสาขา มิคารมาตา, ขุชฺชุตฺตรา อุปาสิกา, สามาวตี อุปาสิกา, สุปฺปวาสา โกลิยธีตา\footnote{โกฬิยธีตา (สฺยา, อิ)}, สุปฺปิยา อุปาสิกา, นกุลมาตา คหปตานี. (1-26)}

\vspace{\tipitakaparskip}
\csromancenter{สามญฺญวคฺโค ปญฺจโม.}
\csromancenter{\textbf{ทุติโย ปณฺณาสโก สมตฺโต.}}
\csromancenter{\textbf{(11) ราคเปยฺยาล} \textbf{(11) ราคเปยฺยาล}}
\proseitem{117}{\csromanfolio{161}ราคสฺส ภิกฺขเว อภิญฺญาย อฏฺฐ ธมฺมา ภาเวตพฺพา.\csromanspacer{} กตเม อฏฺฐ? สมฺมาทิฏฺฐิ สมฺมาสงฺกปฺโป สมฺมาวาจา สมฺมากมฺมนฺโต สมฺมา-อาชีโว สมฺมาวายาโม สมฺมาสติ สมฺมาสมาธิ.\csromanspacer{} ราคสฺส ภิกฺขเว อภิญฺญาย อิเม อฏฺฐ ธมฺมา ภาเวตพฺพาติ. (1)}

\proseitem{118}{ราคสฺส ภิกฺขเว อภิญฺญาย อฏฺฐ ธมฺมา ภาเวตพฺพา.\csromanspacer{} กตเม อฏฺฐ? อชฺฌตฺตํ รูปสญฺญี พหิทฺธา รูปานิ ปสฺสติ ปริตฺตานิ สุวณฺณทุพฺพณฺณานิ, ตานิ อภิภุยฺย “ชานามิ ปสฺสามี”ติ เอวํสญฺญี โหติ.\csromanspacer{} อชฺฌตฺตํ รูปสญฺญี พหิทฺธา รูปานิ ปสฺสติ อปฺปมาณานิ สุวณฺณทุพฺพณฺณานิ, ตานิ อภิภุยฺย “ชานามิ ปสฺสามี”ติ เอวํสญฺญี โหติ.\csromanspacer{} อชฺฌตฺตํ อรูปสญฺญี พหิทฺธา รูปานิ ปสฺสติ ปริตฺตานิ สุวณฺณทุพฺพณฺณนิ, ตานิ อภิภุยฺย “ชานามิ ปสฺสามี”ติ เอวํสญฺญี โหติ.\csromanspacer{} อชฺฌตฺตํ อรูปสญฺญี พหิทฺธา รูปานิ ปสฺสติ อปฺปมาณานิ สุวณฺณทุพฺพณฺณานิ, ตานิ อภิภุยฺย “ชานามิ ปสฺสามี”ติ เอวํสญฺญี โหติ.\csromanspacer{} อชฺฌตฺตํ อรูปสญฺญี พหิทฺธา รูปานิ ปสฺสติ นีลานิ นีลวณฺณานิ นีลทสฺสนานิ นีลนิภาสานิ, ปีตานิ ปีตวณฺณานิ ฯเปฯ โลหิตกานิ โลหิตกวณฺณานิ ฯเปฯ โอทาตานิ โอทาตวณฺณานิ ฯเปฯ โอทาตนิภาสานิ, ตานิ อภิภุยฺย “ชานามิ ปสฺสามี”ติ เอวํสญฺญี โหติ.\csromanspacer{} ราคสฺส ภิกฺขเว อภิญฺญาย อิเม อฏฺฐ ธมฺมา ภาเวตพฺพา. (2)}

\proseitem{119}{ราคสฺส ภิกฺขเว อภิญฺญาย อฏฺฐ ธมฺมา ภาเวตพฺพา.\csromanspacer{} กตเม อฏฺฐ? รูปี รูปานิ ปสฺสติ.\csromanspacer{} อชฺฌตฺตํ อรูปสญฺญี พหิทฺธา รูปานิ ปสฺสติ. “สุภนฺ”เตว อธิมุตฺโต โหติ.\csromanspacer{} สพฺพโส รูปสญฺญานํ สมติกฺกมา ปฏิฆสญฺญานํ อตฺถงฺคมา นานตฺตสญฺญานํ อมนสิการา “อนนฺโต อากาโส”ติ อากาสานญฺจายตนํ อุปสมฺปชฺช วิหรติ.\csromanspacer{} สพฺพโส อากาสานญฺจายตนํ สมติกฺกมฺม “อนนฺตํ วิญฺญาณนฺ”ติ วิญฺญาณญฺจายตนํ อุปสมฺปชฺช วิหรติ.\csromanspacer{} สพฺพโส วิญฺญาณญฺจายตนํ สมติกฺกมฺม “นตฺถิ กิญฺจี”ติ อากิญฺจญฺญายตนํ อุปสมฺปชฺช วิหรติ.\csromanspacer{} สพฺพโส อากิญฺจญฺญายตนํ สมติกฺกมฺม เนวสญฺญานาสญฺญายตนํ อุปสมฺปชฺช วิหรติ.\csromanspacer{} สพฺพโส เนวสญฺญานาสญฺญายตนํ สมติกฺกมฺม สญฺญาเวทยิตนิโรธํ อุปสมฺปชฺช วิหรติ.\csromanspacer{} ราคสฺส ภิกฺขเว อภิญฺญาย อิเม อฏฺฐ ธมฺมา ภาเวตพฺพา. (3)}

\proseitem{120-146}{\csromanfolio{162}ราคสฺส ภิกฺขเว ปริญฺญาย ฯเปฯ ปริกฺขยาย.\csromanspacer{} ปหานาย.\csromanspacer{} ขยาย.\csromanspacer{} วยาย.\csromanspacer{} วิราคาย.\csromanspacer{} นิโรธาย.\csromanspacer{} จาคาย.\csromanspacer{} ปฏินิสฺสคฺคาย ฯเปฯ.\csromanspacer{} อิเม อฏฺฐ ธมฺมา ภาเวตพฺพา. (4-30)}

\proseitem{147-626}{โทสสฺส ฯเปฯ.\csromanspacer{} โมหสฺส.\csromanspacer{} โกธสฺส.\csromanspacer{} อุปนาหสฺส.\csromanspacer{} มกฺขสฺส.\csromanspacer{} ปฬาสสฺส.\csromanspacer{} อิสฺสาย.\csromanspacer{} มจฺฉริยสฺส.\csromanspacer{} มายาย.\csromanspacer{} สาเฐยฺยสฺส.\csromanspacer{} ถมฺภสฺส.\csromanspacer{} สารมฺภสฺส.\csromanspacer{} มานสฺส.\csromanspacer{} อติมานสฺส.\csromanspacer{} มทสฺส.\csromanspacer{} ปมาทสฺส อภิญฺญาย ฯเปฯ ปริญฺญาย.\csromanspacer{} ปริกฺขยาย.\csromanspacer{} ปหานาย.\csromanspacer{} ขยาย.\csromanspacer{} วยาย.\csromanspacer{} วิราคาย.\csromanspacer{} นิโรธาย.\csromanspacer{} จาคาย.\csromanspacer{} ปฏินิสฺสคฺคาย ฯเปฯ.\csromanspacer{} อิเม อฏฺฐ ธมฺมา ภาเวตพฺพาติ. (31-510)}

\nitthitam{ราคเปยฺยาลํ นิฏฺฐิตํ.}
\nitthitam{\textbf{อฏฺฐกนิปาตปาฬิ นิฏฺฐิตา.}}
\csromanheader{2}{\textbf{องฺคุตฺตรนิกาย}}
\gambhira{\textbf{นวกนิปาตปาฬิ}}
\namakkaram{นโม ตสฺส ภควโต อรหโต สมฺมาสมฺพุทฺธสฺส.}
\csromanheader{2}{1. ปฐมปณฺณาสก}
\chapterhead{1. สมฺโพธิวคฺค}
\csromanheader{2}{1. สมฺโพธิสุตฺต}
\proseitem{1}{\csromanfolio{163}เอวํ เม สุตํ\csromandash{}เอกํ สมยํ ภควา สาวตฺถิยํ วิหรติ เชตวเน อนาถปิณฺฑิกสฺส อาราเม, ตตฺร โข ภควา ภิกฺขู อามนฺเตสิ\csromandash{}}

\prose{สเจ ภิกฺขเว อญฺญติตฺถิยา ปริพฺพาชกา เอวํ ปุจฺเฉยฺยุํ “สมฺโพธิปกฺขิกานํ\footnote{สมฺโพธปกฺขิกานํ (สี, สฺยา, อิ)} อาวุโส ธมฺมานํ กา อุปนิสา ภาวนายา”ติ.\csromanspacer{} เอวํ ปุฏฺฐา ตุเมฺห ภิกฺขเว เตสํ อญฺญติตฺถิยานํ ปริพฺพาชกานํ กินฺติ พฺยากเรยฺยาถาติ.\csromanspacer{} ภควํมูลกา โน ภนฺเต ธมฺมา ฯเปฯ ภควโต สุตฺวา ภิกฺขู ธาเรสฺสนฺตีติ.}

\prose{เตน หิ ภิกฺขเว สุณาถ สาธุกํ มนสิ กโรถ, ภาสิสฺสามีติ. “เอวํ ภนฺเต”ติ โข เต ภิกฺขู ภควโต ปจฺจสฺโสสุํ.\csromanspacer{} ภควา เอตทโวจ\csromandash{}}

\prose{สเจ ภิกฺขเว อญฺญติตฺถิยา ปริพฺพาชกา เอวํ ปุจฺเฉยฺยุํ “สมฺโพธิปกฺขิกานํ อาวุโส ธมฺมานํ กา อุปนิสา ภาวนายา”ติ.\csromanspacer{} เอวํ}

\vspace{\tipitakaparskip}
\csromancenter{นวกนิปาตปาฬิ ปุฏฺฐา ตุเมฺห ภิกฺขเว เตสํ อญฺญติตฺถิยานํ ปริพฺพาชกานํ เอวํ พฺยากเรยฺยาถ\csromandash{}}
\prose{\csromanfolio{164}อิธาวุโส ภิกฺขุ กลฺยาณมิตฺโต โหติ กลฺยาณสหาโย กลฺยาณสมฺปวงฺโก.\csromanspacer{} สมฺโพธิปกฺขิกานํ อาวุโส ธมฺมานํ อยํ ปฐมา อุปนิสา ภาวนาย.}

\prose{ปุน จปรํ อาวุโส ภิกฺขุ สีลวา โหติ, ปาติโมกฺขสํวรสํวุโต วิหรติ อาจารโคจรสมฺปนฺโน, อณุมตฺเตสุ วชฺเชสุ ภยทสฺสาวี สมาทาย สิกฺขติ สิกฺขาปเทสุ.\csromanspacer{} สมฺโพธิปกฺขิกานํ อาวุโส ธมฺมานํ อยํ ทุติยา อุปนิสา ภาวนาย.}

\prose{ปุน จปรํ อาวุโส ภิกฺขุ ยายํ กถา อภิสลฺเลขิกา เจโตวิวรณสปฺปายา.\csromanspacer{} เสยฺยถิทํ, อปฺปิจฺฉกถา สนฺตุฏฺฐิกถา ปวิเวกกถา อสํสคฺคกถา วีริยารมฺภกถา สีลกถา สมาธิกถา ปญฺญากถา วิมุตฺติกถา วิมุตฺติญาณทสฺสนกถา.\csromanspacer{} เอวรูปิยา กถาย นิกามลาภี โหติ อกิจฺฉลาภี อกสิรลาภี.\csromanspacer{} สมฺโพธิปกฺขิกานํ อาวุโส ธมฺมานํ อยํ ตติยา อุปนิสา ภาวนาย.}

\prose{ปุน จปรํ อาวุโส ภิกฺขุ อารทฺธวีริโย วิหรติ อกุสลานํ ธมฺมานํ ปหานาย กุสลานํ ธมฺมานํ อุปสมฺปทาย, ถามวา ทฬฺหปรกฺกโม อนิกฺขิตฺตธุโร กุสเลสุ ธมฺเมสุ.\csromanspacer{} สมฺโพธิปกฺขิกานํ อาวุโส ธมฺมานํ อยํ จตุตฺถี อุปนิสา ภาวนาย.}

\prose{ปุน จปรํ อาวุโส ภิกฺขุ ปญฺญวา โหติ อุทยตฺถคามินิยา ปญฺญาย สมนฺนาคโต อริยาย นิพฺเพธิกาย สมฺมา ทุกฺขกฺขยคามินิยา.\csromanspacer{} สมฺโพธิปกฺขิกานํ อาวุโส ธมฺมานํ อยํ ปญฺจมี อุปนิสา ภาวนาย.}

\prose{กลฺยาณมิตฺตสฺเสตํ ภิกฺขเว ภิกฺขุโน ปาฏิกงฺขํ กลฺยาณสหายสฺส กลฺยาณสมฺปวงฺกสฺส “สีลวา ภวิสฺสติ, ปาติโมกฺขสํวรสํวุโส วิหริสฺสติ อาจารโคจรสมฺปนฺโน, อณุมตฺเตสุ วชฺเชสุ ภยทสฺสาวี สมาทาย สิกฺขิสฺสติ สิกฺขาปเทสุ”.}

\prose{กลฺยาณมิตฺตสฺเสตํ ภิกฺขเว ภิกฺขุโน ปาฏิกงฺขํ กลฺยาณสหายสฺส กลฺยาณสมฺปวงฺกสฺส “ยายํ กถา อภิสลฺเลขิกา เจโตวิวรณสปฺปายา.\csromanspacer{} เสยฺยถิทํ, อปฺปิจฺฉกถา สนฺตุฏฺฐิกถา}

\chapterheadpageread{1. สมฺโพธิวคฺค}
\prose{\csromanfolio{165}ปวิเวกกถา อสํสคฺคกถา วีริยารมฺภกถา สีลกถา สมาธิกถา ปญฺญากถา วิมุตฺติกถา วิมุตฺติญาณทสฺสนกถา, เอวรูปิยา กถาย นิกามลาภี ภวิสฺสติ อกิจฺฉลาภี อกสิรลาภี”.}

\prose{กลฺยาณมิตฺตสฺเสตํ ภิกฺขเว ภิกฺขุโน ปาฏิกงฺขํ กลฺยาณสหายสฺส กลฺยาณสมฺปวงฺกสฺส “อารทฺธวีริโย วิหริสฺสติ อกุสลานํ ธมฺมานํ ปหานาย กุสลานํ ธมฺมานํ อุปสมฺปทาย, ถามวา ทฬฺหปรกฺกโม อนิกฺขิตฺตธุโร กุสเลสุ ธมฺเมสุ”.}

\prose{กลฺยาณมิตฺตสฺเสตํ ภิกฺขเว ภิกฺขุโน ปาฏิกงฺขํ กลฺยาณสหายสฺส กลฺยาณสมฺปวงฺกสฺส “ปญฺญวา ภวิสฺสติ อุทยตฺถคามินิยา ปญฺญาย สมนฺนาคโต อริยาย นิพฺเพธิกาย สมฺมา ทุกฺขกฺขยคามินิยา”.}

\prose{เตน จ ปน ภิกฺขเว ภิกฺขุนา อิเมสุ ปญฺจสุ ธมฺเมสุ ปติฏฺฐาย จตฺตาโร ธมฺมา อุตฺตริ\footnote{อุตฺตริํ (สี, สฺยา, อิ)} ภาเวตพฺโพ, อสุภา ภาเวตพฺพา ราคสฺสปหานาย, เมตฺตา ภาเวตพฺพา พฺยาปาทสฺส ปหานาย, อานาปานสฺสติ\footnote{อานาปานสติ (สี, อิ)} ภาเวตพฺพา วิตกฺกุปจฺเฉทาย, อนิจฺจสญฺญา ภาเวตพฺพา อสฺมิมานสมุคฺฆาตาย.\csromanspacer{} อนิจฺจสญฺญิโน ภิกฺขเว อนตฺตสญฺญา สณฺฐาติ, อนตฺตสญฺญี อสฺมิมานสมุคฺฆาตํ ปาปุณาติ ทิฏฺเฐว ธมฺเม นิพฺพานนฺติ.\csromanspacer{}\csromanspacer{}ปฐมํ.}

\csromanheader{2}{2. นิสฺสยสุตฺต}
\proseitem{2}{อถ โข อญฺญตโร ภิกฺขุ เยน ภควา เตนุปสงฺกมิ, อุปสงฺกมิตฺวา ภควนฺตํ ฯเปฯ เอกมนฺตํ นิสินฺโน โข โส ภิกฺขุ ภควนฺตํ เอตทโวจ “นิสฺสยสมฺปนฺโน นิสฺสยสมฺปนฺโนติ ภนฺเต วุจฺจติ, กิตฺตาวตา นุ โข ภนฺเต ภิกฺขุ นิสฺสยสมฺปนฺโน โหตี”ติ.\csromanspacer{} สทฺธํ เจ ภิกฺขุ ภิกฺขุ นิสฺสาย อกุสลํ ปชหติ, กุสลํ ภาเวติ, ปหีนเมวสฺส ตํ อกุสลํ โหติ.\csromanspacer{} หิริํ เจ ภิกฺขุ ภิกฺขุ นิสฺสาย ฯเปฯ.\csromanspacer{} โอตฺตปฺปํ เจ ภิกฺขุ ภิกฺขุ นิสฺสาย ฯเปฯ.\csromanspacer{} วีริยํ เจ ภิกฺขุ ภิกฺขุ นิสฺสาย ฯเปฯ.\csromanspacer{} ปญฺญํ เจ ภิกฺขุ ภิกฺขุ นิสฺสาย อกุสลํ ปชหติ, กุสลํ ภาเวติ, ปหีนเมวสฺส ตํ อกุลลํ โหติ.\csromanspacer{} ตํ หิสฺส ภิกฺขุโน อกุสลํ ปหีนํ โหติ สุปฺปหีนํ, ยํ’ส อริยาย ปญฺญาย ทิสฺวา ปหีนํ.}

\gambhira{นวกนิปาตปาฬิ}
\prose{\csromanfolio{166}เตน จ ปน ภิกฺขุ ภิกฺขุนา อิเมสุ ปญฺจสุ ธมฺเมสุ ปติฏฺฐาย จตฺตาโร อุปนิสฺสาย วิหาตพฺพา.\csromanspacer{} กตเม จตฺตาโร? อิธ ภิกฺขุ ภิกฺขุ สงฺขาเยกํ ปฏิเสวติ, สงฺขาเยกํ อธิวาเสติ, สงฺขาเยกํ ปริวชฺเชติ, สงฺขาเยกํ วิโนเทติ.\csromanspacer{} เอวํ โข ภิกฺขุ ภิกฺขุ นิสฺสยสมฺปนฺโน โหตีติ.\csromanspacer{}\csromanspacer{}ทุติยํ.}

\csromanheader{2}{3. เมฆิยสุตฺต}
\proseitem{3}{เอกํ สมยํ ภควา จาลิกายํ วิหรติ จาลิกาปพฺพเต, เตน โข ปน สมเยน อายสฺมา เมฆิโย ภควโต อุปฏฺฐาโก โหติ, อถ โข อายสฺมา เมฆิโย เยน ภควา เตนุปสงฺกมิ, อุปสงฺกมิตฺวา ภควนฺตํ อภิวาเทตฺวา เอกมนฺตํ อฏฺฐาสิ, เอกมนฺตํ ฐิโต โข อายสฺมา เมฆิโย ภควนฺตํ เอตทโวจ “อิจฺฉามหํ ภนฺเต ชนฺตุคามํ\footnote{ชตุคามํ (สี-ฏฺฐ, สฺยา-ฏฺฐ), ชตฺตุคามํ (ก, อฏฺฐกถายมฺปิ ปาฐนฺตรํ)} ปิณฺฑาย ปวิสิตุนฺ”ติ.\csromanspacer{} ยสฺส ทานิ ตฺวํ เมฆิย กาลํ มญฺญสีติ.}

\prose{อถ โข อายสฺมา เมฆิโย ปุพฺพณฺหสมยํ นิวาเสตฺวา ปตฺตจีวรมาทาย ชนฺตุคามํ ปิณฺฑาย ปาวิสิ, ชนฺตุคาเม ปิณฺฑาย จริตฺวา ปจฺฉาภตฺตํ ปิณฺฑปาตปฏิกฺกนฺโต เยน กิมิกาฬาย นทิยา ตีรํ เตนุปสงฺกมิ, อทฺทสา โข อายสฺมา เมฆิโย กิมิกาฬาย นทิยา ตีเร ชงฺฆาวิหารํ\footnote{ชงฺฆวิหารํ (สฺยา, ก)} อนุจงฺกมมาโน อนุวิจรมาโน อมฺพวนํ ปาสาทิกํ รมณียํ, ทิสฺวานสฺส เอตทโหสิ “ปาสาทิกํ วติทํ อมฺพวนํ รมณียํ, อลํ วติทํ กุลปุตฺตสฺส ปธานตฺถิกสฺส ปธานาย.\csromanspacer{} สเจ มํ ภควา อนุชาเนยฺย, อาคจฺเฉยฺยาหํ อิมํ อมฺพวนํ ปธานายา”ติ.}

\prose{อถ โข อายสฺมา เมฆิโย เยน ภควา เตนุปสงฺกมิ, อุปสงฺกมิตฺวา ภควนฺตํ อภิวาเทตฺวา เอกมนฺตํ นิสีทิ, เอกมนฺตํ นิสินฺโน โข อายสฺมา เมฆิโย ภควนฺตํ เอตทโวจ “อิธาหํ ภนฺเต ปุพฺพณฺหสมยํ นิวาเสตฺวา ปตฺตจีวรมาทาย ชนฺตุคามํ ปิณฺฑาย ปาวิสิํ, ชนฺตุคาเม ปิณฺฑาย จริตฺวา ปจฺฉาภตฺตํ ปิณฺฑปาตปฏิกฺกนฺโต เยน กิมิกาฬาย นทิยา ตีรํ เตนุปสงฺกมิํ, อทฺทสํ โข อหํ ภนฺเต กิมิกาฬาย นทิยา ตีเร ชงฺฆาวิหารํ อนุจงฺกมมาโน อนุวิจรมาโน อมฺพวนํ ปาสาทิกํ รมณียํ, ทิสฺวาน เม เอตทโหสิ ‘ปาสาทิกํ วติทํ อมฺพวนํ รมณียํ, อลํ วติทํ กุลปุตฺตสฺส ปธานตฺถิกสฺส ปธานาย.}

\chapterheadpageread{1. สมฺโพธิวคฺค}
\prose{\csromanfolio{167}สเจ มํ ภควา อนุชาเนยฺย, อาคจฺเฉยฺยาหํ อิมํ อมฺพวนํ ปธานายา’ติ, สเจ มํ ภควา อนุชาเนยฺย, คจฺเฉยฺยาหํ ตํ อมฺพวนํ ปธานายา”ติ.\csromanspacer{} อาคเมหิ ตาว เมฆิย, เอกกมฺหิ\footnote{เอกกมฺหา (สี, อิ)} ตาว\footnote{วต (ก)}, ยาว อญฺโญปิ โกจิ ภิกฺขุ อาคจฺฉตีติ\footnote{ทิสฺสตูติ (สพฺพตฺถ, ฏีกายมฺปิ ปาฐนฺตรํ), อาคจฺฉตูติ, ทิสฺสตีติ (ฏีกายํ ปาฐนฺตรานิ)}.}

\prose{ทุติยมฺปิ โข อายสฺมา เมฆิโย ภควนฺตํ เอตทโวจ “ภควโต ภนฺเต นตฺถิ กิญฺจิ อุตฺตริ กรณียํ, นตฺถิ กตสฺส ปฏิจโย.\csromanspacer{} มยฺหํ โข ปน ภนฺเต อตฺถิ อุตฺตริ กรณียํ, อตฺถิ กตสฺส ปฏิจโย.\csromanspacer{} สเจ มํ ภควา อนุชาเนยฺย, คจฺเฉยฺยาหํ ตํ อมฺพวนํ ปธานายา”ติ.\csromanspacer{} อาคเมหิ ตาว เมฆิย, เอกกมฺหิ ตาว, ยาว อญฺโญปิ โกจิ ภิกฺขุ อาคจฺฉตีติ.}

\prose{ตติยมฺปิ โข อายสฺมา เมฆิโย ภควนฺตํ เอตทโวจ “ภควโต ภนฺเต นตฺถิ กิญฺจิ อุตฺตริ กรณียํ, นตฺถิ กตสฺส ปฏิจโย.\csromanspacer{} มยฺหํ โข ปน ภนฺเต อตฺถิ อุตฺตริ กรณียํ, อตฺถิ กตสฺส ปฏิจโย.\csromanspacer{} สเจ มํ ภควา อนุชาเนยฺย, คจฺเฉยฺยาหํ ตํ อมฺพวนํ ปธานายา”ติ. “ปธานนฺ”ติ โข เมฆิย วทมานํ กินฺติ วเทยฺยาม, ยสฺส ทานิ ตฺวํ เมฆิย กาลํ มญฺญสีติ.}

\prose{อถ โข อายสฺมา เมฆิโย อุฏฺฐายาสนา ภควนฺตํ อภิวาเทตฺวา ปทกฺขิณํ กตฺวา เยน ตํ อมฺพวนํ เตนุปสงฺกมิ, อุปสงฺกมิตฺวา ตํ อมฺพวนํ อชฺโฌคาเหตฺวา อญฺญตรสฺมิํ รุกฺขมูเล ทิวาวิหารํ นิสีทิ.\csromanspacer{} อถ โข อายสฺมโต เมฆิยสฺส ตสฺมิํ อมฺพวเน วิหรนฺตสฺส เยภุยฺเยน ตโย ปาปกา อกุสลา วิตกฺกา สมุทาจรนฺติ.\csromanspacer{} เสยฺยถิทํ, กามวิตกฺโก พฺยาปาทวิตกฺโก วิหิํสาวิตกฺโก.\csromanspacer{} อถ โข อายสฺมโต เมฆิยสฺส เอตทโหสิ “อจฺฉริยํ วต โภ, อพฺภุตํ วต โภ, สทฺธาย จ วตมฺหา อคารสฺมา อนาคาริยํ ปพฺพชิตา, อถ จ ปนิเมหิ ตีหิ ปาปเกหิ อกุสเลหิ วิตกฺเกหิ อนฺวาสตฺตา กามวิตกฺเกน พฺยาปาทวิตกฺเกน วิหิํสาวิตกฺเกนา”ติ.}

\prose{อถ โข อายสฺมา เมฆิโย เยน ภควา เตนุปสงฺกมิ, อุปสงฺกมิตฺวา ภควนฺตํ อภิวาเทตฺวา เอกมนฺตํ นิสีทิ, เอกมนฺตํ นิสินฺโน โข อายสฺมา เมฆิโย ภควนฺตํ เอตทโวจ\csromandash{}}

\gambhira{นวกนิปาตปาฬิ}
\prose{\csromanfolio{168}อิธ มยฺหํ ภนฺเต ตสฺมิํ อมฺพวเน วิหรนฺตสฺส เยภุยฺเยน ตโย ปาปกา อกุสลา วิตกฺกา สมุทาจรนฺติ.\csromanspacer{} เสยฺยถิทํ, กามวิตกฺโก พฺยาปาทวิตกฺโก วิหิํสาวิตกฺโก.\csromanspacer{} ตสฺส มยฺหํ ภนฺเต เอตทโหสิ “อจฺฉริยํ วต โภ, อพฺภุตํ วต โภ, สทฺธาย จ วตมฺหา อคารสฺมา อนคาริยํ ปพฺพชิตา, อถ จ ปนิเมหิ ตีหิ ปาปเกหิ อกุสเลหิ วิตกฺเกหิ อนฺวาสตฺตา กามวิตกฺเกน พฺยาปาทวิตกฺเกน วิหิํสาวิตกฺเกนา”ติ.}

\prose{อปริปกฺกาย เมฆิย เจโตวิมุตฺติยา ปญฺจ ธมฺมา ปริปกฺกาย สํวตฺตนฺติ.\csromanspacer{} กตเม ปญฺจ? อิธ เมฆิย ภิกฺขุ กลฺยาณมิตฺโต โหติ กลฺยาณสหาโย กลฺยาณสมฺปวงฺโก.\csromanspacer{} อปริปกฺกาย เมฆิย เจโตวิมุตฺติยา อยํ ปฐโม ธมฺโม ปริปกฺกาย สํวตฺตหิ.}

\prose{ปุน จปรํ เมฆิย ภิกฺขุ สีลวา โหติ, ปาติโมกฺขสํวรสํวุโต วิหรติ อาจารโคจรสมฺปนฺโน, อณุมตฺเตสุ วชฺเชสุ ภยทสฺสาวี สมาทาย สิกฺขติ สิกฺขาปเทสุ.\csromanspacer{} อปริปกฺกาย เมฆิย เจโตวิมุตฺติยา อยํ ทุติโย ธมฺโม ปริปกฺกาย สํวตฺตติ.}

\prose{ปุน จปรํ เมฆิย ยายํ กถา อภิสลฺเลขิกา เจโตวิวรณสปฺปายา.\csromanspacer{} เสยฺยถิทํ, อปฺปิจฺฉกถา สนฺตุฏฺฐิกถา ปวิเวกกถา อสํสคฺคกถา วีริยารมฺภกถา สีลกถา สมาธิกถา ปญฺญากถา วิมุตฺติกถา วิมุตฺติญาณทสฺสนกถา, เอวรูปิยา กถาย นิกามลาภี โหติ อกิจฺฉลาภี อกสิรลาภี.\csromanspacer{} อปริปกฺกาย เมฆิย เจโตวิมุตฺติยา อยํ ตติโย ธมฺโม ปริปกฺกาย สํวตฺตติ.}

\prose{ปุน จปรํ เมฆิย ภิกฺขุ อารทฺธวีริโย วิหรติ อกุสลานํ ธมฺมานํ ปหานาย กุสลานํ ธมฺมานํ อุปสมฺปทาย, ถามวา ทฬฺหปรกฺกโม อนิกฺขิตฺตธุโร กุสเลสุ ธมฺเมสุ.\csromanspacer{} อปริปกฺกาย เมฆิย เจโตวิมุตฺติยา อยํ จตุตฺโถ ธมฺโม ปริปกฺกาย สํวตฺตติ.}

\prose{ปุน จปรํ เมฆิย ภิกฺขุ ปญฺญวา โหติ อุทยตฺถคามินิยา ปญฺญาย สมนฺนาคโต อริยาย นิพฺเพธิกาย สมฺมา ทุกฺขกฺขยคามินิยา.\csromanspacer{} อปริปกฺกาย เมฆิย เจโตวิมุตฺติยา อยํ ปญฺจโม ธมฺโม ปริปกฺกาย สํวตฺตติ.}

\chapterheadpageread{1. สมฺโพธิวคฺค}
\prose{\csromanfolio{169}กลฺยาณมิตฺตสฺเสตํ เมฆิย ภิกฺขุโน ปาฏิกงฺขํ กลฺยาณสหายสฺส กลฺยาณสมฺปวงฺกสฺส “สีลวา ภวิสฺสติ ฯเปฯ สมาทาย สิกฺขิสฺสติ สิกฺขาปเทสุ”.}

\prose{กลฺยาณมิตฺตสฺเสตํ เมฆิย ภิกฺขุโน ปาฏิกงฺขํ กลฺยาณสหายสฺส กลฺยาณสมฺปวงฺกสฺส “ยายํ กถา อภิสลฺเลขิกา เจโตวิวรณสปฺปายา.\csromanspacer{} เสยฺยถิทํ, อปฺปิจฺฉกถา ฯเปฯ วิมุตฺติญาณทสฺสนกถา, เอวรูปิยา กถาย นิกามลาภี ภวิสฺสติ อกิจฺฉลาภี อกสิรลาภี”.}

\prose{กลฺยาณมิตฺตสฺเสตํ เมฆิย ภิกฺขุโน ปาฏิกงฺขํ กลฺยาณสหายสฺส กลฺยาณสมฺปวงฺกสฺส “อารทฺธวีริโย วิหริสฺสติ ฯเปฯ อนิกฺขิตฺตธุโร กุสเลสุ ธมฺเมสุ”.}

\prose{กลฺยาณมิตฺตสฺเสตํ เมฆิย ภิกฺขุโน ปาฏิกงฺขํ กลฺยาณสหายสฺส กลฺยาณสมฺปวงฺกสฺส “ปญฺญวา ภวิสฺสติ ฯเปฯ สมฺมา ทุกฺขกฺขยคามินิยา”.}

\prose{เตน จ ปน เมฆิย ภิกฺขุนา อิเมสุ ปญฺจสุ ธมฺเมสุ ปติฏฺฐาย จตฺตาโร ธมฺมา อุตฺตริ ภาเวตพฺพา.\csromanspacer{} อสุภา ภาเวตพฺพา ราคสฺส ปหานาย, เมตฺตา ภาเวตพฺพา พฺยาปาทสฺส ปหานาย, อานาปานสฺสติ ภาเวตพฺพา วิตกฺกุปจฺเฉทาย, อนิจฺจสญฺญา ภาเวตพฺพา อสฺมิมานสมุคฺฆาตาย, อนิจฺจสญฺญิโน เมฆิย อนตฺตสญฺญา สณฺฐาติ, อนตฺตสญฺญี อสฺมิมานสมุคฺฆาตํ ปาปุณาติ ทิฏฺเฐว ธมฺเม นิพฺพานนฺติ.\csromanspacer{}\csromanspacer{}ตติยํ.}

\csromanheader{2}{4. นนฺทกสุตฺต}
\proseitem{4}{เอกํ สมยํ ภควา สาวตฺถิยํ วิหรติ เชตวเน อนาถปิณฺฑิกสฺส อาราเม.\csromanspacer{} เตน โข ปน สมเยน อายสฺมา นนฺทโก อุปฏฺฐานสาลายํ ภิกฺขู ธมฺมิยา กถาย สนฺทสฺเสติ สมาทเปติ สมุตฺเตเชติ สมฺปหํเสติ.\csromanspacer{} อถ โข ภควา สายนฺหสมยํ ปฏิสลฺลานา วุฏฺฐิโต เยนุปฏฺฐานสาลา เตนุปสงฺกมิ, อุปสงฺกมิตฺวา พหิทฺวาร โกฏฺฐเก อฏฺฐาสิ กถาปริโยสานํ อาคมยมาโน.\csromanspacer{} อถ โข ภควา กถาปริโยสานํ วิทิตฺวา อุกฺกาเสตฺวา อคฺคฬํ อาโกเฏสิ, วิวริํสุ โข เต ภิกฺขู ภควโต ทฺวารํ.}

\gambhira{นวกนิปาตปาฬิ}
\prose{\csromanfolio{170}อถ โข ภควา อุปฏฺฐานสาลํ ปวิสิตฺวา ปญฺญตฺตาสเน นิสีทิ, นิสชฺช โข ภควา อายสฺมนฺตํ นนฺทกํ เอตทโวจ “ทีโฆ โข ตฺยายํ นนฺทก ธมฺมปริยาโย ภิกฺขูนํ ปฏิภาสิ, อปิ เม ปิฏฺฐิ อาคิลายติ พหิทฺวารโกฏฺฐเก ฐิตสฺส กถาปริโยสานํ อาคมยมานสฺสา”ติ.}

\prose{เอวํ วุตฺเต อายสฺมา นนฺทโก สารชฺชมานรูโป ภควนฺตํ เอตทโวจ “น โข ปน มยํ ภนฺเต ชานาม ‘ภควา พหิทฺวารโกฏฺฐเก ฐิโต’ติ.\csromanspacer{} สเจ หิ มยํ ภนฺเต ชาเนยฺยาม ‘ภควา พหิทฺวารโกฏฺฐเก ฐิโต’ติ.\csromanspacer{} เอตฺตกมฺปิ ( )\footnote{(ธมฺมํ) กตฺถจิ อุปริ 417 ปิฏฺเฐ ปสฺสิตพฺพํ.} โน นปฺปฏิภาเยยฺยา”ติ.}

\prose{อถ โข ภควา อายสฺมนฺตํ นนฺทกํ สารชฺชมานรูปํ วิทิตฺวา อายสฺมนฺตํ นนฺทกํ เอตทโวจ “สาธุ สาธุ นนฺทก, เอตํ โข นนฺทก ตุมฺหากํ ปติรูปํ กุลปุตฺตานํ สทฺธาย อคารสฺมา อนคาริยํ ปพฺพชิตานํ, ยํ ตุเมฺห ธมฺมิยา กถาย สนฺนิสีเทยฺยาถ, สนฺนิปติตานํ โว นนฺทก ทฺวยํ กรณียํ ธมฺมี วา กถา, อริโย วา ตุณฺหีภาโว. * สทฺโธ จ นนฺทก ภิกฺขุ โหติ, โน จ สีลวา, เอวํ โส เตนงฺเคน อปริปูโร โหติ.\csromanspacer{} เตน ตํ องฺคํ ปริปูเรตพฺพํ “กินฺตาหํ สทฺโธ จ อสฺสํ สีลวา จา”ติ.\csromanspacer{} ยโต จ โข นนฺทก ภิกฺขุ สทฺโธ จ โหติ สีลวา จ, เอวํ โส เตนงฺเคน ปริปูโร โหติ.}

\prose{สทฺโธ จ นนฺทก ภิกฺขุ โหติ สีลวา จ โน จ ลาภี อชฺฌตฺตํ เจโตสมาธิสฺส, เอวํ โส เตนงฺเคน อปริปูโร โหติ.\csromanspacer{} เตน ตํ องฺคํ ปริปูเรตพฺพํ “กินฺตาหํ สทฺโธ จ อสฺสํ สีลวา จ, ลาภี จ อชฺฌตฺตํ เจโตสมาธิสฺสา”ติ.\csromanspacer{} ยโต จ โข นนฺทก ภิกฺขุ สทฺโธ จ โหติ สีลวา จ, ลาภี จ อชฺฌตฺตํ เจโตสมาธิสฺส, เอวํ โส เตนงฺเคน ปริปูโร โหติ.}

\prose{สทฺโธ จ นนฺทก ภิกฺขุ โหติ สีลวา จ, ลาภี จ อชฺฌตฺตํ เจโตสมาธิสฺส, น ลาภี อธิปญฺญาธมฺมวิปสฺสนาย, เอวํ โส เตนงฺเคน อปริปูโร โหติ.\csromanspacer{} เสยฺยถาปิ นนฺทก ปาณโก จตุปฺปาทโก อสฺส, ตสฺส เอโก ปาโท โอมโก ลามโก, เอวํ โส เตนงฺเคน อปริปูโร อสฺส.\csromanspacer{} เอวเมวํ โข นนฺทก ภิกฺขุ สทฺโธ จ โหติ สีลวา จ, ลาภี จ อชฺฌตฺตํ เจโตสมาธิสฺส, น ลาภี อธิปญฺญาธมฺมวิปสฺสนาย,}

\chapterheadpageread{1. สมฺโพธิวคฺค}
\prose{\csromanfolio{171}เอวํ โส เตนงฺเคน อปริปูโร โหติ, เตน ตํ องฺคํ ปริปูเรตพฺพํ “กินฺตาหํ สทฺโธ จ อสฺสํ สีลวา จ, ลาภี จ อชฺฌตฺตํ เจโตสมาธิสฺส, ลาภี จ อธิปญฺญาธมฺมวิปสฺสนายา”ติ.}

\prose{ยโต จ โข นนฺทก ภิกฺขุ สทฺโธ จ โหติ สีลวา จ, ลาภี จ อชฺฌตฺตํ เจโตสมาธิสฺส, ลาภี จ อธิปญฺญาธมฺมวิปสฺสนาย, เอวํ โส เตนงฺเคน ปริปูโร โหตีติ.\csromanspacer{} อิทมโวจ ภควา, อิทํ วตฺวาน สุคโต อุฏฺฐายาสนา วิหารํ ปาวิสิ.}

\prose{อถ โข อายสฺมา นนฺทโก อจิรปกฺกนฺตสฺส ภควโต ภิกฺขู อามนฺเตสิ\csromandash{}อิทานิ อาวุโส ภควา จตูหิ ปเทหิ เกวลปริปุณฺณํ ปริสุทฺธํ พฺรหฺมจริยํ ปกาเสตฺวา อุฏฺฐายาสนา วิหารํ ปวิฏฺโฐ, “สทฺโธ จ นนฺทก ภิกฺขุ โหติ, โน จ สีลวา, เอวํ โส เตนงฺเคน อปริปูโร โหติ, เตน ตํ องฺคํ ปริปูเรตพฺพํ ‘กินฺตาหํ สทฺโธ จ อสฺสํ สีลวา จา’ติ.\csromanspacer{} ยโต จ โข นนฺทก ภิกฺขุ สทฺโธ จ โหติ สีลวา จ, เอวํ โส เตนงฺเคน ปริปูโร โหติ.\csromanspacer{} สทฺโธ จ นนฺทก ภิกฺขุ โหติ สีลวา จ, โน จ ลาภี อชฺฌตฺตํ เจโตสมาธิสฺส ฯเปฯ ลาภี จ อชฺฌตฺตํ เจโตสมาธิสฺส, น ลาภี อธิปญฺญาธมฺมวิปสฺสนาย, เอวํ โส เตนงฺเคน อปริปูโร โหติ.\csromanspacer{} เสยฺยถาปิ นนฺทก ปาณโก จตุปฺปาทโก อสฺส, ตสฺส เอโก ปาโท โอมโก ลามโก, เอวํ โส เตนงฺเคน อปริปูโร อสฺส.\csromanspacer{} เอวเมวํ โข นนฺทก ภิกฺขุ สทฺโธ จ โหติ สีลวา จ, ลาภี จ อชฺฌตฺตํ เจโตสมาธิสฺส, น ลาภี อธิปญฺญาธมฺมวิปสฺสนาย, เอวํ โส เตนงฺเคน อปริปูโร โหติ, เตน ตํ องฺคํ ปริปูเรตพฺพํ ‘กินฺตาหํ สทฺโธ จ อสฺสํ สีลวา จ, ลาภี จ อชฺฌตฺตํ เจโตสมาธิสฺส, ลาภี จ อธิปญฺญาธมฺมวิปสฺสนายา’ติ.\csromanspacer{} ยโต จ โข นนฺทก ภิกฺขุ สทฺโธ จ โหติ สีลวา จ ลาภี จ อชฺฌตฺตํ เจโตสมาธิสฺส, ลาภี จ อธิปญฺญาธมฺมวิปสฺสนาย.\csromanspacer{} เอวํ โส เตนงฺเคน ปริปูโร โหตี”ติ.}

\prose{ปญฺจิเม อาวุโส อานิสํสา กาเลน ธมฺมสฺสวเน กาเลน ธมฺมสากจฺฉาย.\csromanspacer{} กตเม ปญฺจ? อิธาวุโส ภิกฺขุ ภิกฺขูนํ ธมฺมํ เทเสติ อาทิกลฺยาณํ มชฺเฌกลฺยาณํ ปริโยสานกลฺยาณํ สาตฺถํ สพฺยญฺชนํ, เกวลปริปุณฺณํ ปริสุทฺธํ พฺรหฺมจริยํ ปกาเสติ.\csromanspacer{} ยถา ยถา อาวุโส ภิกฺขุ ภิกฺขูนํ ธมฺมํ เทเสติ อาทิกลฺยาณํ มชฺเฌกลฺยาณํ ปริโยสานกลฺยาณํ สาตฺถํ สพฺยญฺชนํ, เกวลปริปุณฺณํ ปริสุทฺธํ พฺรหฺมจริยํ}

\vspace{\tipitakaparskip}
\csromancenter{นวกนิปาตปาฬิ ปกาเสติ.\csromanspacer{} ตถา ตถา โส สตฺถุ ปิโย จ โหติ มนาโป จ ครุ จ ภาวนีโย จ.\csromanspacer{} อยํ อาวุโส ปฐโม อานิสํโส กาเลน ธมฺมสฺสวเน กาเลน ธมฺมสากจฺฉาย.}
\prose{\csromanfolio{172}ปุน จปรํ อาวุโส ภิกฺขุ ภิกฺขูนํ ธมฺมํ เทเสติ อาทิกลฺยาณํ มชฺเฌกลฺยาณํ ปริโยสานกลฺยาณํ สาตฺถํ สพฺยญฺชนํ, เกวลปริปุณฺณํ ปริสุทฺธํ พฺรหฺมจริยํ ปกาเสติ.\csromanspacer{} ยถา ยถา อาวุโส ภิกฺขุ ภิกฺขูนํ ธมฺมํ เทเสติ อาทิกลฺยาณํ ฯเปฯ พฺรหฺมจริยํ ปกาเสติ, ตถา ตถา โส ตสฺมิํ ธมฺเม อตฺถปฺปฏิสํเวที จ โหติ ธมฺมปฺปฏิสํเวที จ.\csromanspacer{} อยํ อาวุโส ทุติโย อานิสํโส กาเลน ธมฺมสฺสวเน กาเลน ธมฺมสากจฺฉาย.}

\prose{ปุน จปรํ อาวุโส ภิกฺขุ ภิกฺขูนํ ธมฺมํ เทเสติ อาทิกลฺยาณํ มชฺเฌกลฺยาณํ ปริโยสานกลฺยาณํ สาตฺถํ สพฺยญฺชนํ, เกวลปริปุณฺณํ ปริสุทฺธํ พฺรหฺมจริยํ ปกาเสติ.\csromanspacer{} ยถา ยถา อาวุโส ภิกฺขุ ภิกฺขูนํ ธมฺมํ เทเสติ อาทิกลฺยาณํ ฯเปฯ พฺรหฺมจริยํ ปกาเสติ, ตถา ตถา โส ตสฺมิํ ธมฺเม คมฺภีรํ อตฺถปทํ ปญฺญาย อติวิชฺฌ ปสฺสติ.\csromanspacer{} อยํ อาวุโส ตติโย อานิสํโส กาเลน ธมฺมสฺสวเน กาเลน ธมฺมสากจฺฉาย.}

\prose{ปุน จปรํ อาวุโส ภิกฺขุ ภิกฺขูนํ ธมฺมํ เทเสติ อาทิกลฺยาณํ ฯเปฯ พฺรหฺมจริยํ ปกาเสติ.\csromanspacer{} ยถา ยถา อาวุโส ภิกฺขุ ภิกฺขูนํ ธมฺมํ เทเสติ อาทิกลฺยาณํ ฯเปฯ พฺรหฺมจริยํ ปกาเสติ, ตถา ตถา นํ สพฺรหฺมจารี อุตฺตริ สมฺภาเวนฺติ “อทฺธา อยมายสฺมา ปตฺโต วา ปชฺชติ วา.\csromanspacer{} อยํ อาวุโส จตุตฺโถ อานิสํโส กาเลน ธมฺมสฺสวเน กาเลน ธมฺมสากจฺฉาย.}

\prose{ปุน จปรํ อาวุโส ภิกฺขุ ภิกฺขูนํ ธมฺมํ เทเสติ อาทิกลฺยาณํ มชฺเฌกลฺยาณํ ปริโยสานกลฺยาณํ สาตฺถํ สพฺยญฺชนํ, เกวลปริปุณฺณํ ปริสุทฺธํ พฺรหฺมจริยํ ปกาเสติ.\csromanspacer{} ยถา ยถา อาวุโส ภิกฺขุ ภิกฺขูนํ ธมฺมํ เทเสติ อาทิกลฺยาณํ มชฺเฌกลฺยาณํ ปริโยสานกลฺยาณํ สาตฺถํ สพฺยญฺชนํ, เกวลปริปุณฺณํ ปริสุทฺธํ พฺรหฺมจริยํ ปกาเสติ, ตตฺถ เย โข ภิกฺขู เสขา อปฺปตฺตมานสา อนุตฺตรํ โยคกฺเขมํ ปตฺถยมานา วิหรนฺติ, เต ตํ ธมฺมํ สุตฺวา วีริยํ อารภนฺติ อปฺปตฺตสฺส ปตฺติยา}

\chapterheadpageread{1. สมฺโพธิวคฺค}
\prose{\csromanfolio{173}อนธิคตสฺส อธิคมาย อสจฺฉิกตสฺส สจฺฉิกิริยาย.\csromanspacer{} เย ปน ตตฺถ ภิกฺขู อรหนฺโต ขีณาสวา วุสิตวนฺโต กตกรณียา โอหิตภารา อนุปฺปตฺตสทตฺถา ปริกฺขีณภวสํโยชนา สมฺมทญฺญาวิมุตฺตา, เต ตํ ธมฺมํ สุตฺวา ทิฏฺฐธมฺมสุขวิหารํเยว อนุยุตฺตา วิหรนฺติ.\csromanspacer{} อยํ อาวุโส ปญฺจโม อานิสํโส กาเลน ธมฺมสฺสวเน กาเลน ธมฺมสากจฺฉาย.\csromanspacer{} อิเม โข อาวุโส ปญฺจ อานิสํสา กาเลน ธมฺมสฺสวเน กาเลน ธมฺมสากจฺฉายาติ.\csromanspacer{}\csromanspacer{}จตุตฺถํ.}

\csromanheader{2}{5. พลสุตฺต}
\proseitem{5}{จตฺตาริมานิ ภิกฺขเว พลานิ.\csromanspacer{} กตมานิ จตฺตาริ? ปญฺญาพลํ วีริยพลํ อนวชฺชพลํ สงฺคาหพลํ.\csromanspacer{} กตมญฺจ ภิกฺขเว ปญฺญาพลํ? เย ธมฺมา กุสลา กุสลสงฺขาตา, เย ธมฺมา อกุสลา อกุสลสงฺขาตา, เย ธมฺมา สาวชฺชา สาวชฺชสงฺขาตา, เย ธมฺมา อนวชฺชา อนวชฺชสงฺขาตา, เย ธมฺมา กณฺหา กณฺหสงฺขาตา, เย ธมฺมา สุกฺกา สุกฺกสงฺขาตา, เย ธมฺมา เสวิตพฺพา เสวิตพฺพสงฺขาตา, เย ธมฺมา อเสวิตพฺพา อเสวิตพฺพสงฺขาตา, เย ธมฺมา นาลมริยา นาลมริยสงฺขาตา, เย ธมฺมา อลมริยา อลมริยสงฺขาตา, ตฺยสฺส ธมฺมา ปญฺญาย โวทิฏฺฐา โหนฺติ โวจริตา.\csromanspacer{} อิทํ วุจฺจติ ภิกฺขเว ปญฺญาพลํ.}

\prose{กตมญฺจ ภิกฺขเว วีริยพลํ? เย ธมฺมา อกุสลา อกุสลสงฺขาตา, เย ธมฺมา สาวชฺชา สาวชฺชสงฺขาตา, เย ธมฺมา กณฺหา กณฺหสงฺขาตา, เย ธมฺมา อเสวิตพฺพา อเสวิตพฺพสงฺขาตา, เย ธมฺมา นาลมริยา นาลมริยสงฺขาตา, เตสํ ธมฺมานํ ปหานาย ฉนฺทํ ชเนติ วายมติ วีริยํ อารภติ จิตฺตํ ปคฺคณฺหาติ ปทหติ.\csromanspacer{} เย ธมฺมา กุสลา กุสลสงฺขาตา, เย ธมฺมา อนวชฺชา อนวชฺชสงฺขาตา, เย ธมฺมา สุกฺกา สุกฺกสงฺขาตา, เย ธมฺมา เสวิตพฺพา เสวิตพฺพสงฺขาตา, เย ธมฺมา อลมริยา อลมริยสงฺขาตา, เตสํ ธมฺมานํ ปฏิลาภาย ฉนฺทํ ชเนติ วายมติ วิริยํ อารภติ จิตฺตํ ปคฺคณฺหาติ ปทหติ.\csromanspacer{} อิทํ วุจฺจติ ภิกฺขเว วีริยพลํ.}

\prose{กตมญฺจ ภิกฺขเว อนวชฺชพลํ? อิธ ภิกฺขเว อริยสาวโก อนวชฺเชน กายกมฺเมน สมนฺนาคโต โหติ, อนวชฺเชน วจีกมฺเมน สมนฺนาคโต โหติ, อนวชฺเชน มโนกมฺเมน สมนฺนาคโต โหติ.\csromanspacer{} อิทํ วุจฺจติ ภิกฺขเว อนวชฺชพลํ.}

\gambhira{นวกนิปาตปาฬิ}
\prose{\csromanfolio{174}กตมญฺจ ภิกฺขเว สงฺคาหพลํ, จตฺตาริมานิ ภิกฺขเว สงฺคหวตฺถูนิ ทานํ เปยฺยวชฺชํ อตฺถจริยา สมานตฺตตา.\csromanspacer{} เอตทคฺคํ ภิกฺขเว ทานานํ ยทิทํ ธมฺมทานํ.\csromanspacer{} เอตทคฺคํ ภิกฺขเว เปยฺยวชฺชานํ ยทิทํ อตฺถิกสฺส โอหิตโสตสฺส ปุนปฺปุนํ ธมฺมํ เทเสติ.\csromanspacer{} เอตทคฺคํ ภิกฺขเว อตฺถจริยานํ ยทิทํ อสฺสทฺธํ สทฺธาสมฺปทาย สมาทเปติ นิเวเสติ ปติฏฺฐาเปติ.\csromanspacer{} ทุสฺสีลํ สีลสมฺปทาย.\csromanspacer{} มจฺฉริํ จาคสมฺปทาย.\csromanspacer{} ทุปฺปญฺญํ ปญฺญาสมฺปทาย สมาทเปติ นิเวเสติ ปติฏฺฐาเปติ.\csromanspacer{} เอตทคฺคํ ภิกฺขเว สมานตฺตตานํ ยทิทํ โสตาปนฺโน โสตาปนฺนสฺส สมานตฺโต, สกทาคามี สกทาคามิสฺส สมานตฺโต.\csromanspacer{} อนาคามี อนาคามิสฺส สมานตฺโต, อรหา อรหโต สมานตฺโต.\csromanspacer{} อิทํ วุจฺจติ ภิกฺขเว สงฺคาหพลํ.\csromanspacer{} อิมานิ โข ภิกฺขเว จตฺตาริ พลานิ.}

\prose{อิเมหิ โข ภิกฺขเว จตูหิ พเลหิ สมนฺนาคโต อริยสาวโก ปญฺจ ภยานิ สมติกฺกนฺโต โหติ.\csromanspacer{} กตมานิ ปญฺจ? อาชีวิกภยํ อสิโลกภยํ ปริสสารชฺชภยํ มรณภยํ ทุคฺคติภยํ.\csromanspacer{} ส โข โส ภิกฺขเว อริยสาวโก อิติ ปฏิสญฺจิกฺขติ “นาหํ อาชีวิกภยสฺส ภายามิ.\csromanspacer{} กิสฺสาหํ อาชีวิกภยสฺส ภายิสฺสามิ.\csromanspacer{} อตฺถิ เม จตฺตาริ พลานิ ปญฺญาพลํ วีริยพลํ อนวชฺชพลํ สงฺคาหพลํ.\csromanspacer{} ทุปฺปญฺโญ โข อาชีวิกภยสฺส ภาเยยฺย, กุสีโต อาชีวิกภยสฺส ภาเยยฺย, สาวชฺชกายกมฺมนฺตวจีกมฺมนฺต- มโนกมฺมนฺโต อาชีวิกภยสฺส ภาเยยฺย, อสงฺคาหโก อาชีวิกภยสฺส ภาเยยฺย, นาหํ อสิโลกภยสฺส ภายามิ ฯเปฯ นาหํ ปริสสารชฺชภยสฺส ภายามิ, นาหํ มรณภยสฺส ภายามิ ฯเปฯ นาหํ ทุคฺคติภยสฺส ภายามิ.\csromanspacer{} กิสฺสาหํ ทุคฺคติภยสฺส ภายิสฺสามิ, อตฺถิ เม จตฺตาริ พลานิ ปญฺญาพลํ วีริยพลํ อนวชฺชพลํ สงฺคาหพลํ.\csromanspacer{} ทุปฺปญฺโญ โข ทุคฺคติภยสฺส ภาเยยฺย, กุสีโต ทุคฺคติภยสฺส ภาเยยฺย, สาวชฺชกายกมฺมนฺตวจีกมฺมนฺตมโนกมฺมนฺโต ทุคฺคติภยสฺส ภาเยยฺย, อสงฺคาหโก ทุคฺคติภยสฺส ภาเยยฺย”.\csromanspacer{} อิเมหิ โข ภิกฺขเว จตูหิ พเลหิ สมนฺนาคโต อริยสาวโก อิมานิ ปญฺจ ภยานิ สมติกฺกนฺโต โหตีติ.\csromanspacer{}\csromanspacer{}ปญฺจมํ.}

\csromanheader{2}{6. เสวนาสุตฺต}
\proseitem{6}{ตตฺร โข อายสฺมา สาริปุตฺโต ภิกฺขู อามนฺเตสิ ฯเปฯ.\csromanspacer{} อายสฺมา สาริปุตฺโต เอตทโวจ\csromandash{}}

\chapterheadpageread{1. สมฺโพธิวคฺค}
\prose{\csromanfolio{175}ปุคฺคโลปิ อาวุโส ทุวิเธน เวทิตพฺโพ เสวิตพฺโพปิ อเสวิตพฺโพปิ, จีวรมฺปิ อาวุโส ทุวิเธน เวทิตพฺพํ เสวิตพฺพมฺปิ อเสวิตพฺพมฺปิ, ปิณฺฑปาโตปิ อาวุโส ทุวิเธน เวทิตพฺโพ เสวิตพฺโพปิ อเสวิตพฺโพปิ, เสนาสนมฺปิ อาวุโส ทุวิเธน เวทิตพฺพํ เสวิตพฺพมฺปิ อเสวิตพฺพมฺปิ, คามนิคโมปิ อาวุโส ทุวิเธน เวทิตพฺโพ เสวิตพฺโพปิ อเสวิตพฺโพปิ, ชนปทปเทโสปิ อาวุโส ทุวิเธน เวทิตพฺโพ เสวิตพฺโพปิ อเสวิตพฺโพปิ.}

\prose{“ปุคฺคโลปิ อาวุโส ทุวิเธน เวทิตพฺโพ เสวิตพฺโพ อเสวิตพฺโพปี”ติ อิติ โข ปเนตํ วุตฺตํ, กิญฺเจตํ ปฏิจฺจ วุตฺตํ.\csromanspacer{} ตตฺถ ยํ ชญฺญา ปุคฺคลํ “อิมํ โข เม ปุคฺคลํ เสวโต อกุสลา ธมฺมา อภิวฑฺฒนฺติ, กุสลา ธมฺมา ปริหายนฺติ.\csromanspacer{} เย จ โข เม ปพฺพชิเตน ชีวิตปริกฺขารา สมุทาเนตพฺพา จีวรปิณฺฑปาต เสนาสน คิลานปจฺจย เภสชฺชปริกฺขารา, เต จ กสิเรน สมุทาคจฺฉนฺติ.\csromanspacer{} ยสฺส จมฺหิ อตฺถาย อคารสฺมา อนคาริยํ ปพฺพชิโต, โส จ เม สามญฺญตฺโถ น ภาวนาปาริปูริํ คจฺฉตี”ติ.\csromanspacer{} เตนาวุโส ปุคฺคเลน โส ปุคฺคโล รตฺติภาคํ วา ทิวสภาคํ วา สงฺขาปิ อนาปุจฺฉา ปกฺกมิตพฺพํ, นานุพนฺธิตพฺโพ.}

\prose{ตตฺถ ยํ ชญฺญา ปุคฺคลํ “อิมํ โข เม ปุคฺคลํ เสวโต อกุสลา ธมฺมา อภิวฑฺฒนฺติ, กุสลา ธมฺมา ปริหายนฺติ.\csromanspacer{} เย จ โข เม ปพฺพชิเตน ชีวิตปริกฺขารา สมุทาเนตพฺพา จีวรปิณฺฑปาตเสนาสนคิลาน- ปจฺจยเภสชฺชปริกฺขารา, เต จ อปฺปกสิเรน สมุทาคจฺฉนฺติ.\csromanspacer{} ยสฺส จมฺหิ อตฺถาย อคารสฺมา อนคาริยํ ปพฺพชิโต, โส จ เม สามญฺญตฺโถ น ภาวนาปาริปูริํ คจฺฉตี”ติ.\csromanspacer{} เตนาวุโส ปุคฺคเลน โส ปุคฺคโล สงฺขาปิ อนาปุจฺฉา ปกฺกมิตพฺพํ, นานุพนฺธิตพฺโพ.}

\prose{ตตฺถ ยํ ชญฺญา ปุคฺคลํ “อิมํ โข เม ปุคฺคลํ เสวโต อกุสลา ธมฺมา ปริหายนฺติ, กุสลา ธมฺมา อภิวฑฺฒนฺติ.\csromanspacer{} เย จ โข เม ปพฺพชิเตน ชีวิตปริกฺขารา สมุทาเนตพฺพา จีวรปิณฺฑปาตเสนาสนคิลาน- ปจฺจยเภสชฺชปริกฺขารา, เต จ กสิเรน สมุทาคจฺฉนฺติ.\csromanspacer{} ยสฺส จมฺหิ อตฺถาย อคารสฺมา อนคาริยํ ปพฺพชิโต, โส จ เม สามญฺญตฺโถ ภาวนาปาริปูริํ คจฺฉตี”ติ.\csromanspacer{} เตนาวุโส ปุคฺคเลน โส ปุคฺคโล สงฺขาปิ อนุพนฺธิตพฺโพ, น ปกฺกมิตพฺพํ.}

\gambhira{นวกนิปาตปาฬิ}
\prose{\csromanfolio{176}ตตฺถ ยํ ชญฺญา ปุคฺคลํ “อิมํ โข เม ปุคฺคลํ เสวโต อกุสลา ธมฺมา ปริหายนฺติ, กุสลา ธมฺมา อภิวฑฺฒนฺติ.\csromanspacer{} เย จ โข เม ปพฺพชิเตน ชีวิตปริกฺขารา สมุทาเนตพฺพา จีวรปิณฺฑปาตเสนาสนคิลาน- ปจฺจยเภสชฺชปริกฺขารา, เต จ อปฺปกสิเรน สมุทาคจฺฉนฺติ.\csromanspacer{} ยสฺส จมฺหิ อตฺถาย อคารสฺมา อนคาริยํ ปพฺพชิโต, โส จ เม สามญฺญตฺโถ ภาวนา ปาริปูริํ คจฺฉตี”ติ.\csromanspacer{} เตนาวุโส ปุคฺคเลน โส ปุคฺคโล ยาวชีวํ อนุพนฺธิตพฺโพ, น ปกฺกมิตพฺพํ อปิ ปนุชฺชมาเนน\footnote{ปณุชฺชมาเนน (?)}. “ปุคฺคโลปิ อาวุโส ทุวิเธน เวทิตพฺโพ เสวิตพฺโพปิ อเสวิตพฺโพปี”ติ อิติ ยํ ตํ วุตฺตํ, อิทเมตํ ปฏิจฺจ วุตฺตํ.}

\prose{“จีวรมฺปิ อาวุโส ทุวิเธน เวทิตพฺพํ เสวิตพฺพมฺปิ อเสวิตพฺพมฺปี”ติ อิติ โข ปเนตํ วุตฺตํ, กิญฺเจตํ ปฏิจฺจ วุตฺตํ.\csromanspacer{} ตตฺถ ยํ ชญฺญา จีวรํ “อิทํ โข เม จีวรํ เสวโต อกุสลา ธมฺมา อภิวฑฺฒนฺติ, กุสลา ธมฺมา ปริหายนฺตี”ติ.\csromanspacer{} เอวรูปํ จีวรํ น เสวิตพฺพํ.\csromanspacer{} ตตฺถ ยํ ชญฺญา จีวรํ “อิทํ โข เม จีวรํ เสวโต อกุสลา ธมฺมา ปริหายนฺติ, กุสลา ธมฺมา อภิวฑฺฒนฺตี”ติ.\csromanspacer{} เอวรูปํ จีวรํ เสวิตพฺพํ. “จีวรมฺปิ อาวุโส ทุวิเธน เวทิตพฺพํ เสวิตพฺพมฺปิ อเสวิตพฺพมฺปี”ติ อิติ ยํ ตํ วุตฺตํ, อิทเมตํ ปฏิจฺจ วุตฺตํ.}

\prose{“ปิณฺฑปาโตปิ อาวุโส ทุวิเธน เวทิตพฺโพ เสวิตพฺโพปิ อเสวิตพฺโพปี”ติ อิติ โข ปเนตํ วุตฺตํ, กิญฺเจตํ ปฏิจฺจ วุตฺตํ.\csromanspacer{} ตตฺถ ยํ ชญฺญา ปิณฺฑปาตํ “อิมํ โข เม ปิณฺฑปาตํ เสวโต อกุสลา ธมฺมา อภิวฑฺฒนฺติ, กุสลา ธมฺมา ปริหายนฺตี”ติ.\csromanspacer{} เอวรูโป ปิณฺฑปาโต น เสวิตพฺโพ.\csromanspacer{} ตตฺถ ยํ ชญฺญา ปิณฺฑปาตํ “อิมํ โข เม ปิณฺฑปาตํ เสวโต อกุสลา ธมฺมา ปริหายนฺติ, กุสลา ธมฺมา อภิวฑฺฒนฺตี”ติ.\csromanspacer{} เอวรูโป ปิณฺฑปาโต เสวิตพฺโพ. “ปิณฺฑปาโตปิ อาวุโส ทุวิเธน เวทิตพฺโพ เสวิตพฺโพปิ อเสวิตพฺโพปี”ติ อิติ ยํ ตํ วุตฺตํ, อิทเมตํ ปฏิจฺจ วุตฺตํ.}

\prose{“เสนาสนมฺปิ อาวุโส ทุวิเธน เวทิตพฺพํ เสวิตพฺพมฺปิ อเสวิตพฺพมฺปี”ติ อิติ โข ปเนตํ วุตฺตํ, กิญฺเจตํ ปฏิจฺจ วุตฺตํ.\csromanspacer{} ตตฺถ ยํ ชญฺญา เสนาสนํ “อิทํ โข เม เสนาสนํ เสวโต อกุสลา ธมฺมา อภิวฑฺฒนฺติ, กุสลา ธมฺมา ปริหายนฺตี”ติ.\csromanspacer{} เอวรูปํ เสนาสนํ น เสวิตพฺพํ.\csromanspacer{} ตตฺถ ยํ ชญฺญา เสนาสนํ “อิทํ โข เม เสนาสนํ เสวโต อกุสลา ธมฺมา}

\chapterheadpageread{1. สมฺโพธิวคฺค}
\prose{\csromanfolio{177}ปริหายนฺติ, กุสลา ธมฺมา อภิวฑฺฒนฺตี”ติ.\csromanspacer{} เอวรูปํ เสนาสนํ เสวิตพฺพํ. “เสนาสนมฺปิ อาวุโส ทุวิเธน เวทิตพฺพํ เสวิตพฺพมฺปิ อเสวิตพฺพมฺปี”ติ อิติ ยํ ตํ วุตฺตํ, อิทเมตํ ปฏิจฺจ วุตฺตํ.}

\prose{“คามนิคโมปิ อาวุโส ทุวิเธน เวทิตพฺโพ เสวิตพฺโพปิ อเสวิตพฺโพปี”ติ อิติ โข ปเนตํ วุตฺตํ, กิญฺเจตํ ปฏิจฺจ วุตฺตํ, ตตฺถ ยํ ชญฺญา คามนิคมํ “อิมํ โข เม คามนิคมํ เสวโต อกุสลา ธมฺมา อภิวฑฺฒนฺติ, กุสลา ธมฺมา ปริหายนฺตี”ติ.\csromanspacer{} เอวรูโป คามนิคโม น เสวิตพฺโพ.\csromanspacer{} ตตฺถ ยํชญฺญา คามนิคมํ “อิมํ โข เม คามนิคมํ เสวโต อกุสลา ธมฺมา ปริหายนฺติ, กุสลา ธมฺมา อภิวฑฺฒนฺตี”ติ.\csromanspacer{} เอวรูโป คามนิคโม เสวิตพฺโพ. “คามนิคโมปิ อาวุโส ทุวิเธน เวทิตพฺโพ เสวิตพฺโพปิ อเสวิตพฺโพปี”ติ อิติ ยํ ตํ วุตฺตํ, อิทเมตํ ปฏิจฺจ วุตฺตํ.}

\prose{“ชนปทปเทโสปิ อาวุโส ทุวิเธน เวทิตพฺโพ เสวิตพฺโพปิ อเสวิตพฺโพปี”ติ อิติ โข ปเนตํ วุตฺตํ, กิญฺเจตํ ปฏิจฺจ วุตฺตํ, ตตฺถ ยํ ชญฺญา ชนปทปเทสํ “อิมํ โข เม ชนปทปเทสํ เสวโต อกุสลา ธมฺมา อภิวฑฺฒนฺติ, กุสลา ธมฺมา ปริหายนฺตี”ติ.\csromanspacer{} เอวรูโป ชนปทปเทโส น เสวิตพฺโพ.\csromanspacer{} ตตฺถ ยํ ชญฺญา ชนปทปเทสํ “อิมํ โข เม ชนปทปเทสํ เสวโต อกุสลา ธมฺมา ปริหายนฺติ, กุสลา ธมฺมา อภิวฑฺฒนฺตี”ติ.\csromanspacer{} เอวรูโป ชนปทปเทโส เสวิตพฺโพ. “ชนปทปเทโสปิ อาวุโส ทุวิเธน เวทิตพฺโพ เสวิตพฺโพปิ อเสวิตพฺโพปี”ติ อิติ ยํ ตํ วุตฺตํ, อิทเมตํ ปฏิจฺจ วุตฺตนฺติ.\csromanspacer{}\csromanspacer{}ฉฏฺฐํ.}

\csromanheader{2}{7. สุตวาสุตฺต}
\proseitem{7}{เอกํ สมยํ ภควา ราชคเห วิหรติ คิชฺฌกูเฏ ปพฺพเต.\csromanspacer{} อถ โข สุตวา ปริพฺพาชโก เยน ภควา เตนุปสงฺกมิ, อุปสงฺกมิตฺวา ภควตา สทฺธิํ สมฺโมทิ, สมฺโมทนียํ กถํ สารณียํ วีติสาเรตฺวา เอกมนฺตํ นิสีทิ, เอกมนฺตํ นิสินฺโน โข สุตวา ปริพฺพาชโก ภควนฺตํ เอตทโวจ\csromandash{}}

\prose{เอกมิทาหํ ภนฺเต สมยํ ภควา อิเธว ราชคเห วิหรามิ คิริพฺพเช, ตตฺร เม ภนฺเต ภควโต สมฺมุขา สุตํ สมฺมุขา ปฏิคฺคหิตํ “โย โส}

\vspace{\tipitakaparskip}
\csromancenter{นวกนิปาตปาฬิ สุตวา\footnote{สุตว (สฺยา)} ภิกฺขุ อรหํ ขีณาสโว วุสิตวา กตกรณีโย โอหิตภาโร อนุปฺปตฺตสทตฺโถ ปริกฺขีณภวสํโยชโน สมฺมทญฺญาวิมุตฺโต, อภพฺโพ โส ปญฺจ ฐานานิ อชฺฌาจริตุํ, อภพฺโพ ขีณาสโว ภิกฺขุ สญฺจิจฺจ ปาณํ ชีวิตา โวโรเปตุํ, อภพฺโพ ขีณาสโว ภิกฺขุ อทินฺนํ เถยฺยสงฺขาตํ อาทาตุํ, อภพฺโพ ขีณาสโว ภิกฺขุ เมถุนํ ธมฺมํ ปฏิเสวิตุํ, อภพฺโพ ขีณาสโว ภิกฺขุ สมฺปชานมุสา\footnote{สมฺปชานํ มุสา (ก-สี)} ภาสิตุํ, อภพฺโพ ขีณาสโว ภิกฺขุ สนฺนิธิการกํ กาเม ปริภุญฺชิตุํ.\csromanspacer{} เสยฺยถาปิ ปุพฺเพ อคาริยภูโต”ติ.\csromanspacer{} กจฺจิ เมตํ ภนฺเต ภควโต สุสฺสุตํ สุคฺคหิตํ สุมนสิกตํ สูปธาริตนฺติ.}
\prose{\csromanfolio{178}ตคฺฆ เต เอตํ สุตวา สุสฺสุตํ สุคฺคหิตํ สุมนสิกตํ สูปธาริตํ, ปุพฺเพ จาหํ สุตวา เอตรหิ จ เอวํ วทามิ\csromandash{}โย โส ภิกฺขุ อรหํ ขีณาสโว วุสิตวา กตกรณีโย โอหิตภาโร อนุปฺปตฺตสทตฺโถ ปริกฺขีณภวสํโยชโน สมฺมทญฺญาวิมุตฺโต, อภพฺโพ โส นว ฐานานิ อชฺฌาจริตุํ, อภพฺโพ ขีณาสโว ภิกฺขุ สญฺจิจฺจ ปาณํ ชีวิตา โวโรเปตุํ, อภพฺโพ ขีณาสโว ภิกฺขุ อทินฺนํ เถยฺยสงฺขาตํ อาทาตุํ, อภพฺโพ ขีณาสโว ภิกฺขุ เมถุนํ ธมฺมํ ปฏิเสวิตุํ, อภพฺโพ ขีณาสโว ภิกฺขุ สมฺปชานมุสา ภาสิตุํ, อภพฺโพ ขีณาสโว ภิกฺขุ สนฺนิธิการกํ กาเม ปริภุญฺชิตุํ เสยฺยถาปิ ปุพฺเพ อคาริยภูโต.\csromanspacer{} อภพฺโพ ขีณาสโว ภิกฺขุ ฉนฺทาคติํ คนฺตุํ, อภพฺโพ ขีณาสโว ภิกฺขุ โทสาคติํ คนฺตุํ, อภพฺโพ ขีณาสโว ภิกฺขุ โมหาคติํ คนฺตุํ, อภพฺโพ ขีณาสโว ภิกฺขุ ภยาคติํ คนฺตุํ.\csromanspacer{} ปุพฺเพ จาหํ สุตวา เอตรหิ จ เอวํ วทามิ\csromandash{}โย โส ภิกฺขุ อรหํ ขีณาสโว วุสิตวา กตกรณีโย โอหิตภาโร อนุปฺปตฺตสทตฺโถ ปริกฺขีณภวสํโยชโน สมฺมทญฺญาวิมุตฺโต, อภพฺโพ โส อิมานิ นว ฐานานิ อชฺฌาจริตุนฺติ.\csromanspacer{}\csromanspacer{}สตฺตมํ.}

\csromanheader{2}{8. สชฺฌสุตฺต}
\proseitem{8}{เอกํ สมยํ ภควา ราชคเห วิหรติ คิชฺฌกูเฏ ปพฺพเต.\csromanspacer{} อถ โข สชฺโฌ ปริพฺพาชโก เยน ภควา เตนุปสงฺกมิ, อุปสงฺกมิตฺวา ภควตา สทฺธิํ สมฺโมทิ, สมฺโมทนียํ กถํ สารณียํ วีติสาเรตฺวา}

\chapterheadpageread{1. สมฺโพธิวคฺค}
\prose{\csromanfolio{179}เอกมนฺตํ นิสีทิ, เอกมนฺตํ นิสินฺโน โข สชฺโฌ ปริพฺพาชโก ภควนฺตํ เอตทโวจ\csromandash{}}

\prose{เอกมิทาหํ ภนฺเต สมยํ ภควา อิเธว ราชคเห วิหรามิ คิริพฺพเช, ตตฺร เม ภนฺเต ภควโต สมฺมุขา สุตํ สมฺมุขา ปฏิคฺคหิตํ “โย โส สชฺฌ ภิกฺขุ อรหํ ขีณาสโว วุสิตวา กตกรณีโย โอหิตภาโร อนุปฺปตฺตสทตฺโถ ปริกฺขีณภวสํโยชโน สมฺมทญฺญาวิมุตฺโต, อภพฺโพ โส ปญฺจ ฐานานิ อชฺฌาจริตุํ, อภพฺโพ ขีณาสโว ภิกฺขุ สญฺจิจฺจ ปาณํ ชีวิตา โวโรเปตุํ, อภพฺโพ ขีณาสโว ภิกฺขุ อทินฺนํ เถยฺยสงฺขาตํ อาทาตุํ, อภพฺโพ ขีณาสโว ภิกฺขุ เมถุนํ ธมฺมํ ปฏิเสวิตุํ, อภพฺโพ ขีณาสโว ภิกฺขุ สมฺปชานมุสา ภาสิตุํ, อภพฺโพ ขีณาสโว ภิกฺขุ สนฺนิธิการกํ กาเม ปริภุญฺชิตุํ เสยฺยถาปิ ปุพฺเพ อคาริยภูโต”ติ.\csromanspacer{} กจฺจิ เมตํ ภนฺเต ภควโต สุสฺสุตํ สุคฺคหิตํ สุมนสิกตํ สูปธาริตนฺติ.}

\prose{ตคฺฆ เต เอตํ สชฺฌ สุสฺสุตํ สุคฺคหิตํ สุมนสิกตํ สูปธาริตํ, ปุพฺเพ จาหํ สชฺฌ เอตรหิ จ เอวํ วทามิ\csromandash{}โย โส ภิกฺขุ อรหํ ขีณาสโว วุสิตวา กตกรณีโย โอหิตภาโร อนุปฺปตฺตสทตฺโถ ปริกฺขีณภวสํโยชโน สมฺมทญฺญาวิมุตฺโต, อภพฺโพ โส นว ฐานานิ อชฺฌาจริตุํ, อภพฺโพ ขีณาสโว ภิกฺขุ สญฺจิจฺจ ปาณํ ชีวิตา โวโรเปตุํ ฯเปฯ อภพฺโพ ขีณาสโว ภิกฺขุ สนฺนิธิการกํ กาเม ปริภุญฺชิตุํ เสยฺยถาปิ ปุพฺเพ อคาริยภูโต.\csromanspacer{} อภพฺโพ ขีณาสโว ภิกฺขุ พุทฺธํ ปจฺจกฺขาตุํ, อภพฺโพ ขีณาสโว ภิกฺขุ ธมฺมํ ปจฺจกฺขาตุํ, อภพฺโพ ขีณาสโว ภิกฺขุ สํฆํ ปจฺจกฺขาตุํ, อภพฺโพ ขีณาสโว ภิกฺขุ สิกฺขํ ปจฺจกฺขาตุํ.\csromanspacer{} ปุพฺเพ จาหํ สชฺฌ เอตรหิ จ เอวํ วทามิ\csromandash{}โย โส ภิกฺขุ อรหํ ขีณาสโว วุสิตวา กตกรณีโย โอหิตภาโร อนุปฺปตฺตสทตฺโถ ปริกฺขีณภวสํโยชโน สมฺมทญฺญาวิมุตฺโต, อภพฺโพ โส อิมานิ นว ฐานานิ อชฺฌาจริตุนฺติ.\csromanspacer{}\csromanspacer{}อฏฺฐมํ.}

\csromanheader{2}{9. ปุคฺคลสุตฺต}
\proseitem{9}{นวยิเม ภิกฺขเว ปุคฺคลา สนฺโต สํวิชฺชมานา โลกสฺมิํ.\csromanspacer{} กตเม นว? อรหา, อรหตฺตาย ปฏิปนฺโน, อนาคามี, อนาคามิผลสจฺฉิกิริยาย ปฏิปนฺโน, สกทาคามี, สกทาคามิผลสจฺฉิกิริยาย}

\vspace{\tipitakaparskip}
\csromancenter{นวกนิปาตปาฬิ ปฏิปนฺโน, โสตาปนฺโน, โสตาปตฺติผลสจฺฉิกิริยาย ปฏิปนฺโน, ปุถุชฺชโน, อิเม โข ภิกฺขเว นว ปุคฺคลา สนฺโต สํวิชฺชมานา โลกสฺมินฺติ.\csromanspacer{}\csromanspacer{}นวมํ.}
\csromanheader{2}{10. อาหุเนยฺยสุตฺต}
\proseitem{10}{\csromanfolio{180}นวยิเม ภิกฺขเว ปุคฺคลา อาหุเนยฺยา ปาหุเนยฺยา ทกฺขิเณยฺยา อญฺชลิกรณียา อนุตฺตรํ ปุญฺญกฺเขตฺตํ โลกสฺส.\csromanspacer{} กตเม นว? อรหา, อรหตฺตาย ปฏิปนฺโน, อนาคามี, อนาคามิผลสจฺฉิกิริยาย ปฏิปนฺโน, สกทาคามี, สกทาคามิผลสจฺฉิกิริยาย ปฏิปนฺโน, โสตาปนฺโน, โสตาปตฺติผลสจฺฉิกิริยาย ปฏิปนฺโน, โคตฺรภู.\csromanspacer{} อิเม โข ภิกฺขเว นว ปุคฺคลา อาหุเนยฺยา ฯเปฯ อนุตฺตรํ ปุญฺญกฺเขตฺตํ โลกสฺสาติ.\csromanspacer{}\csromanspacer{}ทสมํ.\csromanspacer{} สมฺโพธิวคฺโค ปฐโม.}

\titlehead{\textbf{ตสฺสุทฺทานํ}}
\csromangathameasure{สมฺโพธิ นิสฺสโย เจว, เมฆิย นนฺทกํ พลํ. \\ เสวนา สุตวา สชฺโฌ, ปุคฺคโล อาหุเนยฺเยน จาติ.}
\csromangathagroup{\csromangathastanza{สมฺโพธิ นิสฺสโย เจว, เมฆิย นนฺทกํ พลํ. \\ เสวนา สุตวา สชฺโฌ, ปุคฺคโล อาหุเนยฺเยน จาติ.}}

\chapterhead{2. สีหนาทวคฺค}
\csromanheader{2}{1. สีหนาทสุตฺต}
\proseitem{11}{เอกํ สมยํ ภควา สาวตฺถิยํ วิหรติ เชตวเน อนาถปิณฺฑิกสฺส อาราเม.\csromanspacer{} อถ โข อายสฺมา สาริปุตฺโต เยน ภควา เตนุปสงฺกมิ, อุปสงฺกมิตฺวา ภควนฺตํ อภิวาเทตฺวา เอกมนฺตํ นิสีทิ, เอกมนฺตํ นิสินฺโน โข อายสฺมา สาริปุตฺโต ภควนฺตํ เอตทโวจ “วุตฺโถ เม ภนฺเต สาวตฺถิยํ วสฺสาวาโส, อิจฺฉามหํ ภนฺเต ชนปทจาริกํ ปกฺกมิตุนฺ”ติ.\csromanspacer{} ยสฺสทานิ ตฺวํ สาริปุตฺต กาลํ มญฺญสีติ.\csromanspacer{} อถ โข อายสฺมา สาริปุตฺโต อุฏฺฐายาสนา ภควนฺตํ อภิวาเทตฺวา ปทกฺขิณํ กตฺวา ปกฺกามิ.\csromanspacer{} อถ โข อญฺญตโร ภิกฺขุ อจิรปกฺกนฺเต อายสฺมนฺเต สาริปุตฺเต ภควนฺตํ เอตทโวจ “อายสฺมา มํ ภนฺเต สาริปุตฺโต อาสชฺช อปฺปฏินิสฺสชฺช จาริกํ ปกฺกนฺโต”ติ.\csromanspacer{} อถ โข ภควา อญฺญตรํ ภิกฺขุํ อามนฺเตสิ}

\chapterheadpageread{2. สีหนาทวคฺค}
\prose{\csromanfolio{181}“เอหิ ตฺวํ ภิกฺขุ มม วจเนน สาริปุตฺตํ อามนฺเตหิ ‘สตฺถา ตํ อาวุโส สาริปุตฺต อามนฺเตตี’ติ”. “เอวํ ภนฺเต”ติ โข โส ภิกฺขุ ภควโต ปฏิสฺสุตฺวา เยนายสฺมา สาริปุตฺโต เตนุปสงฺกมิ, อุปสงฺกมิตฺวา อายสฺมนฺตํ สาริปุตฺตํ เอตทโวจ “สตฺถา ตํ อาวุโส สาริปุตฺต อามนฺเตตี”ติ. “เอวมาวุโส”ติ โข อายสฺมา สาริปุตฺโต ตสฺส ภิกฺขุโน ปจฺจสฺโสสิ.}

\prose{เตน โข ปน สมเยน อายสฺมา จ มหาโมคฺคลฺลาโน\footnote{มหาโมคฺคลาโน (ก)} อายสฺมา จ อานนฺโท อวาปุรณํ\footnote{อปาปุรณํ (สฺยา, ก)} อาทาย วิหาเร อาหิณฺฑนฺติ\footnote{วิหาเรน วิหารํ อนฺวาหิณฺฑนฺติ (สี, อิ), วิหารํ อาหิณฺฑนฺติ (สฺยา)}. “อภิกฺกมถายสฺมนฺโต อภิกฺกมถายสฺมนฺโต อิทานายสฺมา สาริปุตฺโต ภควโต สมฺมุขา สีหนาทํ นทิสฺสตี”ติ.\csromanspacer{} อถ โข อายสฺมา สาริปุตฺโต เยน ภควา เตนุปสงฺกมิ, อุปสงฺกมิตฺวา ภควนฺตํ อภิวาเทตฺวา เอกมนฺตํ นิสีทิ, เอกมนฺตํ นิสินฺนํ โข อายสฺมนฺตํ สาริปุตฺตํ ภควา เอตทโวจ “อิธ เต สาริปุตฺต อญฺญตโร สพฺรหฺมจารี ขียนธมฺมํ อาปนฺโน ‘อายสฺมา มํ ภนฺเต สาริปุตฺโต อาสชฺช อปฺปฏินิสฺสชฺช จาริกํ ปกฺกนฺโต’ติ”.}

\prose{ยสฺส นูน ภนฺเต กาเย กายคตาสติ อนุปฏฺฐิตา อสฺส, โส อิธ อญฺญตรํ สพฺรหฺมจาริํ อาสชฺช อปฺปฏินิสฺสชฺช จาริกํ ปกฺกเมยฺย.}

\prose{เสยฺยถาปิ ภนฺเต ปถวิยํ สุจิมฺปิ นิกฺขิปนฺติ, อสุจิมฺปิ นิกฺขิปนฺติ, คูถคตมฺปิ นิกฺขิปนฺติ, มุตฺตคตมฺปิ นิกฺขิปนฺติ, เขฬคตมฺปิ นิกฺขิปนฺติ, ปุพฺพคตมฺปิ นิกฺขิปนฺติ, โลหิตคตมฺปิ นิกฺขิปนฺติ.\csromanspacer{} น จ เตน ปถวี อฏฺฏียติ วา หรายติ วา ชิคุจฺฉติ วา.\csromanspacer{} เอวเมวํ โข อหํ ภนฺเต ปถวีสเมน เจตสา วิหรามิ วิปุเลน มหคฺคเตน อปฺปมาเณน อเวเรน อพฺยาปชฺเชน.\csromanspacer{} ยสฺส นูน ภนฺเต กาเย กายคตาสติ อนุปฏฺฐิตา อสฺส, โส อิธ อญฺญตรํ สพฺรหฺมจาริํ อาสชฺช อปฺปฏินิสฺสชฺช จาริกํ ปกฺกเมยฺย. (1)}

\prose{เสยฺยถาปิ ภนฺเต อาปสฺมิํ สุจิมฺปิ โธวนฺติ, อสุจิมฺปิ โธวนฺติ, คูถคตมฺปิ, มุตฺตคตมฺปิ, เขฬคตมฺปิ, ปุพฺพคตมฺปิ, โลหิตคตมฺปิ โธวนฺติ.\csromanspacer{} น จ เตน อาโป อฏฺฏียติ วา หรายติ วา ชิคุจฺฉติ วา.\csromanspacer{} เอวเมวํ โข อหํ ภนฺเต อาโปสเมน เจตสา วิหรามิ วิปุเลน มหคฺคเตน}

\vspace{\tipitakaparskip}
\csromancenter{นวกนิปาตปาฬิ อปฺปมาเณน อเวเรน อพฺยาปชฺเชน.\csromanspacer{} ยสฺส นูน ภนฺเต กาเย กายคตาสติ อนุปฏฺฐิตา อสฺส, โส อิธ อญฺญตรํ สพฺรหฺมจาริํ อาสชฺช อปฺปฏินิสฺสชฺช จาริกํ ปกฺกเมยฺย. (2)}
\prose{\csromanfolio{182}เสยฺยถาปิ ภนฺเต เตโช สุจิมฺปิ ฑหติ, อสุจิมฺปิ ฑหติ, คูถคตมฺปิ, มุตฺตคตมฺปิ, เขฬคตมฺปิ, ปุพฺพคตมฺปิ, โลหิตคตมฺปิ ฑหติ.\csromanspacer{} น จ เตน เตโช อฏฺฏียติ วา หรายติ วา ชิคุจฺฉติ วา.\csromanspacer{} เอวเมวํ โข อหํ ภนฺเต เตโชสเมน เจตสา วิหรามิ วิปุเลน มหคฺคเตน อปฺปมาเณน อเวเรน อพฺยาปชฺเชน.\csromanspacer{} ยสฺส นูน ภนฺเต กาเย กายคตาสติ อนุปฏฺฐิตา อสฺส, โส อิธ อญฺญตรํ สพฺรหฺมจาริํ อาสชฺช อปฺปฏินิสฺสชฺช จาริกํ ปกฺกเมยฺย. (3)}

\prose{เสยฺยถาปิ ภนฺเต วาโย สุจิมฺปิ อุปวายติ, อสุจิมฺปิ อุปวายติ, คูถคตมฺปิ, มุตฺตคตมฺปิ, เขฬคตมฺปิ, ปุพฺพคตมฺปิ, โลหิตคตมฺปิ อุปวายติ.\csromanspacer{} น จ เตน วาโย อฏฺฏียติ วา หรายติ วา ชิคุจฺฉติ วา.\csromanspacer{} เอวเมวํ โข อหํ ภนฺเต วาโยสเมน เจตสา วิหรามิ วิปุเลน มหคฺคเตน อปฺปมาเณน อเวเรน อพฺยาปชฺเชน.\csromanspacer{} ยสฺส นูน ภนฺเต กาเย กายคตาสติ อนุปฏฺฐิตา อสฺส, โส อิธ อญฺญตรํ สพฺรหฺมจาริํ อาสชฺช อปฺปฏินิสฺสชฺช จาริกํ ปกฺกเมยฺย. (4)}

\prose{เสยฺยถาปิ ภนฺเต รโชหรณํ สุจิมฺปิ ปุญฺฉติ, อสุจิมฺปิ ปุญฺฉติ, คูถคตมฺปิ, มุตฺตคตมฺปิ, เขฬคตมฺปิ, ปุพฺพคตมฺปิ, โลหิตคตมฺปิ ปุญฺฉติ.\csromanspacer{} น จ เตน รโชหรณํ อฏฺฏียติ วา หรายติ วา ชิคุจฺฉติ วา.\csromanspacer{} เอวเมวํ โข อหํ ภนฺเต รโชหรณสเมน เจตสา วิหรามิ วิปุเลน มหคฺคเตน อปฺปมาเณน อเวเรน อพฺยาปชฺเชน.\csromanspacer{} ยสฺส นูน ภนฺเต กาเย กายคตาสติ อนุปฏฺฐิตา อสฺส, โส อิธ อญฺญตรํ สพฺรหฺมจาริํ อาสชฺช อปฺปฏินิสฺสชฺช จาริกํ ปกฺกเมยฺย. (5)}

\prose{เสยฺยถาปิ ภนฺเต จณฺฑาลกุมารโก วา จณฺฑาลกุมาริกา วา กโฬปิหตฺโถ นนฺตกวาสี คามํ วา นิคมํ วา ปวิสนฺโต นีจจิตฺตํเยว อุปฏฺฐเปตฺวา ปวิสติ.\csromanspacer{} เอวเมวํ โข อหํ ภนฺเต จณฺฑาลกุมารกจณฺฑาลกุมาริกาสเมน เจตสา วิหรามิ วิปุเลน มหคฺคเตน อปฺปมาเณน อเวเรน อพฺยาปชฺเชน.\csromanspacer{} ยสฺส นูน ภนฺเต กาเย}

\chapterheadpageread{2. สีหนาทวคฺค}
\prose{\csromanfolio{183}กายคตาสติ อนุปฏฺฐิตา อสฺส, โส อิธ อญฺญตรํ สพฺรหฺมจาริํ อาสชฺช อปฺปฏินิสฺสชฺช จาริกํ ปกฺกเมยฺย. (6)}

\prose{เสยฺยถาปิ ภนฺเต อุสโภ ฉินฺนวิสาโณ สูรโต สุทนฺโต สุวินีโต รถิยาย รถิยํ สิงฺฆาฏเกน สิงฺฆาฏกํ อนฺวาหิณฺฑนฺโต น กิญฺจิ หิํสติ ปาเทน วา วิสาเณน วา.\csromanspacer{} เอวเมวํ โข อหํ ภนฺเต อุสภฉินฺนวิสาณสเมน เจตสา วิหรามิ วิปุเลน มหคฺคเตน อปฺปมาเณน อเวเรน อพฺยาปชฺเชน.\csromanspacer{} ยสฺส นูน ภนฺเต กาเย กายคตาสติ อนุปฏฺฐิตา อสฺส, โส อิธ อญฺญตรํ สพฺรหฺมจาริํ อาสชฺช อปฺปฏินิสฺสชฺช จาริกํ ปกฺกเมยฺย. (7)}

\prose{เสยฺยถาปิ ภนฺเต อิตฺถี วา ปุริโส วา ทหโร ยุวา มณฺฑนกชาติโก สีสํนฺหาโต อหิกุณเปน วา กุกฺกุรกุณเปน วา มนุสฺสกุณเปน วา กณฺเฐ อาสตฺเตน อฏฺฏีเยยฺย หราเยยฺย ชิคุจฺเฉยฺย.\csromanspacer{} เอวเมวํ โข อหํ ภนฺเต อิมินา ปูติกาเยน อฏฺฏียามิ หรายามิ ชิคุจฺฉามิ.\csromanspacer{} ยสฺส นูน ภนฺเต กาเย กายคตาสติ อนุปฏฺฐิตา อสฺส, โส อิธ อญฺญตรํ สพฺรหฺมจาริํ อาสชฺช อปฺปฏินิสฺสชฺช จาริกํ ปกฺกเมยฺย. (8)}

\prose{เสยฺยถาปิ ภนฺเต ปุริโส เมทกถาลิกํ ปริหเรยฺย ฉิทฺทาวฉิทฺทํ อุคฺฆรนฺตํ ปคฺฆรนฺตํ.\csromanspacer{} เอวเมวํ โข อหํ ภนฺเต อิมํ กายํ ปริหรามิ ฉิทฺทาวฉิทฺทํ อุคฺฆรนฺตํ ปคฺฆรนฺตํ.\csromanspacer{} ยสฺส นูน ภนฺเต กาเย กายคตาสติ อนุปฏฺฐิตา อสฺส, โส อิธ อญฺญตรํ สพฺรหฺมจาริํ อาสชฺช อปฺปฏินิสฺสชฺช จาริกํ ปกฺกเมยฺยาติ. (9)}

\prose{อถ โข โส ภิกฺขุ อุฏฺฐายาสนา เอกํสํ อุตฺตราสงฺคํ กริตฺวา ภควโต ปาเทสุ สิรสา นิปติตฺวา ภควนฺตํ เอตทโวจ “อจฺจโย มํ ภนฺเต อจฺจคมา ยถาพาลํ ยถามูฬฺหํ ยถา-อกุสลํ, โย อหํ อายสฺมนฺตํ สาริปุตฺตํ อสตา ตุจฺฉา มุสา อภูเตน อพฺภาจิกฺขิํ.\csromanspacer{} ตสฺส เม ภนฺเต ภควา อจฺจยํ อจฺจยโต ปฏิคฺคณฺหตุ อายติํ สํวรายา”ติ.\csromanspacer{} ตคฺฆ ตํ\footnote{ตฺวํ (สี, อิ)} ภิกฺขุ อจฺจโย อจฺจคมา ยถาพาลํ ยถามูฬฺหํ ยถา- อกุสลํ, โย ตฺวํ สาริปุตฺตํ อสตา ตุจฺฉา มุสา อภูเตน อพฺภาจิกฺขิ.\csromanspacer{} ยโต จ โข ตฺวํ ภิกฺขุ อจฺจยํ อจฺจยโต ทิสฺวา}

\vspace{\tipitakaparskip}
\csromancenter{นวกนิปาตปาฬิ ยถาธมฺมํ ปฏิกโรสิ, ตํ เต มยํ ปฏิคฺคณฺหาม.\csromanspacer{} วุฑฺฒิเหสา ภิกฺขุ อริยสฺส วินเย โย อจฺจยํ อจฺจยโต ทิสฺวา ยถาธมฺมํ ปฏิกโรติ, อายติํ สํวรํ อาปชฺชตีติ.}
\prose{\csromanfolio{184}อถ โข ภควา อายสฺมนฺตํ สาริปุตฺตํ อามนฺเตสิ “ขม สาริปุตฺต อิมสฺส โมฆปุริสสฺส ปุรา ตสฺส ตตฺเถว สตฺตธา มุทฺธา ผลตี”ติ\footnote{ผลิสฺสตีติ (ก-สี, สฺยา, อิ, ก) อฏฺฐกถาสุ ปน “ผลตีติ”อิเตฺวว ทิสฺสติ.}.\csromanspacer{} ขมามหํ ภนฺเต ตสฺส อายสฺมโต, สเจ มํ โส อายสฺมา เอวมาห.\csromanspacer{} ขมตุ จ เม โส อายสฺมาติ.\csromanspacer{}\csromanspacer{}ปฐมํ.}

\csromanheader{2}{2. ส-อุปาทิเสสสุตฺต}
\proseitem{12}{เอกํ สมยํ ภควา สาวตฺถิยํ วิหรติ เชตวเน อนาถปิณฺฑิกสฺส อาราเม.\csromanspacer{} อถ โข อายสฺมา สาริปุตฺโต ปุพฺพณฺหสมยํ นิวาเสตฺวา ปตฺตจีวรมาทาย สาวตฺถิํ ปิณฺฑาย ปาวิสิ.\csromanspacer{} อถ โข อายสฺมโต สาริปุตฺตสฺส เอตทโหสิ “อติปฺปโค โข ตาว สาวตฺถิยํ ปิณฺฑาย จริตุํ, ยํนูนาหํ เยน อญฺญติตฺถิยานํ ปริพฺพาชกานํ อาราโม เตนุปสงฺกเมยฺยนฺ”ติ.\csromanspacer{} อถ โข อายสฺมา สาริปุตฺโต เยน อญฺญติตฺถิยานํ ปริพฺพาชกานํ อาราโม เตนุปสงฺกมิ, อุปสงฺกมิตฺวา เตหิ อญฺญติตฺถิเยหิ ปริพฺพาชเกหิ สทฺธิํ สมฺโมทิ, สมฺโมทนียํ กถํ สารณียํ วีติสาเรตฺวา เอกมนฺตํ นิสีทิ.}

\prose{เตน โข ปน สมเยน เตสํ อญฺญติตฺถิยานํ ปริพฺพาชกานํ สนฺนิสินฺนานํ สนฺนิปติตานํ อยมนฺตรากถา อุทปาทิ “โย หิ โกจิ อาวุโส ส-อุปาทิเสโส กาลํ กโรติ.\csromanspacer{} สพฺโพ โส อปริมุตฺโต นิรยา, อปริมุตฺโต ติรจฺฉานโยนิยา, อปริมุตฺโต เปตฺติวิสยา, อปริมุตฺโต อปายทุคฺคติวินิปาตา”ติ.\csromanspacer{} อถ โข อายสฺมา สาริปุตฺโต เตสํ อญฺญติตฺถิยานํ ปริพฺพาชกานํ ภาสิตํ เนว อภินนฺทิ นปฺปฏิกฺโกสิ, อนภินนฺทิตฺวา อปฺปฏิกฺโกสิตฺวา อุฏฺฐายาสนา ปกฺกามิ “ภควโต สนฺติเก เอตสฺส ภาสิตสฺส อตฺถํ อาชานิสฺสามี”ติ.\csromanspacer{} อถ โข อายสฺมา สาริปุตฺโต สาวตฺถิยํ ปิณฺฑาย จริตฺวา ปจฺฉาภตฺตํ ปิณฺฑปาตปฏิกฺกนฺโต เยน ภควา เตนุปสงฺกมิ, อุปสงฺกมิตฺวา ภควนฺตํ}

\chapterheadpageread{2. สีหนาทวคฺค}
\prose{\csromanfolio{185}อภิวาเทตฺวา เอกมนฺตํ นิสีทิ, เอกมนฺตํ นิสินฺโน โข อายสฺมา สาริปุตฺโต ภควนฺตํ เอตทโวจ\csromandash{}}

\prose{อิธาหํ ภนฺเต ปุพฺพณฺหสมยํ นิวาเสตฺวา ปตฺตจีวรมาทาย สาวตฺถิํ ปิณฺฑาย ปาวิสิํ, ตสฺส มยฺหํ ภนฺเต เอตทโหสิ “อติปฺปโค โข ตาว สาวตฺถิยํ ปิณฺฑาย จริตุํ, ยํนูนาหํ เยน อญฺญติตฺถิยานํ ปริพฺพาชกานํ อาราโม เตนุปสงฺกเมยฺยนฺ”ติ.\csromanspacer{} อถ โข อหํ ภนฺเต เยน อญฺญติตฺถิยานํ ปริพฺพาชกานํ อาราโม เตนุปสงฺกมิํ, อุปสงฺกมิตฺวา เตหิ อญฺญติตฺถิเยหิ ปริพฺพาชเกหิ สทฺธิํ สมฺโมทิํ, สมฺโมทนียํ กถํ สรณียํ วีติสาเรตฺวา เอกมนฺตํ นิสีทิํ.\csromanspacer{} เตน โข ปน สมเยน เตสํ อญฺญติตฺถิยานํ ปริพฺพาชกานํ สนฺนิสินฺนานํ สนฺนิปติตานํ อยมนฺตรากถา อุทปาทิ “โย หิ โกจิ อาวุโส ส-อุปาทิเสโส กาลํ กโรติ, สพฺโพ โส อปริมุตฺโต นิรยา, อปริมุตฺโต ติรจฺฉานโยนิยา, อปริมุตฺโต เปตฺติวิสยา, อปริมุตฺโต อปายทุคฺคติวินิปาตา”ติ.\csromanspacer{} อถ โข อหํ ภนฺเต เตสํ อญฺญติตฺถิยานํ ปริพฺพาชกานํ ภาสิตํ เนว อภินนฺทิํ นปฺปฏิกฺโกสิํ, อนภินนฺทิตฺวา อปฺปฏิกฺโกสิตฺวา อุฏฺฐายาสนา ปกฺกมิํ “ภควโต สนฺติเก เอตสฺส ภาสิตสฺส อตฺถํ อาชานิสฺสามี”ติ.}

\prose{เก จ\footnote{เกจิ (สฺยา, อิ), เต จ (ก)} สาริปุตฺต อญฺญติตฺถิยา ปริพฺพาชกา พาลา อพฺยตฺตา, เก จ\footnote{เกจิ (สฺยา, อิ, ก) ม 1. 233; อํ 2. 206, 306 ปิฏฺเฐสุ ปาฬิยา สํสนฺเทตพฺพํ.} ส-อุปาทิเสสํ วา “ส-อุปาทิเสโส”ติ ชานิสฺสนฺติ, อนุปาทิเสสํ วา “อนุปาทิเสโส”ติ ชานิสฺสนฺติ.}

\prose{นวยิเม สาริปุตฺต ปุคฺคลา ส-อุปาทิเสสา กาลํ กุรุมานา ปริมุตฺตา นิรยา, ปริมุตฺตา ติรจฺฉานโยนิยา, ปริมุตฺตา เปตฺติวิสยา, ปริมุตฺตา อปายทุคฺคติวินิปาตา.\csromanspacer{} กตเม นว? อิธ สาริปุตฺต เอกจฺโจ ปุคฺคโล สีเลสุ ปริปูรการี โหติ, สมาธิสฺมิํ ปริปูรการี, ปญฺญาย มตฺตโส การี.\csromanspacer{} โส ปญฺจนฺนํ โอรมฺภาคิยานํ สํโยชนานํ ปริกฺขยา อนฺตราปรินิพฺพายี โหติ.\csromanspacer{} อยํ สาริปุตฺต ปฐโม ปุคฺคโล ส-อุปาทิเสโส กาลํ กุรุมาโน ปริมุตฺโต นิรยา, ปริมุตฺโต}

\vspace{\tipitakaparskip}
\csromancenter{นวกนิปาตปาฬิ ติรจฺฉานโยนิยา, ปริมุตฺโต เปตฺติวิสยา, ปริมุตฺโต อปายทุคฺคติวินิปาตา. (1)}
\prose{\csromanfolio{186}ปุน จปรํ สาริปุตฺต อิเธกจฺโจ ปุคฺคโล สีเลสุ ปริปูรการี โหติ, สมาธิสฺมิํ ปริปูรการี, ปญฺญาย มตฺตโส การี.\csromanspacer{} โส ปญฺจนฺนํ โอรมฺภาคิยานํ สํโยชนานํ ปริกฺขยา อุปหจฺจปรินิพฺพายี โหติ ฯเปฯ อสงฺขารปรินิพฺพายี โหติ ฯเปฯ สสงฺขารปรินิพฺพายี โหติ ฯเปฯ อุทฺธํโสโต โหติ อกนิฏฺฐคามี.\csromanspacer{} อยํ สาริปุตฺต ปญฺจโม ปุคฺคโล ส-อุปาทิเสโส กาลํ กุรุมาโน ปริมุตฺโต นิรยา, ปริมุตฺโต ติรจฺฉานโยนิยา, ปริมุตฺโต เปตฺติวิสยา, ปริมุตฺโต อปายทุคฺคติวินิปาตา. (2-5)}

\prose{ปุน จปรํ สาริปุตฺต อิเธกจฺโจ ปุคฺคโล สีเลสุ ปริปูรการี โหติ, สมาธิสฺมิํ มตฺตโส การี, ปญฺญาย มตฺตโส การี.\csromanspacer{} โส ติณฺณํ สํโยชนานํ ปริกฺขยา ราคโทสโมหานํ ตนุตฺตา สกทาคามี โหติ, สกิเทว อิมํ โลกํ อาคนฺตฺวา ทุกฺขสฺสนฺตํ กโรติ.\csromanspacer{} อยํ สาริปุตฺต ฉฏฺโฐ ปุคฺคโล ส-อุปาทิเสโส กาลํ กุรุมาโน ปริมุตฺโต นิรยา ฯเปฯ ปริมุตฺโต อปายทุคฺคตินิปาตา. (6)}

\prose{ปุน จปรํ สาริปุตฺต อิเธกจฺโจ ปุคฺคโล สีเลสุ ปริปูรการี โหติ, สมาธิสฺมิํ มตฺตโส การี, ปญฺญาย มตฺตโส การี.\csromanspacer{} โส ติณฺณํ สํโยชนานํ ปริกฺขยา เอกพีชี โหติ, เอกํเยว มานุสกํ ภวํ นิพฺพตฺเตตฺวา ทุกฺขสฺสนฺตํ กโรติ.\csromanspacer{} อยํ สาริปุตฺต สตฺตโม ปุคฺคโล ส-อุปาทิเสโส กาลํ กุรุมาโน ปริมุตฺโต นิรยา ฯเปฯ ปริมุตฺโต อปายทุคฺคติวินิปาตา. (7)}

\prose{ปุน จปรํ สาริปุตฺต อิเธกจฺโจ ปุคฺคโล สีเลสุ ปริปูรการี โหติ, สมาธิสฺมิํ มตฺตโส การี, ปญฺญาย มตฺตโส การี.\csromanspacer{} โส ติณฺณํ สํโยชนานํ ปริกฺขยา โกลํโกโล โหติ.\csromanspacer{} เทฺว วา ตีณิ วา กุลานิ สนฺธาวิตฺวา สํสริตฺวา ทุกฺขสฺสนฺตํ กโรติ.\csromanspacer{} อยํ สาริปุตฺต อฏฺฐโม ปุคฺคโล ส-อุปาทิเสโส กาลํ กุรุมาโน ปริมุตฺโต นิรยา ฯเปฯ ปริมุตฺโต อปายทุคฺคติวินิปาตา. (8)}

\prose{ปุน จปรํ สาริปุตฺต อิเธกจฺโจ ปุคฺคโล สีเลสุ ปริปูรการี โหติ, สมาธิสฺมิํ มตฺตโส การี, ปญฺญาย มตฺตโส การี.\csromanspacer{} โส}

\chapterheadpageread{2. สีหนาทวคฺค}
\prose{\csromanfolio{187}ติณฺณํ สํโยชนานํ ปริกฺขยา สตฺตกฺขตฺตุปรโม โหติ, สตฺตกฺขตฺตุปรมํ เทเว จ มนุสฺเส จ สนฺธาวิตฺวา สํสริตฺวา ทุกฺขสฺสนฺตํ กโรติ.\csromanspacer{} อยํ สาริปุตฺต นวโม ปุคฺคโล ส-อุปาทิเสโส กาลํ กุรุมาโน ปริมุตฺโต นิรยา, ปริมุตฺโต ติรจฺฉานโยนิยา, ปริมุตฺโต เปตฺติวิสยา, ปริมุตฺโต อปายทุคฺคติวินิปาตา. (9)}

\prose{เก จ สาริปุตฺต อญฺญติตฺถิยา ปริพฺพาชกา พาลา อพฺยตฺตา.\csromanspacer{} เก จ ส-อุปาทิเสสํ วา “ส-อุปาทิเสโส”ติ ชานิสฺสนฺติ, อนุปาทิเสสํ วา “อนุปาทิเสโส”ติ ชานิสฺสนฺติ.\csromanspacer{} อิเม โข สาริปุตฺต นว ปุคฺคลา ส-อุปาทิเสสา กาลํ กุรุมานา ปริมุตฺตา นิรยา, ปริมุตฺตา ติรจฺฉานโยนิยา, ปริมุตฺตา เปตฺติวิสยา, ปริมุตฺตา อปายทุคฺคติวินิปาตา.\csromanspacer{} น ตาวายํ สาริปุตฺต ธมฺมปริยาโย ปฏิภาสิ ภิกฺขูนํ ภิกฺขุนีนํ อุปาสกานํ อุปาสิกานํ.\csromanspacer{} ตํ กิสฺส เหตุ? “มายิมํ ธมฺมปริยายํ สุตฺวา ปมาทํ อาหริํสู”ติ\footnote{อาหริํสุ (สี, อิ)}.\csromanspacer{} อปิ จ มยา\footnote{อปิ จายํ (?)} สาริปุตฺต ธมฺมปริยาโย ปญฺหาธิปฺปาเยน ภาสิโตติ.\csromanspacer{}\csromanspacer{}ทุติยํ.}

\csromanheader{2}{3. โกฏฺฐิกสุตฺต}
\proseitem{13}{อถ โข อายสฺมา มหาโกฏฺฐิโก\footnote{มหาโกฏฺฐิโต (สี, สฺยา, อิ)} เยนายสฺมา สาริปุตฺโต เตนุปสงฺกมิ, อุปสงฺกมิตฺวา อายสฺมตา สาริปุตฺเตน สทฺธิํ สมฺโมทิ, สมฺโมทนียํ กถํ สารณียํ วีติสาเรตฺวา เอกมนฺตํ นิสีทิ, เอกมนฺตํ นิสินฺโน โข อายสฺมา มหาโกฏฺฐิโก อายสฺมนฺตํ สาริปุตฺตํ เอตทโวจ “กิํ นุ โข อาวุโส สาริปุตฺต ‘ยํ กมฺมํ ทิฏฺฐธมฺมเวทนียํ, ตํ เม กมฺมํ สมฺปรายเวทนียํ โหตู’ติ เอตสฺส อตฺถาย ภควติ พฺรหฺมจริยํ วุสฺสตี”ติ.\csromanspacer{} โน หิทํ อาวุโส.}

\prose{กิํ ปนาวุโส สาริปุตฺต “ยํ กมฺมํ สมฺปรายเวทนียํ, ตํ เม กมฺมํ ทิฏฺฐธมฺมเวทนียํ โหตู”ติ เอตสฺส อตฺถาย ภควติ พฺรหฺมจริยํ วุสฺสตีติ.\csromanspacer{} โน หิทํ อาวุโส.}

\prose{กิํ นุ โข อาวุโส สาริปุตฺต “ยํ กมฺมํ สุขเวทนิยํ\footnote{สุขเวทนิยํ (ก) ม 3. 7 ปิฏฺเฐ ปสฺสิตพฺพํ.}, ตํ เม กมฺมํ ทุกฺขเวทนิยํ\footnote{ทุกฺขเวทนิยํ (ก)} โหตู”ติ เอตสฺส อตฺถาย ภควติ พฺรหฺมจริยํ วุสฺสตีติ.\csromanspacer{} โน หิทํ อาวุโส.}

\gambhira{นวกนิปาตปาฬิ}
\prose{\csromanfolio{188}กิํ ปนาวุโส สาริปุตฺต “ยํ กมฺมํ ทุกฺขเวทนียํ, ตํ เม กมฺมํ สุขเวทนียํ โหตู”ติ เอตสฺส อตฺถาย ภควติ พฺรหฺมจริยํ วุสฺสตีติ.\csromanspacer{} โน หิทํ อาวุโส.}

\prose{กิํ นุ โข อาวุโส สาริปุตฺต “ยํ กมฺมํ ปริปกฺกเวทนียํ, ตํ เม กมฺมํ อปริปกฺกเวทนียํ โหตู”ติ เอตสฺส อตฺถาย ภควติ พฺรหฺมจริยํ วุสฺสตีติ.\csromanspacer{} โน หิทํ อาวุโส.\csromanspacer{} กิํ ปนาวุโส สาริปุตฺต “ยํ กมฺมํ อปริปกฺกเวทนียํ, ตํ เม กมฺมํ ปริปกฺกเวทนียํ โหตู”ติ เอตสฺส อตฺถาย ภควติ พฺรหฺมจริยํ วุสฺสตีติ.\csromanspacer{} โน หิทํ อาวุโส.}

\prose{กิํ นุ โข อาวุโส สาริปุตฺต “ยํ กมฺมํ พหุเวทนียํ, ตํ เม กมฺมํ อปฺปเวทนียํ โหตู”ติ เอตสฺส อตฺถาย ภควติ พฺรหฺมจริยํ วุสฺสตีติ.\csromanspacer{} โน หิทํ อาวุโส.}

\prose{กิํ ปนาวุโส สาริปุตฺต “ยํ กมฺมํ อปฺปเวทนียํ, ตํ เม กมฺมํ พหุเวทนียํ โหตู”ติ เอตสฺส อตฺถาย ภควติ พฺรหฺมจริยํ วุสฺสตีติ.\csromanspacer{} โน หิทํ อาวุโส.}

\prose{กิํ นุ โข อาวุโส สาริปุตฺต “ยํ กมฺมํ เวทนียํ, ตํ เม กมฺมํ อเวทนียํ โหตู”ติ เอตสฺส อตฺถาย ภควติ พฺรหฺมจริยํ วุสฺสตีติ.\csromanspacer{} โน หิทํ อาวุโส.}

\prose{กิํ ปนาวุโส สาริปุตฺต “ยํ กมฺมํ อเวทนียํ, ตํ เม กมฺมํ เวทนียํ โหตู”ติ เอตสฺส อตฺถาย ภควติ พฺรหฺมจริยํ วุสฺสตีติ.\csromanspacer{} โน หิทํ อาวุโส.}

\prose{“กิํ นุ โข อาวุโส สาริปุตฺต ‘ยํ กมฺมํ ทิฏฺฐธมฺมเวทนียํ, ตํ เม กมฺมํ สมฺปรายเวทนียํ โหตู’ติ เอตสฺส อตฺถาย ภควติ พฺรหฺมจริยํ วุสฺสตี”ติ อิติ ปุฏฺโฐ สมาโน “โน หิทํ อาวุโส”ติ วเทสิ.}

\prose{“กิํ ปนาวุโส สาริปุตฺต ‘ยํ กมฺมํ สมฺปรายเวทนียํ, ตํ เม กมฺมํ ทิฏฺฐธมฺมเวทนียํ โหตู’ติ เอตสฺส อตฺถาย ภควติ พฺรหฺมจริยํ วุสฺสตี”ติ อิติ ปุฏฺโฐ สมาโน “โน หิทํ อาวุโส”ติ วเทสิ.}

\chapterheadpageread{2. สีหนาทวคฺค}
\prose{\csromanfolio{189}“กิํ นุ โข อาวุโส สาริปุตฺต ‘ยํ กมฺมํ สุขเวทนียํ, ตํ เม กมฺมํ ทุกฺขเวทนียํ โหตู’ติ เอตสฺส อตฺถาย ภควติ พฺรหฺมจริยํ วุสฺสตี”ติ อิติ ปุฏฺโฐ สมาโน “โน หิทํ อาวุโส”ติ วเทสิ.}

\prose{“กิํ ปนาวุโส สาริปุตฺต ‘ยํ กมฺมํ ทุกฺขเวทนียํ, ตํ เม กมฺมํ สุขเวทนียํ โหตู’ติ เอตสฺส อตฺถาย ภควติ พฺรหฺมจริยํ วุสฺสตี”ติ อิติ ปุฏฺโฐ สมาโน “โน หิทํ อาวุโส”ติ วเทสิ.}

\prose{“กิํ นุ โข อาวุโส สาริปุตฺต ‘ยํ กมฺมํ ปริปกฺกเวทนียํ, ตํ เม กมฺมํ อปริปกฺกเวทนียํ โหตู’ติ เอตสฺส อตฺถาย ภควติ พฺรหฺมจริยํ วุสฺสตี”ติ อิติ ปุฏฺโฐ สมาโน “โน หิทํ อาวุโส”ติ วเทสิ.}

\prose{“กิํ ปนาวุโส สาริปุตฺต ‘ยํ กมฺมํ อปริปกฺกเวทนียํ, ตํ เม กมฺมํ ปริปกฺกเวทนียํ โหตู’ติ เอตสฺส อตฺถาย ภควติ พฺรหฺมจริยํ วุสฺสตี”ติ อิติ ปุฏฺโฐ สมาโน “โน หิทํ อาวุโส”ติ วเทสิ.}

\prose{“กิํ นุ โข อาวุโส สาริปุตฺต ‘ยํ กมฺมํ พหุเวทนียํ, ตํ เม กมฺมํ อปฺปเวทนียํ โหตู’ติ เอตสฺส อตฺถาย ภควติ พฺรหฺมจริยํ วุสฺสตี”ติ อิติ ปุฏฺโฐ สมาโน “โน หิทํ อาวุโส”ติ วเทสิ.}

\prose{“กิํ ปนาวุโส สาริปุตฺต ‘ยํ กมฺมํ อปฺปเวทนียํ, ตํ เม กมฺมํ พหุเวทนียํ โหตู’ติ เอตสฺส อตฺถาย ภควติ พฺรหฺมจริยํ วุสฺสตี”ติ อิติ ปุฏฺโฐ สมาโน “โน หิทํ อาวุโส”ติ วเทสิ.}

\prose{“กิํ นุ โข อาวุโส สาริปุตฺต ‘ยํ กมฺมํ เวทนียํ, ตํ เม กมฺมํ อเวทนียํ โหตู’ติ เอตสฺส อตฺถาย ภควติ พฺรหฺมจริยํ วุสฺสตี”ติ อิติ ปุฏฺโฐ สมาโน “โน หิทํ อาวุโส”ติ วเทสิ.}

\prose{“กิํ ปนาวุโส สาริปุตฺต ‘ยํ กมฺมํ อเวทนียํ, ตํ เม กมฺมํ เวทนียํ โหตู’ติ เอตสฺส อตฺถาย ภควติ พฺรหฺมจริยํ วุสฺสตี”ติ อิติ ปุฏฺโฐ สมาโน “โน หิทํ อาวุโส”ติ วเทสิ, อถ กิมตฺถํ จรหาวุโส ภควติ พฺรหฺมจริยํ วุสฺสตีติ.}

\prose{ยํ ขฺวสฺส\footnote{ยํ โข (ก)} อาวุโส อญฺญาตํ อทิฏฺฐํ อปฺปตฺตํ อสจฺฉิกตํ อนภิสเมตํ, ตสฺส ญาณาย ทสฺสนาย ปตฺติยา สจฺฉิกิริยาย อภิสมยาย ภควติ พฺรหฺมจริยํ วุสฺสตีติ\footnote{วุสฺสติ (สฺยา)}. (กิํ ปนสฺสาวุโส}

\vspace{\tipitakaparskip}
\csromancenter{นวกนิปาตปาฬิ อญฺญาตํ อทิฏฺฐํ อปฺปตฺตํ อสจฺฉิกตํ อนภิสเมตํ, ยสฺส ญาณาย ทสฺสนาย ปตฺติยา สจฺฉิกิริยาย อภิสมยาย ภควติ พฺรหฺมจริยํ วุสฺสตีติ.)\footnote{( ) สฺยา-ก-โปตฺถเกสุ นตฺถิ.} “อิทํ ทุกฺขนฺ”ติ ขฺวสฺส\footnote{โข ยํ (ก)} อาวุโส อญฺญาตํ อทิฏฺฐํ อปฺปตฺตํ อสจฺฉิกตํ อนภิสเมตํ, ตสฺส ญาณาย ทสฺสนาย ปตฺติยา สจฺฉิกิริยาย อภิสมยาย ภควติ พฺรหฺมจริยํ วุสฺสติ. “อยํ ทุกฺขสมุทโย”ติ ขฺวสฺส อาวุโส ฯเปฯ “อยํ ทุกฺขนิโรโธ”ติ ขฺวสฺส อาวุโส ฯเปฯ “อยํ ทุกฺขนิโรธคามินี ปฏิปทา”ติ ขฺวสฺส อาวุโส อญฺญาตํ อทิฏฺฐํ อปฺปตฺตํ อสจฺฉิกตํ อนภิสเมตํ, ตสฺส ญาณาย ทสฺสนาย ปตฺติยา สจฺฉิกิริยาย อภิสมยาย ภควติ พฺรหฺมจริยํ วุสฺสติ.\csromanspacer{} อิทํ ขฺวสฺส\footnote{อิติ โข ยํ (ก)} อาวุโส อญฺญาตํ อทิฏฺฐํ อปฺปตฺตํ อสจฺฉิกตํ อนภิสเมตํ, ตสฺส\footnote{ยสฺส (?)} ญาณาย ทสฺสนาย ปตฺติยา สจฺฉิกิริยาย อภิสมยาย ภควติ พฺรหฺมจริยํ วุสฺสตีติ.\csromanspacer{}\csromanspacer{}ตติยํ.}
\csromanheader{2}{4. สมิทฺธิสุตฺต}
\proseitem{14}{\csromanfolio{190}\csromansymbolfootnote{*}{เหฏฺฐา 152 ปิฏฺเฐปิ.}อถ โข อายสฺมา สมิทฺธิ เยนายสฺมา สาริปุตฺโต เตนุปสงฺกมิ, อุปสงฺกมิตฺวา อายสฺมนฺตํ สาริปุตฺตํ อภิวาเทตฺวา เอกมนฺตํ นิสีทิ, เอกมนฺตํ นิสินฺนํ โข อายสฺมนฺตํ สมิทฺธิํ อายสฺมา สาริปุตฺโต เอตทโวจ “กิมารมฺมณา สมิทฺธิ ปุริสสฺส สงฺกปฺปวิตกฺกา อุปฺปชฺชนฺตี”ติ.\csromanspacer{} นามรูปารมฺมณา ภนฺเตติ.\csromanspacer{} เต ปน สมิทฺธิ กฺว นานตฺตํ คจฺฉนฺตีติ.\csromanspacer{} ธาตูสุ ภนฺเตติ.\csromanspacer{} เต ปน สมทฺธิ กิํสมุทยาติ.\csromanspacer{} ผสฺสสมุทยา ภนฺเตติ.\csromanspacer{} เต ปน สมิทฺธิ กิํสโมสรณาติ.\csromanspacer{} เวทนาสโมสรณา ภนฺเตติ.\csromanspacer{} เต ปน สมิทฺธิ กิํปมุขาติ.\csromanspacer{} สมาธิปฺปมุขา ภนฺเตติ.\csromanspacer{} เต ปน สมิทฺธิ กิํอธิปเตยฺยาติ.\csromanspacer{} สตาธิปเตยฺยา ภนฺเตติ.\csromanspacer{} เต ปน สมิทฺธิ กิํอุตฺตราติ.\csromanspacer{} ปญฺญุตฺตรา ภนฺเตติ.\csromanspacer{} เต ปน สมิทฺธิ กิํสาราติ.\csromanspacer{} วิมุตฺติสารา ภนฺเตติ.\csromanspacer{} เต ปน สมิทฺธิ กิํโอคธาติ.\csromanspacer{} อมโตคธา ภนฺเตติ.}

\prose{“กิมารมฺมณา สมิทฺธิ ปุริสสฺส สงฺกปฺปวิตกฺกา อุปฺปชฺชนฺตี”ติ อิติ ปุฏฺโฐ สมาโน “นามรูปารมฺมณา ภนฺเต”ติ วเทสิ. “เต ปน สมิทฺธิ กฺว นานตฺตํ คจฺฉนฺตี”ติ อิติ ปุฏฺโฐ สมาโน “ธาตูสุ ภนฺเต”ติ วเทสิ. “เต ปน สมิทฺธิ กิํสมุทยา”ติ อิติ ปุฏฺโฐ สมาโน “ผสฺสสมุทยา}

\chapterheadpageread{2. สีหนาทวคฺค}
\prose{\csromanfolio{191}ภนฺเต”ติ วเทสิ. “เต ปน สมิทฺธิ กิํสโมสรณา”ติ อิติ ปุฏฺโฐ สมาโน “เวทนาสโมสรณา ภนฺเต”ติ วเทสิ. “เต ปน สมิทฺธิ กิํปมุขา”ติ อิติ ปุฏฺโฐ สมาโน “สมาธิปฺปมุขา ภนฺเต”ติ วเทสิ. “เต ปน สมิทฺธิ กิํอธิปเตยฺยา”ติ อิติ ปุฏฺโฐ สมาโน “สตาธิปเตยฺยา ภนฺเต”ติ วเทสิ. “เต ปน สมิทฺธิ กิํอุตฺตรา”ติ อิติ ปุฏฺโฐ สมาโน “ปญฺญุตฺตรา ภนฺเต”ติ วเทสิ. “เต ปน สมิทฺธิ กิํสารา”ติ อิติ ปุฏฺโฐ สมาโน “วิมุตฺติสารา ภนฺเต”ติ วเทสิ. “เต ปน สมิทฺธิ กิํโอคธา”ติ อิติ ปุฏฺโฐ สมาโน “อมโตคธา ภนฺเต”ติ วเทสิ.\csromanspacer{} สาธุ สาธุ สมิทฺธิ, สาธุ โข ตฺวํ สมิทฺธิ ปุฏฺโฐ\footnote{ปญฺหํ (สี, สฺยา, อิ)} ปุฏฺโฐ วิสฺสชฺเชสิ, เตน จ มา มญฺญีติ.\csromanspacer{}\csromanspacer{}จตุตฺถํ.}

\csromanheader{2}{5. คณฺฑสุตฺต}
\proseitem{15}{เสยฺยถาปิ ภิกฺขเว คณฺโฑ อเนกวสฺสคณิโก, ตสฺสสฺสุ คณฺฑสฺส นว วณมุขานิ นว อเภทนมุขานิ, ตโต ยํ กิญฺจิ ปคฺฆเรยฺย, อสุจิเยว ปคฺฆเรยฺย, ทุคฺคนฺธํเยว ปคฺฆเรยฺย, เชคุจฺฉิยํเยว\footnote{ธชคุจฺฉิเยว (ก)} ปคฺฆเรยฺย.\csromanspacer{} ยํ กิญฺจิ ปสเวยฺย, อสุจิเยว ปสเวยฺย, ทุคฺคนฺธํเยว ปสเวยฺย, เชคุจฺฉิยํเยว ปสเวยฺย.}

\prose{“คณฺโฑ”ติ โข ภิกฺขเว อิมสฺเสตํ จาตุมหาภูติกสฺส\footnote{จาตุมฺมหาภูติกสฺส (สี, สฺยา, อิ)} กายสฺส อธิวจนํ มาตาเปตฺติกสมฺภวสฺส โอทนกุมฺมาสูปจยสฺส อนิจฺจุจฺฉาทนปริมทฺทนเภทนวิทฺธํสนธมฺมสฺส.\csromanspacer{} ตสฺสสฺสุ คณฺฑสฺส นว วณมุขานิ นว อเภทนมุขานิ.\csromanspacer{} ตโต ยํ กิญฺจิ ปคฺฆรติ, อสุจิเยว ปคฺฆรติ, ทุคฺคนฺธํเยว ปคฺฆรติ, เชคุจฺฉิยํเยว ปคฺฆรติ.\csromanspacer{} ยํ กิญฺจิ ปสวติ, อสุจิเยว ปสวติ, ทุคฺคนฺธํเยว ปสวติ, เชคุจฺฉิยํเยว ปสวติ.\csromanspacer{} ตสฺมาติห ภิกฺขเว อิมสฺมิํ กาเย นิพฺพินฺทถาติ.\csromanspacer{}\csromanspacer{}ปญฺจมํ.}

\csromanheader{2}{6. สญฺญาสุตฺต}
\proseitem{16}{นวยิมา ภิกฺขเว สญฺญา ภาวิตา พหุลีกตา มหปฺผลา โหนฺติ มหานิสํสา อมโตคธา อมตปริโยสานา.\csromanspacer{} กตมา}

\vspace{\tipitakaparskip}
\csromancenter{นวกนิปาตปาฬิ นว? อสุภสญฺญา มรณสญฺญา อาหาเร ปฏิกูลสญฺญา\footnote{ปฏิกฺกูลสญฺญา (สี, สฺยา, อิ)} สพฺพโลเก อนภิรตสญฺญา\footnote{อนภิรติสญฺญา (ก) อํ 2. 126 ปิฏฺเฐปิ.} อนิจฺจสญฺญา อนิจฺเจ ทุกฺขสญฺญา ทุกฺเข อนตฺตสญฺญา ปหานสญฺญา วิราคสญฺญา.\csromanspacer{} อิมา โข ภิกฺขเว นว สญฺญา ภาวิตา พหุลีกตา มหปฺผลา โหนฺติ มหานิสํสา อมโตคธา อมตปริโยสานาติ.\csromanspacer{}\csromanspacer{}ฉฏฺฐํ.}
\csromanheader{2}{7. กุลสุตฺต}
\proseitem{17}{\csromanfolio{192}นวหิ ภิกฺขเว องฺเคหิ สมนฺนาคตํ กุลํ อนุปคนฺตฺวา วา นาลํ อุปคนฺตุํ, อุปคนฺตฺวา วา นาลํ นิสีทิตุํ.\csromanspacer{} กตเมหิ นวหิ? น มนาเปน ปจฺจุฏฺเฐนฺติ, น มนาเปน อภิวาเทนฺติ, น มนาเปน อาสนํ เทนฺติ, สนฺตมสฺส ปริคุหนฺติ, พหุกมฺปิ โถกํ เทนฺติ, ปณีตมฺปิ ลูขํ เทนฺติ, อสกฺกจฺจํ เทนฺติ โน สกฺกจฺจํ, น อุปนิสีทนฺติ ธมฺมสฺสวนาย, ภาสิตมสฺส น สุสฺสูสนฺติ.\csromanspacer{} อิเมหิ โข ภิกฺขเว นวหงฺเคหิ สมนฺนาคตํ กุลํ อนุปคนฺตฺวา วา นาลํ อุปคนฺตุํ, อุปคนฺตฺวา วา นาลํ นิสีทิตุํ.}

\prose{นวหิ ภิกฺขเว องฺเคหิ สมนฺนาคตํ กุลํ อนุปคนฺตฺวา วา อลํ อุปคนฺตุํ, อุปคนฺตฺวา วา อลํ นิสีทิตุํ.\csromanspacer{} กตเมหิ นวหิ? มนาเปน ปจฺจุฏฺเฐนฺติ, มนาเปน อภิวาเทนฺติ, มนาเปน อาสนํ เทนฺติ, สนฺตมสฺส น ปริคุหนฺติ, พหุกมฺปิ พหุกํ เทนฺติ, ปณีตมฺปิ ปณีตํ เทนฺติ, สกฺกจฺจํ เทนฺติ โน อสกฺกจฺจํ, อุปนิสีทนฺติ ธมฺมสฺสวนาย, ภาสิตมสฺส สุสฺสูสนฺติ.\csromanspacer{} อิเมหิ โข ภิกฺขเว นวหงฺเคหิ สมนฺนาคตํ กุลํ อนุปคนฺตฺวา วา อลํ อุปคนฺตุํ, อุปคนฺตฺวา วา อลํ นิสีทิตุนฺติ.\csromanspacer{}\csromanspacer{}สตฺตมํ.}

\csromanheader{2}{8. นวงฺคุโปสถสุตฺต}
\proseitem{18}{นวหิ ภิกฺขเว องฺเคหิ สมนฺนาคโต อุโปสโถ อุปวุตฺโถ มหปฺผโล โหติ มหานิสํโส มหาชุติโก มหาวิปฺผาโร.\csromanspacer{} กถํ อุปวุตฺโถ จ ภิกฺขเว นวหงฺเคหิ สมนฺนาคโต อุโปสโถ มหปฺผโล โหติ มหานิสํโส มหาชุติโก มหาวิปฺผาโร, อิธ ภิกฺขเว อริยสาวโก อิติ ปฏิสญฺจิกฺขภิ “ยาวชีวํ อรหนฺโต ปาณาติปาตํ ปหาย ปาณาติปาตา ปฏิวิรตา นิหิตทณฺฑา นิหิตสตฺถา}

\chapterheadpageread{2. สีหนาทวคฺค}
\prose{\csromanfolio{193}ลชฺชี ทยาปนฺนา สพฺพปาณภูตหิตานุกมฺปิโน วิหรนฺติ, อหมฺปชฺช อิมญฺจ รตฺติํ อิมญฺจ ทิวสํ ปาณาติปาตํ ปหาย ปาณาติปาตา ปฏิวิรโต นิหิตทณฺโฑ นิหิตสตฺโถ ลชฺชี ทยาปนฺโน สพฺพปาณภูตหิตานุกมฺปี วิหรามิ.\csromanspacer{} อิมินาปงฺเคน\footnote{อิมินาปิ องฺเคน (ก-สี)} อรหตํ อนุกโรมิ, อุโปสโถ จ เม อุปวุตฺโถ ภวิสฺสตี”ติ.\csromanspacer{} อิมินา ปฐเมน องฺเคน สมนฺนาคโต โหติ ฯเปฯ.}

\prose{“ยาวชีวํ อรหนฺโต อุจฺจาสยนมหาสยนํ ปหาย อุจฺจาสยนมหาสยนา ปฏิวิรตา นีจเสยฺยํ กปฺเปนฺติ มญฺจเก วา ติณสนฺถารเก วา, อหมฺปชฺช อิมญฺจ รตฺติํ อิมญฺจ ทิวสํ อุจฺจาสยนมหาสยนํ ปหาย อุจฺจาสยนมหาสยนา ปฏิวิรโต นีจเสยฺยํ กปฺเปมิ มญฺจเก วา ติณสนฺถารเก วา.\csromanspacer{} อิมินาปงฺเคน อรหตํ อนุกโรมิ, อุโปสโถ จ เม อุปวุตฺโถ ภวิสฺสตี”ติ.\csromanspacer{} อิมินา อฏฺฐเมน องฺเคน สมนฺนาคโต โหติ.}

\prose{เมตฺตาสหคเตน เจตสา เอกํ ทิสํ ผริตฺวา วิหรติ.\csromanspacer{} ตถา ทุติยํ.\csromanspacer{} ตถา ตติยํ.\csromanspacer{} ตถา จตุตฺถํ.\csromanspacer{} อิติ อุทฺธมโธ ติริยํ สพฺพธิ สพฺพตฺตตาย สพฺพาวนฺตํ โลกํ เมตฺตาสหคเตน เจตสา วิปุเลน มหคฺคเตน อปฺปมาเณน อเวเรน อพฺยาปชฺเชน\footnote{อพฺยาปชฺเฌน (ก), อพฺยาพชฺเฌน (?)} ผริตฺวา วิหรติ.\csromanspacer{} อิมินา นวเมน องฺเคน สมนฺนาคโต โหติ.\csromanspacer{} เอวํ อุปวุตฺโถ โข ภิกฺขเว นวหงฺเคหิ สมนฺนาคโต อุโปสโถ มหปฺผโล โหติ มหานิสํโส มหาชุติโก มหาวิปฺผาโรติ.\csromanspacer{}\csromanspacer{}อฏฺฐมํ.}

\csromanheader{2}{9. เทวตาสุตฺต}
\proseitem{19}{อิมญฺจ ภิกฺขเว รตฺติํ สมฺพหุลา เทวตา อภิกฺกนฺตาย รตฺติยา อภิกฺกนฺตวณฺณา เกวลกปฺปํ เชตวนํ โอภาเสตฺวา เยนาหํ เตนุปสงฺกมิํสุ, อุปสงฺกมิตฺวา มํ อภิวาเทตฺวา เอกมนฺตํ อฏฺฐํสุ, เอกมนฺตํ ฐิตา โข ภิกฺขเว ตา เทวตา มํ เอตทโวจุํ “อุปสงฺกมิํสุ โน ภนฺเต ปุพฺเพ มนุสฺสภูตานํ ปพฺพชิตา อคารานิ, เต มยํ ภนฺเต ปจฺจุฏฺฐิมฺห, โน จ โข อภิวาทิมฺห, ตา มยํ ภนฺเต อปริปุณฺณกมฺมนฺตา วิปฺปฏิสารินิโย ปจฺจานุตาปินิโย หีนํ กายํ อุปปนฺนา”ติ.}

\gambhira{นวกนิปาตปาฬิ}
\prose{\csromanfolio{194}อปราปิ มํ ภิกฺขเว สมฺพหุลา เทวตา อุปสงฺกมิตฺวา เอตทโวจุํ “อุปสงฺกมิํสุ โน ภนฺเต ปุพฺเพ มนุสฺสภูตานํ ปพฺพชิตา อคารานิ, เต มยํ ภนฺเต ปจฺจุฏฺฐิมฺห อภิวาทิมฺห\footnote{ปจฺจุฏฺฐิมฺห จ อภิวาทิมฺห จ (สฺยา)}, โน จ เตสํ อาสนํ อทมฺห, ตา มยํ ภนฺเต อปริปุณฺณกมฺมนฺตา วิปฺปฏิสารินิโย ปจฺจานุตาปินิโย หีนํ กายํ อุปปนฺนา”ติ.}

\prose{อปราปิ มํ ภิกฺขเว สมฺพหุลา เทวตา อุปสงฺกมิตฺวา เอตทโวจุํ “อุปสงฺกมิํสุ โน ภนฺเต ปุพฺเพ มนุสฺสภูตานํ ปพฺพชิตา อคารานิ, เต มยํ ภนฺเต ปจฺจุฏฺฐิมฺห อภิวาทิมฺห1, อาสนํ\footnote{อาสนญฺจ (สี, สฺยา)} อทมฺห, โน จ โข ยถาสตฺติ ยถาพลํ สํวิภชิมฺห ฯเปฯ ยถาสตฺติ ยถาพลํ\footnote{ยถาพลํ จ (?)} สํวิภชิมฺห, โน จ โข อุปนิสีทิมฺห ธมฺมสฺสวนาย ฯเปฯ อุปนิสีทิมฺห\footnote{อุปนิสีทิมฺห จ (สฺยา)} ธมฺมสฺสวนาย, โน จ โข โอหิตโสตา ธมฺมํ สุณิมฺห ฯเปฯ โอหิตโสตา จ ธมฺมํ สุณิมฺห, โน จ โข สุตฺวา ธมฺมํ ธารยิมฺห ฯเปฯ สุตฺวา จ ธมฺมํ ธารยิมฺห, โน จ โข ธาตานํ ธมฺมานํ อตฺถํ อุปปริกฺขิมฺห ฯเปฯ ธาตานญฺจ ธมฺมานํ อตฺถํ อุปปริกฺขิมฺห, โน จ โข อตฺถมญฺญาย ธมฺมมญฺญาย ธมฺมานุธมฺมํ ปฏิปชฺชิมฺห.\csromanspacer{} ตา มยํ ภนฺเต อปริปุณฺณกมฺมนฺตา วิปฺปฏิสารินิโย ปจฺจานุตาปินิโย หีนํ กายํ อุปปนฺนา”ติ.}

\prose{อปราปิ มํ ภิกฺขเว สมฺพหุลา เทวตา อุปสงฺกมิตฺวา เอตทโวจุํ “อุปสงฺกมิํสุ โน ภนฺเต ปุพฺเพ มนุสฺสภูตานํ ปพฺพชิตา อคารานิ, เต มยํ ภนฺเต ปจฺจุฏฺฐิมฺห อภิวาทิมฺห1, อาสนํ2 อทมฺห, ยถาสตฺติ ยถาพลํ3 สํวิภชิมฺห, อุปนิสีทิมฺห4 ธมฺมสฺสวนาย, โอหิตโสตา จ ธมฺมํ สุณิมฺห, สุตฺวา จ ธมฺมํ ธารยิมฺห, ธาตานญฺจ ธมฺมานํ อตฺถํ อุปปริกฺขิมฺห, อตฺถมญฺญาย ธมฺมมญฺญาย ธมฺมานุธมฺมํ\footnote{ธมฺมานุธมฺมญฺจ (?)} ปฏิปชฺชิมฺห.\csromanspacer{} ตา มยํ ภนฺเต ปริปุณฺณกมฺมนฺตา อวิปฺปฏิสารินิโย อปจฺจานุตาปินิโย ปณีตํ กายํ อุปปนฺนา”ติ.\csromanspacer{} เอตานิ ภิกฺขเว รุกฺขมูลานิ, เอตานิ สุญฺญาคารานิ, ฌายถ ภิกฺขเว มา ปมาทตฺถ, มา ปจฺฉา วิปฺปฏิสาริโน อหุวตฺถ เสยฺยถาปิ ตา ปุริมิกา เทวตาติ.\csromanspacer{}\csromanspacer{}นวมํ.}

\chapterheadpageread{2. สีหนาทวคฺค}
\csromanheader{2}{10. เวลามสุตฺต}
\proseitem{20}{\csromanfolio{195}เอกํ สมยํ ภควา สาวตฺถิยํ วิหรติ เชตวเน อนาถปิณฺฑิกสฺส อาราเม.\csromanspacer{} อถ โข อนาถปิณฺฑิโก คหปติ เยน ภควา เตนุปสงฺกมิ, อุปสงฺกมิตฺวา ภควนฺตํ อภิวาเทตฺวา เอกมนฺตํ นิสีทิ, เอกมนฺตํ นิสินฺนํ โข อนาถปิณฺฑิกํ คหปติํ ภควา เอตทโวจ\csromandash{}}

\prose{อปิ นุ เต คหปติ กุเล ทานํ ทียตีติ.\csromanspacer{} ทียติ เม ภนฺเต กุเล ทานํ, ตญฺจ โข ลูขํ กณาชกํ พิฬงฺคทุติยนฺติ.\csromanspacer{} ลูขญฺเจปิ\footnote{ลูขํ วาปิ (สฺยา), ลูขญฺจาปิ (ก)} คหปติ ทานํ เทติ ปณีตํ วา, ตญฺจ อสกฺกจฺจํ เทติ, อจิตฺตีกตฺวา\footnote{อจิตฺติํ กตฺวา (ก), อปจิตฺติํ กตฺวา (สฺยา), อจิตฺติกตฺวา (อิ)} เทติ, อสหตฺถา เทติ, อปวิทฺธํ\footnote{อปวิฏฺฐํ (สฺยา)} เทติ, อนาคมนทิฏฺฐิโก เทติ.\csromanspacer{} ยตฺถ ยตฺถ ตสฺส ตสฺส ทานสฺส วิปาโก นิพฺพตฺตติ, น อุฬาราย ภตฺตโภคาย จิตฺตํ นมติ, น อุฬาราย วตฺถโภคาย จิตฺตํ นมติ, น อุฬาราย ยานโภคาย จิตฺตํ นมติ, น อุฬาเรสุ ปญฺจสุ กามคุเณสุ โภคาย จิตฺตํ นมติ.\csromanspacer{} เยปิสฺส เต โหนฺติ “ปุตฺตา”ติ วา “ทารา”ติ วา “ทาสา”ติ วา “เปสฺสา”ติ วา “กมฺมกรา”ติ วา, เตปิ น สุสฺสูสนฺติ, น โสตํ โอทหนฺติ น อญฺญา จิตฺตํ อุปฏฺฐเปนฺติ.\csromanspacer{} ตํ กิสฺส เหตุ? เอวํ เหตํ\footnote{เอวญฺเจตํ (สฺยา, ก)} คหปติ โหติ อสกฺกจฺจํ กตานํ กมฺมานํ วิปาโก.}

\prose{ลูขญฺเจปิ คหปติ ทานํ เทติ ปณีตํ วา, ตญฺจ สกฺกจฺจํ เทติ, จิตฺตีกตฺวา เทติ, สหตฺถา เทติ, อนปวิทฺธํ เทติ, อาคมนทิฏฺฐิโก เทติ.\csromanspacer{} ยตฺถ ยตฺถ ตสฺส ตสฺส ทานสฺส วิปาโก นิพฺพตฺตติ, อุฬาราย ภตฺตโภคาย จิตฺตํ นมติ, อุฬาราย วตฺถโภคาย จิตฺตํ นมติ, อุฬาราย ยานโภคาย จิตฺตํ นมติ, อุฬาเรสุ ปญฺจสุ กามคุเณสุ โภคาย จิตฺตํ นมติ.\csromanspacer{} เยปิสฺส เต โหนฺติ “ปุตฺตา”ติ วา “ทารา”ติ วา “ทาสา”ติ วา “เปสฺสา”ติ วา “กมฺมกรา”ติ วา, เตปิ สุสฺสูสนฺติ, โสตํ โอทหนฺติ, อญฺญา จิตฺตํ อุปฏฺฐเปนฺติ.\csromanspacer{} ตํ กิสฺส เหตุ? เอวํ เหตํ คหปติ โหติ สกฺกจฺจํ กตานํ กมฺมานํ วิปาโก.}

\prose{ภูตปุพฺพํ คหปติ เวลาโม นาม พฺราหฺมโณ อโหสิ, โส เอวรูปํ ทานํ อทาสิ มหาทานํ.\csromanspacer{} จตุราสีติ สุวณฺณปาติสหสฺสานิ อทาสิ}

\vspace{\tipitakaparskip}
\csromancenter{นวกนิปาตปาฬิ รูปิยปูรานิ, จตุราสีติ รูปิยปาติสหสฺสานิ อทาสิ สุวณฺณปูรานิ, จตุราสีติ กํสปาติสหสฺสานิ อทาสิ หิรญฺญปูรานิ, จตุราสีติ หตฺถิสหสฺสานิ อทาสิ โสวณฺณาลงฺการานิ โสวณฺณธชานิ เหมชาลปฺปฏิจฺฉนฺนานิ\footnote{เหมชาลสญฺฉนฺนานิ (สี, อิ)}, จตุราสีติ รถสหสฺสานิ อทาสิ สีหจมฺมปริวารานิ พฺยคฺฆจมฺมปริวารานิ ทีปิจมฺมปริวารานิ ปณฺฑุกมฺพลปริวารานิ โสวณฺณาลงฺการานิ โสวณฺณธชานิ เหมชาลปฺปฏิจฺฉนฺนานิ, จตุราสีติ เธนุสหสฺสานิ อทาสิ ทุกูลสนฺธนานิ\footnote{ทุกูลสนฺทสฺสนานิ (สี), ทุกูลสณฺฐนานิ (สฺยา), ทุกูลสนฺถานานิ (อิ), ทุหสนฺทนานิ (ที 2. 153), ทุกูลสนฺทนานิ (ตตฺถ ปาฐนฺตรํ)} กํสูปธารณานิ, จตุราสีติ กญฺญาสหสฺสานิ อทาสิ อามุตฺตมณิกุณฺฑลาโย\footnote{อามุกฺกมณิกุณฺฑลาโย (?)}, จตุราสีติ ปลฺลงฺกสหสฺสานิ อทาสิ โคนกตฺถตานิ ปฏิกตฺถตานิ ปฏลิกตฺถตานิ กทลิมิคปวรปจฺจตฺถรณานิ ส-อุตฺตรจฺฉทานิ อุภโตโลหิตกูปธานานิ.\csromanspacer{} จตุราสีติ วตฺถโกฏิสหสฺสานิ อทาสิ โขมสุขุมานํ โกเสยฺยสุขุมานํ กมฺพลสุขุมานํ กปฺปาสิกสุขุมานํ, โก ปน วาโท อนฺนสฺส ปานสฺส ขชฺชสฺส โภชฺชสฺส เลยฺยสฺส เปยฺยสฺส, นชฺโช มญฺเญ วิสฺสนฺทนฺติ\footnote{วิสฺสนฺทติ (สี, อิ)}.}
\prose{\csromanfolio{196}สิยา โข ปน เต คหปติ เอวมสฺส “อญฺโญ นูน เตน สมเยน เวลาโม พฺราหฺมโณ อโหสิ, โส\footnote{โย (?)} ตํ ทานํ อทาสิ มหาทานนฺ”ติ.\csromanspacer{} น โข ปเนตํ คหปติ เอวํ ทฏฺฐพฺพํ.\csromanspacer{} อหํ เตน สมเยน เวลาโม พฺราหฺมโณ อโหสิํ, อหํ ตํ ทานํ อทาสิํ มหาทานํ.\csromanspacer{} ตสฺมิํ โข ปน คหปติ ทาเน น โกจิ ทกฺขิเณยฺโย อโหสิ, น ตํ โกจิ ทกฺขิณํ วิโสเธติ.}

\prose{ยํ คหปติ เวลาโม พฺราหฺมโณ ทานํ อทาสิ มหาทานํ, โย เจกํ ทิฏฺฐิสมฺปนฺนํ โภเชยฺย, อิทํ ตโต มหปฺผลตรํ.}

\prose{( )\footnote{(ยญฺจ คหปติ เวลาโม พฺราหฺมโณ ทานํ อทาสิ มหาทานํ) (สี, อิ)} โย จ สตํ ทิฏฺฐิสมฺปนฺนานํ โภเชยฺย, โย เจกํ สกทาคามิํ โภเชยฺย, อิทํ ตโต มหปฺผลตรํ.}

\prose{( )6 โย จ สตํ สกทาคามีนํ โภเชยฺย, โย เจกํ อนาคามิํ โภเชยฺย ฯเปฯ.\csromanspacer{} โย จ สตํ อนาคามีนํ โภเชยฺย, โย}

\chapterheadpageread{2. สีหนาทวคฺค}
\prose{\csromanfolio{197}เจกํ อรหนฺตํ โภเชยฺย ฯเปฯ.\csromanspacer{} โย จ สตํ อรหนฺตานํ โภเชยฺย, โย เจกํ ปจฺเจกพุทฺธํ โภเชยฺย ฯเปฯ.\csromanspacer{} โย จ สตํ ปจฺเจกพุทฺธานํ โภเชยฺย, โย จ ตถาคตํ อรหนฺตํ สมฺมาสมฺพุทฺธํ โภเชยฺย ฯเปฯ.\csromanspacer{} โย จ พุทฺธปฺปมุขํ ภิกฺขุสํฆํ โภเชยฺย ฯเปฯ.\csromanspacer{} โย จ จาตุทฺทิสํ สํฆํ อุทฺทิสฺส วิหารํ การาเปยฺย ฯเปฯ.\csromanspacer{} โย จ ปสนฺนจิตฺโต พุทฺธญฺจ ธมฺมญฺจ สํฆญฺจ สรณํ คจฺเฉยฺย ฯเปฯ.\csromanspacer{} โย จ ปสนฺนจิตฺโต สิกฺขาปทานิ สมาทิเยยฺย ปาณาติปาตา เวรมณิํ, อทินฺนาทานา เวรมณิํ, กาเมสุมิจฺฉาจารา เวรมณิํ, มุสาวาทา เวรมณิํ, สุราเมรยมชฺชปมาทฏฺฐานา เวรมณิํ.\csromanspacer{} โย จ อนฺตมโส คนฺโธหนมตฺตมฺปิ\footnote{คนฺธูหนมตฺตมฺปิ (สี), คทฺทูหนมตฺตมฺปิ (สฺยา, อิ) ม 3. 167 ปิฏฺเฐ จ.} เมตฺตจิตฺตํ ภาเวยฺย. ( )\footnote{(โย จ อจฺฉราสงฺฆาตมตฺตมฺปิ อนิจฺจสญฺญํ ภาเวยฺย) (ก)} อิทํ ตโต มหปฺผลตรํ.}

\prose{ยญฺจ คหปติ เวลาโม พฺราหฺมโณ ทานํ อทาสิ มหาทานํ, โย เจกํ ทิฏฺฐิสมฺปนฺนํ โภเชยฺย, โย จ สตํ ทิฏฺฐิสมฺปนฺนานํ โภเชยฺย, โย เจกํ สกทาคามิํ โภเชยฺย, โย จ สตํ สกทาคามีนํ โภเชยฺย, โย เจกํ อนาคามิํ โภเชยฺย, โย จ สตํ อนาคามีนํ โภเชยฺย, โย เจกํ อรหนฺตํ โภเชยฺย, โย จ สตํ อรหนฺตานํ โภเชยฺย, โย เจกํ ปจฺเจกพุทฺธํ โภเชยฺย, โย จ สตํ ปจฺเจกพุทฺธานํ โภเชยฺย, โย จ ตถาคตํ อรหนฺตํ สมฺมาสมฺพุทฺธํ โภเชยฺย, โย จ พุทฺธปฺปมุขํ ภิกฺขุสํฆํ โภเชยฺย, โย จ จาตุทฺทิสํ สํฆํ อุทฺทิสฺส วิหารํ การาเปยฺย.\csromanspacer{} โย จ ปสนฺนจิตฺโต พุทฺธญฺจ ธมฺมญฺจ สํฆญฺจ สรณํ คจฺเฉยฺย, โย จ ปสนฺนจิตฺโต สิกฺขาปทานิ สมาทิเยยฺย ปาณาติปาตา เวรมณิํ ฯเปฯ สุราเมรยมชฺชปมาทฏฺฐานา เวรมณิํ, โย จ อนฺตมโส คนฺโธหนมตฺตมฺปิ เมตฺตจิตฺตํ ภาเวยฺย, โย จ อจฺฉราสงฺฆาตมตฺตมฺปิ อนิจฺจสญฺญํ ภาเวยฺย, อิทํ ตโต มหปฺผลตรนฺติ.\csromanspacer{}\csromanspacer{}ทสมํ.\csromanspacer{} สีหนาทวคฺโค ทุติโย.}

\titlehead{\textbf{ตสฺสุทฺทานํ}}
\csromangathameasure{นาโท ส-อุปาทิเสโส จ, โกฏฺฐิเกน สมิทฺธินา. \\ คณฺฑสญฺญา กุลํ เมตฺตา, เทวตา เวลาเมน จาติ.}
\csromangathagroup{\csromangathastanza{นาโท ส-อุปาทิเสโส จ, โกฏฺฐิเกน สมิทฺธินา. \\ คณฺฑสญฺญา กุลํ เมตฺตา, เทวตา เวลาเมน จาติ.}}

\gambhira{นวกนิปาตปาฬิ \textbf{3. สตฺตาวาสวคฺค}}
\csromanheader{2}{1. ติฐานสุตฺต}
\proseitem{21}{\csromanfolio{198}ตีหิ ภิกฺขเว ฐาเนหิ อุตฺตรกุรุกา มนุสฺสา เทเว จ ตาวติํเส อธิคฺคณฺหนฺติ ชมฺพุทีปเก จ มนุสฺเส.\csromanspacer{} กตเมหิ ตีหิ? อมมา อปริคฺคหา นิยตายุกา วิเสสคุณา\footnote{วิเสสภุโน (สี, สฺยา, อิ)}.\csromanspacer{} อิเมหิ โข ภิกฺขเว ตีหิ ฐาเนหิ อุตฺตรกุรุกา มนุสฺสา เทเว จ ตาวติํเส อธิคฺคณฺหนฺติ ชมฺพุทีปเก จ มนุสฺเส.}

\prose{ตีหิ ภิกฺขเว ฐาเนหิ เทวา ตาวติํสา อุตฺตรกุรุเก จ มนุสฺเส อธิคฺคณฺหนฺติ ชมฺพุทีปเก จ มนุสฺเส.\csromanspacer{} กตเมหิ ตีหิ? ทิพฺเพน อายุนา ทิพฺเพน วณฺเณน ทิพฺเพน สุเขน.\csromanspacer{} อิเมหิ โข ภิกฺขเว ตีหิ ฐาเนหิ เทวา ตาวติํสา อุตฺตรกุรุเก จ มนุสฺเส อธิคฺคณฺหนฺติ ชมฺพุทีปเก จ มนุสฺเส.}

\prose{\csromansymbolfootnote{*}{อภิ 4. 84 ปิฏฺเฐปิ.}ตีหิ ภิกฺขเว ฐาเนหิ ชมฺพุทีปกา มนุสฺสา อุตฺตรกุรเก จ มนุสฺเส อธิคฺคณฺหนฺติ เทเว จ ตาวติํเส.\csromanspacer{} กตเมหิ ตีหิ? สูรา สติมนฺโต อิธ พฺรหฺมจริยวาโส.\csromanspacer{} อิเมหิ โข ภิกฺขเว ตีหิ ฐาเนหิ ชมฺพุทีปกา มนุสฺสา อุตฺตรกุรุเก จ มนุสฺเส อธิคฺคณฺหนฺติ เทเว จ ตาวติํเสติ.\csromanspacer{}\csromanspacer{}ปฐมํ.}

\csromanheader{2}{2. อสฺสขฬุงฺกสุตฺต}
\proseitem{22}{\csromansymbolfootnote{+}{อํ 1. 291 ปิฏฺเฐปิ.}ตโย จ ภิกฺขเว อสฺสขฬุงฺเก เทเสสฺสามิ ตโย จ ปุริสขฬุงฺเก ตโย จ อสฺสปรสฺเส\footnote{อสฺสสทสฺเส (สี, สฺยา, อิ) อํ 1. 292 ปิฏฺเฐปิ.} ตโย จ ปุริสปรสฺเส\footnote{ปุริสสทสฺเส (สี, สฺยา, อิ)} ตโย จ ภทฺเท อสฺสาชานีเย ตโย จ ภทฺเท ปุริสาชานีเย, ตํ สุณาถ.}

\prose{กตเม จ ภิกฺขเว ตโย อสฺสขฬุงฺกา? อิธ ภิกฺขเว เอกจฺโจ อสฺสขฬุงฺโก ชวสมฺปนฺโน โหติ, น วณฺณสมฺปนฺโน น อาโรหปริณาหสมฺปนฺโน.\csromanspacer{} อิธ ปน ภิกฺขเว เอกจฺโจ อสฺสขฬุงฺโก ชวสมฺปนฺโน จ โหติ วณฺณสมฺปนฺโน จ, น อาโรหปริณาหสมฺปนฺโน.\csromanspacer{} อิธ ปน ภิกฺขเว เอกจฺโจ อสฺสขฬุงฺโก ชวสมฺปนฺโน จ โหติ วณฺณสมฺปนฺโน จ อาโรหปริณาหสมฺปนฺโน จ.\csromanspacer{} อิเม โข ภิกฺขเว ตโย อสฺสขฬุงฺกา.}

\chapterheadpageread{3. สตฺตาวาสวคฺค}
\prose{\csromanfolio{199}กตเม จ ภิกฺขเว ตโย ปุริสขฬุงฺกา? อิธ ภิกฺขเว เอกจฺโจ ปุริสขฬุงฺโก ชวสมฺปนฺโน โหติ, น วณฺณสมฺปนฺโน น อาโรหปริณาห- สมฺปนฺโน.\csromanspacer{} อิธ ปน ภิกฺขเว เอกจฺโจ ปุริสขฬุงฺโก ชวสมฺปนฺโน จ โหติ วณฺณสมฺปนฺโน จ, น อาโรหปริณาหสมฺปนฺโน.\csromanspacer{} อิธ ปน ภิกฺขเว เอกจฺโจ ปุริสขฬุงฺโก ชวสมฺปนฺโน จ โหติ วณฺณสมฺปนฺโน จ, อาโรหปริณาหสมฺปนฺโน จ.}

\prose{กถญฺจ ภิกฺขเว ปุริสขฬุงฺโก ชวสมฺปนฺโน โหติ น วณฺณสมฺปนฺโน น อาโรหปริณาหสมฺปนฺโน? อิธ ภิกฺขเว ภิกฺขุ “อิทํ ทุกฺขนฺ”ติ ยถาภูตํ ปชานาติ, “อยํ ทุกฺขสมุทโย”ติ ยถาภูตํ ปชานาติ, “อยํ ทุกฺขนิโรโธ”ติ ยถาภูตํ ปชานาติ, “อยํ ทุกฺขนิโรโคมินี ปฏิปทา”ติ ยถาภูตํ ปชานาติ, อิทมสฺส ชวสฺมิํ วทามิ.\csromanspacer{} อภิธมฺเม โข ปน อภิวินเย ปญฺหํ ปุฏฺโฐ สํสาเทติ\footnote{สํสาเรติ (ก) อํ 1. 292 ปิฏฺฐาทีสุปิ.} โน วิสฺสชฺเชติ, อิทมสฺส น วณฺณสฺมิํ วทามิ.\csromanspacer{} น โข ปน ลาภี โหติ จีวรปิณฺฑปาตเสนาสนคิลานปจฺจยเภสชฺชปริกฺขารานํ, อิทมสฺส น อาโรหปริณาหสฺมิํ วทามิ.\csromanspacer{} เอวํ โข ภิกฺขเว ปุริสขฬุงฺโก ชวสมฺปนฺโน โหติ, น วณฺณสมฺปนฺโน น อาโรหปริณาหสมฺปนฺโน.}

\prose{กถญฺจ ภิกฺขเว ปุริสขฬุงฺโก ชวสมฺปนฺโน จ โหติ วณฺณสมฺปนฺโน จ, น อาโรหปริณาหสมฺปนฺโน? อิธ ภิกฺขเว ภิกฺขุ “อิทํ ทุกฺขนฺ”ติ ยถาภูตํ ปชานาติ ฯเปฯ “อยํ ทุกฺขนิโรธคามินี ปฏิปทา”ติ ยถาภูตํ ปชานาติ, อิทมสฺส ชวสฺมิํ วทามิ.\csromanspacer{} อภิธมฺเม โข ปน อภิวินเย ปญฺหํ ปุฏฺโฐ วิสฺสชฺเชติ โน สํสาเทติ, อิทมสฺส วณฺณสฺมิํ วทามิ.\csromanspacer{} น โข ปน ลาภี โหติ จีวรปิณฺฑปาตเสนาสนคิลาน- ปจฺจยเภสชฺชปริกฺขารานํ, อิทมสฺส น อาโรหปริณาหสฺมิํ วทามิ.\csromanspacer{} เอวํ โข ภิกฺขเว ปุริสขฬุงฺโก ชวสมฺปนฺโน จ โหติ วณฺณสมฺปนฺโน จ, น อาโรหปริณาหสมฺปนฺโน.}

\prose{กถญฺจ ภิกฺขเว ปุริสขฬุงฺโก ชวสมฺปนฺโน จ โหติ วณฺณสมฺปนฺโน จ อาโรหปริณาหสมฺปนฺโน จ? อิธ ภิกฺขเว ภิกฺขุ “อิทํ ทุกฺขนฺ”ติ ยถาภูตํ ปชานาติ ฯเปฯ “อยํ ทุกฺขนิโรธคามินี ปฏิปทา”ติ ยถาภูตํ ปชานาติ, อิทมสฺส ชวสฺมิํ วทามิ.\csromanspacer{} อภิธมฺเม โข ปน อภิวินเย ปญฺหํ ปุฏฺโฐ วิสฺสชฺเชติ, โน สํสาเทติ, อิทมสฺส วณฺณสฺมิํ วทามิ.\csromanspacer{} ลาภี โข}

\vspace{\tipitakaparskip}
\csromancenter{นวกนิปาตปาฬิ ปน โหติ, จีวรปิณฺฑปาตเสนาสน คิลานปจฺจย เภสชฺชปริกฺขารานํ, อิทมสฺส อาโรหปริณาหสฺมิํ วทามิ.\csromanspacer{} เอวํ โข ภิกฺขเว ปุริสขฬุงฺโก ชวสมฺปนฺโน จ โหติ วณฺณสมฺปนฺโน จ อาโรหปริณาหสมฺปนฺโน จ.\csromanspacer{} อิเม โข ภิกฺขเว ตโย ปุริสขฬุงฺกา.}
\prose{\csromanfolio{200}กตเม จ ภิกฺขเว ตโย อสฺสปรสฺสา? อิธ ภิกฺขเว เอกจฺโจ อสฺสปรสฺโส ฯเปฯ ชวสมฺปนฺโน จ โหติ วณฺณสมฺปนฺโน จ อาโรหปริณาหสมฺปนฺโน จ.\csromanspacer{} อิเม โข ภิกฺขเว ตโย อสฺสปรสฺสา.}

\prose{กตเม จ ภิกฺขเว ตโย ปุริสปรสฺสา? อิธ ภิกฺขเว เอกจฺโจ ปุริสปรสฺโส ฯเปฯ ชวสมฺปนฺโน จ โหติ วณฺณสมฺปนฺโน จ อาโรหปริณาหสมฺปนฺโน จ.}

\prose{กถญฺจ ภิกฺขเว ปุริสปรสฺโส ฯเปฯ ชวสมฺปนฺโน จ โหติ วณฺณสมฺปนฺโน จ อาโรหปริณาหสมฺปนฺโน จ? อิธ ภิกฺขเว ภิกฺขุ ปญฺจนฺนํ โอรมฺภาคิยานํ สํโยชนานํ ปริกฺขยา โอปปาติโก โหติ ตตฺถ ปรินิพฺพายี อนาวตฺติธมฺโม ตสฺมา โลกา, อิทมสฺส ชวสฺมิํ วทามิ.\csromanspacer{} อภิธมฺเม โข ปน อภิวินเย ปญฺหํ ปุฏฺโฐ วิสฺสชฺเชติ, โน สํสาเทติ, อิทมสฺส วณฺณสฺมิํ วทามิ.\csromanspacer{} ลาภี โข ปน โหติ จีวรปิณฺฑปาตเสนาสน- คิลานปจฺจยเภสชฺชปริกฺขารานํ, อิทมสฺส อาโรหปริณาหสฺมิํ วทามิ.\csromanspacer{} เอวํ โข ภิกฺขเว ปุริสปรสฺโส ชวสมฺปนฺโน จ โหติ วณฺณสมฺปนฺโน จ อาโรหปริณาหสมฺปนฺโน จ.\csromanspacer{} อิเม โข ภิกฺขเว ตโย ปุริสปรสฺสา.}

\prose{กตเม จ ภิกฺขเว ตโย ภทฺทา อสฺสาชานียา? อิธ ภิกฺขเว เอกจฺโจ ภทฺโท อสฺสาชานีโย ฯเปฯ ชวสมฺปนฺโน จ โหติ วณฺณสมฺปนฺโน จ อาโรหปริณาหสมฺปนฺโน จ.\csromanspacer{} อิเม โข ภิกฺขเว ตโย ภทฺทา อสฺสาชานียา.}

\prose{กตเม จ ภิกฺขเว ตโย ภทฺทา ปุริสาชานียา? อิธ ภิกฺขเว เอกจฺโจ ภทฺโท ปุริสาชานีโย ฯเปฯ ชวสมฺปนฺโน จ โหติ วณฺณสมฺปนฺโน จ อาโรหปริณาหสมฺปนฺโน จ.}

\prose{กถญฺจ ภิกฺขเว ภทฺโท ปุริสาชานีโย ฯเปฯ ชวสมฺปนฺโน จ โหติ วณฺณสมฺปนฺโน จ อาโรหปริณาหสมฺปนฺโน จ? อิธ ภิกฺขเว ภิกฺขุ อาสวานํ ขยา อนาสวํ เจโตวิมุตฺติํ ปญฺญาวิมุตฺติํ ทิฏฺเฐว ธมฺเม สยํ}

\chapterheadpageread{3. สตฺตาวาสวคฺค}
\prose{\csromanfolio{201}อภิญฺญา สจฺฉิกตฺวา อุปสมฺปชฺช วิหรติ, อิทมสฺส ชวสฺมิํ วทามิ.\csromanspacer{} อภิธมฺเม โข ปน อภิวินเย ปญฺหํ ปุฏฺโฐ วิสฺสชฺเชติ, โน สํสาเทติ, อิทมสฺส วณฺณสฺมิํ วทามิ.\csromanspacer{} ลาภี โข ปน โหติ จีวรปิณฺฑปาตเสนาสน- คิลานปจฺจยเภสชฺชปริกฺขารานํ, อิทมสฺส อาโรหปริณาหสฺมิํ วทามิ.\csromanspacer{} เอวํ โข ภิกฺขเว ภทฺโท ปุริสาชานีโย ชวสมฺปนฺโน จ โหติ วณฺณสมฺปนฺโน จ อาโรหปริณาหสมฺปนฺโน จ.\csromanspacer{} อิเม โข ภิกฺขเว ตโย ภทฺทา ปุริสาชานียาติ.\csromanspacer{}\csromanspacer{}ทุติยํ.}

\csromanheader{2}{3. ตณฺหามูลกสุตฺต}
\proseitem{23}{\csromansymbolfootnote{*}{ที 2. 50 ปิฏฺเฐปิ.}นว ภิกฺขเว ตณฺหามูลเก ธมฺเม เทเสสฺสามิ, ตํ สุณาถ.\csromanspacer{} กตเม จ ภิกฺขเว นว ตณฺหามูลกา ธมฺมา? ตณฺหํ ปฏิจฺจ ปริเยสนา, ปริเยสนํ ปฏิจฺจ ลาโภ, ลาภํ ปฏิจฺจ วินิจฺฉโย, วินิจฺฉยํ ปฏิจฺจ ฉนฺทราโค, ฉนฺทราคํ ปฏิจฺจ อชฺโฌสานํ, อชฺโฌสานํ ปฏิจฺจ ปริคฺคโห, ปริคฺคหํ ปฏิจฺจ มจฺฉริยํ, มจฺฉริยํ ปฏิจฺจ อารกฺโข, อารกฺขาธิกรณํ ทณฺฑาทานํ สตฺถาทานํ กลห วิคฺคห วิวาท ตุวํตุวํ เปสุญฺญ มุสาวาทา อเนเก ปาปกา อกุสลา ธมฺมา สมฺภวนฺติ.\csromanspacer{} อิเม โข ภิกฺขเว นว ตณฺหามูลกา ธมฺมาติ.\csromanspacer{}\csromanspacer{}ตติยํ.}

\csromanheader{2}{4. สตฺตาวาสสุตฺต}
\proseitem{24}{\csromansymbolfootnote{+}{ที 3. 218 ปิฏฺเฐปิ.}นวยิเม ภิกฺขเว สตฺตาวาสา.\csromanspacer{} กตเม นว? สนฺติ ภิกฺขเว สตฺตา นานตฺตกายา นานตฺตสญฺญิโน, เสยฺยถาปิ มนุสฺสา เอกจฺเจ จ เทวา เอกจฺเจ จ วินิปาติกา.\csromanspacer{} อยํ ปฐโม สตฺตาวาโส.}

\prose{สนฺติ ภิกฺขเว สตฺตา นานตฺตกายา เอกตฺตสญฺญิโน, เสยฺยถาปิ เทวา พฺรหฺมกายิกา ปฐมาภินิพฺพตฺตา.\csromanspacer{} อยํ ทุติโย สตฺตาวาโส.}

\prose{สนฺติ ภิกฺขเว สตฺตา เอกตฺตกายา นานตฺตสญฺญิโน, เสยฺยถาปิ เทวา อาภสฺสรา.\csromanspacer{} อยํ ตติโย สตฺตาวาโส.}

\prose{สนฺติ ภิกฺขเว สตฺตา เอกตฺตกายา เอกตฺตสญฺญิโน, เสยฺยถาปิ เทวา สุภกิณฺหา.\csromanspacer{} อยํ จตุตฺโถ สตฺตาวาโส.}

\prose{สนฺติ ภิกฺขเว สตฺตา อสญฺญิโน อปฺปฏิสํเวทิโน, เสยฺยถาปิ เทวา อสญฺญสตฺตา.\csromanspacer{} อยํ ปญฺจโม สตฺตาวาโส.}

\gambhira{นวกนิปาตปาฬิ}
\prose{\csromanfolio{202}สนฺติ ภิกฺขเว สตฺตา สพฺพโส รูปสญฺญานํ สมติกฺกมา ปฏิฆสญฺญานํ อตฺถงฺคมา นานตฺตสญฺญานํ อมนสิการา “อนนฺโต อากาโส”ติ อากาสานญฺจายตนุปคา.\csromanspacer{} อยํ ฉฏฺโฐ สตฺตาวาโส.}

\prose{สนฺติ ภิกฺขเว สตฺตา สพฺพโส อากาสานญฺจายตนํ สมติกฺกมฺม “อนนฺตํ วิญฺญาณนฺ”ติ วิญฺญาณญฺจายตนูปคา.\csromanspacer{} อยํ สตฺตโม สตฺตาวาโส.}

\prose{สนฺติ ภิกฺขเว สตฺตา สพฺพโส วิญฺญาณญฺจายตนํ สมติกฺกมฺม “นตฺถิ กิญฺจี”ติ อากิญฺจญฺญายตนูปคา.\csromanspacer{} อยํ อฏฺฐโม สตฺตาวาโส.}

\prose{สนฺติ ภิกฺขเว สตฺตา สพฺพโส อากิญฺจญฺญายตนํ สมติกฺกมฺม เนวสญฺญานาสญฺญายตนูปคา.\csromanspacer{} อยํ นวโม สตฺตาวาโส.\csromanspacer{} อิเม โข ภิกฺขเว นว สตฺตาวาสาติ.\csromanspacer{}\csromanspacer{}จตุตฺถํ.}

\csromanheader{2}{5. ปญฺญาสุตฺต}
\proseitem{25}{ยโต โข ภิกฺขเว ภิกฺขุโน ปญฺญาย จิตฺตํ สุปริจิตํ โหติ, ตสฺเสตํ ภิกฺขเว ภิกฺขุโน กลฺลํ วจนาย “ขีณา ชาติ, วุสิตํ พฺรหฺมจริยํ, กตํ กรณียํ, นาปรํ อิตฺถตฺตายาติ ปชานามี”ติ.}

\prose{กถญฺจ ภิกฺขเว ภิกฺขุโน ปญฺญาย จิตฺตํ สุปริจิตํ โหติ, “วีตราคํ เม จิตฺตนฺ”ติ ปญฺญาย จิตฺตํ สุปริจิตํ โหติ. “วีตโทสํ เม จิตฺตนฺ”ติ ปญฺญาย จิตฺตํ สุปริจิตํ โหติ. “วีตโมหํ เม จิตฺตนฺ”ติ ปญฺญาย จิตฺตํ สุปริจิตํ โหติ. “อสราคธมฺมํ เม จิตฺตนฺ”ติ ปญฺญาย จิตฺตํ สุปริจิตํ โหติ. “อสโทสธมฺมํ เม จิตฺตนฺ”ติ ปญฺญาย จิตฺตํ สุปริจิตํ โหติ.\csromanspacer{} อสโมหธมฺมํ เม จิตฺตนฺ”ติ ปญฺญาย จิตฺตํ สุปริจิตํ โหติ. “อนาวตฺติธมฺมํ เม จิตฺตํ กามภวายา”ติ ปญฺญาย จิตฺตํ สุปริจิตํ โหติ. “อนาวตฺติธมฺมํ เม จิตฺตํ รูปภวายา”ติ ปญฺญาย จิตฺตํ สุปริจิตํ โหติ. “อนาวตฺติธมฺมํ เม จิตฺตํ อรูปภวายา”ติ ปญฺญาย จิตฺตํ สุปริจิตํ โหติ.\csromanspacer{} ยโต โข ภิกฺขเว ภิกฺขุโน ปญฺญาย จิตฺตํ สุปริจิตํ โหติ, ตสฺเสตํ ภิกฺขเว ภิกฺขุโน กลฺลํ วจนาย “ขีณา ชาติ, วุสิตํ พฺรหฺมจริยํ, กตํ กรณียํ, นาปรํ อิตฺถตฺตายาติ ปชานามี”ติ.\csromanspacer{}\csromanspacer{}ปญฺจมํ.}

\chapterheadpageread{3. สตฺตาวาสวคฺค}
\csromanheader{2}{6. สิลายูปสุตฺต}
\proseitem{26}{\csromanfolio{203}เอกํ สมยํ อายสฺมา จ สาริปุตฺโต อายสฺมา จ จนฺทิกาปุตฺโต ราชคเห วิหรติ เวฬุวเน กลนฺทกนิวาเป.\csromanspacer{} ตตฺร โข อายสฺมา จนฺทิกาปุตฺโต ภิกฺขู อามนฺเตสิ ( )\footnote{(อาวุโส ฯเปฯ เอตทโวจ) (สี)} “เทวทตฺโต อาวุโส ภิกฺขูนํ เอวํ ธมฺมํ เทเสติ.\csromanspacer{} ยโต โข อาวุโส ภิกฺขุโน เจตสา จิตํ โหติ.\csromanspacer{} ตสฺเสตํ ภิกฺขุโน กลฺลํ เวยฺยากรณาย ‘ขีณา ชาติ, วุสิตํ พฺรหฺมจริยํ, กตํ กรณียํ, นาปรํ อิตฺถตฺตายาติ ปชานามี’ติ”.}

\prose{เอวํ วุตฺเต อายสฺมา สาริปุตฺโต อายสฺมนฺตํ จนฺทิกาปุตฺตํ เอตทโวจ “น โข อาวุโส จนฺทิกาปุตฺต เทวทตฺโต ภิกฺขูนํ เอวํ ธมฺมํ เทเสติ, ยโต โข อาวุโส ภิกฺขุโน เจตสา จิตํ โหติ, ตสฺเสตํ ภิกฺขุโน กลฺลํ เวยฺยากรณาย ‘ขีณา ชาติ, วุสิตํ พฺรหฺมจริยํ, กตํ กรณียํ, นาปรํ อิตฺถตฺตายาติ ปชานามี’ติ.\csromanspacer{} เอวญฺจ โข อาวุโส จนฺทิกาปุตฺต เทวทตฺโต ภิกฺขูนํ ธมฺมํ เทเสติ.\csromanspacer{} ยโต โข อาวุโส ภิกฺขุโน เจตสา จิตฺตํ สุปริจิตํ โหติ, ตสฺเสตํ ภิกฺขุโน กลฺลํ เวยฺยากรณาย ‘ขีณา ชาติ, วุสิตํ พฺรหฺมจริยํ, กตํ กรณียํ, นาปรํ อิตฺถตฺตายาติ ปชานามี’ติ”.}

\prose{ทุติยมฺปิ โข อายสฺมา จนฺทิกาปุตฺโต ภิกฺขู อามนฺเตสิ “เทวทตฺโต อาวุโส ภิกฺขูนํ เอวํ ธมฺมํ เทเสติ.\csromanspacer{} ยโต โข อาวุโส ภิกฺขุโน เจตสา จิตํ โหติ, ตสฺเสตํ ภิกฺขุโน กลฺลํ เวยฺยากรณาย ‘ขีณา ชาติ, วุสิตํ พฺรหฺมจริยํ, กตํ กรณียํ, นาปรํ อิตฺถตฺตายาติ ปชานามี’ติ”.\csromanspacer{} ทุติยมฺปิ โข อายสฺมา สาริปุตฺโต อายสฺมนฺตํ จนฺทิกาปุตฺตํ เอตทโวจ “น โข อาวุโส จนฺทิกาปุตฺต เทวทตฺโต ภิกฺขูนํ เอวํ ธมฺมํ เทเสติ.\csromanspacer{} ยโต โข อาวุโส ภิกฺขุโน เจตสา จิตํ โหติ, ตสฺเสตํ ภิกฺขุโน กลฺลํ เวยฺยากรณาย ‘ขีณา ชาติ วุสิตํ พฺรหฺมจริยํ, กตํ กรณียํ, นาปรํ อิตฺถตฺตายาติ ปชานามี’ติ.\csromanspacer{} เอวญฺจ โข อาวุโส จนฺทีกาปุตฺต เทวทตฺโต ภิกฺขูนํ ธมฺมํ เทเสติ.\csromanspacer{} ยโต โข อาวุโส ภิกฺขุโน เจตสา จิตฺตํ สุปริจิตํ โหติ, ตสฺเสตํ ภิกฺขุโน กลฺลํ เวยฺยากรณาย ‘ขีณา ชาติ, วุสิตํ พฺรหฺมจริยํ, กตํ กรณียํ, นาปรํ อิตฺถตฺตายาติ ปชานามี’ติ”.}

\gambhira{นวกนิปาตปาฬิ}
\prose{\csromanfolio{204}ตติยมฺปิ โข อายสฺมา จนฺทิกาปุตฺโต ภิกฺขู อามนฺเตสิ “เทวทตฺโต อาวุโส ภิกฺขูนํ เอวํ ธมฺมํ เทเสติ.\csromanspacer{} ยโต โข อาวุโส ภิกฺขุโน เจตสา จิตํ โหติ, ตสฺเสตํ ภิกฺขุโน กลฺลํ เวยฺยากรณาย ‘ขีณา ชาติ, วุสิตํ พฺรหฺมจริยํ, กตํ กรณียํ, นาปรํ อิตฺถตฺตายาติ ปชานามี’ติ”.\csromanspacer{} ตติยมฺปิ โข อายสฺมา สาริปุตฺโต อายสฺมนฺตํ จนฺทิกาปุตฺตํ เอตทโวจ “น โข อาวุโส จนฺทิกาปุตฺต เทวทตฺโต ภิกฺขูนํ เอวํ ธมฺมํ เทเสติ.\csromanspacer{} ยโต โข อาวุโส ภิกฺขุโน เจตสา จิตํ โหติ, ตสฺเสตํ ภิกฺขุโน กลฺลํ เวยฺยากรณาย ‘ขีณา ชาติ, วุสิตํ พฺรหฺมจริยํ, กตํ กรณียํ, นาปรํ อิตฺถตฺตายาติ ปชานามี’ติ.\csromanspacer{} เอวญฺจ โข อาวุโส จนฺทิกาปุตฺต เทวทตฺโต ภิกฺขูนํ ธมฺมํ เทเสติ.\csromanspacer{} ยโต โข อาวุโส ภิกฺขุโน เจตสา จิตฺตํ สุปริจิตํ โหติ, ตสฺเสตํ ภิกฺขุโน กลฺลํ เวยฺยากรณาย ‘ขีณา ชาติ, วุสิตํ พฺรหฺมจริยํ, กตํ กรณียํ, นาปรํ อิตฺถตฺตายาติ ปชานามี’ติ”.}

\prose{กถญฺจ อาวุโส ภิกฺขุโน เจตสา จิตฺตํ สุปริจิตํ โหติ, “วีตราคํ เม จิตฺตนฺ”ติ เจตสา จิตฺตํ สุปริจิตํ โหติ. “วีตโทสํ เม จิตฺตนฺ”ติ เจตสา จิตฺตํ สุปริจิตํ โหติ. “วีตโมหํ เม จิตฺตนฺ”ติ เจตสา จิตฺตํ สุปริจิตํ โหติ. “อสราคธมฺมํ เม จิตฺตนฺ”ติ เจตสา จิตฺตํ สุปริจิตํ โหติ. “อสโทสธมฺมํ เม จิตฺตนฺ”ติ เจตสา จิตฺตํ สุปริจิตํ โหติ. “อสโมหธมฺมํ เม จิตฺตนฺ”ติ เจตสา จิตฺตํ สุปริจิตํ โหติ. “อนาวตฺติธมฺมํ เม จิตฺตํ กามภวายา”ติ เจตสา จิตฺตํ สุปริจิตํ โหติ.\csromanspacer{} อนาวตฺติธมฺมํ เม จิตฺตํ รูปภวายา”ติ เจตสา จิตฺตํ สุปริจิตํ โหติ. “อนาวตฺติธมฺมํ เม จิตฺตํ อรูปภวายา”ติ เจตสา จิตฺตํ สุปริจิตํ โหติ.\csromanspacer{} เอวํ สมฺมา วิมุตฺตจิตฺตสฺส โข อาวุโส ภิกฺขุโน ภุสา เจปิ จกฺขุวิญฺเญยฺยา รูปา จกฺขุสฺส อาปาถํ อาคจฺฉนฺติ, เนวสฺส จิตฺตํ ปริยาทิยนฺติ, อมิสฺสีกตเมวสฺส จิตฺตํ โหติ ฐิตํ อาเนญฺชปฺปตฺตํ, วยํ จสฺสานุปสฺสติ.}

\prose{เสยฺยถาปิ อาวุโส สิลายูโป โสฬสกุกฺกุโก, ตสฺสสฺสุ อฏฺฐ กุกฺกู เหฏฺฐา เนมงฺคมา อฏฺฐ กุกฺกู อุปริ เนมสฺส.\csromanspacer{} อถ ปุรตฺถิมาย เจปิ ทิสาย อาคจฺเฉยฺย ภุสา วาตวุฏฺฐิ, เนว นํ สงฺกมฺเปยฺย น สมฺปเวเธยฺย.\csromanspacer{} อถ ปจฺฉิมาย.\csromanspacer{} อถ อุตฺตราย.\csromanspacer{} อถ ทกฺขิณาย เจปิ}

\chapterheadpageread{3. สตฺตาวาสวคฺค}
\prose{\csromanfolio{205}ทิสาย อาคจฺเฉยฺย ภุสา วาตวุฏฺฐิ, เนว นํ สงฺกมฺเปยฺย น สมฺปเวเธยฺย.\csromanspacer{} ตํ กิสฺส เหตุ? คมฺภีรตฺตา อาวุโส เนมสฺส สุนิขาตตฺตา สิลายูปสฺส.\csromanspacer{} เอวเมวํ โข อาวุโส สมฺมา วิมุตฺตจิตฺตสฺส ภิกฺขุโน ภุสา เจปิ จกฺขุวิญฺเญยฺยา รูปา จกฺขุสฺส อาปาถํ อาคจฺฉนฺติ, เนวสฺส จิตฺตํ ปริยาทิยนฺติ, อมิสฺสีกตเมวสฺส จิตฺตํ โหติ ฐิตํ อาเนญฺชปฺปตฺตํ, วยํ จสฺสานุปสฺสติ.}

\prose{ภุสา เจปิ โสตวิญฺเญยฺยา สทฺทา.\csromanspacer{} ฆานวิญฺเญยฺยา คนฺธา.\csromanspacer{} ชิวฺหาวิญฺเญยฺยา รสา.\csromanspacer{} กายวิญฺเญยฺยา โผฏฺฐพฺพา.\csromanspacer{} มโนวิญฺเญยฺยา ธมฺมา มนสฺส อาปาถํ อาคจฺฉนฺติ, เนวสฺส จิตฺตํ ปริยาทิยนฺติ, อมิสฺสีกตเมวสฺส จิตฺตํ โหติ ฐิตํ อาเนญฺชปฺปตฺตํ, วยํ จสฺสานุปสฺสตีติ.\csromanspacer{}\csromanspacer{}ฉฏฺฐํ. \textbf{7.}\csromanspacer{}\textbf{ ปฐมเวรสุตฺต}}

\proseitem{27}{\csromansymbolfootnote{*}{อุปริ 404; สํ 3. 338 ปิฏฺเฐปิ.}อถ โข อนาถปิณฺฑิโก คหปติ เยน ภควา เตนุปสงฺกมิ, อุปสงฺกมิตฺวา ภควนฺตํ อภิวาเทตฺวา เอกมนฺตํ นิสีทิ, เอกมนฺตํ นิสินฺนํ โข อนาถปิณฺฑิกํ คหปติํ ภควา เอตทโวจ\csromandash{}}

\prose{ยโต โข คหปติ อริยสาวกสฺส ปญฺจ ภยานิ เวรานิ วูปสนฺตานิ โหนฺติ, จตูหิ จ โสตาปตฺติยงฺเคหิ สมนฺนาคโต โหติ.\csromanspacer{} โส อากงฺขมาโน อตฺตนาว อตฺตานํ พฺยากเรยฺย “ขีณนิรโยมฺหิ ขีณติรจฺฉานโยนิ ขีณเปตฺติวิสโย ขีณาปายทุคฺคติวินิปาโต, โสตาปนฺโนหมสฺมิ อวินิปาตธมฺโม นิยโต สมฺโพธิปรายโณ”ติ.}

\prose{กตมานิ ปญฺจ ภยานิ เวรานิ วูปสนฺตานิ โหนฺติ? ยํ คหปติ ปาณาติปาตี ปาณาติปาตปจฺจยา ทิฏฺฐธมฺมิกมฺปิ ภยํ เวรํ ปสวติ, สมฺปรายิกมฺปิ ภยํ เวรํ ปสวติ, เจตสิกมฺปิ ทุกฺขํ โทมนสฺสํ ปฏิสํเวเทติ, ปาณาติปาตา ปฏิวิรโต เนว ทิฏฺฐธมฺมิกมฺปิ ภยํ เวรํ ปสวติ, น สมฺปรายิกมฺปิ ภยํ เวรํ ปสวติ, น เจตสิกมฺปิ ทุกฺขํ โทมนสฺสํ ปฏิสํเวเทติ.\csromanspacer{} ปาณาติปาตา ปฏิวิรตสฺส เอวํ ตํ ภยํ เวรํ วูปสนฺตํ โหติ.}

\prose{ยํ คหปติ อทินฺนาทายี ฯเปฯ กาเมสุมิจฺฉาจารี.\csromanspacer{} มุสาวาที.\csromanspacer{} สุราเมรยมชฺชปมาทฏฺฐายี สุราเมรยมชฺชปมาทฏฺฐานปจฺจยา ทิฏฺฐธมฺมิกมฺปิ ภยํ เวรํ ปสวติ, สมฺปรายิกมฺปิ ภยํ เวรํ ปสวติ, เจตสิกมฺปิ ทุกฺขํ}

\vspace{\tipitakaparskip}
\csromancenter{นวกนิปาตปาฬิ โทมนสฺสํ ปฏิสํเวเทติ, สุราเมรยมชฺชปมาทฏฺฐานา ปฏิวิรโต เนว ทิฏฺฐธมฺมิกมฺปิ ภยํ เวรํ ปสวติ, น สมฺปรายิกมฺปิ ภยํ เวรํ ปสวติ, น เจตสิกมฺปิ ทุกฺขํ โทมนสฺสํ ปฏิสํเวเทติ.\csromanspacer{} สุราเมรยมชฺชปมาทฏฺฐานา ปฏิวิรตสฺส เอวํ ตํ ภยํ เวรํ วูปสนฺตํ โหติ.\csromanspacer{} อิมานิ ปญฺจ ภยานิ เวรานิ วูปสนฺตานิ โหนฺติ.}
\prose{\csromanfolio{206}กตเมหิ จตูหิ โสตาปตฺติยงฺเคหิ สมนฺนาคโต โหติ? อิธ คหปติ อริยสาวโก พุทฺเธ อเวจฺจปฺปสาเทน สมนฺนาคโต โหติ “อิติปิ โส ภควา อรหํ สมฺมาสมฺพุทฺโธ วิชฺชาจรณสมฺปนฺโน สุคโต โลกวิทู อนุตฺตโร ปุริสทมฺมสารถิ สตฺถา เทวมนุสฺสานํ พุทฺโธ ภควา”ติ.}

\prose{ธมฺเม อเวจฺจปฺปสาเทน สมนฺนาคโต โหติ “สฺวากฺขาโต ภควตา ธมฺโม สนฺทิฏฺฐิโก อกาลิโก เอหิปสฺสิโก โอปเนยฺยิโก ปจฺจตฺตํ เวทิตพฺโพ วิญฺญูหี”ติ.}

\prose{สํเฆ อเวจฺจปฺปสาเทน สมนฺนาคโต โหติ “สุปฺปฏิปนฺโน ภควโต สาวกสํโฆ, อุชุปฺปฏิปนฺโน ภควโต สาวกสํโฆ, ญายปฺปฏิปนฺโน ภควโต สาวกสํโฆ, สามีจิปฺปฏิปนฺโน ภควโต สาวกสํโฆ, ยทิทํ จตฺตาริ ปุริสยุคานิ อฏฺฐ ปุริสปุคฺคลา เอส ภควโต สาวกสํโฆ อาหุเนยฺโย ปาหุเนยฺโย ทกฺขิเณยฺโย อญฺชลิกรณีโย อนุตฺตรํ ปุญฺญกฺเขตฺตํ โลกสฺสา”ติ.}

\prose{อริยกนฺเตหิ สีเลหิ สมนฺนาคโต โหติ อขณฺเฑหิ อจฺฉิทฺเทหิ อสพเลหิ อกมฺมาเสหิ ภุชิสฺเสหิ วิญฺญุปฺปสตฺเถหิ อปรามฏฺเฐหิ สมาธิสํวตฺตนิเกหิ.\csromanspacer{} อิเมหิ จตูหิ โสตาปตฺติยงฺเคหิ สมนฺนาคโต โหติ.}

\prose{ยโต โข คหปติ อริยสาวกสฺส อิมานิ ปญฺจ ภยานิ เวรานิ วูปสนฺตานิ โหนฺติ, อิเมหิ จ จตูหิ โสตปตฺติยงฺเคหิ สมนฺนาคโต โหติ.\csromanspacer{} โส อากงฺขมาโน อตฺตนาว อตฺตานํ พฺยากเรยฺย “ขีณนิรโยมฺหิ ขีณติรจฺฉานโยนิ ขีณเปตฺติวิสโย ขีณาปายทุคฺคติวินิปาโต, โสตาปนฺโนหมสฺมิ อวินิปาตธมฺโม นิยโต สมฺโพธิปรายโณ”ติ.\csromanspacer{}\csromanspacer{}สตฺตมํ.}

\chapterheadpageread{3. สตฺตาวาสวคฺค}
\csromanheader{2}{8. ทุติยเวรสุตฺต}
\proseitem{28}{\csromanfolio{207}\csromansymbolfootnote{*}{สํ 3. 340 ปิฏฺเฐปิ.}ยโต โข ภิกฺขเว อริยสาวกสฺส ปญฺจ ภยานิ เวรานิ วูปสนฺตานิ โหนฺติ, จตูหิ จ โสตาปตฺติยงฺเคหิ สมนฺนาคโต โหติ.\csromanspacer{} โส อากงฺขมาโน อตฺตนาว อตฺตานํ พฺยากเรยฺย “ขีณนิรโยมฺหิ ขีณติรจฺฉานโยนิ ขีณเปตฺติวิสโย ขีณาปายทุคฺคติวินิปาโต, โสตาปนฺโนหมสฺมิ อวินิปาตธมฺโม นิยโต สมฺโพธิปรายโณ”ติ.}

\prose{กตมานิ ปญฺจ ภยานิ เวรานิ วูปสนฺตานิ โหนฺติ? ยํ ภิกฺขเว ปาณาติปาตี ปาณาติปาตปจฺจยา ทิฏฺฐธมฺมิกมฺปิ ภยํ เวรํ ปสวติ, สมฺปรายิกมฺปิ ภยํ เวรํ ปสวติ, เจตสิกมฺปิ ทุกฺขํ โทมนสฺสํ ปฏิสํเวเทติ.\csromanspacer{} ปาณาติปาตา ปฏิวิรโต ฯเปฯ เอวํ ตํ ภยํ เวรํ วูปสนฺตํ โหติ.}

\prose{ยํ ภิกฺขเว อทินฺนาทายี ฯเปฯ สุราเมรยมชฺชปมาทฏฺฐายี สุราเมรยมชฺชปมาทฏฺฐานปจฺจยา ทิฏฺฐธมฺมิกมฺปิ ภยํ เวรํ ปสวติ, สมฺปรายิกมฺปิ ภยํ เวรํ ปสวติ, เจตสิกมฺปิ ทุกฺขํ โทมนสฺสํ ปฏิสํเวเทติ, สุราเมรยมชฺชปมาทฏฺฐานา ปฏิวิรโต เนว ทิฏฺฐธมฺมิกมฺปิ ภยํ เวรํ ปสวติ, น สมฺปรายิกมฺปิ ภยํ เวรํ ปสวติ, น เจตสิกมฺปิ ทุกฺขํ โทมนสฺสํ ปฏิสํเวเทติ.\csromanspacer{} สุราเมรยมชฺชปมาทฏฺฐานา ปฏิวิรตสฺส เอวํ ตํ ภยํ เวรํ วูปสนฺตํ โหติ.\csromanspacer{} อิมานิ ปญฺจ ภยานิ เวรานิ วูปสนฺตานิ โหนฺติ.}

\prose{กตเมหิ จตูหิ โสตาปตฺติยงฺเคหิ สมนฺนาคโต โหติ? อิธ ภิกฺขเว อริยสาวโก พุทฺเธ อเวจฺจปฺปสาเทน สมนฺนาคโต โหติ “อิติปิ โส ภควา ฯเปฯ สตฺถา เทวมนุสฺสานํ พุทฺโธ ภควา”ติ.\csromanspacer{} ธมฺเม ฯเปฯ.\csromanspacer{} สํเฆ.\csromanspacer{} อริยกนฺเตหิ สีเลหิ สมนฺนาคโต โหติ อขณฺเฑหิ อจฺฉิทฺเทหิ อสพเลหิ อกมฺมาเสหิ ภุชิสฺเสหิ วิญฺญุปฺปสตฺเถหิ อปรามฏฺเฐหิ สมาธิสํวตฺตนิเกหิ.\csromanspacer{} อิเมหิ จตูหิ โสตาปตฺติยงฺเคหิ สมนฺนาคโต โหติ.}

\prose{ยโต โข ภิกฺขเว อริยสาวกสฺส อิมานิ ปญฺจ ภยานิ เวรานิ วูปสนฺตานิ โหนฺติ, อิเมหิ จ จตูหิ โสตาปตฺติยงฺเคหิ สมนฺนาคโต โหติ.\csromanspacer{} โส อากงฺขมาโน อตฺตนาว อตฺตานํ พฺยากเรยฺย “ขีณนิรโยมฺหิ ขีณติรจฺฉานโยนิ ขีณเปตฺติวิสโย}

\vspace{\tipitakaparskip}
\csromancenter{นวกนิปาตปาฬิ ขีณาปายทุคฺคติวินิปาโต, โสตาปนฺโนหมสฺมิ อวินิปาตธมฺโม นิยโต สมฺโพธิปรายโณ”ติ.\csromanspacer{}\csromanspacer{}อฏฺฐมํ.}
\csromanheader{2}{9. อาฆาตวตฺถุสุตฺต}
\proseitem{29}{\csromanfolio{208}\csromansymbolfootnote{*}{อภิ 2. 403; ที 3. 218, อุปริ 375 ปิฏฺเฐสุปิ.}นวยิมานิ ภิกฺขเว อาฆาตวตฺถูนิ.\csromanspacer{} กตมานิ นว? “อนตฺถํ เม อจรี”ติ อาฆาตํ พนฺธติ, “อนตฺถํ เม จรตี”ติ อาฆาตํ พนฺธติ, “อนตฺถํ เม จริสฺสตี”ติ อาฆาตํ พนฺธติ, “ปิยสฺส เม มนาปสฺส อนตฺถํ อจรี”ติ ฯเปฯ อนตฺถํ จรตี”ติ ฯเปฯ อนตฺถํ จริสฺสตี”ติ อาฆาตํ พนฺธติ, “อปฺปิยสฺส เม อมนาปสฺส อตฺถํ อจรี”ติ ฯเปฯ อตฺถํ จรตี”ติ ฯเปฯ อตฺถํ จริสฺสตี”ติ อาฆาตํ พนฺธติ.\csromanspacer{} อิมานิ โข ภิกฺขเว นว อาฆาตวตฺถูนีติ.\csromanspacer{}\csromanspacer{}นวมํ.}

\csromanheader{2}{10. อาฆาตปฏิวินยสุตฺต}
\proseitem{30}{\csromansymbolfootnote{+}{ที 3. 218, 255 ปิฏฺเฐสุ.}นวยิเม ภิกฺขเว อาฆาตปฏิวินยา.\csromanspacer{} กตเม นว? “อนตฺถํ เม อจริ1, ตํ กุเตตฺถ ลพฺภา”ติ อาฆาตํ ปฏิวิเนติ, “อนตฺถํ เม จรติ, ตํ กุเตตฺถ ลพฺภา”ติ อาฆาตํ ปฏิวิเนติ, “อนตฺถํ เม จริสฺสติ, ตํ กุเตตฺถ ลพฺภา”ติ อาฆาตํ ปฏิวิเนติ, ปิยสฺส เม มนาปสฺส อนตฺถํ อจริ ฯเปฯ อนตฺถํ จรติ ฯเปฯ อนตฺตํ จริสฺสติ, ตํ กุเตตฺถ ลพฺภา”ติ อาฆาตํ ปฏิวิเนติ, “อปฺปิยสฺส เม อมนาปสฺส อตฺถํ อจริ ฯเปฯ อตฺถํ จรติ ฯเปฯ อตฺถํ จริสฺสติ, ตํ กุเตตฺถ ลพฺภา”ติ อาฆาตํ ปฏิวิเนติ.\csromanspacer{} อิเม โข ภิกฺขเว นว อาฆาตปฏิวินยาติ.\csromanspacer{}\csromanspacer{}ทสมํ.}

\csromanheader{2}{11. อนุปุพฺพนิโรธสุตฺต}
\proseitem{31}{นวยิเม ภิกฺขเว อนุปุพฺพนิโรธา, กตเม นว? ปฐมํ ฌานํ สมาปนฺนสฺส กามสญฺญา\footnote{อามิสฺสสญฺญา (สฺยา)} นิรุทฺธา โหติ, ทุติยํ ฌานํ สมาปนฺนสฺส วิตกฺกวิจารา นิรุทฺธา โหติ, ตติยํ ฌานํ สมาปนฺนสฺส ปีติ นิรุทฺธา โหติ, จตุตฺถํ ฌานํ สมาปนฺนสฺส อสฺสาสปสฺสาสา นิรุทฺธา โหนฺติ, อากาสานญฺจายตนํ สมาปนฺนสฺส รูปสญฺญา นิรุทฺธา โหติ, วิญฺญาณญฺจายตนํ สมาปนฺนสฺส อากาสานญฺจายตนสญฺญา นิรุทฺธา โหติ, อากิญฺจญฺญายตนํ}

\chapterheadpageread{4. มหาวคฺค}
\prose{\csromanfolio{209}สมาปนฺนสฺส วิญฺญาณญฺจายตนสญฺญา นิรุทฺธา โหติ, เนวสญฺญานาสญฺญายตนํ สมาปนฺนสฺส อากิญฺจญฺญายตนสญฺญา นิรุทฺธา โหติ, สญฺญาเวทยิตนิโรธํ สมาปนฺนสฺส สญฺญา จ เวทนา จ นิรุทฺธา โหนฺติ.\csromanspacer{} อิเม โข ภิกฺขเว นว อนุปุพฺพนิโรธาติ * .\csromanspacer{}\csromanspacer{}เอกาทสมํ.\csromanspacer{} สตฺตาวาสวคฺโค ตติโย.}

\titlehead{\textbf{ตสฺสุทฺทานํ}}
\csromangathameasure{ติฐานํ ขฬุงฺโก ตณฺหา, สตฺตปญฺญา สิลายุโป. \\ เทฺว เวรา เทฺว อาฆาตานิ, อนุปุพฺพนิโรเธน จาติ.}
\csromangathagroup{\csromangathastanza{ติฐานํ ขฬุงฺโก ตณฺหา, สตฺตปญฺญา สิลายุโป. \\ เทฺว เวรา เทฺว อาฆาตานิ, อนุปุพฺพนิโรเธน จาติ.}}

\chapterhead{4. มหาวคฺค}
\csromanheader{2}{1. อนุปุพฺพวิหารสุตฺต}
\proseitem{32}{\csromansymbolfootnote{*}{ที 3. 221, 254 ปิฏฺเฐสุปิ.}นวยิเม ภิกฺขเว อนุปุพฺพวิหารา.\csromanspacer{} กตเม นว?\footnote{เอตฺถ สี-อิ-โปตฺถเกสุ “อิธ ภิกฺขเว ภิกฺขุ วิวิจฺเจว กาเมหี”ติ-อาทินา วิตฺถาเรน ปาฏฺโฐ ทิสฺสติ. 2. นว ภิกฺขเว (?)} ปฐมํ ฌานํ ทุติยํ ฌานํ ตติยํ ฌานํ จตุตฺถํ ฌานํ อากาสานญฺจายตนํ วิญฺญาณญฺจายตนํ อากิญฺจญฺญายตนํ เนวสญฺญานาสญฺญายตนํ สญฺญาเวทยิตนิโรโธ.\csromanspacer{} อิเม โข ภิกฺขเว นว อนุปุพฺพวิหาราติ.\csromanspacer{}\csromanspacer{}ปฐมํ.}

\csromanheader{2}{2. อนุปุพฺพวิหารสมาปตฺติสุตฺต}
\proseitem{33}{นวยิมา ภิกฺขเว2 อนุปุพฺพวิหารสมาปตฺติโย เทเสสฺสามิ, ตํ สุณาถ ฯเปฯ.\csromanspacer{} กตมา จ ภิกฺขเว นว อนุปุพฺพวิหารสมาปตฺติโย? ยตฺถ กามา นิรุชฺฌนฺติ, เย จ กาเม นิโรเธตฺวา นิโรเธตฺวา วิหรนฺติ, อทฺธา เต อายสฺมนฺโต นิจฺฉาตา นิพฺพุตา ติณฺณา ปารงฺคตา ตทงฺเคนาติ วทามิ. “กตฺถ กามา นิรุชฺฌนฺติ, เก จ กาเม นิโรเธตฺวา นิโรเธตฺวา วิหรนฺติ, อหเมตํ น ชานามิ, อหเมตํ น ปสฺสามี”ติ อิติ โย เอวํ วเทยฺย, โส เอวมสฺส วจนีโย “อิธาวุโส ภิกฺขุ วิวิจฺเจว กาเมหิ วิวิจฺจ อกุสเลหิ ธมฺเมหิ สวิตกฺกํ สวิจารํ วิเวกชํ ปีติสุขํ ปฐมํ ฌานํ}

\vspace{\tipitakaparskip}
\csromancenter{นวกนิปาตปาฬิ อุปสมฺปชฺช วิหรติ, เอตฺถ กามา นิรุชฺฌนฺติ, เต จ กาเม นิโรเธตฺวา นิโรเธตฺวา วิหรนฺตี”ติ.\csromanspacer{} อทฺธา ภิกฺขเว อสโฐ อมายาวี “สาธู”ติ ภาสิตํ อภินนฺเทยฺย อนุโมเทยฺย, “สาธู”ติ ภาสิตํ อภินนฺทิตฺวา อนุโมทิตฺวา นมสฺสมาโน ปญฺชลิโก ปยิรุปาเสยฺย.}
\prose{\csromanfolio{210}ยตฺถ วิตกฺกวิจารา นิรุชฺฌนฺติ, เย จ วิตกฺกวิจาเร นิโรเธตฺวา นิโรเธตฺวา วิหรนฺติ, อทฺธา เต อายสฺมนฺโต นิจฺฉาตา นิพฺพุตา ติณฺณา ปารงฺคตา ตทงฺเคนาติ วทามิ. “กตฺถ วิตกฺกวิจารา นิรุชฺฌนฺติ, เก จ วิตกฺกวิจาเร นิโรเธตฺวา นิโรเธตฺวา วิหรนฺติ, อหเมตํ น ชานามิ อหเมตํ น ปสฺสามี”ติ อิติ โย เอวํ วเทยฺย, โส เอวมสฺส วจนีโย “อิธาวุโส ภิกฺขุ วิตกฺกวิจารานํ วูปสมา ฯเปฯ ทุติยํ ฌานํ อุปสมฺปชฺช วิหรติ, เอตฺถ วิตกฺกวิจารา นิรุชฺฌนฺติ, เต จ วิตกฺกวิจาเร นิโรเธตฺวา นิโรเธตฺวา วิหรนฺตี”ติ.\csromanspacer{} อทฺธา ภิกฺขเว อสโฐ อมายาวี “สาธู”ติ ภาสิตํ อภินนฺเทยฺย อนุโมเทยฺย, “สาธู”ติ ภาสิตํ อภินนฺทิตฺวา อนุโมทิตฺวา นมสฺสมาโน ปญฺชลิโก ปยิรุปาเสยฺย.}

\prose{ยตฺถ ปีติ นิรุชฺฌติ, เย จ ปีติํ นิโรเธตฺวา นิโรเธตฺวา วิหรนฺติ, อทฺธา เต อายสฺมนฺโต นิจฺฉาตา นิพฺพุตา ติณฺณา ปารงฺคตา ตทงฺเคนาติ วทามิ. “กตฺถ ปีติ นิรุชฺฌติ, เก จ ปีติํ นิโรเธตฺวา นิโรเธตฺวา วิหรนฺติ, อหเมตํ น ชานามิ, อหเมตํ น ปสฺสามี”ติ อิติ โย เอวํ วเทยฺย, โส เอวมสฺส วจนีโย “อิธาวุโส ภิกฺขุ ปีติยา จ วิราคา ฯเปฯ ตติยํ ฌานํ อุปสมฺปชฺช วิหรติ, เอตฺถ ปีติ นิรุชฺฌติ, เต จ ปีติํ นิโรเธตฺวา นิโรเธตฺวา วิหรนฺตี”ติ.\csromanspacer{} อทฺธา ภิกฺขเว อสโฐ อมายาวี “สาธู”ติ ภาสิตํ อภินนฺเทยฺย อนุโมเทยฺย, “สาธู”ติ ภาสิตํ อภินนฺทิตฺวา อนุโมทิตฺวา นมสฺสมาโน ปญฺชลิโก ปยิรุปาเสยฺย.}

\prose{ยตฺถ อุเปกฺขาสุขํ นิรุชฺฌติ, เย จ อุเปกฺขาสุขํ นิโรเธตฺวา นิโรเธตฺวา วิหรนฺติ, อทฺธา เต อายสฺมนฺโต นิจฺฉาตา นิพฺพุตา ติณฺณา ปารงฺคตา ตทงฺเคนาติ วทามิ. “กตฺถ อุเปกฺขาสุขํ นิรุชฺฌติ, เก จ อุเปกฺขาสุขํ นิโรเธตฺวา นิโรเธตฺวา วิหรนฺติ, อหเมตํ น ชานามิ, อหเมตํ น ปสฺสามี”ติ อิติ โย เอวํ วเทยฺย, โส เอวมสฺส วจนีโย “อิธาวุโส ภิกฺขุ สุขสฺส จ ปหานา ฯเปฯ จตุตฺถํ ฌานํ อุปสมฺปชฺช วิหรติ, เอตฺถ อุเปกฺขาสุขํ นิรุชฺฌติ, เต จ อุเปกฺขาสุขํ นิโรเธตฺวา นิโรเธตฺวา วิหรนฺตี”ติ.\csromanspacer{} อทฺธา ภิกฺขเว อสโฐ อมายาวี “สาธู”ติ ภาสิตํ}

\chapterheadpageread{4. มหาวคฺค}
\prose{\csromanfolio{211}อภินนฺเทยฺย อนุโมเทยฺย, “สาธู”ติ ภาสิตํ อภินนฺทิตฺวา อนุโมทิตฺวา นมสฺสมาโน ปญฺชลิโก ปยิรุปาเสยฺย.}

\prose{ยตฺถ รูปสญฺญา นิรุชฺฌติ, เย จ รูปสญฺญํ\footnote{ยตฺถ รูปสญฺญา นิรุชฺฌนฺติ, เย จ รูปสญฺญา (สี, สฺยา, อิ)} นิโรเธตฺวา นิโรเธตฺวา วิหรนฺติ, อทฺธา เต อายสฺมนฺโต นิจฺฉาตา นิพฺพุตา ติณฺณา ปารงฺคตา ตทงฺเคนาติ วทามิ. “กตฺถ รูปสญฺญา นิรุชฺฌติ, เก จ รูปสญฺญํ นิโรเธตฺวา นิโรเธตฺวา วิหรนฺติ, อหเมตํ น ชานามิ, อหเมตํ น ปสฺสามี”ติ อิติ โย เอวํ วเทยฺย, โส เอวมสฺส วจนีโย “อิธาวุโส ภิกฺขุ สพฺพโส รูปสญฺญานํ สมติกฺกมา ปฏิฆสญฺญานํ อตฺถงฺคมา นานตฺตสญฺญานํ อมนสิการา ‘อนนฺโต อากาโส’ติ อากาสานญฺจายตนํ อุปสมฺปชฺช วิหรติ, เอตฺถ รูปสญฺญา นิรุชฺฌติ, เต จ รูปสญฺญํ นิโรเธตฺวา นิโรเธตฺวา วิหรนฺตี”ติ.\csromanspacer{} อทฺธา ภิกฺขเว อสโฐ อมายาวี “สาธู”ติ ภาสิตํ อภินนฺเทยฺย อนุโมเทยฺย, “สาธู”ติ ภาสิตํ อภินนฺทิตฺวา อนุโมทิตฺวา นมสฺสมาโน ปญฺชลิโก ปยิรุปาเสยฺย.}

\prose{ยตฺถ อากาสานญฺจายตนสญฺญา นิรุชฺฌติ, เย จ อากาสานญฺจายตนสญฺญํ นิโรเธตฺวา นิโรเธตฺวา วิหรนฺติ, อทฺธา เต อายสฺมนฺโต นิจฺฉาตา นิพฺพุตา ติณฺณา ปารงฺคตา ตทงฺเคนาติ วทามิ. “กตฺถ อากาสานญฺจายตนสญฺญา นิรุชฺฌติ, เก จ อากาสานญฺจายตนสญฺญํ นิโรเธตฺวา นิโรเธตฺวา วิหรนฺติ, อหเมตํ น ชานามิ, อหเมตํ น ปสฺสามี”ติ อิติ โย เอวํ วเทยฺย, โส เอวมสฺส วจนีโย “อิธาวุโส ภิกฺขุ สพฺพโส อากาสานญฺจายตนํ สมติกฺกมฺม ‘อนนฺตํ วิญฺญาณนฺ’ติ วิญฺญาณญฺจายตนํ อุปสมฺปชฺช วิหรติ, เอตฺถ อากาสานญฺจายตนสญฺญา นิรุชฺฌติ, เต จ อากาสานญฺจายตนสญฺญํ นิโรเธตฺวา นิโรเธตฺวา วิหรนฺตี”ติ.\csromanspacer{} อทฺธา ภิกฺขเว อสโฐ อมายาวี “สาธู”ติ ภาสิตํ อภินนฺเทยฺย อนุโมเทยฺย, “สาธู”ติ ภาสิตํ อภินนฺทิตฺวา อนุโมทิตฺวา นมสฺสมาโน ปญฺชลิโก ปยิรุปาเสยฺย.}

\prose{ยตฺถ วิญฺญาณญฺจายตนสญฺญา นิรุชฺฌติ, เย จ วิญฺญาณญฺจายตน- สญฺญํ นิโรเธตฺวา นิโรเธตฺวา วิหรนฺติ, อทฺธา เต อายสฺมนฺโต นิจฺฉาตา}

\vspace{\tipitakaparskip}
\csromancenter{นวกนิปาตปาฬิ นิพฺพุตา ติณฺณา ปารงฺคตา ตทงฺเคนาติ วทามิ. “กตฺถ วิญฺญาณญฺจายตนสญฺญา นิรุชฺฌติ, เก จ วิญฺญาณญฺจายตนสญฺญํ นิโรเธตฺวา นิโรเธตฺวา วิหรนฺติ, อหเมตํ น ชานามิ, อหเมตํ น ปสฺสามี”ติ อิติ โย เอวํ วเทยฺย, โส เอวมสฺส วจนีโย “อิธาวุโส ภิกฺขุ สพฺพโส วิญฺญาณญฺจายตนํ สมติกฺกมฺม ‘นตฺถิ กิญฺจี’ติ อากิญฺจญฺญายตนํ อุปสมฺปชฺช วิหรติ, เอตฺถ วิญฺญาณญฺจายตนสญฺญา นิรุชฺฌติ, เต จ วิญฺญาณญฺจายตนสญฺญํ นิโรเธตฺวา นิโรเธตฺวา วิหรนฺตี”ติ.\csromanspacer{} อทฺธา ภิกฺขเว อสโฐ อมายาวี “สาธู”ติ ภาสิตํ อภินนฺเทยฺย อนุโมเทยฺย, “สาธู”ติ ภาสิตํ อภินนฺทิตฺวา อนุโมทิตฺวา นมสฺสมาโน ปญฺชลิโก ปยิรุปาเสยฺย.}
\prose{\csromanfolio{212}ยตฺถ อากิญฺจญฺญายตนสญฺญา นิรุชฺฌติ, เย จ อากิญฺจญฺญายตนสญฺญํ นิโรเธตฺวา นิโรเธตฺวา วิหรนฺติ, อทฺธา เต อายสฺมนฺโต นิจฺฉาตา นิพฺพุตา ติณฺณา ปารงฺคตา ตทงฺเคนาติ วทามิ. “กตฺถ อากิญฺจญฺญายตนสญฺญา นิรุชฺฌติ, เก จ อากิญฺจญฺญายตนสญฺญํ นิโรเธตฺวา นิโรเธตฺวา วิหรนฺติ, อหเมตํ น ชานามิ.\csromanspacer{} อหเมตํ น ปสฺสามี”ติ อิติ โย เอวํ วเทยฺย, โส เอวมสฺส วจนีโย “อิธาวุโส ภิกฺขุ สพฺพโส อากิญฺจญฺญายตนํ สมติกฺกมฺม เนวสญฺญานาสญฺญายตนํ อุปสมฺปชฺช วิหรติ, เอตฺถ อากิญฺจญฺญายตนสญฺญา นิรุชฺฌติ, เต จ อากิญฺจญฺญายตนสญฺญํ นิโรเธตฺวา นิโรเธตฺวา วิหรนฺตี”ติ.\csromanspacer{} อทฺธา ภิกฺขเว อสโฐ อมายาวี “สาธู”ติ ภาสิตํ อภินนฺเทยฺย อนุโมเทยฺย, “สาธู”ติ ภาสิตํ อภินนฺทิตฺวา อนุโมทิตฺวา นมสฺสมาโน ปญฺชลิโก ปยิรุปาเสยฺย.}

\prose{ยตฺถ เนวสญฺญานาสญฺญายตนสญฺญา นิรุชฺฌติ, เย จ เนวสญฺญานาสญฺญายตนสญฺญํ นิโรเธตฺวา นิโรเธตฺวา วิหรนฺติ, อทฺธา เต อายสฺมนฺโต นิจฺฉาตา นิพฺพุตา ติณฺณา ปารงฺคตา ตทงฺเคนาติ วทามิ. “กตฺถ เนวสญฺญานาสญฺญายตนสญฺญา นิรุชฺฌติ, เก จ เนวสญฺญานาสญฺญายตน- สญฺญํ นิโรเธตฺวา นิโรเธตฺวา วิหรนฺติ, อหเมตํ น ชานามิ, อหเมตํ น ปสฺสามี”ติ อิติ โย เอวํ วเทยฺย, โส เอวมสฺส วจนีโย “อิธาวุโส ภิกฺขุ สพฺพโส เนวสญฺญานาสญฺญายตนํ สมติกฺกมฺม สญฺญาเวทยิตนิโรธํ อุปสมฺปชฺช วิหรติ, เอตฺถ เนวสญฺญานาสญฺญายตนสญฺญา นิรุชฺฌติ, เต จ เนวสญฺญานาสญฺญายตนสญฺญํ}

\chapterheadpageread{4. มหาวคฺค}
\prose{\csromanfolio{213}นิโรเธตฺวา นิโรเธตฺวา วิหรนฺตี”ติ.\csromanspacer{} อทฺธา ภิกฺขเว อสโฐ อมายาวี “สาธู”ติ ภาสิตํ อภินนฺเทยฺย อนุโมเทยฺย, “สาธู”ติ ภาสิตํ อภินนฺทิตฺวา อนุโมทิตฺวา นมสฺสมาโน ปญฺชลิโก ปยิรุปาเสยฺย.\csromanspacer{} อิมา โข ภิกฺขเว นว อนุปุพฺพวิหารสมาปตฺติโยติ.\csromanspacer{}\csromanspacer{}ทุติยํ.}

\csromanheader{2}{3. นิพฺพานสุขสุตฺต}
\proseitem{34}{เอกํ สมยํ อายสฺมา สาริปุตฺโต ราชคเห วิหรติ เวฬุวเน กลนฺทกนิวาเป, ตตฺร โข อายสฺมา สาริปุตฺโต ภิกฺขู อามนฺเตสิ “สุขมิทํ อาวุโส นิพฺพานํ, สุขมิทํ อาวุโส นิพฺพานนฺ”ติ.\csromanspacer{} เอวํ วุตฺเต อายสฺมา อุทายี อายสฺมนฺตํ สาริปุตฺตํ เอตทโวจ “กิํ ปเนตฺถ อาวุโส สาริปุตฺต สุขํ, ยเทตฺถ นตฺถิ เวทยิตนฺ”ติ.\csromanspacer{} เอตเทว เขฺวตฺถ อาวุโส สุขํ, ยเทตฺถ นตฺถิ เวทยิตํ.\csromanspacer{} ปญฺจิเม อาวุโส กามคุณา, กตเม ปญฺจ? จกฺขุวิญฺเญยฺยา รูปา อิฏฺฐา กนฺตา มนาปา ปิยรูปา กามูปสํหิตา รชนียา.\csromanspacer{} โสตวิญฺเญยฺยา สทฺทา ฯเปฯ.\csromanspacer{} ฆานวิญฺเญยฺยา คนฺธา.\csromanspacer{} ชิวฺหาวิญฺเญยฺยา รสา.\csromanspacer{} กายวิญฺเญยฺยา โผฏฺฐพฺพา อิฏฺฐา กนฺตา มนาปา ปิยรูปา กามูปสํหิตา รชนียา.\csromanspacer{} อิเม โข อาวุโส ปญฺจ กามคุณา, ยํ โข อาวุโส อิเม ปญฺจ กามคุเณ ปฏิจฺจ อุปฺปชฺชติ สุขํ โสมนสฺสํ.\csromanspacer{} อิทํ วุจฺจตาวุโส กามสุขํ.}

\prose{อิธาวุโส ภิกฺขุ วิวิจฺเจว กาเมหิ ฯเปฯ ปฐมํ ฌานํ อุปสมฺปชฺช วิหรติ.\csromanspacer{} ตสฺส เจ อาวุโส ภิกฺขุโน อิมินา วิหาเรน วิหรโต กามสหคตา สญฺญามนสิการา สมุทาจรนฺติ, สฺวสฺส โหติ อาพาโธ.\csromanspacer{} เสยฺยถาปิ อาวุโส สุขิโน ทุกฺขํ อุปฺปชฺเชยฺย ยาวเทว อาพาธาย.\csromanspacer{} เอวเมวสฺส เต กามสหคตา สญฺญามนสิการา สมุทาจรนฺติ, สฺวสฺส โหติ อาพาโธ.\csromanspacer{} โย โข ปนาวุโส อาพาโธ, ทุกฺขเมตํ วุตฺตํ ภควตา.\csromanspacer{} อิมินาปิ โข เอตํ อาวุโส ปริยาเยน เวทิตพฺพํ ยถา สุขํ นิพฺพานํ.}

\prose{ปุน จปรํ อาวุโส ภิกฺขุ วิตกฺกวิจารานํ วูปสมา ฯเปฯ ทุติยํ ฌานํ อุปสมฺปชฺช วิหารติ.\csromanspacer{} ตสฺส เจ อาวุโส ภิกฺขุโน อิมินา วิหาเรน วิหรโต วิตกฺกสหคตา สญฺญามนสิการา สมุทาจรนฺติ, สฺวสฺส โหติ อาพาโธ.\csromanspacer{} เสยฺยถาปิ อาวุโส สุขิโน ทุกฺขํ อุปฺปชฺเชยฺย ยาวเทว อาพาธาย.\csromanspacer{} เอวเมวสฺส เต วิตกฺกสหคตา สญฺญามนสิการา}

\vspace{\tipitakaparskip}
\csromancenter{นวกนิปาตปาฬิ สมุทาจรนฺติ, สฺวสฺส โหติ อาพาโธ.\csromanspacer{} โย โข ปนาวุโส อาพาโธ, ทุกฺขเมตํ วุตฺตํ ภควตา.\csromanspacer{} อิมินาปิ โข เอตํ อาวุโส ปริยาเยน เวทิตพฺพํ ยถา สุขํ นิพฺพานํ.}
\prose{\csromanfolio{214}ปุน จปรํ อาวุโส ภิกฺขุ ปีติยา จ วิราคา ฯเปฯ ตติยํ ฌานํ อุปสมฺปชฺช วิหรติ.\csromanspacer{} ตสฺส เจ อาวุโส ภิกฺขุโน อิมินา วิหาเรน วิหรโต ปีติสหคตา สญฺญามนสิการา สมุทาจรนฺติ, สฺวสฺส โหติ อาพาโธ.\csromanspacer{} เสยฺยถาปิ อาวุโส สุขิโน ทุกฺขํ อุปฺปชฺเชยฺย ยาวเทว อาพาธาย.\csromanspacer{} เอวเมวสฺส เต ปีติสหคตา สญฺญามนสิการา สมุทาจรนฺติ, สฺวสฺส โหติ อาพาโธ.\csromanspacer{} โย โข ปนาวุโส อาพาโธ, ทุกฺขเมตํ วุตฺตํ ภควตา.\csromanspacer{} อิมินาปิ โข เอตํ อาวุโส ปริยาเยน เวทิตพฺพํ ยถา สุขํ นิพฺพานํ.}

\prose{ปุน จปรํ อาวุโส ภิกฺขุ สุขสฺส จ ปหานา ฯเปฯ จตุตฺถํ ฌานํ อุปสมฺปชฺช วิหรติ.\csromanspacer{} ตสฺส เจ อาวุโส ภิกฺขุโน อิมินา วิหาเรน วิหรโต อุเปกฺขาสหคตา สญฺญามนสิการา สมุทาจรนฺติ, สฺวสฺส โหติ อาพาโธ.\csromanspacer{} เสยฺยถาปิ อาวุโส สุขิโน ทุกฺขํ อุปฺปชฺเชยฺย ยาวเทว อาพาธาย.\csromanspacer{} เอวเมวสฺส เต อุเปกฺขาสหคตา สญฺญามนสิการา สมุทาจรนฺติ, สฺวสฺส โหติ อาพาโธ.\csromanspacer{} โย โข ปนาวุโส อาพาโธ, ทุกฺขเมตํ วุตฺตํ ภควตา.\csromanspacer{} อิมินาปิ โข เอตํ อาวุโส ปริยาเยน เวทิตพฺพํ ยถา สุขํ นิพฺพานํ.}

\prose{ปุน จปรํ อาวุโส ภิกฺขุ สพฺพโส รูปสญฺญานํ สมติกฺกมา ปฏิฆสญฺญานํ อตฺถงฺคมา นานตฺตสญฺญานํ อมนสิการา “อนนฺโต อากาโส”ติ อากาสานญฺจายตนํ อุปสมฺปชฺช วิหรติ.\csromanspacer{} ตสฺส เจ อาวุโส ภิกฺขุโน อิมินา วิหาเรน วิหรโต รูปสหคตา สญฺญามนสิการา สมุทาจรนฺติ, สฺวสฺส โหติ อาพาโธ.\csromanspacer{} เสยฺยถาปิ อาวุโส สุขิโน ทุกฺขํ อุปฺปชฺเชยฺย ยาวเทว อาพาธาย.\csromanspacer{} เอวเมวสฺส เต รูปสหคตา สญฺญามนสิการา สมุทาจรนฺติ, สฺวสฺส โหติ อาพาโธ.\csromanspacer{} โย โข ปนาวุโส อาพาโธ, ทุกฺขเมตํ วุตฺตํ ภควตา.\csromanspacer{} อิมินาปิ โข เอตํ อาวุโส ปริยาเยน เวทิตพฺพํ ยถา สุขํ นิพฺพานํ.}

\prose{ปุน จปรํ อาวุโส ภิกฺขุ สพฺพโส อากาสานญฺจายตนํ สมติกฺกมฺม “อนนฺตํ วิญฺญาณนฺ”ติ วิญฺญาณญฺจายตนํ อุปสมฺปชฺช วิหรติ.\csromanspacer{} ตสฺส}

\chapterheadpageread{4. มหาวคฺค}
\prose{\csromanfolio{215}เจ อาวุโส ภิกฺขุโน อิมินา วิหาเรน วิหรโต อากาสานญฺจายตนสหคตา สญฺญามนสิการา สมุทาจรนฺติ, สฺวสฺส โหติ อาพาโธ.\csromanspacer{} เสยฺยถาปิ อาวุโส สุขิโน ทุกฺขํ อุปฺปชฺเชยฺย ยาวเทว อาพาธาย.\csromanspacer{} เอวเมวสฺส เต อากาสานญฺจายตนสหคตา สญฺญามนสิการา สมุทาจรนฺติ, สฺวสฺส โหติ อาพาโธ.\csromanspacer{} โย โข ปนาวุโส อาพาโธ, ทุกฺขเมตํ วุตฺตํ ภควตา.\csromanspacer{} อิมินาปิ โข เอตํ อาวุโส ปริยาเยน เวทิตพฺพํ ยถา สุขํ นิพฺพานํ.}

\prose{ปุน จปรํ อาวุโส ภิกฺขุ สพฺพโส วิญฺญาณญฺจายตนํ สมติกฺกมฺม “นตฺถิ กิญฺจี”ติ อากิญฺจญฺญายตนํ อุปสมฺปชฺช วิหรติ.\csromanspacer{} ตสฺส เจ อาวุโส ภิกฺขุโน อิมินา วิหาเรน วิหรโต วิญฺญาณญฺจายตนสหคตา สญฺญามนสิการา สมุทาจรนฺติ, สฺวสฺส โหติ อาพาโธ.\csromanspacer{} เสยฺยถาปิ อาวุโส สุขิโน ทุกฺขํ อุปฺปชฺเชยฺย ยาวเทว อาพาธาย.\csromanspacer{} เอวเมวสฺส เต วิญฺญาณญฺจายตนสหคตา สญฺญามนสิการา สมุทาจรนฺติ, สฺวสฺส โหติ อาพาโธ.\csromanspacer{} โย โข ปนาวุโส อาพาโธ, ทุกฺขเมตํ วุตฺตํ ภควตา.\csromanspacer{} อิมินาปิ โข เอตํ อาวุโส ปริยาเยน เวทิตพฺพํ ยถา สุขํ นิพฺพานํ.}

\prose{ปุน จปรํ อาวุโส ภิกฺขุ สพฺพโส อากิญฺจญฺญายตนํ สมติกฺกมฺม เนวสญฺญานาสญฺญายตนํ อุปสมฺปชฺช วิหรติ.\csromanspacer{} ตสฺส เจ อาวุโส ภิกฺขุโน อิมินา วิหาเรน วิหรโต อากิญฺจญฺญายตนสหคตา สญฺญามนสิการา สมุทาจรนฺติ, สฺวสฺส โหติ อาพาโธ.\csromanspacer{} เสยฺยถาปิ อาวุโส สุขิโน ทุกฺขํ อุปฺปชฺเชยฺย ยาวเทว อาพาธาย.\csromanspacer{} เอวเมวสฺส เต อากิญฺจญฺญายตนสหคตา สญฺญามนสิการา สมุทาจรนฺติ, สฺวสฺส โหติ อาพาโธ.\csromanspacer{} โย โข ปนาวุโส อาพาโธ, ทุกฺขเมตํ วุตฺตํ ภควตา.\csromanspacer{} อิมินาปิ โข เอตํ อาวุโส ปริยาเยน เวทิตพฺพํ ยถา สุขํ นิพฺพานํ.}

\prose{ปุน จปรํ อาวุโส ภิกฺขุ สพฺพโส เนวสญฺญานาสญฺญายตนํ สมติกฺกมฺม สญฺญาเวทยิตนิโรธํ อุปสมฺปชฺช วิหรติ, ปญฺญาย จสฺส ทิสฺวา อาสวา ปริกฺขีณา โหนฺติ.\csromanspacer{} อิมินาปิ โข เอตํ อาวุโส ปริยาเยน เวทิตพฺพํ ยถา สุขํ นิพฺพานนฺติ.\csromanspacer{}\csromanspacer{}ตติยํ.}

\gambhira{นวกนิปาตปาฬิ \textbf{4. คาวี-อุปมาสุตฺต}}
\proseitem{35}{\csromanfolio{216}เสยฺยถาปิ ภิกฺขเว คาวี ปพฺพเตยฺยา พาลา อพฺยตฺตา อเขตฺตญฺญู อกุสลา วิสเม ปพฺพเต จริตุํ.\csromanspacer{} ตสฺสา เอวมสฺส “ยํนูนาหํ อคตปุพฺพญฺเจว ทิสํ คจฺเฉยฺยํ, อขาทิตปุพฺพานิ จ ติณานิ ขาเทยฺยํ, อปีตปุพฺพานิ จ ปานียานิ ปิเวยฺยนฺ”ติ.\csromanspacer{} สา ปุริมํ ปาทํ น สุปฺปติฏฺฐิตํ ปติฏฺฐาเปตฺวา ปจฺฉิมํ ปาทํ อุทฺธเรยฺย, สา น เจว อคตปุพฺพํ ทิสํ คจฺเฉยฺย, น จ อขาทิตปุพฺพานิ ติณานิ ขาเทยฺย, น จ อปีตปุพฺพานิ ปานียานิ ปิเวยฺย.\csromanspacer{} ยสฺมิํ จสฺสา ปเทสา ฐิตาย เอวมสฺส “ยํนูนาหํ อคตปุพฺพญฺเจว ทิสํ คจฺเฉยฺยํ, อขาทิตปุพฺพานิ จ ติณานิ ขาเทยฺยํ, อปีตปุพฺพานิ จ ปานียานิ ปิเวยฺยนฺ”ติ, ตญฺจ ปเทสํ น โสตฺถินา ปจฺจาคจฺเฉยฺย.\csromanspacer{} ตํ กิสฺส เหตุ? ตถา หิ สา ภิกฺขเว คาวี ปพฺพเตยฺยา พาลา อพฺยตฺตา อเขตฺตญฺญู อกุสลา วิสเม ปพฺพเต จริตุํ.\csromanspacer{} เอวเมวํ โข ภิกฺขเว อิเธกจฺโจ ภิกฺขุ พาโล อพฺยตฺโต อเขตฺตญฺญู อกุสโล วิวิจฺเจว กาเมหิ วิวิจฺจ อกุสเลหิ ธมฺเมหิ สวิตกฺกํ สวิจารํ วิเวกชํ ปีติสุขํ ปฐมํ ฌานํ อุปสมฺปชฺช วิหรติ.\csromanspacer{} โส ตํ นิมิตฺตํ น อาเสวติ น ภาเวติ น พหุลีกโรติ, น สฺวาธิฏฺฐิตํ อธิฏฺฐาติ.}

\prose{ตสฺส เอวํ โหติ “ยํนูนาหํ วิตกฺกวิจารานํ วูปสมา อชฺฌตฺตํ สมฺปสาทนํ เจตโส เอโกทิภาวํ อวิตกฺกํ อวิจารํ สมาธิชํ ปีติสุขํ ทุติยํ ฌานํ อุปสมฺปชฺช วิหเรยฺยนฺ”ติ.\csromanspacer{} โส น สกฺโกติ วิตกฺกวิจารานํ วูปสมา ฯเปฯ ทุติยํ ฌานํ อุปสมฺปชฺช วิหริตุํ.\csromanspacer{} ตสฺส เอวํ โหติ “ยํนูนาหํ วิวิจฺเจว กาเมหิ วิวิจฺจ อกุสเลหิ ธมฺเมหิ สวิตกฺกํ สวิจารํ วิเวกชํ ปีติสุขํ ปฐมํ ฌานํ อุปสมฺปชฺช วิหเรยฺยนฺ”ติ.\csromanspacer{} โส น สกฺโกติ วิวิจฺเจว กาเมหิ ฯเปฯ ปฐม ฌานํ อุปสมฺปชฺช วิหเรตุํ.\csromanspacer{} อยํ วุจฺจติ ภิกฺขเว ภิกฺขุ อุภโต ภฏฺโฐ อุภโต ปริหีโน.\csromanspacer{} เสยฺยถาปิ สา คาวี ปพฺพเตยฺยา พาลา อพฺยตฺตา อเขตฺตญฺญู อกุสลา วิสเม ปพฺพเต จริตุํ.}

\prose{เสยฺยถาปิ ภิกฺขเว คาวี ปพฺพเตยฺยา ปณฺฑิตา พฺยตฺตา เขตฺตญฺญู กุสลา วิสเม ปพฺพเต จริตุํ.\csromanspacer{} ตสฺสา เอวมสฺส “ยํนูนาหํ อคตปุพฺพญฺเจว ทิสํ คจฺเฉยฺยํ, อขาทิตปุพฺพานิ จ ติณานิ ขาเทยฺยํ, อปีตปุพฺพานิ จ ปานียานิ ปิเวยฺยนฺ”ติ.\csromanspacer{} สา ปุริมํ ปาทํ สุปฺปติฏฺฐิตํ ปติฏฺฐาเปตฺวา ปจฺฉิมํ ปาทํ อุทฺธเรยฺย, สา อคตปุพฺพญฺเจว ทิสํ คจฺเฉยฺย, อขาทิตปุพฺพานิ จ}

\chapterheadpageread{4. มหาวคฺค}
\prose{\csromanfolio{217}ติณานิ ขาเทยฺย, อปีตปุพฺพานิ จ ปานียานิ ปิเวยฺย.\csromanspacer{} ยสฺมิํ จสฺสา ปเทเส ฐิตาย เอวมสฺส “ยํนูนาหํ อคตปุพฺพญฺเจว ทิสํ คจฺเฉยฺยํ, อขาทิตปุพฺพานิ จ ติณานิ ขาเทยฺยํ, อปีตปุพฺพานิ จ ปานียานิ ปิเวยฺยนฺ”ติ, ตญฺจ ปเทสํ โสตฺถินา ปจฺจาคจฺเฉยฺย.\csromanspacer{} ตํ กิสฺส เหตุ? ตถา หิ สา ภิกฺขเว คาวี ปพฺพเตยฺยา ปณฺฑิตา พฺยตฺตา เขตฺตญฺญู กุสลา วิสเม ปพฺพเต จริตุํ.\csromanspacer{} เอวเมวํ โข ภิกฺขเว อิเธกจฺโจ ภิกฺขุ ปณฺฑิโต พฺยตฺโต เขตฺตญฺญู กุสโล วิวิจฺเจว กาเมหิ วิวิจฺจ อกุสเลหิ ธมฺเมหิ สวิตกฺกํ สวิจารํ วิเวกชํ ปีติสุขํ ปฐมํ ฌานํ อุปสมฺปชฺช วิหรติ.\csromanspacer{} โส ตํ นิมิตฺตํ อาเสวติ ภาเวติ พหุลีกโรติ, สฺวาธิฏฺฐิตํ อธิฏฺฐาติ.}

\prose{ตสฺส เอวํ โหติ “ยํนูนาหํ วิตกฺกวิจารานํ วูปสมา อชฺฌตฺตํ สมฺปสาทนํ เจตโส เอโกทิภาวํ อวิตกฺกํ อวิจารํ สมาธิชํ ปีติสุขํ ทุติยํ ฌานํ อุปสมฺปชฺช วิหเรยฺยนฺ”ติ.\csromanspacer{} โส ทุติยํ ฌานํ อนภิหิํสมาโน วิตกฺกวิจารานํ วูปสมา ทุติยํ ฌานํ อุปสมฺปชฺช วิหรติ.\csromanspacer{} โส ตํ นิมิตฺตํ อาเสวติ ภาเวติ พหุลีกโรติ, สฺวาธิฏฺฐิตํ อธิฏฺฐาติ.}

\prose{ตสฺส เอวํ โหติ “ยํนูนาหํ ปีติยา จ วิราคา อุเปกฺขโก จ วิหเรยฺยํ, สโต จ สมฺปชาโน สุขญฺจ กาเยน ปฏิสํเวเทยฺยํ, ยํ ตํ อริยา อาจิกฺขนฺติ ‘อุเปกฺขโก สติมา สุขวิหารี’ติ ตติยํ ฌานํ อุปสมฺปชฺช วิหเรยฺยนฺ”ติ.\csromanspacer{} โส ตติยํ ฌานํ อนภิหิํสมาโน ปีติยา จ วิราคา ตติยํ ฌานํ อุปสมฺปชฺช วิหรติ.\csromanspacer{} โส ตํ นิมิตฺตํ อาเสวติ ภาเวติ พหุลีกโรติ, สฺวาธิฏฺฐิตํ อธิฏฺฐาติ.}

\prose{ตสฺส เอวํ โหติ “ยํนูนาหํ สุขสฺส จ ปหานา ทุกฺขสฺส จ ปหานา ปุพฺเพว โสมนสฺสโทมนสฺสานํ อตฺถงฺคมา อทุกฺขมสุขํ อุเปกฺขาสติปาริสุทฺธิํ จตุตฺถํ ฌานํ อุปสมฺปชฺช วิหเรยฺยนฺ”ติ.\csromanspacer{} โส จตุตฺถํ ฌานํ อนภิหิํสมาโน สุขสฺส จ ปหานา ฯเปฯ จตุตฺถํ ฌานํ อุปสมฺปชฺช วิหรติ.\csromanspacer{} โส ตํ นิมิตฺตํ อาเสวติ ภาเวติ พหุลีกโรติ, สฺวาธิฏฺฐิตํ อธิฏฺฐาติ.}

\prose{ตสฺส เอวํ โหติ “ยํนูนาหํ สพฺพโส รูปสญฺญานํ สมติกฺกมา ปฏิฆสญฺญานํ อตฺถงฺคมา นานตฺตสญฺญานํ อมนสิการา ‘อนนฺโต อากาโส’ติ อากาสานญฺจายตนํ อุปสมฺปชฺช วิหเรยฺยนฺ”ติ.\csromanspacer{} โส}

\vspace{\tipitakaparskip}
\csromancenter{นวกนิปาตปาฬิ อากาสานญฺจายตนํ อนภิหิํสมาโน สพฺพโส รูปสญฺญานํ สมติกฺกมา ฯเปฯ อากาสานญฺจายตนํ อุปสมฺปชฺช วิหรติ.\csromanspacer{} โส ตํ นิมิตฺตํ อาเสวติ ภาเวติ พหุลีกโรติ, สฺวาธิฏฺฐิตํ อธิฏฺฐาติ.}
\prose{\csromanfolio{218}ตสฺส เอวํ โหติ “ยํนูนาหํ สพฺพโส อากาสานญฺจายตนํ สมติกฺกมฺม ‘อนนฺตํ วิญฺญาณนฺ’ติ วิญฺญาณญฺจายตนํ อุปสมฺปชฺช วิหเรยฺยนฺ”ติ.\csromanspacer{} โส วิญฺญาณญฺจายตนํ อนภิหิํสมาโน สพฺพโส อากาสานญฺจายตนํ สมติกฺกมฺม “อนนฺตํ วิญฺญาณนฺ”ติ วิญฺญาณญฺจายตนํ อุปสมฺปชฺช วิหรติ.\csromanspacer{} โส ตํ นิมิตฺตํ อาเสวติ ภาเวติ พหุลีกโรติ, สฺวาธิฏฺฐิตํ อธิฏฺฐาติ.}

\prose{ตสฺส เอวํ โหติ “ยํนูนาหํ สพฺพโส วิญฺญาณญฺจายตนํ สมติกฺกมฺม ‘นตฺถิ กิญฺจี’ติ อากิญฺจญฺญายตนํ อุปสมฺปชฺช วิหเรยฺยนฺ”ติ.\csromanspacer{} โส อากิญฺจญฺญายตนํ อนภิหิํสมาโน สพฺพโส วิญฺญาณญฺจายตนํ สมติกฺกมฺม “นตฺถิ กิญฺจี”ติ อากิญฺจญฺญายตนํ อุปสมฺปชฺช วิหรติ.\csromanspacer{} โส ตํ นิมิตฺตํ อาเสวติ ภาเวติ พหุลีกโรติ, สฺวาธิฏฺฐิตํ อธิฏฺฐาติ.}

\prose{ตสฺส เอวํ โหติ “ยํนูนาหํ สพฺพโส อากิญฺจญฺญายตนํ สมติกฺกมฺม เนวสญฺญานาสญฺญายตนํ อุปสมฺปชฺช วิหเรยฺยนฺ”ติ.\csromanspacer{} โส เนวสญฺญานาสญฺญายตนํ อนภิหิํสมาโน สพฺพโส อากิญฺจญฺญายตนํ สมติกฺกมฺม เนวสญฺญานาสญฺญายตนํ อุปสมฺปชฺช วิหรติ, โส ตํ นิมิตฺตํ อาเสวติ ภาเวติ พหุลีกโรติ, สฺวาธิฏฺฐิตํ อธิฏฺฐาติ.}

\prose{ตสฺส เอวํ โหติ “ยํนูนาหํ สพฺพโส เนวสญฺญานาสญฺญายตนํ สมติกฺกมฺม สญฺญาเวทยิตนิโรธํ อุปสมฺปชฺช วิหเรยฺยนฺ”ติ.\csromanspacer{} โส สญฺญาเวทยิตนิโรธํ อนภิหิํสมาโน สพฺพโส เนวสญฺญานาสญฺญายตนํ สมติกฺกมฺม สญฺญาเวทยิตนิโรธํ อุปสมฺปชฺช วิหรติ.}

\prose{ยโต โข ภิกฺขเว ภิกฺขุ ตํ ตเทว สมาปตฺติํ สมาปชฺชติปิ วุฏฺฐาติปิ, ตสฺส มุทุ จิตฺตํ โหติ กมฺมญฺญํ, มุทุนา กมฺมญฺเญน จิตฺเตน อปฺปมาโณ สมาธิ โหติ สุภาวิโต.\csromanspacer{} โส อปฺปมาเณน สมาธินา สุภาวิเตน ยสฺส ยสฺส อภิญฺญาสจฺฉิกรณียสฺส ธมฺมสฺส}

\chapterheadpageread{4. มหาวคฺค}
\prose{\csromanfolio{219}จิตฺตํ อภินินฺนาเมติ อภิญฺญาสจฺฉิกิริยาย.\csromanspacer{} ตตฺร ตเตฺรว สกฺขิภพฺพตํ ปาปุณาติ สติ สติ อายตเน.}

\prose{โส สเจ อากงฺขติ “อเนกวิหิตํ อิทฺธิวิธํ ปจฺจนุภเวยฺยํ.\csromanspacer{} เอโกปิ หุตฺวา พหุธา อสฺสํ, พหุธาปิ หุตฺวา เอโก อสฺสํ ฯเปฯ ยาว พฺรหฺมโลกาปิ กาเยน วสํ วตฺเตยฺยนฺ”ติ.\csromanspacer{} ตตฺร ตเตฺรว สกฺขิภพฺพตํ ปาปุณาติ สติ สติ อายตเน.}

\prose{โส สเจ อากงฺขติ ทิพฺพาย โสตธาตุยา ฯเปฯ สติ สติ อายตเน.}

\prose{โส สเจ อากงฺขติ “ปรสตฺตานํ ปรปุคฺคลานํ เจตสา เจโต ปริจฺจ ปชาเนยฺยํ.\csromanspacer{} สราคํ วา จิตฺตํ “สราคํ จิตฺตนฺ”ติ ปชาเนยฺยํ.\csromanspacer{} วีตราคํ วา จิตฺตํ “วีตราคํ จิตฺตนฺ”ติ ปชาเนยฺยํ.\csromanspacer{} สโทสํ วา จิตฺตํ “สโทสํ จิตฺตนฺ”ติ ปชาเนยฺยํ.\csromanspacer{} วีตโทสํ วา จิตฺตํ “วีตโทสํ จิตฺตนฺ”ติ ปชาเนยฺยํ.\csromanspacer{} สโมหํ วา จิตฺตํ “สโมหํ จิตฺตนฺ”ติ ปชาเนยฺยํ.\csromanspacer{} วีตโมหํ วา จิตฺตํ.\csromanspacer{} สํขิตฺตํ วา จิตฺตํ.\csromanspacer{} วิกฺขิตฺตํ วา จิตฺตํ.\csromanspacer{} มหคฺคตํ วา จิตฺตํ.\csromanspacer{} อมหคฺคตํ วา จิตฺตํ.\csromanspacer{} ส-อุตฺตรํ วา จิตฺตํ.\csromanspacer{} อนุตฺตรํ วา จิตฺตํ.\csromanspacer{} สมาหิตํ วา จิตฺตํ.\csromanspacer{} อสมาหิตํ วา จิตฺตํ.\csromanspacer{} วิมุตฺตํ วา จิตฺตํ.\csromanspacer{} อวิมุตฺตํ วา จิตฺตํ “อวิมุตฺตํ จิตฺตนฺ”ติ ปชาเนยฺยนฺ”ติ.\csromanspacer{} ตตฺร ตเตฺรว สกฺขิภพฺพตํ ปาปุณาติ สติ สติ อายตเน.}

\prose{โส สเจ อากงฺขติ “อเนกวิหิตํ ปุพฺเพนิวาสํ อนุสฺสเรยฺยํ.\csromanspacer{} เสยฺยถิทํ, เอกมฺปิ ชาติํ เทฺวปิ ชาติโย ฯเปฯ อิติ สาการํ ส-อุทฺเทสํ อเนกวิหิตํ ปุพฺเพนิวาสํ อนุสฺสเรยฺยนฺ”ติ.\csromanspacer{} ตตฺร ตเตฺรว สกฺขิภพฺพตํ ปาปุณาติ สติ สติ อายตเน.}

\prose{โส สเจ อากงฺขติ “ทิพฺเพน จกฺขุนา วิสุทฺเธน อติกฺกนฺตมานุสเกน ฯเปฯ ยถากมฺมูปเค สตฺเต ปชาเนยฺยนฺ”ติ.\csromanspacer{} ตตฺร ตเตฺรว สกฺขิภพฺพตํ ปาปุณาติ สติ สติ อายตเน.}

\prose{โส สเจ อากงฺขติ “อาสวานํ ขยา อนาสวํ เจโตวิมุตฺติํ ปญฺญาวิมุตฺติํ ทิฏฺเฐว ธมฺเม สยํ อภิญฺญา สจฺฉิกตฺวา อุปสมฺปชฺช วิหเรยฺยนฺ”ติ.\csromanspacer{} ตตฺร ตเตฺรว สกฺขิภพฺพตํ ปาปุณาติ สติ สติ อายตเนติ.\csromanspacer{}\csromanspacer{}จตุตฺถํ.}

\gambhira{นวกนิปาตปาฬิ \textbf{5. ฌานสุตฺต}}
\proseitem{36}{\csromanfolio{220}ปฐมมฺปาหํ ภิกฺขเว ฌานํ นิสฺสาย อาสวานํ ขยํ วทามิ.\csromanspacer{} ทุติยมฺปาหํ ภิกฺขเว ฌานํ นิสฺสาย อาสวานํ ขยํ วทามิ.\csromanspacer{} ตติยมฺปาหํ ภิกฺขเว ฌานํ นิสฺสาย อาสวานํ ขยํ วทามิ.\csromanspacer{} จตุตฺถมฺปาหํ ภิกฺขเว ฌานํ นิสฺสาย อาสวานํ ขยํ วทามิ.\csromanspacer{} อากาสานญฺจายตนมฺปาหํ ภิกฺขเว นิสฺสาย อาสวานํ ขยํ วทามิ.\csromanspacer{} วิญฺญาณญฺจายตนมฺปาหํ ภิกฺขเว นิสฺสาย อาสวานํ ขยํ วทามิ.\csromanspacer{} อากิญฺจญฺญายตนมฺปาหํ ภิกฺขเว นิสฺสาย อาสวานํ ขยํ วทามิ.\csromanspacer{} เนวสญฺญานาสญฺญายตนมฺปาหํ ภิกฺขเว นิสฺสาย อาสวานํ ขยํ วทามิ.\csromanspacer{} สญฺญาเวทยิตนิโรธมฺปาหํ ภิกฺขเว นิสฺสาย อาสวานํ ขยํ วทามิ.}

\prose{“ปฐมมฺปาหํ ภิกฺขเว ฌานํ นิสฺสาย อาสวานํ ขยํ วทามี”ติ อิติ โข ปเนตํ วุตฺตํ, กิญฺเจตํ ปฏิจฺจ วุตฺตํ.\csromanspacer{} อิธ ภิกฺขเว ภิกฺขุ วิวิจฺเจว กาเมหิ ฯเปฯ ปฐมํ ฌานํ อุปสมฺปชฺช วิหรติ.\csromanspacer{} โส ยเทว ตตฺถ โหติ รูปคตํ เวทนาคตํ สญฺญาคตํ สงฺขารคตํ วิญฺญาณคตํ, เต ธมฺเม อนิจฺจโต ทุกฺขโต โรคโต คณฺฑโต สลฺลโต อฆโต อาพาธโต ปรโต ปโลกโต สุญฺญโต อนตฺตโต สมนุปสฺสติ.\csromanspacer{} โส เตหิ ธมฺเมหิ จิตฺตํ ปฏิวาเปติ\footnote{ปติฏฺฐาเปติ (สฺยา), ปฏิปาเทติ (ก) ม 2. 99 ปิฏฺเฐปิ ปสฺสิตพฺพํ.}.\csromanspacer{} โส เตหิ ธมฺเมหิ จิตฺตํ ปฏิวาเปตฺวา\footnote{ปติฏฺฐาเปตฺวา (สฺยา), ปฏิปาเทตฺวา (ก) 3. ปทาลิตา (ก) อํ 1. 288, 490 ปิฏฺเฐสุปิ.} อมตาย ธาตุยา จิตฺตํ อุปสํหรติ, “เอตํ สนฺตํ เอตํ ปณีตํ, ยทิทํ สพฺพสงฺขารสมโถ สพฺพูปธิปฏินิสฺสคฺโค ตณฺหากฺขโย วิราโค นิโรโธ นิพฺพานนฺ”ติ.\csromanspacer{} โส ตตฺถ ฐิโต อาสวานํ ขยํ ปาปุณาติ.\csromanspacer{} โน เจ อาสวานํ ขยํ ปาปุณาติ.\csromanspacer{} เตเนว ธมฺมราเคน ตาย ธมฺมนนฺทิยา ปญฺจนฺนํ โอรมฺภาคิยานํ สํโยชนานํ ปริกฺขยา โอปปาติโก โหติ, ตตฺถ ปรินิพฺพายี อนาวตฺติธมฺโม ตสฺมา โลกา.}

\prose{เสยฺยถาปิ ภิกฺขเว อิสฺสาโส วา อิสฺสาสนฺเตวาสี วา ติณปุริสรูปเก วา มตฺติกาปุญฺเช วา โยคฺคํ กริตฺวา โส อปเรน สมเยน ทูเรปาตี จ โหติ อกฺขณเวธี จ มหโต จ กายสฺส ปทาเลตา3.\csromanspacer{} เอวเมวํ โข ภิกฺขเว ภิกฺขุ วิวิจฺเจว กาเมหิ ฯเปฯ ปฐมํ ฌานํ}

\chapterheadpageread{4. มหาวคฺค}
\prose{\csromanfolio{221}อุปสมฺปชฺช วิหรติ.\csromanspacer{} โส ยเทว ตตฺถ โหติ รูปคตํ เวทนาคตํ สญฺญาคตํ สงฺขารคตํ วิญฺญาณคตํ, เต ธมฺเม อนิจฺจโต ทุกฺขโต โรคโต คณฺฑโต สลฺลโต อฆโต อาพาธโต ปรโต ปโลกโต สุญฺญโต อนตฺตโต สมนุปสฺสติ.\csromanspacer{} โส เตหิ ธมฺเมหิ จิตฺตํ ปฏิวาเปติ.\csromanspacer{} โส เตหิ ธมฺเมหิ จิตฺตํ ปฏิวาเปตฺวา อมตาย ธาตุยา จิตฺตํ อุปสํหรติ “เอตํ สนฺตํ เอตํ ปณีตํ, ยทิทํ สพฺพสงฺขารสมโถ สพฺพูปธิปฏินิสฺสคฺโค ตณฺหากฺขโย วิราโค นิโรโธ นิพฺพานนฺ”ติ.\csromanspacer{} โส ตตฺถ ฐิโต อาสวานํ ขยํ ปาปุณาติ.\csromanspacer{} โน จ อาสวานํ ขยํ ปาปุณาติ.\csromanspacer{} เตเนว ธมฺมราเคน ตาย ธมฺมนนฺทิยา ปญฺจนฺนํ โอรมฺภาคิยานํ สํโยชนานํ ปริกฺขยา โอปปาติโก โหติ.\csromanspacer{} ตตฺถ ปรินิพฺพายี อนาวตฺติธมฺโม ตสฺมา โลกา. “ปฐมมฺปาหํ ภิกฺขเว ฌานํ นิสฺสาย อาสวานํ ขยํ วทามี”ติ อิติ ยํ ตํ วุตฺตํ, อิทเมตํ ปฏิจฺจ วุตฺตํ.}

\prose{“ทุติยมฺปาหํ ภิกฺขเว ฌานํ นิสฺสาย ฯเปฯ.\csromanspacer{} ตติยมฺปาหํ ภิกฺขเว ฌานํ นิสฺสาย ฯเปฯ.\csromanspacer{} จตุตฺถมฺปาหํ ภิกฺขเว ฌานํ นิสฺสาย อาสวานํ ขยํ วทามี”ติ อิติ โข ปเนตํ วุตฺตํ, กิญฺเจตํ ปฏิจฺจ วุตฺตํ.\csromanspacer{} อิธ ภิกฺขเว ภิกฺขุ สุขสฺส จ ปหานา ทุกฺขสฺส จ ปหานา ปุพฺเพว โสมนสฺสโทมนสฺสานํ อตฺถงฺคมา อทุกฺขมสุขํ อุเปกฺขาสติปาริสุทฺธิํ จตุตฺถํ ฌานํ อุปสมฺปชฺช วิหรติ.\csromanspacer{} โส ยเทว ตตฺถ โหติ รูปคตํ เวทนาคตํ สญฺญาคตํ สงฺขารคตํ วิญฺญาณคตํ, เต ธมฺเม อนิจฺจโต ทุกฺขโต โรคโต คณฺฑโต สลฺลโต อฆโต อาพาธโต ปรโต ปโลกโต สุญฺญโต อนตฺตโต สมนุปสฺสติ.\csromanspacer{} โส เตหิ ธมฺเมหิ จิตฺตํ ปฏิวาเปติ.\csromanspacer{} โส เตหิ ธมฺเมหิ จิตฺตํ ปฏิวาเปตฺวา อมตาย ธาตุยา จิตฺตํ อุปสํหรติ “เอตํ สนฺตํ เอตํ ปณีตํ, ยทิทํ สพฺพสงฺขารสมโถ สพฺพูปธิปฏินิสฺสคฺโค ตณฺหากฺขโย วิราโค นิโรโธ นิพฺพานนฺ”ติ.\csromanspacer{} โส ตตฺถ ฐิโต อาสวานํ ขยํ ปาปุณาติ.\csromanspacer{} โน เจ อาสวานํ ขยํ ปาปุณาติ.\csromanspacer{} เตนว ธมฺมราเคน ตาย ธมฺมนนฺทิยา ปญฺจนฺนํ โอรมฺภาคิยานํ สํโยชนานํ ปริกฺขยา โอปปาติโก โหติ.\csromanspacer{} ตตฺถ ปรินิพฺพายี อนาวตฺติธมฺโม ตสฺมา โลกา.}

\prose{เสยฺยถาปิ ภิกฺขเว อิสฺสาโส วา อิสฺสาสนฺเตวาสี วา ติณปุริสรูปเก วา มตฺติกาปุญฺเช วา โยคฺคํ กริตฺวา โส อปเรน สมเยน ทูเรปาตี จ โหติ อกฺขณเวธี จ มหโต จ กายสฺส}

\vspace{\tipitakaparskip}
\csromancenter{นวกนิปาตปาฬิ ปทาเลตา.\csromanspacer{} เอวเมวํ โข ภิกฺขเว ภิกฺขุ สุขสฺส จ ปหานา ฯเปฯ จตุตฺถํ ฌานํ อุปสมฺปชฺช วิหรติ.\csromanspacer{} โส ยเทว ตตฺถ โหติ รูปคตํ เวทนาคตํ ฯเปฯ อนาวตฺติธมฺโม ตสฺมา โลกา. “จตุตฺถมฺปาหํ ภิกฺขเว ฌานํ นิสฺสาย อาสวานํ ขยํ วทามี”ติ อิติ ยํ ตํ วุตฺตํ, อิทเมตํ ปฏิจฺจ วุตฺตํ.}
\prose{\csromanfolio{222}“อากาสานญฺจายตนมฺปาหํ ภิกฺขเว ฌานํ นิสฺสาย อาสวานํ ขยํ วทามี”ติ อิติ โข ปเนตํ วุตฺตํ, กิญฺเจตํ ปฏิจฺจ วุตฺตํ.\csromanspacer{} อิธ ภิกฺขเว ภิกฺขุ สพฺพโส รูปสญฺญานํ สมติกฺกมา ปฏิฆสญฺญานํ อตฺถงฺคมา นานตฺตสญฺญานํ อมนสิการา “อนนฺโต อากาโส”ติ อากาสานญฺจายตนํ อุปสมฺปชฺช วิหรติ.\csromanspacer{} โส ยเทว ตตฺถ โหติ เวทนาคตํ สญฺญาคตํ สงฺขารคตํ วิญฺญาณคตํ.\csromanspacer{} เต ธมฺเม อนิจฺจโต ทุกฺขโต โรคโต คณฺฑโต สลฺลโต อฆโต อาพาธโต ปรโต ปโลกโต สุญฺญโต อนตฺตโต สมนุปสฺสติ.\csromanspacer{} โส เตหิ ธมฺเมหิ จิตฺตํ ปฏิวาเปติ โส เตหิ ธมฺเมหิ จิตฺตํ ปฏิวาเปตฺวา อมตาย ธาตุยา จิตฺตํ อุปสํหรติ “เอตํ สนฺตํ เอตํ ปณีตํ, ยทิทํ สพฺพสงฺขารสโถ สพฺพูปธิปฏินิสฺสคฺโค ตณฺหากฺขโย วิราโค นิโรโธ นิพฺพานนฺ”ติ.\csromanspacer{} โส ตตฺถ ฐิโต อาสวานํ ขยํ ปาปุณาติ.\csromanspacer{} โน เจ อาสวานํ ขยํ ปาปุณาติ.\csromanspacer{} เตเนว ธมฺมราเคน ตาย ธมฺมนนฺทิยา ปญฺจนฺนํ โอรมฺภาคิยานํ สํโยชนานํ ปริกฺขยา โอปปาติโก โหติ.\csromanspacer{} ตตฺถ ปรินิพฺพายี อนาวตฺติธมฺโม ตสฺมา โลกา.}

\prose{เสยฺยถาปิ ภิกฺขเว อิสฺสาโส วา อิสฺสาสนฺเตวาสี วา ติณปุริสรูปเก วา มตฺติกาปุญฺเช วา โยคฺคํ กริตฺวา โส อปเรน สมเยน ทูเรปาตี จ โหติ อกฺขณเวธี จ มหโต จ กายสฺส ปทาเลตา.\csromanspacer{} เอวเมวํ โข ภิกฺขเว ภิกฺขุ สพฺพโส รูปสญฺญานํ สมติกฺกมา ปฏิฆสญฺญานํ อตฺถงฺคมา นานตฺตสญฺญานํ อมนสิการา “อนนฺโต อากาโส”ติ อากาสานญฺจายตนํ อุปสมฺปชฺช วิหรติ.\csromanspacer{} โส ยเทว ตตฺถ โหติ เวทนาคตํ สญฺญาคตํ ฯเปฯ อนาวตฺติธมฺโม ตสฺมา โลกา. “อากาสานญฺจายตนมฺปาหํ ภิกฺขเว นิสฺสาย อาสวานํ ขยํ วทามี”ติ อิติ ยํ ตํ วุตฺตํ, อิทเมตํ ปฏิจฺจ วุตฺตํ.}

\prose{“วิญฺญาณญฺจายตนมฺปาหํ ภิกฺขเว นิสฺสาย ฯเปฯ.\csromanspacer{} อากิญฺจญฺญายตนมฺปาหํ ภิกฺขเว นิสฺสาย อาสวานํ ขยํ วทามี”ติ อิติ โข ปเนตํ วุตฺตํ, กิญฺเจตํ ปฏิจฺจ วุตฺตํ, อิธ ภิกฺขเว ภิกฺขุ สพฺพโส วิญฺญาณญฺจายตนํ}

\chapterheadpageread{4. มหาวคฺค}
\prose{\csromanfolio{223}สมติกฺกมฺม “นตฺถิ กิญฺจี”ติ อากิญฺจญฺญายตนํ อุปสมฺปชฺช วิหรติ.\csromanspacer{} โส ยเทว ตตฺถ โหติ เวทนาคตํ สญฺญาคตํ สงฺขารคตํ วิญฺญาณคตํ เต ธมฺเม อนิจฺจโต ทุกฺขโต โรคโต คณฺฑโต สลฺลโต อฆโต อาพาธโต ปรโต ปโลกโต สุญฺญโต อนตฺตโต สมนุปสฺสติ.\csromanspacer{} โส เตหิ ธมฺเมหิ จิตฺตํ ปฏิวาเปติ.\csromanspacer{} โส เตหิ ธมฺเมหิ จิตฺตํ ปฏิวาเปตฺวา อมตาย ธาตุยา จิตฺตํ อุปสํหรติ “เอตํ สนฺตํ เอตํ ปณีตํ, ยทิทํ สพฺพสงฺขารสมโถ สพฺพูปธิปฏินิสฺสคฺโค ตณฺหากฺขโย วิราโค นิโรโธ นิพฺพานนฺ”ติ.\csromanspacer{} โส ตตฺถ ฐิโต อาสวานํ ขยํ ปาปุณาติ.\csromanspacer{} โน เจ อาสวานํ ขยํ ปาปุณาติ.\csromanspacer{} เตเนว ธมฺมราเคน ตาย ธมฺมนนฺทิยา ปญฺจนฺนํ โอรมฺภาคิยานํ สํโยชนานํ ปริกฺขยา โอปปาติโก โหติ.\csromanspacer{} ตตฺถ ปรินิพฺพายี อนาวตฺติธมฺโม ตสฺมา โลกา.}

\prose{เสยฺยถาปิ ภิกฺขเว อิสฺสาโส วา อิสฺสาสนฺเตวาสี วา ติณปุริสรูปเก วา มตฺติกาปุญฺเช วา โยคฺคํ กริตฺวา โส อปเรน สมเยน ทูเรปาตี จ โหติ อกฺขณเวธี จ มหโต จ กายสฺส ปทาเลตา.\csromanspacer{} เอวเมวํ โข ภิกฺขเว ภิกฺขุ สพฺพโส วิญฺญาณญฺจายตนํ สมติกฺกมฺม “นตฺถิ กิญฺจี”ติ อากิญฺจญฺญายตนํ อุปสมฺปชฺช วิหรติ.\csromanspacer{} โส ยเทว ตตฺถ โหติ เวทนาคตํ สญฺญาคตํ สงฺขารคตํ วิญฺญาณคตํ, เต ธมฺเม อนิจฺจโต ทุกฺขโต โรคโต คณฺฑโต สลฺลโต อฆโต อาพาธโต ปรโต ปโลกโต สุญฺญโต อนตฺตโต สมนุปสฺสติ.\csromanspacer{} โส เตหิ ธมฺเมหิ จิตฺตํ ปฏิวาเปติ.\csromanspacer{} โส เตหิ ธมฺเมหิ จิตฺตํ ปฏิวาเปตฺวา อมตาย ธาตุยา จิตฺตํ อุปสํหรติ “เอตํ สนฺตํ เอตํ ปณีตํ, ยทิทํ สพฺพสงฺขารสมโถ สพฺพูปธิปฏินิสฺสคฺโค ตณฺหากฺขโย วิราโค นิโรโธ นิพฺพานนฺ”ติ.\csromanspacer{} โส ตตฺถ ฐิโต อาสวานํ ขยํ ปาปุณาติ.\csromanspacer{} โน เจ อาสวานํ ขยํ ปาปุณาติ.\csromanspacer{} เตเนว ธมฺมราเคน ตาย ธมฺมนนฺทิยา ปญฺจนฺนํ โอรมฺภาคิยานํ สํโยชนานํ ปริกฺขยา โอปปาติโก โหติ.\csromanspacer{} ตตฺถ ปรินิพฺพายี อนาวตฺติธมฺโม ตสฺมา โลกา. “อากิญฺจญฺญา- ยตนมฺปาหํ ภิกฺขเว นิสฺสาย อาสวานํ ขยํ วทามี”ติ อิติ ยํ ตํ วุตฺตํ, อิทเมตํ ปฏิจฺจ วุตฺตํ.}

\prose{อิติ โข ภิกฺขเว ยาวตา สญฺญาสมาปตฺติ, ตาวตา อญฺญาปฏิเวโธ.\csromanspacer{} ยานิ จ โข อิมานิ ภิกฺขเว นิสฺสาย เทฺว อายตนานิ เนวสญฺญานาสญฺญายตน- สมาปตฺติ จ สญฺญาเวทยิตนิโรโธ จ.}

\vspace{\tipitakaparskip}
\csromancenter{นวกนิปาตปาฬิ ฌายีเหเต ภิกฺขเว สมาปตฺติกุสเลหิ สมาปตฺติวุฏฺฐานกุสเลหิ สมาปชฺชิตฺวา วุฏฺฐหิตฺวา สมฺมา อกฺขาตพฺพานีติ วทามีติ.\csromanspacer{}\csromanspacer{}ปญฺจมํ.}
\csromanheader{2}{6. อานนฺทสุตฺต}
\proseitem{37}{\csromanfolio{224}เอกํ สมยํ อายสฺมา อานนฺโท โกสมฺพิยํ วิหรติ โฆสิตาราเม.\csromanspacer{} ตตฺร โข อายสฺมา อานนฺโท ภิกฺขู อามนฺเตสิ “อาวุโส ภิกฺขเว”ติ. “อาวุโส”ติ โข เต ภิกฺขู อายสฺมโต อานนฺทสฺส ปจฺจสฺโสสุํ.\csromanspacer{} อายสฺมา อานนฺโท เอตทโวจ\csromandash{}}

\prose{อจฺฉริยํ อาวุโส อพฺภุตํ อาวุโส, ยาวญฺจิทํ เตน ภควตา ชานตา ปสฺสตา อรหตา สมฺมาสมฺพุทฺเธน สมฺพาเธ โอกาสาธิคโม อนุพุทฺโธ สตฺตานํ วิสุทฺธิยา โสกปริเทวานํ สมติกฺกมาย ทุกฺขโทมนสฺสานํ อตฺถงฺคมาย ญายสฺส อธิคมาย นิพฺพานสฺส สจฺฉิกิริยาย, ตเทว นาม จกฺขุํ ภวิสฺสติ เต รูปา ตญฺจายตนํ, โน ปฏิสํเวทิสฺสติ\footnote{ปฏิสํเวทยติ. (ก)}.\csromanspacer{} ตเทว นาม โสตํ ภวิสฺสติ เต สทฺทา ตญฺจายตนํ, โน ปฏิสํเวทิสฺสติ.\csromanspacer{} ตเทว นาม ฆานํ ภวิสฺสติ เต คนฺธา ตญฺจายตนํ, โน ปฏิสํเวทิสฺสติ.\csromanspacer{} สาว นาม ชิวฺหา ภวิสฺสติ เต รสา ตญฺจายตนํ, โน ปฏิสํเวทิสฺสติ.\csromanspacer{} โสว นาม กาโย ภวิสฺสติ เต โผฏฺฐพฺพา ตญฺจายตนํ, โน ปฏิสํเวทิสฺสตีติ.}

\prose{เอวํ วุตฺเต อายสฺมา อุทายี อายสฺมนฺตํ อานนฺทํ เอตทโวจ “สญฺญีเมว นุ โข อาวุโส อานนฺท ตทายตนํ โน ปฏิสํเวเทติ อุทาหุ อสญฺญี”ติ.\csromanspacer{} สญฺญีเมว โข อาวุโส ตทายตนํ โน ปฏิสํเวเทติ โน อสญฺญีติ.}

\prose{กิํสญฺญี ปนาวุโส ตทายตนํ โน ปฏิสํเวเทตีติ.\csromanspacer{} อิธาวุโส ภิกฺขุ สพฺพโส รูปสญฺญานํ สมติกฺกมา ปฏิฆสญฺญานํ อตฺถงฺคมา นานตฺตสญฺญานํ อมนสิการา “อนนฺโต อากาโส”ติ อากาสานญฺจายตนํ อุปสมฺปชฺช วิหรติ.\csromanspacer{} เอวํสญฺญีปิ โข อาวุโส ตทายตนํ โน ปฏิสํเวเทติ.}

\chapterheadpageread{4. มหาวคฺค}
\prose{\csromanfolio{225}ปุน จปรํ อาวุโส ภิกฺขุ สพฺพโส อากาสานญฺจายตนํ สมติกฺกมฺม “อนนฺตํ วิญฺญาณนฺ”ติ วิญฺญาณญฺจายตนํ อุปสมฺปชฺช วิหรติ.\csromanspacer{} เอวํสญฺญีปิ โข อาวุโส ตทายตนํ โน ปฏิสํเวเทติ.}

\prose{ปุน จปรํ อาวุโส ภิกฺขุ สพฺพโส วิญฺญาณญฺจายตนํ สมติกฺกมฺม “นตฺถิ กิญฺจี”ติ อากิญฺจญฺญายตนํ อุปสมฺปชฺช วิหรติ.\csromanspacer{} เอวํสญฺญีปิ โข อาวุโส ตทายตนํ โน ปฏิสํเวเทตีติ.}

\prose{เอกมิทาหํ อาวุโส สมยํ สาเกเต วิหรามิ อญฺชนวเน มิคทาเย.\csromanspacer{} อถ โข อาวุโส ชฏิลวาสิกา\footnote{ชฏิลคาหิยา (สี, อิ), ชฏิลภาคิกา (สฺยา)} ภิกฺขุนี เยนาหํ เตนุปสงฺกมิ, อุปสงฺกมิตฺวา มํ อภิวาเทตฺวา เอกมนฺตํ อฏฺฐาสิ, เอกมนฺตํ ฐิตา โข อาวุโส ชฏิลวาสิกา ภิกฺขุนี มํ เอตทโวจ “ยายํ ภนฺเต อานนฺท สมาธิ น จาภินโต น จาปนโต น จ สสงฺขารนิคฺคยฺหวาริตคโต\footnote{สสงฺขารนิคฺคยฺหวาริตวโต (สี, สฺยา, กํ, อิ), สสงฺขารนิคฺคยฺหวาริวาวโฏ (ก) อํ 1. 256, อํ 2. 20; ที 3. 234 ปิฏฺเฐสุปิ.}, วิมุตฺตตฺตา ฐิโต, ฐิตตฺตา สนฺตุสิโต, สนฺตุสิตตฺตา โน ปริตสฺสติ.\csromanspacer{} อยํ ภนฺเต อานนฺท สมาธิ กิํผโล วุตฺโต ภควตา”ติ.}

\prose{เอวํ วุตฺเต โสหํ อาวุโส ชฏิลวาสิกํ ภิกฺขุนิํ เอตทโวจํ “ยายํ ภคินิ สมาธิ น จาภินโต น จาปนโต น จ สสงฺขารนิคฺคยฺหวาริตคโต.\csromanspacer{} วิมุตฺตตฺตา ฐิโต, ฐิตตฺตา สนฺตุสิโต, สนฺตุสิตตฺตา โน ปริตสฺสติ.\csromanspacer{} อยํ ภคินิ สมาธิ อญฺญาผโล วุตฺโต ภควตา”ติ.\csromanspacer{} เอวํสญฺญีปิ โข อาวุโส ตทายตนํ โน ปฏิสํเวเทตีติ.\csromanspacer{}\csromanspacer{}ฉฏฺฐํ.}

\csromanheader{2}{7. โลกายติกสุตฺต}
\proseitem{38}{อถ โข เทฺว โลกายติกา พฺราหฺมณา เยน ภควา เตนุปสงฺกมิํสุ, อุปสงฺกมิตฺวา ภควตา สทฺธิํ สมฺโมทิํสุ, สมฺโมทนียํ กถํ สารณียํ วีติสาเรตฺวา เอกมนฺตํ นิสีทิํสุ, เอกมนฺตํ นิสินฺนา โข เต พฺราหฺมณา ภควนฺตํ เอตทโวจุํ\csromandash{}}

\prose{ปูรโณ โภ โคตม กสฺสโป สพฺพญฺญู สพฺพทสฺสาวี อปริเสสํ ญาณทสฺสนํ ปฏิชานาติ “จรโต จ เม ติฏฺฐโต จ สุตฺตสฺส จ}

\vspace{\tipitakaparskip}
\csromancenter{นวกนิปาตปาฬิ ชาครสฺส จ สตตํ สมิตํ ญาณทสฺสนํ ปจฺจุปฏฺฐิตนฺ”ติ.\csromanspacer{} โส เอวมาห “อหํ อนนฺเตน ญาเณน อนนฺตํ โลกํ ชานํ ปสฺสํ วิหรามี”ติ.\csromanspacer{} อยมฺปิ\footnote{อยมฺปิ หิ (สฺยา, ก)} โภ โคตม นิคณฺโฐ นาฏปุตฺโต สพฺพญฺญู สพฺพทสฺสาวี อปริเสสํ ญาณทสฺสนํ ปฏิชานาติ “จรโต จ เม ติฏฺฐโต จ สุตฺตสฺส จ ชาครสฺส จ สตตํ สมิตํ ญาณทสฺสนํ ปจฺจุปฏฺฐิตนฺ”ติ.\csromanspacer{} โส\footnote{โสปิ (?) 3. ทฬฺหธมฺโม (สพฺพตฺถ) อํ 1. 357; ม 1. 117 ปิฏฺเฐสุ จ, ตํสํวณฺณนาฏีกาโย จ โมคฺคลฺลานพฺยากรณญฺจ โอโลเกตพฺพา. 4. ตาลจฺฉาติํ (สี, สฺยา, อิ), ตาลจฺฉาทิํ (ก) อํ 1. 357; ม 1. 117 ปิฏฺเฐสุ ปสฺสิตพฺพํ.} เอวมาห “อหํ อนนฺเตน ญาเณน อนนฺตํ โลกํ ชานํ ปสฺสํ วิหรามี”ติ.\csromanspacer{} อิเมสํ โภ โคตม อุภินฺนํ ญาณวาทานํ อุภินฺนํ อญฺญมญฺญํ วิปจฺจนีกวาทานํ โก สจฺจํ อาห, โก มุสาติ.}
\prose{\csromanfolio{226}อลํ พฺราหฺมณา ติฏฺฐเตตํ “อิเมสํ อุภินฺนํ ญาณวาทานํ อุภินฺนํ อญฺญมญฺญํ วิปจฺจนีกวาทานํ โก สจฺจํ อาห, โก มุสา”ติ.\csromanspacer{} ธมฺมํ โว พฺราหฺมณา เทเสสฺสามิ, ตํ สุณาถ สาธุกํ มนสิ กโรถ, ภาสิสฺสามีติ. “เอวํ โภ”ติ โข เต พฺราหฺมณา ภควโต ปจฺจสฺโสสุํ.\csromanspacer{} ภควา เอตทโวจ\csromandash{}}

\prose{เสยฺยถาปิ พฺราหฺมณา จตฺตาโร ปุริสา จตุทฺทิสา ฐิตา ปรเมน ชเวน จ สมนฺนาคตา ปรเมน จ ปทวีติหาเรน, เต เอวรูเปน ชเวน สมนฺนาคตา อสฺสุ.\csromanspacer{} เสยฺยถาปิ นาม ทฬฺหธมฺมา3 ธนุคฺคโห สิกฺขิโต กตหตฺโถ กตูปาสโน ลหุเกน อสเนน อปฺปกสิเรน ติริยํ ตาลจฺฉายํ4 อติปาเตยฺย, เอวรูเปน จ ปทวีติหาเรน.\csromanspacer{} เสยฺยถาปิ นาม ปุรตฺถิมา สมุทฺทา ปจฺฉิโม สมุทฺโท, อถ ปุรตฺถิมาย ทิสาย ฐิโต ปุริโส เอวํ วเทยฺย “อหํ คมเนน โลกสฺส อนฺตํ ปาปุณิสฺสามี”ติ.\csromanspacer{} โส อญฺญเตฺรว อสิตปีตขายิตสายิตา อญฺญตฺร อุจฺจารปสฺสาวกมฺมา อญฺญตฺร นิทฺทากิลมถปฏิวิโนทนา วสฺสสตายุโก วสฺสสตชีวี วสฺสสตํ คนฺตฺวา อปฺปตฺวาว โลกสฺส อนฺตํ อนฺตรา กาลํ กเรยฺย.\csromanspacer{} อถ ปจฺฉิมาย ทิสาย ฯเปฯ.\csromanspacer{} อถ อุตฺตราย ทิสาย.\csromanspacer{} อถ ทกฺขิณาย ทิสาย ฐิโต ปุริโส เอวํ วเทยฺย “อหํ คมเนน โลกสฺส อนฺตํ ปาปุณิสฺสามี”ติ.\csromanspacer{} โส อญฺญเตฺรว อสิตปีตขายิตสายิตา อญฺญตฺร อุจฺจารปสฺสาวกมฺมา}

\chapterheadpageread{4. มหาวคฺค}
\prose{\csromanfolio{227}อญฺญตฺร นิทฺทากิลมถปฏิวิโนทนา วสฺสสตายุโก วสฺสสตชีวี วสฺสสตํ คนฺตฺวา อปฺปตฺวาว โลกสฺส อนฺตํ อนฺตรา กาลํ กเรยฺย.\csromanspacer{} ตํ กิสฺส เหตุ? นาหํ พฺราหฺมณา “เอวรูปาย สนฺธาวนิกาย โลกสฺส อนฺตํ ญาเตยฺยํ ทฏฺเฐยฺยํ ปตฺเตยฺยนฺ”ติ วทามิ, น จาหํ พฺราหฺมณา อปฺปตฺวาว โลกสฺส อนฺตํ ทุกฺขสฺส อนฺตกิริยํ วทามิ.}

\prose{ปญฺจิเม พฺราหฺมณา กามคุณา อริยสฺส วินเย “โลโก”ติ วุจฺจติ.\csromanspacer{} กตเม ปญฺจ, จกฺขุวิญฺเญยฺยา รูปา อิฏฺฐา กนฺตา มนาปา ปิยรูปา กามูปสํหิตา รชนียา, โสตวิญฺเญยฺยา สทฺทา ฯเปฯ.\csromanspacer{} ฆานวิญฺเญยฺยา คนฺธา.\csromanspacer{} ชิวฺหาวิญฺเญยฺยา รสา.\csromanspacer{} กายวิญฺเญยฺยา โผฏฺฐพฺพา อิฏฺฐา กนฺตา มนาปา ปิยรูปา กามูปสํหิตา รชนียา.\csromanspacer{} อิเม โข พฺราหฺมณา ปญฺจ กามคุณา อริยสฺส วินเย “โลโก”ติ วุจฺจติ.}

\prose{อิธ พฺราหฺมณา ภิกฺขุ วิวิจฺเจว กาเมหิ วิวิจฺจ อกุสเลหิ ธมฺเมหิ สวิตกฺกํ สวิจารํ วิเวกชํ ปีติสุขํ ปฐมํ ฌานํ อุปสมฺปชฺช วิหรติ.\csromanspacer{} อยํ วุจฺจติ พฺราหฺมณา ภิกฺขุ โลกสฺส อนฺตมาคมฺม โลกสฺส อนฺเต วิหรติ.\csromanspacer{} ตมญฺเญ เอวมาหํสุ “อยมฺปิ โลกปริยาปนฺโน, อยมฺปิ อนิสฺสโฏ โลกมฺหา”ติ.\csromanspacer{} อหมฺปิ หิ\footnote{อหมฺปิ (สี, อิ)} พฺราหฺมณา เอวํ วทามิ “อยมฺปิ โลกปริยาปนฺโน, อยมฺปิ อนิสฺสโฏ โลกมฺหา”ติ.}

\prose{ปุน จปรํ พฺราหฺมณา ภิกฺขุ วิตกฺกวิจารานํ วูปสมา ฯเปฯ ทุติยํ ฌานํ.\csromanspacer{} ตติยํ ฌานํ.\csromanspacer{} จตุตฺถํ ฌานํ อุปสมฺปชฺช วิหรติ.\csromanspacer{} อยํ วุจฺจติ พฺราหฺมณา ภิกฺขุ โลกสฺส อนฺตมาคมฺม โลกสฺส อนฺเต วิหรติ.\csromanspacer{} ตมญฺเญ เอวมาหํสุ “อยมฺปิ โลกปริยาปนฺโน, อยมฺปิ อนิสฺสโฏ โลกมฺหา”ติ.\csromanspacer{} อหมฺปิ หิ พฺราหฺมณา เอวํ วทามิ “อยมฺปิ โลกปริยาปนฺโน, อยมฺปิ อนิสฺสโฏ โลกมฺหา”ติ.}

\prose{ปุน จปรํ พฺราหฺมณา ภิกฺขุ สพฺพโส รูปสญฺญานํ สมติกฺกมา ปฏิฆสญฺญานํ อตฺถงฺคมา นานตฺตสญฺญานํ อมนสิการา “อนนฺโต อากาโส”ติ อากาสานญฺจายตนํ อุปสมฺปชฺช วิหรติ.\csromanspacer{} อยํ วุจฺจติ พฺราหฺมณา ภิกฺขุ โลกสฺส อนฺตมาคมฺม โลกสฺส อนฺเต วิหรติ.\csromanspacer{} ตมญฺเญ เอวมาหํสุ “อยมฺปิ โลกปริยาปนฺโน, อยมฺปิ อนิสฺสโฏ}

\vspace{\tipitakaparskip}
\csromancenter{นวกนิปาตปาฬิ โลกมฺหา”ติ.\csromanspacer{} อหมฺปิ หิ พฺราหฺมณา เอวํ วทามิ “อยมฺปิ โลกปริยาปนฺโน, อยมฺปิ อนิสฺสโฏ โลกมฺหา”ติ.}
\prose{\csromanfolio{228}ปุน จปรํ พฺราหฺมณา ภิกฺขุ สพฺพโส อากาสานญฺจายตนํ สมติกฺกมฺม “อนนฺตํ วิญฺญาณนฺ”ติ วิญฺญาณญฺจายตนํ อุปสมฺปชฺช วิหรติ ฯเปฯ สพฺพโส วิญฺญาณญฺจายตนํ สมติกฺกมฺม “นตฺถิ กิญฺจี”ติ อากิญฺจญฺญายตนํ อุปสมฺปชฺช วิหรติ ฯเปฯ สพฺพโส อากิญฺจญฺญายตนํ สมติกฺกมฺม เนวสญฺญานาสญฺญายตนํ อุปสมฺปชฺช วิหรติ.\csromanspacer{} อยํ วุจฺจติ พฺราหฺมณา ภิกฺขุ โลกสฺส อนฺตมาคมฺม โลกสฺส อนฺเต วิหรติ.\csromanspacer{} ตมญฺเญ เอวมาหํสุ “อยมฺปิ โลกปริยาปนฺโน, อยมฺปิ อนิสฺสโฏ โลกมฺหา”ติ.\csromanspacer{} อหมฺปิ หิ พฺราหฺมณา เอวํ วทามิ “อยมฺปิ โลกปริยาปนฺโน, อยมฺปิ อนิสฺสโฏ โลกมฺหา”ติ.}

\prose{ปุน จปรํ พฺราหฺมณา ภิกฺขุ สพฺพโส เนวสญฺญานาสญฺญายตนํ สมติกฺกมฺม สญฺญาเวทยิตนิโรธํ อุปสมฺปชฺช วิหรติ, ปญฺญาย จสฺส ทิสฺวา อาสวา ปริกฺขีณา โหนฺติ.\csromanspacer{} อยํ วุจฺจติ พฺราหฺมณา ภิกฺขุ “โลกสฺส อนฺตมาคมฺม โลกสฺส อนฺเต วิหรติ, ติณฺโณ โลเก วิสตฺติกนฺ”ติ.\csromanspacer{}\csromanspacer{}สตฺตมํ.}

\csromanheader{2}{8. เทวาสุรสงฺคามสุตฺต}
\proseitem{39}{ภูตปุพฺพํ ภิกฺขเว เทวาสุรสงฺคาโม สมุปพฺยูโฬฺห\footnote{สมุปพฺพูโฬฺห (สี, อิ)} อโหสิ.\csromanspacer{} ตสฺมิํ โข ปน ภิกฺขเว สงฺคาเม อสุรา ชินิํสุ, เทวา ปราชยิํสุ\footnote{ปราชิยิํสุ (สี, สฺยา, ก)}.\csromanspacer{} ปราชิตา จ ภิกฺขเว เทวา\footnote{เทวา ภีตา (อิ)} อปยิํสุเยว\footnote{อปยํเสฺวว (สี)} อุตฺตเรนาภิมุขา, อภิยิํสุ\footnote{อภิยํสุ (สี)} อสุรา.\csromanspacer{} อถ โข ภิกฺขเว เทวานํ เอตทโหสิ “อภิยนฺเตว โข อสุรา, ยํนูน มยํ ทุติยมฺปิ อสุเรหิ สงฺคาเมยฺยามา”ติ.\csromanspacer{} ทุติยมฺปิ โข ภิกฺขเว เทวา อสุเรหิ สงฺคาเมสุํ.\csromanspacer{} ทุติยมฺปิ โข ภิกฺขเว อสุราว ชินิํสุ, เทวา ปราชยิํสุ.\csromanspacer{} ปราชิตา จ ภิกฺขเว เทวา อปยิํสุเยว อุตฺตเรนาภิมุขา, อภิยิํสุ อสุรา. อถ โข ภิกฺขเว เทวานํ เอตทโหสิ “อภิยนฺเตว โข อสุรา, ยํนูน มยํ ตติยมฺปิ อสุเรหิ สงฺคาเมยฺยามา”ติ.\csromanspacer{} ตติยมฺปิ โข}

\chapterheadpageread{4. มหาวคฺค}
\prose{\csromanfolio{229}ภิกฺขเว เทวา อสุเรหิ สงฺคาเมสุํ.\csromanspacer{} ตติยมฺปิ โข ภิกฺขเว อสุราว ชินิํสุ, เทวา ปราชยิํสุ.\csromanspacer{} ปราชิตา จ ภิกฺขเว เทวา ภีตา เทวปุรํเยว ปวิสิํสุ.\csromanspacer{} เทวปุรคตานญฺจ ปน\footnote{ปุน (ก)} ภิกฺขเว เทวานํ เอตทโหสิ “ภีรุตฺตานคเตน โข ทานิ มยํ เอตรหิ อตฺตนา วิหราม, อกรณียา อสุเรหี”ติ.\csromanspacer{} อสุรานมฺปิ ภิกฺขเว เอตทโหสิ “ภีรุตฺตานคเตน โข ทานิ เทวา เอตรหิ อตฺตนา วิหรนฺติ, อกรณียา อเมฺหหี”ติ.}

\prose{ภูตปุพฺพํ ภิกฺขเว เทวาสุรสงฺคาโม สมุปพฺยูโฬฺห อโหสิ.\csromanspacer{} ตสฺมิํ โข ปน ภิกฺขเว สงฺคาเม เทวา ชินิํสุ, อสุรา ปราชยิํสุ.\csromanspacer{} ปราชิตา จ ภิกฺขเว อสุรา อปยิํสุเยว ทกฺขิเณนาภิมุขา, อภิยิํสุ เทวา.\csromanspacer{} อถ โข ภิกฺขเว อสุรานํ เอตทโหสิ “อภิยนฺเตว โข เทวา, ยํนูน มยํ ทุติยมฺปิ เทเวหิ สงฺคาเมยฺยามา”ติ.\csromanspacer{} ทุติยมฺปิ โข ภิกฺขเว อสุรา เทเวหิ สงฺคาเมสุํ.\csromanspacer{} ทุติยมฺปิ โข ภิกฺขเว เทวา ชินิํสุ, อสุรา ปราชยิํสุ.\csromanspacer{} ปราชิตา จ ภิกฺขเว อสุรา อปยิํสุเยว ทกฺขิเณนาภิมุขา, อภิยิํสุ เทวา.}

\prose{อถ โข ภิกฺขเว อสุรานํ เอตทโหสิ “อภิยนฺเตว โข เทวา, ยํนูน มยํ ตติยมฺปิ เทเวหิ สงฺคาเมยฺยามา”ติ.\csromanspacer{} ตติยมฺปิ โข ภิกฺขเว อสุรา เทเวหิ สงฺคาเมสุํ.\csromanspacer{} ตติยมฺปิ โข ภิกฺขเว เทวา ชินิํสุ, อสุรา ปราชยิํสุ.\csromanspacer{} ปราชิตา จ ภิกฺขเว อสุรา ภีตา อสุรปุรํเยว ปวิสิํสุ.\csromanspacer{} อสุรปุรคตานญฺจ ปน ภิกฺขเว อสุรานํ เอตทโหสิ “ภีรุตฺตานคเตน โข ทานิ มยํ เอตรหิ อตฺตนา วิหราม, อกรณียา เทเวหี”ติ.\csromanspacer{} เทวานมฺปิ ภิกฺขเว เอตทโหสิ “ภีรุตฺตานคเตน โข ทานิ อสุรา เอตรหิ อตฺตนา วิหรนฺติ, อกรณียา อเมฺหหี”ติ.}

\prose{เอวเมวํ โข ภิกฺขเว ยสฺมิํ สมเย ภิกฺขุ วิวิจฺเจว กาเมหิ วิวิจฺจ อกุสเลหิ ธมฺเมหิ สวิตกฺกํ สวิจารํ วิเวกชํ ปีติสุขํ ปฐมํ ฌานํ อุปสมฺปชฺช วิหรติ, ตสฺมิํ ภิกฺขเว สมเย ภิกฺขุสฺส เอวํ โหติ “ภีรุตฺตานคเตน โข ทานาหํ เอตรหิ อตฺตนา วิหรามิ, อกรณีโย}

\vspace{\tipitakaparskip}
\csromancenter{นวกนิปาตปาฬิ มารสฺสา”ติ.\csromanspacer{} มารสฺสาปิ ภิกฺขเว ปาปิมโต เอวํ โหติ “ภีรุตฺตานคเตน โข ทานิ ภิกฺขุ เอตรหิ อตฺตนา วิหรติ, อกรณีโย มยฺหนฺ”ติ.}
\prose{\csromanfolio{230}ยสฺมิํ ภิกฺขเว สมเย ภิกฺขุ วิตกฺกวิจารานํ วูปสมา ฯเปฯ ทุติยํ ฌานํ.\csromanspacer{} ตติยํ ฌานํ.\csromanspacer{} จตุตฺถํ ฌานํ อุปสมฺปชฺช วิหรติ, ตสฺมิํ ภิกฺขเว สมเย ภิกฺขุสฺส เอวํ โหติ “ภีรุตฺตานคเตน โข ทานาหํ เอตรหิ อตฺตนา วิหรามิ, อกรณีโย มารสฺสา”ติ.\csromanspacer{} มารสฺสาปิ ภิกฺขเว ปาปิมโต เอวํ โหติ “ภีรุตฺตานคเตน โข ทานิ ภิกฺขุ เอตรหิ อตฺตนา วิหรติ, อกรณีโย มยฺหนฺ”ติ.}

\prose{ยสฺมิํ ภิกฺขเว สมเย ภิกฺขุ สพฺพโส รูปสญฺญานํ สมติกฺกมา ปฏิฆสญฺญานํ อตฺถงฺคมา นานตฺตสญฺญานํ อมนสิการา “อนนฺโต อากาโส”ติ อากาสานญฺจายตนํ อุปสมฺปชฺช วิหรติ.\csromanspacer{} อยํ วุจฺจติ ภิกฺขเว ภิกฺขุ “อนฺตมกาสิ มารํ อปทํ, วธิตฺวา มารจกฺขุํ อทสฺสนํ คโต ปาปิมโต, ติณฺโณ โลเก วิสตฺติกนฺ”ติ.}

\prose{ยสฺมิํ ภิกฺขเว สมเย ภิกฺขุ สพฺพโส อากาสานญฺจายตนํ สมติกฺกมฺม “อนนฺตํ วิญฺญาณนฺ”ติ วิญฺญาณญฺจายตนํ อุปสมฺปชฺช วิหรติ.\csromanspacer{} สพฺพโส วิญฺญาณญฺจายตนํ สมติกฺกมฺม “นตฺถิ กิญฺจี”ติ อากิญฺจญฺญายตนํ อุปสมฺปชฺช วิหรติ.\csromanspacer{} สพฺพโส อากิญฺจญฺญายตนํ สมติกฺกมฺม เนวสญฺญานาสญฺญายตนํ อุปสมฺปชฺช วิหรติ.\csromanspacer{} สพฺพโส เนวสญฺญานาสญฺญายตนํ สมติกฺกมฺม สญฺญาเวทยิตนิโรธํ อุปสมฺปชฺช วิหรติ, ปญฺญาย จสฺส ทิสฺวา อาสวา ปริกฺขีณา โหนฺติ.\csromanspacer{} อยํ วุจฺจติ ภิกฺขเว ภิกฺขุ “อนฺตมกาสิ มารํ อปทํ, วธิตฺวา มารจกฺขุํ อทสฺสนํ คโต ปาปิมโต, ติณฺโณ โลเก วิสตฺติกนฺ”ติ.\csromanspacer{}\csromanspacer{}อฏฺฐมํ.}

\csromanheader{2}{9. นาคสุตฺต}
\proseitem{40}{ยสฺมิํ ภิกฺขเว สมเย อารญฺญิกสฺส นาคสฺส โคจรปสุตสฺส หตฺถีปิ หตฺถินิโยปิ หตฺถิกลภาปิ หตฺถิจฺฉาปาปิ ปุรโต ปุรโต คนฺตฺวา ติณคฺคานิ ฉินฺทนฺติ, เตน ภิกฺขเว อารญฺญิโก นาโค อฏฺฏียติ หรายติ ชิคุจฺฉติ.\csromanspacer{} ยสฺมิํ ภิกฺขเว สมเย อารญฺญิกสฺส นาคสฺส โคจรปสุตสฺส หตฺถีปิ หตฺถินิโยปิ หตฺถิกลภาปิ หตฺถิจฺฉาปาปิ โอภคฺโคภคฺคํ สาขาภงฺคํ ขาทนฺติ, เตน ภิกฺขเว อารญฺญิโก นาโค อฏฺฏียติ หรายติ ชิคุจฺฉติ.}

\chapterheadpageread{4. มหาวคฺค}
\prose{\csromanfolio{231}ยสฺมิํ ภิกฺขเว สมเย อารญฺญิกสฺส นาคสฺส โอคาหํ โอติณฺณสฺส หตฺถีปิ หตฺถินิโยปิ หตฺถิกลภาปิ หตฺถิจฺฉาปาปิ ปุรโต ปุรโต คนฺตฺวา โสณฺฑาย อุทกํ อาโลเฬนฺติ, เตน ภิกฺขเว อารญฺญิโก นาโค อฏฺฏียติ หรายติ ชิคุจฺฉติ.\csromanspacer{} ยสฺมิํ ภิกฺขเว สมเย อารญฺญิกสฺส นาคสฺส โอคาหา อุตฺติณฺณสฺส หตฺถินิโย กายํ อุปนิฆํสนฺติโย คจฺฉนฺติ, เตน ภิกฺขเว อารญฺญิโก นาโค อฏฺฏียติ หรายติ ชิคุจฺฉติ.}

\prose{ตสฺมิํ ภิกฺขเว สมเย อารญฺญิกสฺส นาคสฺส เอวํ โหติ “อหํ โข เอตรหิ อากิณฺโณ วิหรามิ หตฺถีหิ หตฺถินีหิ หตฺถิกลเภหิ หตฺถิจฺฉาเปหิ, ฉินฺนคฺคานิ เจว ติณานิ ขาทามิ, โอภคฺโคภคฺคญฺจ เม สาขาภงฺคํ ขาทนฺติ\footnote{ขาทิตํ (สฺยา, ก) วิ 3. 500 ปิฏฺเฐ ปสฺสิตพฺพํ. 2. โอคาหาปิ จ (สฺยา, ก) วิ 3. 500 ปิฏฺเฐ ปสฺสิตพฺพํ. 3. น โอภคฺโคภคฺคญฺจ สาขาภงฺคํ ขาทติ (สฺยา, ก) 4. เอตฺถ ปิสทฺโท สพฺพตฺถปิ นตฺถิ. 5. กณฺฑุํ สํหนฺติ (สี, อิ), กณฺฑุํ สํหนติ (สฺยา), เอตฺถ กณฺฑุวนทุกฺขํ วิเนตีติ อตฺโถ.}, อาวิลานิ จ ปานียานิ ปิวามิ, โอคาหา จ2 เม อุตฺติณฺณสฺส หตฺถินิโย กายํ อุปนิฆํสนฺติโย คจฺฉนฺติ.\csromanspacer{} ยํนูนาหํ เอโก คณสฺมา วูปกฏฺโฐ วิหเรยฺยนฺ”ติ.\csromanspacer{} โส อปเรน สมเยน เอโก คณสฺมา วูปกฏฺโฐ วิหรติ, อจฺฉินฺนคฺคานิ เจว ติณานิ ขาทติ, โอภคฺโคภคฺคญฺจสฺส สาขาภงฺคํ น ขาทนฺติ3, อนาวิลานิ จ ปานียานิ ปิวติ, โอคาหา จสฺส อุตฺติณฺณสฺส น หตฺถินิโย กายํ อุปนิฆํสนฺติโย คจฺฉนฺติ.}

\prose{ตสฺมิํ ภิกฺขเว สมเย อารญฺญิกสฺส นาคสฺส เอวํ โหติ “อหํ โข ปุพฺเพ อากิณฺโณ วิหาสิํ หตฺถีหิ หตฺถินีหิ หตฺถิกลเภหิ หตฺถิจฺฉาเปหิ, ฉินฺนคฺคานิ เจว ติณานิ ขาทิํ, โอภคฺโคภคฺคญฺจ เม สาขาภงฺคํ ขาทิํสุ, อาวิลานิ จ ปานียานิ อปายิํ, โอคาหา4 จ เม อุตฺติณฺณสฺส หตฺถินิโย กายํ อุปนิฆํสนฺติโย อคมํสุ.\csromanspacer{} โสหํ เอตรหิ เอโก คณสฺมา วูปกฏฺโฐ วิหรามิ, อจฺฉินฺนคฺคานิ เจว ติณานิ ขาทามิ, โอภคฺโคภคฺคญฺจ เม สาขาภงฺคํ น ขาทนฺติ, อนาวิลานิ จ ปานียานิ ปิวามิ, โอคาหา จ เม อุตฺติณฺณสฺส น หตฺถินิโย กายํ อุปนิฆํสนฺติโย คจฺฉนฺตี”ติ.\csromanspacer{} โส โสณฺฑาย สาขาภงฺคํ ภญฺชิตฺวา สาขาภงฺเคน กายํ ปริมชฺชิตฺวา อตฺตมโน โสณฺฑํ สํหรติ5.}

\prose{เอวเมวํ โข ภิกฺขเว ยสฺมิํ สมเย ภิกฺขุ อากิณฺโณ วิหรติ ภิกฺขูหิ ภิกฺขุนีหิ อุปาสเกหิ อุปาสิกาหิ รญฺญา ราชมหามตฺเตหิ}

\vspace{\tipitakaparskip}
\csromancenter{นวกนิปาตปาฬิ ติตฺถิเยหิ ติตฺถิยสาวเกหิ, ตสฺมิํ ภิกฺขเว สมเย ภิกฺขุสฺส เอวํ โหติ “อหํ โข เอตรหิ อากิณฺโณ วิหรามิ ภิกฺขูหิ ภิกฺขุนีหิ อุปาสเกหิ อุปาสิกาหิ รญฺญา ราชมหามตฺเตหิ ติตฺถิเยติ ติตฺถิยสาวเกหิ, ยํนูนาหํ เอโก คณสฺมา วูปกฏฺโฐ วิหเรยฺยนฺ”ติ.\csromanspacer{} โส วิวิตฺตํ เสนาสนํ ภชติ อรญฺญํ รุกฺขมูลํ ปพฺพตํ กนฺทรํ คิริคุหํ สุสานํ วนปตฺถํ อพฺโภกาสํ ปลาลปุญฺชํ, โส อรญฺญคโต วา รุกฺขมูลคโต วา สุญฺญาคารคโต วา นิสีทติ ปลฺลงฺกํ อาภุชิตฺวา อุชุํ กายํ ปณิธาย ปริมุขํ สติํ อุปฏฺฐเปตฺวา.}
\prose{\csromanfolio{232}โส อภิชฺฌํ โลเก ปหาย วิคตาภิชฺเฌน เจตสา วิหรติ, อภิชฺฌาย จิตฺตํ ปริโสเธติ.\csromanspacer{} พฺยาปาทปโทสํ ปหาย อพฺยาปนฺนจิตฺโต วิหรติ สพฺพปาณภูตหิตานุกมฺปี, พฺยาปาทปโทสา จิตฺตํ ปริโสเธติ.\csromanspacer{} ถินมิทฺธํ ปหาย วิคตถินมิทฺโธ วิหรติ อาโลกสญฺญี สโต สมฺปาชาโน, ถินมิทฺธา จิตฺตํ ปริโสเธติ.\csromanspacer{} อุทฺธจฺจกุกฺกุจฺจํ ปหาย อนุทฺธโต วิหรติ อชฺฌตฺตํ วูปสนฺตจิตฺโต, อุทฺธจฺจกุกฺกุจฺจา จิตฺตํ ปริโสเธติ.\csromanspacer{} วิจิกิจฺฉํ ปหาย ติณฺณวิจิกิจฺโฉ วิหรติ อกถํกถี กุสเลสุ ธมฺเมสุ, วิจิกิจฺฉาย จิตฺตํ ปริโสเธติ.\csromanspacer{} โส อิเม ปญฺจ นีวรเณ ปหาย เจตโส อุปกฺกิเลเส ปญฺญาย ทุพฺพลีกรเณ วิวิจฺเจว กาเมหิ วิวิจฺจ อกุสเลหิ ธมฺเมหิ สวิตกฺกํ สวิจารํ วิเวกชํ ปีติสุขํ ปฐมํ ฌานํ อุปสมฺปชฺช วิหรติ.\csromanspacer{} โส อตฺตมโน โสณฺฑํ สํหรติ\footnote{กณฺฑุํ สํหนฺติ (สี, อิ), กณฺฑุํ สํหนติ (สฺยา), เอตฺถ กณฺฑุวนสทิสํ ฌานปฏิปกฺขํ กิเลสทุกฺขํ วิเนตีติ อตฺโถ.}.\csromanspacer{} วิตกฺกวิจารานํ วูปสมา ฯเปฯ ทุติยํ ฌานํ.\csromanspacer{} ตติยํ ฌานํ.\csromanspacer{} จตุตฺถํ ฌานํ อุปสมฺปชฺช วิหรติ.\csromanspacer{} โส อตฺตมโน โสณฺฑํ สํหรติ.}

\prose{สพฺพโส รูปสญฺญานํ สมติกฺกมา ปฏิฆสญฺญานํ อตฺถงฺคมา นานตฺตสญฺญานํ อมนสิการา “อนนฺโต อากาโส”ติ อากาสานญฺจายตนํ อุปสมฺปชฺช วิหรติ.\csromanspacer{} โส อตฺตมโน โสณฺฑํ สํหรติ.\csromanspacer{} สพฺพโส อากาสานญฺจายตนํ สมติกฺกมฺม “อนนฺตํ วิญฺญาณนฺ”ติ วิญฺญาณญฺจายตนํ อุปสมฺปชฺช วิหรติ.\csromanspacer{} สพฺพโส วิญฺญาณญฺจายตนํ สมติกฺกมฺม “นตฺถิ กิญฺจี”ติ อากิญฺจญฺญายตนํ อุปสมฺปชฺช วิหรติ.\csromanspacer{} สพฺพโส อากิญฺจญฺญายตนํ สมติกฺกมฺม เนวสญฺญานาสญฺญายตนํ อุปสมฺปชฺช วิหรติ.\csromanspacer{} สพฺพโส}

\chapterheadpageread{4. มหาวคฺค}
\prose{\csromanfolio{233}เนวสญฺญานาสญฺญายตนํ สมติกฺกมฺม สญฺญาเวทยิตนิโรธํ อุปสมฺปชฺช วิหรติ, ปญฺญาย จสฺส ทิสฺวา อาสวา ปริกฺขีณา โหนฺติ.\csromanspacer{} โส อตฺตมโน โสณฺฑํ สํหรตีติ.\csromanspacer{}\csromanspacer{}นวมํ.}

\csromanheader{2}{10. ตปุสฺสสุตฺต}
\proseitem{41}{เอกํ สมยํ ภควา มลฺเลสุ วิหรติ อุรุเวลกปฺปํ นาม มลฺลานํ นิคโม.\csromanspacer{} อถ โข ภควา ปุพฺพณฺหสมยํ นิวาเสตฺวา ปตฺตจีวรมาทาย อุรุเวลกปฺปํ ปิณฺฑาย ปาวิสิ, อุรุเวลกปฺเป ปิณฺฑาย จริตฺวา ปจฺฉาภตฺตํ ปิณฺฑปาตปฏิกฺกนฺโต อายสฺมนฺตํ อานนฺทํ อามนฺเตสิ “อิเธว ตาว ตฺวํ อานนฺท โหติ, ยาวาหํ มหาวนํ อชฺโฌคาหามิ ทิวาวิหารายา”ติ. “เอวํ ภนฺเต”ติ โข อายสฺมา อานนฺโท ภควโต ปจฺจสฺโสสิ.\csromanspacer{} อถ โข ภควา มหาวนํ อชฺโฌคาเหตฺวา อญฺญตรสฺมิํ รุกฺขมูเล ทิวาวิหารํ นิสีทิ.}

\prose{อถ โข ตปุสฺโส คหปติ เยนายสฺมา อานนฺโท เตนุปสงฺกมิ, อุปสงฺกมิตฺวา อายสฺมนฺตํ อานนฺทํ อภิวาเทตฺวา เอกมนฺตํ นิสีทิ, เอกมนฺตํ นิสินฺโน โข ตปุสฺโส คหปติ อายสฺมนฺตํ อานนฺทํ เอตทโวจ\csromandash{}}

\prose{มยํ ภนฺเต อานนฺท คิหี กามโภคิโน กามารามา กามรตา กามสมฺมุทิตา, เตสํ โน ภนฺเต อมฺหากํ คิหีนํ กามโภคีนํ กามารามานํ กามรตานํ กามสมฺมุทิตานํ ปปาโต วิย ขายติ, ยทิทํ เนกฺขมฺมํ.\csromanspacer{} สุตํ เมตํ ภนฺเต “อิมสฺมิํ ธมฺมวินเย ทหรานํ ทหรานํ ภิกฺขูนํ เนกฺขมฺเม จิตฺตํ ปกฺขนฺทติ ปสีทติ สนฺติฏฺฐติ วิมุจฺจติ ‘เอตํ สนฺตนฺ’ติ ปสฺสโต\footnote{ปสฺสตํ (?)}”.\csromanspacer{} ตยิทํ ภนฺเต อิมสฺมิํ ธมฺมวินเย ภิกฺขูนํ พหุนา ชเนน วิสภาโค, ยทิทํ เนกฺขมฺมนฺติ.}

\prose{อตฺถิ โข เอตํ คหปติ กถาปาภตํ ภควนฺตํ ทสฺสนาย, อายาม คหปติ เยน ภควา เตนุปสงฺกมิสฺสาม, อุปสงฺกมิตฺวา ภควโต เอตมตฺถํ อาโรเจสฺสาม, ยถา โน ภควา พฺยากริสฺสติ, ตถา นํ ธาเรสฺสามาติ.}

\prose{“เอวํ ภนฺเต”ติ โข ตปุสฺโส คหปติ อายสฺมโต อานนฺทสฺส ปจฺจสฺโสสิ.\csromanspacer{} อถ โข อายสฺมา อานนฺโท ตปุสฺเสน คหปตินา}

\vspace{\tipitakaparskip}
\csromancenter{นวกนิปาตปาฬิ สทฺธิํ เยน ภควา เตนุปสงฺกมิ, อุปสงฺกมิตฺวา ภควนฺตํ อภิวาเทตฺวา เอกมนฺตํ นิสีทิ, เอกมนฺตํ นิสินฺโน โข อายสฺมา อานนฺโท ภควนฺตํ เอตทโวจ\csromandash{}}
\prose{\csromanfolio{234}อยํ ภนฺเต ตปุสฺโส คหปติ เอวมาห “มยํ ภนฺเต อานนฺท คิหี กามโภคิโน กามารามา กามรตา กามสมฺมุทิตา, เตสํ โน ภนฺเต อมฺหากํ คิหีนํ กามโภคีนํ กามารามานํ กามรตานํ กามสมฺมุทิตานํ ปปาโต วิย ขายติ, ยทิทํ เนกฺขมฺมํ.\csromanspacer{} สุตํ เมตํ ภนฺเต ‘อิมสฺมิํ ธมฺมวินเย ทหรานํ ทหรานํ ภิกฺขูนํ เนกฺขมฺเม จิตฺตํ ปกฺขนฺทติ ปสีทติ สนฺติฏฺฐติ วิมุจฺจติ เอตํ สนฺตนฺติ ปสฺสโต’.\csromanspacer{} ตยิทํ ภนฺเต อิมสฺมิํ ธมฺมวินเย ภิกฺขูนํ พหุนา ชเนน วิสภาโค, ยทิทํ เนกฺขมฺมนฺ”ติ.}

\prose{เอวเมตํ อานนฺท, เอวเมตํ อานนฺท, มยฺหมฺปิ โข อานนฺท ปุพฺเพว สมฺโพธา อนภิสมฺพุทฺธสฺส โพธิสตฺตสฺเสว สโต เอตทโหสิ “สาธุ เนกฺขมฺมํ สาธุ ปวิเวโก”ติ.\csromanspacer{} ตสฺส มยฺหํ อานนฺท เนกฺขมฺเม จิตฺตํ น ปกฺขนฺทติ นปฺปสีทติ น สนฺติฏฺฐติ น วิมุจฺจติ “เอตํ สนฺตนฺ”ติ ปสฺสโต.\csromanspacer{} ตสฺส มยฺหํ อานนฺท เอตทโหสิ “โก นุโข เหตุ โก ปจฺจโย, เยน เม เนกฺขมฺเม จิตฺตํ น ปกฺขนฺทติ นปฺปสีทติ น สนฺติฏฺฐติ น วิมุจฺจติ ‘เอตํ สนฺตนฺ’ติ ปสฺสโต”.\csromanspacer{} ตสฺส มยฺหํ อานนฺท เอตทโหสิ “กาเมสุ โข เม อาทีนโว อทิฏฺโฐ, โส จ เม อพหุลีกโต.\csromanspacer{} เนกฺขมฺเม จ อานิสํโส อนธิคโต, โส จ เม อนาเสวิโต.\csromanspacer{} ตสฺมา เม เนกฺขมฺเม จิตฺตํ น ปกฺขนฺทติ นปฺปสีทติ น สนฺติฏฺฐติ น วิมุจฺจติ ‘เอตํ สนฺตนฺ’ติ ปสฺสโต”.\csromanspacer{} ตสฺส มยฺหํ อานนฺท เอตทโหสิ “สเจ โข อหํ กาเมสุ อาทีนวํ ทิสฺวา ตํ พหุลํ กเรยฺยํ\footnote{พหุลีกเรยฺยํ (สี, สฺยา, อิ)}, เนกฺขมฺเม อานิสํสํ อธิคมฺม ตมาเสเวยฺยํ.\csromanspacer{} ฐานํ โข ปเนตํ วิชฺชติ, ยํ เม เนกฺขมฺเม จิตฺตํ ปกฺขนฺเทยฺย ปสีเทยฺย สนฺติฏฺเฐยฺย วิมุจฺเจยฺย ‘เอตํ สนฺตนฺ’ติ ปสฺสโต”.\csromanspacer{} โส โข อหํ อานนฺท อปเรน สมเยน กาเมสุ อาทีนวํ ทิสฺวา ตํ พหุลมกาสิํ, เนกฺขมฺเม อานิสํสํ อธิคมฺม ตมาเสวิํ.\csromanspacer{} ตสฺส มยฺหํ อานนฺท เนกฺขมฺเม จิตฺตํ ปกฺขนฺทติ ปสีทติ สนฺติฏฺฐติ วิมุจฺจติ “เอตํ สนฺตนฺ”ติ ปสฺสโต, โส โข อหํ อานนฺท วิวิจฺเจว กาเมหิ วิวิจฺจ อกุสเลหิ ธมฺเมหิ สวิตกฺกํ สวิจารํ วิเวกชํ ปีติสุขํ ปฐมํ ฌานํ อุปสมฺปชฺช วิหรามิ.\csromanspacer{} ตสฺส มยฺหํ อานนฺท อิมินา วิหาเรน วิหรโต กามสหคตา สญฺญามนสิการา สมุทาจรนฺติ, สฺวสฺส เม โหติ}

\chapterheadpageread{4. มหาวคฺค}
\prose{\csromanfolio{235}อาพาโธ.\csromanspacer{} เสยฺยถาปิ อานนฺท สุขิโน ทุกฺขํ อุปฺปชฺเชยฺย ยาวเทว อาพาธาย.\csromanspacer{} เอวเมวสฺส เม กามสหคตา สญฺญามนสิการา สมุทาจรนฺติ, สฺวสฺส เม โหติ อาพาโธ.}

\prose{ตสฺส มยฺหํ อานนฺท เอตทโหสิ “ยํนูนาหํ วิตกฺกวิจารานํ วูปสมา ฯเปฯ ทุติยํ ฌานํ อุปสมฺปชฺช วิหเรยฺยนฺ”ติ.\csromanspacer{} ตสฺส มยฺหํ อานนฺท อวิตกฺเก จิตฺตํ น ปกฺขนฺทติ นปฺปสีทติ น สนฺติฏฺฐติ น วิมุจฺจติ “เอตํ สนฺตนฺ”ติ ปสฺสโต.\csromanspacer{} ตสฺส มยฺหํ อานนฺท เอตทโหสิ “โก นุ โข เหตุ โก ปจฺจโย, เยน เม อวิตกฺเก จิตฺตํ น ปกฺขนฺทติ นปฺปสีทติ น สนฺติฏฺฐติ น วิมุจฺจติ ‘เอตํ สนฺตนฺ’ติ ปสฺสโต”.\csromanspacer{} ตสฺส มยฺหํ อานนฺท เอตทโหสิ “วิตกฺเกสุ โข เม อาทีนโว อทิฏฺโฐ, โส จ เม อพหุลีกโต.\csromanspacer{} อวิตกฺเก จ อานิสํโส อนธิคโต, โส จ เม อนาเสวิโต.\csromanspacer{} ตสฺมา เม อวิตกฺเก จิตฺตํ น ปกฺขนฺทติ นปฺปสีทติ น สนฺติฏฺฐติ น วิมุจฺจติ ‘เอตํ สนฺตนฺ’ติ ปสฺสโต”.\csromanspacer{} ตสฺส มยฺหํ อานนฺท เอตทโหสิ “สเจ โข อหํ วิตกฺเกสุ อาทีนวํ ทิสฺวา ตํ พหุลํ กเรยฺยํ, อวิตกฺเก อานิสํสํ อธิคมฺม ตมาเสเวยฺยํ.\csromanspacer{} ฐานํ โข ปเนตํ วิชฺชติ, ยํ เม อวิตกฺเก จิตฺตํ ปกฺขนฺเทยฺย ปสีเทยฺย สนฺติฏฺเฐยฺย วิมุจฺเจยฺย ‘เอตํ สนฺตนฺ’ติ ปสฺสโต”.\csromanspacer{} โส โข อหํ อานนฺท อปเรน สมเยน วิตกฺเกสุ อาทีนวํ ทิสฺวา ตํ พหุลมกาสิํ, อวิตกฺเก อานิสํสํ อธิคมฺม ตมาเสวิํ.\csromanspacer{} ตสฺส มยฺหํ อานนฺท อวิตกฺเก จิตฺตํ ปกฺขนฺทติ ปสีทติ สนฺติฏฺฐติ วิมุจฺจติ “เอตํ สนฺตนฺ”ติ ปสฺสโต.\csromanspacer{} โส โข อหํ อานนฺท วิตกฺกวิจารานํ วูปสมา ฯเปฯ ทุติยํ ฌานํ อุปสมฺปชฺช วิหรามิ.\csromanspacer{} ตสฺส มยฺหํ อานนฺท อิมินา วิหาเรน วิหรโต วิตกฺกสหคตา สญฺญามนสิการา สมุทาจรนฺติ.\csromanspacer{} สฺวสฺส เม โหติ อาพาโธ.\csromanspacer{} เสยฺยถาปิ อานนฺท สุขิโน ทุกฺขํ อุปฺปชฺเชยฺย ยาวเทว อาพาธาย.\csromanspacer{} เอวเมวสฺส เม วิตกฺกสหคตา สญฺญามนสิการา สมุทาจรนฺติ, สฺวสฺส เม โหติ อาพาโธ.}

\prose{ตสฺส มยฺหํ อานนฺท เอตทโหสิ “ยํนูนาหํ ปีติยา จ วิราคา อุเปกฺขโก จ วิหเรยฺยํ สโต จ สมฺปชาโน, สุขญฺจ กาเยน ปฏิสํเวเทยฺยํ, ยํ ตํ อริยา อาจิกฺขนฺติ ‘อุเปกฺขโก สติมา สุขวิหารี’ติ ตติยํ ฌานํ อุปสมฺปชฺช วิหเรยฺยนฺ”ติ.\csromanspacer{} ตสฺส มยฺหํ อานนฺท นิปฺปีติเก จิตฺตํ น ปกฺขนฺทติ นปฺปสีทติ น สนฺติฏฺฐติ น วิมุจฺจติ “เอตํ สนฺตนฺ”ติ ปสฺสโต.\csromanspacer{} ตสฺส มยฺหํ อานนฺท เอตทโหสิ “โก นุ โข เหตุ โก}

\vspace{\tipitakaparskip}
\csromancenter{นวกนิปาตปาฬิ ปจฺจโย, เยน เม นิปฺปีติเก จิตฺตํ น ปกฺขนฺทติ นปฺปสีทติ น สนฺติฏฺฐติ น วิมุจฺจติ ‘เอตํ สนฺตนฺ’ติ ปสฺสโต”.\csromanspacer{} ตสฺส มยฺหํ อานนฺท เอตทโหสิ “ปีติยา โข เม อาทีนโว อทิฏฺโฐ, โส จ เม อพหุลีกโต.\csromanspacer{} นิปฺปีติเก จ อานิสํโส อนธิคโต, โส จ เม อนาเสวิโต.\csromanspacer{} ตสฺมา เม นิปฺปีติเก จิตฺตํ น ปกฺขนฺทติ นปฺปสีทติ น สนฺติฏฺฐติ น วิมุจฺจติ ‘เอตํ สนฺตนฺ’ติ ปสฺสโต”.\csromanspacer{} ตสฺส มยฺหํ อานนฺท เอตทโหสิ “สเจ โข อหํ ปีติยา อาทีนวํ ทิสฺวา ตํ พหุลํ กเรยฺยํ, นิปฺปีติเก อานิสํสํ อธิคมฺม ตมาเสเวยฺยํ, ฐานํ โข ปเนตํ วิชฺชติ, ยํ เม นิปฺปีติเก จิตฺตํ ปกฺขนฺเทยฺย ปสีเทยฺย สนฺติฏฺเฐยฺย วิมุจฺเจยฺย ‘เอตํ สนฺตนฺ’ติ ปสฺสโต”.\csromanspacer{} โส โข อหํ อานนฺท อปเรน สมเยน ปีติยา อาทีนวํ ทิสฺวา ตํ พหุลมกาสิํ, นิปฺปีติเก อานิสํสํ อธิคมฺม ตมาเสวิํ.\csromanspacer{} ตสฺส มยฺหํ อานนฺท นิปฺปีติเก จิตฺตํ ปกฺขนฺทติ ปสีทติ สนฺติฏฺฐติ วิมุจฺจติ “เอตํ สนฺตนฺ”ติ ปสฺสโต.\csromanspacer{} โส โข อหํ อานนฺท ปีติยา จ วิราคา ฯเปฯ ตติยํ ฌานํ อุปสมฺปชฺช วิหรามิ.\csromanspacer{} ตสฺส มยฺหํ อานนฺท อิมินา วิหาเรน วิหรโต ปีติสหคตา สญฺญามนสิการา สมุทาจรนฺติ, สฺวสฺส เม โหติ อาพาโธ.\csromanspacer{} เสยฺยถาปิ อานนฺท สุขิโน ทุกฺขํ อุปฺปชฺเชยฺย ยาวเทว อาพาธาย.\csromanspacer{} เอวเมวสฺส เม ปีติสหคตา สญฺญามนสิการา สมุทาจรนฺติ, สฺวสฺส เม โหติ อาพาโธ.}
\prose{\csromanfolio{236}ตสฺส มยฺหํ อานนฺท เอตทโหสิ “ยํนูนาหํ สุขสฺส จ ปหานา ทุกฺขสฺส จ ปหานา ปุพฺเพว โสมนสฺสโทมนสฺสานํ อตฺถงฺคมา อทุกฺขมสุขํ อุเปกฺขาสติปาริสุทฺธิํ จตุตฺถํ ฌานํ อุปสมฺปชฺช วิหเรยฺยนฺ”ติ.\csromanspacer{} ตสฺส มยฺหํ อานนฺท อทุกฺขมสุเข จิตฺตํ น ปกฺขนฺทติ นปฺปสีทติ น สนฺติฏฺฐติ น วิมุจฺจติ “เอตํ สนฺตนฺ”ติ ปสฺสโต.\csromanspacer{} ตสฺส มยฺหํ อานนฺท เอตทโหสิ “โก นุ โข เหตุ โก ปจฺจโย, เยน เม อทุกฺขมสุเข จิตฺตํ น ปกฺขนฺทติ นปฺปสีทติ น สนฺติฏฺฐติ น วิมุจฺจติ ‘เอตํ สนฺตนฺ’ติ ปสฺสโต”.\csromanspacer{} ตสฺส มยฺหํ อานนฺท เอตทโหสิ “อุเปกฺขาสุเข โข เม อาทีนโว อทิฏฺโฐ, โส จ เม อพหุลีกโต.\csromanspacer{} อทุกฺขมสุเข จ อานิสํโส อนธิคโต, โส จ เม อนาเสวิโต.\csromanspacer{} ตสฺมา เม อทุกฺขมสุเข จิตฺตํ น ปกฺขนฺทติ นปฺปสีทติ น สนฺติฏฺฐติ น วิมุจฺจติ ‘เอตํ สนฺตนฺ’ติ ปสฺสโต”.\csromanspacer{} ตสฺส มยฺหํ อานนฺท เอตทโหสิ “สเจ โข อหํ อุเปกฺขาสุเข อาทีนวํ ทิสฺวา ตํ พหุลํ กเรยฺยํ, อทุกฺขมสุเข อานิสํสํ อธิคมฺม ตมาเสเวยฺยํ.}

\chapterheadpageread{4. มหาวคฺค}
\prose{\csromanfolio{237}ฐานํ โข ปเนตํ วิชฺชติ, ยํ เม อทุกฺขมสุเข จิตฺตํ ปกฺขนฺเทยฺย ปสีเทยฺย สนฺติฏฺเฐยฺย วิมุจฺเจยฺย ‘เอตํ สนฺตนฺ’ติ ปสฺสโต”.\csromanspacer{} โส โข อหํ อานนฺท อปเรน สมเรน อุเปกฺขาสุเข อาทีนวํ ทิสฺวา ตํ พหุลมกาสิํ, อทุกฺขมสุเข อานิสํสํ อธิคมฺม ตมาเสวิํ.\csromanspacer{} ตสฺส มยฺหํ อานนฺท อทุกฺขมสุเข จิตฺตํ ปกฺขนฺทติ ปสีทติ สนฺติฏฺฐติ วิมุจฺจติ “เอตํ สนฺตนฺ”ติ ปสฺสโต.\csromanspacer{} โส โข อหํ อานนฺท สุขสฺส จ ปหานา ฯเปฯ จตุตฺถํ ฌานํ อุปสมฺปชฺช วิหรามิ.\csromanspacer{} ตสฺส มยฺหํ อานนฺท อิมินา วิหาเรน วิหรโต อุเปกฺขาสหคตา สญฺญามนสิการา สมุทาจรนฺติ, สฺวสฺส เม โหติ อาพาโธ.\csromanspacer{} เสยฺยถาปิ อานนฺท สุขิโน ทุกฺขํ อุปฺปชฺเชยฺย ยาวเทว อาพาธาย.\csromanspacer{} เอวเมวสฺส เม อุเปกฺขาสหคตา สญฺญามนสิการา สมุทาจรนฺติ, สฺวสฺส เม โหติ อาพาโธ.}

\prose{ตสฺส มยฺหํ อานนฺท เอตทโหสิ “ยํนูนาหํ สพฺพโส รูปสญฺญานํ สมติกฺกมา ปฏิฆสญฺญานํ อตฺถงฺคมา นานตฺตสญฺญานํ อมนสิการา ‘อนนฺโต อากาโส’ติ อากาสานญฺจายตนํ อุปสมฺปชฺช วิหเรยฺยนฺ”ติ.\csromanspacer{} ตสฺส มยฺหํ อานนฺท อากาสานญฺจายตเน จิตฺตํ น ปกฺขนฺทติ นปฺปสีทติ น สนฺติฏฺฐติ น วิมุจฺจติ “เอตํ สนฺตนฺ”ติ ปสฺสโต.\csromanspacer{} ตสฺส มยฺหํ อานนฺท เอตทโหสิ “โก นุโข เหตุ โก ปจฺจโย, เยน เม อากาสานญฺจายตเน จิตฺตํ น ปกฺขนฺทติ นปฺปสีทติ น สนฺติฏฺฐติ น วิมุจฺจติ ‘เอตํ สนฺตนฺ’ติ ปสฺสโต”.\csromanspacer{} ตสฺส มยฺหํ อานนฺท เอตทโหสิ “รูเปสุ โข เม อาทีนโว อทิฏฺโฐ, โส จ อพหุลีกโต.\csromanspacer{} อากาสานญฺจายตเน จ อานิสํโส อนธิคโต, โส จ เม อนาเสวิโต.\csromanspacer{} ตสฺมา เม อากาสานญฺจายตเน จิตฺตํ น ปกฺขนฺทติ นปฺปสีทติ น สนฺติฏฺฐติ น วิมุจฺจติ ‘เอตํ สนฺตนฺ’ติ ปสฺสโต”.\csromanspacer{} ตสฺส มยฺหํ อานนฺท เอตทโหสิ “สเจ โข อหํ รูเปสุ อาทีนวํ ทิสฺวา ตํ พหุลํ กเรยฺยํ, อากาสานญฺจายตเน อานิสํสํ อธิคมฺม ตมาเสเวยฺยํ.\csromanspacer{} ฐานํ โข ปเนตํ วิชฺชติ, ยํ เม อากาสานญฺจายตเน จิตฺตํ ปกฺขนฺเทยฺย ปสีเทยฺย สนฺติฏฺเฐยฺย วิมุจฺเจยฺย ‘เอตํ สนฺตนฺ’ติ ปสฺสโต.\csromanspacer{} โส โข อหํ อานนฺท อปเรน สมเยน รูเปสุ อาทีนวํ ทิสฺวา ตํ พหุลมกาสิํ, อากาสานญฺจายตเน อานิสํสํ อธิคมฺม ตมาเสวิํ.\csromanspacer{} ตสฺส มยฺหํ อานนฺท อากาสานญฺจายตเน จิตฺตํ ปกฺขนฺทติ ปสีทติ สนฺติฏฺฐติ วิมุจฺจติ ‘เอตํ สนฺตนฺ’ติ ปสฺสโต”.\csromanspacer{} โส โข อหํ อานนฺท สพฺพโส รูปสญฺญานํ}

\vspace{\tipitakaparskip}
\csromancenter{นวกนิปาตปาฬิ สมติกฺกมา ปฏิฆสญฺญานํ อตฺถงฺคมา นานตฺตสญฺญานํ อมนสิการา “อนนฺโต อากาโส”ติ อากาสานญฺจายตนํ อุปสมฺปชฺช วิหรามิ.\csromanspacer{} ตสฺส มยฺหํ อานนฺท อิมินา วิหาเรน วิหรโต รูปสหคตา สญฺญามนสิการา สมุทาจรนฺติ, สฺวสฺส เม โหติ อาพาโธ.\csromanspacer{} เสยฺยถาปิ อานนฺท สุขิโน ทุกฺขํ อุปฺปชฺเชยฺย ยาวเทว อาพาธาย.\csromanspacer{} เอวเมวสฺส เม รูปสหคตา สญฺญามนสิการา สมุทาจรนฺติ, สฺวสฺส เม โหติ อาพาโธ.}
\prose{\csromanfolio{238}ตสฺส มยฺหํ อานนฺท เอตทโหสิ “ยํนูนาหํ สพฺพโส อากาสานญฺจายตนํ สมติกฺกมฺม ‘อนนฺตํ วิญฺญาณนฺ’ติ วิญฺญาณญฺจายตนํ อุปสมฺปชฺช วิหเรยฺยนฺ”ติ.\csromanspacer{} ตสฺส มยฺหํ อานนฺท วิญฺญาณญฺจายตเน จิตฺตํ น ปกฺขนฺทติ นปฺปสีทติ น สนฺติฏฺฐติ น วิมุจฺจติ “เอตํ สนฺตนฺ”ติ ปสฺสโต.\csromanspacer{} ตสฺส มยฺหํ อานนฺท เอตทโหสิ “โก นุ โข เหตุ โก ปจฺจโย, เยน เม วิญฺญาณญฺจายตเน จิตฺตํ น ปกฺขนฺทติ นปฺปสีทติ น สนฺติฏฺฐติ น วิมุจฺจติ ‘เอตํ สนฺตนฺ’ติ ปสฺสโต”.\csromanspacer{} ตสฺส มยฺหํ อานนฺท เอตทโหสิ “อากาสานญฺจายตเน โข เม อาทีนโว อทิฏฺโฐ, โส จ อพหุลีกโต.\csromanspacer{} วิญฺญาณญฺจายตเน จ อานิสํโส อนธิคโต, โส จ เม อนาเสวิโต.\csromanspacer{} ตสฺมา เม วิญฺญาณญฺจายตเน จิตฺตํ น ปกฺขนฺทติ นปฺปสีทติ น สนฺติฏฺฐติ น วิมุจฺจติ ‘เอตํ สนฺตนฺ’ติ ปสฺสโต”.\csromanspacer{} ตสฺส มยฺหํ อานนฺท เอตทโหสิ “สเจ โข อหํ อากาสานญฺจายตเน อาทีนวํ ทิสฺวา ตํ พหุลํ กเรยฺยํ, วิญฺญาณญฺจายตเน อานิสํสํ อธิคมฺม ตมาเสเวยฺยํ.\csromanspacer{} ฐานํ โข ปเนตํ วิชฺชติ, ยํ เม วิญฺญาณญฺจายตเน จิตฺตํ ปกฺขนฺเทยฺย ปสีเทยฺย สนฺติฏฺเฐยฺย วิมุจฺเจยฺย ‘เอตํ สนฺตนฺ’ติ ปสฺสโต”.\csromanspacer{} โส โข อหํ อานนฺท อปเรน สมเยน อากาสานญฺจายตเน อาทีนวํ ทิสฺวา ตํ พหุลมกาสิํ, วิญฺญาณญฺจายตเน อานิสํสํ อธิคมฺม ตมาเสวิํ.\csromanspacer{} ตสฺส มยฺหํ อานนฺท วิญฺญาณญฺจายตเน จิตฺตํ ปกฺขนฺทติ ปสีทติ สนฺติฏฺฐติ วิมุจฺจติ “เอตํ สนฺตนฺ’ติ ปสฺสโต”.\csromanspacer{} โส โข อหํ อานนฺท สพฺพโส อากาสานญฺจายตนํ สมติกฺกมฺม “อนนฺตํ วิญฺญาณ นฺ”ติ วิญฺญาณญฺจายตนํ อุปสมฺปชฺช วิหรามิ.\csromanspacer{} ตสฺส มยฺหํ อานนฺท อิมินา วิหาเรน วิหรโต อากาสานญฺจายตนสหคตา สญฺญามนสิการา สมุทาจรนฺติ, สฺวสฺส เม โหติ อาพาโธ.\csromanspacer{} เสยฺยถาปิ อานนฺท สุขิโน ทุกฺขํ อุปฺปชฺเชยฺย ยาวเทว}

\chapterheadpageread{4. มหาวคฺค}
\prose{\csromanfolio{239}อาพาธาย.\csromanspacer{} เอวเมวสฺส เม อากาสานญฺจายตนสหคตา สญฺญามนสิการา สมุทาจรนฺติ, สฺวสฺส เม โหติ อาพาโธ.}

\prose{ตสฺส มยฺหํ อานนฺท เอตทโหสิ “ยํนูนาหํ สพฺพโส วิญฺญาณญฺจายตนํ สมติกฺกมฺม “นตฺถิ กิญฺจี”ติ อากิญฺจญฺญายตนํ อุปสมฺปชฺช วิหเรยฺยนฺ”ติ.\csromanspacer{} ตสฺส มยฺหํ อานนฺท อากิญฺจญฺญายตเน จิตฺตํ น ปกฺขนฺทติ นปฺปสีทติ น สนฺติฏฺฐติ น วิมุจฺจติ “เอตํ สนฺตนฺ”ติ ปสฺสโต.\csromanspacer{} ตสฺส มยฺหํ อานนฺท เอตทโหสิ “โก นุ โข เหตุ โก ปจฺจโย, เยน เม อากิญฺจญฺญายตเน จิตฺตํ น ปกฺขนฺทติ นปฺปสีทติ น สนฺติฏฺฐติ น วิมุจฺจติ ‘เอตํ สนฺตนฺ’ติ ปสฺสโต”.\csromanspacer{} ตสฺส มยฺหํ อานนฺท เอตทโหสิ “วิญฺญาณญฺจายตเน โข เม อาทีนโว อทิฏฺโฐ, โส จ เม อพหุลีกโต.\csromanspacer{} อากิญฺจญฺญายตเน จ อานิสํโส อนธิคโต, โส จ เม อนาเสวิโต.\csromanspacer{} ตสฺมา เม อากิญฺจญฺญายตเน จิตฺตํ น ปกฺขนฺทติ นปฺปสีทติ น สนฺติฏฺฐติ น วิมุจฺจติ ‘เอตํ สนฺตนฺ’ติ ปสฺสโต”.\csromanspacer{} ตสฺส มยฺหํ อานนฺท เอตทโหสิ “สเจ โข อหํ วิญฺญาณญฺจายตเน อาทีนวํ ทิสฺวา ตํ พหุลํ กเรยฺยํ, อากิญฺจญฺญายตเน อานิสํสํ อธิคมฺม ตมาเสเวยฺยํ.\csromanspacer{} ฐานํ โข ปเนตํ วิชฺชติ, ยํ เม อากิญฺจญฺญายตเน จิตฺตํ ปกฺขนฺเทยฺย ปสีเทยฺย สนฺติฏฺเฐยฺย วิมุจฺเจยฺย ‘เอตํ สนฺตนฺ’ติ ปสฺสโต”.\csromanspacer{} โส โข อหํ อานนฺท อปเรน สมเยน วิญฺญาณญฺจายตเน อาทีนวํ ทิสฺวา ตํ พหุลมกาสิํ, อากิญฺจญฺญายตเน อานิสํสํ อธิคมฺม ตมาเสวิํ.\csromanspacer{} ตสฺส มยฺหํ อานนฺท อากิญฺจญฺญายตเน จิตฺตํ ปกฺขนฺทติ ปสีทติ สนฺติฏฺฐติ วิมุจฺจติ “เอตํ สนฺตนฺ”ติ ปสฺสโต.\csromanspacer{} โส โข อหํ อานนฺท สพฺพโส วิญฺญาณญฺจายตนํ สมติกฺกมฺม “นตฺถิ กิญฺจี”ติ อากิญฺจญฺญายตนํ อุปสมฺปชฺช วิหรามิ.\csromanspacer{} ตสฺส มยฺหํ อานนฺท อิมินา วิหาเรน วิหรโต วิญฺญาณญฺจายตนสหคตา สญฺญามนสิการา สมุทาจรนฺติ, สฺวสฺส เม โหติ อาพาโธ.\csromanspacer{} เสยฺยถาปิ อานนฺท สุขิโน ทุกฺขํ อุปฺปชฺเชยฺย ยาวเทว อาพาธาย.\csromanspacer{} เอวเมวสฺส เม วิญฺญาณญฺจายตนสหคตา สญฺญามนสิการา สมุทาจรนฺติ, สฺวสฺส เม โหติ อาพาโธ.}

\prose{ตสฺส มยฺหํ อานนฺท เอตทโหสิ “ยํนูนาหํ สพฺพโส อากิญฺจญฺญายตนํ สมติกฺกมฺม เนวสญฺญานาสญฺญายตนํ อุปสมฺปชฺช วิหเรยฺยนฺ”ติ.\csromanspacer{} ตสฺส มยฺหํ อานนฺท เนวสญฺญานาสญฺญายตเน จิตฺตํ น ปกฺขนฺทติ นปฺปสีทติ น สนฺติฏฺฐติ น วิมุจฺจติ “เอตํ สนฺตนฺ”ติ ปสฺสโต.\csromanspacer{} ตสฺส มยฺหํ อานนฺท}

\vspace{\tipitakaparskip}
\csromancenter{นวกนิปาตปาฬิ เอตทโหสิ “โก นุ โข เหตุ โก ปจฺจโย, เยน เม เนวสญฺญานาสญฺญายตเน จิตฺตํ น ปกฺขนฺทติ นปฺปสีทติ น สนฺติฏฺฐติ น วิมุจฺจติ ‘เอตํ สนฺตนฺ’ติ ปสฺสโต”.\csromanspacer{} ตสฺส มยฺหํ อานนฺท เอตทโหสิ “อากิญฺจญฺญายตเน โข เม อาทีนโว อทิฏฺโฐ, โส จ เม อพหุลีกโต.\csromanspacer{} เนวสญฺญานาสญฺญายตเน จ อานิสํโส อนธิคโต, โส จ เม อนาเสวิโต.\csromanspacer{} ตสฺมา เม เนวสญฺญานาสญฺญายตเน จิตฺตํ น ปกฺขนฺทติ นปฺปสีทติ น สนฺติฏฺฐติ น วิมุจฺจติ ‘เอตํ สนฺตนฺ’ติ ปสฺสโต”.\csromanspacer{} ตสฺส มยฺหํ อานนฺท เอตทโหสิ “สเจ โข อหํ อากิญฺจญฺญายตเน อาทีนวํ ทิสฺวา ตํ พหุลํ กเรยฺยํ, เนวสญฺญานาสญฺญายตเน อานิสํสํ อธิคมฺม ตมาเสเวยฺยํ.\csromanspacer{} ฐานํ โข ปเนตํ วิชฺชติ, ยํ เม เนวสญฺญานาสญฺญายตเน จิตฺตํ ปกฺขนฺเทยฺย ปสีเทยฺย สนฺติฏฺเฐยฺย วิมุจฺเจยฺย ‘เอตํ สนฺตนฺ’ติ ปสฺสโต”.\csromanspacer{} โส โข อหํ อานนฺท อปเรน สมเยน อากิญฺจญฺญายตเน อาทีนวํ ทิสฺวา ตํ พหุลมกาสิํ, เนวสญฺญานาสญฺญายตเน อานิสํสํ อธิคมฺม ตมาเสวิํ.\csromanspacer{} ตสฺส มยฺหํ อานนฺท เนวสญฺญานาสญฺญายตเน จิตฺตํ ปกฺขนฺทติ ปสีทติ สนฺติฏฺฐติ วิมุจฺจติ “เอตํ สนฺตนฺ”ติ ปสฺสโต.\csromanspacer{} โส โข อหํ อานนฺท สพฺพโส อากิญฺจญฺญายตนํ สมติกฺกมฺม เนวสญฺญานาสญฺญายตนํ อุปสมฺปชฺช วิหรามิ.\csromanspacer{} ตสฺส มยฺหํ อานนฺท อิมินา วิหาเรน วิหรโต อากิญฺจญฺญายตนสหคตา สญฺญามนสิการา สมุทาจรนฺติ, สฺวสฺส เม โหติ อาพาโธ.\csromanspacer{} เสยฺยถาปิ อานนฺท สุขิโน ทุกฺขํ อุปฺปชฺเชยฺย ยาวเทว อาพาธาย.\csromanspacer{} เอวเมวสฺส เม อากิญฺจญฺญายตนสหคตา สญฺญามนสิการา สมุทาจรนฺติ, สฺวสฺส เม โหติ อาพาโธ.}
\prose{\csromanfolio{240}ตสฺส มยฺหํ อานนฺท เอตทโหสิ “ยํนูนาหํ เนวสญฺญานาสญฺญายตนํ สมติกฺกมฺม สญฺญาเวทยิตนิโรธํ อุปสมฺปชฺช วิหเรยฺยนฺ”ติ.\csromanspacer{} ตสฺส มยฺหํ อานนฺท สญฺญาเวทยิตนิโรเธ จิตฺตํ น ปกฺขนฺทติ นปฺปสีทติ น สนฺติฏฺฐติ น วิมุจฺจติ “เอตํ สนฺตนฺ”ติ ปสฺสโต.\csromanspacer{} ตสฺส มยฺหํ อานนฺท เอตทโหสิ “โก นุ โข เหตุ โก ปจฺจโย, เยน เม สญฺญาเวทยิตนิโรเธ จิตฺตํ น ปกฺขนฺทติ นปฺปสีทติ น สนฺติฏฺฐติ น วิมุจฺจติ ‘เอตํ สนฺตนฺ’ติ ปสฺสโต”.\csromanspacer{} ตสฺส มยฺหํ อานนฺท เอตทโหสิ “เนวสญฺญานาสญฺญายตเน โข เม อาทีนโว อทิฏฺโฐ, โส จ เม อพหุลีกโต.\csromanspacer{} สญฺญาเวทยิตนิโรเธ จ อานิสํโส อนธิคโต,}

\chapterheadpageread{4. มหาวคฺค}
\prose{\csromanfolio{241}โส จ เม อนาเสวิโต.\csromanspacer{} ตสฺมา เม สญฺญาเวทยิตนิโรเธ จิตฺตํ น ปกฺขนฺทติ นปฺปสีทติ น สนฺติฏฺฐติ น วิมุจฺจติ ‘เอตํ สนฺตนฺ’ติ ปสฺสโต”.\csromanspacer{} ตสฺส มยฺหํ อานนฺท เอตทโหสิ “สเจ โข อหํ เนวสญฺญานาสญฺญายตเน อาทีนวํ ทิสฺวา ตํ พหุลํ กเรยฺยํ, สญฺญาเวทยิตนิโรเธ อานิสํสํ อธิคมฺม ตมาเสเวยฺยํ.\csromanspacer{} ฐานํ โข ปเนตํ วิชฺชติ, ยํ เม สญฺญาเวทยิตนิโรเธ จิตฺตํ ปกฺขนฺเทยฺย ปสีเทยฺย สนฺติฏฺเฐยฺย วิมุจฺเจยฺย ‘เอตํ สนฺตนฺ’ติ ปสฺสโต”.\csromanspacer{} โส โข อหํ อานนฺท อปเรน สมเยน เนวสญฺญานาสญฺญายตเน อาทีนวํ ทิสฺวา ตํ พหุลมกาสิํ, สญฺญาเวทยิตนิโรเธ อานิสํสํ อธิคมฺม ตมาเสวิํ.\csromanspacer{} ตสฺส มยฺหํ อานนฺท สญฺญาเวทยิตนิโรเธ จิตฺตํ ปกฺขนฺทติ ปสีทติ สนฺติฏฺฐติ วิมุจฺจติ “เอตํ สนฺตนฺ”ติ ปสฺสโต.\csromanspacer{} โส โข อหํ อานนฺท สพฺพโส เนวสญฺญานาสญฺญายตนํ สมติกฺกมฺม สญฺญาเวทยิตนิโรธํ อุปสมฺปชฺช วิหรามิ, ปญฺญาย จ เม ทิสฺวา อาสวา ปริกฺขยํ อคมํสุ.}

\prose{ยาวกีวญฺจาหํ อานนฺท อิมา นว อนุปุพฺพวิหารสมาปตฺติโย น เอวํ อนุโลมปฏิโลมํ สมาปชฺชีมฺปิ วุฏฺฐหิมฺปิ, เนว ตาวาหํ อานนฺท สเทวเก โลเก สมารเก สพฺรหฺมเก สสฺสมณพฺราหฺมณิยา ปชาย สเทวมนุสฺสาย อนุตฺตรํ สมฺมาสมฺโพธิํ อภิสมฺพุทฺโธติ ปจฺจญฺญาสิํ.\csromanspacer{} ยโต จ โข อหํ อานนฺท อิมา นว อนุปุพฺพวิหารสมาปตฺติโย เอวํ อนุโลมปฏิโลมํ สมาปชฺชิมฺปิ วุฏฺฐหิมฺปิ, อถาหํ อานนฺท สเทวเก โลเก สมารเก สพฺรหฺมเก สสฺสมณพฺราหฺมณิยา ปชาย สเทวมนุสฺสาย อนุตฺตรํ สมฺมาสมฺโพธิํ อภิสมฺพุทฺโธติ ปจฺจญฺญาสิํ, ญาณญฺจ ปน เม ทสฺสนํ อุทปาทิ “อกุปฺปา เม เจโตวิมุตฺติ\footnote{วิมุตฺติ (ก-สี, ก)}, อยมนฺติมา ชาติ, นตฺถิ ทานิ ปุนพฺภโว”ติ.\csromanspacer{}\csromanspacer{}ทสฺสมํ.\csromanspacer{} มหาวคฺโค จตุตฺโถ.}

\titlehead{\textbf{ตสฺสุทฺทานํ}}
\csromangathameasure{เทฺว วิหารา จ นิพฺพานํ, คาวี ฌาเนน ปญฺจมํ. \\ อานนฺโท พฺราหฺมณา เทโว, นาเคน ตปุสฺเสน จาติ.}
\csromangathagroup{\csromangathastanza{เทฺว วิหารา จ นิพฺพานํ, คาวี ฌาเนน ปญฺจมํ. \\ อานนฺโท พฺราหฺมณา เทโว, นาเคน ตปุสฺเสน จาติ.}}

\gambhira{นวกนิปาตปาฬิ \textbf{5. สามญฺญวคฺค}}
\csromanheader{2}{1. สมฺพาธสุตฺต}
\proseitem{42}{\csromanfolio{242}เอกํ สมยํ อายสฺมา อานนฺโท โกสมฺพิยํ วิหรติ โฆสิตาราเม.\csromanspacer{} อถ โข อายสฺมา อุทายี เยนายสฺมา อานนฺโท เตนุปสงฺกมิ, อุปสงฺกมิตฺวา อายสฺมตา อานนฺเทน สทฺธิํ สมฺโมทิ, สมฺโมทนียํ กถํ สารณียํ วีติสาเรตฺวา เอกมนฺตํ นิสีทิ, เอกมนฺตํ นิสินฺโน โข อายสฺมา อุทายี อายสฺมนฺตํ อานนฺทํ เอตทโวจ\csromandash{}วุตฺตมิทํ อาวุโส ปญฺจาลจณฺเฑน เทวปุตฺเตน.}

\csromangathameasure{“สมฺพาเธ คตํ โอกาสํ, อวิทฺวา ภูริเมธโส. \\ โย ฌานมพุชฺฌิ พุทฺโธ, ปฏิลีนนิสโภ มุนี”ติ.}
\csromangathagroup{\csromangathastanza{“สมฺพาเธ คตํ\footnote{สมฺพาเธ วต (สี)} โอกาสํ, อวิทฺวา ภูริเมธโส. \\ โย ฌานมพุชฺฌิ พุทฺโธ, ปฏิลีนนิสโภ มุนี”ติ.}}

\prose{กตโม อาวุโส สมฺพาโธ, กตโม สมฺพาเธ โอกาสาธิคโม วุตฺโต ภควตาติ.\csromanspacer{} ปญฺจิเม อาวุโส กามคุณา สมฺพาโธ วุตฺโต ภควตา.\csromanspacer{} กตเม ปญฺจ? จกฺขุวิญฺเญยฺยา รูปา อิฏฺฐา กนฺตา มนาปา ปิยรูปา กามูปสํหิตา รชนียา.\csromanspacer{} โสตวิญฺเญยฺยา สทฺทา ฯเปฯ.\csromanspacer{} ฆานวิญฺเญยฺยา คนฺธา.\csromanspacer{} ชิวฺหาวิญฺเญยฺยา รสา.\csromanspacer{} กายวิญฺเญยฺยา โผฏฺฐพฺพา อิฏฺฐา กนฺตา มนาปา ปิยรูปา กามูปสํหิตา รชนียา.\csromanspacer{} อิเม โข อาวุโส ปญฺจ กามคุณา สมฺพาโธ วุตฺโต ภควตา.}

\prose{อิธาวุโส ภิกฺขุ วิวิจฺเจว กาเมหิ ฯเปฯ ปฐมํ ฌานํ อุปสมฺปชฺช วิหรติ.\csromanspacer{} เอตฺตาวตาปิ โข อาวุโส สมฺพาเธ โอกาสาธิคโม วุตฺโต ภควตา ปริยาเยน.\csromanspacer{} ตตฺราป’ตฺถิ สมฺพาโธ, กิญฺจ ตตฺถ สมฺพาโธ, ยเทว ตตฺถ วิตกฺกวิจารา อนิรุทฺธา โหนฺติ.\csromanspacer{} อยเมตฺถ สมฺพาโธ.}

\prose{ปุน จปรํ อาวุโส ภิกฺขุ วิตกฺกวิจารานํ วูปสมา ฯเปฯ ทุติยํ ฌานํ อุปสมฺปชฺช วิหรติ.\csromanspacer{} เอตฺตาวตาปิ โข อาวุโส สมฺพาเธ โอกาสาธิคโม วุตฺโต ภควตา ปริยาเยน.\csromanspacer{} ตตฺรป’ตฺถิ สมฺพาโธ, กิญฺจ ตตฺถ สมฺโพโธ, ยเทว ตตฺถ ปีติ อนิรุทฺธา โหติ.\csromanspacer{} อยเมตฺถ สมฺพาโธ.}

\chapterheadpageread{5. สามญฺญวคฺค}
\prose{\csromanfolio{243}ปุน จปรํ อาวุโส ภิกฺขุ ปีติยา จ วิราคา ฯเปฯ ตติยํ ฌานํ อุปสมฺปชฺช วิหรติ.\csromanspacer{} เอตฺตาวตาปิ โข อาวุโส สมฺพาเธ โอกาสาธิคโม วุตฺโต ภควตา ปริยาเยน.\csromanspacer{} ตตฺราป’ตฺถิ สมฺพาโธ, กิญฺจ ตตฺถ สมฺพาโธ, ยเทว ตตฺถ อุเปกฺขาสุขํ อนิรุทฺธํ โหติ.\csromanspacer{} อยเมตฺถ สมฺพาโธ.}

\prose{ปุน จปรํ อาวุโส ภิกฺขุ สุขสฺส จ ปหานา ฯเปฯ จตุตฺถํ ฌานํ อุปสมฺปชฺช วิหรติ.\csromanspacer{} เอตฺตาวตาปิ โข อาวุโส สมฺพาเธ โอกาสาธิคโม วุตฺโต ภควตา ปริยาเยน.\csromanspacer{} ตตฺราปตฺถิ สมฺพาโธ, กิญฺจ ตตฺถ สมฺพาโธ, ยเทว ตตฺถ รูปสญฺญา อนิรุทฺธา โหติ.\csromanspacer{} อยเมตฺถ สมฺพาโธ.}

\prose{ปุน จปรํ อาวุโส ภิกฺขุ สพฺพโส รูปสญฺญานํ สมติกฺกมา ปฏิฆสญฺญานํ อตฺถงฺคมา นานตฺตสญฺญานํ อมนสิการา “อนนฺโต อากาโส”ติ อากาสานญฺจายตนํ อุปสมฺปชฺช วิหรติ.\csromanspacer{} เอตฺตาวตาปิ โข อาวุโส สมฺพาเธ โอกาสาธิคโม วุตฺโต ภควตา ปริยาเยน.\csromanspacer{} ตตฺราป’ตฺถิ สมฺพาโธ, กิญฺจ ตตฺถ สมฺพาโธ, ยเทว ตตฺถ อากาสานญฺจายตนสญฺญา อนิรุทฺธา โหติ.\csromanspacer{} อยเมตฺถ สมฺพาโธ.}

\prose{ปุน จปรํ อาวุโส ภิกฺขุ สพฺพโส อากาสานญฺจายตนํ สมติกฺกมฺม “อนนฺตํ วิญฺญาณนฺ”ติ วิญฺญาณญฺจายตนํ อุปสมฺปชฺช วิหรติ.\csromanspacer{} เอตฺตาวตาปิ โข อาวุโส สมฺพาเธ โอกาสาธิคโม วุตฺโต ภควตา ปริยาเยน.\csromanspacer{} ตตฺราป’ตฺถิ สมฺโพโธ, กิญฺจ ตตฺถ สมฺพาโธ ยเทว ตตฺถ วิญฺญาณญฺจายตนสญฺญา อนิรุทฺธา โหติ.\csromanspacer{} อยเมตฺถ สมฺพาโธ.}

\prose{ปุน จปรํ อาวุโส ภิกฺขุ สพฺพโส วิญฺญาณญฺจายตนํ สมติกฺกมฺม “นตฺถิ กิญฺจี”ติ อากิญฺจญฺญายตนํ อุปสมฺปชฺช วิหรติ.\csromanspacer{} เอตฺตาวตาปิ โข อาวุโส สมฺพาเธ โอกาสาธิคโม วุตฺโต ภควตา ปริยาเยน.\csromanspacer{} ตตฺราป’ตฺถิ สมฺโพโธ, กิญฺจ ตตฺถ สมฺพาโธ, ยเทว ตตฺถ อากิญฺจญฺญายตนสญฺญา อนิรุทฺธา โหติ.\csromanspacer{} อยเมตฺถ สมฺพาโธ.}

\prose{ปุน จปรํ อาวุโส ภิกฺขุ สพฺพโส อากิญฺจญฺญายตนํ สมติกฺกมฺม เนวสญฺญานาสญฺญายตนํ อุปสมฺปชฺช วิหรติ.\csromanspacer{} เอตฺตาวตาปิ โข อาวุโส สมฺพาเธ โอกาสาธิคโม วุตฺโต ภควตา ปริยาเยน.}

\vspace{\tipitakaparskip}
\csromancenter{นวกนิปาตปาฬิ ตตฺราป’ตฺถิ สมฺพาโธ, กิญฺจ ตตฺถ สมฺพาโธ, ยเทว ตตฺถ เนวสญฺญานาสญฺญายตนสญฺญา อนิรุทฺธา โหติ.\csromanspacer{} อยเมตฺถ สมฺโพโธ.}
\prose{\csromanfolio{244}ปุน จปรํ อาวุโส ภิกฺขุ สพฺพโส เนวสญฺญานาสญฺญายตนํ สมติกฺกมฺม สญฺญาเวทยิตนิโรธํ อุปสมฺปชฺช วิหรติ, ปญฺญาย จสฺส ทิสฺวา อาสวา ปริกฺขีณา โหนฺติ.\csromanspacer{} เอตฺตาวตาปิ โข อาวุโส สมฺพาเธ โอกาสาธิคโม วุตฺโต ภควตา นิปฺปริยาเยนาติ.\csromanspacer{}\csromanspacer{}ปฐมํ.}

\csromanheader{2}{2. กายสกฺขีสุตฺต}
\proseitem{43}{“กายสกฺขี กายสกฺขี”ติ อาวุโส วุจฺจติ, กิตฺตาวตา นุ โข อาวุโส กายสกฺขี วุตฺโต ภควตาติ.\csromanspacer{} อิธาวุโส ภิกฺขุ วิวิจฺเจว กาเมหิ ฯเปฯ ปฐมํ ฌานํ อุปสมฺปชฺช วิหรติ.\csromanspacer{} ยถา ยถา จ ตทายตนํ, ตถา ตถา นํ กาเยน ผุสิตฺวา วิหรติ.\csromanspacer{} เอตฺตาวตาปิ โข อาวุโส กายสกฺขี วุตฺโต ภควตา ปริยาเยน. (1)}

\prose{ปุน จปรํ อาวุโส ภิกฺขุ วิตกฺกวิจารานํ วูปสมา ฯเปฯ ทุติยํ ฌานํ.\csromanspacer{} ตติยํ ฌานํ.\csromanspacer{} จตุตฺถํ ฌานํ อุปสมฺปชฺช วิหรติ.\csromanspacer{} ยถา ยถา จ ตทายตนํ, ตถา ตถา นํ กาเยน ผุสิตฺวา วิหรติ.\csromanspacer{} เอตฺตาวตาปิ โข อาวุโส กายสกฺขี วุตฺโต ภควตา ปริยาเยน. (2-4)}

\prose{ปุน จปรํ อาวุโส ภิกฺขุ สพฺพโส รูปสญฺญานํ สมติกฺกมา ปฏิฆสญฺญานํ อตฺถงฺคมา นานตฺตสญฺญานํ อมนสิการา “อนนฺโต อากาโส”ติ อากาสานญฺจายตนํ อุปสมฺปชฺช วิหรติ.\csromanspacer{} ยถา ยถา จ ตทายตนํ, ตถา ตถา นํ กาเยน ผุสิตฺวา วิหรติ.\csromanspacer{} เอตฺตาวตาปิ โข อาวุโส กายสกฺขี วุตฺโต ภควตา ปริยาเยน ฯเปฯ. (5-8)}

\prose{ปุน จปรํ อาวุโส ภิกฺขุ สพฺพโส เนวสญฺญานาสญฺญายตนํ สมติกฺกมฺม สญฺญาเวทยิตนิโรธํ อุปสมฺปชฺช วิหรติ, ปญฺญาย จสฺส ทิสฺวา อาสวา ปริกฺขีณา โหนฺติ.\csromanspacer{} ยถา ยถา จ ตทายตนํ, ตถา ตถา นํ กาเยน ผุสิตฺวา วิหรติ.\csromanspacer{} เอตฺตาวตาปิ โข อาวุโส กายสกฺขี วุตฺโต ภควตา นิปฺปริยาเยนาติ. (9).\csromanspacer{} ทุติยํ.}

\chapterheadpageread{5. สามญฺญวคฺค}
\csromanheader{2}{3. ปญฺญาวิมุตฺตสุตฺต}
\proseitem{44}{\csromanfolio{245}“ปญฺญาวิมุตฺโต ปญฺญาวิมุตฺโต”ติ อาวุโส วุจฺจติ, กิตฺตาวตา นุ โข อาวุโส ปญฺญาวิมุตฺโต วุตฺโต ภควตาติ.\csromanspacer{} อิธาวุโส ภิกฺขุ วิวิจฺเจว กาเมหิ ฯเปฯ ปฐมํ ฌานํ อุปสมฺปชฺช วิหรติ, ปญฺญาย จ นํ ปชานาติ.\csromanspacer{} เอตฺตาวตาปิ โข อาวุโส ปญฺญาวิมุตฺโต วุตฺโต ภควตา ปริยาเยน ฯเปฯ.}

\prose{ปุน จปรํ อาวุโส ภิกฺขุ สพฺพโส เนวสญฺญานาสญฺญายตนํ สมติกฺกมฺม สญฺญาเวทยิตนิโรธํ อุปสมฺปชฺช วิหรติ, ปญฺญาย จสฺส ทิสฺวา อาสวา ปริกฺขีณา โหนฺติ, ปญฺญาย จ นํ ปชานาติ.\csromanspacer{} เอตฺตาวตาปิ โข อาวุโส ปญฺญาวิมุตฺโต วุตฺโต ภควตา นิปฺปริยาเยนาติ.\csromanspacer{}\csromanspacer{}ตติยํ.}

\csromanheader{2}{4. อุภโตภาควิมุตฺตสุตฺต}
\proseitem{45}{“อุภโตภาควิมุตฺโต อุภโตภาควิมุตฺโต”ติ อาวุโส วุจฺจติ, กิตฺตาวตา นุ โข อาวุโส อุภโตภาควิมุตฺโต วุตฺโต ภควตาติ.\csromanspacer{} อิธาวุโส ภิกฺขุ วิวิจฺเจว กาเมหิ ฯเปฯ ปฐมํ ฌานํ อุปสมฺปชฺช วิหรติ, ยถา ยถา จ ตทายตนํ, ตถา ตถา นํ กาเยน ผุสิตฺวา วิหรติ, ปญฺญาย จ นํ ปชานาติ.\csromanspacer{} เอตฺตาวตาปิ โข อาวุโส อุภโตภาควิมุตฺโต วุตฺโต ภควตา ปริยาเยน ฯเปฯ.}

\prose{ปุน จปรํ อาวุโส ภิกฺขุ สพฺพโส เนวสญฺญานาสญฺญายตนํ สมติกฺกมฺม สญฺญาเวทยิตนิโรธํ อุปสมฺปชฺช วิหรติ, ปญฺญาย จสฺส ทิสฺวา อาสวา ปริกฺขีณา โหนฺติ, ยถา ยถา จ ตทายตนํ, ตถา ตถา นํ กาเยน ผุสิตฺวา วิหรติ, ปญฺญาย จ นํ ปชานาติ.\csromanspacer{} เอตฺตาวตาปิ โข อาวุโส อุภโตภาควิมุตฺโต วุตฺโต ภควตา นิปฺปริยาเยนาติ.\csromanspacer{}\csromanspacer{}จตุตฺถํ.}

\csromanheader{2}{5. สนฺทิฏฺฐิกธมฺมสุตฺต}
\proseitem{46}{“สนฺทิฏฺฐิโก ธมฺโม สนฺทิฏฺฐิโก ธมฺโม”ติ อาวุโส วุจฺจติ, กิตฺตาวตา นุ โข อาวุโส สนฺทิฏฺฐิโก ธมฺโม วุตฺโต ภควตาติ.\csromanspacer{} อิธาวุโส ภิกฺขุ วิวิจฺเจว กาเมหิ ฯเปฯ ปฐมํ ฌานํ อุปสมฺปชฺช วิหรติ.}

\vspace{\tipitakaparskip}
\csromancenter{นวกนิปาตปาฬิ เอตฺตาวภาปิ โข อาวุโส สนฺทิฏฺฐิโก ธมฺโม วุตฺโต ภควตา ปริยาเยน ฯเปฯ.}
\prose{\csromanfolio{246}ปุน จปรํ อาวุโส ภิกฺขุ สพฺพโส เนวสญฺญานาสญฺญายตนํ สมติกฺกมฺม สญฺญาเวทยิตนิโรธํ อุปสมฺปชฺช วิหรติ, ปญฺญาย จสฺส ทิสฺวา อาสวา ปริกฺขีณา โหนฺติ.\csromanspacer{} เอตฺตาวตาปิ โข อาวุโส สนฺทิฏฺฐิโก ธมฺโม วุตฺโต ภควตา นิปฺปริยาเยนาติ.\csromanspacer{}\csromanspacer{}ปญฺจมํ.}

\csromanheader{2}{6. สนฺทิฏฺฐิกนิพฺพานสุตฺต}
\proseitem{47}{“สนฺทิฏฺฐิกํ นิพฺพานํ สนฺทิฏฺฐิกํ นิพฺพานนฺ”ติ อาวุโส วุจฺจติ, กิตฺตาวตา นุ โข อาวุโส สนฺทิฏฺฐิกํ นิพฺพานํ วุตฺตํ ภควตาติ.\csromanspacer{} อิธาวุโส ภิกฺขุ วิวิจฺเจว กาเมหิ ฯเปฯ ปฐมํ ฌานํ อุปสมฺปชฺช วิหรติ.\csromanspacer{} เอตฺตาวตาปิ โข อาวุโส สนฺทิฏฺฐิกํ นิพฺพานํ วุตฺตํ ภควตา ปริยาเยน ฯเปฯ.}

\prose{ปุน จปรํ อาวุโส ภิกฺขุ สพฺพโส เนวสญฺญานาสญฺญายตนํ สมติกฺกมฺม สญฺญาเวทยิตนิโรธํ อุปสมฺปชฺช วิหรติ, ปญฺญาย จสฺส ทิสฺวา อาสวา ปริกฺขีณา โหนฺติ.\csromanspacer{} เอตฺตาวตาปิ โข อาวุโส สนฺทิฏฺฐิกํ นิพฺพานํ วุตฺตํ ภควตา นิปฺปริยาเยนาติ.\csromanspacer{}\csromanspacer{}ฉฏฺฐํ.}

\csromanheader{2}{7. นิพฺพานสุตฺต}
\vspace{\tipitakaparskip}
\csromancenter{“นิพฺพานํ นิพฺพานนฺ”ติ อาวุโส วุจฺจติ ฯเปฯ.\csromanspacer{}\csromanspacer{}สตฺตมํ.}
\csromanheader{2}{8. ปรินิพฺพานสุตฺต}
\vspace{\tipitakaparskip}
\csromancenter{“ปรินิพฺพานํ ปรินิพฺพานนฺ”ติ ฯเปฯ.\csromanspacer{}\csromanspacer{}อฏฺฐมํ.}
\csromanheader{2}{9. ตทงฺคนิพฺพานสุตฺต}
\vspace{\tipitakaparskip}
\csromancenter{“ตทงฺคนิพฺพานํ ตทงฺคนิพฺพานนฺ”ติ อาวุโส วุจฺจติ ฯเปฯ.\csromanspacer{}\csromanspacer{}นวมํ.}
\csromanheader{2}{10. ทิฏฺฐธมฺมนิพฺพานสุตฺต}
\proseitem{51}{“ทิฏฺฐธมฺมนิพฺพานํ ทิฏฺฐธมฺมนิพฺพานนฺ”ติ อาวุโส วุจฺจติ, กิตฺตาวตา นุ โข อาวุโส ทิฏฺฐธมฺมนิพฺพานํ วุตฺตํ ภควตาติ.\csromanspacer{} อิธาวุโส ภิกฺขุ วิวิจฺเจว}

\chapterheadpageread{5. สามญฺญวคฺค}
\prose{\csromanfolio{247}กาเมหิ ฯเปฯ ปฐมํ ฌานํ อุปสมฺปชฺช วิหรติ.\csromanspacer{} เอตฺตาวตาปิ โข อาวุโส ทิฏฺฐธมฺมนิพฺพานํ วุตฺตํ ภควตา ปริยาเยน ฯเปฯ.}

\prose{ปุน จปรํ อาวุโส ภิกฺขุ สพฺพโส เนวสญฺญานาสญฺญายตนํ สมติกฺกมฺม สญฺญาเวทยิตนิโรธํ อุปสมฺปชฺช วิหรติ, ปญฺญาย จสฺส ทิสฺวา อาสวา ปริกฺขีณา โหนฺติ.\csromanspacer{} เอตฺตาวตาปิ โข อาวุโส ทิฏฺฐธมฺมนิพฺพานํ วุตฺตํ ภควตา นิปฺปริยาเยนาติ.\csromanspacer{}\csromanspacer{}ทสมํ.}

\vspace{\tipitakaparskip}
\csromancenter{สามญฺญวคฺโค ปญฺจโม.}
\titlehead{\textbf{ตสฺสุทฺทานํ}}
\csromangathameasure{สมฺพาโธ กายสกฺขี ปญฺญา, อุภโตภาโค สนฺทิฏฺฐิกา เทฺว. \\ นิพฺพานํ ปรินิพฺพานํ, ตทงฺคทิฏฺฐธมฺมิเกน จาติ. \\ \textbf{ปฐโม ปณฺณาสโก สมตฺโต.}}
\csromangathagroup{\csromangathastanza{สมฺพาโธ กายสกฺขี ปญฺญา, อุภโตภาโค สนฺทิฏฺฐิกา เทฺว. \\ นิพฺพานํ ปรินิพฺพานํ, ตทงฺคทิฏฺฐธมฺมิเกน จาติ. \\ \textbf{ปฐโม ปณฺณาสโก สมตฺโต.}}}

\csromanheader{2}{2. ทุติยปณฺณาสก}
\chapterhead{\textbf{(6) 1. เขมวคฺค}}
\csromanheader{2}{1. เขมสุตฺต}
\proseitem{52}{\csromanfolio{248}“เขมํ เขมนฺ”ติ อาวุโส วุจฺจติ, กิตฺตาวตา นุ โข อาวุโส เขมํ วุตฺตํ ภควตาติ.\csromanspacer{} อิธาวุโส ภิกฺขุ วิวิจฺเจว กาเมหิ ฯเปฯ ปฐมํ ฌานํ อุปสมฺปชฺช วิหรติ.\csromanspacer{} เอตฺตาวตาปิ โข อาวุโส เขมํ วุตฺตํ ภควตา ปริยาเยน ฯเปฯ.}

\prose{ปุน จปรํ อาวุโส ภิกฺขุ สพฺพโส เนวสญฺญานาสญฺญายตนํ สมติกฺกมฺม สญฺญาเวทยิตนิโรธํ อุปสมฺปชฺช วิหรติ, ปญฺญาย จสฺส ทิสฺวา อาสวา ปริกฺขีณา โหนฺติ.\csromanspacer{} เอตฺตาวตาปิ โก อาวุโส เขมํ วุตฺตํ ภควตา นิปฺปริยาเยนาติ.\csromanspacer{}\csromanspacer{}ปฐมํ.}

\csromanheader{2}{2. เขมปฺปตฺตสุตฺต}
\vspace{\tipitakaparskip}
\csromancenter{“เขมปฺปตฺโต เขมปฺปตฺโต”ติ อาวุโส วุจฺจติ.\csromanspacer{}\csromanspacer{}ทุติยํ.}
\csromanheader{2}{3. อมตสุตฺต}
\vspace{\tipitakaparskip}
\csromancenter{“อมตํ อมตนฺ”ติ อาวุโส วุจฺจติ.\csromanspacer{}\csromanspacer{}ตติยํ.}
\csromanheader{2}{4. อมตปฺปตฺตสุตฺต}
\vspace{\tipitakaparskip}
\csromancenter{“อมตปฺปตฺโต อมตปฺปตฺโต”ติ อาวุโส วุจฺจติ.\csromanspacer{}\csromanspacer{}จตุตฺถํ.}
\csromanheader{2}{5. อภยสุตฺต}
\vspace{\tipitakaparskip}
\csromancenter{“อภยํ อภยนฺ”ติ อาวุโส วุจฺจติ.\csromanspacer{}\csromanspacer{}ปญฺจมํ.}
\csromanheader{2}{6. อภยปฺปตฺตสุตฺต}
\vspace{\tipitakaparskip}
\csromancenter{“อภยปฺปตฺโต อภยปฺปตฺโต”ติ อาวุโส วุจฺจติ.\csromanspacer{}\csromanspacer{}ฉฏฺฐํ.}
\csromanheader{2}{7. ปสฺสทฺธิสุตฺต}
\proseitem{58}{“ปสฺสทฺธิ ปสฺสทฺธี”ติ อาวุโส วุจฺจติ.\csromanspacer{}\csromanspacer{}สตฺตมํ.}

\chapterheadpageread{(6) 1. เขมวคฺค}
\csromanheader{2}{8. อนุปุพฺพปสฺสทฺธิสุตฺต}
\vspace{\tipitakaparskip}
\csromancenter{“อนุปุพฺพปสฺสทฺธิ อนุปุพฺพปสฺสทฺธี”ติ อาวุโส วุจฺจติ.\csromanspacer{}\csromanspacer{}อฏฺฐมํ.}
\csromanheader{2}{9. นิโรธสุตฺต}
\vspace{\tipitakaparskip}
\csromancenter{“นิโรโธ นิโรโธ”ติ อาวุโส วุจฺจติ.\csromanspacer{}\csromanspacer{}นวมํ.}
\csromanheader{2}{10. อนุปุพฺพนิโรธสุตฺต}
\proseitem{61}{\csromanfolio{249}“อนุปุพฺพนิโรโธ อนุปุพฺพนิโรโธ”ติ อาวุโส วุจฺจติ, กิตฺตาวตา นุ โข อาวุโส อนุปุพฺพนิโรโธ วุตฺโต ภควตาติ.\csromanspacer{} อิธาวุโส ภิกฺขุ วิวิจฺเจว กาเมหิ ฯเปฯ ปฐมํ ฌานํ อุปสมฺปชฺช วิหรติ.\csromanspacer{} เอตฺตาวตาปิ โข อาวุโส อนุปุพฺพนิโรโธ วุตฺโต ภควตา ปริยาเยน ฯเปฯ.}

\prose{ปุน จปรํ อาวุโส ภิกฺขุ สพฺพโส เนวสญฺญานาสญฺญายตนํ สมติกฺกมฺม สญฺญาเวทยิตนิโรธํ อุปสมฺปชฺช วิหรติ, ปญฺญาย จสฺส ทิสฺวา อาสวา ปริกฺขีณา โหนฺติ.\csromanspacer{} เอตฺตาวตาปิ โข อาวุโส อนุปุพฺพนิโรโธ วุตฺโต ภควตา นิปฺปริยาเยนาติ.\csromanspacer{}\csromanspacer{}ทสมํ.}

\csromanheader{2}{11. อภพฺพสุตฺต}
\proseitem{62}{นว ภิกฺขเว ธมฺเม อปฺปหาย อภพฺโพ อรหตฺตํ สจฺฉิกาตุํ.\csromanspacer{} กตเม นว? ราคํ โทสํ โมหํ โกธํ อุปนาหํ มกฺขํ ปฬาสํ อิสฺสํ มจฺฉริยํ.\csromanspacer{} อิเม โข ภิกฺขเว นว ธมฺเม อปฺปหาย อภพฺโพ อรหตฺตํ สจฺฉิกาตุํ.\csromanspacer{} นว ภิกฺขเว ธมฺเม ปหาย ภพฺโพ อรหตฺตํ สจฺฉิกาตุํ.\csromanspacer{} กตเม นว? ราคํ โทสํ โมหํ โกธํ อุปนาหํ มกฺขํ ปฬาสํ อิสํ มจฺฉริยํ.\csromanspacer{} อิเม โข ภิกฺขเว นว ธมฺเม ปหาย ภพฺโพ อรหตฺตํ สจฺฉิกาตุนฺติ.\csromanspacer{}\csromanspacer{}เอกาทสมํ.\csromanspacer{} เขมวคฺโค ปฐโม.}

\titlehead{\textbf{ตสฺสุทฺทานํ}}
\csromangathameasure{เขโม จ อมตํ เจว, อภยํ ปสฺสทฺธีเยน จ. \\ นิโรโธ อนุปุพฺโพ จ, ธมฺมํ ปหาย ภพฺเพน จาติ.}
\csromangathagroup{\csromangathastanza{เขโม จ อมตํ เจว, อภยํ ปสฺสทฺธีเยน จ. \\ นิโรโธ อนุปุพฺโพ จ, ธมฺมํ ปหาย ภพฺเพน จาติ.}}

\chapterheadpageread{นวกนิปาตปาฬิ \textbf{(7) 2. สติปฏฺฐานวคฺค}}
\csromanheader{2}{1. สิกฺขาทุพฺพลฺยสุตฺต}
\proseitem{63}{\csromanfolio{250}ปญฺจิมานิ ภิกฺขเว สิกฺขาทุพฺพลฺยานิ.\csromanspacer{} กตมานิ ปญฺจ? ปาณาติปาโต อทินฺนาทานํ กาเมสุมิจฺฉาจาโร มุสาวาโท สุรเมรยมชฺชปมาทฏฺฐานํ.\csromanspacer{} อิมานิ โข ภิกฺขเว ปญฺจ สิกฺขาทุพฺพลฺยานิ.}

\prose{อิเมสํ โข ภิกฺขเว ปญฺจนฺนํ สิกฺขาทุพฺพลฺยานํ ปหานาย จตฺตาโร สติปฏฺฐานา ภาเวตพฺพา.\csromanspacer{} กตเม จตฺตาโร? อิธ ภิกฺขเว ภิกฺขุ กาเย กายานุปสฺสี วิหรติ อาตาปี สมฺปชาโน สติมา วิเนยฺย โลเก อภิชฺฌาโทมนสฺสํ.\csromanspacer{} เวทนาสุ ฯเปฯ จิตฺเต.\csromanspacer{} ธมฺเมสุ ธมฺมานุปสฺสี วิหรติ อาตาปี สมฺปชาโน สติมา วิเนยฺย โลเก อภิชฺฌาโทมนสฺสํ.\csromanspacer{} อิเมสํ โข ภิกฺขเว ปญฺจนฺนํ สิกฺขาทุพฺพลฺยานํ ปหานาย อิเม จตฺตาโร สติปฏฺฐานา ภาเวตพฺพาติ.\csromanspacer{}\csromanspacer{}ปฐมํ.}

\csromanheader{2}{2. นีวรณสุตฺต}
\proseitem{64}{ปญฺจิมานิ ภิกฺขเว นีวรณานิ.\csromanspacer{} กตมานิ ปญฺจ? กามจฺฉนฺทนีวรณํ พฺยาปาทนีวรณํ ถินมิทฺธนีวรณํ อุทฺธจฺจกุกฺกุจฺจนีวรณํ วิจิกิจฺฉานีวรณํ, อิมานิ โข ภิกฺขเว ปญฺจ นีวรณานิ.}

\prose{อิเมสํ โข ภิกฺขเว ปญฺจนฺนํ นีวรณานํ ปหานาย จตฺตาโร สติปฏฺฐานา ภาเวตพฺพา.\csromanspacer{} กตเม จตฺตาโร? อิธ ภิกฺขเว ภิกฺขุ กาเย กายานุปสฺสี วิหรติ อาตาปี สมฺปชาโน สติมา วิเนยฺย โลเก อภิชฺฌาโทมนสฺสํ.\csromanspacer{} เวทนาสุ ฯเปฯ จิตฺเต.\csromanspacer{} ธมฺเมสุ ธมฺมานุปสฺสี วิหรติ อาตาปี สมฺปชาโน สติมา วิเนยฺย โลเก อภิชฺฌาโทมนสฺสํ.\csromanspacer{} อิเมสํ โข ภิกฺขเว ปญฺจนฺนํ นีวรณานํ ปหานาย อิเม จตฺตาโร สติปฏฺฐานา ภาเวตพฺพาติ.\csromanspacer{}\csromanspacer{}ทุติยํ.}

\csromanheader{2}{3. กามคุณสุตฺต}
\proseitem{65}{ปญฺจิเม ภิกฺขเว กามคุณา.\csromanspacer{} กตเม ปญฺจ? จกฺขุวิญฺเญยฺยา รูปา อิฏฺฐา กนฺตา มนาปา ปิยรูปา กามูปสํหิตา รชนียา.\csromanspacer{} โสตวิญฺเญยฺยา สทฺทา ฯเปฯ.\csromanspacer{} ฆานวิญฺเญยฺยา คนฺธา.\csromanspacer{} ชิวฺหาวิญฺเญยฺยา รสา.\csromanspacer{} กายวิญฺเญยฺยา}

\csromanheader{1}{(7) 2. สติปฏฺฐานวคฺค}
\prose{\csromanfolio{251}โผฏฺฐพฺพา อิฏฺฐา กนฺตา มนาปา ปิยรูปา กามูปสํหิตา รชนียา.\csromanspacer{} อิเม โข ภิกฺขเว ปญฺจ กามคุณา.}

\prose{อิเมสํ โข ภิกฺขเว ปญฺจนฺนํ กามคุณานํ ปหานาย ฯเปฯ.\csromanspacer{} อิเม จตฺตาโร สติปฏฺฐานา ภาเวตพฺพาติ.\csromanspacer{}\csromanspacer{}ตติยํ.}

\csromanheader{2}{4. อุปาทานกฺขนฺธสุตฺต}
\csromanmatikaanchor{mtk.3}
\csromantocmarkref{subsection}{2. นีวรณสุตฺต}{mtk.3}
\bookmark[dest={mtk.3},level=2]{2. นีวรณสุตฺต}
\proseitem{66}{ปญฺจิเม ภิกฺขเว อุปาทานกฺขนฺธา.\csromanspacer{} กตเม ปญฺจ? รูปุปาทานกฺขนฺโธ เวทนุปาทานกฺขนฺโธ สญฺญุปาทานกฺขนฺโธ สงฺขารุปาทานกฺขนฺโธ วิญฺญาณุปาทานกฺขนฺโธ.\csromanspacer{} อิเม โข ภิกฺขเว ปญฺจุปาทานกฺขนฺธา.}

\prose{อิเมสํ โข ภิกฺขเว ปญฺจนฺนํ อุปาทานกฺขนฺธานํ ปหานาย ฯเปฯ.\csromanspacer{} อิเม จตฺตาโร สติปฏฺฐานา ภาเวตพฺพาติ.\csromanspacer{}\csromanspacer{}จตุตฺถํ.}

\csromanheader{2}{5. โอรมฺภาคิยสุตฺต}
\csromanmatikaanchor{mtk.4}
\csromantocmarkref{subsection}{3. กามคุณสุตฺต}{mtk.4}
\bookmark[dest={mtk.4},level=2]{3. กามคุณสุตฺต}
\proseitem{67}{ปญฺจิมานิ ภิกฺขเว โอรมฺภาคิยานิ สํโยชนานิ.\csromanspacer{} กตมานิ ปญฺจ? สกฺกายทิฏฺฐิ วิจิกิจฺฉา สีลพฺพตปรามาโส กามจฺฉนฺโท พฺยาปาโท.\csromanspacer{} อิมานิ โข ภิกฺขเว ปญฺโจรมฺภาคิยานิ สํโยชนานิ.}

\prose{อิเมสํ โข ภิกฺขเว ปญฺจนฺนํ โอรมฺภาคิยานํ สํโยชนานํ ปหานาย ฯเปฯ.\csromanspacer{} อิเม จตฺตาโร สติปฏฺฐานา ภาเวตพฺพาติ.\csromanspacer{}\csromanspacer{}ปญฺจมํ.}

\csromanheader{2}{6. คติสุตฺต}
\proseitem{68}{ปญฺจิมา ภิกฺขเว คติโย.\csromanspacer{} กตมา ปญฺจ? นิรโย ติรจฺฉานโยนิ เปตฺติวิสโย มนุสฺสา เทวา.\csromanspacer{} อิมา โข ภิกฺขเว ปญฺจ คติโย.}

\prose{อิมาสํ โข ภิกฺขเว ปญฺจนฺนํ คตีนํ ปหานาย ฯเปฯ.\csromanspacer{} อิเม จตฺตาโร สติปฏฺฐานา ภาเวตพฺพาติ.\csromanspacer{}\csromanspacer{}ฉฏฺฐํ.}

\csromanmatikaanchor{mtk.5}
\csromantocmarkref{subsection}{4. อุปาทานกฺขนฺธสุตฺต}{mtk.5}
\bookmark[dest={mtk.5},level=2]{4. อุปาทานกฺขนฺธสุตฺต}
\csromanheader{2}{7. มจฺฉริยสุตฺต}
\proseitem{69}{ปญฺจิมานิ ภิกฺขเว มจฺฉริยานิ.\csromanspacer{} กตมานิ ปญฺจ? อาวาสมจฺฉริยํ กุลมจฺฉริยํ ลาภมจฺฉริยํ วณฺณมจฺฉริยํ ธมฺมมจฺฉริยํ.\csromanspacer{} อิมานิ โข ภิกฺขเว ปญฺจ มจฺฉริยานิ.}

\gambhira{นวกนิปาตปาฬิ}
\csromanmatikaanchor{mtk.6}
\csromantocmarkref{subsection}{5. โอรมฺภาคิยสุตฺต}{mtk.6}
\bookmark[dest={mtk.6},level=2]{5. โอรมฺภาคิยสุตฺต}
\prose{\csromanfolio{252}อิเมสํ โข ภิกฺขเว ปญฺจนฺนํ มจฺฉริยานํ ปหานาย ฯเปฯ.\csromanspacer{} อิเม จตฺตาโร สติปฏฺฐานา ภาเวตพฺพาติ.\csromanspacer{}\csromanspacer{}สตฺตมํ.}

\csromanheader{2}{8. อุทฺธมฺภาคิยสุตฺต}
\proseitem{70}{ปญฺจิมานิ ภิกฺขเว อุทฺธมฺภาคิยานิ สํโยชนานิ.\csromanspacer{} กตมานิ ปญฺจ? รูปราโค อรูปราโค มาโน อุทฺธจฺจํ อวิชฺชา.\csromanspacer{} อิมานิ โข ภิกฺขเว ปญฺจุทฺธมฺภาคิยานิ สํโยชนานิ.}

\csromanmatikaanchor{mtk.7}
\csromantocmarkref{subsection}{6. คติสุตฺต}{mtk.7}
\bookmark[dest={mtk.7},level=2]{6. คติสุตฺต}
\prose{อิเมสํ โข ภิกฺขเว ปญฺจนฺนํ อุทฺธมฺภาคิยานํ สํโยชนานํ ปหานาย ฯเปฯ.\csromanspacer{} อิเม จตฺตาโร สติปฏฺฐานา ภาเวตพฺพาติ.\csromanspacer{}\csromanspacer{}อฏฺฐมํ.}

\csromanheader{2}{9. เจโตขิลสุตฺต}
\proseitem{71}{\csromansymbolfootnote{*}{อํ 2. 217; ที 3. 198; ม 1. 145 ปิฏฺเฐสุปิ.}ปญฺจิเม ภิกฺขเว เจโตขิลา\footnote{เจโตขีลา (ก)}.\csromanspacer{} กตเม ปญฺจ? อิธ ภิกฺขเว ภิกฺขุ สตฺถริ กงฺขติ วิจิกิจฺฉติ นาธิมุจฺจติ น สมฺปสีทติ.\csromanspacer{} โย โส ภิกฺขเว ภิกฺขุ สตฺถริ กงฺขติ วิจิกิจฺฉติ นาธิมุจฺจติ น สมฺปสีทติ, ตสฺส จิตฺตํ น นมติ อาตปฺปาย อนุโยคาย สาตจฺจาย ปธานาย.\csromanspacer{} ยสฺส จิตฺตํ น นมติ อาตปฺปาย อนุโยคาย สาตจฺจาย ปธานาย.\csromanspacer{} อยํ ปฐโม เจโตขิโล.}

\csromanmatikaanchor{mtk.8}
\csromantocmarkref{subsection}{7. มจฺฉริยสุตฺต}{mtk.8}
\bookmark[dest={mtk.8},level=2]{7. มจฺฉริยสุตฺต}
\prose{ปุน จปรํ ภิกฺขเว ภิกฺขุ ธมฺเม กงฺขติ ฯเปฯ สํเฆ กงฺขติ สิกฺขาย กงฺขติ สพฺรหฺมจารีสุ กุปิโต โหติ อนตฺตมโน อาหตจิตฺโต ขิลชาโต.\csromanspacer{} โย โส ภิกฺขเว ภิกฺขุ สพฺรหฺมจารีสุ กุปิโต โหติ อนตฺตมโน อาหตจิตฺโต ขิลชาโต, ตสฺส จิตฺตํ น นมติ อาตปฺปาย อนุโยคาย สาตจฺจาย ปธานาย.\csromanspacer{} ยสฺส จิตฺตํ น นมติ อาตปฺปาย อนุโยคาย สาตจฺจาย ปธานาย.\csromanspacer{} อยํ ปญฺจโม เจโตขิโล.}

\prose{อิเมสํ โข ภิกฺขเว ปญฺจนฺนํ เจโตขิลานํ ปหานาย ฯเปฯ.\csromanspacer{} อิเม จตฺตาโร สติปฏฺฐานา ภาเวตพฺพาติ.\csromanspacer{}\csromanspacer{}นวมํ.}

\csromanheader{2}{10. เจตโสวินิพนฺธสุตฺต}
\proseitem{72}{ปญฺจิเม ภิกฺขเว เจตโสวินิพนฺธา\footnote{เจโตวินิพทฺธา (สารตฺถทีปนีฏีกา) อํ 2. 218; ที 3. 199 ปิฏฺเฐสุปิ.}.\csromanspacer{} กตเม ปญฺจ? อิธ ภิกฺขเว ภิกฺขุ กาเมสุ อวีตราโค โหติ อวิคตจฺฉนฺโท อวิคตเปโม}

\csromanmatikaanchor{mtk.9}
\csromantocmarkref{subsection}{8. อุทฺธมฺภาคิยสุตฺต}{mtk.9}
\bookmark[dest={mtk.9},level=2]{8. อุทฺธมฺภาคิยสุตฺต}
\chapterheadpageread{(7) 2. สติปฏฺฐานวคฺค}
\csromanmatikaanchor{mtk.10}
\csromanmatikaanchor{mtk.12}
\csromantocmarkref{subsection}{9. เจโตขิลสุตฺต}{mtk.10}
\bookmark[dest={mtk.10},level=2]{9. เจโตขิลสุตฺต}
\prose{\csromanfolio{253}อวิคตปิปาโส อวิคตปริฬาโห อวิคตตโณฺห.\csromanspacer{} โย โส ภิกฺขเว ภิกฺขุ กาเมสุ อวีตราโค โหติ อวิคตจฺฉนฺโท อวิคตเปโม อวิคตปิปาโส อวิคตปริฬาโห อวิคตตโณฺห, ตสฺส จิตฺตํ น นมติ อาตปฺปาย อนุโยคาย สาตจฺจาย ปธานาย.\csromanspacer{} ยสฺส จิตฺตํ น นมติ อาตปฺปาย อนุโยคาย สาตจฺจาย ปธานาย.\csromanspacer{} อยํ ปฐโม เจตโสวินิพนฺโธ.}

\prose{ปุน จปรํ ภิกฺขเว ภิกฺขุ กาเย อวีตราโค โหติ ฯเปฯ รูเป อวีตราโค โหติ.\csromanspacer{} ยาวทตฺถํ อุทราวเทหกํ ภุญฺชิตฺวา เสยฺยสุขํ ปสฺสสุขํ มิทฺธสุขํ อนุยุตฺโต วิหรติ.\csromanspacer{} อญฺญตรํ เทวนิกายํ ปณิธาย พฺรหฺมจริยํ จรติ “อิมินาหํ สีเลน วา วเตน วา ตเปน วา พฺรหฺมจริเยน วา เทโว วา ภวิสฺสามิ เทวญฺญตโร วา”ติ.\csromanspacer{} โย โส ภิกฺขเว ภิกฺขุ อญฺญตรํ เทวนิกายํ ปณิธาย พฺรหฺมจริยํ จรติ “อิมินาหํ สีเลน วา วเตน วา ตเปน วา พฺรหฺมจริเยน วา เทโว วา ภวิสฺสามิ เทวญฺญตโร วา”ติ.\csromanspacer{} ตสฺส จิตฺตํ น นมติ อาตปฺปาย อนุโยคาย สาตจฺจาย ปธานาย.\csromanspacer{} ยสฺส จิตฺตํ น นมติ อาตปฺปาย อนุโยคาย สาตจฺจาย ปธานาย.\csromanspacer{} อยํ ปญฺจโม เจตโสวินิพนฺโธ.\csromanspacer{} อิเม โข ภิกฺขเว ปญฺจ เจตโสวินิพนฺธา.}

\prose{อิเมสํ โข ภิกฺขเว ปญฺจนฺนํ เจตโสวินิพนฺธานํ ปหานาย จตฺตาโร สติปฏฺฐานา ภาเวตพฺพา.\csromanspacer{} กตเม จตฺตาโร? อิธ ภิกฺขเว ภิกฺขุ กาเย กายานุปสฺสี วิหรติ อาตาปี สมฺปชาโน สติมา วิเนยฺยโลเก อภิชฺฌาโทมนสฺสํ.\csromanspacer{} เวทนาสุ ฯเปฯ จิตฺเต.\csromanspacer{} ธมฺเมสุ ธมฺมานุปสฺสี วิหรติ อาตาปี สมฺปชาโน สติมา วิเนยฺย โลเก อภิชฺฌาโทมนสฺสํ.\csromanspacer{} อิเมสํ โข ภิกฺขเว ปญฺจนฺนํ เจตโสวินิพนฺธานํ ปหานาย อิเม จตฺตาโร สติปฏฺฐานา ภาเวตพฺพาติ.\csromanspacer{}\csromanspacer{}ทสมํ.\csromanspacer{} สติปฏฺฐานวคฺโค ทุติโย.}

\titlehead{\textbf{ตสฺสุทฺทานํ}}
\csromangathameasure{สิกฺขา นีวรณากามา, ขนฺธา จ โอรมฺภาคิยา คติ. \\ มจฺเฉรํ อุทฺธมฺภาคิยา อฏฺฐมํ, เจโตขิลา วินิพนฺธาติ.}
\csromangathagroup{\csromangathastanza{สิกฺขา นีวรณากามา, ขนฺธา จ โอรมฺภาคิยา คติ. \\ มจฺเฉรํ อุทฺธมฺภาคิยา อฏฺฐมํ, เจโตขิลา วินิพนฺธาติ.}}

\csromanheader{1}{นวกนิปาตปาฬิ \textbf{(8) 3. สมฺมปฺปธานวคฺค}}
\csromanmatikaanchor{mtk.11}
\csromantocmarkref{subsection}{10. เจตโสวินิพนฺธสุตฺต}{mtk.11}
\bookmark[dest={mtk.11},level=2]{10. เจตโสวินิพนฺธสุตฺต}
\csromantocmarkref{section}{อุทฺทานคาถา}{mtk.12}
\bookmark[dest={mtk.12},level=1]{อุทฺทานคาถา}
\csromanheader{2}{1. สิกฺขสุตฺต}
\proseitem{73}{\csromanfolio{254}ปญฺจิมานิ ภิกฺขเว สิกฺขาทุพฺพลฺยานิ.\csromanspacer{} กตมานิ ปญฺจ? ปาณาติปาโต ฯเปฯ สุราเมรยมชฺชปมาทฏฺฐานํ.\csromanspacer{} อิมานิ โข ภิกฺขเว ปญฺจ สิกฺขาทุพฺพลฺยานิ.}

\prose{อิเมสํ โข ภิกฺขเว ปญฺจนฺนํ สิกฺขาทุพฺพลฺยานํ ปหานาย จตฺตาโร สมฺมปฺปธานา ภาเวตพฺพา.\csromanspacer{} กตเม จตฺตาโร? อิธ ภิกฺขเว ภิกฺขุ อนุปฺปนฺนานํ ปาปกานํ อกุสลานํ ธมฺมานํ อนุปฺปาทาย ฉนฺทํ ชเนติ วายมติ วีริยํ อารภติ จิตฺตํ ปคฺคณฺหาติ ปทหติ.\csromanspacer{} อุปฺปนฺนานํ ปาปกานํ อกุสลานํ ธมฺมานํ ปหานาย ฉนฺทํ ชเนติ วายมติ วีริยํ อารภติ จิตฺตํ ปคฺคณฺหาติ ปทหติ.\csromanspacer{} อนุปฺปนฺนานํ กุสลานํ ธมฺมานํ อุปฺปาทาย ฉนฺทํ ชเนติ วายมติ วีริยํ อารภติ จิตฺตํ ปคฺคณฺหาติ ปทหติ.\csromanspacer{} อุปฺปนฺนานํ กุสลานํ ธมฺมานํ ฐิติยา อสมฺโมสาย ภิยฺโยภาวาย เวปุลฺลาย ภาวนาย ปาริปูริยา ฉนฺทํ ชเนติ วายมติ วีริยํ อารภติ จิตฺตํ ปคฺคณฺหาติ ปทหติ.\csromanspacer{} อิเมสํ โข ภิกฺขเว ปญฺจนฺนํ สิกฺขาทุพฺพลฺยานํ ปหานาย อิเม จตฺตาโร สมฺมปฺปธานา ภาเวตพฺพาติ.\csromanspacer{}\csromanspacer{}ปฐมํ.}

\proseitem{74-81}{(ยถา สติปฏฺฐานวคฺเค, ตถา สมฺมปฺปธานวเสน วิตฺถาเรตพฺพา.)}

\csromanheader{2}{10. เจตโสวินิพนฺธสุตฺต}
\proseitem{82}{ปญฺจิเม ภิกฺขเว เจตโสวินิพนฺธา.\csromanspacer{} กตเม ปญฺจ? อิธ ภิกฺขเว ภิกฺขุ กาเมสุ อวีตราโค โหติ ฯเปฯ.\csromanspacer{} อิเม โข ภิกฺขเว ปญฺจ เจตโสวินิพนฺธา.}

\prose{อิเมสํ โข ภิกฺขเว ปญฺจนฺนํ เจตโสวินิพนฺธานํ ปหานาย จตฺตาโร สมฺมปฺปธานา ภาเวตพฺพา.\csromanspacer{} กตเม จตฺตาโร? อิธ ภิกฺขเว ภิกฺขุ อนุปฺปนฺนานํ ปาปกานํ อกุสลานํ ธมฺมานํ อนุปฺปาทาย ฉนฺทํ ชเนติ วายมติ วีริยํ อารภติ จิตฺตํ ปคฺคณฺหาติ ปทหติ.\csromanspacer{} อุปฺปนฺนานํ ปาปกานํ อกุสลานํ ธมฺมานํ ปหานาย.\csromanspacer{} อนุปฺปนฺนานํ กุสลานํ ธมฺมานํ อุปฺปาทาย.\csromanspacer{} อุปฺปนฺนานํ กุสลานํ ธมฺมานํ ฐิติยา อสมฺโมสาย ภิยฺโยภาวาย เวปุลฺลาย ภาวนาย ปาริปูริยา ฉนฺทํ ชเนติ วายมติ วีริยํ อารภติ จิตฺตํ ปคฺคณฺหาติ}

\csromanheader{1}{\textbf{(9) 4. อิทฺธิปาทวคฺค}}
\csromanmatikaanchor{mtk.13}
\csromantocmarknopage{section}{(8) 3. สมฺมปฺปธานวคฺค}{mtk.13}
\bookmark[dest={mtk.13},level=1]{(8) 3. สมฺมปฺปธานวคฺค}
\prose{\csromanfolio{255}ปทหติ.\csromanspacer{} อิเมสํ โข ภิกฺขเว ปญฺจนฺนํ เจตโสวินิพนฺธานํ ปหานาย อิเม จตฺตาโร สมฺมปฺปธานา ภาเวตพฺพาติ.\csromanspacer{}\csromanspacer{}ทสมํ.}

\csromanmatikaanchor{mtk.14}
\csromantocmarkref{subsection}{1. สิกฺขสุตฺต}{mtk.14}
\bookmark[dest={mtk.14},level=2]{1. สิกฺขสุตฺต}
\csromanheader{2}{สมฺมปฺปธานวคฺโค ตติโย.}
\csromansectionrule
\chapterhead{\textbf{(9) 4. อิทฺธิปาทวคฺค}}
\csromanheader{2}{1. สิกฺขสุตฺต}
\csromanmatikaanchor{mtk.15}
\csromantocmarkref{subsection}{10. เจตโสวินิพนฺธสุตฺต}{mtk.15}
\bookmark[dest={mtk.15},level=2]{10. เจตโสวินิพนฺธสุตฺต}
\proseitem{83}{ปญฺจิมานิ ภิกฺขเว สิกฺขาทุพฺพลฺยานิ.\csromanspacer{} กตมานิ ปญฺจ? ปาณาติปาโต ฯเปฯ.\csromanspacer{} สุราเมรยมชฺชปมาทฏฺฐานํ.\csromanspacer{} อิมานิ โข ภิกฺขเว ปญฺจ สิกฺขาทุพฺพลฺยานิ. อิเมสํ โข ภิกฺขเว ปญฺจนฺนํ สิกฺขาทุพฺพลฺยานํ ปหานาย จตฺตาโร อิทฺธิปาทา ภาเวตพฺพา.\csromanspacer{} กตเม จตฺตาโร? อิธ ภิกฺขเว ภิกฺขุ ฉนฺทสมาธิปธานสงฺขารสมนฺนาคตํ อิทฺธิปาทํ ภาเวติ, วีริยสมาธิ.\csromanspacer{} จิตฺตสมาธิ.\csromanspacer{} วีมํสาสมาธิปธานสงฺขารสมนฺนาคตํ อิทฺธิปาทํ ภาเวติ.\csromanspacer{} อิเมสํ โข ภิกฺขเว ปญฺจนฺนํ สิกฺขาทุพฺพลฺยานํ ปหานาย อิเม จตฺตาโร อิทฺธิปาทา ภาเวตพฺพาติ.\csromanspacer{}\csromanspacer{}ปฐมํ.}

\proseitem{84-91}{(ยถา สติปฏฺฐานวคฺเค, ตถา อิทฺธิปาทวเสน วิตฺถาเรตพฺพา.)}

\csromanheader{2}{10. เจตโสวินิพนฺธสุตฺต}
\csromanmatikaanchor{mtk.16}
\csromantocmarknopage{section}{(9) 4. อิทฺธิปาทวคฺค}{mtk.16}
\bookmark[dest={mtk.16},level=1]{(9) 4. อิทฺธิปาทวคฺค}
\proseitem{92}{ปญฺจิเม ภิกฺขเว เจตโสวินิพนฺธา.\csromanspacer{} กตเม ปญฺจ? อิธ ภิกฺขเว ภิกฺขุ กาเมสุ อวีตราโค โหติ ฯเปฯ.\csromanspacer{} อิเม โข ภิกฺขเว ปญฺจ เจตโสวินิพนฺธา.}

\prose{อิเมสํ โข ภิกฺขเว ปญฺจนฺนํ เจตโสวินิพนฺธานํ ปหานาย อิเม จตฺตาโร อิทฺธิปาทา ภาเวตพฺพา.\csromanspacer{} กตเม จตฺตาโร? อิธ ภิกฺขเว ภิกฺขุ ฉนฺทสมาธิปธานสงฺขารสมนฺนาคตํ อิทฺธิปาทํ ภาเวติ, วีริยสมาธิ.\csromanspacer{} จิตฺตสมาธิ.\csromanspacer{} วีมํสาสมาธิปธานสงฺขารสมนฺนาคตํ อิทฺธิปาทํ ภาเวติ.\csromanspacer{} อิเมสํ โข ภิกฺขเว ปญฺจนฺนํ เจตโสวินิพนฺธานํ ปหานาย อิเม จตฺตาโร อิทฺธิปาทา ภาเวตพฺพาติ.\csromanspacer{}\csromanspacer{}ทสมํ.\csromanspacer{} อิทฺธิปาทวคฺโค จตุตฺโถ.}

\gambhira{นวกนิปาตปาฬิ}
\csromanmatikaanchor{mtk.17}
\csromanmatikaanchor{mtk.19}
\csromantocmarkref{subsection}{1. สิกฺขสุตฺต}{mtk.17}
\bookmark[dest={mtk.17},level=2]{1. สิกฺขสุตฺต}
\csromangathameasure{ยเถว สติปฏฺฐานา, ปธานา จตุโรปิ จ. \\ จตฺตาโร อิทฺธิปาทา จ, ตเถว สมฺปโยชเยติ.}
\csromangathagroup{\csromangathastanza{\csromanfolio{256}ยเถว สติปฏฺฐานา, ปธานา จตุโรปิ จ.}\csromangathastanza{จตฺตาโร อิทฺธิปาทา จ, ตเถว สมฺปโยชเยติ.}}

\vspace{\tipitakaparskip}
\csromancenter{\textbf{(10) 5.}\csromanspacer{}\textbf{ ราคเปยฺยาล}}
\proseitem{93}{ราคสฺส ภิกฺขเว อภิญฺญาย นว ธมฺมา ภาเวตพฺพา.\csromanspacer{} กตเม นว? อสุภสญฺญา มรณสญฺญา อาหาเร ปฏิกูลสญฺญา สพฺพโลเก อนภิรตสญฺญา อนิจฺจสญฺญา อนิจฺเจ ทุกฺขสญฺญา ทุกฺเข อนตฺตสญฺญา ปหานสญฺญา วิราคสญฺญา, ราคสฺส ภิกฺขเว อภิญฺญาย อิเม นว ธมฺมา ภาเวตพฺพาติ. (1)}

\proseitem{94}{ราคสฺส ภิกฺขเว อภิญฺญาย นว ธมฺมา ภาเวตพฺพา.\csromanspacer{} กตเม นว? ปฐมํ ฌานํ.\csromanspacer{} ทุติยํ ฌานํ.\csromanspacer{} ตติยํ ฌานํ.\csromanspacer{} จตุตฺถํ ฌานํ.\csromanspacer{} อากาสานญฺจายตนํ.\csromanspacer{} วิญฺญาณญฺจายตนํ.\csromanspacer{} อากิญฺจญฺญายตนํ.\csromanspacer{} เนวสญฺญานาสญฺญายตนํ.\csromanspacer{} สญฺญาเวทยิตนิโรโธ.\csromanspacer{} ราคสฺส ภิกฺขเว อภิญฺญาย อิเม นว ธมฺมา ภาเวตพฺพาติ. (2)}

\csromanmatikaanchor{mtk.18}
\csromantocmarkref{subsection}{10. เจตโสวินิพนฺธสุตฺต}{mtk.18}
\bookmark[dest={mtk.18},level=2]{10. เจตโสวินิพนฺธสุตฺต}
\csromantocmarkref{section}{(10) 5. ราคเปยฺยาล}{mtk.19}
\bookmark[dest={mtk.19},level=1]{(10) 5. ราคเปยฺยาล}
\proseitem{95-112}{ราคสฺส ภิกฺขเว ปริญฺญาย ฯเปฯ ปริกฺขยาย ฯเปฯ ปหานาย ฯเปฯ ขยาย ฯเปฯ วยาย ฯเปฯ วิราคาย ฯเปฯ นิโรธาย ฯเปฯ จาคาย ฯเปฯ ปฏินิสฺสคฺคาย ฯเปฯ.\csromanspacer{} อิเม นว ธมฺมา ภาเวตพฺพา. (3-20)}

\proseitem{113-432}{โทสสฺส ฯเปฯ.\csromanspacer{} โมหสฺส.\csromanspacer{} โกธสฺส.\csromanspacer{} อุปนาหสฺส.\csromanspacer{} มกฺขสฺส.\csromanspacer{} ปฬาสสฺส.\csromanspacer{} อิสฺสาย.\csromanspacer{} มจฺฉริยสฺส.\csromanspacer{} มายาย.\csromanspacer{} สาฏฺเฐยฺยสฺส.\csromanspacer{} ถมฺภสฺส.\csromanspacer{} สารมฺภสฺส.\csromanspacer{} มานสฺส.\csromanspacer{} อติมานสฺส.\csromanspacer{} มทสฺส.\csromanspacer{} ปมาทสฺส.\csromanspacer{} อภิญฺญาย ฯเปฯ ปริญฺญาย.\csromanspacer{} ปริกฺขยาย.\csromanspacer{} ปหานาย.\csromanspacer{} ขยาย.\csromanspacer{} วยาย.\csromanspacer{} วิราคาย.\csromanspacer{} นิโรธาย.\csromanspacer{} จาคาย.\csromanspacer{} ปฏินิสฺสคฺคาย ฯเปฯ.\csromanspacer{} อิเม นว ธมฺมา ภาเวตพฺพาติ. (21-340)}

\nitthitam{ราคเปยฺยาลํ นิฏฺฐิตํ.}
\nitthitam{\textbf{นวกนิปาตปาฬิ นิฏฺฐิตา.}}
\csromanheader{2}{\textbf{องฺคุตฺตรนิกาย}}
\gambhira{\textbf{ทสกนิปาตปาฬิ}}
\namakkaram{นโม ตสฺส ภควตา อรหโต สมฺมาสมฺพุทฺธสฺส.}
\csromanheader{2}{1. ปฐมปณฺณาสก}
\csromanheader{1}{1. อานิสํสวคฺค}
\csromanheader{2}{1. กิมตฺถิยสุตฺต}
\proseitem{1}{\csromanfolio{257}เอวํ เม สุตํ\csromandash{}เอกํ สมยํ ภควา สาวตฺถิยํ วิหรติ เชตวเน อนาถปิณฺฑิกสฺส อาราเม.\csromanspacer{} อถ โข อายสฺมา อานนฺโท เยน ภควา เตนุปสงฺกมิ, อุปสงฺกมิตฺวา ภควนฺตํ อภิวาเทตฺวา เอกมนฺตํ นิสีทิ, เอกมนฺตํ นิสินฺโน โข อายสฺมา อานนฺโท ภควนฺตํ เอตทโวจ\csromandash{}}

\prose{กิมตฺถิยานิ ภนฺเต กุสลานิ สีลานิ กิมานิสํสานีติ.\csromanspacer{} อวิปฺปฏิสารตฺถานิ โข อานนฺท กุสลานิ สีลานิ อวิปฺปฏิสารานิสํสานีติ.}

\prose{อวิปฺปฏิสาโร ปน ภนฺเต กิมตฺถิโย กิมานิสํโสติ.\csromanspacer{} อวิปฺปฏิสาโร โข อานนฺท ปาโมชฺชตฺโถ ปาโมชฺชานิสํโสติ\footnote{ปามุชฺชตฺโถ ปามุชฺชานิสํโสติ (สี, สฺยา, อิ) อุปริ 515 ปิฏฺเฐปิ.}.}

\prose{ปาโมชฺชํ ปน ภนฺเต กิมตฺถิยํ กิมานิสํสนฺติ.\csromanspacer{} ปาโมชฺชํ โข อานนฺท ปีตตฺถํ ปีตานิสํสนฺติ.}

\csromanmatikaanchor{mtk.20}
\csromantocmarknopage{chapter}{ทสกนิปาตปาฬิ}{mtk.20}
\bookmark[dest={mtk.20},level=0]{ทสกนิปาตปาฬิ}
\csromantocmarknopage{chapter}{1. ปฐมปณฺณาสก}{mtk.21}
\bookmark[dest={mtk.21},level=0]{1. ปฐมปณฺณาสก}
\prose{ปีติ ปน ภนฺเต กิมตฺถิยา กิมานิสํสาติ.\csromanspacer{} ปีติ โข อานนฺท ปสฺสทฺธตฺถา ปสฺสทฺธานิสํสาติ.}

\gambhira{ทสกนิปาตปาฬิ}
\prose{\csromanfolio{258}ปสฺสทฺธิ ปน ภนฺเต กิมตฺถิยา กิมานิสํสาติ.\csromanspacer{} ปสฺสทฺธิ โข อานนฺท สุขตฺถา สุขานิสํสาติ.}

\csromanmatikaanchor{mtk.22}
\csromantocmarknopage{section}{1. อานิสํสวคฺค}{mtk.22}
\bookmark[dest={mtk.22},level=1]{1. อานิสํสวคฺค}
\prose{สุขํ ปน ภนฺเต กิมตฺถิยํ กิมานิสํสนฺติ.\csromanspacer{} สุขํ โข อานนฺท สมาธตฺถํ สมาธานิสํสนฺติ.}

\csromanmatikaanchor{mtk.23}
\csromantocmarkref{subsection}{1. กิมตฺถิยสุตฺต}{mtk.23}
\bookmark[dest={mtk.23},level=2]{1. กิมตฺถิยสุตฺต}
\prose{สมาธิ ปน ภนฺเต กิมตฺถิโย กิมานิสํโสติ.\csromanspacer{} สมาธิ โข อานนฺท ยถาภูตญาณทสฺสนตฺโถ ยถาภูตญาณทสฺสนานิสํโสติ.}

\prose{ยถาภูตญาณทสฺสนํ ปน ภนฺเต กิมตฺถิยํ กิมานิสํสนฺติ.\csromanspacer{} ยถาภูตญาณทสฺสนํ โข อานนฺท นิพฺพิทาวิราคตฺถํ นิพฺพิทาวิราคานิสํสนฺติ.}

\prose{นิพฺพิทาวิราโค ปน ภนฺเต กิมตฺถิโย กิมานิสํโสติ.\csromanspacer{} นิพฺพิทาวิราโค โข อานนฺท วิมุตฺติญาณทสฺสนตฺโถ วิมุตฺติญาณทสฺสนานิสํโส\footnote{...นิสํโสติ (สี, ก)}.}

\prose{อิติ โข อานนฺท กุสลานิ สีลานิ อวิปฺปฏิสารตฺถานิ อวิปฺปฏิสารานิสํสานิ, อวิปฺปฏิสาโร ปาโมชฺชตฺโถ ปาโมชฺชานิสํโส, ปาโมชฺชํ ปีตตฺถํ ปีตานิสํสํ, ปีติ ปสฺสทฺธตฺถา ปสฺสทฺธานิสํสา, ปสฺสทฺธิ สุขตฺถา สุขานิสํสา, สุขํ สมาธตฺถํ สมาธานิสํสํ, สมาธิ ยถาภูตญาณทสฺสนตฺโถ ยถาภูตญาณทสฺสนานิสํโส, ยถาภูตญาณทสฺสนํ นิพฺพิทาวิราคตฺถํ นิพฺพิทาวิราคานิสํสํ, นิพฺพิทาวิราโค วิมุตฺติญาณทสฺสนตฺโถ วิมุตฺติญาณทสฺสนานิสํโส.\csromanspacer{} อิติ โข อานนฺท กุสลานิ สีลานิ อนุปุพฺเพน อคฺคาย ปเรนฺตีติ\footnote{อรหตฺตาย ปูเรนฺตีติ (สฺยา)}.\csromanspacer{}\csromanspacer{}ปฐมํ.}

\csromanheader{2}{2. เจตนากรณียสุตฺต}
\proseitem{2}{\csromansymbolfootnote{*}{อุปริ 516 ปิฏฺเฐปิ.}สีลวโต ภิกฺขเว สีลสมฺปนฺนสฺส น เจตนาย กรณียํ “อวิปฺปฏิสาโร เม อุปฺปชฺชตู”ติ, ธมฺมตา เอสา ภิกฺขเว ยํ สีลวโต สีลสมฺปนฺนสฺส อวิปฺปฏิสาโร อุปฺปชฺชติ.\csromanspacer{} อวิปฺปฏิสาริสฺส ภิกฺขเว น เจตนาย กรณียํ “ปาโมชฺชํ เม อุปฺปชฺชตู”ติ, ธมฺมตา เอสา ภิกฺขเว ยํ อวิปฺปฏิสาริสฺส ปาโมชฺชํ ชายติ.\csromanspacer{} ปมุทิตสฺส ภิกฺขเว น เจตนาย กรณียํ “ปีติ เม อุปชฺชตู”ติ, ธมฺมตา เอสา ภิกฺขเว ยํ ปมุทิตสฺส ปีติ อุปฺปชฺชติ.\csromanspacer{} ปีติมนสฺส ภิกฺขเว น เจตนาย กรณียํ “กาโย เม ปสฺสมฺภตู”ติ, ธมฺมตา เอสา ภิกฺขเว ยํ}

\csromanmatikaanchor{mtk.24}
\csromanmatikaanchor{mtk.250}
\csromantocmarkref{subsection}{2. เจตนากรณียสุตฺต}{mtk.24}
\bookmark[dest={mtk.24},level=2]{2. เจตนากรณียสุตฺต}
\chapterheadpageread{1. อานิสํสวคฺค}
\prose{\csromanfolio{259}ปีติมนสฺส กาโย ปสฺสมฺภติ.\csromanspacer{} ปสฺสทฺธกายสฺส ภิกฺขเว น เจตนาย กรณียํ “สุขํ เวทิยามี”ติ, ธมฺมตา เอสา ภิกฺขเว ยํ ปสฺสทฺธกาโย สุขํ เวทิยติ.\csromanspacer{} สุขิโน ภิกฺขเว น เจตนาย กรณียํ “จิตฺตํ เม สมาธิยตู”ติ, ธมฺมตา เอสา ภิกฺขเว ยํ สุขิโน จิตฺตํ สมาธิยติ.\csromanspacer{} สมาหิตสฺส ภิกฺขเว น เจตนาย กรณียํ “ยถาภูตํ ชานามิ ปสฺสามี”ติ, ธมฺมตา เอสา ภิกฺขเว ยํ สมาหิโต ยถาภูตํ ชานาติ ปสฺสติ.\csromanspacer{} ยถาภูตํ ภิกฺขเว ชานโต ปสฺสโต น เจตนาย กรณียํ “นิพฺพินฺทามิ วิรชฺชามี”ติ, ธมฺมตา เอสา ภิกฺขเว ยํ ยถาภูตํ ชานํ ปสฺสํ นิพฺพินฺทติ วิรชฺชติ.\csromanspacer{} นิพฺพินฺนสฺส\footnote{นิพฺพินฺทสฺส (สี, ก)} ภิกฺขเว วิรตฺตสฺส น เจตนาย กรณียํ “วิมุตฺติญาณทสฺสนํ สจฺฉิกโรมี”ติ, ธมฺมตา เอสา ภิกฺขเว ยํ นิพฺพินฺโน\footnote{นิพฺพินฺโท (สี, ก)} วิรตฺโต วิมุตฺติญาณทสฺสนํ สจฺฉิกโรติ.}

\prose{อิติ โข ภิกฺขเว นิพฺพิทาวิราโค วิมุตฺติญาณทสฺสนตฺโถ วิมุตฺติญาณทสฺสนานิสํโส, ยถาภูตญาณทสฺสนํ นิพฺพิทาวิราคตฺถํ นิพฺพิทาวิราคานิสํสํ, สมาธิ ยถาภูตญาณทสฺสนตฺโถ ยถาภูตญาณทสฺสนานิสํโส, สุขํ สมาธตฺถํ สมาธานิสํสํ, ปสฺสทฺธิ สุขตฺถา สุขานิสํสา, ปีติ ปสฺสทฺธตฺถา ปสฺสทฺธานิสํสา, ปาโมชฺชํ ปีตตฺถํ ปีตานิสํสํ, อวิปฺปฏิสาโร ปาโมชฺชตฺโถ ปาโมชฺชานิสํโส, กุสลานิ สีลานิ อวิปฺปฏิสารตฺถานิ อวิปฺปฏิสารานิสํสานิ.\csromanspacer{} อิติ โข ภิกฺขเว ธมฺมา ธมฺเม อภิสนฺเทนฺติ, ธมฺมา ธมฺเม ปริปูเรนฺติ อปารา ปารํ คมนายาติ.\csromanspacer{}\csromanspacer{}ทุติยํ.}

\csromanheader{2}{3. ปฐม-อุปนิสสุตฺต}
\proseitem{3}{\csromansymbolfootnote{*}{อํ 2. 16, อุปริ 517 ปิฏฺเฐสุปิ.}ทุสฺสีลสฺส ภิกฺขเว สีลวิปนฺนสฺส หตูปนิโส โหติ อวิปฺปฏิสาโร, อวิปฺปฏิสาเร อสติ อวิปฺปฏิสารวิปนฺนสฺส หตูปนิสํ โหติ ปาโมชฺชํ, ปาโมชฺเช อสติ ปาโมชฺชวิปนฺนสฺส หตูปนิสา โหติ ปีติ, ปีติยา อสติ ปีติวิปนฺนสฺส หตูปนิสา โหติ ปสฺสทฺธิ, ปสฺสทฺธิยา อสติ ปสฺสทฺธิวิปนฺนสฺส หตูปนิสํ โหติ สุขํ, สุเข อสติ สุขวิปนฺนสฺส หตูปนิโส โหติ สมฺมาสมาธิ.\csromanspacer{} สมฺมาสมาธิมฺหิ อสติ สมฺมาสมาธิวิปนฺนสฺส หตูปนิสํ โหติ ยถาภูตญาณทสฺสนํ, ยถาภูตญาณทสฺสเน อสติ ยถาภูตญาณทสฺสนวิปนฺนสฺส หตูปนิโส โหติ}

\vspace{\tipitakaparskip}
\csromancenter{ทสกนิปาตปาฬิ นิพฺพิทาวิราโค, นิพฺพิทาวิราเค อสติ นิพฺพิทาวิราควิปนฺนสฺส หตูปนิสํ โหติ วิมุตฺติญาณทสฺสนํ.\csromanspacer{} เสยฺยถาปิ ภิกฺขเว รุกฺโข สาขาปลาสวิปนฺโน, ตสฺส ปปฏิกาปิ น ปาริปูริํ คจฺฉติ, ตโจปิ.\csromanspacer{} เผคฺคุปิ.\csromanspacer{} สาโรปิ น ปาริปูริํ คจฺฉติ.\csromanspacer{} เอวเมวํ โข ภิกฺขเว ทุสฺสีลสฺส สีลวิปนฺนสฺส หตูปนิโส โหติ อวิปฺปฏิสโร, อวิปฺปฏิสาเร อสติ อวิปฺปฏิสารวิปนฺนสฺส หตูปนิสํ โหติ ฯเปฯ วิมุตฺติญาณทสฺสนํ.}
\prose{\csromanfolio{260}สีลวโต ภิกฺขเว สีลสมฺปนฺนสฺส อุปนิสสมฺปนฺโน โหติ อวิปฺปฏิสาโร, อวิปฺปฏิสาเร สติ อวิปฺปฏิสารสมฺปนฺนสฺส อุปนิสสมฺปนฺนํ โหติ ปาโมชฺชํ, ปาโมชฺเช สติ ปาโมชฺชสมฺปนฺนสฺส อุปนิสสมฺปนฺนา โหติ ปีติ, ปีติยา สติ ปีติสมฺปนฺนสฺส อุปนิสสมฺปนฺนา โหติ ปสฺสทฺธิ, ปสฺสทฺธิยา สติ ปสฺสทฺธิสมฺปนฺนสฺส อุปนิสสมฺปนฺนํ โหติ สุขํ, สุเข สติ สุขสมฺปนฺนสฺส อุปนิสสมฺปนฺโน โหติ สมฺมาสมาธิ, สมฺมาสมาธิมฺหิ สติ สมฺมาสมาธิสมฺปนฺนสฺส อุปนิสสมฺปนฺนํ โหติ ยถาภูตญาณทสฺสนํ, ยถาภูตญาณทสฺสเน สติ ยถาภูตญาณ- ทสฺสนสมฺปนฺนสฺส อุปนิสสมฺปนฺโน โหติ นิพฺพิทาวิราโค, นิพฺพิทาวิราเค สติ นิพฺพิทาวิราคสมฺปนฺนสฺส อุปนิสสมฺปนฺนํ โหติ วิมุตฺติญาณทสฺสนํ.\csromanspacer{} เสยฺยถาปิ ภิกฺขเว รุกฺโข สาขาปลาสสมฺปนฺโน, ตสฺส ปปฏิกาปิ ปาริปูริํ คจฺฉติ, ตโจปิ.\csromanspacer{} เผคฺคุปิ.\csromanspacer{} สาโรปิ ปาริปูริํ คจฺฉติ.\csromanspacer{} เอวเมวํ โข ภิกฺขเว สีลวโต สีลสมฺปนฺนสฺส อุปนิสสมฺปนฺโน โหติ อวิปฺปฏิสาโร, อวิปฺปฏิสาเร สติ อวิปฺปฏิสาร- สมฺปนฺนสฺส อุปนิสสมฺปนฺนํ โหติ ฯเปฯ วิมุตฺติญาณทสฺสนนฺติ.\csromanspacer{}\csromanspacer{}ตติยํ.}

\csromanheader{2}{4. ทุติย-อุปนิสสุตฺต}
\proseitem{4}{\csromansymbolfootnote{*}{อุปริ 519 ปิฏฺเฐปิ.}ตตฺร โข อายสฺมา สาติปุตฺโต ภิกฺขู อามนฺเตสิ ฯเปฯ.\csromanspacer{} ทุสฺสีลสฺส อาวุโส สีลวิปนฺนสฺส หตูปนิโส โหติ อวิปฺปฏิสาโร, อวิปฺปฏิสาเร อสติ อวิปฺปฏิสารวิปนฺนสฺส หตูปนิสํ โหติ ฯเปฯ วิมุตฺติญาณทสฺสนํ.\csromanspacer{} เสยฺยถาปิ อาวุโส รุกฺโข สาขาปลาสวิปนฺโน, ตสฺส ปปฏิกาปิ น ปาริปูริํ คจฺฉติ, ตโจปิ.\csromanspacer{} เผคฺคุปิ.\csromanspacer{} สาโรปิ น ปาริปูริํ คจฺฉติ.\csromanspacer{} เอวเมวํ โข อาวุโส ทุสฺสีลสฺส สีลวิปนฺนสฺส หตูปนิโส โหติ อวิปฺปฏิสาโร, อวิปฺปฏิสาเร อสติ อวิปฺปฏิสารวิปนฺนสฺส หตูปนิสํ โหติ ฯเปฯ วิมุตฺติญาณทสฺสนํ.}

\chapterheadpageread{1. อานิสํสวคฺค}
\prose{\csromanfolio{261}สีลวโต อาวุโส สีลสมฺปนฺนสฺส อุปนิสสมฺปนฺโน โหติ อวิปฺปฏิสาโร, อวิปฺปฏิสาเร สติ อวิปฺปฏิสารสมฺปนฺนสฺส อุปนิสสมฺปนฺนํ โหติ ฯเปฯ วิมุตฺติญาณทสฺสนํ.\csromanspacer{} เสยฺยถาปิ อาวุโส รุกฺโข สาขาปลาสสมฺปนฺโน, ตสฺส ปปฏิกาปิ ปาริปูริํ คจฺฉติ, ตโจปิ.\csromanspacer{} เผคฺคุปิ.\csromanspacer{} สาโรปิ ปาริปูริํ คจฺฉติ.\csromanspacer{} เอวเมวํ โข อาวุโส สีลวโต สีลสมฺปนฺนสฺส อุปนิสสมฺปนฺโน โหติ อวิปฺปฏิสาโร, อวิปฺปฏิสาเร สติ อวิปฺปฏิสารสมฺปนฺนสฺส อุปนิสสมฺปนฺนํ โหติ ฯเปฯ วิมุตฺติญาณทสฺสนนฺติ.\csromanspacer{}\csromanspacer{}จตุตฺถํ.}

\csromanheader{2}{5. ตติย-อุปนิสสุตฺต}
\csromanmatikaanchor{mtk.25}
\csromantocmarkref{subsection}{3. ปฐม-อุปนิสสุตฺต}{mtk.25}
\bookmark[dest={mtk.25},level=2]{3. ปฐม-อุปนิสสุตฺต}
\proseitem{5}{\csromansymbolfootnote{*}{อุปริ 520 ปิฏฺเฐปิ.}ตตฺร โข อายสฺมา อานนฺโท ภิกฺขู อามนฺเตสิ ฯเปฯ.\csromanspacer{} ทุสฺสีลสฺส อาวุโส สีลวิปนฺนสฺส หตูปนิโส โหติ อวิปฺปฏิสาโร, อวิปฺปฏิสาเร อสติ อวิปฺปฏิสารวิปนฺนสฺส หตูปนิสํ โหติ ปาโมชฺชํ, ปาโมชฺเช อสติ ปาโมชฺชวิปนฺนสฺส หตูปนิสา โหติ ปีติ, ปีติยา อสติ ปีติวิปนฺนสฺส หตูปนิสา โหติ ปสฺสทฺธิ, ปสฺสทฺธิยา อสติ ปสฺสทฺธิวิปนฺนสฺส หตูปนิสํ โหติ สุขํ, สุเข อสติ สุขวิปนฺนสฺส หตูปนิโส โหติ สมฺมาสมาธิ, สมฺมาสมาธิมฺหิ อสติ สมฺมาสมาธิวิปนฺนสฺส หตูปนิสํ โหติ ยถาภูตญาณทสฺสนํ, ยถาภูตญาณทสฺสเน อสติ ยถาภูตญาณทสฺสนวิปนฺนสฺส หตูปนิโส โหติ นิพฺพิทาวิราโค, นิพฺพิทาวิราเค อสติ นิพฺพิทาวิราควิปนฺนสฺส หตูปนิสํ โหติ วิมุตฺติญาณทสฺสนํ.\csromanspacer{} เสยฺยถาปิ อาวุโส รุกฺโข สาขาปลาสวิปนฺโน, ตสฺส ปปฏิกาปิ น ปาริปูริํ คจฺฉติ, ตโจปิ.\csromanspacer{} เผคฺคุปิ.\csromanspacer{} สาโรปิ น ปาริปูริํ คจฺฉติ.\csromanspacer{} เอวเมวํ โข อาวุโส ทุสีลสฺส สีลวิปนฺนสฺส หตูปนิโส โหติ อวิปฺปฏิสาโร, อวิปฺปฏิสาเร อสติ อวิปฺปฏิสารวิปนฺนสฺส หตูปนิสํ โหติ ฯเปฯ.\csromanspacer{} วิมุตฺติญาณทสฺสนํ.}

\prose{สีลวโต อาวุโส สีลสมฺปนฺนสฺส อุปนิสสมฺปนฺโน โหติ อวิปฺปฏิสาโร, อวิปฺปฏิสาเร สติ อวิปฺปฏิสารสมฺปนฺนสฺส อุปนิสสมฺปนฺนํ โหติ ปาโมชฺชํ, ปาโมชฺเช สติ ปาโมชฺชสมฺปนฺนสฺส อุปนิสสมฺปนฺนา โหติ ปีติ, ปีติยา สติ ปีติสมฺปนฺนสฺส อุปนิสสมฺปนฺนา โหติ ปสฺสทฺธิ, ปสฺสทฺธิยา สติ ปสฺสทฺธิสมฺปนฺนสฺส อุปนิสสมฺปนฺนํ โหติ สุขํ, สุเข สติ สุขสมฺปนฺนสฺส อุปนิสสมฺปนฺโน โหติ สมฺมาสมาธิ, สมฺมาสมาธิมฺหิ สติ สมฺมาสมาธิสมฺปนฺนสฺส อุปนิสสมฺปนฺนํ โหติ ยถาภูตญาณทสฺสนํ, ยถาภูตญาณทสกนิปาตปาฬิ \csromanfolio{262}ทสฺสเน สติ ยถาภูตญาณทสฺสนสมฺปนฺนสฺส อุปนิสสมฺปนฺโน โหติ นิพฺพิทาวิราโค, นิพฺพิทาวิราเค สติ นิพฺพิทาวิราคสมฺปนฺนสฺส อุปนิสสมฺปนฺนํ โหติ วิมุตฺติญาณทสฺสนํ.\csromanspacer{} เสยฺยถาปิ อาวุโส รุกฺโข สาขาปลาสสมฺปนฺโน, ตสฺส ปปฏิกาปิ ปาริปูริํ คจฺฉติ, ตโจปิ.\csromanspacer{} เผคฺคุปิ.\csromanspacer{} สาโรปิ ปาริปูริํ คจฺฉติ.\csromanspacer{} เอวเมวํ โข อาวุโส สีลวโต สีลสมฺปนฺนสฺส อุปนิสสมฺปนฺโน โหติ อวิปฺปฏิสาโร, อวิปฺปฏิสาเร สติ อวิปฺปฏิสารสมฺปนฺนสฺส อุปนิสสมฺปนฺนํ โหติ ฯเปฯ วิมุตฺติญาณทสฺสนนฺติ.\csromanspacer{}\csromanspacer{}ปญฺจมํ.}

\csromanheader{2}{6. สมาธิสุตฺต}
\csromanmatikaanchor{mtk.26}
\csromantocmarkref{subsection}{4. ทุติย-อุปนิสสุตฺต}{mtk.26}
\bookmark[dest={mtk.26},level=2]{4. ทุติย-อุปนิสสุตฺต}
\proseitem{6}{\csromansymbolfootnote{*}{อุปริ 551 ปิฏฺเฐปิ.}อถ โข อายสฺมา อานนฺโท เยน ภควา เตนุปสงฺกมิ ฯเปฯ เอกมนฺตํ นิสินฺโน โข อายสฺมา อานนฺโท ภควนฺตํ เอตทโวจ “สิยา นุ โข ภนฺเต ภิกฺขุโน ตถารูโป สมาธิปฏิลาโภ, ยถา เนว ปถวิยํ ปถวิสญฺญี\footnote{ปฐวิสญฺญี (สี), ปฐวีสญฺญี (สฺยา)} อสฺส, น อาปสฺมิํ อาโปสญฺญี อสฺส, น เตชสฺมิํ เตโชสญฺญี อสฺส, น วายสฺมิํ วาโยสญฺญี อสฺส, น อากาสานญฺจายตเน อากาสานญฺจายตนสญฺญี อสฺส, น วิญฺญาณญฺจายตเน วิญฺญาณญฺจายตนสญฺญี อสฺส, น อากิญฺจญฺญายตเน อากิญฺจญฺญายตนสญฺญี อสฺส, น เนวสญฺญานาสญฺญายตเน เนวสญฺญานาสญฺญา- ยตนสญฺญี อสฺส, น อิธโลเก อิธโลกสญฺญี อสฺส, น ปรโลเก ปรโลกสญฺญี อสฺส, สญฺญี จ ปน อสฺสา”ติ.\csromanspacer{} สิยา อานนฺท ภิกฺขุโน ตถารูโป สมาธิปฏิลาโภ, ยถา เนว ปถวิยํ ปถวิสญฺญี อสฺส, น อาปสฺมิํ อาโปสญฺญี อสฺส, น เตชสฺมิํ เตโชสญฺญี อสฺส, น วายสฺมิํ วาโยสญฺญี อสฺส, น อากาสานญฺจายตเน อากาสานญฺจายตนสญฺญี อสฺส, น วิญฺญาณญฺจายตเน วิญฺญาณญฺจายตนสญฺญี อสฺส, น อากิญฺจญฺญายตเน อากิญฺจญฺญายตนสญฺญี อสฺส, น เนวสญฺญานาสญฺญายตเน เนวสญฺญานาสญฺญายตนสญฺญี อสฺส, น อิธโลเก อิธโลกสญฺญี อสฺส, น ปรโลเก ปรโลกสญฺญี อสฺส, สญฺญี จ ปน อสฺสาติ.}

\prose{ยถา กถํ ปน ภนฺเต สิยา ภิกฺขุโน ตถารูโป สมาธิปฏิลาโภ, ยถา เนว ปถวิยํ ปถวิสญฺญี อสฺส, น อาปสฺมิํ อาโปสญฺญี อสฺส, น เตชสฺมิํ เตโชสญฺญี อสฺส, น วายสฺมิํ วาโยสญฺญี อสฺส, น อากาสานญฺจายตเน อากาสานญฺจายตนสญฺญี อสฺส,}

\chapterheadpageread{1. อานิสํสวคฺค}
\prose{\csromanfolio{263}น วิญฺญาณญฺจายตเน วิญฺญาณญฺจายตนสญฺญี อสฺส, น อากิญฺจญฺญายตเน อากิญฺจญฺญายตนสญฺญี อสฺส, น เนวสญฺญานาสญฺญายตเน เนวสญฺญานาสญฺญา- ยตนสญฺญี อสฺส, น อิธโลเก อิธโลกสญฺญี อสฺส, น ปรโลเก ปรโลกสญฺญี อสฺส, สญฺญี จ ปน อสฺสาติ.}

\csromanmatikaanchor{mtk.27}
\csromantocmarkref{subsection}{5. ตติย-อุปนิสสุตฺต}{mtk.27}
\bookmark[dest={mtk.27},level=2]{5. ตติย-อุปนิสสุตฺต}
\prose{อิธานนฺท ภิกฺขุ เอวํ สญฺญี โหติ “เอตํ สนฺตํ เอตํ ปณีตํ, ยทิทํ สพฺพสงฺขารสมโถ สพฺพูปธิปฏินิสฺสคฺโค ตณฺหากฺขโย วิราโค นิโรโธ นิพฺพานนฺ”ติ.\csromanspacer{} เอวํ โข อานนฺท สิยา ภิกฺขุโน ตถารูโป สมาธิปฏิลาโภ, ยถา เนว ปถวิยํ ปถวิสญฺญี อสฺส, น อาปสฺมิํ อาโปสญฺญี อสฺส, น เตชสฺมิํ เตโชสญฺญี อสฺส, น วายสฺมิํ วาโยสญฺญี อสฺส, น อากาสานญฺจายตเน อากาสานญฺจายตนสญฺญี อสฺส, น วิญฺญาณญฺจายตเน วิญฺญาณญฺจายตนสญฺญี อสฺส, น อากิญฺจญฺญายตเน อากิญฺจญฺญายตนสญฺญี อสฺส, น เนวสญฺญานาสญฺญายตเน เนวสญฺญานาสญฺญายตนสญฺญี อสฺส, น อิธโลเก อิธโลกสญฺญี อสฺส, น ปรโลเก ปรโลกสญฺญี อสฺส, สญฺญี จ ปน อสฺสาติ.\csromanspacer{}\csromanspacer{}ฉฏฺฐํ.}

\csromanheader{2}{7. สาริปุตฺตสุตฺต}
\proseitem{7}{อถ โข อายสฺมา อานนฺโท เยนายสฺมา สาริปุตฺโต เตนุปสงฺกมิ, อุปสงฺกมิตฺวา อายสฺมตา สาริปุตฺเตน สทฺธิํ สมฺโมทิ, สมฺโมทนียํ กถํ สารณียํ วีติสาเรตฺวา เอกมนฺตํ นิสีทิ, เอกมนฺตํ นิสินฺโน โข อายสฺมา อานนฺโท อายสฺมนฺตํ สาริปุตฺตํ เอตทโวจ\csromandash{}}

\prose{สิยา นุ โข อาวุโส สาริปุตฺต ภิกฺขุโน ตถารูโป สมาธิปฏิลาโภ, ยถา เนว ปถวิยํ ปถวิสญฺญี อสฺส, น อาปสฺมิํ อาโปสญฺญี อสฺส, น เตชสฺมิํ เตโชสญฺญี อสฺส, น วายสฺมิํ วาโยสญฺญี อสฺส, น อากาสานญฺจายตเน อากาสานญฺจายตนสญฺญี อสฺส, น วิญฺญาณญฺจายตเน วิญฺญาณญฺจายตนสญฺญี อสฺส, น อากิญฺจญฺญายตเน อากิญฺจญฺญายตนสญฺญี อสฺส, น เนวสญฺญานาสญฺญา- ยตเน เนวสญฺญานาสญฺญายตนสญฺญี อสฺส, น อิธโลเก อิธโลกสญฺญี อสฺส, น ปรโลเก ปรโลกสญฺญี อสฺส, สญฺญี จ ปน อสฺสาติ.}

\csromanmatikaanchor{mtk.28}
\csromantocmarkref{subsection}{6. สมาธิสุตฺต}{mtk.28}
\bookmark[dest={mtk.28},level=2]{6. สมาธิสุตฺต}
\gambhira{ทสกนิปาตปาฬิ}
\prose{\csromanfolio{264}สิยา อาวุโส อานนฺท ภิกฺขุโน ตถารูโป สมาธิปฏิลาโภ, ยถา เนว ปถวิยํ ปถวิสญฺญี อสฺส ฯเปฯ น ปรโลเก ปรโลกสญฺญี อสฺส, สญฺญี จ ปน อสฺสาติ.}

\prose{ยถา กถํ ปนาวุโส สาริปุตฺต สิยา ภิกฺขุโน ตถารูโป สมาธิปฏิลาโภ, ยถา เนว ปถวิยํ ปถวิสญฺญี อสฺส ฯเปฯ สญฺญี จ ปน อสฺสาติ.\csromanspacer{} เอกมิทาหํ อาวุโส อานนฺท สมยํ อิเธว สาวตฺถิยํ วิหรามิ อนฺธวนสฺมิํ, ตตฺถาหํ\footnote{อถาหํ (ก)} ตถารูปํ สมาธิํ สมาปชฺชิํ\footnote{ปฏิลภามิ (ก)}, ยถา เนว ปถวิยํ ปถวิสญฺญี อโหสิํ, น อาปสฺมิํ อาโปสญฺญี อโหสิํ, น เตชสฺมิํ เตโชสญฺญี อโหสิํ, น วายสฺมิํ วาโยสญฺญี อโหสิํ, น อากาสานญฺจายตเน อากาสานญฺจายตนสญฺญี อโหสิํ, น วิญฺญาณญฺจายตเน วิญฺญาณญฺจายตนสญฺญี อโหสิํ, น อากิญฺจญฺญายตเน อากิญฺจญฺญายตนสญฺญี อโหสิํ, น เนวสญฺญานาสญฺญายตเน เนวสญฺญานาสญฺญายตนสญฺญี อโหสิํ, น อิธโลเก อิธโลกสญฺญี อโหสิํ, น ปรโลเก ปรโลกสญฺญี อโหสิํ, สญฺญี จ ปน อโหสินฺติ.}

\prose{กิํ สญฺญี ปนายสฺมา สาริปุตฺโต\footnote{กิํสญฺญี ปนาวุโส สาริปุตฺต (ก)} ตสฺมิํ สมเย อโหสีติ. “ภวนิโรโธ นิพฺพานํ ภวนิโรโธ นิพฺพานนฺ”ติ โข เม อาวุโส อญฺญาว สญฺญา อุปฺปชฺชติ, อญฺญาว สญฺญา นิรุชฺฌติ.\csromanspacer{} เสยฺยถาปิ อาวุโส สกลิกคฺคิสฺส ฌายมานสฺส อญฺญาว อจฺจิ อุปฺปชฺชติ, อญฺญาว อจฺจิ นิรุชฺฌติ.\csromanspacer{} เอวเมวํ โข อาวุโส “ภวนิโรโธ นิพฺพานํ ภวนิโรโธ นิพฺพานนฺ”ติ อญฺญาว สญฺญา อุปฺปชฺชติ, อญฺญาว สญฺญา นิรุชฺฌติ.\csromanspacer{} ภวนิโรโธ นิพฺพานนฺติ\footnote{นิพฺพานํ (สี, ก)} สญฺญี จ ปนาหํ อาวุโส ตสฺมิํ สมเย อโหสินฺติ.\csromanspacer{}\csromanspacer{}สตฺตมํ.}

\csromanheader{2}{8. ฌานสุตฺต}
\proseitem{8}{สทฺโธ จ\footnote{อิมสฺมิํ วาเกฺย อยํ จ-กาโร นตฺถิ สฺยามโปตฺถเก.} ภิกฺขเว ภิกฺขุ โหติ โน จ\footnote{โน (สฺยา) เอวมุปริปิ. เหฏฺฐา 133 ปิฏฺเฐปิ.} สีลวา, เอวํ โส เตนงฺเคน อปริปูโร โหติ.\csromanspacer{} เตน ตํ องฺคํ ปริปูเรตพฺพํ “กินฺตาหํ สทฺโธ จ อสฺสํ สีลวา จา”ติ.\csromanspacer{} ยโต จ โข ภิกฺขเว ภิกฺขุ สทฺโธ จ โหติ สีลวา จ, เอวํ โส เตนงฺเคน ปริปูโร โหติ.}

\csromanmatikaanchor{mtk.29}
\csromantocmarkref{subsection}{7. สาริปุตฺตสุตฺต}{mtk.29}
\bookmark[dest={mtk.29},level=2]{7. สาริปุตฺตสุตฺต}
\chapterheadpageread{1. อานิสํสวคฺค}
\prose{\csromanfolio{265}สทฺโธ จ ภิกฺขเว ภิกฺขุ โหติ สีลวา จ โน จ พหุสฺสุโต ฯเปฯ พหุสฺสุโต จ โน จ ธมฺมกถิโก.\csromanspacer{} ธมฺมกถิโก จ โน จ ปริสาวจโร.\csromanspacer{} ปริสาวจโร จ โน จ วิสารโท ปริสาย ธมฺมํ เทเสติ.\csromanspacer{} วิสารโท จ ปริสาย ธมฺมํ เทเสติ โน จ วินยธโร.\csromanspacer{} วินยธโร จ โน จ อารญฺญิโก\footnote{อารญฺญโก (ก)} ปนฺตเสนาสโน.\csromanspacer{} อารญฺญิโก จ ปนฺตเสนาสโน โน จ จตุนฺนํ ฌานานํ อาภิเจตสิกานํ ทิฏฺฐธมฺมสุขวิหารานํ นิกามลาภี โหติ อกิจฺฉลาภี อกสิรลาภี.\csromanspacer{} จตุนฺนญฺจ ฌานานํ อาภิเจตสิกานํ ทิฏฺฐธมฺมสุขวิหารานํ นิกามลาภี โหติ อกิจฺฉลาภี อกสิรลาภี โน จ อาสวานํ ขยา อนาสวํ เจโตวิมุตฺติํ ปญฺญาวิมุตฺติํ ทิฏฺเฐว ธมฺเม สยํ อภิญฺญา สจฺฉิกตฺวา อุปสมฺปชฺช วิหรติ, เอวํ โส เตนงฺเคน อปริปูโร โหติ.\csromanspacer{} เตน ตํ องฺคํ ปริปูเรตพฺพํ “กินฺตาหํ สทฺโธ จ อสฺสํ สีลวา จ พหุสฺสุโต จ ธมฺมกถิโก จ ปริสาวจโร จ, วิสารโท จ ปริสาย ธมฺมํ เทเสยฺยํ, วินยธโร จ อารญฺญิโก จ ปนฺตเสนาสโน, จตุนฺนญฺจ ฌานานํ อาภิเจตสิกานํ ทิฏฺฐธมฺมสุขวิหารานํ นิกามลาภี อสฺสํ อกิจฺฉลาภี อกสิรลาภี, อาสวานญฺจ ขยา อนาสวํ เจโตวิมุตฺติํ ปญฺญาวิมุตฺติํ ทิฏฺเฐว ธมฺเม สยํ อภิญฺญา สจฺฉิกตฺวา อุปสมฺปชฺช วิหเรยฺยนฺ”ติ.}

\prose{ยโต จ โข ภิกฺขเว ภิกฺขุ สทฺโธ จ โหติ สีลวา จ พหุสฺสุโต จ ธมฺมกถิโก จ ปริสาวจโร จ, วิสารโท จ ปริสาย ธมฺมํ เทเสติ, วินยธโร จ อารญฺญิโก จ ปนฺตเสนาสโน, จตุนฺนญฺจ ฌานานํ อาภิเจตสิกานํ ทิฏฺฐธมฺมสุขวิหารานํ นิกามลาภี โหติ อกิจฺฉลาภี อกสิรลาภี, อาสวานญฺจ ขยา อนาสวํ เจโตวิมุตฺติํ ปญฺญาวิมุตฺติํ ทิฏฺเฐว ธมฺเม สยํ อภิญฺญา สจฺฉิกตฺวา อุปสมฺปชฺช วิหรติ.\csromanspacer{} เอวํ โส เตนงฺเคน ปริปูโร โหติ.\csromanspacer{} อิเมหิ โข ภิกฺขเว ทสหิ ธมฺเมหิ สมนฺนาคโต ภิกฺขุ สมนฺตปาสาทิโก จ โหติ สพฺพาการปริปูโร จาติ.\csromanspacer{}\csromanspacer{}อฏฺฐมํ.}

\csromanheader{2}{9. สนฺตวิโมกฺขสุตฺต}
\proseitem{9}{สทฺโธ จ ภิกฺขเว ภิกฺขุ โหติ โน จ สีลวา ฯเปฯ สีลวา จ โน จ พหุสฺสุโต.\csromanspacer{} พหุสฺสุโต จ โน จ ธมฺมกถิโก, ธมฺมกถิโก จ}

\vspace{\tipitakaparskip}
\csromancenter{ทสกนิปาตปาฬิ โน จ ปริสาวจโร.\csromanspacer{} ปริสาวจโร จ โน จ วิสารโท ปริสาย ธมฺมํ เทเสติ.\csromanspacer{} วิสารโท จ ปริสาย ธมฺมํ เทเสติ โน จ วินยธโร.\csromanspacer{} วินยธโร จ โน จ อารญฺญิโก ปนฺตเสนาสโน.\csromanspacer{} อารญฺญิโก จ ปนฺตเสนาสโน โน จ เย เต สนฺตา วิโมกฺขา อติกฺกมฺม รูเป อารุปฺปา, เต กาเยน ผุสิตฺวา วิหรติ.\csromanspacer{} เย เต สนฺตา วิโมกฺขา อติกฺกมฺม รูเป อารุปฺปา, เต จ กาเยน ผุสิตฺวา วิหรติ, โน จ อาสวานํ ขยา อนาสวํ เจโตวิมุตฺติํ ปญฺญาวิมุตฺติํ ทิฏฺเฐว ธมฺเม สยํ อภิญฺญา สจฺฉิกตฺวา อุปสมฺปชฺช วิหรติ, เอวํ โส เตนงฺเคน อปริปูโร โหติ.\csromanspacer{} เตน ตํ องฺคํ ปริปูเรตพฺพํ “กินฺตาหํ สทฺโธ จ อสฺสํ สีลวา จ พหุสฺสุโต จ ธมฺมกถิโก จ ปริสาวจโร จ, วิสารโท จ ปริสาย ธมฺมํ เทเสยฺยํ, วินยธโร จ อารญฺญิโก จ ปนฺตเสนาสโน, เย เต สนฺตา วิโมกฺขา อติกฺกมฺม รูเป อารุปฺปา, เต จ กาเยน ผุสิตฺวา วิหเรยฺยํ, อาสวานญฺจ ขยา อนาสวํ เจโตวิมุตฺติํ ปญฺญาวิมุตฺติํ ทิฏฺเฐว ธมฺเม สยํ อภิญฺญา สจฺฉิกตฺวา อุปสมฺปชฺช วิหเรยฺยนฺ”ติ.}
\prose{\csromanfolio{266}ยโต จ โข ภิกฺขเว ภิกฺขุ สทฺโธ จ โหติ สีลวา จ พหุสฺสุโต จ ธมฺมกถิโก จ ปริสาวจโร จ, วิสารโท จ ปริสาย ธมฺมํ เทเสติ, วินยธโร จ อารญฺญิโก จ ปนฺตเสนาสโน, เย เต สนฺตา วิโมกฺขา อติกฺกมฺม รูเป อารุปฺปา, เต จ กาเยน ผุสิตฺวา วิหรติ, อาสวานญฺจ ขยา อนาสวํ เจโตวิมุตฺติํ ปญฺญาวิมุตฺติํ ทิฏฺเฐว ธมฺเม สยํ อภิญฺญา สจฺฉิกตฺวา อุปสมฺปชฺช วิหรติ.\csromanspacer{} เอวํ โส เตนงฺเคน ปริปูโร โหติ.\csromanspacer{} อิเมหิ โข ภิกฺขเว ทสหิ ธมฺเมหิ สมนฺนาคโต ภิกฺขุ สมนฺตปาสาทิโก จ โหติ สพฺพาการปริปูโร จาติ.\csromanspacer{}\csromanspacer{}นวมํ.}

\csromanmatikaanchor{mtk.30}
\csromantocmarkref{subsection}{8. ฌานสุตฺต}{mtk.30}
\bookmark[dest={mtk.30},level=2]{8. ฌานสุตฺต}
\csromanheader{2}{10. วิชฺชาสุตฺต}
\proseitem{10}{สทฺโธ จ ภิกฺขเว ภิกฺขุ โหติ โน จ สีลวา, เอวํ โส เตนงฺเคน อปริปูโร โหติ.\csromanspacer{} เตน ตํ องฺคํ ปริปูเรตพฺพํ “กินฺตาหํ สทฺโธ จ อสฺสํ สีลวา จา”ติ.\csromanspacer{} ยโต จ โข ภิกฺขเว ภิกฺขุ สทฺโธ จ โหติ สีลวา จ, เอวํ โส เตนงฺเคน ปริปูโร โหติ.}

\prose{สทฺโธ จ ภิกฺขเว ภิกฺขุ โหติ สีลวา จ โน จ พหุสฺสุโต, พหุสฺสุโต จ โน จ ธมฺมกถิโก ฯเปฯ ธมฺมกถิโก จ โน จ ปริสาวจโร.\csromanspacer{} ปริสาวจโร จ โน จ วิสารโท ปริสาย ธมฺมํ เทเสติ.}

\chapterheadpageread{1. อานิสํสวคฺค}
\csromanmatikaanchor{mtk.31}
\csromanmatikaanchor{mtk.33}
\csromantocmarkref{subsection}{9. สนฺตวิโมกฺขสุตฺต}{mtk.31}
\bookmark[dest={mtk.31},level=2]{9. สนฺตวิโมกฺขสุตฺต}
\prose{\csromanfolio{267}วิสารโท จ ปริสาย ธมฺมํ เทเสติ โน จ วินยธโร.\csromanspacer{} วินยธโร จ โน จ อเนกวิหิตํ ปุพฺเพนิวาสํ อนุสฺสรติ, เสยฺยถิทํ, เอกมฺปิ ชาติํ เทฺวปิ ชาติโย ฯเปฯ อิติ สาการํ ส-อุทฺเทสํ อเนกวิหิตํ ปุพฺเพนิวาสํ อนุสฺสรติ.\csromanspacer{} อเนกวิหิตญฺจ ฯเปฯ ปุพฺเพนิวาสํ.\csromanspacer{} อนุสฺสรติ โน จ ทิพฺเพน จกฺขุนา วิสุทฺเธน อติกฺกนฺตมานุสเกน ฯเปฯ ยถากมฺมูปเค สตฺเต ปชานาติ.\csromanspacer{} ทิพฺเพน จ จกฺขุนา วิสุทฺเธน อติกฺกนฺตมานุสเกน ฯเปฯ ยถากมฺมูปเค สตฺเต ปชานาติ, โน จ อาสวานํ ขยา ฯเปฯ สจฺฉิกตฺวา อุปสมฺปชฺช วิหรติ, เอวํ โส เตนงฺเคน อปริปูโร โหติ.\csromanspacer{} เตน ตํ องฺคํ ปริปูเรตพฺพํ “กินฺตาหํ สทฺโธ จ อสฺสํ สีลวา จ พหุสฺสุโต จ ธมฺมกถิโก จ ปริสาวจโร จ, วิสารโท จ ปริสาย ธมฺมํ เทเสยฺยํ, วินยธโร จ, อเนกวิหิตญฺจ ปุพฺเพนิวาสํ อนุสฺสเรยฺยํ, เสยฺยถิทํ, เอกมฺปิ ชาติํ เทฺวปิ ชาติโย ฯเปฯ อิติ สาการํ ส-อุทฺเทสํ อเนกวิหิตํ ปุพฺเพนิวาสํ อนุสฺสเรยฺยํ, ทิพฺเพน จ จกฺขุนา วิสุทฺเธน อติกฺกนฺตมานุสเกน ฯเปฯ ยถากมฺมูปเค สตฺเต ปชาเนยฺยํ, อาสวานญฺจ ขยา ฯเปฯ สจฺฉิกตฺวา อุปสมฺปชฺช วิหเรยฺยนฺ”ติ.}

\prose{ยโต จ โข ภิกฺขเว ภิกฺขุ สทฺโธ จ โหติ สีลวา จ พหุสฺสุโต จ ธมฺมกถิโก จ ปริสาวจโร จ, วิสารโท จ ปริสาย ธมฺมํ เทเสติ, วินยธโร จ อเนกวิหิตญฺจ ปุพฺเพนิวาสํ อนุสฺสรติ, เสยฺยถิทํ, เอกมฺปิ ชาติํ เทฺวปิ ชาติโย ฯเปฯ อิติ สาการํ ส-อุทฺเทสํ อเนกวิหิตํ ปุพฺเพนิวาสํ อนุสฺสรติ, ทิพฺเพน จ จกฺขุนา วิสุทฺเธน อติกฺกนฺตมานุสเกน ฯเปฯ ยถากมฺมูปเค สตฺเต ปชานาติ, อาสวานญฺจ ขยา อนาสวํ เจโตวิมุตฺติํ ปญฺญาวิมุตฺติํ ทิฏฺเฐว ธมฺเม สยํ อภิญฺญา สจฺฉิกตฺวา อุปสมฺปชฺช วิหรติ, เอวํ โส เตนงฺเคน ปริปูโร โหติ.\csromanspacer{} อิเมหิ โข ภิกฺขเว ทสหิ ธมฺเมหิ สมนฺนาคโต ภิกฺขุ สมนฺตปาสาทิโก จ โหติ สพฺพาการปริปูโร จาติ.\csromanspacer{}\csromanspacer{}ทสมํ.}

\vspace{\tipitakaparskip}
\csromancenter{อานิสํสวคฺโค ปฐโม.}
\titlehead{\textbf{ตสฺสุทฺทานํ}}
\csromangathameasure{กิมตฺถิยํ เจตนา จ, ตโย อุปนิสาปิ จ. \\ สมาธิ สาริปุตฺโต จ, ฌานํ สนฺเตน วิชฺชยาติ.}
\csromangathagroup{\csromangathastanza{กิมตฺถิยํ เจตนา จ, ตโย อุปนิสาปิ จ. \\ สมาธิ สาริปุตฺโต จ, ฌานํ สนฺเตน วิชฺชยาติ.}}

\csromanmatikaanchor{mtk.32}
\csromantocmarkref{subsection}{10. วิชฺชาสุตฺต}{mtk.32}
\bookmark[dest={mtk.32},level=2]{10. วิชฺชาสุตฺต}
\csromantocmarkref{section}{อุทฺทานคาถา}{mtk.33}
\bookmark[dest={mtk.33},level=1]{อุทฺทานคาถา}
\csromanheader{1}{ทสกนิปาตปาฬิ \textbf{2. นาถวคฺค}}
\csromanheader{2}{1. เสนาสนสุตฺต}
\proseitem{11}{\csromanfolio{268}ปญฺจงฺคสมนฺนาคโต ภิกฺขเว ภิกฺขุ ปญฺจงฺคสมนฺนาคตํ เสนาสนํ เสวมาโน ภชมาโน นจิรสฺเสว อาสวานํ ขยา อนาสวํ เจโตวิมุตฺติํ ปญฺญาวิมุตฺติํ ทิฏฺเฐว ธมฺเม สยํ อภิญฺญา สจฺฉิกตฺวา อุปสมฺปชฺช วิหเรยฺย.}

\prose{กถญฺจ ภิกฺขเว ภิกฺขุ ปญฺจงฺคสมนฺนาคโต โหติ? อิธ ภิกฺขเว ภิกฺขุ สทฺโธ โหติ, สทฺทหติ ตถาคตสฺส โพธิํ “อิติปิ โส ภควา ฯเปฯ ภควา”ติ, อปฺปาพาโธ โหติ อปฺปาตงฺโก สมเวปากินิยา คหณิยา สมนฺนาคโต นาติสีตาย นาจฺจุณฺหาย มชฺฌิมาย ปธานกฺขมาย, อสโฐ โหติ อมายาวี ยถาภูตํ อตฺตานํ อาวิกตฺตา สตฺถริ วา วิญฺญูสุ วา สพฺรหฺมจารีสุ, อารทฺธวีริโย วิหรติ อกุสลานํ ธมฺมานํ ปหานาย กุสลานํ ธมฺมานํ อุปสมฺปทาย, ถามวา ทฬฺหปรกฺกโม อนิกฺขิตฺตธุโร กุสเลสุ ธมฺเมสุ, ปญฺญวา โหติ อุทยตฺถคามินิยา ปญฺญาย สมนฺนาคโต อริยาย นิพฺเพธิกาย สมฺมา ทุกฺขกฺขยคามินิยา.\csromanspacer{} เอวํ โข ภิกฺขเว ภิกฺขุ ปญฺจงฺคสมนฺนาคโต โหติ.}

\prose{กถญฺจ ภิกฺขเว เสนาสนํ ปญฺจงฺคสมนฺนาคตํ โหติ? อิธ ภิกฺขเว เสนาสนํ นาติทูรํ โหติ นาจฺจาสนฺนํ คมนาคมนสมฺปนฺนํ ทิวา อปฺปากิณฺณํ รตฺติํ อปฺปสทฺทํ อปฺปนิคฺโฆสํ อปฺปฑํสมกสวาตาตปสรีสปสมฺผสฺสํ\footnote{อปฺปฑํส...สิริํสปสมฺผสฺสํ (สี, สฺยา, อิ)}, ตสฺมิํ โข ปน เสนาสเน วิหรนฺตสฺส อปฺปกสิเรน อุปฺปชฺชนฺติ จีวรปิณฺฑปาตเสนาสนคิลานปจฺจยเภสชฺชปริกฺขารา, ตสฺมิํ โข ปน เสนาสเน เถรา ภิกฺขู วิหรนฺติ พหุสฺสุตา อาคตาคมา ธมฺมธรา วินยธรา มาติกาธรา, เต กาเลน กาลํ อุปสงฺกมิตฺวา ปริปุจฺฉติ ปริปญฺหติ “อิทํ ภนฺเต กถํ อิมสฺส โก อตฺโถ”ติ.\csromanspacer{} ตสฺส เต อายสฺมนฺโต อวิวฏญฺเจว วิวรนฺติ, อนุตฺตานีกตญฺจ อุตฺตานิํ กโรนฺติ, อเนกวิหิเตสุ จ กงฺขาฐานิเยสุ ธมฺเมสุ กงฺขํ ปฏิวิโนเทนฺติ.\csromanspacer{} เอวํ โข ภิกฺขเว เสนาสนํ ปญฺจงฺคสมนฺนาคตํ โหติ.\csromanspacer{} ปญฺจงฺคสมนฺนาคโต โข ภิกฺขเว ภิกฺขุ ปญฺจงฺคสมนฺนาคตํ เสนาสนํ เสวมาโน ภชมาโน}

\vspace{\tipitakaparskip}
\csromancenter{นาถวคฺค นจิรสฺเสว อาสวานํ ขยา ฯเปฯ สจฺฉิกตฺวา อุปสมฺปชฺช วิหเรยฺยาติ.\csromanspacer{}\csromanspacer{}ปฐมํ.}
\csromanheader{2}{2. ปญฺจงฺคสุตฺต}
\proseitem{12}{\csromanfolio{269}ปญฺจงฺควิปฺปหีโน ภิกฺขเว ภิกฺขุ ปญฺจงฺคสมนฺนาคโต อิมสฺมิํ ธมฺมวินเย เกวลี วุสิตวา “อุตฺตมปุริโส”ติ วุจฺจติ.\csromanspacer{} กถญฺจ ภิกฺขเว ภิกฺขุ ปญฺจงฺควิปฺปหีโน โหติ? อิธ ภิกฺขเว ภิกฺขุโน กามจฺฉนฺโท ปหีโน โหติ, พฺยาปาโท ปหีโน โหติ, ถินมิทฺธํ\footnote{ถีนมิทฺธํ (สี, สฺยา, อิ)} ปหีนํ โหติ, อุทฺธจฺจกุกฺกุจฺจํ ปหีนํ โหติ, วิจิกิจฺฉา ปหีนา โหติ.\csromanspacer{} เอวํ โข ภิกฺขเว ภิกฺขุ ปญฺจงฺควิปฺปหีโน โหติ.}

\prose{กถญฺจ ภิกฺขเว ภิกฺขุ ปญฺจงฺคสมนฺนาคโต โหติ? อิธ ภิกฺขเว ภิกฺขุ อเสเขน สีลกฺขนฺเธน สมนฺนาคโต โหติ, อเสเขน สมาธิกฺขนฺเธน สมนฺนาคโต โหติ, อเสเขน ปญฺญากฺขนฺเธน สมนฺนาคโต โหติ, อเสเขน วิมุตฺติกฺขนฺเธน สมนฺนาคโต โหติ, อเสเขน วิมุตฺติญาณทสฺสนกฺขนฺเธน สมนฺนาคโต โหติ.\csromanspacer{} เอวํ โข ภิกฺขเว ภิกฺขุ ปญฺจงฺคสมนฺนาคโต โหติ.}

\csromanmatikaanchor{mtk.34}
\csromantocmarknopage{section}{2. นาถวคฺค}{mtk.34}
\bookmark[dest={mtk.34},level=1]{2. นาถวคฺค}
\prose{ปญฺจงฺควิปฺปหีโน โข ภิกฺขเว ภิกฺขุ ปญฺจงฺคสมนฺนาคโต อิมสฺมิํ ธมฺมวินเย เกวลี วุสิตวา “อุตฺตมปุริโส”ติ วุจฺจติ.}

\csromanmatikaanchor{mtk.35}
\csromantocmarkref{subsection}{1. เสนาสนสุตฺต}{mtk.35}
\bookmark[dest={mtk.35},level=2]{1. เสนาสนสุตฺต}
\csromangathameasure{กามจฺฉนฺโท จ พฺยาปาโท, ถินมิทฺธญฺจ ภิกฺขุโน. \\ อุทฺธจฺจํ วิจิกิจฺฉา จ, สพฺพโสว น วิชฺชติ. \\ อเสเขน จ สีเลน, อเสเขน สมาธินา. \\ วิมุตฺติยา จ สมฺปนฺโน, ญาเณน จ ตถาวิโธ. \\ ส เว ปญฺจงฺคสมฺปนฺโน, ปญฺจ องฺเค วิวชฺชยํ. \\ อิมสฺมิํ ธมฺมวินเย, “เกวลี” อิติ วุจฺจตีติ.ทุติยํ.}
\csromangathagroup{\csromangathastanza{กามจฺฉนฺโท จ พฺยาปาโท, ถินมิทฺธญฺจ ภิกฺขุโน. \\ อุทฺธจฺจํ วิจิกิจฺฉา จ, สพฺพโสว น วิชฺชติ.}\csromangathastanza{อเสเขน จ สีเลน, อเสเขน สมาธินา. \\ วิมุตฺติยา จ สมฺปนฺโน, ญาเณน จ ตถาวิโธ.}\csromangathastanza{ส เว ปญฺจงฺคสมฺปนฺโน, ปญฺจ องฺเค\footnote{ปญฺจงฺคานิ (สฺยา)} วิวชฺชยํ\footnote{วิวชฺชิย (ก)}. \\ อิมสฺมิํ ธมฺมวินเย, “เกวลี” อิติ วุจฺจตีติ.\csromanspacer{}\csromanspacer{}ทุติยํ.}}

\csromanheader{2}{3. สํโยชนสุตฺต}
\proseitem{13}{ทสยิมานิ ภิกฺขเว สํโยชนานิ.\csromanspacer{} กตมานิ ทส? ปญฺโจรมฺภาคิยานิ สํโยชนานิ ปญฺจุทฺธมฺภาคิยานิ สํโยชนานิ.\csromanspacer{} กตมานิ}

\csromanmatikaanchor{mtk.36}
\csromantocmarkref{subsection}{2. ปญฺจงฺคสุตฺต}{mtk.36}
\bookmark[dest={mtk.36},level=2]{2. ปญฺจงฺคสุตฺต}
\vspace{\tipitakaparskip}
\csromancenter{ทสกนิปาตปาฬิ ปญฺโจรมฺภาคิยานิ สํโยชนานิ? สกฺกายทิฏฺฐิ วิจิกิจฺฉา สีลพฺพตปรามาโส กามจฺฉนฺโท พฺยาปาโท.\csromanspacer{} อิมานิ ปญฺโจรมฺภาคิยานิ สํโยชนานิ.}
\csromanmatikaanchor{mtk.37}
\csromanmatikaanchor{mtk.38}
\csromantocmarkref{subsection}{3. สํโยชนสุตฺต}{mtk.37}
\bookmark[dest={mtk.37},level=2]{3. สํโยชนสุตฺต}
\csromantocmarkref{subsection}{4. เจโตขีลสุตฺต}{mtk.38}
\bookmark[dest={mtk.38},level=2]{4. เจโตขีลสุตฺต}
\prose{\csromanfolio{270}กตมานิ ปญฺจุทฺธมฺภาคิยานิ สํโยชนานิ? รูปราโค อรูปราโค มาโน อุทฺธจฺจํ อวิชฺชา.\csromanspacer{} อิมานิ ปญฺจุทฺธมฺภาคิยานิ สํโยชนานิ.\csromanspacer{} อิมานิ โข ภิกฺขเว ทส สํโยชนานีติ.\csromanspacer{}\csromanspacer{}ตติยํ.}

\csromanheader{2}{4. เจโตขิลสุตฺต}
\proseitem{14}{ยสฺส กสฺสจิ ภิกฺขเว ภิกฺขุสฺส วา ภิกฺขุนิยา วา ปญฺจ เจโตขิลา อปฺปหีนา ปญฺจ เจตโสวินิพนฺธา อสมุจฺฉินฺนา, ตสฺส ยา รตฺติ วา ทิวโส วา อาคจฺฉติ, หานิเยว ปาฏิกงฺขา กุสเลสุ ธมฺเมสุ, โน วุทฺธิ.}

\prose{กตมสฺส ปญฺจ เจโตขิลา อปฺปหีนา โหนฺติ? อิธ ภิกฺขเว ภิกฺขุ สตฺถริ กงฺขติ วิจิกิจฺฉติ นาธิมุจฺจติ น สมฺปสีทติ, โย โส ภิกฺขเว ภิกฺขุ สตฺถริ กงฺขติ วิจิกิจฺฉติ นาธิมุจฺจติ น สมฺปสีทติ, ตสฺส จิตฺตํ น นมติ อาตปฺปาย อนุโยคาย สาตจฺจาย ปธานาย.\csromanspacer{} ยสฺส จิตฺตํ น นมติ อาตปฺปาย อนุโยคาย สาตจฺจาย ปธานาย, เอวมสฺสายํ ปฐโม เจโตขิโล อปฺปหีโน โหติ.}

\prose{ปุน จปรํ ภิกฺขเว ภิกฺขุ ธมฺเม กงฺขติ ฯเปฯ สํเฆ กงฺขติ.\csromanspacer{} สิกฺขาย กงฺขติ.\csromanspacer{} สพฺรหฺมจารีสุ กุปิโต โหติ อนตฺตมโน อาหตจิตฺโต ขิลชาโต.\csromanspacer{} โย โส ภิกฺขเว ภิกฺขุ สพฺรหฺมจารีสุ กุปิโต โหติ อนตฺตมโน อาหตจิตฺโต ขิลชาโต, ตสฺส จิตฺตํ น นมติ อาตปฺปาย อนุโยคาย สาตจฺจาย ปธานาย.\csromanspacer{} ยสฺส จิตฺตํ น นมติ อาตปฺปาย อนุโยคาย สาตจฺจาย ปธานาย, เอวมสฺสายํ ปญฺจโม เจโตขิโล อปฺปหีโน โหติ.\csromanspacer{} อิมสฺส ปญฺจ เจโตขิลา อปฺปหีนา โหนฺติ.}

\prose{กตมสฺส ปญฺจ เจตโสวินิพนฺธา อสมุจฺฉินฺนา โหนฺติ? อิธ ภิกฺขเว ภิกฺขุ กาเมสุ อวีตราโค โหติ อวิคตจฺฉนฺโท อวิคตเปโม อวิคตปิปาโส อวิคตปริฬาโห อวิคตตโณฺห.\csromanspacer{} โย โส ภิกฺขเว ภิกฺขุ กาเมสุ อวีตราโค โหติ อวิคตจฺฉนฺโท อวิคตเปโม อวิคตปิปาโส อวิคตปริฬาโห}

\vspace{\tipitakaparskip}
\csromancenter{นาถวคฺค อวิคตตโณฺห, ตสฺส จิตฺตํ น นมติ อาตปฺปาย อนุโยคาย สาตจฺจาย ปธานาย.\csromanspacer{} ยสฺส จิตฺตํ น นมติ อาตปฺปาย อนุโยคาย สาตจฺจาย ปธานาย, เอวมสฺสายํ ปฐโม เจตโสวินิพนฺโธ อสมุจฺฉินฺโน โหติ.}
\prose{\csromanfolio{271}ปุน จปรํ ภิกฺขเว ภิกฺขุ กาเย อวีตราโค โหติ ฯเปฯ รูเป อวีตราโค โหติ ฯเปฯ ยาวทตฺถํ อุทราวเทหกํ ภุญฺชิตฺวา เสยฺยสุขํ ปสฺสสุขํ มิทฺธสุขํ อนุยุตฺโต วิหรติ.\csromanspacer{} อญฺญตรํ เทวนิกายํ ปณิธาย พฺรหฺมจริยํ จรติ “อิมินาหํ สีเลน วา วเตน วา ตเปน วา พฺรหฺมจริเยน วา เทโว วา ภวิสฺสามิ เทวญฺญตโร วา”ติ.\csromanspacer{} โย โส ภิกฺขเว ภิกฺขุ อญฺญตรํ เทวนิกายํ ปณิธาย พฺรหฺมจริยํ จรติ “อิมินาหํ สีเลน วา วเตน วา ตเปน วา พฺรหฺมจริเยน วา เทโว วา ภวิสฺสามิ เทวญฺญตโร วา”ติ, ตสฺส จิตฺตํ น นมติ อาตปฺปาย อนุโยคาย สาตจฺจาย ปธานาย.\csromanspacer{} ยสฺส จิตฺตํ น นมติ อาตปฺปาย อนุโยคาย สาตจฺจาย ปธานาย, เอวมสฺสายํ ปญฺจโม เจตโสวินิพนฺโธ อสมุจฺฉินฺโน โหติ.\csromanspacer{} อิมสฺส ปญฺจ เจตโสวินิพนฺธา อสมุจฺฉินฺนา โหนฺติ.}

\prose{ยสฺส กสฺสจิ ภิกฺขเว ภิกฺขุสฺส วา ภิกฺขุนิยา วา อิเม ปญฺจ เจโตขิลา อปฺปหีนา อิเม ปญฺจ เจตโสวินิพนฺธา อสมุจฺฉินฺนา, ตสฺส ยา รตฺติ วา ทิวโส วา อาคจฺฉติ, หานิเยว ปาฏิกงฺขา กุสเลสุ ธมฺเมสุ โน วุทฺธิ.}

\prose{เสยฺยถาปิ ภิกฺขเว กาฬปกฺเข จนฺทสฺส ยา รตฺติ วา ทิวโส วา อาคจฺฉติ, หายเตว วณฺเณน, หายติ มณฺฑเลน, หายติ อาภาย, หายติ อาโรหปริณาเหน.\csromanspacer{} เอวเมวํ โข ภิกฺขเว ยสฺส กสฺสจิ ภิกฺขุสฺส วา ภิกฺขุนิยา วา อิเม ปญฺจ เจโตขิลา อปฺปหีนา, อิเม ปญฺจ เจตโสวินิพนฺธา อสมุจฺฉินฺนา, ตสฺส ยา รตฺติ วา ทิวโส วา อาคจฺฉติ, หานิเยว ปาฏิกงฺขา กุสเลสุ ธมฺเมสุ โน วุทฺธิ.}

\prose{ยสฺส กสฺสจิ ภิกฺขเว ภิกฺขุสฺส วา ภิกฺขุนิยา วา ปญฺจ เจโตขิลา ปหีนา ปญฺจ เจตโสวินิพนฺธา สุสมุจฺฉินฺนา, ตสฺส ยา รตฺติ วา ทิวโส วา อาคจฺฉติ, วุทฺธิเยว ปาฏิกงฺขา กุสเลสุ ธมฺเมสุ โน ปริหานิ.}

\gambhira{ทสกนิปาตปาฬิ}
\prose{\csromanfolio{272}กตมสฺส ปญฺจ เจโตขิลา ปหีนา โหนฺติ? อิธ ภิกฺขเว ภิกฺขุ สตฺถริ น กงฺขติ น วิจิกิจฺฉติ อธิมุจฺจติ สมฺปสีทติ.\csromanspacer{} โย โส ภิกฺขเว ภิกฺขุ สตฺถริ น กงฺขติ น วิจิกิจฺฉติ อธิมุจฺจติ สมฺปสีทติ, ตสฺส จิตฺตํ นมติ อาตปฺปาย อนุโยคาย สาตจฺจาย ปธานาย.\csromanspacer{} ยสฺส จิตฺตํ นมติ อาตปฺปาย อนุโยคาย สาตจฺจาย ปธานาย.\csromanspacer{} เอวมสฺสายํ ปฐโม เจโตขิโล ปหีโน โหติ.}

\prose{ปุน จปรํ ภิกฺขเว ภิกฺขุ ธมฺเม น กงฺขติ ฯเปฯ สํเฆ น กงฺขติ.\csromanspacer{} สิกฺขาย น กงฺขติ.\csromanspacer{} สพฺรหฺมจารีสุ น กุปิโต โหติ อตฺตมโน น อาหตจิตฺโต น ขิลชาโต.\csromanspacer{} โย โส ภิกฺขเว ภิกฺขุ สพฺรหฺมจารีสุ น กุปิโต โหติ อตฺตมโน น อาหตจิตฺโต น ขิลชาโต ตสฺส จิตฺตํ นมติ อาตปฺปาย อนุโยคาย สาตจฺจาย ปธานาย.\csromanspacer{} ยสฺส จิตฺตํ นมติ อาตปฺปาย อนุโยคาย สาตจฺจาย ปธานาย, เอวมสฺสายํ ปญฺจโม เจโตขิโล ปหีโน โหติ.\csromanspacer{} อิมสฺส ปญฺจ เจโตขิลา ปหีนา โหนฺติ.}

\prose{กตมสฺส ปญฺจ เจตโสวินิพนฺธา สุสมุจฺฉินฺนา โหนฺติ, อิธ ภิกฺขเว ภิกฺขุ กาเมสุ วีตราโค โหติ วิคตจฺฉนฺโท วิคตเปโม วิคตปิปาโส วิคตปริฬาโห วิคตตโณฺห.\csromanspacer{} โย โส ภิกฺขเว ภิกฺขุ กาเมสุ วีตราโค โหติ วิคตจฺฉนฺโท วิคตเปโม วิคตปิปาโส วิคตปริฬาโห วิคตตโณฺห, ตสฺส จิตฺตํ นมติ อาตปฺปาย อนุโยคาย สาตจฺจาย ปธานาย.\csromanspacer{} ยสฺส จิตฺตํ นมติ อาตปฺปาย อนุโยคาย สาตจฺจาย ปธานาย, เอวมสฺสายํ ปฐโม เจตโสวินิพนฺโธ สุสมุจฺฉินฺโน โหติ.}

\prose{ปุน จปรํ ภิกฺขเว ภิกฺขุ กาเย วีตราโค โหติ ฯเปฯ รูเป วีตราโค โหติ ฯเปฯ น ยาวทตฺถํ อุทราวเทหกํ ภุญฺชิตฺวา เสยฺยสุขํ ปสฺสสุขํ มิทฺธสุขํ อนุยุตฺโต วิหรติ, น อญฺญตรํ เทวนิกายํ ปณิธาย พฺรหฺมจริยํ จรติ “อิมินาหํ สีเลน วา วเตน วา ตเปน วา พฺรหฺมจริเยน วา เทโว วา ภวิสฺสามิ เทวญฺญตโร วา”ติ.\csromanspacer{} โย โส ภิกฺขเว ภิกฺขุ น อญฺญตรํ เทวนิกายํ ปณิธาย ฯเปฯ เทวญฺญตโร วาติ, ตสฺส จิตฺตํ นมติ อาตปฺปาย อนุโยคาย สาตจฺจาย ปธานาย.\csromanspacer{} ยสฺส จิตฺตํ นมติ อาตปฺปาย อนุโยคาย สาตจฺจาย ปธานาย, เอวมสฺสายํ ปญฺจโม}

\vspace{\tipitakaparskip}
\csromancenter{นาถวคฺค เจตโสวินิพนฺโธ สุสมุจฺฉินฺโน โหติ.\csromanspacer{} อิมสฺส ปญฺจ เจตโสวินิพนฺธา สุสมุจฺฉินฺนา โหนฺติ.}
\prose{\csromanfolio{273}ยสฺส กสฺสจิ ภิกฺขเว ภิกฺขุสฺส วา ภิกฺขุนิยา วา อิเม ปญฺจ เจโตขิลา ปหีนา อิเม ปญฺจ เจตโสวินิพนฺธา สุสมุจฺฉินฺนา, ตสฺส ยา รตฺติ วา ทิวโส วา อาคจฺฉติ, วุทฺธิเยว ปาฏิกงฺขา กุสเลสุ ธมฺเมสุ โน ปริหานิ.}

\prose{เสยฺยถาปิ ภิกฺขเว ชุณฺหปกฺเข จนฺทสฺส ยา รตฺติ วา ทิวโส วา อาคจฺฉติ, วฑฺฒเตว วณฺเณน, วฑฺฒติ มณฺฑเลน, วฑฺฒติ อาภาย, วฑฺฒติ อาโรหปริณาเหน.\csromanspacer{} เอวเมวํ โข ภิกฺขเว ยสฺส กสฺสจิ ภิกฺขุสฺส วา ภิกฺขุนิยา วา อิเม ปญฺจ เจโตขิลา ปหีนา อิเม ปญฺจ เจตโสวินิพนฺธา สุสมุจฺฉินฺนา, ตสฺส ยา รตฺติ วา ทิวโส วา อาคจฺฉติ, วุทฺธิเยว ปาฏิกงฺขา กุสเลสุ ธมฺเมสุ โน ปริหานีติ.\csromanspacer{}\csromanspacer{}จตุตฺถํ.}

\csromanheader{2}{5. อปฺปมาทสุตฺต}
\proseitem{15}{ยาวตา ภิกฺขเว สตฺตา อปทา วา ทฺวิปทา วา จตุปฺปทา วา พหุปฺปทา วา รูปิโน วา อรูปิโน วา สญฺญิโน วา อสญฺญิโน วา เนวสญฺญินาสญฺญิโน วา, ตถาคโต เตสํ อคฺคมกฺขายติ อรหํ สมฺมาสมฺพุทฺโธ.\csromanspacer{} เอวเมวํ โข ภิกฺขเว เย เกจิ กุสลา ธมฺมา, สพฺเพ เต อปฺปมาทมูลกา อปฺปมาทสโมสรณา, อปฺปมาโท เตสํ\footnote{เตสํ ธมฺมานํ (สี, ก) สํ 3. 38 ปิฏฺเฐปิ.} อคฺคมกฺขายติ.}

\prose{เสยฺยถาปิ ภิกฺขเว ยานิ กานิจิ ชงฺคลานํ\footnote{ชงฺคมานํ (สี, อิ) สํ 3. 40 ปิฏฺเฐปิ.} ปาณานํ ปทชาตานิ, สพฺพานิ ตานิ หตฺถิปเท สโมธานํ คจฺฉนฺติ, หตฺถิปทํ เตสํ อคฺคมกฺขายติ ยทิทํ มหนฺตตฺเตน.\csromanspacer{} เอวเมวํ โข ภิกฺขเว เย เกจิ กุสลา ธมฺมา, สพฺเพ เต อปฺปมาทมูลกา อปฺปมาทสโมสรณา, อปฺปมาโท เตสํ1 อคฺคมกฺขายติ.}

\prose{เสยฺยถาปิ ภิกฺขเว กูฏาคารสฺส ยา กาจิ โคปานสิโย, สพฺพา ตา กูฏงฺคมา กูฏนินฺนา กูฏสโมสรณา, กูโฏ ตาสํ อคฺคมกฺขายติ.\csromanspacer{} เอวเมวํ โข ภิกฺขเว เย เกจิ กุสลา ธมฺมา, สพฺเพ เต อปฺปมาทมูลกา อปฺปมาทสโมสรณา, อปฺปมาโท เตสํ อคฺคมกฺขายติ.}

\gambhira{ทสกนิปาตปาฬิ}
\prose{\csromanfolio{274}เสยฺยถาปิ ภิกฺขเว เย เกจิ มูลคนฺธา, กาฬานุสาริยํ เตสํ อคฺคมกฺขายติ.\csromanspacer{} เอวเมวํ โข ภิกฺขเว ฯเปฯ.}

\prose{เสยฺยถาปิ ภิกฺขเว เย เกจิ สารคนฺธา, โลหิตจนฺทนํ เตสํ อคฺคมกฺขายติ.\csromanspacer{} เอวเมวํ โข ภิกฺขเว ฯเปฯ.}

\prose{เสยฺยถาปิ ภิกฺขเว เย เกจิ ปุปฺผคนฺธา, วสฺสิกํ เตสํ อคฺคมกฺขายติ.\csromanspacer{} เอวเมวํ โข ภิกฺขเว ฯเปฯ.}

\prose{เสยฺยถาปิ ภิกฺขเว เย เกจิ ขุทฺทราชาโน\footnote{กุฑฺฑราชาโน (สี, สฺยา, อิ), กุฏฺฏราชาโน, กูฏราชาโน (ก) อํ 2. 320 ปิฏฺเฐปิ.}, สพฺเพ เต รญฺโญ จกฺกวตฺติสฺส อนุยนฺตา ภวนฺติ, ราชา เตสํ จกฺกวตฺตี อคฺคมกฺขายติ.\csromanspacer{} เอวเมวํ โข ภิกฺขเว ฯเปฯ.}

\csromanmatikaanchor{mtk.39}
\csromantocmarkref{subsection}{5. อปฺปมาทสุตฺต}{mtk.39}
\bookmark[dest={mtk.39},level=2]{5. อปฺปมาทสุตฺต}
\prose{เสยฺยถาปิ ภิกฺขเว ยา กาจิ ตารกรูปานํ ปภา, สพฺพา ตา จนฺทปฺปภาย กลํ นาคฺฆนฺติ โสฬสิํ, จนฺทปฺปภา ตาสํ อคฺคมกฺขายติ.\csromanspacer{} เอวเมวํ โข ภิกฺขเว ฯเปฯ.}

\prose{เสยฺยถาปิ ภิกฺขเว สรทสมเย วิทฺเธ วิคตวลาหเก เทเว อาทิจฺโจ นภํ อพฺภุสฺสกฺกมาโน\footnote{อพฺภุสฺสุกฺกมาโน(สี) สํ 3. 41 ปิฏฺเฐปิ.} สพฺพํ อากาสคตํ ตมคตํ อภิวิหจฺจ ภาสเต จ ตปเต จ วิโรจติ จ.\csromanspacer{} เอวเมวํ โข ภิกฺขเว ฯเปฯ.}

\prose{เสยฺยถาปิ ภิกฺขเว ยา กาจิ มหานทิโย.\csromanspacer{} เสยฺยถิทํ, คงฺคา ยมุนา อจิรวตี สรภู มหี, สพฺพา ตา สมุทฺทงฺคมา สมุทฺทนินฺนา สมุทฺทโปณา สมุทฺทปพฺภารา, มหาสมุทฺโท ตาสํ อคฺคมกฺขายติ.\csromanspacer{} เอวเมวํ โข ภิกฺขเว เย เกจิ กุสลา ธมฺมา, สพฺเพ เต อปฺปมาทมูลกา อปฺปมาทสโมสรณา, อปฺปมาโท เตสํ อคฺคมกฺขายตีติ.}

\prose{.\csromanspacer{} ปญฺจมํ.}

\csromanheader{2}{6. อาหุเนยฺยสุตฺต}
\proseitem{16}{ทสยิเม ภิกฺขเว ปุคฺคลา อาหุเนยฺยา ปาหุเนยฺยา ทกฺขิเณยฺยา อญฺชลิกรณียา อนุตฺตรํ ปุญฺญกฺเขตฺตํ โลกสฺส.\csromanspacer{} กตเม ทส? ตถาคโต อรหํ สมฺมาสมฺพุทฺโธ ปจฺเจกพุทฺโธ อุภโตภาควิมุตฺโต ปญฺญาวิมุตฺโต กายสกฺขี ทิฏฺฐิปฺปตฺโต สทฺธาวิมุตฺโต สทฺธานุสารี ธมฺมานุสารี โคตฺรภู.\csromanspacer{} อิเม โข ภิกฺขเว ทส ปุคฺคลา อาหุเนยฺยา ฯเปฯ อนุตฺตรํ ปุญฺญกฺเขตฺตํ โลกสฺสาติ.\csromanspacer{}\csromanspacer{}ฉฏฺฐํ.}

\csromanheader{2}{2. นาถวคฺค \textbf{7. ปฐมนาถสุตฺต}}
\proseitem{17}{\csromanfolio{275}\csromansymbolfootnote{*}{ที 3. 221, 225 ปิฏฺเฐสุปิ.}สนาถา ภิกฺขเว วิหรถ มา อนาถา, ทุกฺขํ ภิกฺขเว อนาโถ วิหรติ.\csromanspacer{} ทสยิเม ภิกฺขเว นาถกรณา ธมฺมา.\csromanspacer{} กตเม ทส? อิธ ภิกฺขเว ภิกฺขุ สีลวา โหติ, ปาติโมกฺขสํวรสํวุโต วิหรติ อาจารโคจรสมฺปนฺโน, อณุมตฺเตสุ วชฺเชสุ ภยทสฺสาวี สมาทาย สิกฺขติ สิกฺขาปเทสุ.\csromanspacer{} ยมฺปิ ภิกฺขเว ภิกฺขุ สีลวา โหติ ฯเปฯ สมาทาย สิกฺขติ สิกฺขาปเทสุ.\csromanspacer{} อยมฺปิ ธมฺโม นาถกรโณ. (1)}

\prose{ปุน จปรํ ภิกฺขเว ภิกฺขุ พหุสฺสุโต โหติ สุตธโร สุตสนฺนิจโย, เย เต ธมฺมา อาทิกลฺยาณา มชฺเฌกลฺยาณา ปริโยสานกลฺยาณา สาตฺถํ สพฺยญฺชนํ\footnote{สาตฺถา สพฺยญฺชนา (สี)} เกวลปริปุณฺณํ ปริสุทฺธํ พฺรหฺมจริยํ อภิวทนฺติ, ตถารูปาสฺส ธมฺมา พหุสฺสุตา\footnote{พหู สุตา (?)} โหนฺติ ธาตา วจสา ปริจิตา มนสานุเปกฺขิตา ทิฏฺฐิยา สุปฺปฏิวิทฺธา.\csromanspacer{} ยมฺปิ ภิกฺขเว ภิกฺขุ พหุสฺสุโต โหติ ฯเปฯ ทิฏฺฐิยา สุปฺปฏิวิทฺธา, อยมฺปิ ธมฺโม นาถกรโณ. (2)}

\prose{ปุน จปรํ ภิกฺขเว ภิกฺขุ กลฺยาณมิตฺโต โหติ กลฺยาณสหาโย ลลฺยาณสมฺปวงฺโก.\csromanspacer{} ยมฺปิ ภิกฺขเว ภิกฺขุ กลฺยาณมิตฺโต โหติ กลฺยาณสหาโย กลฺยาณสมฺปวงฺโก, อยมฺปิ ธมฺโม นาถกรโณ. (3)}

\prose{ปุน จปรํ ภิกฺขเว ภิกฺขุ สุวโจ โหติ โสวจสฺสกรเณหิ ธมฺเมหิ สมนฺนาคโต, ขโม ปทกฺขิณคฺคาหี อนุสาสนิํ.\csromanspacer{} ยมฺปิ ภิกฺขเว ภิกฺขุ สุวโจ โหติ ฯเปฯ อนุสาสนิํ, อยมฺปิ ธมฺโม นาถกรโณ. (4)}

\prose{ปุน จปรํ ภิกฺขเว ภิกฺขุ ยานิ ตานิ สพฺรหฺมจารีนํ อุจฺจาวจานิ กิํกรณียานิ, ตตฺถ ทกฺโข โหติ อนลโส, ตตฺรูปายาย วีมํสาย สมนฺนาคโต, อลํ กาตุํ อลํ สํวิธาตุํ.\csromanspacer{} ยมฺปิ ภิกฺขเว ภิกฺขุ ยานิ ตานิ สพฺรหฺมจารีนํ ฯเปฯ อลํ กาตุํ อลํ สํวิธาตุํ, อยมฺปิ ธมฺโม นาถกรโณ. (5)}

\prose{ปุน จปรํ ภิกฺขเว ภิกฺขุ ธมฺมกาโม โหติ ปิยสมุทาหาโร, อภิธมฺเม อภิวินเย อุฬารปาโมชฺโช.\csromanspacer{} ยมฺปิ ภิกฺขเว ภิกฺขุ ธมฺมกาโม}

\csromanmatikaanchor{mtk.40}
\csromantocmarkref{subsection}{6. อาหุเนยฺยสุตฺต}{mtk.40}
\bookmark[dest={mtk.40},level=2]{6. อาหุเนยฺยสุตฺต}
\vspace{\tipitakaparskip}
\csromancenter{ทสกนิปาตปาฬิ โหติ ปิยสมุทาหาโร, อภิธมฺเม อภิวินเย อุฬารปาโมชฺโช, อยมฺปิ ธมฺโม นาถกรโณ. (6)}
\prose{\csromanfolio{276}ปุน จปรํ ภิกฺขเว ภิกฺขุ อารทฺธวีริโย วิหรติ อกุสลานํ ธมฺมานํ ปหานาย กุสลานํ ธมฺมานํ อุปสมฺปทาย, ถามวา ทฬฺหปรกฺกโม อนิกฺขิตฺตธุโร กุสเลสุ ธมฺเมสุ.\csromanspacer{} ยมฺปิ ภิกฺขเว ภิกฺขุ อารทฺธวีริโย วิหรติ อกุสลานํ ธมฺมานํ ปหานาย กุสลานํ ธมฺมานํ อุปสมฺปทาย, ถามวา ทฬฺหปรกฺกโม อนิกฺขิตฺตธุโร กุสเลสุ ธมฺเมสุ, อยมฺปิ ธมฺโม นาถกรโณ. (7)}

\csromanmatikaanchor{mtk.41}
\csromantocmarkref{subsection}{7. ปฐมนาถสุตฺต}{mtk.41}
\bookmark[dest={mtk.41},level=2]{7. ปฐมนาถสุตฺต}
\prose{ปุน จปรํ ภิกฺขเว ภิกฺขุ สนฺตุฏฺโฐ โหติ อิตรีตร จีวร ปิณฺฑปาตเสนาสน คิลานปจฺจย เภสชฺช ปริกฺขาเรน.\csromanspacer{} ยมฺปิ ภิกฺขเว ภิกฺขุ สนฺตุฏฺโฐ โหติ อิตรีตรจีวรปิณฺฑปาตเสนาสนคิลานปจฺจย- เภสชฺชปริกฺขาเรน, อยมฺปิ ธมฺโม นาถกรโณ. (8)}

\prose{ปุน จปรํ ภิกฺขเว ภิกฺขุ สติมา โหติ ปรเมน สติเนปกฺเกน สมนฺนาคโต, จิรกตมฺปิ จิรภาสิตมฺปิ สริตา อนุสฺสริตา.\csromanspacer{} ยมฺปิ ภิกฺขเว ภิกฺขุ สติมา โหติ ปรเมน สติเนปกฺเกน สมนฺนาคโต, จิรกตมฺปิ จิรภาสิตมฺปิ สริตา อนุสฺสริตา, อยมฺปิ ธมฺโม นาถกรโณ. (9)}

\prose{ปุน จปรํ ภิกฺขเว ภิกฺขุ ปญฺญวา โหติ อุทยตฺถคามินิยา ปญฺญาย สมนฺนาคโต อริยาย นิพฺเพธิกาย สมฺมา ทุกฺขกฺขยคามินิยา.\csromanspacer{} ยมฺปิ ภิกฺขเว ภิกฺขุ ปญฺญวา โหติ อุทยตฺถคามินิยา ปญฺญาย สมนฺนาคโต อริยาย นิพฺเพธิกาย สมฺมา ทุกฺขกฺขยคามินิยา, อยมฺปิ ธมฺโม นาถกรโณ. (10)}

\prose{สนาถา ภิกฺขเว วิหรถ, มา อนาถา, ทุกฺขํ ภิกฺขเว อนาโถ วิหรติ.\csromanspacer{} อิเม โข ภิกฺขเว ทส นาถกรณา ธมฺมาติ.\csromanspacer{}\csromanspacer{}สตฺตมํ.}

\csromanheader{2}{8. ทุติยนาถสุตฺต}
\proseitem{18}{เอวํ เม สุตํ\csromandash{}เอกํ สมยํ ภควา สาวตฺถิยํ วิหรติ เชตวเน อนาถปิณฺฑิกสฺส อาราเม.\csromanspacer{} ตตฺร โข ภควา ภิกฺขู อามนฺเตสิ “ภิกฺขโว”ติ. “ภทนฺเต”ติ เต ภิกฺขู ภควโต ปจฺจสฺโสสุํ.\csromanspacer{} ภควา เอตทโวจ\csromandash{}}

\chapterheadpageread{2. นาถวคฺค}
\prose{\csromanfolio{277}สนาถา ภิกฺขเว วิหรถ, มา อนาถา, ทุกฺขํ ภิกฺขเว อนาโถ วิหรติ.\csromanspacer{} ทสยิเม ภิกฺขเว นาถกรณา ธมฺมา.\csromanspacer{} กตเม ทส? อิธ ภิกฺขเว ภิกฺขุ สีลวา โหติ ฯเปฯ สมาทาย สิกฺขติ สิกฺขาปเทสุ. “สีลวา วตายํ ภิกฺขุ ปาติโมกฺขสํวรสํวุโต วิหรติ อาจารโคจรสมฺปนฺโน, อณุมตฺเตสุ วชฺเชสุ ภยทสฺสาวี สมาทาย สิกฺขติ สิกฺขาปเทสู”ติ เถราปิ นํ ภิกฺขู วตฺตพฺพํ อนุสาสิตพฺพํ มญฺญนฺติ, มชฺฌิมาปิ ภิกฺขู.\csromanspacer{} นวาปิ ภิกฺขู วตฺตพฺพํ อนุสาสิตพฺพํ มญฺญนฺติ.\csromanspacer{} ตสฺส เถรานุกมฺปิตสฺส มชฺฌิมานุกมฺปิตสฺส นวานุกมฺปิตสฺส วุทฺธิเยว ปาฏิกงฺขากุสเลสุ ธมฺเมสุ, โน ปริหานิ.\csromanspacer{} อยมฺปิ ธมฺโม นาถกรโณ. (1)}

\prose{ปุน จปรํ ภิกฺขเว ภิกฺขุ พหุสฺสุโต โหติ ฯเปฯ ทิฏฺฐิยา สุปฺปฏิวิทฺธา. “พหุสฺสุโต วตายํ ภิกฺขุ สุตธโร สุตสนฺนิจโย, เย เต ธมฺมา อาทิกลฺยาณา มชฺเฌกลฺยาณา ปริโยสานกลฺยาณา สาตฺถํ สพฺยญฺชนํ เกวลปริปุณฺณํ ปริสุทฺธํ พฺรหฺมจริยํ อภิวทนฺติ, ตถารูปาสฺส ธมฺมา พหุสฺสุตา โหนฺติ ธาตา วจสา ปริจิตามนสานุเปกฺขิตา ทิฏฺฐิยา สุปฺปฏิวิทฺธา”ติ เถราปิ นํ ภิกฺขู วตฺตพฺพํ อนุสาสิตพฺพํ มญฺญนฺติ, มชฺฌิมาปิ ภิกฺขู.\csromanspacer{} นวาปิ ภิกฺขู วตฺตพฺพํ อนุสาสิตพฺพํ มญฺญนฺติ, ตสฺส เถรานุกมฺปิตสฺส มชฺฌิมานุกมฺปิตสฺส นวานุกมฺปิตสฺส วุทฺธิเยว ปาฏิกงฺขา กุสเลสุ ธมฺเมสุ, โน ปริหานิ.\csromanspacer{} อยมฺปิ ธมฺโม นาถกรโณ. (2)}

\prose{ปุน จปรํ ภิกฺขเว ภิกฺขุ กลฺยาณมิตฺโต โหติ กลฺยาณสหาโย กลฺยาณสมฺปวงฺโก. “กลฺยาณมิตฺโต วตายํ ภิกฺขุ กลฺยาณสหาโย กลฺยาณสมฺปวงฺโก”ติ เถราปิ นํ ภิกฺขู วตฺตพฺพํ อนุสาสิตพฺพํ มญฺญนฺติ, มชฺฌิมาปิ ภิกฺขู.\csromanspacer{} นวาปิ ภิกฺขู วตฺตพฺพํ อนุสาสิตพฺพํ มญฺญนฺติ.\csromanspacer{} ตสฺส เถรานุกมฺปิตสฺส มชฺฌิมานุกมฺปิตสฺส นวานุกมฺปิตสฺส วุทฺธิเยว ปาฏิกงฺขา กุสเลสุ ธมฺเมสุ, โน ปริหานิ.\csromanspacer{} อยมฺปิ ธมฺโม นาถกรโณ. (3)}

\prose{ปุน จปรํ ภิกฺขเว ภิกฺขุ สุวโจ โหติ โสวจสฺสกรเณหิ ธมฺเมหิ สมนฺนาคโต ขโม ปทกฺขิณคฺคาหี อนุสาสนิํ. “สุวโจ วตายํ ภิกฺขุ โสวจสฺสกรเณหิ ธมฺเมหิ สมนฺนาคโต ขโม ปทกฺขิณคฺคาหี อนุสาสนินฺ”ติ เถราปิ นํ ภิกฺขู วตฺตพฺพํ อนุสาสิตพฺพํ มญฺญนฺติ, มชฺฌิมาปิ ภิกฺขู.\csromanspacer{} นวาปิ ภิกฺขู วตฺตพฺพํ อนุสาสิตพฺพํ มญฺญนฺติ.\csromanspacer{} ตสฺส เถรานุกมฺปิตสฺส}

\vspace{\tipitakaparskip}
\csromancenter{ทสกนิปาตปาฬิ มชฺฌิมานุกมฺปิตสฺส นวานุกมฺปิตสฺส วุทฺธิเยว ปาฏิกงฺขา กุสเลสุ ธมฺเมสุ, โน ปริหานิ.\csromanspacer{} อยมฺปิ ธมฺโม นาถกรโณ. (4)}
\prose{\csromanfolio{278}ปุน จปรํ ภิกฺขเว ภิกฺขุ ยานิ ตานิ สพฺรหฺมจารีนํ อุจฺจาวจานิ กิํกรณียานิ, ตตฺถ ทกฺโข โหติ อนลโส ตตฺรูปายาย วีมํสาย สมนฺนาคโต, อลํ กาตุํ อลํ สํวิธาตุํ. “ยานิ ตานิ สพฺรหฺมจารีนํ อุจฺจาวจานิ กิํกรณียานิ, ตตฺถ ทกฺโข วตายํ ภิกฺขุ อนลโส ตตฺรูปายาย วีมํสาย สมนฺนาคโต, อลํ กาตุํ อลํ สํวิธาตุนฺ”ติ เถราปิ นํ ภิกฺขู วตฺตพฺพํ อนุสาสิตพฺพํ มญฺญนฺติ, มชฺฌิมาปิ ภิกฺขู.\csromanspacer{} นวาปิ ภิกฺขู วตฺตพฺพํ อนุสาสิตพฺพํ มญฺญนฺติ.\csromanspacer{} ตสฺส เถรานุกมฺปิตสฺส มชฺฌิมานุกมฺปิตสฺส นวานุกมฺปิตสฺส วุทฺธิเยว ปาฏิกงฺขา กุสเลสุ ธมฺเมสุ, โน ปริหานิ.\csromanspacer{} อยมฺปิ ธมฺโม นาถกรโณ. (5)}

\csromanmatikaanchor{mtk.42}
\csromantocmarkref{subsection}{8. ทุติยนาถสุตฺต}{mtk.42}
\bookmark[dest={mtk.42},level=2]{8. ทุติยนาถสุตฺต}
\prose{ปุน จปรํ ภิกฺขเว ภิกฺขุ ธมฺมกาโม โหติ ปิยสมุทาหาโร, อภิธมฺเม อภิวินเย อุฬารปาโมชฺโช. “ธมฺมกาโม วตายํ ภิกฺขุ ปิยสมุทาหาโร, อภิธมฺเม อภิวินเย อุฬารปาโมชฺโช”ติ เถราปิ นํ ภิกฺขู วตฺตพฺพํ อนุสาสิตพฺพํ มญฺญนฺติ, มชฺฌิมาปิ ภิกฺขู.\csromanspacer{} นวาปิ ภิกฺขู วตฺตพฺพํ อนุสาสิตพฺพํ มญฺญนฺติ.\csromanspacer{} ตสฺส เถรานุกมฺปิตสฺส มชฺฌิมานุกมฺปิตสฺส นวานุกมฺปิตสฺส วุทฺธิเยว ปาฏิกงฺขา กุสเลสุ ธมฺเมสุ, โน ปริหานิ.\csromanspacer{} อยมฺปิ ธมฺโม นาถกรโณ. (6)}

\prose{ปุน จปรํ ภิกฺขเว ภิกฺขุ อารทฺธวีริโย วิหรติ อกุสลานํ ธมฺมานํ ปหานาย, กุสลานํ ธมฺมานํ อุปสมฺปทาย, ถามวา ทฬฺหปรกฺกโม อนิกฺขิตฺตธุโร กุสเลสุ ธมฺเมสุ. “อารทฺธวีริโย วตายํ ภิกฺขุ วิหรติ อกุสลานํ ธมฺมานํ ปหานาย, กุสลานํ ธมฺมานํ อุปสมฺปทาย, ถามวา ทฬฺหปรกฺกโม อนิกฺขิตฺตธุโร กุสเลสุ ธมฺเมสู”ติ เถราปิ นํ ภิกฺขู วตฺตพฺพํ อนุสาสิตพฺพํ มญฺญนฺติ, มชฺฌิมาปิ ภิกฺขู.\csromanspacer{} นวาปิ ภิกฺขู วตฺตพฺพํ อนุสาสิตพฺพํ มญฺญนฺติ.\csromanspacer{} ตสฺส เถรานุกมฺปิตสฺส มชฺฌิมานุกมฺปิตสฺส นวานุกมฺปิตสฺส วุทฺธิเยว ปาฏิกงฺขา กุสเลสุ ธมฺเมสุ, โน ปริหานิ.\csromanspacer{} อยมฺปิ ธมฺโม นาถกรโณ. (7)}

\prose{ปุน จปรํ ภิกฺขเว ภิกฺขุ สนฺตุฏฺโฐ โหติ อิตรีตรจีวร ปิณฺฑปาตเสนาสนคิลานปจฺจยเภสชฺชปริกฺขาเรน. “สนฺตุฏฺโฐ วตายํ ภิกฺขุ อิตรีตรจีวรปิณฺฑปาตเสนาสนคิลานปจฺจยเภสชฺช ปริกฺขาเรนา”ติ}

\vspace{\tipitakaparskip}
\csromancenter{นาถวคฺค เถราปิ นํ ภิกฺขู วตฺตพฺพํ อนุสาสิตพฺพํ มญฺญนฺติ, มชฺฌิมาปิ ภิกฺขู.\csromanspacer{} นวาปิ ภิกฺขู วตฺตพฺพํ อนุสาสิตพฺพํ มญฺญนฺติ.\csromanspacer{} ตสฺส เถรานุกมฺปิตสฺส มชฺฌิมานุกมฺปิตสฺส นวานุกมฺปิตสฺส วุทฺธิเยว ปาฏิกงฺขา กุสเลสุ ธมฺเมสุ, โน ปริหานิ.\csromanspacer{} อยมฺปิ ธมฺโม นาถกรโณ. (8)}
\prose{\csromanfolio{279}ปุน จปรํ ภิกฺขเว ภิกฺขุ สติมา โหติ ปรเมน สติเนปกฺเกน สมนฺนาคโต, จิรกตมฺปิ จิรภาสิตมฺปิ สริตา อนุสฺสริตา. “สติมา วตายํ ภิกฺขุ ปรเมน สติเนปกฺเกน สมนฺนาคโต, จิรกตมฺปิ จิรภาสิตมฺปิ สริตา อนุสฺสริตา”ติ เถราปิ นํ ภิกฺขู วตฺตพฺพํ อนุสาสิตพฺพํ มญฺญนฺติ, มชฺฌิมาปิ ภิกฺขู.\csromanspacer{} นวาปิ ภิกฺขู วตฺตพฺพํ อนุสาสิตพฺพํ มญฺญนฺติ.\csromanspacer{} ตสฺส เถรานุกมฺปิตสฺส มชฺฌิมานุกมฺปิตสฺส นวานุกมฺปิตสฺส วุทฺธิเยว ปาฏิกงฺขา กุสเลสุ ธมฺเมสุ, โน ปริหานิ.\csromanspacer{} อยมฺปิ ธมฺโม นาถกรโณ. (9)}

\prose{ปุน จปรํ ภิกฺขเว ภิกฺขุ ปญฺญวา โหติ อุทยตฺถคามินิยา ปญฺญาย สมนฺนาคโต อริยาย นิพฺเพธิกาย สมฺมา ทุกฺขกฺขยคามินิยา. “ปญฺญวา วตายํ ภิกฺขุ อุทยตฺถคามินิยา ปญฺญาย สมนฺนาคโต อริยาย นิพฺเพธิกาย สมฺมา ทุกฺขกฺขยคามินิยา”ติ เถราปิ นํ ภิกฺขู วตฺตพฺพํ อนุสาสิตพฺพํ มญฺญนฺติ, มชฺฌิมาปิ ภิกฺขู.\csromanspacer{} นวาปิ ภิกฺขู วตฺตพฺพํ อนุสาสิตพฺพํ มญฺญนฺติ.\csromanspacer{} ตสฺส เถรานุกมฺปิตสฺส ฯเปฯ โน ปริหานิ.\csromanspacer{} อยมฺปิ ธมฺโม นาถกรโณ. (10)}

\prose{สนาถา ภิกฺขเว วิหรถ, มา อนาถา, ทุกฺขํ ภิกฺขเว อนาโถ วิหรติ.\csromanspacer{} อิเม โข ภิกฺขเว ทส นาถกรณา ธมฺมาติ.\csromanspacer{} อิทมโวจ ภควา.\csromanspacer{} อตฺตมนา เต ภิกฺขู ภควโต ภาสิตํ อภินนฺทุนฺติ.\csromanspacer{}\csromanspacer{}อฏฺฐมํ.}

\csromanheader{2}{9. ปฐม-อริยาวาสสุตฺต}
\proseitem{19}{\csromansymbolfootnote{*}{ที 3. 224, 257 ปิฏฺเฐสุปิ.}ทสยิเม ภิกฺขเว อริยาวาสา, เย อริยา อาวสิํสุ วา อาวสนฺติ วา อาวสิสฺสนฺติ วา.\csromanspacer{} กตเม ทส? อิธ ภิกฺขเว ภิกฺขุ ปญฺจงฺควิปฺปหีโน โหติ ฉฬงฺคสมนฺนาคโต เอการกฺโข จตุราปสฺเสโน ปณุนฺนปจฺเจกสจฺโจ\footnote{ปนุณฺณปจฺเจกสจฺโจ (ก)} สมวยสฏฺเฐสโน อนาวิลสงฺกปฺโป ปสฺสทฺธกายสงฺขาโร สุวิมุตฺตจิตฺโต สุวิมุตฺตปญฺโญ.\csromanspacer{} อิเม โข ภิกฺขเว ทส}

\vspace{\tipitakaparskip}
\csromancenter{ทสกนิปาตปาฬิ อริยาวาสา, เย อริยา อาวสิํสุ วา อาวสนฺติ วา อาวสิสฺสนฺติ วาติ.\csromanspacer{}\csromanspacer{}นวมํ.}
\csromanheader{2}{10. ทุติย-อริยาวาสสุตฺต}
\proseitem{20}{\csromanfolio{280}เอกํ สมยํ ภควา กุรูสุ วิหรติ กมฺมาสธมฺมํ นาม กุรูนํ นิคโม.\csromanspacer{} ตตฺร โข ภควา ภิกฺขู อามนฺเตสิ ฯเปฯ.}

\prose{ทสยิเม ภิกฺขเว อริยาวาสา, เย อริยา อาวสิํสุ วา อาวสนฺติ วา อาวสิสฺสนฺติ วา.\csromanspacer{} กตเม ทส? อิธ ภิกฺขเว ภิกฺขุ ปญฺจงฺควิปฺปหีโน โหติ ฉฬงฺคสมนฺนาคโต เอการกฺโข จตุราปสฺเสโน ปณุนฺนปจฺเจกสจฺโจ สมวยสฏฺเฐสโน อนาวิลสงฺกปฺโป ปสฺสทฺธกายสงฺขาโร สุวิมุตฺตจิตฺโต สุวิมุตฺตปญฺโญ.}

\prose{กถญฺจ ภิกฺขเว ภิกฺขุ ปญฺจงฺควิปฺปหีโน โหติ? อิธ ภิกฺขเว ภิกฺขุโน กามจฺฉนฺโท ปหีโน โหติ, พฺยาปาโท ปหีโน โหติ, ถินมิทฺธํ ปหีนํ โหติ, อุทฺธจฺจกุกฺกุจฺจํ ปหีนํ โหติ, วิจิกิจฺฉา ปหีนา โหติ.\csromanspacer{} เอวํ โข ภิกฺขเว ภิกฺขุ ปญฺจงฺควิปฺปหีโน โหติ. (1)}

\prose{กถญฺจ ภิกฺขเว ภิกฺขุ ฉฬงฺคสมนฺนาคโต โหติ? อิธ ภิกฺขเว ภิกฺขุ จกฺขุนา รูปํ ทิสฺวา เนว สุมโน โหติ, น ทุมฺมโน, อุเปกฺขโก วิหรติ สโต สมฺปชาโน.\csromanspacer{} โสเตน สทฺทํ สุตฺวา ฯเปฯ.\csromanspacer{} ฆาเนน คนฺธํ ฆายิตฺวา.\csromanspacer{} ชิวฺหาย รสํ สายิตฺวา.\csromanspacer{} กาเยน โผฏฺฐพฺพํ ผุสิตฺวา.\csromanspacer{} มนสา ธมฺมํ วิญฺญาย เนว สุมโน โหติ, น ทุมฺมโน, อุเปกฺขโก วิหรติ สโต สมฺปชาโน.\csromanspacer{} เอวํ โข ภิกฺขเว ภิกฺขุ ฉฬงฺคสมนฺนาคโต โหติ. (2)}

\prose{กถญฺจ ภิกฺขเว ภิกฺขุ เอการกฺโข โหติ? อิธ ภิกฺขเว ภิกฺขุ สตารกฺเขน เจตสา สมนฺนาคโต โหติ.\csromanspacer{} เอวํ โข ภิกฺขเว ภิกฺขุ เอการกฺโข โหติ. (3)}

\csromanmatikaanchor{mtk.43}
\csromantocmarkref{subsection}{9. ปฐม-อริยาวาสสุตฺต}{mtk.43}
\bookmark[dest={mtk.43},level=2]{9. ปฐม-อริยาวาสสุตฺต}
\prose{กถญฺจ ภิกฺขเว ภิกฺขุ จตุราปสฺเสโน โหติ? อิธ ภิกฺขเว ภิกฺขุ สงฺขาเยกํ ปฏิเสวติ, สงฺขาเยกํ อธิวาเสติ, สงฺขาเยกํ ปริวชฺเชติ, สงฺขาเยกํ วิโนเทติ.\csromanspacer{} เอวํ โข ภิกฺขเว ภิกฺขุ จตุราปสฺเสโน โหติ. (4)}

\chapterheadpageread{2. นาถวคฺค}
\prose{\csromanfolio{281}กถญฺจ ภิกฺขเว ภิกฺขุ ปณุนฺนปจฺเจกสจฺโจ โหติ? อิธ ภิกฺขเว ภิกฺขุโน ยานิ ตานิ ปุถุสมณพฺราหฺมณานํ ปุถุปจฺเจกสจฺจานิ.\csromanspacer{} เสยฺยถิทํ, “สสฺสโต โลโก”ติ วา “อสสฺสโต โลโก”ติ วา “อนฺตวา โลโก”ติ วา “อนนฺตวา โลโก”ติ วา “ตํ ชีวํ ตํ สรีรนฺ”ติ วา “อญฺญํ ชีวํ อญฺญํ สรีรนฺ”ติ วา “โหติ ตถาคโต ปรํ มรณา”ติ วา “น โหติ ตถาคโต ปรํ มรณา”ติ วา “โหติ จ น จ โหติ ตถาคโต ปรํ มรณา”ติ วา “เนว โหติ น น โหติ ตถาคโต ปรํ มรณา”ติ วา, สพฺพานิ ตานิ นุนฺนานิ โหนฺติ ปณุนฺนานิ\footnote{นุณฺณานิ โหนฺติ ปนุณฺณานิ (?)} จตฺตานิ วนฺตานิ มุตฺตานิ ปหีนานิ ปฏินิสฺสฏฺฐานิ.\csromanspacer{} เอวํ โข ภิกฺขเว ภิกฺขุ ปณุนฺนปจฺเจกสจฺโจ โหติ. (5)}

\csromanmatikaanchor{mtk.44}
\csromantocmarkref{subsection}{10. ทุติย-อริยาวาสสุตฺต}{mtk.44}
\bookmark[dest={mtk.44},level=2]{10. ทุติย-อริยาวาสสุตฺต}
\prose{กถญฺจ ภิกฺขเว ภิกฺขุ สมวยสฏฺเฐสโน โหติ? อิธ ภิกฺขเว ภิกฺขุโน กาเมสนา ปหีนา โหติ, ภเวสนา ปหีนา โหติ, พฺรหฺมจริเยสนา ปฏิปฺปสฺสทฺธา.\csromanspacer{} เอวํ โข ภิกฺขเว ภิกฺขุ สมวยสฏฺเฐสโน โหติ. (6)}

\prose{กถญฺจ ภิกฺขเว ภิกฺขุ อนาวิลสงฺกปฺโป โหติ? อิธ ภิกฺขเว ภิกฺขุโน กามสงฺกปฺโป ปหีโน โหติ, พฺยาปาทสงฺกปฺโป ปหีโน โหติ, วิหิํสาสงฺกปฺโป ปหีโน โหติ.\csromanspacer{} เอวํ โข ภิกฺขเว ภิกฺขุ อนาวิลสงฺกปฺโป โหติ. (7)}

\prose{กถญฺจ ภิกฺขเว ภิกฺขุ ปสฺสทฺธกายสงฺขาโร โหติ? อิธ ภิกฺขเว ภิกฺขุ สุขสฺส จ ปหานา ทุกฺขสฺส จ ปหานา ปุพฺเพว โสมนสฺสโทมนสฺสานํ อตฺถงฺคมา อทุกฺขมสุขํ อุเปกฺขาสติปาริสุทฺธิํ จตุตฺถํ ฌานํ อุปสมฺปชฺช วิหรติ.\csromanspacer{} เอวํ โข ภิกฺขเว ภิกฺขุ ปสฺสทฺธกายสงฺขาโร โหติ. (8)}

\prose{กถญฺจ ภิกฺขเว ภิกฺขุ สุวิมุตฺตจิตฺโต โหติ? อิธ ภิกฺขเว ภิกฺขุโน ราคา จิตฺตํ วิมุตฺตํ โหติ, โทสา จิตฺตํ วิมุตฺตํ โหติ, โมหา จิตฺตํ วิมุตฺตํ โหติ.\csromanspacer{} เอวํ โข ภิกฺขเว ภิกฺขุ สุวิมุตฺตจิตฺโต โหติ. (9)}

\prose{กถญฺจ ภิกฺขเว ภิกฺขุ สุวิมุตฺตปญฺโญ โหติ? อิธ ภิกฺขเว ภิกฺขุ “ราโค เม ปหีโน อุจฺฉินฺนมูโล ตาลาวตฺถุกโต อนภาวํกโต อายติํ อนุปฺปาทธมฺโม”ติ ปชานาติ. “โทโส เม ปหีโน ฯเปฯ.}

\vspace{\tipitakaparskip}
\csromancenter{ทสกนิปาตปาฬิ โมโห เม ปหีโน อุจฺฉินฺนมูโล ตาลาวตฺถุกโต อนภาวํ กโต อายติํ อนุปฺปาทธมฺโม”ติ ปชานาติ.\csromanspacer{} เอวํ โข ภิกฺขเว ภิกฺขุ สุวิมุตฺตปญฺโญ โหติ. (10)}
\csromanmatikaanchor{mtk.45}
\csromantocmarkref{section}{อุทฺทานคาถา}{mtk.45}
\bookmark[dest={mtk.45},level=1]{อุทฺทานคาถา}
\csromantocmarknopage{section}{3. มหาวคฺค}{mtk.46}
\bookmark[dest={mtk.46},level=1]{3. มหาวคฺค}
\prose{\csromanfolio{282}เย หิ เกจิ ภิกฺขเว อตีตมทฺธานํ อริยา อริยาวาเส อาวสิํสุ, สพฺเพ เต อิเมว ทส อริยาวาเส อาวสิํสุ.\csromanspacer{} เย หิ\footnote{เยปิ (?)} เกจิ ภิกฺขเว อนาคตมทฺธานํ อริยา อริยาวาเส อาวสิสฺสนฺติ, สพฺเพ เต อิเมว ทส อริยาวาเส อาวสิสฺสนฺติ.\csromanspacer{} เย หิ1 เกจิ ภิกฺขเว เอตรหิ อริยา อริยาวาเส อาวสนฺติ, สพฺเพ เต อิเมว ทส อริยาวาเส อาวสนฺติ.\csromanspacer{} อิเม โข ภิกฺขเว ทส อริยาวาสา, เย อริยา อาวสิํสุ วา อาวสนฺติ วา อาวสิสฺสนฺติ วาติ.\csromanspacer{}\csromanspacer{}ทสมํ.\csromanspacer{} นาถวคฺโค ทุติโย.}

\titlehead{\textbf{ตสฺสุทฺทานํ}}
\csromangathameasure{เสนาสนํ จ ปญฺจงฺคํ, สํโยชนาขิเลน จ. \\ อปฺปมาโท อาหุเนยฺโย, เทฺว นาถา เทฺว อริยาวาสาติ.}
\csromangathagroup{\csromangathastanza{เสนาสนํ จ ปญฺจงฺคํ, สํโยชนาขิเลน จ. \\ อปฺปมาโท อาหุเนยฺโย, เทฺว นาถา เทฺว อริยาวาสาติ.}}

\chapterhead{3. มหาวคฺค}
\csromanheader{2}{1. สีหนาทสุตฺต}
\proseitem{21}{สีโห ภิกฺขเว มิคราชา สายนฺหสมยํ อาสยา นิกฺขมติ, อาสยา นิกฺขมิตฺวา วิชมฺภติ, วิชมฺภิตฺวา สมนฺตา จตุทฺทิสํ\footnote{จตุทฺทิสา (สฺยา, ก) อํ 2. 365 ปิฏฺเฐปิ.} อนุวิโลเกติ, สมนฺตา จตุทฺทิสํ2 อนุวิโลเกตฺวา ติกฺขตฺตุํ สีหนาทํ นทติ, ติกฺขตฺตุํ สีหนาทํ นทิตฺวา โคจราย ปกฺกมติ.\csromanspacer{} ตํ กิสฺส เหตุ? “มาหํ ขุทฺทเก ปาเณ วิสมคเต สงฺฆาตํ อาปาเทสินฺ”ติ.}

\prose{“สีโห”ติ โข ภิกฺขเว ตถาคตสฺเสตํ อธิวจนํ อรหโต สมฺมาสมฺพุทฺธสฺส.\csromanspacer{} ยํ โข ภิกฺขเว ตถาคโต ปริสาย ธมฺมํ เทเสติ, อิทมสฺส โหติ สีหนาทสฺมิํ.}

\chapterheadpageread{3. มหาวคฺค}
\prose{\csromanfolio{283}\csromansymbolfootnote{*}{ม 1. 99; อภิ 2. 328; ขุ 9. 356 ปิฏฺเฐสุปิ.}ทสยิมานิ ภิกฺขเว ตถาคตสฺส ตถาคตพลานิ, เยหิ พเลหิ สมนฺนาคโต ตถาคโต อาสภํ ฐานํ ปฏิชานาติ, ปริสาสุ สีหนาทํ นทติ, พฺรหฺมจกฺกํ ปวตฺเตติ.\csromanspacer{} กตมานิ ทส? อิธ ภิกฺขเว ตถาคโต ฐานญฺจ ฐานโต อฏฺฐานญฺจ อฏฺฐานโต ยถาภูตํ ปชานาติ.\csromanspacer{} ยมฺปิ ภิกฺขเว ตถาคโต ฐานญฺจ ฐานโต อฏฺฐานญฺจ อฏฺฐานโต ยถาภูตํ ปชานาติ.\csromanspacer{} อิทมฺปิ ภิกฺขเว ตถาคตสฺส ตถาคตพลํ โหติ, ยํ พลํ อาคมฺม ตถาคโต อาสภํ ฐานํ ปฏิชานาติ, ปริสาสุ สีหนาทํ นทติ, พฺรหฺมจกฺกํ ปวตฺเตติ. (1)}

\prose{ปุน จปรํ ภิกฺขเว ตถาคโต อตีตานาคตปจฺจุปฺปนฺนานํ กมฺมสมาทานานํ ฐานโส เหตุโส วิปากํ ยถาภูตํ ปชานาติ.\csromanspacer{} ยมฺปิ ภิกฺขเว ตถาคโต อตีตานาคตปจฺจุปฺปนฺนานํ กมฺมสมาทานานํ ฐานโส เหตุโส วิปากํ ยถาภูตํ ปชานาติ.\csromanspacer{} อิทมฺปิ ภิกฺขเว ตถาคตสฺส ตถาคตพลํ โหติ, ยํ พลํ อาคมฺม ตถาคโต อาสภํ ฐานํ ปฏิชานาติ, ปริสาสุ สีหนาทํ นทติ, พฺรหฺมจกฺกํ ปวตฺเตติ. (2)}

\prose{ปุน จปรํ ภิกฺขเว ตถาคโต สพฺพตฺถคามินิํ ปฏิปทํ ยถาภูตํ ปชานาติ.\csromanspacer{} ยมฺปิ ภิกฺขเว ตถาคโต สพฺพตฺถคามินิํ ปฏิปทํ ยถาภูตํ ปชานาติ.\csromanspacer{} อิทมฺปิ ภิกฺขเว ตถาคตสฺส ตถาคตพลํ โหติ, ยํ พลํ อาคมฺม ตถาคโต อาสภํ ฐานํ ปฏิชานาติ, ปริสาสุ สีหนาทํ นทติ, พฺรหฺมจกฺกํ ปวตฺเตติ. (3)}

\prose{ปุน จปรํ ภิกฺขเว ตถาคโต อเนกธาตุํ นานาธาตุํ โลกํ\footnote{อเนกธาตุนานาธาตุโลกํ (สี, ก)} ยถาภูตํ ปชานาติ.\csromanspacer{} ยมฺปิ ภิกฺขเว ตถาคโต อเนกธาตุํ นานาธาตุํ โลกํ ยถาภูตํ ปชานาติ.\csromanspacer{} อิทมฺปิ ภิกฺขเว ตถาคตสฺส ตถาคตพลํ โหติ ฯเปฯ พฺรหฺมจกฺกํ ปวตฺเตติ. (4)}

\prose{ปุน จปรํ ภิกฺขเว ตถาคโต สตฺตานํ นานาธิมุตฺติกตํ ยถาภูตํ ปชานาติ.\csromanspacer{} ยมฺปิ ภิกฺขเว ตถาคโต สตฺตานํ นานาธิมุตฺติกตํ ยถาภูตํ ปชานาติ.\csromanspacer{} อิทมฺปิ ภิกฺขเว ตถาคตสฺส ตถาคตพลํ โหติ ฯเปฯ พฺรหฺมจกฺกํ ปวตฺเตติ. (5)}

\csromanmatikaanchor{mtk.47}
\csromantocmarkref{subsection}{1. สีหนาทสุตฺต}{mtk.47}
\bookmark[dest={mtk.47},level=2]{1. สีหนาทสุตฺต}
\prose{ปุน จปรํ ภิกฺขเว ตถาคโต ปรสตฺตานํ ปรปุคฺคลานํ อินฺทฺริยปโรปริยตฺตํ ยถาภูตํ ปชานาติ.\csromanspacer{} ยมฺปิ ภิกฺขเว ตถาคโต ปรสตฺตานํ}

\vspace{\tipitakaparskip}
\csromancenter{ทสกนิปาตปาฬิ ปรปุคฺคลานํ อินฺทฺริยปโรปริยตฺตํ ยถาภูตํ ปชานาติ.\csromanspacer{} อิทมฺปิ ภิกฺขเว ตถาคตสฺส ตถาคตพลํ โหติ ฯเปฯ พฺรหฺมจกฺกํ ปวตฺเตติ. (6)}
\prose{\csromanfolio{284}ปุน จปรํ ภิกฺขเว ตถาคโต ฌานวิโมกฺขสมาธิ- สมาปตฺตีนํ สํกิเลสํ โวทานํ วุฏฺฐานํ ยถาภูตํ ปชานาติ.\csromanspacer{} ยมฺปิ ฯเปฯ ปชานาติ.\csromanspacer{} อิทมฺปิ ภิกฺขเว ตถาคตสฺส ตถาคตพลํ โหติ ฯเปฯ พฺรหฺมจกฺกํ ปวตฺเตติ. (7)}

\prose{ปุน จปรํ ภิกฺขเว ตถาคโต อเนกวิหิตํ ปุพฺเพนิวาสํ อนุสฺสรติ.\csromanspacer{} เสยฺยถิทํ, เอกมฺปิ ชาติํ เทฺวปิ ชาติโย ติสฺโสปิ ชาติโย จตสฺโสปิ ชาติโย ปญฺจปิ ชาติโย ทสปิ ชาติโย วีสมฺปิ ชาติโย ติํสมฺปิ ชาติโย จตฺตาลีสมฺปิ ชาติโย ปญฺญาสมฺปิ ชาติโย ชาติสตมฺปิ ชาติสหสฺสมฺปิ ชาติสตสหสฺสมฺปิ อเนเกปิ สํวฏฺฏกปฺเป อเนเกปิ วิวฏฺฏกปฺเป อเนเกปิ สํวฏฺฏวิวฏฺฏกปฺเป อมุตฺราสิํ “เอวํนาโม เอวํโคตฺโต เอวํวณฺโณ เอวมาหาโร เอวํสุขทุกฺขปฺปฏิสํเวที เอวมายุปริยนฺโต, โส ตโต จุโต อมุตฺร อุทปาทิํ, ตตฺราปาสิํ “เอวํนาโม เอวํโคตฺโต เอวํวณฺโณ เอวมาหาโร เอวํสุขทุกฺขปฺปฏิสํเวที เอวมายุปริยนฺโต, โส ตโต จุโต อิธูปปนฺโน”ติ อิติ สาการํ ส-อุทฺเทสํ อเนกวิหิตํ ปุพฺเพนิวาสํ อนุสฺสรติ.\csromanspacer{} ยมฺปิ ภิกฺขเว ตถาคโต อเนกวิหิตํ ปุพฺเพนิวาสํ อนุสฺสรติ.\csromanspacer{} เสยฺยถิทํ, เอกมฺปิ ชาติํ เทฺวปิ ชาติโย ฯเปฯ อิติ สาการํ ส-อุทฺเทสํ อเนกวิหิตํ ปุพฺเพนิวาสํ อนุสฺสรติ.\csromanspacer{} อิทมฺปิ ภิกฺขเว ตถาคตสฺส ตถาคตพลํ โหติ, ยํ พลํ อาคมฺม ตถาคโต อาสภํ ฐานํ ปฏิชานาติ, ปริสาสุ สีหนาทํ นทติ, พฺรหฺมจกฺกํ ปวตฺเตติ. (8)}

\prose{ปุน จปรํ ภิกฺขเว ตถาคโต ทิพฺเพน จกฺขุนา วิสุทฺเธน อติกฺกนฺตมานุสเกน สตฺเต ปสฺสติ จวมาเน อุปปชฺชมาเน หีเน ปณีเต สุวณฺเณ ทุพฺพณฺเณ สุคเต ทุคฺคเต, ยถากมฺมูปเค สตฺเต ปชานาติ “อิเม วต โภนฺโต สตฺตา กายทุจฺจริเตน สมนฺนาคตา วจีทุจฺจริเตน สมนฺนาคตา มโนทุจฺจริเตน สมนฺนาคตา อริยานํ อุปวาทกา มิจฺฉาทิฏฺฐิกา มิจฺฉาทิฏฺฐิกมฺมสมาทานา, เต กายสฺส เภทา ปรํ มรณา อปายํ ทุคฺคติํ วินิปาตํ นิรยํ อุปปนฺนา.\csromanspacer{} อิเม วา ปน โภนฺโต สตฺตา}

\vspace{\tipitakaparskip}
\csromancenter{มหาวคฺค กายสุจริเตน สมนฺนาคตา วจีสุจริเตน สมนฺนาคตา มโนสุจริเตน สมนฺนาคตา อริยานํ อนุปวาทกา สมฺมาทิฏฺฐิกา สมฺมาทิฏฺฐิกมฺมสมาทานา, เต กายสฺส เภทา ปรํ มรณา สุคติํ สคฺคํ โลกํ อุปปนฺนา”ติ.\csromanspacer{} อิติ ทิพฺเพน จกฺขุนา วิสุทฺเธน อติกฺกนฺตมานุสเกน สตฺเต ปสฺสติ จวมาเน อุปปชฺชมาเน หีเน ปณีเต สุวณฺเณ ทุพฺพณฺเณ สุคเต ทุคฺคเต, ยถากมฺมูปเค สตฺเต ปชานาติ.\csromanspacer{} ยมฺปิ ภิกฺขเว ตถาคโต ทิพฺเพน จกฺขุนา วิสุทฺเธน อติกฺกนฺตมานุสเกน ฯเปฯ ยถากมฺมูปเค สตฺเต ปชานาติ.\csromanspacer{} อิทมฺปิ ภิกฺขเว ตถาคตสฺส ตถาคตพลํ โหติ, ยํ พลํ อาคมฺม ตถาคโต อาสภํ ฐานํ ปฏิชานาติ, ปริสาสุ สีหนาทํ นทติ, พฺรหฺมจกฺกํ ปวตฺเตติ. (9)}
\prose{\csromanfolio{285}ปุน จปรํ ภิกฺขเว ตถาคโต อาสวานํ ขยา อนาสวํ เจโตวิมุตฺติํ ปญฺญาวิมุตฺติํ ทิฏฺเฐว ธมฺเม สยํ อภิญฺญา สจฺฉิกตฺวา อุปสมฺปชฺช วิหรติ.\csromanspacer{} ยมฺปิ ภิกฺขเว ตถาคโต อาสวานํ ขยา อนาสวํ เจโตวิมุตฺติํ ปญฺญาวิมุตฺติํ ทิฏฺเฐว ธมฺเม สยํ อภิญฺญา สจฺฉิกตฺวา อุปสมฺปชฺช วิหรติ.\csromanspacer{} อิทมฺปิ ภิกฺขเว ตถาคตสฺส ตถาคตพลํ โหติ, ยํ พลํ อาคมฺม ตถาคโต อาสภํ ฐานํ ปฏิชานาติ, ปริสาสุ สีหนาทํ นทติ, พฺรหฺมจกฺกํ ปวตฺเตติ. (10)}

\prose{อิมานิ โข ภิกฺขเว ทส ตถาคตสฺส ตถาคตพลานิ, เยหิ พเลหิ สมนฺนาคโต ตถาคโต อาสภํ ฐานํ ปฏิชานาติ, ปริสาสุ สีหนาทํ นทติ, พฺรหฺมจกฺกํ ปวตฺเตตีติ.\csromanspacer{}\csromanspacer{}ปฐมํ.}

\csromanheader{2}{2. อธิวุตฺติปทสุตฺต}
\proseitem{22}{อถ โข อายสฺมา อานนฺโท เยน ภควา เตนุปสงฺกมิ, อุปสงฺกมิตฺวา ภควนฺตํ อภิวาเทตฺวา เอกมนฺตํ นิสีทิ, เอกมนฺตํ นิสินฺนํ โข อายสฺมนฺตํ อานนฺทํ ภควา เอตทโวจ\csromandash{}}

\prose{เย เต อานนฺท ธมฺมา เตสํ เตสํ อธิวุตฺติปทานํ\footnote{อธิมุตฺติปทานํ (ก)} อภิญฺญา สจฺฉิกิริยาย สํวตฺตนฺติ.\csromanspacer{} วิสารโท อหํ อานนฺท ตตฺถ ปฏิชานามิ เตสํ เตสํ ตถา ตถา ธมฺมํ เทเสตุํ, ยถา ยถา ปฏิปนฺโน สนฺตํ วา “อตฺถี”ติ ญสฺสติ, อสนฺตํ วา “นตฺถี”ติ ญสฺสติ, หีนํ วา “หีนนฺ”ติ}

\vspace{\tipitakaparskip}
\csromancenter{ทสกนิปาตปาฬิ ญสฺสติ, ปณีตํ วา “ปณีตนฺ”ติ ญสฺสติ, ส-อุตฺตรํ วา “ส-อุตฺตรนฺ”ติ ญสฺสติ, อนุตฺตรํ วา “อนุตฺตรนฺ”ติ ญสฺสติ.\csromanspacer{} ยถา ยถา วา ปน ตํ ญาเตยฺยํ วา ทฏฺเฐยฺยํ วา สจฺฉิกเรยฺยํ วา ตถา ตถา ญสฺสติ วา ทกฺขติ วา สจฺฉิกริสฺสติ วาติ ฐานเมตํ วิชฺชติ.\csromanspacer{} เอตทานุตฺตริยํ อานนฺท ญาณานํ, ยทิทํ ตตฺถ ตตฺถ ยถาภูตญาณํ.\csromanspacer{} เอตสฺมา จาหํ อานนฺท ญาณา อญฺญํ ญาณํ อุตฺตริตรํ วา ปณีตตรํ วา นตฺถีติ วทามิ.}
\prose{\csromanfolio{286}ทสยิมานิ อานนฺท ตถาคตสฺส ตถาคตพลานิ, เยหิ พเลหิ สมนฺนาคโต ตถาคโต อาสภํ ฐานํ ปฏิชานาติ, ปริสาสุ สีหนาทํ นทติ, พฺรหฺมจกฺกํ ปวตฺเตติ.\csromanspacer{} กตมานิ ทส? อิธานนฺท ตถาคโต ฐานญฺจ ฐานโต อฏฺฐานญฺจ อฏฺฐานโต ยถาภูตํ ปชานาติ.\csromanspacer{} ยมฺปานนฺท ตถาคโต ฐานญฺจ ฐานโต อฏฺฐานญฺจ อฏฺฐานโต ยถาภูตํ ปชานาติ.\csromanspacer{} อิทมฺปานนฺท ตถาคตสฺส ตถาคตพลํ โหติ, ยํ พลํ อาคมฺม ตถาคโต อาสภํ ฐานํ ปฏิชานาติ, ปริสาสุ สีหนาทํ นทติ, พฺรหฺมจกฺกํ ปวตฺเตติ. (1)}

\prose{ปุน จปรํ อานนฺท ตถาคโต อตีตานาคตปจฺจุปฺปนฺนานํ กมฺมสมาทานานํ ฐานโส เหตุโส วิปากํ ยถาภูตํ ปชานาติ.\csromanspacer{} ยมฺปานนฺท ฯเปฯ.\csromanspacer{} อิทมฺปานนฺท ฯเปฯ. (2)}

\prose{ปุน จปรํ อานนฺท ตถาคโต สพฺพตฺถคามินิํ ปฏิปทํ ยถาภูตํ ปชานาติ.\csromanspacer{} ยมฺปานนฺท ฯเปฯ.\csromanspacer{} อิทมฺปานนฺท ฯเปฯ. (3)}

\prose{ปุน จปรํ อานนฺท ตถาคโต อเนกธาตุํ นานาธาตุํ โลกํ ยถาภูตํ ปชานาติ.\csromanspacer{} ยมฺปานนฺท ฯเปฯ.\csromanspacer{} อิทมฺปานนฺท ฯเปฯ. (4)}

\prose{ปุน จปรํ อานนฺท ตถาคโต สตฺตานํ นานาธิมุตฺติกตํ ยถาภูตํ ปชานาติ.\csromanspacer{} ยมฺปานนฺท ฯเปฯ.\csromanspacer{} อิทมฺปานนฺท ฯเปฯ. (5)}

\csromanmatikaanchor{mtk.48}
\csromantocmarkref{subsection}{2. อธิวุตฺติปทสุตฺต}{mtk.48}
\bookmark[dest={mtk.48},level=2]{2. อธิวุตฺติปทสุตฺต}
\prose{ปุน จปรํ อานนฺท ตถาคโต ปรสตฺตานํ ปรปุคฺคลานํ อินฺทฺริยปโรปริยตฺตํ ยถาภูตํ ปชานาติ.\csromanspacer{} ยมฺปานนฺท ฯเปฯ.\csromanspacer{} อิทมฺปานนฺท ฯเปฯ. (6)}

\prose{ปุน จปรํ อานนฺท ตถาคโต ฌานวิโมกฺขสมาธิสมาปตฺตีนํ สํกิเลสํ โวทานํ วุฏฺฐานํ ยถาภูตํ ปชานาติ.\csromanspacer{} ยมฺปานนฺท ฯเปฯ.\csromanspacer{} อิทมฺปานนฺท ฯเปฯ. (7)}

\chapterheadpageread{3. มหาวคฺค}
\prose{\csromanfolio{287}ปุน จปรํ อานนฺท ตถาคโต อเนกวิหิตํ ปุพฺเพนิวาสํ อนุสฺสรติ.\csromanspacer{} เสยฺยถิทํ, เอกมฺปิ ชาติํ เทฺวปิ ชาติโย ฯเปฯ อิติ สาการํ ส-อุทฺเทสํ อเนกวิหิตํ ปุพฺเพนิวาสํ อนุสฺสรติ.\csromanspacer{} ยมฺปานนฺท ฯเปฯ.\csromanspacer{} อิทมฺปานนฺท ฯเปฯ. (8)}

\prose{ปุน จปรํ อานนฺท ตถาคโต ทิพฺเพน จกฺขุนา วิสุทฺเธน อติกฺกนฺตมานุสเกน ฯเปฯ ยถากมฺมูปเค สตฺเต ปชานาติ.\csromanspacer{} ยมฺปานนฺท ฯเปฯ.\csromanspacer{} อิทมฺปานนฺท ฯเปฯ. (9)}

\prose{ปุน จปรํ อานนฺท ตถาคโต อาสวานํ ขยา อนาสวํ เจโตวิมุตฺติํ ปญฺญาวิมุตฺติํ ทิฏฺเฐว ธมฺเม สยํ อภิญฺญา สจฺฉิกตฺวา อุปสมฺปชฺช วิหรติ.\csromanspacer{} ยมฺปานนฺท ตถาคโต อาสวานํ ขยา อนาสวํ เจโตวิมุตฺติํ ฯเปฯ สจฺฉิกตฺวา อุปสมฺปชฺช วิหรติ.\csromanspacer{} อิทมฺปานนฺท ตถาคตสฺส ตถาคตพลํ โหติ, ยํ พลํ อาคมฺม ตถาคโต อาสภํ ฐานํ ปฏิชานาติ, ปริสาสุ สีหนาทํ นทติ, พฺรหฺมจกฺกํ ปวตฺเตติ. (10)}

\prose{อิมานิ โข อานนฺท ทส ตถาคตสฺส ตถาคตพลานิ, เยหิ พเลหิ สมนฺนาคโต ตถาคโต อาสภํ ฐานํ ปฏิชานาติ, ปริสาสุ สีหนาทํ นทติ, พฺรหฺมจกฺกํ ปวตฺเตตีติ.\csromanspacer{}\csromanspacer{}ทุติยํ.}

\csromanheader{2}{3. กายสุตฺต}
\proseitem{23}{อตฺถิ ภิกฺขเว ธมฺมา กาเยน ปหาตพฺพา โน วาจาย, อตฺถิ ภิกฺขเว ธมฺมา วาจาย ปหาตพฺพา โน กาเยน, อตฺถิ ภิกฺขเว ธมฺมา เนว กาเยน ปหาตพฺพา โน วาจาย, ปญฺญาย ทิสฺวา\footnote{ทิสฺวา ทิสฺวา (สี, สฺยา)} ปหาตพฺพา.}

\prose{กตเม จ ภิกฺขเว ธมฺมา กาเยน ปหาตพฺพา โน วาจาย? อิธ ภิกฺขเว ภิกฺขุ อกุสลํ อาปนฺโน โหติ กิญฺจิ เทสํ\footnote{กญฺจิ เทว เทสํ (สี, สฺยา)} กาเยน.\csromanspacer{} ตเมนํ อนุวิจฺจ วิญฺญู สพฺรหฺมจารี เอวมาหํสุ “อายสฺมา โข อกุสลํ อาปนฺโน กิญฺจิ เทสํ กาเยน, สาธุ วตายสฺมา กายทุจฺจริตํ ปหาย กายสุจริตํ ภาเวตู”ติ.\csromanspacer{} โส อนุวิจฺจ วิญฺญูหิ สพฺรหฺมจารีหิ วุจฺจมาโน กายทุจฺจริตํ ปหาย กายสุจริตํ ภาเวติ.\csromanspacer{} อิเม วุจฺจนฺติ ภิกฺขเว ธมฺมา กาเยน ปหาตพฺพา โน วาจาย.}

\gambhira{ทสกนิปาตปาฬิ}
\prose{\csromanfolio{288}กตเม จ ภิกฺขเว ธมฺมา วาจาย ปหาตพฺพา โน กาเยน? อิธ ภิกฺขเว ภิกฺขุ อกุสลํ อาปนฺโน โหติ กิญฺจิ เทสํ วาจาย.\csromanspacer{} ตเมนํ อนุวิจฺจ วิญฺญู สพฺรหฺมจารี เอวมาหํสุ “อายสฺมา โข อกุสลํ อาปนฺโน กิญฺจิ เทสํ วาจาย, สาธุ วตายสฺมา วจีทุจฺจริตํ ปหาย วจีสุจริตํ ภาเวตู”ติ.\csromanspacer{} โส อนุวิจฺจ วิญฺญูหิ สพฺรหฺมจารีหิ วุจฺจมาโน วจีทุจฺจริตํ ปหาย วจีสุจริตํ ภาเวติ.\csromanspacer{} อิเม วุจฺจนฺติ ภิกฺขเว ธมฺมา วาจาย ปหาตพฺพา โน กาเยน.}

\prose{กตเม จ ภิกฺขเว ธมฺมา เนว กาเยน ปหาตพฺพา โน วาจาย, ปญฺญาย ทิสฺวา ปหาตพฺพา? โลโภ ภิกฺขเว เนว กาเยน ปหาตพฺโพ โน วาจาย, ปญฺญาย ทิสฺวา ปหาตพฺโพ.\csromanspacer{} โทโส ภิกฺขเว ฯเปฯ.\csromanspacer{} โมโห.\csromanspacer{} โกโธ.\csromanspacer{} อุปนาโห.\csromanspacer{} มกฺโข.\csromanspacer{} ปฬาโส.\csromanspacer{} มจฺฉริยํ ภิกฺขเว เนว กาเยน ปหาตพฺพํ โน วาจาย, ปญฺญาย ทิสฺวา ปหาตพฺพํ.}

\prose{ปาปิกา ภิกฺขเว อิสฺสา เนว กาเยน ปหาตพฺพา โน วาจาย, ปญฺญาย ทิสฺวา ปหาตพฺพา.\csromanspacer{} กตมา จ ภิกฺขเว ปาปิกา อิสฺสา? อิธ ภิกฺขเว อิชฺฌติ คหปติสฺส วา คหปติปุตฺตสฺส วา ธเนน วา ธญฺเญน วา รชเตน วา ชาตรูเปน วา.\csromanspacer{} ตตฺราญฺญตรสฺส ทาสสฺส วา อุปวาสสฺส วา เอวํ โหติ “อโห วติมสฺส คหปติสฺส วา คหปติปุตฺตสฺส วา น อิชฺเฌยฺย ธเนน วา ธญฺเญน วา รชเตน วา ชาตรูเปน วา”ติ.\csromanspacer{} สมโณ วา ปน พฺราหฺมโณ วา ลาภี โหติ จีวรปิณฺฑปาตเสนาสนคิลานปจฺจยเภสชฺชปริกฺขารานํ, ตตฺราญฺญตรสฺส สมณสฺส วา พฺราหฺมณสฺส วา เอวํ โหติ “อโห วต อยมายสฺมา น ลาภี อสฺส จีวรปิณฺฑปาตเสนาสนคิลานปจฺจยเภสชฺช- ปริกฺขารานนฺ”ติ.\csromanspacer{} อยํ วุจฺจติ ภิกฺขเว ปาปิกา อิสฺสา.}

\prose{ปาปิกา ภิกฺขเว อิจฺฉา เนว กาเยน ปหาตพฺพา โน วาจาย, ปญฺญาย ทิสฺวา ปหาตพฺพา.\csromanspacer{} กตมา จ ภิกฺขเว ปาปิกา อิจฺฉา? * อิธ ภิกฺขเว เอกจฺโจ อสฺสทฺโธ สมาโน “สทฺโธติ มํ ชาเนยฺยุนฺ”ติ อิจฺฉติ, ทุสฺสีโล สมาโน “สีลวาติ มํ ชาเนยฺยุนฺ”ติ อิจฺฉติ, อปฺปสฺสุโต สมาโน “พหุสฺสุโตติ มํ ชาเนยฺยุนฺ”ติ อิจฺฉติ, สงฺคณิการาโม สมาโน “ปวิวิตฺโตติ มํ ชาเนยฺยุนฺ”ติ อิจฺฉติ, กุสีโต สมาโน “อารทฺธวีริโยติ มํ ชาเนยฺยุนฺ”ติ อิจฺฉติ, มุฏฺฐสฺสติ สมาโน “อุปฏฺฐิตสฺสตีติ}

\vspace{\tipitakaparskip}
\csromancenter{มหาวคฺค มํ ชาเนยฺยุนฺ”ติ อิจฺฉติ, อสมาหิโต สมาโน “สมาหิโตติ มํ ชาเนยฺยุนฺ”ติ อิจฺฉติ, ทุปฺปญฺโญ สมาโน “ปญฺญวาติ มํ ชาเนยฺยุนฺ”ติ อิจฺฉติ, อขีณาสโว สมาโน “ขีณาสโวติ มํ ชาเนยฺยุนฺ”ติ อิจฺฉติ, อยํ วุจฺจติ ภิกฺขเว ปาปิกา อิจฺฉา, อิเม วุจฺจนฺติ ภิกฺขเว ธมฺมา เนว กาเยน ปหาตพฺพา โน วาจาย, ปญฺญาย ทิสฺวา ปหาตพฺพา.}
\csromanmatikaanchor{mtk.49}
\csromantocmarkref{subsection}{3. กายสุตฺต}{mtk.49}
\bookmark[dest={mtk.49},level=2]{3. กายสุตฺต}
\prose{\csromanfolio{289}ตญฺเจ ภิกฺขเว ภิกฺขุํ โลโภ อภิภุยฺย อิริยติ, โทโส.\csromanspacer{} โมโห.\csromanspacer{} โกโธ.\csromanspacer{} อุปนาโห.\csromanspacer{} มกฺโข.\csromanspacer{} ปฬาโส.\csromanspacer{} มจฺฉริยํ.\csromanspacer{} ปาปิกา อิสฺสา.\csromanspacer{} ปาปิกา อิจฺฉา อภิภุยฺย อิริยติ, โส เอวมสฺส เวทิตพฺโพ “นายมายสฺมา ตถา ปชานาติ, ยถา ปชานโต โลโภ น โหติ, ตถาหิมํ อายสฺมนฺตํ โลโภ อภิภุยฺย อิริยติ.\csromanspacer{} นายมายสฺมา ตถา ปชานาติ, ยถา ปชานโต โทโส น โหติ, โมโห.\csromanspacer{} โกโธ.\csromanspacer{} อุปนาโห.\csromanspacer{} มกฺโข.\csromanspacer{} ปฬาโส.\csromanspacer{} มจฺฉริยํ.\csromanspacer{} ปาปิกา อิสฺสา.\csromanspacer{} ปาปิกา อิจฺฉา น โหติ, ตถาหิมํ อายสฺมนฺตํ ปาปิกา อิจฺฉา อภิภุยฺย อิริยตี”ติ.}

\prose{ตญฺเจ ภิกฺขเว ภิกฺขุํ โลโภ นาภิภุยฺย อิริยติ, โทโส.\csromanspacer{} โมโห.\csromanspacer{} โกโธ.\csromanspacer{} อุปนาโห.\csromanspacer{} มกฺโข.\csromanspacer{} ปฬาโส.\csromanspacer{} มจฺฉริยํ.\csromanspacer{} ปาปิกา อิสฺสา.\csromanspacer{} ปาปิกา อิจฺฉา นาภิภุยฺย อิริยติ, โส เอวมสฺส เวทิตพฺโพ “ตถา อยมายสฺมา ปชานาติ, ยถา ปชานโต โลโภ น โหติ, ตถาหิมํ อายสฺมนฺตํ โลโภ นาภิภุยฺย อิริยติ.\csromanspacer{} ตถา อยมายสฺมา ปชานาติ, ยถา ปชานโต โทโส น โหติ, โมโห.\csromanspacer{} โกโธ.\csromanspacer{} อุปนาโห.\csromanspacer{} มกฺโข.\csromanspacer{} ปฬาโส.\csromanspacer{} มจฺฉริยํ.\csromanspacer{} ปาปิกา อิสฺสา.\csromanspacer{} ปาปิกา อิจฺฉา น โหติ, ตถาหิมํ อายสฺมนฺตํ ปาปิกา อิจฺฉา นาภิภุยฺย อิริยตี”ติ.\csromanspacer{}\csromanspacer{}ตติยํ.}

\csromanheader{2}{4. มหาจุนฺทสุตฺต}
\proseitem{24}{เอกํ สมยํ อายสฺมา มหาจุนฺโท เจตีสุ วิหรติ สหชาติยํ.\csromanspacer{} ตตฺร โข อายสฺมา มหาจุนฺโท ภิกฺขู อามนฺเตสิ “อาวุโส ภิกฺขเว”ติ. “อาวุโส”ติ โข เต ภิกฺขู อายสฺมโต มหาจุนฺทสฺส ปจฺจสฺโสสุํ.\csromanspacer{} อายสฺมา มหาจุนฺโท เอตทโวจ\csromandash{}}

\prose{ญาณวาทํ อาวุโส ภิกฺขุ วทมาโน “ชานามิมํ ธมฺมํ ปสฺสามิมํ ธมฺมนฺ”ติ.\csromanspacer{} ตญฺเจ อาวุโส ภิกฺขุํ โลโภ อภิภุยฺย ติฏฺฐติ, โทโส.}

\vspace{\tipitakaparskip}
\csromancenter{ทสกนิปาตปาฬิ โมโห.\csromanspacer{} โกโธ.\csromanspacer{} อุปนาโห.\csromanspacer{} มกฺโข.\csromanspacer{} ปฬาโส.\csromanspacer{} มจฺฉริยํ.\csromanspacer{} ปาปิกา อิสฺสา.\csromanspacer{} ปาปิกา อิจฺฉา อภิภุยฺย ติฏฺฐติ, โส เอวมสฺส เวทิตพฺโพ “นายมายสฺมา ตถา ปชานาติ, ยถา ปชานโต โลโภ น โหติ, ตถาหิมํ อายสฺมนฺตํ โลโภ อภิภุยฺย ติฏฺฐติ.\csromanspacer{} นายมายสฺมา ตถา ปชานาติ, ยถา ปชานโต โทโส น โหติ, โมโห.\csromanspacer{} โกโธ.\csromanspacer{} อุปนาโห.\csromanspacer{} มกฺโข.\csromanspacer{} ปฬาโส.\csromanspacer{} มจฺฉริยํ.\csromanspacer{} ปาปิกา อิสฺสา.\csromanspacer{} ปาปิกา อิจฺฉา น โหติ, ตถาหิมํ อายสฺมนฺตํ ปาปิกา อิจฺฉา อภิภุยฺย ติฏฺฐตี”ติ.}
\prose{\csromanfolio{290}ภาวนาวาทํ อาวุโส ภิกฺขุ วทมาโน “ภาวิตกาโยมฺหิ ภาวิตสีโล ภาวิตจิตฺโต ภาวิตปญฺโญ”ติ.\csromanspacer{} ตญฺเจ อาวุโส ภิกฺขุํ โลโภ อภิภุยฺย ติฏฺฐติ, โทโส.\csromanspacer{} โมโห.\csromanspacer{} โกโธ.\csromanspacer{} อุปนาโห.\csromanspacer{} มกฺโข.\csromanspacer{} ปฬาโส.\csromanspacer{} มจฺฉริยํ.\csromanspacer{} ปาปิกา อิสฺสา.\csromanspacer{} ปาปิกา อิจฺฉา อภิภุยฺย ติฏฺฐติ, โส เอวมสฺส เวทิตพฺโพ “นายมายสฺมา ตถา ปชานาติ, ยถา ปชานโต โลโภ น โหติ, ตถาหิมํ อายสฺมนฺตํ โลโภ อภิภุยฺย ติฏฺฐติ.\csromanspacer{} นายมายสฺมา ตถา ปชานาติ, ยถา ปชานโต โทโส น โหติ, โมโห.\csromanspacer{} โกโธ.\csromanspacer{} อุปนาโห.\csromanspacer{} มกฺโข.\csromanspacer{} ปฬาโส.\csromanspacer{} มจฺฉริยํ.\csromanspacer{} ปาปิกา อิสฺสา.\csromanspacer{} ปาปิกา อิจฺฉา น โหติ, ตถาหิมํ อายสฺมนฺตํ ปาปิกา อิจฺฉา อภิภุยฺย ติฏฺฐตี”ติ.}

\prose{ญาณวาทญฺจ อาวุโส ภิกฺขุ วทมาโน ภาวนาวาทญฺจ “ชานามิมํ ธมฺมํ ปสฺสามิมํ ธมฺมํ, ภาวิตกาโยมฺหิ ภาวิตสีโล ภาวิตจิตฺโต ภาวิตปญฺโญ”ติ.\csromanspacer{} ตญฺเจ อาวุโส ภิกฺขุํ โลโภ อภิภุยฺย ติฏฺฐติ, โทโส.\csromanspacer{} โมโห.\csromanspacer{} โกโธ.\csromanspacer{} อุปนาโห.\csromanspacer{} มกฺโข.\csromanspacer{} ปฬาโส.\csromanspacer{} มจฺฉริยํ.\csromanspacer{} ปาปิกา อิสฺสา.\csromanspacer{} ปาปิกา อิจฺฉา อภิภุยฺย ติฏฺฐติ, โส เอวมสฺส เวทิตพฺโพ “นายมายสฺมา ตถา ปชานาติ, ยถา ปชานโต โลโภ น โหติ, ตถาหิมํ อายสฺมนฺตํ โลโภ อภิภุยฺย ติฏฺฐติ.\csromanspacer{} นายมายสฺมา ตถา ปชานาติ, ยถา ปชานโต โทโส น โหติ, โมโห.\csromanspacer{} โกโธ.\csromanspacer{} อุปนาโห.\csromanspacer{} มกฺโข.\csromanspacer{} ปฬาโส.\csromanspacer{} มจฺฉริยํ.\csromanspacer{} ปาปิกา อิสฺสา.\csromanspacer{} ปาปิกา อิจฺฉา น โหติ, ตถาหิมํ อายสฺมนฺตํ ปาปิกา อิจฺฉา อภิภุยฺย ติฏฺฐตี”ติ.}

\chapterheadpageread{3. มหาวคฺค}
\prose{\csromanfolio{291}เสยฺยถาปิ อาวุโส ปุริโส ทลิทฺโทว สมาโน อฑฺฒวาทํ วเทยฺย, อธโนว สมาโน ธนวาวาทํ วเทยฺย, อโภโคว สมาโน โภควาวาทํ วเทยฺย, โส กิสฺมิญฺจิเทว ธนกรณีเย สมุปฺปนฺเน น สกฺกุเณยฺย อุปนีหาตุํ ธนํ วา ธญฺญํ วา รชตํ วา ชาตรูปํ วา.\csromanspacer{} ตเมนํ เอวํ ชาเนยฺยุํ “ทลิทฺโทว อยมายสฺมา สมาโน อฑฺฒวาทํ วเทติ, อธโนว อยมายสฺมา สมาโน ธนวาวาทํ วเทติ, อโภควาว อยมายสฺมา สมาโน โภควาวาทํ วเทติ.\csromanspacer{} ตํ กิสฺส เหตุ? ตถา หิ อยมายสฺมา กิสฺมิญฺจิเทว ธนกรณีเย สมุปฺปนฺเน น สกฺโกติ อุปนีหาตุํ ธนํ วา ธญฺญํ วา รชตํ วา ชาตรูปํ วา”ติ.}

\prose{เอวเมวํ โข อาวุโส ญาณวาทญฺจ ภิกฺขุ วทมาโน ภาวนาวาทญฺจ “ชานามิมํ ธมฺมํ ปสฺสามิมํ ธมฺมํ, ภาวิตกาโยมฺหิ ภาวิตสีโล ภาวิตจิตฺโต ภาวิตปญฺโญ”ติ.\csromanspacer{} ตํ เจ อาวุโส ภิกฺขุํ โลโภ อภิภุยฺย ติฏฺฐติ, โทโส.\csromanspacer{} โมโห.\csromanspacer{} โกโธ.\csromanspacer{} อุปนาโห.\csromanspacer{} มกฺโข.\csromanspacer{} ปฬาโส.\csromanspacer{} มจฺฉริยํ.\csromanspacer{} ปาปิกา อิสฺสา.\csromanspacer{} ปาปิกา อิจฺฉา อภิภุยฺย ติฏฺฐติ, โส เอวมสฺส เวทิตพฺโพ “นายมายสฺมา ตถา ปชานาติ, ยถา ปชานโต โลโภ น โหติ, ตถาหิมํ อายสฺมนฺตํ โลโภ อภิภุยฺย ติฏฺฐติ.\csromanspacer{} นายมายสฺมา ตถา ปชานาติ, ยถา ปชานโต โทโส น โหติ, โมโห.\csromanspacer{} โกโธ.\csromanspacer{} อุปนาโห.\csromanspacer{} มกฺโข.\csromanspacer{} ปฬาโส.\csromanspacer{} มจฺฉริยํ.\csromanspacer{} ปาปิกา อิสฺสา.\csromanspacer{} ปาปิกา อิจฺฉา น โหติ.\csromanspacer{} ตถาหิมํ อายสฺมนฺตํ ปาปิกา อิจฺฉา อภิภุยฺย ติฏฺฐตี”ติ.}

\csromanmatikaanchor{mtk.50}
\csromantocmarkref{subsection}{4. มหาจุนฺทสุตฺต}{mtk.50}
\bookmark[dest={mtk.50},level=2]{4. มหาจุนฺทสุตฺต}
\prose{ญาณวาทํ อาวุโส ภิกฺขุ วทมาโน “ชานามิมํ ธมฺมํ ปสฺสามิมํ ธมฺมนฺ”ติ.\csromanspacer{} ตํ เจ อาวุโส ภิกฺขุํ โลโภ นาภิภุยฺย ติฏฺฐติ, โทโส.\csromanspacer{} โมโห.\csromanspacer{} โกโธ.\csromanspacer{} อุปนาโห.\csromanspacer{} มกฺโข.\csromanspacer{} ปฬาโส.\csromanspacer{} มจฺฉริยํ.\csromanspacer{} ปาปิกา อิสฺสา.\csromanspacer{} ปาปิกา อิจฺฉา นาภิภุยฺย ติฏฺฐติ.\csromanspacer{} โส เอวมสฺส เวทิตพฺโพ “อยมายสฺมา ตถา ปชานาติ, ยถา ปชานโต โลโภ น โหติ, ตถาหิมํ อายสฺมนฺตํ โลโภ นาภิภุยฺย ติฏฺฐติ.\csromanspacer{} ตถา อยมายสฺมา ปชานาติ, ยถา ปชานโต โทโส น โหติ, โมโห.\csromanspacer{} โกโธ.\csromanspacer{} อุปนาโห.\csromanspacer{} มกฺโข.\csromanspacer{} ปฬาโส.\csromanspacer{} มจฺฉริยํ.\csromanspacer{} ปาปิกา อิสฺสา.\csromanspacer{} ปาปิกา อิจฺฉา น โหติ, ตถาหิมํ อายสฺมนฺตํ ปาปิกา อิจฺฉา นาภิภุยฺย ติฏฺฐตี”ติ.}

\gambhira{ทสกนิปาตปาฬิ}
\prose{\csromanfolio{292}ภาวนาวาทํ อาวุโส ภิกฺขุ วทมาโน “ภาวิตกาโยมฺหิ ภาวิตสีโล ภาวิตจิตฺโต ภาวิตปญฺโญ”ติ.\csromanspacer{} ตํ เจ อาวุโส ภิกฺขุํ โลโภ นาภิภุยฺย ติฏฺฐติ, โทโส.\csromanspacer{} โมโห.\csromanspacer{} โกโธ.\csromanspacer{} อุปนาโห.\csromanspacer{} มกฺโข.\csromanspacer{} ปฬาโส.\csromanspacer{} มจฺฉริยํ.\csromanspacer{} ปาปิกา อิสฺสา.\csromanspacer{} ปาปิกา อิจฺฉา นาภิภุยฺย ติฏฺฐติ, โส เอวมสฺส เวทิตพฺโพ “ตถา อยมายสฺมา ปชานาติ, ยถา ปชานโต โลโภ น โหติ, ตถาหิมํ อายสฺมนฺตํ โลโภ นาภิภุยฺย ติฏฺฐติ.\csromanspacer{} ตถา อยมายสฺมา ปชานาติ, ยถา ปชานโต โทโส น โหติ, โมโห.\csromanspacer{} โกโธ.\csromanspacer{} อุปนาโห.\csromanspacer{} มกฺโข.\csromanspacer{} ปฬาโส.\csromanspacer{} มจฺฉริยํ.\csromanspacer{} ปาปิกา อิสฺสา.\csromanspacer{} ปาปิกา อิจฺฉา น โหติ, ตถาหิมํ อายสฺมนฺตํ ปาปิกา อิจฺฉา นาภิภุยฺย ติฏฺฐตี”ติ.}

\prose{ญาณวาทญฺจ อาวุโส ภิกฺขุ วทมาโน ภาวนาวาทญฺจ “ชานามิมํ ธมฺมํ ปสฺสามิมํ ธมฺมํ, ภาวิตกาโยมฺหิ ภาวิตสีโล ภาวิตจิตฺโต ภาวิตปญฺโญ”ติ.\csromanspacer{} ตํ เจ อาวุโส ภิกฺขุํ โลโภ นาภิภุยฺย ติฏฺฐติ, โทโส.\csromanspacer{} โมโห.\csromanspacer{} โกโธ.\csromanspacer{} อุปนาโห.\csromanspacer{} มกฺโข.\csromanspacer{} ปฬาโส.\csromanspacer{} มจฺฉริยํ.\csromanspacer{} ปาปิกา อิสฺสา.\csromanspacer{} ปาปิกา อิจฺฉา นาภิภุยฺย ติฏฺฐติ, โส เอวมสฺส เวทิตพฺโพ “ตถา อยมายสฺมา ปชานาติ, ยถา ปชานาโต โลโภ น โหติ, ตถาหิมํ อายสฺมนฺตํ โลโภ นาภิภุยฺย ติฏฺฐติ.\csromanspacer{} ตถา อยมายสฺมา ปชานาติ, ยถา ปชานโต โทโส น โหติ.\csromanspacer{} โมโห.\csromanspacer{} โกโธ.\csromanspacer{} อุปนาโห.\csromanspacer{} มกฺโข.\csromanspacer{} ปฬาโส.\csromanspacer{} มจฺฉริยํ.\csromanspacer{} ปาปิกา อิสฺสา.\csromanspacer{} ปาปิกา อิจฺฉา น โหติ, ตถาหิมํ อายสฺมนฺตํ ปาปิกา อิจฺฉา นาภิภุยฺย ติฏฺฐตี”ติ.}

\prose{เสยฺยถาปิ อาวุโส ปุริโส อฑฺโฒว สมาโน อฑฺฒวาทํ วเทยฺย, ธนวาว สมาโน ธนวาวาทํ วเทยฺย, โภควาว สมาโน โภควาวาทํ วเทยฺย, โส กิสฺมิญฺจิเทว ธนกรณีเย สมุปฺปนฺเน สกฺกุเณยฺย อุปนีหาตุํ ธนํ วา ธญฺญ วา รชตํ วา ชาตรูปํ วา.\csromanspacer{} ตเมนํ เอวํ ชาเนยฺยุํ “อฑฺโฒว อยมายสฺมา สมาโน อฑฺฒวาทํ วเทติ, ธนวาว อยมายสฺมา สมาโน ธนวาวาทํ วเทติ, โภควาว อยมายสฺมา สมาโน โภควาวาทํ วเทติ.\csromanspacer{} ตํ กิสฺส เหตุ? ตถา หิ อยมายสฺมา กิสฺมิญฺจิเทว ธนกรณีเย สมุปฺปนฺเน สกฺโกติ อุปนีหาตุํ ธนํ วา ธญฺญํ วา รชตํ วา ชาตรูปํ วา”ติ.}

\chapterheadpageread{3. มหาวคฺค}
\prose{\csromanfolio{293}เอวเมวํ โข อาวุโส ญาณวาทญฺจ ภิกฺขุ วทมาโน ภาวนาวาทญฺจ “ชานามิมํ ธมฺมํ ปสฺสามิมํ ธมฺมํ, ภาวิตกาโยมฺหิ ภาวิตสีโล ภาวิตจิตฺโต ภาวิตปญฺโญ”ติ.\csromanspacer{} ตํ เจ อาวุโส ภิกฺขุํ โลโภ นาภิภุยฺย ติฏฺฐติ, โทโส.\csromanspacer{} โมโห.\csromanspacer{} โกโธ.\csromanspacer{} อุปนาโห.\csromanspacer{} มกฺโข.\csromanspacer{} ปฬาโส.\csromanspacer{} มจฺฉริยํ.\csromanspacer{} ปาปิกา อิสฺสา.\csromanspacer{} ปาปิกา อิจฺฉา นาภิภุยฺย ติฏฺฐติ.\csromanspacer{} โส เอวมสฺส เวทิตพฺโพ “ตถา อยมายสฺมา ปชานาติ, ยถา ปชานโต โลโภ น โหติ, ตถาหิมํ อายสฺมนฺตํ โลโภ นาภิภุยฺย ติฏฺฐติ.\csromanspacer{} ตถา อยมายสฺมา ปชานาติ, ยถา ปชานโต โทโส น โหติ.\csromanspacer{} โมโห.\csromanspacer{} โกโธ.\csromanspacer{} อุปนาโห.\csromanspacer{} มกฺโข.\csromanspacer{} ปฬาโส.\csromanspacer{} มจฺฉริยํ.\csromanspacer{} ปาปิกา อิสฺสา.\csromanspacer{} ปาปิกา อิจฺฉา น โหติ, ตถาหิมํ อายสฺมนฺตํ ปาปิกา อิจฺฉา นาภิภุยฺย ติฏฺฐตี”ติ.\csromanspacer{}\csromanspacer{}จตุตฺถํ.}

\csromanheader{2}{5. กสิณสุตฺต}
\proseitem{25}{\csromansymbolfootnote{*}{ที 3. 223, 257; อุปริ 304 ปิฏฺเฐสุปิ.}ทสยิมานิ ภิกฺขเว กสิณายตนานิ.\csromanspacer{} กตมานิ ทส? ปถวีกสิณเมโก สญฺชานาติ อุทฺธํ อโธ ติริยํ อทฺวยํ อปฺปมาณํ, อาโปกสิณเมโก สญฺชานาติ ฯเปฯ.\csromanspacer{} เตโชกสิณเมโก สญฺชานาติ.\csromanspacer{} วาโยกสิณเมโก สญฺชานาติ.\csromanspacer{} นีลกสิณเมโก สญฺชานาติ.\csromanspacer{} ปีตกสิณเมโก สญฺชานาติ.\csromanspacer{} โลหิตกสิณเมโก สญฺชานาติ.\csromanspacer{} โอทาตกสิณเมโก สญฺชานาติ.\csromanspacer{} อากาสกสิณเมโก สญฺชานาติ.\csromanspacer{} วิญฺญาณกสิณเมโก สญฺชานาติ อุทฺธํ อโธ ติริยํ อทฺวยํ อปฺปมาณํ.\csromanspacer{} อิมานิ โข ภิกฺขเว ทส กสิณายตนานีติ.\csromanspacer{}\csromanspacer{}ปญฺจมํ.}

\csromanheader{2}{6. กาฬีสุตฺต}
\proseitem{26}{เอกํ สมยํ อายสฺมา มหากจฺจาโน อวนฺตีสุ วิหรติ กุรรฆเร\footnote{กุลฆเร (ก), กุรุรฆเร (วิ 3. 283 ปิฏฺเฐ.)} ปวตฺเต ปพฺพเต.\csromanspacer{} อถ โข กาฬี อุปาสิกา กุรรฆริกา เยนายสฺมา มหากจฺจาโน เตนุปสงฺกมิ, อุปสงฺกมิตฺวา อายสฺมนฺตํ มหากจฺจานํ อภิวาเทตฺวา เอกมนฺตํ นิสีทิ, เอกมนฺตํ นิสินฺโน โข กาฬี}

\vspace{\tipitakaparskip}
\csromancenter{ทสกนิปาตปาฬิ อุปาสิกา กุรรฆริกา อายสฺมนฺตํ มหากจฺจานํ เอตทโวจ “วุตฺตมิทํ ภนฺเต ภวคตา กุมาริปเญฺหสุ\csromandash{}}
\vspace{0.5\baselineskip}
\csromangathameasure{‘อตฺถสฺส ปตฺติํ หทยสฺส สนฺติํ, \\ เชตฺวาน เสนํ ปิยสาตรูปํ. \\ เอโกหํ ฌายํ สุขมนุโพธิํ, \\ ตสฺมา ชเนน น กโรมิ สกฺขิํ. \\ สกฺขี น สมฺปชฺชติ เกนจิ เม’ติ.}
\csromangathagroup{\csromangathastanza{\csromanfolio{294}‘อตฺถสฺส ปตฺติํ หทยสฺส สนฺติํ, \\ เชตฺวาน เสนํ ปิยสาตรูปํ. \\ เอโกหํ\footnote{เอกาหํ (ก)} ฌายํ สุขมนุโพธิํ, \\ ตสฺมา ชเนน น กโรมิ สกฺขิํ\footnote{เอกาหํ (ก)}.}\csromangathastanza{สกฺขี\footnote{สขี (ก)} น สมฺปชฺชติ เกนจิ เม’ติ.}}

\prose{อิมสฺส โข ภนฺเต ภควตา สํขิตฺเตน ภาสิตสฺส กถํ วิตฺถาเรน อตฺโถ ทฏฺฐพฺโพ”ติ.}

\prose{ปถวีกสิณสมาปตฺติปรมา โข ภคินิ เอเก สมณพฺราหฺมณา อตฺโถติ อภินิพฺพตฺเตสุํ\footnote{อตฺถาภินิพฺพตฺเตสุํ (สี, สฺยา)}.\csromanspacer{} ยาวตา โข ภคินิ ปถวีกสิณสมาปตฺติปรมตา, ตทภิญฺญาสิ ภควา, ตทภิญฺญาย ภควา อสฺสาทมทฺทส\footnote{อาทิมทฺทส (สี, สฺยา)}, อาทีนวมทฺทส.\csromanspacer{} นิสฺสรณมทฺทส.\csromanspacer{} มคฺคามคฺคญาณมทฺทส.\csromanspacer{} ตสฺส อสฺสาททสฺสนเหตุ อาทีนวทสฺสนเหตุ นิสฺสรณทสฺสนเหตุ มคฺคามคฺคญาณทสฺสนเหตุ อตฺถสฺส ปตฺติ หทยสฺส สนฺติ วิทิตา โหติ.}

\csromanmatikaanchor{mtk.51}
\csromantocmarkref{subsection}{5. กสิณสุตฺต}{mtk.51}
\bookmark[dest={mtk.51},level=2]{5. กสิณสุตฺต}
\prose{อาโปกสิณสมาปตฺติปรมา โข ภคินิ ฯเปฯ.\csromanspacer{} เตโชกสิณสมาปตฺติปรมา โข ภคินิ.\csromanspacer{} วาโยกสิณสมาปตฺติปรมา โข ภคินิ.\csromanspacer{} นีลกสิณสมาปตฺติปรมา โข ภคินิ.\csromanspacer{} ปีตกสิณสมาปตฺติปรมา โข ภคินิ.\csromanspacer{} โลหิตกสิณสมาปตฺติปรมา โข ภคินิ.\csromanspacer{} โอทาตกสิณสมาปตฺติปรมา โข ภคินิ.\csromanspacer{} อากาสกสิณสมาปตฺติปรมา โข ภคินิ.\csromanspacer{} วิญฺญาณกสิณสมาปตฺติปรมา โข ภคินิ เอเก สมณพฺราหฺมณา อตฺโถติ อภินิพฺพตฺเตสุํ.\csromanspacer{} ยาวตา โข ภคินิ วิญฺญาณกสิณสมาปตฺติปรมตา, ตทภิญฺญาสิ ภควา, ตทภิญฺญาย ภควา อสฺสาทมทฺทส.\csromanspacer{} อาทีนวมทฺทส.\csromanspacer{} นิสฺสรณมทฺทส.\csromanspacer{} มคฺคามคฺคญาณทสฺสนมทฺทส.\csromanspacer{} ตสฺส อสฺสาททสฺสนเหตุ อาทีนวทสฺสนเหตุ นิสฺสรณทสฺสนเหตุ มคฺคามคฺคญาณทสฺสนเหตุ อตฺถสฺส ปตฺติ หทยสฺส สนฺติ วิทิตา โหติ.\csromanspacer{} อิติ โข ภคินิ ยํ ตํ วุตฺตํ ภควตา กุมาริปเญฺหสุ\csromandash{}}

\chapterheadpageread{3. มหาวคฺค}
\csromanmatikaanchor{mtk.52}
\csromantocmarkref{subsection}{6. กาฬีสุตฺต}{mtk.52}
\bookmark[dest={mtk.52},level=2]{6. กาฬีสุตฺต}
\csromangathameasure{“อตฺถสฺส ปตฺติํ หทยสฺส สนฺติํ, \\ เชตฺวาน เสนํ ปิยสาตรูปํ. \\ เอโกหํ ฌายํ สุขมนุโพธิํ, \\ ตสฺมา ชเนน น กโรมิ สกฺขิํ. \\ สกฺขี น สมฺปชฺชติ เกนจิ เม”ติ.}
\csromangathagroup{\csromangathastanza{\csromanfolio{295}“อตฺถสฺส ปตฺติํ หทยสฺส สนฺติํ, \\ เชตฺวาน เสนํ ปิยสาตรูปํ. \\ เอโกหํ ฌายํ สุขมนุโพธิํ, \\ ตสฺมา ชเนน น กโรมิ สกฺขิํ.}\csromangathastanza{สกฺขี น สมฺปชฺชติ เกนจิ เม”ติ.}}

\prose{อิมสฺส โข ภคินิ ภควตา สํขิตฺเตน ภาสิตสฺส เอวํ วิตฺถาเรน อตฺโถ ทฏฺฐพฺโพติ.\csromanspacer{}\csromanspacer{}ฉฏฺฐํ.}

\csromanheader{2}{7. ปฐมมหาปญฺหาสุตฺต}
\proseitem{27}{เอกํ สมยํ ภควา สาวตฺถิยํ วิหรติ เชตวเน อนาถปิณฺฑิกสฺส อาราเม.\csromanspacer{} อถ โข สมฺพหุลา ภิกฺขู ปุพฺพณฺหสมยํ นิวาเสตฺวา ปตฺตจีวรมาทาย สาวตฺถิํ ปิณฺฑาย ปวิสิํสุ.\csromanspacer{} อถ โข เตสํ ภิกฺขูนํ เอตทโหสิ “อติปฺปโค โข ตาว สาวตฺถิยํ ปิณฺฑาย จริตุํ, ยํนูน มยํ เยน อญฺญติตฺถิยานํ ปริพฺพาชกานํ อาราโม เตนุปสงฺกเมยฺยามา”ติ.}

\prose{อถ โข เต ภิกฺขู เยน อญฺญติตฺถิยานํ ปริพฺพาชกานํ อาราโม เตนุปสงฺกมิํสุ, อุปสงฺกมิตฺวา เตหิ อญฺญติตฺถิเยหิ ปริพฺพาชเกหิ สทฺธิํ สมฺโมทิํสุ, สมฺโมทนียํ กถํ สารณียํ วีติสาเรตฺวา เอกมนฺตํ นิสีทิํสุ, เอกมนฺตํ นิสินฺเน โข เต ภิกฺขู เต อญฺญติตฺถิยา ปริพฺพาชกา เอตทโวจุํ\csromandash{}}

\prose{สมโณ อาวุโส โคตโม สาวกานํ เอวํ ธมฺมํ เทเสติ “เอถ ตุเมฺห ภิกฺขเว สพฺพํ ธมฺมํ อภิชานาถ, สพฺพํ ธมฺมํ อภิญฺญาย วิหรถา”ติ.\csromanspacer{} มยมฺปิ โข อาวุโส สาวกานํ เอวํ ธมฺมํ เทเสม “เอถ ตุเมฺห อาวุโส สพฺพํ ธมฺมํ อภิชานาถ, สพฺพํ ธมฺมํ อภิญฺญาย วิหรถา”ติ.\csromanspacer{} อิธ โน อาวุโส โก วิเสโส โก อธิปฺปยาโส กิํ นานากรณํ สมณสฺส วา โคตมสฺส อมฺหากํ วา ยทิทํ ธมฺมเทสนาย วา ธมฺมเทสนํ อนุสาสนิยา วา อนุสาสนินฺติ.}

\prose{อถ โข เต ภิกฺขู เตสํ อญฺญติตฺถิยานํ ปริพฺพาชกานํ ภาสิตํ เนว อภินนฺทิํสุ นปฺปฏิกฺโกสิํสุ, อนภินนฺทิตฺวา อปฺปฏิกฺโกสิตฺวา}

\vspace{\tipitakaparskip}
\csromancenter{ทสกนิปาตปาฬิ อุฏฺฐายาสนา ปกฺกมิํสุ “ภควโต สนฺติเก เอตสฺส ภาสิตสฺส อตฺถํ อาชานิสฺสามา”ติ.}
\prose{\csromanfolio{296}อถ โข เต ภิกฺขู สาวตฺถิยํ ปิณฺฑาย จริตฺวา ปจฺฉาภตฺตํ ปิณฺฑปาตปฏิกฺกนฺตา เยน ภควา เตนุปสงฺกมิํสุ, อุปสงฺกมิตฺวา ภควนฺตํ อภิวาเทตฺวา เอกมนฺตํ นิสีทิํสุ, เอกมนฺตํ นิสินฺนา โข เต ภิกฺขู ภควนฺตํ เอตทโวจุํ\csromandash{}}

\prose{อิธ มยํ ภนฺเต ปุพฺพณฺหสมยํ นิวาเสตฺวา ปตฺตจีวรมาทาย สาวตฺถิํ ปิณฺฑาย ปวิสิมฺหา, เตสํ โน ภนฺเต อมฺหากํ เอตทโหสิ “อติปฺปโค โข ตาว สาวตฺถิยํ ปิณฺฑาย จริตุํ, ยํนูน มยํ เยน อญฺญติตฺถิยานํ ปริพฺพาชกานํ อาราโม เตนุสงฺกเมยฺยามา”ติ.\csromanspacer{} อถ โข มยํ ภนฺเต เยน อญฺญติตฺถิยานํ ปริพฺพาชกานํ อาราโม เตนุปสงฺกมิมฺหา, อุปสงฺกมิตฺวา เตหิ อญฺญติตฺถิเยหิ ปริพฺพาชเกหิ สทฺธิํ สมฺโมทิมฺหา, สมฺโมทนียํ กถํ สารณียํ วีติสาเรตฺวา เอกมนฺตํ นิสีทิมฺหา, เอกมนฺตํ นิสินฺเน โข ภนฺเต อญฺญติตฺถิยา ปริพฺพาชกา อเมฺห เอตทโวจุํ\csromandash{}}

\prose{สมโณ อาวุโส โคตโม สาวกานํ เอวํ ธมฺมํ เทเสติ “เอถ ตุเมฺห ภิกฺขเว สพฺพํ ธมฺมํ อภิชานาถ, สพฺพํ ธมฺมํ อภิญฺญาย วิหรถา”ติ.\csromanspacer{} มยมฺปิ โข อาวุโส สาวกานํ เอวํ ธมฺมํ เทเสม “เอถ ตุเมฺห อาวุโส สพฺพํ ธมฺมํ อภิชานาถ, สพฺพํ ธมฺมํ อภิญฺญาย วิหรถา”ติ.\csromanspacer{} อิธ โน อาวุโส โก วิเสโส โก อธิปฺปยาโส กิํ นานากรณํ สมณสฺส วา โคตมสฺส อมฺหากํ วา ยทิทํ ธมฺมเทสนาย วา ธมฺมเทสนํ อนุสาสนิยา วา อนุสาสนินฺติ.}

\csromanmatikaanchor{mtk.53}
\csromantocmarkref{subsection}{7. ปฐมมหาปญฺหาสุตฺต}{mtk.53}
\bookmark[dest={mtk.53},level=2]{7. ปฐมมหาปญฺหาสุตฺต}
\prose{อถ โข มยํ ภนฺเต เตสํ อญฺญติตฺถิยานํ ปริพฺพาชกานํ ภาสิตํ เนว อภินนฺทิมฺหา นปฺปฏิกฺโกสิมฺหา, อนภินนฺทิตฺวา อปฺปฏิกฺโกสิตฺวา อุฏฺฐายาสนา ปกฺกมิมฺหา “ภควโต สนฺติเก เอตสฺส ภาสิตสฺส อตฺถํ อาชานิสฺสมา”ติ.}

\prose{เอวํวาทิโน ภิกฺขเว อญฺญติตฺถิยา ปริพฺพาชกา เอวมสฺสุ วจนียา “เอโก อาวุโส ปโญฺห เอโก อุทฺเทโส เอกํ เวยฺยากรณํ, เทฺว ปญฺหา เทฺว อุทฺเทสา เทฺว เวยฺยากรณานิ, ตโย ปญฺหา ตโย อุทฺเทสา ตีณิ เวยฺยากรณานิ, จตฺตาโร ปญฺหา จตฺตาโร อุทฺเทสา จตฺตาริ เวยฺยากรณานิ, ปญฺจ ปญฺหา ปญฺจุทฺเทสา ปญฺจ เวยฺยากรณานิ,}

\vspace{\tipitakaparskip}
\csromancenter{มหาวคฺค ฉ ปญฺหา ฉ อุทฺเทสา ฉ เวยฺยากรณานิ, สตฺต ปญฺหา สตฺตุทฺเทสา สตฺต เวยฺยากรณานิ, อฏฺฐ ปญฺหา อฏฺฐุทฺเทสา อฏฺฐ เวยฺยากรณานิ, นว ปญฺหา นวุทฺเทสา นว เวยฺยากรณานิ, ทส ปญฺหา ทสุทฺเทสา ทส เวยฺยากรณานี”ติ.\csromanspacer{} เอวํ ปุฏฺฐา ภิกฺขเว อญฺญติตฺถิยา ปริพฺพาชกา น เจว สมฺมายิสฺสนฺติ, อุตฺตริ จ วิฆาตํ อาปชฺชิสฺสนฺติ.\csromanspacer{} ตํ กิสฺส เหตุ? ยถา ตํ ภิกฺขเว อวิสยสฺมิํ.\csromanspacer{} นาหํ ตํ ภิกฺขเว ปสฺสามิ สเทวเก โลเก สมารเก สพฺรหฺมเก สสฺสมณพฺราหฺมณิยา ปชาย สเทวมนุสฺสาย, โย อิเมสํ ปญฺหานํ เวยฺยากรเณน จิตฺตํ อาราเธยฺย อญฺญตฺร ตถาคเตน วา ตถาคตสาวเกน วา, อิโต วา ปน สุตฺวา.}
\prose{\csromanfolio{297}“เอโก ปโญฺห เอโก อุทฺเทโส เอกํ เวยฺยากรณนฺ”ติ อิติ โข ปเนตํ วุตฺตํ, กิญฺเจตํ ปฏิจฺจ วุตฺตํ.\csromanspacer{} เอกธมฺเม ภิกฺขเว ภิกฺขุ สมฺมา นิพฺพินฺทมาโน สมฺมา วิรชฺชมาโน สมฺมา วิมุจฺจมาโน สมฺมา ปริยนฺตทสฺสาวี สมฺมทตฺถํ อภิสเมจฺจ ทิฏฺเฐว ธมฺเม ทุกฺขสฺสนฺตกโร โหติ.\csromanspacer{} กตมสฺมิํ เอกธมฺเม? สพฺเพ สตฺตา อาหารฏฺฐิติกา.\csromanspacer{} อิมสฺมิํ โข ภิกฺขเว เอกธมฺเม ภิกฺขุ สมฺมา นิพฺพินฺทมาโน สมฺมา วิรชฺชมาโน สมฺมา วิมุจฺจมาโน สมฺมา ปริยนฺตทสฺสาวี สมฺมทตฺถํ อภิสเมจฺจ ทิฏฺเฐว ธมฺเม ทุกฺขสฺสนฺตกโร โหติ. “เอโก ปโญฺห เอโก อุทฺเทโส เอกํ เวยฺยากรณนฺ”ติ อิติ ยํ ตํ วุตฺตํ, อิทเมตํ ปฏิจฺจ วุตฺตํ.}

\prose{“เทฺว ปญฺหา เทฺว อุทฺเทสา เทฺว เวยฺยากรณานี”ติ อิติ โข ปเนตํ วุตฺตํ, กิญฺเจตํ ปฏิจฺจ วุตฺตํ.\csromanspacer{} ทฺวีสุ ภิกฺขเว ธมฺเมสุ ภิกฺขุ สมฺมา นิพฺพินฺทมาโน สมฺมา วิรชฺชมาโน สมฺมา วิมุจฺจมาโน สมฺมา ปริยนฺตทสฺสาวี สมฺมทตฺถํ อภิสเมจฺจ ทิฏฺเฐว ธมฺเม ทุกฺขสฺสนฺตกโร โหติ.\csromanspacer{} กตเมสุ ทฺวีสุ? นาเม จ รูเป จ.\csromanspacer{} อิเมสุ โข ภิกฺขเว ทฺวีสุ ธมฺเมสุ ภิกฺขุ สมฺมา นิพฺพินฺทมาโน สมฺมา วิรชฺชมาโน สมฺมา วิมุจฺจมาโน สมฺมา ปริยนฺตทสฺสาวี สมฺมทตฺถํ อภิสเมจฺจ ทิฏฺเฐว ธมฺเม ทุกฺขสฺสนฺตกโร โหติ. “เทฺว ปญฺหา เทฺว อุทฺเทสา เทฺว เวยฺยากรณานี”ติ อิติ ยํ ตํ วุตฺตํ, อิทเมตํ ปฏิจฺจ วุตฺตํ.}

\prose{“ตโย ปญฺหา ตโย อุทฺเทสา ตีณิ เวยฺยากรณานี”ติ อิติ โข ปเนตํ วุตฺตํ, กิญฺเจตํ ปฏิจฺจ วุตฺตํ.\csromanspacer{} ตีสุ ภิกฺขเว ธมฺเมสุ ภิกฺขุ สมฺมา นิพฺพินฺทมาโน สมฺมา วิรชฺชมาโน สมฺมา วิมุจฺจมาโน สมฺมา ปริยนฺตทสฺสาวี สมฺมทตฺถํ อภิสเมจฺจ ทิฏฺเฐว ธมฺเม ทุกฺขสฺสนฺตกโร}

\vspace{\tipitakaparskip}
\csromancenter{ทสกนิปาตปาฬิ โหติ.\csromanspacer{} กตเมสุ ตีสุ? ตีสุ เวทนาสุ.\csromanspacer{} อิเมสุ โข ภิกฺขเว ตีสุ ธมฺเมสุ ภิกฺขุ สมฺมา นิพฺพินฺทมาโน สมฺมา วิรชฺชมาโน สมฺมา วิมุจฺจมาโน สมฺมา ปริยนฺตทสฺสาวี สมฺมทตฺถํ อภิสเมจฺจ ทิฏฺเฐว ธมฺเม ทุกฺขสฺสนฺตกโร โหติ. “ตโย ปญฺหา ตโย อุทฺเทสา ตีณิ เวยฺยากรณานี”ติ อิติ ยํ ตํ วุตฺตํ, อิทเมตํ ปฏิจฺจ วุตฺตํ.}
\prose{\csromanfolio{298}“จตฺตาโร ปญฺหา จตฺตาโร อุทฺเทสา จตฺตาริ เวยฺยากรณานี”ติ อิติ โข ปเนตํ วุตฺตํ, กิญฺเจตํ ปฏิจฺจ วุตฺตํ.\csromanspacer{} จตูสุ ภิกฺขเว ธมฺเมสุ ภิกฺขุ สมฺมา นิพฺพินฺทมาโน สมฺมา วิรชฺชมาโน สมฺมา วิมุจฺจมาโน สมฺมา ปริยนฺตทสฺสาวี สมฺมทตฺถํ อภิสเมจฺจ ทิฏฺเฐว ธมฺเม ทุกฺขสฺสนฺตกโร โหติ.\csromanspacer{} กตเมสุ จตูสุ? จตูสุ อาหาเรสุ.\csromanspacer{} อิเมสุ โข ภิกฺขเว จตูสุ ธมฺเมสุ ภิกฺขุ สมฺมา นิพฺพินฺทมาโน สมฺมา วิรชฺชมาโน สมฺมา วิมุจฺจมาโน สมฺมา ปริยนฺตทสฺสาวี สมฺมทตฺถํ อภิสเมจฺจ ทิฏฺเฐว ธมฺเม ทุกฺขสฺสนฺตกโร โหติ. “จตฺตาโร ปญฺหา จตฺตาโร อุทฺเทสา จตฺตาริ เวยฺยากรณานี”ติ อิติ ยํ ตํ วุตฺตํ, อิทเมตํ ปฏิจฺจ วุตฺตํ.}

\prose{“ปญฺจ ปญฺหา ปญฺจุทฺเทสา ปญฺจ เวยฺยากรณานี”ติ อิติ โข ปเนตํ วุตฺตํ, กิญฺเจตํ ปฏิจฺจ วุตฺตํ.\csromanspacer{} ปญฺจสุ ภิกฺขเว ธมฺเมสุ ภิกฺขุ สมฺมา นิพฺพินฺทมาโน สมฺมา วิรชฺชมาโน สมฺมา วิมุจฺจมาโน สมฺมา ปริยนฺตทสฺสาวี สมฺมทตฺถํ อภิสเมจฺจ ทิฏฺเฐว ธมฺเม ทุกฺขสฺสนฺตกโร โหติ.\csromanspacer{} กตเมสุ ปญฺจสุ? ปญฺจสุ อุปาทานกฺขนฺเธสุ.\csromanspacer{} อิเมสุ โข ภิกฺขเว ปญฺจสุ ธมฺเมสุ ภิกฺขุ สมฺมา นิพฺพินฺทมาโน สมฺมา วิรชฺชมาโน สมฺมา วิมุจฺจมาโน สมฺมา ปริยนฺตทสฺสาวี สมฺมทตฺถํ อภิสเมจฺจ ทิฏฺเฐว ธมฺเม ทุกฺขสฺสนฺตกโร โหติ. “ปญฺจ ปญฺหา ปญฺจุทฺเทสา ปญฺจ เวยฺยากรณานี”ติ อิติ ยํ ตํ วุตฺตํ, อิทเมตํ ปฏิจฺจ วุตฺตํ.}

\prose{“ฉ ปญฺหา ฉ อุทฺเทสา ฉ เวยฺยากรณานี”ติ อิติ โข ปเนตํ วุตฺตํ, กิญฺเจตํ ปฏิจฺจ วุตฺตํ.\csromanspacer{} ฉสุ ภิกฺขเว ธมฺเมสุ ภิกฺขุ สมฺมา นิพฺพินฺทมาโน สมฺมา วิรชฺชมาโน สมฺมา วิมุจฺจมาโน สมฺมา ปริยนฺตทสฺสาวี สมฺมทตฺถํ อภิสเมจฺจ ทิฏฺเฐว ธมฺเม ทุกฺขสฺสนฺตกโร โหติ.\csromanspacer{} กตเมสุ ฉสุ? ฉสุ อชฺฌตฺติเกสุ อายตเนสุ.\csromanspacer{} อิเมสุ โข ภิกฺขเว ฉสุ ธมฺเมสุ ภิกฺขุ สมฺมา นิพฺพินฺทมาโน สมฺมา วิรชฺชมาโน สมฺมา วิมุจฺจมาโน สมฺมา ปริยนฺตทสฺสาวี สมฺมทตฺถํ อภิสเมจฺจ ทิฏฺเฐว ธมฺเม ทุกฺขสฺสนฺตกโร โหติ. “ฉ ปญฺหา ฉ อุทฺเทสา ฉ เวยฺยากรณานี”ติ อิติ ยํ ตํ วุตฺตํ, อิทเมตํ ปฏิจฺจ วุตฺตํ.}

\chapterheadpageread{3. มหาวคฺค}
\prose{\csromanfolio{299}“สตฺต ปญฺหา สตฺตุทฺเทสา สตฺต เวยฺยากรณานี”ติ อิติ โข ปเนตํ วุตฺตํ, กิญฺเจตํ ปฏิจฺจ วุตฺตํ.\csromanspacer{} สตฺตสุ ภิกฺขเว ธมฺเมสุ ภิกฺขุ สมฺมา นิพฺพินฺทมาโน สมฺมา วิรชฺชมาโน สมฺมา วิมุจฺจมาโน สมฺมา ปริยนฺตทสฺสาวี สมฺมทตฺถํ อภิสเมจฺจ ทิฏฺเฐว ธมฺเม ทุกฺขสฺสนฺตกโร โหติ.\csromanspacer{} กตเมสุ สตฺตสุ? สตฺตสุ วิญฺญาณฏฺฐิตีสุ.\csromanspacer{} อิเมสุ โข ภิกฺขเว สตฺตสุ ธมฺเมสุ ภิกฺขุ สมฺมา นิพฺพินฺทมาโน สมฺมา วิรชฺชมาโน สมฺมา วิมุจฺจมาโน สมฺมา ปริยนฺตทสฺสาวี สมฺมทตฺถํ อภิสเมจฺจ ทิฏฺเฐว ธมฺเม ทุกฺขสฺสนฺตกโร โหติ. “สตฺต ปญฺหา สตฺตุทฺเทสา สตฺต เวยฺยากรณานี”ติ อิติ ยํ ตํ วุตฺตํ, อิทเมตํ ปฏิจฺจ วุตฺตํ.}

\prose{“อฏฺฐ ปญฺหา อฏฺฐุทฺเทสา อฏฺฐ เวยฺยากรณานี”ติ อิติ โข ปเนตํ วุตฺตํ, กิญฺเจตํ ปฏิจฺจ วุตฺตํ.\csromanspacer{} อฏฺฐสุ ภิกฺขเว ธมฺเมสุ ภิกฺขุ สมฺมา นิพฺพินฺทมาโน สมฺมา วิรชฺชมาโน สมฺมา วิมุจฺจมาโน สมฺมา ปริยนฺตทสฺสาวี สมฺมทตฺถํ อภิสเมจฺจ ทิฏฺเฐว ธมฺเม ทุกฺขสฺสนฺตกโร โหติ.\csromanspacer{} กตเมสุ อฏฺฐสุ? อฏฺฐสุ โลกธมฺเมสุ.\csromanspacer{} อิเมสุ โข ภิกฺขเว อฏฺฐสุ ธมฺเมสุ ภิกฺขุ สมฺมา นิพฺพินฺทมาโน ฯเปฯ ทุกฺขสฺสนฺตกโร โหติ. “อฏฺฐ ปญฺหา อฏฺฐุทฺเทสา อฏฺฐ เวยฺยากรณานี”ติ อิติ ยํ ตํ วุตฺตํ, อิทเมตํ ปฏิจฺจ วุตฺตํ.}

\prose{“นว ปญฺหา นวุทฺเทสา นว เวยฺยากรณานี”ติ อิติ โข ปเนตํ วุตฺตํ, กิญฺเจตํ ปฏิจฺจ วุตฺตํ.\csromanspacer{} นวสุ ภิกฺขเว ธมฺเมสุ ภิกฺขุ สมฺมา นิพฺพินฺทมาโน สมฺมา วิรชฺชมาโน สมฺมา วิมุจฺจมาโน สมฺมา ปริยนฺตทสฺสาวี สมฺมทตฺถํ อภิสเมจฺจ ทิฏฺเฐว ธมฺเม ทุกฺขสฺสนฺตกโร โหติ.\csromanspacer{} กตเมสุ นวสุ? นวสุ สตฺตาวาเสสุ.\csromanspacer{} อิเมสุ โข ภิกฺขเว นวสุ ธมฺเมสุ ภิกฺขุ สมฺมา นิพฺพินฺทมาโน สมฺมา วิรชฺชมาโน สมฺมา วิมุจฺจมาโน สมฺมา ปริยนฺตทสฺสาวี สมฺมทตฺถํ อภิสเมจฺจ ทิฏฺเฐว ธมฺเม ทุกฺขสฺสนฺตกโร โหติ. “นว ปญฺหา นวุทฺเทสา นว เวยฺยากรณานี”ติ อิติ ยํ ตํ วุตฺตํ, อิทเมตํ ปฏิจฺจ วุตฺตํ.}

\prose{“ทส ปญฺหา ทสุทฺเทสา ทส เวยฺยากรณานี”ติ อิติ โข ปเนตํ วุตฺตํ, กิญฺเจตํ ปฏิจฺจ วุตฺตํ.\csromanspacer{} ทสสุ ภิกฺขเว ธมฺเมสุ ภิกฺขุ สมฺมา นิพฺพินฺทมาโน สมฺมา วิรชฺชมาโน สมฺมา วิมุจฺจมาโน สมฺมา ปริยนฺตทสฺสาวี สมฺมทตฺถํ อภิสเมจฺจ ทิฏฺเฐว ธมฺเม ทุกฺขสฺสนฺตกโร โหติ.\csromanspacer{} กตเมสุ ทสสุ? ทสสุ อกุสเลสุ กมฺมปเถสุ.\csromanspacer{} อิเมสุ โข ภิกฺขเว ทสสุ ธมฺเมสุ ภิกฺขุ สมฺมา นิพฺพินฺทมาโน สมฺมา วิรชฺชมาโน สมฺมา วิมุจฺจมาโน สมฺมา ปริยนฺตทสฺสาวี สมฺมทตฺถํ อภิสเมจฺจ ทิฏฺเฐว ธมฺเม ทุกฺขสฺสนฺตกโร}

\vspace{\tipitakaparskip}
\csromancenter{ทสกนิปาตปาฬิ โหติ. “ทส ปญฺหา ทสุทฺเทสา ทส เวยฺยากรณานี”ติ อิติ ยํ ตํ วุตฺตํ, อิทเมตํ ปฏิจฺจ วุตฺตนฺติ.\csromanspacer{}\csromanspacer{}สตฺตมํ.}
\csromanheader{2}{8. ทุติยมหาปญฺหาสุตฺต}
\proseitem{28}{\csromanfolio{300}เอกํ สมยํ ภควา กชงฺคลายํ วิหรติ เวฬุวเน.\csromanspacer{} อถ โข สมฺพหุลา กชงฺคลกา อุปาสกา เยน กชงฺคลิกา ภิกฺขุนี เตนุปสงฺกมิํสุ, อุปสงฺกมิตฺวา กชงฺคลิกํ ภิกฺขุนิํ อภิวาเทตฺวา เอกมนฺตํ นิสีทิํสุ, เอกมนฺตํ นิสินฺนา โข กชงฺคลกา อุปาสกา กชงฺคลิกํ ภิกฺขุนิํ เอตทโวจุํ\csromandash{}}

\prose{วุตฺตมิทํ อยฺเย ภควตา มหาปเญฺหสุ “เอโก ปโญฺห เอโก อุทฺเทโส เอกํ เวยฺยากรณํ, เทฺว ปญฺหา เทฺว อุทฺเทสา เทฺว เวยฺยากรณานิ, ตโย ปญฺหา ตโย อุทฺเทสา ตีณิ เวยฺยากรณานิ, จตฺตาโร ปญฺหา จตฺตาโร อุทฺเทสา จตฺตาริ เวยฺยากรณานิ, ปญฺจ ปญฺหา ปญฺจุทฺเทสา ปญฺจ เวยฺยากรณานิ, ฉ ปญฺหา ฉ อุทฺเทสา ฉ เวยฺยากรณานิ, สตฺต ปญฺหา สตฺตุทฺเทสา สตฺต เวยฺยากรณานิ, อฏฺฐ ปญฺหา อฏฺฐุทฺเทสา อฏฺฐ เวยฺยากรณานิ, นว ปญฺหา นวุทฺเทสา นว เวยฺยากรณานิ, ทส ปญฺหา ทสุทฺเทสา ทส เวยฺยากรณานี”ติ.\csromanspacer{} อิมสฺส นุ โข อยฺเย ภควตา สํขิตฺเตน ภาสิตสฺส กถํ วิตฺตาเรน อตฺโถ ทฏฺฐพฺโพติ.}

\prose{น โข ปเนตํ อาวุโส ภควโต สมฺมุขา สุตํ สมฺมุขา ปฏิคฺคหิตํ, นปิ มโนภาวนียานํ ภิกฺขูนํ สมฺมุขา สุตํ สมฺมุขา ปฏิคฺคหิตํ.\csromanspacer{} อปิ จ ยถา เมตฺถ ขายติ, ตํ สุณาถ สาธุกํ มนสิ กโรถ, ภาสิสฺสามีติ. “เอวํ อยฺเย”ติ โข กชงฺคลกา อุปาสกา กชงฺคลิกาย ภิกฺขุนิยา ปจฺจสฺโสสุํ.\csromanspacer{} กชงฺคลิกา ภิกฺขุนี เอตทโวจ\csromandash{}}

\prose{“เอโก ปโญฺห เอโก อุทฺเทโส เอกํ เวยฺยากรณนฺ”ติ อิติ โข ปเนตํ วุตฺตํ ภควตา, กิญฺเจตํ ปฏิจฺจ วุตฺตํ.\csromanspacer{} เอกธมฺเม อาวุโส ภิกฺขุ สมฺมา นิพฺพินฺทมาโน สมฺมา วิรชฺชมาโน สมฺมา วิมุจฺจมาโน สมฺมา ปริยนฺตทสฺสาวี สมฺมทตฺถํ อภิสเมจฺจ ทิฏฺเฐว ธมฺเม ทุกฺขสฺสนฺตกโร โหติ.\csromanspacer{} กตมสฺมิํ เอกธมฺเม? สพฺเพ สตฺตา อาหารฏฺฐิติกา.\csromanspacer{} อิมสฺมิํ โข อาวุโส เอกธมฺเม ภิกฺขุ สมฺมา นิพฺพินฺทมาโน สมฺมา วิรชฺชมาโน สมฺมา วิมุจฺจมาโน สมฺมา ปริยนฺตทสฺสาวี สมฺมทตฺถํ อภิสเมจฺจ ทิฏฺเฐว}

\vspace{\tipitakaparskip}
\csromancenter{มหาวคฺค ธมฺเม ทุกฺขสฺสนฺตกโร โหติ. “เอโก ปโญฺห เอโก อุทฺเทโส เอกํ เวยฺยากรณนฺ”ติ อิติ ยํ ตํ วุตฺตํ ภควตา, อิทเมตํ ปฏิจฺจ วุตฺตํ.}
\prose{\csromanfolio{301}“เทฺว ปญฺหา เทฺว อุทฺเทสา เทฺว เวยฺยากรณานี”ติ อิติ โข ปเนตํ วุตฺตํ ภควตา, กิญฺเจตํ ปฏิจฺจ วุตฺตํ.\csromanspacer{} ทฺวีสุ อาวุโส ธมฺเมสุ ภิกฺขุ สมฺมา นิพฺพินฺทมาโน สมฺมา วิรชฺชมาโน สมฺมา วิมุจฺจมาโน สมฺมา ปริยนฺตทสฺสาวี สมฺมทตฺถํ อภิสเมจฺจ ทิฏฺเฐว ธมฺเม ทุกฺขสฺสนฺตกโร โหติ.\csromanspacer{} กตเมสุ ทฺวีสุ? นาเม จ รูเป จ ฯเปฯ.\csromanspacer{} กตเมสุ ตีสุ? ตีสุ เวทนาสุ.\csromanspacer{} อิเมสุ โข อาวุโส ตีสุ ธมฺเมสุ ภิกฺขุ สมฺมา นิพฺพินฺทมาโน สมฺมา วิรชฺชมาโน สมฺมา วิมุจฺจมาโน สมฺมา ปริยนฺตทสฺสาวี สมฺมทตฺถํ อภิสเมจฺจ ทิฏฺเฐว ธมฺเม ทุกฺขสฺสนฺตกโร โหติ. “ตโย ปญฺหา ตโย อุทฺเทสา ตีณิ เวยฺยากรณานี”ติ อิติ ยํ ตํ วุตฺตํ ภควตา, อิทเมตํ ปฏิจฺจ วุตฺตํ.}

\prose{“จตฺตาโร ปญฺหา จตฺตาโร อุทฺเทสา จตฺตาริ เวยฺยากรณานี”ติ อิติ โข ปเนตํ วุตฺตํ ภควตา, กิญฺเจตํ ปฏิจฺจ วุตฺตํ.\csromanspacer{} จตูสุ อาวุโส ธมฺเมสุ ภิกฺขุ สมฺมา สุภาวิตจิตฺโต สมฺมา ปริยนฺตทสฺสาวี สมฺมทตฺถํ อภิสเมจฺจ ทิฏฺเฐว ธมฺเม ทุกฺขสฺสนฺตกโร โหติ.\csromanspacer{} กตเมสุ จตูสุ? จตูสุ สติปฏฺฐาเนสุ.\csromanspacer{} อิเมสุ โข อาวุโส จตูสุ ธมฺเมสุ ภิกฺขุ สมฺมา สุภาวิตจิตฺโต สมฺมา ปริยนฺตทสฺสาวี สมฺมทตฺถํ อภิสเมจฺจ ทิฏฺเฐว ธมฺเม ทุกฺขสฺสนฺตกโร โหติ. “จตฺตาโร ปญฺหา จตฺตาโร อุทฺเทสา จตฺตาริ เวยฺยากรณานี”ติ อิติ ยํ ตํ วุตฺตํ ภควตา, อิทเมตํ ปฏิจฺจ วุตฺตํ.}

\prose{“ปญฺจ ปญฺหา ปญฺจุทฺเทสา ปญฺจ เวยฺยากรณานี”ติ อิติ โข ปเนตํ วุตฺตํ ภควตา, กิญฺเจตํ ปฏิจฺจ วุตฺตํ.\csromanspacer{} ปญฺจสุ อาวุโส ธมฺเมสุ ภิกฺขุ สมฺมา สุภาวิตจิตฺโต สมฺมา ปริยนฺตทสฺสาวี สมฺมทตฺถํ อภิสเมจฺจ ทิฏฺเฐว ธมฺเม ทุกฺขสฺสนฺตกโร โหติ.\csromanspacer{} กตเมสุ ปญฺจสุ? ปญฺจสุ อินฺทฺริเยสุ ฯเปฯ.\csromanspacer{} กตเมสุ ฉสุ? ฉสุ นิสฺสรณียาสุ ธาตูสุ ฯเปฯ.\csromanspacer{} กตเมสุ สตฺตสุ? สตฺตสุ โพชฺฌงฺเคสุ ฯเปฯ.\csromanspacer{} กตเมสุ อฏฺฐสุ? อฏฺฐสุ อริย-อฏฺฐงฺคิกมคฺเคสุ.\csromanspacer{} อิเมสุ โข อาวุโส อฏฺฐสุ ธมฺเมสุ ภิกฺขุ สมฺมา สุภาวิตจิตฺโต สมฺมา ปริยนฺตทสฺสาวี สมฺมทตฺถํ อภิสเมจฺจ ทิฏฺเฐว ธมฺเม ทุกฺขสฺสนฺตกโร โหติ. “อฏฺฐ ปญฺหา อฏฺฐุทฺเทสา อฏฺฐ เวยฺยากรณานี”ติ อิติ ยํ ตํ วุตฺตํ ภควตา, อิทเมตํ ปฏิจฺจ วุตฺตํ.}

\csromanmatikaanchor{mtk.54}
\csromantocmarkref{subsection}{8. ทุติยมหาปญฺหาสุตฺต}{mtk.54}
\bookmark[dest={mtk.54},level=2]{8. ทุติยมหาปญฺหาสุตฺต}
\gambhira{ทสกนิปาตปาฬิ}
\prose{\csromanfolio{302}“นว ปญฺหา นวุทฺเทสา นวเวยฺยากรณานี”ติ อิติ โข ปเนตํ วุตฺตํ ภควตา, กิญฺเจตํ ปฏิจฺจ วุตฺตํ.\csromanspacer{} นวสุ อาวุโส ธมฺเมสุ ภิกฺขุ สมฺมา นิพฺพินฺทมาโน สมฺมา วิรชฺชมาโน สมฺมา วิมุจฺจมาโน สมฺมา ปริยนฺตทสฺสาวี สมฺมทตฺถํ อภิสเมจฺจ ทิฏฺเฐว ธมฺเม ทุกฺขสฺสนฺตกโร โหติ.\csromanspacer{} กตเมสุ นวสุ? นวสุ สตฺตาวาเสสุ.\csromanspacer{} อิเมสุ โข อาวุโส นวสุ ธมฺเมสุ ภิกฺขุ สมฺมา นิพฺพินฺทมาโน สมฺมา วิรชฺชมาโน สมฺมา วิมุจฺจมาโน สมฺมา ปริยนฺตทสฺสาวี สมฺมทตฺถํ อภิสเมจฺจ ทิฏฺเฐว ธมฺเม ทุกฺขสฺสนฺตกโร โหติ. “นว ปญฺหา นวุทฺเทสา นว เวยฺยากรณานี”ติ อิติ ยํ ตํ วุตฺตํ ภควตา อิทเมตํ ปฏิจฺจ วุตฺตํ.}

\prose{“ทส ปญฺหา ทสุทฺเทสา ทส เวยฺยากรณานี”ติ อิติ โข ปเนตํ วุตฺตํ ภควตา, กิญฺเจตํ ปฏิจฺจ วุตฺตํ.\csromanspacer{} ทสสุ อาวุโส ธมฺเมสุ ภิกฺขุ สมฺมา สุภาวิตจิตฺโต สมฺมา ปริยนฺตทสฺสาวี สมฺมทตฺถํ อภิสเมจฺจ ทิฏฺเฐว ธมฺเม ทุกฺขสฺสนฺตกโร โหติ.\csromanspacer{} กตเมสุ ทสสุ? ทสสุ กุสเลสุ กมฺมปเถสุ.\csromanspacer{} อิเมสุ โข อาวุโส ทสสุ ธมฺเมสุ ภิกฺขุ สมฺมา สุภาวิตจิตฺโต สมฺมา ปริยนฺตทสฺสาวี สมฺมทตฺถํ อภิสเมจฺจ ทิฏฺเฐว ธมฺเม ทุกฺขสฺสนฺตกโร โหติ. “ทส ปญฺหา ทสุทฺเทสา ทส เวยฺยากรณานี”ติ อิติ ยํ ตํ วุตฺตํ ภควตา, อิทเมตํ ปฏิจฺจ วุตฺตํ.}

\prose{อิติ โข อาวุโส ยํ ตํ วุตฺตํ ภควตา สํขิตฺเตน ภาสิตาสุ มหาปญฺหาสุ “เอโก ปโญฺห เอโก อุทฺเทโส เอกํ เวยฺยากรณํ ฯเปฯ ทส ปญฺหา ทสุทฺเทสา ทส เวยฺยากรณานี”ติ.\csromanspacer{} อิมสฺส โข อหํ อาวุโส ภควตา สํขิตฺเตน ภาสิตสฺส เอวํ วิตฺถาเรน อตฺถํ อาชานามิ.\csromanspacer{} อากงฺขมานา จ ปน ตุเมฺห อาวุโส ภควนฺตญฺเญว อุปสงฺกมิตฺวา เอตมตฺถํ ปฏิปุจฺเฉยฺยาถ.\csromanspacer{} ยถา โว\footnote{ยถา โข (ก), ยถา โน (พหูสุ) อุปริ 441-2, 466-8 ปิฏฺเฐสุ ปน ปาฐเภโท นตฺถิ.} ภควา พฺยากโรติ, ตถา นํ ธาเรยฺยาถาติ. “เอวํ อยฺเย”ติ โข กชงฺคลกา อุปาสกา กชงฺคลิกาย โข ภิกฺขุนิยา ภาสิตํ อภินนฺทิตฺวา อนุโมทิตฺวา อุฏฺฐายาสนา กชงฺคลิกํ ภิกฺขุนิํ อภิวาเทตฺวา ปทกฺขิณํ กตฺวา เยน ภควา เตนุปสงฺกมิํสุ, อุปสงฺกมิตฺวา ภควนฺตํ อภิวาเทตฺวา เอกมนฺตํ นิสีทิํสุ, เอกมนฺตํ นิสินฺนา โข กชงฺคลกา อุปาสกา ยาวตโก อโหสิ กชงฺคลิกาย ภิกฺขุนิยา สทฺธิํ กถาสลฺลาโป, ตํ สพฺพํ ภควโต อาโรเจสุํ.}

\chapterheadpageread{3. มหาวคฺค}
\prose{\csromanfolio{303}สาธุ สาธุ คหปตโย, ปณฺฑิตา คหปตโย กชงฺคลิกา ภิกฺขุนี, มหาปญฺญา คหปตโย กชงฺคลิกา ภิกฺขุนี, มญฺเจปิ ตุเมฺห คหปตโย อุปสงฺกมิตฺวา เอตมตฺถํ ปฏิปุจฺเฉยฺยาถ.\csromanspacer{} อหมฺปิ เจตํ เอวเมวํ\footnote{เอวเมว (ก) ม 1. 161 ปิฏฺฐาทีสุ ปสฺสิตพฺพํ.} พฺยากเรยฺยํ, ยถา ตํ กชงฺคลิกาย ภิกฺขุนิยา พฺยากตํ.\csromanspacer{} เอโส เจว ตสฺส\footnote{เอโส เจเวตสฺส (ม 161 ปิฏฺฐาทีสุ.)} อตฺโถ, เอวญฺจ นํ ธาเรยฺยาถาติ.\csromanspacer{}\csromanspacer{}อฏฺฐมํ.}

\csromanheader{2}{9. ปฐมโกสลสุตฺต}
\proseitem{29}{ยาวตา ภิกฺขเว กาสิโกสลา, ยาวตา รญฺโญ ปเสนทิสฺส โกสลสฺส วิชิตํ\footnote{วิชิเต (สี, ก)}, ราชา ตตฺถ ปเสนทิ โกสโล อคฺคมกฺขายติ.\csromanspacer{} รญฺโญปิ โข ภิกฺขเว ปเสนทิสฺส โกสลสฺส อตฺเถว อญฺญถตฺตํ, อตฺถิ วิปริณาโม.\csromanspacer{} เอวํ ปสฺสํ ภิกฺขเว สุตวา อริยสาวโก ตสฺมิมฺปิ นิพฺพินฺทติ, ตสฺมิํ นิพฺพินฺทนฺโต อคฺเค วิรชฺชติ, ปเคว หีนสฺมิํ. (1)}

\prose{ยาวตา ภิกฺขเว จนฺทิมสูริยา ปริหรนฺติ ทิสา ภนฺติ วิโรจมานา, ตาว สหสฺสธา โลโก.\csromanspacer{} ตสฺมิํ สหสฺสธา โลเก สหสฺสํ จนฺทานํ สหสฺสํ สูริยานํ\footnote{สุริยานํ (สี, สฺยา, กํ, อิ)} สหสฺสํ สิเนรุปพฺพตราชานํ สหสฺสํ ชมฺพุทีปานํ สหสฺสํ อปรโคยานานํ สหสฺสํ อุตฺตรกุรูนํ สหสฺสํ ปุพฺพวิเทหานํ จตฺตาริ มหาสมุทฺทสหสฺสานิ จตฺตาริ มหาราชสหสฺสานิ สหสฺสํ จาตุมหาราชิกานํ สหสฺสํ ตาวติํสานํ สหสฺสํ ยามานํ สหสฺสํ ตุสิตานํ สหสฺสํ นิมฺมานรตีนํ สหสฺสํ ปรนิมฺมิตวสวตฺตีนํ สหสฺสํ พฺรหฺมโลกานํ.\csromanspacer{} ยาวตา ภิกฺขเว สหสฺสี โลกธาตุ, มหาพฺรหฺมา ตตฺถ อคฺคมกฺขายติ.\csromanspacer{} มหาพฺรหฺมุโนปิ โข ภิกฺขเว อตฺเถว อญฺญถตฺตํ, อตฺถิ วิปริณาโม.\csromanspacer{} เอวํ ปสฺสํ ภิกฺขเว สุตวา อริยสาวโก ตสฺมิมฺปิ นิพฺพินฺทติ, ตสฺมิํ นิพฺพินฺทนฺโต อคฺเค วิรชฺชติ, ปเคว หีนสฺมิํ. (2)}

\prose{โหติ โส ภิกฺขเว สมโย, ยํ อยํ โลโก สํวฏฺฏติ, สํวฏฺฏมาเน ภิกฺขเว โลเก เยภุยฺเยน สตฺตา อาภสฺสรสํวตฺตนิกา\footnote{อาภสฺสรวตฺตนิกา (สี, สฺยา)} ภวนฺติ.\csromanspacer{} เต ตตฺถ โหนฺติ มโนมยา ปีติภกฺขา สยํปภา อนฺตลิกฺเขจรา สุภฏฺฐายิโน จิรํ ทีฆมทฺธานํ ติฏฺฐนฺติ.\csromanspacer{} สํวฏฺฏมาเน ภิกฺขเว โลเก}

\vspace{\tipitakaparskip}
\csromancenter{ทสกนิปาตปาฬิ อาภสฺสรา เทวา อคฺคมกฺขายนฺติ, อาภสฺสรานมฺปิ โข ภิกฺขเว เทวานํ อตฺเถว อญฺญถตฺตํ, อตฺถิ วิปริณาโม.\csromanspacer{} เอวํ ปสฺสํ ภิกฺขเว สุตวา อริยสาวโก ตสฺมิมฺปิ นิพฺพินฺทติ, ตสฺมิํ นิพฺพินฺทนฺโต อคฺเค วิรชฺชติ, ปเคว หีนสฺมิํ. (3)}
\prose{\csromanfolio{304}\csromansymbolfootnote{+}{เหฏฺฐา 293 ปิฏฺเฐปิ.}ทสยิมานิ ภิกฺขเว กสิณายตนานิ.\csromanspacer{} กตมานิ ทส? ปถวีกสิณเมโก สญฺชานาติ อุทฺธํ อโธ ติริยํ อทฺวยํ อปฺปมาณํ.\csromanspacer{} อาโปกสิณเมโก สญฺชานาติ ฯเปฯ.\csromanspacer{} เตโชกสิณเมโก สญฺชานาติ.\csromanspacer{} วาโยกสิณเมโก สญฺชานาติ.\csromanspacer{} นีลกสิณเมโก สญฺชานาติ.\csromanspacer{} ปีตกสิณเมโก สญฺชานาติ.\csromanspacer{} โลหิตกสิณเมโก สญฺชานาติ.\csromanspacer{} โอทาตกสิณเมโก สญฺชานาติ.\csromanspacer{} อากาสกสิณเมโก สญฺชานาติ.\csromanspacer{} วิญฺญาณกสิณเมโก สญฺชานาติ.\csromanspacer{} อุทฺธํ อโธ ติริยํ อทฺวยํ อปฺปมาณํ.\csromanspacer{} อิมานิ โข ภิกฺขเว ทส กสิณายตนานิ.}

\prose{เอตทคฺคํ ภิกฺขเว อิเมสํ ทสนฺนํ กสิณายตนานํ, ยทิทํ วิญฺญาณกสิณํ เอโก สญฺชานาติ อุทฺธํ อโธ ติริยํ อทฺวยํ อปฺปมาณํ.\csromanspacer{} เอวํสญฺญิโนปิ โข ภิกฺขเว สนฺติ สตฺตา.\csromanspacer{} เอวํสญฺญีนมฺปิ โข ภิกฺขเว สตฺตานํ อตฺเถว อญฺญถตฺตํ, อตฺถิ วิปริณาโม.\csromanspacer{} เอวํ ปสฺสํ ภิกฺขเว สุตวา อริยสาวโก ตสฺมิมฺปิ นิพฺพินฺทติ, ตสฺมิํ นิพฺพินฺทนฺโต อคฺเค วิรชฺชติ, ปเคว หีนสฺมิํ. (4)}

\prose{\csromansymbolfootnote{*}{ที 3. 215, 250, เหฏฺฐา 125 ปิฏฺเฐสุปิ.}อฏฺฐิมานิ ภิกฺขเว อภิภายตนานิ.\csromanspacer{} กตมานิ อฏฺฐ? อชฺฌตฺตํ รูปสญฺญี เอโก พหิทฺธา รูปานิ ปสฺสติ ปริตฺตานิ สุวณฺณทุพฺพณฺณานิ, ตานิ อภิภุยฺย “ชานามิ ปสฺสามี”ติ เอวํสญฺญี โหติ.\csromanspacer{} อิทํ ปฐมํ อภิภายตนํ.}

\prose{อชฺฌตฺตํ รูปสญฺญี เอโก พหิทฺธา รูปานิ ปสฺสติ อปฺปมาณานิ สุวณฺณทุพฺพณฺณานิ, ตานิ อภิภุยฺย “ชานามิ ปสฺสามี”ติ เอวํสญฺญี โหติ.\csromanspacer{} อิทํ ทุติยํ อภิภายตนํ.}

\csromanmatikaanchor{mtk.55}
\csromantocmarkref{subsection}{9. ปฐมโกสลสุตฺต}{mtk.55}
\bookmark[dest={mtk.55},level=2]{9. ปฐมโกสลสุตฺต}
\prose{อชฺฌตฺตํ อรูปสญฺญี เอโก พหิทฺธา รูปานิ ปสฺสติ ปริตฺตานิ สุวณฺณทุพฺพณฺณานิ, ตานิ อภิภุยฺย “ชานามิ ปสฺสามี”ติ เอวํสญฺญี โหติ.\csromanspacer{} อิทํ ตติยํ อภิภายตนํ.}

\chapterheadpageread{3. มหาวคฺค}
\prose{\csromanfolio{305}อชฺฌตฺตํ อรูปสญฺญี เอโก พหิทฺธา รูปานิ ปสฺสติ อปฺปมาณานิ สุวณฺณทุพฺพณฺณานิ, ตานิ อภิภุยฺย “ชานามิ ปสฺสามี”ติ เอวํสญฺญี โหติ.\csromanspacer{} อิทํ จตุตฺถํ อภิภายตนํ.}

\prose{อชฺฌตฺตํ อรูปสญฺญี เอโก พหิทฺธา รูปานิ ปสฺสติ นีลานิ นีลวณฺณานิ นีลนิทสฺสนานิ นีลนิภาสานิ.\csromanspacer{} เสยฺยถาปิ นาม อุมาปุปฺผํ นีลํ นีลวณฺณํ นีลนิทสฺสนํ นีลนิภาสํ.\csromanspacer{} เสยฺยถา วา ปน ตํ วตฺถํ พาราณเสยฺยกํ อุภโตภาควิมฏฺฐํ นีลํ นีลวณฺณํ นีลนิทสฺสนํ นีลนิภาสํ.\csromanspacer{} เอวเมวํ อชฺฌตฺตํ อรูปสญฺญี เอโก พหิทฺธา รูปานิ ปสฺสติ นีลานิ นีลวณฺณานิ นีลนิทสฺสนานิ นีลนิภาสานิ, ตานิ อภิภุยฺย “ชานามิ ปสฺสามี”ติ เอวํสญฺญี โหติ.\csromanspacer{} อิทํ ปญฺจมํ อภิภายตนํ.}

\prose{อชฺฌตฺตํ อรูปสญฺญี เอโก พหิทฺธา รูปานิ ปสฺสติ ปีตานิ ปีตวณฺณานิ ปีตนิทสฺสนานิ ปีตนิภาสานิ.\csromanspacer{} เสยฺยถาปิ นาม กณิการปุปฺผํ ปีตํ ปีตวณฺณํ ปีตนิทสฺสนํ ปีตนิภาสํ.\csromanspacer{} เสยฺยถา วา ปน ตํ วตฺถํ พาราณเสยฺยกํ อุภโตภาควิมฏฺฐํ ปีตํ ปีตวณฺณํ ปีตนิทสฺสนํ ปีตนิภาสํ.\csromanspacer{} เอวเมวํ อชฺฌตฺตํ อรูปสญฺญี เอโก พหิทฺธา รูปานิ ปสฺสติ ปีตานิ ปีตวณฺณานิ ปีตนิทสฺสนานิ ปีตนิภาสานิ, ตานิ อภิภุยฺย “ชานามิ ปสฺสามี”ติ เอวํสญฺญี โหติ.\csromanspacer{} อิทํ ฉฏฺฐํ อภิภายตนํ.}

\prose{อชฺฌตฺตํ อรูปสญฺญี เอโก พหิทฺธา รูปานิ ปสฺสติ โลหิตกานิ โลหิตกวณฺณานิ โลหิตกนิทสฺสนานิ โลหิตกนิภาสานิ.\csromanspacer{} เสยฺยถาปิ นาม พนฺธุชีวกปุปฺผํ โลหิตกํ โลหิตกวณฺณํ โลหิตกนิทสฺสนํ โลหิตกนิภาสํ.\csromanspacer{} เสยฺยถา วา ปน ตํ วตฺถํ พาราณเสยฺยกํ อุภโตภาควิมฏฺฐํ โลหิตกํ โลหิตกวณฺณํ โลหิตกนิทสฺสนํ โลหิตกนิภาสํ.\csromanspacer{} เอวเมวํ อชฺฌตฺตํ อรูปสญฺญี เอโก พหิทฺธา รูปานิ ปสฺสติ โลหิตกานิ โลหิตกวณฺณานิ โลหิตกนิทสฺสนานิ โลหิตกนิภาสานิ ตานิ อภิภุยฺย “ชานามิ ปสฺสามี”ติ เอวํสญฺญี โหติ.\csromanspacer{} อิทํ สตฺตมํ อภิภายตนํ.}

\prose{อชฺฌตฺตํ อรูปสญฺญี เอโก พหิทฺธา รูปานิ ปสฺสติ โอทาตานิ โอทาตวณฺณานิ โอทาตนิทสฺสนานิ โอทาตนิภาสานิ.\csromanspacer{} เสยฺยถาปิ นาม โอสธิตารกา โอทาตา โอทาตวณฺณา โอทาตนิทสฺสนา โอทาตนิภาสา.}

\vspace{\tipitakaparskip}
\csromancenter{ทสกนิปาตปาฬิ เสยฺยถา วา ปน ตํ วตฺถํ พาราณเสยฺยกํ อุภโตภาควิมฏฺฐํ โอทาตํ โอทาตวณฺณํ โอทาตนิทสฺสนํ โอทาตนิภาสํ.\csromanspacer{} เอวเมวํ อชฺฌตฺตํ อรูปสญฺญี เอโก พหิทฺธา รูปานิ ปสฺสติ โอทาตานิ โอทาตวณฺณานิ โอทาตนิทสฺสนานิ โอทาตนิภาสานิ, ตานิ อภิภุยฺย “ชานามิ ปสฺสามี”ติ เอวํสญฺญี โหติ.\csromanspacer{} อิทํ อฏฺฐมํ อภิภายตนํ.\csromanspacer{} อิมานิ โข ภิกฺขเว อฏฺฐ อภิภายตนานิ.}
\prose{\csromanfolio{306}เอตทคฺคํ ภิกฺขเว อิเมสํ อฏฺฐนฺนํ อภิภายตนานํ, ยทิทํ อชฺฌตฺตํ อรูปสญฺญี เอโก พหิทฺธา รูปานิ ปสฺสติ โอทาตานิ โอทาตวณฺณานิ โอทาตนิทสฺสนานิ โอทาตนิภาสานิ, ตานิ อภิภุยฺย “ชานามิ ปสฺสามี”ติ เอวํสญฺญี โหติ.\csromanspacer{} เอวํสญฺญิโนปิ โข ภิกฺขเว สนฺติ สตฺตา.\csromanspacer{} เอวํสญฺญีนมฺปิ โข ภิกฺขเว สตฺตานํ อตฺเถว อญฺญถตฺตํ, อตฺถิ วิปริณาโม.\csromanspacer{} เอวํ ปสฺสํ ภิกฺขเว สุตวา อริยสาวโก ตสฺมิมฺปิ นิพฺพินฺทติ, ตสฺมิํ นิพฺพินฺทนฺโต อคฺเค วิรชฺชติ, ปเคว หีนสฺมิํ. (5)}

\prose{จตสฺโส อิมา ภิกฺขเว ปฏิปทา.\csromanspacer{} กตมา จตสฺโส? ทุกฺขา ปฏิปทา ทนฺธาภิญฺญา ทุกฺขา ปฏิปทา ขิปฺปาภิญฺญา สุขา ปฏิปทา ทนฺธาภิญฺญา สุขา ปฏิปทา ขิปฺปาภิญฺญา.\csromanspacer{} อิมา โข ภิกฺขเว จตสฺโส ปฏิปทา.}

\prose{เอตทคฺคํ ภิกฺขเว อิมาสํ จตุนฺนํ ปฏิปทานํ, ยทิทํ สุขา ปฏิปทา ขิปฺปาภิญฺญา.\csromanspacer{} เอวํปฏิปนฺนาปิ โข ภิกฺขเว สนฺติ สตฺตา.\csromanspacer{} เอวํปฏิปนฺนานมฺปิ โข ภิกฺขเว สตฺตานํ อตฺเถว อญฺญถตฺตํ, อตฺถิ วิปริณาโม.\csromanspacer{} เอวํ ปสฺสํ ภิกฺขเว สุตวา อริยสาวโก ตสฺมิมฺปิ นิพฺพินฺทติ, ตสฺมิํ นิพฺพินฺทนฺโต อคฺเค วิรชฺชติ, ปเคว หีนสฺมิํ. (6)}

\prose{จตสฺโส อิมา ภิกฺขเว สญฺญา.\csromanspacer{} กตมา จตสฺโส? ปริตฺตเมโก สญฺชานาติ, มหคฺคตเมโก สญฺชานาติ, อปฺปมาณเมโก สญฺชานาติ, “นตฺถิ กิญฺจี”ติ อากิญฺจญฺญายตนเมโก สญฺชานาติ.\csromanspacer{} อิมา โข ภิกฺขเว จตสฺโส สญฺญา.}

\prose{เอตทคฺคํ ภิกฺขเว อิมาสํ จตุนฺนํ สญฺญานํ, ยทิทํ “นตฺถิ กิญฺจี”ติ อากิญฺจญฺญายตนเมโก สญฺชานาติ.\csromanspacer{} เอวํสญฺญิโนปิ โข ภิกฺขเว สนฺติ สตฺตา.\csromanspacer{} เอวํสญฺญีนมฺปิ โข ภิกฺขเว สตฺตานํ อตฺเถว อญฺญถตฺตํ, อตฺถิ วิปริณาโม.\csromanspacer{} เอวํ ปสฺสํ ภิกฺขเว สุตวา อริยสาวโก ตสฺมิมฺปิ นิพฺพินฺทติ, ตสฺมิํ นิพฺพินฺทนฺโต อคฺเค วิรชฺชติ, ปเคว หีนสฺมิํ. (7)}

\chapterheadpageread{3. มหาวคฺค}
\prose{\csromanfolio{307}เอตทคฺคํ ภิกฺขเว พาหิรกานํ ทิฏฺฐิคตานํ “ยทิทํ โน จสฺสํ โน จ เม สิยา, น ภวิสฺสามิ น เม ภวิสฺสตี”ติ.\csromanspacer{} เอวํทิฏฺฐิโน ภิกฺขเว เอตํ ปาฏิกงฺขํ “ยา จายํ ภเว อปฺปฏิกุลฺยตา, สา จสฺส น ภวิสฺสติ, ยา จายํ ภวนิโรเธ ปาฏิกุลฺยตา, สา จสฺส น ภวิสฺสตี”ติ.\csromanspacer{} เอวํทิฏฺฐิโนปิ โข ภิกฺขเว สนฺติ สตฺตา.\csromanspacer{} เอวํทิฏฺฐีนมฺปิ โข ภิกฺขเว สตฺตานํ อตฺเถว อญฺญถตฺตํ, อตฺถิ วิปริณาโม.\csromanspacer{} เอวํ ปสฺสํ ภิกฺขเว สุตวา อริยสาวโก ตสฺมิมฺปิ นิพฺพินฺทติ, ตสฺมิํ นิพฺพินฺทนฺโต อคฺเค วิรชฺชติ, ปเคว หีนสฺมิํ. (8)}

\prose{สนฺติ ภิกฺขเว เอเก สมณพฺราหฺมณา ปรมตฺถวิสุทฺธิํ ปญฺญาเปนฺติ.\csromanspacer{} เอตทคฺคํ ภิกฺขเว ปรมตฺถวิสุทฺธิํ ปญฺญเปนฺตานํ, ยทิทํ สพฺพโส อากิญฺจญฺญายตนํ สมติกฺกมฺม เนวสญฺญานาสญฺญายตนํ อุปสมฺปชฺช วิหรติ.\csromanspacer{} เต ตทภิญฺญาย ตสฺส สจฺฉิกิริยาย ธมฺมํ เทเสนฺติ.\csromanspacer{} เอวํวาทิโนปิ โข ภิกฺขเว สนฺติ สตฺตา.\csromanspacer{} เอวํวาทีนมฺปิ โข ภิกฺขเว สตฺตานํ อตฺเถว อญฺญถตฺตํ, อตฺถิ วิปริณาโม.\csromanspacer{} เอวํ ปสฺสํ ภิกฺขเว สุตวา อริยสาวโก ตสฺมิมฺปิ นิพฺพินฺทติ, ตสฺมิํ นิพฺพินฺทนฺโต อคฺเค วิรชฺชติ, ปเคว หีนสฺมิํ. (9)}

\prose{สนฺติ ภิกฺขเว เอเก สมณพฺราหฺมณา ปรมทิฏฺฐธมฺมนิพฺพานํ ปญฺญาเปนฺติ.\csromanspacer{} เอตทคฺคํ ภิกฺขเว ปรมทิฏฺฐธมฺมนิพฺพานํ ปญฺญเปนฺตานํ, ยทิทํ ฉนฺนํ ผสฺสายตนานํ สมุทยญฺจ อตฺถงฺคมญฺจ อสฺสาทญฺจ อาทีนวญฺจ นิสฺสรณญฺจ ยถาภูตํ วิทิตฺวา อนุปาทา วิโมกฺโข.\csromanspacer{} เอวํวาทิํ โข มํ ภิกฺขเว เอวมกฺขายิํ เอเก สมณพฺราหฺมณา อสตา ตุจฺฉา มุสา อภูเตน อพฺภาจิกฺขนฺติ “สมโณ โคตโม น กามานํ ปริญฺญํ ปญฺญาเปติ, น รูปานํ ปริญฺญํ ปญฺญาเปติ, น เวทนานํ ปริญฺญํ ปญฺญาเปตี”ติ.\csromanspacer{} กามานญฺจาหํ ภิกฺขเว ปริญฺญํ ปญฺญาเปมิ, รูปานญฺจ ปริญฺญํ ปญฺญาเปมิ, เวทนานญฺจ ปริญฺญํ ปญฺญาเปมิ, ทิฏฺเฐว ธมฺเม นิจฺฉาโต นิพฺภุโต สีติภูโต อนุปาทา ปรินิพฺพานํ ปญฺญาเปมี”ติ. (10).\csromanspacer{}\csromanspacer{}นวมํ.}

\csromanheader{2}{10. ทุติยโกสลสุตฺต}
\proseitem{30}{เอกํ สมยํ ภควา สาวตฺถิยํ วิหรติ เชตวเน อนาถปิณฺฑิกสฺส อาราเม.\csromanspacer{} เตน โข ปน สมเยน ราชา ปเสนทิ โกสโล อุยฺโยธิกา นิวตฺโต โหติ วิชิตสงฺคาโม ลทฺธาธิปฺปาโย.\csromanspacer{} อถ โข ราชา ปเสนทิ โกสโล เยน อาราโม เตน ปายาสิ, ยาวติกา ยานสฺส ภูมิ ยาเนน คนฺตฺวา ยานา ปจฺโจโรหิตฺวา ปตฺติโกว อารามํ ปาวิสิ.\csromanspacer{} เตน โข ปน สมเยน สมฺพหุลา}

\vspace{\tipitakaparskip}
\csromancenter{ทสกนิปาตปาฬิ ภิกฺขู อพฺโภกาเส จงฺกมนฺติ.\csromanspacer{} อถ โข ราชา ปเสนทิ โกสโล เยน เต ภิกฺขู เตนุปสงฺกมิ, อุปสงฺกมิตฺวา เต ภิกฺขู เอตทโวจ “กหํ นุ โข ภนฺเต ภควา เอตรหิ วิหรติ อรหํ สมฺมาสมฺพุทฺโธ, ทสฺสนกามา หิ มยํ ภนฺเต ตํ ภควนฺตํ อรหนฺตํ สมฺมาสมฺพุทฺธนฺ”ติ.\csromanspacer{} เอโส มหาราช วิหาโร สํวุตทฺวาโร, เตน อปฺปสทฺโท อุปสงฺกมิตฺวา อตรมาโน อาลินฺทํ ปวิสิตฺวา อุกฺกาสิตฺวา อคฺคฬํ อาโกเฏหิ, วิวริสฺสติ เต ภควา ทฺวารนฺติ.}
\prose{\csromanfolio{308}อถ โข ราชา ปเสนทิ โกสโล เยน โส วิหาโร สํวุตทฺวาโร, เตน อปฺปสทฺโท อุปสงฺกมิตฺวา อตรมาโน อาลินฺทํ ปวิสิตฺวา อุกฺกาสิตฺวา อคฺคฬํ อาโกเฏสิ, วิวริ ภควา ทฺวารํ.\csromanspacer{} อถ โข ราชา ปเสนทิ โกสโล วิหารํ ปวิสิตฺวา ภควโต ปาเทสุ สิรสา นิปติตฺวา ภควโต ปาทานิ มุเขน จ ปริจุมฺพติ, ปาณีหิ จ ปริสมฺพาหติ, นามญฺจ สาเวติ “ราชาหํ ภนฺเต ปเสนทิ โกสโล, ราชาหํ ภนฺเต ปเสนทิ โกสโล”ติ.}

\prose{กํ ปน ตฺวํ มหาราช อตฺถวสํ สมฺปสฺสมาโน อิมสฺมิํ สรีเร เอวรูปํ ปรมนิปจฺจการํ กโรสิ, เมตฺตูปหารํ อุปทํเสสีติ.\csromanspacer{} กตญฺญุตํ โข อหํ ภนฺเต กตเวทิตํ สมฺปสฺสมาโน ภควติ เอวรูปํ ปรมนิปจฺจการํ กโรมิ, เมตฺตูปหารํ อุปทํเสมิ.}

\prose{ภควา หิ ภนฺเต พหุชนหิตาย ปฏิปนฺโน พหุชนสุขาย, พหุโน ชนสฺส อริเย ญาเย ปติฏฺฐาปิตา ยทิทํ กลฺยาณธมฺมตาย กุสลธมฺมตาย.\csromanspacer{} ยมฺปิ ภนฺเต ภควา พหุชนหิตาย ปฏิปนฺโน พหุชนสุขาย, พหุโน ชนสฺส อริเย ญาเย ปติฏฺฐาปิตา ยทิทํ กลฺยาณธมฺมตาย กุสลธมฺมตาย, อิทมฺปิ โข อหํ ภนฺเต อตฺถวสํ สมฺปสฺสมาโน ภควติ เอวรูปํ ปรมนิปจฺจการํ กโรมิ, เมตฺตูปหารํ อุปทํเสมิ. (1)}

\prose{ปุน จปรํ ภนฺเต ภควา สีลวา วุทฺธสีโล อริยสีโล กุสลสีโล กุสลสีเลน สมนฺนาคโต.\csromanspacer{} ยมฺปิ ภนฺเต ภควา สีลวา วุทฺธสีโล อริยสีโล กุสลสีโล กุสลสีเลน สมนฺนาคโต, อิทมฺปิ โข อหํ ภนฺเต อตฺถวสํ สมฺปสฺสมาโน ภควติ เอวรูปํ ปรมนิปจฺจการํ กโรมิ, เมตฺตูปหารํ อุปทํเสมิ. (2)}

\chapterheadpageread{3. มหาวคฺค}
\prose{\csromanfolio{309}ปุน จปรํ ภนฺเต ภควา ทีฆรตฺตํ อารญฺญิโก อรญฺญวนปตฺถานิ ปนฺตานิ เสนาสนานิ ปฏิเสวติ.\csromanspacer{} ยมฺปิ ภนฺเต ภควา ทีฆรตฺตํ อารญฺญิโก อรญฺญวนปตฺถานิ ปนฺตานิ เสนาสนานิ ปฏิเสวติ, อิทมฺปิ โข อหํ ภนฺเต อตฺถวสํ สมฺปสฺสมาโน ภควติ เอวรูปํ ปรมนิปจฺจการํ กโรมิ, เมตฺตูปหารํ อุปทํเสมิ. (3)}

\csromanmatikaanchor{mtk.56}
\csromantocmarkref{subsection}{10. ทุติยโกสลสุตฺต}{mtk.56}
\bookmark[dest={mtk.56},level=2]{10. ทุติยโกสลสุตฺต}
\prose{ปุน จปรํ ภนฺเต ภควา สนฺตุฏฺโฐ อิตรีตรจีวรปิณฺฑปาต- เสนาสนคิลานปจฺจยเภสชฺชปริกฺขาเรน.\csromanspacer{} ยมฺปิ ภนฺเต ภควา สนฺตุฏฺโฐ อิตรีตรจีวรปิณฺฑปาตเสนาสนคิลานปจฺจยเภสชฺช- ปริกฺขาเรน, อิทมฺปิ โข อหํ ภนฺเต อตฺถวสํ สมฺปสฺสมาโน ภควติ เอวรูปํ ปรมนิปจฺจการํ กโรมิ, เมตฺตูปหารํ อุปทํเสมิ. (4)}

\prose{ปุน จปรํ ภนฺเต ภควา อาหุเนยฺโย ปาหุเนยฺโย ทกฺขิเณยฺโย อญฺชลิกรณีโย อนุตฺตรํ ปุญฺญกฺเขตฺตํ โลกสฺส.\csromanspacer{} ยมฺปิ ภนฺเต ภควา อาหุเนยฺโย ปาหุเนยฺโย ทกฺขิเณยฺโย อญฺชลิกรณีโย อนุตฺตรํ ปุญฺญกฺเขตฺตํ โลกสฺส, อิทมฺปิ โข อหํ ภนฺเต อตฺถวสํ สมฺปสฺสมาโน ภควติ เอวรูปํ ปรมนิปจฺจการํ กโรมิ, เมตฺตูปหารํ อุปทํเสมิ. (5)}

\prose{ปุน จปรํ ภนฺเต ภควา ยายํ กถา อภิสลฺเลขิกา เจโตวิวรณสปฺปายา.\csromanspacer{} เสยฺยถิทํ, อปฺปิจฺฉกถา สนฺตุฏฺฐิกถา ปวิเวกกถา อสํสคฺคกถา วีริยารมฺภกถา สีลกถา สมาธิกถา ปญฺญากถา วิมุตฺติกถา วิมุตฺติญาณทสฺสนกถา, เอวรูปาย กถาย นิกามลาภี อกิจฺฉลาภี อกสิรลาภี.\csromanspacer{} ยมฺปิ ภนฺเต ภควา ยายํ กถา อภิสลฺเลขิกา เจโตวิวรณสปฺปายา.\csromanspacer{} เสยฺยถิทํ, อปฺปิจฺฉกถา ฯเปฯ วิมุตฺติญาณทสฺสนกถา, เอวรูปาย กถาย นิกามลาภี อกิจฺฉลาภี อกสิรลาภี, อิทมฺปิ โข อหํ ภนฺเต อตฺถวสํ สมฺปสฺสมาโน ภควติ เอวรูปํ ปรมนิปจฺจการํ กโรมิ, เมตฺตูปหารํ อุปทํเสมิ. (6)}

\prose{ปุน จปรํ ภนฺเต ภควา จตุนฺนํ ฌานานํ อาภิเจตสิกานํ ทิฏฺฐธมฺมสุขวิหารานํ นิกามลาภี อกิจฺฉลาภี อกสิรลาภี.\csromanspacer{} ยมฺปิ ภนฺเต ภควา จตุนฺนํ ฌานานํ อาภิเจตสิกานํ ทิฏฺฐธมฺมสุขวิหารานํ นิกามลาภี อกิจฺฉลาภี อกสิรลาภี, อิทมฺปิ โข อหํ ภนฺเต อตฺถวสํ สมฺปสฺสมาโน ภควติ เอวรูปํ ปรมนิปจฺจการํ กโรมิ, เมตฺตูปหารํ อุปทํเสมิ. (7)}

\gambhira{ทสกนิปาตปาฬิ}
\prose{\csromanfolio{310}ปุน จปรํ ภนฺเต ภควา อเนกวิหิตํ ปุพฺเพนิวาสํ อนุสฺสรติ.\csromanspacer{} เสยฺยถิทํ, เอกมฺปิ ชาติํ เทฺวปิ ชาติโย ติสฺโสปิ ชาติโย จตสฺโสปิ ชาติโย ปญฺจปิ ชาติโย ทสปิ ชาติโย วีสมฺปิ ชาติโย ติํสมฺปิ ชาติโย จตฺตาลีสมฺปิ ชาติโย ปญฺญาสมฺปิ ชาติโย ชาติสตมฺปิ ชาติสหสฺสมฺปิ ชาติสตสหสฺสมฺปิ อเนเกปิ สํวฏฺฏกปฺเป อเนเกปิ วิวฏฺฏกปฺเป อเนเกปิ สํวฏฺฏวิวฏฺฏกปฺเป “อมุตฺราสิํ เอวํนาโม เอวํโคตฺโต เอวํวณฺโณ เอวมาหาโร เอวํสุขทุกฺขปฺปฏิสํเวที เอวมายุปริยนฺโต, โส ตโต จุโต อมุตฺร อุทปาทิํ, ตตฺราปาสิํ เอวํนาโม เอวํโคตฺโต เอวํวณฺโณ เอวมาหาโร เอวํสุขทุกฺขปฺปฏิสํเวที เอวมายุปริยนฺโต, โส ตโต จุโต อิธูปปนฺโน”ติ.\csromanspacer{} อิติ สาการํ ส-อุทฺเทสํ อเนกวิหิตํ ปุพฺเพนิวาสํ อนุสฺสรติ.\csromanspacer{} ยมฺปิ ภนฺเต ภควา อเนกวิหิตํ ปุพฺเพนิวาสํ อนุสฺสรติ.\csromanspacer{} เสยฺยถิทํ, เอกมฺปิ ชาติํ เทฺวปิ ชาติโย ฯเปฯ อิติ สาการํ ส-อุทฺเทสํ อเนกวิหิตํ ปุพฺเพนิวาสํ อนุสฺสรติ.\csromanspacer{} อิทมฺปิ โข อหํ ภนฺเต อตฺถวสํ สมฺปสฺสมาโน ภควติ เอวรูปํ ปรมนิปจฺจการํ กโรมิ, เมตฺตูปหารํ อุปทํเสมิ. (8)}

\prose{ปุน จปรํ ภนฺเต ภควา ทิพฺเพน จกฺขุนา วิสุทฺเธน อติกฺกนฺตมานุสเกน สตฺเต ปสฺสติ จวมาเน อุปปชฺชมาเน หีเน ปณีเต สุวณฺเณ ทุพฺพณฺเณ สุคเต ทุคฺคเต, ยถากมฺมูปเค สตฺเต ปชานาติ “อิเม วต โภนฺโต สตฺตา กายทุจฺจริเตน สมนฺนาคตา วจีทุจฺจริเตน สมนฺนาคตา มโนทุจฺจริเตน สมนฺนาคตา อริยานํ อุปวาทกา มิจฺฉาทิฏฺฐิกา มิจฺฉาทิฏฺฐิกมฺมสมาทานา, เต กายสฺส เภทา ปรํ มรณา อปายํ ทุคฺคติํ วินิปาตํ นิรยํ อุปปนฺนา.\csromanspacer{} อิเม วา ปน โภนฺโต สตฺตา กายสุจริเตน สมนฺนาคตา วจีสุจริเตน สมนฺนาคตา มโนสุจริเตน สมนฺนาคตา อริยานํ อนุปวาทกา สมฺมาทิฏฺฐิกา สมฺมาทิฏฺฐิกมฺม- สมาทานา, เต กายสฺส เภทา ปรํ มรณา สุคติํ สคฺคํ โลกํ อุปปนฺนา”ติ.\csromanspacer{} อิติ ทิพฺเพน จกฺขุนา วิสุทฺเธน อติกฺกนฺตมานุสเกน สตฺเต ปสฺสติ ฯเปฯ ยถากมฺมูปเค สตฺเต ปชานาติ.\csromanspacer{} ยมฺปิ ภนฺเต ภควา ทิพฺเพน จกฺขุนา วิสุทฺเธน อติกฺกนฺตมานุสเกน ฯเปฯ ยถากมฺมูปเค สตฺเต ปชานาติ.\csromanspacer{} อิทมฺปิ โข อหํ ภนฺเต อตฺถวสํ สมฺปสฺสมาโน ภควติ เอวรูปํ ปรมนิปจฺจการํ กโรมิ, เมตฺตูปหารํ อุปทํเสมิ. (9)}

\chapterheadpageread{4. อุปาลิวคฺค}
\csromanmatikaanchor{mtk.57}
\csromantocmarkref{section}{อุทฺทานคาถา}{mtk.57}
\bookmark[dest={mtk.57},level=1]{อุทฺทานคาถา}
\prose{\csromanfolio{311}ปุน จปรํ ภนฺเต ภควา อาสวานํ ขยา อนาสวํ เจโตวิมุตฺติํ ปญฺญาวิมุตฺติํ ทิฏฺเฐว ธมฺเม สยํ อภิญฺญา สจฺฉิกตฺวา อุปสมฺปชฺช วิหรติ.\csromanspacer{} ยมฺปิ ภนฺเต ภควา อาสวานํ ขยา อนาสวํ เจโตวิมุตฺติํ ฯเปฯ สจฺฉิกตฺวา อุปสมฺปชฺช วิหรติ, อิทมฺปิ โข อหํ ภนฺเต อตฺถวสํ สมฺปสฺสมาโน ภควติ เอวรูปํ ปรมนิปจฺจการํ กโรมิ, เมตฺตูปหารํ อุปทํเสมิ. (10)}

\prose{หนฺท จ ทานิ มยํ ภนฺเต คจฺฉาม, พหุกิจฺจา มยํ พหุกรณียาติ.\csromanspacer{} ยสฺส ทานิ ตฺวํ มหาราช กาลํ มญฺญสีติ.\csromanspacer{} อถ โข ราชา ปเสนทิ โกสโล อุฏฺฐายาสนา ภควนฺตํ อภิวาเทตฺวา ปทกฺขิณํ กตฺวา ปกฺกามีติ.\csromanspacer{}\csromanspacer{}ทสมํ.\csromanspacer{} มหาวคฺโค ตติโย.}

\titlehead{\textbf{ตสฺสุทฺทานํ}}
\csromangathameasure{สีหาธิวุตฺติ กาเยน, จุนฺเทน กสิเณน จ. \\ กาฬี จ เทฺว มหาปญฺหา, โกสเลหิ ปเร ทุเวติ.}
\csromangathagroup{\csromangathastanza{สีหาธิวุตฺติ กาเยน, จุนฺเทน กสิเณน จ. \\ กาฬี จ เทฺว มหาปญฺหา, โกสเลหิ ปเร ทุเวติ.}}

\chapterhead{4. อุปาลิวคฺค}
\csromanheader{2}{1. อุปาลิสุตฺต}
\proseitem{31}{อถ โข อายสฺมา อุปาลิ เยน ภควา เตนุปสงฺกมิ, อุปสงฺกมิตฺวา ภควนฺตํ อภิวาเทตฺวา เอกมนฺตํ นิสีทิ, เอกมนฺตํ นิสินฺโน โข อายสฺมา อุปาลิ ภควนฺตํ เอตทโวจ “กติ นุ โข ภนฺเต อตฺถวเส ปฏิจฺจ ตถาคเตน สาวกานํ สิกฺขาปทํ ปญฺญตฺตํ ปาติโมกฺขํ อุทฺทิฏฺฐนฺ”ติ.}

\prose{ทส โข อุปาลิ อตฺถวเส ปฏิจฺจ ตถาคเตน สาวกานํ สิกฺขาปทํ ปญฺญตฺตํ ปาติโมกฺขํ อุทฺทิฏฺฐํ.\csromanspacer{} กตเม ทส? สํฆสุฏฺฐุตาย สํฆผาสุตาย ทุมฺมงฺกูนํ ปุคฺคลานํ นิคฺคหาย เปสลานํ ภิกฺขูนํ ผาสุวิหาราย ทิฏฺฐธมฺมิกานํ อาสวานํ สํวราย สมฺปรายิกานํ อาสวานํ ปฏิฆาตาย อปฺปสนฺนานํ ปสาทาย ปสนฺนานํ ภิยฺโยภาวาย สทฺธมฺมฏฺฐิติยา วินยานุคฺคหาย.\csromanspacer{} อิเม โข อุปาลิ ทส อตฺถวเส ปฏิจฺจ ตถาคเตน สาวกานํ สิกฺขาปทํ ปญฺญตฺตํ ปาโมกฺขํ อุทฺทิฏฺฐนฺติ.\csromanspacer{}\csromanspacer{}ปฐมํ.}

\csromanheader{2}{ทสกนิปาตปาฬิ \textbf{2. ปาติโมกฺขฏฺฐปนาสุตฺต}}
\proseitem{32}{\csromanfolio{312}กติ นุ โข ภนฺเต ปาติโมกฺขฏฺฐปนาติ.\csromanspacer{} ทส โข อุปาลิ ปาติโมกฺขฏฺฐปนา.\csromanspacer{} กตเม ทส? ปาราชิโก ตสฺสํ ปริสายํ นิสินฺโน โหติ, ปาราชิกกถา วิปฺปกตา โหติ, อนุปสมฺปนฺโน ตสฺสํ ปริสายํ นิสินฺโน โหติ, อนุปสมฺปนฺนกถา วิปฺปกตา โหติ, สิกฺขํ ปจฺจกฺขาตโก ตสฺสํ ปริสายํ นิสินฺโน โหติ, สิกฺขํ ปจฺจกฺขาตกกถา วิปฺปกตา โหติ, ปณฺฑโก ตสฺสํ ปริสายํ นิสินฺโน โหติ, ปณฺฑกกถา วิปฺปกตา โหติ, ภิกฺขุนิทูสโก ตสฺสํ ปริสายํ นิสินฺโน โหติ, ภิกฺขุนิทูสกกถา วิปฺปกตา โหติ.\csromanspacer{} อิเม โข อุปาลิ ทส ปาติโมกฺขฏฺฐปนาติ.\csromanspacer{}\csromanspacer{}ทุติยํ.}

\csromanheader{2}{3. อุพฺพาหิกาสุตฺต}
\proseitem{33}{\csromansymbolfootnote{*}{วิ 4. 223 ปิฏฺเฐปิ.}กติหิ นุ โข ภนฺเต ธมฺเมหิ สมนฺนาคโต ภิกฺขุ อุพฺพาหิกาย สมฺมนฺนิตพฺโพติ.\csromanspacer{} ทสหิ โข อุปาลิ ธมฺเมหิ สมนฺนาคโต ภิกฺขุ อุพฺพาหิกาย สมฺมนฺนิตพฺโพ.\csromanspacer{} กตเมหิ ทสหิ? อิธุปาลิ ภิกฺขุ สีลวา โหติ, ปาติโมกฺขสํวรสํวุโต วิหรติ อาจารโคจรสมฺปนฺโน, อณุมตฺเตสุ วชฺเชสุ ภยทสฺสาวี สมาทาย สิกฺขติ สิกฺขาปเทสุ.\csromanspacer{} พหุสฺสุโต โหติ สุตธโร สุตสนฺนิจโย, เย เต ธมฺมา อาทิกลฺยาณา มชฺเฌกลฺยาณา ปริโยสานกลฺยาณา สาตฺถํ สพฺยญฺชนํ\footnote{สตฺถา สพฺยญฺชนา (สี) เอวมุปริปิ.} เกวลปริปุณฺณํ ปริสุทฺธํ พฺรหฺมจริยํ อภิวทนฺติ, ตถารูปาสฺส ธมฺมา พหุสฺสุตา โหนฺติ ธาตา วจสา ปริจิตา มนสานุเปกฺขิตา ทิฏฺฐิยา สุปฺปฏิวิทฺธา.\csromanspacer{} อุภยานิ โข ปนสฺส ปาติโมกฺขานิ วิตฺถาเรน สฺวาคตานิ โหนฺติ สุวิภตฺตานิ สุปฺปวตฺตีนิ สุวินิจฺฉิตานิ สุตฺตโส อนุพฺยญฺชนโส, วินเย โข ปน ฐิโต โหติ อสํหีโร.\csromanspacer{} ปฏิพโล โหติ อุโภ อตฺถปจฺจตฺถิเก สญฺญาเปตุํ ปญฺญาเปตุํ นิชฺฌาเปตุํ เปกฺเขตุํ ปสาเทตุํ.\csromanspacer{} อธิกรณสมุปฺปาทวูปสมกุสโล โหติ.\csromanspacer{} อธิกรณํ ชานาติ.\csromanspacer{} อธิกรณสมุทยํ ชานาติ.\csromanspacer{} อธิกรณนิโรธํ ชานาติ.\csromanspacer{} อธิกรณนิโรธคามินิํ ปฏิปทํ ชานาติ.\csromanspacer{} อิเมหิ โข อุปาลิ ทสหิ ธมฺเมหิ สมนฺนาคโต ภิกฺขุ อุพฺพาหิกาย สมฺมนฺนิตพฺโพติ.\csromanspacer{}\csromanspacer{}ตติยํ.}

\csromanheader{2}{4. อุปาลิวคฺค \textbf{4. อุปสมฺปทาสุตฺต}}
\proseitem{34}{\csromanfolio{313}กติหิ นุ โข ภนฺเต ธมฺเมหิ สมนฺนาคเตน ภิกฺขุนา อุปสมฺปาเทตพฺพนฺติ.\csromanspacer{} ทสหิ โข อุปาลิ ธมฺเมหิ สมนฺนาคเตน ภิกฺขุนา อุปสมฺปาเทตพฺพํ.\csromanspacer{} กตเมหิ ทสหิ? อิธุปาลิ ภิกฺขุ สีลวา โหติ, ปาติโมกฺขสํวรสํวุโต วิหรติ อาจารโคจรสมฺปนฺโน, อณุมตฺเตสุ วชฺเชสุ ภยทสฺสาวี สมาทาย สิกฺขติ สิกฺขาปเทสุ.\csromanspacer{} พหุสฺสุโต โหติ สุตธโร สุตสนฺนิจโย, เย เต ธมฺมา อาทิกลฺยาณา มชฺเฌกลฺยาณา ปริโยสานกลฺยาณา สาตฺถํ สพฺยญฺชนํ เกวลปริปุณฺณํ ปริสุทฺธํ พฺรหฺมจริยํ อภิวทนฺติ, ตถารูปาสฺส ธมฺมา พหุสฺสุตา โหนฺติ ธาตา วจสา ปริจิตา มนสานุเปกฺขิตา ทิฏฺฐิยา สุปฺปฏิวิทฺธา.\csromanspacer{} ปาติโมกฺขํ โข ปนสฺส วิตฺถาเรน สฺวาคตํ โหติ สุวิภตฺตํ สุปฺปวตฺตํ สุวินิจฺฉิตํ สุตฺตโส อนุพฺยญฺชนโส.\csromanspacer{} ปฏิพโล โหติ คิลานํ อุปฏฺฐาตุํ วา อุปฏฺฐาเปตุํ วา.\csromanspacer{} ปฏิพโล โหติ อนภิรติํ วูปกาเสตุํ วา วูปกาสาเปตุํ วา.\csromanspacer{} ปฏิพโล โหติ อุปฺปนฺนํ กุกฺกุจฺจํ ธมฺมโต วิโนเทตุํ.\csromanspacer{} ปฏิพโล โหติ อุปฺปนฺนํ ทิฏฺฐิคตํ ธมฺมโต วิเวเจตุํ.\csromanspacer{} ปฏิพโล โหติ อธิสีเล สมาทเปตุํ.\csromanspacer{} ปฏิพโล โหติ อธิจิตฺเต สมาทเปตุํ.\csromanspacer{} ปฏิพโล โหติ อธิปญฺญาย สมาทเปตุํ.\csromanspacer{} อิเมหิ โข อุปาลิ ทสหิ ธมฺเมหิ สมนฺนาคเตน ภิกฺขุนา อุปสมฺปาเทตพฺพนฺติ.\csromanspacer{}\csromanspacer{}จตุตฺถํ.}

\csromanmatikaanchor{mtk.58}
\csromanmatikaanchor{mtk.59}
\csromantocmarknopage{section}{4. อุปาลิวคฺค}{mtk.58}
\bookmark[dest={mtk.58},level=1]{4. อุปาลิวคฺค}
\csromantocmarkref{subsection}{1. อุปาลิสุตฺต}{mtk.59}
\bookmark[dest={mtk.59},level=2]{1. อุปาลิสุตฺต}
\csromanheader{2}{5. นิสฺสยสุตฺต}
\proseitem{35}{กติหิ นุ โข ภนฺเต ธมฺเมหิ สมนฺนาคเตน ภิกฺขุนา นิสฺสโย ทาตพฺโพติ.\csromanspacer{} ทสหิ โข อุปาลิ ธมฺเมหิ สมนฺนาคเตน ภิกฺขุนา นิสฺสโย ทาตพฺโพ.\csromanspacer{} กตเมหิ ทสหิ? อิธุปาลิ ภิกฺขุ สีลวา โหติ ฯเปฯ สมาทาย สิกฺขติ สิกฺขาปเทสุ.\csromanspacer{} พหุสฺสุโต โหติ ฯเปฯ ทิฏฺฐิยา สุปฺปฏิวิทฺธา.\csromanspacer{} ปาติโมกฺขํ โข ปนสฺส วิตฺถาเรน สฺวาคตํ โหติ สุวิภตฺตํ สุปฺปวตฺตํ สุวินิจฺฉิตํ สุตฺตโส อนุพฺยญฺชนโส.\csromanspacer{} ปฏิพโล โหติ คิลานํ อุปฏฺฐาตุํ วา อุปฏฺฐาเปตุํ วา.\csromanspacer{} ปฏิพโล โหติ อนภิรติํ วูปกาเสตุํ วา วูปกาสาเปตุํ วา.\csromanspacer{} ปฏิพโล โหติ อุปฺปนฺนํ กุกฺกุจฺจํ ธมฺมโต วิโนเทตุํ.\csromanspacer{} ปฏิพโล โหติ อุปฺปนฺนํ ทิฏฺฐิคตํ ธมฺมโต วิเวเจตุํ.\csromanspacer{} ปฏิพโล โหติ อธิสีเล ฯเปฯ อธิจิตฺเต.\csromanspacer{} อธิปญฺญาย สมาทเปตุํ.\csromanspacer{} อิเมหิ โข อุปาลิ ทสหิ ธมฺเมหิ สมนฺนาคเตน ภิกฺขุนา นิสฺสโย ทาตพฺโพติ.\csromanspacer{}\csromanspacer{}ปญฺจมํ.}

\csromanheader{2}{ทสกนิปาตปาฬิ \textbf{6. สามเณรสุตฺต}}
\csromanmatikaanchor{mtk.60}
\csromantocmarkref{subsection}{2. ปาติโมกฺขฏฺฐปนาสุตฺต}{mtk.60}
\bookmark[dest={mtk.60},level=2]{2. ปาติโมกฺขฏฺฐปนาสุตฺต}
\proseitem{36}{\csromanfolio{314}กติหิ นุ โข ภนฺเต ธมฺเมหิ สมนฺนาคเตน ภิกฺขุนา สามเณโร อุปฏฺฐาเปตพฺโพติ.\csromanspacer{} ทสหิ โข อุปาลิ ธมฺเมหิ สมนฺนาคเตน ภิกฺขุนา สามเณโร อุปฏฺฐาเปตพฺโพ.\csromanspacer{} กตเมหิ ทสหิ? อิธุปาลิ ภิกฺขุ สีลวา โหติ ฯเปฯ สมาทาย สิกฺขติ สิกฺขาปเทสุ.\csromanspacer{} พหุสฺสุโต โหติ ฯเปฯ ทิฏฺฐิยา สุปฺปฏิวิทฺธา.\csromanspacer{} ปาติโมกฺขํ โข ปนสฺส วิตฺถาเรน สฺวาคตํ โหติ สุวิภตฺตํ สุปฺปวตฺตํ สุวินิจฺฉิตํ สุตฺตโส อนุพฺยญฺชนโส.\csromanspacer{} ปฏิพโล โหติ คิลานํ อุปฏฺฐาตุํ วา อุปฏฺฐาเปตุํ วา.\csromanspacer{} ปฏิพโล โหติ อนภิรติํ วูปกาเสตุํ วา วูปกาสาเปตุํ วา.\csromanspacer{} ปฏิพโล โหติ อุปฺปนฺนํ กุกฺกุจฺจํ ธมฺมโต วิโนเทตุํ.\csromanspacer{} ปฏิพโล โหติ อุปฺปนฺนํ ทิฏฺฐิคตํ ธมฺมโต วิเวเจตุํ.\csromanspacer{} ปฏิพโล โหติ อธิสีเล สมาทเปตุํ.\csromanspacer{} ปฏิพโล โหติ อธิจิตฺเต สมาทเปตุํ.\csromanspacer{} ปฏิพโล โหติ อธิปญฺญาย สมาทเปตุํ.\csromanspacer{} อิเมหิ โข อุปาลิ ทสหิ ธมฺเมหิ สมนฺนาคเตน ภิกฺขุนา สามเณโร อุปฏฺฐาเปตพฺโพติ.\csromanspacer{}\csromanspacer{}ฉฏฺฐํ.}

\csromanheader{2}{7. สํฆเภทสุตฺต}
\csromanmatikaanchor{mtk.61}
\csromantocmarkref{subsection}{3. อุพฺพาหิกาสุตฺต}{mtk.61}
\bookmark[dest={mtk.61},level=2]{3. อุพฺพาหิกาสุตฺต}
\proseitem{37}{“สํฆเภโท สํฆเภโท”ติ ภนฺเต วุจฺจติ, กิตฺตาวตา นุ โข ภนฺเต สํโฆ ภินฺโน โหตีติ.\csromanspacer{} อิธุปาลิ ภิกฺขู อธมฺมํ “ธมฺโม”ติ ทีเปนฺติ, ธมฺมํ “อธมฺโม”ติ ทีเปนฺติ, อวินยํ “วินโย”ติ ทีเปนฺติ, วินยํ “อวินโย”ติ ทีเปนฺติ, อภาสิตํ อลปิตํ ตถาคเตน “ภาสิตํ ลปิตํ ตถาคเตนา”ติ ทีเปนฺติ, ภาสิตํ ลปิตํ ตถาคเตน “อภาสิตํ อลปิตํ ตถาคเตนา”ติ ทีเปนฺติ, อนาจิณฺณํ ตถาคเตน “อาจิณฺณํ ตถาคเตนา”ติ ทีเปนฺติ, อาจิณฺณํ ตถาคเตน “อนาจิณฺณํ ตถาคเตนา”ติ ทีเปนฺติ, อปญฺญตฺตํ ตถาคเตน “ปญฺญตฺตํ ตถาคเตนา”ติ ทีเปนฺติ, ปญฺญตฺตํ ตถาคเตน “อปญฺญตฺตํ ตถาคเตนา”ติ ทีเปนฺติ.\csromanspacer{} เต อิเมหิ ทสหิ วตฺถูหิ อวกสฺสนฺติ อปกสฺสนฺติ, อาเวนิ\footnote{อาเวนิํ (วิ 4. 368 ปิฏฺเฐ) อาเวณิ, อาเวณิกํ (ตตฺเถว อโธลิปิ)} กมฺมานิ กโรนฺติ, อาเวนิ ปาติโมกฺขํ อุทฺทิสนฺติ.\csromanspacer{} เอตฺตาวตา โข อุปาลิ สํโฆ ภินฺโน โหตีติ.\csromanspacer{}\csromanspacer{}สตฺตมํ.}

\csromanheader{2}{4. อุปาลิวคฺค \textbf{8. สํฆสามคฺคีสุตฺต}}
\csromanmatikaanchor{mtk.62}
\csromantocmarkref{subsection}{4. อุปสมฺปทาสุตฺต}{mtk.62}
\bookmark[dest={mtk.62},level=2]{4. อุปสมฺปทาสุตฺต}
\proseitem{38}{\csromanfolio{315}\csromansymbolfootnote{*}{วิ 4. 368 ปิฏฺเฐปิ.}สํฆสามคฺคี สํฆสามคฺคี”ติ ภนฺเต วุจฺจติ, กิตฺตาวตา นุ โข ภนฺเต สํโฆ สมคฺโค โหตีติ.\csromanspacer{} อิธุปาลิ ภิกฺขู อธมฺมํ “อธมฺโม”ติ ทีเปนฺติ, ธมฺมํ “ธมฺโม”ติ ทีเปนฺติ, อวินยํ “อวินโย”ติ ทีเปนฺติ, วินยํ “วินโย”ติ ทีเปนฺติ, อภาสิตํ อลปิตํ ตถาคเตน “อภาสิตํ อลปิตํ ตถาคเตนา”ติ ทีเปนฺติ, ภาสิตํ ลปิตํ ตถาคเตน “ภาสิตํ ลปิตํ ตถาคเตนา”ติ ทีเปนฺติ, อนาจิณฺณํ ตถาคเตน “อนาจิณฺณํ ตถาคเตนา”ติ ทีเปนฺติ, อาจิณฺณํ ตถาคเตน “อาจิณฺณํ ตถาคเตนา”ติ ทีเปนฺติ, อปญฺญตฺตํ ตถาคเตน “อปญฺญตฺตํ ตถาคเตนา”ติ ทีเปนฺติ, ปญฺญตฺตํ ตถาคเตน “ปญฺญตฺตํ ตถาคเตนา”ติ ทีเปนฺติ.\csromanspacer{} เต อิเมหิ ทสหิ วตฺถูหิ น อวกสฺสนฺติ น อปกสฺสนฺติ, น อาเวนิ กมฺมานิ กโรนฺติ, น อาเวนิ ปาติโมกฺขํ อุทฺทิสนฺติ.\csromanspacer{} เอตฺตาวตา โข อุปาลิ สํโฆ สมคฺโค โหตีติ.\csromanspacer{}\csromanspacer{}อฏฺฐมํ.}

\csromanheader{2}{9. ปฐม-อานนฺทสุตฺต}
\csromanmatikaanchor{mtk.63}
\csromantocmarkref{subsection}{5. นิสฺสยสุตฺต}{mtk.63}
\bookmark[dest={mtk.63},level=2]{5. นิสฺสยสุตฺต}
\proseitem{39}{อถ โข อายสฺมา อานนฺโท เยน ภควา เตนุปสงฺกมิ, อุปสงฺกมิตฺวา ภควนฺตํ อภิวาเทตฺวา เอกมนฺตํ นิสีทิ, เอกมนฺตํ นิสินฺโน โข อายสฺมา อานนฺโท ภควนฺตํ เอตทโวจ “สํฆเภโท สํฆเภโท”ติ ภนฺเต วุจฺจติ, กิตฺตาวตา นุ โข ภนฺเต สํโฆ ภินฺโน โหตีติ.\csromanspacer{} อิธานนฺท ภิกฺขู อธมฺมํ “ธมฺโม”ติ ทีเปนฺติ, ธมฺมํ “อธมฺโม”ติ ทีเปนฺติ, อวินยํ “วินโย”ติ ทีเปนฺติ ฯเปฯ ปญฺญตฺตํ ตถาคเตน “อปญฺญตฺตํ ตถาคเตนา”ติ ทีเปนฺติ.\csromanspacer{} เต อิเมหิ ทสหิ วตฺถูหิ อวกสฺสนฺติ อปกสฺสนฺติ, อาเวนิ กมฺมานิ กโรนฺติ, อาเวนิ ปาติโมกฺขํ อุทฺทิสนฺติ.\csromanspacer{} เอตฺตาวตา โข อานนฺท สํโฆ ภินฺโน โหตีติ.}

\prose{สมคฺคํ ปน ภนฺเต สํฆํ ภินฺทิตฺวา กิํ โส ปสวตีติ.\csromanspacer{} กปฺปฏฺฐิกํ อานนฺท กิพฺพิสํ ปสวตีติ.\csromanspacer{} กิํ ปน ภนฺเต กปฺปฏฺฐิกํ กิพฺพิสนฺติ.\csromanspacer{} กปฺปํ อานนฺท นิรยมฺหิ ปจฺจตีติ.}

\csromanmatikaanchor{mtk.64}
\csromantocmarkref{subsection}{6. สามเณรสุตฺต}{mtk.64}
\bookmark[dest={mtk.64},level=2]{6. สามเณรสุตฺต}
\csromangathameasure{อาปายิโก เนรยิโก, กปฺปฏฺโฐ สํฆเภทโก. \\ วคฺครโต อธมฺมฏฺโฐ, โยคกฺเขมา ปธํสติ.}
\csromangathagroup{\csromangathastanza{อาปายิโก เนรยิโก, กปฺปฏฺโฐ สํฆเภทโก. \\ วคฺครโต อธมฺมฏฺโฐ, โยคกฺเขมา ปธํสติ.}}

\prose{สํฆํ สมคฺคํ ภินฺทิตฺวา\footnote{เภตฺวาน (สี, สฺยา), ภิตฺวาน (ก) วิ 4. 369; ขุ 1. 203; อภิ 4. 346 ปิฏฺเฐสุปิ.},}

\csromanmatikaanchor{mtk.65}
\csromantocmarkref{subsection}{7. สํฆเภทสุตฺต}{mtk.65}
\bookmark[dest={mtk.65},level=2]{7. สํฆเภทสุตฺต}
\prose{กปฺปํ นิรยมฺหิ ปจฺจตีติ.\csromanspacer{}\csromanspacer{}นวมํ.}

\csromanheader{2}{ทสกนิปาตปาฬิ \textbf{10. ทุติย-อานนฺทสุตฺต}}
\csromanmatikaanchor{mtk.66}
\csromanmatikaanchor{mtk.69}
\csromantocmarkref{subsection}{8. สํฆสามคฺคีสุตฺต}{mtk.66}
\bookmark[dest={mtk.66},level=2]{8. สํฆสามคฺคีสุตฺต}
\proseitem{40}{\csromanfolio{316}“สํฆสามคฺคี สํฆสามคฺคี”ติ ภนฺเต วุจฺจติ, กิตฺตาวตา นุ โข ภนฺเต สํโฆ สมคฺโค โหตีติ.\csromanspacer{} อิธานนฺท ภิกฺขู อธมฺมํ “อธมฺโม”ติ ทีเปนฺติ, ธมฺมํ “ธมฺโม”ติ ทีเปนฺติ, อวินยํ “อวินโย”ติ ทีเปนฺติ, วินยํ “วินโย”ติ ทีเปนฺติ, อภาสิตํ อลปิตํ ตถาคเตน “อภาสิตํ อลปิตํ ตถาคเตนา”ติ ทีเปนฺติ, ภาสิตํ ลปิตํ ตถาคเตน “ภาสิตํ ลปิตํ ตถาคเตนา”ติ ทีเปนฺติ, อนาจิณฺณํ ตถาคเตน “อนาจิณฺณํ ตถาคเตนา”ติ ทีเปนฺติ, อาจิณฺณํ ตถาคเตน “อาจิณฺณํ ตถาคเตนา”ติ ทีเปนฺติ, อปญฺญตฺตํ ตถาคเตน “อปญฺญตฺตํ ตถาคเตนา”ติ ทีเปนฺติ, ปญฺญตฺตํ ตถาคเตน “ปญฺญตฺตํ ตถาคเตนา”ติ ทีเปนฺติ.\csromanspacer{} เต อิเมหิ ทสหิ วตฺถูหิ น อวกสฺสนฺติ น อปกสฺสนฺติ, น อาเวนิ กมฺมานิ กโรนฺติ, น อาเวนิ ปาติโมกฺขํ อุทฺทิสนฺติ.\csromanspacer{} เอตฺตาวตา โข อานนฺท สํโฆ สมคฺโค โหตีติ.}

\prose{ภินฺนํ ปน ภนฺเต สํฆํ สมคฺคํ กตฺวา กิํ โส ปสวตีติ.\csromanspacer{} พฺรหฺมํ อานนฺท ปุญฺญํ ปสวตีติ.\csromanspacer{} กิํ ปน ภนฺเต พฺรหฺมํ ปุญฺญนฺติ.\csromanspacer{} กปฺปํ อานนฺท สคฺคมฺหิ โมทตีติ.}

\csromanmatikaanchor{mtk.67}
\csromantocmarkref{subsection}{9. ปฐม-อานนฺทสุตฺต}{mtk.67}
\bookmark[dest={mtk.67},level=2]{9. ปฐม-อานนฺทสุตฺต}
\prose{สุขา สํฆสฺส สามคฺคี, สมคฺคานญฺจ อนุคฺคโห.}

\prose{สมคฺครโต ธมฺมฏฺโฐ, โยคกฺเขมา น ธํสติ.}

\prose{สํฆํ สมคฺคํ กตฺวาน, กปฺปํ สคฺคมฺหิ โมทตีติ.\csromanspacer{}\csromanspacer{}ทสมํ.\csromanspacer{} อุปาลิวคฺโค จตุตฺโถ.}

\titlehead{\textbf{ตสฺสุทฺทานํ}}
\csromangathameasure{อุปาลิ ฐปนา อุพฺพาโห, อุปสมฺปทนิสฺสยา. \\ สามเณโร จ เทฺว เภทา, อานนฺเทหิ ปเร ทุเวติ.}
\csromangathagroup{\csromangathastanza{อุปาลิ ฐปนา อุพฺพาโห, อุปสมฺปทนิสฺสยา. \\ สามเณโร จ เทฺว เภทา, อานนฺเทหิ ปเร ทุเวติ.}}

\csromanheader{1}{5. อกฺโกสวคฺค}
\csromanmatikaanchor{mtk.68}
\csromantocmarkref{subsection}{10. ทุติย-อานนฺทสุตฺต}{mtk.68}
\bookmark[dest={mtk.68},level=2]{10. ทุติย-อานนฺทสุตฺต}
\csromantocmarkref{section}{อุทฺทานคาถา}{mtk.69}
\bookmark[dest={mtk.69},level=1]{อุทฺทานคาถา}
\csromanheader{2}{1. วิวาทสุตฺต}
\proseitem{41}{อถ โข อายสฺมา อุปาลิ เยน ภควา เตนุปสงฺกมิ, อุปสงฺกมิตฺวา ภควนฺตํ อภิวาเทตฺวา เอกมนฺตํ นิสีทิ, เอกมนฺตํ นิสินฺโน}

\vspace{\tipitakaparskip}
\csromancenter{อกฺโกสวคฺค โข อายสฺมา อุปาลิ ภควนฺตํ เอตทโวจ “โก นุ โข ภนฺเต เหตุ โก ปจฺจโย, เยน สํเฆ ภณฺฑนกลหวิคฺคหวิวาทา อุปฺปชฺชนฺติ, ภิกฺขู จ น ผาสุ\footnote{ผาสุํ (?)} วิหรนฺตี”ติ.\csromanspacer{} อิธุปาลิ ภิกฺขู อธมฺมํ “ธมฺโม”ติ ทีเปนฺติ, ธมฺมํ “อธมฺโม”ติ ทีเปนฺติ, อวินยํ “วินโย”ติ ทีเปนฺติ, วินยํ “อวินโย”ติ ทีเปนฺติ, อภาสิตํ อลปิตํ ตถาคเตน “ภาสิตํ ลปิตํ ตถาคเตนา”ติ ทีเปนฺติ, ภาสิตํ ลปิตํ ตถาคเตน “อภาสิตํ อลปิตํ ตถาคเตนา”ติ ทีเปนฺติ, อนาจิณฺณํ ตถาคเตน “อาจิณฺณํ ตถาคเตนา”ติ ทีเปนฺติ, อาจิณฺณํ ตถาคเตน “อนาจิณฺณํ ตถาคเตนา”ติ ทีเปนฺติ, อปญฺญตฺตํ ตถาคเตน “ปญฺญตฺตํ ตถาคเตนา”ติ ทีเปนฺติ, ปญฺญตฺตํ ตถาคเตน “อปญฺญตฺตํ ตถาคเตนา”ติ ทีเปนฺติ.\csromanspacer{} อยํ โข อุปาลิ เหตุ อยํ ปจฺจโย, เยน สํเฆ ภณฺฑนกลหวิคฺคหวิวาทา อุปฺปชฺชนฺติ, ภิกฺขู จ น ผาสุ วิหรนฺตีติ.\csromanspacer{}\csromanspacer{}ปฐมํ.}
\csromanheader{2}{2. ปฐมวิวาทมูลสุตฺต}
\proseitem{42}{\csromanfolio{317}กติ นุ โข ภนฺเต วิวาทมูลานีติ.\csromanspacer{} ทส โข อุปาลิ วิวาทมูลานิ.\csromanspacer{} กตมานิ ทส? อิธุปาลิ ภิกฺขู อธมฺมํ “ธมฺโม”ติ ทีเปนฺติ, ธมฺมํ “อธมฺโม”ติ ทีเปนฺติ, อวินยํ “วินโย”ติ ทีเปนฺติ, วินยํ “อวินโย”ติ ทีเปนฺติ, อภาสิตํ อลปิตํ ตถาคเตน “ภาสิตํ ลปิตํ ตถาคเตนา”ติ ทีเปนฺติ, ภาสิตํ ลปิตํ ตถาคเตน “อภาสิตํ อลปิตํ ตถาคเตนา”ติ ทีเปนฺติ, อนาจิณฺณํ ตถาคเตน “อาจิณฺณํ ตถาคเตนา”ติ ทีเปนฺติ, อาจิณฺณํ ตถาคเตน “อนาจิณฺณํ ตถาคเตนา”ติ ทีเปนฺติ, อปญฺญตฺตํ ตถาคเตน “ปญฺญตฺตํ ตถาคเตนา”ติ ทีเปนฺติ, ปญฺญตฺตํ ตถาคเตน “อปญฺญตฺตํ ตถาคเตนา”ติ ทีเปนฺติ.\csromanspacer{} อิมานิ โข อุปาลิ ทส วิวาทมูลานีติ.\csromanspacer{}\csromanspacer{}ทุติยํ.}

\csromanheader{2}{3. ทุติยวิวาทมูลสุตฺต}
\proseitem{43}{กติ นุ โข ภนฺเต วิวาทมูลานีติ.\csromanspacer{} ทส โข อุปาลิ วิวาทมูลานิ.\csromanspacer{} กตมานิ ทส? อิธุปาลิ ภิกฺขู อนาปตฺติํ “อาปตฺตี”ติ ทีเปนฺติ, อาปตฺติํ “อนาปตฺตี”ติ ทีเปนฺติ, ลหุกํ อาปตฺติํ “ครุกาปตฺตี”ติ}

\vspace{\tipitakaparskip}
\csromancenter{ทสกนิปาตปาฬิ ทีเปนฺติ, ครุกํ อาปตฺติํ “ลหุกาปตฺตี”ติ ทีเปนฺติ, ทุฏฺฐุลฺลํ อาปตฺติํ “อทุฏฺฐุลฺลาปตฺตี”ติ ทีเปนฺติ, อทุฏฺฐุลฺลํ อาปตฺติํ “ทุฏฺฐุลฺลาปตฺตี”ติ ทีเปนฺติ, สาวเสสํ อาปตฺติํ “อนวเสสาปตฺตี”ติ ทีเปนฺติ, อนวเสสํ อาปตฺติํ “สาวเสสาปตฺตี”ติ ทีเปนฺติ, สปฺปฏิกมฺมํ อาปตฺติํ “อปฺปฏิกมฺมาปตฺตี”ติ ทีเปนฺติ, อปฺปฏิกมฺมํ อาปตฺติํ “สปฺปฏิกมฺมาปตฺตี”ติ ทีเปนฺติ.\csromanspacer{} อิมานิ โข อุปาลิ ทส วิวาทมูลานีติ.\csromanspacer{}\csromanspacer{}ตติยํ.}
\csromanmatikaanchor{mtk.70}
\csromantocmarknopage{section}{5. อกฺโกสวคฺค}{mtk.70}
\bookmark[dest={mtk.70},level=1]{5. อกฺโกสวคฺค}
\csromanheader{2}{4. กุสินารสุตฺต}
\csromanmatikaanchor{mtk.71}
\csromantocmarkref{subsection}{1. วิวาทสุตฺต}{mtk.71}
\bookmark[dest={mtk.71},level=2]{1. วิวาทสุตฺต}
\proseitem{44}{\csromanfolio{318}เอกํ สมยํ ภควา กุสินารายํ วิหรติ พลิหรเณ วนสณฺเฑ.\csromanspacer{} ตตฺร โข ภควา ภิกฺขู อามนฺเตสิ “ภิกฺขโว”ติ. “ภทนฺเต”ติ เต ภิกฺขู ภควโต ปจฺจสฺโสสุํ.\csromanspacer{} ภควา เอตทโวจ\csromandash{}}

\prose{\csromansymbolfootnote{*}{วิ 4. 436; วิ 5. 328 ปิฏฺเฐสุปิ.}โจทเกน ภิกฺขเว ภิกฺขุนา ปรํ โจเทตุกาเมน ปญฺจ ธมฺเม อชฺฌตฺตํ ปจฺจเวกฺขิตฺวา ปญฺจ ธมฺเม อชฺฌตฺตํ อุปฏฺฐาเปตฺวา ปโร โจเทตพฺโพ.\csromanspacer{} กตเม ปญฺจ ธมฺมา อชฺฌตฺตํ ปจฺจเวกฺขิตพฺพา? โจทเกน ภิกฺขเว ภิกฺขุนา ปรํ โจเทตุกาเมน เอวํ ปจฺจเวกฺขิตพฺพํ “ปริสุทฺธกายสมาจาโร นุ โขมฺหิ, ปริสุทฺเธนมฺหิ กายสมาจาเรน สมนฺนาคโต อจฺฉิทฺเทน อปฺปฏิมํเสน, สํวิชฺชติ นุ โข เม เอโส ธมฺโม อุทาหุ โน”ติ.\csromanspacer{} โน เจ ภิกฺขเว ภิกฺขุ ปริสุทฺธกายสมาจาโร โหติ ปริสุทฺเธน กายสมาจาเรน สมนฺนาคโต อจฺฉิทฺเทน อปฺปฏิมํเสน, ตสฺส ภวนฺติ วตฺตาโร “อิงฺฆ ตาว อายสฺมา กายิกํ สิกฺขสฺสู”ติ.\csromanspacer{} อิติสฺส ภวนฺติ วตฺตาโร.}

\prose{ปุน จปรํ ภิกฺขเว โจทเกน ภิกฺขุนา ปรํ โจเทตุกาเมน เอวํ ปจฺจเวกฺขิตพฺพํ “ปริสุทฺธวจีสมาจาโร นุ โขมฺหิ, ปริสุทฺเธนมฺหิ วจีสมาจาเรน สมนฺนาคโต อจฺฉิทฺเทน อปฺปฏิมํเสน, สํวิชฺชติ นุ โข เม เอโส ธมฺโม อุทาหุ โน”ติ.\csromanspacer{} โน เจ ภิกฺขเว ภิกฺขุ ปริสุทฺธวจีสมาจาโร โหติ ปริสุทฺเธน วจีสมาจาเรน สมนฺนาคโต อจฺฉิทฺเทน อปฺปฏิมํเสน, ตสฺส ภวนฺติ วตฺตาโร “อิงฺฆ ตาว อายสฺมา วาจสิกํ สิกฺขสฺสู”ติ.\csromanspacer{} อิติสฺส ภวนฺติ วตฺตาโร.}

\csromanmatikaanchor{mtk.72}
\csromantocmarkref{subsection}{2. ปฐมวิวาทมูลสุตฺต}{mtk.72}
\bookmark[dest={mtk.72},level=2]{2. ปฐมวิวาทมูลสุตฺต}
\prose{ปุน จปรํ ภิกฺขเว โจทเกน ภิกฺขุนา ปรํ โจเทตุกาเมน เอวํ ปจฺจเวกฺขิตพฺพํ “เมตฺตํ นุ โข เม จิตฺตํ ปจฺจุปฏฺฐิตํ สพฺรหฺมจารีสุ อนาฆาตํ สํวิชฺชติ นุ โข เม เอโส ธมฺโม อุทาหุ โน”ติ.\csromanspacer{} โน เจ ภิกฺขเว}

\vspace{\tipitakaparskip}
\csromancenter{อกฺโกสวคฺค ภิกฺขุโน เมตฺตํ จิตฺตํ ปจฺจุปฏฺฐิตํ โหติ สพฺรหฺมจารีสุ อนาฆาตํ, ตสฺส ภวนฺติ วตฺตาโร “อิงฺฆ ตาว อายสฺมา สพฺรหฺมจารีสุ เมตฺตํ จิตฺตํ อุปฏฺฐาเปหี”ติ.\csromanspacer{} อิติสฺส ภวนฺติ วตฺตาโร.}
\csromanmatikaanchor{mtk.73}
\csromantocmarkref{subsection}{3. ทุติยวิวาทมูลสุตฺต}{mtk.73}
\bookmark[dest={mtk.73},level=2]{3. ทุติยวิวาทมูลสุตฺต}
\prose{\csromanfolio{319}ปุน จปรํ ภิกฺขเว โจทเกน ภิกฺขุนา ปรํ โจเทตุกาเมน เอวํ ปจฺจเวกฺขิตพฺพํ “พหุสฺสุโต นุ โขมฺหิ สุตธโร สุตสนฺนิจโย, เย เต ธมฺมา อาทิกลฺยาณา มชฺเฌกลฺยาณา ปริโยสานกลฺยาณา สาตฺถํ สพฺยญฺชนํ เกวลปริปุณฺณํ ปริสุทฺธํ พฺรหฺมจริยํ อภิวทนฺติ, ตถารูปา เม ธมฺมา พหุสฺสุตา โหนฺติ ธาตา วจสา ปริจิตา มนสานุเปกฺขิตา ทิฏฺฐิยา สุปฺปฏิวิทฺธา, สํวิชฺชติ นุ โข เม เอโส ธมฺโม อุทาหุ โน”ติ.\csromanspacer{} โน เจ ภิกฺขเว ภิกฺขุ พหุสฺสุโต โหติ สุตธโร สุตสนฺนิจโย.\csromanspacer{} เย เต ธมฺมา อาทิกลฺยาณา มชฺเฌกลฺยาณา ปริโยสานกลฺยาณา สาตฺถํ สพฺยญฺชนํ เกวลปริปุณฺณํ ปริสุทฺธํ พฺรหฺมจริยํ อภิวทนฺติ, ตถารูปาสฺส ธมฺมา พหุสฺสุตา โหนฺติ ธาตา วจสา ปริจิตา มนสานุเปกฺขิตา ทิฏฺฐิยา สุปฺปฏิวิทฺธา, ตสฺส ภวนฺติ วตฺตาโร “อิงฺฆ ตาว อายสฺมา อาคมํ ปริยาปุณสฺสู”ติ.\csromanspacer{} อิติสฺส ภวนฺติ วตฺตาโร.}

\prose{ปุน จปรํ ภิกฺขเว โจทเกน ภิกฺขุนา ปรํ โจเทตุกาเมน เอวํ ปจฺจเวกฺขิตพฺพํ “อุภยานิ โข ปน เม ปาติโมกฺขานิ วิตฺถาเรน สฺวาคตานิ โหนฺติ สุวิภตฺตานิ สุปฺปวตฺตีนิ สุวินิจฺฉิตานิ สุตฺตโส อนุพฺยญฺชนโส, สํวิชฺชติ นุ โข เม เอโส ธมฺโม อุทาหุ โน”ติ.\csromanspacer{} โน เจ ภิกฺขเว ภิกฺขุโน อุภยานิ ปาติโมกฺขานิ วิตฺถาเรน สฺวาคตานิ โหนฺติ สุวิภตฺตานิ สุปฺปวตฺตีนิ สุวินิจฺฉิตานิ สุตฺตโส อนุพฺยญฺชนโส.\csromanspacer{} อิทํ ปนายสฺมา “กตฺถ วุตฺตํ ภควตา”ติ อิติ ปุฏฺโฐ น สมฺปายิสฺสติ, ตสฺส ภวนฺติ วตฺตาโร “อิงฺฆ ตาว อายสฺมา วินยํ สิกฺขสฺสู”ติ.\csromanspacer{} อิติสฺส ภวนฺติ วตฺตาโร.\csromanspacer{} อิเม ปญฺจ ธมฺมา อชฺฌตฺตํ ปจฺจเวกฺขิตพฺพา.}

\prose{กตเม ปญฺจ ธมฺมา อชฺฌตฺตํ อุปฏฺฐาเปตพฺพา? กาเลน วกฺขามิ โน อกาเลน, ภูเตน วกฺขามิ โน อภูเตน, สเณฺหน วกฺขามิ โน ผรุเสน, อตฺถสํหิเตน วกฺขามิ โน อนตฺถสํหิเตน, เมตฺตจิตฺโต วกฺขามิ โน โทสนฺตโรติ.\csromanspacer{} อิเม ปญฺจ ธมฺมา อชฺฌตฺตํ อุปฏฺฐาเปตพฺพา.\csromanspacer{} โจทเกน ภิกฺขเว ภิกฺขุนา ปรํ โจเทตุกาเมน อิเม ปญฺจ ธมฺเม อชฺฌตฺตํ ปจฺจเวกฺขิตฺวา อิเม ปญฺจ ธมฺเม อชฺฌตฺตํ อุปฏฺฐาเปตฺวา ปโร โจเทตพฺโพติ.\csromanspacer{}\csromanspacer{}จตุตฺถํ.}

\csromanmatikaanchor{mtk.74}
\csromantocmarkref{subsection}{4. กุสินารสุตฺต}{mtk.74}
\bookmark[dest={mtk.74},level=2]{4. กุสินารสุตฺต}
\csromanheader{2}{ทสกนิปาตปาฬิ \textbf{5. ราชนฺเตปุรปฺปเวสนสุตฺต}}
\proseitem{45}{\csromanfolio{320}\csromansymbolfootnote{*}{วิ 2. 207 ปิฏฺเฐปิ.}ทสยิเม ภิกฺขเว อาทีนวา ราชนฺเตปุรปฺปเวสเน.\csromanspacer{} กตเม ทส? อิธ ภิกฺขเว ราชา มเหสิยา สทฺธิํ นิสินฺโน โหติ, ตตฺร ภิกฺขุ ปวิสติ, มเหสี วา ภิกฺขุํ ทิสฺวา สิตํ ปาตุ กโรติ, ภิกฺขุ วา มเหสิํ ทิสฺวา สิตํ ปาตุ กโรติ.\csromanspacer{} ตตฺถ รญฺโญ เอวํ โหติ “อทฺธา อิเมสํ กตํ วา กริสฺสนฺติ วา”ติ.\csromanspacer{} อยํ ภิกฺขเว ปฐโม อาทีนโว ราชนฺเตปุรปฺปเวสเน.}

\prose{ปุน จปรํ ภิกฺขเว ราชา พหุกิจฺโจ พหุกรณีโย อญฺญตรํ อิตฺถิํ คนฺตฺวา น สรติ, สา เตน คพฺภํ คณฺหาติ.\csromanspacer{} ตตฺถ รญฺโญ เอวํ โหติ “น โข อิธ อญฺโญ โกจิ ปวิสติ อญฺญตฺร ปพฺพชิเตน, สิยา นุ โข ปพฺพชิตสฺส กมฺมนฺ”ติ.\csromanspacer{} อยํ ภิกฺขเว ทุติโย อาทีนโว ราชนฺเตปุรปฺปเวสเน.}

\prose{ปุน จปรํ ภิกฺขเว รญฺโญ อนฺเตปุเร อญฺญตรํ รตนํ นสฺสติ.\csromanspacer{} ตตฺถ รญฺโญ เอวํ โหติ “น โข อิธ อญฺโญ โกจิ ปวิสติ อญฺญตฺร ปพฺพชิเตน, สิยา นุ โข ปพฺพชิตสฺส กมฺมนฺ”ติ.\csromanspacer{} อยํ ภิกฺขเว ตติโย อาทีนโว ราชนฺเตปุรปฺปเวสเน.}

\prose{ปุน จปรํ ภิกฺขเว รญฺโญ อนฺเตปุเร อพฺภนฺตรา คุยฺหมนฺตา พหิทฺธา สมฺเภทํ คจฺฉนฺติ.\csromanspacer{} ตตฺถ รญฺโญ เอวํ โหติ “น โข อิธ อญฺโญ โกจิ ปวิสติ อญฺญตฺร ปพฺพชิเตน, สิยา นุ โข ปพฺพชิตสฺส กมฺมนฺ”ติ.\csromanspacer{} อยํ ภิกฺขเว จตุตฺโถ อาทีนโว ราชนฺเตปุรปฺปเวสเน.}

\prose{ปุน จปรํ ภิกฺขเว รญฺโญ อนฺเตปุเร ปิตา วา ปุตฺตํ ปตฺเถติ, ปุตฺโต วา ปิตรํ ปตฺเถติ.\csromanspacer{} เตสํ เอวํ โหติ “น โข อิธ อญฺโญ โกจิ ปวิสติ อญฺญตฺร ปพฺพชิเตน, สิยา นุ โข ปพฺพชิตสฺส กมฺมนฺ”ติ.\csromanspacer{} อยํ ภิกฺขเว ปญฺจโม อาทีนโว ราชนฺเตปุรปฺปเวสเน.}

\prose{ปุน จปรํ ภิกฺขเว ราชา นีจฏฺฐานิยํ อุจฺเจ ฐาเน ฐเปติ.\csromanspacer{} เยสํ ตํ อมนาปํ, เตสํ เอวํ โหติ “ราชา โข ปพฺพชิเตน สํสฏฺโฐ, สิยา นุ โข ปพฺพชิตสฺส กมฺมนฺ”ติ.\csromanspacer{} อยํ ภิกฺขเว ฉฏฺโฐ อาทีนโว ราชนฺเตปุรปฺปเวสเน.}

\chapterheadpageread{5. อกฺโกสวคฺค}
\prose{\csromanfolio{321}ปุน จปรํ ภิกฺขเว ราชา อุจฺจฏฺฐานิยํ นีเจ ฐาเน ฐเปติ.\csromanspacer{} เยสํ ตํ อมนาปํ, เตสํ เอวํ โหติ “ราชา โข ปพฺพชิเตน สํสฏฺโฐ, สิยา นุ โข ปพฺพชิตสฺส กมฺมนฺ”ติ.\csromanspacer{} อยํ ภิกฺขเว สตฺตโม อาทีนโว ราชนฺเตปุรปฺปเวสเน.}

\csromanmatikaanchor{mtk.75}
\csromantocmarkref{subsection}{5. ราชนฺเตปุรปฺปเวสนสุตฺต}{mtk.75}
\bookmark[dest={mtk.75},level=2]{5. ราชนฺเตปุรปฺปเวสนสุตฺต}
\prose{ปุน จปรํ ภิกฺขเว ราชา อกาเล เสนํ อุยฺโยเชติ.\csromanspacer{} เยสํ ตํ อมนาปํ, เตสํ เอวํ โหติ “ราชา โข ปพฺพชิเตน สํสฏฺโฐ, สิยา นุ โข ปพฺพชิตสฺส กมฺมนฺ”ติ.\csromanspacer{} อยํ ภิกฺขเว อฏฺฐโม อาทีนโว ราชนฺเตปุรปฺปเวสเน.}

\prose{ปุน จปรํ ภิกฺขเว ราชา กาเล เสนํ อุยฺโยเชตฺวา อนฺตรามคฺคโต นิวตฺตาเปติ.\csromanspacer{} เยสํ ตํ อมนาปํ, เตสํ เอวํ โหติ “ราชา โข ปพฺพชิเตน สํสฏฺโฐ, สิยา นุ โข ปพฺพชิตสฺส กมฺมนฺ”ติ.\csromanspacer{} อยํ ภิกฺขเว นวโม อาทีนโว ราชนฺเตปุรปฺปเวสเน.}

\prose{ปุน จปรํ ภิกฺขเว รญฺโญ อนฺเตปุรํ หตฺถิสมฺมทฺทํ อสฺสสมฺมทฺทํ รถสมฺมทฺทํ รชนียานิ รูปสทฺทคนฺธรสโผฏฺฐพฺพานิ.\csromanspacer{} ยานิ น ปพฺพชิตสฺส สารุปฺปานิ.\csromanspacer{} อยํ ภิกฺขเว ทสโม อาทีนโว ราชนฺเตปุรปฺปเวสเน.\csromanspacer{} อิเม โข ภิกฺขเว ทส อาทีนวา ราชนฺเตปุรปฺปเวสเนติ.\csromanspacer{}\csromanspacer{}ปญฺจมํ.}

\csromanheader{2}{6. สกฺกสุตฺต}
\proseitem{46}{เอกํ สมยํ ภควา สกฺเกสุ วิหรติ กปิลวตฺถุสฺมิํ นิโคฺรธาราเม.\csromanspacer{} อถ โข สมฺพหุลา สกฺกา อุปาสกา ตทหุโปสเถ เยน ภควา เตนุปสงฺกมิํสุ, อุปสงฺกมิตฺวา ภควนฺตํ อภิวาเทตฺวา เอกมนฺตํ นิสีทิํสุ, เอกมนฺตํ นิสินฺเน โข สกฺเก อุปาสเก ภควา เอตทโวจ “อปิ นุ ตุเมฺห สกฺกา อฏฺฐงฺคสมนฺนาคตํ อุโปสถํ อุปวสถา”ติ.\csromanspacer{} อปฺเปกทา มยํ ภนฺเต อฏฺฐงฺคสมนฺนาคตํ อุโปสถํ อุปวสาม, อปฺเปกทา น อุปวสามาติ.\csromanspacer{} เตสํ โว สกฺกา อลาภา เตสํ ทุลฺลทฺธํ, เย ตุเมฺห เอวํ โสกสภเย ชีวิเต มรณสภเย ชีวิเต อปฺเปกทา อฏฺฐงฺคสมนฺนาคตํ อุโปสถํ อุปวสถ, อปฺเปกทา น อุปวสถ.}

\prose{ตํ กิํ มญฺญถ สกฺกา, อิธ ปุริโส เยน เกนจิ กมฺมฏฺฐาเนน อนาปชฺช อกุสลํ ทิวสํ อฑฺฒกหาปณํ นิพฺพิเสยฺย, “ทกฺโข ปุริโส อุฏฺฐานสมฺปนฺโน”ติ อลํ วจนายาติ.\csromanspacer{} เอวํ ภนฺเต.}

\gambhira{ทสกนิปาตปาฬิ}
\prose{\csromanfolio{322}ตํ กิํ มญฺญถ สกฺกา, อิธ ปุริโส เยน เกนจิ กมฺมฏฺฐาเนน อนาปชฺช อกุสลํ ทิวสํ กหาปณํ นิพฺพิเสยฺย, “ทกฺโข ปุริโส อุฏฺฐานสมฺปนฺโน”ติ อลํ วจนายาติ.\csromanspacer{} เอวํ ภนฺเต.}

\prose{ตํ กิํ มญฺญถ สกฺกา, อิธ ปุริโส เยน เกนจิ กมฺมฏฺฐาเนน อนาปชฺช อกุสลํ ทิวสํ เทฺว กหาปเณ นิพฺพิเสยฺย.\csromanspacer{} ตโย กหาปเณ นิพฺพิเสยฺย.\csromanspacer{} จตฺตาโร กหาปเณ นิพฺพิเสยฺย.\csromanspacer{} ปญฺจ กหาปเณ นิพฺพิเสยฺย.\csromanspacer{} ฉ กหาปเณ นิพฺพิเสยฺย.\csromanspacer{} สตฺต กหาปเณ นิพฺพิเสยฺย.\csromanspacer{} อฏฺฐ กหาปเณ นิพฺพิเสยฺย.\csromanspacer{} นว กหาปเณ นิพฺพิเสยฺย.\csromanspacer{} ทส กหาปเณ นิพฺพิเสยฺย.\csromanspacer{} วีส กหาปเณ นิพฺพิเสยฺย.\csromanspacer{} ติํส กหาปเณ นิพฺพิเสยฺย.\csromanspacer{} จตฺตารีสํ กหาปเณ นิพฺพิเสยฺย.\csromanspacer{} ปญฺญาสํ กหาปเณ นิพฺเพเสยฺย.\csromanspacer{} กหาปณสตํ นิพฺพิเสยฺย, “ทกฺโข ปุริโส อุฏฺฐานสมฺปนฺโน”ติ อลํ วจนายาติ.\csromanspacer{} เอวํ ภนฺเต.}

\prose{ตํ กิํ มญฺญถ สกฺกา, อปิ นุ โส ปุริโส ทิวเส ทิวเส กหาปณสตํ กหาปณสหสฺสํ นิพฺพิสมาโน ลทฺธํ ลทฺธํ นิกฺขิปนฺโต วสฺสสตายุโก วสฺสสตชีวี มหนฺตํ โภคกฺขนฺธํ อธิคจฺเฉยฺยาติ.\csromanspacer{} เอวํ ภนฺเต.}

\prose{ตํ กิํ มญฺญถ สกฺกา, อปิ นุ โส ปุริโส โภคเหตุ โภคนิทานํ โภคาธิกรณํ เอกํ วา รตฺติํ เอกํ วา ทิวสํ อุปฑฺฒํ วา รตฺติํ อุปฑฺฒํ วา ทิวสํ เอกนฺตสุขปฺปฏิสํเวที วิหเรยฺยาติ.\csromanspacer{} โน เหตํ ภนฺเต.\csromanspacer{} ตํ กิสฺส เหตุ? กามา หิ ภนฺเต อนิจฺจา ตุจฺฉา มุสา โมสธมฺมาติ.}

\prose{อิธ ปน โว สกฺกา มม สาวโก ทส วสฺสานิ อปฺปมตฺโต อาตาปี ปหิตตฺโต วิหรนฺโต ยถา มยานุสิฏฺฐํ, ตถา ปฏิปชฺชมาโน สตมฺปิ วสฺสานิ สตมฺปิ วสฺสสตานิ สตมฺปิ วสฺสสหสฺสานิ เอกนฺตสุขปฺปฏิสํเวที วิหเรยฺย, โส จ ขฺวสฺส สกทาคามี วา อนาคามี วา อปณฺณกํ วา โสตาปนฺโน.\csromanspacer{} ติฏฺฐนฺตุ สกฺกา ทส วสฺสานิ.}

\csromanmatikaanchor{mtk.76}
\csromantocmarkref{subsection}{6. สกฺกสุตฺต}{mtk.76}
\bookmark[dest={mtk.76},level=2]{6. สกฺกสุตฺต}
\prose{อิธ มม สาวโก นว วสฺสานิ.\csromanspacer{} อฏฺฐ วสฺสานิ.\csromanspacer{} สตฺต วสฺสานิ.\csromanspacer{} ฉ วสฺสานิ.\csromanspacer{} ปญฺจ วสฺสานิ.\csromanspacer{} จตฺตาริ วสฺสานิ.\csromanspacer{} ตีณิ วสฺสานิ.\csromanspacer{} เทฺว วสฺสานิ.\csromanspacer{} เอกํ วสฺสํ อปฺปมตฺโต อาตาปี ปหิตตฺโต วิหรนฺโต ยถา มยานุสิฏฺฐํ, ตถา ปฏิปชฺชมาโน สตมฺปิ วสฺสานิ สตมฺปิ วสฺสสตานิ สตมฺปิ วสฺสสหสฺสานิ เอกนฺตสุขปฺปฏิสํเวที วิหเรยฺย, โส จ ขฺวสฺส}

\vspace{\tipitakaparskip}
\csromancenter{อกฺโกสวคฺค สกทาคามี วา อนาคามี วา อปณฺณกํ วา โสตาปนฺโน.\csromanspacer{} ติฏฺฐตุ สกฺกา เอกํ วสฺสํ.}
\prose{\csromanfolio{323}อิธ มม สาวโก ทส มาเส อปฺปมตฺโต อาตาปี ปหิตตฺโต วิหรนฺโต ยถา มยานุสิฏฺฐํ, ตถา ปฏิปชฺชมาโน สตมฺปิ วสฺสานิ สตมฺปิ วสฺสสตานิ สตมฺปิ วสฺสสหสฺสานิ เอกนฺตสุขปฺปฏิสํเวที วิหเรยฺย, โส จ ขฺวสฺส สกทาคามี วา อนาคามี วา อปณฺณกํ วา โสตาปนฺโน.\csromanspacer{} ติฏฺฐนฺตุ สกฺกา ทส มาสา.}

\prose{อิธ มม สาวโก นว มาเส.\csromanspacer{} อฏฺฐ มาเส.\csromanspacer{} สตฺต มาเส.\csromanspacer{} ฉ มาเส.\csromanspacer{} ปญฺจ มาเส.\csromanspacer{} จตฺตาโร มาเส.\csromanspacer{} ตโย มาเส.\csromanspacer{} เทฺว มาเส.\csromanspacer{} เอกํ มาสํ.\csromanspacer{} อฑฺฒมาสํ อปฺปมตฺโต อาตาปี ปหิตตฺโต วิหรนฺโต ยถา มยานุสิฏฺฐํ, ตถา ปฏิปชฺชมาโน สตมฺปิ วสฺสานิ สตมฺปิ วสฺสสตานิ สตมฺปิ วสฺสสหสฺสานิ เอกนฺตสุขปฺปฏิสํเวที วิหเรยฺย, โส จ ขฺวสฺส สกทาคามี วา อนาคามี วา อปณฺณกํ วา โสตาปนฺโน, ติฏฺฐตุ สกฺกา อฑฺฒมาโส.}

\prose{อิธ มม สาวโก ทส รตฺตินฺทิเว\footnote{รตฺติทิเว (ก)} อปฺปมตฺโต อาตาปี ปหิตตฺโต วิหรนฺโต ยถา มยานุสิฏฺฐํ, ตถา ปฏิปชฺชมาโน สตมฺปิ วสฺสานิ สตมฺปิ วสฺสสตานิ สตมฺปิ วสฺสสหสฺสานิ เอกนฺตสุขปฺปฏิสํเวที วิหเรยฺย, โส จ ขฺวสฺส สกทาคามี วา อนาคามี วา อปณฺณกํ วา โสตาปนฺโน.\csromanspacer{} ติฏฺฐนฺตุ สกฺกา ทส รตฺตินฺทิวา.}

\prose{อิธ มม สาวโก นว รตฺตินฺทิเว.\csromanspacer{} อฏฺฐ รตฺตินฺทิเว.\csromanspacer{} สตฺต รตฺตินฺทิเว.\csromanspacer{} ฉ รตฺตินฺทิเว.\csromanspacer{} ปญฺจ รตฺตินฺทิเว.\csromanspacer{} จตฺตาโร รตฺตินฺทิเว.\csromanspacer{} ตโย รตฺตินฺทิเว.\csromanspacer{} เทฺว รตฺตินฺทิเว.\csromanspacer{} เอกํ รตฺตินฺทิวํ อปฺปมตฺโต อาตาปี ปหิตตฺโต วิหรนฺโต ยถา มยานุสิฏฺฐํ, ตถา ปฏิปชฺชมาโน สตมฺปิ วสฺสานิ สตมฺปิ วสฺสสตานิ สตมฺปิ วสฺสสหสฺสานิ เอกนฺตสุขปฺปฏิสํเวที วิหเรยฺย, โส จ ขฺวสฺส สกทาคามี วา อนาคามี วา อปณฺณกํ วา โสตาปนฺโน.\csromanspacer{} เตสํ โว สกฺกา อลาภา เตสํ ทุลฺลทฺธํ, เย ตุเมฺห เอวํ โสกสภเย ชีวิเต มรณสภเย ชีวิเต อปฺเปกทา อฏฺฐงฺคสมนฺนาคตํ อุโปสถํ อุปวสถ, อปฺเปกทา น อุปวสถาติ.\csromanspacer{} เอเต มยํ ภนฺเต อชฺชตคฺเค อฏฺฐงฺคสมนฺนาคตํ อุโปสถํ อุปวสิสฺสามาติ.\csromanspacer{}\csromanspacer{}ฉฏฺฐํ.}

\csromanheader{2}{ทสกนิปาตปาฬิ \textbf{7. มหาลิสุตฺต}}
\proseitem{47}{\csromanfolio{324}เอกํ สมยํ ภควา เวสาลิยํ วิหรติ มหาวเน กูฏาคารสาลายํ.\csromanspacer{} อถ โข มหาลิ ลิจฺฉวิ เยน ภควา เตนุปสงฺกมิ, อุปสงฺกมิตฺวา ภควนฺตํ อภิวาเทตฺวา เอกมนฺตํ นิสีทิ, เอกมนฺตํ นิสินฺโน โข มหาลิ ลิจฺฉวิ ภควนฺตํ เอตทโวจ “โก นุ โข ภนฺเต เหตุ, โก ปจฺจโย ปาปสฺส กมฺมสฺส กิริยาย ปาปสฺส กมฺมสฺส ปวตฺติยา”ติ.\csromanspacer{} โลโภ โข มหาลิ เหตุ, โลโภ ปจฺจโย ปาปสฺส กมฺมสฺส กิริยาย ปาปสฺส กมฺมสฺส ปวตฺติยา, โทโส โข มหาลิ เหตุ, โทโส ปจฺจโย ปาปสฺส กมฺมสฺส กิริยาย ปาปสฺส กมฺมสฺส ปวตฺติยา, โมโห โข มหาลิ เหตุ, โมโห ปจฺจโย ปาปสฺส กมฺมสฺส กิริยาย ปาปสฺส กมฺมสฺส ปวตฺติยา, อโยนิโส มนสิกาโร โข มหาลิ เหตุ, อโยนิโส มนสิกาโร ปจฺจโย ปาปสฺส กมฺมสฺส กิริยาย ปาปสฺส กมฺมสฺส ปวตฺติยา, มิจฺฉาปณิหิตํ โข มหาลิ จิตฺตํ เหตุ, มิจฺฉาปณิหิตํ จิตฺตํ ปจฺจโย ปาปสฺส กมฺมสฺส กิริยาย ปาปสฺส กมฺมสฺส ปวตฺติยาติ.\csromanspacer{} อยํ โข มหาลิ เหตุ, อยํ ปจฺจโย ปาปสฺส กมฺมสฺส กิริยาย ปาปสฺส กมฺมสฺส ปวตฺติยาติ.}

\prose{โก ปน ภนฺเต เหตุ, โก ปจฺจโย กลฺยาณสฺส กมฺมสฺส กิริยาย กลฺยาณสฺส กมฺมสฺส ปวตฺติยาติ.\csromanspacer{} อโลโภ โข มหาลิ เหตุ, อโลโภ ปจฺจโย กลฺยาณสฺส กมฺมสฺส กิริยาย กลฺยาณสฺส กมฺมสฺส ปวตฺติยา, อโทโส โข มหาลิ เหตุ, อโทโส ปจฺจโย กลฺยาณสฺส กมฺมสฺส กิริยาย กลฺยาณสฺส กมฺมสฺส ปวตฺติยา, อโมโห โข มหาลิ เหตุ; อโมโห ปจฺจโย กลฺยาณสฺส กมฺมสฺส กิริยาย กลฺยาณสฺส กมฺมสฺส ปวตฺติยา, โยนิโส มนสิกาโร โข มหาลิ เหตุ, โยนิโส มนสิกาโร ปจฺจโย กลฺยาณสฺส กมฺมสฺส กิริยาย กลฺยาณสฺส กมฺมสฺส ปวตฺติยา, สมฺมาปณิหิตํ โข มหาลิ จิตฺตํ เหตุ, สมฺมาปณิหิตํ จิตฺตํ ปจฺจโย กลฺยาณสฺส กมฺมสฺส กิริยาย กลฺยาณสฺส กมฺมสฺส ปวตฺติยา.\csromanspacer{} อยํ โข มหาลิ เหตุ, อยํ ปจฺจโย กลฺยาณสฺส กมฺมสฺส กิริยาย กลฺยาณสฺส กมฺมสฺส ปวตฺติยา.\csromanspacer{} อิเม จ มหาลิ ทส ธมฺมา โลเก น สํวิชฺเชยฺยุํ, นยิธ ปญฺญาเยถ “อธมฺมจริยาวิสมจริยา”ติ วา “ธมฺมจริยาสมจริยา”ติ วา.\csromanspacer{} ยสฺมา จ โข มหาลิ อิเม ทส ธมฺมา}

\vspace{\tipitakaparskip}
\csromancenter{อกฺโกสวคฺค โลเก สํวิชฺชนฺติ, ตสฺมา ปญฺญายติ “อธมฺมจริยาวิสมจริยา”ติ วา “ธมฺมจริยาสมจริยา”ติ วาติ.\csromanspacer{}\csromanspacer{}สตฺตมํ.}
\csromanheader{2}{8. ปพฺพชิต-อภิณฺหสุตฺต}
\proseitem{48}{\csromanfolio{325}ทสยิเม ภิกฺขเว ธมฺมา ปพฺพชิเตน อภิณฺหํ ปจฺจเวกฺขิตพฺพา.\csromanspacer{} กตเม ทส? “เววณฺณิยมฺหิ อชฺฌุปคโต”ติ ปพฺพชิเตน อภิณฺหํ ปจฺจเวกฺขิตพฺพํ. “ปรปฏิพทฺธา เม ชีวิกา”ติ ปพฺพชิเตน อภิณฺหํ ปจฺจเวกฺขิตพฺพํ. “อญฺโญ เม อากปฺโป กรณีโย”ติ ปพฺพชิเตน อภิณฺหํ ปจฺจเวกฺขิตพฺพํ. “กจฺจิ นุ โข เม อตฺตา สีลโต น อุปวทตี”ติ ปพฺพชิเตน อภิณฺหํ ปจฺจเวกฺขิตพฺพํ. “กจฺจิ นุ โข มํ อนุวิจฺจ วิญฺญู สพฺรหฺมจารี สีลโต น อุปวทนฺตี”ติ ปพฺพชิเตน อภิณฺหํ ปจฺจเวกฺขิตพฺพํ. “สพฺเพหิ เม ปิเยหิ มนาเปหิ นานาภาโว วินาภาโว”ติ ปพฺพชิเตน อภิณฺหํ ปจฺจเวกฺขิตพฺพํ. “กมฺมสฺสโกมฺหิ กมฺมทายาโท กมฺมโยนิ กมฺมพนฺธุ กมฺมปฏิสรโณ, ยํ กมฺมํ กริสฺสามิ กลฺยาณํ วา ปาปกํ วา, ตสฺส ทายาโท ภวิสฺสามี”ติ ปพฺพชิเตน อภิณฺหํ ปจฺจเวกฺขิตพฺพํ. “กถํภูตสฺส เม รตฺตินฺทิวา วีติวตฺตนฺตี”ติ ปพฺพชิเตน อภิณฺหํ ปจฺจเวกฺขิตพฺพํ. “กจฺจิ นุ โข อหํ สุญฺญาคาเร อภิรมามี”ติ ปพฺพชิเตน อภิณฺหํ ปจฺจเวกฺขิตพฺพํ. “อตฺถิ นุ โข เม อุตฺตริ มนุสฺสธมฺโม อลมริยญาณทสฺสนวิเสโส อธิคโต, เยนาหํ\footnote{โยหํ (สี, อิ, ก), โสหํ (สฺยา)} ปจฺฉิเม กาเล สพฺรหฺมจารีหิ ปุฏฺโฐ น มงฺกุ ภวิสฺสามี”ติ ปพฺพชิเตน อภิณฺหํ ปจฺจเวกฺขิตพฺพํ.\csromanspacer{} อิเม โข ภิกฺขเว ทส ธมฺมา ปพฺพชิเตน อภิณฺหํ ปจฺจเวกฺขิตพฺพาติ.\csromanspacer{}\csromanspacer{}อฏฺฐมํ.}

\csromanheader{2}{9. สรีรฏฺฐธมฺมสุตฺต}
\proseitem{49}{ทสยิเม ภิกฺขเว ธมฺมา สรีรฏฺฐา.\csromanspacer{} กตเม ทส? สีตํ อุณฺหํ ชิฆจฺฉา ปิปาสา อุจฺจาโร ปสฺสาโว กายสํวโร วจีสํวโร อาชีวสํวโร โปโนภวิโก\footnote{โปโนพฺภวิโก (ก)} ภวสงฺขาโร.\csromanspacer{} อิเม โข ภิกฺขเว ทส ธมฺมา สรีรฏฺฐาติ.\csromanspacer{}\csromanspacer{}นวมํ.}

\csromanheader{2}{10. ภณฺฑนสุตฺต}
\csromanmatikaanchor{mtk.77}
\csromantocmarkref{subsection}{7. มหาลิสุตฺต}{mtk.77}
\bookmark[dest={mtk.77},level=2]{7. มหาลิสุตฺต}
\proseitem{50}{เอกํ สมยํ ภควา สาวตฺถิยํ วิหรติ เชตวเน อนาถปิณฺฑิกสฺส อาราเม.\csromanspacer{} เตน โข ปน สมเยน สมฺพหุลา ภิกฺขู ปจฺฉาภตฺตํ}

\vspace{\tipitakaparskip}
\csromancenter{ทสกนิปาตปาฬิ ปิณฺฑปาตปฏิกฺกนฺตา อุปฏฺฐานสาลายํ สนฺนิสินฺนา สนฺนิปติตา ภณฺฑนชาตา กลหชาตา วิวาทาปนฺนา อญฺญมญฺญํ มุขสตฺตีหิ วิตุทนฺตา วิหรนฺติ.}
\prose{\csromanfolio{326}อถ โข ภควา สายนฺหสมยํ ปฏิสลฺลานา วุฏฺฐิโต เยน อุปฏฺฐานสาลา เตนุปสงฺกมิ, อุปสงฺกมิตฺวา ปญฺญตฺเต อาสเน นิสีทิ, นิสชฺช โข ภควา ภิกฺขู อามนฺเตสิ “กาย นุตฺถ ภิกฺขเว เอตรหิ กถาย สนฺนิสินฺนา สนฺนิปติตา, กา จ ปน โว อนฺตรากถา วิปฺปกตา”ติ.}

\prose{อิธ มยํ ภนฺเต ปจฺฉาภตฺตํ ปิณฺฑปาตปฏิกฺกนฺตา อุปฏฺฐานสาลายํ สนฺนิสินฺนา สนฺนิปติตา ภณฺฑนชาตา กลหชาตา วิวาทาปนฺนา อญฺญมญฺญํ มุขสตฺตีหิ วิตุทนฺตา วิหรามาติ.\csromanspacer{} น โข ปเนตํ ภิกฺขเว ตุมฺหากํ ปติรูปํ กุลปุตฺตานํ สทฺธาย อคารสฺมา อนคาริยํ ปพฺพชิตานํ, ยํ ตุเมฺห ภณฺฑนชาตา กลหชาตา วิวาทาปนฺนา อญฺญมญฺญํ มุขสตฺตีหิ วิตุทนฺตา วิหเรยฺยาถ.}

\csromanmatikaanchor{mtk.78}
\csromantocmarkref{subsection}{8. ปพฺพชิต-อภิณฺหสุตฺต}{mtk.78}
\bookmark[dest={mtk.78},level=2]{8. ปพฺพชิต-อภิณฺหสุตฺต}
\prose{ทสยิเม ภิกฺขเว ธมฺมา สารณียา ปิยกรณา ครุกรณา\footnote{ปิยกราณา ครุกราณา (?)} สงฺคหาย อวิวาทาย สามคฺคิยา เอกีภาวาย สํวตฺตนฺติ.\csromanspacer{} กตเม ทส? อิธ ภิกฺขเว ภิกฺขุ สีลวา โหติ, ปาติโมกฺขสํวรสํวุโต วิหรติ อาจารโคจรสมฺปนฺโน, อณุมตฺเตสุ วชฺเชสุ ภยทสฺสาวี สมาทาย สิกฺขติ สิกฺขาปเทสุ.\csromanspacer{} ยมฺปิ ภิกฺขเว ภิกฺขุ สีลวา โหติ ฯเปฯ สมาทาย สิกฺขติ สิกฺขาปเทสุ, อยมฺปิ ธมฺโม สารณีโย ปิยกรโณ ครุกรโณ สงฺคหาย อวิวาทาย สามคฺคิยา เอกีภาวาย สํวตฺตติ.}

\prose{ปุน จปรํ ภิกฺขเว ภิกฺขุ พหุสฺสุโต โหติ สุตธโร สุตสนฺนิจโย, เย เต ธมฺมา อาทิกลฺยาณา มชฺเฌกลฺยาณา ปริโยสานกลฺยาณา สาตฺถํ สพฺยญฺชนํ เกวลปริปุณฺณํ ปริสุทฺธํ พฺรหฺมจริยํ อภิวทนฺติ, ตถารูปาสฺส ธมฺมา พหุสฺสุตา โหนฺติ ธาตา วจสา ปริจิตา มนสานุเปกฺขิตา ทิฏฺฐิยา สุปฺปฏิวิทฺธา.\csromanspacer{} ยมฺปิ ภิกฺขเว ภิกฺขุ พหุสฺสุโต โหติ ฯเปฯ ทิฏฺฐิยา สุปฺปฏิวิทฺธา, อยมฺปิ ธมฺโม สารณีโย ปิยกรโณ ครุกรโณ สงฺคหาย อวิวาทาย สามคฺคิยา เอกีภาวาย สํวตฺตติ.}

\csromanmatikaanchor{mtk.79}
\csromantocmarkref{subsection}{9. สรีรฏฺฐธมฺมสุตฺต}{mtk.79}
\bookmark[dest={mtk.79},level=2]{9. สรีรฏฺฐธมฺมสุตฺต}
\chapterheadpageread{5. อกฺโกสวคฺค}
\prose{\csromanfolio{327}ปุน จปรํ ภิกฺขเว ภิกฺขุ กลฺยาณมิตฺโต โหติ กลฺยาณสหาโย กลฺยาณสมฺปวงฺโก.\csromanspacer{} ยมฺปิ ภิกฺขเว ภิกฺขุ กลฺยาณมิตฺโต โหติ กลฺยาณสหาโย กลฺยาณสมฺปวงฺโก, อยมฺปิ ธมฺโม สารณีโย ปิยกรโณ ครุกรโณ สงฺคหาย อวิวาทาย สามคฺคิยา เอกีภาวาย สํวตฺตติ.}

\csromanmatikaanchor{mtk.80}
\csromantocmarkref{subsection}{10. ภณฺฑนสุตฺต}{mtk.80}
\bookmark[dest={mtk.80},level=2]{10. ภณฺฑนสุตฺต}
\prose{ปุน จปรํ ภิกฺขเว ภิกฺขุ สุวโจ โหติ โสวจสฺสกรเณหิ ธมฺเมหิ สมนฺนาคโต, ขโม ปทกฺขิณคฺคาหี อนุสาสนิํ.\csromanspacer{} ยมฺปิ ภิกฺขเว ภิกฺขุ สุวโจ โหติ โสวจสฺสกรเณหิ ธมฺเมหิ สมนฺนาคโต, ขโม ปทกฺขิณคฺคาหี อนุสาสนิํ, อยมฺปิ ธมฺโม สารณีโย ปิยกรโณ ครุกรโณ สงฺคหาย อวิวาทาย สามคฺคิยา เอกีภาวาย สํวตฺตติ.}

\prose{ปุน จปรํ ภิกฺขเว ภิกฺขุ ยานิ ตานิ สพฺรหฺมจารีนํ อุจฺจาวจานิ กิํกรณียานิ, ตตฺถ ทกฺโข โหติ อนลโส, ตตฺรูปายาย วีมํสาย สมนฺนาคโต อลํ กาตุํ อลํ สํวิธาตุํ.\csromanspacer{} ยมฺปิ ภิกฺขเว ภิกฺขุ ยานิ ตานิ สพฺรหฺมจารีนํ อุจฺจาวจานิ กิํกรณียานิ, ตตฺถ ทกฺโข โหติ อนลโส, ตตฺรูปายาย วีมํสาย สมนฺนาคโต อลํ กาตุํ อลํ สํวิธาตุํ, อยมฺปิ ธมฺโม สารณีโย ปิยกรโณ ครุกรโณ สงฺคหาย อวิวาทาย สามคฺคิยา เอกีภาวาย สํวตฺตติ.}

\prose{ปุน จปรํ ภิกฺขเว ภิกฺขุ ธมฺมกาโม โหติ ปิยสมุทาหาโร, อภิธมฺเม อภิวินเย อุฬารปาโมชฺโช.\csromanspacer{} ยมฺปิ ภิกฺขเว ภิกฺขุ ธมฺมกาโม โหติ ปิยสมุทาหาโร, อภิธมฺเม อภิวินเย อุฬารปาโมชฺโช, อยมฺปิ ธมฺโม สารณีโย ปิยกรโณ ครุกรโณ สงฺคหาย อวิวาทาย สามคฺคิยา เอกีภาวาย สํวตฺตติ.}

\prose{ปุน จปรํ ภิกฺขเว ภิกฺขุ อารทฺธวีริโย วิหรติ อกุสลานํ ธมฺมานํ ปหานาย กุสลานํ ธมฺมานํ อุปสมฺปทาย, ถามวา ทฬฺหปรกฺกโม อนิกฺขิตฺตธุโร กุสเลสุ ธมฺเมสุ.\csromanspacer{} ยมฺปิ ภิกฺขเว ภิกฺขุ อารทฺธวีริโย วิหรติ อกุสลานํ ธมฺมานํ ปหานาย กุสลานํ ธมฺมานํ อุปสมฺปทาย, ถามวา ทฬฺหปรกฺกโม อนิกฺขิตฺตธุโร กุสเลสุ ธมฺเมสุ, อยมฺปิ ธมฺโม สารณีโย ปิยกรโณ ครุกรโณ สงฺคหาย อวิวาทาย สามคฺคิยา เอกีภาวาย สํวตฺตติ.}

\gambhira{ทสกนิปาตปาฬิ}
\csromanmatikaanchor{mtk.81}
\csromantocmarkref{section}{อุทฺทานคาถา}{mtk.81}
\bookmark[dest={mtk.81},level=1]{อุทฺทานคาถา}
\csromantocmarknopage{chapter}{2. ทุติยปณฺณาสก}{mtk.82}
\bookmark[dest={mtk.82},level=0]{2. ทุติยปณฺณาสก}
\prose{\csromanfolio{328}ปุน จปรํ ภิกฺขเว ภิกฺขุ สนฺตุฏฺโฐ โหติ อิตรีตรจีวร- ปิณฺฑปาตเสนาสนคิลานปจฺจย เภสชฺช ปริกฺขาเรน.\csromanspacer{} ยมฺปิ ภิกฺขเว ภิกฺขุ สนฺตุฏฺโฐ โหติ อิตรีตรจีวรปิณฺฑปาตเสนาสนคิลานปจฺจย- เภสชฺชปริกฺขาเรน, อยมฺปิ ธมฺโม สารณีโย ฯเปฯ สํวตฺตติ.}

\prose{ปุน จปรํ ภิกฺขเว ภิกฺขุ สติมา โหติ ปรเมน สติเนปกฺเกน สมนฺนาคโต, จิรกตมฺปิ จิรภาสิตมฺปิ สริตา อนุสฺสริตา.\csromanspacer{} ยมฺปิ ภิกฺขเว ภิกฺขุ สติมา โหติ ปรเมน สติเนปกฺเกน สมนฺนาคโต, จิรกตมฺปิ จิรภาสิตมฺปิ สริตา อนุสฺสริตา, อยมฺปิ ธมฺโม สารณีโย ฯเปฯ สํวตฺตติ.}

\prose{ปุน จปรํ ภิกฺขเว ภิกฺขุ ปญฺญวา โหติ อุทยตฺถคามินิยา ปญฺญาย สมนฺนาคโต อริยาย นิพฺเพธิกาย สมฺมา ทุกฺขกฺขยคามินิยา.\csromanspacer{} ยมฺปิ ภิกฺขเว ภิกฺขุ ปญฺญวา โหติ อุทยตฺถคามินิยา, ปญฺญาย สมนฺนาคโต อริยาย นิพฺเพธิกาย สมฺมา ทุกฺขกฺขยคามินิยา, อยมฺปิ ธมฺโม สารณีโย ฯเปฯ สํวตฺตติ.\csromanspacer{} อิเม โข ภิกฺขเว ทส ธมฺมา สารณียา ปิยกรณา ครุกรณา สงฺคหาย อวิวาทาย สามคฺคิยา เอกีภาวาย สํวตฺตนฺตีติ.\csromanspacer{}\csromanspacer{}ทสมํ.\csromanspacer{} อกฺโกสวคฺโค ปญฺจโม.}

\titlehead{\textbf{ตสฺสุทฺทานํ}}
\csromangathameasure{วิวาทา เทฺว จ มูลานิ, กุสินารปเวสเน. \\ สกฺโก มหาลิ อภิณฺหํ, สรีรฏฺฐา จ ภณฺฑนาติ. \\ \textbf{ปฐโม ปณฺณาสโก สมตฺโต.}}
\csromangathagroup{\csromangathastanza{วิวาทา เทฺว จ มูลานิ, กุสินารปเวสเน. \\ สกฺโก มหาลิ อภิณฺหํ, สรีรฏฺฐา จ ภณฺฑนาติ. \\ \textbf{ปฐโม ปณฺณาสโก สมตฺโต.}}}

\csromanheader{2}{2. ทุติยปณฺณาสก}
\csromanheader{1}{\textbf{(6) 1. สจิตฺตวคฺค}}
\csromanheader{2}{1. สจิตฺตสุตฺต}
\proseitem{51}{\csromanfolio{329}เอกํ สมยํ ภควา สาวตฺถิยํ วิหรติ เชตวเน อนาถปิณฺฑิกสฺส อาราเม.\csromanspacer{} ตตฺร โข ภควา ภิกฺขู อามนฺเตสิ “ภิกฺขโว”ติ. “ภทนฺเต”ติ เต สิกฺขู ภควโต ปจฺจสฺโสสุํ.\csromanspacer{} ภควา เอตทโวจ\csromandash{}}

\prose{โน เจ ภิกฺขเว ภิกฺขุ ปรจิตฺตปริยายกุสโล โหติ, อถ “สจิตฺตปริยายกุสโล ภวิสฺสามี”ติ\footnote{ภวิสฺสามาติ (สฺยา)} เอวํ หิ โว ภิกฺขเว สิกฺขิตพฺพํ.}

\prose{กถญฺจ ภิกฺขเว ภิกฺขุ สจิตฺตปริยายกุสโล โหติ? เสยฺยถาปิ ภิกฺขเว อิตฺถี วา ปุริโส วา ทหโร ยุวา มณฺฑนกชาติโก อาทาเส วา ปริสุทฺเธ ปริโยทาเต อจฺเฉ วา อุทกปตฺเต สกํ มุขนิมิตฺตํ ปจฺจเวกฺขมาโน สเจ ตตฺถ ปสฺสติ รชํ วา องฺคณํ วา, ตสฺเสว รชสฺส วา องฺคณสฺส วา ปหานาย วายมติ.\csromanspacer{} โน เจ ตตฺถ ปสฺสติ รชํ วา องฺคณํ วา, เตเนวตฺตมโน โหติ ปริปุณฺณสงฺกปฺโป “ลาภา วต เม ปริสุทฺธํ วต เม”ติ.\csromanspacer{} เอวเมวํ โข ภิกฺขเว ภิกฺขุโน ปจฺจเวกฺขณา พหุการา\footnote{ภิกฺขุ ปจฺจเวกฺขมาโน พหุกาโร (ก)} โหติ กุสเลสุ ธมฺเมสุ “อภิชฺฌาลุ นุ โข พหุลํ วิหรามิ, อนภิชฺฌาลุ นุ โข พหุลํ วิหรามิ.\csromanspacer{} พฺยาปนฺนจิตฺโต นุ โข พหุลํ วิหรามิ, อพฺยาปนฺนจิตฺโต นุ โข พหุลํ วิหรามิ.\csromanspacer{} ถินมิทฺธปริยุฏฺฐิโต นุ โข พหุลํ วิหรามิ, วิคตถินมิทฺโธ นุ โข พหุลํ วิหรามิ.\csromanspacer{} อุทฺธโต นุ โข พหุลํ วิหรามิ, อนุทฺธโต นุ โข พหุลํ วิหรามิ.\csromanspacer{} วิจิกิจฺโฉ นุ โข พหุลํ วิหรามิ, ติณฺณวิจิกิจฺโฉ นุ โข พหุลํ วิหรามิ.\csromanspacer{} โกธโน นุ โข พหุลํ วิหรามิ, อกฺโกธโน นุ โข พหุลํ วิหรามิ.\csromanspacer{} สํกิลิฏฺฐจิตฺโต นุ โข พหุลํ วิหรามิ, อสํกิลิฏฺฐจิตฺโต นุ โข พหุลํ วิหรามิ.\csromanspacer{} สารทฺธกาโย นุ โข พหุลํ วิหรามิ, อสารทฺธกาโย นุ โข พหุลํ วิหรามิ.\csromanspacer{} กุสีโต นุ โข พหุลํ วิหรามิ, อารทฺธวีริโย นุ โข พหุลํ วิหรามิ.\csromanspacer{} อสมาหิโต นุ โข พหุลํ วิหรามิ, สมาหิโต นุ โข พหุลํ วิหรามี”ติ.}

\gambhira{ทสกนิปาตปาฬิ}
\prose{\csromanfolio{330}สเจ ภิกฺขเว ภิกฺขุ ปจฺจเวกฺขมาโน เอวํ ชานาติ “อภิชฺฌาลุ พหุลํ วิหรามิ, พฺยาปนฺนจิตฺโต พหุลํ วิหรามิ, ถินมิทฺธปริยุฏฺฐิโต พหุลํ วิหรามิ, อุทฺธโต พหุลํ วิหรามิ, วิจิกิจฺโฉ พหุลํ วิหรามิ, โกธโน พหุลํ วิหรามิ, สํกิลิฏฺฐจิตฺโต พหุลํ วิหรามิ, สารทฺธกาโย พหุลํ วิหรามิ, กุสีโต พหุลํ วิหรามิ, อสมาหิโต พหุลํ วิหรามี”ติ.\csromanspacer{} เตน ภิกฺขเว ภิกฺขุนา เตสํเยว ปาปกานํ อกุสลานํ ธมฺมานํ ปหานาย อธิมตฺโต ฉนฺโท จ วายาโม จ อุสฺสาโห จ อุสฺโสฬฺหี จ อปฺปฏิวานี จ สติ จ สมฺปชญฺญญฺจ กรณียํ.\csromanspacer{} เสยฺยถาปิ ภิกฺขเว อาทิตฺตเจโล วา อาทิตฺตสีโส วา ตสฺเสว เจลสฺส วา สีสสฺส วา นิพฺพาปนาย อธิมตฺตํ ฉนฺทญฺจ วายามญฺจ อุสฺสาหญฺจ อุสฺโสฬฺหิญฺจ อปฺปฏิวานิญฺจ สติญฺจ สมฺปชญฺญญฺจ กเรยฺย.\csromanspacer{} เอวเมวํ โข เตน ภิกฺขเว ภิกฺขุนา เตสํเยว ปาปกานํ อกุสลานํ ธมฺมานํ ปหานาย อธิมตฺโต ฉนฺโท จ วายาโม จ อุสฺสาโห จ อุสฺโสฬฺหี จ อปฺปฏิวานี จ สติ จ สมฺปชญฺญญฺจ กรณียํ.}

\prose{สเจ ปน ภิกฺขเว ภิกฺขุ ปจฺจเวกฺขมาโน เอวํ ชานาติ “อนภิชฺฌาลุ พหุลํ วิหรามิ, อพฺยาปนฺนจิตฺโต พหุลํ วิหรามิ, วิคตถินมิทฺโธ พหุลํ วิหรามิ.\csromanspacer{} อนุทฺธโต พหุลํ วิหรามิ, ติณฺณวิจิกิจฺโฉ พหุลํ วิหรามิ, อกฺโกธโน พหุลํ วิหรามิ, อสํกิลิฏฺฐจิตฺโต พหุลํ วิหรามิ, อสารทฺธกาโย พหุลํ วิหรามิ, อารทฺธวีริโย พหุลํ วิหรามิ, สมาหิโต พหุลํ วิหรามี”ติ.\csromanspacer{} เตน ภิกฺขเว ภิกฺขุนา เตสุเยว กุสเลสุ ธมฺเมสุ ปติฏฺฐาย อุตฺตริ อาสวานํ ขยาย โยโค กรณีโยติ.\csromanspacer{}\csromanspacer{}ปฐมํ.}

\csromanheader{2}{2. สาริปุตฺตสุตฺต}
\csromanmatikaanchor{mtk.83}
\csromantocmarknopage{section}{(6) 1. สจิตฺตวคฺค}{mtk.83}
\bookmark[dest={mtk.83},level=1]{(6) 1. สจิตฺตวคฺค}
\proseitem{52}{ตตฺร โข อายสฺมา สาริปุตฺโต ภิกฺขู อามนฺเตสิ “อาวุโส ภิกฺขเว”ติ. “อาวุโส”ติ โข เต ภิกฺขู อายสฺมโต สาริปุตฺตสฺส ปจฺจสฺโสสุํ.\csromanspacer{} อายสฺมา สาริปุตฺโต เอตทโวจ\csromandash{}}

\csromanmatikaanchor{mtk.84}
\csromantocmarkref{subsection}{1. สจิตฺตสุตฺต}{mtk.84}
\bookmark[dest={mtk.84},level=2]{1. สจิตฺตสุตฺต}
\prose{โน เจ อาวุโส ภิกฺขุ ปรจิตฺตปริยายกุสโล โหติ, อถ “สจิตฺตปริยายกุสโล ภวิสฺสามี”ติ เอวํ หิ โว อาวุโส สิกฺขิตพฺพํ.}

\chapterheadpageread{(6) 1. สจิตฺตวคฺค}
\prose{\csromanfolio{331}กถญฺจาวุโส ภิกฺขุ สจิตฺตปริยายกุสโล โหติ, เสยฺยถาปิ อาวุโส อิตฺถี วา ปุริโส วา ทหโร ยุวา มณฺฑนกชาติโก อาทาเส วา ปริสุทฺเธ ปริโยทาเต อจฺเฉ วา อุทปตฺเต สกํ มุขนิมิตฺตํ ปจฺจเวกฺขมาโน สเจ ตตฺถ ปสฺสติ รชํ วา องฺคณํ วา, ตสฺเสว รชสฺส วา องฺคณสฺส วา ปหานาย วายมติ.\csromanspacer{} โน เจ ตตฺถ ปสฺสติ รชํ วา องฺคณํ วา, เตเนวตฺตมโน โหติ ปริปุณฺณสงฺกปฺโป “ลาภา วต เม ปริสุทฺธํ วต เม”ติ.}

\prose{เอวเมวํ โข อาวุโส ภิกฺขุโน ปจฺจเวกฺขณา พหุการา โหติ กุสเลสุ ธมฺเมสุ “อภิชฺฌาลุ นุ โข พหุลํ วิหรามิ, อนภิชฺฌาลุ นุ โข พหุลํ วิหรามิ.\csromanspacer{} พฺยาปนฺนจิตฺโต นุ โข พหุลํ วิหรามิ, อพฺยาปนฺนจิตฺโต นุ โข พหุลํ วิหรามิ.\csromanspacer{} ถินมิทฺธปริยุฏฺฐิโต นุ โข พหุลํ วิหรามิ, วิคตถินมิทฺโธ นุ โข พหุลํ วิหรามิ.\csromanspacer{} อุทฺธโต นุ โข พหุลํ วิหรามิ, อนุทฺธโต นุ โข พหุลํ วิหรามิ.\csromanspacer{} วิจิกิจฺโฉ นุ โข พหุลํ วิหรามิ, ติณฺณวิจิกิจฺโฉ นุ โข พหุลํ วิหรามิ.\csromanspacer{} โกธโน นุ โข พหุลํ วิหรามิ, อกฺโกธโน นุ โข พหุลํ วิหรามิ.\csromanspacer{} สํกิลิฏฺฐจิตฺโต นุ โข พหุลํ วิหรามิ, อสํกิลิฏฺฐจิตฺโต นุ โข พหุลํ วิหรามิ.\csromanspacer{} สารทฺธกาโย นุ โข พหุลํ วิหรามิ, อสารทฺธกาโย นุ โข พหุลํ วิหรามิ.\csromanspacer{} กุสีโต นุ โข พหุลํ วิหรามิ, อารทฺธวีริโย นุ โข พหุลํ วิหรามิ.\csromanspacer{} สมาหิโต นุ โข พหุลํ วิหรามิ, อสมาหิโต นุ โข พหุลํ วิหรามี”ติ.}

\prose{สเจ อาวุโส ภิกฺขุ ปจฺจเวกฺขมาโน เอวํ ชานาติ “อภิชฺฌาลุ พหุลํ วิหรามิ ฯเปฯ อสมาหิโต พหุลํ วิหรามี”ติ.\csromanspacer{} เตนาวุโส ภิกฺขุนา เตสํเยว ปาปกานํ อกุสลานํ ธมฺมานํ ปหานาย อธิมตฺโต ฉนฺโท จ วายาโม จ อุสฺสาโห จ อุสฺโสฬฺหี จ อปฺปฏิวานี จ สติ จ สมฺปชญฺญญฺจ กรณียํ.\csromanspacer{} เสยฺยถาปิ อาวุโส อาทิตฺตเจโล วา อาทิตฺตสีโส วา ตสฺเสว เจลสฺส วา สีสสฺส วา นิพฺพาปนาย อธิมตฺตํ ฉนฺทญฺจ วายามญฺจ อุสฺสาหญฺจ อุสฺโสฬฺหิญฺจ อปฺปฏิวานิญฺจ สติญฺจ สมฺปชญฺญญฺจ กเรยฺย.\csromanspacer{} เอวเมวํ โข อาวุโส เตน ภิกฺขุนา เตสํเยว ปาปกานํ อกุสลานํ ธมฺมานํ ปหานาย อธิมตฺโต ฉนฺโท จ วายาโม จ อุสฺสาโห จ อุสฺโสฬฺหี จ อปฺปฏิวานี จ สติ จ สมฺปชญฺญญฺจ กรณียํ.}

\gambhira{ทสกนิปาตปาฬิ}
\prose{\csromanfolio{332}สเจ ปนาวุโส ภิกฺขุ ปจฺจเวกฺขมาโน เอวํ ชานาติ “อนภิชฺฌาลุ พหุลํ วิหรามิ ฯเปฯ สมาหิโต พหุลํ วิหรามี”ติ.\csromanspacer{} เตนาวุโส ภิกฺขุนา เตสุเยว กุสเลสุ ธมฺเมสุ ปติฏฺฐาย อุตฺตริ อาสวานํ ขยาย โยโค กรณีโยติ.\csromanspacer{}\csromanspacer{}ทุติยํ.}

\csromanmatikaanchor{mtk.85}
\csromantocmarkref{subsection}{2. สาริปุตฺตสุตฺต}{mtk.85}
\bookmark[dest={mtk.85},level=2]{2. สาริปุตฺตสุตฺต}
\csromanheader{2}{3. ฐิติสุตฺต}
\proseitem{53}{ฐิติมฺปาหํ ภิกฺขเว น วณฺณยามิ กุสเลสุ ธมฺเมสุ, ปเคว ปริหานิํ.\csromanspacer{} วุฑฺฒิญฺจ โข อหํ ภิกฺขเว วณฺณยามิ กุสเลสุ ธมฺเมสุ, โน ฐิติํ โน หานิํ.}

\prose{กถญฺจ ภิกฺขเว หานิ โหติ กุสเลสุ ธมฺเมสุ, โน ฐิติ โน วุฑฺฒิ? อิธ ภิกฺขเว ภิกฺขุ ยตฺตโก โหติ สทฺธาย สีเลน สุเตน จาเคน ปญฺญาย ปฏิภาเนน, ตสฺส เต ธมฺมา เนว ติฏฺฐนฺติ โน วฑฺฒนฺติ.\csromanspacer{} หานิเมตํ ภิกฺขเว วทามิ กุสเลสุ ธมฺเมสุ, โน ฐิติํ โน วุฑฺฒิํ.\csromanspacer{} เอวํ โข ภิกฺขเว หานิ โหติ กุสเลสุ ธมฺเมสุ โน ฐิติ โน วุฑฺฒิ.}

\prose{กถญฺจ ภิกฺขเว ฐิติ โหติ กุสเลสุ ธมฺเมสุ, โน หานิ โน วุฑฺฒิ? อิธ ภิกฺขเว ภิกฺขุ ยตฺตโก โหติ สทฺธาย สีเลน สุเตน จาเคน ปญฺญาย ปฏิภาเนน, ตสฺส เต ธมฺมา เนว หายนฺติ โน วฑฺฒนฺติ.\csromanspacer{} ฐิติเมตํ ภิกฺขเว วทามิ กุสเลสุ ธมฺเมสุ, โน หานิํ โน วุฑฺฒิํ.\csromanspacer{} เอวํ โข ภิกฺขเว ฐิติ โหติ กุสเลสุ ธมฺเมสุ โน วุฑฺฒิ โน หานิ.}

\prose{กถญฺจ ภิกฺขเว วุฑฺฒิ โหติ กุสเลสุ ธมฺเมสุ, โน ฐิติ โน หานิ? อิธ ภิกฺขเว ภิกฺขุ ยตฺตโก โหติ สทฺธาย สีเลน สุเตน จาเคน ปญฺญาย ปฏิภาเนน, ตสฺส เต ธมฺมา เนว ติฏฺฐนฺติ โน หายนฺติ.\csromanspacer{} วุฑฺฒิเมตํ ภิกฺขเว วทามิ กุสเลสุ ธมฺเมสุ, โน ฐิติํ โน หานิํ.\csromanspacer{} เอวํ โข ภิกฺขเว วุฑฺฒิ โหติ กุสเลสุ ธมฺเมสุ โน ฐิติ โน หานิ.}

\prose{โน เจ ภิกฺขเว ภิกฺขุ ปรจิตฺตปริยายกุสโล โหติ, อถ “สจิตฺตปริยายกุสโล ภวิสฺสามี”ติ เอวํ หิ โว ภิกฺขเว สิกฺขิตพฺพํ.}

\prose{กถญฺจ ภิกฺขเว ภิกฺขุ สจิตฺตปริยายกุสโล โหติ? เสยฺยถาปิ ภิกฺขเว อิตฺถี วา ปุริโส วา ทหโร ยุวา มณฺฑนกชาติโก อาทาเส}

\vspace{\tipitakaparskip}
\csromancenter{(6) 1.\csromanspacer{} สจิตฺตวคฺค วา ปริสุทฺเธ ปริโยทาเต อจฺเฉ วา อุทปตฺเต สกํ มุขนิมิตฺตํ ปจฺจเวกฺขมาโน สเจ ตตฺถ ปสฺสติ รชํ วา องฺคณํ วา, ตสฺเสว รชสฺส วา องฺคณสฺส วา ปหานาย วายมติ, โน เจ ตตฺถ ปสฺสติ รชํ วา องฺคณํ วา, เตเนวตฺตมโน โหติ ปริปุณฺณสงฺกปฺโป “ลาภา วต เม ปริสุทฺธํ วต เม”ติ.\csromanspacer{} เอวเมวํ โข ภิกฺขเว ภิกฺขุโน ปจฺจเวกฺขณา พหุการา โหติ กุสเลสุ ธมฺเมสุ “อภิชฺฌาลุ นุ โข พหุลํ วิหรามิ, อนภิชฺฌาลุ นุ โข พหุลํ วิหรามิ.\csromanspacer{} พฺยาปนฺนจิตฺโต นุ โข พหุลํ วิหรามิ, อพฺยาปนฺนจิตฺโต นุ โข พหุลํ วิหรามิ.\csromanspacer{} ถินมิทฺธปริยุฏฺฐิโต นุ โข พหุลํ วิหรามิ, วิคตถินมิทฺโธ นุ โข พหุลํ วิหรามิ.\csromanspacer{} อุทฺธโต นุ โข พหุลํ วิหรามิ, อนุทฺธโต นุ โข พหุลํ วิหรามิ.\csromanspacer{} วิจิกิจฺโฉ นุ โข พหุลํ วิหรามิ, ติณฺณวิจิกิจฺโฉ นุ โข พหุลํ วิหรามิ.\csromanspacer{} โกธโน นุ โข พหุลํ วิหรามิ, อกฺโกธโน นุ โข พหุลํ วิหรามิ.\csromanspacer{} สํกิลิฏฺฐจิตฺโต นุ โข พหุลํ วิหรามิ, อสํกิลิฏฺฐจิตฺโต นุ โข พหุลํ วิหรามิ.\csromanspacer{} สารทฺธกาโย นุ โข พหุลํ วิหรามิ, อสารทฺธกาโย นุ โข พหุลํ วิหรามิ.\csromanspacer{} กุสีโต นุ โข พหุลํ วิหรามิ, อารทฺธวีริโย นุ โข พหุลํ วิหรามิ.\csromanspacer{} สมาหิโต นุ โข พหุลํ วิหรามิ, อสมาหิโต นุ โข พหุลํ วิหรามี”ติ.}
\prose{\csromanfolio{333}สเจ ภิกฺขเว ภิกฺขุ ปจฺจเวกฺขมาโน เอวํ ชานาติ “อภิชฺฌาลุ พหุลํ วิหรามิ.\csromanspacer{} พฺยาปนฺนจิตฺโต พหุลํ วิหรามิ.\csromanspacer{} ถินมิทฺธปริยุฏฺฐิโต พหุลํ วิหรามิ.\csromanspacer{} อุทฺธโต พหุลํ วิหรามิ.\csromanspacer{} วิจิกิจฺโฉ พหุลํ วิหรามิ.\csromanspacer{} โกธโน พหุลํ วิหรามิ.\csromanspacer{} สํกิลิฏฺฐจิตฺโต พหุลํ วิหรามิ.\csromanspacer{} สารทฺธกาโย พหุลํ วิหรามิ.\csromanspacer{} กุสีโต พหุลํ วิหรามิ.\csromanspacer{} อสมาหิโต พหุลํ วิหรามี”ติ.\csromanspacer{} เตน ภิกฺขเว ภิกฺขุนา เตสํเยว ปาปกานํ อกุสลานํ ธมฺมานํ ปหานาย อธิมตฺโต ฉนฺโท จ วายาโม จ อุสฺสาโห จ อุสฺโสฬฺหี จ อปฺปฏิวานี จ สติ จ สมฺปชญฺญญฺจ กรณียํ.\csromanspacer{} เสยฺยถาปิ ภิกฺขเว อาทิตฺตเจโล วา อาทิตฺตสีโส วา ตสฺเสว เจลสฺส วา สีสสฺส วา นิพฺพาปนาย อธิมตฺตํ ฉนฺทญฺจ วายามญฺจ อุสฺสาหญฺจ อุสฺโสฬฺหิญฺจ อปฺปฏิวานิญฺจ สติญฺจ สมฺปชญฺญญฺจ กเรยฺย.\csromanspacer{} เอวเมวํ โข ภิกฺขเว เตน ภิกฺขุนา เตสํเยว ปาปกานํ อกุสลานํ ธมฺมานํ ปหานาย อธิมตฺโต ฉนฺโท จ วายาโม จ อุสฺสาโห จ อุสฺโสฬฺหี จ อปฺปฏิวานี จ สติ จ สมฺปชญฺญญฺจ กรณียํ.}

\csromanmatikaanchor{mtk.86}
\csromantocmarkref{subsection}{3. ฐิติสุตฺต}{mtk.86}
\bookmark[dest={mtk.86},level=2]{3. ฐิติสุตฺต}
\gambhira{ทสกนิปาตปาฬิ}
\prose{\csromanfolio{334}สเจ ปน ภิกฺขเว ภิกฺขุ ปจฺจเวกฺขมาโน เอวํ ชานาติ “อนภิชฺฌาลุ พหุลํ วิหรามิ.\csromanspacer{} อพฺยาปนฺนจิตฺโต พหุลํ วิหรามิ.\csromanspacer{} วิคตถินมิทฺโธ พหุลํ วิหรามิ.\csromanspacer{} อนุทฺธโต พหุลํ วิหรามิ.\csromanspacer{} ติณฺณวิจิกิจฺโฉ พหุลํ วิหรามิ.\csromanspacer{} อกฺโกธโน พหุลํ วิหรามิ.\csromanspacer{} อสํกิลิฏฺฐจิตฺโต พหุลํ วิหรามิ.\csromanspacer{} อสารทฺธกาโย พหุลํ วิหรามิ.\csromanspacer{} อารทฺธวีริโย พหุลํ วิหรามิ.\csromanspacer{} สมาหิโต พหุลํ วิหรามี”ติ.\csromanspacer{} เตน ภิกฺขเว ภิกฺขุนา เตสุเยว กุสเลสุ ธมฺเมสุ ปติฏฺฐาย อุตฺตริ อาสวานํ ขยาย โยโค กรณีโยติ.\csromanspacer{}\csromanspacer{}ตติยํ.}

\csromanheader{2}{4. สมถสุตฺต}
\proseitem{54}{โน เจ ภิกฺขเว ภิกฺขุ ปรจิตฺตปริยายกุสโล โหติ? อถ “สจิตฺตปริยายกุสโล ภวิสฺสามี”ติ เอวํ หิ โว ภิกฺขเว สิกฺขิตพฺพํ.}

\prose{กถญฺจ ภิกฺขเว ภิกฺขุ สจิตฺตปริยายกุสโล โหติ, เสยฺยถาปิ ภิกฺขเว อิตฺถี วา ปุริโส วา ทหโร ยุวา มณฺฑนกชาติโก อาทาเส วา ปริสุทฺเธ ปริโยทาเต อจฺเฉ วา อุทปตฺเต สกํ มุขนิมิตฺตํ ปจฺจเวกฺขมาโน สเจ ตตฺถ ปสฺสติ รชํ วา องฺคณํ วา, ตสฺเสว รชสฺส วา องฺคณสฺส วา ปหานาย วายมติ.\csromanspacer{} โน เจ ตตฺถ ปสฺสติ รชํ วา องฺคณํ วา, เตเนวตฺตมโน โหติ ปริปุณฺณสงฺกปฺโป “ลาภา วต เม ปริสุทฺธํ วต เม”ติ.\csromanspacer{} เอวเมวํ โข ภิกฺขเว ภิกฺขุโน ปจฺจเวกฺขณา พหุการา โหติ กุสเลสุ ธมฺเมสุ. “ลาภี นุ โขมฺหิ อชฺฌตฺตํ เจโตสมถสฺส, น นุ โขมฺหิ ลาภี อชฺฌตฺตํ เจโตสมถสฺส, ลาภี นุ โขมฺหิ อธิปญฺญา- ธมฺมวิปสฺสนาย, น นุ โขมฺหิ ลาภี อธิปญฺญาธมฺมวิปสฺสนายา”ติ.}

\prose{สเจ ภิกฺขเว ภิกฺขุ ปจฺจเวกฺขมาโน เอวํ ชานาติ “ลาภีมฺหิ อชฺฌตฺตํ เจโตสมถสฺส น ลาภี อธิปญฺญาธมฺมวิปสฺสนายา”ติ.\csromanspacer{} เตน ภิกฺขเว ภิกฺขุนา อชฺฌตฺตํ เจโตสมเถ ปติฏฺฐาย อธิปญฺญาธมฺม- วิปสฺสนาย โยโค กรณีโย.\csromanspacer{} โส อปเรน สมเยน ลาภี เจว โหติ อชฺฌตฺตํ เจโตสมถสฺส ลาภี จ อธิปญฺญาธมฺมวิปสฺสนาย.}

\prose{สเจ ปน ภิกฺขเว ภิกฺขุ ปจฺจเวกฺขมาโน เอวํ ชานาติ “ลาภีมฺหิ อธิปญฺญาธมฺมวิปสฺสนาย น ลาภี อชฺฌตฺตํ เจโตสมถสฺสา”ติ.\csromanspacer{} เตน}

\vspace{\tipitakaparskip}
\csromancenter{(6) 1.\csromanspacer{} สจิตฺตวคฺค ภิกฺขเว ภิกฺขุนา อธิปญฺญาธมฺมวิปสฺสนาย ปติฏฺฐาย อชฺฌตฺตํ เจโตสมเถ โยโค กรณีโย.\csromanspacer{} โส อปเรน สมเยน ลาภี เจว โหติ อธิปญฺญาธมฺมวิปสฺสนาย ลาภี จ อชฺฌตฺตํ เจโตสมถสฺส.}
\prose{\csromanfolio{335}สเจ ปน ภิกฺขเว ภิกฺขุ ปจฺจเวกฺขมาโน เอวํ ชานาติ “น ลาภี อชฺฌตฺตํ เจโตสมถสฺส น ลาภี อธิปญฺญาธมฺมวิปสฺสนายา”ติ.\csromanspacer{} เตน ภิกฺขเว ภิกฺขุนา เตสํเยว กุสลานํ ธมฺมานํ ปฏิลาภาย อธิมตฺโต ฉนฺโท จ วายาโม จ อุสฺสาโห จ อุสฺโสฬฺหี จ อปฺปฏิวานี จ สติ จ สมฺปชญฺญญฺจ กรณียํ.\csromanspacer{} เสยฺยถาปิ ภิกฺขเว อาทิตฺตเจโล วา อาทิตฺตสีโส วา ตสฺเสว เจลสฺส วา สีสสฺส วา นิพฺพาปนาย อธิมตฺตํ ฉนฺทญฺจ วายามญฺจ อุสฺสาหญฺจ อุสฺโสฬฺหิญฺจ อปฺปฏิวานิญฺจ สติญฺจ สมฺปชญฺญญฺจ กเรยฺย.\csromanspacer{} เอวเมวํ โข ภิกฺขเว เตน ภิกฺขุนา เตสํเยว กุสลานํ ธมฺมานํ ปฏิลาภาย อธิมตฺโต ฉนฺโท จ วายาโม จ อุสฺสาโห จ อุสฺโสฬฺหี จ อปฺปฏิวานี จ สติ จ สมฺปชญฺญญฺจ กรณียํ.\csromanspacer{} โส อปเรน สมเยน ลาภี เจว โหติ อชฺฌตฺตํ เจโตสมถสฺส ลาภี จ อธิปญฺญาธมฺม- วิปสฺสนาย.}

\prose{สเจ ปน ภิกฺขเว ภิกฺขุ ปจฺจเวกฺขมาโน เอวํ ชานาติ “ลาภีมฺหิ อชฺฌตฺตํ เจโตสมถสฺส ลาภี อธิปญฺญาธมฺมวิปสฺสนายา”ติ.\csromanspacer{} เตน ภิกฺขเว ภิกฺขุนา เตสุเยว กุสเลสุ ธมฺเมสุ ปติฏฺฐาย อุตฺตริ อาสวานํ ขยาย โยโค กรณีโย.}

\prose{จีวรมฺปาหํ ภิกฺขเว ทุวิเธน วทามิ เสวิตพฺพมฺปิ อเสวิตพฺพมฺปิ.\csromanspacer{} ปิณฺฑปาตมฺปาหํ ภิกฺขเว ทุวิเธน วทามิ เสวิตพฺพมฺปิ อเสวิตพฺพมฺปิ.\csromanspacer{} เสนาสนมฺปาหํ ภิกฺขเว ทุวิเธน วทามิ เสวิตพฺพมฺปิ อเสวิตพฺพมฺปิ.\csromanspacer{} คามนิคมมฺปาหํ ภิกฺขเว ทุวิเธน วทามิ เสวิตพฺพมฺปิ อเสวิตพฺพมฺปิ.\csromanspacer{} ชนปทปเทสมฺปาหํ ภิกฺขเว ทุวิเธน วทามิ เสวิตพฺพมฺปิ อเสวิตพฺพมฺปิ.\csromanspacer{} ปุคฺคลมฺปาหํ ภิกฺขเว ทุวิเธน วทามิ เสวิตพฺพมฺปิ อเสวิตพฺพมฺปิ.}

\csromanmatikaanchor{mtk.87}
\csromantocmarkref{subsection}{4. สมถสุตฺต}{mtk.87}
\bookmark[dest={mtk.87},level=2]{4. สมถสุตฺต}
\prose{“จีวรมฺปาหํ ภิกฺขเว ทุวิเธน วทามิ เสวิตพฺพมฺปิ อเสวิตพฺพมฺปี”ติ อิติ โข ปเนตํ วุตฺตํ, กิญฺเจตํ ปฏิจฺจ วุตฺตํ.\csromanspacer{} ตตฺถ ยํ ชญฺญา จีวรํ “อิทํ โข เม จีวรํ เสวโต อกุสลา ธมฺมา อภิวฑฺฒนฺติ, กุสลา ธมฺมา ปริยายนฺตี”ติ.\csromanspacer{} เอวรูปํ จีวรํ น เสวิตพฺพํ.\csromanspacer{} ตตฺถ ยํ ชญฺญา จีวรํ “อิทํ โข}

\vspace{\tipitakaparskip}
\csromancenter{ทสกนิปาตปาฬิ เม จีวรํ เสวโต อกุสลา ธมฺมา ปริหายนฺติ, กุสลา ธมฺมา อภิวฑฺฒนฺตี”ติ.\csromanspacer{} เอวรูปํ จีวรํ เสวิตพฺพํ. “จีวรมฺปาหํ ภิกฺขเว ทุวิเธน วทามิ เสวิตพฺพมฺปิ อเสวิตพฺพมฺปี”ติ อิติ ยํ ตํ วุตฺตํ, อิทเมตํ ปฏิจฺจ วุตฺตํ.}
\prose{\csromanfolio{336}“ปิณฺฑปาตมฺปาหํ ภิกฺขเว ทุวิเธน วทามิ เสวิตพฺพมฺปิ อเสวิตพฺพมฺปี”ติ อิติ โข ปเนตํ วุตฺตํ, กิญฺเจตํ ปฏิจฺจ วุตฺตํ.\csromanspacer{} ตตฺถ ยํ ชญฺญา ปิณฺฑปาตํ “อิมํ โข เม ปิณฺฑปาตํ เสวโต อกุสลา ธมฺมา อภิวฑฺฒนฺติ, กุสลา ธมฺมา ปริหายนฺตี”ติ.\csromanspacer{} เอวรูโป ปิณฺฑปาโต น เสวิตพฺโพ.\csromanspacer{} ตตฺถ ยํ ชญฺญา ปิณฺฑปาตํ “อิมํ โข เม ปิณฺฑปาตํ เสวโต อกุสลา ธมฺมา ปริหายนฺติ, กุสลา ธมฺมา อภิวฑฺฒนฺตี”ติ.\csromanspacer{} เอวรูโป ปิณฺฑปาโต เสวิตพฺโพ. “ปิณฺฑปาตมฺปาหํ ภิกฺขเว ทุวิเธน วทามิ เสวิตพฺพมฺปิ อเสวิตพฺพมฺปี”ติ อิติ ยํ ตํ วุตฺตํ, อิทเมตํ ปฏิจฺจ วุตฺตํ.}

\prose{“เสนาสนมฺปาหํ ภิกฺขเว ทุวิเธน วทามิ เสวิตพฺพมฺปิ อเสวิตพฺพมฺปี”ติ อิติ โข ปเนตํ วุตฺตํ, กิญฺเจตํ ปฏิจฺจ วุตฺตํ.\csromanspacer{} ตตฺถ ยํ ชญฺญา เสนาสนํ “อิทํ โข เม เสนาสนํ เสวโต อกุสลา ธมฺมา อภิวฑฺฒนฺติ, กุสลา ธมฺมา ปริหายนฺตี”ติ.\csromanspacer{} เอวรูปํ เสนาสนํ น เสวิตพฺพํ.\csromanspacer{} ตตฺถ ยํ ชญฺญา เสนาสนํ “อิทํ โข เม เสนาสนํ เสวโต อกุสลา ธมฺมา ปริหายนฺติ, กุสลา ธมฺมา อภิวฑฺฒนฺตี”ติ.\csromanspacer{} เอวรูปํ เสนาสนํ เสวิตพฺพํ. “เสนาสนมฺปาหํ ภิกฺขเว ทุวิเธน วทามิ เสวิตพฺพมฺปิ อเสวิตพฺพมฺปี”ติ อิติ ยํ ตํ วุตฺตํ, อิทเมตํ ปฏิจฺจ วุตฺตํ.}

\prose{“คามนิคมมฺปาหํ ภิกฺขเว ทุวิเธน วทามิ เสวิตพฺพมฺปิ อเสวิตพฺพมฺปี”ติ.\csromanspacer{} อิติ โข ปเนตํ วุตฺตํ, กิญฺเจตํ ปฏิจฺจ วุตฺตํ.\csromanspacer{} ตตฺถ ยํ ชญฺญา คามนิคมํ “อิมํ โข เม คามนิคมํ เสวโต อกุสลา ธมฺมา อภิวฑฺฒนฺติ, กุสลา ธมฺมา ปริหายนฺตี”ติ.\csromanspacer{} เอวรูโป คามนิคโม น เสวิตพฺโพ.\csromanspacer{} ตตฺถ ยํ ชญฺญา คามนิคมํ “อิมํ โข เม คามนิคมํ เสวโต อกุสลา ธมฺมา ปริหายนฺติ, กุสลา ธมฺมา อภิวฑฺฒนฺตี”ติ.\csromanspacer{} เอวรูโป คามนิคโม เสวิตพฺโพ. “คามนิคมมฺปาหํ ภิกฺขเว ทุวิเธน วทามิ เสวิตพฺพมฺปิ อเสวิตพฺพมฺปี”ติ อิติ ยํ ตํ วุตฺตํ, อิทเมตํ ปฏิจฺจ วุตฺตํ.}

\prose{ชนปทปเทสมฺปาหํ ภิกฺขเว ทุวิเธน วทามิ เสวิตพฺพมฺปิ อเสวิตพฺพมฺปี”ติ อิติ โข ปเนตํ วุตฺตํ, กิญฺเจตํ ปฏิจฺจ วุตฺตํ.\csromanspacer{} ตตฺถ ยํ ชญฺญา ชนปทปเทสํ “อิมํ โข เม ชนปทปเทสํ เสวโต อกุสลา ธมฺมา อภิวฑฺฒนฺติ,}

\vspace{\tipitakaparskip}
\csromancenter{(6) 1.\csromanspacer{} สจิตฺตวคฺค กุสลา ธมฺมา ปริหายนฺตี”ติ.\csromanspacer{} เอวรูโป ชนปทปเทโส น เสวิตพฺโพ.\csromanspacer{} ตตฺถ ยํ ชญฺญา ชนปทปเทสํ “อิมํ โข เม ชนปทปเทสํ เสวโต อกุสลา ธมฺมา ปริหายนฺติ, กุสลา ธมฺมา อภิวฑฺฒนฺตี”ติ.\csromanspacer{} เอวรูโป ชนปทปเทโส เสวิตพฺโพ. “ชนปทปเทสมฺปาหํ ภิกฺขเว ทุวิเธน วทามิ เสวิตพฺพมฺปิ อเสวิตพฺพมฺปี”ติ อิติ ยํ ตํ วุตฺตํ, อิทเมตํ ปฏิจฺจ วุตฺตํ.}
\prose{\csromanfolio{337}“ปุคฺคลมฺปาหํ ภิกฺขเว ทุวิเธน วทามิ เสวิตพฺพมฺปิ อเสวิตพฺพมฺปี”ติ อิติ โข ปเนตํ วุตฺตํ, กิญฺเจตํ ปฏิจฺจ วุตฺตํ.\csromanspacer{} ตตฺถ ยํ ชญฺญา ปุคฺคลํ “อิมํ โข เม ปุคฺคลํ เสวโต อกุสลา ธมฺมา อภิวฑฺฒนฺติ, กุสลา ธมฺมา ปริหายนฺตี”ติ.\csromanspacer{} เอวรูโป ปุคฺคโล น เสวิตพฺโพ.\csromanspacer{} ตตฺถ ยํ ชญฺญา ปุคฺคลํ “อิมํ โข เม ปุคฺคลํ เสวโต อกุสลา ธมฺมา ปริหายนฺติ, กุสลา ธมฺมา อภิวฑฺฒนฺตี”ติ.\csromanspacer{} เอวรูโป ปุคฺคโล เสวิตพฺโพ. “ปุคฺคลมฺปาหํ ภิกฺขเว ทุวิเธน วทามิ เสวิตพฺพมฺปิ อเสวิตพฺพมฺปี”ติ อิติ ยํ ตํ วุตฺตํ, อิทเมตํ ปฏิจฺจ วุตฺตนฺติ.\csromanspacer{}\csromanspacer{}จตุตฺถํ.}

\csromanheader{2}{5. ปริหานสุตฺต}
\proseitem{55}{ตตฺร โข อายสฺมา สาริปุตฺโต ภิกฺขู อามนฺเตสิ “อาวุโส ภิกฺขเว”ติ\footnote{ภิกฺขโวติ (สี, สฺยา)}. “อาวุโส”ติ โข เต ภิกฺขู อายสฺมโต สาริปุตฺตสฺส ปจฺจสฺโสสุํ.\csromanspacer{} อายสฺมา สาริปุตฺโต เอตทโวจ\csromandash{}}

\prose{“ปริหานธมฺโม ปุคฺคโล ปริหานธมฺโม ปุคฺคโล”ติ อาวุโส วุจฺจติ, “อปริหานธมฺโม ปุคฺคโล อปริหานธมฺโม ปุคฺคโล”ติ อาวุโส วุจฺจติ.\csromanspacer{} กิตฺตาวตา นุ โข อาวุโส ปริหานธมฺโม ปุคฺคโล วุตฺโต ภควตา, กิตฺตาวตา จ ปน อปริหานธมฺโม ปุคฺคโล วุตฺโต ภควตาติ.\csromanspacer{} ทูรโตปิ โข มยํ อาวุโส อาคจฺฉาม อายสฺมโต สาริปุตฺตสฺส สนฺติเก เอตสฺส ภาสิตสฺส อตฺถมญฺญาตุํ.\csromanspacer{} สาธุ วตายสฺมนฺตํเยว สาริปุตฺตํ ปฏิภาตุ เอตสฺส ภาสิตสฺส อตฺโถ, อายสฺมโต สาริปุตฺตสฺส สุตฺวา ภิกฺขู ธาเรสฺสนฺตีติ.}

\gambhira{ทสกนิปาตปาฬิ}
\prose{\csromanfolio{338}เตนาหาวุโส สุณาถ สาธุกํ มนสิ กโรถ, ภาสิสฺสามีติ. “เอวมาวุโส”ติ โข เต ภิกฺขู อายสฺมโต สาริปุตฺตสฺส ปจฺจสฺโสสุํ.\csromanspacer{} อายสฺมา สาริปุตฺโต เอตทโวจ\csromandash{}}

\prose{กิตฺตาวตา นุ โข อาวุโส ปริหานธมฺโม ปุคฺคโล วุตฺโต ภควตา, อิธาวุโส ภิกฺขุ อสฺสุตญฺเจว ธมฺมํ น สุณาติ, สุตา จสฺส ธมฺมา สมฺโมสํ คจฺฉนฺติ.\csromanspacer{} เย จสฺส ธมฺมา ปุพฺเพ เจตโส อสมฺผุฏฺฐปุพฺพา, เต จสฺส น สมุทาจรนฺติ, อวิญฺญาตญฺเจว น วิชานาติ.\csromanspacer{} เอตฺตาวตา โข อาวุโส ปริหานธมฺโม ปุคฺคโล วุตฺโต ภควตา.}

\prose{กิตฺตาวตา จ ปนาวุโส อปริหานธมฺโม ปุคฺคโล วุตฺโต ภควตา, อิธาวุโส ภิกฺขุ อสฺสุตญฺเจว ธมฺมํ สุณาติ, สุตา จสฺส ธมฺมา น สมฺโมสํ คจฺฉนฺติ.\csromanspacer{} เย จสฺส ธมฺมา ปุพฺเพ เจตโส อสมฺผุฏฺฐปุพฺพา, เต จสฺส สมุทาจรนฺติ, อวิญฺญาตญฺเจว วิชานาติ.\csromanspacer{} เอตฺตาวตา โข อาวุโส อปริหานธมฺโม ปุคฺคโล วุตฺโต ภควตา.}

\prose{โน เจ อาวุโส ภิกฺขุ ปรจิตฺตปริยายกุสโล โหติ, อถ “สจิตฺตปริยายกุสโล ภวิสฺสามี”ติ เอวํ หิ โว อาวุโส สิกฺขิตพฺพํ.}

\prose{กถญฺจาวุโส ภิกฺขุ สจิตฺตปริยายกุสโล โหติ, เสยฺยถาปิ อาวุโส อิตฺถี วา ปุริโส วา ทหโร ยุวา มณฺฑนกชาติโก อาทาเส วา ปริสุทฺเธ ปริโยทาเต อจฺเฉ วา อุทปตฺเต สกํ มุขนิมิตฺตํ ปจฺจเวกฺขมาโน สเจ ตตฺถ ปสฺสติ รชํ วา องฺคณํ วา, ตสฺเสว รชสฺส วา องฺคณสฺส วา ปหานาย วายมติ.\csromanspacer{} โน เจ ตตฺถ ปสฺสติ รชํ วา องฺคณํ วา, เตเนวตฺตมโน โหติ ปริปุณฺณสงฺกปฺโป “ลาภา วต เม ปริสุทฺธํ วต เม”ติ.\csromanspacer{} เอวเมว โข อาวุโส ภิกฺขุโน ปจฺจเวกฺขณา พหุการา โหติ กุสเลสุ ธมฺเมสุ “อนภิชฺฌาลุ นุ โข พหุลํ วิหรามิ, สํวิชฺชติ นุ โข เม เอโส ธมฺโม อุทาหุ โน.\csromanspacer{} อพฺยาปนฺนจิตฺโต นุ โข พหุลํ วิหรามิ, สํวิชฺชติ นุ โข เม เอโส ธมฺโม อุทาหุ โน.\csromanspacer{} วิคตถินมิทฺโธ นุ โข พหุลํ วิหรามิ, สํวิชฺชติ นุ โข เม เอโส ธมฺโม อุทาหุ โน.\csromanspacer{} อนุทฺธโต นุ โข พหุลํ วิหรามิ, สํวิชฺชติ นุ โข เม เอโส ธมฺโม อุทาหุ โน.\csromanspacer{} ติณฺณวิจิกิจฺโฉ นุ โข พหุลํ วิหรามิ, สํวิชฺชติ นุ โข เม เอโส ธมฺโม อุทาหุ โน.\csromanspacer{} อกฺโกธโน นุ โข พหุลํ วิหรามิ,}

\csromanmatikaanchor{mtk.88}
\csromantocmarkref{subsection}{5. ปริหานสุตฺต}{mtk.88}
\bookmark[dest={mtk.88},level=2]{5. ปริหานสุตฺต}
\vspace{\tipitakaparskip}
\csromancenter{(6) 1.\csromanspacer{} สจิตฺตวคฺค สํวิชฺชติ นุ โข เม เอโส ธมฺโม อุทาหุ โน.\csromanspacer{} อสํกิลิฏฺฐจิตฺโต นุ โข พหุลํ วิหรามิ, สํวิชฺชติ นุ โข เม เอโส ธมฺโม อุทาหุ โน.\csromanspacer{} ลาภี นุ โขมฺหิ อชฺฌตฺตํ ธมฺมปาโมชฺชสฺส, สํวิชฺชติ นุ โข เม เอโส ธมฺโม อุทาหุ โน.\csromanspacer{} ลาภี นุ โขมฺหิ อชฺฌตฺตํ เจโตสมถสฺส, สํวิชฺชติ นุ โข เม เอโส ธมฺโม อุทาหุ โน.\csromanspacer{} ลาภี นุ โขมฺหิ อธิปญฺญาธมฺมวิปสฺสนาย, สํวิชฺชติ นุ โข เม เอโส ธมฺโม อุทาหุ โน”ติ.}
\prose{\csromanfolio{339}สเจ ปน อาวุโส ภิกฺขุ ปจฺจเวกฺขมาโน สพฺเพปิเม กุสเล ธมฺเม อตฺตนิ น สมนุปสฺสติ, เตนาวุโส ภิกฺขุนา สพฺเพสํเยว อิเมสํ กุสลานํ ธมฺมานํ ปฏิลาภาย อธิมตฺโต ฉนฺโท จ วายาโม จ อุสฺสาโห จ อุสฺโสฬฺหี จ อปฺปฏิวานี จ สติ จ สมฺปชญฺญญฺจ กรณียํ.\csromanspacer{} เสยฺยาถาปิ อาวุโส อาทิตฺตเจโล วา อาทิตฺตสีโส วา ตสฺเสว เจลสฺส วา สีสสฺส วา นิพฺพาปนาย อธิมตฺตํ ฉนฺทญฺจ วายามญฺจ อุสฺสาหญฺจ อุสฺโสฬฺหิญฺจ อปฺปฏิวานิญฺจ สติญฺจ สมฺปชญฺญญฺจ กเรยฺย.\csromanspacer{} เอวเมวํ โข อาวุโส เตน ภิกฺขุนา สพฺเพสํเยว กุสลานํ ธมฺมานํ ปฏิลาภาย อธิมตฺโต ฉนฺโท จ วายาโม จ อุสฺสาโห จ อุสฺโสฬฺหี จ อปฺปฏิวานี จ สติ จ สมฺปชญฺญญฺจ กรณียํ.}

\prose{สเจ ปนาวุโส ภิกฺขุ ปจฺจเวกฺขมาโน เอกจฺเจ กุสเล ธมฺเม อตฺตนิ สมนุปสฺสติ, เอกจฺเจ กุสเล ธมฺเม อตฺตนิ น สมนุปสฺสติ.\csromanspacer{} เตนาวุโส ภิกฺขุนา เย กุสเล ธมฺเม อตฺตนิ สมนุปสฺสติ, เตสุ กุสเลสุ ธมฺเมสุ ปติฏฺฐาย.\csromanspacer{} เย กุสเล ธมฺเม อตฺตนิ น สมนุปสฺสติ, เตสํ กุสลานํ ธมฺมานํ ปฏิลาภาย อธิมตฺโต ฉนฺโท จ วายาโม จ อุสฺสาโห จ อุสฺโสฬฺหี จ อปฺปฏิวานี จ สติ จ สมฺปชญฺญญฺจ กรณียํ.\csromanspacer{} เสยฺยถาปิ อาวุโส อาทิตฺตเจโล วา อาทิตฺตสีโส วา ตสฺเสว เจลสฺส วา สีสสฺส วา นิพฺพาปนาย อธิมตฺตํ ฉนฺทญฺจ วายมญฺจ อุสฺสาหญฺจ อุสฺโสฬฺหิญฺจ อปฺปฏิวานิญฺจ สติญฺจ สมฺปชญฺญญฺจ กเรยฺย.\csromanspacer{} เอวเมวํ โข อาวุโส เตน ภิกฺขุนา เย กุสเล ธมฺเม อตฺตนิ สมนุปสฺสติ, เตสุ กุสเลสุ ธมฺเมสุ ปติฏฺฐาย.\csromanspacer{} เย กุสเล ธมฺเม อตฺตนิ น สมนุปสฺสติ, เตสํ กุสลานํ ธมฺมานํ ปฏิลาภาย อธิมตฺโต ฉนฺโท จ วายาโม จ อุสฺสาโห จ อุสฺโสฬฺหี จ อปฺปฏิวานี จ สติ จ สมฺปชญฺญญฺจ กรณียํ.}

\prose{สเจ ปนาวุโส ภิกฺขุ ปจฺจเวกฺขมาโน สพฺเพปิเม กุสเล ธมฺเม อตฺตนิ สมนุปสฺสติ, เตนาวุโส ภิกฺขุนา สพฺเพเสฺวว อิเมสุ}

\vspace{\tipitakaparskip}
\csromancenter{ทสกนิปาตปาฬิ กุสเลสุ ธมฺเมสุ ปติฏฺฐาย อุตฺตริ อาสวานํ ขยาย โยโค กรณีโยติ.\csromanspacer{}\csromanspacer{}ปญฺจมํ.}
\csromanheader{2}{6. ปฐมสญฺญาสุตฺต}
\proseitem{56}{\csromanfolio{340}ทสยิมา ภิกฺขเว สญฺญา ภาวิตา พหุลีกตา มหปฺผลา โหนฺติ มหานิสํสา อมโตคธา อมตปริโยสานา.\csromanspacer{} กตมา ทส? อสุภสญฺญา มรณสญฺญา อาหาเร ปฏิกูลสญฺญา สพฺพโลเก อนภิรตสญฺญา อนิจฺจสญฺญา อนิจฺเจ ทุกฺขสญฺญา ทุกฺเข อนตฺตสญฺญา ปหานสญฺญา วิราคสญฺญา นิโรธสญฺญา.\csromanspacer{} อิมา โข ภิกฺขเว ทส สญฺญา ภาวิตา พหุลีกตา มหปฺผลา โหนฺติ มหานิสํสา อมโตคธา อมตปริโยสานาติ.\csromanspacer{}\csromanspacer{}ฉฏฺฐํ.}

\csromanheader{2}{7. ทุติยสญฺญาสุตฺต}
\proseitem{57}{ทสยิมา ภิกฺขเว สญฺญา ภาวิตา พหุลีกตา มหปฺผลา โหนฺติ มหานิสํสา อมโตคธา อมตปริโยสานา.\csromanspacer{} กตมา ทส? อนิจฺจสญฺญา อนตฺตสญฺญา มรณสญฺญา อาหาเร ปฏิกูลสญฺญา สพฺพโลเก อนภิรตสญฺญา อฏฺฐิกสญฺญา ปุฬวกสญฺญา\footnote{ปุลวกสญฺญา (สี, อิ) ปุฬุวกสญฺญา (ก), อํ 1. 43 ปิฏฺเฐปิ.} วินีลกสญฺญา วิจฺฉิทฺทกสญฺญา อุทฺธุมาตกสญฺญา.\csromanspacer{} อิมา โข ภิกฺขเว ทส สญฺญา ภาวิตา พหุลีกตา มหปฺผลา โหนฺติ มหานิสํสา อมโตคธา อมตปริโยสานาติ.\csromanspacer{}\csromanspacer{}สตฺตมํ.}

\csromanheader{2}{8. มูลกสุตฺต}
\proseitem{58}{\csromansymbolfootnote{*}{เหฏฺฐา 152 ปิฏฺเฐปิ.}สเจ ภิกฺขเว อญฺญติตฺถิยา ปริพฺพาชกา เอวํ ปุจฺเฉยฺยุํ “กิํ มูลกา อาวุโส สพฺเพ ธมฺมา.\csromanspacer{} กิํ สมฺภวา สพฺเพ ธมฺมา.\csromanspacer{} กิํ สมุทยา สพฺเพ ธมฺมา.\csromanspacer{} กิํ สโมสรณา สพฺเพ ธมฺมา.\csromanspacer{} กิํ ปมุขา สพฺเพ ธมฺมา.\csromanspacer{} กิํ อธิปเตยฺยา สพฺเพ ธมฺมา.\csromanspacer{} กิํ อุตฺตรา สพฺเพ ธมฺมา.\csromanspacer{} กิํ สารา สพฺเพ ธมฺมา.\csromanspacer{} กิํ โอคธา สพฺเพ ธมฺมา.\csromanspacer{} กิํ ปริโยสานา สพฺเพ ธมฺมา”ติ.\csromanspacer{} เอวํ ปุฏฺฐา ตุเมฺห ภิกฺขเว เตสํ อญฺญติตฺถิยานํ ปริพฺพาชกานํ กินฺติ พฺยากเรยฺยาถาติ.\csromanspacer{} ภควํมูลกา โน ภนฺเต ธมฺมา ภควํเนตฺติกา ภควํปฏิสรณา, สาธุ วต ภนฺเต ภควนฺตํเยว ปฏิภาตุ เอตสฺส ภาสิตสฺส อตฺโถ, ภควโต สุตฺวา ภิกฺขู ธาเรสฺสนฺตีติ.}

\chapterheadpageread{(6) 1. สจิตฺตวคฺค}
\prose{\csromanfolio{341}เตน หิ ภิกฺขเว สุณาถ สาธุกํ มนสิ กโรถ, ภาสิสฺสามีติ. “เอวํ ภนฺเต”ติ โข เต ภิกฺขู ภควโต ปจฺจสฺโสสุํ.\csromanspacer{} ภควา เอตทโวจ\csromandash{}}

\prose{สเจ ภิกฺขเว อญฺญติตฺถิยา ปริพฺพาชกา เอวํ ปุจฺเฉยฺยุํ “กิํ มูลกา อาวุโส สพฺเพ ธมฺมา.\csromanspacer{} กิํ สมฺภวา สพฺเพ ธมฺมา.\csromanspacer{} กิํ สมุทยา สพฺเพ ธมฺมา.\csromanspacer{} กิํ สโมสรณา สพฺเพ ธมฺมา.\csromanspacer{} กิํ ปมุขา สพฺเพ ธมฺมา.\csromanspacer{} กิํ อธิปเตยฺยา สพฺเพ ธมฺมา.\csromanspacer{} กิํ อุตฺตรา สพฺเพ ธมฺมา.\csromanspacer{} กิํ สารา สพฺเพ ธมฺมา.\csromanspacer{} กิํ โอคธา สพฺเพ ธมฺมา.\csromanspacer{} กิํ ปริโยสานา สพฺเพ ธมฺมา”ติ.\csromanspacer{} เอวํ ปุฏฺฐา ตุเมฺห ภิกฺขเว เตสํ อญฺญติตฺถิยานํ ปริพฺพาชกานํ เอวํ พฺยากเรยฺยาถ. “ฉนฺทมูลกา อาวุโส สพฺเพ ธมฺมา.\csromanspacer{} มนสิการสมฺภวา สพฺเพ ธมฺมา.\csromanspacer{} ผสฺสสมุทยา สพฺเพ ธมฺมา.\csromanspacer{} เวทนาสโมสรณา สพฺเพ ธมฺมา.\csromanspacer{} สมาธิปฺปมุขา สพฺเพ ธมฺมา.\csromanspacer{} สตาธิปเตยฺยา สพฺเพ ธมฺมา.\csromanspacer{} ปญฺญุตฺตรา สพฺเพ ธมฺมา.\csromanspacer{} วิมุตฺติสารา สพฺเพ ธมฺมา.\csromanspacer{} อมโตคธา สพฺเพ ธมฺมา.\csromanspacer{} นิพฺพานปริโยสานา สพฺเพ ธมฺมา”ติ.\csromanspacer{} เอวํ ปุฏฺฐา ตุเมฺห ภิกฺขเว เตสํ อญฺญติตฺถิยานํ ปริพฺพาชกานํ เอวํ พฺยากเรยฺยาถาติ.\csromanspacer{}\csromanspacer{}อฏฺฐมํ.}

\csromanmatikaanchor{mtk.89}
\csromantocmarkref{subsection}{6. ปฐมสญฺญาสุตฺต}{mtk.89}
\bookmark[dest={mtk.89},level=2]{6. ปฐมสญฺญาสุตฺต}
\csromanheader{2}{9. ปพฺพชฺชาสุตฺต}
\proseitem{59}{ตสฺมาติห ภิกฺขเว เอวํ สิกฺขิตพฺพํ “ยถาปพฺพชฺชาปริจิตญฺจ โน จิตฺตํ ภวิสฺสติ, น จุปฺปนฺนา ปาปกา อกุสลา ธมฺมา จิตฺตํ ปริยาทาย ฐสฺสนฺติ, อนิจฺจสญฺญาปริจิตญฺจ โน จิตฺตํ ภวิสฺสติ, อนตฺตสญฺญาปริจิตญฺจ โน จิตฺตํ ภวิสฺสติ, อสุภสญฺญาปริจิตญฺจ โน จิตฺตํ ภวิสฺสติ, อาทีนวสญฺญาปริจิตญฺจ โน จิตฺตํ ภวิสฺสติ, โลกสฺส สมญฺจ วิสมญฺจ ญตฺวา ตํสญฺญาปริจิตญฺจ โน จิตฺตํ ภวิสฺสติ, โลกสฺส ภวญฺจ\footnote{สมฺภวญฺจ (สี, สฺยา)} วิภวญฺจ ญตฺวา ตํสญฺญาปริจิตญฺจ โน จิตฺตํ ภวิสฺสติ, โลกสฺส สมุทยญฺจ อตฺถงฺคมญฺจ ญตฺวา ตํสญฺญาปริจิตญฺจ โน จิตฺตํ ภวิสฺสติ, ปหานสญฺญาปริจิตญฺจ โน จิตฺตํ ภวิสฺสติ, วิราคสญฺญาปริจิตญฺจ โน จิตฺตํ ภวิสฺสติ, นิโรธสญฺญาปริจิตญฺจ โน จิตฺตํ ภวิสฺสตี”ติ.\csromanspacer{} เอวํ หิ โว ภิกฺขเว สิกฺขิตพฺพํ.}

\csromanmatikaanchor{mtk.90}
\csromantocmarkref{subsection}{7. ทุติยสญฺญาสุตฺต}{mtk.90}
\bookmark[dest={mtk.90},level=2]{7. ทุติยสญฺญาสุตฺต}
\gambhira{ทสกนิปาตปาฬิ}
\prose{\csromanfolio{342}ยโต โข ภิกฺขเว ภิกฺขุโน ยถาปพฺพชฺชาปริจิตญฺจ จิตฺตํ โหติ, น จุปฺปนฺนา ปาปกา อกุสลา ธมฺมา จิตฺตํ ปริยาทาย ติฏฺฐนฺติ, อนิจฺจสญฺญาปริจิตญฺจ จิตฺตํ โหติ, อนตฺตสญฺญาปริจิตญฺจ จิตฺตํ โหติ, อสุภสญฺญาปริจิตญฺจ จิตฺตํ โหติ, อาทีนวสญฺญาปริจิตญฺจ จิตฺตํ โหติ, โลกสฺส สมญฺจ วิสมญฺจ ญตฺวา ตํสญฺญาปริจิตญฺจ จิตฺตํ โหติ, โลกสฺส ภวญฺจ วิภวญฺจ ญตฺวา ตํสญฺญาปริจิตญฺจ จิตฺตํ โหติ, โลกสฺส สมุทยญฺจ อตฺถงฺคมญฺจ ญตฺวา ตํสญฺญาปริจิตญฺจ จิตฺตํ โหติ, ปหานสญฺญาปริจิตญฺจ จิตฺตํ โหติ, วิราคสญฺญาปริจิตญฺจ จิตฺตํ โหติ, นิโรธสญฺญาปริจิตญฺจ จิตฺตํ โหติ, ตสฺส ทฺวินฺนํ ผลานํ อญฺญตรํ ผลํ ปาฏิกงฺขํ ทิฏฺเฐว ธมฺเม อญฺญา สติ วา อุปาทิเสเส อนาคามิตาติ.\csromanspacer{}\csromanspacer{}นวมํ.}

\csromanmatikaanchor{mtk.91}
\csromantocmarkref{subsection}{8. มูลกสุตฺต}{mtk.91}
\bookmark[dest={mtk.91},level=2]{8. มูลกสุตฺต}
\csromanheader{2}{10. คิริมานนฺทสุตฺต}
\proseitem{60}{เอกํ สมยํ ภควา สาวตฺถิยํ วิหรติ เชตวเน อนาถปิณฺฑิกสฺส อาราเม.\csromanspacer{} เตน โข ปน สมเยน อายสฺมา คิริมานนฺโท อาพาธิโก โหติ ทุกฺขิโต พาฬฺหคิลาโน.\csromanspacer{} อถ โข อายสฺมา อานนฺโท เยน ภควา เตนุปสงฺกมิ, อุปสงฺกมิตฺวา ภควนฺตํ อภิวาเทตฺวา เอกมนฺตํ นิสีทิ, เอกมนฺตํ นิสินฺโน โข อายสฺมา อานนฺโท ภควนฺตํ เอตทโวจ\csromandash{}}

\prose{“อายสฺมา ภนฺเต คิริมานนฺโท อาพาธิโก โหติ ทุกฺขิโต พาฬฺหคิลาโน, สาธุ ภนฺเต ภควา เยนายสฺมา คิริมานนฺโท เตนุปสงฺกมตุ อนุกมฺปํ อุปาทายา”ติ.\csromanspacer{} สเจ โข ตฺวํ อานนฺท คิริมานนฺทสฺส ภิกฺขุโน ทส สญฺญา ภาเสยฺยาสิ, ฐานํ โข ปเนตํ วิชฺชติ, ยํ คิริมานนฺทสฺส ภิกฺขุโน ทส สญฺญา สุตฺวา โส อาพาโธ ฐานโส ปฏิปฺปสฺสมฺเภยฺย.}

\prose{กตมา ทส? อนิจฺจสญฺญา อนตฺตสญฺญา อสุภสญฺญา อาทีนวสญฺญา ปหานสญฺญา วิราคสญฺญา นิโรธสญฺญา สพฺพโลเก อนภิรตสญฺญา\footnote{อนภิรติสญฺญา (ก)} สพฺพสงฺขาเรสุ อนิจฺฉาสญฺญา อานาปานสฺสติ.}

\prose{กตมา จานนฺท อนิจฺจสญฺญา? อิธานนฺท ภิกฺขุ อรญฺญคโต วา รุกฺขมูลคโต วา สุญฺญาคารคโต วา อิติ ปฏิสญฺจิกฺขติ “รูปํ อนิจฺจํ, เวทนา}

\csromanmatikaanchor{mtk.92}
\csromantocmarkref{subsection}{9. ปพฺพชฺชาสุตฺต}{mtk.92}
\bookmark[dest={mtk.92},level=2]{9. ปพฺพชฺชาสุตฺต}
\vspace{\tipitakaparskip}
\csromancenter{(6) 1.\csromanspacer{} สจิตฺตวคฺค อนิจฺจา, สญฺญา อนิจฺจา, สงฺขารา อนิจฺจา, วิญฺญาณํ อนิจฺจนฺ”ติ.\csromanspacer{} อิติ อิเมสุ ปญฺจสุ อุปาทานกฺขนฺเธสุ อนิจฺจานุปสฺสี วิหรติ.\csromanspacer{} อยํ วุจฺจตานนฺท อนิจฺจสญฺญา. (1)}
\prose{\csromanfolio{343}กตมา จานนฺท อนตฺตสญฺญา? อิธานนฺท ภิกฺขุ อรญฺญคโต วา รุกฺขมูลคโต วา สุญฺญาคารคโต วา อิติ ปฏิสญฺจิกฺขติ “จกฺขุ อนตฺตา, รูปา อนตฺตา, โสตํ อนตฺตา, สทฺทา อนตฺตา, ฆานํ อนตฺตา, คนฺธา อนตฺตา, ชิวฺหา อนตฺตา, รสา อนตฺตา, กายา อนตฺตา, โผฏฺฐพฺพา อนตฺตา, มโน อนตฺตา, ธมฺมา อนตฺตา”ติ.\csromanspacer{} อิติ อิเมสุ ฉสุ อชฺฌตฺติกพาหิเรสุ อายตเนสุ อนตฺตานุปสฺสี วิหรติ.\csromanspacer{} อยํ วุจฺจตานนฺท อนตฺตสญฺญา. (2)}

\prose{กตมา จานนฺท อสุภสญฺญา? อิธานนฺท ภิกฺขุ อิมเมว กายํ อุทฺธํ ปาทตลา อโธ เกสมตฺถกา ตจปริยนฺตํ ปูรํ นานาปฺปการสฺส อสุจิโน ปจฺจเวกฺขติ “อตฺถิ อิมสฺมิํ กาเย เกสา โลมา นขา ทนฺตา ตโจ มํสํ นฺหารุ อฏฺฐิ อฏฺฐิมิญฺชํ วกฺกํ หทยํ ยกนํ กิโลมกํ ปิหกํ ปปฺผาสํ อนฺตํ อนฺตคุณํ อุทริยํ กรีสํ ปิตฺตํ เสมฺหํ ปุพฺโพ โลหิตํ เสโท เมโท อสฺสุ วสา เขโฬ สิงฺฆาณิกา ลสิกา มุตฺตนฺ”ติ.\csromanspacer{} อิติ อิมสฺมิํ กาเย อสุภานุปสฺสี วิหรติ.\csromanspacer{} อยํ วุจฺจตานนฺท อสุภสญฺญา. (3)}

\prose{กตมา จานนฺท อาทีนวสญฺญา? อิธานนฺท ภิกฺขุ อรญฺญคโต วา รุกฺขมูลคโต วา สุญฺญาคารคโต วา อิติ ปฏิสญฺจิกฺขติ “พหุทุกฺโข โข อยํ กาโย พหุ-อาทีนโว, อิติ อิมสฺมิํ กาเย วิวิธา อาพาธา อุปฺปชฺชนฺติ.\csromanspacer{} เสยฺยถิทํ, จกฺขุโรโค โสตโรโค ฆานโรโค ชิวฺหาโรโค กายโรโค สีสโรโค กณฺณโรโค มุขโรโค ทนฺตโรโค โอฏฺฐโรโค กาโส สาโส ปินาโส ฑาโห\footnote{qอโห (สี, สฺยา)} ชโร กุจฺฉิโรโค มุจฺฉา ปกฺขนฺทิกา สูลา วิสูจิกา กุฏฺฐํ คณฺโฑ กิลาโส โสโส อปมาโร ททฺทุ กณฺฑุ กจฺฉุ นขสา วิตจฺฉิกา โลหิตํ ปิตฺตํ\footnote{โลหิตปิตฺตํ (สี)} มธุเมโห อํสา ปิฬกา ภคนฺทลา ปิตฺตสมุฏฺฐานา อาพาธา เสมฺหสมุฏฺฐานา อาพาธา วาตสมุฏฺฐานา อาพาธา สนฺนิปาติกา อาพาธา อุตุปริณามชา อาพาธา วิสมปริหารชา อาพาธา โอปกฺกมิกา อาพาธา กมฺมวิปากชา อาพาธา สีตํ อุณฺหํ ชิฆจฺฉา ปิปาสา อุจฺจาโร ปสฺสาโว”ติ.}

\csromanmatikaanchor{mtk.93}
\csromantocmarkref{subsection}{10. คิริมานนฺทสุตฺต}{mtk.93}
\bookmark[dest={mtk.93},level=2]{10. คิริมานนฺทสุตฺต}
\vspace{\tipitakaparskip}
\csromancenter{ทสกนิปาตปาฬิ อิติ อิมสฺมิํ กาเย อาทีนวานุปสฺสี วิหรติ.\csromanspacer{} อยํ วุจฺจตานนฺท อาทีนวสญฺญา. (4)}
\prose{\csromanfolio{344}กตมา จานนฺท ปหานสญฺญา? อิธานนฺท ภิกฺขุ อุปฺปนฺนํ กามวิตกฺกํ นาธิวาเสติ ปชหติ วิโนเทติ พฺยนฺตีกโรติ อนภาวํ คเมติ, อุปฺปนฺนํ พฺยาปาทวิตกฺกํ นาธิวาเสติ ปชหติ วิโนเทติ พฺยนฺตีกโรติ อนภาวํ คเมติ, อุปฺปนฺนํ วิหิํสาวิตกฺกํ นาธิวาเสติ ปชหติ วิโนเทติ พฺยนฺตีกโรติ อนภาวํ คเมติ, อุปฺปนฺนุปฺปนฺเน ปาปเก อกุสเล ธมฺเม นาธิวาเสติ ปชหติ วิโนเทติ พฺยนฺตีกโรติ อนภาวํ คเมติ.\csromanspacer{} อยํ วุจฺจตานนฺท ปหานสญฺญา. (5)}

\prose{กตมา จานนฺท วิราคสญฺญา? อิธานนฺท ภิกฺขุ อรญฺญคโต วา รุกฺขมูลคโต วา สุญฺญาคารคโต วา อิติ ปฏิสญฺจิกฺขติ “เอตํ สนฺตํ เอตํ ปณีตํ, ยทิทํ สพฺพสงฺขารสมโถ สพฺพูปธิปฺปฏินิสฺสคฺโค ตณฺหากฺขโย วิราโค นิพฺพานนฺ”ติ.\csromanspacer{} อยํ วุจฺจตานนฺท วิราคสญฺญา. (6)}

\prose{กตมา จานนฺท นิโรธสญฺญา? อิธานนฺท ภิกฺขุ อรญฺญคโต วา รุกฺขมูลคโต วา สุญฺญาคารคโต วา อิติ ปฏิสญฺจิกฺขติ “เอตํ สนฺตํ เอตํ ปณีตํ, ยทิทํ สพฺพสงฺขารสมโถ สพฺพูปธิปฺปฏินิสฺสคฺโค ตณฺหากฺขโย นิโรโธ นิพฺพานนฺ”ติ.\csromanspacer{} อยํ วุจฺจตานนฺท นิโรธสญฺญา. (7)}

\prose{กตมา จานนฺท สพฺพโลเก อนภิรตสญฺญา? อิธานนฺท ภิกฺขุ เย โลเก อุปาทานา เจตโส อธิฏฺฐานาภินิเวสานุสยา, เต ปชหนฺโต วิหรติ อนุปาทิยนฺโต.\csromanspacer{} อยํ วุจฺจตานนฺท สพฺพโลเก อนภิรตสญฺญา. (8)}

\prose{กตมา จานนฺท สพฺพสงฺขาเรสุ อนิจฺฉาสญฺญา? อิธานนฺท ภิกฺขุ สพฺพสงฺขาเรสุ อฏฺฏียติ หรายติ ชิคุจฺฉติ.\csromanspacer{} อยํ วุจฺจตานนฺท สพฺพสงฺขาเรสุ อนิจฺฉาสญฺญา. (9)}

\prose{กตมา จานนฺท อานาปานสฺสติ? อิธานนฺท ภิกฺขุ อรญฺญคโต วา รุกฺขมูลคโต วา สุญฺญาคารคโต วา นิสีทติ ปลฺลงฺกํ อาภุชิตฺวา อุชุํ กายํ ปนิธาย ปริมุขํ สติํ อุปฏฺฐเปตฺวา, โส สโตว อสฺสสติ, สโตว ปสฺสสติ, ทีฆํ วา อสฺสสนฺโต “ทีฆํ อสฺสสามี”ติ ปชานาติ, ทีฆํ วา ปสฺสสนฺโต “ทีฆํ ปสฺสสามี”ติ ปชานาติ, รสฺสํ วา อสฺสสนฺโต “รสฺสํ อสฺสสามี”ติ ปชานาติ, รสฺสํ วา ปสฺสสนฺโต}

\vspace{\tipitakaparskip}
\csromancenter{(6) 1.\csromanspacer{} สจิตฺตวคฺค “รสฺสํ ปสฺสสามี”ติ ปชานาติ, “สพฺพกายปฏิสํเวที อสฺสสิสฺสามี”ติ สิกฺขติ, “สพฺพกายปฏิสํเวที ปสฺสสิสฺสามี”ติ สิกฺขติ, “ปสฺสมฺภยํ กายสงฺขารํ อสฺสสิสฺสามี”ติ สิกฺขติ, “ปสฺสมฺภยํ กายสงฺขารํ ปสฺสสิสฺสามี”ติ สิกฺขติ, “ปีติปฏิสํเวที อสฺสสิสฺสามี”ติ สิกฺขติ “ปีติปฏิสํเวที ปสฺสสิสฺสามี”ติ สิกฺขติ, “สุขปฏิสํเวที อสฺสสิสฺสามี”ติ สิกฺขติ, “สุขปฏิสํเวที ปสฺสสิสฺสามี”ติ สิกฺขติ, “จิตฺตสงฺขาร- ปฏิสํเวที อสฺสสิสฺสามี”ติ สิกฺขติ, “จิตฺตสงฺขารปฏิสํเวที ปสฺสสิสฺสามี”ติ สิกฺขติ, “ปสฺสมฺภยํ จิตฺตสงฺขารํ อสฺสสิสฺสามี”ติ สิกฺขติ, “ปสฺสมฺภยํ จิตฺตสงฺขารํ ปสฺสสิสฺสามี”ติ สิกฺขติ, “จิตฺตปฏิสํเวที อสฺสสิสฺสามี”ติ สิกฺขติ, “จิตฺตปฏิสํเวที ปสฺสสิสฺสามี”ติ สิกฺขติ, อภิปฺปโมทยํ จิตฺตํ ฯเปฯ สมาทหํ จิตฺตํ ฯเปฯ วิโมจยํ จิตฺตํ ฯเปฯ อนิจฺจานุปสฺสี ฯเปฯ วิราคานุปสฺสี ฯเปฯ นิโรธานุปสฺสี ฯเปฯ “ปฏินิสฺสคฺคานุปสฺสี อสฺสสิสฺสามี”ติ สิกฺขติ, “ปฏินิสฺสคฺคานุปสฺสี ปสฺสสิสฺสามี”ติ สิกฺขติ.\csromanspacer{} อยํ วุจฺจตานนฺท อานาปานสฺสติ. (10)}
\csromanmatikaanchor{mtk.94}
\csromantocmarkref{section}{อุทฺทานคาถา}{mtk.94}
\bookmark[dest={mtk.94},level=1]{อุทฺทานคาถา}
\prose{\csromanfolio{345}สเจ โข ตฺวํ อานนฺท คิริมานนฺทสฺส ภิกฺขุโน อิมา ทส สญฺญา ภาเสยฺยาสิ, ฐานํ โข ปเนตํ วิชฺชติ, ยํ คิริมานนฺทสฺส ภิกฺขุโน อิมา ทส สญฺญา สุตฺวา โส อาพาโธ ฐานโส ปฏิปฺปสฺสมฺเภยฺยาติ.}

\prose{อถ โข อายสฺมา อานนฺโท ภควโต สนฺติเก อิมา ทส สญฺญา อุคฺคเหตฺวา เยนายสฺมา คิริมานนฺโท เตนุปสงฺกมิ, อุปสงฺกมิตฺวา อายสฺมโต คิริมานนฺทสฺส อิมา ทส สญฺญา อภาสิ.\csromanspacer{} อถ โข อายสฺมโต คิริมานนฺทสฺส ทส สญฺญา สุตฺวา โส อาพาโธ ฐานโส ปฏิปฺปสฺสมฺภิ, วุฏฺฐหิ จายสฺมา คิริมานนฺโท ตมฺหา อาพาธา, ตถา ปหีโน จ ปนายสฺมโต คิริมานนฺทสฺส โส อาพาโธ อโหสีติ.\csromanspacer{}\csromanspacer{}ทสมํ.\csromanspacer{} สจิตฺตวคฺโค ปฐโม.}

\titlehead{\textbf{ตสฺสุทฺทานํ}}
\csromangathameasure{สจิตฺตญฺจ สาริปุตฺต, ฐิติ จ สมเถน จ. \\ ปริหาโน จ เทฺว สญฺญา, มูลา ปพฺพชิตํ คิรีติ.}
\csromangathagroup{\csromangathastanza{สจิตฺตญฺจ สาริปุตฺต, ฐิติ จ สมเถน จ. \\ ปริหาโน จ เทฺว สญฺญา, มูลา ปพฺพชิตํ คิรีติ.}}

\gambhira{ทสกนิปาตปาฬิ \textbf{(7)} \textbf{2. ยมกวคฺค}}
\csromanheader{2}{1. อวิชฺชาสุตฺต}
\proseitem{61}{\csromanfolio{346}“ปุริมา ภิกฺขเว โกฏิ น ปญฺญายติ อวิชฺชาย ‘อิโต ปุพฺเพ อวิชฺชา นาโหสิ, อถ ปจฺฉา สมภวี’ติ” เอวญฺเจตํ ภิกฺขเว วุจฺจติ, อถ จ ปน ปญฺญายติ “อิทปฺปจฺจยา อวิชฺชา”ติ.}

\prose{อวิชฺชมฺปาหํ\footnote{อวิชฺชมฺปหํ (สี, สฺยา)} ภิกฺขเว สาหารํ วทามิ, โน อนาหารํ.\csromanspacer{} โก จาหาโร อวิชฺชาย, “ปญฺจ นีวรณา”ติสฺส วจนียํ.\csromanspacer{} ปญฺจปาหํ ภิกฺขเว นีวรเณ สาหาเร วทามิ, โน อนาหาเร.\csromanspacer{} โก จาหาโร ปญฺจนฺนํ นีวรณานํ, “ตีณิ ทุจฺจริตานี”ติสฺส วจนียํ.\csromanspacer{} ตีณิปาหํ ภิกฺขเว ทุจฺจริตานิ สาหารานิ วทามิ, โน อนาหารานิ.\csromanspacer{} โก จาหาโร ติณฺณํ ทุจฺจริตานํ, “อินฺทฺริย-อสํวโร”ติสฺส วจนียํ.\csromanspacer{} อินฺทฺริย-อสํวรมฺปาหํ ภิกฺขเว สาหารํ วทามิ, โน อนาหารํ.\csromanspacer{} โก จาหาโร อินฺทฺริย-อสํวรสฺส, “อสตาสมฺปชญฺญนฺ”ติสฺส วจนียํ.\csromanspacer{} อสตาสมฺปชญฺญมฺปาหํ ภิกฺขเว สาหารํ วทามิ, โน อนาหารํ.\csromanspacer{} โก จาหาโร อสตาสมฺปชญฺญสฺส, “อโยนิโสมนสิกาโร”ติสฺส วจนียํ.\csromanspacer{} อโยนิโสมนสิการมฺปหํ ภิกฺขเว สาหารํ วทามิ, โน อนาหารํ.\csromanspacer{} โก จาหาโร อโยนิโสมนสิการสฺส, “อสฺสทฺธิยนฺ”ติสฺส วจนียํ.\csromanspacer{} อสฺสทฺธิยมฺปาหํ ภิกฺขเว สาหารํ วทามิ, โน อนาหารํ.\csromanspacer{} โก จาหาโร อสฺสทฺธิยสฺส, “อสทฺธมฺมสฺสวนนฺ” ติสฺส วจนียํ.\csromanspacer{} อสทฺธมฺมสฺสวนมฺปาหํ ภิกฺขเว สาหารํ วทามิ, โน อนาหารํ.\csromanspacer{} โก จาหาโร อสทฺธมฺมสฺสวนสฺส, “อสปฺปุริสสํเสโว”ติสฺส วจนียํ.}

\prose{อิติ โข ภิกฺขเว อสปฺปุริสสํเสโว ปริปูโร อสทฺธมฺมสฺสวนํ ปริปูเรติ, อสทฺธมฺมสฺสวนํ ปริปูรํ อสฺสทฺธิยํ ปริปูเรติ, อสฺสทฺธิยํ ปริปูรํ อโยนิโสมนสิการํ ปริปูเรติ, อโยนิโสมนสิกาโร ปริปูโร อสตาสมฺปชญฺญํ ปริปูเรติ, อสตาสมฺปชญฺญํ ปริปูรํ อินฺทฺริย-อสํวรํ ปริปูเรติ, อินฺทฺริย-อสํวโร ปริปูโร ตีณิ ทุจฺจริตานิ ปริปูเรติ, ตีณิ ทุจฺจริตานิ ปริปูรานิ ปญฺจ นีวรเณ ปริปูเรนฺติ, ปญฺจ นีวรณา ปริปูรา อวิชฺชํ ปริปูเรนฺติ.\csromanspacer{} เอวเมติสฺสา อวิชฺชาย อาหาโร โหติ, เอวญฺจ ปาริปูริ.}

\csromanheader{1}{(7) 2. ยมกวคฺค}
\prose{\csromanfolio{347}เสยฺยถาปิ ภิกฺขเว อุปริปพฺพเต ถุลฺลผุสิตเก เทเว วสฺสนฺเต ( )\footnote{(คลคลายนฺเต) (สี), (คฬคฬายนฺเต) (สฺยา)} ตํ อุทกํ ยถานินฺนํ ปวตฺตมานํ ปพฺพตกนฺทรปทรสาขา ปริปูเรติ, ปพฺพตกนฺทรปทรสาขา ปริปูรา, กุโสพฺเภ\footnote{กุสฺสุพฺเภ (สี), กุสุพฺเภ (สฺยา), กุโสมฺเภ (ก) อํ 1. 245 ปิฏฺเฐปิ.} ปริปูเรนฺติ, กุโสพฺภา ปริปูรา มหาโสพฺเภ\footnote{มหาโสมฺเภ (ก)} ปริปูเรนฺติ, มหาโสพฺภา ปริปูรา กุนฺนทิโย ปริปูเรนฺติ, กุนฺนทิโย ปริปูรา มหานทิโย ปริปูเรนฺติ, มหานทิโย ปริปูรา มหาสมุทฺทํ สาครํ ปริปูเรนฺติ.\csromanspacer{} เอวเมตสฺส มหาสมุทฺทสฺส สาครสฺส อาหาโร โหติ, เอวญฺจ ปาริปูริ.}

\prose{เอวเมวํ โข ภิกฺขเว อสปฺปุริสสํเสโว ปริปูโร อสทฺธมฺมสฺสวนํ ปริปูเรติ, อสทฺธมฺมสฺสวนํ ปริปูรํ อสฺสทฺธิยํ ปริปูเรติ, อสฺสทฺธิยํ ปริปูรํ อโยนิโสมนสิการํ ปริปูเรติ, อโยนิโสมนสิกาโร ปริปูโร อสตาสมฺปชญฺญํ ปริปูเรติ, อสตาสมฺปชญฺญํ ปริปูรํ อินฺทฺริย-อสํวรํ ปริปูเรติ, อินฺทฺริย-อสํวโร ปริปูโร ตีณิ ทุจฺจริตานิ ปริปูเรติ, ตีณิ ทุจฺจริตานิ ปริปูรานิ ปญฺจ นีวรเณ ปริปูเรนฺติ, ปญฺจ นีวรณา ปริปูรา อวิชฺชํ ปริปูเรนฺติ.\csromanspacer{} เอวเมติสฺสา อวิชฺชาย อาหาโร โหติ, เอวญฺจ ปาริปูริ.}

\prose{วิชฺชาวิมุตฺติมฺปาหํ ภิกฺขเว สาหารํ วทามิ, โน อนาหารํ.\csromanspacer{} โก จาหาโร วิชฺชาวิมุตฺติยา, “สตฺต โพชฺฌงฺคา”ติสฺส วจนียํ.\csromanspacer{} สตฺตปาหํ ภิกฺขเว โพชฺฌงฺเค สาหาเร วทามิ, โน อนาหาเร.\csromanspacer{} โก จาหาโร สตฺตนฺนํ โพชฺฌงฺคานํ, “จตฺตาโร สติปฏฺฐานา”ติสฺส วจนียํ.\csromanspacer{} จตฺตาโรปาหํ ภิกฺขเว สติปฏฺฐาเน สาหาเร วทามิ, โน อนาหาเร.\csromanspacer{} โก จาหาโร จตุนฺนํ สติปฏฺฐานานํ, “ตีณิ สุจริตานี”ติสฺส วจนียํ.\csromanspacer{} ตีณิปาหํ ภิกฺขเว สุจริตานิ สาหารานิ วทามิ, โน อนาหารานิ.\csromanspacer{} โก จาหาโร ติณฺณํ สุจริตานํ, “อินฺทฺริยสํวโร”ติสฺส วจนียํ.\csromanspacer{} อินฺทฺริยสํวรมฺปาหํ ภิกฺขเว สาหารํ วทามิ, โน อนาหารํ.\csromanspacer{} โก จาหาโร อินฺทฺริยสํวรสฺส, “สติสมฺปชญฺญนฺ”ติสฺส วจนียํ.\csromanspacer{} สติสมฺปชญฺญมฺปาหํ ภิกฺขเว สาหารํ วทามิ, โน อนาหารํ.\csromanspacer{} โก จาหาโร สติสมฺปชญฺญสฺส, “โยนิโสมนสิกาโร”ติสฺส วจนียํ.\csromanspacer{} โยนิโสมนสิการมฺปาหํ ภิกฺขเว สาหารํ วทามิ, โน อนาหารํ.\csromanspacer{} โก จาหาโร โยนิโสมนสิการสฺส, “สทฺธา”ติสฺส วจนียํ.\csromanspacer{} สทฺธมฺปาหํ ภิกฺขเว สาหารํ วทามิ, โน อนาหารํ.\csromanspacer{} โก จาหาโร}

\vspace{\tipitakaparskip}
\csromancenter{ทสกนิปาตปาฬิ สทฺธาย, “สทฺธมฺมสฺสวนนฺ”ติสฺส วจนียํ.\csromanspacer{} สทฺธมฺมสฺสวนมฺปาหํ ภิกฺขเว สาหารํ วทามิ, โน อนาหารํ.\csromanspacer{} โก จาหาโร สทฺธมฺมสฺสวนสฺส. “สปฺปุริสสํเสโว”ติสฺส วจนียํ.}
\csromanmatikaanchor{mtk.95}
\csromanmatikaanchor{mtk.96}
\csromantocmarknopage{section}{(7) 2. ยมกวคฺค}{mtk.95}
\bookmark[dest={mtk.95},level=1]{(7) 2. ยมกวคฺค}
\csromantocmarkref{subsection}{1. อวิชฺชาสุตฺต}{mtk.96}
\bookmark[dest={mtk.96},level=2]{1. อวิชฺชาสุตฺต}
\prose{\csromanfolio{348}อิติ โข ภิกฺขเว สปฺปุริสสํเสโว ปริปูโร สทฺธมฺมสฺสวนํ ปริปูเรติ, สทฺธมฺมสฺสวนํ ปริปูรํ สทฺธํ ปริปูเรติ, สทฺธา ปริปูรา โยนิโสมนสิการํ ปริปูเรติ, โยนิโสมนสิกาโร ปริปูโร สติสมฺปชญฺญํ ปริปูเรติ, สติสมฺปชญฺญํ ปริปูรํ อินฺทฺริยสํวรํ ปริปูเรติ, อินฺทฺริยสํวโร ปริปูโร ตีณิ สุจริตานิ ปริปูเรติ, ตีณิ สุจริตานิ ปริปูรานิ จตฺตาโร สติปฏฺฐาเน ปริปูเรนฺติ, จตฺตาโร สติปฏฺฐานา ปริปูรา สตฺต โพชฺฌงฺเค ปริปูเรนฺติ, สตฺต โพชฺฌงฺคา ปริปูรา วิชฺชาวิมุตฺติํ ปริปูเรนฺติ.\csromanspacer{} เอวเมติสฺสา วิชฺชาวิมุตฺติยา อาหาโร โหติ, เอวญฺจ ปาริปูริ.}

\prose{เสยฺยถาปิ ภิกฺขเว อุปริปพฺพเต ถุลฺลผุสิตเก เทเว วสฺสนฺเต ตํ อุทกํ ยถานินฺนํ ปวตฺตมานํ ปพฺพตกนฺทรปทรสาขา ปริปูเรติ, ปพฺพตกนฺทรปทรสาขา ปริปูรา กุโสพฺเภ ปริปูเรนฺติ, กุโสพฺภา ปริปูรา มหาโสพฺเภ ปริปูเรนฺติ, มหาโสพฺภา ปริปูรา กุนฺนทิโย ปริปูเรนฺติ, กุนฺนทิโย ปริปูรา มหานทิโย ปริปูเรนฺติ, มหานทิโย ปริปูรา มหาสมุทฺทํ สาครํ ปริปูเรนฺติ.\csromanspacer{} เอวเมตสฺส มหาสมุทฺทสฺส สาครสฺส อาหาโร โหติ, เอวญฺจ ปาริปูริ.}

\prose{เอวเมวํ โข ภิกฺขเว สปฺปุริสสํเสโว ปริปูโร สทฺธมฺมสฺสวนํ ปริปูเรติ, สทฺธมฺมสฺสวนํ ปริปูรํ สทฺธํ ปริปูเรติ, สทฺธา ปริปูรา โยนิโสมนสิการํ ปริปูเรติ, โยนิโสมนสิกาโร ปริปูโร สติสมฺปชญฺญํ ปริปูเรติ, สติสมฺปชญฺญํ ปริปูรํ อินฺทฺริยสํวรํ ปริปูเรติ, อินฺทฺริยสํวโร ปริปูโร ตีณิ สุจริตานิ ปริปูเรติ, ตีณิ สุจริตานิ ปริปูรานิ จตฺตาโร สติปฏฺฐาเน ปริปูเรนฺติ, จตฺตาโร สติปฏฺฐานา ปริปูรา สตฺต โพชฺฌงฺเค ปริปูเรนฺติ, สตฺต โพชฺฌงฺคา ปริปูรา วิชฺชาวิมุตฺติํ ปริปูเรนฺติ.\csromanspacer{} เอวเมติสฺสา วิชฺชาวิมุตฺติยา อาหาโร โหติ, เอวญฺจ ปาริปูรีติ.\csromanspacer{}\csromanspacer{}ปฐมํ.}

\csromanheader{2}{2. ตณฺหาสุตฺต}
\proseitem{62}{“ปุริมา ภิกฺขเว โกฏิ น ปญฺญายติ ภวตณฺหาย ‘อิโต ปุพฺเพ ภวตณฺหา นาโหสิ, อถ ปจฺฉา สมภวี’ติ” เอวญฺเจตํ ภิกฺขเว วุจฺจติ, อถ จ ปน ปญฺญายติ “อิทปฺปจฺจยา ภวตณฺหา”ติ.}

\chapterheadpageread{(7) 2. ยมกวคฺค}
\prose{\csromanfolio{349}ภวตณฺหามฺปาหํ ภิกฺขเว สาหารํ วทามิ, โน อนาหารํ.\csromanspacer{} โก จาหาโร ภวตณฺหาย, “อวิชฺชา”ติสฺส วจนียํ.\csromanspacer{} อวิชฺชมฺปาหํ ภิกฺขเว สาหารํ วทามิ, โน อนาหารํ.\csromanspacer{} โก จหาโร อวิชฺชาย, “ปญฺจ นีวรณา”ติสฺส วจนียํ.\csromanspacer{} ปญฺจ นีวรเณปาหํ ภิกฺขเว สาหาเร วทามิ, โน อนาหาเร.\csromanspacer{} โก จาหาโร ปญฺจนฺนํ นีวรณานํ, “ตีณิ ทุจฺจริตานี”ติสฺส วจนียํ.\csromanspacer{} ตีณิปาหํ ภิกฺขเว ทุจฺจริตานิ สาหารานิ วทามิ, โน อนาหารานิ.\csromanspacer{} โก จาหาโร ติณฺณนฺนํ ทุจฺจริตานํ, “อินฺทฺริย-อสํวโร”ติสฺส วจนียํ.\csromanspacer{} อินฺทฺริย-อสํวรมฺปาหํ ภิกฺขเว สาหารํ วทามิ, โน อนาหารํ.\csromanspacer{} โก จาหาโร อินฺทฺริย-อสํวรสฺส “อสตาสมฺปชญฺญนฺ”ติสฺส วจนียํ.\csromanspacer{} อสตาสมฺปชญฺญมฺปาหํ ภิกฺขเว สาหารํ วทามิ, โน อนาหารํ.\csromanspacer{} โก จาหาโร อสตา สมฺปชญฺญสฺส, “อโยนิโสมนสิกาโร”ติสฺส วจนียํ.\csromanspacer{} อโยนิโสมนสิการมฺปาหํ ภิกฺขเว สาหารํ วทามิ, โน อนาหารํ.\csromanspacer{} โก จาหาโร อโยนิโสมนสิการสฺส, “อสฺสทฺธิยนฺ”ติสฺส วจนียํ.\csromanspacer{} อสฺสทฺธิยมฺปาหํ ภิกฺขเว สาหารํ วทามิ, โน อนาหารํ.\csromanspacer{} โก จาหาโร อสฺสทฺธิยสฺส, “อสฺสทฺธมฺมสฺสวนนฺ”ติสฺส วจนียํ.\csromanspacer{} อสฺสทฺธมฺมสฺสวนมฺปาหํ ภิกฺขเว สาหารํ วทามิ, โน อนาหารํ.\csromanspacer{} โก จาหาโร อสฺสทฺธมฺมสฺสวนสฺส, “อสปฺปุริสสํเสโว”ติสฺส วจนียํ.}

\prose{อิติ โข ภิกฺขเว อสปฺปุริสฺสสํเสโว ปริปูโร อสฺสทฺธมฺมสฺสวนํ ปริปูเรติ, อสฺสทฺธมฺมสฺสวนํ ปริปูรํ อสฺสทฺธิยํ ปริปูเรติ, อสฺสทฺธิยํ ปริปูรํ อโยนิโสมนสิการํ ปริปูเรติ, อโยนิโสมนสิกาโร ปริปูโร อสตาสมฺปชญฺญํ ปริปูเรติ, อสตาสมฺปชญฺญํ ปริปูรํ อินฺทฺริย-อสํวรํ ปริปูเรติ, อินฺทฺริย-อสํวโร ปริปูโร ตีณิ ทุจฺจริตานิ ปริปูเรติ, ตีณิ ทุจฺจริตานิ ปริปูรานิ ปญฺจ นีวรเณ ปริปูเรนฺติ, ปญฺจ นีวรณา ปริปูรา อวิชฺชํ ปริปูเรนฺติ, อวิชฺชา ปริปูรา ภวตณฺหํ ปริปูเรติ.\csromanspacer{} เอวเมติสฺสา ภวตณฺหาย อาหาโร โหติ, เอวญฺจ ปาริปูริ.}

\prose{เสยฺยถาปิ ภิกฺขเว อุปริปพฺพเต ถุลฺลผุสิตเก เทเว วสฺสนฺเต ตํ อุทกํ ยถานินฺนํ ปวตฺตมานํ ปพฺพตกนฺทรปทรสาขา ปริปูเรติ, ปพฺพตกนฺทรปทรสาขา ปริปูรา กุโสพฺเภ ปริปูเรนฺติ, กุโสพฺภา ปริปูรา มหาโสพฺเภ ปริปูเรนฺติ, มหาโสพฺภา ปริปูรา กุนฺนทิโย ปริปูเรนฺติ, กุนฺนทิโย ปริปูรา มหานทิโย ปริปูเรนฺติ, มหานทิโย ปริปูรา}

\vspace{\tipitakaparskip}
\csromancenter{ทสกนิปาตปาฬิ มหาสมุทฺทํ สาครํ ปริปูเรนฺติ.\csromanspacer{} เอวเมตสฺส มหาสมุทฺทสฺส สาครสฺส อาหาโร โหติ, เอวญฺจ ปาริปูริ.}
\prose{\csromanfolio{350}เอวเมวํ โข ภิกฺขเว อสปฺปุริสสํเสโว ปริปูโร อสฺสทฺธมฺมสฺสวนํ ปริปูเรติ, อสฺสทฺธมฺมสฺสวนํ ปริปูรํ อสฺสทฺธิยํ ปริปูเรติ, อสฺสทฺธิยํ ปริปูรํ อโยนิโสมนสิการํ ปริปูเรติ, อโยนิโสมนสิกาโร ปริปูโร อสตาสมฺปชญฺญํ ปริปูเรติ, อสตาสมฺปชญฺญํ ปริปูรํ อินฺทฺริย-อสํวรํ ปริปูเรติ, อินฺทฺริย-อสํวโร ปริปูโร ตีณิ ทุจฺจริตานิ ปริปูเรติ, ตีณิ ทุจฺจริตานิ ปริปูรานิ ปญฺจ นีวรเณ ปริปูเรนฺติ, ปญฺจ นีวรณา ปริปูรา อวิชฺชํ ปริปูเรนฺติ, อวิชฺชา ปริปูรา ภวตณฺหํ ปริปูเรติ.\csromanspacer{} เอวเมติสฺสา ภวตณฺหาย อาหาโร โหติ, เอวญฺจ ปาริปูริ.}

\prose{วิชฺชาวิมุตฺติมฺปาหํ ภิกฺขเว สาหารํ วทามิ, โน อนาหารํ.\csromanspacer{} โก จาหาโร วิชฺชา วิมุตฺติยา, “สตฺต โพชฺฌงฺคา”ติสฺส วจนียํ.\csromanspacer{} สตฺตปาหํ ภิกฺขเว โพชฺฌงฺเค สาหาเร วทามิ, โน อนาหาเร.\csromanspacer{} โก จาหาโร สตฺตนฺนํ โพชฺฌงฺคานํ, “จตฺตาโร สติปฏฺฐานา”ติสฺส วจนียํ.\csromanspacer{} จตฺตาโรปาหํ ภิกฺขเว สติปฏฺฐาเน สาหาเร วทามิ โน อนาหาเร.\csromanspacer{} โก จาหาโร จตุนฺนํ สติปฏฺฐานานํ, “ตีณิ สุจริตานี”ติสฺส วจนียํ.\csromanspacer{} ตีณิปาหํ ภิกฺขเว สุจริตานิ สาหารานิ วทามิ, โน อนาหารานิ.\csromanspacer{} โก จาหาโร ติณฺณนฺนํ สุจริตานํ, “อินฺทฺริยสํวโร”ติสฺส วจนียํ.\csromanspacer{} อินฺทฺริยสํวรมฺปาหํ ภิกฺขเว สาหารํ วทามิ, โน อนาหารํ.\csromanspacer{} โก จาหาโร อินฺทฺริยสํวรสฺส, “สติสมฺปชญฺญนฺ”ติสฺส วจนียํ.\csromanspacer{} สติสมฺปชญฺญมฺปาหํ ภิกฺขเว สาหารํ วทามิ, โน อนาหารํ.\csromanspacer{} โก จาหาโร สติสมฺปชญฺญสฺส, “โยนิโสมนสิกาโร”ติสฺส วจนียํ.\csromanspacer{} โยนิโสมนสิการมฺปาหํ ภิกฺขเว สาหารํ วทามิ, โน อนาหารํ.\csromanspacer{} โก จาหาโร โยนิโสมนสิการสฺส, “สทฺธา”ติสฺส วจนียํ.\csromanspacer{} สทฺธมฺปาหํ ภิกฺขเว สาหารํ วทามิ, โน อนาหารํ.\csromanspacer{} โก จาหาโร สทฺธาย, “สทฺธมฺมสฺสวนนฺ”ติสฺส วจนียํ.\csromanspacer{} สทฺธมฺมสฺสวนมฺปาหํ ภิกฺขเว สาหารํ วทามิ, โน อนาหารํ.\csromanspacer{} โก จาหาโร สทฺธมฺมสฺสวนสฺส, “สปฺปุริสสํเสโว”ติสฺส วจนียํ.}

\csromanmatikaanchor{mtk.97}
\csromantocmarkref{subsection}{2. ตณฺหาสุตฺต}{mtk.97}
\bookmark[dest={mtk.97},level=2]{2. ตณฺหาสุตฺต}
\prose{อิติ โข ภิกฺขเว สปฺปุริสสํเสโว ปริปูโร สทฺธมฺมสฺสวนํ ปริปูเรติ, สทฺธมฺมสฺสวนํ ปริปูรํ สทฺธํ ปริปูเรติ, สทฺธา ปริปูรา โยนิโสมนสิการํ ปริปูเรติ, โยนิโสมนสิกาโร ปริปูโร สติสมฺปชญฺญํ ปริปูเรติ,}

\vspace{\tipitakaparskip}
\csromancenter{(7) 2.\csromanspacer{} ยมกวคฺค สติสมฺปชญฺญํ ปริปูรํ อินฺทฺริยสํวรํ ปริปูเรติ, อินฺทฺริยสํวโร ปริปูโร ตีณิ สุจริตานิ ปริปูเรติ, ตีณิ สุจริตานิ ปริปูรานิ จตฺตาโร สติปฏฺฐาเน ปริปูเรนฺติ, จตฺตาโร สติปฏฺฐานา ปริปูรา สตฺต โพชฺฌงฺเค ปริปูเรนฺติ, สตฺต โพชฺฌงฺคา ปริปูรา วิชฺชาวิมุตฺติํ ปริปูเรนฺติ.\csromanspacer{} เอวเมติสฺสา วิชฺชาวิมุตฺติยา อาหาโร โหติ, เอวญฺจ ปาริปูริ.}
\prose{\csromanfolio{351}เสยฺยถาปิ ภิกฺขเว อุปริปพฺพเต ถุลฺลผุสิตเก เทเว วสฺสนฺเต ตํ อุทกํ ยถานินฺนํ ปวตฺตมานํ ฯเปฯ.\csromanspacer{} เอวเมตสฺส มหาสมุทฺทสฺส สาครสฺส อาหาโร โหติ, เอวญฺจ ปาริปูริ.\csromanspacer{} เอวเมวํ โข ภิกฺขเว สปฺปุริสสํเสโว ปริปูโร สทฺธมฺมสฺสวนํ ปริปูเรติ ฯเปฯ.\csromanspacer{} เอวเมติสฺสา วิชฺชาวิมุตฺติยา อาหาโร โหติ, เอวญฺจ ปาริปูรีติ.\csromanspacer{}\csromanspacer{}ทุติยํ.}

\csromanheader{2}{3. นิฏฺฐงฺคตสุตฺต}
\proseitem{63}{เย เกจิ ภิกฺขเว มยิ นิฏฺฐํ คตา, สพฺเพ เต ทิฏฺฐิสมฺปนฺนา.\csromanspacer{} เตสํ ทิฏฺฐิสมฺปนฺนานํ ปญฺจนฺนํ อิธ นิฏฺฐา, ปญฺจนฺนํ อิธ วิหาย นิฏฺฐา.\csromanspacer{} กตเมสํ ปญฺจนฺนํ อิธ นิฏฺฐา? สตฺตกฺขตฺตุปรมสฺส โกลํโกลสฺส เอกพีชิสฺส สกทาคามิสฺส, โย จ ทิฏฺเฐว ธมฺเม อรหา, อิเมสํ ปญฺจนฺนํ อิธ นิฏฺฐา.\csromanspacer{} กตเมสํ ปญฺจนฺนํ อิธ วิหาย นิฏฺฐา? อนฺตราปรินิพฺพายิสฺส อุปหจฺจปรินิพฺพายิสฺส อสงฺขารปรินิพฺพายิสฺส สสงฺขารปรินิพฺพายิสฺส อุทฺธํโสตสฺส อกนิฏฺฐคามิโน.\csromanspacer{} อิเมสํ ปญฺจนฺนํ อิธ วิหาย นิฏฺฐา.\csromanspacer{} เย เกจิ ภิกฺขเว มยิ นิฏฺฐํ คตา, สพฺเพ เต ทิฏฺฐิสมฺปนฺนา.\csromanspacer{} เตสํ ทิฏฺฐิสมฺปนฺนานํ อิเมสํ ปญฺจนฺนํ อิธ นิฏฺฐา, อิเมสํ ปญฺจนฺนํ อิธ วิหาย นิฏฺฐาติ.\csromanspacer{}\csromanspacer{}ตติยํ.}

\csromanheader{2}{4. อเวจฺจปฺปสนฺนสุตฺต}
\proseitem{64}{เย เกจิ ภิกฺขเว มยิ อเวจฺจปฺปสนฺนา, สพฺเพ เต โสตาปนฺนา.\csromanspacer{} เตสํ โสตาปนฺนานํ ปญฺจนฺนํ อิธ นิฏฺฐา, ปญฺจนฺนํ อิธ วิหาย นิฏฺฐา.\csromanspacer{} กตเมสํ ปญฺจนฺนํ อิธ นิฏฺฐา? สตฺตกฺขตฺตุปรมสฺส โกลํโกลสฺส เอกพีชิสฺส สกทาคามิสฺส, โย จ ทิฏฺเฐว ธมฺเม อรหา, อิเมสํ ปญฺจนฺนํ อิธ นิฏฺฐา.\csromanspacer{} กตเมสํ ปญฺจนฺนํ อิธ วิหาย นิฏฺฐา? อนฺตราปรินิพฺพายิสฺส อุปหจฺจปริพฺพายิสฺส อสงฺขารปรินิพฺพายิสฺส สสงฺขารปรินิพฺพายิสฺส อุทฺธํโสตสฺส อกนิฏฺฐคามิโน.\csromanspacer{} อิเมสํ ปญฺจนฺนํ อิธ วิหาย นิฏฺฐา.\csromanspacer{} เย เกจิ ภิกฺขเว มยิ อเวจฺจปฺปสนฺนา, สพฺเพ เต โสตาปนฺนา.}

\vspace{\tipitakaparskip}
\csromancenter{ทสกนิปาตปาฬิ เตสํ โสตาปนฺนานํ อิเมสํ ปญฺจนฺนํ อิธ นิฏฺฐา, อิเมสํ ปญฺจนฺนํ อิธ วิหาย นิฏฺฐาติ.\csromanspacer{}\csromanspacer{}จตุตฺถํ.}
\csromanheader{2}{5. ปฐมสุขสุตฺต}
\proseitem{65}{\csromanfolio{352}เอกํ สมยํ อายสฺมา สาริปุตฺโต มคเธสุ วิหรติ นาลกคามเก.\csromanspacer{} อถ โข สามณฺฑกานิ ปริพฺพาชโก เยนายสฺมา สาริปุตฺโต เตนุปสงฺกมิ, อุปสงฺกมิตฺวา อายสฺมตา สาริปุตฺเตน สทฺธิํ สมฺโมทิ, สมฺโมทนียํ กถํ สารณียํ วีติสาเรตฺวา เอกมนฺตํ นิสีทิ, เอกมนฺตํ นิสินฺโน โข สามณฺฑกานิ ปริพฺพาชโก อายสฺมนฺตํ สาริปุตฺตํ เอตทโวจ\csromandash{}}

\prose{กิํ นุ โข อาวุโส สาริปุตฺต สุขํ กิํ ทุกฺขนฺติ.\csromanspacer{} อภินิพฺพตฺติ โข อาวุโส ทุกฺขา, อนภินิพฺพตฺติ สุขา.\csromanspacer{} อภินิพฺพตฺติยา อาวุโส สติ อิทํ ทุกฺขํ ปาฏิกงฺขํ\csromandash{}สีตํ อุณฺหํ ชิฆจฺฉา ปิปาสา อุจฺจาโร ปสฺสาโว อคฺคิสมฺผสฺโส ทณฺฑสมฺผสฺโส สตฺถสมฺผสฺโส ญาตีปิ มิตฺตาปิ สงฺคมฺม สมาคมฺม โรเสนฺติ.\csromanspacer{} อภินิพฺพตฺติยา อาวุโส สติ อิทํ ทุกฺขํ ปาฏิกงฺขํ.\csromanspacer{} อนภินิพฺพตฺติยา อาวุโส สติ อิทํ สุขํ ปาฏิกงฺขํ\csromandash{}น สีตํ น อุณฺหํ น ชิฆจฺฉา น ปิปาสา น อุจฺจาโร น ปสฺสาโว น อคฺคิสมฺผสฺโส น ทณฺฑสมฺผสฺโส น สตฺถสมฺผสฺโส ญาตีปิ มิตฺตาปิ สงฺคมฺม สมาคมฺม น โรเสนฺติ.\csromanspacer{} อนภินิพฺพตฺติยา อาวุโส สติ อิทํ สุขํ ปาฏิกงฺขนฺติ.\csromanspacer{}\csromanspacer{}ปญฺจมํ.}

\csromanheader{2}{6. ทุติยสุขสุตฺต}
\csromanmatikaanchor{mtk.98}
\csromantocmarkref{subsection}{3. นิฏฺฐงฺคตสุตฺต}{mtk.98}
\bookmark[dest={mtk.98},level=2]{3. นิฏฺฐงฺคตสุตฺต}
\proseitem{66}{เอกํ สมยํ อายสฺมา สาริปุตฺโต มคเธสุ วิหรติ นาลกคามเก.\csromanspacer{} อถ โข สามณฺฑกานิ ปริพฺพาชโก เยนายสฺมา สาริปุตฺโต เตนุปสงฺกมิ, อุปสงฺกมิตฺวา อายสฺมตา สาริปุตฺเตน สทฺธิํ สมฺโมทิ, สมฺโมทนียํ กถํ สารณียํ วีติสาเรตฺวา เอกมนฺตํ นิสีทิ, เอกมนฺตํ นิสินฺโน โข สามณฺฑกานิ ปริพฺพาชโก อายสฺมนฺตํ สาริปุตฺตํ เอตทโวจ\csromandash{}}

\prose{กิํ นุ โข อาวุโส สาริปุตฺต อิมสฺมิํ ธมฺมวินเย สุขํ กิํ ทุกฺขนฺติ.\csromanspacer{} อนภิรติ โข อาวุโส อิมสฺมิํ ธมฺมวินเย ทุกฺขา, อภิรติ สุขา.\csromanspacer{} อนภิรติยา อาวุโส สติ อิทํ ทุกฺขํ ปาฏิกงฺขํ\csromandash{}คจฺฉนฺโตปิ สุขํ}

\csromanmatikaanchor{mtk.99}
\csromantocmarkref{subsection}{4. อเวจฺจปฺปสนฺนสุตฺต}{mtk.99}
\bookmark[dest={mtk.99},level=2]{4. อเวจฺจปฺปสนฺนสุตฺต}
\vspace{\tipitakaparskip}
\csromancenter{(7) 2.\csromanspacer{} ยมกวคฺค สาตํ นาธิคจฺฉติ, ฐิโตปิ.\csromanspacer{} นิสินฺโนปิ.\csromanspacer{} สยาโนปิ.\csromanspacer{} คามคโตปิ.\csromanspacer{} อรญฺญคโตปิ.\csromanspacer{} รุกฺขมูลคโตปิ.\csromanspacer{} สุญฺญาคารคโตปิ.\csromanspacer{} อพฺโภกาสคโตปิ.\csromanspacer{} ภิกฺขุมชฺฌคโตปิ สุขํ สาตํ นาธิคจฺฉติ.\csromanspacer{} อนภิรติยา อาวุโส สติ อิทํ ทุกฺขํ ปาฏิกงฺขํ.}
\prose{\csromanfolio{353}อภิรติยา อาวุโส สติ อิทํ สุขํ ปาฏิกงฺขํ\csromandash{}คจฺฉนฺโตปิ สุขํ สาตํ อธิคจฺฉติ.\csromanspacer{} ฐิโตปิ.\csromanspacer{} นิสินฺโนปิ.\csromanspacer{} สยาโนปิ.\csromanspacer{} คามคโตปิ.\csromanspacer{} อรญฺญคโตปิ.\csromanspacer{} รุกฺขมูลคโตปิ.\csromanspacer{} สุญฺญาคารคโตปิ.\csromanspacer{} อพฺโภกาสคโตปิ.\csromanspacer{} ภิกฺขุมชฺฌคโตปิ สุขํ สาตํ อธิคจฺฉติ.\csromanspacer{} อภิรติยา อาวุโส สติ อิทํ สุขํ ปาฏิกงฺขนฺติ.\csromanspacer{}\csromanspacer{}ฉฏฺฐํ.}

\csromanheader{2}{7. ปฐมนฬกปานสุตฺต}
\csromanmatikaanchor{mtk.100}
\csromantocmarkref{subsection}{5. ปฐมสุขสุตฺต}{mtk.100}
\bookmark[dest={mtk.100},level=2]{5. ปฐมสุขสุตฺต}
\proseitem{67}{เอกํ สมยํ ภควา โกสเลสุ จาริกํ จรมาโน มหตา ภิกฺขุสํเฆน สทฺธิํ เยน นฬกปานํ นาม โกสลานํ นิคโม, ตทวสริ.\csromanspacer{} ตตฺร สุทํ ภควา นฬกปาเน วิหรติ ปลาสวเน.\csromanspacer{} เตน โข ปน สมเยน ภควา ตทหุโปสเถ ภิกฺขุสํฆปริวุโต นิสินฺโน โหติ.\csromanspacer{} อถ โข ภควา พหุเทว รตฺติํ ภิกฺขูนํ ธมฺมิยา กถาย สนฺทสฺเสตฺวา สมาทเปตฺวา สมุตฺเตเชตฺวา สมฺปหํเสตฺวา ตุณฺหีภูตํ ตุณฺหีภูตํ ภิกฺขุสํฆํ อนุวิโลเกตฺวา อายสฺมนฺตํ สาริปุตฺตํ อามนฺเตสิ\csromandash{}}

\prose{วิคตถินมิทฺโธ\footnote{วิคตถีนมิทฺโธ (สี, สฺยา, กํ, อิ)} โข สาริปุตฺต ภิกฺขุสํโฆ, ปฏิภาตุ ตํ สาริปุตฺต ภิกฺขูนํ ธมฺมี กถา, ปิฏฺฐิ เม อาคิลายติ, ตมหํ อายมิสฺสามีติ. “เอวํ ภนฺเต”ติ โข อายสฺมา สาริปุตฺโต ภควโต ปจฺจสฺโสสิ.}

\prose{อถ โข ภควา จตุคฺคุณํ สงฺฆาฏิํ ปญฺญาเปตฺวา ทกฺขิเณน ปสฺเสน สีหเสยฺยํ กปฺเปสิ ปาเท ปาทํ อจฺจาธาย สโต สมฺปชาโน อุฏฺฐานสญฺญํ มนสิ กริตฺวา.\csromanspacer{} ตตฺร โข อายสฺมา สาริปุตฺโต ภิกฺขู อามนฺเตสิ “อาวุโส ภิกฺขเว”ติ. “อาวุโส”ติ โข เต ภิกฺขู อายสฺมโต สาริปุตฺตสฺส ปจฺจสฺโสสุํ.\csromanspacer{} อายสฺมา สาริปุตฺโต เอตทโวจ\csromandash{}}

\csromanmatikaanchor{mtk.101}
\csromantocmarkref{subsection}{6. ทุติยสุขสุตฺต}{mtk.101}
\bookmark[dest={mtk.101},level=2]{6. ทุติยสุขสุตฺต}
\gambhira{ทสกนิปาตปาฬิ}
\prose{\csromanfolio{354}ยสฺส กสฺสจิ อาวุโส สทฺธา นตฺถิ กุสเลสุ ธมฺเมสุ, หิรี\footnote{หิริ (สี, สฺยา, กํ, อิ)} นตฺถิ.\csromanspacer{} โอตฺตปฺปํ นตฺถิ.\csromanspacer{} วีริยํ\footnote{วิริยํ (สี, สฺยา, กํ, อิ)} นตฺถิ.\csromanspacer{} ปญฺญา นตฺถิ กุสเลสุ ธมฺเมสุ, ตสฺส ยา รตฺติ วา ทิวโส วา อาคจฺฉติ, หานิเยว ปาฏิกงฺขา กุสเลสุ ธมฺเมสุ, โน วุทฺธิ.\csromanspacer{} เสยฺยถาปิ อาวุโส กาฬปกฺเข จนฺทสฺส ยา รตฺติ วา ทิวโส วา อาคจฺฉติ, หายเตว วณฺเณน, หายติ มณฺฑเลน, หายติ อาภาย, หายติ อาโรหปริณาเหน.\csromanspacer{} เอวเมวํ โข อาวุโส ยสฺส กสฺสจิ สทฺธา นตฺถิ กุสเลสุ ธมฺเมสุ, หิรี นตฺถิ.\csromanspacer{} โอตฺตปฺปํ นตฺถิ.\csromanspacer{} วีริยํ นตฺถิ.\csromanspacer{} ปญฺญา นตฺถิ กุสเลสุ ธมฺเมสุ, ตสฺส ยา รตฺติ วา ทิวโส วา อาคจฺฉติ, หานิเยว ปาฏิกงฺขา กุสเลสุ ธมฺเมสุ, โน วุทฺธิ.}

\prose{“อสฺสทฺโธ ปุริสปุคฺคโล”ติ อาวุโส ปริหานเมตํ. “อหิริโก ปุริสปุคฺคโล”ติ อาวุโส ปริหานเมตํ. “อโนตฺตปฺปี ปุริสปุคฺคโล”ติ อาวุโส ปริหานเมตํ. “กุสีโต ปุริสปุคฺคโล”ติ อาวุโส ปริหานเมตํ. “ทุปฺปญฺโญ ปุริสปุคฺคโล”ติ อาวุโส ปริหานเมตํ. “โกธโน ปุริสปุคฺคโล”ติ อาวุโส ปริหานเมตํ. “อุปนาหี ปุริสปุคฺคโล”ติ อาวุโส ปริหานเมตํ. “ปาปิจฺโฉ ปุริสปุคฺคโล”ติ อาวุโส ปริหานเมตํ. “ปาปมิตฺโต ปุริสปุคฺคโล”ติ อาวุโส ปริหานเมตํ. “มิจฺฉาทิฏฺฐิโก ปุริสปุคฺคโล”ติ อาวุโส ปริหานเมตํ.}

\prose{ยสฺส กสฺสจิ อาวุโส สทฺธา อตฺถิ กุสเลสุ ธมฺเมสุ, หิรี อตฺถิ.\csromanspacer{} โอตฺตปฺปํ อตฺถิ.\csromanspacer{} ปญฺญา อตฺถิ กุสเลสุ ธมฺเมสุ, ตสฺส ยา รตฺติ วา ทิวโส วา อาคจฺฉติ, วุทฺธิเยว ปาฏิกงฺขา กุสเลสุ ธมฺเมสุ, โน ปริหานิ.\csromanspacer{} เสยฺยถาปิ อาวุโส ชุณฺหปกฺเข จนฺทสฺส ยา รตฺติ วา ทิวโส วา อาคจฺฉติ, วฑฺฒเตว วณฺเณน, วฑฺฒติ มณฺฑเลน, วฑฺฒติ อาภาย, วฑฺฒติ อาโรหปริณาเหน.\csromanspacer{} เอวเมวํ โข อาวุโส ยสฺส กสฺสจิ สทฺธา อตฺถิ กุสเลสุ ธมฺเมสุ, หิรี อตฺถิ.\csromanspacer{} โอตฺตปฺปํ อตฺถิ.\csromanspacer{} วีริยํ อตฺถิ.\csromanspacer{} ปญฺญา อตฺถิ กุสเลสุ ธมฺเมสุ, ตสฺส ยา รตฺติ วา ทิวโส วา อาคจฺฉติ, วุทฺธิเยว ปาฏิกงฺขา กุสเลสุ ธมฺเมสุ, โน ปริหานิ.}

\chapterheadpageread{(7) 2. ยมกวคฺค}
\csromanmatikaanchor{mtk.102}
\csromantocmarkref{subsection}{7. ปฐมนฬกปานสุตฺต}{mtk.102}
\bookmark[dest={mtk.102},level=2]{7. ปฐมนฬกปานสุตฺต}
\prose{\csromanfolio{355}“สทฺโธ ปุริสปุคฺคโล”ติ อาวุโส อปริหานเมตํ. “หิรีมา ปุริสปุคฺคโล”ติ อาวุโส อปริหานเมตํ. “โอตฺตปฺปี ปุริสปุคฺคโล”ติ อาวุโส อปริหานเมตํ. “อารทฺธวีริโย ปุริสปุคฺคโล”ติ อาวุโส อปริหานเมตํ. “ปญฺญวา ปุริสปุคฺคโล”ติ อาวุโส อปริหานเมตํ. “อกฺโกธโน ปุริสปุคฺคโล”ติ อาวุโส อปริหานเมตํ. “อนุปนาหี ปุริสปุคฺคโล”ติ อาวุโส อปริหานเมตํ. “อปฺปิจฺโฉ ปุริสปุคฺคโล”ติ อาวุโส อปริหานเมตํ. “กลฺยาณมิตฺโต ปุริสปุคฺคโล”ติ อาวุโส อปริหานเมตํ. “สมฺมาทิฏฺฐิโก ปุริสปุคฺคโล”ติ อาวุโส อปริหานเมตนฺติ.}

\prose{อถ โข ภควา ปจฺจุฏฺฐาย อายสฺมนฺตํ สาริปุตฺตํ อามนฺเตสิ\csromandash{} สาธุ สาธุ สาริปุตฺต, ยสฺส กสฺสจิ สาริปุตฺต สทฺธา นตฺถิ กุสเลสุ ธมฺเมสุ, หิรี นตฺถิ.\csromanspacer{} โอตฺตปฺปํ นตฺถิ.\csromanspacer{} วีริยํ นตฺถิ.\csromanspacer{} ปญฺญา นตฺถิ กุสเลสุ ธมฺเมสุ, ตสฺส ยา รตฺถิ วา ทิวโส วา อาคจฺฉติ, หานิเยว ปาฏิกงฺขา กุสเลสุ ธมฺเมสุ, โน วุทฺธิ.\csromanspacer{} เสยฺยถาปิ สาริปุตฺต กาฬปกฺเข จนฺทสฺส ยา รตฺติ วา ทิวโส วา อาคจฺฉติ, หายเตว วณฺเณน, หายติ มณฺฑเลน, หายติ อาภาย, หายติ อาโรหปริณาเหน.\csromanspacer{} เอวเมวํ โข สาริปุตฺต ยสฺส กสฺสจิ สทฺธา นตฺถิ กุสเลสุ ธมฺเมสุ ฯเปฯ.\csromanspacer{} ปญฺญา นตฺถิ กุสเลสุ ธมฺเมสุ, ตสฺส ยา รตฺติ วา ทิวโส วา ฯเปฯ โน วุทฺธิ.}

\prose{“อสฺสทฺโธ ปุริสปุคฺคโล”ติ สาริปุตฺต ปริหานเมตํ.\csromanspacer{} อหิริโก.\csromanspacer{} อโนตฺตปฺปี.\csromanspacer{} กุสีโต.\csromanspacer{} ทุปฺปญฺโญ.\csromanspacer{} โกธโน.\csromanspacer{} อุปนาหี.\csromanspacer{} ปาปิจฺโฉ.\csromanspacer{} ปาปมิตฺโต. “มิจฺฉาทิฏฺฐิโก ปุริสปุคฺคโล”ติ สาริปุตฺต ปริหานเมตํ.}

\prose{ยสฺส กสฺสจิ สาริปุตฺต สทฺธา อตฺถิ กุสเลสุ ธมฺเมสุ, หิรี อตฺถิ.\csromanspacer{} โอตฺตปฺปํ อตฺถิ.\csromanspacer{} วีริยํ อตฺถิ.\csromanspacer{} ปญฺญา อตฺถิ กุสเลสุ ธมฺเมสุ, ตสฺส ยา รตฺติ วา ทิวโส วา อาคจฺฉติ, วุทฺธิเยว ปาฏิกงฺขา กุสเลสุ ธมฺเมสุ, โน ปริหานิ.\csromanspacer{} เสยฺยถาปิ สาริปุตฺต ชุณฺหปกฺเข จนฺทสฺส ยา รตฺติ วา ทิวโส วา อาคจฺฉติ, วฑฺฒเตว วณฺเณน, วฑฺฒติ มณฺฑเลน, วฑฺฒติ อาภาย, วฑฺฒติ อาโรหปริณาเหน.\csromanspacer{} เอวเมวํ โข สาริปุตฺต ยสฺส กสฺสจิ สทฺธา อตฺถิ กุสเลสุ ธมฺเมสุ, หิรี อตฺถิ.\csromanspacer{} โอตฺตปฺปํ อตฺถิ.\csromanspacer{} วีริยํ อตฺถิ.\csromanspacer{} ปญฺญา อตฺถิ กุสเลสุ ธมฺเมสุ, ตสฺส ยา รตฺติ วา ทิวโส วา อาคจฺฉติ, วุทฺธิเยว ปาฏิกงฺขา กุสเลสุ ธมฺเมสุ, โน ปริหานิ.}

\gambhira{ทสกนิปาตปาฬิ}
\prose{\csromanfolio{356}“สทฺโธ ปุริสปุคฺคโล”ติ สาริปุตฺต อปริหานเมตํ.\csromanspacer{} หิรีมา.\csromanspacer{} โอตฺตปฺปี.\csromanspacer{} อารทฺธวีริโย.\csromanspacer{} ปญฺญวา.\csromanspacer{} อกฺโกธโน.\csromanspacer{} อนุปนาหี.\csromanspacer{} อปฺปิจฺโฉ.\csromanspacer{} กลฺยาณมิตฺโต. “สมฺมาทิฏฺฐิโก ปุริสปุคฺคโล”ติ สาริปุตฺต อปริหานเมตนฺติ.\csromanspacer{}\csromanspacer{}สตฺตมํ.}

\csromanheader{2}{8. ทุติยนฬกปานสุตฺต}
\proseitem{68}{เอกํ สมยํ ภควา นฬกปาเน วิหรติ ปลาสวเน.\csromanspacer{} เตน โข ปน สมเยน ภควา ตทหุโปสเถ ภิกฺขุสํฆปริวุโต นิสินฺโน โหติ.\csromanspacer{} อถ โข ภควา พหุเทว รตฺติํ ภิกฺขูนํ ธมฺมิยา กถาย สนฺทสฺเสตฺวา สมาทเปตฺวา สมุตฺเตเชตฺวา สมฺปหํเสตฺวา ตุณฺหีภูตํ ตุณฺหีภูตํ ภิกฺขุสํฆํ อนุวิโลเกตฺวา อายสฺมนฺตํ สาริปุตฺตํ อามนฺเตสิ\csromandash{}}

\prose{วิคตถินมิทฺโธ โข สาริปุตฺต ภิกฺขุสํโฆ, ปฏิภาตุ ตํ สาริปุตฺต ภิกฺขูนํ ธมฺมี กถา, ปิฏฺฐิ เม อาคิลายติ, ตมหํ อายมิสฺสามีติ. “เอวํ ภนฺเต”ติ โข อายสฺมา สาริปุตฺโต ภควโต ปจฺจสฺโสสิ.}

\prose{อถ โข ภควา จตุคฺคุณํ สงฺฆาฏิํ ปญฺญาเปตฺวา ทกฺขิเณน ปสฺเสน สีหเสยฺยํ กปฺเปสิ ปาเท ปาทํ อจฺจาธาย สโต สมฺปชาโน อุฏฺฐานสญฺญํ มนสิ กริตฺวา.\csromanspacer{} ตตฺร โข อายสฺมา สาริปุตฺโต ภิกฺขู อามนฺเตสิ “อาวุโส ภิกฺขเว”ติ. “อาวุโส”ติ โข เต ภิกฺขู อายสฺมโต สาริปุตฺตสฺส ปจฺจสฺโสสุํ.\csromanspacer{} อายสฺมา สาริปุตฺโต เอตทโวจ\csromandash{}}

\prose{ยสฺส กสฺสจิ อาวุโส สทฺธา นตฺถิ กุสเลสุ ธมฺเมสุ, หิรี นตฺถิ.\csromanspacer{} โอตฺตปฺปํ นตฺถิ.\csromanspacer{} วีริยํ นตฺถิ.\csromanspacer{} ปญฺญา นตฺถิ.\csromanspacer{} โสตาวธานํ นตฺถิ.\csromanspacer{} ธมฺมธารณา นตฺถิ.\csromanspacer{} อตฺถูปปริกฺขา นตฺถิ.\csromanspacer{} ธมฺมานุธมฺมปฺ- ปฏิปตฺติ นตฺถิ.\csromanspacer{} อปฺปมาโท นตฺถิ กุสเลสุ ธมฺเมสุ, ตสฺส ยา รตฺติ วา ทิวโส วา อาคจฺฉติ, หานิเยว ปาฏิกงฺขา กุสเลสุ ธมฺเมสุ โน วุทฺธิ.\csromanspacer{} เสยฺยถาปิ อาวุโส กาฬปกฺเข จนฺทสฺส ยา รตฺติ วา ทิวโส วา อาคจฺฉติ, หายเตว วณฺเณน, หายติ มณฺฑเลน, หายติ อาภาย, หายติ อาโรหปริณาเหน.\csromanspacer{} เอวเมวํ โข อาวุโส ยสฺส กสฺสจิ สทฺธา นตฺถิ กุสเลสุ ธมฺเมสุ, หิรี นตฺถิ.\csromanspacer{} โอตฺตปฺปํ นตฺถิ.\csromanspacer{} วีริยํ นตฺถิ.\csromanspacer{} ปญฺญา นตฺถิ.\csromanspacer{} โสตาวธานํ นตฺถิ.\csromanspacer{} ธมฺมธารณา นตฺถิ.\csromanspacer{} อตฺถูปปริกฺขา นตฺถิ.\csromanspacer{} ธมฺมานุธมฺมปฺปฏิ- ปตฺติ นตฺถิ.\csromanspacer{} อปฺปมาโท นตฺถิ กุสเลสุ ธมฺเมสุ, ตสฺส ยา รตฺติ}

\vspace{\tipitakaparskip}
\csromancenter{(7) 2.\csromanspacer{} ยมกวคฺค วา ทิวโส วา อาคจฺฉติ, หานิเยว ปาฏิกงฺขา กุสเลสุ ธมฺเมสุ โน วุทฺธิ.}
\prose{\csromanfolio{357}ยสฺส กสฺสจิ อาวุโส สทฺธา อตฺถิ กุสเลสุ ธมฺเมสุ, หิรี อตฺถิ.\csromanspacer{} โอตฺตปฺปํ อตฺถิ.\csromanspacer{} วีริยํ อตฺถิ.\csromanspacer{} ปญฺญา อตฺถิ.\csromanspacer{} โสตาวธานํ อตฺถิ.\csromanspacer{} ธมฺมธารณา อตฺถิ.\csromanspacer{} อตฺถูปปริกฺขา อตฺถิ.\csromanspacer{} ธมฺมานุธมฺมปฺปฏิปตฺติ อตฺถิ.\csromanspacer{} อปฺปมาโท อตฺถิ กุสเลสุ ธมฺเมสุ, ตสฺส ยา รตฺติ วา ทิวโส วา อาคจฺฉติ, วุทฺธิเยว ปาฏิกงฺขา กุสเลสุ ธมฺเมสุ, โน ปริหานิ.\csromanspacer{} เสยฺยถาปิ อาวุโส ชุณฺหปกฺเข จนฺทสฺส ยา รตฺติ วา ทิวโส วา อาคจฺฉติ, วฑฺฒเตว วณฺเณน, วฑฺฒติ มณฺฑเลน, วฑฺฒติ อาภาย, วฑฺฒติ อาโรหปริณาเหน.\csromanspacer{} เอวเมวํ โข อาวุโส ยสฺส กสฺสจิ สทฺธา อตฺถิ กุสเลสุ ธมฺเมสุ ฯเปฯ.\csromanspacer{} อปฺปมาโท อตฺถิ กุสเลสุ ธมฺเมสุ, ตสฺส ยา รตฺติ วา ทิวโส วา อาคจฺฉติ, วุทฺธิเยว ปาฏิกงฺขา กุสเลสุ ธมฺเมสุ, โน ปริหานีติ.}

\prose{อถ โข ภควา ปจฺจุฏฺฐาย อายสฺมนฺตํ สาริปุตฺตํ อามนฺเตสิ\csromandash{} สาธุ สาธุ สาริปุตฺต, ยสฺส กสฺสจิ สาริปุตฺต สทฺธา นตฺถิ กุสเลสุ ธมฺเมสุ, หิรี นตฺถิ.\csromanspacer{} โอตฺตปฺปํ นตฺถิ.\csromanspacer{} ปญฺญา นตฺถิ.\csromanspacer{} วีริยํ นตฺถิ.\csromanspacer{} โสตาวธานํ นตฺถิ.\csromanspacer{} ธมฺมธารณา นตฺถิ.\csromanspacer{} อตฺถูปปริกฺขา นตฺถิ.\csromanspacer{} ธมฺมานุธมฺมปฺปฏิปตฺติ นตฺถิ.\csromanspacer{} อปฺปมาโท นตฺถิ กุสเลสุ ธมฺเมสุ, ตสฺส ยา รตฺติ วา ทิวโส วา อาคจฺฉติ, หานิเยว ปาฏิกงฺขา กุสเลสุ ธมฺเมสุ, โน วุทฺธิ.\csromanspacer{} เสยฺยถาปิ สาริปุตฺต กาฬปกฺเข จนฺทสฺส ยา รตฺติ วา ทิวโส วา อาคจฺฉติ, หายเตว วณฺเณน, หายติ มณฺฑเลน, หายติ อาภาย, หายติ อาโรหปริณาเหน.\csromanspacer{} เอวเมวํ โข สาริปุตฺต ยสฺส กสฺสจิ สทฺธา นตฺถิ กุสเลสุ ธมฺเมสุ ฯเปฯ.\csromanspacer{} อปฺปมาโท นตฺถิ กุสเลสุ ธมฺเมสุ, ตสฺส ยา รตฺติ วา ทิวโส วา อาคจฺฉติ, หานิเยว ปาฏิกงฺขา กุสเลสุ ธมฺเมสุ, โน วุทฺธิ.}

\prose{ยสฺส กสฺสจิ สาริปุตฺต สทฺธา อตฺถิ กุสเลสุ ธมฺเมสุ, หิรี อตฺถิ.\csromanspacer{} โอตฺตปฺปํ อตฺถิ.\csromanspacer{} วีริยํ อตฺถิ.\csromanspacer{} ปญฺญา อตฺถิ.\csromanspacer{} โสตาวธานํ อตฺถิ.\csromanspacer{} ธมฺมธารณา อตฺถิ.\csromanspacer{} อตฺถูปปริกฺขา อตฺถิ.\csromanspacer{} ธมฺมานุธมฺมปฺปฏิปตฺติ อตฺถิ.\csromanspacer{} อปฺปมาโท อตฺถิ กุสเลสุ ธมฺเมสุ, ตสฺส ยา รตฺติ วา ทิวโส วา อาคจฺฉติ, วุทฺธิเยว ปาฏิกงฺขา กุสเลสุ ธมฺเมสุ, โน ปริหานิ.\csromanspacer{} เสยฺยถาปิ สาริปุตฺต ชุณฺหปกฺเข จนฺทสฺส ยา รตฺติ วา ทิวโส วา อาคจฺฉติ, วฑฺฒเตว วณฺเณน, วฑฺฒติ มณฺฑเลน, วฑฺฒติ อาภาย,}

\csromanmatikaanchor{mtk.103}
\csromantocmarkref{subsection}{8. ทุติยนฬกปานสุตฺต}{mtk.103}
\bookmark[dest={mtk.103},level=2]{8. ทุติยนฬกปานสุตฺต}
\vspace{\tipitakaparskip}
\csromancenter{ทสกนิปาตปาฬิ วฑฺฒติ อาโรหปริณาเหน.\csromanspacer{} เอวเมวํ โข สาริปุตฺต ยสฺส กสฺสจิ สทฺธา อตฺถิ กุสเลสุ ธมฺเมสุ ฯเปฯ.\csromanspacer{} อปฺปมาโท อตฺถิ กุสเลสุ ธมฺเมสุ, ตสฺส ยา รตฺติ วา ทิวโส วา อาคจฺฉติ, วุทฺธิเยว ปาฏิกงฺขา กุสเลสุ ธมฺเมสุ, โน ปริหานีติ.\csromanspacer{}\csromanspacer{}อฏฺฐมํ.}
\csromanheader{2}{9. ปฐมกถาวตฺถุสุตฺต}
\proseitem{69}{\csromanfolio{358}เอกํ สมยํ ภควา สาวตฺถิยํ วิหรติ เชตวเน อนาถปิณฺฑิกสฺส อาราเม.\csromanspacer{} เตน โข ปน สมเยน สมฺพหุลา ภิกฺขู ปจฺฉาภตฺตํ ปิณฺฑปาตปฏิกฺกนฺตา อุปฏฺฐานสาลายํ สนฺนิสินฺนา สนฺนิปติตา อเนกวิหิตํ ติรจฺฉานกถํ อนุยุตฺตา วิหรนฺติ.\csromanspacer{} เสยฺยถิทํ, * ราชกถํ โจรกถํ มหามตฺตกถํ เสนากถํ ภยกถํ ยุทฺธกถํ อนฺนกถํ ปานกถํ วตฺถกถํ สยนกถํ มาลากถํ คนฺธกถํ ญาติกถํ ยานกถํ คามกถํ นิคมกถํ นครกถํ ชนปทกถํ อิตฺถิกถํ\footnote{อิตฺถิกถํ ปุริสกถํ (ก) ม ฏฺฐ 3. 156 ปิฏฺเฐ ปสฺสิตพฺพํ.} สูรกถํ วิสิขากถํ กุมฺภฏฺฐานกถํ ปุพฺพเปตกถํ นานตฺตกถํ โลกกฺขายิกํ สมุทฺทกฺขายิกํ อิติภวาภวกถํ อิติ วาติ.}

\prose{อถ โข ภควา สายนฺหสมยํ ปฏิสลฺลานา วุฏฺฐิโต เยน อุปฏฺฐานสาลา เตนุปสงฺกมิ, อุปสงฺกมิตฺวา ปญฺญตฺเต อาสเน นิสีทิ, นิสชฺช โข ภควา ภิกฺขู อามนฺเตสิ “กาย นุตฺถ ภิกฺขเว เอตรติ กถาย สนฺนิสินฺนา สนฺนิปติตา, กา จ ปน โว อนฺตรากถา วิปฺปกตา”ติ.}

\prose{อิธ มยํ ภนฺเต ปจฺฉาภตฺตํ ปิณฺฑปาตปฏิกฺกนฺตา อุปฏฺฐานสาลายํ สนฺนิสินฺนา สนฺนิปติตา, อเนกวิหิตํ ติรจฺฉานกถํ อนุยุตฺตา วิหราม.\csromanspacer{} เสยฺยถิทํ, ราชกถํ โจรกถํ ฯเปฯ อิติภวาภวกถํ อิติ วาติ.\csromanspacer{} น โข ปเนตํ ภิกฺขเว ตุมฺหากํ ปติรูปํ กุลปุตฺตานํ สทฺธาย อคารสฺมา อนคาริยํ ปพฺพชิตานํ, ยํ ตุเมฺห อเนกวิหิตํ ติรจฺฉานกถํ อนุยุตฺตา วิหเรยฺยาถ.\csromanspacer{} เสยฺยถิทํ, ราชกถํ โจรกถํ มหามตฺตกถํ เสนากถํ ภยกถํ ยุทฺธกถํ อนฺนกถํ ปานกถํ วตฺถกถํ สยนกถํ มาลากถํ คนฺธกถํ ญาติกถํ ยานกถํ คามกถํ นิคมกถํ นครกถํ ชนปทกถํ อิตฺถิกถํ สูรกถํ วิสิขากถํ กุมฺภฏฺฐานกถํ ปุพฺพเปตกถํ นานตฺตกถํ โลกกฺขายิกํ สมุทฺทกฺขายิกํ อิติภวาภวกถํ อิติ วาติ.}

\chapterheadpageread{(7) 2. ยมกวคฺค}
\prose{\csromanfolio{359}ทสยิมานิ ภิกฺขเว กถาวตฺถูนิ.\csromanspacer{} กตมานิ ทส? อปฺปิจฺฉกถา สนฺตุฏฺฐิกถา ปวิเวกกถา อสํสคฺคกถา วีริยารมฺภกถา สีลกถา สมาธิกถา ปญฺญากถา วิมุตฺติกถา วิมุตฺติญาณทสฺสนกถาติ, อิมานิ โข ภิกฺขเว ทส กถาวตฺถูนิ.}

\prose{อิเมสํ เจ ตุเมฺห ภิกฺขเว ทสนฺนํ กถาวตฺถูนํ อุปาทายุปาทาย กถํ กเถยฺยาถ, อิเมสมฺปิ จนฺทิมสูริยานํ เอวํมหิทฺธิกานํ เอวํมหานุภาวานํ เตชสา เตชํ ปริยาทิเยยฺยาถ, โก ปน วาโท อญฺญติตฺถิยานํ ปริพฺพาชกานนฺติ.\csromanspacer{}\csromanspacer{}นวมํ.}

\csromanheader{2}{10. ทุติยกถาวตฺถุสุตฺต}
\proseitem{70}{เอกํ สมยํ ภควา สาวตฺถิยํ วิหรติ เชตวเน อนาถปิณฺฑิกสฺส อาราเม.\csromanspacer{} เตน โข ปน สมเยน สมฺพหุลา ภิกฺขู ปจฺฉาภตฺตํ ปิณฺฑปาตปฏิกฺกนฺตา อุปฏฺฐานสาลายํ สนฺนิสินฺนา สนฺนิปติตา อเนกวิหิตํ ติรจฺฉานกถํ อนุยุตฺตา วิหรนฺติ.\csromanspacer{} เสยฺยถิทํ, ราชกถํ โจรกถํ มหามตฺตกถํ ฯเปฯ อิติภวาภวกถํ อิติ วาติ.}

\csromanmatikaanchor{mtk.104}
\csromantocmarkref{subsection}{9. ปฐมกถาวตฺถุสุตฺต}{mtk.104}
\bookmark[dest={mtk.104},level=2]{9. ปฐมกถาวตฺถุสุตฺต}
\prose{ทสยิมานิ ภิกฺขเว ปาสํสานิ ฐานานิ.\csromanspacer{} กตมานิ ทส? อิธ ภิกฺขเว ภิกฺขุ อตฺตนา จ อปฺปิจฺโฉ โหติ, อปฺปิจฺฉกถญฺจ ภิกฺขูนํ กตฺตา โหติ. “อปฺปิจฺโฉ ภิกฺขุ อปฺปิจฺฉกถญฺจ ภิกฺขูนํ กตฺตา”ติ ปาสํสเมตํ ฐานํ.}

\prose{อตฺตนา จ สนฺตุฏฺโฐ โหติ, สนฺตุฏฺฐิกถญฺจ ภิกฺขูนํ กตฺตา โหติ. “สนฺตุฏฺโฐ ภิกฺขุ สนฺตุฏฺฐิกถญฺจ ภิกฺขูนํ กตฺตา”ติ ปาสํสเมตํ ฐานํ.}

\prose{อตฺตนา จ ปวิวิตฺโต โหติ, ปวิเวกกถญฺจ ภิกฺขูนํ กตฺตา โหติ. “ปวิวิตฺโต ภิกฺขุ ปวิเวกกถญฺจ ภิกฺขูนํ กตฺตา”ติ ปาสํสเมตํ ฐานํ.}

\prose{อตฺตนา จ อสํสฏฺโฐ โหติ, อสํสฏฺฐกถญฺจ ภิกฺขูนํ กตฺตา โหติ. “อสํสฏฺโฐ ภิกฺขุ อสํสฏฺฐกถญฺจ ภิกฺขูนํ กตฺตา”ติ ปาสํสเมตํ ฐานํ.}

\prose{อตฺตนา จ อารทฺธวีริโย โหติ, วีริยารมฺภกถญฺจ ภิกฺขูนํ กตฺตา โหติ. “อารทฺธวีริโย ภิกฺขุ วีริยารมฺภกถญฺจ ภิกฺขูนํ กตฺตา”ติ ปาสํสเมตํ ฐานํ.}

\gambhira{ทสกนิปาตปาฬิ}
\csromanmatikaanchor{mtk.105}
\csromanmatikaanchor{mtk.106}
\csromantocmarkref{subsection}{10. ทุติยกถาวตฺถุสุตฺต}{mtk.105}
\bookmark[dest={mtk.105},level=2]{10. ทุติยกถาวตฺถุสุตฺต}
\csromantocmarkref{section}{อุทฺทานคาถา}{mtk.106}
\bookmark[dest={mtk.106},level=1]{อุทฺทานคาถา}
\prose{\csromanfolio{360}อตฺตนา จ สีลสมฺปนฺโน โหติ, สีลสมฺปทากถญฺจ ภิกฺขูนํ กตฺตา โหติ. “สีลสมฺปนฺโน ภิกฺขุ สีลสมฺปทากถญฺจ ภิกฺขูนํ กตฺตา”ติ ปาสํสเมตํ ฐานํ.}

\prose{อตฺตนา จ สมาธิสมฺปนฺโน โหติ, สมาธิสมฺปทากถญฺจ ภิกฺขูนํ กตฺตา โหติ. “สมาธิสมฺปนฺโน ภิกฺขุ สมาธิสมฺปทากถญฺจ ภิกฺขูนํ กตฺตา”ติ ปาสํสเมตํ ฐานํ.}

\prose{อตฺตนา จ ปญฺญาสมฺปนฺโน โหติ, ปญฺญาสมฺปทากถญฺจ ภิกฺขูนํ กตฺตา โหติ. “ปญฺญาสมฺปนฺโน ภิกฺขุ ปญฺญาสมฺปทากถญฺจ ภิกฺขูนํ กตฺตา”ติ ปาสํสเมตํ ฐานํ.}

\prose{อตฺตนา จ วิมุตฺติสมฺปนฺโน โหติ, วิมุตฺติสมฺปทากถญฺจ ภิกฺขูนํ กตฺตา โหติ. “วิมุตฺติสมฺปนฺโน ภิกฺขุ วิมุตฺติสมฺปทากถญฺจ ภิกฺขูนํ กตฺตา”ติ ปาสํสเมตํ ฐานํ.}

\prose{อตฺตนา จ วิมุตฺติญาณทสฺสนสมฺปนฺโน โหติ, วิมุตฺติญาณทสฺสนสมฺปทากถญฺจ ภิกฺขูนํ กตฺตา โหติ. “วิมุตฺติญาณทสฺสนสมฺปนฺโน ภิกฺขุ วิมุตฺติญาณทสฺสนสมฺปทากถญฺจ ภิกฺขูนํ กตฺตา”ติ ปาสํสเมตํ ฐานํ.\csromanspacer{} อิมานิ โข ภิกฺขเว ทส ปาสํสานิ ฐานานีติ.\csromanspacer{}\csromanspacer{}ทสมํ.\csromanspacer{} ยมกวคฺโค ทุติโย.}

\titlehead{\textbf{ตสฺสุทฺทานํ}}
\csromangathameasure{อวิชฺชา ตณฺหา นิฏฺฐา จ, อเวจฺจ เทฺว สุขานิ จ. \\ นฬกปาเน เทฺว วุตฺตา, กถาวตฺถูปเร ทุเวติ.}
\csromangathagroup{\csromangathastanza{อวิชฺชา ตณฺหา นิฏฺฐา จ, อเวจฺจ เทฺว สุขานิ จ. \\ นฬกปาเน เทฺว วุตฺตา, กถาวตฺถูปเร ทุเวติ.}}

\csromanheader{1}{\textbf{(8) 3. อากงฺขวคฺค}}
\csromanheader{2}{1. อากงฺขสุตฺต}
\proseitem{71}{\csromansymbolfootnote{*}{อํ 1. 321; ขุ 1. 274 ปิฏฺเฐสุปิ.}เอกํ สมยํ ภควา สาวตฺถิยํ วิหรติ เชตวเน อนาถปิณฺฑิกสฺส อาราเม.\csromanspacer{} ตตฺร โข ภควา ภิกฺขู อามนฺเตสิ “ภิกฺขโว”ติ. “ภทนฺเต”ติ เต ภิกฺขู ภควโต ปจฺจสฺโสสุํ.\csromanspacer{} ภควา เอตทโวจ\csromandash{}}

\chapterheadpageread{(8) 3. อากงฺขวคฺค}
\prose{\csromanfolio{361}สมฺปนฺนสีลา ภิกฺขเว วิหรถ สมฺปนฺนปาติโมกฺขา, ปาติโมกฺขสํวรสํวุตา วิหรถ อาจารโคจรสมฺปนฺนา, อณุมตฺเตสุ วชฺเชสุ ภยทสฺสาวิโน สมาทาย สิกฺขถ สิกฺขาปเทสุ.}

\prose{อากงฺเขยฺย เจ ภิกฺขเว ภิกฺขุ “สพฺรหฺมจารีนํ ปิโย จสฺสํ มนาโป จ ครุ จ ภาวนีโย จา”ติ, สีเลเสฺววสฺส ปริปูรการี อชฺฌตฺตํ เจโตสมถมนุยุตฺโต อนิรากตชฺฌาโน วิปสฺสนาย สมนฺนาคโต พฺรูเหตา สุญฺญาคารานํ.}

\prose{อากงฺเขยฺย เจ ภิกฺขเว ภิกฺขุ “ลาภี อสฺสํ จีวรปิณฺฑปาตเสนาสนคิลานปจฺจยเภสชฺชปริกฺขารานนฺ”ติ, สีเลเสฺววสฺส ปริปูรการี อชฺฌตฺตํ เจโตสมถมนุยุตฺโต อนิรากตชฺฌาโน วิปสฺสนาย สมนฺนาคโต พฺรูเหตา สุญฺญาคารานํ.}

\prose{อากงฺเขยฺย เจ ภิกฺขเว ภิกฺขุ “เยสาหํ ปริภุญฺชามิ จีวรปิณฺฑปาตเสนาสนคิลานปจฺจยเภสชฺชปริกฺขารานํ, เตสํ เต การา มหปฺผลา อสฺสุ มหานิสํสา”ติ, สีเลเสฺววสฺส ฯเปฯ พฺรูเหตา สุญฺญาคารานํ.}

\prose{อากงฺเขยฺย เจ ภิกฺขเว ภิกฺขุ “เย เม\footnote{เย มํ (ม 1. 40 ปิฏฺเฐ)} เปตา ญาตี สาโลหิตา กาลงฺกตา\footnote{กาลกตา (สี, สฺยา, กํ, อิ)} ปสนฺนจิตฺตา อนุสฺสรนฺติ, เตสํ ตํ มหปฺผลํ อสฺส มหานิสํสนฺ”ติ, สีเลเสฺววสฺส ฯเปฯ พฺรูเหตา สุญฺญาคารานํ.}

\csromanmatikaanchor{mtk.107}
\csromantocmarknopage{section}{(8) 3. อากงฺขวคฺค}{mtk.107}
\bookmark[dest={mtk.107},level=1]{(8) 3. อากงฺขวคฺค}
\prose{อากงฺเขยฺย เจ ภิกฺขเว ภิกฺขุ “สนฺตุฏฺโฐ อสฺสํ อิตรีตรจีวรปิณฺฑปาต เสนาสน คิลานปจฺจย เภสชฺช ปริกฺขาเรนา”ติ, สีเลเสฺววสฺส ฯเปฯ พฺรูเหตา สุญฺญาคารานํ.}

\csromanmatikaanchor{mtk.108}
\csromantocmarkref{subsection}{1. อากงฺขสุตฺต}{mtk.108}
\bookmark[dest={mtk.108},level=2]{1. อากงฺขสุตฺต}
\prose{อากงฺเขยฺย เจ ภิกฺขเว ภิกฺขุ “ขโม อสฺสํ สีตสฺส อุณฺหสฺส ชิฆจฺฉาย ปิปาสาย ฑํส มกส วาตาตป สรีสป สมฺผสฺสานํ ทุรุตฺตานํ ทุราคตานํ วจนปถานํ, อุปฺปนฺนานํ สารีริกานํ เวทนานํ ทุกฺขานํ ติพฺพานํ\footnote{ติปฺปานํ (สี, สฺยา, กํ, อิ)} ขรานํ กฏุกานํ อสาตานํ อมนาปานํ ปาณหรานํ อธิวาสกชาติโก อสฺสนฺ”ติ, สีเลเสฺววสฺส ฯเปฯ พฺรูเหตา สุญฺญาคารานํ.}

\prose{อากงฺเขยฺย เจ ภิกฺขเว ภิกฺขุ “อรติรติสโห อสฺสํ, น จ มํ อรติรติ สเหยฺย, อุปฺปนฺนํ อรติรติํ อภิภุยฺย อภิภุยฺย วิหเรยฺยนฺ”ติ, สีเลเสฺววสฺส ฯเปฯ พฺรูเหตา สุญฺญาคารานํ.}

\gambhira{ทสกนิปาตปาฬิ}
\prose{\csromanfolio{362}อากงฺเขยฺย เจ ภิกฺขเว ภิกฺขุ “ภยเภรวสโห อสฺสํ, น จ มํ ภยเภรโว สเหยฺย, อุปฺปนฺนํ ภยเภรวํ อภิภุยฺย อภิภุยฺย วิหเรยฺยนฺ”ติ, สีเลเสฺววสฺส ฯเปฯ พฺรูเหตา สุญฺญาคารานํ.}

\prose{อากงฺเขยฺย เจ ภิกฺขเว ภิกฺขุ “จตุนฺนํ ฌานานํ อาภิเจตสิกานํ ทิฏฺฐธมฺมสุขวิหารานํ นิกามลาภี อสฺสํ อกิจฺฉลาภี อกสิรลาภี”ติ, สีเลเสฺววสฺส ฯเปฯ พฺรูเหตา สุญฺญาคารานํ.}

\prose{อากงฺเขยฺย เจ ภิกฺขเว ภิกฺขุ “อาสวานํ ขยา อนาสวํ เจโตวิมุตฺติํ ปญฺญาวิมุตฺติํ ทิฏฺเฐว ธมฺเม สยํ อภิญฺญา สจฺฉิกตฺวา อุปสมฺปชฺช วิหเรยฺยนฺ”ติ, สีเลเสฺววสฺส ปริปูรการี อชฺฌตฺตํ เจโตสมถมนุยุตฺโต อนิรากตชฺฌาโน วิปสฺสนาย สมนฺนาคโต พฺรูเหตา สุญฺญาคารานํ.}

\prose{“สมฺปนฺนสีลา ภิกฺขเว วิหรถ สมฺปนฺนปาติโมกฺขา, ปาติโมกฺขสํวรสํวุตา วิหรถ อาจารโคจรสมฺปนฺนา, อณุมตฺเตสุ วชฺเชสุ ภยทสฺสาวิโน สมาทาย สิกฺขถ สิกฺขาปเทสู”ติ อิติ ยํตํ วุตฺตํ, อิทเมตํ ปฏิจฺจ วุตฺตนฺติ.\csromanspacer{}\csromanspacer{}ปฐมํ.}

\csromanheader{2}{2. กณฺฏกสุตฺต}
\proseitem{72}{เอกํ สมยํ ภควา เวสาลิยํ วิหรติ มหาวเน กูฏาคารสาลายํ สมฺพหุเลหิ อภิญฺญาเตหิ อภิญฺญาเตหิ เถเรหิ สาวเกหิ สทฺธิํ อายสฺมตา จ จาเลน\footnote{ปาเลน (สฺยา)} อายสฺมตา จ อุปจาเลน\footnote{อุปฺปาเลน (สฺยา)} อายสฺมตา จ กุกฺกุเฏน\footnote{กกฺกเฏน (สี, สฺยา)} อายสฺมตา จ กฬิมฺเภน\footnote{กวิมฺเภน (สี)} อายสฺมตา จ นิกเฏน\footnote{กเฏน (สี)} อายสฺมตา จ กฏิสฺสเหน อญฺเญหิ จ อภิญฺญาเตหิ อภิญฺญาเตหิ เถเรหิ สาวเกหิ สทฺธิํ.}

\prose{เตน โข ปน สมเยน สมฺพหุลา อภิญฺญาตา อภิญฺญาตา ลิจฺฉวี ภเทฺรหิ ภเทฺรหิ ยาเนหิ ปรปุราย\footnote{ปรํปุราย (สฺยา-ฏฺฐ)} อุจฺจาสทฺทา มหาสทฺทา มหาวนํ อชฺโฌคาหนฺติ ภควนฺตํ ทสฺสนาย.\csromanspacer{} อถ โข เตสํ อายสฺมนฺตานํ เอตทโหสิ “อิเม โข สมฺพหุลา อภิญฺญาตา อภิญฺญาตา ลิจฺฉวี ภเทฺรหิ ภเทฺรหิ ยาเนหิ ปรปุราย อุจฺจาสทฺทา มหาสทฺทา มหาวนํ}

\vspace{\tipitakaparskip}
\csromancenter{(8) 3.\csromanspacer{} อากงฺขวคฺค อชฺโฌคาหนฺติ ภควนฺตํ ทสฺสนาย.\csromanspacer{} สทฺทกณฺฏกา โข ปน ฌานา วุตฺตา ภควตา, ยํนูน มยํ เยน โคสิงฺคสาลวนทาโย เตนุปสงฺกเมยฺยาม, ตตฺถ มยํ อปฺปสทฺทา อปฺปากิณฺณา ผาสุํ\footnote{ผาสุ (สฺยา, ก)} วิหเรยฺยามา”ติ.\csromanspacer{} อถ โข เต อายสฺมนฺโต เยน โคสิงฺคสาลวนทาโย เตนุปสงฺกมิํสุ, ตตฺถ เต อายสฺมนฺโต อปฺปสทฺทา อปฺปากิณฺณา ผาสุํ วิหรนฺติ.}
\prose{\csromanfolio{363}อถ โข ภควา ภิกฺขู อามนฺเตสิ “กหํ นุ โข ภิกฺขเว จาโล กหํ อุปจาโล, กหํ กุกฺกุโฏ, กหํ กฬิมฺโภ, กหํ นิกโฏ, กหํ กฏิสฺสโห, กหํ นุ โข เต ภิกฺขเว เถรา สาวกา คตา”ติ.}

\prose{อิธ ภนฺเต เตสํ อายสฺมนฺตานํ เอตทโหสิ “อิเม โข สมฺพหุลา อภิญฺญาตา อภิญฺญาตา ลิจฺฉวี ภเทฺรหิ ภเทฺรหิ ยาเนหิ ปรปุราย อุจฺจาสทฺทา มหาสทฺทา มหาวนํ อชฺโฌคาหนฺติ ภควนฺตํ ทสฺสนาย.\csromanspacer{} สทฺทกณฺฏกา โข ปน ฌานา วุตฺตา ภควตา, ยํนูน มยํ เยน โคสิงฺคสาลวนทาโย เตนุปสงฺกเมยฺยาม, ตตฺถ มยํ อปฺปสทฺทา อปฺปากิณฺณา ผาสุํ วิหเรยฺยามา”ติ.\csromanspacer{} อถ โข เต ภนฺเต อายสฺมนฺโต เยน โคสิงฺคสาลวนทาโย เตนุปสงฺกมิํสุ, ตตฺถ เต อายสฺมนฺโต อปฺปสทฺทา อปฺปากิณฺณา ผาสุํ วิหรนฺตีติ.}

\prose{สาธุ สาธุ ภิกฺขเว ยถา เต มหาสาวกา สมฺมา พฺยากรมานา พฺยากเรยฺยุํ, สทฺทกณฺฏกา หิ ภิกฺขเว ฌานา วุตฺตา มยา.}

\prose{ทสยิเม ภิกฺขเว กณฺฏกา.\csromanspacer{} กตเม ทส? ปวิเวการามสฺส สงฺคณิการามตา กณฺฏโก, อสุภนิมิตฺตานุโยคํ อนุยุตฺตสฺส สุภนิมิตฺตานุโยโค กณฺฏโก, อินฺทฺริเยสุ คุตฺตทฺวารสฺส วิสูกทสฺสนํ กณฺฏโก, พฺรหฺมจริยสฺส มาตุคามูปจาโร\footnote{มาตุคาโมปวิจาโร (สี), มาตุคามูปวิจโร (ก)} กณฺฏโก, * ปฐมสฺส ฌานสฺส สทฺโท กณฺฏโก, ทุติยสฺส ฌานสฺส วิตกฺกวิจารา กณฺฏกา, ตติยสฺส ฌานสฺส ปีติ กณฺฏโก, จตุตฺถสฺส ฌานสฺส อสฺสาสปสฺสาโส กณฺฏโก, สญฺญาเวทยิต- นิโรธสมาปตฺติยา สญฺญา จ เวทนา จ กณฺฏโก, ราโค กณฺฏโก, โทโส กณฺฏโก, โมโห กณฺฏโก.}

\gambhira{ทสกนิปาตปาฬิ}
\csromanmatikaanchor{mtk.109}
\csromantocmarkref{subsection}{2. กณฺฏกสุตฺต}{mtk.109}
\bookmark[dest={mtk.109},level=2]{2. กณฺฏกสุตฺต}
\prose{\csromanfolio{364}อกณฺฏกา ภิกฺขเว วิหรถ, นิกฺกณฺฏกา ภิกฺขเว วิหรถ, อกณฺฏกนิกฺกณฺฏกา ภิกฺขเว วิหรถ, อกณฺฏกา ภิกฺขเว อรหนฺโต, นิกฺกณฺฏกา ภิกฺขเว อรหนฺโต, อกณฺฏกนิกฺกณฺฏกา ภิกฺขเว อรหนฺโตติ.\csromanspacer{}\csromanspacer{}ทุติยํ.}

\csromanheader{2}{3. อิฏฺฐธมฺมสุตฺต}
\proseitem{73}{ทสยิเม ภิกฺขเว ธมฺมา อิฏฺฐา กนฺตา มนาปา ทุลฺลภา โลกสฺมิํ.\csromanspacer{} กตเม ทส? โภคา อิฏฺฐา กนฺตา มนาปา ทุลฺลภา โลกสฺมิํ, วณฺโณ อิฏฺโฐ กนฺโต มนาโป ทุลฺลโภ โลกสฺมิํ, อาโรคฺยํ อิฏฺฐํ กนฺตํ มนาปํ ทุลฺลภํ โลกสฺมิํ, สีลํ อิฏฺฐํ กนฺตํ มนาปํ ทุลฺลภํ โลกสฺมิํ, พฺรหฺมจริยํ อิฏฺฐํ กนฺตํ มนาปํ ทุลฺลภํ โลกสฺมิํ, มิตฺตา อิฏฺฐา กนฺตา มนาปา ทุลฺลภา โลกสฺมิํ, พาหุสจฺจํ อิฏฺฐํ กนฺตํ มนาปํ ทุลฺลภํ โลกสฺมิํ, ปญฺญา อิฏฺฐา กนฺตา มนาปา ทุลฺลภา โลกสฺมิํ, ธมฺมา อิฏฺฐา กนฺตา มนาปา ทุลฺลภา โลกสฺมิํ, สคฺคา อิฏฺฐา กนฺตา มนาปา ทุลฺลภา โลกสฺมิํ.}

\prose{อิเมสํ โข ภิกฺขเว ทสนฺนํ ธมฺมานํ อิฏฺฐานํ กนฺตานํ มนาปานํ ทุลฺลภานํ โลกสฺมิํ ทส ธมฺมา ปริปนฺถา\footnote{ปริพนฺธา(ก)}.\csromanspacer{} อาลสฺยํ อนุฏฺฐานํ โภคานํ ปริปนฺโถ, อมณฺฑนา อวิภูสนา วณฺณสฺส ปริปนฺโถ, อสปฺปายกิริยา อาโรคฺยสฺส ปริปนฺโถ, ปาปมิตฺตตา สีลานํ ปริปนฺโถ, อินฺทฺริย-อสํวโร พฺรหฺมจริยสฺส ปริปนฺโถ, วิสํวาทนา มิตฺตานํ ปริปนฺโถ, อสชฺฌายกิริยา พาหุสจฺจสฺส ปริปนฺโถ, อสุสฺสูสา อปริปุจฺฉา ปญฺญาย ปริปนฺโถ, อนนุโยโค อปจฺจเวกฺขณา ธมฺมานํ ปริปนฺโถ, มิจฺฉาปฏิปตฺติ สคฺคานํ ปริปนฺโถ.\csromanspacer{} อิเมสํ โข ภิกฺขเว ทสนฺนํ อิฏฺฐานํ กนฺตานํ มนาปานํ ทุลฺลภานํ โลกสฺมิํ อิเม ทส ธมฺมา ปริปนฺถา.}

\prose{อิเมสํ โข ภิกฺขเว ทสนฺนํ ธมฺมานํ อิฏฺฐานํ กนฺตานํ มนาปานํ ทุลฺลภานํ โลกสฺมิํ ทส ธมฺมา อาหารา.\csromanspacer{} อุฏฺฐานํ อนาลสฺยํ โภคานํ อาหาโร, มณฺฑนา วิภูสนา วณฺณสฺส อาหาโร, สปฺปายกิริยา อาโรคฺยสฺส อาหาโร, กลฺยาณมิตฺตตา สีลานํ อาหาโร, อินฺทฺริยสํวโร พฺรหฺมจริยสฺส อาหาโร, อวิสํวาทนา มิตฺตานํ อาหาโร, สชฺฌายกิริยา พาหุสจฺจสฺส อาหาโร, สุสฺสูสา ปริปุจฺฉา ปญฺญาย อาหาโร, อนุโยโค ปจฺจเวกฺขณา ธมฺมานํ อาหาโร, สมฺมาปฏิปตฺติ สคฺคานํ}

\vspace{\tipitakaparskip}
\csromancenter{(8) 3.\csromanspacer{} อากงฺขวคฺค อาหาโร.\csromanspacer{} อิเมสํ โข ภิกฺขเว ทสนฺนํ ธมฺมานํ อิฏฺฐานํ กนฺตานํ มนาปานํ ทุลฺลภานํ โลกสฺมิํ อิเม ทส ธมฺมา อาหาราติ.\csromanspacer{}\csromanspacer{}ตติยํ.}
\csromanheader{2}{4. วฑฺฒิสุตฺต}
\csromanmatikaanchor{mtk.110}
\csromanmatikaanchor{mtk.111}
\csromantocmarkref{subsection}{3. อิฏฺฐธมฺมสุตฺต}{mtk.110}
\bookmark[dest={mtk.110},level=2]{3. อิฏฺฐธมฺมสุตฺต}
\csromantocmarkref{subsection}{4. วุฑฺฒิสุตฺต}{mtk.111}
\bookmark[dest={mtk.111},level=2]{4. วุฑฺฒิสุตฺต}
\proseitem{74}{\csromanfolio{365}ทสหิ ภิกฺขเว วฑฺฒีหิ วฑฺฒมาโน อริยสาวโก อริยาย วฑฺฒิยา วฑฺฒติ, สาราทายี จ โหติ วราทายี กายสฺส.\csromanspacer{} กตเมหิ ทสหิ? เขตฺตวตฺถูหิ วฑฺฒติ, ธนธญฺเญน วฑฺฒติ, ปุตฺตทาเรหิ วฑฺฒติ, ทาสกมฺมกรโปริเสหิ วฑฺฒติ, จตุปฺปเทหิ วฑฺฒติ สทฺธาย วฑฺฒติ, สีเลน วฑฺฒติ, สุเตน วฑฺฒติ, จาเคน วฑฺฒติ, ปญฺญาย วฑฺฒติ.\csromanspacer{} อิเมหิ โข ภิกฺขเว ทสหิ วฑฺฒีหิ วฑฺฒมาโน อริยสาวโก อริยาย วฑฺฒิยา วฑฺฒติ, สาราทายี จ โหติ วราทายี กายสฺสาติ.}

\csromangathameasure{ธเนน ธญฺเญน จ โยธ วฑฺฒติ, \\ ปุตฺเตหิ ทาเรหิ จตุปฺปเทหิ จ. \\ ส โภควา โหติ ยสสฺสิ ปูชิโต, \\ ญาตีหิ มิตฺเตหิ อโถปิ ราชุภิ. \\ สทฺธาย สีเลน จ โยธ วฑฺฒติ, \\ ปญฺญาย จาเคน สุเตน จูภยํ. \\ โส ตาทิโส สปฺปุริโส วิจกฺขโณ, \\ ทิฏฺเฐว ธมฺเม อุภเยน วฑฺฒตีติ.จตุตฺถํ.}
\csromangathagroup{\csromangathastanza{ธเนน ธญฺเญน จ โยธ วฑฺฒติ, \\ ปุตฺเตหิ ทาเรหิ จตุปฺปเทหิ จ. \\ ส โภควา โหติ ยสสฺสิ ปูชิโต, \\ ญาตีหิ มิตฺเตหิ อโถปิ ราชุภิ.}\csromangathastanza{สทฺธาย สีเลน จ โยธ วฑฺฒติ, \\ ปญฺญาย จาเคน สุเตน จูภยํ. \\ โส ตาทิโส สปฺปุริโส วิจกฺขโณ, \\ ทิฏฺเฐว ธมฺเม อุภเยน วฑฺฒตีติ.\csromanspacer{}\csromanspacer{}จตุตฺถํ.}}

\csromanheader{2}{5. มิคสาลาสุตฺต}
\proseitem{75}{เอกํ สมยํ ภควา สาวตฺถิยํ วิหรติ เชตวเน อนาถปิณฺฑิกสฺส อาราเม.\csromanspacer{} อถ โข อายสฺมา อานนฺโท ปุพฺพณฺหสมยํ นิวาเสตฺวา ปตฺตจีวรมาทาย เยน มิคสาลาย อุปาสิกาย นิเวสนํ เตนุปสงฺกมิ, อุปสงฺกมิตฺวา ปญฺญตฺเต อาสเน นิสีทิ.\csromanspacer{} อถ โข มิคสาลา อุปาสิกา เยนายสฺมา อานนฺโท เตนุปสงฺกมิ, อุปสงฺกมิตฺวา อายสฺมนฺตํ อานนฺทํ อภิวาเทตฺวา เอกมนฺตํ นิสีทิ, เอกมนฺตํ นิสินฺนา โข มิคสาลา อุปาสิกา อายสฺมนฺตํ อานนฺทํ เอตทโวจ\csromandash{}}

\prose{กถํ กถํ นามายํ ภนฺเต อานนฺท ภควตา ธมฺโม เทสิโต อญฺเญยฺโย.\csromanspacer{} ยตฺร หิ นาม พฺรหฺมจารี จ อพฺรหฺมจารี จ อุโภ สมสมคติกา ภวิสฺสนฺติ อภิสมฺปรายํ, ปิตา เม ภนฺเต ปุราโณ พฺรหฺมจารี}

\vspace{\tipitakaparskip}
\csromancenter{ทสกนิปาตปาฬิ โหติ อาราจารี\footnote{อนาจารี (ก)} วิรโต เมถุนา คามธมฺมา, โส กาลงฺกโต ภควตา พฺยากโต “สกทาคามี สตฺโต\footnote{สกทาคามิสตฺโต (สี, สฺยา, อิ)} ตุสิตํ กายํ อุปปนฺโน”ติ.\csromanspacer{} ปิตามโห เม\footnote{เปตฺตาปิ โย เม (สี), ปิตุ ปิโย เม (สฺยา) อํ 2. 305 ปิฏฺเฐปิ.} ภนฺเต อิสิทตฺโต อพฺรหฺมจารี อโหสิ สทารสนฺตุฏฺโฐ, โสปิ กาลงฺกโต ภควตา พฺยากโต “สกทาคามี สตฺโต ตุสิตํ กายํ อุปปนฺโน”ติ.}
\prose{\csromanfolio{366}กถํ กถํ นามายํ ภนฺเต อานนฺท ภควตา ธมฺมา เทสิโต อญฺเญยฺโย, ยตฺร หิ นาม พฺรหฺมจารี จ อพฺรหฺมจารี จ อุโภ สมสมคติกา ภวิสฺสนฺติ อภิสมฺปรายนฺติ.\csromanspacer{} เอวํ โข ปเนตํ ภคินิ ภควตา พฺยากตนฺติ.}

\prose{อถ โข อายสฺมา อานนฺโท มิคสาลาย อุปาสิกาย นิเวสเน ปิณฺฑปาตํ คเหตฺวา อุฏฺฐายาสนา ปกฺกามิ.\csromanspacer{} อถ โข อายสฺมา อานนฺโท ปจฺฉาภตฺตํ ปิณฺฑปาตปฏิกฺกนฺโต เยน ภควา เตนุปสงฺกมิ, อุปสงฺกมิตฺวา ภควนฺตํ อภิวาเทตฺวา เอกมนฺตํ นิสีทิ, เอกมนฺตํ นิสินฺโน โข อายสฺมา อานนฺโท ภควนฺตํ เอตทโวจ\csromandash{}}

\prose{อิธาหํ ภนฺเต ปุพฺพณฺหสมยํ นิวาเสตฺวา ปตฺตจีวรมาทาย เยน มิคสาลาย อุปาสิกาย นิเวสนํ เตนุปสงฺกมิํ, อุปสงฺกมิตฺวา ปญฺญตฺเต อาสเน นิสีทิํ.\csromanspacer{} อถ โข ภนฺเต มิคสาลา อุปาสิกา เยนาหํ เตนุปสงฺกมิ, อุปสงฺกมิตฺวา มํ อภิวาเทตฺวา เอกมนฺตํ นิสีทิ, เอกมนฺตํ นิสินฺนา โข ภนฺเต มิคสาลา อุปาสิกา มํ เอตทโวจ\csromandash{}}

\prose{กถํ กถํ นามายํ ภนฺเต อานนฺท ภควตา ธมฺโม เทสิโต อญฺเญยฺโย, ยตฺร หิ นาม พฺรหฺมจารี จ อพฺรหฺมจารี จ อุโภ สมสมคติกา ภวิสฺสนฺติ อภิสมฺปรายํ, ปิตา เม ภนฺเต ปุราโณ พฺรหฺมจารี อโหสิ อาราจารี วิรโต เมถุนา คามธมฺมา, โส กาลงฺกโต ภควตา พฺยากโต “สกทาคามี สตฺโต ตุสิตํ กายํ อุปปนฺโน”ติ.\csromanspacer{} ปิตามโห เม ภนฺเต อิสิทตฺโต อพฺรหฺมจารี อโหสิ สทารสนฺตุฏฺโฐ, โสปิ กาลงฺกโต ภควตา พฺยากโต “สกทาคามี สตฺโต ตุสิตํ กายํ อุปปนฺโน”ติ.}

\chapterheadpageread{(8) 3. อากงฺขวคฺค}
\csromanmatikaanchor{mtk.112}
\csromantocmarkref{subsection}{5. มิคสาลาสุตฺต}{mtk.112}
\bookmark[dest={mtk.112},level=2]{5. มิคสาลาสุตฺต}
\prose{\csromanfolio{367}กถํ กถํ นามายํ ภนฺเต อานนฺท ภควตา ธมฺโม เทสิโต อญฺเญยฺโย, ยตฺร หิ นาม พฺรหฺมจารี จ อพฺรหฺมจารี จ อุโภ สมสมคติกา ภวิสฺสนฺติ อภิสมฺปรายนฺติ.\csromanspacer{} เอวํ วุตฺเต อหํ ภนฺเต มิคสาลํ อุปาสิกํ เอตทโวจํ “เอวํ โข ปเนตํ ภคินิ ภควตา พฺยากตนฺ”ติ.}

\prose{กา จานนฺท มิคสาลา อุปาสิกา พาลา อพฺยตฺตา อมฺมกา อมฺมกปญฺญา\footnote{อมฺพกา อมฺพกปญฺญา (สี, อิ), อนฺธกา อนฺธกปญฺญา (สฺยา)}, เก จ ปุริสปุคฺคลปโรปริเย ญาเณ.}

\prose{ทสยิเม อานนฺท ปุคฺคลา สนฺโต สํวิชฺชมานา โลกสฺมิํ.\csromanspacer{} กตเม ทส? อิธานนฺท เอกจฺโจ ทุสฺสีโล โหติ, ตญฺจ เจโตวิมุตฺติํ ปญฺญาวิมุตฺติํ ยถาภูตํ นปฺปชานาติ, ยตฺถสฺส ตํ ทุสฺสิลฺยํ อปริเสสํ นิรุชฺฌติ.\csromanspacer{} ตสฺส สวเนนปิ อกตํ โหติ, พาหุสจฺเจนปิ อกตํ โหติ, ทิฏฺฐิยาปิ อปฺปฏิวิทฺธํ โหติ สามายิกมฺปิ วิมุตฺติํ น ลภติ.\csromanspacer{} โส กายสฺส เภทา ปรํ มรณา หานาย ปเรติ โน วิเสสาย, หานคามีเยว โหติ โน วิเสสคามี. (1)}

\prose{อิธ ปนานนฺท เอกจฺโจ ปุคฺคโล ทุสฺสีโล โหติ, ตญฺจ เจโตวิมุตฺติํ ปญฺญาวิมุตฺติํ ยถาภูตํ ปชานาติ, ยตฺถสฺส ตํ ทุสฺสิลฺยํ อปริเสสํ นิรุชฺฌติ.\csromanspacer{} ตสฺส สวเนนปิ กตํ โหติ, พาหุสจฺเจนปิ กตํ โหติ, ทิฏฺฐิยาปิ ปฏิวิทฺธํ\footnote{สุปฺปฏิวิทฺธํ (สฺยา)} โหติ, สามายิกมฺปิ วิมุตฺติํ ลภติ.\csromanspacer{} โส กายสฺส เภทา ปรํ มรณา วิเสสาย ปเรติ โน หานาย, วิเสสคามีเยว โหติ โน หานคามี. (2)}

\prose{ตตฺรานนฺท ปมาณิกา ปมิณนฺติ “อิมสฺสปิ เตว ธมฺมา อปรสฺสปิ เตว ธมฺมา, กสฺมา เนสํ เอโก หีโน เอโก ปณีโต”ติ.\csromanspacer{} ตํ หิ เตสํ อานนฺท โหติ ทีฆรตฺตํ อหิตาย ทุกฺขาย.}

\prose{ตตฺรานนฺท ยฺวายํ ปุคฺคโล ทุสฺสีโล โหติ, ตญฺจ เจโตวิมุตฺติํ ปญฺญาวิมุตฺติํ ยถาภูตํ ปชานาติ, ยตฺถสฺส ตํ ทุสฺสิลฺยํ อปริเสสํ นิรุชฺฌติ.\csromanspacer{} ตสฺส สวเนนปิ กตํ โหติ, พาหุสจฺเจนปิ กตํ โหติ, ทิฏฺฐิยาปิ ปฏิวิทฺธํ โหติ, สามายิกมฺปิ วิมุตฺติํ ลภติ.\csromanspacer{} อยํ อานนฺท ปุคฺคโล อมุนา ปุริเมน ปุคฺคเลน อภิกฺกนฺตตโร จ ปณีตตโร จ.\csromanspacer{} ตํ กิสฺส เหตุ? อิมํ หานนฺท ปุคฺคลํ ธมฺมโสโต นิพฺพหติ, ตทนฺตรํ โก ชาเนยฺย อญฺญตฺร ตถาคเตน.\csromanspacer{} ตสฺมาติหานนฺท มา ปุคฺคเลสุ ปมาณิกา}

\vspace{\tipitakaparskip}
\csromancenter{ทสกนิปาตปาฬิ อหุวตฺถ, มา ปุคฺคเลสุ ปมาณํ คณฺหิตฺถ, ขญฺญติ หานนฺท ปุคฺคเลสุ ปมาณํ คณฺหนฺโต, อหํ วา อานนฺท\footnote{อหญฺจานนฺท (สี, สฺยา, ก) อํ 2. 307 ปิฏฺเฐ ปสฺสิตพฺพํ.} ปุคฺคเลสุ ปมาณํ คเณฺหยฺยํ, โย วา ปนสฺส มาทิโส.}
\prose{\csromanfolio{368}อิธ ปนานนฺท เอกจฺโจ ปุคฺคโล สีลวา โหติ, ตญฺจ เจโตวิมุตฺติํ ปญฺญาวิมุตฺติํ ยถาภูตํ นปฺปชานาติ, ยตฺถสฺส ตํ สีลํ อปริเสสํ นิรุชฺฌติ.\csromanspacer{} ตสฺส สวเนนปิ อกตํ โหติ, พาหุสจฺเจนปิ อกตํ โหติ, ทิฏฺฐิยาปิ อปฺปฏิวิทฺธํ โหติ, สามายิกมฺปิ วิมุตฺติํ น ลภติ.\csromanspacer{} โส กายสฺส เภทา ปรํ มรณา หานาย ปเรติ โน วิเสสาย, หานคามีเยว โหติ โน วิเสสคามี. (3)}

\prose{อิธ ปนานนฺท เอกจฺโจ ปุคฺคโล สีลวา โหติ, ตญฺจ เจโตวิมุตฺติํ ปญฺญาวิมุตฺติํ ยถาภูตํ ปชานาติ, ยตฺถสฺส ตํ สีลํ อปริเสสํ นิรุชฺฌติ.\csromanspacer{} ตสฺส สวเนนปิ กตํ โหติ, พาหุสจฺเจนปิ กตํ โหติ, ทิฏฺฐิยาปิ ปฏิวิทฺธํ โหติ, สามายิกมฺปิ วิมุตฺติํ ลภติ.\csromanspacer{} โส กายสฺส เภทา ปรํ มรณา วิเสสาย ปเรติ โน หานาย, วิเสสคามิเยว โหติ โน หานคามี. (4)}

\prose{ตตฺรานนฺท ปมาณิกา ปมิณนฺติ ฯเปฯ อหํ วา อานนฺท ปุคฺคเลสุ ปมาณํ คเณฺหยฺยํ, โย วา ปนสฺส มาทิโส.}

\prose{อิธ ปนานนฺท เอกจฺโจ ปุคฺคโล ติพฺพราโค โหติ, ตญฺจ เจโตวิมุตฺติํ ปญฺญาวิมุตฺติํ ยถาภูตํ นปฺปชานาติ, ยตฺถสฺส โส ราโค อปริเสโส นิรุชฺฌติ.\csromanspacer{} ตสฺส สเวเนนปิ อกตํ โหติ, พาหุสจฺเจนปิ อกตํ โหติ, ทิฏฺฐิยาปิ อปฺปฏิวิทฺธํ โหติ, สามายิกมฺปิ วิมุตฺติํ น ลภติ.\csromanspacer{} โส กายสฺส เภทา ปรํ มรณา หานาย ปเรติ โน วิเสสาย, หานคามีเยว โหติ โน วิเสสคามี. (5)}

\prose{อิธ ปนานนฺท เอกจฺโจ ปุคฺคโล ติพฺพราโค โหติ, ตญฺจ เจโตวิมุตฺติํ ปญฺญาวิมุตฺติํ ยถาภูตํ ปชานาติ, ยตฺถสฺส โส ราโค อปริเสโส นิรุชฺฌติ.\csromanspacer{} ตสฺส สวเนนปิ กตํ โหติ, พาหุสจฺเจนปิ กตํ โหติ, ทิฏฺฐิยาปิ ปฏิวิทฺธํ โหติ, สามายิกมฺปิ วิมุตฺติํ ลภติ.\csromanspacer{} โส กายสฺส เภทา ปรํ มรณา วิเสสาย ปเรติ โน หานาย, วิเสสคามีเยว โหติ โน หานคามี. (6)}

\chapterheadpageread{(8) 3. อากงฺขวคฺค}
\prose{\csromanfolio{369}ตตฺรานนฺท ปมาณิกา ปมิณนฺติ ฯเปฯ อหํ วา อานนฺท ปุคฺคเลสุ ปมาณํ คเณฺหยฺยํ.\csromanspacer{} โย วา ปนสฺส มาทิโส.}

\prose{อิธ ปนานนฺท เอกจฺโจ ปุคฺคโล โกธโน โหติ, ตญฺจ เจโตวิมุตฺติํ ปญฺญาวิมุตฺติํ ยถาภูตํ นปฺปชานาติ, ยตฺถสฺส โส โกโธ อปริเสโส นิรุชฺฌติ.\csromanspacer{} ตสฺส สวเนนปิ อกตํ โหติ, พาหุสจฺเจนปิ อกตํ โหติ, ทิฏฺฐิยาปิ อปฺปฏิวิทฺธํ โหติ, สามายิกมฺปิ วิมุตฺติํ น ลภติ.\csromanspacer{} โส กายสฺส เภทา ปรํ มรณา หานาย ปเรติ โน วิเสสาย, หานคามีเยว โหติ โน วิเสสคามี. (7)}

\prose{อิธ ปนานนฺท เอกจฺโจ ปุคฺคโล โกธโน โหติ, ตญฺจ เจโตวิมุตฺติํ ปญฺญาวิมุตฺติํ ยถาภูตํ ปชานาติ, ยตฺถสฺส โส โกโธ อปริเสโส นิรุชฺฌติ.\csromanspacer{} ตสฺส สวเนนปิ กตํ โหติ, พาหุสจฺเจนปิ กตํ โหติ, ทิฏฺฐิยาปิ ปฏิวิทฺธํ โหติ, สามายิกมฺปิ วิมุตฺติํ ลภติ.\csromanspacer{} โส กายสฺส เภทา ปรํ มรณา วิเสสาย ปเรติ โน หานาย, วิเสสคามีเยว โหติ โน หานคามี. (8)}

\prose{ตตฺรานนฺท ปมาณิกา ปมิณนฺติ ฯเปฯ อหํ วา อานนฺท ปุคฺคเลสุ ปมาณํ คเณฺหยฺยํ โย วา ปนสฺส มาทิโส.}

\prose{อิธ ปนานนฺท เอกจฺโจ ปุคฺคโล อุทฺธโต โหติ, ตญฺจ เจโตวิมุตฺติํ ปญฺญาวิมุตฺติํ ยถาภูตํ นปฺปชานาติ, ยตฺถสฺส ตํ อุทฺธจฺจํ อปริเสสํ นิรุชฺฌติ.\csromanspacer{} ตสฺส สวเนนปิ อกตํ โหติ, พาหุสจฺเจนปิ อกตํ โหติ, ทิฏฺฐิยา อปฺปฏิวิทฺธํ โหติ, สามายิกมฺปิ วิมุตฺติํ น ลภติ.\csromanspacer{} โส กายสฺส เภทา ปรํ มรณา หานาย ปเรติ โน วิเสสาย, หานคามีเยว โหติ โน วิเสสคามี. (9)}

\prose{อิธ ปนานนฺท เอกจฺโจ ปุคฺคโล อุทฺธโต โหติ, ตญฺจ เจโตวิมุตฺติํ ปญฺญาวิมุตฺติํ ยถาภูตํ ปชานาติ, ยตฺถสฺส ตํ อุทฺธจฺจํ อปริเสสํ นิรุชฺฌติ.\csromanspacer{} ตสฺส สวเนนปิ กตํ โหติ, พาหุสจฺเจนปิ กตํ โหติ, ทิฏฺฐิยาปิ ปฏิวิทฺธํ โหติ, สามายิกมฺปิ วิมุตฺติํ ลภติ.\csromanspacer{} โส กายสฺส เภทา ปรํ มรณา วิเสสาย ปเรติ โน หานาย, วิเสสคามีเยว โหติ โน หานคามี. (10)}

\gambhira{ทสกนิปาตปาฬิ}
\prose{\csromanfolio{370}ตตฺรานนฺท ปมาณิกา ปมิณนฺติ “อิมสฺสปิ เตว ธมฺมา อปรสฺสปิ เตว ธมฺมา, กสฺมา เนสํ เอโก หีโน เอโก ปณีโต”ติ.\csromanspacer{} ตํ หิ เตสํ อานนฺท โหติ ทีฆรตฺตํ อหิตาย ทุกฺขาย.}

\prose{ตตฺรานนฺท ยฺวายํ ปุคฺคโล อุทฺธโต โหติ, ตญฺจ เจโตวิมุตฺติํ ปญฺญาวิมุตฺติํ ยถาภูตํ ปชานาติ, ยตฺถสฺส ตํ อุทฺธจฺจํ อปริเสสํ นิรุชฺฌติ.\csromanspacer{} ตสฺส สวเนนปิ กตํ โหติ, พาหุสจฺเจนปิ กตํ โหติ, ทิฏฺฐิยาปิ ปฏิวิทฺธํ โหติ, สามายิกมฺปิ วิมุตฺติํ ลภติ.\csromanspacer{} อยํ อานนฺท ปุคฺคโล อมุนา ปุริเมน ปุคฺคเลน อภิกฺกนฺตตโร จ ปณีตตโร จ.\csromanspacer{} ตํ กิสฺส เหตุ? อิมํ หานนฺท ปุคฺคลํ ธมฺมโสโต นิพฺพหติ, ตทนฺตรํ โก ชาเนยฺย อญฺญตฺร ตถาคเตน.\csromanspacer{} ตสฺมาติหานนฺท มา ปุคฺคเลสุ ปมาณิกา อหุวตฺถ, มา ปุคฺคเลสุ ปมาณํ คณฺหิตฺถ, ขญฺญติ หานนฺท ปุคฺคเลสุ ปมาณํ คณฺหนฺโต.\csromanspacer{} อหํ วา อานนฺท ปุคฺคเลสุ ปมาณํ คเณฺหยฺยํ, โย วา ปนสฺส มาทิโส.}

\prose{กา จานนฺท มิคสาลา อุปาสิกา พาลา อพฺยตฺตา อมฺมกา อมฺมกปญฺญา, เก จ ปุริสปุคฺคลปโรปริเย ญาเณ.\csromanspacer{} อิเม โข อานนฺท ทส ปุคฺคลา สนฺโต สํวิชฺชมานา โลกสฺมิํ.}

\prose{ยถารูเปน อานนฺท สีเลน ปุราโณ สมนฺนาคโต อโหสิ, ตถารูเปน สีเลน อิสิทตฺโต สมนฺนาคโต อภวิสฺส, นยิธ ปุราโณ อิสิทตฺตสฺส คติมฺปิ อญฺญสฺส.\csromanspacer{} ยถารูปาย จ อานนฺท ปญฺญาย อิสิทตฺโต สมนฺนาคโต อโหสิ, ตถารูปาย ปญฺญาย ปุราโณ สมนฺนาคโต อภวิสฺส, นยิธ อิสิทตฺโต ปุราณสฺส คติมฺปิ อญฺญสฺส.\csromanspacer{} อิติ โข อานนฺท อิเม ปุคฺคลา อุโภ เอกงฺคหีนาติ.\csromanspacer{}\csromanspacer{}ปญฺจมํ.}

\csromanheader{2}{6. ตโยธมฺมสุตฺต}
\proseitem{76}{ตโยเม ภิกฺขเว ธมฺมา โลเก น สํวิชฺเชยฺยุํ, น ตถาคโต โลเก อุปฺปชฺเชยฺย อรหํ สมฺมาสมฺพุทฺโธ, น ตถาคตปฺปเวทิโต ธมฺมวินโย โลเก ทิพฺเพยฺย.\csromanspacer{} กตเม ตโย? ชาติ จ ชรา จ มรณญฺจ.\csromanspacer{} อิเม โข ภิกฺขเว ตโย ธมฺมา โลเก น สํวิชฺเชยฺยุํ, น ตถาคโต โลเก อุปฺปชฺเชยฺย อรหํ สมฺมาสมฺพุทฺโธ, น ตถาคตปฺปเวทิโต ธมฺมวินโย โลเก ทิพฺเพยฺย.\csromanspacer{} ยสฺมา จ โข ภิกฺขเว อิเม ตโย ธมฺมา โลเก สํวิชฺชนฺติ, ตสฺมา ตถาคโต โลเก อุปฺปชฺชติ อรหํ สมฺมาสมฺพุทฺโธ, ตสฺมา ตถาคตปฺปเวทิโต ธมฺมวินโย โลเก ทิพฺพติ.}

\chapterheadpageread{(8) 3. อากงฺขวคฺค}
\prose{\csromanfolio{371}ตโยเม ภิกฺขเว ธมฺเม อปฺปหาย อภพฺโพ ชาติํ ปหาตุํ ชรํ ปหาตุํ มรณํ ปหาตุํ.\csromanspacer{} กตเม ตโย? ราคํ อปฺปหาย โทสํ อปฺปหาย โมหํ อปฺปหาย.\csromanspacer{} อิเม โข ภิกฺขเว ตโย ธมฺเม อปฺปหาย อภพฺโพ ชาติํ ปหาตุํ ชรํ ปหาตุํ มรณํ ปหาตุํ. (1)}

\prose{ตโยเม ภิกฺขเว ธมฺเม อปฺปหาย อภพฺโพ ราคํ ปหาตุํ โทสํ ปหาตุํ โมหํ ปหาตุํ.\csromanspacer{} กตเม ตโย? สกฺกายทิฏฺฐิํ อปฺปหาย, วิจิกิจฺฉํ อปฺปหาย, สีลพฺพตปรามาสํ อปฺปหาย.\csromanspacer{} อิเม โข ภิกฺขเว ตโย ธมฺเม อปฺปหาย อภพฺโพ ราคํ ปหาตุํ โทสํ ปหาตุํ โมหํ ปหาตุํ. (2)}

\prose{ตโยเม ภิกฺขเว ธมฺเม อปฺปหาย อภพฺโพ สกฺกายทิฏฺฐิํ ปหาตุํ วิจิกิจฺฉํ ปหาตุํ สีลพฺพตรามาสํ ปหาตุํ.\csromanspacer{} กตเม ตโย? อโยนิโสมนสิการํ อปฺปหาย กุมฺมคฺคเสวนํ อปฺปหาย เจตโส ลีนตฺตํ อปฺปหาย.\csromanspacer{} อิเม โข ภิกฺขเว ตโย ธมฺเม อปฺปหาย อภพฺโพ สกฺกายทิฏฺฐิํ ปหาตุํ วิจิกิจฺฉํ ปหาตุํ สีลพฺพตปรามาสํ ปหาตุํ. (3)}

\prose{ตโยเม ภิกฺขเว ธมฺเม อปฺปหาย อภพฺโพ อโยนิโสมนสิการํ ปหาตุํ กุมฺมคฺคเสวนํ ปหาตุํ เจตโส ลีนตฺตํ ปหาตุํ.\csromanspacer{} กตเม ตโย? มุฏฺฐสฺสจฺจํ อปฺปหาย อสมฺปชญฺญํ อปฺปหาย เจตโส วิกฺเขปํ อปฺปหาย.\csromanspacer{} อิเม โข ภิกฺเขเว ตโย ธมฺเม อปฺปหาย อภพฺโพ อโยนิโสมนสิการํ ปหาตุํ กุมฺมคฺคเสวนํ ปหาตุํ เจตโส ลีนตฺตํ ปหาตุํ. (4)}

\prose{ตโยเม ภิกฺขเว ธมฺเม อปฺปหาย อภพฺโพ มุฏฺฐสฺสจฺจํ ปหาตุํ อสมฺปชญฺญํ ปหาตุํ เจตโส วิกฺเขปํ ปหาตุํ.\csromanspacer{} กตเม ตโย? อริยานํ อทสฺสนกมฺยตํ อปฺปหาย อริยธมฺมสฺส\footnote{อริยธมฺมํ (สฺยา)} อโสตุกมฺยตํ อปฺปหาย อุปารมฺภจิตฺตตํ อปฺปหาย.\csromanspacer{} อิเม โข ภิกฺขเว ตโย ธมฺเม อปฺปหาย อภพฺโพ มุฏฺฐสฺสจฺจํ ปหาตุํ อสมฺปชญฺญํ ปหาตุํ เจตโส วิกฺเขปํ ปหาตุํ. (5)}

\prose{ตโยเม ภิกฺขเว ธมฺเม อปฺปหาย อภพฺโพ อริยานํ อทสฺสนกมฺยตํ ปหาตุํ อริยธมฺมสฺส อโสตุกมฺยตํ ปหาตุํ อุปารมฺภจิตฺตตํ ปหาตุํ.\csromanspacer{} กตเม ตโย? อุทฺธจฺจํ อปฺปหาย อสํวรํ อปฺปหาย ทุสฺสิลฺยํ}

\csromanmatikaanchor{mtk.113}
\csromantocmarkref{subsection}{6. ตโยธมฺมสุตฺต}{mtk.113}
\bookmark[dest={mtk.113},level=2]{6. ตโยธมฺมสุตฺต}
\vspace{\tipitakaparskip}
\csromancenter{ทสกนิปาตปาฬิ อปฺปหาย.\csromanspacer{} อิเม โข ภิกฺขเว ตโย ธมฺเม อปฺปหาย อภพฺโพ อริยานํ อทสฺสนกมฺยตํ ปหาตุํ อริยธมฺมสฺส อโสตุกมฺยตํ ปหาตุํ อุปารมฺภจิตฺตตํ ปหาตุํ. (6)}
\prose{\csromanfolio{372}ตโยเม ภิกฺขเว ธมฺเม อปฺปหาย อภพฺโพ อุทฺธจฺจํ ปหาตุํ อสํวรํ ปหาตุํ ทุสฺสิลฺยํ ปหาตุํ.\csromanspacer{} กตเม ตโย? อสฺสทฺธิยํ อปฺปหาย อวทญฺญุตํ อปฺปหาย โกสชฺชํ อปฺปหาย.\csromanspacer{} อิเม โข ภิกฺขเว ตโยธมฺเม อปฺปหาย อภพฺโพ อุทฺธจฺจํ ปหาตุํ อสํวรํ ปหาตุํ ทุสฺสิลฺยํ ปหาตุํ. (7)}

\prose{ตโยเม ภิกฺขเว ธมฺเม อปฺปหาย อภพฺโพ อสฺสทฺธิยํ ปหาตุํ อวทญฺญุตํ ปหาตุํ โกสชฺชํ ปหาตุํ.\csromanspacer{} กตเม ตโย? อนาทริยํ อปฺปหาย โทวจสฺสตํ อปฺปหาย ปาปมิตฺตตํ อปฺปหาย.\csromanspacer{} อิเม โข ภิกฺขเว ตโย ธมฺเม อปฺปหาย อภพฺโพ อสฺสทฺธิยํ ปหาตุํ อวทญฺญุตํ ปหาตุํ โกสชฺชํ ปหาตุํ. (8)}

\prose{ตโยเม ภิกฺขเว ธมฺเม อปฺปหาย อภพฺโพ อนาทริยํ ปหาตุํ โทวจสฺสตํ ปหาตุํ ปาปมิตฺตตํ ปหาตุํ.\csromanspacer{} กตเม ตโย? อหิริกํ อปฺปหาย อโนตฺตปฺปํ อปฺปหาย ปมาทํ อปฺปหาย.\csromanspacer{} อิเม โข ภิกฺขเว ตโย ธมฺเม อปฺปหาย อภพฺโพ อนาทริยํ ปหาตุํ โทวจสฺสตํ ปหาตุํ ปาปมิตฺตตํ ปหาตุํ. (9)}

\prose{อหิริโกยํ ภิกฺขเว อโนตฺตาปี ปมตฺโต โหติ, โส ปมตฺโต สมาโน อภพฺโพ อนาทริยํ ปหาตุํ โทวจสฺสตํ ปหาตุํ ปาปมิตฺตตํ ปหาตุํ.\csromanspacer{} โส ปาปมิตฺโต สมาโน อภพฺโพ อสฺสทฺธิยํ ปหาตุํ อวทญฺญุตํ ปหาตุํ โกสชฺชํ ปหาตุํ.\csromanspacer{} โส กุสีโต สมาโน อภพฺโพ อุทฺธจฺจํ ปหาตุํ อสํวรํ ปหาตุํ ทุสฺสิลฺยํ ปหาตุํ.\csromanspacer{} โส ทุสฺสีโล สมาโน อภพฺโพ อริยานํ อทสฺสนกมฺยตํ ปหาตุํ อริยธมฺมสฺส อโสตุกมฺยตํ ปหาตุํ อุปารมฺภจิตฺตตํ ปหาตุํ.\csromanspacer{} โส อุปารมฺภจิตฺโต สมาโน อภพฺโพ มุฏฺฐสฺสจฺจํ ปหาตุํ อสมฺปชญฺญํ ปหาตุํ เจตโส วิกฺเขปํ ปหาตุํ.\csromanspacer{} โส วิกฺขิตฺตจิตฺโต สมาโน อภพฺโพ อโยนิโสมนสิการํ ปหาตุํ กุมฺมคฺคเสวนํ ปหาตุํ เจตโส ลีนตฺตํ ปหาตุํ.\csromanspacer{} โส ลีนจิตฺโต สมาโน อภพฺโพ สกฺกายทิฏฺฐิํ ปหาตุํ วิจิกิจฺฉํ ปหาตุํ สีลพฺพตปรามาสํ ปหาตุํ.\csromanspacer{} โส วิจิกิจฺโฉ สมาโน อภพฺโพ ราคํ ปหาตุํ}

\vspace{\tipitakaparskip}
\csromancenter{(8) 3.\csromanspacer{} อากงฺขวคฺค โทสํ ปหาตุํ โมหํ ปหาตุํ.\csromanspacer{} โส ราคํ อปฺปหาย โทสํ อปฺปหาย โมหํ อปฺปหาย อภพฺโพ ชาติํ ปหาตุํ ชรํ ปหาตุํ มรณํ ปหาตุํ. (10)}
\prose{\csromanfolio{373}ตโยเม ภิกฺขเว ธมฺเม ปหาย ภพฺโพ ชาติํ ปหาตุํ ชรํ ปหาตุํ มรณํ ปหาตุํ.\csromanspacer{} กตเม ตโย? ราคํ ปหาย โทสํ ปหาย โมหํ ปหาย.\csromanspacer{} อิเม โข ภิกฺขเว ตโย ธมฺเม ปหาย ภพฺโพ ชาติํ ปหาตุํ ชรํ ปหาตุํ มรณํ ปหาตุํ. (1)}

\prose{ตโยเม ภิกฺขเว ธมฺเม ปหาย ภพฺโพ ราคํ ปหาตุํ โทสํ ปหาตุํ โมหํ ปหาตุํ.\csromanspacer{} กตเม ตโย? สกฺกายทิฏฺฐิํ ปหาย วิจิกิจฺฉํ ปหาย สีลพฺพตปรามาสํ ปหาย.\csromanspacer{} อิเม โข ภิกฺขเว ตโย ธมฺเม ปหาย ภพฺโพ ราคํ ปหาตุํ โทสํ ปหาตุํ โมหํ ปหาตุํ. (2)}

\prose{ตโยเม ภิกฺขเว ธมฺเม ปหาย ภพฺโพ สกฺกายทิฏฺฐิํ ปหาตุํ วิจิกิจฺฉํ ปหาตุํ สีลพฺพตปรามาสํ ปหาตุํ.\csromanspacer{} กตเม ตโย? อโยนิโสมนสิการํ ปหาย กุมฺมคฺคเสวนํ ปหาย เจตโส ลีนตฺตํ ปหาย.\csromanspacer{} อิเม โข ภิกฺขเว ตโย ธมฺเม ปหาย ภพฺโพ สกฺกายทิฏฺฐิํ ปหาตุํ วิจิกิจฺฉํ ปหาตุํ สีลพฺพตปรามาสํ ปหาตุํ. (3)}

\prose{ตโยเม ภิกฺขเว ธมฺเม ปหาย ภพฺโพ อโยนิโสมนสิการํ ปหาตุํ กุมฺมคฺคเสวนํ ปหาตุํ เจตโส ลีนตฺตํ ปหาตุํ.\csromanspacer{} กตเม ตโย? มุฏฺฐสจฺจํ ปหาย อสมฺปชญฺญํ ปหาย เจตโส วิกฺเขปํ ปหาย.\csromanspacer{} อิเม โข ภิกฺขเว ตโย ธมฺเม ปหาย ภพฺโพ อโยนิโสมนสิการํ ปหาตุํ กุมฺมคฺคเสวนํ ปหาตุํ เจตโส ลีนตฺตํ ปหาตุํ. (4)}

\prose{ตโยเม ภิกฺขเว ธมฺเม ปหาย ภพฺโพ มุฏฺฐสฺสจฺจํ ปหาตุํ อสมฺปชญฺญํ ปหาตุํ เจตโส วิกฺเขปํ ปหาตุํ.\csromanspacer{} กตเม ตโย? อริยานํ อทสฺสนกมฺยตํ ปหาย อริยธมฺมสฺส อโสตุกมฺยตํ ปหาย อุปารมฺภจิตฺตตํ ปหาย.\csromanspacer{} อิเม โข ภิกฺขเว ตโย ธมฺเม ปหาย ภพฺโพ มุฏฺฐสฺสจฺจํ ปหาตุํ อสมฺปชญฺญํ ปหาตุํ เจตโส วิกฺเขปํ ปหาตุํ. (5)}

\prose{ตโยเม ภิกฺขเว ธมฺเม ปหาย ภพฺโพ อริยานํ อทสฺสนกมฺยตํ ปหาตุํ อริยธมฺมสฺส อโสตุกมฺยตํ ปหาตุํ อุปารมฺภจิตฺตตํ ปหาตุํ.\csromanspacer{} กตเม ตโย? อุทฺธจฺจํ ปหาย อสํวรํ ปหาย ทุสฺสิลฺยํ ปหาย.}

\vspace{\tipitakaparskip}
\csromancenter{ทสกนิปาตปาฬิ อิเม โข ภิกฺขเว ตโย ธมฺเม ปหาย ภพฺโพ อริยานํ อทสฺสนกมฺยตํ ปหาตุํ อริยธมฺมสฺส อโสตุกมฺยตํ ปหาตุํ อุปารมฺภจิตฺตตํ ปหาตุํ. (6)}
\prose{\csromanfolio{374}ตโยเม ภิกฺขเว ธมฺเม ปหาย ภพฺโพ อุทฺธจฺจํ ปหาตุํ อสํวรํ ปหาตุํ ทุสฺสิลฺยํ ปหาตุํ.\csromanspacer{} กตเม ตโย? อสฺสทฺธิยํ ปหาย อวทญฺญุตํ ปหาย โกสชฺชํ ปหาย.\csromanspacer{} อิเม โข ภิกฺขเว ตโย ธมฺเม ปหาย ภพฺโพ อุทฺธจฺจํ ปหาตุํ อสํวรํ ปหาตุํ ทุสฺสิลฺยํ ปหาตุํ. (7)}

\prose{ตโยเม ภิกฺขเว ธมฺเม ปหาย ภพฺโพ อสฺสทฺธิยํ ปหาตุํ อวทญฺญุตํ ปหาตุํ โกสชฺชํ ปหาตุํ.\csromanspacer{} กตเม ตโย? อนาทริยํ ปหาย โทวจสฺสตํ ปหาย ปาปมิตฺตตํ ปหาย.\csromanspacer{} อิเม โข ภิกฺขเว ตโย ธมฺเม ปหาย ภพฺโพ อสฺสทฺธิยํ ปหาตุํ อวทญฺญุตํ ปหาตุํ โกสชฺชํ ปหาตุํ. (8)}

\prose{ตโยเม ภิกฺขเว ธมฺเม ปหาย ภพฺโพ อนาทริยํ ปหาตุํ โทวจสฺสตํ ปหาตุํ ปาปมิตฺตตํ ปหาตุํ.\csromanspacer{} กตเม ตโย? อหิริกํ ปหาย อโนตฺตปฺปํ ปหาย ปมาทํ ปหาย.\csromanspacer{} อิเม โข ภิกฺขเว ตโย ธมฺเม ปหาย ภพฺโพ อนาทริยํ ปหาตุํ โทวจสฺสตํ ปหาตุํ ปาปมิตฺตตํ ปหาตุํ. (9)}

\prose{หิรีมายํ ภิกฺขเว โอตฺตาปี อปฺปมตฺโต โหติ, โส อปฺปมตฺโต สมาโน ภพฺโพ อนาทริยํ ปหาตุํ โทวจสฺสตํ ปหาตุํ ปาปมิตฺตตํ ปหาตุํ.\csromanspacer{} โส กลฺยาณมิตฺโต สมาโน ภพฺโพ อสฺสทฺธิยํ ปหาตุํ อวทญฺญุตํ ปหาตุํ โกสชฺชํ ปหาตุํ.\csromanspacer{} โส อารทฺธวีริโย สมาโน ภพฺโพ อุทฺธจฺจํ ปหาตุํ อสํวรํ ปหาตุํ ทุสฺสิลฺยํ ปหาตุํ.\csromanspacer{} โส สีลวา สมาโน ภพฺโพ อริยานํ อทสฺสนกมฺยตํ ปหาตุํ อริยธมฺมสฺส อโสตุกมฺยตํ ปหาตุํ อุปารมฺภจิตฺตตํ ปหาตุํ.\csromanspacer{} โส อนุปารมฺภจิตฺโต สมาโน ภพฺโพ มุฏฺฐสฺสจฺจํ ปหาตุํ อสมฺปชญฺญํ ปหาตุํ เจตโส วิกฺเขปํ ปหาตุํ.\csromanspacer{} โส อวิกฺขิตฺตจิตฺโต สมาโน ภพฺโพ อโยนิโสมนสิการํ ปหาตุํ กุมฺมคฺคเสวนํ ปหาตุํ เจตโส ลีนตฺตํ ปหาตุํ.\csromanspacer{} โส อลีนจิตฺโต สมาโน ภพฺโพ สกฺกายทิฏฺฐิํ ปหาตุํ วิจิกิจฺฉํ ปหาตุํ สีลพฺพตปรามาสํ ปหาตุํ.\csromanspacer{} โส อวิจิกิจฺโฉ สมาโน ภพฺโพ ราคํ ปหาตุํ โทสํ ปหาตุํ โมหํ ปหาตุํ.\csromanspacer{} โส ราคํ}

\vspace{\tipitakaparskip}
\csromancenter{(8) 3.\csromanspacer{} อากงฺขวคฺค ปหาย โทสํ ปหาย โมหํ ปหาย ภพฺโพ ชาติํ ปหาตุํ ชรํ ปหาตุํ มรณํ ปหาตุนฺติ. (10).\csromanspacer{}\csromanspacer{}ฉฏฺฐํ.}
\csromanheader{2}{7. กากสุตฺต}
\proseitem{77}{\csromanfolio{375}ทสหิ ภิกฺขเว อสทฺธมฺเมหิ สมนฺนาคโต กาโก.\csromanspacer{} กตเมหิ ทสหิ? ธํสี จ ปคพฺโภ จ ตินฺติโณ\footnote{นิลฺลชฺโช (ก) ตินฺติโณติ ตินฺติณํ วุจฺจติ ตณฺหา ... (สี, สฺยา-ฏฺฐ) อภิธมฺเม ขุทฺทกวตฺถุวิภงฺเค (364) ปิฏฺเฐ ตินฺติณปทนิทฺเทเส ปสฺสิตพฺพํ.} จ มหคฺฆโส จ ลุทฺโท จ อการุณิโก จ ทุพฺพโล จ โอรวิตา จ มุฏฺฐสฺสติ จ เนจยิโก\footnote{เนรสิโก (สี) ตทฏฺฐกถายํ ปน “เนจยิโก”เตฺวว ทิสฺสติ.} จ.\csromanspacer{} อิเมหิ โข ภิกฺขเว ทสหิ อสทฺธมฺเมหิ สมนฺนาคโต กาโก.\csromanspacer{} เอวเมวํ โข ภิกฺขเว ทสหิ อสทฺธมฺเมหิ สมนฺนาคโต ปาปภิกฺขุ.\csromanspacer{} กตเมหิ ทสหิ? ธํสี จ ปคพฺโภ จ ตินฺติโณ จ มหคฺฆโส จ ลุทฺโท จ อการุณิโก จ ทุพฺพโล จ โอรวิตา จ มุฏฺฐสฺสติ จ เนจยิโก จ.\csromanspacer{} อิเมหิ โข ภิกฺขเว ทสหิ อสทฺธมฺเมหิ สมนฺนาคโต ปาปภิกฺขูติ.\csromanspacer{}\csromanspacer{}สตฺตมํ.}

\csromanheader{2}{8. นิคณฺฐสุตฺต}
\proseitem{78}{ทสหิ ภิกฺขเว อสทฺธมฺเมหิ สมนฺนาคตา นิคณฺฐา.\csromanspacer{} กตเมหิ ทสหิ? อสฺสทฺธา ภิกฺขเว นิคณฺฐา, ทุสฺสีลา ภิกฺขเว นิคณฺฐา, อหิริกา ภิกฺขเว นิคณฺฐา, อโนตฺตปฺปิโน ภิกฺขเว นิคณฺฐา, อสปฺปุริสสมฺภตฺติโน ภิกฺขเว นิคณฺฐา, อตฺตุกฺกํสกปรวมฺภกา ภิกฺขเว นิคณฺฐา, สนฺทิฏฺฐิปรามาสา อาธานคฺคาหี ทุปฺปฏินิสฺสคฺคิโน ภิกฺขเว นิคณฺฐา, กุหกา ภิกฺขเว นิคณฺฐา, ปาปิจฺฉา ภิกฺขเว นิคณฺฐา, ปาปมิตฺตา ภิกฺขเว นิคณฺฐา.\csromanspacer{} อิเมหิ โข ภิกฺขเว ทสหิ อสทฺธมฺเมหิ สมนฺนาคตา นิคณฺฐาติ.\csromanspacer{}\csromanspacer{}อฏฺฐมํ.}

\csromanheader{2}{9. อาฆาตวตฺถุสุตฺต}
\proseitem{79}{\csromansymbolfootnote{*}{เหฏฺฐา 208 ปิฏฺเฐปิ.}ทสยิมานิ ภิกฺขเว อาฆาตวตฺถูนิ.\csromanspacer{} กตมานิ ทส? “อนตฺถํ เม อจรี”ติ อาฆาตํ พนฺธติ, “อนตฺถํ เม จรตี”ติ อาฆาตํ พนฺธติ, “อนตฺถํ เม จริสฺสตี”ติ อาฆาตํ พนฺธติ, “ปิยสฺส เม มนาปสฺส อนตฺถํ อจรีติ ฯเปฯ อนตฺถํ จรตีติ ฯเปฯ อนตฺถํ จริสฺสตี”ติ อาฆาตํ พนฺธติ, “อปฺปิยสฺส เม อมนาปสฺส อตฺถํ อจรีติ ฯเปฯ อตฺถํ จรตีติ ฯเปฯ อตฺถํ}

\vspace{\tipitakaparskip}
\csromancenter{ทสกนิปาตปาฬิ จริสฺสตี”ติ อาฆาตํ พนฺธติ, อฏฺฐาเน จ กุปฺปติ.\csromanspacer{} อิมานิ โข ภิกฺขเว ทส อาฆาตวตฺถูนีติ.\csromanspacer{}\csromanspacer{}นวมํ.}
\csromanheader{2}{10. อาฆาตปฏิวินยสุตฺต}
\csromanmatikaanchor{mtk.114}
\csromanmatikaanchor{mtk.118}
\csromantocmarkref{subsection}{7. กากสุตฺต}{mtk.114}
\bookmark[dest={mtk.114},level=2]{7. กากสุตฺต}
\proseitem{80}{\csromanfolio{376}ทสยิเม ภิกฺขเว อาฆาตปฏิวินยา.\csromanspacer{} กตเม ทส? “อนตฺถํ เม อจริ, ตํ กุเตตฺถ ลพฺภา”ติ อาฆาตํ ปฏิวิเนติ, “อนตฺถํ เม จรติ, ตํ กุเตตฺถ ลพฺภา”ติ อาฆาตํ ปฏิวิเนติ, “อนตฺถํ เม จริสฺสติ, ตํ กุเตตฺถ ลพฺภา”ติ อาฆาตํ ปฏิวิเนติ, “ปิยสฺส เม มนาปสฺส อนตฺถํ อจริ ฯเปฯ จรติ ฯเปฯ จริสฺสติ, ตํ กุเตตฺถ ลพฺภา”ติ อาฆาตํ ปฏิวิเนติ, “อปฺปิยสฺส เม อมนาปสฺส อตฺถํ อจริ ฯเปฯ อตฺถํ จรติ ฯเปฯ อตฺถํ จริสฺสติ, ตํ กุเตตฺถ ลพฺภา”ติ อาฆาตํ ปฏิวิเนติ, อฏฺฐาเน จ น กุปฺปติ.\csromanspacer{} อิเม โข ภิกฺขเว ทส อาฆาตปฏิวินยาติ.\csromanspacer{}\csromanspacer{}ทสมํ.\csromanspacer{} อากงฺขวคฺโค ตติโย.}

\titlehead{\textbf{ตสฺสุทฺทานํ}}
\csromangathameasure{อากงฺโข กณฺฏโก อิฏฺฐา, วฑฺฒิ จ มิคสาลาย. \\ ตโยธมฺมา จ กาโก จ, นิคณฺฐา เทฺว จ อาฆาตาติ.}
\csromangathagroup{\csromangathastanza{อากงฺโข กณฺฏโก อิฏฺฐา, วฑฺฒิ จ มิคสาลาย. \\ ตโยธมฺมา จ กาโก จ, นิคณฺฐา เทฺว จ อาฆาตาติ.}}

\csromanmatikaanchor{mtk.115}
\csromantocmarkref{subsection}{8. นิคณฺฐสุตฺต}{mtk.115}
\bookmark[dest={mtk.115},level=2]{8. นิคณฺฐสุตฺต}
\csromanheader{1}{\textbf{(9) 4. เถรวคฺค}}
\csromanheader{2}{1. วาหนสุตฺต}
\csromanmatikaanchor{mtk.116}
\csromantocmarkref{subsection}{9. อาฆาตวตฺถุสุตฺต}{mtk.116}
\bookmark[dest={mtk.116},level=2]{9. อาฆาตวตฺถุสุตฺต}
\proseitem{81}{เอกํ สมยํ ภควา จมฺปายํ วิหรติ คคฺคราย โปกฺขรณิยา ตีเร.\csromanspacer{} อถ โข อายสฺมา วาหโน เยน ภควา เตนุปสงฺกมิ, อุปสงฺกมิตฺวา ภควนฺตํ อภิวาเทตฺวา เอกมนฺตํ นิสีทิ, เอกมนฺตํ นิสินฺโน โข อายสฺมา วาหโน ภควนฺตํ เอตทโวจ “กติหิ นุ โข ภนฺเต ธมฺเมหิ ตถาคโต นิสฺสโฏ วิสํยุตฺโต วิปฺปมุตฺโต วิมริยาทีกเตน เจตสา วิหรตี”ติ.}

\prose{ทสหิ โข วาหน ธมฺเมหิ ตถาคโต นิสฺสโฏ วิสํยุตฺโต วิปฺปมุตฺโต วิมริยาทีกเตน เจตสา วิหรติ.\csromanspacer{} กตเมหิ ทสหิ? รูเปน โข วาหน ตถาคโต นิสฺสโฏ วิสํยุตฺโต วิปฺปมุตฺโต วิมริยาทีกเตน}

\vspace{\tipitakaparskip}
\csromancenter{(9) 4.\csromanspacer{} เถรวคฺค เจตสา วิหรติ, เวทนาย โข วาหน ฯเปฯ สญฺญาย โข วาหน.\csromanspacer{} สงฺขาเรหิ โข วาหน.\csromanspacer{} วิญฺญาเณน โข วาหน.\csromanspacer{} ชาติยา โข วาหน.\csromanspacer{} ชราย โข วาหน.\csromanspacer{} มรเณน โข วาหน.\csromanspacer{} ทุกฺเขหิ โข วาหน.\csromanspacer{} กิเลเสหิ โข วาหน ตถาคโต นิสฺสโฏ วิสํยุตฺโต วิปฺปมุตฺโต วิมริยาทีกเตน เจตสา วิหรติ.\csromanspacer{} เสยฺยถาปิ วาหน อุปฺปลํ วา ปทุมํ วา ปุณฺฑรีกํ วา อุทเก ชาตํ อุทเก สํวฑฺฒํ อุทกา ปจฺจุคฺคมฺม ฐิตํ อนุปลิตฺตํ อุทเกน.\csromanspacer{} เอวเมวํ โข วาหน อิเมหิ ทสหิ ธมฺเมหิ ตถาคโต นิสฺสโฏ วิสํยุตฺโต วิปฺปมุตฺโต วิมริยาทีกเตน เจตสา วิหรตีติ.\csromanspacer{}\csromanspacer{}ปฐมํ.}
\csromanmatikaanchor{mtk.117}
\csromantocmarkref{subsection}{10. อาฆาตปฏิวินยสุตฺต}{mtk.117}
\bookmark[dest={mtk.117},level=2]{10. อาฆาตปฏิวินยสุตฺต}
\csromantocmarkref{section}{อุทฺทานคาถา}{mtk.118}
\bookmark[dest={mtk.118},level=1]{อุทฺทานคาถา}
\csromanheader{2}{2. อานนฺทสุตฺต}
\proseitem{82}{\csromanfolio{377}อถ โข อายสฺมา อานนฺโท เยน ภควา เตนุปสงฺกมิ, อุปสงฺกมิตฺวา ภควนฺตํ อภิวาเทตฺวา เอกมนฺตํ นิสีทิ, เอกมนฺตํ นิสินฺนํ โข อายสฺมนฺตํ อานนฺทํ ภควา เอตทโวจ\csromandash{}}

\prose{โส วตานนฺท ภิกฺขุ อสฺสทฺโธ สมาโน อิมสฺมิํ ธมฺมวินเย วุทฺธิํ วิรูฬฺหิํ เวปุลฺลํ อาปชฺชิสฺสตีติ เนตํ ฐานํ วิชฺชติ. (1)}

\prose{โส วตานนฺท ภิกฺขุ ทุสฺสีโล สมาโน อิมสฺมิํ ธมฺมวินเย วุทฺธิํ วิรูฬฺหิํ เวปุลฺลํ อาปชฺชิสฺสตีติ เนตํ ฐานํ วิชฺชติ. (2)}

\csromanmatikaanchor{mtk.119}
\csromantocmarknopage{section}{(9) 4. เถรวคฺค}{mtk.119}
\bookmark[dest={mtk.119},level=1]{(9) 4. เถรวคฺค}
\prose{โส วตานนฺท ภิกฺขุ อปฺปสฺสุโต สมาโน อิมสฺมิํ ธมฺมวินเย วุทฺธิํ วิรูฬฺหิํ เวปุลฺลํ อาปชฺชิสฺสตีติ เนตํ ฐานํ วิชฺชติ. (3)}

\csromanmatikaanchor{mtk.120}
\csromantocmarkref{subsection}{1. วาหนสุตฺต}{mtk.120}
\bookmark[dest={mtk.120},level=2]{1. วาหนสุตฺต}
\prose{โส วตานนฺท ภิกฺขุ ทุพฺพโจ สมาโน อิมสฺมิํ ธมฺมวินเย วุทฺธิํ วิรูฬฺหิํ เวปุลฺลํ อาปชฺชิสฺสตีติ เนตํ ฐานํ วิชฺชติ. (4)}

\prose{โส วตานนฺท ภิกฺขุ ปาปมิตฺโต สมาโน อิมสฺมิํ ธมฺมวินเย วุทฺธิํ วิรูฬฺหิํ เวปุลฺลํ อาปชฺชิสฺสตีติ เนตํ ฐานํ วิชฺชติ. (5)}

\prose{โส วตานนฺท ภิกฺขุ กุสีโต สมาโน อิมสฺมิํ ธมฺมวินเย วุทฺธิํ วิรูฬฺหิํ เวปุลฺลํ อาปชฺชิสฺสตีติ เนตํ ฐานํ วิชฺชติ. (6)}

\prose{โส วตานนฺท ภิกฺขุ มุฏฺฐสฺสติ สมาโน อิมสฺมิํ ธมฺมวินเย วุทฺธิํ วิรูฬฺหิํ เวปุลฺลํ อาปชฺชิสฺสตีติ เนตํ ฐานํ วิชฺชติ. (7)}

\csromanmatikaanchor{mtk.121}
\csromantocmarkref{subsection}{2. อานนฺทสุตฺต}{mtk.121}
\bookmark[dest={mtk.121},level=2]{2. อานนฺทสุตฺต}
\gambhira{ทสกนิปาตปาฬิ}
\prose{\csromanfolio{378}โส วตานนฺท ภิกฺขุ อสนฺตุฏฺโฐ สมาโน อิมสฺมิํ ธมฺมวินเย วุทฺธิํ วิรูฬฺหิํ เวปุลฺลํ อาปชฺชิสฺสตีติ เนตํ ฐานํ วิชฺชติ. (8)}

\prose{โส วตานนฺท ภิกฺขุ ปาปิจฺโฉ สมาโน อิมสฺมิํ ธมฺมวินเย วุทฺธิํ วิรูฬฺหิํ เวปุลฺลํ อาปชฺชิสฺสตีติ เนตํ ฐานํ วิชฺชติ. (9)}

\prose{โส วตานนฺท ภิกฺขุ มิจฺฉาทิฏฺฐิโก สมาโน อิมสฺมิํ ธมฺมวินเย วุทฺธิํ วิรูฬฺหิํ เวปุลฺลํ อาปชฺชิสฺสตีติ เนตํ ฐานํ วิชฺชติ. (10)}

\prose{โส วตานนฺท ภิกฺขุ อิเมหิ ทสหิ ธมฺเมหิ สมนฺนาคโต อิมสฺมิํ ธมฺมวินเย วุทฺธิํ วิรูฬฺหิํ เวปุลฺลํ อาปชฺชิสฺสตีติ เนตํ ฐานํ วิชฺชติ.}

\prose{โส วตานนฺท ภิกฺขุ สทฺโธ สมาโน อิมสฺมิํ ธมฺมวินเย วุทฺธิํ วิรูฬฺหิํ เวปุลฺลํ อาปชฺชิสฺสตีติ ฐานเมตํ วิชฺชติ. (1)}

\prose{โส วตานนฺท ภิกฺขุ สีลวา สมาโน อิมสฺมิํ ธมฺมวินเย วุทฺธิํ วิรูฬฺหิํ เวปุลฺลํ อาปชฺชิสฺสตีติ ฐานเมตํ วิชฺชติ. (2)}

\prose{โส วตานนฺท ภิกฺขุ พหุสฺสุโต สุตธโร สมาโน อิมสฺมิํ ธมฺมวินเย วุทฺธิํ วิรูฬฺหิํ เวปุลฺลํ อาปชฺชิสฺสตีติ ฐานเมตํ วิชฺชติ. (3)}

\prose{โส วตานนฺท ภิกฺขุ สุวโจ สมาโน อิมสฺมิํ ธมฺมวินเย วุทฺธิํ วิรูฬฺหิํ เวปุลฺลํ อาปชฺชิสฺสตีติ ฐานเมตํ วิชฺชติ. (4)}

\prose{โส วตานนฺท ภิกฺขุ กลฺยาณมิตฺโต สมาโน อิมสฺมิํ ธมฺมวินเย วุทฺธิํ วิรูฬฺหิํ เวปุลฺลํ อาปชฺชิสฺสตีติ ฐานเมตํ วิชฺชติ. (5)}

\prose{โส วตานนฺท ภิกฺขุ อารทฺธวีริโย สมาโน อิมสฺมิํ ธมฺมวินเย วุทฺธิํ วิรูฬฺหิํ เวปุลฺลํ อาปชฺชิสฺสตีติ ฐานเมตํ วิชฺชติ. (6)}

\prose{โส วตานนฺท ภิกฺขุ อุปฏฺฐิตสฺสติ สมาโน อิมสฺมิํ ธมฺมวินเย วุทฺธิํ วิรูฬฺหิํ เวปุลฺลํ อาปชฺชิสฺสตีติ ฐานเมตํ วิชฺชติ. (7)}

\prose{โส วตานนฺท ภิกฺขุ สนฺตุฏฺโฐ สมาโน อิมสฺมิํ ธมฺมวินเย วุทฺธิํ วิรูฬฺหิํ เวปุลฺลํ อาปชฺชิสฺสตีติ ฐานเมตํ วิชฺชติ. (8)}

\prose{โส วตานนฺท ภิกฺขุ อปฺปิจฺโฉ สมาโน อิมสฺมิํ ธมฺมวินเย วุทฺธิํ วิรูฬฺหิํ เวปุลฺลํ อาปชฺชิสฺสตีติ ฐานเมตํ วิชฺชติ. (9)}

\chapterheadpageread{(9) 4. เถรวคฺค}
\prose{\csromanfolio{379}โส วตานนฺท ภิกฺขุ สมฺมาทิฏฺฐิโก สมาโน อิมสฺมิํ ธมฺมวินเย วุทฺธิํ วิรูฬฺหิํ เวปุลฺลํ อาปชฺชิสฺสตีติ ฐานเมตํ วิชฺชติ. (10)}

\prose{โส วตานนฺท ภิกฺขุ อิเมหิ ทสหิ ธมฺเมหิ สมนฺนาคโต อิมสฺมิํ ธมฺมวินเย วุทฺธิํ วิรูฬฺหิํ เวปุลฺลํ อาปชฺชิสฺสตีติ ฐานเมตํ วิชฺชตีติ.\csromanspacer{}\csromanspacer{}ทุติยํ.}

\csromanheader{2}{3. ปุณฺณิยสุตฺต}
\proseitem{83}{อถ โข อายสฺมา ปุณฺณิโย เยน ภควา เตนุปสงฺกมิ, อุปสงฺกมิตฺวา ภควนฺตํ อภิวาเทตฺวา เอกมนฺตํ นิสีทิ, เอกมนฺตํ นิสินฺโน โข อายสฺมา ปุณฺณิโย ภควนฺตํ เอตทโวจ “โก นุ โข ภนฺเต เหตุ โก ปจฺจโย, เยน อปฺเปกทา ตถาคตํ ธมฺมเทสนา ปฏิภาติ, อปฺเปกทา นปฺปฏิภาตี”ติ.}

\prose{สทฺโธ จ ปุณฺณิย ภิกฺขุ โหติ, โน จ อุปสงฺกมิตา, เนว ตาว ตถาคตํ ธมฺมเทสนา ปฏิภาติ.\csromanspacer{} ยโต จ โข ปุณฺณิย ภิกฺขุ สทฺโธ จ โหติ อุปสงฺกมิตา จ, เอวํ ตถาคตํ ธมฺมเทสนา ปฏิภาติ.}

\prose{สทฺโธ จ ปุณฺณิย ภิกฺขุ โหติ อุปสงฺกมิตา จ, โน จ ปยิรุปาสิตา ฯเปฯ ปยิรุปาสิตา จ, โน จ ปริปุจฺฉิตา.\csromanspacer{} ปริปุจฺฉิตา จ, โน จ โอหิตโสโต ธมฺมํ สุณาติ.\csromanspacer{} โอหิตโสโต จ ธมฺมํ สุณาติ, โน จ สุตฺวา ธมฺมํ ธาเรติ.\csromanspacer{} สุตฺวา จ ธมฺมํ ธาเรติ, โน จ ธาตานํ ธมฺมานํ อตฺถํ อุปปริกฺขติ.\csromanspacer{} ธาตานญฺจ ธมฺมานํ อตฺถํ อุปปริกฺขติ, โน จ อตฺถมญฺญาย ธมฺมมญฺญาย ธมฺมานุธมฺมปฺปฏิปนฺโน โหติ.\csromanspacer{} อตฺถมญฺญาย ธมฺมมญฺญาย ธมฺมานุธมฺมปฺปฏิปนฺโน จ โหติ, โน จ กลฺยาณวาโจ โหติ กลฺยาณวากฺกรโณ โปริยา วาจาย สมนฺนาคโต วิสฺสฏฺฐาย อเนลคฬาย อตฺถสฺส วิญฺญาปนิยา.\csromanspacer{} กลฺยาณวาโจ จ โหติ กลฺยาณวากฺกรโณ โปริยา วาจาย สมนฺนาคโต วิสฺสฏฺฐาย อเนลคฬาย อตฺถสฺส วิญฺญาปนิยา, โน จ สนฺทสฺสโก โหติ สมาทปโก สมุตฺเตชโก สมฺปหํสโก สพฺรหฺมจารีนํ.\csromanspacer{} เนว ตาว ตถาคตํ ธมฺมเทสนา ปฏิภาติ.}

\prose{ยโต จ โข ปุณฺณิย ภิกฺขุ สทฺโธ จ โหติ อุปสงฺกมิตา จ ปยิรุปาสิตา จ ปริปุจฺฉิตา จ, โอหิตโสโต จ ธมฺมํ สุณาติ, สุตฺวา จ ธมฺมํ ธาเรติ, ธาตานญฺจ ธมฺมานํ อตฺถํ อุปปริกฺขติ, อตฺถมญฺญาย ธมฺมมญฺญาย}

\vspace{\tipitakaparskip}
\csromancenter{ทสกนิปาตปาฬิ ธมฺมานุธมฺมปฺปฏิปนฺโน จ โหติ, กลฺยาณวาโจ จ โหติ กลฺยาณวากฺกรโณ โปริยา วาจาย สมนฺนาคโต วิสฺสฏฺฐาย อเนลคฬาย อตฺถสฺส วิญฺญาปนิยา, สนฺทสฺสโก จ โหติ สมาทปโก สมุตฺเตชโก สมฺปหํสโก สพฺรหฺมจารีนํ.\csromanspacer{} เอวํ ตถาคตํ ธมฺมเทสนา ปฏิภาติ.\csromanspacer{} อิเมหิ โข ปุณฺณิย ทสหิ ธมฺเมหิ สมนฺนาคตา\footnote{สมนฺนาคโต (ก)} 2เอกนฺตปฏิภานา\footnote{เอกนฺตปฏิภานํ (สี)} ตถาคตํ ธมฺมเทสนา โหตีติ\footnote{เอกนฺตํ ตถาคตํ ธมฺมเทสนา ปฏิภาตีติ (สฺยา)}.\csromanspacer{}\csromanspacer{}ตติยํ.}
\csromanheader{2}{4. พฺยากรณสุตฺต}
\proseitem{84}{\csromanfolio{380}ตตฺร โข อายสฺมา มหาโมคฺคลฺลาโน ภิกฺขู อามนฺเตสิ “อาวุโส ภิกฺขเว”ติ. “อาวุโส”ติ โข เต ภิกฺขู อายสฺมโต มหาโมคฺคลฺลานสฺส ปจฺจสฺโสสุํ.\csromanspacer{} อายสฺมา มหาโมคฺคลฺลาโน เอตทโวจ\csromandash{}}

\prose{อิธาวุโส ภิกฺขุ อญฺญํ พฺยากโรติ “ขีณา ชาติ, วุสิตํ พฺรหฺมจริยํ, กตํ กรณียํ, นาปรํ อิตฺถตฺตายาติ ปชานามี”ติ.\csromanspacer{} ตเมนํ ตถาคโต วา ตถาคตสาวโก วา ฌายี สมาปตฺติกุสโล ปรจิตฺตกุสโล ปรจิตฺตปริยายกุสโล สมนุยุญฺชติ สมนุคฺคาหติ สมนุภาสติ.\csromanspacer{} โส ตถาคเตน วา ตถาคตสาวเกน วา ฌายินา สมาปตฺติกุสเลน ปรจิตฺตกุสเลน ปรจิตฺตปริยายกุสเลน สมนุยุญฺชิยมาโน สมนุคฺคาหิยมาโน สมนุภาสิยมาโน อิรีณํ อาปชฺชติ, วิจินํ\footnote{วิสินํ (สี-ฏฺฐ)} อาปชฺชติ, อนยํ อาปชฺชติ, พฺยสนํ อาปชฺชติ, อนยพฺยสนํ อาปชฺชติ.}

\csromanmatikaanchor{mtk.122}
\csromantocmarkref{subsection}{3. ปุณฺณิยสุตฺต}{mtk.122}
\bookmark[dest={mtk.122},level=2]{3. ปุณฺณิยสุตฺต}
\prose{ตเมนํ ตถาคโต วา ตถาคตสาวโก วา ฌายี สมาปตฺติกุสโล ปรจิตฺตกุสโล ปรจิตฺตปริยายกุสโล เอวํ เจตสา เจโต ปริจฺจ มนสิ กโรติ “กิํ นุ โข อยมายสฺมา อญฺญํ พฺยากโรติ ‘ขีณา ชาติ, วุสิตํ พฺรหฺมจริยํ, กตํ กรณียํ, นาปรํ อิตฺถตฺตายาติ ปชานามี’ติ”.}

\prose{ตเมนํ ตถาคโต วา ตถาคตสาวโก วา ฌายี สมาปตฺติกุสโล ปรจิตฺตกุสโล ปรจิตฺตปริยายกุสโล เอวํ เจตสา เจโต ปริจฺจ ปชานาติ\csromandash{}}

\chapterheadpageread{(9) 4. เถรวคฺค}
\prose{\csromanfolio{381}โกธโน โข อยมายสฺมา, โกธปริยุฏฺฐิเตน เจตสา พหุลํ วิหรติ, โกธปริยุฏฺฐานํ โข ปน ตถาคตปฺปเวทิเต ธมฺมวินเย ปริหานเมตํ.}

\prose{อุปนาหี โข ปน อยมายสฺมา, อุปนาหปริยุฏฺฐิเตน เจตสา พหุลํ วิหรติ, อุปนาหปริยุฏฺฐานํ โข ปน ตถาคตปฺปเวทิเต ธมฺมวินเย ปริหานเมตํ.}

\prose{มกฺขี โข ปน อยมายสฺมา, มกฺขปริยุฏฺฐิเตน เจตสา พหุลํ วิหรติ, มกฺขปริยุฏฺฐานํ โข ปน ตถาคตปฺปเวทิเต ธมฺมวินเย ปริหานเมตํ.}

\csromanmatikaanchor{mtk.123}
\csromantocmarkref{subsection}{4. พฺยากรณสุตฺต}{mtk.123}
\bookmark[dest={mtk.123},level=2]{4. พฺยากรณสุตฺต}
\prose{ปฬาสี โข ปน อยมายสฺมา, ปฬาสปริยุฏฺฐิเตน เจตสา พหุลํ วิหรติ, ปฬาสปริยุฏฺฐานํ โข ปน ตถาคตปฺปเวทิเต ธมฺมวินเย ปริหานเมตํ.}

\prose{อิสฺสุกี โข ปน อยมายสฺมา, อิสฺสาปริยุฏฺฐิเตน เจตสา พหุลํ วิหรติ, อิสฺสาปริยุฏฺฐานํ โข ปน ตถาคตปฺปเวทิเต ธมฺมวินเย ปริหานเมตํ.}

\prose{มจฺฉรี โข ปน อยมายสฺมา, มจฺเฉรปริยุฏฺฐิเตน เจตสา พหุลํ วิหรติ, มจฺเฉรปริยุฏฺฐานํ โข ปน ตถาคตปฺปเวทิเต ธมฺมวินเย ปริหานเมตํ.}

\prose{สโฐ โข ปน อยมายสฺมา, สาเฐยฺยปริยุฏฺฐิเตน เจตสา พหุลํ วิหรติ, สาเฐยฺยปริยุฏฺฐานํ โข ปน ตถาคตปฺปเวทิเต ธมฺมวินเย ปริหานเมตํ.}

\prose{มายาวี โข ปน อยมายสฺมา, มายาปริยุฏฺฐิเตน เจตสา พหุลํ วิหรติ, มายาปริยุฏฺฐานํ โข ปน ตถาคตปฺปเวทิเต ธมฺมวินเย ปริหานเมตํ.}

\prose{ปาปิจฺโฉ โข ปน อยมายสฺมา, อิจฺฉาปริยุฏฺฐิเตน เจตสา พหุลํ วิหรติ, อิจฺฉาปริยุฏฺฐานํ โข ปน ตถาคตปฺปเวทิเต ธมฺมวินเย ปริหานเมตํ.}

\gambhira{ทสกนิปาตปาฬิ}
\prose{\csromanfolio{382}สติ\footnote{มุฏฺฐสฺสติ (สี, สฺยา)} โข ปน อยมายสฺมา อุตฺตริ กรณีเย โอรมตฺตเกน วิเสสาธิคเมน อนฺตรา โวสานํ อาปนฺโน, อนฺตรา โวสานคมนํ โข ปน ตถาคตปฺปเวทิเต ธมฺมวินเย ปริหานเมตํ.}

\prose{โส วตาวุโส ภิกฺขุ อิเม ทส ธมฺเม อปฺปหาย อิมสฺมิํ ธมฺมวินเย วุทฺธิํ วิรูฬฺหิํ เวปุลฺลํ อาปชฺชิสฺสตีติ เนตํ ฐานํ วิชฺชติ.\csromanspacer{} โส วตาวุโส ภิกฺขุ อิเม ทส ธมฺเม ปหาย อิมสฺมิํ ธมฺมวินเย วุทฺธิํ วิรูฬฺหิํ เวปุลฺลํ อาปชฺชิสฺสตีติ ฐานเมตํ วิชฺชตีติ.\csromanspacer{}\csromanspacer{}จตุตฺถํ.}

\csromanheader{2}{5. กตฺถีสุตฺต}
\proseitem{85}{เอกํ สมยํ อายสฺมา มหาจุนฺโท เจตีสุ วิหรติ สหชาติยํ.\csromanspacer{} ตตฺร โข อายสฺมา มหาจุนฺโท ภิกฺขู อามนฺเตสิ “อาวุโส ภิกฺขเว”ติ. “อาวุโส”ติ โข เต ภิกฺขู อายสฺมโต มหาจุนฺทสฺส ปจฺจโสสุํ.\csromanspacer{} อายสฺมา มหาจุนฺโท เอตทโวจ\csromandash{}}

\prose{อิธาวุโส ภิกฺขุ กตฺถี โหติ วิกตฺถี อธิคเมสุ “อหํ ปฐมํ ฌานํ สมาปชฺชามิปิ วุฏฺฐหามิปิ, อหํ ทุติยํ ฌานํ สมาปชฺชามิปิ วุฏฺฐหามิปิ.\csromanspacer{} อหํ ตติยํ ฌานํ สมาปชฺชามิปิ วุฏฺฐหามิปิ, อหํ จตุตฺถํ ฌานํ สมาปชฺชามิปิ วุฏฺฐหามิปิ, อหํ อากาสานญฺจายตนํ สมาปชฺชามิปิ วุฏฺฐหามิปิ, อหํ วิญฺญาณญฺจายตนํ สมาปชฺชามิปิ วุฏฺฐหามิปิ, อหํ อากิญฺจญฺญายตนํ สมาปชฺชามิปิ วุฏฺฐหามิปิ, อหํ เนวสญฺญานาสญฺญายตนํ สมาปชฺชามิปิ วุฏฺฐหามิปิ, อหํ สญฺญาเวทยิตนิโรธํ สมาปชฺชามิปิ วุฏฺฐหามิปี”ติ.}

\prose{ตเมนํ ตถาคโต วา ตถาคตสาวโก วา ฌายี สมาปตฺติกุสโล ปรจิตฺตกุสโล ปรจิตฺตปริยายกุสโล สมนุยุญฺชติ สมนุคฺคาหติ สมนุภาสติ.\csromanspacer{} โส ตถาคเตน วา ตถาคตสาวเกน วา ฌายินา สมาปตฺติกุสเลน ปรจิตฺตกุสเลน ปรจิตฺตปริยายกุสเลน สมนุยุญฺชิยมาโน สมนุคฺคาหิยมาโน สมนุภาสิยมาโน อิรีณํ อาปชฺชติ, วิจินํ อาปชฺชติ, อนยํ อาปชฺชติ, พฺยสนํ อาปชฺชติ, อนยพฺยสนํ อาปชฺชติ.}

\chapterheadpageread{(9) 4. เถรวคฺค}
\prose{\csromanfolio{383}ตเมนํ ตถาคโต วา ตถาคตสาวโก วา ฌายี สมาปตฺติกุสโล ปรจิตฺตกุสโล ปรจิตฺตปริยายกุสโล เอวํ เจตสา เจโต ปริจฺจ มนสิ กโรติ “กิํ นุ โข อยมายสฺมา กตฺถี โหติ วิกตฺถี อธิคเมสุ “อหํ ปฐมํ ฌานํ สมาปชฺชามิปิ วุฏฺฐหามิปิ ฯเปฯ อหํ สญฺญาเวทยิตนิโรธํ สมาปชฺชามิปิ วุฏฺฐหามิปี”ติ.}

\prose{ตเมนํ ตถาคโต วา ตถาคตสาวโก วา ฌายี สมาปตฺติกุสโล ปรจิตฺตกุสโล ปรจิตฺตปริยายกุสโล เอวํ เจตสา เจโต ปริจฺจ ปชานาติ\csromandash{}}

\prose{ทีฆรตฺตํ โข อยมายสฺมา ขณฺฑการี ฉิทฺทการี สพลการี กมฺมาสการี น สนฺตตการี น สนฺตตวุตฺติ สีเลสุ.\csromanspacer{} ทุสฺสีโล โข อยมายสฺมา, ทุสฺสิลฺยํ โข ปน ตถาคตปฺปเวทิเต ธมฺมวินเย ปริหานเมตํ.}

\prose{อสฺสทฺโธ โข ปน อยมายสฺมา, อสฺสทฺธิยํ โข ปน ตถาคตปฺปเวทิเต ธมฺมวินเย ปริหานเมตํ.}

\csromanmatikaanchor{mtk.124}
\csromantocmarkref{subsection}{5. กตฺถีสุตฺต}{mtk.124}
\bookmark[dest={mtk.124},level=2]{5. กตฺถีสุตฺต}
\prose{อปฺปสฺสุโต โข ปน อยมายสฺมา อนาจาโร, อปฺปสจฺจํ โข ปน ตถาคตปฺปเวทิเต ธมฺมวินเย ปริหานเมตํ.}

\prose{ทุพฺพโจ โข ปน อยมายสฺมา, โทวจสฺสตา โข ปน ตถาคตปฺปเวทิเต ธมฺมวินเย ปริหานเมตํ.}

\prose{ปาปมิตฺโต โข ปน อยมายสฺมา, ปาปมิตฺตตา โข ปน ตถาคตปฺปเวทิเต ธมฺมวินเย ปริหานเมตํ.}

\prose{กุสีโต โข ปน อยมายสฺมา, โกสชฺชํ โข ปน ตถาคตปฺปเวทิเต ธมฺมวินเย ปริหานเมตํ.}

\prose{มุฏฺฐสฺสติ โข ปน อยมายสฺมา, มุฏฺฐสฺสจฺจํ โข ปน ตถาคตปฺปเวทิเต ธมฺมวินเย ปริหานเมตํ.}

\prose{กุหโก โข ปน อยมายสฺมา, โกหญฺญํ โข ปน ตถาคตปฺปเวทิเต ธมฺมวินเย ปริหานเมตํ.}

\prose{ทุพฺภโร โข ปน อยมายสฺมา, ทุพฺภรตา โข ปน ตถาคตปฺปเวทิเต ธมฺมวินเย ปริหานเมตํ.}

\gambhira{ทสกนิปาตปาฬิ}
\prose{\csromanfolio{384}ทุปฺปญฺโญ โข ปน อยมายสฺมา, ทุปฺปญฺญตา โข ปน ตถาคตปฺปเวทิเต ธมฺมวินเย ปริหานเมตํ.}

\prose{เสยฺยถาปิ อาวุโส สหายโก สหายกํ เอวํ วเทยฺย “ยทา เต สมฺม ธเนน\footnote{พนฺโธ (ก)} ธนกรณียํ อสฺส, ยาเจยฺยาสิ มํ\footnote{ยาจิสฺสสิ มํ (สี), ปเวเทยฺยาสิ มํ (สฺยา), ปราเชยฺยาปิ มํ (ก)} ธนํ, ทสฺสามิ เต ธนนฺ”ติ.\csromanspacer{} โส กิญฺจิเทว ธนกรณีเย สมุปฺปนฺเน สหายโก สหายกํ เอวํ วเทยฺย “อตฺโถ เม สมฺมธเนน, เทหิ เม ธนนฺ”ติ.\csromanspacer{} โส เอวํ วเทยฺย “เตน หิ สมฺม อิธ ขนาหี”ติ.\csromanspacer{} โส ตตฺร ขนนฺโต นาธิคจฺเฉยฺย, โส เอวํ วเทยฺย “อลิกํ มํ สมฺม อวจ, ตุจฺฉกํ มํ สมฺม อวจ ‘อิธ ขนาหี’ติ”.\csromanspacer{} โส เอวํ วเทยฺย “นาหํ ตํ สมฺม อลิกํ อวจํ, ตุจฺฉกํ อวจํ, เตน หิ สมฺม อิธ ขนาหี”ติ.\csromanspacer{} โส ตตฺรปิ ขนนฺโต นาธิคจฺเฉยฺย, โส เอวํ วเทยฺย “อลิกํ มํ สมฺม อวจ, ตุจฺฉกํ มํ สมฺม อวจ ‘อิธ ขนาหี’ติ”.\csromanspacer{} โส เอวํ วเทยฺย “นาหํ ตํ สมฺม อลิกํ อวจํ, ตุจฺฉกํ อวจํ, เตน หิ สมฺม อิธ ขนาหี”ติ.\csromanspacer{} โส ตตฺรปิ ขนนฺโต นาธิคจฺเฉยฺย, โส เอวํ วเทยฺย “อลิกํ มํ สมฺม อวจ, ตุจฺฉกํ มํ สมฺม อวจ ‘อิธ ขนาหี’ติ.\csromanspacer{} โส เอวํ วเทยฺย “นาหํ ตํ สมฺม อลิกํ อวจํ, ตุจฺฉกํ อวจํ, อปิ จ อหเมว อุมฺมาทํ ปาปุณิํ เจตโส วิปริยายนฺ”ติ.}

\prose{เอวเมวํ โข อาวุโส ภิกฺขุ กตฺถี โหติ วิกตฺถี อธิคเมสุ “อหํ ปฐมํ ฌานํ สมาปชฺชามิปิ วุฏฺฐหามิปิ, อหํ ทุติยํ ฌานํ สมาปชฺชามิปิ วุฏฺฐหามิปิ, อหํ ตติยํ ฌานํ สมาปชฺชามิปิ วุฏฺฐหามิปิ, อหํ จตุตฺถํ ฌานํ สมาปชฺชามิปิ วุฏฺฐหามิปิ, อหํ อากาสานญฺชายตนํ สมาปชฺชามิปิ วุฏฺฐหามิปิ, อหํ วิญฺญาณญฺจายตนํ สมาปชฺชามิปิ วุฏฺฐหามิปิ, อหํ อากิญฺจญฺญายตนํ สมาปชฺชามิปิ วุฏฺฐหามิปิ, อหํ เนวสญฺญานาสญฺญายตนํ สมาปชฺชามิปิ วุฏฺฐหามิปิ, อหํ สญฺญาเวทยิตนิโรธํ สมาปชฺชามิปิ วุฏฺฐหามิปี”ติ.}

\prose{ตเมนํ ตถาคโต วา ตถาคตสาวโก วา ฌายี สมาปตฺติกุสโล ปรจิตฺตกุสโล ปรจิตฺตปริยายกุสโล สมนุยุญฺชติ สมนุคฺคาหติ สมนุภาสติ.\csromanspacer{} โส ตถาคเตน วา ตถาคตสาวเกน วา ฌายินา สมาปตฺติกุสเลน ปรจิตฺตกุสเลน ปรจิตฺตปริยายกุสเลน สมนุยุญฺชิยมาโน สมนุคฺคาหิยมาโน สมนุภาสิยมาโน}

\vspace{\tipitakaparskip}
\csromancenter{(9) 4.\csromanspacer{} เถรวคฺค อิรีณํ อาปชฺชติ, วิจินํ อาปชฺชติ, อนยํ อาปชฺชติ, พฺยสนํ อาปชฺชติ, อนยพฺยสนํ อาปชฺชติ.}
\prose{\csromanfolio{385}ตเมนํ ตถาคโต วา ตถาคตสาวโก วา ฌายี สมาปตฺติกุสโล ปรจิตฺตกุสโล ปรจิตฺตปริยายกุสโล เอวํ เจตสา เจโต ปริจฺจ มนสิ กโรติ “กิํ นุ โข อยมายสฺมา กตฺถี โหติ วิกตฺถี อธิคเมสุ “อหํ ปฐมํ ฌานํ สมาปชฺชามิปิ ฯเปฯ อหํ สญฺญาเวทยิตนิโรธํ สมาปชฺชามิปิ วุฏฺฐหามิปี”ติ.}

\prose{ตเมนํ ตถาคโต วา ตถาคตสาวโก วา ฌายี สมาปตฺติกุสโล ปรจิตฺตปริยายกุสโล เจตสา เจโต ปริจฺจ ปชานาติ.}

\prose{ทีฆรตฺตํ โข อยมายสฺมา ขณฺฑการี ฉิทฺทการี สพลการี กมฺมาสการี น สนฺตตการี น สนฺตตวุตฺติ สีเลสุ.\csromanspacer{} ทุสฺสีโล โข อยมายสฺมา, ทุสฺสิลฺยํ โข ปน ตถาคตปฺปเวทิเต ธมฺมวินเย ปริหานเมตํ.}

\prose{อสฺสทฺโธ โข ปน อยมายสฺมา, อสฺสทฺธิยํ โข ปน ตถาคตปฺปเวทิเต ธมฺมวินเย ปริหานเมตํ.}

\prose{อปฺปสฺสุโต โข ปน อยมายสฺมา อนาจาโร, อปฺปสจฺจํ โข ปน ตถาคตปฺปเวทิเต ธมฺมวินเย ปริหานเมตํ.}

\prose{ทุพฺพโจ โข ปน อยมายสฺมา, โทวจสฺสตา โข ปน ตถาคตปฺปเวทิเต ธมฺมวินเย ปริหานเมตํ.}

\prose{ปาปมิตฺโต โข ปน อยมายสฺมา, ปาปมิตฺตตา โข ปน ตถาคตปฺปเวทิเต ธมฺมวินเย ปริหานเมตํ.}

\prose{กุสีโต โข ปน อยมายสฺมา, โกสชฺชํ โข ปน ตถาคตปฺปเวทิเต ธมฺมวินเย ปริหานเมตํ.}

\prose{มุฏฺฐสฺสติํ โข ปน อยมายสฺมา, มุฏฺฐสฺสจฺจํ โข ปน ตถาคตปฺปเวทิเต ธมฺมวินเย ปริหานเมตํ.}

\prose{กุหโก โข ปน อยมายสฺมา, โกหญฺญํ โข ปน ตถาคตปฺปเวทิเต ธมฺมวินเย ปริหานเมตํ.}

\gambhira{ทสกนิปาตปาฬิ}
\prose{\csromanfolio{386}ทุพฺภโร โข ปน อยมายสฺมา, ทุพฺภรตา โข ปน ตถาคตปฺปเวทิเต ธมฺมวินเย ปริหานเมตํ.}

\prose{ทุปฺปญฺโญ โข ปน อยมายสฺมา, ทุปฺปญฺญตา โข ปน ตถาคตปฺปเวทิเต ธมฺมวินเย ปริหานเมตํ.}

\prose{โส วตาวุโส ภิกฺขุ อิเม ทส ธมฺเม อปฺปหาย อิมสฺมิํ ธมฺมวินเย วุทฺธิํ วิรูฬฺหิํ เวปุลฺลํ อาปชฺชิสฺสตีติ เนตํ ฐานํ วิชฺชติ.\csromanspacer{} โส วตาวุโส ภิกฺขุ อิเม ทส ธมฺเม ปหาย อิมสฺมิํ ธมฺมวินเย วุทฺธิํ วิรูฬฺหิํ เวปุลฺลํ อาปชฺชิสฺสตีติ ฐานเมตํ วิชฺชตีติ.\csromanspacer{}\csromanspacer{}ปญฺจมํ.}

\csromanheader{2}{6. อธิมานสุตฺต}
\proseitem{86}{เอกํ สมยํ อายสฺมา มหากสฺสโป ราชคเห วิหรติ เวฬุวเน กลนฺทกนิวาเป.\csromanspacer{} ตตฺร โข อายสฺมา มหากสฺสโป ภิกฺขู อามนฺเตสิ “อาวุโส ภิกฺขเว”ติ. “อาวุโส”ติ โข เต ภิกฺขู อายสฺมโต มหากสฺสปสฺส ปจฺจสฺโสสุํ.\csromanspacer{} อายสฺมา มหากสฺสโป เอตทโวจ\csromandash{}}

\prose{อิธาวุโส ภิกฺขุ อญฺญํ พฺยากโรติ “ขีณา ชาติ, วุสิตํ พฺรหฺมจริยํ, กตํ กรณียํ, นาปรํ อิตฺถตฺตายาติ ปชานามี”ติ.\csromanspacer{} ตเมนํ ตถาคโต วา ตถาคตสาวโก วา ฌายี สมาปตฺติกุสโล ปรจิตฺตกุสโล ปรจิตฺตปริยายกุสโล สมนุยุญฺชติ สมนุคฺคาหติ สมนุภาสติ.\csromanspacer{} โส ตถาคเตน วา ตถาคตสาวเกน วา ฌายินา สมาปตฺติกุสเลน ปรจิตฺตกุสเลน ปรจิตฺตปริยายกุสเลน สมนุยุญฺชิยมาโน สมนุคฺคาหิยมาโน สมนุภาสิยมาโน อิรีณํ อาปชฺชติ, วิจินํ อาปชฺชติ, อนยํ อาปชฺชติ, พฺยสนํ อาปชฺชติ, อนยพฺยสนํ อาปชฺชติ.}

\prose{ตเมนํ ตถาคโต วา ตถาคตสาวโก วา ฌายี สมาปตฺติกุสโล ปรจิตฺตกุสโล ปรจิตฺตปริยายกุสโล เอวํ เจตสา เจโต ปริจฺจ มนสิ กโรติ “กิํ นุ โข อยมายสฺมา อญฺญํ พฺยากโรติ ‘ขีณา ชาติ, วุสิตํ พฺรหฺมจริยํ, กตํ กรณียํ, นาปรํ อิตฺถตฺตายาติ ปชานามี’ติ”.}

\chapterheadpageread{(9) 4. เถรวคฺค}
\prose{\csromanfolio{387}ตเมนํ ตถาคโต วา ตถาคตสาวโก วา ฌายี สมาปตฺติกุสโล ปรจิตฺตกุสโล ปรจิตฺตปริยายกุสโล เอวํ เจตสา เจโต ปริจฺจ ปชานาติ\csromandash{}}

\prose{อธิมานิโก โข อยมายสฺมา อธิมานสจฺโจ อปฺปตฺเต ปตฺตสญฺญี อกเต กตสญฺญี อนธิคเต อธิคตสญฺญี อธิมาเนน อญฺญํ พฺยากโรติ “ขีณา ชาติ, วุสิตํ พฺราหฺมจริยํ, กตํ กรณียํ, นาปรํ อิตฺถตฺตายาติ ปชานามี”ติ.}

\prose{ตเมนํ ตถาคโต วา ตถาคตสาวโก วา ฌายี สมาปตฺติกุสโล ปรจิตฺตกุสโล ปรจิตฺตปริยายกุสโล เอวํ เจตสา เจโต ปริจฺจ มนสิ กโรติ “กิํ นุ โข อยมายสฺมา นิสฺสาย อธิมานิโก อธิมานสจฺโจ อปฺปตฺเต ปตฺตสญฺญี อกเต กตสญฺญี อนธิคเต อธิคตสญฺญี อธิมาเนน อญฺญํ พฺยากโรติ ‘ขีณา ชาติ, วุสิตํ พฺรหฺมจริยํ, กตํ กรณียํ, นาปรํ อิตฺถตฺตายาติ ปชานามี’ติ”.}

\prose{ตเมนํ ตถาคโต วา ตถาคตสาวโก วา ฌายี สมาปตฺติกุสโล ปรจิตฺตกุสโล ปรจิตฺตปริยายกุสโล เอวํ เจตสา เจโต ปริจฺจ ปชานาติ\csromandash{}}

\csromanmatikaanchor{mtk.125}
\csromantocmarkref{subsection}{6. อธิมานสุตฺต}{mtk.125}
\bookmark[dest={mtk.125},level=2]{6. อธิมานสุตฺต}
\prose{พหุสฺสุโต โข ปน อยมายสฺมา สุตธโร สุตสนฺนิจโย, เย เต ธมฺมา อาทิกลฺยาณา มชฺเฌกลฺยาณา ปริโยสานกลฺยาณา สาตฺถํ สพฺยญฺชนํ เกวลปริปุณฺณํ ปริสุทฺธํ พฺรหฺมจริยํ อภิวทนฺติ, ตถารูปาสฺส ธมฺมา พหุสฺสุตา โหนฺติ ธาตา วจสา ปริจิตา มนสานุเปกฺขิตา ทิฏฺฐิยา สุปฺปฏิวิทฺธา.\csromanspacer{} ตสฺมา อยมายสฺมา อธิมานิโก อธิมานสจฺโจ อปฺปตฺเต ปตฺตสญฺญี อกเต กตสญฺญี อนธิคเต อธิคตสญฺญี อธิมาเนน อญฺญํ พฺยากโรติ “ขีณา ชาติ, วุสิตํ พฺรหฺมจริยํ, กตํ กรณียํ, นาปรํ อิตฺถตฺตายาติ ปชานามี”ติ.}

\prose{ตเมนํ ตถาคโต วา ตถาคตสาวโก วา ฌายี สมาปตฺติกุสโล ปรจิตฺตกุสโล ปรจิตฺตปริยายกุสโล เอวํ เจตสา เจโต ปริจฺจ ปชานาติ\csromandash{}}

\prose{อภิชฺฌาลุ โข ปน อยมายสฺมา อภิชฺฌาปริยุฏฺฐิเตน เจตสา พหุลํ วิหรติ, อภิชฺฌาปริยุฏฺฐานํ โข ปน ตถาคตปฺปเวทิเต ธมฺมวินเย ปริหานเมตํ.}

\gambhira{ทสกนิปาตปาฬิ}
\prose{\csromanfolio{388}พฺยาปนฺโน โข ปน อยมายสฺมา พฺยาปาทปริยุฏฺฐิเตน เจตสา พหุลํ วิหรติ, พฺยาปาทปริยุฏฺฐานํ โข ปน ตถาคตปฺปเวทิเต ธมฺมวินเย ปริหานเมตํ.}

\prose{ถินมิทฺโธ โข ปน อยมายสฺมา ถินมิทฺธปริยุฏฺฐิเตน เจตสา พหุลํ วิหรติ, ถินมิทฺธปริยุฏฺฐานํ โข ปน ตถาคตปฺปเวทิเต ธมฺมวินเย ปริหานเมตํ.}

\prose{อุทฺธโต โข ปน อยมาสฺมา อุทฺธจฺจปริยุฏฺฐิเตน เจตสา พหุลํ วิหรติ, อุทฺธจฺจปริยุฏฺฐานํ โข ปน ตถาคตปฺปเวทิเต ธมฺมวินเย ปริหานเมตํ.}

\prose{วิจิกิจฺโฉ โข ปน อยมายสฺมา วิจิกิจฺฉาปริยุฏฺฐิเตน เจตสา พหุลํ วิหรติ, วิจิกิจฺฉาปริยุฏฺฐานํ โข ปน ตถาคตปฺปเวทิเต ธมฺมวินเย ปริหานเมตํ.}

\prose{กมฺมาราโม โข ปน อยมายสฺมา กมฺมรโต กมฺมารามตํ อนุยุตฺโต, กมฺมารามตา โข ปน ตถาคตปฺปเวทิเต ธมฺมวินเย ปริหานเมตํ.}

\prose{ภสฺสาราโม โข ปน อยมายสฺมา ภสฺสรโต ภสฺสารามตํ อนุยุตฺโต, ภสฺสารามตา โข ปน ตถาคตปฺปเวทิเต ธมฺมวินเย ปริหานเมตํ.}

\prose{นิทฺทาราโม โข ปน อยมายสฺมา นิทฺทารโต นิทฺทารามตํ อนุยุตฺโต, นิทฺทารามตา โข ปน ตถาคตปฺปเวทิเต ธมฺมวินเย ปริหานเมตํ.}

\prose{สงฺคณิการาโม โข ปน อยมายสฺมา สงฺคณิกรโต สงฺคณิการามตํ อนุยุตฺโต, สงฺคณิการามตา โข ปน ตถาคตปฺปเวทิเต ธมฺมวินเย ปริหานเมตํ.}

\prose{สติ โข ปน อยมายสฺมา อุตฺตริ กรณีเย โอรมตฺตเกน วิเสสาธิคเมน อนฺตรา โวสานํ อาปนฺโน, อนฺตรา โวสานคมนํ โข ปน ตถาคตปฺปเวทิเต ธมฺมวินเย ปริหานเมตํ.}

\prose{โส วตาวุโส ภิกฺขุ อิเม ทส ธมฺเม อปฺปหาย อิมสฺมิํ ธมฺมวินเย วุทฺธิํ วิรูฬฺหิํ เวปุลฺลํ อาปชฺชิสฺสตีติ เนตํ ฐานํ วิชฺชติ.\csromanspacer{} โส วตาวุโส}

\vspace{\tipitakaparskip}
\csromancenter{(9) 4.\csromanspacer{} เถรวคฺค ภิกฺขุ อิเม ทส ธมฺเม ปหาย อิมสฺมิํ ธมฺมวินเย วุทฺธิํ วิรูฬฺหิํ เวปุลฺลํ อาปชฺชิสฺสตีติ ฐานเมตํ วิชฺชตีติ.\csromanspacer{}\csromanspacer{}ฉฏฺฐํ.}
\csromanheader{2}{7. นปฺปิยสุตฺต}
\proseitem{87}{\csromanfolio{389}ตตฺร โข ภควา กาลงฺกตํ ภิกฺขุํ\footnote{กลนฺทกํ ภิกฺขุํ (สี), กาฬกภิกฺขุํ (สฺยา)} อารพฺภ ภิกฺขู อามนฺเตสิ “ภิกฺขโว”ติ. “ภทนฺเต”ติ เต ภิกฺขู ภควโต ปจฺจสฺโสสุํ.\csromanspacer{} ภควา เอตทโวจ\csromandash{}}

\prose{อิธ ภิกฺขเว ภิกฺขุ อธิกรณิโก โหติ อธิกรณสมถสฺส น วณฺณวาที.\csromanspacer{} ยมฺปิ ภิกฺขเว ภิกฺขุ อธิกรณิโก โหติ อธิกรณสมถสฺส น วณฺณวาที, อยมฺปิ ธมฺโม น ปิยตาย น ครุตาย น ภาวนาย น สามญฺญาย น เอกีภาวาย สํวตฺตติ.}

\prose{ปุน จปรํ ภิกฺขเว ภิกฺขุ น สิกฺขากาโม โหติ สิกฺขาสมาทานสฺส\footnote{สิกฺขากามสฺส (ก)} น วณฺณวาที.\csromanspacer{} ยมฺปิ ภิกฺขเว ภิกฺขุ น สิกฺขากาโม โหติ สิกฺขาสมาทานสฺส น วณฺณวาที, อยมฺปิ ธมฺโม น ปิยตาย น ครุตาย น ภาวนาย น สามญฺญาย น เอกีภาวาย สํวตฺตติ.}

\prose{ปุน จปรํ ภิกฺขเว ภิกฺขุ ปาปิจฺโฉ โหติ อิจฺฉาวินยสฺส น วณฺณวาที.\csromanspacer{} ยมฺปิ ภิกฺขเว ภิกฺขุ ปาปิจฺโฉ โหติ อิจฺฉาวินยสฺส น วณฺณวาที, อยมฺปิ ธมฺโม น ปิยตาย น ครุตาย น ภาวนาย น สามญฺญาย น เอกีภาวาย สํวตฺตติ.}

\prose{ปุน จปรํ ภิกฺขเว ภิกฺขุ โกธโน โหติ โกธวินยสฺส น วณฺณวาที.\csromanspacer{} ยมฺปิ ภิกฺขเว ภิกฺขุ โกธโน โหติ โกธวินยสฺส น วณฺณวาที, อยมฺปิ ธมฺโม น ปิยตาย น ครุตาย น ภาวนาย น สามญฺญาย น เอกีภาวาย สํวตฺตติ.}

\prose{ปุน จปรํ ภิกฺขเว ภิกฺขุ มกฺขี โหติ มกฺขวินยสฺส น วณฺณวาที.\csromanspacer{} ยมฺปิ ภิกฺขเว ภิกฺขุ มกฺขี โหติ มกฺขวินยสฺส น วณฺณวาที, อยมฺปิ ธมฺโม น ปิยตาย น ครุตาย น ภาวนาย น สามญฺญาย น เอกีภาวาย สํวตฺตติ.}

\gambhira{ทสกนิปาตปาฬิ}
\prose{\csromanfolio{390}ปุน จปรํ ภิกฺขเว ภิกฺขุ สโฐ โหติ สาเฐยฺยวินยสฺส น วณฺณวาที.\csromanspacer{} ยมฺปิ ภิกฺขเว ภิกฺขุ สโฐ โหติ สาเฐยฺยวินยสฺส น วณฺณวาที, อยมฺปิ ธมฺโม น ปิยตาย น ครุตาย น ภาวนาย น สามญฺญาย น เอกีภาวาย สํวตฺตติ.}

\csromanmatikaanchor{mtk.126}
\csromantocmarkref{subsection}{7. นปฺปิยสุตฺต}{mtk.126}
\bookmark[dest={mtk.126},level=2]{7. นปฺปิยสุตฺต}
\prose{ปุน จปรํ ภิกฺขเว ภิกฺขุ มายาวี โหติ มายาวินยสฺส น วณฺณวาที.\csromanspacer{} ยมฺปิ ภิกฺขเว ภิกฺขุ มายาวี โหติ มายาวินยสฺส น วณฺณวาที, อยมฺปิ ธมฺโม น ปิยตาย น ครุตาย น ภาวนาย น สามญฺญาย น เอกีภาวาย สํวตฺตติ.}

\prose{ปุน จปรํ ภิกฺขเว ภิกฺขุ ธมฺมานํ น นิสามกชาติโก โหติ ธมฺมนิสนฺติยา น วณฺณวาที.\csromanspacer{} ยมฺปิ ภิกฺขเว ภิกฺขุ ธมฺมานํ น นิสามกชาติโก โหติ ธมฺมนิสนฺติยา น วณฺณวาที, อยมฺปิ ธมฺโม น ปิยตาย น ครุตาย น ภาวนาย น สามญฺญาย น เอกีภาวาย สํวตฺตติ.}

\prose{ปุน จปรํ ภิกฺขเว ภิกฺขุ น ปฏิสลฺลีโน โหติ ปฏิสลฺลานสฺส น วณฺณวาที.\csromanspacer{} ยมฺปิ ภิกฺขเว ภิกฺขุ น ปฏิสลฺลีโน โหติ ปฏิสลฺลานสฺส น วณฺณวาที, อยมฺปิ ธมฺโม น ปิยตาย น ครุตาย น ภาวนาย น สามญฺญาย น เอกีภาวาย สํวตฺตติ.}

\prose{ปุน จปรํ ภิกฺขเว ภิกฺขุ สพฺรหฺมจารีนํ น ปฏิสนฺถารโก\footnote{ปฏิสนฺธารโก (ก)} โหติ ปฏิสนฺถารกสฺส น วณฺณวาที.\csromanspacer{} ยมฺปิ ภิกฺขเว ภิกฺขุ สพฺรหฺมจารีนํ น ปฏิสนฺถารโก โหติ ปฏิสนฺถารกสฺส น วณฺณวาที, อยมฺปิ ธมฺโม น ปิยตาย น ครุตาย น ภาวนาย น สามญฺญาย น เอกีภาวาย สํวตฺตติ.}

\prose{เอวรูปสฺส ภิกฺขเว ภิกฺขุโน กิญฺจาปิ เอวํ อิจฺฉา อุปฺปชฺเชยฺย “อโห วต มํ สพฺรหฺมจารี สกฺกเรยฺยุํ ครุํ กเรยฺยุํ\footnote{ครุกเรยฺยุํ (สี, สฺยา)} มาเนยฺยุํ ปูเชยฺยุนฺ”ติ.\csromanspacer{} อถ โข นํ สพฺรหฺมจารี น เจว สกฺกโรนฺติ น ครุํ กโรนฺติ\footnote{ครุกโรนฺติ (สี, สฺยา)} น มาเนนฺติ น ปูเชนฺติ.\csromanspacer{} ตํ กิสฺส เหตุ? ถถาหิสฺส ภิกฺขเว วิญฺญู สพฺรหฺมจารี เต ปาปเก อกุสเล ธมฺเม อปฺปหีเน สมนุปสฺสนฺติ.}

\prose{เสยฺยถาปิ ภิกฺขเว อสฺสขฬุงฺกสฺส กิญฺจาปิ เอวํ อิจฺฉา อุปฺปชฺเชยฺย “อโห วต มํ มนุสฺสา อาชานียฏฺฐาเน ฐเปยฺยุํ, อาชานียโภชนญฺจ โภเชยฺยุํ, อาชานียปริมชฺชนญฺจ ปริมชฺเชยฺยุนฺ”ติ.\csromanspacer{} อถ โข นํ มนุสฺสา}

\vspace{\tipitakaparskip}
\csromancenter{(9) 4.\csromanspacer{} เถรวคฺค น เจว อาชานียฏฺฐาเน ฐเปนฺติ, น จ อาชานียโภชนํ โภเชนฺติ, น จ อาชานียปริมชฺชนํ ปริมชฺชนฺติ.\csromanspacer{} ตํ กิสฺส เหตุ? ตถาหิสฺส ภิกฺขเว วิญฺญู มนุสฺสา ตานิ สาเฐยฺยานิ กูเฏยฺยานิ ชิเมฺหยฺยานิ วงฺเกยฺยานิ อปฺปหีนานิ สมนุปสฺสนฺติ.\csromanspacer{} เอวเมวํ โข ภิกฺขเว เอวรูปสฺส ภิกฺขุโน กิญฺจาปิ เอวํ อิจฺฉา อุปฺปชฺเชยฺย “อโห วต มํ สพฺรหฺมจารี สกฺกเรยฺยุํ ครุํ กเรยฺยุํ มาเนยฺยุํ ปูเชยฺยุนฺ”ติ.\csromanspacer{} อถ โข นํ สพฺรหฺมจารี น เจว สกฺกโรนฺติ น ครุํ กโรนฺติ น มาเนนฺติ น ปูเชนฺติ.\csromanspacer{} ตํ กิสฺส เหตุ? ตถาหิสฺส ภิกฺขเว วิญฺญู สพฺรหฺมจารี เต ปาปเก อกุสเล ธมฺเม อปฺปหีเน สมนุปสฺสนฺติ.}
\prose{\csromanfolio{391}อิธ ปน ภิกฺขเว ภิกฺขุ น อธิกรณิโก โหติ อธิกรณสมถสฺส วณฺณวาที.\csromanspacer{} ยมฺปิ ภิกฺขเว ภิกฺขุ น อธิกรณิโก โหติ อธิกรณสมถสฺส วณฺณวาที, อยมฺปิ ธมฺโม ปิยตาย ครุตาย ภาวนาย สามญฺญาย เอกีภาวาย สํวตฺตติ.}

\prose{ปุน จปรํ ภิกฺขเว ภิกฺขุ สิกฺขากาโม โหติ สิกฺขาสมาทานสฺส วณฺณวที.\csromanspacer{} ยมฺปิ ภิกฺขเว ภิกฺขุ สิกฺขากาโม โหติ สิกฺขาสมาทานสฺส วณฺณวาที, อยมฺปิ ธมฺโม ปิยตาย ครุตาย ภาวนาย สามญฺญาย เอกีภาวาย สํวตฺตติ.}

\prose{ปุน จปรํ ภิกฺขเว ภิกฺขุ อปฺปิจฺโฉ โหติ อิจฺฉาวินยสฺส วณฺณวาที.\csromanspacer{} ยมฺปิ ภิกฺขเว ภิกฺขุ อปฺปิจฺโฉ โหติ อิจฺฉาวินยสฺส วณฺณวาที, อยมฺปิ ธมฺโม ฯเปฯ เอกีภาวาย สํวตฺตติ.}

\prose{ปุน จปรํ ภิกฺขเว ภิกฺขุ อกฺโกธโน โหติ โกธวินยสฺส วณฺณวาที.\csromanspacer{} ยมฺปิ ภิกฺขเว ภิกฺขุ อกฺโกธโน โหติ โกธวินยสฺส วณฺณวาที, อยมฺปิ ธมฺโม ฯเปฯ เอกีภาวาย สํวตฺตติ.}

\prose{ปุน จปรํ ภิกฺขเว ภิกฺขุ อมกฺขี โหติ มกฺขวินยสฺส วณฺณวาที.\csromanspacer{} ยมฺปิ ภิกฺขเว ภิกฺขุ อมกฺขี โหติ มกฺขวินยสฺส วณฺณวาที, อยมฺปิ ธมฺโม ฯเปฯ เอกีภาวาย สํวตฺตติ.}

\prose{ปุน จปรํ ภิกฺขเว ภิกฺขุ อสโฐ โหติ สาเฐยฺยวินยสฺส วณฺณวาที.\csromanspacer{} ยมฺปิ ภิกฺขเว ภิกฺขุ อสโฐ โหติ สาเฐยฺยวินยสฺส วณฺณวาที, อยมฺปิ ธมฺโม ฯเปฯ เอกีภาวาย สํวตฺตติ.}

\gambhira{ทสกนิปาตปาฬิ}
\prose{\csromanfolio{392}ปุน จปรํ ภิกฺขเว ภิกฺขุ อมายาวี โหติ มายาวินยสฺส วณฺณวาที.\csromanspacer{} ยมฺปิ ภิกฺขเว ภิกฺขุ อมายาวี โหติ มายาวินยสฺส วณฺณวาที, อยมฺปิ ธมฺโม ฯเปฯ เอกีภาวาย สํวตฺตติ.}

\prose{ปุน จปรํ ภิกฺขเว ภิกฺขุ ธมฺมานํ นิสามกชาติโก โหติ ธมฺมนิสนฺติยา วณฺณวาที.\csromanspacer{} ยมฺปิ ภิกฺขเว ภิกฺขุ ธมฺมานํ นิสามกชาติโก โหติ ธมฺมนิสนฺติยา วณฺณวาที, อยมฺปิ ธมฺโม ฯเปฯ เอกีภาวาย สํวตฺตติ.}

\prose{ปุน จปรํ ภิกฺขเว ภิกฺขุ ปฏิสลฺลีโน โหติ ปฏิสลฺลานสฺส วณฺณวาที.\csromanspacer{} ยมฺปิ ภิกฺขเว ภิกฺขุ ปฏิสลฺลีโน โหติ ปฏิสลฺลานสฺส วณฺณวาที, อยมฺปิ ธมฺโม ฯเปฯ เอกีภาวาย สํวตฺตติ.}

\prose{ปุน จปรํ ภิกฺขเว ภิกฺขุ สพฺรหฺมจารีนํ ปฏิสนฺถารโก โหติ ปฏิสนฺถารกสฺส วณฺณวาที.\csromanspacer{} ยมฺปิ ภิกฺขเว ภิกฺขุ สพฺรหฺมจารีนํ ปฏิสนฺถารโก โหติ ปฏิสนฺถารกสฺส วณฺณวาที, อยมฺปิ ธมฺโม ปิยตาย ครุตาย ภาวนาย สามญฺญาย เอกีภาวาย สํวตฺตติ.}

\prose{เอวรูปสฺส ภิกฺขเว ภิกฺขุโน กิญฺจาปิ น เอวํ อิจฺฉา อุปฺปชฺเชยฺย “อโห วต มํ สพฺรหฺมจารี สกฺกเรยฺยุํ ครุํ กเรยฺยุํ มาเนยฺยุํ ปูเชยฺยุนฺ”ติ.\csromanspacer{} อถ โข นํ สพฺรหฺมจารี สกฺกโรนฺติ ครุํ กโรนฺติ มาเนนฺติ ปูเชนฺติ.\csromanspacer{} ตํ กิสฺส เหตุ? ตถาหิสฺส ภิกฺขเว วิญฺญู สพฺรหฺมจารี เต ปาปเก อกุสเล ธมฺเม ปหีเน สมนุปสฺสนฺติ.}

\prose{เสยฺยถาปิ ภิกฺขเว ภทฺทสฺส อสฺสาชานียสฺส กิญฺจาปิ น เอวํ อิจฺฉา อุปฺปชฺเชยฺย “อโห วต มํ มนุสฺสา อาชานียฏฺฐาเน ฐเปยฺยุํ, อาชานียโภชนญฺจ โภเชยฺยุํ, อาชานียปริมชฺชนญฺจ ปริมชฺเชยฺยุนฺ”ติ.\csromanspacer{} อถ โข นํ มนุสฺสา อาชานียฏฺฐาเน จ ฐเปนฺติ, อาชานียโภชนญฺจ โภเชนฺติ อาชานียปริมชฺชนญฺจ ปริมชฺชนฺติ.\csromanspacer{} ตํ กิสฺส เหตุ? ตถาหิสฺส ภิกฺขเว วิญฺญู มนุสฺสา ตานิ สาเฐยฺยานิ กูเฏยฺยานิ ชิเมฺหยฺยานิ วงฺเกยฺยานิ ปหีนานิ สมนุปสฺสนฺติ.}

\prose{เอวเมวํ โข ภิกฺขเว เอวรูปสฺส ภิกฺขุโน กิญฺจาปิ น เอวํ อิจฺฉา อุปฺปชฺเชยฺย “อโห วต มํ สพฺรหฺมจารี สกฺกเรยฺยุํ ครุํ กเรยฺยุํ มาเนยฺยุํ ปูเชยฺยุนฺ”ติ.\csromanspacer{} อถ โข นํ สพฺรหฺมจารี สกฺกโรนฺติ ครุํ กโรนฺติ มาเนนฺติ ปูเชนฺติ.\csromanspacer{} ตํ กิสฺส เหตุ? ตถาหิสฺส ภิกฺขเว วิญฺญู สพฺรหฺมจารี เต ปาปเก อกุสเล ธมฺเม ปหีเน สมนุปสฺสนฺตีติ.\csromanspacer{}\csromanspacer{}สตฺตมํ.}

\csromanheader{2}{(9) 4. เถรวคฺค \textbf{8. อกฺโกสกสุตฺต}}
\proseitem{88}{\csromanfolio{393}โย โส ภิกฺขเว ภิกฺขุ อกฺโกสกปริภาสโก อริยูปวาที สพฺรหฺมจารีนํ.\csromanspacer{} ฐานเมตํ อวกาโส\footnote{อฏฺฐานเมตํ อนวกาโส (สี, สฺยา, อิ)}, ยํ โส ทสนฺนํ พฺยสนานํ อญฺญตรํ พฺยสนํ นิคจฺเฉยฺย\footnote{น นิคจฺเฉยฺย (สี, สฺยา, อิ)}.\csromanspacer{} กตเมสํ ทสนฺนํ? อนธิคตํ นาธิคจฺฉติ, อธิคตา ปริหายติ, สทฺธมฺมสฺส น โวทายนฺติ, สทฺธมฺเมสุ วา อธิมานิโก โหติ, อนภิรโต วา พฺรหฺมจริยํ จรติ, อญฺญตรํ วา สํกิลิฏฺฐํ อาปตฺติํ อาปชฺชติ, คาฬฺหํ วา โรคาตงฺกํ ผุสติ, อุมฺมาทํ วา ปาปุณาติ จิตฺตกฺเขปํ, สมฺมูโฬฺห กาลํ กโรติ, กายสฺส เภทา ปรํ มรณา อปายํ ทุคฺคติํ วินิปาตํ นิรยํ อุปปชฺชติ.\csromanspacer{} โย โส ภิกฺขเว ภิกฺขุ อกฺโกสกปริภาสโก อริยูปวาที สพฺรหฺมจารีนํ.\csromanspacer{} ฐานเมตํ อวกาโส, ยํ โส อิเมสํ ทสนฺนํ พฺยสนานํ อญฺญตรํ พฺยสนํ นิคจฺเฉยฺยาติ.\csromanspacer{}\csromanspacer{}อฏฺฐมํ.}

\csromanheader{2}{9. โกกาลิกสุตฺต}
\proseitem{89}{\csromansymbolfootnote{*}{สํ 1. 152; ขุ 1. 379 ปิฏฺเฐสุปิ.}อถ โข โกกาลิโก ภิกฺขุ เยน ภควา เตนุปสงฺกมิ, อุปสงฺกมิตฺวา ภควนฺตํ อภิวาเทตฺวา เอกมนฺตํ นิสีทิ, เอกมนฺตํ นิสินฺโน โข โกกาลิโก ภิกฺขุ ภควนฺตํ เอตทโวจ “ปาปิจฺฉา ภนฺเต สาริปุตฺตโมคฺคลฺลานา, ปาปิกานํ อิจฺฉานํ วสํ คตา”ติ.\csromanspacer{} มา เหวํ โกกาลิก, มา เหวํ โกกาลิก, ปสาเทหิ โกกาลิก สาริปุตฺตโมคฺคลฺลาเนสุ จิตฺตํ, เปสลา สาริปุตฺตโมคฺคลฺลานาติ.}

\prose{ทุติยมฺปิ โข โกกาลิโก ภิกฺขุ ภควนฺตํ เอตทโวจ “กิญฺจาปิ เม ภนฺเต ภควา สทฺธายิโก ปจฺจยิโก, อถ โข ปาปิจฺฉาว สาริปุตฺตโมคฺคลฺลานา ปาปิกานํ อิจฺฉานํ วสํ คตา”ติ.\csromanspacer{} มา เหวํ โกกาลิก, มา เหวํ โกกาลิก, ปสาเทหิ โกกาลิก สาริปุตฺตโมคฺคลฺลาเนสุ จิตฺตํ, เปสลา สาริปุตฺตโมคฺคลฺลานาติ.}

\prose{ตติยมฺปิ โข ภิกฺขเว โกกาลิโก ภิกฺขุ ภควนฺตํ เอตทโวจ “กิญฺจาปิ เม ภนฺเต ภควา สทฺธายิโก ปจฺจยิโก, อถ โข ปาปิจฺฉาว สาริปุตฺตโมคฺคลฺลานา ปาปิกานํ อิจฺฉานํ วสํ คตา”ติ.\csromanspacer{} มา เหวํ โกกาลิก, มา เหวํ โกกาลิก, ปสาเทหิ}

\vspace{\tipitakaparskip}
\csromancenter{ทสกนิปาตปาฬิ โกกาลิก สาริปุตฺตโมคฺคลฺลาเนสุ จิตฺตํ, เปสลา สาริปุตฺตโมคฺคลฺลานาติ.}
\prose{\csromanfolio{394}อถ โข โกกาลิโก ภิกฺขุ อุฏฺฐายาสนา ภควนฺตํ อภิวาเทตฺวา ปทกฺขิณํ กตฺวา ปกฺกามิ, อจิรปกฺกนฺตสฺส จ โกกาลิกสฺส ภิกฺขุโน สาสปมตฺตีหิ ปีฬกาหิ สพฺโพ กาโย ผุโฏ อโหสิ, สาสปมตฺติโย หุตฺวา มุคฺคมตฺติโย อเหสุํ, มุคฺคมตฺติโย หุตฺวา กลายมตฺติโย อเหสุํ, กลายมตฺติโย หุตฺวา โกลฏฺฐิมตฺติโย อเหสุํ, โกลฏฺฐิมตฺติโย หุตฺวา โกลมตฺติโย อเหสุํ, โกลมตฺติโย หุตฺวา อามลกมตฺติโย อเหสุํ, อามลกมตฺติโย หุตฺวา (ติณฺฑุกมตฺติโย อเหสุํ, ติณฺฑุกมตฺติโย หุตฺวา)\footnote{สํ 1. 152 ปิฏฺเฐ; ขุ 1. 380 ปิฏฺเฐสุ นตฺถิ.} เพฬุวสลาฏุกมตฺติโย อเหสุํ, เพฬุวสลาฏุกมตฺติโย หุตฺวา พิลฺลมตฺติโย อเหสุํ, พิลฺลมตฺติโย หุตฺวา ปภิชฺชิํสุ, ปุพฺพญฺจ โลหิตญฺจ ปคฺฆริํสุ, โส สุทํ กทลิปตฺเตสุ เสติ มจฺโฉว วิสคิลิโต.}

\prose{อถ โข ตุรู ปจฺเจกพฺรหฺมา\footnote{ตุทุปฺปจฺเจกพฺรหฺมา (สี, อิ), ตุทิ ปจฺเจกพฺรหฺมา (สฺยา), ตุริ ปจฺเจกพฺรหฺมา (ก) สํ 1. 151 ปิฏฺเฐปิ.} เยน โกกาลิโก ภิกฺขุ เตนุปสงฺกมิ, อุปสงฺกมิตฺวา เวหาเส ฐตฺวา โกกาลิกํ ภิกฺขุํ เอตทโวจ “ปสาเทหิ โกกาลิก สาริปุตฺตโมคฺคลฺลาเนสุ จิตฺตํ, เปสลา สาริปุตฺตโมคฺคลฺลานา”ติ.\csromanspacer{} โกสิ ตฺวํ อาวุโสติ.\csromanspacer{} อหํ ตุรู ปจฺเจกพฺรหฺมาติ.\csromanspacer{} นนุ ตฺวํ อาวุโส ภควตา อนาคามี พฺยากโต, อถ กิญฺจรหิ อิธาคโต, ปสฺส ยาวญฺจ เต อิทํ อปรทฺธนฺติ.}

\csromanmatikaanchor{mtk.127}
\csromantocmarkref{subsection}{8. อกฺโกสกสุตฺต}{mtk.127}
\bookmark[dest={mtk.127},level=2]{8. อกฺโกสกสุตฺต}
\prose{อถ โข ตุรู ปจฺเจกพฺรหฺมา โกกาลิกํ ภิกฺขุํ คาถาหิ อชฺฌภาสิ\csromandash{}}

\prose{“ปุริสสฺส หิ ชาตสฺส, กุฐารี ชายเต มุเข.}

\csromanmatikaanchor{mtk.128}
\csromantocmarkref{subsection}{9. โกกาลิกสุตฺต}{mtk.128}
\bookmark[dest={mtk.128},level=2]{9. โกกาลิกสุตฺต}
\prose{ยาย ฉินฺทติ อตฺตานํ, พาโล ทุพฺภาสิตํ ภณํ.}

\prose{โย นินฺทิยํ ปสํสติ, ตํ วา นินฺทติ โย ปสํสิโย.}

\prose{วิจินาติ มุเขน โส กลิํ, กลินา เตน สุขํ น วินฺทติ.}

\prose{อปฺปมตฺตโก อยํ กลิ, โย อกฺเขสุ ธนปราชโย.}

\prose{สพฺพสฺสาปิ สหาปิ อตฺตนา, อยเมว มหตฺตโร กลิ.}

\prose{โย สุคเตสุ มนํ ปทูสเย.}

\chapterheadpageread{(9) 4. เถรวคฺค}
\prose{\csromanfolio{395}สตํ สหสฺสานํ นิรพฺพุทานํ, ฉตฺติํสติ ปญฺจ จ อพฺพุทานิ.}

\prose{ยมริยครหี นิรยํ อุเปติ, วาจํ มนญฺจ ปณิธาย ปาปกนฺ”ติ.}

\prose{อถ โข โกกาลิโก ภิกฺขุ เตเนว อาพาเธน กาลมกาสิ, กาลงฺกโต จ โกกาลิโก ภิกฺขุ ปทุมํ นิรยํ อุปปชฺชติ สาริปุตฺตโมคฺคลฺลาเนสุ จิตฺตํ อาฆาเตตฺวา.}

\prose{อถ โข พฺรหฺมา สหมฺปติ อภิกฺกนฺตาย รตฺติยา อภิกฺกนฺตวณฺโณ เกวลกปฺปํ เชตวนํ โอภาเสตฺวา เยน ภควา เตนุปสงฺกมิ, อุปสงฺกมิตฺวา ภควนฺตํ อภิวาเทตฺวา เอกมนฺตํ อฏฺฐาสิ, เอกมนฺตํ ฐิโต โข พฺรหฺมา สหมฺปติ ภควนฺตํ เอตทโวจ “โกกาลิโก ภนฺเต ภิกฺขุ กาลงฺกโต, กาลงฺกโต จ ภนฺเต โกกาลิโก ภิกฺขุ ปทุมํ นิรยํ อุปปนฺโน สาริปุตฺตโมคฺคลฺลาเนสุ จิตฺตํ อาฆาเตตฺวา”ติ.\csromanspacer{} อิทมโวจ พฺรหฺมา สหมฺปติ, อิทํ วตฺวา ภควนฺตํ อภิวาเทตฺวา ปทกฺขิณํ กตฺวา ตตฺเถวนฺตรธายิ.}

\prose{อถ โข ภควา ตสฺสา รตฺติยา อจฺจเยน ภิกฺขู อามนฺเตสิ\csromandash{}อิมํ ภิกฺขเว รตฺติํ พฺรหฺมา สหมฺปติ อภิกฺกนฺตาย รตฺติยา อภิกฺกนฺตวณฺโณ เกวลกปฺปํ เชตวนํ โอภาเสตฺวา เยนาหํ เตนุปสงฺกมิ, อุปสงฺกมิตฺวา มํ อภิวาเทตฺวา เอกมนฺตํ อฏฺฐาสิ, เอกมนฺตํ ฐิโต โข ภิกฺขเว พฺรหฺมา สหมฺปติ มํ เอตทโวจ “โกกาลิโก ภนฺเต ภิกฺขุ กาลงฺกโต, กาลงฺกโต จ ภนฺเต โกกาลิโก ภิกฺขุ ปทุมํ นิรยํ อุปปนฺโน สาริปุตฺตโมคฺคลฺลาเนสุ จิตฺตํ อาฆาเตตฺวา”ติ.\csromanspacer{} อิทมโวจ ภิกฺขเว พฺรหฺมา สหมฺปติ, อิทํ วตฺวา มํ อภิวาเทตฺวา ปทกฺขิณํ กตฺวา ตตฺเถวนฺตรธายีติ.}

\prose{เอวํ วุตฺเต อญฺญตโร ภิกฺขุ ภควนฺตํ เอตทโวจ “กีว ทีฆํ นุ โข ภนฺเต ปทุเม นิรเย อายุปฺปมาณนฺ”ติ.\csromanspacer{} ทีฆํ โข ภิกฺขุ ปทุเม นิรเย อายุปฺปมาณํ, น ตํ สุกรํ สงฺขาตุํ “เอตฺตกานิ วสฺสานี”ติ วา “เอตฺตกานิ วสฺสสตานี”ติ วา “เอตฺตกานิ วสฺสสหสฺสานี”ติ วา “เอตฺตกานิ วสฺสสตสหสฺสานี”ติ วาติ.}

\prose{สกฺกา ปน ภนฺเต อุปมํ กาตุนฺติ. “สกฺกา ภิกฺขู”ติ ภควา อโวจ.\csromanspacer{} เสยฺยถาปิ ภิกฺขุ วีสติ ขาริโก โกสลโก ติลวาโห, ตโต ปุริโส วสฺสสตสฺส วสฺสสตสฺส อจฺจเยน เอกเมกํ ติลํ}

\vspace{\tipitakaparskip}
\csromancenter{ทสกนิปาตปาฬิ อุทฺธเรยฺย, ขิปฺปตรํ โข โส ภิกฺขุ วีสติขาริโก โกสลโก ติลวาโห อิมินา อุปกฺกเมน ปริกฺขยํ ปริยาทานํ คจฺเฉยฺย, น เตฺวว เอโก อพฺพุโท นิรโย.\csromanspacer{} เสยฺยถาปิ ภิกฺขุ วีสติ อพฺพุทา นิรยา, เอวเมโก นิรพฺพุโท นิรโย.\csromanspacer{} เสยฺยถาปิ ภิกฺขุ วีสติ นิรพฺพุทา นิรยา, เอวเมโก อพโพ นิรโย.\csromanspacer{} เสยฺยถาปิ ภิกฺขุ วีสติ อพพา นิรยา, เอวเมโก อฏโฏ นิรโย.\csromanspacer{} เสยฺยถาปิ ภิกฺขุ วีสติ อฏฏา นิรยา, เอวเมโก อหโห นิรโย.\csromanspacer{} เสยฺยถาปิ ภิกฺขุ วีสติ อหหา นิรยา, เอวเมโก กุมุโท นิรโย.\csromanspacer{} เสยฺยถาปิ ภิกฺขุ วีสติ กุมุทา นิรยา, เอวเมโก โสคนฺธิโก นิรโย.\csromanspacer{} เสยฺยถาปิ ภิกฺขุ วีสติ โสคนฺธิกา นิรยา, เอวเมโก อุปฺปลโก นิรโย.\csromanspacer{} เสยฺยถาปิ ภิกฺขุ วีสติ อุปฺปลกา นิรยา, เอวเมโก ปุณฺฑรีโก นิรโย.\csromanspacer{} เสยฺยถาปิ ภิกฺขุ วีสติ ปุณฺฑรีกา นิรยา, เอวเมโก ปทุโม นิรโย.\csromanspacer{} ปทุมํ โข ปน ภิกฺขุ นิรยํ โกกาลิโก ภิกฺขุ อุปปนฺโน สาริปุตฺตโมคฺคลฺลาเนสุ จิตฺตํ อาฆาเตตฺวาติ.\csromanspacer{} อิทมโวจ ภควา, อิทํ วตฺวาน สุคโต อถาปรํ เอตทโวจ สตฺถา\csromandash{}}
\vspace{0.5\baselineskip}
\csromangathameasure{“ปุริสสฺส หิ ชาตสฺส, กุฐารี ชายเต มุเข. \\ ยาย ฉินฺทติ อตฺตานํ, พาโล ทุพฺภาสิตํ ภณํ. \\ โย นินฺทิยํ ปสํสติ, ตํ วา นินฺทติ โย ปสํสิโย. \\ วิจินาติ มุเขน โส กลิํ, กลินา เตน สุขํ น วินฺทติ. \\ อปฺปมตฺตโก อยํ กลิ, โย อกฺเขสุ ธนปราชโย. \\ สพฺพสฺสาปิ สหาปิ อตฺตนา, อยเมว มหตฺตโร กลิ.}
\csromangathagroup{\csromangathastanza{\csromanfolio{396}“ปุริสสฺส หิ ชาตสฺส, กุฐารี ชายเต มุเข. \\ ยาย ฉินฺทติ อตฺตานํ, พาโล ทุพฺภาสิตํ ภณํ.}\csromangathastanza{โย นินฺทิยํ ปสํสติ, ตํ วา นินฺทติ โย ปสํสิโย. \\ วิจินาติ มุเขน โส กลิํ, กลินา เตน สุขํ น วินฺทติ.}\csromangathastanza{อปฺปมตฺตโก อยํ กลิ, โย อกฺเขสุ ธนปราชโย. \\ สพฺพสฺสาปิ สหาปิ อตฺตนา, อยเมว มหตฺตโร กลิ.}}

\prose{โย สุคฺคเตสุ มนํ ปทูสเย.}

\csromangathameasure{สตํ สหสฺสานํ นิรพฺพุทานํ, \\ ฉตฺติํสติ ปญฺจ จ อพฺพุทานิ. \\ ยมริยครหี นิรยํ อุเปติ, \\ วาจํ มนญฺจ ปณิธาย ปาปกนฺ”ติ.นวมํ.}
\csromangathagroup{\csromangathastanza{สตํ สหสฺสานํ นิรพฺพุทานํ, \\ ฉตฺติํสติ ปญฺจ จ อพฺพุทานิ. \\ ยมริยครหี นิรยํ อุเปติ, \\ วาจํ มนญฺจ ปณิธาย ปาปกนฺ”ติ.\csromanspacer{}\csromanspacer{}นวมํ.}}

\csromanheader{2}{10. ขีณาสวพลสุตฺต}
\proseitem{90}{อถ โข อายสฺมา สาริปุตฺโต เยน ภควา เตนุปสงฺกมิ, อุปสงฺกมิตฺวา ภควนฺตํ อภิวาเทตฺวา เอกมนฺตํ นิสีทิ, เอกมนฺตํ นิสินฺนํ โข อายสฺมนฺตํ สาริปุตฺตํ ภควา เอตทโวจ “กติ นุ โข สาริปุตฺต}

\vspace{\tipitakaparskip}
\csromancenter{(9) 4.\csromanspacer{} เถรวคฺค ขีณาสวสฺส ภิกฺขุโน พลานิ, เยหิ พเลหิ สมนฺนาคโต ขีณาสโว ภิกฺขุ อาสวานํ ขยํ ปฏิชานาติ ‘ขีณา เม อาสวา’ติ”.}
\prose{\csromanfolio{397}ทส ภนฺเต ขีณาสวสฺส ภิกฺขุโน พลานิ, เยหิ พเลหิ สมนฺนาคโต ขีณาสโว ภิกฺขุ อาสวานํ ขยํ ปฏิชานาติ “ขีณา เม อาสวา”ติ.\csromanspacer{} กตมานิ ทส? * อิธ ภนฺเต ขีณาสวสฺส ภิกฺขุโน อนิจฺจโต สพฺเพ สงฺขารา ยถาภูตํ สมฺมปฺปญฺญาย สุทิฏฺฐา โหนฺติ.\csromanspacer{} ยมฺปิ ภนฺเต ขีณาสวสฺส ภิกฺขุโน อนิจฺจโต สพฺเพ สงฺขารา ยถาภูตํ สมฺมปฺปญฺญาย สุทิฏฺฐา โหนฺติ, อิทมฺปิ ภนฺเต ขีณาสวสฺส ภิกฺขุโน พลํ โหติ, ยํ พลํ อาคมฺม ขีณาสโว ภิกฺขุ อาสวานํ ขยํ ปฏิชานาติ “ขีณา เม อาสวา”ติ. (1)}

\prose{ปุน จปรํ ภนฺเต ขีณาสวสฺส ภิกฺขุโน องฺคารกาสูปมา กามา ยถาภูตํ สมฺมปฺปญฺญาย สุทิฏฺฐา โหนฺติ.\csromanspacer{} ยมฺปิ ภนฺเต ขีณาสวสฺส ภิกฺขุโน องฺคารกาสูปมา กามา ยถาภูตํ สมฺมปฺปญฺญาย สุทิฏฺฐา โหนฺติ, อิทมฺปิ ภนฺเต ขีณาสวสฺส ภิกฺขุโน พลํ โหติ, ยํ พลํ อาคมฺม ขีณาสโว ภิกฺขุ อาสวานํ ขยํ ปฏิชานาติ “ขีณา เม อาสวา”ติ. (2)}

\prose{ปุน จปรํ ภนฺเต ขีณาสวสฺส ภิกฺขุโน วิเวกนินฺนํ จิตฺตํ โหติ วิเวกโปณํ วิเวกปพฺภารํ วิเวกฏฺฐํ เนกฺขมฺมาภิรตํ พฺยนฺตีภูตํ สพฺพโส อาสวฏฺฐานิเยหิ ธมฺเมหิ.\csromanspacer{} ยมฺปิ ภนฺเต ขีณาสวสฺส ภิกฺขุโน วิเวกนินฺนํ จิตฺตํ โหติ วิเวกโปณํ วิเวกปพฺภารํ วิเวกฏฺฐํ เนกฺขมฺมาภิรตํ พฺยนฺตีภูตํ สพฺพโส อาสวฏฺฐานิเยหิ ธมฺเมหิ.\csromanspacer{} อิทมฺปิ ภนฺเต ขีณาสวสฺส ภิกฺขุโน พลํ โหติ, ยํ พลํ อาคมฺม ขีณาสโว ภิกฺขุ อาสวานํ ขยํ ปฏิชานาติ “ขีณา เม อาสวา”ติ. (3)}

\prose{ปุน จปรํ ภนฺเต ขีณาสวสฺส ภิกฺขุโน จตฺตาโร สติปฏฺฐานา ภาวิตา โหนฺติ สุภาวิตา.\csromanspacer{} ยมฺปิ ภนฺเต ขีณาสวสฺส ภิกฺขุโน จตฺตาโร สติปฏฺฐานา ภาวิตา โหนฺติ สุภาวิตา.\csromanspacer{} อิทมฺปิ ภนฺเต ขีณาสวสฺส ภิกฺขุโน พลํ โหติ, ยํ พลํ อาคมฺม ขีณาสโว ภิกฺขุ อาสวานํ ขยํ ปฏิชานาติ “ขีณา เม อาสวา”ติ. (4)}

\prose{ปุน จปรํ ภนฺเต ขีณาสวสฺส ภิกฺขุโน จตฺตาโร สมฺมปฺปธานา ภาวิตา โหนฺติ สุภาวิตา ฯเปฯ จตฺตาโร อิทฺธิปาทา ภาวิตา โหนฺติ}

\vspace{\tipitakaparskip}
\csromancenter{ทสกนิปาตปาฬิ สุภาวิตา ฯเปฯ ปญฺจินฺทฺริยานิ.\csromanspacer{} ปญฺจ พลานิ ภาวิตานิ โหนฺติ สุภาวิตานิ.\csromanspacer{} สตฺต โพชฺฌงฺคา ภาวิตา โหนฺติ สุภาวิตา.\csromanspacer{} อริโย อฏฺฐงฺคิโก มคฺโค ภาวิโต โหติ สุภาวิโต.\csromanspacer{} ยมฺปิ ภนฺเต ขีณาสวสฺส ภิกฺขุโน อริโย อฏฺฐงฺคิโก มคฺโค ภาวิโต โหติ สุภาวิโต.\csromanspacer{} อิทมฺปิ ภนฺเต ขีณาสวสฺส ภิกฺขุโน พลํ โหติ, ยํ พลํ อาคมฺม ขีณาสโว ภิกฺขุ อาสวานํ ขยํ ปฏิชานาติ “ขีณา เม อาสวา”ติ. (5-10)}
\csromanmatikaanchor{mtk.129}
\csromanmatikaanchor{mtk.130}
\csromantocmarkref{subsection}{10. ขีณาสวพลสุตฺต}{mtk.129}
\bookmark[dest={mtk.129},level=2]{10. ขีณาสวพลสุตฺต}
\csromantocmarkref{section}{อุทฺทานคาถา}{mtk.130}
\bookmark[dest={mtk.130},level=1]{อุทฺทานคาถา}
\prose{\csromanfolio{398}อิมานิ โข ภนฺเต ทส ขีณาสวสฺส ภิกฺขุโน พลานิ, เยหิ พเลหิ สมนฺนาคโต ขีณาสโว ภิกฺขุ อาสวานํ ขยํ ปฏิชานาติ “ขีณา เม อาสวา”ติ.\csromanspacer{}\csromanspacer{}ทสมํ.\csromanspacer{} เถรวคฺโค จตุตฺโถ.}

\titlehead{\textbf{ตสฺสุทฺทานํ}}
\csromangathameasure{วาหนานนฺโท ปุณฺณิโย, พฺยากรํ กตฺถิมานิโก. \\ นปิยกฺโกสโกกาลิ, ขีณาสวพเลน จาติ.}
\csromangathagroup{\csromangathastanza{วาหนานนฺโท ปุณฺณิโย, พฺยากรํ กตฺถิมานิโก. \\ นปิยกฺโกสโกกาลิ, ขีณาสวพเลน จาติ.}}

\csromanheader{1}{\textbf{(10) 5. อุปาลิวคฺค}}
\csromanheader{2}{1. กามโภคีสุตฺต}
\proseitem{91}{เอกํ สมยํ ภควา สาวตฺถิยํ วิหรติ เชตวเน อนาถปิณฺฑิกสฺส อาราเม.\csromanspacer{} อถ โข อนาถปิณฺฑิโก คหปติ เยน ภควา เตนุปสงฺกมิ, อุปสงฺกมิตฺวา ภควนฺตํ อภิวาเทตฺวา เอกมนฺตํ นิสีทิ, เอกมนฺตํ นิสินฺนํ โข อนาถปิณฺฑิกํ คหปติํ ภควา เอตทโวจ\csromandash{}}

\prose{ทสยิเม คหปติ กามโภคี สนฺโต สํวิชฺชมานา โลกสฺมิํ.\csromanspacer{} กตเม ทส? อิธ คหปติ เอกจฺโจ กามโภคี อธมฺเมน โภเค ปริเยสติ สาหเสน, อธมฺเมน โภเค ปริเยสิตฺวา สาหเสน น อตฺตานํ สุเขติ น ปีเณติ\footnote{น อตฺตานํ สุเขติ ปีเณติ (สี, สฺยา, อิ) เอวมุปริปิ.} น สํวิภชติ น ปุญฺญานิ กโรติ. (1)}

\chapterheadpageread{(10) 5. อุปาลิวคฺค}
\prose{\csromanfolio{399}อิธ ปน คหปติ เอกจฺโจ กามโภคี อธมฺเมน โภเค ปริเยสติ สาหเสน, อธมฺเมน โภเค ปริเยสิตฺวา สาหเสน อตฺตานํ สุเขติ ปีเณติ น สํวิภชติ น ปุญฺญานิ กโรติ. (2)}

\prose{อิธ ปน คหปติ เอกจฺโจ กามโภคี อธมฺเมน โภเค ปริเยสติ สาหเสน, อธมฺเมน โภเค ปริเยสิตฺวา สาหเสน อตฺตานํ สุเขติ ปีเณติ สํวิภชติ ปุญฺญานิ กโรติ. (3)}

\prose{อิธ ปน คหปติ เอกจฺโจ กามโภคี ธมฺมาธมฺเมน โภเค ปริเยสติ สาหเสนปิ อสาหเสนปิ, ธมฺมาธมฺเมน โภเค ปริเยสิตฺวา สาหเสนปิ อสาหเสนปิ น อตฺตานํ สุเขติ น ปีเณติ น สํวิภชติ น ปุญฺญานิ กโรติ. (4)}

\prose{อิธ ปน คหปติ เอกจฺโจ กามโภคี ธมฺมาธมฺเมน โภเค ปริเยสติ สาหเสนปิ อสาหเสนปิ, ธมฺมาธมฺเมน โภเค ปริเยสิตฺวา สาหเสนปิ อสาหเสนปิ อตฺตานํ สุเขติ ปีเณติ น สํวิภชติ น ปุญฺญานิ กโรติ. (5)}

\prose{อิธ ปน คหปติ เอกจฺโจ กามโภคี ธมฺมาธมฺเมน โภเค ปริเยสติ สาหเสนปิ อสาหเสนปิ, ธมฺมาธมฺเมน โภเค ปริเยสิตฺวา สาหเสนปิ อสาหเสนปิ อตฺตานํ สุเขติ ปีเณติ สํวิภชติ ปุญฺญานิ กโรติ. (6)}

\prose{อิธ ปน คหปติ เอกจฺโจ กามโภคี ธมฺเมน โภเค ปริเยสติ อสาหเสน, ธมฺเมน โภเค ปริเยสิตฺวา อสาหเสน น อตฺตานํ สุเขติ น ปีเณติ น สํวิภชติ น ปุญฺญานิ กโรติ. (7)}

\prose{อิธ ปน คหปติ เอกจฺโจ กามโภคี ธมฺเมน โภเค ปริเยสติ อสาหเสน, ธมฺเมน โภเค ปริเยสิตฺวา อสาหเสน อตฺตานํ สุเขติ ปีเณติ น สํวิภชติ น ปุญฺญานิ กโรติ. (8)}

\csromanmatikaanchor{mtk.131}
\csromantocmarknopage{section}{(10) 5. อุปาลิวคฺค}{mtk.131}
\bookmark[dest={mtk.131},level=1]{(10) 5. อุปาลิวคฺค}
\prose{อิธ ปน คหปติ เอกจฺโจ กามโภคี ธมฺเมน โภเค ปริเยสติ อสาหเสน, ธมฺเมน โภเค ปริเยสิตฺวา อสาหเสน อตฺตานํ สุเขติ ปีเณติ สํวิภชติ ปุญฺญานิ กโรติ, เต จ โภเค คถิโต\footnote{คมิโต (ก) อํ 1. 278 ปิฏฺเฐปิ ปสฺสิตพฺพํ.}}

\csromanmatikaanchor{mtk.132}
\csromantocmarkref{subsection}{1. กามโภคีสุตฺต}{mtk.132}
\bookmark[dest={mtk.132},level=2]{1. กามโภคีสุตฺต}
\vspace{\tipitakaparskip}
\csromancenter{ทสกนิปาตปาฬิ มุจฺฉิโต อชฺโฌสนฺโน\footnote{อชฺฌาปนฺโน (สพฺพตฺถ) อํ 1. 278 ปิฏฺเฐ สุตฺตวณฺณนา ฏีกา โอโลเกตพฺพา.} อนาทีนวทสฺสาวี อนิสฺสรณปญฺโญ ปริภุญฺชติ. (9)}
\prose{\csromanfolio{400}อิธ ปน คหปติ เอกจฺโจ กามโภคี ธมฺเมน โภเค ปริเยสติ อสาหเสน, ธมฺเมน โภเค ปริเยสิตฺวา อสาหเสน อตฺตานํ สุเขติ ปีเณติ สํวิภชติ ปุญฺญานิ กโรติ, เต จ โภเค อคถิโต อมุจฺฉิโต อนชฺโฌสนฺโน อาทีนวทสฺสาวี นิสฺสรณปญฺโญ ปริภุญฺชติ. (10)}

\prose{ตตฺร คหปติ ยฺวายํ กามโภคี อธมฺเมน โภเค ปริเยสติ สาหเสน, อธมฺเมน โภเค ปริเยสิตฺวา สาหเสน น อตฺตานํ สุเขติ น ปีเณติ น สํวิภชติ น ปุญฺญานิ กโรติ.\csromanspacer{} อยํ คหปติ กามโภคี ตีหิ ฐาเนหิ คารโยฺห. “อธมฺเมน โภเค ปริเยสติ สาหเสนา”ติ อิมินา ปฐเมน ฐาเนน คารโยฺห, “น อตฺตานํ สุเขติ น ปีเณตี”ติ อิมินา ทุติเยน ฐาเนน คารโยฺห, “น สํวิภชติ น ปุญฺญานิ กโรตี”ติ อิมินา ตติเยน ฐาเนน คารโยฺห.\csromanspacer{} อยํ คหปติ กามโภคี อิเมหิ ตีหิ ฐาเนหิ คารโยฺห. (1)}

\prose{ตตฺร คหปติ ยฺวายํ กามโภคี อธมฺเมน โภเค ปริเยสติ สาหเสน, อธมฺเมน โภเค ปริเยสิตฺวา สาหเสน อตฺตานํ สุเขติ ปีเณติ น สํวิภชติ น ปุญฺญานิ กโรติ.\csromanspacer{} อยํ คหปติ กามโภคี ทฺวีหิ ฐาเนหิ คารโยฺห, เอเกน ฐาเนน ปาสํโส. “อธมฺเมน โภเค ปริเยสติ สาหเสนา”ติ อิมินา ปฐเมน ฐาเนน คารโยฺห, “อตฺตานํ สุเขติ ปีเณตี”ติ อิมินา เอเกน ฐาเนน ปาสํโส, “น สํวิภชติ น ปุญฺญานิ กโรตี”ติ อิมินา ทุติเยน ฐาเนน คารโยฺห.\csromanspacer{} อยํ คหปติ กามโภคี อิเมหิ ทฺวีหิ ฐาเนหิ คารโยฺห, อิมินา เอเกน ฐาเนน ปาสํโส. (2)}

\prose{ตตฺร คหปติ ยฺวายํ กามโภคี อธมฺเมน โภเค ปริเยสติ สาหเสน, อธมฺเมน โภเค ปริเยสิตฺวา สาหเสน อตฺตานํ สุเขติ ปีเณติ สํวิภชติ ปุญฺญานิ กโรติ.\csromanspacer{} อยํ คหปติ กามโภคี เอเกน ฐาเนน คารโยฺห, ทฺวีหิ ฐาเนหิ ปาสํโส. “อธมฺเมน โภเค ปริเยสติ สาหเสนา”ติ อิมินา เอเกน ฐาเนน คารโยฺห, “อตฺตานํ}

\vspace{\tipitakaparskip}
\csromancenter{(10) 5.\csromanspacer{} อุปาลิวคฺค สุเขติ ปีเณตี”ติ อิมินา ปฐเมน ฐาเนน ปาสํโส, “สํวิภชติ ปุญฺญานิ กโรตี”ติ อิมินา ทุติเยน ฐาเนน ปาสํโส.\csromanspacer{} อยํ คหปติ กามโภคี อิมินา เอเกน ฐาเนน คารโยฺห, อิเมหิ ทฺวีหิ ฐาเนหิ ปาสํโส. (3)}
\prose{\csromanfolio{401}ตตฺร คหปติ ยฺวายํ กามโภคี ธมฺมาธมฺเมน โภเค ปริเยสติ สาหเสน อสาหเสนปิ, ธมฺมาธมฺเมน โภเค ปริเยสิตฺวา สาหเสนปิ อสาหเสนปิ น อตฺตานํ สุเขติ น ปีเณติ น สํวิภชติ น ปุญฺญานิ กโรติ.\csromanspacer{} อยํ คหปติ กามโภคี เอเกน ฐาเนน ปาสํโส, ตีหิ ฐาเนหิ คารโยฺห. “ธมฺเมน โภเค ปริเยสติ อสาหเสนา”ติ อิมินา เอเกน ฐาเนน ปาสํโส. “อธมฺเมน โภเค ปริเยสติ สาหเสนา”ติ อิมินา ปฐเมน ฐาเนน คารโยฺห, “น อตฺตานํ สุเขติ น ปีเณตี”ติ อิมินา ทุติเยน ฐาเนน คารโยฺห, “น สํวิภชติ น ปุญฺญานิ กโรตี”ติ อิมินา ตติเยน ฐาเนน คารโยฺห.\csromanspacer{} อยํ คหปติ กามโภคี อิมินา เอเกน ฐาเนน ปาสํโส, อิเมหิ ตีหิ ฐาเนหิ คารโยฺห. (4)}

\prose{ตตฺร คหปติ ยฺวายํ กามโภคี ธมฺมาธมฺเมน โภเค ปริเยสติ สาหเสนปิ อสาหเสนปิ, ธมฺมาธมฺเมน โภเค ปริเยสิตฺวา สาหเสนปิ อสาหเสนปิ อตฺตานํ สุเขติ ปีเณติ น สํวิภชติ น ปุญฺญานิ กโรติ.\csromanspacer{} อยํ คหปติ กามโภคี ทฺวีหิ ฐาเนหิ ปาสํโส, ทฺวีหิ ฐาเนหิ คารโยฺห. “ธมฺเมน โภเค ปริเยสติ อสาหเสนา”ติ อิมินา ปฐเมน ฐาเนน ปาสํโส, “อธมฺเมน โภเค ปริเยสติ สาหเสนา”ติ อิมินา ปฐเมน ฐาเนน คารโยฺห, “อตฺตานํ สุเขติ ปีเณตี”ติ อิมินา ทุติเยน ฐาเนน ปาสํโส, “น สํวิภชติ น ปุญฺญานิ กโรตี”ติ อิมินา ทุติเยน ฐาเนน คารโยฺห.\csromanspacer{} อยํ คหปติ กามโภคี อิเมหิ ทฺวีหิ ฐาเนหิ ปาสํโส, อิเมหิ ทฺวีหิ ฐาเนหิ คารโยฺห. (5)}

\prose{ตตฺร คหปติ ยฺวายํ กามโภคี ธมฺมาธมฺเมน โภเค ปริเยสติ สาหเสนปิ อสาหเสนปิ, ธมฺมาธมฺเมน โภเค ปริเยสิตฺวา สาหเสนปิ อสาหเสนปิ อตฺตานํ สุเขติ ปีเณติ สํวิภชติ ปุญฺญานิ กโรติ.\csromanspacer{} อยํ คหปติ กามโภคี ตีหิ ฐาเนหิ ปาสํโส,}

\vspace{\tipitakaparskip}
\csromancenter{ทสกนิปาตปาฬิ เอเกน ฐาเนน คารโยฺห. “ธมฺเมน โภเค ปริเยสติ อสาหเสนา”ติ อิมินา ปฐเมน ฐาเนน ปาสํโส, “อธมฺเมน โภเค ปริเยสติ สาหเสนา”ติ อิมินา เอเกน ฐาเนน คารโยฺห. “อตฺตานํ สุเขติ ปีเณตี”ติ อิมินา ทุติเยน ฐาเนน ปาสํโส. “สํวิภชติ ปุญฺญานิ กโรตี”ติ อิมินา ตติเยน ฐาเนน ปาสํโส, อยํ คหปติ กามโภคี อิเมหิ ตีหิ ฐาเนหิ ปาสํโส, อิมินา เอเกน ฐาเนน คารโยฺห. (6)}
\prose{\csromanfolio{402}ตตฺร คหปติ ยฺวายํ กามโภคี ธมฺเมน โภเค ปริเยสติ อสาหเสน, ธมฺเมน โภเค ปริเยสิตฺวา อสาหเสน น อตฺตานํ สุเขติ น ปีเณติ น สํวิภชติ น ปุญฺญานิ กโรติ.\csromanspacer{} อยํ คหปติ กามโภคี เอเกน ฐาเนน ปาสํโส, ทฺวีหิ ฐาเนหิ คารโยฺห. “ธมฺเมน โภเค ปริเยสติ อสาหเสนา”ติ อิมินา เอเกน ฐาเนน ปาสํโส, “น อตฺตานํ สุเขติ น ปีเณตี”ติ อิมินา ปฐเมน ฐาเนน คารโยฺห, “น สํวิภชติ น ปุญฺญานิ กโรตี”ติ อิมินา ทุติเยน ฐาเนน คารโยฺห.\csromanspacer{} อยํ คหปติ กามโภคี อิมินา เอเกน ฐาเนน ปาสโส, อิเมหิ ทฺวีหิ ฐาเนหิ คารโยฺห. (7)}

\prose{ตตฺร คหปติ ยฺวายํ กามโภคี ธมฺเมน โภเค ปริเยสติ อสาหเสน, ธมฺเมน โภเค ปริเยสิตฺวา อสาหเสน อตฺตานํ สุเขติ ปีเณติ น สํวิภชติ น ปุญฺญานิ กโรติ.\csromanspacer{} อยํ คหปติ กามโภคี ทฺวีหิ ฐาเนหิ ปาสํโส, เอเกน ฐาเนน คารโยฺห. “ธมฺเมน โภเค ปริเยสติ อสาหเสนา”ติ อิมินา ปฐเมน ฐาเนน ปาสํโส, “อตฺตานํ สุเขติ ปีเณตี”ติ อิมินา ทุติเยน ฐาเนน ปาสํโส, “น สํวิภชติ น ปุญฺญานิ กโรตี”ติ อิมินา เอเกน ฐาเนน คารโยฺห.\csromanspacer{} อยํ คหปติ กามโภคี อิเมหิ ทฺวีหิ ฐาเนหิ ปาสํโส, อิมินา เอเกน ฐาเนน คารโยฺห. (8)}

\prose{ตตฺร คหปติ ยฺวายํ กามโภคี ธมฺเมน โภเค ปริเยสติ อสาหเสน, ธมฺเมน โภเค ปริเยสิตฺวา อสาหเสน อตฺตานํ สุเขติ ปีเณติ สํวิภชติ ปุญฺญานิ กโรติ, เต จ โภเค คถิโต มุจฺฉิโต อชฺโฌสนฺโน อนาทีนวทสฺสาวี อนิสฺสรณปญฺโญ ปริภุญฺชติ.\csromanspacer{} อยํ คหปติ กามโภคี ตีหิ ฐาเนหิ ปาสํโส, เอเกน ฐาเนน คารโยฺห. “ธมฺเมน โภเค ปริเยสติ อสาหเสนา”ติ อิมินา}

\vspace{\tipitakaparskip}
\csromancenter{(10) 5.\csromanspacer{} อุปาลิวคฺค ปฐเมน ฐาเนน ปาสํโส, “อตฺตานํ สุเขติ ปีเณตี”ติ อิมินา ทุติเยน ฐาเนน ปาสํโส, “สํวิภชติ ปุญฺญานิ กโรตี”ติ อิมินา ตติเยน ฐาเนน ปาสํโส, “เต จ โภเค คถิโต มุจฺฉิโต อชฺโฌสนฺโน อนาทีนวทสฺสาวี อนิสฺสรณปญฺโญ ปริภุญฺชตี”ติ อิมินา เอเกน ฐาเนน คารโยฺห.\csromanspacer{} อยํ คหปติ กามโภคี อิเมหิ ตีหิ ฐาเนหิ ปาสํโส.\csromanspacer{} อิมินา เอเกน ฐาเนน คารโยฺห. (9)}
\prose{\csromanfolio{403}ตตฺร คหปติ ยฺวายํ กามโภคี ธมฺเมน โภเค ปริเยสติ อสาหเสน, ธมฺเมน โภเค ปริเยสิตฺวา อสาหเสน อตฺตานํ สุเขติ ปีเณติ สํวิภชติ ปุญฺญานิ กโรติ, เต จ โภเค อคถิโต อมุจฺฉิโต อนชฺโฌสนฺโน อาทีนวทสฺสาวี นิสฺสรณปญฺโญ ปริภุญฺชติ.\csromanspacer{} อยํ คหปติ กามโภคี จตูหิ ฐาเนหิ ปาสํโส. “ธมฺเมน โภเค ปริเยสติ อสาหเสนา”ติ อิมินา ปฐเมน ฐาเนน ปาสํโส, “อตฺตานํ สุเขติ ปีเณตี”ติ อิมินา ทุติเยน ฐาเนน ปาสํโส, “สํวิภชติ ปุญฺญานิ กโรตี”ติ อิมินา ตติเยน ฐาเนน ปาสํโส, “เต จ โภเค อคถิโต อมุจฺฉิโต อนชฺโฌสนฺโน อาทีนวทสฺสาวี นิสฺสรณปญฺโญ ปริภุญฺชตี”ติ อิมินา จตุตฺเถน ฐาเนน ปาสํโส.\csromanspacer{} อยํ คหปติ กามโภคี อิเมหิ จตูหิ ฐาเนหิ ปาสํโส. (10)}

\prose{อิเม โข คหปติ ทส กามโภคี สนฺโต สํวิชฺชมานา โลกสฺมิํ.\csromanspacer{} อิเมสํ โข คหปติ ทสนฺนํ กามโภคีนํ ยฺวายํ กามโภคี ธมฺเมน โภเค ปริเยสติ อสาหเสน, ธมฺเมน โภเค ปริเยสิตฺวา อสาหเสน อตฺตานํ สุเขติ ปีเณติ สํวิภชติ ปุญฺญานิ กโรติ, เต จ โภเค อคถิโต อมุจฺฉิโต อนชฺโฌสนฺโน อาทีนวทสฺสาวี นิสฺสรณปญฺโญ ปริภุญฺชติ.\csromanspacer{} อยํ อิเมสํ ทสนฺนํ กามโภคีนํ อคฺโค จ เสฏฺโฐ จ ปาโมกฺโข\footnote{โมกฺโข (ก-สี) อํ 1. 408; อํ 2. 193; สํ 2. 224 ปิฏฺเฐสุปิ.} จ อุตฺตโม จ ปวโร จ.\csromanspacer{} เสยฺยถาปิ คหปติ ควา ขีรํ, ขีรมฺหา ทธิ, ทธิมฺหา นวนีตํ, นวนีตมฺหา สปฺปิ, สปฺปิมฺหา สปฺปิมณฺโฑ, สปฺปิมณฺโฑ ตตฺถ อคฺคมกฺขายติ.}

\prose{เอวเมวํ โข คหปติ อิเมสํ ทสนฺนํ กามโภคีนํ ยฺวายํ กามโภคี ธมฺเมน โภเค ปริเยสติ อสาหเสน, ธมฺเมน โภเค ปริเยสิตฺวา อสาหเสน อตฺตานํ สุเขติ ปีเณติ}

\vspace{\tipitakaparskip}
\csromancenter{ทสกนิปาตปาฬิ สํวิภชติ ปุญฺญานิ กโรติ, เต จ โภเค อคถิโต อมุจฺฉิโต อนชฺโฌสนฺโน อาทีนวทสฺสาวี นิสฺสรณปญฺโญ ปริภุญฺชติ.\csromanspacer{} อยํ อิเมสํ ทสนฺนํ กามโภคีนํ อคฺโค จ เสฏฺโฐ จ ปาโมกฺโข\footnote{โมกฺโข (ก-สี) อํ 2. 193 ปิฏฺเฐปิ.} จ อุตฺตโม จ ปวโร จาติ.\csromanspacer{}\csromanspacer{}ปฐมํ.}
\csromanheader{2}{2. ภยสุตฺต}
\proseitem{92}{\csromanfolio{404}\csromansymbolfootnote{*}{เหฏฺฐา 205; สํ 3. 338 ปิฏฺเฐสุปิ.}อถ โข อนาถปิณฺฑิโก คหปติ เยน ภควา เตนุปสงฺกมิ, อุปสงฺกมิตฺวา ภควนฺตํ อภิวาเทตฺวา เอกมนฺตํ นิสีทิ.\csromanspacer{} เอกมนฺตํ นิสินฺนํ โข อนาถปิณฺฑิกํ คหปติํ ภควา เอตทโวจ\csromandash{}}

\prose{ยโต โข คหปติ อริยสาวกสฺส ปญฺจ ภยานิ เวรานิ วูปสนฺตานิ โหนฺติ, จตูหิ จ โสตาปตฺติยงฺเคหิ สมนฺนาคโต โหติ, อริโย จสฺส ญาโย ปญฺญาย สุทิฏฺโฐ โหติ สุปฺปฏิวิทฺโธ.\csromanspacer{} โส อากงฺขมาโน อตฺตนาว อตฺตานํ พฺยากเรยฺย “ขีณนิรโยมฺหิ ขีณติรจฺฉานโยนิ ขีณเปตฺติวิสโย ขีณาปายทุคฺคติวินิปาโต โสตาปนฺโนหมสฺมิ อวินิปาตธมฺโม นิยโต สมฺโพธิปรายโณ”ติ.}

\prose{กตมานิ ปญฺจ ภยานิ เวรานิ วูปสนฺตานิ โหนฺติ? ยํ คหปติ ปาณาติปาตี ปาณาติปาตปจฺจยา ทิฏฺฐธมฺมิกมฺปิ ภยํ เวรํ ปสวติ, สมฺปรายิกมฺปิ ภยํ เวรํ ปสวติ, เจตสิกมฺปิ ทุกฺขํ โทมนสฺสํ ปฏิสํเวเทติ.\csromanspacer{} ปาณาติปาตา ปฏิวิรโต เนว ทิฏฺฐธมฺมิกมฺปิ\footnote{เนว ทิฏฺฐธมฺมิกํ.} ภยํ เวรํ ปสวติ, น สมฺปรายิกมฺปิ\footnote{น สมฺปรายิกํ.} ภยํ เวรํ ปสวติ, น เจตสิกมฺปิ\footnote{น เจตสิกํ (สี, สฺยา, อิ)} ทุกฺขํ โทมนสฺสํ ปฏิสํเวเทติ.\csromanspacer{} ปาณาติปาตา ปฏิวิรตสฺส เอวํ ตํ ภยํ เวรํ วูปสนฺตํ โหติ.}

\prose{ยํ คหปติ อทินฺนาทายี ฯเปฯ กาเมสุมิจฺฉาจารี.\csromanspacer{} มุสาวาที.\csromanspacer{} สุราเมรยมชฺชปมาทฏฺฐายี สุราเมรยมชฺชปมาทฏฺฐานปจฺจยา ทิฏฺฐธมฺมิกมฺปิ ภยํ เวรํ ปสวติ, สมฺปรายิกมฺปิ ภยํ เวรํ ปสวติ, เจตสิกมฺปิ ทุกฺขํ โทมนสฺสํ ปฏิสํเวเทติ.\csromanspacer{} สุราเมรยมชฺชปมาทฏฺฐานา ปฏิวิรโต เนว ทิฏฺฐธมฺมิกมฺปิ ภยํ เวรํ ปสวติ, น สมฺปรายิกมฺปิ ภยํ เวรํ ปสวติ, น เจตสิกมฺปิ ทุกฺขํ โทมนสฺสํ ปฏิสํเวเทติ.\csromanspacer{} สุราเมรยมชฺชปมาทฏฺฐานา ปฏิวิรตสฺส}

\vspace{\tipitakaparskip}
\csromancenter{(10) 5.\csromanspacer{} อุปาลิวคฺค เอวํ ตํ ภยํ เวรํ วูปสนฺตํ โหติ.\csromanspacer{} อิมานิ ปญฺจ ภยานิ เวรานิ วูปสนฺตานิ โหนฺติ.}
\prose{\csromanfolio{405}กตเมหิ จตูหิ โสตาปตฺติยงฺเคหิ สมนฺนาคโต โหติ? อิธ คหปติ อริยสาวโก พุทฺเธ อเวจฺจปฺปสาเทน สมนฺนาคโต โหติ “อิติปิ โส ภควา ฯเปฯ พุทฺโธ ภควา”ติ.\csromanspacer{} ธมฺเม อเวจฺจปฺปสาเทน สมนฺนาคโต โหติ “สฺวากฺขาโต ภควตา ธมฺโม สนฺทิฏฺฐิโก อกาลิโก เอหิปสฺสิโก โอปเนยฺยิโก ปจฺจตฺตํ เวทิตพฺโพ วิญฺญูหี”ติ.\csromanspacer{} สํเฆ อเวจฺจปฺปสาเทน สมนฺนาคโต โหติ “สุปฺปฏิปนฺโน ภควโต สาวกสํโฆ อุชุปฺปฏิปนฺโน ภควโต สาวกสํโฆ ญายปฺปฏิปนฺโน ภควโต สาวกสํโฆ สามีจิปฺปฏิปนฺโน ภควโต สาวกสํโฆ, ยทิทํ จตฺตาริ ปุริสยุคานิ อฏฺฐ ปุริสปุคฺคลา, เอส ภควโต สาวกสํโฆ อาหุเนยฺโย ปาหุเนยฺโย ทกฺขิเณยฺโย อญฺชลิกรณีโย อนุตฺตรํ ปุญฺญกฺเขตฺตํ โลกสฺสา”ติ.\csromanspacer{} อริยกนฺเตหิ สีเลหิ สมนฺนาคโต โหติ อขณฺเฑหิ อจฺฉิทฺเทหิ อสพเลหิ อกมฺมาเสหิ ภุชิสฺเสหิ วิญฺญุปฺปสตฺเถหิ อปรามฏฺเฐหิ สมาธิสํวตฺตนิเกหิ.\csromanspacer{} อิเมหิ จตูหิ โสตาปตฺติยงฺเคหิ สมนฺนาคโต โหติ.}

\prose{กตโม จสฺส อริโย ญาโย ปญฺญาย สุทิฏฺโฐ โหติ สุปฺปฏิวิทฺโธ? อิธ คหปติ อริยสาวโก อิติ ปฏิสญฺจิกฺขติ “อิติ อิมสฺมิํ สติ อิทํ โหติ, อิมสฺสุปฺปาทา อิทํ อุปฺปชฺชติ.\csromanspacer{} อิมสฺมิํ อสติ อิทํ น โหติ, อิมสฺส นิโรธา อิทํ นิรุชฺฌติ\csromandash{}ยทิทํ อวิชฺชาปจฺจยา สงฺขารา, สงฺขารปจฺจยา วิญฺญาณํ, วิญฺญาณปจฺจยา นามรูปํ, นามรูปปจฺจยา สฬายตนํ, สฬายตนปจฺจยา ผสฺโส, ผสฺสปจฺจยา เวทนา, เวทนาปจฺจยา ตณฺหา, ตณฺหาปจฺจยา อุปาทานํ, อุปาทานปจฺจยา ภโว, ภวปจฺจยา ชาติ, ชาติปจฺจยา ชรา มรณํ โสก ปริเทว ทุกฺข โทมนสฺสุปายาสา สมฺภวนฺติ, เอวเมตสฺส เกวลสฺส ทุกฺขกฺขนฺธสฺส สมุทโย โหติ.\csromanspacer{} อวิชฺชาย เตฺวว อเสสวิราคนิโรธา สงฺขารนิโรโธ ฯเปฯ เอวเมตสฺส เกวลสฺส ทุกฺขกฺขนฺธสฺส นิโรโธ โหตี”ติ.\csromanspacer{} อยญฺจสฺส อริโย ญาโย ปญฺญาย สุทิฏฺโฐ โหติ สุปฺปฏิวิทฺโธ.}

\prose{ยโต โข คหปติ อริยสาวกสฺส อิมานิ ปญฺจ ภยานิ เวรานิ วูปสนฺตานิ โหนฺติ, อิเมหิ จ จตูหิ โสตาปตฺติยงฺเคหิ สมนฺนาคโต โหติ, อยญฺจสฺส อริโย ญาโย ปญฺญาย สุทิฏฺโฐ โหติ}

\vspace{\tipitakaparskip}
\csromancenter{ทสกนิปาตปาฬิ สุปฺปฏิวิทฺโธ.\csromanspacer{} โส อากงฺขมาโน อตฺตนาว อตฺตานํ พฺยากเรยฺย “ขีณนิรโยมฺหิ ขีณติรจฺฉานโยนิ ขีณเปตฺติวิสโย ขีณาปายทุคฺคติวินิปาโต โสตาปนฺโนหมสฺมิ อวินิปาตธมฺโม นิยโต สมฺโพธิปรายโณ”ติ.\csromanspacer{}\csromanspacer{}ทุติยํ.}
\csromanheader{2}{3. กิํทิฏฺฐิกสุตฺต}
\csromanmatikaanchor{mtk.133}
\csromantocmarkref{subsection}{2. ภยสุตฺต}{mtk.133}
\bookmark[dest={mtk.133},level=2]{2. ภยสุตฺต}
\proseitem{93}{\csromanfolio{406}เอกํ สมยํ ภควา สาวตฺถิยํ วิหรติ เชตวเน อนาถปิณฺฑิกสฺส อาราเม.\csromanspacer{} อถ โข อนาถปิณฺฑิโก คหปติ ทิวา ทิวสฺส สาวตฺถิยา นิกฺขมิ ภควนฺตํ ทสฺสนาย.\csromanspacer{} อถ โข อนาถปิณฺฑิกสฺส คหปติสฺส เอตทโหสิ “อกาโล โข ตาว ภควนฺตํ ทสฺสนาย, ปฏิสลฺลีโน ภควา, มโนภาวนียานมฺปิ ภิกฺขูนํ อกาโล ทสฺสนาย, ปฏิสลฺลีนา มโนภาวนียา ภิกฺขู, ยํนูนาหํ เยน อญฺญติตฺถิยานํ ปริพฺพาชกานํ อาราโม, เตนุปสงฺกเมยฺยนฺ”ติ. อถ โข อนาถปิณฺฑิโก คหปติ เยน อญฺญติตฺถิยานํ ปริพฺพาชกานํ อาราโม เตนุปสงฺกมิ.\csromanspacer{} เตน โข ปน สมเยน อญฺญติตฺถิยา ปริพฺพาชกา สงฺคมฺม สมาคมฺม อุนฺนาทิโน อุจฺจาสทฺทมหาสทฺทา อเนกวิหิตํ ติรจฺฉานกถํ กเถนฺตา นิสินฺนา โหนฺติ.\csromanspacer{} อทฺทสํสุ โข เต อญฺญติตฺถิยา ปริพฺพาชกา อนาถปิณฺฑิกํ คหปติํ ทูรโตว อาคจฺฉนฺตํ, ทิสฺวาน อญฺญมญฺญํ สณฺฐาเปสุํ “อปฺปสทฺทา โภนฺโต โหนฺตุ, มา โภนฺโต สทฺทมกตฺถ.\csromanspacer{} อยํ อนาถปิณฺฑิโก คหปติ อารามํ อาคจฺฉติ สมณสฺส โคตมสฺส สาวโก, ยาวตา โข ปน สมณสฺส โคตมสฺส สาวกา คิหี โอทาตวสนา สาวตฺถิยํ ปฏิวสนฺติ, อยํ เตสํ อญฺญตโร อนาถปิณฺฑิโก คหปติ.\csromanspacer{} อปฺปสทฺทกาเม โข ปน เต อายสฺมนฺโต อปฺปสทฺทวินีตา อปฺปสทฺทสฺส วณฺณวาทิโน, อปฺเปว นาม อปฺปสทฺทํ ปริสํ วิทิตฺวา อุปสงฺกมิตพฺพํ มญฺเญยฺยา”ติ.}

\prose{อถ โข เต อญฺญติตฺถิยา ปริพฺพาชกา ตุณฺหี อเหสุํ, อถ โข อนาถปิณฺฑิโก คหปติ เยน เต อญฺญติตฺถิยา ปริพฺพาชกา เตนุปสงฺกมิ, อุปสงฺกมิตฺวา เตหิ อญฺญติตฺถิเยหิ ปริพฺพาชเกหิ สทฺธิํ สมฺโมทิ, สมฺโมทนียํ กถํ สารณียํ วีติสาเรตฺวา เอกมนฺตํ นิสีทิ, เอกมนฺตํ นิสินฺนํ โข อนาถปิณฺฑิกํ คหปติํ เต อญฺญติตฺถิยา}

\vspace{\tipitakaparskip}
\csromancenter{(10) 5.\csromanspacer{} อุปาลิวคฺค ปริพฺพาชกา เอตทโวจุํ “วเทหิ คหปติ, กิํทิฏฺฐิโก สมโณ โคตโม”ติ.\csromanspacer{} น โข อหํ ภนฺเต ภควโต สพฺพํ ทิฏฺฐิํ ชานามีติ.}
\prose{\csromanfolio{407}อิติ กิร ตฺวํ คหปติ น สมณสฺส โคตมสฺส สพฺพํ ทิฏฺฐิํ ชานาสิ วเทหิ คหปติ กิํทิฏฺฐิกา ภิกฺขูติ.\csromanspacer{} ภิกฺขูนมฺปิ โข อหํ ภนฺเต น สพฺพํ ทิฏฺฐิํ ชานามีติ.}

\prose{อิติ กิร ตฺวํ คหปติ น สมณสฺส โคตมสฺส สพฺพํ ทิฏฺฐิํ ชานาสิ, นปิ ภิกฺขูนํ สพฺพํ ทิฏฺฐิํ ชานาสิ, วเทหิ คหปติ กิํทิฏฺฐิโกสิ ตุวนฺติ.\csromanspacer{} เอตํ โข ภนฺเต อเมฺหหิ น ทุกฺกรํ พฺยากาตุํ, ยํทิฏฺฐิกา มยํ.\csromanspacer{} อิงฺฆ ตาว อายสฺมนฺโต ยถาสกานิ ทิฏฺฐิคตานิ พฺยากโรนฺตุ, ปจฺฉาเปตํ อเมฺหหิ น ทุกฺกรํ ภวิสฺสติ พฺยากาตุํ, ยํทิฏฺฐิกา มยนฺติ.}

\prose{เอวํ วุตฺเต อญฺญตโร ปริพฺพาชโก อนาถปิณฺฑิกํ คหปติํ เอตทโวจ “สสฺสโต โลโก, อิทเมว สจฺจํ โมฆมญฺญนฺติ เอวํทิฏฺฐิโก อหํ คหปตี”ติ.}

\prose{อญฺญตโรปิ โข ปริพฺพาชโก อนาถปิณฺฑิกํ คหปติํ เอตทโวจ “อสสฺสโต โลโก, อิทเมว สจฺจํ โมฆมญฺญนฺติ เอวํทิฏฺฐิโก อหํ คหปตี”ติ.}

\prose{อญฺญตโรปิ โข ปริพฺพาชโก อนาถปิณฺฑิกํ คหปติํ เอตทโวจ “อนฺตวา โลโก ฯเปฯ อนนฺตวา โลโก.\csromanspacer{} ตํ ชีวํ ตํ สรีรํ.\csromanspacer{} อญฺญํ ชีวํ อญฺญํ สรีรํ.\csromanspacer{} โหติ ตถาคโต ปรํ มรณา.\csromanspacer{} น โหติ ตถาคโต ปรํ มรณา.\csromanspacer{} โหติ จ น โหติ ตถาคโต ปรํ มรณา.\csromanspacer{} เนว โหติ น น โหติ ตถาคโต ปรํ มรณา, อิทเมว สจฺจํ โมฆมญฺญนฺติ เอวํทิฏฺฐิโก อหํ คหปตี”ติ.}

\prose{เอวํ วุตฺเต อนาถปิณฺฑิโก คหปติ เต ปริพฺพาชเก เอตทโวจ\csromandash{}ยฺวายํ ภนฺเต อายสฺมา เอวมาห “สสฺสโต โลโก อิทเมว สจฺจํ โมฆมญฺญนฺติ เอวํทิฏฺฐิโก อหํ คหปตี”ติ, อิมสฺส อยมายสฺมโต ทิฏฺฐิ อตฺตโน วา อโยนิโสมนสิการเหตุ อุปฺปนฺนา ปรโตโฆสปจฺจยา วา.\csromanspacer{} สา โข ปเนสา ทิฏฺฐิ ภูตา สงฺขตา เจตยิตา ปฏิจฺจสมุปฺปนฺนา, ยํ โข ปน กิญฺจิ ภูตํ สงฺขตํ เจตยิตํ ปฏิจฺจสมุปฺปนฺนํ ตทนิจฺจํ, ยทนิจฺจํ ตํ ทุกฺขํ, ยํ ทุกฺขํ ตเทเวโส อายสฺมา อลฺลีโน ตเทเวโส อายสฺมา อชฺฌุปคโต.}

\gambhira{ทสกนิปาตปาฬิ}
\csromanmatikaanchor{mtk.134}
\csromantocmarkref{subsection}{3. กิํทิฏฺฐิกสุตฺต}{mtk.134}
\bookmark[dest={mtk.134},level=2]{3. กิํทิฏฺฐิกสุตฺต}
\prose{\csromanfolio{408}โยปายํ ภนฺเต อายสฺมา เอวมาห “อสสฺสโต โลโก, อิทเมว สจฺจํ โมฆมญฺญนฺติ เอวํทิฏฺฐิโก อหํ คหปตี”ติ, อิมสฺสาปิ อยมายสฺมโต ทิฏฺฐิ อตฺตโน วา อโยนิโสมนสิการเหตุ อุปฺปนฺนา ปรโตโฆสปจฺจยา วา.\csromanspacer{} สา โข ปเนสา ทิฏฺฐิ ภูตา สงฺขตา เจตยิตา ปฏิจฺจสมุปฺปนฺนา, ยํ โข ปน กิญฺจิ ภูตํ สงฺขตํ เจตยิตํ ปฏิจฺจสมุปฺปนฺนํ ตทนิจฺจํ, ยทนิจฺจํ ตํ ทุกฺขํ, ยํ ทุกฺขํ ตเทเวโส อายสฺมา อลฺลีโน ตเทเวโส อายสฺมา อชฺฌุปคโต.}

\prose{โยปายํ ภนฺเต อายสฺมา เอวมาห “อนฺตวา โลโก ฯเปฯ อนนฺตวา โลโก.\csromanspacer{} ตํ ชีวํ ตํ สรีรํ.\csromanspacer{} อญฺญํ ชีวํ อญฺญํ สรีนํ.\csromanspacer{} โหติ ตถาคโต ปรํ มรณา.\csromanspacer{} น โหติ ตถาคโต ปรํ มรณา.\csromanspacer{} โหติ จ น จ โหติ ตถาคโต ปรํ มรณา.\csromanspacer{} เนว โหติ น น โหติ ตถาคโต ปรํ มรณา, อิทเมว สจฺจํ โมฆมญฺญนฺติ เอวํทิฏฺฐิโก อหํ คหปตี”ติ.\csromanspacer{} อิมสฺสาปิ อยมายสฺมโต ทิฏฺฐิ อตฺตโน วา อโยนิโสมนสิการเหตุ อุปฺปนฺนา ปรโตโฆสปจฺจยา วา.\csromanspacer{} สา โข ปเนสา ทิฏฺฐิ ภูตา สงฺขตา เจตยิตา ปฏิจฺจสมุปฺปนฺนา, ยํ โข ปน กิญฺจิ ภูตํ สงฺขตํ เจตยิตํ ปฏิจฺจสมุปฺปนฺนํ ตทนิจฺจํ, ยทนิจฺจํ ตํ ทุกฺขํ, ยํ ทุกฺขํ ตเทเวโส อายสฺมา อลฺลีโน ตเทเวโส อายสฺมา อชฺฌุปคโตติ.}

\prose{เอวํ วุตฺเต เต ปริพฺพาชกา อนาถปิณฺฑิกํ คหปติํ เอตทโวจุํ “พฺยากตานิ โข คหปติ อเมฺหหิ สพฺเพเหว ยถาสกา ทิฏฺฐิคตานิ, วเทหิ คหปติ กิํทิฏฺฐิโกสิ ตุวนฺ”ติ.\csromanspacer{} ยํ โข ภนฺเต กิญฺจิ ภูตํ สงฺขตํ เจตยิตํ ปฏิจฺจสมุปฺปนฺนํ ตทนิจฺจํ, ยทนิจฺจํ ตํ ทุกฺขํ, ยํ ทุกฺขํ ตํ “เนตํ มม, เนโสหมสฺมิ, น เมโส อตฺตา”ติ เอวํทิฏฺฐิโก อหํ ภนฺเตติ.}

\prose{ยํ โข คหปติ กิญฺจิ ภูตํ สงฺขตํ เจตยิตํ ปฏิจฺจสมุปฺปนฺนํ ตทนิจฺจํ, ยทนิจฺจํ ตํ ทุกฺขํ ตเทว ตฺวํ คหปติ อลฺลีโน ตเทว ตฺวํ คหปติ อชฺฌุปคโตติ.}

\prose{ยํ โข ภนฺเต กิญฺจิ ภูตํ สงฺขตํ เจตยิตํ ปฏิจฺจสมุปฺปนฺนํ ตทนิจฺจํ, ยทนิจฺจํ ตํ ทุกฺขํ, ยํ ทุกฺขํ ตํ “เนตํ มม, เนโสหมสฺมิ, นเมโส อตฺตา”ติ เอวเมตํ ยถาภูตํ สมฺมปฺปญฺญาย สุทิฏฺฐํ, ตสฺส จ อุตฺตริ นิสฺสรณํ ยถาภูตํ ปชานามีติ.}

\chapterheadpageread{(10) 5. อุปาลิวคฺค}
\prose{\csromanfolio{409}เอวํ วุตฺเต เต ปริพฺพาชกา ตุณฺหีภูตา มงฺกุภูตา ปตฺตกฺขนฺธา อโธมุขา ปชฺฌายนฺตา อปฺปฏิภานา นิสีทิํสุ.\csromanspacer{} อถ โข อนาถปิณฺฑิโก คหปติ เต ปริพฺพาชเก ตุณฺหีภูเต มงฺกุภูเต ปตฺตกฺขนฺเธ อโธมุเข ปชฺฌายนฺเต อปฺปฏิภาเน วิทิตฺวา อุฏฺฐายาสนา เยน ภควา เตนุปสงฺกมิ, อุปสงฺกมิตฺวา ภควนฺตํ อภิวาเทตฺวา เอกมนฺตํ นิสีทิ, เอกมนฺตํ นิสินฺโน โข อนาถปิณฺฑิโก คหปติ ยาวตโก อโหสิ เตหิ อญฺญติตฺถิเยหิ ปริพฺพาชเกหิ สทฺธิํ กถาสลฺลาโป, ตํ สพฺพํ ภควโต อาโรเจสิ.\csromanspacer{} สาธุ สาธุ คหปติ, เอวํ โข เต คหปติ โมฆปุริสา กาเลน กาลํ สหธมฺเมน สุนิคฺคหิตํ นิคฺคเหตพฺพาติ.}

\prose{อถ โข ภควา อนาถปิณฺฑิกํ คหปติํ ธมฺมิยา กถาย สนฺทสฺเสสิ สมาทเปสิ สมุตฺเตเชสิ สมฺปหํเสสิ.\csromanspacer{} อถ โข อนาถปิณฺฑิโก คหปติ ภควตา ธมฺมิยา กถาย สนฺทสฺสิโต สมาทปิโต สมุตฺเตชิโต สมฺปหํสิโต อุฏฺฐายาสนา ภควนฺตํ อภิวาเทตฺวา ปทกฺขิณํ กตฺวา ปกฺกามิ.}

\prose{อถ โข ภควา อจิรปกฺกนฺเต อนาถปิณฺฑิเก คหปติมฺหิ ภิกฺขู อามนฺเตสิ “โยปิ โส ภิกฺขเว ภิกฺขุ วสฺสสตุปสมฺปนฺโน\footnote{ภิกฺขุ ทีฆรตฺตํ อเวธิ ธมฺโม (สฺยา)} อิมสฺมิํ ธมฺมวินเย, โสปิ เอวเมวํ อญฺญติตฺถิเย ปริพฺพาชเก สหธมฺเมน สุนิคฺคหิตํ นิคฺคเณฺหยฺย, ยถา ตํ อนาถปิณฺฑิเกน คหปตินา นิคฺคหิตา”ติ.\csromanspacer{}\csromanspacer{}ตติยํ.}

\csromanheader{2}{4. วชฺชิยมาหิตสุตฺต}
\proseitem{94}{เอกํ สมยํ ภควา จมฺปายํ วิหรติ คคฺคราย โปกฺขรณิยา ตีเร.\csromanspacer{} อถ โข วชฺชิยมาหิโต คหปติ ทิวา ทิวสฺส จมฺปาย นิกฺขมิ ภควนฺตํ ทสฺสนาย.\csromanspacer{} อถ โข วชฺชิยมาหิตสฺส คหปติสฺส เอตทโวจ “อกาโล โข ตาว ภควนฺตํ ทสฺสนาย, ปฏิสลฺลีโน ภควา, มโนภาวนียานมฺปิ ภิกฺขูนํ อกาโล ทสฺสนาย, ปฏิสลฺลีนา มโนภาวนียาปิ ภิกฺขู, ยํนูนาหํ เยน อญฺญติตฺถิยานํ ปริพฺพาชกานํ อาราโม เตนุปสงฺกเมยฺยนฺ”ติ.}

\gambhira{ทสกนิปาตปาฬิ}
\prose{\csromanfolio{410}อถ โข วชฺชิยมาหิโต คหปติ เยน อญฺญติตฺถิยานํ ปริพฺพาชกานํ อาราโม เตนุปสงฺกมิ.\csromanspacer{} เตน โข ปน สมเยน เต อญฺญติตฺถิยา ปริพฺพาชกา สงฺคมฺม สมาคมฺม อุนฺนาทิโน อุจฺจาสทฺทมหาสทฺทา อเนกวิหิตํ ติรจฺฉานกถํ กเถนฺตา นิสินฺนา โหนฺติ.}

\prose{อทฺทสํสุ โข เต อญฺญติตฺถิยา ปริพฺพาชกา วชฺชิยมาหิตํ คหปติํ ทูรโตว อาคจฺฉนฺตํ, ทิสฺวาน อญฺญมญฺญํ สณฺฐาเปสุํ “อปฺปสทฺทา โภนฺโต โหนฺตุ, มา โภนฺโต สทฺทมกตฺถ.\csromanspacer{} อยํ วชฺชิยมาหิโต คหปติ อาคจฺฉติ สมณสฺส โคตมสฺส สาวโก.\csromanspacer{} ยาวตา โข ปน สมณสฺส โคตมสฺส สาวกา คิหี โอทาตวสนา จมฺปายํ ปฏิวสนฺติ, อยํ เตสํ อญฺญตโร วชฺชิยมาหิโต คหปติ.\csromanspacer{} อปฺปสทฺทกามา โข ปน เต อายสฺมนฺโต อปฺปสทฺทวินีตา อปฺปสทฺทสฺส วณฺณวาทิโน, อปฺเปว นาม อปฺปสทฺทํ ปริสํ วิทิตฺวา อุปสงฺกมิตพฺพํ มญฺเญยฺยา”ติ.}

\prose{อถ โข เต อญฺญติตฺถิยา ปริพฺพาชกา ตุณฺหี อเหสุํ.\csromanspacer{} อถ โข วชฺชิยมาหิโต คหปติ เยน เต อญฺญติตฺถิยา ปริพฺพาชกา เตนุปสงฺกมิ, อุปสงฺกมิตฺวา เตหิ อญฺญติตฺถิเยหิ ปริพฺพาชเกหิ สทฺธิํ สมฺโมทิ, สมฺโมทนียํ กถํ สารณียํ วีติสาเรตฺวา เอกมนฺตํ นิสีทิ, เอกมนฺตํ นิสินฺนํ โข วชฺชิยมาหิตํ คหปติํ เต อญฺญติตฺถิยา ปริพฺพาชกา เอตทโวจุํ “สจฺจํ กิร คหปติ สมโณ โคตโม สพฺพํ ตปํ ครหติ, สพฺพํ ตปสฺสิํ ลูขาชีวิํ เอกํเสน อุปกฺโกสติ อุปวทตี”ติ.\csromanspacer{} น โข ภนฺเต ภควา สพฺพํ ตปํ ครหติ, นปิ สพฺพํ ตปสฺสิํ ลูขาชีวิํ เอกํเสน อุปกฺโกสติ อุปวทติ.\csromanspacer{} คารยฺหํ โข ภนฺเต ภควา ครหติ, ปสํสิตพฺพํ ปสํสติ.\csromanspacer{} คารยฺหํ โข ปน ภนฺเต ภควา ครหนฺโต ปสํสิตพฺพํ ปสํสนฺโต วิภชฺชวาโท ภควา, น โส ภควา เอตฺถ เอกํสวาโทติ.}

\prose{เอวํ วุตฺเต อญฺญตโร ปริพฺพาชโก วชฺชิยมาหิตํ คหปติํ เอตทโวจ “อาคเมหิ ตฺวํ คหปติ, ยสฺส ตฺวํ สมณสฺส โคตมสฺส วณฺณํ ภาสติ, สมโณ โคตโม เวนยิโก อปฺปญฺญตฺติโก”ติ.\csromanspacer{} เอตฺถปาหํ ภนฺเต อายสฺมนฺเต วกฺขามิ สหธมฺเมน, “อิทํ กุสลนฺ”ติ ภนฺเต ภควตา ปญฺญตฺตํ, “อิทํ อกุสลนฺ”ติ ภนฺเต ภควตา ปญฺญตฺตํ, อิติ กุสลากุสลํ ภควา ปญฺญาปยมาโน สปญฺญตฺติโก ภควา, น โส ภควา เวนยิโก อปฺปญฺญตฺติโกติ.}

\chapterheadpageread{(10) 5. อุปาลิวคฺค}
\prose{\csromanfolio{411}เอวํ วุตฺเต เต ปริพฺพาชกา ตุณฺหีภูตา มงฺกุภูตา ปตฺตกฺขนฺธา อโธมุขา ปชฺฌายนฺตา อปฺปฏิภานา นิสีทิํสุ.\csromanspacer{} อถ โข วชฺชิยมาหิโต คหปติ เต ปริพฺพาชเก ตุณฺหีภูเต มงฺกุภูเต ปตฺตกฺขนฺเธ อโธมุเข ปชฺฌายนฺเต อปฺปฏิภาเน วิทิตฺวา อุฏฺฐายาสนา เยน ภควา เตนุปสงฺกมิ, อุปสงฺกมิตฺวา ภควนฺตํ อภิวาเทตฺวา เอกมนฺตํ นิสีทิ, เอกมนฺตํ นิสินฺโน โข วชฺชิยมาหิโต คหปติ ยาวตโก อโหสิ เตหิ อญฺญติตฺถิเยหิ ปริพฺพาชเกหิ สทฺธิํ กถาสลฺลาโป, ตํ สพฺพํ ภควโต อาโรเจสิ.}

\prose{สาธุ สาธุ คหปติ, เอวํ โข เต คหปติ โมฆปุริสา กาเลน กาลํ สหธมฺเมน สุนิคฺคหิตํ นิคฺคเหตพฺพา.\csromanspacer{} นาหํ คหปติ สพฺพํ ตปํ “ตปิตพฺพนฺ”ติ วทามิ, น จ ปนาหํ คหปติ สพฺพํ ตปํ “น ตปิตพฺพนฺ”ติ วทามิ.\csromanspacer{} นาหํ คหปติ สพฺพํ สมาทานํ “สมาทิตพฺพนฺ”ติ วทามิ, น ปนาหํ คหปติ สพฺพํ สมาทานํ “น สมาทิตพฺพนฺ”ติ วทามิ.\csromanspacer{} นาหํ คหปติ สพฺพํ ปธานํ “ปทหิตพฺพนฺ”ติ วทามิ, น ปนาหํ คหปติ สพฺพํ ปธานํ “น ปทหิตพฺพนฺ”ติ วทามิ.\csromanspacer{} นาหํ คหปติ สพฺโพ ปฏินิสฺสคฺโค “ปฏินิสฺสชฺชิตพฺโพ”ติ วทามิ, น ปนาหํ คหปติ สพฺโพ ปฏินิสฺสคฺโค “น ปฏินิสฺสชฺชิตพฺโพ”ติ วทามิ.\csromanspacer{} นาหํ คหปติ สพฺพา วิมุตฺติ “วิมุจฺจิตพฺพา”ติ วทามิ, น ปนาหํ คหปติ สพฺพา วิมุตฺติ “น วิมุจฺจิตพฺพา”ติ วทามิ.}

\prose{ยํ หิ คหปติ ตปํ ตปโต อกุสลา ธมฺมา อภิวฑฺฒนฺติ, กุสลา ธมฺมา ปริหายนฺติ.\csromanspacer{} เอวรูปํ ตปํ “น ตปิตพฺพนฺ”ติ วทามิ.\csromanspacer{} ยญฺจ ขฺวสฺส คหปติ ตปํ ตปโต อกุสลา ธมฺมา ปริหายนฺติ, กุสลา ธมฺมา อภิวฑฺฒนฺติ.\csromanspacer{} เอวรูปํ ตปํ “ตปิตพฺพนฺ”ติ วทามิ.}

\prose{ยํ หิ คหปติ สมาทานํ สมาทิยโต อกุสลา ธมฺมา อภิวฑฺฒนฺติ, กุสลา ธมฺมา ปริหายนฺติ.\csromanspacer{} เอวรูปํ สมาทานํ “น สมาทิตพฺพนฺ”ติ วทามิ.\csromanspacer{} ยญฺจ ขฺวสฺส คหปติ สมาทานํ สมาทิยโต อกุสลา ธมฺมา ปริหายนฺติ, กุสลา ธมฺมา อภิวฑฺฒนฺติ.\csromanspacer{} เอวรูปํ สมาทานํ “สมาทิตพฺพนฺ”ติ วทามิ.}

\csromanmatikaanchor{mtk.135}
\csromantocmarkref{subsection}{4. วชฺชิยมาหิตสุตฺต}{mtk.135}
\bookmark[dest={mtk.135},level=2]{4. วชฺชิยมาหิตสุตฺต}
\prose{ยํ หิ คหปติ ปธานํ ปทหโต อกุสลา ธมฺมา อภิวฑฺฒนฺติ, กุสลา ธมฺมา ปริหายนฺติ.\csromanspacer{} เอวรูปํ ปธานํ “น ปทหิตพฺพนฺ”ติ วทามิ.\csromanspacer{} ยญฺจ ขฺวสฺส คหปติ ปธานํ ปทหโต อกุสลา ธมฺมา ปริหายนฺติ, กุสลา ธมฺมา อภิวฑฺฒนฺติ.\csromanspacer{} เอวรูปํ ปธานํ “ปทหิตพฺพนฺ”ติ วทามิ.}

\gambhira{ทสกนิปาตปาฬิ}
\prose{\csromanfolio{412}ยํ หิ คหปติ ปฏินิสฺสคฺคํ ปฏินิสฺสชฺชโต อกุสลา ธมฺมา อภิวฑฺฒนฺติ, กุสลา ธมฺมา ปริหายนฺติ.\csromanspacer{} เอวรูโป ปฏินิสฺสคฺโค “น ปฏินิสฺสชฺชิตพฺโพ”ติ วทามิ.\csromanspacer{} ยญฺจ ขฺวสฺส คหปติ ปฏินิสฺสคฺคํ ปฏินิสฺสชฺชโต อกุสลา ธมฺมา ปริหายนฺติ, กุสลา ธมฺมา อภิวฑฺฒนฺติ.\csromanspacer{} เอวรูโป ปฏินิสฺสคฺโค “ปฏินิสฺสชฺชิตพฺโพ”ติ วทามิ.}

\prose{ยํ หิ คหปติ วิมุตฺติํ วิมุจฺจโต อกุสลา ธมฺมา อภิวฑฺฒนฺติ, กุสลา ธมฺมา ปริหายนฺติ.\csromanspacer{} เอวรูปา วิมุตฺติ “น วิมุจฺจิตพฺพา”ติ วทามิ.\csromanspacer{} ยญฺจ ขฺวสฺส คหปติ วิมุตฺติํ วิมุจฺจโต อกุสลา ธมฺมา ปริหายนฺติ, กุสลา ธมฺมา อภิวฑฺฒนฺติ.\csromanspacer{} เอวรูปา วิมุตฺติ “วิมุจฺจิตพฺพา”ติ วทามีติ.}

\prose{อถ โข วชฺชิยมาหิโต คหปติ ภควตา ธมฺมิยา กถาย สนฺทสฺสิโต สมาทปิโต สมุตฺเตชิโต สมฺปหํสิโต อุฏฺฐายาสนา ภควนฺตํ อภิวาเทตฺวา ปทกฺขิณํ กตฺวา ปกฺกามิ.}

\prose{อถ โข ภควา อจิรปกฺกนฺเต วชฺชิยมาหิเต คหปติมฺหิ ภิกฺขู อามนฺเตสิ “โยปิ โส ภิกฺขเว ภิกฺขุ ทีฆรตฺตํ อปฺปรชกฺโข อิมสฺมิํ ธมฺมวินเย, โสปิ เอวเมวํ อญฺญติตฺถิเย ปริพฺพาชเก สหธมฺเมน สุนิคฺคหิตํ นิคฺคเณฺหยฺย, ยถา ตํ วชฺชิยมาหิเตน คหปตินา นิคฺคหิตาติ.\csromanspacer{}\csromanspacer{}จตุตฺถํ.}

\csromanheader{2}{5. อุตฺติยสุตฺต}
\proseitem{95}{อถ โข อุตฺติโย ปริพฺพาชโก เยน ภควา เตนุปสงฺกมิ, อุปสงฺกมิตฺวา ภควตา สทฺธิํ สมฺโมทิ, สมฺโมทนียํ กถํ สารณียํ วีติสาเรตฺวา เอกมนฺตํ นิสีทิ, เอกมนฺตํ นิสินฺโน โข อุตฺติโย ปริพฺพาชโก ภควนฺตํ เอตทโวจ “กิํ นุ โข โภ โคตม สสฺสโต โลโก, อิทเมว สจฺจํ โมฆมญฺญนฺ”ติ.\csromanspacer{} อพฺยากตํ โข เอตํ อุตฺติย มยา “สสฺสโต โลโก, อิทเมว สจฺจํ โมฆมญฺญนฺ”ติ.}

\prose{กิํ ปน โภ โคตม อสสฺสโต โลโก, อิทเมว สจฺจํ โมฆมญฺญนฺติ.\csromanspacer{} เอตมฺปิ โข อุตฺติย อพฺยากตํ มยา “อสสฺสโต โลโก, อิทเมว สจฺจํ โมฆมญฺญนฺ”ติ.}

\prose{กิํ นุ โข โภ โคตม อนฺตวา โลโก ฯเปฯ อนนฺตวา โลโก.\csromanspacer{} ตํ ชีวํ ตํ สรีรํ.\csromanspacer{} อญฺญํ ชีวํ อญฺญํ สรีรํ.\csromanspacer{} โหติ ตถาคโต ปรํ}

\vspace{\tipitakaparskip}
\csromancenter{(10) 5.\csromanspacer{} อุปาลิวคฺค มรณา.\csromanspacer{} น โหติ ตถาคโต ปรํ มรณา.\csromanspacer{} โหติ จ น จ โหติ ตถาคโต ปรํ มรณา.\csromanspacer{} เนว โหติ น น โหติ ตถาคโต ปรํ มรณา, อิทเมว สจฺจํ โมฆมญฺญนฺติ.\csromanspacer{} เอตมฺปิ โข อุตฺติย อพฺยากตํ มยา “เนว โหติ น น โหติ ตถาคโต ปรํ มรณา, อิทเมว สจฺจํ โมฆมญฺญนฺ”ติ.}
\prose{\csromanfolio{413}“กิํ นุ โข โภ โคตม สสฺสโต โลโก, อิทเมว สจฺจํ โมฆมญฺญนฺ”ติ อิติ ปุฏฺโฐ สมาโน “อพฺยากตํ โข เอตํ อุตฺติย มยา สสฺสโต โลโก, อิทเมว สจฺจํ โมฆมญฺญนฺ”ติ วเทสิ.}

\prose{“กิํ ปน โภ โคตม อสสฺสโต โลโก, อิทเมว สจฺจํ โมฆมญฺญนฺ”ติ อิติ ปุฏฺโฐ สมาโน เอตมฺปิ โข อุตฺติย อพฺยากตํ มยา “อสสฺสโต โลโก, อิทเมว สจฺจํ โมฆมญฺญนฺ”ติ วเทสิ.}

\prose{“กิํ นุ โข โภ โคตม อนฺตวา โลโก ฯเปฯ อนนฺตวา โลโก.\csromanspacer{} ตํ ชีวํ ตํ สรีรํ.\csromanspacer{} อญฺญํ ชีวํ อญฺญํ สรีรํ.\csromanspacer{} โหติ ตถาคโต ปรํ มรณา.\csromanspacer{} น โหติ ตถาคโต ปรํ มรณา.\csromanspacer{} โหติ จ น จ โหติ ตถาคโต ปรํ มรณา, เนว โหติ น น โหติ ตถาคโต ปรํ มรณา, อิทเมว สจฺจํ โมฆมญฺญนฺ”ติ อิติ ปุฏฺโฐ สมาโน “เอตมฺปิ โข อุตฺติย อพฺยากตํ มยา เนว โหติ น น โหติ ตถาคโต ปรํ มรณา, อิทเมว สจฺจํ โมฆมญฺญนฺ”ติ วเทสิ.\csromanspacer{} อถ กิญฺจรหิ โภตา โคตเมน พฺยากตนฺติ.}

\prose{อภิญฺญาย โข อหํ อุตฺติย สาวกานํ ธมฺมํ เทเสมิ สตฺตานํ วิสุทฺธิยา โสกปริเทวานํ สมติกฺกมาย ทุกฺขโทมนสฺสานํ อตฺถงฺคมาย ญายสฺส อธิคมาย นิพฺพานสฺส สจฺฉิกิริยายาติ.}

\prose{ยํ ปเนตํ ภวํ โคตโม อภิญฺญาย สาวกานํ ธมฺมํ เทเสสิ สตฺตานํ วิสุทฺธิยา โสกปริเทวานํ สมติกฺกมาย ทุกฺขโทมนสฺสานํ อตฺถงฺคมาย ญายสฺส อธิคมาย นิพฺพานสฺส สจฺฉิกิริยาย.\csromanspacer{} สพฺโพ วา\footnote{สพฺโพ จ (ก)} เตน โลโก นียติ\footnote{นียิสฺสติ (สี), นิยฺยาสฺสติ (สฺยา), นิยฺยิสฺสติ (อิ)} อุปฑฺโฒ วา ติภาโค วาติ\footnote{ติภาโค วาติ วเทหิ (ก)}.\csromanspacer{} เอวํ วุตฺเต ภควา ตุณฺหี อโหสิ.}

\gambhira{ทสกนิปาตปาฬิ}
\prose{\csromanfolio{414}อถ โข อายสฺมโต อานนฺทสฺส เอตทโหสิ “มา เหวํ โข อุตฺติโย ปริพฺพาชโก ปาปกํ ทิฏฺฐิคตํ ปฏิลภิ ‘สพฺพสามุกฺกํสิกํ วต เม สมโณ โคตโม ปญฺหํ ปุฏฺโฐ สํสาเทติ, โน วิสฺสชฺเชติ, น นูน วิสหตี’ติ, ตทสฺส อุตฺติยสฺส ปริพฺพาชกสฺส ทีฆรตฺตํ อหิตาย ทุกฺขายา”ติ.}

\csromanmatikaanchor{mtk.136}
\csromantocmarkref{subsection}{5. อุตฺติยสุตฺต}{mtk.136}
\bookmark[dest={mtk.136},level=2]{5. อุตฺติยสุตฺต}
\prose{อถ โข อายสฺมา อานนฺโท อุตฺติยํ ปริพฺพาชกํ เอตทโวจ“เตนหาวุโส อุตฺติย อุปมํ เต กริสฺสามิ, อุปมาย มิเธกจฺเจ วิญฺญู ปุริสา ภาสิตสฺส อตฺถํ อาชานนฺติ.\csromanspacer{} เสยฺยถาปิ อาวุโส อุตฺติย รญฺโญ ปจฺจนฺติมํ นครํ ทฬฺหุทฺธาปํ\footnote{ทฬฺหุทฺทาปํ (สี, อิ)} ทฬฺหปาการโตรณํ เอกทฺวารํ, ตตฺรสฺส โทวาริโก ปณฺฑิโต พฺยตฺโต เมธาวี อญฺญาตานํ นิวาเรตา ญาตานํ ปเวเสตา, โส ตสฺส นครสฺส สมนฺตา อนุปริยายปถํ อนุกฺกมติ, อนุปริยายปถํ อนุกฺกมมาโน น ปสฺเสยฺย ปาการสนฺธิํ วา ปาการวิวรํ วา อนฺตมโส พิฬารนิกฺขมนมตฺตมฺปิ.\csromanspacer{} โน จ ขฺวสฺส เอวํ ญาณํ โหติ “เอตฺตกา ปาณา อิมํ นครํ ปวิสนฺติ วา นิกฺขมนฺติ วา”ติ.\csromanspacer{} อถ โข ขฺวสฺส เอวเมตฺถ โหติ “เย โข เกจิ โอฬาริกา ปาณา อิมํ นครํ ปวิสนฺติ วา นิกฺขมนฺติ วา, สพฺเพ เต อิมินา ทฺวาเรน ปวิสนฺติ วา นิกฺขมนฺติ วา”ติ.}

\prose{เอวเมวํ โข อาวุโส อุตฺติย น ตถาคตสฺส เอวํ อุสฺสุกฺกํ โหติ “สพฺโพ วา เตน โลโก นียติ อุปฑฺโฒ วา ติภาโค วา”ติ.\csromanspacer{} อถ โข เอวเมตฺถ ตถาคตสฺส โหติ “เย โข เกจิ โลกมฺหา นียิํสุ วา นียนฺติ วา นียิสฺสนฺติ วา, สพฺเพ เต ปญฺจ นีวรเณ ปหาย เจตโส อุปกฺกิเลเส ปญฺญาย ทุพฺพลีกรเณ จตูสุ สติปฏฺฐาเนสุ สุปฺปติฏฺฐิตจิตฺตา สตฺต โพชฺฌงฺเค ยถาภูตํ ภาเวตฺวา เอวเมเต\footnote{เอวเมเตน (ก)} โลกมฺหา นียิํสุ วา นียนฺติ วา นียิสฺสนฺติ วา”ติ.\csromanspacer{} ยเทว โข ตฺวํ\footnote{ยเทว เขฺวตฺถ (ก)} อาวุโส อุตฺติย ภควนฺตํ ปญฺหํ\footnote{อิมํ ปญฺหํ (สฺยา, ก)} อปุจฺฉิ, ตเทเวตํ ปญฺหํ ภควนฺตํ อญฺเญน ปริยาเยน อปุจฺฉิ, ตสฺมา เต ตํ ภควา น พฺยากาสี”ติ.\csromanspacer{}\csromanspacer{}ปญฺจมํ.}

\csromanheader{2}{(10) 5. อุปาลิวคฺค \textbf{6. โกกนุทสุตฺต}}
\proseitem{96}{\csromanfolio{415}เอกํ สมยํ อายสฺมา อานนฺโท ราชคเห วิหรติ ตโปทาราเม.\csromanspacer{} อถ โข อายสฺมา อานนฺโท รตฺติยา ปจฺจูสสมยํ ปจฺจุฏฺฐาย เยน ตโปทา เตนุปสงฺกมิ คตฺตานิ ปริสิญฺจิตุํ, ตโปทาย\footnote{ตโปเท (ก)} คตฺตานิ ปริสิญฺจิตฺวา ปจฺจุตฺตริตฺวา เอกจีวโร อฏฺฐาสิ คตฺตานิ ปุพฺพาปยมาโน.\csromanspacer{} โกกนุโทปิ โข ปริพฺพาชโก รตฺติยา ปจฺจูสสมยํ ปจฺจุฏฺฐาย เยน ตโปทา เตนุปสงฺกมิ คตฺตานิ ปริสิญฺจิตุํ.}

\prose{อทฺทสา โข โกกนุโท ปริพฺพาชโก อายสฺมนฺตํ อานนฺทํ ทูรโตว อาคจฺฉนฺตํ, ทิสฺวาน อายสฺมนฺตํ อานนฺทํ เอตทโวจ “เกฺวตฺถ\footnote{โก เตตฺถ (สี), กฺวตฺถ (อิ, ก)} อาวุโส”ติ. “อหมาวุโส ภิกฺขู”ติ.}

\prose{กตเมสํ อาวุโส ภิกฺขูนนฺติ.\csromanspacer{} สมณานํ อาวุโส สกฺยปุตฺติยานนฺติ.}

\prose{ปุจฺเฉยฺยาม มยํ อายสฺมนฺตํ กิญฺจิเทว เทสํ, สเจ อายสฺมา โอกาสํ กโรติ ปญฺหสฺส เวยฺยากรณายาติ.\csromanspacer{} ปุจฺฉาวุโส, สุตฺวา เวทิสฺสามาติ.}

\prose{กิํ นุ โข โภ “สสฺสโต โลโก, อิทเมว สจฺจํ โมฆมญฺญนฺ”ติ เอวํทิฏฺฐิ\footnote{เอวํทิฏฺฐิโก (สฺยา)} ภวนฺติ.\csromanspacer{} น โข อหํ อาวุโส เอวํทิฏฺฐิ “สสฺสโต โลโก, อิทเมว สจฺจํ โมฆมญฺญนฺ”ติ.}

\prose{กิํ ปน โภ “อสสฺสโต โลโก, อิทเมว สจฺจํ โมฆมญฺญนฺ”ติ เอวํทิฏฺฐิ ภวนฺติ.\csromanspacer{} น โข อหํ อาวุโส เอวํทิฏฺฐิ “อสสฺสโต โลโก, อิทเมว สจฺจํ โมฆมญฺญนฺ”ติ.}

\prose{กิํ นุ โข โภ “อนฺตวา โลโก ฯเปฯ อนนฺตวา โลโก.\csromanspacer{} ตํ ชีวํ ตํ สรีรํ.\csromanspacer{} อญฺญํ ชีวํ อญฺญํ สรีรํ.\csromanspacer{} โหติ ตถาคโต ปรํ มรณา.\csromanspacer{} น โหติ ตถาคโต ปรํ มรณา.\csromanspacer{} โหติ จ น จ โหติ ตถาคโต ปรํ มรณา.\csromanspacer{} เนว โหติ น น โหติ ตถาคโต ปรํ มรณา, อิทเมว สจฺจํ โมฆมญฺญนฺติ เอวํทิฏฺฐิ ภวนฺ”ติ.\csromanspacer{} น โข อหํ อาวุโส เอวํทิฏฺฐิ “เนว โหติ น น โหติ ตถาคโต ปรํ มรณา, อิทเมว สจฺจํ โมฆมญฺญนฺ”ติ.}

\gambhira{ทสกนิปาตปาฬิ}
\prose{\csromanfolio{416}เตน หิ ภวํ น ชานาติ น ปสฺสตีติ.\csromanspacer{} น โข อหํ อาวุโส น ชานามิ น ปสฺสามิ, ชานามหํ อาวุโส ปสฺสามีติ.}

\prose{“กิํ นุ โข โภ สสฺสโต โลโก, อิทเมว สจฺจํ โมฆมญฺญนฺติ เอวํทิฏฺฐิ ภวนฺ”ติ อิติ ปุฏฺโฐ สมาโน “น โข อหํ อาวุโส เอวํทิฏฺฐิ สสฺสโต โลโก, อิทเมว สจฺจํ โมฆมญฺญนฺ”ติ วเทสิ.}

\prose{“กิํ ปน โภ อสสฺสโต โลโก, อิทเมว สจฺจํ โมฆมญฺญนฺติ เอวํทิฏฺฐิ ภวนฺ”ติ อิติ ปุฏฺโฐ สมาโน “น โข อหํ อาวุโส เอวํทิฏฺฐิ อสสฺสโต โลโก, อิทเมว สจฺจํ โมฆมญฺญนฺ”ติ วเทสิ.}

\csromanmatikaanchor{mtk.137}
\csromantocmarkref{subsection}{6. โกกนุทสุตฺต}{mtk.137}
\bookmark[dest={mtk.137},level=2]{6. โกกนุทสุตฺต}
\prose{“กิํ นุ โข โภ อนฺตวา โลโก ฯเปฯ อนนฺตวา โลโก.\csromanspacer{} ตํ ชีวํ ตํ สรีรํ.\csromanspacer{} อญฺญํ ชีวํ อญฺญํ สรีรํ.\csromanspacer{} โหติ ตถาคโต ปรํ มรณา.\csromanspacer{} น โหติ ตถาคโต ปรํ มรณา.\csromanspacer{} โหติ จ น จ โหติ ตถาคโต ปรํ มรณา.\csromanspacer{} เนว โหติ น น โหติ ตถาคโต ปรํ มรณา, อิทเมว สจฺจํ โมฆมญฺญนฺติ เอวํทิฏฺฐิ ภวนฺ”ติ อิติ ปุฏฺโฐ สมาโน “น โข อหํ อาวุโส เอวํทิฏฺฐิ เนว โหติ น น โหติ ตถาคโต ปรํ มรณา, อิทเมว สจฺจํ โมฆมญฺญนฺ”ติ วเทสิ.}

\prose{“เตน หิ ภวํ น ชานาติ น ปสฺสตี”ติ อิติ ปุฏฺโฐ สมาโน “น โข อหํ อาวุโส น ชานามิ น ปสฺสามิ, ชานามหํ อาวุโส ปสฺสามี”ติ วเทสิ.\csromanspacer{} ยถา กถํ ปนาวุโส อิมสฺส ภาสิตสฺส อตฺโถ ทฏฺฐพฺโพติ.}

\prose{“สสฺสโต โลโก, อิทเมว สจฺจํ โมฆมญฺญนฺ”ติ โข อาวุโส ทิฏฺฐิคตเมตํ. “อสสฺสโต โลโก, อิทเมว สจฺจํ โมฆมญฺญนฺ”ติ โข อาวุโส ทิฏฺฐิคตเมตํ. “อนฺตวา โลโก ฯเปฯ อนนฺตวา โลโก.\csromanspacer{} ตํ ชีวํ ตํ สรีรํ.\csromanspacer{} อญฺญํ ชีวํ อญฺญํ สรีรํ.\csromanspacer{} โหติ ตถาคโต ปรํ มรณา.\csromanspacer{} น โหติ ตถาคโต ปรํ มรณา.\csromanspacer{} โหติ จ น จ โหติ ตถาคโต ปรํ มรณา.\csromanspacer{} เนว โหติ น น โหติ ตถาคโต ปรํ มรณา, อิทเมว สจฺจํ โมฆ- มญฺญนฺ”ติ โข อาวุโส ทิฏฺฐิคตเมตํ.}

\chapterheadpageread{(10) 5. อุปาลิวคฺค}
\prose{\csromanfolio{417}ยาวตา อาวุโส ทิฏฺฐิ\footnote{ทิฏฺฐิคตา (สพฺพตฺถ)} ยาวตา ทิฏฺฐิฏฺฐานํ ทิฏฺฐิ-อธิฏฺฐานํ ทิฏฺฐิปริยุฏฺฐานํ ทิฏฺฐิสมุฏฺฐานํ ทิฏฺฐิสมุคฺฆาโต\footnote{ยาวตา ทิฏฺฐิฏฺฐาน อธิฏฺฐาน ปริยุฏฺฐาน สมุฏฺฐาน สมุคฺฆาโต (สี, อิ)}, ตมหํ ชานามิ ตมหํ ปสฺสามิ.\csromanspacer{} ตมหํ ชานนฺโต ตมหํ ปสฺสนฺโต กฺยาหํ วกฺขามิ “น ชานามิ น ปสฺสามี”ติ.\csromanspacer{} ชานามหํ อาวุโส ปสฺสามีติ.}

\prose{โก นาโม อายสฺมา, กถญฺจ ปนายสฺมนฺตํ สพฺรหฺมจารี ชานนฺตีติ. “อานนฺโท”ติ โข เม อาวุโส นามํ, “อานนฺโท”ติ จ ปน มํ สพฺรหฺมจารี ชานนฺตีติ.\csromanspacer{} มหาจริเยน วต กิร โภ สทฺธิํ มนฺตยมานา น ชานิมฺห “อายสฺมา อานนฺโท”ติ.\csromanspacer{} สเจ หิ มยํ ชาเนยฺยาม “อยํ อายสฺมา อานนฺโท”ติ, เอตฺตกมฺปิ โน นปฺปฏิภาเยยฺย\footnote{นปฺปฏิภาเสยฺยาม (ก) นปฺปฏิภาเสยฺย (พหูสุ) ม 3. 172 ปิฏฺเฐ ปสฺสิตพฺพํ.}, ขมตุ จ เม อายสฺมา อานนฺโทติ.\csromanspacer{}\csromanspacer{}ฉฏฺฐํ.}

\csromanheader{2}{7. อาหุเนยฺยสุตฺต}
\proseitem{97}{ทสหิ ภิกฺขเว ธมฺเมหิ สมนฺนาคโต ภิกฺขุ อาหุเนยฺโย โหติ ปาหุเนยฺโย ทกฺขิเณยฺโย อญฺชลิกรณีโย อนุตฺตรํ ปุญฺญกฺเขตฺตํ โลกสฺส.}

\prose{กตเมหิ ทสหิ? อิธ ภิกฺขเว ภิกฺขุ สีลวา โหติ, ปาติโมกฺขสํวรสํวุโต วิหรติ อาจารโคจรสมฺปนฺโน, อณุมตฺเตสุ วชฺเชสุ ภยทสฺสาวี สมาทาย สิกฺขติ สิกฺขาปเทสุ.}

\prose{พหุสฺสุโต โหติ สุตธโร สุตสนฺนิจโย, เย เต ธมฺมา อาทิกลฺยาณา มชฺเฌกลฺยาณา ปริโยสานกลฺยาณา สาตฺถํ สพฺยญฺชนํ เกวลปริปุณฺณํ ปริสุทฺธํ พฺรหฺมจริยํ อภิวทนฺติ, ตถารูปาสฺส ธมฺมา พหุสฺสุตา โหนฺติ ธาตา วจสา ปริจิตา มนสานุเปกฺขิตา ทิฏฺฐิยา สุปฺปฏิวิทฺธา.}

\prose{กลฺยาณมิตฺโต โหติ กลฺยาณสหาโย กลฺยาณสมฺปวงฺโก.}

\prose{สมฺมาทิฏฺฐิโก โหติ สมฺมาทสฺสเนน สมนฺนาคโต.}

\prose{อเนกวิหิตํ อิทฺธิวิธํ ปจฺจนุโภติ, เอโกปิ หุตฺวา พหุธา โหติ, พหุธาปิ หุตฺวา เอโก โหติ, อาวิภาวํ ติโรภาวํ ติโรกุฏฺฏํ}

\vspace{\tipitakaparskip}
\csromancenter{ทสกนิปาตปาฬิ ติโรปาการํ ติโรปพฺพตํ อสชฺชมาโน คจฺฉติ เสยฺยถาปิ อากาเส, ปถวิยาปิ อุมฺมุชฺชนิมุชฺชํ กโรติ เสยฺยถาปิ อุทเก, อุทเกปิ อภิชฺชมาเน คจฺฉติ เสยฺยถาปิ ปถวิยํ, อากาเสปิ ปลฺลงฺเกน กมติ เสยฺยถาปิ ปกฺขี สกุโณ, อิเมปิ จนฺทิมสูริเย เอวํมหิทฺธิเก เอวํมหานุภาเว ปาณินา ปรามสติ\footnote{ปริมสติ (สี)} ปริมชฺชติ, ยาว พฺราหฺมโลกาปิ กาเยน วสํ วตฺเตติ.}
\prose{\csromanfolio{418}ทิพฺพาย โสตธาตุยา วิสุทฺธาย อติกฺกนฺตมานุสิกาย อุโภ สทฺเท สุณาติ ทิพฺเพ จ มานุเส จ เย ทูเร สนฺติเก จ.}

\prose{ปรสตฺตานํ ปรปุคฺคลานํ เจตสา เจโต ปริจฺจ ปชานาติ, สราคํ วา จิตฺตํ “สราคํ จิตฺตนฺ”ติ ปชานาติ.\csromanspacer{} วีตราคํ วา จิตฺตํ “วีตราคํ จิตฺตนฺ”ติ ปชานาติ, สโทสํ วา จิตฺตํ ฯเปฯ วีตโทสํ วา จิตฺตํ ฯเปฯ สโมหํ วา จิตฺตํ.\csromanspacer{} วีตโมหํ วา จิตฺตํ.\csromanspacer{} สํขิตฺตํ วา จิตฺตํ.\csromanspacer{} วิกฺขิตฺตํ วา จิตฺตํ.\csromanspacer{} มหคฺคตํ วา จิตฺตํ.\csromanspacer{} อมหคฺคตํ วา จิตฺตํ.\csromanspacer{} ส-อุตฺตรํ วา จิตฺตํ.\csromanspacer{} อนุตฺตรํ วา จิตฺตํ.\csromanspacer{} สมาหิตํ วา จิตฺตํ.\csromanspacer{} อสมาหิตํ วา จิตฺตํ.\csromanspacer{} วิมุตฺตํ วา จิตฺตํ.\csromanspacer{} อวิมุตฺตํ วา จิตฺตํ “อวิมุตฺตํ จิตฺตนฺ”ติ ปชานาติ.}

\prose{อเนกวิหิตํ ปุพฺเพนิวาสํ อนุสฺสรติ.\csromanspacer{} เสยฺยถิทํ, เอกมฺปิ ชาติํ เทฺวปิ ชาติโย ติสฺโสปิ ชาติโย จตสฺโสปิ ชาติโย ปญฺจปิ ชาติโย ทสปิ ชาติโย วีสมฺปิ ชาติโย ติํสมฺปิ ชาติโย จตฺตาลีสมฺปิ ชาติโย ปญฺญาสมฺปิ ชาติโย ชาติสตมฺปิ ชาติสหสฺสมฺปิ ชาติสตสหสฺสมฺปิ อเนเกปิ สํวฏฺฏกปฺเป อเนเกปิ วิวฏฺฏกปฺเป อเนเกปิ สํวฏฺฏวิวฏฺฏกปฺเป “อมุตฺราสิํ เอวํนาโม เอวํโคตฺโต เอวํวณฺโณ เอวมาหาโร เอวํสุขทุกฺขปฏิสํเวที เอวมายุปริยนฺโต, โส ตโต จุโต อมุตฺร อุทปาทิํ, ตตฺราปาสิํ เอวํนาโม เอวํโคตฺโต เอวํวณฺโณ เอวมาหาโร เอวํสุขทุกฺขปฏิสํเวที เอวมายุปริยนฺโต, โส ตโต จุโต อิธูปปนฺโน”ติ อิติ สาการํ ส-อุทฺเทสํ อเนกวิหิตํ ปุพฺเพนิวาสํ อนุสฺสรติ.}

\prose{ทิพฺเพน จกฺขุนา วิสุทฺเธน อติกฺกนฺตมานุสเกน สตฺเต ปสฺสติ จวมาเน อุปปชฺชมาเน หีเน ปณีเต สุวณฺเณ ทุพฺพณฺเณ สุคเต ทุคฺคเต, ยถากมฺมูปเค สตฺเต ปชานาติ, “อิเม วต โข โภนฺโต สตฺตา}

\csromanmatikaanchor{mtk.138}
\csromantocmarkref{subsection}{7. อาหุเนยฺยสุตฺต}{mtk.138}
\bookmark[dest={mtk.138},level=2]{7. อาหุเนยฺยสุตฺต}
\vspace{\tipitakaparskip}
\csromancenter{(10) 5.\csromanspacer{} อุปาลิวคฺค กายทุจฺจริเตน สมนฺนาคตา วจีทุจฺจริเตน สมนฺนาคตา มโนทุจฺจริเตน สมนฺนาคตา, อริยานํ อุปวาทกา, มิจฺฉาทิฏฺฐิกา มิจฺฉาทิฏฺฐิกมฺมสมาทานา.\csromanspacer{} เต กายสฺส เภทา ปรํ มรณา อปายํ ทุคฺคติํ วินิปาตํ นิรยํ อุปปนฺนา.\csromanspacer{} อิเม วา ปน โภนฺโต สตฺตา กายสุจริเตน สมนฺนาคตา วจีสุจริเตน สมนฺนาคตา มโนสุจริเตน สมนฺนาคตา, อริยานํ อนุปวาทกา, สมฺมาทิฏฺฐิกา สมฺมาทิฏฺฐิกมฺมสมาทานา.\csromanspacer{} เต กายสฺส เภทา ปรํ มรณา สุคติํ สคฺคํ โลกํ อุปปนฺนา”ติ.\csromanspacer{} อิติ ทิพฺเพน จกฺขุนา วิสุทฺเธน อติกฺกนฺตมานุสเกน สตฺเต ปสฺสติ จวมาเน อุปปชฺชมาเน หีเน ปณีเต สุวณฺเณ ทุพฺพณฺเณ สุคเต ทุคฺคเต, ยถากมฺมูปเค สตฺเต ปชานาติ.}
\prose{\csromanfolio{419}อาสวานํ ขยา อนาสวํ เจโตวิมุตฺติํ ปญฺญาวิมุตฺติํ ทิฏฺเฐว ธมฺเม สยํ อภิญฺญา สจฺฉิกตฺวา อุปสมฺปชฺช วิหรติ.\csromanspacer{} อิเมหิ โข ภิกฺขเว ทสหิ ธมฺเมหิ สมนฺนาคโต ภิกฺขุ อาหุเนยฺโย โหติ ปาหุเนโย ทกฺขิเณยฺโย อญฺชลิกรณีโย อนุตฺตรํ ปุญฺญกฺเขตฺตํ โลกสฺสาติ.\csromanspacer{}\csromanspacer{}สตฺตมํ.}

\csromanheader{2}{8. เถรสุตฺต}
\proseitem{98}{ทสหิ ภิกฺขเว ธมฺเมหิ สมนฺนาคโต เถโร ภิกฺขุ ยสฺสํ ยสฺสํ ทิสายํ วิหรติ, ผาสุเยว วิหรติ.\csromanspacer{} กตเมหิ ทสหิ? เถโร โหติ รตฺตญฺญู จิรปพฺพชิโต.\csromanspacer{} สีลวา โหติ ฯเปฯ สมาทาย สิกฺขติ สิกฺขาปเทสุ.\csromanspacer{} พหุสฺสุโต โหติ ฯเปฯ ทิฏฺฐิยา สุปฺปฏิวิทฺโธ.\csromanspacer{} อุภยานิ โข ปนสฺส ปาติโมกฺขานิ วิตฺถาเรน สฺวาคตานิ โหนฺติ สุวิภตฺตานิ สุปฺปวตฺตีนิ สุวินิจฺฉิตานิ สุตฺตโส อนุพฺยญฺชนโส.\csromanspacer{} อธิกรณสมุปฺปาท- วูปสมกุสโล โหติ.\csromanspacer{} ธมฺมกาโม โหติ ปิยสมุทาหาโร อภิธมฺเม อภิวินเย อุฬารปาโมชฺโช.\csromanspacer{} สนฺตุฏฺโฐ โหติ อิตรีตรจีวร ปิณฺฑปาต เสนาสน คิลานปฺปจฺจย เภสชฺชปริกฺขาเรน.\csromanspacer{} ปาสาทิโก โหติ อภิกฺกนฺตปฏิกฺกนฺเต\footnote{อภิกฺกนฺตปฏิกฺกนฺโต (ก)} สุสํวุโต อนฺตรฆเร นิสชฺชาย.\csromanspacer{} จตุนฺนํ ฌานานํ อาภิเจตสิกานํ ทิฏฺฐธมฺมสุขวิหารานํ นิกามลาภี โหติ อกิจฺฉลาภี อกสิรลาภี.\csromanspacer{} อาสวานญฺจ ขยา อนาสวํ เจโตวิมุตฺติํ ปญฺญาวิมุตฺติํ ทิฏฺเฐว ธมฺเม สยํ อภิญฺญา สจฺฉิกตฺวา อุปสมฺปชฺช วิหรติ.}

\vspace{\tipitakaparskip}
\csromancenter{ทสกนิปาตปาฬิ อิเมหิ โข ภิกฺขเว ทสหิ ธมฺเมหิ สมนฺนาคโต เถโร ภิกฺขุ ยสฺสํ ยสฺสํ ทิสายํ วิหรติ, ผาสุเยว วิหรตีติ.\csromanspacer{}\csromanspacer{}อฏฺฐมํ.}
\csromanheader{2}{9. อุปาลิสุตฺต}
\proseitem{99}{\csromanfolio{420}อถ โข อายสฺมา อุปาลิ เยน ภควา เตนุปสงฺกมิ, อุปสงฺกมิตฺวา ภควนฺตํ อภิวาเทตฺวา เอกมนฺตํ นิสีทิ, เอกมนฺตํ นิสินฺโน โข อายสฺมา อุปาลิ ภควนฺตํ เอตทโวจ “อิจฺฉามหํ ภนฺเต อรญฺญวนปตฺถานิ ปนฺตานิ เสนาสนานิ ปฏิเสวิตุนฺ”ติ.}

\prose{ทุรภิสมฺภวานิ หิ โข\footnote{ทุรภิสมฺภวานิ โข (สี, อิ)} อุปาลิ อรญฺญวนปตฺถานิ ปนฺตานิ เสนาสนานิ, ทุกฺกรํ ปวิเวกํ, ทุรภิรมํ เอกตฺเต, หรนฺติ มญฺเญ มโน วนานิ สมาธิํ อลภมานสฺส ภิกฺขุโน.\csromanspacer{} โย โข อุปาลิ เอวํ วเทยฺย “อหํ สมาธิํ อลภมาโน อรญฺญวนปตฺถานิ ปนฺตานิ เสนาสนานิ ปฏิเสวิสฺสามี”ติ, ตสฺเสตํ ปาฏิกงฺขํ “สํสีทิสฺสติ วา อุปฺลวิสฺสติ วา”ติ\footnote{อุปฺปิลวิสฺสติ วา (สี, สฺยา, อิ)}.}

\prose{เสยฺยถาปิ อุปาลิ มหา-อุทกรหโท, อถ อาคจฺเฉยฺย หตฺถินาโค สตฺตรตโน วา อฑฺฒฏฺฐรตโน\footnote{อฏฺฐรตโน (สี, อิ)} วา, ตสฺส เอวมสฺส “ยํนูนาหํ อิมํ อุทกรหทํ โอคาเหตฺวา กณฺณสํโธวิกมฺปิ ขิฑฺฑํ กีเฬยฺยํ ปิฏฺฐิสํโธวิกมฺปิ ขิฑฺฑํ กีเฬยฺยํ, กณฺณสํโธวิกมฺปิ ขิฑฺฑํ กีฬิตฺวา ปิฏฺฐิสํโธวิกมฺปิ ขิฑฺฑํ กีฬิตฺวา นฺหตฺวา\footnote{นหาตฺวา (สี, อิ), นฺหตฺวา (สฺยา)} จ ปิวิตฺวา จ ปจฺจุตฺตริตฺวา เยน กามํ ปกฺกเมยฺยนฺ”ติ, โส ตํ อุทกรหทํ โอคาเหตฺวา กณฺณสํโธวิกมฺปิ ขิฑฺฑํ กีเฬยฺย ปิฏฺฐิสํโธวิกมฺปิ ขิฑฺฑํ กีเฬยฺย, กณฺณสํโธวิกมฺปิ ขิฑฺฑํ กีฬิตฺวา ปิฏฺฐิสํโธวิกมฺปิ ขิฑฺฑํ กีฬิตฺวา นฺหตฺวา จ ปิวิตฺวา จ ปจฺจุตฺตริตฺวา เยน กามํ ปกฺกเมยฺย.\csromanspacer{} ตํ กิสฺส เหตุ? มหา อุปาลิ\footnote{มหา หุปาลิ (สี, อิ)} อตฺตภาโว คมฺภีเร คาธํ วินฺทติ.}

\prose{อถ อาคจฺเฉยฺย สโส วา พิฬาโร วา, ตสฺส เอวมสฺส “โก จาหํ, โก จ หตฺถินาโค, ยํนูนาหํ อิมํ อุทกรหทํ โอคาเหตฺวา กณฺณสํโธวิกมฺปิ ขิฑฺฑํ กีเฬยฺยํ, ปิฏฺฐิสํโธวิกมฺปิ ขิฑฺฑํ กีเฬยฺยํ, กณฺณสํโธวิกมฺปิ ขิฑฺฑํ กีฬิตฺวา ปิฏฺฐิสํโธวิกมฺปิ ขิฑฺฑํ กีฬิตฺวา นฺหตฺวา จ ปิวิตฺวา จ ปจฺจุตฺตริตฺวา เยน กามํ ปกฺกเมยฺยนฺ”ติ.\csromanspacer{} โส ตํ อุทกรหทํ สหสา อปฺปฏิสงฺขา ปกฺขนฺเทยฺย, ตสฺเสตํ ปาฏิกงฺขํ “สํสีทิสฺสติ วา อุปฺลวิสฺสติ}

\vspace{\tipitakaparskip}
\csromancenter{(10) 5.\csromanspacer{} อุปาลิวคฺค วา”ติ.\csromanspacer{} ตํ กิสฺส เหตุ? ปริตฺโต อุปาลิ อตฺตภาโว คมฺภีเร คาธํ น วินฺทติ.\csromanspacer{} เอวเมวํ โข อุปาลิ โย เอวํ วเทยฺย “อหํ สมาธิํ อลภมาโน อรญฺญวนปตฺถานิ ปนฺตานิ เสนาสนานิ ปฏิเสวิสฺสามี”ติ, ตสฺเสตํ ปาฏิกงฺขํ “สํสีทิสฺสติ วา อุปฺลวิสฺสติ วา”ติ.}
\prose{\csromanfolio{421}เสยฺยถาปิ อุปาลิ ทหโร กุมาโร มนฺโท อุตฺตานเสยฺยโก สเกน มุตฺตกรีเสน กีฬติ.\csromanspacer{} ตํ กิํ มญฺญสิ อุปาลิ, นนฺวายํ เกวลา ปริปูรา พาลขิฑฺฑาติ.\csromanspacer{} เอวํ ภนฺเต.}

\prose{ส โข โส อุปาลิ กุมาโร อปเรน สมเยน วุทฺธิมนฺวาย อินฺทฺริยานํ ปริปากมนฺวาย ยานิ กานิจิ กุมารกานํ กีฬาปนกานิ ภวนฺติ.\csromanspacer{} เสยฺยถิทํ, วงฺกกํ\footnote{วงฺกํ (สี, อิ)} ฆฏิกํ โมกฺขจิกํ จิงฺคุลกํ\footnote{ปิงฺคุลิกํ (สฺยา), จิงฺกุลกํ (ก)} ปตฺตาฬฺหกํ รถกํ ธนุกํ, เตหิ กีฬติ.\csromanspacer{} ตํ กิํ มญฺญสิ อุปลิ, นนฺวายํ ขิฑฺฑา ปุริมาย ขิฑฺฑาย อภิกฺกนฺตตรา จ ปณีตตรา จาติ.\csromanspacer{} เอวํ ภนฺเต.}

\prose{ส โข โส อุปาลิ กุมาโร อปเรน สมเยน วุทฺธิมนฺวาย อินฺทฺริยานํ ปริปากมนฺวาย ปญฺจหิ กามคุเณหิ สมปฺปิโต สมงฺคิภูโต ปริจาเรติ จกฺขุวิญฺเญยฺเยหิ รูเปหิ อิฏฺเฐหิ กนฺเตหิ มนาเปหิ ปิยรูเปหิ กามูปสํหิเตหิ รชนีเยหิ.\csromanspacer{} โสตวิญฺเญยฺเยหิ สทฺเทหิ.\csromanspacer{} ฆานวิญฺเญยฺเยหิ คนฺเธหิ.\csromanspacer{} ชิวฺหาวิญฺเญยฺเยหิ รเสหิ.\csromanspacer{} กายวิญฺเญยฺเยหิ โผฏฺฐพฺเพหิ อิฏฺเฐหิ กนฺเตหิ มนาเปหิ ปิยรูเปหิ กามูปสํหิเตหิ รชนีเยหิ.\csromanspacer{} ตํ กิํ มญฺญสิ อุปาลิ, นนฺวายํ ขิฑฺฑา ปุริมาหิ ขิฑฺฑาหิ อภิกฺกนฺตตรา จ ปณีตตรา จาติ.\csromanspacer{} เอวํ ภนฺเต.}

\csromanmatikaanchor{mtk.139}
\csromantocmarkref{subsection}{8. เถรสุตฺต}{mtk.139}
\bookmark[dest={mtk.139},level=2]{8. เถรสุตฺต}
\prose{\csromansymbolfootnote{*}{ที 1. 58; ม 2. 190 ปิฏฺเฐสุปิ.}อิธ โข ปน โว\footnote{โวติ นิปาตมตฺตํ (ฏฺฐ)} อุปาลิ ตถาคโต โลเก อุปฺปชฺชติ อรหํ สมฺมาสมฺพุทฺโธ วิชฺชาจรณสมฺปนฺโน สุคโต โลกวิทู อนุตฺตโร ปุริสทมฺมสารถิ สตฺถา เทวมนุสฺสานํ พุทฺโธ ภควา.\csromanspacer{} โส อิมํ โลกํ สเทวกํ สมารกํ สพฺรหฺมกํ สสฺสมณพฺราหฺมณิํ ปชํ สเทวมนุสฺสํ สยํ อภิญฺญา สจฺฉิกตฺวา ปเวเทติ, โส ธมฺมํ เทเสติ อาทิกลฺยาณํ มชฺเฌกลฺยาณํ ปริโยสานกลฺยาณํ สาตฺถํ สพฺยญฺชนํ เกวลปริปุณฺณํ ปริสุทฺธํ พฺรหฺมจริยํ ปกาเสติ.}

\prose{ตํ ธมฺมํ สุณาติ คหปติ วา คหปติปุตฺโต วา อญฺญตรสฺมิํ วา กุเล ปจฺจาชาโต.\csromanspacer{} โส ตํ ธมฺมํ สุตฺวา ตถาคเต สทฺธํ ปฏิลภติ,}

\vspace{\tipitakaparskip}
\csromancenter{ทสกนิปาตปาฬิ โส เตน สทฺธาปฏิลาเภน สมนฺนาคโต อิติ ปฏิสญฺจิกฺขติ “สมฺพาโธ ฆราวาโส รชาปโถ, อพฺโภกาโส ปพฺพชฺชา, นยิทํ สุกรํ อคารํ อชฺฌาวสตา เอกนฺตปริปุณฺณํ เอกนฺตปริสุทฺธํ สงฺขลิขิตํ พฺรหฺมจริยํ จริตุํ, ยํนูนาหํ เกสมสฺสุํ โอหาเรตฺวา กาสายานิ วตฺถานิ อจฺฉาเทตฺวา อคารสฺมา อนคาริยํ ปพฺพเชยฺยนฺ”ติ.}
\csromanmatikaanchor{mtk.140}
\csromantocmarkref{subsection}{9. อุปาลิสุตฺต}{mtk.140}
\bookmark[dest={mtk.140},level=2]{9. อุปาลิสุตฺต}
\prose{\csromanfolio{422}โส อปเรน สมเยน อปฺปํ วา โภคกฺขนฺธํ ปหาย มหนฺตํ วา โภคกฺขนฺตํ ปหาย อปฺปํ วา ญาติปริวฏฺฏํ ปหาย มหนฺตํ วา ญาติปริวฑฺฑํ ปหาย เกสมสฺสุํ โอหาเรตฺวา กาสายานิ วตฺถานิ อจฺฉาเทตฺวา อคารสฺมา อนคาริยํ ปพฺพชติ.}

\prose{โส เอวํ ปพฺพชิโต สมาโน ภิกฺขูนํ สิกฺขาสาชีวสมาปนฺโน ปาณาติปาตํ ปหาย ปาณาติปาตา ปฏิวิรโต โหติ นิหิตทณฺโฑ นิหิตสตฺโถ ลชฺชี ทยาปนฺโน สพฺพปาณภูตหิตานุกมฺปี วิหรติ.}

\prose{อทินฺนาทานํ ปหาย อทินฺนาทานา ปฏิวิรโต โหติ ทินฺนาทายี ทินฺนปาฏิกงฺขี อเถเนน สุจิภูเตน อตฺตานา วิหรติ.}

\prose{อพฺรหฺมจริยํ ปหาย พฺรหฺมจารี โหติ อาราจารี วิรโต เมถุนา คามธมฺมา.}

\prose{มุสาวาทํ ปหาย มุสาวาทา ปฏิวิรโต โหติ สจฺจวาที สจฺจสนฺโธ เถโต ปจฺจยิโก อวิสํวาทโก โลกสฺส.}

\prose{ปิสุณํ วาจํ ปหาย ปิสุณาย วาจาย ปฏิวิรโต โหติ, อิโต สุตฺวา น อมุตฺร อกฺขาตา อิเมสํ เภทาย, อมุตฺร วา สุตฺวา น อิเมสํ อกฺขาตา อมูสํ เภทาย, อิติ ภินฺนานํ วา สนฺธาตา สหิตานํ วา อนุปฺปทาตา สมคฺคาราโม สมคฺครโต สมคฺคนนฺที สมคฺคกรณิํ วาจํ ภาสิตา โหติ.}

\prose{ผรุสํ วาจํ ปหาย ผรุสาย วาจาย ปฏิวิรโต โหติ, ยา สา วาจา เนลา กณฺณสุขา เปมนียา หทยงฺคมา โปรี พหุชนกนฺตา พหุชนมนาปา, ตถารูปิํ วาจํ ภาสิตา โหติ.}

\prose{สมฺผปฺปลาปํ ปหาย สมฺผปฺปลาปา ปฏิวิรโต โหติ กาลวาที ภูตวาที อตฺถวาที ธมฺมวาที วินยวาที นิธานวติํ วาจํ ภาสิตา กาเลน สาปเทสํ ปริยนฺตวติํ อตฺถสํหิตํ.}

\chapterheadpageread{(10) 5. อุปาลิวคฺค}
\prose{\csromanfolio{423}โส พีชคามภูตคามสมารมฺภา ปฏิวิรโต โหติ.\csromanspacer{} เอกภตฺติโก โหติ รตฺตูปรโต, วิรโต วิกาลโภชนา.\csromanspacer{} นจฺจคีตวาทิตวิสูกทสฺสนา ปฏิวิรโต โหติ.\csromanspacer{} มาลาคนฺธวิเลปนธารณมณฺฑนวิภูสนฏฺฐานา ปฏิวิรโต โหติ.\csromanspacer{} อุจฺจาสยนมหาสยนา ปฏิวิรโต โหติ.\csromanspacer{} ชาตรูปรชตปฏิคฺคหณา ปฏิวิรโต โหติ.\csromanspacer{} อามกธญฺญปฏิคฺคหณา ปฏิวิรโต โหติ.\csromanspacer{} อามกมํสปฏิคฺคหณา ปฏิวิรโต โหติ.\csromanspacer{} อิตฺถิกุมาริกปฏิคฺคหณา ปฏิวิรโต โหติ.\csromanspacer{} ทาสิทาส- ปฏิคฺคหณา ปฏิวิรโต โหติ.\csromanspacer{} อเชฬกปฏิคฺคหณา ปฏิวิรโต โหติ.\csromanspacer{} กุกฺกุฏสูกรปฏิคฺคหณา ปฏิวิรโต โหติ.\csromanspacer{} หตฺถิควสฺสวฬวปฏิคฺคหณา ปฏิวิรโต โหติ.\csromanspacer{} เขตฺตวตฺถุปฏิคฺคหณา ปฏิวิรโต โหติ.\csromanspacer{} ทูเตยฺยปหิณค- มนานุโยคา ปฏิวิรโต โหติ.\csromanspacer{} กยวิกฺกยา ปฏิวิรโต โหติ.\csromanspacer{} ตุลากูฏ- กํสกูฏมานกูฏา ปฏิวิรโต โหติ.\csromanspacer{} อุกฺโกฏนวญฺจนนิกติสาจิโยคา ปฏิวิรโต โหติ.\csromanspacer{} เฉทน วธ พนฺธน วิปราโมส อาโลป สหสาการา ปฏิวิรโต โหติ.}

\prose{โส สนฺตุฏฺโฐ โหติ กายปริหาริเกน จีวเรน กุจฺฉิปริหาริเกน ปิณฺฑปาเตน, เยน เยเนว ปกฺกมติ, สมาทาเยว ปกฺกมติ.\csromanspacer{} เสยฺยถาปิ นาม ปกฺขี สกุโณ เยน เยเนว เฑติ สปตฺตภาโรว เฑติ.\csromanspacer{} เอวเมวํ ภิกฺขุ สนฺตุฏฺโฐ โหติ กายปริหาริเกน จีวเรน กุจฺฉิปริหาริเกน ปิณฺฑปาเตน, เยน เยเนว ปกฺกมติ, สมาทาเยว ปกฺกมติ.\csromanspacer{} โส อิมินา อริเยน สีลกฺขนฺเธน สมนฺนาคโต อชฺฌตฺตํ อนวชฺชสุขํ ปฏิสํเวเทติ.}

\prose{โส จกฺขุนา รูปํ ทิสฺวา น นิมิตฺตคฺคาหี โหติ นานุพฺยญฺชนคฺคาหี, ยตฺวาธิกรณเมนํ จกฺขุนฺทฺริยํ อสํวุตํ วิหรนฺตํ อภิชฺฌาโทมนสฺสา ปาปกา อกุสลา ธมฺมา อนฺวาสฺสเวยฺยุํ.\csromanspacer{} ตสฺส สํวราย ปฏิปชฺชติ, รกฺขติ จกฺขุนฺทฺริยํ, จกฺขุนฺทฺริเย สํวรํ อาปชฺชติ.\csromanspacer{} โสเตน สทฺทํ สุตฺวา.\csromanspacer{} ฆาเนน คนฺธํ ฆายิตฺวา.\csromanspacer{} ชิวฺหาย รสํ สายิตฺวา.\csromanspacer{} กาเยน โผฏฺฐพฺพํ ผุสิตฺวา.\csromanspacer{} มนสา ธมฺมํ วิญฺญาย น นิมิตฺตคฺคาหี โหติ นานุพฺยญฺชนคฺคาหี, ยตฺวาธิกรณเมนํ มนินฺทฺริยํ อสํวุตํ วิหรนฺตํ อภิชฺฌาโทมนสฺสา ปาปกา อกุสลา ธมฺมา อนฺวาสฺสเวยฺยุํ, ตสฺส สํวราย ปฏิปชฺชติ, รกฺขติ มนินฺทฺริยํ, มนินฺทฺริเย สํวรํ อาปชฺชติ.\csromanspacer{} โส อิมินา อริเยน อินฺทฺริยสํวเรน สมนฺนาคโต อชฺฌตฺตํ อพฺยาเสกสุขํ ปฏิสํเวเทติ.}

\gambhira{ทสกนิปาตปาฬิ}
\prose{\csromanfolio{424}โส อภิกฺกนฺเต ปฏิกฺกนฺเต สมฺปชานการี โหติ, อาโลกิเต วิโลกิเต สมฺปชานการี โหติ, สมิญฺชิเต ปสาริเต สมฺปชานการี โหติ, สํฆาฏิปตฺตจีวรธารเณ สมฺปชานการี โหติ, อสิเต ปีเต ขายิเต สายิเต สมฺปชานการี โหติ, อุจฺจารปสฺสาวกมฺเม สมฺปชานการี โหติ, คเต ฐิเต นิสินฺเน สุตฺเต ชาคริเต ภาสิเต ตุณฺหีภาเว สมฺปชานการี โหติ.}

\prose{โส อิมินา จ อริเยน สีลกฺขนฺเธน สมนฺนาคโต, อิมินา จ อริเยน อินฺทฺริยสํวเรน สมนฺนาคโต; อิมินา จ อริเยน สติสมฺปชญฺเญน สมนฺนาคโต วิวิตฺตํ เสนาสนํ ภชติ อรญฺญํ รุกฺขมูลํ ปพฺพตํ กนฺทรํ คิริคุหํ สุสานํ วนปตฺถํ อพฺโภกาสํ ปลาลปุญฺชํ.\csromanspacer{} โส อรญฺญคโต วา รุกฺขมูลคโต วา สุญฺญาคารคโต วา นิสีทติ ปลฺลงฺกํ อาภุชิตฺวา อุชุํ กายํ ปณิธาย ปริมุขํ สติํ อุปฏฺฐเปตฺวา.}

\prose{โส อภิชฺฌํ โลเก ปหาย วิคตาภิชฺเฌน เจตสา วิหรติ, อภิชฺฌาย จิตฺตํ ปริโสเธติ.\csromanspacer{} พฺยาปาทปโทสํ ปหาย อพฺยาปนฺนจิตฺโต วิหรติ สพฺพปาณภูตหิตานุกมฺปี, พฺยาปาทปโทสา จิตฺตํ ปริโสเธติ.\csromanspacer{} ถินมิทํ ปหาย วิคตถินมิทฺโธ วิหรติ อาโลกสญฺญี สโต สมฺปชาโน, ถินมิทฺธา จิตฺตํ ปริโสเธติ.\csromanspacer{} อุทฺธจฺจกุกฺกุจฺจํ ปหาย อนุทฺธโต วิหรติ อชฺฌตฺตํ วูปสนฺตจิตฺโต, อุทฺธจฺจกุกฺกุจฺจา จิตฺตํ ปริโสเธติ.\csromanspacer{} วิจิกิจฺฉํ ปหาย ติณฺณวิจิกิจฺโฉ วิหรติ อกถํกถี กุสเลสุ ธมฺเมสุ, วิจิกิจฺฉาย จิตฺตํ ปริโสเธติ.}

\prose{โส อิเม ปญฺจ นีวรเณ ปหาย เจตโส อุปกฺกิเลเส ปญฺญาย ทุพฺพลีกรเณ วิวิจฺเจว กาเมหิ วิวิจฺจ อกุสเลหิ ธมฺเมหิ สวิตกฺกํ สวิจารํ วิเวกชํ ปีติสุขํ ปฐมํ ฌานํ อุปสมฺปชฺช วิหรติ.\csromanspacer{} ตํ กิํ มญฺญสิ อุปาลิ, นนฺวายํ วิหาโร ปุริเมหิ วิหาเรหิ อภิกฺกนฺตตโร จ ปณีตตโร จาติ.\csromanspacer{} เอวํ ภนฺเต.}

\prose{อิมมฺปิ โข อุปาลิ มม สาวกา อตฺตนิ ธมฺมํ สมฺปสฺสมานา อรญฺญวนปตฺถานิ ปนฺตานิ เสนาสนานิ ปฏิเสวนฺติ, โน จ โข ตาว อนุปฺปตฺตสทตฺถา วิหรนฺติ.}

\chapterheadpageread{(10) 5. อุปาลิวคฺค}
\prose{\csromanfolio{425}ปุน จปรํ อุปาลิ ภิกฺขุ วิตกฺกวิจารานํ วูปสมา ฯเปฯ ทุติยํ ฌานํ อุปสมฺปชฺช วิหรติ.\csromanspacer{} ตํ กิํ มญฺญสิ อุปาลิ, นนฺวายํ วิหาโร ปุริเมหิ วิหเรติ อภิกฺกนฺตโร จ ปณีตตโร จาติ.\csromanspacer{} เอวํ ภนฺเต.}

\prose{อิมมฺปิ โข อุปาลิ มม สาวกา อตฺตนิ ธมฺมํ สมฺปสฺสมานา อรญฺญวนปตฺถานิ ปนฺตานิ เสนาสนานิ ปฏิเสวนฺติ, โน จ โข ตาว อนุปฺปตฺตทสตฺถา วิหรนฺติ.}

\prose{ปุนจปรํ อุปาลิ ภิกฺขุ ปีติยา จ วิราคา ฯเปฯ ตติยํ ฌานํ อุปสมฺปชฺช วิหรติ.\csromanspacer{} ตํ กิํ มญฺญสิ อุปาลิ, นนฺวายํ วิหาโร ปุริเมหิ วิหเรหิ อภิกฺกนฺตตโร จ ปณีตตโร จาติ.\csromanspacer{} เอวํ ภนฺเต.}

\prose{อิมมฺปิ โข อุปาลิ มม สาวกา อตฺตนิ ธมฺมํ สมฺปสฺสมานา อรญฺญวนปตฺถานิ ปนฺตานิ เสนาสนานิ ปฏิเสวนฺติ, โน จ โข ตาว อนุปฺปตฺตสทตฺถา วิหรนฺติ.}

\prose{ปุน จปรํ อุปาลิ ภิกฺขุ สุขสฺส จ ปหานา ฯเปฯ จตุตฺถํ ฌานํ ฯเปฯ.}

\prose{ปุน จปรํ อุปาลิ ภิกฺขุ สพฺพโส รูปสญฺญานํ สมติกฺกมา ปฏิฆสญฺญานํ อตฺถงฺคมา นานตฺตสญฺญานํ อมนสิการา “อนนฺโต อากาโส”ติ อากาสานญฺจายตนํ อุปสมฺปชฺช วิหรติ.\csromanspacer{} ตํ กิํ มญฺญสิ อุปาลิ, นนฺวายํ วิหาโร ปุริเมหิ วิหาเรหิ อภิกฺกนฺตตโร จ ปณีตตโร จาติ.\csromanspacer{} เอวํ ภนฺเต.}

\prose{อิมมฺปิ โข อุปาลิ มม สาวกา อตฺตนิ ธมฺมํ สมฺปสฺสมานา อรญฺญวนปตฺถานิ ปนฺตานิ เสนาสนานิ ปฏิเสวนฺติ, โน จ โข ตาว อนุปฺปตฺตสทตฺถา วิหรนฺติ.}

\prose{ปุน จปรํ อุปาลิ ภิกฺขุ สพฺพโส อากาสานญฺจายตนํ สมติกฺกมฺม “อนนฺตํ วิญฺญาณนฺ”ติ วิญฺญาณญฺจายตนํ อุปสมฺปชฺช วิหรติ ฯเปฯ.}

\prose{สพฺพโส วิญฺญาณญฺจายตนํ สมติกฺกมฺม “นตฺถิ กิญฺจี”ติ อากิญฺจญฺญายตนํ อุปสมฺปชฺช วิหรติ ฯเปฯ.}

\prose{สพฺพโส อากิญฺจญฺญายตนํ สมติกฺกมฺม “สนฺตเมตํ ปณีตเมตนฺ”ติ เนวสญฺญานาสญฺญายตนํ อุปสมฺปชฺช วิหรติ.\csromanspacer{} ตํ กิํ มญฺญสิ อุปาลิ, นนฺวายํ วิหาโร ปุริเมหิ วิหาเรหิ อภิกฺกนฺตตโร จ ปณีตตโร จาติ.\csromanspacer{} เอวํ ภนฺเต.}

\gambhira{ทสกนิปาตปาฬิ}
\csromanmatikaanchor{mtk.141}
\csromanmatikaanchor{mtk.142}
\csromantocmarkref{subsection}{10. อภพฺพสุตฺต}{mtk.141}
\bookmark[dest={mtk.141},level=2]{10. อภพฺพสุตฺต}
\csromantocmarkref{section}{อุทฺทานคาถา}{mtk.142}
\bookmark[dest={mtk.142},level=1]{อุทฺทานคาถา}
\prose{\csromanfolio{426}อิมมฺปิ โข อุปาลิ มม สาวกา อตฺตนิ ธมฺมํ สมฺปสฺสมานา อรญฺญวนปตฺถานิ ปนฺตานิ เสนาสนานิ ปฏิเสวนฺติ, โน จ โข ตาว อนุปฺปตฺตสทตฺถา วิหรนฺติ.}

\prose{ปุน จปรํ อุปาลิ ภิกฺขุ สพฺพโส เนวสญฺญานาสญฺญายตนํ สมติกฺกมฺม สญฺญาเวทยิตนิโรธํ อุปสมฺปชฺช วิหรติ, ปญฺญาย จสฺส ทิสฺวา อาสวา ปริกฺขีณา โหนฺติ.\csromanspacer{} ตํ กิํ มญฺญสิ อุปาลิ, นนฺวายํ วิหาโร ปุริเมติ วิหาเรหิ อภิกฺกนฺตตโร จ ปณีตตโร จาติ.\csromanspacer{} เอวํ ภนฺเต.}

\prose{อิมมฺปิ โข อุปาลิ มม สาวกา อตฺตนิ ธมฺมํ สมฺปสฺสมานา อรญฺญวนปตฺถานิ ปนฺตานิ เสนาสนานิ ปฏิเสวนฺติ, อนุปฺปตฺตสทตฺถา จ วิหรนฺติ.\csromanspacer{} อิงฺฆ ตฺวํ อุปาลิ สํเฆ วิหราหิ, สํเฆ เต วิหรโต ผาสุ ภวิสฺสตีติ.\csromanspacer{}\csromanspacer{}นวมํ.}

\csromanheader{2}{10. อภพฺพสุตฺต}
\proseitem{100}{ทสยิเม ภิกฺขเว ธมฺเม อปฺปหาย อภพฺโพ อรหตฺตํ สจฺฉิกาตุํ.\csromanspacer{} กตเม ทส, ราคํ โทสํ โมหํ โกธํ อุปนาหํ มกฺขํ ปฬาสํ อิสฺสํ มจฺฉริยํ มานํ.\csromanspacer{} อิเม โข ภิกฺขเว ทส ธมฺเม อปฺปหาย อภพฺโพ อรหตฺตํ สจฺฉิกาตุํ.}

\prose{ทสยิเม ภิกฺขเว ธมฺเม ปหาย ภพฺโพ อรหตฺตํ สจฺฉิกาตุํ.\csromanspacer{} กตเม ทส, ราคํ โทสํ โมหํ โกธํ อุปนาหํ มกฺขํ ปฬาสํ อิสฺสํ มจฺฉริยํ มานํ.\csromanspacer{} อิเม โข ภิกฺขเว ทส ธมฺเม ปหาย ภพฺโพ อรหตฺตํ สจฺฉิกาตุนฺติ.\csromanspacer{}\csromanspacer{}ทสมํ.}

\vspace{\tipitakaparskip}
\csromancenter{อุปาลิวคฺโค ปญฺจโม.}
\titlehead{\textbf{ตสฺสุทฺทานํ}}
\csromangathameasure{กามโภคี ภยํ ทิฏฺฐิ, วชฺชิยมาหิตุตฺติยา. \\ โกกนุโท อาหุเนยฺโย, เถโร อุปาลิ อภพฺโพติ. \\ \textbf{ทุติโย ปณฺณาสโก สมตฺโต.}}
\csromangathagroup{\csromangathastanza{กามโภคี ภยํ ทิฏฺฐิ, วชฺชิยมาหิตุตฺติยา. \\ โกกนุโท อาหุเนยฺโย, เถโร อุปาลิ อภพฺโพติ. \\ \textbf{ทุติโย ปณฺณาสโก สมตฺโต.}}}

\chapterheadpageread{3. ตติยปณฺณาสก}
\csromanheader{1}{\textbf{(11) 1. สมณสญฺญาวคฺค}}
\csromanheader{2}{1. สมณสญฺญาสุตฺต}
\proseitem{101}{\csromanfolio{427}ติสฺโส อิมา ภิกฺขเว สมณสญฺญา ภาวิตา พหุลีกตา สตฺต ธมฺเม ปริปูเรนฺติ.\csromanspacer{} กตมา ติสฺโส? เววณฺณิยมฺหิ อชฺฌุปคโต, ปรปฏิพทฺธา เม ชีวิกา, อญฺโญ เม อากปฺโป กรณีโยติ.\csromanspacer{} อิมา โข ภิกฺขเว ติสฺโส สมณสญฺญา ภาวิตา พหุลีกตา สตฺต ธมฺเม ปริปูเรนฺติ.}

\prose{กตเม สตฺต? สนฺตตการี\footnote{สตตการี (สฺยา, อิ, ก)} โหติ สนฺตตวุตฺติ\footnote{สตตวุตฺติ (สฺยา, อิ)} สีเลสุ, อนภิชฺฌาลุ โหติ, อพฺยาปชฺโช โหติ, อนติมานี โหติ, สิกฺขากาโม โหติ, อิทมตฺถํติสฺส โหติ ชีวิตปริกฺขาเรสุ, อารทฺธวีริโย จ\footnote{อารทฺธวิริโย จ (สี, อิ), อารทฺธวิริโย (สฺยา)} วิหรติ.\csromanspacer{} อิมา โข ภิกฺขเว ติสฺโส สมณสญฺญา ภาวิตา พหุลีกตา อิเม สตฺต ธมฺเม ปริปูเรนฺตีติ.\csromanspacer{}\csromanspacer{}ปฐมํ.}

\csromanheader{2}{2. โพชฺฌงฺคสุตฺต}
\proseitem{102}{สตฺติเม ภิกฺขเว โพชฺฌงฺคา ภาวิตา พหุลีกตา ติสฺโส วิชฺชา ปริปูเรนฺติ.\csromanspacer{} กตเม สตฺต? สติสมฺโพชฺฌงฺโค ธมฺมวิจยสมฺโภชฺฌงฺโค วีริยสมฺโพชฺฌงฺโค ปีติสมฺโพชฺฌงฺโค ปสฺสทฺธิสมฺโพชฺฌงฺโค สมาธิสมฺโพชฺฌงฺโค อุเปกฺขาสมฺโพชฺฌงฺโค.\csromanspacer{} อิเม โข ภิกฺขเว สตฺต โพชฺฌงฺคา ภาวิตา พหุลีกตา ติสฺโส วิชฺชา ปริปูเรนฺติ.\csromanspacer{} กตมา ติสฺโส? อิธ ภิกฺขเว ภิกฺขุ อเนกวิหิตํ ปุพฺเพนิวาสํ อนุสฺสรติ.\csromanspacer{} เสยฺยถิทํ, เอกมฺปิ ชาติํ เทฺวปิ ชาติโย ติสฺโสปิ ชาติโย ฯเปฯ อิติ สาการํ ส-อุทฺเทสํ อเนกวิหิตํ ปุพฺเพนิวาสํ อนุสฺสรติ.\csromanspacer{} ทิพฺเพน จกฺขุนา วิสุทฺเธน อติกฺกนฺตมานุสเกน วฺ ยถากมฺมูปเค สตฺเต ปชานาติ.\csromanspacer{} อาสวานํ ขยา ฯเปฯ สจฺฉิกตฺวา อุปสมฺปชฺช วิหรติ.\csromanspacer{} อิเม โข ภิกฺขเว สตฺต โพชฺฌงฺคา ภาวิตา พหุลีกตา อิมา ติสฺโส วิชฺชา ปริปูเรนฺตีติ.\csromanspacer{}\csromanspacer{}ทุติยํ.}

\csromanheader{2}{ทสกนิปาตปาฬิ \textbf{3. มิจฺฉตฺตสุตฺต}}
\proseitem{103}{\csromanfolio{428}มิจฺฉตฺตํ ภิกฺขเว อาคมฺม วิราธนา โหติ โน อาราธนา.\csromanspacer{} กถญฺจ ภิกฺขเว มิจฺฉตฺตํ อาคมฺม วิราธนา โหติ โน อาราธนา? มิจฺฉาทิฏฺฐิกสฺส ภิกฺขเว มิจฺฉาสงฺกปฺโป ปโหติ, มิจฺฉาสงฺกปฺปสฺส มิจฺฉาวาจา ปโหติ, มิจฺฉาวาจสฺส มิจฺฉากมฺมนฺโต ปโหติ, มิจฺฉากมฺมนฺตสฺส มิจฺฉา-อาชีโว ปโหติ, มิจฺฉา-อาชีวสฺส มิจฺฉาวายาโม ปโหติ, มิจฺฉาวายามสฺส มิจฺฉาสติ ปโหติ, มิจฺฉาสติสฺส มิจฺฉาสมาธิ ปโหติ, มิจฺฉาสมาธิสฺส มิจฺฉาญาณํ ปโหติ, มิจฺฉาญาณิสฺส\footnote{มิจฺฉาญาณสฺส (อิ, ก)} มิจฺฉาวิมุตฺติ ปโหติ.\csromanspacer{} เอวํ โข ภิกฺขเว มิจฺฉตฺตํ อาคมฺม วิราธนา โหติ โน อาราธนา.}

\prose{สมฺมตฺตํ ภิกฺขเว อาคมฺม อาราธนา โหติ โน วิราธนา.\csromanspacer{} กถญฺจ ภิกฺขเว สมฺมตฺตํ อาคมฺม อาราธนา โหติ โน วิราธนา? สมฺมาทิฏฺฐิกสฺส ภิกฺขเว สมฺมาสงฺกปฺโป ปโหติ, สมฺมาสงฺกปฺปสฺส สมฺมาวาจา ปโหติ, สมฺมาวาจสฺส สมฺมากมฺมนฺโต ปโหติ, สมฺมากมฺมนฺตสฺส สมฺมา อาชีโว ปโหติ, สมฺมา-อาชีวสฺส สมฺมาวายาโม ปโหติ, สมฺมาวายามสฺส สมฺมาสติ ปโหติ, สมฺมาสติสฺส สมฺมาสมฺมาธิ ปโหติ, สมฺมาสมาธิสฺส สมฺมาญาณํ ปโหติ, สมฺมาญาณิสฺส\footnote{สมฺมาญาณสฺส (อิ, ก)} สมฺมาวิมุตฺติ ปโหติ.\csromanspacer{} เอวํ โข ภิกฺขเว สมฺมตฺตํ อาคมฺม อาราธนา โหติ โน วิราธนาติ.\csromanspacer{}\csromanspacer{}ตติยํ.}

\csromanheader{2}{4. พีชสุตฺต}
\proseitem{104}{\csromansymbolfootnote{*}{อํ 1. 33; อภิ 4. 367 ปิฏฺเฐสุปิ.}มิจฺฉาทิฏฺฐิกสฺส ภิกฺขเว ปุริสปุคฺคลสฺส มิจฺฉาสงฺกปฺปสฺส มิจฺฉาวาจสฺส มิจฺฉากมฺมนฺตสฺส มิจฺฉา-อาชีวสฺส มิจฺฉาวายามสฺส มิจฺฉาสติสฺส มิจฺฉาสมาธิสฺส มิจฺฉาญาณิสฺส มิจฺฉาวิมุตฺติสฺส ยญฺจ กายกมฺมํ ยถาทิฏฺฐิ สมตฺตํ สมาทินฺนํ\footnote{สมาทิณฺณํ (อิ, ก)}, ยญฺจ วจีกมฺมํ.\csromanspacer{} ยญฺจ มโนกมฺมํ ยถาทิฏฺฐิ สมตฺตํ สมาทินฺนํ, ยา จ เจตนา ยา จ ปตฺถนา โย จ ปณิธิ เย จ สงฺขารา, สพฺเพ เต ธมฺมา อนิฏฺฐาย อกนฺตาย อมนาปาย อหิตาย ทุกฺขาย สํวตฺตนฺติ.\csromanspacer{} ตํ กิสฺส เหตุ? ทิฏฺฐิ หิสฺส\footnote{ทิฏฺฐิ หิ (สี, สฺยา, อิ)} ภิกฺขเว ปาปิกา.}

\csromanmatikaanchor{mtk.143}
\csromantocmarknopage{chapter}{3. ตติยปณฺณาสก}{mtk.143}
\bookmark[dest={mtk.143},level=0]{3. ตติยปณฺณาสก}
\prose{เสยฺยถาปิ ภิกฺขเว นิมฺพพีชํ วา โกสาตกิพีชํ วา ติตฺตกาลาพุพีชํ วา อลฺลาย ปถวิยา นิกฺขิตฺตํ ยญฺเจว ปถวิรสํ อุปาทิยติ, ยญฺจ อาโปรสํ}

\csromanmatikaanchor{mtk.144}
\csromantocmarknopage{section}{(11) 1. สมณสญฺญาวคฺค}{mtk.144}
\bookmark[dest={mtk.144},level=1]{(11) 1. สมณสญฺญาวคฺค}
\vspace{\tipitakaparskip}
\csromancenter{(11) 1.\csromanspacer{} สมณสญฺญาวคฺค อุปาทิยติ, สพฺพํ ตํ ติตฺตกตฺตาย กฏุกตฺตาย อสาตตฺตาย สํวตฺตติ.\csromanspacer{} ตํ กิสฺส เหตุ? พีชํ หิ ภิกฺขเว ปาปกํ.\csromanspacer{} เอวเมวํ โข ภิกฺขเว มิจฺฉาทิฏฺฐิกสฺส ปุริสปุคฺคลสฺส มิจฺฉาสงฺกปฺปสฺส มิจฺฉาวาจสฺส มิจฺฉากมฺมนฺตสฺส มิจฺฉา-อาชีวสฺส มิจฺฉาวายามสฺส มิจฺฉาสติสฺส มิจฺฉาสมาธิสฺส มิจฺฉาญาณิสฺส มิจฺฉาวิมุตฺติสฺส ยญฺเจว กายกมฺมํ ยถาทิฏฺฐิ สมตฺตํ สมาทินฺนํ, ยญฺจ วจีกมฺมํ.\csromanspacer{} ยญฺจ มโนกมฺมํ ยถาทิฏฺฐิ สมตฺตํ สมาทินฺนํ, ยา จ เจตนา ยา จ ปตฺถนา โย จ ปณิธิ เย จ สงฺขารา, สพฺเพ เต ธมฺมา อนิฏฺฐาย อกนฺตาย อมนาปาย อหิตาย ทุกฺขาย สํวตฺตนฺติ.\csromanspacer{} ตํ กิสฺส เหตุ? ทิฏฺฐิ หิสฺส ภิกฺขเว ปาปิกา.}
\csromanmatikaanchor{mtk.145}
\csromantocmarkref{subsection}{1. สมณสญฺญาสุตฺต}{mtk.145}
\bookmark[dest={mtk.145},level=2]{1. สมณสญฺญาสุตฺต}
\prose{\csromanfolio{429}สมฺมาทิฏฺฐิกสฺส ภิกฺขเว ปุริสปุคฺคลสฺส สมฺมาสงฺกปฺปสฺส สมฺมาวาจสฺส สมฺมากมฺมนฺตสฺส สมฺมา-อาชีวสฺส สมฺมาวายมสฺส สมฺมาสติสฺส สมฺมา สมาธิสฺส สมฺมาญาณิสฺส สมฺมาวิมุตฺติสฺส ยญฺเจว กายกมฺมํ ยถาทิฏฺฐิ สมตฺตํ สมาทินฺนํ, ยญฺจ วจีกมฺมํ ยถาทิฏฺฐิ สมตฺตํ สมาทินฺนํ.\csromanspacer{} ยญฺจ มโนกมฺมํ ยถาทิฏฺฐิ สมตฺตํ สมาทินฺนํ, ยา จ เจตนา ยา จ ปตฺถนา โย จ ปณิธิ เย จ สงฺขารา, สพฺเพ เต ธมฺมา อิฏฺฐาย กนฺตาย มนาปาย หิตาย สุขาย สํวตฺตนฺติ.\csromanspacer{} ตํ กิสฺส เหตุ? ทิฏฺฐิ หิสฺส ภิกฺขเว ภทฺทิกา.}

\prose{เสยฺยถาปิ ภิกฺขเว อุจฺฉุพีชํ วา สาลิพีชํ วา มุทฺทิกาพีชํ วา อลฺลาย ปถวิยา นิกฺขิตฺตํ ยญฺจ ปถวิรสํ อุปาทิยติ, ยญฺจ อาโปรสํ อุปาทิยติ, สพฺพํ ตํ สาตตฺตาย มธุรตฺตาย อเสจนกตฺตาย สํวตฺตติ.\csromanspacer{} ตํ กิสฺส เหตุ? พีชํ หิ ภิกฺขเว ภทฺทกํ.\csromanspacer{} เอวเมวํ โข ภิกฺขเว สมฺมาทิฏฺฐิกสฺส ฯเปฯ สมฺมาวิมุตฺติสฺส ยญฺเจว กายกมฺมํ ยถาทิฏฺฐิ สมตฺตํ สมาทินฺนํ, ยญฺจ วจีกมฺมํ.\csromanspacer{} ยญฺจ มโนกมฺมํ ยถาทิฏฺฐิ สมตฺตํ สมาทินฺนํ, ยา จ เจตนา ยา จ ปตฺถนา โย จ ปณิธิ เย จ สงฺขารา, สพฺเพ เต ธมฺมา อิฏฺฐาย กนฺตาย มนาปาย หิตาย สุขาย สํวตฺตนฺติ.\csromanspacer{} ตํ กิสฺส เหตุ? ทิฏฺฐิ หิสฺส ภิกฺขเว ภทฺทิกาติ.\csromanspacer{}\csromanspacer{}จตุตฺถํ.}

\csromanheader{2}{5. วิชฺชาสุตฺต}
\csromanmatikaanchor{mtk.146}
\csromantocmarkref{subsection}{2. โพชฺฌงฺคสุตฺต}{mtk.146}
\bookmark[dest={mtk.146},level=2]{2. โพชฺฌงฺคสุตฺต}
\proseitem{105}{อวิชฺชา ภิกฺขเว ปุพฺพงฺคมา อกุสลานํ ธมฺมานํ สมาปตฺติยา อนฺวเทว อหิริกํ อโนตฺตปํ.\csromanspacer{} อวิชฺชาคตสฺส ภิกฺขเว อวิทฺทสุโน มิจฺฉาทิฏฺฐิ ปโหติ, มิจฺฉาทิฏฺฐิกสฺส มิจฺฉาสงฺกปฺโป ปโหติ, มิจฺฉาสงฺกปฺปสฺส มิจฺฉาวาจา ปโหติ, มิจฺฉาวาจสฺส มิจฺฉากมฺมนฺโต ปโหติ,}

\vspace{\tipitakaparskip}
\csromancenter{ทสกนิปาตปาฬิ มิจฺฉากมฺมนฺตสฺส มิจฺฉา-อาชีโว ปโหติ, มิจฺฉา-อาชีวสฺส มิจฺฉาวายาโม ปโหติ, มิจฺฉาวายามสฺส มิจฺฉาสติ ปโหติ, มิจฺฉาสติสฺส มิจฺฉาสมาธิ ปโหติ, มิจฺฉาสมาธิสฺส มิจฺฉาญาณํ ปโหติ, มิจฺฉาญาณิสฺส มิจฺฉาวิมุตฺติ ปโหติ.}
\csromanmatikaanchor{mtk.147}
\csromantocmarkref{subsection}{3. มิจฺฉตฺตสุตฺต}{mtk.147}
\bookmark[dest={mtk.147},level=2]{3. มิจฺฉตฺตสุตฺต}
\prose{\csromanfolio{430}วิชฺชา ภิกฺขเว ปุพฺพงฺคมา กุสลานํ ธมฺมานํ สมาปตฺติยา, อนฺวเทว หิโรตฺตปฺปํ.\csromanspacer{} วิชฺชาคตสฺส ภิกฺขเว วิทฺทสุโน สมฺมาทิฏฺฐิ ปโหติ, สมฺมาทิฏฺฐิกสฺส สมฺมาสงฺกปฺโป ปโหติ, สมฺมาสงฺกปฺปสฺส สมฺมาวาจา ปโหติ, สมฺมาวาจสฺส สมฺมากมฺมนฺโต ปโหติ, สมฺมากมฺมนฺตสฺส สมฺมา-อาชีโว ปโหติ, สมฺมา-อาชีวสฺส สมฺมาวายาโม ปโหติ, สมฺมาวายามสฺส สมฺมาสติ ปโหติ, สมฺมาสติสฺส สมฺมาสมาธิ ปโหติ, สมฺมาสมาธิสฺส สมฺมาญาณํ ปโหติ, สมฺมาญาณิสฺส สมฺมาวิมุตฺติ ปโหตีติ.\csromanspacer{}\csromanspacer{}ปญฺจมํ.}

\csromanheader{2}{6. นิชฺชรสุตฺต}
\proseitem{106}{\csromansymbolfootnote{*}{ที 3. 259 ปิฏฺเฐ.}ทสยิมานิ ภิกฺขเว นิชฺชรวตฺถูนิ.\csromanspacer{} กตมานิ ทส? สมฺมาทิฏฺฐิกสฺส ภิกฺขเว มิจฺฉาทิฏฺฐิ นิชฺชิณฺณา โหติ, เย จ มิจฺฉาทิฏฺฐิปจฺจยา อเนเก ปาปกา อกุสลา ธมฺมา สมฺภวนฺติ, เต จสฺส นิชฺชิณฺณา โหนฺติ, สมฺมาทิฏฺฐิปจฺจยา จ อเนเก กุสลา ธมฺมา ภาวนาปาริปูริํ คจฺฉนฺติ.}

\csromanmatikaanchor{mtk.148}
\csromantocmarkref{subsection}{4. พีชสุตฺต}{mtk.148}
\bookmark[dest={mtk.148},level=2]{4. พีชสุตฺต}
\prose{สมฺมาสงฺกปฺปสฺส ภิกฺขเว มิจฺฉาสงฺกปฺโป นิชฺชิณฺโณ โหติ, เย จ มิจฺฉาสงฺกปฺปปจฺจยา อเนเก ปาปกา อกุสลา ธมฺมา สมฺภวนฺติ เต จสฺส นิชฺชิณฺณา โหนฺติ, สมฺมาสงฺกปฺปปจฺจยา จ อเนเก กุสลา ธมฺมา ภาวนาปาริปูริํ คจฺฉนฺติ.}

\prose{สมฺมาวาจสฺส ภิกฺขเว มิจฺฉาวาจา นิชฺชิณฺณา โหติ, เย จ มิจฺฉาวาจาปจฺจยา อเนเก ปาปกา อกุสลา ธมฺมา สมฺภวนฺติ, เต จสฺส นิชฺชิณฺณา โหนฺติ, สมฺมาวาจาปจฺจยา จ อเนเก กุสลา ธมฺมา ภาวนาปาริปูริํ คจฺฉนฺติ.}

\prose{สมฺมากมฺมนฺตสฺส ภิกฺขเว มิจฺฉากมฺมนฺโต นิชฺชิณฺโณ โหติ, เย จ มิจฺฉากมฺมนฺตปจฺจยา อเนเก ปาปกา อกุสลา ธมฺมา สมฺภวนฺติ, เต จสฺส นิชฺชิณฺณา โหนฺติ, สมฺมากมฺมนฺตปจฺจยา จ อเนเก กุสลา ธมฺมา ภาวนาปาริปูริํ คจฺฉนฺติ.}

\chapterheadpageread{(11) 1. สมณสญฺญาวคฺค}
\prose{\csromanfolio{431}สมฺมา-อาชีวสฺส ภิกฺขเว มิจฺฉา-อาชีโว นิชฺชิณฺโณ โหติ.\csromanspacer{} เย จ มิจฺฉา-อาชีวปจฺจยา อเนเก ปาปกา อกุสลา ธมฺมา สมฺภวนฺติ, เต จสฺส นิชฺชิณฺณา โหนฺติ.\csromanspacer{} สมฺมา-อาชีวปจฺจยา จ อเนเก กุสลา ธมฺมา ภาวนาปาริปูริํ คจฺฉนฺติ.}

\prose{สมฺมาวายามสฺส ภิกฺขเว มิจฺฉาวายาโม นิชฺชิณฺโณ โหติ.\csromanspacer{} เย จ มิจฺฉาวายามปจฺจยา อเนเก ปาปกา อกุสลา ธมฺมา สมฺภวนฺติ, เต จสฺส นิชฺชิณฺณา โหนฺติ.\csromanspacer{} สมฺมาวายามปจฺจยา จ อเนเก กุสลา ธมฺมา ภาวนาปาริปูริํ คจฺฉนฺติ.}

\csromanmatikaanchor{mtk.149}
\csromantocmarkref{subsection}{5. วิชฺชาสุตฺต}{mtk.149}
\bookmark[dest={mtk.149},level=2]{5. วิชฺชาสุตฺต}
\prose{สมฺมาสติสฺส ภิกฺขเว มิจฺฉาสติ นิชฺชิณฺณา โหติ.\csromanspacer{} เย จ มิจฺฉาสติปจฺจยา อเนเก ปาปกา อกุสลา ธมฺมา สมฺภวนฺติ, เต จสฺส นิชฺชิณฺณา โหนฺติ.\csromanspacer{} สมฺมาสติปจฺจยา จ อเนเก กุสลา ธมฺมา ภาวนาปาริปูริํ คจฺฉนฺติ.}

\prose{สมฺมาสมาธิสฺส ภิกฺขเว มิจฺฉาสมาธิ นิชฺชิณฺโณ โหติ.\csromanspacer{} เย จ มิจฺฉาสมาธิปจฺจยา อเนเก ปาปกา อกุสลา ธมฺมา สมฺภวนฺติ, เต จสฺส นิชฺชิณฺณา โหนฺติ.\csromanspacer{} สมฺมาสมาธิปจฺจยา จ อเนเก กุสลา ธมฺมา ภาวนาปาริปูริํ คจฺฉนฺติ.}

\prose{สมฺมาญาณิสฺส ภิกฺขเว มิจฺฉาญาณํ นิชฺชิณฺณํ โหติ.\csromanspacer{} เย จ มิจฺฉาญาณปจฺจยา อเนเก ปาปกา อกุสลา ธมฺมา สมฺภวนฺติ, เต จสฺส นิชฺชิณฺณา โหนฺติ.\csromanspacer{} สมฺมาญาณปจฺจยา จ อเนเก กุสลา ธมฺมา ภาวนาปาริปูริํ คจฺฉนฺติ.}

\prose{สมฺมาวิมุตฺติสฺส ภิกฺขเว มิจฺฉาวิมุตฺติ นิชฺชิณฺณา โหติ.\csromanspacer{} เย จ มิจฺฉาวิมุตฺติปจฺจยา อเนเก ปาปกา อกุสลา ธมฺมา สมฺภวนฺติ, เต จสฺส นิชฺชิณฺณา โหนฺติ.\csromanspacer{} สมฺมาวิมุตฺติปจฺจยา จ อเนเก กุสลา ธมฺมา ภาวนาปาริปูริํ คจฺฉนฺติ.\csromanspacer{} อิมานิ โข ภิกฺขเว ทส นิชฺชรวตฺถูนีติ.\csromanspacer{}\csromanspacer{}ฉฏฺฐํ.}

\csromanmatikaanchor{mtk.150}
\csromantocmarkref{subsection}{6. นิชฺชรสุตฺต}{mtk.150}
\bookmark[dest={mtk.150},level=2]{6. นิชฺชรสุตฺต}
\csromanheader{2}{7. โธวนสุตฺต}
\proseitem{107}{อตฺถิ ภิกฺขเว ทกฺขิเณสุ ชนปเทสุ โธวนํ นาม, ตตฺถ โหติ อนฺนมฺปิ ปานมฺปิ ขชฺชมฺปิ โภชฺชมฺปิ เลยฺยมฺปิ เปยฺยมฺปิ นจฺจมฺปิ คีตมฺปิ วาทิตมฺปิ.\csromanspacer{} อตฺเถตํ ภิกฺขเว โธวนํ, เนตํ นตฺถีติ วทามิ.\csromanspacer{} ตญฺจ โข เอตํ ภิกฺขเว โธวนํ หีนํ คมฺมํ โปถุชฺชนิกํ อนริยํ อนตฺถสํหิตํ น นิพฺพิทาย น}

\vspace{\tipitakaparskip}
\csromancenter{ทสกนิปาตปาฬิ วิราคาย น นิโรธาย น อุปสมาย น อภิญฺญาย น สมฺโพธาย น นิพฺพานาย สํวตฺตติ.}
\prose{\csromanfolio{432}อหญฺจ โข ภิกฺขเว อริยํ โธวนํ เทเสสฺสามิ, ยํ โธวนํ เอกนฺตนิพฺพิทาย วิราคาย นิโรธาย อุปสมาย อภิญฺญาย สมฺโพธาย นิพฺพานาย สํวตฺตติ, ยํ โธวนํ อาคมฺม ชาติธมฺมา สตฺตา ชาติยา ปริมุจฺจนฺติ, ชราธมฺมา สตฺตา ชราย ปริมุจฺจนฺติ, มรณธมฺมา สตฺตา มรเณน ปริมุจฺจนฺติ, โสกปริเทวทุกฺขโทมนสฺสุปายาสธมฺมา สตฺตา โสกปริเทวทุกฺขโทมนสฺสุปายาเสหิ ปริมุจฺจนฺติ.\csromanspacer{} ตํ สุณาถ สาธุกํ มนสิ กโรถ, ภาสิสฺสามีติ. “เอวํ ภนฺเต”ติ โข เต ภิกฺขู ภควโต ปจฺจสฺโสสุํ.\csromanspacer{} ภควา เอตทโวจ\csromandash{}}

\prose{กตมญฺจ ตํ ภิกฺขเว อริยํ โธวนํ, (ยํ โธวนํ)\footnote{( ) นตฺถิ สฺยามโปตฺถเก.} เอกนฺตนิพฺพิทาย วิราคาย นิโรธาย อุปสมาย อภิญฺญาย สมฺโพธาย นิพฺพานาย สํวตฺตติ, ยํ โธวนํ อาคมฺม ชาติธมฺมา สตฺตา ชาติยา ปริมุจฺจนฺติ, ชราธมฺมา สตฺตา ชราย ปริมุจฺจนฺติ, มรณธมฺมา สตฺตา มรเณน ปริมุจฺจนฺติ, โสกปริเทวทุกฺขโทมนสฺสุปายาสธมฺมา สตฺตา โสกปริเทวทุกฺขโทมนสฺสุปายาเสหิ ปริมุจฺจนฺติ.}

\prose{สมฺมาทิฏฺฐิกสฺส ภิกฺขเว มิจฺฉาทิฏฺฐิ นิทฺโธตา โหติ.\csromanspacer{} เย จ มิจฺฉาทิฏฺฐิปจฺจยา อเนเก ปาปกา อกุสลา ธมฺมา สมฺภวนฺติ, เต จสฺส นิทฺโธตา โหนฺติ.\csromanspacer{} สมฺมาทิฏฺฐิปจฺจยา จ อเนเก กุสลา ธมฺมา ภาวนาปาริปูริํ คจฺฉนฺติ.}

\prose{สมฺมาสงฺกปฺปสฺส ภิกฺขเว มิจฺฉาสงฺกปฺโป นิทฺโธโต โหติ ฯเปฯ.\csromanspacer{} สมฺมาวาจสฺส ภิกฺขเว มิจฺฉาวาจา นิทฺโธตา โหติ.\csromanspacer{} สมฺมากมฺมนฺตสฺส ภิกฺขเว มิจฺฉากมฺมนฺโต นิทฺโธโต โหติ.\csromanspacer{} สมฺมา-อาชีวสฺส ภิกฺขเว มิจฺฉา-อาชีโว นิทฺโธโต โหติ.\csromanspacer{} สมฺมาวายามสฺส ภิกฺขเว มิจฺฉาวายาโม นิทฺโธโต โหติ.\csromanspacer{} สมฺมาสติสฺส ภิกฺขเว มิจฺฉาสติ นิทฺโธตา โหติ.\csromanspacer{} สมฺมาสมาธิสฺส ภิกฺขเว มิจฺฉาสมาธิ นิทฺโธโต โหติ.\csromanspacer{} สมฺมาญาณิสฺส ภิกฺขเว มิจฺฉาญาณํ นิทฺโธตํ โหติ ฯเปฯ.}

\chapterheadpageread{(11) 1. สมณสญฺญาวคฺค}
\prose{\csromanfolio{433}สมฺมาวิมุตฺติสฺส ภิกฺขเว มิจฺฉาวิมุตฺติ นิทฺโธตา โหติ.\csromanspacer{} เย จ มิจฺฉาวิมุตฺติปจฺจยา อเนเก ปาปกา อกุสลา ธมฺมา สมฺภวนฺติ, เต จสฺส นิทฺโธตา โหนฺติ.\csromanspacer{} สมฺมาวิมุตฺติปจฺจยา จ อเนเก กุสลา ธมฺมา ภาวนาปาริปูริํ คจฺฉนฺติ.\csromanspacer{} อิทํ โข ตํ ภิกฺขเว อริยํ โธวนํ เอกนฺตนิพฺพิทาย วิราคาย นิโรธาย อุปสมาย อภิญฺญาย สมฺโพธาย นิพฺพานาย สํวตฺตติ.\csromanspacer{} ยํ โธวนํ อาคมฺม ชาติธมฺมา สตฺตา ชาติยา ปริมุจฺจนฺติ, ชราธมฺมา สตฺตา ชราย ปริมุจฺจนฺติ, มรณธมฺมา สตฺตา มรเณน ปริมุจฺจนฺติ, โสก ปริเทว ทุกฺข โทมนสฺสุปายาสธมฺมา สตฺตา โสก ปริเทว ทุกฺขโทมนสฺสุปายาเสหิ ปริมุจฺจนฺตีติ.\csromanspacer{}\csromanspacer{}สตฺตมํ.}

\csromanheader{2}{8. ติกิจฺฉกสุตฺต}
\proseitem{108}{ติกิจฺฉกา ภิกฺขเว วิเรจนํ เทนฺติ ปิตฺตสมุฏฺฐานานมฺปิ อาพาธานํ ปฏิฆาตาย เสมฺหสมุฏฺฐานานมฺปิ อาพาธานํ ปฏิฆาตาย วาตสมุฏฺฐานานมฺปิ อาพาธานํ ปฏิฆาตาย.\csromanspacer{} อตฺเถตํ ภิกฺขเว วิเรจนํ, เนตํ นตฺถีติ วทามิ.\csromanspacer{} ตญฺจ โข เอตํ ภิกฺขเว วิเรจนํ สมฺปชฺชติปิ วิปชฺชติปิ.}

\prose{อหญฺจ โข ภิกฺขเว อริยํ วิเรจนํ เทเสสฺสามิ, ยํ วิเรจนํ สมฺปชฺชติเยว, โน วิปชฺชติ.\csromanspacer{} ยํ วิเรจนํ อาคมฺม ชาติธมฺมา สตฺตา ชาติยา ปริมุจฺจนฺติ, ชราธมฺมา สตฺตา ชราย ปริมุจฺจนฺติ, มรณธมฺมา สตฺตา มรเณน ปริมุจฺจนฺติ, โสก ปริเทว ทุกฺข โทมนสฺสุปายาสธมฺมา สตฺตา โสก ปริเทว ทุกฺข โทมนสฺสุปายาเสหิ ปริมุจฺจนฺติ.\csromanspacer{} ตํ สุณาถ สาธุกํ มนสิ กโรถ, ภาสิสฺสามีติ. “เอวํ ภนฺเต”ติ โข เต ภิกฺขู ภควโต ปจฺจสฺโสสุํ.\csromanspacer{} ภควา เอตทโวจ\csromandash{}}

\csromanmatikaanchor{mtk.151}
\csromantocmarkref{subsection}{7. โธวนสุตฺต}{mtk.151}
\bookmark[dest={mtk.151},level=2]{7. โธวนสุตฺต}
\prose{กตมญฺจ ตํ ภิกฺขเว อริยํ วิเรจนํ, ยํ วิเรจนํ สมฺปชฺชติเยว, โน วิปชฺชติ.\csromanspacer{} ยํ วิเรจนํ อาคมฺม ชาติธมฺมา สตฺตา ชาติยา ปริมุจฺจนฺติ, ชราธมฺมา สตฺตา ชราย ปริมุจฺจนฺติ, มรณธมฺมา สตฺตา มรเณน ปริมุจฺจนฺติ, โสก ปริเทว ทุกฺข โทมนสฺสุปายาสธมฺมา สตฺตา โสกปริเทวทุกฺขโทมนสฺสุปายาเสหิ ปริมุจฺจนฺติ.}

\prose{สมฺมาทิฏฺฐิกสฺส ภิกฺขเว มิจฺฉาทิฏฺฐิ วิริตฺตา โหติ.\csromanspacer{} เย จ มิจฺฉาทิฏฺฐิปจฺจยา อเนเก ปาปกา อกุสลา ธมฺมา สมฺภวนฺติ, เต จสฺส วิริตฺตา โหนฺติ.\csromanspacer{} สมฺมาทิฏฺฐิปจฺจยา จ อเนเก กุสลา ธมฺมา ภาวนาปาริปูริํ คจฺฉนฺติ.}

\gambhira{ทสกนิปาตปาฬิ}
\prose{\csromanfolio{434}สมฺมาสงฺกปฺปสฺส ภิกฺขเว มิจฺฉาสงฺกปฺโป วิริตฺโต โหติ ฯเปฯ.\csromanspacer{} สมฺมาวาจสฺส ภิกฺขเว มิจฺฉาวาจา วิริตฺตา โหติ.\csromanspacer{} สมฺมากมฺมนฺตสฺส ภิกฺขเว มิจฺฉากมฺมนฺโต วิริตฺโต โหติ.\csromanspacer{} สมฺมา-อาชีวสฺส ภิกฺขเว มิจฺฉา-อาชีโว วิริตฺโต โหติ.\csromanspacer{} สมฺมาวายามสฺส ภิกฺขเว มิจฺฉาวายาโม วิริตฺโต โหติ.\csromanspacer{} สมฺมาสติสฺส ภิกฺขเว มิจฺฉาสติ วิริตฺตา โหติ.\csromanspacer{} สมฺมาสมาธิสฺส ภิกฺขเว มิจฺฉาสมาธิ วิริตฺโต โหติ.\csromanspacer{} สมฺมาญาณิสฺส ภิกฺขเว มิจฺฉาญาณํ วิริตฺตํ โหติ ฯเปฯ.}

\prose{สมฺมาวิมุตฺติสฺส ภิกฺขเว มิจฺฉาวิมุตฺติ วิริตฺตา โหติ.\csromanspacer{} เย จ มิจฺฉาวิมุตฺติปจฺจยา อเนเก ปาปกา อกุสลา ธมฺมา สมฺภวนฺติ, เต จสฺส วิริตฺตา โหนฺติ.\csromanspacer{} สมฺมาวิมุตฺติปจฺจยา จ อเนเก กุสลา ธมฺมา ภาวนาปาริปูริํ คจฺฉนฺติ.\csromanspacer{} อิทํ โข ตํ ภิกฺขเว อริยํ วิเรจนํ, ยํ วิเรจนํ สมฺปชฺชติเยว, โน วิปชฺชติ.\csromanspacer{} ยํ วิเรจนํ อาคมฺม ชาติธมฺมา สตฺตา ชาติยา ปริมุจฺจนฺติ ฯเปฯ โสก ปริเทว ทุกฺข โทมนสฺสุปายาเสหิ ปริมุจฺจนฺตี.\csromanspacer{}\csromanspacer{}อฏฺฐมํ.}

\csromanheader{2}{9. วมนสุตฺต}
\proseitem{109}{ติกิจฺฉกา ภิกฺขเว วมนํ เทนฺติ ปิตฺตสมุฏฺฐานานมฺปิ อาพาธานํ ปฏิฆาตาย เสมฺหสมุฏฺฐานานมฺปิ อาพาธานํ ปฏิฆาตาย วาตสมุฏฺฐานานมฺปิ อาพาธานํ ปฏิฆาตาย.\csromanspacer{} อตฺเถตํ ภิกฺขเว วมนํ, เนตํ นตฺถีติ วทามิ.\csromanspacer{} ตญฺจ โข เอตํ ภิกฺขเว วมนํ สมฺปชฺชติปิ วิปชฺชติปิ.}

\prose{อหญฺจ โข ภิกฺขเว อริยํ วมนํ เทเสสฺสามิ, ยํ วมนํ สมฺปชฺชติเยว, โน วิปชฺชติ.\csromanspacer{} ยํ วมนํ อาคมฺม ชาติธมฺมา สตฺตา ชาติยา ปริมุจฺจนฺติ, ชราธมฺมา สตฺตา ชราย ปริมุจฺจนฺติ, มรณธมฺมา สตฺตา มรเณน ปริมุจฺจนฺติ, โสก ปริเทว ทุกฺข โทมนสฺสุปายาสธมฺมา สตฺตา โสกปริเทว ทุกฺขโทมนสฺสุปายาเสหิ ปริมุจฺจนฺติ.\csromanspacer{} ตํ สุณาถ ฯเปฯ.}

\prose{กตมญฺจ ตํ ภิกฺขเว อริยํ วมนํ, ยํ วมนํ สมฺปชฺชติเยว, โน วิปชฺชติ.\csromanspacer{} ยํ วมนํ อาคมฺม ชาติธมฺมา สตฺตา ชาติยา ปริมุจฺจนฺตินฺ ฯเปฯ โสกปริเทว ทุกฺขโทมนสฺสุปายาสธมฺมา สตฺตา โสกปริเทวทุกฺขโทมนสฺสุปายาเสหิ ปริมุจฺจนฺติ.}

\csromanmatikaanchor{mtk.152}
\csromantocmarkref{subsection}{8. ติกิจฺฉกสุตฺต}{mtk.152}
\bookmark[dest={mtk.152},level=2]{8. ติกิจฺฉกสุตฺต}
\chapterheadpageread{(11) 1. สมณสญฺญาวคฺค}
\prose{\csromanfolio{435}สมฺมาทิฏฺฐิกสฺส ภิกฺขเว มิจฺฉาทิฏฺฐิ วนฺตา โหติ.\csromanspacer{} เย จ มิจฺฉาทิฏฺฐิปจฺจยา อเนเก ปาปกา อกุสลา ธมฺมา สมฺภวนฺติ, เต จสฺส วนฺตา โหนฺติ.\csromanspacer{} สมฺมาทิฏฺฐิปจฺจยา จ อเนเก กุสลา ธมฺมา ภาวนาปาริปูริํ คจฺฉนฺติ.}

\prose{สมฺมาสงฺกปฺปสฺส ภิกฺขเว มิจฺฉาสงฺกปฺโป วนฺโต โหติ ฯเปฯ.\csromanspacer{} สมฺมาวาจสฺส ภิกฺขเว มิจฺฉาวาจา วนฺตา โหติ.\csromanspacer{} สมฺมากมฺมนฺตสฺส ภิกฺขเว มิจฺฉากมฺมนฺโต วนฺโต โหติ.\csromanspacer{} สมฺมา-อาชีวสฺส ภิกฺขเว มิจฺฉา-อาชีโว วนฺโต โหติ.\csromanspacer{} สมฺมาวายามสฺส ภิกฺขเว มิจฺฉาวายาโม วนฺโต โหติ.\csromanspacer{} สมฺมาสติสฺส ภิกฺขเว มิจฺฉาสติ วนฺตา โหติ.\csromanspacer{} สมฺมาสมาธิสฺส ภิกฺขเว มิจฺฉาสมาธิ วนฺโต โหติ.\csromanspacer{} สมฺมาญาณิสฺส ภิกฺขเว มิจฺฉาญาณํ วนฺตํ โหติ ฯเปฯ.}

\prose{สมฺมาวิมุตฺติสฺส ภิกฺขเว มิจฺฉาวิมุตฺติ วนฺตา โหติ.\csromanspacer{} เย จ มิจฺฉาวิมุตฺติปจฺจยา อเนเก ปาปกา อกุสลา ธมฺมา สมฺภวนฺติ, เต จสฺส วนฺตา โหนฺติ.\csromanspacer{} สมฺมาวิมุตฺติปจฺจยา จ อเนเก กุสลา ธมฺมา ภาวนาปาริปูริํ คจฺฉนฺติ.\csromanspacer{} อิทํ โข ตํ ภิกฺขเว อริยํ วมนํ, ยํ วมนํ สมฺปชฺชติเยว, โน วิปชฺชติ.\csromanspacer{} ยํ วมนํ อาคมฺม ชาติธมฺมา สตฺตา ชาติยา ปริมุจฺจนฺติ ฯเปฯ โสกปริเทวทุกฺขโทมนสฺสุปายาเสหิ ปริมุจฺจนฺตีติ.\csromanspacer{}\csromanspacer{}นวมํ.}

\csromanheader{2}{10. นิทฺธมนียสุตฺต}
\proseitem{110}{ทสยิเม ภิกฺขเว นิทฺธมนียา ธมฺมา.\csromanspacer{} กตเม ทส? สมฺมาทิฏฺฐิกสฺส ภิกฺขเว มิจฺฉาทิฏฺฐิ นิทฺธนฺตา โหติ.\csromanspacer{} เย จ มิจฺฉาทิฏฺฐิปจฺจยา อเนเก ปาปกา อกุสลา ธมฺมา สมฺภวนฺติ, เต จสฺส นิทฺธนฺตา โหนฺติ.\csromanspacer{} สมฺมาทิฏฺฐิปจฺจยา จ อเนเก กุสลา ธมฺมา ภาวนาปาริปูริํ คจฺฉนฺติ.}

\prose{สมฺมาสงฺกปฺปสฺส ภิกฺขเว มิจฺฉาสงฺกปฺโป นิทฺธนฺโต โหติ ฯเปฯ สมฺมาวาจสฺส ภิกฺขเว มิจฺฉาวาจา นิทฺธนฺตา โหติ.\csromanspacer{} สมฺมากมฺมนฺตสฺส ภิกฺขเว มิจฺฉากมฺมนฺโต นิทฺธนฺโต โหติ.\csromanspacer{} สมฺมา-อาชีวสฺส ภิกฺขเว มิจฺฉา-อาชีโว นิทฺธนฺโต โหติ.\csromanspacer{} สมฺมาวายามสฺส ภิกฺขเว มิจฺฉาวายาโม นิทฺธนฺโต โหติ.\csromanspacer{} สมฺมาสติสฺส ภิกฺขเว มิจฺฉาสติ นิทฺธนฺโต โหติ.\csromanspacer{} สมฺมาสมาธิสฺส ภิกฺขเว มิจฺฉาสมาธิ นิทฺธนฺโต โหติ.\csromanspacer{} สมฺมาญาณิสฺส ภิกฺขเว มิจฺฉาญาณํ นิทฺธนฺตํ โหติ.}

\vspace{\tipitakaparskip}
\csromancenter{ทสกนิปาตปาฬิ สมฺมาวิมุตฺติสฺส ภิกฺขเว มิจฺฉาวิมุตฺติ นิทฺธนฺตา โหติ.\csromanspacer{} เย จ มิจฺฉาวิมุตฺติปจฺจยา อเนเก ปาปกา อกุสลา ธมฺมา สมฺภวนฺติ, เต จสฺส นิทฺธนฺตา โหนฺติ.\csromanspacer{} สมฺมาวิมุตฺติปจฺจยา จ อเนเก กุสลา ธมฺมา ภาวนาปาริปูริํ คจฺฉนฺติ.\csromanspacer{} อิเม โข ภิกฺขเว ทส นิทฺธมนียา ธมฺมาติ.\csromanspacer{}\csromanspacer{}ทสมํ.}
\csromanmatikaanchor{mtk.153}
\csromantocmarkref{subsection}{9. วมนสุตฺต}{mtk.153}
\bookmark[dest={mtk.153},level=2]{9. วมนสุตฺต}
\csromanheader{2}{11. ปฐม-อเสขสุตฺต}
\csromanmatikaanchor{mtk.154}
\csromanmatikaanchor{mtk.157}
\csromantocmarkref{subsection}{10. นิทฺธมนียสุตฺต}{mtk.154}
\bookmark[dest={mtk.154},level=2]{10. นิทฺธมนียสุตฺต}
\proseitem{111}{\csromanfolio{436}อถ โข อญฺญตโร ภิกฺขุ เยน ภควา เตนุปสงฺกมิ, อุปสงฺกมิตฺวา ภควนฺตํ อภิวาเทตฺวา เอกมนฺตํ นิสีทิ, เอกมนฺตํ นิสินฺโน โข โส ภิกฺขุ ภควนฺตํ เอตทโวจ\csromandash{}}

\prose{“อเสโข อเสโข”ติ ภนฺเต วุจฺจติ, กิตฺตาวตา ภนฺเต ภิกฺขุ อเสโข โหตีติ.\csromanspacer{} อิธ ภิกฺขุ ภิกฺขุ อเสขาย สมฺมาทิฏฺฐิยา สมนฺนาคโต โหติ, อเสเขน สมฺมาสงฺกปฺเปน สมนฺนาคโต โหติ, อเสขาย สมฺมาวาจาย สมนฺนาคโต โหติ, อเสเขน สมฺมา กมฺมนฺเตน สมนฺนาคโต โหติ, อเสเขน สมฺมา-อาชีเวน สมนฺนาคโต โหติ, อเสเขน สมฺมาวายาเมน สมนฺนาคโต โหติ, อเสขาย สมฺมาสติยา สมนฺนาคโต โหติ, อเสเขน สมฺมาสมาธินา สมนฺนาคโต โหติ, อเสเขน สมฺมาญาเณน สมนฺนาคโต โหติ, อเสขาย สมฺมาวิมุตฺติยา สมนฺนาคโต โหติ.\csromanspacer{} เอวํ โข ภิกฺขุ ภิกฺขุ อเสโข โหตีติ.\csromanspacer{}\csromanspacer{}เอกาทสมํ.}

\csromanheader{2}{12. ทุติย-อเสขสุตฺต}
\proseitem{112}{ทสยิเม ภิกฺขเว อเสขิยา ธมฺมา.\csromanspacer{} กตเม ทส? อเสขา สมฺมาทิฏฺฐิ, อเสโข สมฺมาสงฺกปฺโป, อเสขา สมฺมาวาจา, อเสโข สมฺมากมฺมนฺโต, อเสโข สมฺมา-อาชีโว, อเสโข สมฺมาวายาโม, อเสขา สมฺมาสติ, อเสโข สมฺมาสมาธิ, อเสขํ สมฺมาญาณํ, อเสขา สมฺมาวิมุตฺติ.\csromanspacer{} อิเม โข ภิกฺขเว ทส อเสขิยา ธมฺมาติ.\csromanspacer{}\csromanspacer{}ทฺวาทสมํ.\csromanspacer{} สมณสญฺญาวคฺโค ปฐโม.}

\titlehead{\textbf{ตสฺสุทฺทานํ}}
\csromangathameasure{สญฺญา โพชฺฌงฺคา มิจฺฉตฺตํ, พีชํ วิชฺชาย นิชฺชรํ. \\ โธวนํ ติกิจฺฉา วมนํ, นิทฺธวนํ เทฺว อเสขาติ.}
\csromangathagroup{\csromangathastanza{สญฺญา โพชฺฌงฺคา มิจฺฉตฺตํ, พีชํ วิชฺชาย นิชฺชรํ. \\ โธวนํ ติกิจฺฉา วมนํ, นิทฺธวนํ เทฺว อเสขาติ.}}

\chapterheadpageread{\textbf{(12) 2. ปจฺโจโรหณิวคฺค} \textbf{(12) 2. ปจฺโจโรหณิวคฺค}}
\csromanheader{2}{1. ปฐม-อธมฺมสุตฺต}
\proseitem{113}{\csromanfolio{437}\csromansymbolfootnote{*}{อุปริ 464 ปิฏฺเฐปิ.}อธมฺโม จ ภิกฺขเว เวทิตพฺโพ อนตฺโถ จ, ธมฺโม จ เวทิตพฺโพ อตฺโถ จ, อธมฺมญฺจ วิทิตฺวา อนตฺถญฺจ, ธมฺมญฺจ วิทิตฺวา อตฺถญฺจ ยถา ธมฺโม ยถา อตฺโถ, ตถา ปฏิปชฺชิตพฺพํ.}

\prose{กตโม จ ภิกฺขเว อธมฺโม จ อนตฺโถ จ? มิจฺฉาทิฏฺฐิ มิจฺฉาสงฺกปฺโป มิจฺฉาวาจา มิจฺฉากมฺมนฺโต มิจฺฉา-อาชีโว มิจฺฉาวายาโม มิจฺฉาสติ มิจฺฉาสมาธิ มิจฺฉาญาณํ มิจฺฉาวิมุตฺติ.\csromanspacer{} อยํ วุจฺจติ ภิกฺขเว อธมฺโม จ อนตฺโถ จ.}

\prose{กตโม จ ภิกฺขเว ธมฺโม จ อตฺโถ จ? สมฺมาทิฏฺฐิ สมฺมาสงฺกปฺโป สมฺมาวาจา สมฺมากมฺมนฺโต สมฺมา-อาชีโว สมฺมาวายาโม สมฺมาสติ สมฺมาสมาธิ สมฺมาญาณํ สมฺมาวิมุตฺติ.\csromanspacer{} อยํ วุจฺจติ ภิกฺขเว ธมฺโม จ อตฺโถ จ.}

\csromanmatikaanchor{mtk.155}
\csromantocmarkref{subsection}{11. ปฐม-อเสขสุตฺต}{mtk.155}
\bookmark[dest={mtk.155},level=2]{11. ปฐม-อเสขสุตฺต}
\prose{“อธมฺโม จ ภิกฺขเว เวทิตพฺโพ อนตฺโถ จ, ธมฺโม จ เวทิตพฺโพ อตฺโถ จ, อธมฺมญฺจ วิทิตฺวา อนตฺถญฺจ, ธมฺมญฺจ วิทิตฺวา อตฺถญฺจ ยถา ธมฺโม ยถา อตฺโถ, ตถา ปฏิปชฺชิตพฺพนฺ”ติ อิติ ยํ ตํ วุตฺตํ, อิทเมตํ ปฏิจฺจ วุตฺตนฺติ.\csromanspacer{}\csromanspacer{}ปฐมํ.}

\csromanheader{2}{2. ทุติย-อธมฺมสุตฺต}
\proseitem{114}{อธมฺโม จ ภิกฺขเว เวทิตพฺโพ ธมฺโม จ, อนตฺโถ จ เวทิตพฺโพ อตฺโถ จ, อธมฺมญฺจ วิทิตฺวา ธมฺมญฺจ, อนตฺถญฺจ วิทิตฺวา อตฺถญฺจ ยถา ธมฺโม ยถา อตฺโถ, ตถา ปฏิปชฺชิตพฺพํ.}

\csromanmatikaanchor{mtk.156}
\csromantocmarkref{subsection}{12. ทุติย-อเสขสุตฺต}{mtk.156}
\bookmark[dest={mtk.156},level=2]{12. ทุติย-อเสขสุตฺต}
\csromantocmarkref{section}{อุทฺทานคาถา}{mtk.157}
\bookmark[dest={mtk.157},level=1]{อุทฺทานคาถา}
\prose{กตโม จ ภิกฺขเว อธมฺโม, กตโม จ ธมฺโม, กตโม จ อนตฺโถ, กตโม จ อตฺโถ.}

\prose{มิจฺฉาทิฏฺฐิ ภิกฺขเว อธมฺโม, สมฺมาทิฏฺฐิ ธมฺโม.\csromanspacer{} เย จ มิจฺฉาทิฏฺฐิปจฺจยา อเนเก ปาปกา อกุสลา ธมฺมา สมฺภวนฺติ, อยํ อนตฺโถ.\csromanspacer{} สมฺมาทิฏฺฐิปจฺจยา จ อเนเก กุสลา ธมฺมา ภาวนาปาริปูริํ คจฺฉนฺติ, อยํ อตฺโถ.}

\gambhira{ทสกนิปาตปาฬิ}
\prose{\csromanfolio{438}มิจฺฉาสงฺกปฺโป ภิกฺขเว อธมฺโม, สมฺมาสงฺกปฺโป ธมฺโม.\csromanspacer{} เย จ มิจฺฉาสงฺกปฺปปจฺจยา อเนเก ปาปกา อกุสลา ธมฺมา สมฺภวนฺติ, อยํ อนตฺโถ.\csromanspacer{} สมฺมาสงฺกปฺปปจฺจยา จ อเนเก กุสลา ธมฺมา ภาวนาปาริปูริํ คจฺฉนฺติ, อยํ อตฺโถ.}

\prose{มิจฺฉาวาจา ภิกฺขเว อธมฺโม, สมฺมาวาจา ธมฺโม.\csromanspacer{} เย จ มิจฺฉาวาจาปจฺจยา อเนเก ปาปกา อกุสลา ธมฺมา สมฺภวนฺติ, อยํ อนตฺโถ.\csromanspacer{} สมฺมาวาจาปจฺจยา จ อเนเก กุสลา ธมฺมา ภาวนาปาริปูริํ คจฺฉนฺติ, อยํ อตฺโถ.}

\csromanmatikaanchor{mtk.158}
\csromanmatikaanchor{mtk.159}
\csromantocmarknopage{section}{(12) 2. ปจฺโจโรหณิวคฺค}{mtk.158}
\bookmark[dest={mtk.158},level=1]{(12) 2. ปจฺโจโรหณิวคฺค}
\csromantocmarkref{subsection}{1. ปฐม-อธมฺมสุตฺต}{mtk.159}
\bookmark[dest={mtk.159},level=2]{1. ปฐม-อธมฺมสุตฺต}
\prose{มิจฺฉากมฺมนฺโต ภิกฺขเว อธมฺโม, สมฺมากมฺมนฺโต ธมฺโม.\csromanspacer{} เย จ มิจฺฉากมฺมนฺตปจฺจยา อเนเก ปาปกา อกุสลา ธมฺมา สมฺภวนฺติ, อยํ อนตฺโถ.\csromanspacer{} สมฺมากมฺมนฺตปจฺจยา จ อเนเก กุสลา ธมฺมา ภานาปาริปูริํ คจฺฉนฺติ, อยํ อตฺโถ.}

\prose{มิจฺฉา-อาชีโว ภิกฺขเว อธมฺโม, สมฺมา-อาชีโว ธมฺโม.\csromanspacer{} เย จ มิจฺฉา-อาชีวปจฺจยา อเนเก ปาปกา อกุสลา ธมฺมา สมฺภวนฺติ, อยํ อนตฺโถ.\csromanspacer{} สมฺมา-อาชีวปจฺจยา จ อเนเก กุสลา ธมฺมา ภาวนาปาริปูริํ คจฺฉนฺติ, อยํ อตฺโถ.}

\prose{มิจฺฉาวายาโม ภิกฺขเว อธมฺโม, สมฺมาวายาโม ธมฺโม.\csromanspacer{} เย จ มิจฺฉาวายามปจฺจยา อเนเก ปาปกา อกุสลา ธมฺมา สมฺภวนฺติ, อยํ อนตฺโถ.\csromanspacer{} สมฺมาวายามปจฺจยา จ อเนเก กุสลา ธมฺมา ภาวนาปาริปูริํ คจฺฉนฺติ, อยํ อตฺโถ.}

\prose{มิจฺฉาสติ ภิกฺขเว อธมฺโม, สมฺมาสติ ธมฺโม.\csromanspacer{} เย จ มิจฺฉาสติปจฺจยา อเนเก ปาปกา อกุสลา ธมฺมา สมฺภวนฺติ, อยํ อนตฺโถ.\csromanspacer{} สมฺมาสติปจฺจยา จ อเนเก กุสลา ธมฺมา ภาวนาปาริปูริํ คจฺฉนฺติ, อยํ อตฺโถ.}

\prose{มิจฺฉาสมาธิ ภิกฺขเว อธมฺโม, สมฺมาสมาธิ ธมฺโม.\csromanspacer{} เย จ มิจฺฉาสมาธิปจฺจยา อเนเก ปาปกา อกุสลา ธมฺมา สมฺภวนฺติ, อยํ อนตฺโถ.\csromanspacer{} สมฺมาสมาธิปจฺจยา จ อเนเก กุสลา ธมฺมา ภาวนาปาริปูริํ คจฺฉนฺติ, อยํ อตฺโถ.}

\csromanmatikaanchor{mtk.160}
\csromantocmarkref{subsection}{2. ทุติย-อธมฺมสุตฺต}{mtk.160}
\bookmark[dest={mtk.160},level=2]{2. ทุติย-อธมฺมสุตฺต}
\prose{มิจฺฉาญาณํ ภิกฺขเว อธมฺโม, สมฺมาญาณํ ธมฺโม.\csromanspacer{} เย จ มิจฺฉาญาณปจฺจยา อเนเก ปาปกา อกุสลา ธมฺมา สมฺภวนฺติ, อยํ อนตฺโถ.}

\vspace{\tipitakaparskip}
\csromancenter{(12) 2.\csromanspacer{} ปจฺโจโรหณิวคฺค สมฺมาญาณปจฺจยา จ อเนเก กุสลา ธมฺมา ภาวนาปาริปูริํ คจฺฉนฺติ, อยํ อตฺโถ.}
\prose{\csromanfolio{439}มิจฺฉาวิมุตฺติ ภิกฺขเว อธมฺโม, สมฺมาวิมุตฺติ ธมฺโม.\csromanspacer{} เย จ มิจฺฉาวิมุตฺติปจฺจยา อเนเก ปาปกา อกุสลา ธมฺมา สมฺภวนฺติ, อยํ อนตฺโถ.\csromanspacer{} สมฺมาวิมุตฺติปจฺจยา จ อเนเก กุสลา ธมฺมา ภาวนาปาริปูริํ คจฺฉนฺติ อยํ อตฺโถ.}

\prose{“อธมฺโม จ ภิกฺขเว เวทิตพฺโพ ธมฺโม จ, อนตฺโถ จ เวทิตพฺโพ อตฺโถ จ, อธมฺมญฺจ วิทิตฺวา ธมฺมญฺจ, อนตฺถญฺจ วิทิตฺวา อตฺถญฺจ ยถา ธมฺโม ยถา อตฺโถ, ตถา ปฏิปชฺชิตพฺพนฺ”ติ อิติ ยํ ตํ วุตฺตํ, อิทเมตํ ปฏิจฺจ วุตฺตนฺติ.\csromanspacer{}\csromanspacer{}ทุติยํ.}

\csromanheader{2}{3. ตติย-อธมฺมสุตฺต}
\proseitem{115}{อธมฺโม จ ภิกฺขเว เวทิตพฺโพ ธมฺโม จ, อนตฺโถ จ เวทิตพฺโพ อตฺโถ จ, อธมฺมญฺจ วิทิตฺวา ธมฺมญฺจ, อนตฺถญฺจ วิทิตฺวา อตฺถญฺจ ยถา ธมฺโม ยถา อตฺโถ, ตถา ปฏิปชฺชิตพฺพนฺติ อิทมโวจ ภควา อิทํ วตฺวาน สุคโต อุฏฺฐายาสนา วิหารํ ปาวิสิ.}

\prose{อถ โข เตสํ ภิกฺขูนํ อจิรปกฺกนฺตสฺส ภควโต เอตทโหสิ “อิทํ โข โน อาวุโส ภควา สํขิตฺเตน อุทฺเทสํ อุทฺทิสิตฺวา วิตฺถาเรน อตฺถํ อวิภชิตฺวา อุฏฺฐายาสนา วิหารํ ปวิฏฺโฐ, ‘อธมฺโม จ ภิกฺขเว เวทิตพฺโพ ธมฺโม จ, อนตฺโถ จ เวทิตพฺโพ อตฺโถ จ, อธมฺมญฺจ วิทิตฺวา ธมฺมญฺจ, อนตฺถญฺจ วิทิตฺวา อตฺถญฺจ ยถา ธมฺโม ยถา อตฺโถ, ตถา ปฏิปชฺชิตพฺพนฺ’ติ.\csromanspacer{} โก นุ โข อิมสฺส ภควตา สํขิตฺเตน อุทฺเทสสฺส อุทฺทิฏฺฐสฺส วิตฺถาเรน อตฺถํ อวิภตฺตสฺส วิตฺถาเรน อตฺถํ วิภเชยฺยา”ติ.}

\prose{อถ โข เตสํ ภิกฺขูนํ เอตทโหสิ “อยํ โข อายสฺมา อานนฺโท สตฺถุ เจว สํวณฺณิโต สมฺภาวิโต จ วิญฺญูนํ สพฺรหฺมจารีนํ, ปโหติ จายสฺมา อานนฺโท อิมสฺส ภควตา สํขิตฺเตน อุทฺเทสสฺส อุทฺทิฏฺฐสฺส วิตฺถาเรน อตฺถํ อวิภตฺตสฺส วิตฺถาเรน อตฺถํ วิภชิตุํ, ยํนูน มยํ เยนายสฺมา อานนฺโท เตนุปสงฺกเมยฺยาม, อุปสงฺกมิตฺวา อายสฺมนฺตํ อานนฺทํ เอตมตฺถํ ปฏิปุจฺเฉยฺยาม\footnote{ปุจฺเฉยฺยาม (สี, สฺยา, อิ) ม 1. 156 ปิฏฺฐาทีสุ ปสฺสิตพฺพํ.}.\csromanspacer{} ยถา โน อายสฺมา อานนฺโท พฺยากริสฺสติ, ตถา นํ ธาเรสฺสามา”ติ.}

\gambhira{ทสกนิปาตปาฬิ}
\prose{\csromanfolio{440}อถ โข เต ภิกฺขู เยนายสฺมา อานนฺโท เตนุปสงฺกมิํสุ, อุปสงฺกมิตฺวา อายสฺมตา อานนฺเทน สทฺธิํ สมฺโมทิํสุ, สมฺโมทนียํ กถํ สารณียํ วีติสาเรตฺวา เอกมนฺตํ นิสีทิํสุ, เอกมนฺตํ นิสินฺนา โข เต ภิกฺขู อายสฺมนฺตํ อานนฺทํ เอตทโวจุํ\csromandash{}}

\prose{อิทํ โข โน อาวุโส อานนฺท ภควา สํขิตฺเตน อุทฺเทสํ อุทฺทิสิตฺวา วิตฺถาเรน อตฺถํ อวิภชิตฺวา อุฏฺฐายาสนา วิหารํ ปวิฏฺโฐ “อธมฺโม จ ฯเปฯ ตถา ปฏิปชฺชิตพฺพนฺ”ติ.}

\prose{เตสํ โน อาวุโส อมฺหากํ อจิรปกฺกนฺตสฺส ภควโต เอตทโหสิ “อิทํ โข โน อาวุโส ภควตา สํขิตฺเตน อุทฺเทสํ อุทฺทิสิตฺวา วิตฺถาเรน อตฺถํ อวิภชิตฺวา อุฏฺฐายาสนา วิหารํ ปวิฏฺโฐ, ‘อธมฺโม จ ฯเปฯ ตถา ปฏิปชฺชิตพฺพนฺ’ติ.\csromanspacer{} โก นุ โข อิมสฺส ภควตา สํขิตฺเตน อุทฺเทสสฺส อุทฺทิฏฺฐสฺส วิตฺถาเรน อตฺถํ อวิภตฺตสฺส วิตฺถาเรน อตฺถํ วิภเชยฺยา”ติ.}

\prose{เตสํ โน อาวุโส อมฺหากํ เอตทโหสิ “อยํ โข อายสฺมา อานนฺโท สตฺถุ เจว สํวณฺณิโต, สมฺภาวิโต จ วิญฺญูนํ สพฺรหฺมจารีนํ, ปโหติ จายสฺมา อานนฺโท อิมสฺส ภควตา สํขิตฺเตน อุทฺเทสสฺส อุทฺทิฏฺฐสฺส วิตฺถาเรน อตฺถํ อวิภตฺตสฺส วิตฺถาเรน อตฺถํ วิภชิตุํ, ยํนูน มยํ เยนายสฺมา อานนฺโท เตนุปสงฺกเมยฺยาม, อุปสงฺกมิตฺวา อายสฺมนฺตํ อานนฺทํ เอตมตฺถํ ปฏิปุจฺเฉยฺยาม, ยถา โน อายสฺมา อานนฺโท พฺยากริสฺสติ, ตถา นํ ธาเรสฺสามา”ติ.\csromanspacer{} วิภชตุ อายสฺมา อานนฺโทติ.}

\prose{เสยฺยถาปิ อาวุโส ปุริโส สารตฺถิโก สารคเวสี สารปริเยสนํ จรมาโน มหโต รุกฺขสฺส ติฏฺฐโต สารวโต อติกฺกมฺเมว มูลํ อติกฺกมฺม ขนฺธํ สาขาปลาเส สารํ ปริเยสิตพฺพํ มญฺเญยฺย, เอวํสมฺปทมิทํ อายสฺมนฺตานํ สตฺถริ สมฺมุขีภูเต ตํ ภควนฺตํ อติสิตฺวา อเมฺห เอตมตฺถํ ปฏิปุจฺฉิตพฺพํ มญฺญถ.\csromanspacer{} โส หาวุโส ภควา ชานํ ชานาติ, ปสฺสํ ปสฺสติ, จกฺขุภูโต ญาณภูโต ธมฺมภูโต พฺรหฺมภูโต วตฺตา ปวตฺตา อตฺถสฺส นินฺเนตา อมตสฺส ทาตา ธมฺมสฺสามี ตถาคโต, โส เจว ปเนตสฺส กาโล อโหสิ, ยํ ตุเมฺห}

\vspace{\tipitakaparskip}
\csromancenter{(12) 2.\csromanspacer{} ปจฺโจโรหณิวคฺค ภควนฺตํเยว อุปสงฺกมิตฺวา เอตมตฺถํ ปฏิปุจฺเฉยฺยาถ.\csromanspacer{} ยถา โว ภควา พฺยากเรยฺย, ตถา นํ ธาเรยฺยาถาติ.}
\prose{\csromanfolio{441}อทฺธาวุโส อานนฺท ภควา ชานํ ชานาติ, ปสฺสํ ปสฺสติ, จกฺขุภูโต ญาณภูโต ธมฺมภูโต พฺรหฺมภูโต วตฺตา ปวตฺตา อตฺถสฺส นินฺเนตา อมตสฺส ทาตา ธมฺมสฺสามี ตถาคโต, โส เจว ปเนตสฺส กาโล อโหสิ, ยํ มยํ ภควนฺตํเยว อุปสงฺกมิตฺวา เอตมตฺถํ ปฏิปุจฺเฉยฺยาม.\csromanspacer{} ยถา โน ภควา พฺยากเรยฺย, ตถา นํ ธาเรยฺยาม.\csromanspacer{} อปิ จายสฺมา อานนฺโท สตฺถุ เจว สํวณฺณิโต, สมฺภาวิโต จ วิญฺญูนํ สพฺรหฺมจารีนํ, ปโหติ จายสฺมา อานนฺโท อิมสฺส ภควตา สํขิตฺเตน อุทฺเทสสฺส อุทฺทิฏฺฐสฺส วิตฺถาเรน อตฺถํ อวิภตฺตสฺส วิตฺถาเรน อตฺถํ วิภชิตุํ, วิภชตายสฺมา อานนฺโท อครุํ กตฺวาติ.}

\csromanmatikaanchor{mtk.161}
\csromantocmarkref{subsection}{3. ตติย-อธมฺมสุตฺต}{mtk.161}
\bookmark[dest={mtk.161},level=2]{3. ตติย-อธมฺมสุตฺต}
\prose{เตนหาวุโส สุณาถ สาธุกํ มนสิ กโรถ, ภาสิสฺสามีติ. “เอวมาวุโส”ติ โข เต ภิกฺขู อายสฺมโต อานนฺทสฺส ปจฺจสฺโสสุํ.\csromanspacer{} อถายสฺมา อานนฺโท เอตทโวจ\csromandash{}}

\prose{ยํ โข โน อาวุโส ภควา สํขิตฺเตน อุทฺเทสํ อุทฺทิสิตฺวา วิตฺถาเรน อตฺถํ อวิภชิตฺวา อุฏฺฐายาสนา วิหารํ ปวิฏฺโฐ, “อธมฺโม จ ภิกฺขเว เวทิตพฺโพ ธมฺโม จ, อนตฺโถ จ เวทิตพฺโพ อตฺโถ จ, อธมฺมญฺจ วิทิตฺวา ธมฺมญฺจ, อนตฺถญฺจ วิทิตฺวา อตฺถญฺจ ยถา ธมฺโม ยถา อตฺโถ, ตถา ปฏิปชฺชิตพฺพนฺ”ติ.}

\prose{กตโม จาวุโส อธมฺโม, กตโม จ ธมฺโม, กตโม จ อนตฺโถ, กตโม จ อตฺโถ.}

\prose{มิจฺฉาทิฏฺฐิ อาวุโส อธมฺโม, สมฺมาทิฏฺฐิ ธมฺโม.\csromanspacer{} เย จ มิจฺฉาทิฏฺฐิปจฺจยา อเนเก ปาปกา อกุสลา ธมฺมา สมฺภวนฺติ, อยํ อนตฺโถ.\csromanspacer{} สมฺมาทิฏฺฐิปจฺจยา จ อเนเก กุสลา ธมฺมา ภาวนาปาริปูริํ คจฺฉนฺติ, อยํ อตฺโถ.}

\prose{มิจฺฉาสงฺกปฺโป อาวุโส อธมฺโม, สมฺมาสงฺกปฺโป ธมฺโม.\csromanspacer{} มิจฺฉาวาจา อาวุโส อธมฺโม, สมฺมาวาจา ธมฺโม.\csromanspacer{} มิจฺฉากมฺมนฺโต อาวุโส อธมฺโม, สมฺมากมฺมนฺโต ธมฺโม.\csromanspacer{} มิจฺฉา-อาชีโว อาวุโส อธมฺโม,}

\vspace{\tipitakaparskip}
\csromancenter{ทสกนิปาตปาฬิ สมฺมา อาชีโว ธมฺโม.\csromanspacer{} มิจฺฉาวายาโม อาวุโส อธมฺโม, สมฺมาวายาโม ธมฺโม.\csromanspacer{} มิจฺฉาสติ อาวุโส อธมฺโม, สมฺมาสติ ธมฺโม.\csromanspacer{} มิจฺฉาสมาธิ อาวุโส อธมฺโม, สมฺมาสมาธิ ธมฺโม.\csromanspacer{} มิจฺฉาญาณํ อาวุโส อธมฺโม, สมฺมาญาณํ ธมฺโม.}
\prose{\csromanfolio{442}มิจฺฉาวิมุตฺติ อาวุโส อธมฺโม, สมฺมาวิมุตฺติ ธมฺโม.\csromanspacer{} เย จ มิจฺฉาวิมุตฺติปจฺจยา อเนเก ปาปกา อกุสลา ธมฺมา สมฺภวนฺติ, อยํ อนตฺโถ, สมฺมาวิมุตฺติปจฺจยา จ อเนเก กุสลา ธมฺมา ภาวนาปาริปูริํ คจฺฉนฺติ, อยํ อตฺโถ.}

\prose{อยํ โข โน อาวุโส ภควา สํขิตฺเตน อุทฺเทสํ อุทฺทิสิตฺวา วิตฺถาเรน อตฺถํ อวิภชิตฺวา อุฏฺฐายาสนา วิหารํ ปวิฏฺโฐ, “อธมฺโม จ ภิกฺขเว เวทิตพฺโพ ธมฺโม จ ฯเปฯ ตถา ปฏิปชฺชิตพฺพนฺ”ติ.\csromanspacer{} อิมสฺส โข อหํ อาวุโส ภควตา สํขิตฺเตน อุทฺเทสสฺส อุทฺทิฏฺฐสฺส วิตฺถาเรน อตฺถํ อวิภตฺตสฺส เอวํ วิตฺถาเรน อตฺถํ อาชานามิ, อากงฺขมานา จ ปน ตุเมฺห อาวุโส ภควนฺตํเยว อุปสงฺกมิตฺวา เอตมตฺถํ ปฏิปุจฺเฉยฺยาถ.\csromanspacer{} ยถา โว ภควา พฺยากโรติ\footnote{พฺยากเรยฺย (สฺยา)}, ตถา นํ ธาเรยฺยาถาติ.}

\prose{“เอวมาวุโส”ติ โข เต ภิกฺขู อายสฺมโต อานนฺทสฺส ภาสิตํ อภินนฺทิตฺวา อนุโมทิตฺวา อุฏฺฐายาสนา เยน ภควา เตนุปสงฺกมิํสุ, อุปสงฺกมิตฺวา ภควนฺตํ อภิวาเทตฺวา เอกมนฺตํ นิสีทิํสุ, เอกมนฺตํ นิสินฺนา โข เต ภิกฺขู ภควนฺตํ เอตทโวจุํ\csromandash{}}

\prose{ยํ โข โน ภควา สํขิตฺเตน อุทฺเทสํ อุทฺทิสิตฺวา วิตฺถาเรน อตฺถํ อวิภชิตฺวา อุฏฺฐายาสนา วิหารํ ปวิฏฺโฐ, “อธมฺโม จ ภิกฺขเว เวทิตพฺโพ ฯเปฯ ตถา ปฏิปชฺชิตพฺพนฺ”ติ.}

\prose{เตสํ โน ภนฺเต อมฺหากํ อจิรปกฺกนฺตสฺส ภควโต เอตทโหสิ “อิทํ โข โน อาวุโส ภควา สํขิตฺเตน อุทฺเทสํ อุทฺทิสิตฺวา วิตฺตาเรน อตฺถํ อวิภชิตฺวา อุฏฺฐายาสนา วิหารํ ปวิฏฺโฐ, ‘อธมฺโม จ ภิกฺขเว เวทิตพฺโพ ฯเปฯ ตถา ปฏิปชฺชิตพฺพนฺ’ติ.\csromanspacer{} โก นุ โข อิมสฺส ภควตา สํขิตฺเตน อุทฺเทสสฺส อุทฺทิฏฺฐสฺส วิตฺถาเรน อตฺถํ อวิภตฺตสฺส วิตฺถาเรน อตฺถํ วิภเชยฺยา”ติ.}

\csromanheader{1}{(12) 2. ปจฺโจโรหณิวคฺค}
\prose{\csromanfolio{443}เตสํ โน ภนฺเต อมฺหากํ เอตทโหสิ “อยํ โข อายสฺมา อานนฺโท สตฺถุ เจว สํวณฺณิโต, สมฺภาวิโต จ วิญฺญูนํ สพฺรหฺมจารีนํ, ปโหติ จายสฺมา อานนฺโท อิมสฺส ภควตา สํขิตฺเตน อุทฺเทสสฺส อุทฺทิฏฺฐสฺส วิตฺถาเรน อตฺถํ อวิภตฺตสฺส วิตฺถาเรน อตฺถํ วิภชิตุํ.\csromanspacer{} ยํนูน มยํ เยนายสฺมา อานนฺโท เตนุปสงฺกเมยฺยาม, อุปสงฺกมิตฺวา อายสฺมนฺตํ อานนฺทํ เอตมตฺถํ ปฏิปุจฺเฉยฺยาม.\csromanspacer{} ยถา โน อายสฺมา อานนฺโท พฺยากริสฺสติ, ตถา นํ ธาเรสฺสามา”ติ.}

\prose{อถ โข มยํ ภนฺเต เยนายสฺมา อานนฺโท เตนุปสงฺกมิมฺหา, อุปสงฺกมิตฺวา อายสฺมนฺตํ อานนฺทํ เอตมตฺถํ อปุจฺฉิมฺหา.\csromanspacer{} เตสํ โน ภนฺเต อายสฺมตา อานนฺเทน อิเมหิ อากาเรหิ อิเมหิ ปเทหิ อิเมหิ พฺยญฺชเนหิ อตฺโถ สุวิภตฺโตติ\footnote{วิภตฺโตติ (?) เอวเมว หิ อญฺเญสุ อีทิสสุตฺเตสุ ทิสฺสติ.}.}

\prose{สาธุ สาธุ ภิกฺขเว, ปณฺฑิโต ภิกฺขเว อานนฺโท, มหาปญฺโญ ภิกฺขเว อานนฺโท.\csromanspacer{} มํ เจปิ ตุเมฺห ภิกฺขเว อุปสงฺกมิตฺวา เอตมตฺถํ ปฏิปุจฺเฉยฺยาถ, อหมฺปิ เจตํ เอวเมวํ\footnote{อหมฺปิ ตํ เอวเมวํ (ม 1. 161 ปิฏฺฐาทีสุ)} พฺยากเรยฺยํ, ยถา ตํ อานนฺเทน พฺยากตํ, เอโส เจว ตสฺส\footnote{เอโส เจเวตสฺส (ม 1. 161 ปิฏฺฐาทีสุ)} อตฺโถ, เอวญฺจ นํ ธาเรยฺยาถาติ.\csromanspacer{}\csromanspacer{}ตติยํ.}

\csromanheader{2}{4. อชิตสุตฺต}
\proseitem{116}{อถ โข อชิโต ปริพฺพาชโก เยน ภควา เตนุปสงฺกมิ, อุปสงฺกมิตฺวา ภควตา สทฺธิํ สมฺโมทิ, สมฺโมทนียํ กถํ สารณียํ วีติสาเรตฺวา เอกมนฺตํ นิสีทิ, เอกมนฺตํ นิสินฺโน โข อชิโต ปริพฺพาชโก ภควนฺตํ เอตทโวจ\csromandash{}}

\prose{อมฺหากํ โภ โคตม ปณฺฑิโต นาม สพฺรหฺมจารี, เตน ปญฺจมตฺตานิ จิตฺตฏฺฐานสตานิ จินฺติตานิ, เยหิ อญฺญติตฺถิยา อุปารทฺธาว ชานนฺติ\footnote{อุปารทฺธา ปชานนฺติ (สี)} อุปารทฺธสฺมาติ\footnote{อุปารทฺธมฺหาติ (สี, อิ)}.}

\gambhira{ทสกนิปาตปาฬิ}
\prose{\csromanfolio{444}อถ โข ภควา ภิกฺขู อามนฺเตสิ “ธาเรถ โน ตุเมฺห ภิกฺขเว ปณฺฑิตวตฺถูนี”ติ.\csromanspacer{} เอตสฺส ภควา กาโล เอตสฺส สุคต กาโล.\csromanspacer{} ยํ ภควา ภาเสยฺย, ภควโต สุตฺวา ภิกฺขู ธาเรสฺสนฺตีติ.}

\prose{เตน หิ ภิกฺขเว สุณาถ สาธุกํ มนสิ กโรถ, ภาสิสฺสามีติ. “เอวํ ภนฺเต”ติ โข เต ภิกฺขู ภควโต ปจฺจสฺโสสุํ.\csromanspacer{} ภควา เอตทโวจ\csromandash{}}

\prose{อิธ ภิกฺขเว เอกจฺโจ อธมฺมิเกน วาเทน อธมฺมิกํ วาทํ อภินิคฺคณฺหาติ อภินิปฺปีเฬติ, เตน จ อธมฺมิกํ ปริสํ รญฺเชติ, เตน สา อธมฺมิกา ปริสา อุจฺจาสทฺทมหาสทฺทา โหติ “ปณฺฑิโต วต โภ ปณฺฑิโต วต โภ”ติ.}

\prose{อิธ ปน ภิกฺขเว เอกจฺโจ อธมฺมิเกน วาเทน ธมฺมิกํ วาทํ อภินิคฺคณฺหาติ อภินิปฺปีเฬติ, เตน จ อธมฺมิกํ ปริสํ รญฺเชติ, เตน สา อธมฺมิกา ปริสา อุจฺจาสทฺทมหาสทฺทา โหติ “ปณฺฑิโต วต โภ ปณฺฑิโต วต โภ”ติ.}

\prose{อิธ ปน ภิกฺขเว เอกจฺโจ อธมฺมิเกน วาเทน ธมฺมิกญฺจ วาทํ อธมฺมิกญฺจ วาทํ อภินิคฺคณฺหาติ อภินิปฺปีเฬติ, เตน จ อธมฺมิกํ ปริสํ รญฺเชติ, เตน สา อธมฺมิกา ปริสา อุจฺจาสทฺทมหาสทฺทา โหติ “ปณฺฑิโต วต โภ ปณฺฑิโต วต โภ”ติ.}

\prose{อธมฺโม จ ภิกฺขเว เวทิตพฺโพ ธมฺโม จ, อนตฺโถ จ เวทิตพฺโพ อตฺโถ จ, อธมฺมญฺจ วิทิตฺวา ธมฺมญฺจ อนตฺถญฺจ วิทิตฺวา อตฺถญฺจ ยถา ธมฺโม ยถา อตฺโถ, ตถา ปฏิปชฺชิตพฺพํ.}

\prose{กตโม จ ภิกฺขเว อธมฺโม, กตโม จ ธมฺโม, กตโม จ อนตฺโถ, กตโม จ อตฺโถ.\csromanspacer{} มิจฺฉาทิฏฺฐิ ภิกฺขเว อธมฺโม, สมฺมาทิฏฺฐิ ธมฺโม.\csromanspacer{} เย จ มิจฺฉาทิฏฺฐิปจฺจยา อเนเก ปาปกา อกุสลา ธมฺมา สมฺภวนฺติ, อยํ อนตฺโถ.\csromanspacer{} สมฺมาทิฏฺฐิปจฺจยา จ อเนเก กุสลา ธมฺมา ภาวนาปาริปูริํ คจฺฉนฺติ, อยํ อตฺโถ.}

\prose{มิจฺฉาสงฺกปฺโป ภิกฺขเว อธมฺโม, สมฺมาสงฺกปฺโป ธมฺโม.\csromanspacer{} มิจฺฉาวาจา ภิกฺขเว อธมฺโม, สมฺมาวาจา ธมฺโม.\csromanspacer{} มิจฺฉากมฺมนฺโต ภิกฺขเว อธมฺโม, สมฺมากมฺมนฺโต ธมฺโม.\csromanspacer{} มิจฺฉา-อาชีโว ภิกฺขเว อธมฺโม, สมฺมา-อาชีโว}

\csromanmatikaanchor{mtk.162}
\csromantocmarkref{subsection}{4. อชิตสุตฺต}{mtk.162}
\bookmark[dest={mtk.162},level=2]{4. อชิตสุตฺต}
\vspace{\tipitakaparskip}
\csromancenter{(12) 2.\csromanspacer{} ปจฺโจโรหณิวคฺค ธมฺโม.\csromanspacer{} มิจฺฉาวายาโม ภิกฺขเว อธมฺโม, สมฺมาวายาโม ธมฺโม.\csromanspacer{} มิจฺฉาสติ ภิกฺขเว อธมฺโม, สมฺมาสติ ธมฺโม.\csromanspacer{} มิจฺฉาสมาธิ ภิกฺขเว อธมฺโม, สมฺมาสมาธิ ธมฺโม.\csromanspacer{} มิจฺฉาญาณํ ภิกฺขเว อธมฺโม, สมฺมาญาณํ ธมฺโม.}
\prose{\csromanfolio{445}มิจฺฉาวิมุตฺติ ภิกฺขเว อธมฺโม, สมฺมาวิมุตฺติ ธมฺโม.\csromanspacer{} เย จ มิจฺฉาวิมุตฺติปจฺจยา อเนเก ปาปกา อกุสลา ธมฺมา สมฺภวนฺติ, อยํ อนตฺโถ.\csromanspacer{} สมฺมาวิมุตฺติปจฺจยา จ อเนเก กุสลา ธมฺมา ภาวนาปาริปูริํ คจฺฉนฺติ, อยํ อตฺโถ.}

\prose{“อธมฺโม จ ภิกฺขเว เวทิตพฺโพ ธมฺโม จ, อนตฺโถ จ เวทิตพฺโพ อตฺโถ จ, อธมฺมญฺจ วิทิตฺวา ธมฺมญฺจ, อนตฺถญฺจ วิทิตฺวา อตฺถญฺจ ยถา ธมฺโม ยถา อตฺโถ, ตถา ปฏิปชฺชิตพฺพนฺ”ติ อิติ ยํ ตํ วุตฺตํ, อิทเมตํ ปฏิจฺจ วุตฺตนฺติ.\csromanspacer{}\csromanspacer{}จตุตฺถํ.}

\csromanheader{2}{5. สงฺคารวสุตฺต}
\proseitem{117}{\csromansymbolfootnote{*}{อุปริ 462 ปิฏฺเฐปิ.}อถ โข สงฺคารโว พฺราหฺมโณ เยน ภควา เตนุปสงฺกมิ, อุปสงฺกมิตฺวา ภควตา สทฺธิํ สมฺโมทิ, สมฺโมทนียํ กถํ สารณียํ วีติสาเรตฺวา เอกมนฺตํ นิสีทิ, เอกมนฺตํ นิสินฺโน โข สงฺคารโว พฺราหฺมโณ ภควนฺตํ เอตทโวจ “กิํ นุ โข โภ โคตม โอริมํ ตีรํ, กิํ ปาริมํ ตีรนฺ”ติ.\csromanspacer{} มิจฺฉาทิฏฺฐิ โข พฺราหฺมณ โอริมํ ตีรํ, สมฺมาทิฏฺฐิ ปาริมํ ตีรํ.\csromanspacer{} มิจฺฉาสงฺกปฺโป โอริมํ ตีรํ, สมฺมาสงฺกปฺโป ปาริมํ ตีรํ.\csromanspacer{} มิจฺฉาวาจา โอริมํ ตีรํ, สมฺมาวาจา ปาริมํ ตีรํ.\csromanspacer{} มิจฺฉากมฺมนฺโต โอริมํ ตีรํ, สมฺมากมฺมนฺโต ปาริมํ ตีรํ.\csromanspacer{} มิจฺฉา-อาชีโว โอริมํ ตีรํ, สมฺมา-อาชีโว ปาริมํ ตีรํ.\csromanspacer{} มิจฺฉาวายาโม โอริมํ ตีรํ, สมฺมาวายาโม ปาริมํ ตีรํ.\csromanspacer{} มิจฺฉาสติ โอริมํ ตีรํ, สมฺมาสติ ปาริมํ ตีรํ.\csromanspacer{} มิจฺฉาสมาธิ โอริมํ ตีรํ, สมฺมาสมาธิ ปาริมํ ตีรํ.\csromanspacer{} มิจฺฉาญาณํ โอริมํ ตีรํ, สมฺมาญาณํ ปาริมํ ตีรํ.\csromanspacer{} มิจฺฉาวิมุตฺติ โอริมํ ตีรํ, สมฺมาวิมุตฺติ ปาริมํ ตีรนฺติ.\csromanspacer{} อิทํ โข พฺราหฺมณ โอริมํ ตีรํ, อิทํ ปาริมํ ตีรนฺติ.}

\csromangathameasure{อปฺปกา เต มนุสฺเสสุ, เย ชนา ปารคามิโน. \\ อถายํ อิตรา ปชา, ตีรเมวานุธาวติ. \\ เย จ โข สมฺมทกฺขาเต, ธมฺเม ธมฺมานุวตฺติโน. \\ เต ชนา ปารเมสฺสนฺติ, มจฺจุเธยฺยํ สุทุตฺตรํ.}
\csromangathagroup{\csromangathastanza{อปฺปกา เต มนุสฺเสสุ, เย ชนา ปารคามิโน. \\ อถายํ อิตรา ปชา, ตีรเมวานุธาวติ.}\csromangathastanza{เย จ โข สมฺมทกฺขาเต, ธมฺเม ธมฺมานุวตฺติโน. \\ เต ชนา ปารเมสฺสนฺติ, มจฺจุเธยฺยํ สุทุตฺตรํ.}}

\gambhira{ทสกนิปาตปาฬิ}
\csromanmatikaanchor{mtk.163}
\csromanmatikaanchor{mtk.164}
\csromantocmarkref{subsection}{5. สงฺคารวสุตฺต}{mtk.163}
\bookmark[dest={mtk.163},level=2]{5. สงฺคารวสุตฺต}
\csromantocmarkref{subsection}{6. โอรีมตีรสุตฺต}{mtk.164}
\bookmark[dest={mtk.164},level=2]{6. โอรีมตีรสุตฺต}
\csromangathameasure{กณฺหํ ธมฺมํ วิปฺปหาย, สุกฺกํ ภาเวถ ปณฺฑิโต. \\ โอกา อโนก มาคมฺม, วิเวเก ยตฺถ ทูรมํ. \\ ตตฺราภิรติมิจฺเฉยฺย, หิตฺวา กาเม อกิญฺจโน. \\ ปริโยทเปยฺย อตฺตานํ, จิตฺตเกฺลเสหิ ปณฺฑิโต. \\ เยสํ สมฺโพธิยงฺเคสุ, สมฺมา จิตฺตํ สุภาวิตํ. \\ อาทานปฏินิสฺสคฺเค, อนุปาทาย เย รตา. \\ ขีณาสวา ชุติมนฺโต, เต โลเก ปรินิพฺพุตาติ.ปญฺจมํ.}
\csromangathagroup{\csromangathastanza{\csromanfolio{446}กณฺหํ ธมฺมํ วิปฺปหาย, สุกฺกํ ภาเวถ ปณฺฑิโต. \\ โอกา อโนก มาคมฺม, วิเวเก ยตฺถ ทูรมํ.}\csromangathastanza{ตตฺราภิรติมิจฺเฉยฺย, หิตฺวา กาเม อกิญฺจโน. \\ ปริโยทเปยฺย อตฺตานํ, จิตฺตเกฺลเสหิ ปณฺฑิโต.}\csromangathastanza{เยสํ สมฺโพธิยงฺเคสุ, สมฺมา จิตฺตํ สุภาวิตํ. \\ อาทานปฏินิสฺสคฺเค, อนุปาทาย เย รตา. \\ ขีณาสวา ชุติมนฺโต\footnote{ชุตีมนฺโต (สี)}, เต โลเก ปรินิพฺพุตาติ.\csromanspacer{}\csromanspacer{}ปญฺจมํ.}}

\csromanheader{2}{6. โอริมตีรสุตฺต}
\proseitem{118}{โอริมญฺจ ภิกฺขเว ตีรํ เทเสสฺสามิ ปาริมญฺจ ตีรํ, ตํ สุณาถ สาธุกํ มนสิ กโรถ, ภาสิสฺสามีติ. “เอวํ ภนฺเต”ติ โข เต ภิกฺขู ภควโต ปจฺจสฺโสสุํ.\csromanspacer{} ภควา เอตทโวจ\csromandash{}}

\prose{กตมญฺจ ภิกฺขเว โอริมํ ตีรํ, กตมญฺจ ปาริมํ ตีรํ? มิจฺฉาทิฏฺฐิ โอริมํ ตีรํ, สมฺมาทิฏฺฐิ ปาริมํ ตีรํ ฯเปฯ มิจฺฉาวิมุตฺติ โอริมํ ตีรํ, สมฺมาวิมุตฺติ ปาริมํ ตีรํ.\csromanspacer{} อิทํ โข ภิกฺขเว โอริมํ ตีรํ, อิทํ ปาริมํ ตีรนฺติ.}

\csromangathameasure{อปฺปกา เต มนุสฺเสสุ, เย ชนา ปารคามิโน. \\ อถายํ อิตรา ปชา, ตีรเมวานุธาวติ. \\ เย จ โข สมฺมทกฺขาเต, ธมฺเม ธมฺมานุวตฺติโน. \\ เต ชนา ปารเมสฺสนฺติ, มจฺจุเธยฺยํ สุทุตฺตรํ. \\ กณฺหํ ธมฺมํ วิปฺปหาย, สุกฺกํ ภาเวถ ปณฺฑิโต. \\ โอกา อโนก มาคมฺม, วิเวเก ยตฺถ ทูรมํ. \\ ตตฺราภิรติมิจฺเฉยฺย, หิตฺวา กาเม อกิญฺจโน. \\ ปริโยทเปยฺย อตฺตานํ, จิตฺตเกฺลเสหิ ปณฺฑิโต. \\ เยสํ สมฺโพธิยงฺเคสุ, สมฺมา จิตฺตํ สุภาวิตํ. \\ อาทานปฏินิสฺสคฺเค, อนุปาทาย เย รตา. \\ ขีณาสวา ชุติมนฺโต, เต โลเก ปรินิพฺพุตาติ.ฉฏฺฐํ.}
\csromangathagroup{\csromangathastanza{อปฺปกา เต มนุสฺเสสุ, เย ชนา ปารคามิโน. \\ อถายํ อิตรา ปชา, ตีรเมวานุธาวติ.}\csromangathastanza{เย จ โข สมฺมทกฺขาเต, ธมฺเม ธมฺมานุวตฺติโน. \\ เต ชนา ปารเมสฺสนฺติ, มจฺจุเธยฺยํ สุทุตฺตรํ.}\csromangathastanza{กณฺหํ ธมฺมํ วิปฺปหาย, สุกฺกํ ภาเวถ ปณฺฑิโต. \\ โอกา อโนก มาคมฺม, วิเวเก ยตฺถ ทูรมํ.}\csromangathastanza{ตตฺราภิรติมิจฺเฉยฺย, หิตฺวา กาเม อกิญฺจโน. \\ ปริโยทเปยฺย อตฺตานํ, จิตฺตเกฺลเสหิ ปณฺฑิโต.}\csromangathastanza{เยสํ สมฺโพธิยงฺเคสุ, สมฺมา จิตฺตํ สุภาวิตํ. \\ อาทานปฏินิสฺสคฺเค, อนุปาทาย เย รตา. \\ ขีณาสวา ชุติมนฺโต, เต โลเก ปรินิพฺพุตาติ.\csromanspacer{}\csromanspacer{}ฉฏฺฐํ.}}

\csromanheader{2}{(12) 2. ปจฺโจโรหณิวคฺค \textbf{7. ปฐมปจฺโจโรหณีสุตฺต}}
\proseitem{119}{\csromanfolio{447}เตน โข ปน สมเยน ชาณุสฺโสณิ\footnote{ชานุสฺโสนิ (ก-สี), ชานุสฺโสณิ (ก-สี), ชาณุโสณิ (ก)} พฺราหฺมโณ ตทหุโปสเถ สีสํนฺหาโต\footnote{สีสํนหาโต (สี, อิ), สีสนฺหาโต (สฺยา)} นวํ โขมยุคํ นิวตฺโถ อลฺลกุสมุฏฺฐิํ อาทาย ภควโต อวิทูเร เอกมนฺตํ ฐิโต โหติ.}

\prose{อทฺทสา โข ภควา ชาณุสฺโสณิํ พฺราหฺมณํ ตทหุโปสเถ สีสํนฺหาตํ นวํ โขมยุคํ นิวตฺถํ อลฺลกุสมุฏฺฐิํ อาทาย เอกมนฺตํ ฐิตํ, ทิสฺวาน ชาณุสฺโสณิํ พฺราหฺมณํ เอตทโวจ “กิํ นุ ตฺวํ พฺราหฺมณ ตทหุโปสเถ สีสํนฺหาโต นวํ โขมยุคํ นิวตฺโถ อลฺลกุสมุฏฺฐิํ อาทาย เอกมนฺตํ ฐิโต, กิํ นฺวชฺช\footnote{กิํ นุ อชฺช (สฺยา), กิํ นุ โข อชฺช (อิ), กิํ นุ ขฺวชฺช (ก)} พฺราหฺมณกุลสฺสาติ\footnote{พฺราหฺมณ พฺรหฺมกุสลสฺสาติ (ก)}.\csromanspacer{} ปจฺโจโรหณี โภ โคตม อชฺช พฺราหฺมณกุลสฺสา”ติ\footnote{พฺรหฺมกุสลสฺสาติ (ก)}.}

\prose{ยถา กถํ ปน พฺราหฺมณ พฺราหฺมณานํ ปจฺโจโรหณี โหตีติ, อิธ โภ โคตม พฺราหฺมณา ตทหุโปสเถ สีสํนฺหาตา นวํ โขมยุคํ นิวตฺถา อลฺเลน โคมเยน ปถวิํ โอปุญฺชิตฺวา หริเตหิ กุเสหิ ปตฺถริตฺวา\footnote{ปวิตฺถาเรตฺวา (ก)} อนฺตรา จ เวลํ อนฺตรา จ อคฺยาคารํ เสยฺยํ กปฺเปนฺติ.\csromanspacer{} เต ตํ รตฺติํ ติกฺขตฺตุํ ปจฺจุฏฺฐาย ปญฺชลิกา อคฺคิํ นมสฺสนฺติ “ปจฺโจโรหาม ภวนฺตํ ปจฺโจโรหาม ภวนฺตนฺ”ติ.\csromanspacer{} พหุเกน จ สปฺปิเตลนวนีเตน อคฺคิํ สนฺตปฺเปนฺติ, ตสฺสา จ รตฺติยา อจฺจเยน ปณีเตน ขาทนีเยน โภชนีเยน พฺราหฺมเณ สนฺตปฺเปนฺติ.\csromanspacer{} เอวํ โภ โคตม พฺราหฺมณานํ ปจฺโจโรหณี โหตีติ.}

\prose{อญฺญถา โข พฺราหฺมณ พฺราหฺมณานํ ปจฺโจโรหณี โหติ, อญฺญถา จ ปน อริยสฺส วินเย ปจฺโจโรหณี โหตีติ.\csromanspacer{} ยถา กถํ ปน โภ โคตม อริยสฺส วินเย ปจฺโจโรหณี โหติ, สาธุ เม ภวํ โคตโม ตถา ธมฺมํ เทเสตุ, ยถา อริยสฺส วินเย ปจฺโจโรหณี โหตีติ.}

\gambhira{ทสกนิปาตปาฬิ}
\prose{\csromanfolio{448}เตน หิ พฺราหฺมณ สุณาหิ สาธุกํ มนสิ กโรหิ, ภาสิสฺสามีติ. “เอวํ โภ”ติ โข ชาณุสฺโสณิ พฺราหฺมโณ ภควโต ปจฺจสฺโสสิ, ภควา เอตทโวจ\csromandash{}}

\prose{อิธ พฺราหฺมณ อริยสาวโก อิติ ปฏิสญฺจิกฺขติ “มิจฺฉาทิฏฺฐิยา โข ปาปโก วิปาโก ทิฏฺเฐ เจว ธมฺเม อภิสมฺปรายญฺจา”ติ.\csromanspacer{} โส อิติ ปฏิสงฺขาย มิจฺฉาทิฏฺฐิํ ปชหติ, มิจฺฉาทิฏฺฐิยา ปจฺโจโรหติ.}

\prose{“มิจฺฉาสงฺกปฺปสฺส โข ปาปโก วิปาโก ทิฏฺเฐ เจว ธมฺเม อภิสมฺปรายญฺจา”ติ.\csromanspacer{} โส อิติ ปฏิสงฺขาย มิจฺฉาสงฺกปฺปํ ปชหติ, มิจฺฉาสงฺกปฺปา ปจฺโจโรหติ.}

\prose{“มิจฺฉาวาจาย โข ปาปโก วิปาโก ทิฏฺเฐ เจว ธมฺเม อภิสมฺปรายญฺจา”ติ.\csromanspacer{} โส อิติ ปฏิสงฺขาย มิจฺฉาวาจํ ปชหติ, มิจฺฉาวาจาย ปจฺโจโรหติ.}

\prose{“มิจฺฉากมฺมนฺตสฺส โข ปาปโก วิปาโก ทิฏฺเฐ เจว ธมฺเม อภิสมฺปรายญฺจา”ติ.\csromanspacer{} โส อิติ ปฏิสงฺขาย มิจฺฉากมฺมนฺตํ ปชหติ, มิจฺฉากมฺมนฺตา ปจฺโจโรหติ.}

\prose{“มิจฺฉา-อาชีวสฺส โข ปาปโก วิปาโก ทิฏฺเฐ เจว ธมฺเม อภิสมฺปรายญฺจา”ติ.\csromanspacer{} โส อิติ ปฏิสงฺขาย มิจฺฉา-อาชีวํ ปชหติ, มิจฺฉา-อาชีวา ปจฺโจโรหติ.}

\csromanmatikaanchor{mtk.165}
\csromantocmarkref{subsection}{7. ปฐมปจฺโจโรหณีสุตฺต}{mtk.165}
\bookmark[dest={mtk.165},level=2]{7. ปฐมปจฺโจโรหณีสุตฺต}
\prose{“มิจฺฉาวายามสฺส โข ปาปโก วิปาโก ทิฏฺเฐ เจว ธมฺเม อภิสมฺปรายญฺจา”ติ.\csromanspacer{} โส อิติ ปฏิสงฺขาย มิจฺฉาวายามํ ปชหติ, มิจฺฉาวายามา ปจฺโจโรหติ.}

\prose{“มิจฺฉาสติยา โข ปาปโก วิปาโก ทิฏฺเฐ เจว ธมฺเม อภิสมฺปรายญฺจา”ติ.\csromanspacer{} โส อิติ ปฏิสงฺขาย มิจฺฉาสติํ ปชหติ, มิจฺฉาสติยา ปจฺโจโรหติ.}

\prose{“มิจฺฉาสมาธิสฺส โข ปาปโก วิปาโก ทิฏฺเฐ เจว ธมฺเม อภิสมฺปรายญฺจา”ติ.\csromanspacer{} โส อิติ ปฏิสงฺขาย มิจฺฉาสมาธิํ ปชหติ, มิจฺฉาสมาธิมฺหา ปจฺโจโรหติ.}

\chapterheadpageread{(12) 2. ปจฺโจโรหณิวคฺค}
\prose{\csromanfolio{449}“มิจฺฉาญาณสฺส โข ปาปโก วิปาโก ทิฏฺเฐ เจว ธมฺเม อภิสมฺปรายญฺจา”ติ.\csromanspacer{} โส อิติ ปฏิสงฺขาย มิจฺฉาญาณํ ปชหติ, มิจฺฉาญาณมฺหา ปจฺโจโรหติ.}

\prose{“มิจฺฉาวิมุตฺติยา โข ปาปโก วิปาโก ทิฏฺเฐ เจว ธมฺเม อภิสมฺปรายญฺจา”ติ.\csromanspacer{} โส อิติ ปฏิสงฺขาย มิจฺฉาวิมุตฺติํ ปชหติ, มิจฺฉาวิมุตฺติยา ปจฺโจโรหติ.\csromanspacer{} เอวํ โข พฺราหฺมณ อริยสฺส วินเย ปจฺโจโรหณี โหตีติ.}

\prose{อญฺญถา โภ โคตม พฺราหฺมณานํ ปจฺโจโรหณี, อญฺญถา จ ปน อริยสฺส วินเย ปจฺโจโรหณี โหติ.\csromanspacer{} อิมิสฺสา จ โภ โคตม อริยสฺส วินเย ปจฺโจโรหณิยา พฺราหฺมณานํ ปจฺโจโรหณี กลํ นาคฺฆติ โสฬสิํ.\csromanspacer{} อภิกฺกนฺตํ โภ โคตม ฯเปฯ อุปาสกํ มํ ภวํ โคตโม ธาเรตุ อชฺชตคฺเค ปาณุเปตํ สรณํ คตนฺติ.\csromanspacer{}\csromanspacer{}สตฺตมํ.}

\csromanheader{2}{8. ทุติยปจฺโจโรหณีสุตฺต}
\proseitem{120}{อริยํ โว ภิกฺขเว ปจฺโจโรหณิํ เทเสสฺสามิ, ตํ สุณาถ.\csromanspacer{} กตมา จ ภิกฺขเว อริยา ปจฺโจโรหณี? อิธ ภิกฺขเว อริยสาวโก อิติ ปฏิสญฺจิกฺขติ “มิจฺฉาทิฏฺฐิยา โข ปาปโก วิปาโก ทิฏฺเฐ เจว ธมฺเม อภิสมฺปรายญฺจา”ติ.\csromanspacer{} โส อิติ ปฏิสงฺขาย มิจฺฉาทิฏฺฐิํ ปชหติ.\csromanspacer{} มิจฺฉาทิฏฺฐิยา ปจฺโจโรหติ.\csromanspacer{} มิจฺฉาสงฺกปฺปสฺส โข ปาปโก วิปาโก.\csromanspacer{} มิจฺฉาวาจาย โข.\csromanspacer{} มิจฺฉากมฺมนฺตสฺส โข.\csromanspacer{} มิจฺฉา-อาชีวสฺส โข.\csromanspacer{} มิจฺฉาวายามสฺส โข.\csromanspacer{} มิจฺฉาสติยา โข.\csromanspacer{} มิจฺฉาสมาธิสฺส โข.\csromanspacer{} มิจฺฉาญาณสฺส โข.\csromanspacer{} มิจฺฉาวิมุตฺติยา โข ปาปโก วิปาโก ทิฏฺเฐ เจว ธมฺเม อภิสมฺปรายญฺจาติ.\csromanspacer{} โส อิติ ปฏิสงฺขาย มิจฺฉาวิมุตฺติํ ปชหติ, มิจฺฉาวิมุตฺติยา ปจฺโจโรหติ.\csromanspacer{} อยํ วุจฺจติ ภิกฺขเว อริยา ปจฺโจโรหณีติ.\csromanspacer{}\csromanspacer{}อฏฺฐมํ.}

\csromanheader{2}{9. ปุพฺพงฺคมสุตฺต}
\proseitem{121}{สูริยสฺส ภิกฺขเว อุทยโต เอตํ ปุพฺพงฺคมํ เอตํ ปุพฺพนิมิตฺตํ ยทิทํ อรุณุคฺคํ.\csromanspacer{} เอวเมวํ โข ภิกฺขเว กุสลานํ ธมฺมานํ เอตํ ปุพฺพงฺคมํ เอตํ ปุพฺพนิมิตฺตํ, ยทิทํ สมฺมาทิฏฺฐิ.\csromanspacer{} สมฺมาทิฏฺฐิกสฺส ภิกฺขเว สมฺมาสงฺกปฺโป ปโหติ, สมฺมาสงฺกปฺปสฺส สมฺมาวาจา ปโหติ, สมฺมาวาจสฺส}

\vspace{\tipitakaparskip}
\csromancenter{ทสกนิปาตปาฬิ สมฺมากมฺมนฺโต ปโหติ, สมฺมากมฺมนฺตสฺส สมฺมา-อาชีโว ปโหติ, สมฺมา-อาชีวสฺส สมฺมาวายาโม ปโหติ, สมฺมาวายามสฺส สมฺมาสติ ปโหติ, สมฺมาสติสฺส สมฺมาสมาธิ ปโหติ, สมฺมาสมาธิสฺส สมฺมาญาณํ ปโหติ, สมฺมาญาณิสฺส สมฺมาวิมุตฺติ ปโหตีติ.\csromanspacer{}\csromanspacer{}นวมํ.}
\csromanheader{2}{10. อาสวกฺขยสุตฺต}
\csromanmatikaanchor{mtk.166}
\csromanmatikaanchor{mtk.169}
\csromanmatikaanchor{mtk.172}
\csromantocmarkref{subsection}{8. ทุติยปจฺโจโรหณีสุตฺต}{mtk.166}
\bookmark[dest={mtk.166},level=2]{8. ทุติยปจฺโจโรหณีสุตฺต}
\proseitem{122}{\csromanfolio{450}ทสยิเม ภิกฺขเว ธมฺมา ภาวิตา พหุลีกตา อาสวานํ ขยาย สํวตฺตนฺติ.\csromanspacer{} กตเม ทส? สมฺมาทิฏฺฐิ สมฺมาสงฺกปฺโป สมฺมาวาจา สมฺมากมฺมนฺโต สมฺมา-อาชีโว สมฺมาวายาโม สมฺมาสติ สมฺมาสมาธิ สมฺมาญาณํ สมฺมาวิมุตฺติ.\csromanspacer{} อิเม โข ภิกฺขเว ทส ธมฺมา ภาวิตา พหุลีกตา อาสวานํ ขยาย สํวตฺตนฺตีติ.\csromanspacer{}\csromanspacer{}ทสมํ.\csromanspacer{} ปจฺโจโรหณิวคฺโค ทุติโย.}

\titlehead{\textbf{ตสฺสุทฺทานํ}}
\csromangathameasure{ตโย อธมฺมา อชิโต, สงฺคารโว จ โอริมํ. \\ เทฺว เจว ปจฺโจโรหณี, ปุพฺพงฺคมํ อาสวกฺขโยติ.}
\csromangathagroup{\csromangathastanza{ตโย อธมฺมา อชิโต, สงฺคารโว จ โอริมํ. \\ เทฺว เจว ปจฺโจโรหณี, ปุพฺพงฺคมํ อาสวกฺขโยติ.}}

\csromanheader{1}{\textbf{(13) 3. ปริสุทฺธวคฺค}}
\csromanheader{2}{1. ปฐมสุตฺต}
\proseitem{123}{ทสยิเม ภิกฺขเว ธมฺมา ปริสุทฺธา ปริโยทาตา นาญฺญตฺร สุคตวินยา.\csromanspacer{} กตเม ทส? สมฺมาทิฏฺฐิ สมฺมาสงฺกปฺโป สมฺมาวาจา สมฺมากมฺมนฺโต สมฺมา-อาชีโว สมฺมาวายาโม สมฺมาสติ สมฺมาสมาธิ สมฺมาญาณํ สมฺมาวิมุตฺติ.\csromanspacer{} อิเม โข ภิกฺขเว ทส ธมฺมา ปริสุทฺธา ปริโยทาตา นาญฺญตฺร สุคตวินยาติ.\csromanspacer{}\csromanspacer{}ปฐมํ.}

\csromanheader{2}{2. ทุติยสุตฺต}
\proseitem{124}{ทสยิเม ภิกฺขเว ธมฺมา อนุปฺปนฺนา อุปฺปชฺชนฺติ นาญฺญตฺร สุคตวินยา.\csromanspacer{} กตเม ทส? สมฺมาทิฏฺฐิ ฯเปฯ สมฺมาวิมุตฺติ.\csromanspacer{} อิเม โข ภิกฺขเว ทส ธมฺมา อนุปฺปนฺนา อุปฺปชฺชนฺติ นาญฺญตฺร สุคตวินยาติ.\csromanspacer{}\csromanspacer{}ทุติยํ.}

\csromanmatikaanchor{mtk.167}
\csromantocmarkref{subsection}{9. ปุพฺพงฺคมสุตฺต}{mtk.167}
\bookmark[dest={mtk.167},level=2]{9. ปุพฺพงฺคมสุตฺต}
\chapterheadpageread{(13) 3. ปริสุทฺธวคฺค \textbf{3. ตติยสุตฺต}}
\proseitem{125}{\csromanfolio{451}ทสยิเม ภิกฺขเว ธมฺมา มหปฺผลา มหานิสํสา นาญฺญตฺร สุคตวินยา.\csromanspacer{} กตเม ทส? สมฺมาทิฏฺฐิ ฯเปฯ สมฺมาวิมุตฺติ.\csromanspacer{} อิเม โข ภิกฺขเว ทส ธมฺมา มหปฺผลา มหานิสํสา นาญฺญตฺร สุคตวินยาติ.\csromanspacer{}\csromanspacer{}ตติยํ.}

\csromanheader{2}{4. จตุตฺถสุตฺต}
\csromanmatikaanchor{mtk.168}
\csromantocmarkref{subsection}{10. อาสวกฺขยสุตฺต}{mtk.168}
\bookmark[dest={mtk.168},level=2]{10. อาสวกฺขยสุตฺต}
\csromantocmarkref{section}{อุทฺทานคาถา}{mtk.169}
\bookmark[dest={mtk.169},level=1]{อุทฺทานคาถา}
\proseitem{126}{ทสยิเม ภิกฺขเว ธมฺมา ราควินยปริโยสานา โหนฺติ, โทสวินยปริโยสานา โหนฺติ, โมหวินยปริโยสานา โหนฺติ นาญฺญตฺร สุคตวินยา.\csromanspacer{} กตเม ทส? สมฺมาทิฏฺฐิ ฯเปฯ สมฺมาวิมุตฺติ.\csromanspacer{} อิเม โข ภิกฺขเว ทส ธมฺมา ราควินยปริโยสานา โหนฺติ, โทสวินยปริโยสานา โหนฺติ, โมหวินยปริโยสานา โหนฺติ นาญฺญตฺร สุคตวินยาติ.\csromanspacer{}\csromanspacer{}จตุตฺถํ.}

\csromanheader{2}{5. ปญฺจมสุตฺต}
\proseitem{127}{ทสยิเม ภิกฺขเว ธมฺมา เอกนฺตนิพฺพิทาย วิราคาย นิโรธาย อุปสมาย อภิญฺญาย สมฺโพธาย นิพฺพานาย สํวตฺตนฺติ นาญฺญตฺร สุคตวินยา.\csromanspacer{} กตเม ทส? สมฺมาทิฏฺฐิ ฯเปฯ สมฺมาวิมุตฺติ.\csromanspacer{} อิเม โข ภิกฺขเว ทส ธมฺมา เอกนฺตนิพฺพิทาย วิราคาย นิโรธาย อุปสมาย อภิญฺญาย สมฺโพธาย นิพฺพานาย สํวตฺตนฺติ นาญฺญตฺร สุคตวินยาติ.\csromanspacer{}\csromanspacer{}ปญฺจมํ.}

\csromanheader{2}{6. ฉฏฺฐสุตฺต}
\csromanmatikaanchor{mtk.170}
\csromantocmarknopage{section}{(13) 3. ปริสุทฺธวคฺค}{mtk.170}
\bookmark[dest={mtk.170},level=1]{(13) 3. ปริสุทฺธวคฺค}
\proseitem{128}{ทสยิเม ภิกฺขเว ธมฺมา ภาวิตา พหุลีกตา อนุปฺปนฺนา อุปฺปชฺชนฺติ นาญฺญตฺร สุคตวินยา.\csromanspacer{} กตเม ทส? สมฺมาทิฏฺฐิ ฯเปฯ สมฺมาวิมุตฺติ.\csromanspacer{} อิเม โข ภิกฺขเว ทส ธมฺมา ภาวิตา พหุลีกตา อนุปฺปนฺนา อุปฺปชฺชนฺติ นาญฺญตฺร สุคตวินยาติ.\csromanspacer{}\csromanspacer{}ฉฏฺฐํ.}

\csromanmatikaanchor{mtk.171}
\csromantocmarkref{subsection}{1. ปฐมสุตฺต}{mtk.171}
\bookmark[dest={mtk.171},level=2]{1. ปฐมสุตฺต}
\csromantocmarkref{subsection}{2-11. ทุติยาทิสุตฺตานิ}{mtk.172}
\bookmark[dest={mtk.172},level=2]{2-11. ทุติยาทิสุตฺตานิ}
\csromanheader{2}{7. สตฺตมสุตฺต}
\proseitem{129}{ทสยิเม ภิกฺขเว ธมฺมา ภาวิตา พหุลีกตา มหปฺผลา โหนฺติ มหานิสํสา นาญฺญตฺร สุคตวินยา.\csromanspacer{} กตเม ทส? สมฺมาทิฏฺฐิ ฯเปฯ สมฺมาวิมุตฺติ.\csromanspacer{} อิเม โข ภิกฺขเว ทส ธมฺมา ภาวิตา พหุลีกตา มหปฺผลา โหนฺติ มหานิสํสา นาญฺญตฺร สุคตวินยาติ.\csromanspacer{}\csromanspacer{}สตฺตมํ.}

\gambhira{ทสกนิปาตปาฬิ \textbf{8. อฏฺฐมสุตฺต}}
\proseitem{130}{\csromanfolio{452}ทสยิเม ภิกฺขเว ธมฺมา ภาวิตา มหุลีกตา ราควินยปริโยสานา โหนฺติ, โทสวินยปริโยสานา โหนฺติ, โมหวินยปริโยสานา โหนฺติ นาญฺญตฺร สุคตวินยา.\csromanspacer{} กตเม ทส? สมฺมาทิฏฺฐิ ฯเปฯ สมฺมาวิมุตฺติ.\csromanspacer{} อิเม โข ภิกฺขเว ทส ธมฺมา ภาวิตา พหุลีกตา ราควินยปริโยสานา โหนฺติ, โทสวินยปริโยสานา โหนฺติ, โมหวินยปริโยสานา โหนฺติ นาญฺญตฺร สุคตวินยาติ.\csromanspacer{}\csromanspacer{}อฏฺฐมํ.}

\csromanheader{2}{9. นวมสุตฺต}
\proseitem{131}{ทสยิเม ภิกฺขเว ธมฺมา ภาวิตา พหุลีกตา เอกนฺตนิพฺพิทาย วิราคาย นิโรธาย อุปสมาย อภิญฺญาย สมฺโพธาย นิพฺพานาย สํวตฺตนฺติ นาญฺญาตฺร สุคตวินยา.\csromanspacer{} กตเม ทส? สมฺมาทิฏฺฐิ ฯเปฯ สมฺมาวิมุตฺติ.\csromanspacer{} อิเม โข ภิกฺขเว ทส ธมฺมา ภาวิตา พหุลีกตา เอกนฺตนิพฺพิทาย วิราคาย นิโรธาย อุปสมาย อภิญฺญาย สมฺโพธาย นิพฺพานาย สํวตฺตนฺติ นาญฺญตฺร สุคตวินยาติ.\csromanspacer{}\csromanspacer{}นวมํ.}

\csromanheader{2}{10. ทสมสุตฺต}
\proseitem{132}{ทสยิเม ภิกฺขเว มิจฺฉตฺตา.\csromanspacer{} กตเม ทส? มิจฺฉาทิฏฺฐิ มิจฺฉาสงฺกปฺโป มิจฺฉาวาจา มิจฺฉากมฺมนฺโต มิจฺฉา-อาชีโว มิจฺฉาวายาโม มิจฺฉาสติ มิจฺฉาสมาธิ มิจฺฉาญาณํ มิจฺฉาวิมุตฺติ.\csromanspacer{} อิเม โข ภิกฺขเว ทส มิจฺฉตฺตาติ.\csromanspacer{}\csromanspacer{}ทสมํ.}

\csromanheader{2}{11. เอกาทสมสุตฺต}
\proseitem{133}{ทสยิเม ภิกฺขเว สมฺมตฺตา.\csromanspacer{} กตเม ทส? สมฺมาทิฏฺฐิ สมฺมาสงฺกปฺโป สมฺมาวาจา สมฺมากมฺมนฺโต สมฺมา-อาชีโว สมฺมาวายาโม สมฺมาสติ สมฺมาสมาธิ สมฺมาญาณํ สมฺมาวิมุตฺติ.\csromanspacer{} อิเม โข ภิกฺขเว ทส สมฺมตฺตาติ.\csromanspacer{}\csromanspacer{}เอกาทสมํ.}

\vspace{\tipitakaparskip}
\csromancenter{ปริสุทฺธวคฺโค ตติโย.}
\csromanheader{1}{\textbf{(14) 4. สาธุวคฺค} \textbf{(14) 4. สาธุวคฺค}}
\csromanheader{2}{1. สาธุสุตฺต}
\proseitem{134}{\csromanfolio{453}\csromansymbolfootnote{*}{อุปริ 482 ปิฏฺเฐปิ.}สาธุญฺจ โว ภิกฺขเว เทเสสฺสามิ อสาธุญฺจ, ตํ สุณาถ สาธุกํ มนสิ กโรถ, ภาสิสฺสามีติ. “เอวํ ภนฺเต”ติ โข เต ภิกฺขู ภควโต ปจฺจสฺโสสุํ.\csromanspacer{} ภควา เอตทโวจ\csromandash{}}

\prose{กตมญฺจ ภิกฺขเว อสาธุ? มิจฺฉาทิฏฺฐิ มิจฺฉาสงฺกปฺโป มิจฺฉาวาจา มิจฺฉากมฺมนฺโต มิจฺฉา-อาชีโว มิจฺฉาวายาโม มิจฺฉาสติ มิจฺฉาสมาธิ มิจฺฉาญาณํ มิจฺฉาวิมุตฺติ.\csromanspacer{} อิทํ วุจฺจติ ภิกฺขเว อสาธุ.\csromanspacer{} กตมญฺจ ภิกฺขเว สาธุ? สมฺมาทิฏฺฐิ สมฺมาสงฺกปฺโป สมฺมาวาจา สมฺมากมฺมนฺโต สมฺมา-อาชีโว สมฺมาวายาโม สมฺมาสติ สมฺมาสมาธิ สมฺมาญาณํ สมฺมาวิมุตฺติ.\csromanspacer{} อิทํ วุจฺจติ ภิกฺขเว สาธูติ.\csromanspacer{}\csromanspacer{}ปฐมํ.}

\csromanheader{2}{2. อริยธมฺมสุตฺต}
\proseitem{135}{อริยธมฺมญฺจ โว ภิกฺขเว เทเสสฺสามิ อนริยธมฺมญฺจ, ตํ สุณาถ ฯเปฯ.\csromanspacer{} กตโม จ ภิกฺขเว อนริโย ธมฺโม? มิจฺฉาทิฏฺฐิ ฯเปฯ มิจฺฉาวิมุตฺติ.\csromanspacer{} อยํ วุจฺจติ ภิกฺขเว อนริโย ธมฺโม.\csromanspacer{} กตโม จ ภิกฺขเว อริโย ธมฺโม? สมฺมาทิฏฺฐิ ฯเปฯ สมฺมาวิมุตฺติ.\csromanspacer{} อยํ วุจฺจติ ภิกฺขเว อริโย ธมฺโมติ.\csromanspacer{}\csromanspacer{}ทุติยํ.}

\csromanheader{2}{3. อกุสลสุตฺต}
\proseitem{136}{อกุสลญฺจ โว ภิกฺขเว เทเสสฺสามิ กุสลญฺจ, ตํ สุณาถ ฯเปฯ.\csromanspacer{} กตมญฺจ ภิกฺขเว อกุสลํ? มิจฺฉาทิฏฺฐิ ฯเปฯ มิจฺฉาวิมุตฺติ.\csromanspacer{} อิทํ วุจฺจติ ภิกฺขเว อกุสลํ.\csromanspacer{} กตมญฺจ ภิกฺขเว กุสลํ? สมฺมาทิฏฺฐิ ฯเปฯ สมฺมาวิมุตฺติ.\csromanspacer{} อิทํ วุจฺจติ ภิกฺขเว กุสลนฺติ.\csromanspacer{}\csromanspacer{}ตติยํ.}

\csromanheader{2}{4. อตฺถสุตฺต}
\proseitem{137}{อตฺถญฺจ โว ภิกฺขเว เทเสสฺสามิ อนตฺถญฺจ, ตํ สุณาถ ฯเปฯ.\csromanspacer{} กตโม จ ภิกฺขเว อนตฺโถ? มิจฺฉาทิฏฺฐิ ฯเปฯ มิจฺฉาวิมุตฺติ.\csromanspacer{} อยํ วุจฺจติ ภิกฺขเว อนตฺโถ.\csromanspacer{} กตโม จ ภิกฺขเว อตฺโถ? สมฺมาทิฏฺฐิ ฯเปฯ สมฺมาวิมุตฺติ.\csromanspacer{} อยํ วุจฺจติ ภิกฺขเว อตฺโถติ.\csromanspacer{}\csromanspacer{}จตุตฺถํ.}

\csromanheader{2}{ทสกนิปาตปาฬิ \textbf{5. ธมฺมสุตฺต}}
\proseitem{138}{\csromanfolio{454}ธมฺมญฺจ โว ภิกฺขเว เทเสสฺสามิ อธมฺมญฺจ, ตํ สุณาถ ฯเปฯ.\csromanspacer{} กตโม จ ภิกฺขเว อธมฺโม? มิจฺฉาทิฏฺฐิ ฯเปฯ มิจฺฉาวิมุตฺติ.\csromanspacer{} อยํ วุจฺจติ ภิกฺขเว อธมฺโม.\csromanspacer{} กตโม จ ภิกฺขเว ธมฺโม? สมฺมาทิฏฺฐิ ฯเปฯ สมฺมาวิมุตฺติ.\csromanspacer{} อยํ วุจฺจติ ภิกฺขเว ธมฺโมติ.\csromanspacer{}\csromanspacer{}ปญฺจมํ.}

\csromanmatikaanchor{mtk.173}
\csromantocmarknopage{section}{(14) 4. สาธุวคฺค}{mtk.173}
\bookmark[dest={mtk.173},level=1]{(14) 4. สาธุวคฺค}
\csromanheader{2}{6. สาสวสุตฺต}
\csromanmatikaanchor{mtk.174}
\csromantocmarkref{subsection}{1. สาธุสุตฺต}{mtk.174}
\bookmark[dest={mtk.174},level=2]{1. สาธุสุตฺต}
\proseitem{139}{สาสวญฺจ โว ภิกฺขเว ธมฺมํ เทเสสฺสามิ อนาสวญฺจ, ตํ สุณาถ ฯเปฯ.\csromanspacer{} กตโม จ ภิกฺขเว สาสโว ธมฺโม? มิจฺฉาทิฏฺฐิ ฯเปฯ มิจฺฉาวิมุตฺติ.\csromanspacer{} อยํ วุจฺจติ ภิกฺขเว สาสโว ธมฺโม.\csromanspacer{} กตโม จ ภิกฺขเว อนาสโว ธมฺโม? สมฺมาทิฏฺฐิ ฯเปฯ สมฺมาวิมุตฺติ.\csromanspacer{} อยํ วุจฺจติ ภิกฺขเว อนาสโว ธมฺโมติ.\csromanspacer{}\csromanspacer{}ฉฏฺฐํ.}

\csromanheader{2}{7. สาวชฺชสุตฺต}
\proseitem{140}{สาวชฺชญฺจ โว ภิกฺขเว ธมฺมํ เทเสสฺสามิ อนวชฺชญฺจ, ตํ สุณาถ ฯเปฯ.\csromanspacer{} กตโม จ ภิกฺขเว สาวชฺโช ธมฺโม? มิจฺฉาทิฏฺฐิ ฯเปฯ มิจฺฉาวิมุตฺติ.\csromanspacer{} อยํ วุจฺจติ ภิกฺขเว สาวชฺโช ธมฺโม.\csromanspacer{} กตโม จ ภิกฺขเว อนวชฺโช ธมฺโม? สมฺมาทิฏฺฐิ ฯเปฯ สมฺมาวิมุตฺติ.\csromanspacer{} อยํ วุจฺจติ ภิกฺขเว อนวชฺโช ธมฺโมติ.\csromanspacer{}\csromanspacer{}สตฺตมํ.}

\csromanmatikaanchor{mtk.175}
\csromantocmarkref{subsection}{2. อริยธมฺมสุตฺต}{mtk.175}
\bookmark[dest={mtk.175},level=2]{2. อริยธมฺมสุตฺต}
\csromanheader{2}{8. ตปนียสุตฺต}
\proseitem{141}{ตปนียญฺจ โว ภิกฺขเว ธมฺมํ เทเสสฺสามิ อตปนียญฺจ, ตํ สุณาถ ฯเปฯ.\csromanspacer{} กตโม จ ภิกฺขเว ตปนีโย ธมฺโม? มิจฺฉาทิฏฺฐิ ฯเปฯ มิจฺฉาวิมุตฺติ.\csromanspacer{} อยํ วุจฺจติ ภิกฺขเว ตปนีโย ธมฺโม.\csromanspacer{} กตโม จ ภิกฺขเว อตปนีโย ธมฺโม? สมฺมาทิฏฺฐิ ฯเปฯ สมฺมาวิมุตฺติ.\csromanspacer{} อยํ วุจฺจติ ภิกฺขเว อตปนีโย ธมฺโมติ.\csromanspacer{}\csromanspacer{}อฏฺฐมํ.}

\csromanmatikaanchor{mtk.176}
\csromantocmarkref{subsection}{3. อกุสลสุตฺต}{mtk.176}
\bookmark[dest={mtk.176},level=2]{3. อกุสลสุตฺต}
\csromanheader{2}{9. อาจยคามิสุตฺต}
\proseitem{142}{อาจยคามิญฺจ โว ภิกฺขเว ธมฺมํ เทเสสฺสามิ อปจยคามิญฺจ, ตํ สุณาถ ฯเปฯ.\csromanspacer{} กตโม จ ภิกฺขเว อาจยคามี ธมฺโม? มิจฺฉาทิฏฺฐิ ฯเปฯ มิจฺฉาวิมุตฺติ.\csromanspacer{} อยํ วุจฺจติ ภิกฺขเว อาจยคามี ธมฺโม.\csromanspacer{} กตโม จ}

\csromanmatikaanchor{mtk.177}
\csromantocmarkref{subsection}{4. อตฺถสุตฺต}{mtk.177}
\bookmark[dest={mtk.177},level=2]{4. อตฺถสุตฺต}
\vspace{\tipitakaparskip}
\csromancenter{\textbf{(15) 5.}\csromanspacer{}\textbf{ อริยวคฺค} ภิกฺขเว อปจยคามี ธมฺโม? สมฺมาทิฏฺฐิ ฯเปฯ สมฺมาวิมุตฺติ.\csromanspacer{} อยํ วุจฺจติ ภิกฺขเว อปจยคามี ธมฺโมติ.\csromanspacer{}\csromanspacer{}นวมํ.}
\csromanheader{2}{10. ทุกฺขุทฺรยสุตฺต}
\csromanmatikaanchor{mtk.178}
\csromantocmarkref{subsection}{5. ธมฺมสุตฺต}{mtk.178}
\bookmark[dest={mtk.178},level=2]{5. ธมฺมสุตฺต}
\proseitem{143}{\csromanfolio{455}ทุกฺขุทฺรยญฺจ โว ภิกฺขเว ธมฺมํ เทเสสฺสามิ สุขุทฺรยญฺจ, ตํ สุณาถ ฯเปฯ.\csromanspacer{} กตโม จ ภิกฺขเว ทุกฺขุทฺรโย ธมฺโม? มิจฺฉาทิฏฺฐิ ฯเปฯ มิจฺฉาวิมุตฺติ.\csromanspacer{} อยํ วุจฺจติ ภิกฺขเว ทุกฺขุทฺรโย ธมฺโม.\csromanspacer{} กตโม จ ภิกฺขเว สุขุทฺรโย ธมฺโม? สมฺมาทิฏฺฐิ ฯเปฯ สมฺมาวิมุตฺติ.\csromanspacer{} อยํ วุจฺจติ ภิกฺขเว สุขุทฺรโย ธมฺโมติ.\csromanspacer{}\csromanspacer{}ทสมํ.}

\csromanheader{2}{11. ทุกฺขวิปากสุตฺต}
\csromanmatikaanchor{mtk.179}
\csromantocmarkref{subsection}{6. สาสวสุตฺต}{mtk.179}
\bookmark[dest={mtk.179},level=2]{6. สาสวสุตฺต}
\proseitem{144}{ทุกฺขวิปากญฺจ โว ภิกฺขเว ธมฺมํ เทเสสฺสามิ สุขวิปากญฺจ, ตํ สุณาถฯเปฯ.\csromanspacer{} กตโม จ ภิกฺขเว ทุกฺขวิปาโก ธมฺโม? มิจฺฉาทิฏฺฐิ ฯเปฯ มิจฺฉาวิมุตฺติ.\csromanspacer{} อยํ วุจฺจติ ภิกฺขเว ทุกฺขวิปาโก ธมฺโม.\csromanspacer{} กตโม จ ภิกฺขเว สุขวิปาโก ธมฺโม? สมฺมาทิฏฺฐิ ฯเปฯ สมฺมาวิมุตฺติ.\csromanspacer{} อยํ วุจฺจติ ภิกฺขเว สุขวิปาโก ธมฺโมติ.\csromanspacer{}\csromanspacer{}เอกาทสมํ.\csromanspacer{} สาธุวคฺโค จตุตฺโถ.}
\csromansectionrule

\csromanheader{1}{\textbf{(15) 5. อริยวคฺค}}
\csromanmatikaanchor{mtk.180}
\csromantocmarkref{subsection}{7. สาวชฺชสุตฺต}{mtk.180}
\bookmark[dest={mtk.180},level=2]{7. สาวชฺชสุตฺต}
\csromanheader{2}{1. อริยมคฺคสุตฺต}
\proseitem{145}{อริยมคฺคญฺจ โว ภิกฺขเว ธมฺมํ เทเสสฺสามิ อนริยมคฺคญฺจ, ตํ สุณาถ ฯเปฯ.\csromanspacer{} กตโม จ ภิกฺขเว อนริโย มคฺโค? มิจฺฉาทิฏฺฐิ ฯเปฯ มิจฺฉาวิมุตฺติ.\csromanspacer{} อยํ วุจฺจติ ภิกฺขเว อนริโย มคฺโค.\csromanspacer{} กตโม จ ภิกฺขเว อริโย มคฺโค? สมฺมาทิฏฺฐิ ฯเปฯ สมฺมาวิมุตฺติ.\csromanspacer{} อยํ วุจฺจติ ภิกฺขเว อริโย มคฺโคติ.\csromanspacer{}\csromanspacer{}ปฐมํ.}

\csromanmatikaanchor{mtk.181}
\csromantocmarkref{subsection}{8. ตปนียสุตฺต}{mtk.181}
\bookmark[dest={mtk.181},level=2]{8. ตปนียสุตฺต}
\csromanheader{2}{2. กณฺหมคฺคสุตฺต}
\proseitem{146}{กณฺหมคฺคญฺจ โว ภิกฺขเว ธมฺมํ เทเสสฺสามิ สุกฺกมคฺคญฺจ, ตํ สุณาถ ฯเปฯ.\csromanspacer{} กตโม จ ภิกฺขเว กณฺหมคฺโค? มิจฺฉาทิฏฺฐิ ฯเปฯ มิจฺฉาวิมุตฺติ.}

\csromanmatikaanchor{mtk.182}
\csromantocmarkref{subsection}{9. อาจยคามิสุตฺต}{mtk.182}
\bookmark[dest={mtk.182},level=2]{9. อาจยคามิสุตฺต}
\gambhira{ทสกนิปาตปาฬิ}
\prose{\csromanfolio{456}อยํ วุจฺจติ ภิกฺขเว กณฺหมคฺโค.\csromanspacer{} กตโม จ ภิกฺขเว สุกฺกมคฺโค? สมฺมาทิฏฺฐิ ฯเปฯ สมฺมาวิมุตฺติ.\csromanspacer{} อยํ วุจฺจติ ภิกฺขเว สุกฺกมคฺโคติ.\csromanspacer{}\csromanspacer{}ทุติยํ.}

\csromanheader{2}{3. สทฺธมฺมสุตฺต}
\csromanmatikaanchor{mtk.183}
\csromantocmarkref{subsection}{10. ทุกฺขุทฺรยสุตฺต}{mtk.183}
\bookmark[dest={mtk.183},level=2]{10. ทุกฺขุทฺรยสุตฺต}
\proseitem{147}{สทฺธมฺมญฺจ โว ภิกฺขเว ธมฺมํ เทเสสฺสามิ อสทฺธมฺมญฺจ, ตํ สุณาถ ฯเปฯ.\csromanspacer{} กตโม จ ภิกฺขเว อสทฺธมฺโม? มิจฺฉาทิฏฺฐิ ฯเปฯ มิจฺฉาวิมุตฺติ.\csromanspacer{} อยํ วุจฺจติ ภิกฺขเว อสทฺธมฺโม.\csromanspacer{} กตโม จ ภิกฺขเว สทฺธมฺโม? สมฺมาทิฏฺฐิ ฯเปฯ สมฺมาวิมุตฺติ.\csromanspacer{} อยํ วุจฺจติ ภิกฺขเว สทฺธมฺโมติ.\csromanspacer{}\csromanspacer{}ตติยํ.}

\csromanheader{2}{4. สปฺปุริสธมฺมสุตฺต}
\csromanmatikaanchor{mtk.184}
\csromantocmarkref{subsection}{11. ทุกฺขวิปากสุตฺต}{mtk.184}
\bookmark[dest={mtk.184},level=2]{11. ทุกฺขวิปากสุตฺต}
\proseitem{148}{สปฺปุริสธมฺมญฺจ โว ภิกฺขเว เทเสสฺสามิ อสปฺปุริสธมฺมญฺจ, ตํ สุณาถ ฯเปฯ.\csromanspacer{} กตโม จ ภิกฺขเว อสปฺปุริสธมฺโม? มิจฺฉาทิฏฺฐิ ฯเปฯ มิจฺฉาวิมุตฺติ.\csromanspacer{} อยํ วุจฺจติ ภิกฺขเว อสปฺปุริสธมฺโม.\csromanspacer{} กตโม จ ภิกฺขเว สปฺปุริสธมฺโม? สมฺมาทิฏฺฐิ ฯเปฯ สมฺมาวิมุตฺติ.\csromanspacer{} อยํ วุจฺจติ ภิกฺขเว สปฺปุริสธมฺโมติ.\csromanspacer{}\csromanspacer{}จตุตฺถํ.}

\csromanheader{2}{5. อุปฺปาเทตพฺพสุตฺต}
\csromanmatikaanchor{mtk.185}
\csromantocmarknopage{section}{(15) 5. อริยวคฺค}{mtk.185}
\bookmark[dest={mtk.185},level=1]{(15) 5. อริยวคฺค}
\proseitem{149}{อุปฺปาเทตพฺพญฺจ โว ภิกฺขเว ธมฺมํ เทเสสฺสามิ น อุปฺปาเทตพฺพญฺจ, ตํ สุณาถ ฯเปฯ.\csromanspacer{} กตโม จ ภิกฺขเว น อุปฺปาเทตพฺโพ ธมฺโม? มิจฺฉาทิฏฺฐิ ฯเปฯ มิจฺฉาวิมุตฺติ.\csromanspacer{} อยํ วุจฺจติ ภิกฺขเว น อุปฺปาเทตพฺโพ ธมฺโม.\csromanspacer{} กตโม จ ภิกฺขเว อุปฺปาเทตพฺโพ ธมฺโม? สมฺมาทิฏฺฐิ ฯเปฯ สมฺมาวิมุตฺติ.\csromanspacer{} อยํ วุจฺจติ ภิกฺขเว อุปฺปาเทตพฺโพ ธมฺโมติ.\csromanspacer{}\csromanspacer{}ปญฺจมํ.}

\csromanmatikaanchor{mtk.186}
\csromantocmarkref{subsection}{1. อริยมคฺคสุตฺต}{mtk.186}
\bookmark[dest={mtk.186},level=2]{1. อริยมคฺคสุตฺต}
\csromanheader{2}{6. อาเสวิตพฺพสุตฺต}
\proseitem{150}{อาเสวิตพฺพญฺจ โว ภิกฺขเว ธมฺมํ เทเสสฺสามิ น อาเสวิตพฺพญฺจ, ตํ สุณาถ ฯเปฯ.\csromanspacer{} กตโม จ ภิกฺขเว น อาเสวิตพฺโพ ธมฺโม? มิจฺฉาทิฏฺฐิ ฯเปฯ มิจฺฉาวิมุตฺติ.\csromanspacer{} อยํ วุจฺจติ ภิกฺขเว น อาเสวิตพฺโพ ธมฺโม.\csromanspacer{} กตโม จ ภิกฺขเว อาเสวิตพฺโพ ธมฺโม? สมฺมาทิฏฺฐิ ฯเปฯ สมฺมาวิมุตฺติ.\csromanspacer{} อยํ วุจฺจติ ภิกฺขเว อาเสวิตพฺโพ ธมฺโมติ.\csromanspacer{}\csromanspacer{}ฉฏฺฐํ.}

\csromanmatikaanchor{mtk.187}
\csromantocmarkref{subsection}{2. กณฺหมคฺคสุตฺต}{mtk.187}
\bookmark[dest={mtk.187},level=2]{2. กณฺหมคฺคสุตฺต}
\csromanheader{2}{(15) 5. อริยวคฺค \textbf{7. ภาเวตพฺพสุตฺต}}
\proseitem{151}{\csromanfolio{457}ภาเวตพฺพญฺจ โว ภิกฺขเว ธมฺมํ เทเสสฺสามิ น ภาเวตพฺพญฺจ, ตํ สุณาถ ฯเปฯ.\csromanspacer{} กตโม จ ภิกฺขเว น ภาเวตพฺโพ ธมฺโม? มิจฺฉาทิฏฺฐิ ฯเปฯ มิจฺฉาวิมุตฺติ.\csromanspacer{} อยํ วุจฺจติ ภิกฺขเว น ภาเวตพฺโพ ธมฺโม.\csromanspacer{} กตโม จ ภิกฺขเว ภาเวตพฺโพ ธมฺโม? สมฺมาทิฏฺฐิ ฯเปฯ สมฺมาวิมุตฺติ.\csromanspacer{} อยํ วุจฺจติ ภิกฺขเว ภาเวตพฺโพ ธมฺโมติ.\csromanspacer{}\csromanspacer{}สตฺตมํ.}

\csromanheader{2}{8. พหุลีกาตพฺพสุตฺต}
\proseitem{152}{พหุลีกาตพฺพญฺจ โว ภิกฺขเว ธมฺมํ เทเสสฺสามิ น พหุลีกาตพฺพญฺจ, ตํ สุณาถ ฯเปฯ.\csromanspacer{} กตโม จ ภิกฺขเว น พหุลีกาตพฺโพ ธมฺโม? มิจฺฉาทิฏฺฐิ ฯเปฯ มิจฺฉาวิมุตฺติ.\csromanspacer{} อยํ วุจฺจติ ภิกฺขเว น พหุลีกาตพฺโพ ธมฺโม.\csromanspacer{} กตโม จ ภิกฺขเว พหุลีกาตพฺโพ ธมฺโม? สมฺมาทิฏฺฐิ ฯเปฯ สมฺมาวิมุตฺติ.\csromanspacer{} อยํ วุจฺจติ ภิกฺขเว พหุลีกาตพฺโพ ธมฺโมติ.\csromanspacer{}\csromanspacer{}อฏฺฐมํ.}

\csromanmatikaanchor{mtk.188}
\csromantocmarkref{subsection}{3. สทฺธมฺมสุตฺต}{mtk.188}
\bookmark[dest={mtk.188},level=2]{3. สทฺธมฺมสุตฺต}
\csromanheader{2}{9. อนุสฺสริตพฺพสุตฺต}
\proseitem{153}{อนุสฺสริตพฺพญฺจ โว ภิกฺขเว ธมฺมํ เทเสสฺสามิ น อนุสฺสริตพฺพญฺจ, ตํ สุณาถ ฯเปฯ.\csromanspacer{} กตโม จ ภิกฺขเว น อนุสฺสริตพฺโพ ธมฺโม? มิจฺฉาทิฏฺฐิ ฯเปฯ มิจฺฉาวิมุตฺติ.\csromanspacer{} อยํ วุจฺจติ ภิกฺขเว น อนุสฺสริตพฺโพ ธมฺโม.\csromanspacer{} กตโม จ ภิกฺขเว อนุสฺสริตพฺโพ ธมฺโม? สมฺมาทิฏฺฐิ ฯเปฯ สมฺมาวิมุตฺติ.\csromanspacer{} อยํ วุจฺจติ ภิกฺขเว อนุสฺสริตพฺโพ ธมฺโมติ.\csromanspacer{}\csromanspacer{}นวมํ.}

\csromanmatikaanchor{mtk.189}
\csromantocmarkref{subsection}{4. สปฺปุริสธมฺมสุตฺต}{mtk.189}
\bookmark[dest={mtk.189},level=2]{4. สปฺปุริสธมฺมสุตฺต}
\csromanheader{2}{10. สจฺฉิกาตพฺพสุตฺต}
\proseitem{154}{สจฺฉิกาตพฺพญฺจ โว ภิกฺขเว ธมฺมํ เทเสสฺสามิ น สจฺฉิกาตพฺพญฺจ, ตํ สุณาถ ฯเปฯ.\csromanspacer{} กตโม จ ภิกฺขเว น สจฺฉิกาตพฺโพ ธมฺโม? มิจฺฉาทิฏฺฐิ ฯเปฯ มิจฺฉาวิมุตฺติ.\csromanspacer{} อยํ วุจฺจติ ภิกฺขเว น สจฺฉิกาตพฺโพ ธมฺโม.\csromanspacer{} กตโม จ ภิกฺขเว สจฺฉิกาตพฺโพ ธมฺโม? สมฺมาทิฏฺฐิ ฯเปฯ สมฺมาวิมุตฺติ.\csromanspacer{} อยํ วุจฺจติ ภิกฺขเว สจฺฉิกาตพฺโพ ธมฺโมติ.\csromanspacer{}\csromanspacer{}ทสมํ.\csromanspacer{} อริยวคฺโค ปญฺจโม.}

\csromanmatikaanchor{mtk.190}
\csromantocmarkref{subsection}{5. อุปฺปาเทตพฺพสุตฺต}{mtk.190}
\bookmark[dest={mtk.190},level=2]{5. อุปฺปาเทตพฺพสุตฺต}
\vspace{\tipitakaparskip}
\csromancenter{\textbf{ตติโย ปณฺณาสโก สมตฺโต.}}
\chapterheadpageread{4. จตุตฺถปณฺณาสก}
\csromanmatikaanchor{mtk.191}
\csromantocmarkref{subsection}{6. อาเสวิตพฺพสุตฺต}{mtk.191}
\bookmark[dest={mtk.191},level=2]{6. อาเสวิตพฺพสุตฺต}
\csromanheader{1}{\textbf{(16) 1. ปุคฺคลวคฺค}}
\csromanheader{2}{1. เสวิตพฺพสุตฺต}
\csromanmatikaanchor{mtk.192}
\csromantocmarkref{subsection}{7. ภาเวตพฺพสุตฺต}{mtk.192}
\bookmark[dest={mtk.192},level=2]{7. ภาเวตพฺพสุตฺต}
\proseitem{155}{\csromanfolio{458}ทสหิ ภิกฺขเว ธมฺเมหิ สมนฺนาคโต ปุคฺคโล น เสวิตพฺโพ.\csromanspacer{} กตเมหิ ทสหิ? มิจฺฉาทิฏฺฐิโก โหติ, มิจฺฉาสงฺกปฺโป โหติ, มิจฺฉาวาโจ โหติ, มิจฺฉากมฺมนฺโต โหติ, มิจฺฉา-อาชีโว โหติ, มิจฺฉาวายาโม โหติ, มิจฺฉาสติ โหติ, มิจฺฉาสมาธิ โหติ, มิจฺฉาญาณี โหติ, มิจฺฉาวิมุตฺติ โหติ.\csromanspacer{} อิเมหิ โข ภิกฺขเว ทสหิ ธมฺเมหิ สมนฺนาคโต ปุคฺคโล น เสวิตพฺโพ.}

\prose{ทสหิ ภิกฺขเว ธมฺเมหิ สมนฺนาคโต ปุคฺคโล เสวิตพฺโพ.\csromanspacer{} กตเมหิ ทสหิ? สมฺมาทิฏฺฐิโก โหติ, สมฺมาสงฺกปฺโป โหติ, สมฺมาวาโจ โหติ, สมฺมากมฺมนฺโต โหติ, สมฺมา-อาชีโว โหติ, สมฺมาวายาโม โหติ, สมฺมาสติ โหติ, สมฺมาสมาธิ โหติ, สมฺมาญาณี โหติ, สมฺมาวิมุตฺติ โหติ.\csromanspacer{} อิเมหิ โข ภิกฺขเว ทสหิ ธมฺเมหิ สมนฺนาคโต ปุคฺคโล เสวิตพฺโพติ.}

\csromanmatikaanchor{mtk.193}
\csromantocmarkref{subsection}{8. พหุลีกาตพฺพสุตฺต}{mtk.193}
\bookmark[dest={mtk.193},level=2]{8. พหุลีกาตพฺพสุตฺต}
\csromanheader{2}{2-12. ภชิตพฺพาทิสุตฺต}
\proseitem{156-166}{ทสหิ ภิกฺขเว ธมฺเมหิ สมนฺนาคโต ปุคฺคโล น ภชิตพฺโพ ฯเปฯ.\csromanspacer{} ภชิตพฺโพ ฯเปฯ.\csromanspacer{} น ปยิรุปาสิตพฺโพ.\csromanspacer{} ปยิรุปาสิตพฺโพ ฯเปฯ.\csromanspacer{} น ปุชฺโช โหติ.\csromanspacer{} ปุชฺโช โหติ ฯเปฯ.\csromanspacer{} น ปาสํโส โหติ.\csromanspacer{} ปาสํโส โหติ ฯเปฯ.\csromanspacer{} อคารโว โหติ.\csromanspacer{} สคารโว โหติ ฯเปฯ.\csromanspacer{} อปฺปติสฺโส โหติ.\csromanspacer{} สปฺปติสฺโส โหติ ฯเปฯ.\csromanspacer{} น อาราธโก โหติ.\csromanspacer{} อาราธโก โหติ ฯเปฯ.\csromanspacer{} น วิสุชฺฌติ.\csromanspacer{} วิสุชฺฌติ ฯเปฯ.\csromanspacer{} มานํ นาธิโภติ.\csromanspacer{} มานํ อธิโภติ ฯเปฯ.\csromanspacer{} ปญฺญาย น วฑฺฒติ.\csromanspacer{} ปญฺญาย วฑฺฒติ ฯเปฯ.}

\csromanmatikaanchor{mtk.194}
\csromantocmarkref{subsection}{9. อนุสฺสริตพฺพสุตฺต}{mtk.194}
\bookmark[dest={mtk.194},level=2]{9. อนุสฺสริตพฺพสุตฺต}
\prose{พหุํ อปุญฺญํ ปสวติ.\csromanspacer{} พหุํ ปุญฺญํ ปสวติ, กตเมหิ ทสหิ? สมฺมาทิฏฺฐิโก โหติ, สมฺมาสงฺกปฺโป โหติ, สมฺมาวาโจ โหติ, สมฺมากมฺมนฺโต โหติ, สมฺมา-อาชีโว โหติ, สมฺมาวายาโม}

\vspace{\tipitakaparskip}
\csromancenter{\textbf{(17) 2.}\csromanspacer{}\textbf{ ชาณุสฺโสณิวคฺค} โหติ, สมฺมาสติ โหติ, สมฺมาสมาธิ โหติ สมฺมาญาณี โหติ, สมฺมาวิมุตฺติ โหติ.\csromanspacer{} อิเมหิ โข ภิกฺขเว ทสหิ ธมฺเมหิ สมนฺนาคโต ปุคฺคโล พหุํ ปุญฺญํ ปสวตีติ.}
\csromanmatikaanchor{mtk.195}
\csromantocmarkref{subsection}{10. สจฺฉิกาตพฺพสุตฺต}{mtk.195}
\bookmark[dest={mtk.195},level=2]{10. สจฺฉิกาตพฺพสุตฺต}
\csromanheader{2}{ปุคฺคลวคฺโค ปฐโม.}
\csromansectionrule
\csromanheader{1}{\textbf{(17) 2. ชาณุสฺโสณิวคฺค}}
\csromanheader{2}{1. พฺราหฺมณปจฺโจโรหณีสุตฺต}
\csromanmatikaanchor{mtk.196}
\csromantocmarknopage{chapter}{4. จตุตฺถปณฺณาสก}{mtk.196}
\bookmark[dest={mtk.196},level=0]{4. จตุตฺถปณฺณาสก}
\proseitem{167}{\csromanfolio{459}เตน โข ปน สมเยน ชาณุสฺโสณิ พฺราหฺมโณ ตทหุโปสเถ สีสํนฺหาโต นวํ โขมยุคํ นิวตฺโถ อลฺลกุสมุฏฺฐิํ อาทาย ภควโต อวิทูเร เอกมนฺตํ ฐิโต โหติ.}

\csromanmatikaanchor{mtk.197}
\csromantocmarknopage{section}{(16) 1. ปุคฺคลวคฺค}{mtk.197}
\bookmark[dest={mtk.197},level=1]{(16) 1. ปุคฺคลวคฺค}
\prose{อทฺทสา โข ภควา ชาณุสฺโสณิํ พฺราหฺมณํ ตทหุโปสเถ สีสํนฺหาตํ นวํ โขมยุคํ นิวตฺถํ อลฺลกุสมุฏฺฐิํ อาทาย เอกมนฺตํ ฐิตํ, ทิสฺวาน ชาณุสฺโสณิํ พฺราหฺมณํ เอตทโวจ “กิํ นุ ตฺวํ พฺราหฺมณ ตทหุโปสเถ สีสํนฺหาโต นวํ โขมยุคํ นิวตฺโถ อลฺลกุสมุฏฺฐิํ อาทาย เอกมนฺตํ ฐิโต, กิํ นฺวชฺช พฺราหฺมณกุลสฺสา”ติ.\csromanspacer{} ปจฺโจโรหณี โภ โคตม อชฺช พฺราหฺมณกุลสฺสาติ.}

\csromanmatikaanchor{mtk.198}
\csromantocmarkref{subsection}{1. เสวิตพฺพสุตฺต}{mtk.198}
\bookmark[dest={mtk.198},level=2]{1. เสวิตพฺพสุตฺต}
\prose{ยถา กถํ ปน พฺราหฺมณ พฺราหฺมณานํ ปจฺโจโรหณี โหตีติ, อิธ โภ โคตม พฺราหฺมณา ตทหุโปสเถ สีสํนฺหาตา นวํ โขมยุคํ นิวตฺถา อลฺเลน โคมเยน ปถวิํ โอปุญฺชิตฺวา หริเตหิ กุเสหิ ปตฺถริตฺวา อนฺตรา จ เวลํ อนฺตรา จ อคฺยาคารํ เสยฺยํ กปฺเปนฺติ, เต ตํ รตฺติํ ติกฺขตฺตุํ ปจฺจุฏฺฐาย ปญฺชลิกา อคฺคิํ นมสฺสนฺติ “ปจฺโจโรหาม ภวนฺตํ ปจฺโจโรหาม ภวนฺตนฺ”ติ, พหุเกน จ สปฺปิเตลนวนีเตน อคฺคิํ สนฺตปฺเปนฺติ, ตสฺสา จ รตฺติยา อจฺจเยน ปณีเตน ขาทนีเยน โภชนีเยน พฺราหฺมเณ สนฺตปฺเปนฺติ, เอวํ โภ โคตม พฺราหฺมณานํ ปจฺโจโรหณี โหตีติ.}

\prose{อญฺญถา โข พฺราหฺมณ พฺราหฺมณานํ ปจฺโจโรหณี โหติ, อญฺญถา จ ปน อริยสฺส วินเย ปจฺโจโรหณี โหตีติ.\csromanspacer{} ยถา}

\gambhira{ทสกนิปาตปาฬิ}
\csromanmatikaanchor{mtk.199}
\csromantocmarkref{subsection}{2-12. ภชิตพฺพาทิสุตฺต}{mtk.199}
\bookmark[dest={mtk.199},level=2]{2-12. ภชิตพฺพาทิสุตฺต}
\prose{\csromanfolio{460}กถํ ปน โภ โคตม อริยสฺส วินเย ปจฺโจโรหณี โหติ, สาธุ เม ภวํ โคตโม ตถา ธมฺมํ เทเสตุ, ยถา อริยสฺส วินเย ปจฺโจโรหณี โหตีติ.}

\prose{เตน หิ พฺราหฺมณ สุณาหิ สาธุกํ มนสิ กโรหิ, ภาสิสฺสามีติ. “เอวํ โภ”ติ โข ชาณุสฺโสณิ พฺราหฺมโณ ภควโต ปจฺจสฺโสสิ.\csromanspacer{} ภควา เอตทโวจ\csromandash{}}

\prose{อิธ พฺราหฺมณ อริยสาวโก อิติ ปฏิสญฺจิกฺขติ “ปาณาติปาตสฺส โข ปาปโก วิปาโก ทิฏฺเฐ เจว ธมฺเม อภิสมฺปรายญฺจา”ติ.\csromanspacer{} โส อิติ ปฏิสงฺขาย ปาณาติปาตํ ปชหติ, ปาณาติปาตา ปจฺโจโรหติ.}

\prose{“อทินฺนาทานสฺส โข ปาปโก วิปาโก ทิฏฺเฐ เจว ธมฺเม อภิสมฺปรายญฺจา”ติ.\csromanspacer{} โส อิติ ปฏิสงฺขาย อทินฺนาทานํ ปชหติ, อทินฺนาทานา ปจฺโจโรหติ.}

\prose{“กาเมสุมิจฺฉาจารสฺส โข ปาปโก วิปาโก ทิฏฺเฐ เจว ธมฺเม อภิสมฺปรายญฺจา”ติ.\csromanspacer{} โส อิติ ปฏิสงฺขาย กาเมสุมิจฺฉาจารํ ปชหติ, กาเมสุมิจฺฉาจารา ปจฺโจโรหติ.}

\csromanmatikaanchor{mtk.200}
\csromantocmarknopage{section}{(17) 2. ชาณุสฺโสณิวคฺค}{mtk.200}
\bookmark[dest={mtk.200},level=1]{(17) 2. ชาณุสฺโสณิวคฺค}
\prose{“มุสาวาทสฺส โข ปาปโก วิปาโก ทิฏฺเฐ เจว ธมฺเม อภิสมฺปรายญฺจา”ติ.\csromanspacer{} โส อิติ ปฏิสงฺขาย มุสาวาทํ ปชหติ, มุสาวาทา ปจฺโจโรหติ.}

\csromanmatikaanchor{mtk.201}
\csromantocmarkref{subsection}{1. พฺราหฺมณปจฺโจโรหณีสุตฺต}{mtk.201}
\bookmark[dest={mtk.201},level=2]{1. พฺราหฺมณปจฺโจโรหณีสุตฺต}
\prose{“ปิสุณาย วาจาย โข ปาปโก วิปาโก ทิฏฺเฐ เจว ธมฺเม อภิสมฺปรายญฺจา”ติ.\csromanspacer{} โส อิติ ปฏิสงฺขาย ปิสุณํ วาจํ ปชหติ, ปิสุณาย วาจาย ปจฺโจโรหติ.}

\prose{“ผรุสาย วาจาย โข ปาปโก วิปาโก ทิฏฺเฐ เจว ธมฺเม อภิสมฺปรายญฺจา”ติ.\csromanspacer{} โส อิติ ปฏิสงฺขาย ผรุสํ วาจํ ปชหติ, ผรุสาย วาจาย ปจฺโจโรหติ.}

\prose{“สมฺผปฺปลาปสฺส โข ปาปโก วิปาโก ทิฏฺเฐ เจว ธมฺเม อภิสมฺปรายญฺจา”ติ.\csromanspacer{} โส อิติ ปฏิสงฺขาย สมฺผปฺปลาปํ ปชหติ, สมฺผลาปา ปจฺโจโรหติ.}

\chapterheadpageread{(17) 2. ชาณุสฺโสณิวคฺค}
\prose{\csromanfolio{461}“อภิชฺฌาย โข ปาปโก วิปาโก ทิฏฺเฐ เจว ธมฺเม อภิสมฺปรายญฺจา”ติ.\csromanspacer{} โส อิติ ปฏิสงฺขาย อภิชฺฌํ ปชหติ, อภิชฺฌาย ปจฺโจโรหติ.}

\prose{“พฺยาปาทสฺส โข ปาปโก วิปาโก ทิฏฺเฐ เจว ธมฺเม อภิสมฺปรายญฺจา”ติ.\csromanspacer{} โส อิติ ปฏิสงฺขาย พฺยาปาทํ ปชหติ, พฺยาปาทา ปจฺโจโรหติ.}

\prose{“มิจฺฉาทิฏฺฐิยา โข ปาปโก วิปาโก ทิฏฺเฐ เจว ธมฺเม อภิสมฺปรายญฺจา”ติ, โส อิติ ปฏิสงฺขาย มิจฺฉาทิฏฺฐิํ ปชหติ, มิจฺฉาทิฏฺฐิยา ปจฺโจโรหติ.\csromanspacer{} เอวํ โข พฺราหฺมณ อริยสฺส วินเย ปจฺโจโรหณี โหตีติ.}

\prose{อญฺญถา โข โภ โคตม พฺราหฺมณานํ ปจฺโจโรหณี โหติ, อญฺญถา จ ปน อริยสฺส วินเย ปจฺโจโรหณี โหติ.\csromanspacer{} อิมิสฺสา โภ โคตม อริยสฺส วินเย ปจฺโจโรหณิยา พฺราหฺมณานํ ปจฺโจโรหณี กลํ นาคฺฆติ โสฬสิํ.\csromanspacer{} อภิกฺกนฺตํ โภ โคตม ฯเปฯ อุปาสกํ มํ ภวํ โคตโม ธาเรตุ อชฺชตคฺเค ปาณุเปตํ สรณํ คตนฺติ.\csromanspacer{}\csromanspacer{}ปฐมํ.}

\csromanheader{2}{2. อริยปจฺโจโรหณีสุตฺต}
\proseitem{168}{อริยํ โว ภิกฺขเว ปจฺโจโรหณิํ เทเสสฺสามิ, ตํ สุณาถ สาธุกํ มนสิ กโรถ, ภาสิสฺสามีติ. “เอวํ ภนฺเต”ติ โข เต ภิกฺขู ภควโต ปจฺจสฺโสสุํ.\csromanspacer{} ภควา เอตทโวจ\csromandash{}}

\prose{กตมา จ ภิกฺขเว อริยา ปจฺโจโรหณี, อิธ ภิกฺขเว อริยสาวโก อิติ ปฏิสญฺจิกฺขติ “ปาณาติปาตสฺส โข ปาปโก วิปาโก ทิฏฺเฐ เจว ธมฺเม อภิสมฺปรายญฺจา”ติ.\csromanspacer{} โส อิติ ปฏิสงฺขาย ปาณาติปาตํ ปชหติ, ปาณาติปาตา ปจฺโจโรหติ.}

\prose{“อทินฺนาทานสฺส โข ปาปโก วิปาโก ทิฏฺเฐ เจว ธมฺเม อภิสมฺปรายญฺจา”ติ, โส อิติ ปฏิสงฺขาย อทินฺนาทานํ ปชหติ, อทินฺนาทานา ปจฺโจโรหติ.}

\prose{กาเมสุมิจฺฉาจารสฺส โข ปาปโก วิปาโก ฯเปฯ กาเมสุมิจฺฉาจารา ปจฺโจโรหติ.}

\gambhira{ทสกนิปาตปาฬิ}
\prose{\csromanfolio{462}มุสาวาทสฺส โข ปาปโก วิปาโก ฯเปฯ มุสาวาทา ปจฺโจโรหติ.}

\prose{ปิสุณาย วาจาย โข ปาปโก วิปาโก ฯเปฯ ปิสุณาย วาจาย ปจฺโจโรหติ.}

\prose{ผรุสาย วาจาย โข ปาปโก วิปาโก ฯเปฯ ผรุสาย วาจาย ปจฺโจโรหติ.}

\prose{สมฺผปฺปลาปสฺส โข ปาปโก วิปาโก ฯเปฯ สมฺผปฺปลาปา ปจฺโจโรหติ.}

\prose{อภิชฺฌาย โข ปาปโก วิปาโก ฯเปฯ อภิชฺฌาย ปจฺโจโรหติ.}

\prose{พฺยาปาทสฺส โข ปาปโก วิปาโก ฯเปฯ พฺยาปาทา ปจฺโจโรหติ.}

\csromanmatikaanchor{mtk.202}
\csromantocmarkref{subsection}{2. อริยปจฺโจโรหณีสุตฺต}{mtk.202}
\bookmark[dest={mtk.202},level=2]{2. อริยปจฺโจโรหณีสุตฺต}
\prose{“มิจฺฉาทิฏฺฐิยา โข ปาปโก วิปาโก ทิฏฺเฐ เจว ธมฺเม อภิสมฺปรายญฺจา”ติ.\csromanspacer{} โส อิติ ปฏิสงฺขาย มิจฺฉาทิฏฺฐิํ ปชหติ, มิจฺฉาทิฏฺฐิยา ปจฺโจโรหติ.\csromanspacer{} อยํ วุจฺจติ ภิกฺขเว อริยา ปจฺโจโรหณีติ.\csromanspacer{}\csromanspacer{}ทุติยํ.}

\csromanheader{2}{3. สงฺคารวสุตฺต}
\proseitem{169}{\csromansymbolfootnote{*}{เหฏฺฐา 445 ปิฏฺเฐปิ.}อถ โข สงฺคารโว พฺราหฺมโณ เยน ภควา เตนุปสงฺกมิ, อุปสงฺกมิตฺวา ภควตา สทฺธิํ สมฺโมทิ, สมฺโมทนียํ กถํ สารณียํ วีติสาเรตฺวา เอกมนฺตํ นิสีทิ, เอกมนฺตํ นิสินฺโน โข สงฺคารโว พฺราหฺมโณ ภควนฺตํ เอตทโวจ\csromandash{}}

\prose{กิํ นุ โข โภ โคตม โอริมํ ตีรํ, กิํ ปาริมํ ตีรนฺติ.\csromanspacer{} ปาณาติปาโต โข พฺราหฺมณ โอริมํ ตีรํ, ปาณาติปาตา เวรมณี ปาริมํ ตีรํ.\csromanspacer{} อทินฺนาทานํ โข พฺราหฺมณ โอริมํ ตีรํ, อทินฺนาทานา เวรมณี ปาริมํ ตีรํ.\csromanspacer{} กาเมสุมิจฺฉาจาโร โอริมํ ตีรํ, กาเมสุมิจฺฉาจารา เวรมณี ปาริมํ ตีรํ.\csromanspacer{} มุสาวาโท โอริมํ ตีรํ, มุสาวาทา เวรมณี ปาริมํ ตีรํ.\csromanspacer{} ปิสุณา วาจา โอริมํ ตีรํ, ปิสุณาย วาจาย เวรมณี ปาริมํ ตีรํ.\csromanspacer{} ผรุสา วาจา โอริมํ ตีรํ, ผรุสาย วาจาย เวรมณี ปาริมํ ตีรํ.\csromanspacer{} สมฺผปฺปลาโป โอริมํ ตีรํ, สมฺผปฺปลาปา เวรมณี ปาริมํ ตีรํ.\csromanspacer{} อภิชฺฌา โอริมํ ตีรํ, อนภิชฺฌา ปาริมํ ตีรํ.\csromanspacer{} พฺยาปาโท โอริมํ ตีรํ, อพฺยาปาโท ปาริมํ ตีรํ.\csromanspacer{} มิจฺฉาทิฏฺฐิ โอริมํ ตีรํ, สมฺมาทิฏฺฐิ ปาริมํ ตีรํ.\csromanspacer{} อิทํ โข พฺราหฺมณ โอริมํ ตีรํ, อิทํ ปาริมํ ตีรนฺติ.}

\chapterheadpageread{(17) 2. ชาณุสฺโสณิวคฺค}
\csromangathameasure{อปฺปกา เต มนุสฺเสสุ, เย ชนา ปารคามิโน. \\ อถายํ อิตรา ปชา, ตีรเมวานุธาวติ. \\ เย จ โข สมฺมทกฺขาเต, ธมฺเม ธมฺมานุวตฺติโน. \\ เต ชนา ปารเมสฺสนฺติ, มจฺจุเธยฺยํ สุทุตฺตรํ. \\ กณฺหํ ธมฺมํ วิปฺปหาย, สุกฺกํ ภาเวถ ปณฺฑิโต. \\ โอกา อโนก มาคมฺม, วิเวเก ยตฺถ ทูรมํ. \\ ตตฺราภิรติมิจฺเฉยฺย, หิตฺวา กาเม อกิญฺจโน. \\ ปริโยทเปยฺย อตฺตานํ, จิตฺตเกฺลเสหิ ปณฺฑิโต. \\ เยสํ สมฺโพธิยงฺเคสุ, สมฺมา จิตฺตํ สุภาวิตํ. \\ อาทานปฏินิสฺสคฺเค, อนุปาทาย เย รตา. \\ ขีณาสวา ชุติมนฺโต, เต โลเก ปรินิพฺพุตาติ.ตติยํ.}
\csromangathagroup{\csromangathastanza{\csromanfolio{463}อปฺปกา เต มนุสฺเสสุ, เย ชนา ปารคามิโน. \\ อถายํ อิตรา ปชา, ตีรเมวานุธาวติ.}\csromangathastanza{เย จ โข สมฺมทกฺขาเต, ธมฺเม ธมฺมานุวตฺติโน. \\ เต ชนา ปารเมสฺสนฺติ, มจฺจุเธยฺยํ สุทุตฺตรํ.}\csromangathastanza{กณฺหํ ธมฺมํ วิปฺปหาย, สุกฺกํ ภาเวถ ปณฺฑิโต. \\ โอกา อโนก มาคมฺม, วิเวเก ยตฺถ ทูรมํ.}\csromangathastanza{ตตฺราภิรติมิจฺเฉยฺย, หิตฺวา กาเม อกิญฺจโน. \\ ปริโยทเปยฺย อตฺตานํ, จิตฺตเกฺลเสหิ ปณฺฑิโต.}\csromangathastanza{เยสํ สมฺโพธิยงฺเคสุ, สมฺมา จิตฺตํ สุภาวิตํ. \\ อาทานปฏินิสฺสคฺเค, อนุปาทาย เย รตา. \\ ขีณาสวา ชุติมนฺโต, เต โลเก ปรินิพฺพุตาติ.\csromanspacer{}\csromanspacer{}ตติยํ.}}

\csromanheader{2}{4. โอริมสุตฺต}
\proseitem{170}{โอริมญฺจ ภิกฺขเว ตีรํ เทเสสฺสามิ ปาริมญฺจ ตีรํ, ตํ สุณาถ ฯเปฯ.\csromanspacer{} กตมญฺจ ภิกฺขเว โอริมํ ตีรํ, กตมญฺจ ปาริมํ ตีรํ, ปาณาติปาโต ภิกฺขเว โอริมํ ตีรํ, ปาณาติปาตา เวรมณี ปาริมํ ตีรํ.\csromanspacer{} อทินฺนาทานํ โอริมํ ตีรํ, อทินฺนาทานา เวรมณี ปาริมํ ตีรํ.\csromanspacer{} กาเมสุมิจฺฉาจาโร โอริมํ ตีรํ, กาเมสุมิจฺฉาจารา เวรมณี ปาริมํ ตีรํ.\csromanspacer{} มุสาวาโท โอริมํ ตีรํ, มุสาวาทา เวรมณี ปาริมํ ตีรํ.\csromanspacer{} ปิสุณา วาจา โอริมํ ตีรํ, ปิสุณาย วาจาย เวรมณี ปาริมํ ตีรํ.\csromanspacer{} ผรุสา วาจา โอริมํ ตีรํ, ผรุสาย วาจาย เวรมณี ปาริมํ ตีรํ, สมฺผปฺปลาโป โอริมํ ตีรํ, สมฺผปฺปลาปา เวรมนี ปาริมํ ตีรํ.\csromanspacer{} อภิชฺฌา โอริมํ ตีรํ, อนภิชฺฌา ปาริมํ ตีรํ.\csromanspacer{} พฺยาปาโท โอริมํ ตีรํ, อพฺยาปาโท ปาริมํ ตีรํ.\csromanspacer{} มิจฺฉาทิฏฺฐิ โอริมํ ตีรํ, สมฺมาทิฏฺฐิ ปาริมํ ตีรํ.\csromanspacer{} อิทํ โข ภิกฺขเว โอริมํ ตีรํ, อิทํ ปาริมํ ตีรนฺติ.}

\csromanmatikaanchor{mtk.203}
\csromantocmarkref{subsection}{3. สงฺคารวสุตฺต}{mtk.203}
\bookmark[dest={mtk.203},level=2]{3. สงฺคารวสุตฺต}
\csromangathameasure{อปฺปกา เต มนุสฺเสสุ, เย ชนา ปารคามิโน. \\ อถายํ อิตรา ปชา, ตีรเมวานุธาวติ. \\ เย จ โข สมฺมทกฺขาเต, ธมฺเม ธมฺมานุวตฺติโน. \\ เต ชนา ปารเมสฺสนฺติ, มจฺจุเธยฺยํ สุทุตฺตรํ.}
\csromangathagroup{\csromangathastanza{อปฺปกา เต มนุสฺเสสุ, เย ชนา ปารคามิโน. \\ อถายํ อิตรา ปชา, ตีรเมวานุธาวติ.}\csromangathastanza{เย จ โข สมฺมทกฺขาเต, ธมฺเม ธมฺมานุวตฺติโน. \\ เต ชนา ปารเมสฺสนฺติ, มจฺจุเธยฺยํ สุทุตฺตรํ.}}

\gambhira{ทสกนิปาตปาฬิ}
\csromangathameasure{กณฺหํ ธมฺมํ วิปฺปหาย, สุกฺกํ ภาเวถ ปณฺฑิโต. \\ โอกา อโนก มาคมฺม, วิเวเก ยตฺถ ทูรมํ. \\ ตตฺราภิรติมิจฺเฉยฺย, หิตฺวา กาเม อกิญฺจโน. \\ ปริโยทเปยฺย อตฺตานํ, จิตฺตเกฺลเสหิ ปณฺฑิโต. \\ เยสํ สมฺโพธิยงฺเคสุ, สมฺมา จิตฺตํ สุภาวิตํ. \\ อาทานปฏินิสฺสคฺเค, อนุปาทาย เย รตา. \\ ขีณาสวา ชุติมนฺโต, เต โลเก ปรินิพฺพุตาติ.จตุตฺถํ.}
\csromangathagroup{\csromangathastanza{\csromanfolio{464}กณฺหํ ธมฺมํ วิปฺปหาย, สุกฺกํ ภาเวถ ปณฺฑิโต. \\ โอกา อโนก มาคมฺม, วิเวเก ยตฺถ ทูรมํ.}\csromangathastanza{ตตฺราภิรติมิจฺเฉยฺย, หิตฺวา กาเม อกิญฺจโน. \\ ปริโยทเปยฺย อตฺตานํ, จิตฺตเกฺลเสหิ ปณฺฑิโต.}\csromangathastanza{เยสํ สมฺโพธิยงฺเคสุ, สมฺมา จิตฺตํ สุภาวิตํ. \\ อาทานปฏินิสฺสคฺเค, อนุปาทาย เย รตา. \\ ขีณาสวา ชุติมนฺโต, เต โลเก ปรินิพฺพุตาติ.\csromanspacer{}\csromanspacer{}จตุตฺถํ.}}

\csromanheader{2}{5. ปฐม-อธมฺมสุตฺต}
\proseitem{171}{\csromansymbolfootnote{*}{เหฏฺฐา 437 ปิฏฺเฐปิ.}อธมฺโม จ ภิกฺขเว เวทิตพฺโพ อนตฺโถ จ, ธมฺโม จ เวทิตพฺโพ อตฺโถ จ.\csromanspacer{} อธมฺมญฺจ วิทิตฺวา อนตฺถญฺจ, ธมฺมญฺจ วิทิตฺวา อตฺถญฺจ ยถา ธมฺโม ยถา อตฺโถ, ตถา ปฏิปชฺชิตพฺพํ.}

\prose{กตโม จ ภิกฺขเว อธมฺโม จ อนตฺโถ จ? ปาณาติปาโต อทินฺนาทานํ กาเมสุมิจฺฉาจาโร มุสาวาโท ปิสุณา วาจา ผรุสา วาจา สมฺผปฺปลาโป อภิชฺฌา พฺยาปาโท มิจฺฉาทิฏฺฐิ.\csromanspacer{} อยํ วุจฺจติ ภิกฺขเว อธมฺโม จ อนตฺโถ จ.}

\prose{กตโม จ ภิกฺขเว ธมฺโม จ อตฺโถ จ? ปาณาติปาตา เวรมณี, อทินฺนาทานา เวรมณี, กาเมสุมิจฺฉาจารา เวรมณี, มุสาวาทา เวรมนี, ปิสุณาย วาจาย เวรมณี, ผรุสาย วาจาย เวรมณี, สมฺผปฺปลาปา เวรมณี, อนภิชฺฌา อพฺยาปาโท สมฺมาทิฏฺฐิ.\csromanspacer{} อยํ วุจฺจติ ภิกฺขเว ธมฺโม จ อตฺโถ จ.}

\csromanmatikaanchor{mtk.204}
\csromantocmarkref{subsection}{4. โอริมสุตฺต}{mtk.204}
\bookmark[dest={mtk.204},level=2]{4. โอริมสุตฺต}
\prose{“อธมฺโม จ ภิกฺขเว เวทิตพฺโพ อนตฺโถ จ, ธมฺโม จ เปทิตพฺโพ อตฺโถ จ.\csromanspacer{} อธมฺมญฺจ วิทิตฺวา อนตฺถญฺจ, ธมฺมญฺจ วิทิตฺวา อตฺถญฺจ ยถา ธมฺโม ยถา อตฺโถ, ตถา ปฏิปชฺชิตพฺพนฺ”ติ อิติ ยํ ตํ วุตฺตํ, อิทเมตํ ปฏิจฺจ วุตฺตนฺติ.\csromanspacer{}\csromanspacer{}ปญฺจมํ.}

\csromanheader{2}{6. ทุติย-อธมฺมสุตฺต}
\proseitem{172}{อธมฺโม จ ภิกฺขเว เวทิตพฺโพ ธมฺโม จ, อนตฺโถ จ เวทิตพฺโพ อตฺโถ จ.\csromanspacer{} อธมฺมญฺจ วิทิตฺวา ธมฺมญฺจ, อนตฺถญฺจ วิทิตฺวา อตฺถญฺจ ยถา}

\vspace{\tipitakaparskip}
\csromancenter{(17) 2.\csromanspacer{} ชาณุสฺโสณิวคฺค ธมฺโม ยถา อตฺโถ, ตถา ปฏิปชฺชิตพฺพนฺติ อิทมโวจ ภควา, อิทํ วตฺวาน สุคโต อุฏฺฐายาสนา วิหารํ ปาวิสิ.}
\prose{\csromanfolio{465}อถ โข เตสํ ภิกฺขูนํ อจิรปกฺกนฺตสฺส ภควโต เอตทโหสิ “อิทํ โข โน อาวุโส ภควา สํขิตฺเตน อุทฺเทสํ อุทฺทิสิตฺวา วิตฺถาเรน อตฺถํ อวิภชิตฺวา อุฏฺฐายาสนา วิหารํ ปวิฏฺโฐ ‘อธมฺโม จ ภิกฺขเว เวทิตพฺโพ ธมฺโม จ, อนตฺโถ จ เวทิตพฺโพ อตฺโถ จ.\csromanspacer{} อธมฺมญฺจ วิทิตฺวา ธมฺมญฺจ, อนตฺถญฺจ วิทิตฺวา อตฺถญฺจ ยถา ธมฺโม ยถา อตฺโถ, ตถา ปฏิปชฺชิตพฺพนฺ’ติ.\csromanspacer{} โก นุ โข อิมสฺส ภควาตา สํขิตฺเตน อุทฺเทสสฺส อุทฺทิฏฺฐสฺส วิตฺถาเรน อตฺถํ อวิภตฺตสฺส วิตฺถาเรน อตฺถํ วิภเชยฺยา”ติ.}

\prose{อถ โข เตสํ ภิกฺขูนํ เอตทโหสิ “อยํ โข อายสฺมา มหากจฺจาโน สตฺถุ เจว สํวณฺณิโต สมฺภาวิโต จ วิญฺญูนํ สพฺรหฺมจารีนํ, ปโหติ จายสฺมา มหากจฺจาโน อิมสฺส ภควตา สํขิตฺเตน อุทฺเทสสฺส อุทฺทิฏฺฐสฺส วิตฺถาเรน อตฺถํ อวิภตฺตสฺส วิตฺถาเรน อตฺถํ วิภชิตุํ, ยํนูน มยํ เยนายสฺมา มหากจฺจาโน เตนุปสงฺกเมยฺยาม, อุปสงฺกมิตฺวา อายสฺมนฺตํ มหากจฺจานํ เอตมตฺถํ ปุจฺเฉยฺยาม.\csromanspacer{} ยถา โน อายสฺมา มหากจฺจาโน พฺยากริสฺสติ, ตถา นํ ธาเรสฺสามา”ติ.}

\prose{อถ โข เต ภิกฺขู เยนายสฺมา มหากจฺจาโน เตนุปสงฺกมิํสุ, อุปสงฺกมิตฺวา อายสฺมตา มหากจฺจาเนน สทฺธิํ สมฺโมทิํสุ, สมฺโมทนียํ กถํ สารณียํ วีติสาเรตฺวา เอกมนฺตํ นิสีทิํสุ, เอกมนฺตํ นิสินฺนา โข เต ภิกฺขู อายสฺมนฺตํ มหากจฺจานํ เอตทโวจุํ\csromandash{}}

\prose{อิทํ โข โน อาวุโส กจฺจาน ภควา สํขิตฺเตน อุทฺเทสํ อุทฺทิสิตฺวา วิตฺถาเรน อตฺถํ อวิภชิตฺวา อุฏฺฐายาสนา วิหารํ ปวิฏฺโฐ, “อธมฺโม จ ภิกฺขเว เวทิตพฺโพ ธมฺโม จ, อนตฺโถ จ เวทิตพฺโพ อตฺโถ จ, อธมฺมญฺจ วิทิตฺวา ธมฺมญฺจ, อนตฺถญฺจ วิทิตฺวา อตฺถญฺจ ยถา ธมฺโม ยถา อตฺโถ, ตถา ปฏิปชฺชิตพฺพนฺ”ติ.}

\csromanmatikaanchor{mtk.205}
\csromantocmarkref{subsection}{5. ปฐม-อธมฺมสุตฺต}{mtk.205}
\bookmark[dest={mtk.205},level=2]{5. ปฐม-อธมฺมสุตฺต}
\prose{เตสํ โน อาวุโส อมฺหากํ อจิรปกฺกนฺตสฺส ภควโต เอตทโหสิ “อิทํ โข โน อาวุโส ภควา สํขิตฺเตน อุทฺเทสํ อุทฺทิสิตฺวา วิตฺถาเรน อตฺถํ อวิภชิตฺวา อุฏฺฐายาสนา วิหารํ ปวิฏฺโฐ, ‘อธมฺโม จ ภิกฺขเว ฯเปฯ ตถา ปฏิปชฺชิตพฺพนฺ’ติ, โก นุ โข อิมสฺส ภควตา}

\gambhira{ทสกนิปาตปาฬิ}
\prose{\csromanfolio{466}สํขิตฺเตน อุทฺเทสสฺส อุทฺทิฏฺฐสฺส วิตฺถาเรน อตฺถํ อวิภตฺตสฺส วิตฺถาเรน อตฺถํ วิภเชยฺยา”ติ.}

\prose{เตสํ โน อาวุโส อมฺหากํ เอตทโหสิ “อยํ โข อายสฺมา มหากจฺจาโน สตฺถุ เจว สํวณฺณิโต สมฺภาวิโต จ วิญฺญูนํ สพฺรหฺมจารีนํ, ปโหติ จายสฺมา มหากจฺจาโน อิมสฺส ภควตา สํขิตฺเตน อุทฺเทสสฺส อุทฺทิฏฺฐสฺส วิตฺถาเรน อตฺถํ อวิภตฺตสฺส วิตฺถาเรน อตฺถํ วิภชิตุํ, ยํนูน มยํ เยนายสฺมา มหากจฺจาโน เตนุปสงฺกเมยฺยาม, อุปสงฺกมิตฺวา อายสฺมนฺตํ มหากจฺจานํ เอตมตฺถํ ปฏิปุจฺเฉยฺยาม.\csromanspacer{} ยถา โน อายสฺมา มหากจฺจาโน พฺยากริสฺสติ, ตถา นํ ธาเรสฺสามา”ติ.\csromanspacer{} วิภชตุ อายสฺมา มหากจฺจาโนติ.}

\prose{เสยฺยถาปิ อาวุโส ปุริโส สารตฺถิโก สารํ คเวสี สารปริเยสนํ จรมาโน มหโต รุกฺขสฺส ติฏฺฐโต สารวโต อติกฺกมฺเมว มูลํ อติกฺกมฺม ขนฺธํ สาขาปลาเส สารํ ปริเยสิตพฺพํ มญฺเญยฺย, เอวํสมฺปทมิทํ อายสฺมนฺตานํ สตฺถริ สมฺมุขีภูเต ตํ ภควนฺตํ อติสิตฺวา อเมฺห เอตมตฺถํ ปฏิปุจฺฉิตพฺพํ มญฺญถ\footnote{มญฺเญถ (สี), มญฺเญยฺยาถ (ก)}.\csromanspacer{} โส หาวุโส ภควา ชานํ ชานาติ, ปสฺสํ ปสฺสติ, จกฺขุภูโต ญาณภูโต ธมฺมภูโต พฺรหฺมภูโต วตฺตา ปวตฺตา อตฺถสฺส นินฺเนตา อมตสฺส ทาตา ธมฺมสฺสามี ตถาคโต, โส เจว ปเนตสฺส กาโล อโหสิ.\csromanspacer{} ยํ ตุเมฺห ภควนฺตํเยว อุปสงฺกมิตฺวา เอตมตฺถํ ปฏิปุจฺเฉยฺยาถ.\csromanspacer{} ยถา โว ภควา พฺยากเรยฺย, ตถา นํ ธาเรยฺยาถาติ.}

\csromanmatikaanchor{mtk.206}
\csromantocmarkref{subsection}{6. ทุติย-อธมฺมสุตฺต}{mtk.206}
\bookmark[dest={mtk.206},level=2]{6. ทุติย-อธมฺมสุตฺต}
\prose{อทฺธา อาวุโส กจฺจาน ภควา ชานํ ชานาติ, ปสฺสํ ปสฺสติ, จกฺขุภูโต ญาณภูโต ธมฺมภูโต พฺรหฺมภูโต วตฺตา ปวตฺตา อตฺถสฺส นินฺเนตา อมตสฺส ทาตา ธมฺมสฺสามี ตถาคโต, โส เจว ปเนตสฺส กาโล อโหสิ.\csromanspacer{} ยํ มยํ ภควนฺตํเยว อุปสงฺกมิตฺวา เอตมตฺถํ ปฏิปุจฺเฉยฺยาม.\csromanspacer{} ยถา โน ภควา พฺยากเรยฺย, ตถา นํ ธาเรยฺยาม.\csromanspacer{} อปิ จายสฺมา มหากจฺจาโน สตฺถุ เจว สํวณฺณิโต สมฺภาวิโต จ วิญฺญูนํ สพฺรหฺมจารีนํ, ปโหติ จายสฺมา มหากจฺจาโน อิมสฺส ภควตา สํขิตฺเตน อุทฺเทสสฺส อุทฺทิฏฺฐสฺส วิตฺถาเรน อตฺถํ อวิภตฺตสฺส วิตฺถาเรน อตฺถํ วิภชิตุํ, วิภชตายสฺมา มหากจฺจาโน อครุํ กริตฺวาติ.}

\chapterheadpageread{(17) 2. ชาณุสฺโสณิวคฺค}
\prose{\csromanfolio{467}เตน หาวุโส สุณาถ สาธุกํ มนสิ กโรถ, ภาสิสฺสามีติ. “เอวํ อาวุโส”ติ โข เต ภิกฺขู อายสฺมโต มหากจฺจานสฺส ปจฺจสฺโสสุํ.\csromanspacer{} อถายสฺมา มหากจฺจาโน เอตทโวจ\csromandash{}}

\prose{ยํ โข โน อาวุโส ภควา สํขิตฺเตน อุทฺเทสํ อุทฺทิสิตฺวา วิตฺถาเรน อตฺถํ อวิภชิตฺวา อุฏฺฐายาสนา วิหารํ ปวิฏฺโฐ, “อธมฺโม จ ภิกฺขเว เวทิตพฺโพ ฯเปฯ ตถา ปฏิปชฺชิตพฺพนฺ”ติ.}

\prose{กตโม จาวุโส อธมฺโม, กตโม จ ธมฺโม, กตโม จ อนตฺโถ, กตโม จ อตฺโถ? ปาณาติปาโต อาวุโส อธมฺโม, ปาณาติปาตา เวรมณี ธมฺโม.\csromanspacer{} เย จ ปาณาติปาตปจฺจยา อเนเก ปาปกา อกุสลา ธมฺมา สมฺภวนฺติ, อยํ อนตฺโถ.\csromanspacer{} ปาณาติปาตา เวรมณิปจฺจยา จ อเนเก กุสลา ธมฺมา ภาวนาปาริปูริํ คจฺฉนฺติ, อยํ อตฺโถ.}

\prose{อทินฺนาทานํ อาวุโส อธมฺโม, อทินฺนาทานา เวรมณี ธมฺโม.\csromanspacer{} เย จ อทินฺนาทานปจฺจยา อเนเก ปาปกา อกุสลา ธมฺมา สมฺภวนฺติ, อยํ อนตฺโถ.\csromanspacer{} อทินฺนาทานา เวรมณิปจฺจยา จ อเนเก กุสลา ธมฺมา ภาวนาปาริปูริํ คจฺฉนฺติ, อยํ อตฺโถ.}

\prose{กาเมสุมิจฺฉาจาโร อาวุโส อธมฺโม, กาเมสุมิจฺฉาจารา เวรมณี ธมฺโม.\csromanspacer{} เย จ กาเมสุมิจฺฉาจารปจฺจยา อเนเก ปาปกา อกุสลา ธมฺมา สมฺภวนฺติ.\csromanspacer{} อยํ อนตฺโถ, กาเมสุมิจฺฉาจารา เวรมณิปจฺจยา จ อเนเก กุสลา ธมฺมา ภาวนาปาริปูริํ คจฺฉนฺติ, อยํ อตฺโถ.}

\prose{มุสาวาโท อาวุโส อธมฺโม, มุสาวาทา เวรมณี ธมฺโม.\csromanspacer{} เย จ มุสาวาทปจฺจยา อเนเก ปาปกา อกุสลา ธมฺมา สมฺภวนฺติ.\csromanspacer{} อยํ อนตฺโถ.\csromanspacer{} มุสาวาทาเวรมณิปจฺจยา จ อเนเก กุสลา ธมฺมา ภาวนาปาริปูริํ คจฺฉนฺติ, อยํ อตฺโถ.}

\prose{ปิสุณา วาจา อาวุโส อธมฺโม, ปิสุณาย วาจาย เวรมณี ธมฺโม.\csromanspacer{} เย จ ปิสุณาวาจาปจฺจยา อเนเก ปาปกา อกุสลา ธมฺมา สมฺภวนฺติ, อยํ อนตฺโถ.\csromanspacer{} ปิสุณาย วาจาย เวรมณิปจฺจยา จ อเนเก กุสลา ธมฺมา ภาวนาปาริปูริํ คจฺฉนฺติ, อยํ อตฺโถ.}

\gambhira{ทสกนิปาตปาฬิ}
\prose{\csromanfolio{468}ผรุสา วาจา อาวุโส อธมฺโม, ผรุสาย วาจาย เวรมณี ธมฺโม.\csromanspacer{} เย จ ผรุสาวาจาปจฺจยา อเนเก ปาปกา อกุสลา ธมฺมา สมฺภวนฺติ, อยํ อนตฺโถ.\csromanspacer{} ผรุสาย วาจาย เวรมณิปจฺจยา จ อเนเก กุสลา ธมฺมา ภาวนาปาริปูริํ คจฺฉนฺติ, อยํ อตฺโถ.}

\prose{สมฺผปฺปลาโป อาวุโส อธมฺโม, สมฺผปฺปลาปา เวรมณี ธมฺโม.\csromanspacer{} เย จ สมฺผปฺปลาปปจฺจยา อเนเก ปาปกา อกุสลา ธมฺมา สมฺภวนฺติ, อยํ อนตฺโถ.\csromanspacer{} สมฺผปฺปลาปา เวรมณิปจฺจยา จ อเนเก กุสลา ธมฺมา ภาวนาปาริปูริํ คจฺฉนฺติ, อยํ อตฺโถ.}

\prose{อภิชฺฌา อาวุโส อธมฺโม, อนภิชฺฌา ธมฺโม.\csromanspacer{} เย จ อภิชฺฌาปจฺจยา อเนเก ปาปกา อกุสลา ธมฺมา สมฺภวนฺติ, อยํ อนตฺโถ.\csromanspacer{} อนภิชฺฌาปจฺจยา จ อเนเก กุสลา ธมฺมา ภาวนาปาริปูริํ คจฺฉนฺติ, อยํ อตฺโถ.}

\prose{พฺยาปาโท อาวุโส อธมฺโม, อพฺยาปาโท ธมฺโม.\csromanspacer{} เย จ พฺยาปาทปจฺจยา อเนเก ปาปกา อกุสลา ธมฺมา สมฺภวนฺติ, อยํ อนตฺโถ, อพฺยาปาทปจฺจยา จ อเนเก กุสลา ธมฺมา ภาวนาปาริปูริํ คจฺฉนฺติ, อยํ อตฺโถ.}

\prose{มิจฺฉาทิฏฺฐิ อาวุโส อธมฺโม.\csromanspacer{} สมฺมาทิฏฺฐิ ธมฺโม.\csromanspacer{} เย จ มิจฺฉาทิฏฺฐิปจฺจยา อเนเก ปาปกา อกุสลา ธมฺมา สมฺภวนฺติ, อยํ อนตฺโถ, สมฺมาทิฏฺฐิปจฺจยา จ อเนเก กุสลา ธมฺมา ภาวนาปาริปูริํ คจฺฉนฺติ, อยํ อตฺโถ.}

\prose{ยํ โข โน อาวุโส ภควา สํขิตฺเตน อุทฺเทสํ อุทฺทิสิตฺวา วิตฺถาเรน อตฺถํ อวิภชิตฺวา อุฏฺฐายาสนา วิหารํ ปวิฏฺโฐ, “อธมฺโม จ ภิกฺขเว เวทิตพฺโพ ฯเปฯ ตถา ปฏิปชฺชิตพฺพนฺ”ติ.\csromanspacer{} อิมสฺส โข อหํ อาวุโส ภควตา สํขิตฺเตน อุทฺเทสสฺส อุทฺทิฏฺฐสฺส วิตฺถาเรน อตฺถํ อวิภตฺตสฺส เอวํ วิตฺถาเรน อตฺถํ อาชานามิ.\csromanspacer{} อากงฺขมานา จ ปน ตุเมฺห อาวุโส ภควนฺตํเยว อุปสงฺกมิตฺวา เอตมตฺถํ ปฏิปุจฺเฉยฺยาถ.\csromanspacer{} ยถา โว ภควา พฺยากโรติ, ตถา นํ ธาเรยฺยาถาติ.}

\prose{“เอวมาวุโส”ติ โข เต ภิกฺขู อายสฺมโต มหากจฺจานสฺส ภาสิตํ อภินนฺทิตฺวา อนุโมทิตฺวา อุฏฺฐายาสนา เยน ภควา}

\vspace{\tipitakaparskip}
\csromancenter{(17) 2.\csromanspacer{} ชาณุสฺโสณิวคฺค เตนุปสงฺกมิํสุ, อุปสงฺกมิตฺวา ภควนฺตํ อภิวาเทตฺวา เอกมนฺตํ นิสีทิํสุ, เอกมนฺตํ นิสินฺนา โข เต ภิกฺขู ภควนฺตํ เอตทโวจุํ\csromandash{}}
\prose{\csromanfolio{469}ยํ โข โน ภนฺเต ภควา สํขิตฺเตน อุทฺเทสํ อุทฺทิสิตฺวา วิตฺถาเรน อตฺถํ อวิภชิตฺวา อุฏฺฐายาสนา วิหารํ ปวิฏฺโฐ, “อธมฺโม จ ภิกฺขเว เวทิตพฺโพ ฯเปฯ ตถา ปฏิปชฺชิตพฺพนฺ”ติ.}

\prose{เตสํ โน ภนฺเต อมฺหากํ อจิรปกฺกนฺตสฺส ภควโต เอตทโหสิ “อิทํ โข โน อาวุโส ภควา สํขิตฺเตน อุทฺเทสํ อุทฺทิสิตฺวา วิตฺถาเรน อตฺถํ อวิภชิตฺวา อุฏฺฐายาสนา วิหารํ ปวิฏฺโฐ ‘อธมฺโม จ ภิกฺขเว เวทิตพฺโพ ฯเปฯ ตถา ปฏิปชฺชิตพฺพนฺ’ติ.\csromanspacer{} โก นุ โข อิมสฺส ภควตา สํขิตฺเตน อุทฺเทสสฺส อุทฺทิฏฺฐสฺส วิตฺถาเรน อตฺถํ อวิภตฺตสฺส วิตฺถาเรน อตฺถํ วิภเชยฺยา”ติ.}

\prose{เตสํ โน ภนฺเต อมฺหากํ เอตทโหสิ “อยํ โข อายสฺมา มหากจฺจาโน สตฺถุ เจว สํวณฺณิโต, สมฺภาวิโต จ วิญฺญูนํ สพฺรหฺมจารีนํ, ปโหติ จายสฺมา มหากจฺจาโน อิมสฺส ภควตา สํขิตฺเตน อุทฺเทสสฺส อุทฺทิฏฺฐสฺส วิตฺถาเรน อตฺถํ อวิภตฺตสฺส วิตฺถาเรน อตฺถํ วิภชิตุํ.\csromanspacer{} ยํนูน มยํ เยนายสฺมา มหากจฺจาโน เตนุปสงฺกเมยฺยาม, อุปสงฺกมิตฺวา อายสฺมนฺตํ มหากจฺจานํ เอตมตฺถํ ปฏิปุจฺเฉยฺยาม.\csromanspacer{} ยถา โน อายสฺมา มหากจฺจาโน พฺยากริสฺสติ, ตถา นํ ธาเรสฺสามา”ติ.}

\prose{อถ โข มยํ ภนฺเต เยนายสฺมา มหากจฺจาโน เตนุปสงฺกมิมฺหา, อุปสงฺกมิตฺวา อายสฺมนฺตํ มหากจฺจานํ เอตมตฺถํ อปุจฺฉิมฺหา.\csromanspacer{} เตสํ โน ภนฺเต อายสฺมตา มหากจฺจาเนน อิเมหิ อกฺขเรหิ อิเมหิ ปเทหิ อิเมหิ พฺยญฺชเนหิ อตฺโถ สุวิภตฺโตติ.}

\prose{สาธุ สาธุ ภิกฺขเว, ปณฺฑิโต ภิกฺขเว มหากจฺจาโน, มหาปญฺโญ ภิกฺขเว มหากจฺจาโน.\csromanspacer{} มํ เจปิ ตุเมฺห ภิกฺขเว อุปสงฺกมิตฺวา เอตมตฺถํ ปฏิปุจฺเฉยฺยาถ, อหมฺปิ เจตํ เอวเมวํ พฺยากเรยฺยํ, ยถา ตํ มหากจฺจาเนน พฺยากตํ, เอโส เจว ตสฺส อตฺโถ, เอวญฺจ นํ ธาเรยฺยาถาติ.\csromanspacer{}\csromanspacer{}ฉฏฺฐํ.}

\gambhira{ทสกนิปาตปาฬิ}
\csromanheader{2}{7. ตติย-อธมฺมสุตฺต}
\proseitem{173}{\csromanfolio{470}อธมฺโม จ ภิกฺขเว เวทิตพฺโพ ธมฺโม จ, อนตฺโถ จ เวทิตพฺโพ อตฺโถ จ.\csromanspacer{} อธมฺมญฺจ วิทิตฺวา ธมฺมญฺจ, อนตฺถญฺจ วิทิตฺวา อตฺถญฺจ ยถา ธมฺโม ยถา อตฺโถ, ตถา ปฏิปชฺชิตพฺพํ.}

\prose{กตโม จ ภิกฺขเว อธมฺโม, กตโม จ ธมฺโม.\csromanspacer{} กตโม จ อนตฺโถ, กตโม จ อตฺโถ? ปาณาติปาโต ภิกฺขเว อธมฺโม, ปาณาติปาตา เวรมณี ธมฺโม.\csromanspacer{} เย จ ปาณาติปาตปจฺจยา อเนเก ปาปกา อกุสลา ธมฺมา สมฺภวนฺติ, อยํ อนตฺโถ.\csromanspacer{} ปาณาติปาตา เวรมณิปจฺจยา จ อเนเก กุสลา ธมฺมา ภาวนาปาริปูริํ คจฺฉนฺติ, อยํ อตฺโถ.}

\prose{อทินฺนาทานํ ภิกฺขเว อธมฺโม, อทินฺนาทานา เวรมณี ธมฺโม, กาเมสุมิจฺฉาจาโร ภิกฺขเว อธมฺโม, กาเมสุมิจฺฉาจารา เวรมณี ธมฺโม.\csromanspacer{} มุสาวาโท ภิกฺขเว อธมฺโม, มุสาวาทา เวรมณี ธมฺโม.\csromanspacer{} ปิสุณา วาจา ภิกฺขเว อธมฺโม, ปิสุณาย วาจาย เวรมณี ธมฺโม.\csromanspacer{} ผรุสา วาจา ภิกฺขเว อธมฺโม, ผรุสาย วาจาย เวรมณี ธมฺโม.\csromanspacer{} สมฺผปฺปลาโป ภิกฺขเว อธมฺโม.\csromanspacer{} สมฺผปฺปลาปา เวรมณี ธมฺโม.\csromanspacer{} อภิชฺฌา ภิกฺขเว อธมฺโม.\csromanspacer{} อนภิชฺฌา ธมฺโม.\csromanspacer{} พฺยาปาโท ภิกฺขเว อธมฺโม, อพฺยาปาโท ธมฺโม.}

\prose{มิจฺฉาทิฏฺฐิ ภิกฺขเว อธมฺโม, สมฺมาทิฏฺฐิ ธมฺโม.\csromanspacer{} เย จ มิจฺฉาทิฏฺฐิปจฺจยา อเนเก ปาปกา อกุสลา ธมฺมา สมฺภวนฺติ, อยํ อนตฺโถ.\csromanspacer{} สมฺมาทิฏฺฐิปจฺจยา จ อเนเก กุสลา ธมฺมา ภาวนาปาริปูริํ คจฺฉนฺติ, อยํ อตฺโถ.}

\prose{“อธมฺโม จ ภิกฺขเว เวทิตพฺโพ ธมฺโม จ, อนตฺโถ จ เวทิตพฺโพ อตฺโถ จ.\csromanspacer{} อธมฺมญฺจ วิทิตฺวา ธมฺมญฺจ, อนตฺถญฺจ วิทิตฺวา อตฺถญฺจ ยถา ธมฺโม ยถา อตฺโถ, ตถา ปฏิปชฺชิตพฺพนฺ”ติ อิติ ยํ ตํ วุตฺตํ, อิทเมตํ ปฏิจฺจ วุตฺตนฺติ.\csromanspacer{}\csromanspacer{}สตฺตมํ.}

\csromanheader{2}{8. กมฺมนิทานสุตฺต}
\proseitem{174}{ปาณาติปาตมฺปาหํ ภิกฺขเว ติวิธํ วทามิ โลภเหตุกมฺปิ โทสเหตุกมฺปิ โมหเหตุกมฺปิ.}

\chapterheadpageread{(17) 2. ชาณุสฺโสณิวคฺค}
\prose{\csromanfolio{471}อทินฺนาทานมฺปาหํ ภิกฺขเว ติวิธํ วทามิ โลภเหตุกมฺปิ โทสเหตุกมฺปิ โมหเหตุกมฺปิ.}

\prose{กาเมสุมิจฺฉาจารมฺปาหํ ภิกฺขเว ติวิธํ วทามิ โลภเหตุกมฺปิ โทสเหตุกมฺปิ โมหเหตุกมฺปิ.}

\prose{มุสาวาทมฺปาหํ ภิกฺขเว ติวิธํ วทามิ โลภเหตุกมฺปิ โทสเหตุกมฺปิ โมหเหตุกมฺปิ.}

\csromanmatikaanchor{mtk.207}
\csromantocmarkref{subsection}{7. ตติย-อธมฺมสุตฺต}{mtk.207}
\bookmark[dest={mtk.207},level=2]{7. ตติย-อธมฺมสุตฺต}
\prose{ปิสุณวาจมฺปาหํ ภิกฺขเว ติวิธํ วทามิ โลภเหตุกมฺปิ โทสเหตุกมฺปิ โมหเหตุกมฺปิ.}

\prose{ผรุสวาจมฺปาหํ ภิกฺขเว ติวิธํ วทามิ โลภเหตุกมฺปิ โทสเหตุกมฺปิ โมหเหตุกมฺปิ.}

\prose{สมฺผปฺปลาปมฺปาหํ ภิกฺขเว ติวิธํ วทามิ โลภเหตุกมฺปิ โทสเหตุกมฺปิ โมหเหตุกมฺปิ.}

\prose{อภิชฺฌมฺปาหํ ภิกฺขเว ติวิธํ วทามิ โลภเหตุกมฺปิ โทสเหตุกมฺปิ โมหเหตุกมฺปิ.}

\prose{พฺยาปาทมฺปาหํ ภิกฺขเว ติวิธํ วทามิ โลภเหตุกมฺปิ โทสเหตุกมฺปิ โมหเหตุกมฺปิ.}

\prose{มิจฺฉาทิฏฺฐิมฺปาหํ ภิกฺขเว ติวิธํ วทามิ โลภเหตุกมฺปิ โทสเหตุกมฺปิ โมหเหตุกมฺปิ.\csromanspacer{} อิติ โข ภิกฺขเว โลโภ กมฺมนิทานสมฺภโว, โทโส กมฺมนิทานสมฺภโว, โมโห กมฺมนิทานสมฺภโว, โลภกฺขยา กมฺมนิทานสงฺขโย, โทสกฺขยา กมฺมนิทานสงฺขโย, โมหกฺขยา กมฺมนิทานสงฺขโยติ.\csromanspacer{}\csromanspacer{}อฏฺฐมํ.}

\csromanmatikaanchor{mtk.208}
\csromantocmarkref{subsection}{8. กมฺมนิทานสุตฺต}{mtk.208}
\bookmark[dest={mtk.208},level=2]{8. กมฺมนิทานสุตฺต}
\csromanheader{2}{9. ปริกฺกมนสุตฺต}
\proseitem{175}{สปริกฺกมโน อยํ ภิกฺขเว ธมฺโม, นายํ ธมฺโม อปริกฺกมโน.\csromanspacer{} กถญฺจ ภิกฺขเว สปริกฺกมโน อยํ ธมฺโม, นายํ ธมฺโม อปริกฺกมโน? ปาณาติปาติสฺส ภิกฺขเว ปาณาติปาตา เวรมณี ปริกฺกมนํ โหติ, อทินฺนาทายิสฺส ภิกฺขเว อทินฺนาทานา เวรมณี ปริกฺกมนํ โหติ, กาเมสุมิจฺฉาจาริสฺส ภิกฺขเว กาเมสุมิจฺฉาจารา เวรมณี ปริกฺกมนํ โหติ, มุสาวาทิสฺส ภิกฺขเว มุสาวาทา เวรมณี ปริกฺกมนํ โหติ,}

\gambhira{ทสกนิปาตปาฬิ}
\prose{\csromanfolio{472}ปิสุณวาจสฺส ภิกฺขเว ปิสุณาย วาจาย เวรมณี ปริกฺกมนํ โหติ, ผรุสวาจสฺส ภิกฺขเว ผรุสาย วาจาย เวรมณี ปริกฺกมนํ โหติ, สมฺผปฺปลาปิสฺส ภิกฺขเว สมฺผปฺปลาปา เวรมณี ปริกฺกมนํ โหติ, อภิชฺฌาลุสฺส ภิกฺขเว อนภิชฺฌา ปริกฺกมนํ โหติ, พฺยาปนฺนจิตฺตสฺส\footnote{พฺยาปาทสฺส (สี, อิ, ก), พฺยาปนฺนสฺส (สฺยา)} ภิกฺขเว อพฺยาปาโท ปริกฺกมนํ โหติ.\csromanspacer{} มิจฺฉาทิฏฺฐิสฺส ภิกฺขเว สมฺมาทิฏฺฐิ ปริกฺกมนํ โหติ.\csromanspacer{} เอวํ โข ภิกฺขเว สปริกฺกมโน อยํ ธมฺโม, นายํ ธมฺโม อปริกฺกมโนติ.\csromanspacer{}\csromanspacer{}นวมํ.}

\csromanheader{2}{10. จุนฺทสุตฺต}
\proseitem{176}{เอวํ เม สุตํ\csromandash{}เอกํ สมยํ ภควา ปาวายํ\footnote{จมฺปายํ (ก-สี) ที 2. 105 ปิฏฺเฐ ปสฺสิตพฺพํ.} วิหรติ จุนฺทสฺส กมฺมารปุตฺตสฺส อมฺพวเน.\csromanspacer{} อถ โข จุนฺโท กมฺมารปุตฺโต เยน ภควา เตนุปสงฺกมิ, อุปสงฺกมิตฺวา ภควนฺตํ อภิวาเทตฺวา เอกมนฺตํ นิสีทิ, เอกมนฺตํ นิสินฺนํ โข จุนฺทํ กมฺมารปุตฺตํ ภควา เอตทโวจ “กสฺส โน ตฺวํ จุนฺท โสเจยฺยานิ โรเจสี”ติ.\csromanspacer{} พฺราหฺมณา ภนฺเต ปจฺฉาภูมกา กมณฺฑลุกา เสวาลมาลิกา\footnote{เสวาลมาลกา (สี, สฺยา, อิ)} อคฺคิปริจาริกา อุทโกโรหกา โสเจยฺยานิ ปญฺญเปนฺติ, เตสาหํ โสเจยฺยานิ โรเจมีติ.}

\prose{ยถา กถํ ปน จุนฺท พฺราหฺมณา ปจฺฉาภูมกา กมณฺฑลุกา เสวาลมาลิกา อคฺคิปริจาริกา อุทโกโรหกา โสเจยฺยานิ ปญฺญเปนฺตีติ.\csromanspacer{} อิธ ภนฺเต พฺราหฺมณา ปจฺฉาภูมกา กมณฺฑลุกา เสวาลมาลิกา อคฺคิปริจาริกา อุทโกโรหกา, เต สาวกํ\footnote{สาวเก (สฺยา, ก)} เอวํ สมาทเปนฺติ. “เอหิ ตฺวํ อมฺโภ ปุริส กาลสฺเสว\footnote{สกาลสฺเสว (สฺยา)} อุฏฺฐหนฺโตว\footnote{อุฏฺฐหนฺโต (สฺยา), วุฏฺฐหนฺโต (อิ, ก)} สยนมฺหา ปถวิํ อามเสยฺยาสิ.\csromanspacer{} โน เจ ปถวิํ อามเสยฺยาสิ, อลฺลานิ โคมยานิ อามเสยฺยาสิ.\csromanspacer{} โน เจ อลฺลานิ โคมยานิ อามเสยฺยาสิ, หริตานิ ติณานิ อามเสยฺยาสิ.\csromanspacer{} โน เจ หริตานิ ติณานิ อามเสยฺยาสิ, อคฺคิํ ปริจเรยฺยาสิ.\csromanspacer{} โน เจ อคฺคิํ ปริจเรยฺยาสิ, ปญฺชลิโก อาทิจฺจํ นมสฺเสยฺยาสิ.\csromanspacer{} โน เจ ปญฺชลิโก อาทิจฺจํ นมสฺเสยฺยาสิ, สายตติยกํ อุทกํ โอโรเหยฺยาสี”ติ.\csromanspacer{} เอวํ โข ภนฺเต พฺราหฺมณา}

\vspace{\tipitakaparskip}
\csromancenter{(17) 2.\csromanspacer{} ชาณุสฺโสณิวคฺค ปจฺฉาภูมกา กมณฺฑลุกา เสวาลมาลิกา อคฺคิปริจาริกา อุทโกโรหกา โสเจยฺยานิ ปญฺญเปนฺติ, เตสาหํ โสเจยฺยานิ โรเจมีติ.}
\prose{\csromanfolio{473}อญฺญถา โข จุนฺท พฺราหฺมณา ปจฺฉาภูมกา กมณฺฑลุกา เสวาลมาลิกา อคฺคิปริจาริกา อุทโกโรหกา โสเจยฺยานิ ปญฺญเปนฺติ, อญฺญถา จ ปน อริยสฺส วินเย โสเจยฺยํ โหตีติ.\csromanspacer{} ยถา กถํ ปน ภนฺเต อริยสฺส วินเย โสเจยฺยํ โหติ, สาธุ เม ภนฺเต ภควา ตถา ธมฺมํ เทเสตุ, ยถา อริยสฺส วินเย โสเจยฺยํ โหตีติ.}

\prose{เตน หิ จุนฺท สุณาหิ สาธุกํ มนสิ กโรหิ, ภาสิสฺสามีติ. “เอวํ ภนฺเต”ติ โข จุนฺโท กมฺมารปุตฺโต ภควโต ปจฺจสฺโสสิ.\csromanspacer{} ภควา เอตทโวจ\csromandash{}}

\prose{ติวิธํ โข จุนฺท กาเยน อโสเจยฺยํ โหติ, จตุพฺพิธํ วาจาย อโสเจยฺยํ โหติ, ติวิธํ มนสา อโสเจยฺยํ โหติ.}

\prose{กถญฺจ จุนฺท ติวิธํ กาเยน อโสเจยฺยํ โหติ? อิธ จุนฺท เอกจฺโจ ปาณาติปาตี โหติ, ลุทฺโท โลหิตปาณิ หตปหเต นิวิฏฺโฐ, อทยาปนฺโน สพฺพปาณภูเตสุ\footnote{ปาณภูเตสุ (ก)}.}

\csromanmatikaanchor{mtk.209}
\csromantocmarkref{subsection}{9. ปริกฺกมนสุตฺต}{mtk.209}
\bookmark[dest={mtk.209},level=2]{9. ปริกฺกมนสุตฺต}
\prose{อทินฺนาทายี โหติ, ยํ ตํ ปรสฺส ปรวิตฺตูปกรณํ คามคตํ วา อรญฺญคตํ วา, ตํ อทินฺนํ เถยฺยสงฺขาตํ อาทาตา โหติ.}

\prose{กาเมสุมิจฺฉาจารี โหติ, ยา ตา มาตุรกฺขิตา ปิตุรกฺขิตา มาตาปิตุรกฺขิตา\footnote{นตฺถิ สี-สฺยา-อิ-โปตฺถเกสุ.} ภาตุรกฺขิตา ภคินิรกฺขิตา ญาติรกฺขิตา โคตฺตรกฺขิตา2 ธมฺมรกฺขิตา สสามิกา สปริทณฺฑา อนฺตมโส มาลาคุฬปริกฺขิตฺตาปิ, ตถารูปาสุ จาริตฺตํ อาปชฺชิตา โหติ.\csromanspacer{} เอวํ โข จุนฺท ติวิธํ กาเยน อโสเจยฺยํ โหติ.}

\prose{กถญฺจ จุนฺท จตุพฺพิธํ วาจาย อโสเจยฺยํ โหติ? อิธ จุนฺท เอกจฺโจ มุสาวาที โหติ, สภคฺคโต วา ปริสคฺคโต วา ญาติมชฺฌคโต วา ปูคมชฺฌคโต วา ราชกุลมชฺฌคโต วา อภินีโต สกฺขิปุฏฺโฐ}

\gambhira{ทสกนิปาตปาฬิ}
\csromanmatikaanchor{mtk.210}
\csromantocmarkref{subsection}{10. จุนฺทสุตฺต}{mtk.210}
\bookmark[dest={mtk.210},level=2]{10. จุนฺทสุตฺต}
\prose{\csromanfolio{474}“เอหมฺโภ ปุริส ยํ ชานาสิ, ตํ วเทหี”ติ, 1โส อชานํ วา อาห “ชานามี”ติ, ชานํ วา อาห “น ชานามี”ติ, อปสฺสํ วา อาห “ปสฺสามี”ติ, ปสฺสํ วา อาห “น ปสฺสามี”ติ1, อิติ อตฺตเหตุ วา ปรเหตุ วา อามิสกิญฺจิกฺขเหตุ วา สมฺปชานมุสา ภาสิตา โหติ.}

\prose{ปิสุณวาโจ โหติ.\csromanspacer{} อิโต สุตฺวา อมุตฺร อกฺขาตา อิเมสํ เภทาย, อมุตฺร วา สุตฺวา อิเมสํ อกฺขาตา อมูสํ เภทาย.\csromanspacer{} อิติ สมคฺคานํ วา เภตฺตา\footnote{เภทาตา (ก)} ภินฺนานํ วา อนุปฺปทาตา, วคฺคาราโม วคฺครโต วคฺคนนฺที วคฺคกรณิํ วาจํ ภาสิตา โหติ.}

\prose{ผรุสวาโจ โหติ, ยา สา วาจา อณฺฑกา กกฺกสา ปรกฏุกา ปราภิสชฺชนี โกธสามนฺตา อสมาธิสํวตฺตนิกา, ตถารูปิํ วาจํ ภาสิตา โหติ.}

\prose{สมฺผปฺปลาปี โหติ, อกาลวาที อภูตวาที อนตฺถวาที อธมฺมวาที อวินยวาที อนิธานวติํ วาจํ ภาสิตา โหติ อกาเลน อนปเทสํ อปริยนฺตวติํ อนตฺถสํหิตํ.\csromanspacer{} เอวํ โข จุนฺท จตุพฺพิธํ วาจาย อโสเจยฺยํ โหติ.}

\prose{กถญฺจ จุนฺท ติวิธํ มนสา อโสเจยฺยํ โหติ? อิธ จุนฺท เอกจฺโจ อภิชฺฌาลุ โหติ.\csromanspacer{} ยํ ตํ ปรสฺส ปรวิตฺตูปกรณํ, ตํ อภิชฺฌาตา\footnote{อภิชฺฌิตา (ก) ม 1. 356 ปิฏฺเฐปิ ปสฺสิตพฺพํ.} โหติ “อโห วต ยํ ปรสฺส, ตํ มมสฺสา”ติ.}

\prose{พฺยาปนฺนจิตฺโต โหติ ปทุฏฺฐมนสงฺกปฺโป “อิเม สตฺตา หญฺญนฺตุ วา พชฺฌนฺตุ วา อุจฺฉิชฺชนฺตุ วา วินสฺสนฺตุ วา มา วา อเหสุนฺ”ติ\footnote{มาวา อเหสุํ อิติ วาติ (สี, อิ, ก)}.}

\prose{มิจฺฉาทิฏฺฐิโก โหติ วิปรีตทสฺสโน “นตฺถิ ทินฺนํ, นตฺถิ ยิฏฺฐํ, นตฺถิ หุตํ, นตฺถิ สุกฏทุกฺกฏานํ\footnote{นตฺเถตฺถ ปาฐเภโท.} กมฺมานํ ผลํ วิปาโก, นตฺถิ อยํ โลโก, นตฺถิ ปโร โลโก, นตฺถิ มาตา, นตฺถิ ปิตา, นตฺถิ สตฺตา โอปปาติกา, นตฺถิ โลเก สมณพฺราหฺมณา สมฺมคฺคตา สมฺมาปฏิปนฺนา, เย อิมญฺจ โลกํ ปรญฺจ โลกํ สยํ อภิญฺญา สจฺฉิกตฺวา ปเวเทนฺตี”ติ.\csromanspacer{} เอวํ โข จุนฺท มนสา ติวิธํ อโสเจยฺยํ โหติ.}

\chapterheadpageread{(17) 2. ชาณุสฺโสณิวคฺค}
\prose{\csromanfolio{475}อิเม โข จุนฺท ทส อกุสลกมฺมปถา\footnote{อกุสลา กมฺมปถา (?)}.\csromanspacer{} อิเมหิ โข จุนฺท ทสหิ อกุสเลหิ กมฺมปเถหิ สมนฺนาคโต กาลสฺเสว อุฏฺฐหนฺโตว สยนมฺหา ปถวิํ เจปิ อามสติ, อสุจิเยว โหติ, โน เจปิ ปถวิํ อามสติ, อสุจิเยว โหติ.}

\prose{อลฺลานิ เจปิ โคมยานิ อามสติ, อสุจิเยว โหติ.\csromanspacer{} โน เจปิ อลฺลานิ โคมยานิ อามสติ, อสุจิเยว โหติ.}

\prose{หริตานิ เจปิ ติณานิ อามสติ, อสุจิเยว โหติ.\csromanspacer{} โน เจปิ หริตานิ ติณานิ อามสติ, อสุจิเยว โหติ.}

\prose{อคฺคิํ เจวิ ปริจรติ, อสุจิเยว โหติ.\csromanspacer{} โน เจปิ อคฺคิํ ปริจรติ, อสุจิเยว โหติ.}

\prose{ปญฺชลิโก เจปิ อาทิจฺจํ นมสฺสติ, อสุจิเยว โหติ.\csromanspacer{} โน เจปิ ปญฺชลิโก อาทิจฺจํ นมสฺสติ, อสุจิเยว โหติ.}

\prose{สายตติยกํ เจปิ อุทกํ โอโรหติ, อสุจิเยว โหติ.\csromanspacer{} โน เจปิ สายตติยกํ อุทกํ โอโรหติ, อสุจิเยว โหติ.\csromanspacer{} ตํ กิสฺส เหตุ? อิเม จุนฺท ทส อกุสลกมฺมปถา อสุจีเยว\footnote{อสุจิจฺเจว (สฺยา)} โหนฺติ อสุจิกรณา จ.}

\prose{อิเมสํ ปน จุนฺท ทสนฺนํ อกุสลานํ กมฺมปถานํ สมนฺนาคมนเหตุ นิรโย ปญฺญายติ, ติรจฺฉานโยนิ ปญฺญายติ, เปตฺติวิสโย ปญฺญายติ, ยา วา\footnote{ยา จ (ก)} ปนญฺญาปิ กาจิ ทุคฺคติโย\footnote{ทุคฺคติ โหติ (สฺยา, ก)}.}

\prose{ติวิธํ โข จุนฺท กาเยน โสเจยฺยํ โหติ, จตุพฺพิธํ วาจาย โสเจยฺยํ โหติ, ติวิธํ มนสา โสเจยฺยํ โหติ.}

\prose{กถํ จุนฺท ติวิธํ กาเยน โสเจยฺยํ โหติ? อิธ จุนฺท เอกจฺโจ ปาณาติปาตํ ปหาย ปาณาติปาตา ปฏิวิรโต โหติ, นิหิตทณฺโฑ นิหิตสตฺโถ ลชฺชี ทยาปนฺโน สพฺพปาณภูตหิตานุกมฺปี วิหรติ.}

\gambhira{ทสกนิปาตปาฬิ}
\prose{\csromanfolio{476}อทินฺนาทานํ ปหาย อทินฺนาทานา ปฏิวิรโต โหติ, ยํ ตํ ปรสฺส ปรวิตฺตูปกรณํ คามคตํ วา อรญฺญคตํ วา, น ตํ อทินฺนํ\footnote{ตํ นาทินฺนํ (ก-สี, ม 1. 357 ปิฏฺเฐปิ)} เถยฺยสงฺขาตํ อาทาตา โหติ.}

\prose{กาเมสุมิจฺฉาจารํ ปหาย กาเมสุมิจฺฉาจารา ปฏิวิรโต โหติ, ยา ตา มาตุรกฺขิตา ปิตุรกฺขิตา มาตาปิตุรกฺขิตา ภาตุรกฺขิตา ภคินิรกฺขิตา ญาติรกฺขิตา โคตฺตรกฺขิตา ธมฺมรกฺขิตา สสามิกา สปริทณฺฑา อนฺตมโส มาลาคุฬปริกฺขิตฺตาปิ, ตถารูปาสุ น จาริตฺตํ อาปชฺชิตา โหติ.\csromanspacer{} เอวํ โข จุนฺท ติวิธํ กาเยน โสเจยฺยํ โหติ.}

\prose{กถญฺจ จุนฺท จตุพฺพิธํ วาจาย โสเจยฺยํ โหติ? อิธ จุนฺท เอกจฺโจ มุสาวาทํ ปหาย มุสาวาทา ปฏิวิรโต โหติ, สภคฺคโต วา ปริสคฺคโต วา ญาติมชฺฌคโต วา ปูคมชฺฌคโต วา ราชกุลมชฺฌคโต วา อภินีโต สกฺขิปุฏฺโฐ “เอหมฺโภ ปุริส ยํ ชานาสิ, ตํ วเทหี”ติ, โส อชานํ วา อาห “น ชานามี”ติ, ชานํ วา อาห “ชานามี”ติ, อปสฺสํ วา อาห “น ปสฺสามี”ติ, ปสฺสํ วา อาห “ปสฺสามี”ติ, อิติ อตฺตเหตุ วา ปรเหตุ วา อามิสกิญฺจิกฺขเหตุ วา น สมฺปชานมุสา ภาสิตา โหติ.}

\prose{ปิสุณํ วาจํ ปหาย ปิสุณาย วาจาย ปฏิวิรโต โหติ, น อิโต สุตฺวา อมุตฺร อกฺขาตา อิเมสํ เภทาย, น อมุตฺร วา สุตฺวา อิเมสํ อกฺขาตา อมูสํ เภทาย.\csromanspacer{} อิติ ภินฺนานํ วา สนฺธาตา สหิตานํ วา อนุปฺปทาตา, สมคฺคาราโม สมคฺครโต สมคฺคนนฺที สมคฺคกรณิํ วาจํ ภาสิตา โหติ.}

\prose{ผรุสํ วาจํ ปหาย ผรุสาย วาจาย ปฏิวิรโต โหติ, ยา สา วาจา เนลา กณฺณสุขา เปมนียา หทยงฺคมา โปรี พหุชนกนฺตา พหุชนมนาปา, ตถารูปิํ วาจํ ภาสิตา โหติ.}

\prose{สมฺผปฺปลาปํ ปหาย สมฺผปฺปลาปา ปฏิวิรโต โหติ, กาลวาที ภูตวาที อตฺถวาที ธมฺมวาที วินยวาที นิธานวติํ วาจํ ภาสิตา โหติ กาเลน สาปเทสํ ปริยนฺตวติํ อตฺถสํหิตํ.\csromanspacer{} เอวํ โข จุนฺท จตุพฺพิธํ วาจาย โสเจยฺยํ โหติ.}

\chapterheadpageread{(17) 2. ชาณุสฺโสณิวคฺค}
\prose{\csromanfolio{477}กถญฺจ จุนฺท ติวิธํ มนสา โสเจยฺยํ โหติ? อิธ จุนฺท เอกจฺโจ อนภิชฺฌาลุ โหติ, ยํ ตํ ปรสฺส ปรวิตฺตูปกรณํ, ตํ อนภิชฺฌิตา โหติ “อโห วต ยํ ปรสฺส ตํ มมสฺสา”ติ.}

\prose{อพฺยาปนฺนจิตฺโต โหติ อปฺปทุฏฺฐมนสงฺกปฺโป “อิเม สตฺตา อเวรา โหนฺตุ\footnote{อิทํ ปทํ สี-สฺยา-อิ-โปตฺถเกสุ นตฺถิ, ตถา ม 1. 358 ปิฏฺฐาทีสุปิ.} อพฺยาปชฺชา อนีฆา, สุขี อตฺตานํ ปริหรนฺตู”ติ.}

\prose{สมฺมาทิฏฺฐิโก โหติ อวิปรีตทสฺสโน, “อตฺถิ ทินฺนํ, อตฺถิ ยิฏฺฐํ อตฺถิ หุตํ, อตฺถิ สุกฏทุกฺกฏานํ กมฺมานํ ผลํ วิปาโก, อตฺถิ อยํ โลโก, อตฺถิ ปโร โลโก, อตฺถิ มาตา, อตฺถิ ปิตา, อตฺถิ สตฺตา โอปปาติกา, อตฺถิ โลเก สมณพฺราหฺมณา สมฺมคฺคตา สมฺมาปฏิปนฺนา, เย อิมญฺจ โลกํ ปรญฺจ โลกํ สยํ อภิญฺญา สจฺฉิกตฺวา ปเวเทนฺตี”ติ.\csromanspacer{} เอวํ โข จุนฺท ติวิธํ มนสา โสเจยฺยํ โหติ.}

\prose{อิเม โข จุนฺท ทส กุสลกมฺมปถา.\csromanspacer{} อิเมหิ โข จุนฺท ทสหิ กุสเลหิ กมฺมปเถหิ สมนฺนาคโต กาลสฺเสว อุฏฺฐหนฺโตว สยนมฺหา ปถวิํ เจปิ อามสติ, สุจิเยว โหติ.\csromanspacer{} โน เจปิ ปถวิํ อามสติ, สุจิเยว โหติ.}

\prose{อลฺลานิ เจปิ โคมยานิ อามสติ, สุจิเยว โหติ.\csromanspacer{} โน เจปิ อลฺลานิ โคมยานิ อามสติ, สุจิเยว โหติ.}

\prose{หริตานิ เจปิ ติณานิ อามสติ, สุจิเยว โหติ.\csromanspacer{} โน เจปิ หริตานิ ติณานิ อามสติ, สุจิเยว โหติ.}

\prose{อคฺคิํ เจปิ ปริจรติ, สุจิเยว โหติ.\csromanspacer{} โน เจปิ อคฺคิํ ปริจรติ, สุจิเยว โหติ.}

\prose{ปญฺจลิโก เจปิ อาทิจฺจํ นมสฺสติ, สุจิเยว โหติ.\csromanspacer{} โน เจปิ ปญฺชลิโก อาทิจฺจํ นมสฺสติ, สุจิเยว โหติ.}

\prose{สายตติยกํ เจปิ อุทกํ โอโรหติ, สุจิเยว โหติ.\csromanspacer{} โน เจปิ สายตติยกํ อุทกํ โอโรหติ, สุจิเยว โหติ.\csromanspacer{} ตํ กิสฺส เหตุ? อิเม จุนฺท ทส กุสลกมฺมปถา สุจีเยว โหนฺติ สุจิกรณา จ.}

\gambhira{ทสกนิปาตปาฬิ}
\prose{\csromanfolio{478}อิเมสํ ปน จุนฺท ทสนฺนํ กุสลานํ กมฺมปถานํ สมนฺนาคมนเหตุ เทวา ปญฺญายนฺติ, มนุสฺสา ปญฺญายนฺติ, ยา วา ปนญฺญาปิ กาจิ สุคติโยติ\footnote{สุคติ โหตีติ (สฺยา), สุคติ โหติ (ก)}.}

\prose{เอวํ วุตฺเต จุนฺโท กมฺมารปุตฺโต ภควนฺตํ เอตทโวจ “อภิกฺกนฺตํ ภนฺเต ฯเปฯ อุปาสกํ มํ ภนฺเต ภควา ธาเรตุ อชฺชตคฺเค ปาณุเปตํ สรณํ คตนฺ”ติ.\csromanspacer{}\csromanspacer{}ทสมํ.}

\csromanheader{2}{11. ชาณุสฺโสณิสุตฺต}
\proseitem{177}{อถ โข ชาณุสฺโสณิ พฺราหฺมโณ เยน ภควา เตนุปสงฺกมิ, อุปสงฺกมิตฺวา ภควตา สทฺธิํ สมฺโมทิ, สมฺโมทนียํ กถํ สารณียํ วีติสาเรตฺวา เอกมนฺตํ นิสีทิ, เอกมนฺตํ นิสินฺโน โข ชาณุสฺโสณิ พฺราหฺมโณ ภควนฺตํ เอตทโวจ\csromandash{}}

\prose{มยมสฺสุ โภ โคตม พฺราหฺมณา นาม ทานานิ เทม, สทฺธานิ กโรม “อิทํ ทานํ เปตานํ ญาติสาโลหิตานํ อุปกปฺปตุ, อิทํ ทานํ เปตา ญาติสาโลหิตา ปริภุญฺชนฺตู”ติ.\csromanspacer{} กจฺจิ ตํ โภ โคตม ทานํ เปตานํ ญาติสาโลหิตานํ อุปกปฺปติ, กจฺจิ เต เปตา ญาติสาโลหิตา ตํ ทานํ ปริภุญฺชนฺตีติ.\csromanspacer{} ฐาเน โข พฺราหฺมณ อุปกปฺปติ, โน อฏฺฐาเนติ.}

\prose{กตมํ ปน โภ\footnote{กตมญฺจ ปน โภ (สี, อิ) กตมํ (สฺยา)} โคตม ฐานํ, กตมํ อฏฺฐานนฺติ.\csromanspacer{} อิธ พฺราหฺมณ เอกจฺโจ ปาณาติปาตี โหติ, อทินฺนาทายี โหติ, กาเมสุมิจฺฉาจารี โหติ, มุสาวาที โหติ, ปิสุณวาโจ โหติ, ผรุสวาโจ โหติ, สมฺผปฺปลาปี โหติ, อภิชฺฌาลุ โหติ, พฺยาปนฺนจิตฺโต โหติ, มิจฺฉาทิฏฺฐิโก โหติ.\csromanspacer{} โส กายสฺส เภทา ปรํ มรณา นิรยํ อุปปชฺชติ.\csromanspacer{} โย เนรยิกานํ สตฺตานํ อาหาโร, เตน โส ตตฺถ ยาเปติ, เตน โส ตตฺถ ติฏฺฐติ.\csromanspacer{} อิทมฺปิ\footnote{อิทํ (สฺยา)} โข พฺราหฺมณ อฏฺฐานํ, ยตฺถ ฐิตสฺส ตํ ทานํ น อุปกปฺปติ.}

\prose{อิธ ปน พฺราหฺมณ เอกจฺโจ ปาณาติปาตี ฯเปฯ มิจฺฉาทิฏฺฐิโก โหติ.\csromanspacer{} โส กายสฺส เภทา ปรํ มรณา ติรจฺฉานโยนิํ อุปปชฺชติ.\csromanspacer{} โย ติรจฺฉานโยนิกานํ สตฺตานํ อาหาโร, เตน โส ตตฺถ ยาเปติ,}

\vspace{\tipitakaparskip}
\csromancenter{(17) 2.\csromanspacer{} ชาณุสฺโสณิวคฺค เตน โส ตตฺถ ติฏฺฐติ.\csromanspacer{} อิทมฺปิ โข พฺราหฺมณ อฏฺฐานํ, ยตฺถ ฐิตสฺส ตํ ทานํ น อุปกปฺปติ.}
\prose{\csromanfolio{479}อิธ ปน พฺราหฺมณ เอกจฺโจ ปาณาติปาตา ปฏิวิรโต โหติ, อทินฺนาทานา ปฏิวิรโต โหติ, กาเมสุมิจฺฉาจารา ปฏิวิรโต โหติ, มุสาวาทา ปฏิวิรโต โหติ, ปิสุณาย วาจาย ปฏิวิรโต โหติ, ผรุสาย วาจาย ปฏิวิรโต โหติ, สมฺผปฺปลาปา ปฏิวิรโต โหติ, อนภิชฺฌาลุ โหติ, อพฺยาปนฺนจิตฺโต โหติ, สมฺมาทิฏฺฐิโก โหติ.\csromanspacer{} โส กายสฺส เภทา ปรํ มรณา มนุสฺสานํ สหพฺยตํ อุปปชฺชติ.\csromanspacer{} โย มนุสฺสานํ อาหาโร, เตน โส ตตฺถ ยาเปติ, เตน โส ตตฺถ ติฏฺฐติ.\csromanspacer{} อิทมฺปิ โข พฺราหฺมณ อฏฺฐานํ, ยตฺถ ฐิตสฺส ตํ ทานํ น อุปกปฺปติ.}

\prose{อิธ ปน พฺราหฺมณ เอกจฺโจ ปาณาติปาตา ปฏิวิรโต โหติ ฯเปฯ สมฺมาทิฏฺฐิโก โหติ.\csromanspacer{} โส กายสฺส เภทา ปรํ มรณา เทวานํ สหพฺยตํ อุปปชฺชติ.\csromanspacer{} โย เทวานํ อาหาโร, เตน โส ตตฺถ ยาเปติ, เตน โส ตตฺถ ติฏฺฐติ.\csromanspacer{} อิทมฺปิ พฺราหฺมณ อฏฺฐานํ, ยตฺถ ฐิตสฺส ตํ ทานํ อุปกปฺปติ.}

\prose{อิธ ปน พฺราหฺมณ เอกจฺโจ ปาณาติปาตี โหติ ฯเปฯ มิจฺฉาทิฏฺฐิโก โหติ.\csromanspacer{} โส กายสฺส เภทา ปรํ มรณา เปตฺติวิสยํ อุปปชฺชติ.\csromanspacer{} โย เปตฺติเวสยิกานํ สตฺตานํ อาหาโร, เตน โส ตตฺถ ยาเปติ, เตน โส ตตฺถ ติฏฺฐติ.\csromanspacer{} ยํ วา ปนสฺส อิโต อนุปฺปเวจฺฉนฺติ มิตฺตามจฺจา วา ญาติสาโลหิตา วา\footnote{มิตฺตา วา อมจฺจา วา ญาตี วา สาโลหิตา วา (สี, อิ)}, เตน โส ตตฺถ ยาเปติ, เตน โส ตตฺถ ติฏฺฐติ.\csromanspacer{} อิทํ โข พฺราหฺมณ ฐานํ, ยตฺถ ฐิตสฺส ตํ ทานํ อุปกปฺปตีติ.}

\prose{สเจ ปน โภ โคตม โส เปโต ญาติสาโลหิโต ตํ ฐานํ อนุปปนฺโน โหติ, โก ตํ ทานํ ปริภุญฺชตีติ.\csromanspacer{} อญฺเญปิสฺส พฺราหฺมณ เปตา ญาติสาโลหิตา ตํ ฐานํ อุปปนฺนา โหนฺติ, เต ตํ ทานํ ปริภุญฺชนฺตีติ.}

\prose{สเจ ปน โภ โคตม โส เจว เปโต ญาติสาโลหิโต ตํ ฐานํ อนุปปนฺโน โหติ, อญฺเญปิสฺส ญาติสาโลหิตา เปตา}

\gambhira{ทสกนิปาตปาฬิ}
\csromanmatikaanchor{mtk.211}
\csromantocmarkref{subsection}{11. ชาณุสฺโสณิสุตฺต}{mtk.211}
\bookmark[dest={mtk.211},level=2]{11. ชาณุสฺโสณิสุตฺต}
\prose{\csromanfolio{480}ตํ ฐานํ อนุปปนฺนา โหนฺติ, โก ตํ ทานํ ปริภุญฺชตีติ.\csromanspacer{} อฏฺฐานํ โข เอตํ พฺราหฺมณ อนวกาโส, ยํ ตํ ฐานํ วิวิตฺตํ อสฺส อิมินา ทีเฆน อทฺธุนา ยทิทํ เปเตหิ ญาติสาโลหิเตหิ.\csromanspacer{} อปิ จ พฺราหฺมณ ทายโกปิ อนิปฺผโลติ.}

\prose{อฏฺฐาเนปิ ภวํ โคตโม ปริกปฺปํ วทตีติ.\csromanspacer{} อฏฺฐาเนปิ โข อหํ พฺราหฺมณ ปริกปฺปํ วทามิ.\csromanspacer{} อิธ พฺราหฺมณ เอกจฺโจ ปาณาติปาตี โหติ, อทินฺนาทายี โหติ, กาเมสุมิจฺฉาจารี โหติ, มุสาวาที โหติ, ปิสุณวาโจ โหติ, ผรุสวาโจ โหติ, สมฺผปฺปลาปี โหติ, อภิชฺฌาลุ โหติ, พฺยาปนฺนจิตฺโต โหติ, มิจฺฉาทิฏฺฐิโก โหติ, โส ทาตา โหติ สมณสฺส วา พฺราหฺมณสฺส วา อนฺนํ ปานํ วตฺถํ ยานํ มาลาคนฺธวิเลปนํ เสยฺยาวสถปทีเปยฺยํ, โส กายสฺส เภทา ปรํ มรณา หตฺถีนํ สหพฺยตํ อุปปชฺชติ, โส ตตฺถ ลาภี โหติ อนฺนสฺส ปานสฺส มาลานานาลงฺการสฺส\footnote{มาลาคนฺธวิเลปนสฺส นานาลงฺการสฺส (ก)}.}

\prose{ยํ โข พฺราหฺมณ อิธ ปาณาติปาตี อทินฺนาทายี กาเมสุมิจฺฉาจารี มุสาวาที ปิสุณวาโจ ผรุสวาโจ สมฺผปฺปลาปี อภิชฺฌาลุ พฺยาปนฺนจิตฺโต มิจฺฉาทิฏฺฐิโก, เตน โส กายสฺส เภทา ปรํ มรณา หตฺถีนํ สหพฺยตํ อุปปชฺชติ.\csromanspacer{} ยญฺจ โข โส ทาตา โหติ สมณสฺส วา พฺราหฺมณสฺส วา อนฺนํ ปานํ วตฺถํ ยานํ มาลาคนฺธวิเลปนํ เสยฺยาวสถปทีเปยฺยํ, เตน โส ตตฺถ ลาภี โหติ อนฺนสฺส ปานสฺส มาลานานาลงฺการสฺส.}

\prose{อิธ ปน พฺราหฺมณ เอกจฺโจ ปาณาติปาตี โหติ ฯเปฯ มิจฺฉาทิฏฺฐิโก โหติ.\csromanspacer{} โส ทาตา โหติ สมณสฺส วา พฺราหฺมณสฺส วา อนฺนํ ปานํ วตฺถํ ยานํ มาลาคนฺธวิเลปนํ เสยฺยาวสถปทีเปยฺยํ, โส กายสฺส เภทา ปรํ มรณา อสฺสานํ สหพฺยตํ อุปปชฺชติ ฯเปฯ คุนฺนํ สหพฺยตํ อุปปชฺชติ ฯเปฯ กุกฺกุรานํ สหพฺยตํ อุปปชฺชติ\footnote{“กุกฺกุรานํ สหพฺยตํ อุปปชฺชตี”ติ อยํ วาโร เกสุจิ สีหฬโปตฺถเกสุ น ทิสฺสตีติ อิงฺคลิสโปตฺถเก อโธลิปิ. ตํ ทสวารคณนาย สเมติ.}.\csromanspacer{} โส ตตฺถ ลาภี โหติ อนฺนสฺส ปานสฺส มาลานานาลงฺการสฺส.}

\chapterheadpageread{(17) 2. ชาณุสฺโสณิวคฺค}
\prose{\csromanfolio{481}ยํ โข พฺราหฺมณ อิธ ปาณาติปาตี ฯเปฯ มิจฺฉาทิฏฺฐิโก.\csromanspacer{} เตน โส กายสฺส เภทา ปรํ มรณา กุกฺกุรานํ สหพฺยตํ อุปปชฺชติ.\csromanspacer{} ยญฺจ โข โส ทาตา โหติ สมณสฺส วา พฺราหฺมณสฺส วา อนฺนํ ปานํ วตฺถํ ยานํ มาลาคนฺธวิเลปนํ เสยฺยาวสถปทีเปยฺยํ, เตน โส ตตฺถ ลาภี โหติ อนฺนสฺส ปานสฺส มาลานานาลงฺการสฺส.}

\prose{อิธ ปน พฺราหฺมณ เอกจฺโจ ปาณาติปาตา ปฏิวิรโต โหติ ฯเปฯ สมฺมาทิฏฺฐิโก โหติ.\csromanspacer{} โส ทาตา โหติ สมณสฺส วา พฺราหฺมณสฺส วา อนฺนํ ปานํ วตฺถํ ยานํ มาลาคนฺธวิเลปนํ เสยฺยาวสถปทีเปยฺยํ, โส กายสฺส เภทา ปรํ มรณา มนุสฺสานํ สหพฺยตํ อุปปชฺชติ.\csromanspacer{} โส ตตฺถ ลาภี โหติ มานุสกานํ ปญฺจนฺนํ กามคุณานํ.}

\prose{ยํ โข พฺราหฺมณ อิธ ปาณาติปาตา ปฏิวิรโต โหติ ฯเปฯ สมฺมาทิฏฺฐิโก.\csromanspacer{} เตน โส กายสฺส เภทา ปรํ มรณา มนุสฺสานํ สหพฺยตํ อุปปชฺชติ.\csromanspacer{} ยญฺจ โข โส ทาตา โหติ สมณสฺส วา พฺราหฺมณสฺส วา อนฺนํ ปานํ วตฺถํ ยานํ มาลาคนฺธวิเลปนํ เสยฺยาวสถปทีเปยฺยํ, เตน โส ตตฺถ ลาภี โหติ มานุสกานํ ปญฺจนฺนํ กามคุณานํ.}

\prose{อิธ ปน พฺราหฺมณ เอกจฺโจ ปาณาติปาตา ปฏิวิรโต โหติ ฯเปฯ สมฺมาทิฏฺฐิโก โหติ.\csromanspacer{} โส ทาตา โหติ สมณสฺส วา พฺราหฺมณสฺส วา อนฺนํ ปานํ วตฺถํ ยานํ มาลาคนฺธวิเลปนํ เสยฺยาวสถปทีเปยฺยํ, โส กายสฺส เภทา ปรํ มรณา เทวานํ สหพฺยตํ อุปปชฺชติ.\csromanspacer{} โส ตตฺถ ลาภี โหติ ทิพฺพานํ ปญฺจนฺนํ กามคุณานํ.}

\prose{ยํ โข พฺราหฺมณ อิธ ปาณาติปาตา ปฏิวิรโต โหติ ฯเปฯ สมฺมาทิฏฺฐิโก.\csromanspacer{} เตน โส กายสฺส เภทา ปรํ มรณา เทวานํ สหพฺยตํ อุปปชฺชติ.\csromanspacer{} ยญฺจ โข โส ทาตา โหติ สมณสฺส วา พฺราหฺมณสฺส วา อนฺนํ ปานํ วตฺถํ ยานํ มาลาคนฺธวิเลปนํ เสยฺยาวสถปทีเปยฺยํ, เตน โส ตตฺถ ลาภี โหติ ทิพฺพานํ ปญฺจนฺนํ กามคุณานํ, อปิ จ พฺราหฺมณ ทายโกปิ อนิปฺผโลติ.}

\prose{อจฺฉริยํ โภ โคตม, อพฺภุตํ โภ โคตม, ยาวญฺจิทํ โภ โคตม อลเมว ทานานิ ทาตุํ, อลํ สทฺธานิ กาตุํ.\csromanspacer{} ยตฺร หิ นาม}

\gambhira{ทสกนิปาตปาฬิ}
\prose{\csromanfolio{482}ทายโกปิ อนิปฺผโลติ.\csromanspacer{} เอวเมตํ พฺราหฺมณ\footnote{เอวเมตํ พฺราหฺมณ เอวเมตํ พฺราหฺมณ (สี, สฺยา)}, ทายโกปิ หิ พฺราหฺมณ อนิปฺผโลติ.}

\prose{อภิกฺกนฺตํ โภ โคตม, อภิกฺกนฺตํ โภ โคตม ฯเปฯ อุปาสกํ มํ ภวํ โคตโม ธาเรตุ อชฺชตคฺเค ปาณุเปตํ สรณํ คตนฺติ.\csromanspacer{}\csromanspacer{}เอกาทสมํ.}

\csromanheader{2}{ชาณุสฺโสณิวคฺโค\footnote{ยมกวคฺโค (ก)} ทุติโย.}
\csromansectionrule
\csromanheader{1}{\textbf{(18) 3. สาธุวคฺค}}
\csromanheader{2}{1. สาธุสุตฺต}
\proseitem{178}{\csromansymbolfootnote{*}{เหฏฺฐา 453 ปิฏฺเฐปิ.}สาธุญฺจ โว ภิกฺขเว เทเสสฺสามิ อสาธุญฺจ, ตํ สุณาถ สาธุกํ มนสิ กโรถ, ภาสิสฺสามี”ติ. “เอวํ ภนฺเต”ติ โข เต ภิกฺขู ภควโต ปจฺจสฺโสสุํ.\csromanspacer{} ภควา เอตทโวจ\csromandash{}}

\prose{กตมญฺจ ภิกฺขเว อสาธุ? ปาณาติปาโต อทินฺนาทานํ กาเมสุมิจฺฉาจาโร มุสาวาโท ปิสุณา วาจา ผรุสา วาจา สมฺผปฺปลาโป อภิชฺฌา พฺยาปาโท มิจฺฉาทิฏฺฐิ.\csromanspacer{} อิทํ วุจฺจติ ภิกฺขเว อสาธุ.}

\prose{กตมญฺจ ภิกฺขเว สาธุ? ปาณาติปาตา เวรมณี อทินฺนาทานา เวรมณี กาเมสุมิจฺฉาจารา เวรมณี มุสาวาทา เวรมณี ปิสุณาย วาจาย เวรมณี ผรุสาย วาจาย เวรมณี สมฺผปฺปลาปา เวรมณี อนภิชฺฌา อพฺยาปาโท สมฺมาทิฏฺฐิ.\csromanspacer{} อิทํ วุจฺจติ ภิกฺขเว สาธูติ.\csromanspacer{}\csromanspacer{}ปฐมํ.}

\csromanheader{2}{2. อริยธมฺมสุตฺต}
\proseitem{179}{อริยธมฺมญฺจ โว ภิกฺขเว เทเสสฺสามิ อนริยธมฺมญฺจ, ตํ สุณาถ ฯเปฯ.\csromanspacer{} กตโม จ ภิกฺขเว อนริโย ธมฺโม? ปาณาติปาโต ฯเปฯ มิจฺฉาทิฏฺฐิ.\csromanspacer{} อยํ วุจฺจติ ภิกฺขเว อนริโย ธมฺโม.}

\prose{กตโม จ ภิกฺขเว อริโย ธมฺโม? ปาณาติปาตา เวรมณี ฯเปฯ สมฺมาทิฏฺฐิ.\csromanspacer{} อยํ วุจฺจติ ภิกฺขเว อริโย ธมฺโมติ.\csromanspacer{}\csromanspacer{}ทุติยํ.}

\csromanheader{2}{(18) 3. สาธุวคฺค \textbf{3. กุสลสุตฺต}}
\proseitem{180}{\csromanfolio{483}กุสลญฺจ โว ภิกฺขเว เทเสสฺสามิ อกุสลญฺจ, ตํ สุณาถ ฯเปฯ.\csromanspacer{} กตมญฺจ ภิกฺขเว อกุสลํ? ปาณาติปาโต ฯเปฯ มิจฺฉาทิฏฺฐิ.\csromanspacer{} อิทํ วุจฺจติ ภิกฺขเว อกุสลํ.}

\prose{กตมญฺจ ภิกฺขเว กุสลํ? ปาณาติปาตา เวรมณี ฯเปฯ สมฺมาทิฏฺฐิ.\csromanspacer{} อิทํ วุจฺจติ ภิกฺขเว กุสลนฺติ.\csromanspacer{}\csromanspacer{}ตติยํ.}

\csromanheader{2}{4. อตฺถสุตฺต}
\csromanmatikaanchor{mtk.212}
\csromantocmarknopage{section}{(18) 3. สาธุวคฺค}{mtk.212}
\bookmark[dest={mtk.212},level=1]{(18) 3. สาธุวคฺค}
\proseitem{181}{อตฺถญฺจ โว ภิกฺขเว เทเสสฺสามิ อนตฺถญฺจ, ตํ สุณาถ ฯเปฯ.\csromanspacer{} กตโม จ ภิกฺขเว อนตฺโถ? ปาณาติปาโต ฯเปฯ มิจฺฉาทิฏฺฐิ.\csromanspacer{} อยํ วุจฺจติ ภิกฺขเว อนตฺโถ.}

\csromanmatikaanchor{mtk.213}
\csromantocmarkref{subsection}{1. สาธุสุตฺต}{mtk.213}
\bookmark[dest={mtk.213},level=2]{1. สาธุสุตฺต}
\prose{กตโม จ ภิกฺขเว อตฺโถ? ปาณาติปาตา เวรมณี ฯเปฯ สมฺมาทิฏฺฐิ.\csromanspacer{} อยํ วุจฺจติ ภิกฺขเว อตฺโถติ.\csromanspacer{}\csromanspacer{}จตุตฺถํ.}

\csromanheader{2}{5. ธมฺมสุตฺต}
\proseitem{182}{ธมฺมญฺจ โว ภิกฺขเว เทเสสฺสามิ อธมฺมญฺจ, ตํ สุณาถ ฯเปฯ.\csromanspacer{} กตโม จ ภิกฺขเว อธมฺโม? ปาณาติปาโต ฯเปฯ มิจฺฉาทิฏฺฐิ.\csromanspacer{} อยํ วุจฺจติ ภิกฺขเว อธมฺโม.}

\prose{กตโม จ ภิกฺขเว ธมฺโม? ปาณาติปาตา เวรมณี ฯเปฯ สมฺมาทิฏฺฐิ.\csromanspacer{} อยํ วุจฺจติ ภิกฺขเว ธมฺโมติ.\csromanspacer{}\csromanspacer{}ปญฺจมํ.}

\csromanmatikaanchor{mtk.214}
\csromantocmarkref{subsection}{2. อริยธมฺมสุตฺต}{mtk.214}
\bookmark[dest={mtk.214},level=2]{2. อริยธมฺมสุตฺต}
\csromanheader{2}{6. อาสวสุตฺต}
\proseitem{183}{สาสวญฺจ โว ภิกฺขเว ธมฺมํ เทเสสฺสามิ อนาสวญฺจ, ตํ สุณาถ ฯเปฯ.\csromanspacer{} กตโม จ ภิกฺขเว สาสโว ธมฺโม? ปาณาติปาโต ฯเปฯ มิจฺฉาทิฏฺฐิ.\csromanspacer{} อยํ วุจฺจติ ภิกฺขเว สาสโว ธมฺโม.}

\prose{กตโม จ ภิกฺขเว อนาสโว ธมฺโม? ปาณาติปาตา เวรมณี ฯเปฯ สมฺมาทิฏฺฐิ.\csromanspacer{} อยํ วุจฺจติ ภิกฺขเว อนาสโว ธมฺโมติ.\csromanspacer{}\csromanspacer{}ฉฏฺฐํ.}

\csromanmatikaanchor{mtk.215}
\csromantocmarkref{subsection}{3. กุสลสุตฺต}{mtk.215}
\bookmark[dest={mtk.215},level=2]{3. กุสลสุตฺต}
\gambhira{ทสกนิปาตปาฬิ}
\csromanheader{2}{7. วชฺชสุตฺต}
\proseitem{184}{\csromanfolio{484}สาวชฺชญฺจ โว ภิกฺขเว ธมฺมํ เทเสสฺสามิ อนวชฺชญฺจ, ตํ สุณาถ ฯเปฯ.\csromanspacer{} กตโม จ ภิกฺขเว สาวชฺโช ธมฺโม? ปาณาติปาโต ฯเปฯ มิจฺฉาทิฏฺฐิ.\csromanspacer{} อยํ วุจฺจติ ภิกฺขเว สาวชฺโช ธมฺโม.}

\csromanmatikaanchor{mtk.216}
\csromantocmarkref{subsection}{4. อตฺถสุตฺต}{mtk.216}
\bookmark[dest={mtk.216},level=2]{4. อตฺถสุตฺต}
\prose{กตโม จ ภิกฺขเว อนวชฺโช ธมฺโม? ปาณาติปาตา เวรมณี ฯเปฯ สมฺมาทิฏฺฐิ.\csromanspacer{} อยํ วุจฺจติ ภิกฺขเว อนวชฺโช ธมฺโมติ.\csromanspacer{}\csromanspacer{}สตฺตมํ.}

\csromanheader{2}{8. ตปนียสุตฺต}
\proseitem{185}{ตปนียญฺจ โว ภิกฺขเว ธมฺมํ เทเสสฺสามิ อตปนียญฺจ, ตํ สุณาถ ฯเปฯ.\csromanspacer{} กตโม จ ภิกฺขเว ตปนีโย ธมฺโม? ปาณาติปาโต ฯเปฯ มิจฺฉาทิฏฺฐิ.\csromanspacer{} อยํ วุจฺจติ ภิกฺขเว ตปนีโย ธมฺโม.}

\csromanmatikaanchor{mtk.217}
\csromantocmarkref{subsection}{5. ธมฺมสุตฺต}{mtk.217}
\bookmark[dest={mtk.217},level=2]{5. ธมฺมสุตฺต}
\prose{กตโม จ ภิกฺขเว อตปนีโย ธมฺโม? ปาณาติปาตา เวรมณี ฯเปฯ สมฺมาทิฏฺฐิ.\csromanspacer{} อยํ วุจฺจติ ภิกฺขเว อตปนีโย ธมฺโมติ.\csromanspacer{}\csromanspacer{}อฏฺฐมํ.}

\csromanheader{2}{9. อาจยคามิสุตฺต}
\proseitem{186}{อาจยคามิญฺจ โว ภิกฺขเว ธมฺมํ เทเสสฺสามิ อปจยคามิญฺจ, ตํ สุณาถ ฯเปฯ.\csromanspacer{} กตโม จ ภิกฺขเว อาจยคามี ธมฺโม? ปาณาติปาโต ฯเปฯ มิจฺฉาทิฏฺฐิ.\csromanspacer{} อยํ วุจฺจติ ภิกฺขเว อาจยคามี ธมฺโม.}

\csromanmatikaanchor{mtk.218}
\csromantocmarkref{subsection}{6. อาสวสุตฺต}{mtk.218}
\bookmark[dest={mtk.218},level=2]{6. อาสวสุตฺต}
\prose{กตโม จ ภิกฺขเว อปจยคามี ธมฺโม? ปาณาติปาตา เวรมณี ฯเปฯ สมฺมาทิฏฺฐิ.\csromanspacer{} อยํ วุจฺจติ ภิกฺขเว อปจยคามี ธมฺโมติ.\csromanspacer{}\csromanspacer{}นวมํ.}

\csromanheader{2}{10. ทุกฺขุทฺรยสุตฺต}
\proseitem{187}{ทุกฺขุทฺรยญฺจ โว ภิกฺขเว ธมฺมํ เทเสสฺสามิ สุขุทฺรยญฺจ, ตํ สุณาถ ฯเปฯ.\csromanspacer{} กตโม จ ภิกฺขเว ทุกฺขุทฺรโย ธมฺโม? ปาณาติปาโต ฯเปฯ มิจฺฉาทิฏฺฐิ.\csromanspacer{} อยํ วุจฺจติ ภิกฺขเว ทุกฺขุทฺรโย ธมฺโม.}

\prose{กตโม จ ภิกฺขเว สุขุทฺรโย ธมฺโม? ปาณาติปาตา เวรมณี ฯเปฯ สมฺมาทิฏฺฐิ.\csromanspacer{} อยํ วุจฺจติ ภิกฺขเว สุขุทฺรโย ธมฺโมติ.\csromanspacer{}\csromanspacer{}ทสมํ.}

\csromanmatikaanchor{mtk.219}
\csromantocmarkref{subsection}{7. วชฺชสุตฺต}{mtk.219}
\bookmark[dest={mtk.219},level=2]{7. วชฺชสุตฺต}
\csromanheader{2}{\textbf{(19) 4. อริยมคฺควคฺค} \textbf{11. วิปากสุตฺต}}
\proseitem{188}{\csromanfolio{485}ทุกฺขวิปากญฺจ โว ภิกฺขเว ธมฺมํ เทเสสฺสามิ สุขวิปากญฺจ, ตํ สุณาถ ฯเปฯ.\csromanspacer{} กตโม จ ภิกฺขเว ทุกฺขวิปาโก ธมฺโม? ปาณาติปาโต ฯเปฯ มิจฺฉาทิฏฺฐิ.\csromanspacer{} อยํ วุจฺจติ ภิกฺขเว ทุกฺขวิปาโก ธมฺโม.}

\prose{กตโม จ ภิกฺขเว สุขวิปาโก ธมฺโม? ปาณาติปาตา เวรมณี ฯเปฯ สมฺมาทิฏฺฐิ.\csromanspacer{} อยํ วุจฺจติ ภิกฺขเว สุขวิปาโก ธมฺโมติ.\csromanspacer{}\csromanspacer{}เอกาทสมํ.}

\csromanmatikaanchor{mtk.220}
\csromantocmarkref{subsection}{8. ตปณียสุตฺต}{mtk.220}
\bookmark[dest={mtk.220},level=2]{8. ตปณียสุตฺต}
\csromanheader{2}{สาธุวคฺโค ตติโย.}
\csromansectionrule
\csromanheader{1}{\textbf{(19) 4. อริยมคฺควคฺค}}
\csromanheader{2}{1. อริยมคฺคสุตฺต}
\csromanmatikaanchor{mtk.221}
\csromantocmarkref{subsection}{9. อาจยคามิสุตฺต}{mtk.221}
\bookmark[dest={mtk.221},level=2]{9. อาจยคามิสุตฺต}
\proseitem{189}{อริยมคฺคญฺจ โว ภิกฺขเว เทเสสฺสามิ อนริยมคฺคญฺจ, ตํ สุณาถ ฯเปฯ.\csromanspacer{} กตโม จ ภิกฺขเว อนริโย มคฺโค? ปาณาติปาโต ฯเปฯ มิจฺฉาทิฏฺฐิ.\csromanspacer{} อยํ วุจฺจติ ภิกฺขเว อนริโย มคฺโค.}

\prose{กตโม จ ภิกฺขเว อริโย มคฺโค? ปาณาติปาตา เวรมณี ฯเปฯ สมฺมาทิฏฺฐิ.\csromanspacer{} อยํ วุจฺจติ ภิกฺขเว อริโย มคฺโคติ.\csromanspacer{}\csromanspacer{}ปฐมํ.}

\csromanheader{2}{2. กณฺหมคฺคสุตฺต}
\csromanmatikaanchor{mtk.222}
\csromantocmarkref{subsection}{10. ทุกฺขุทฺรยสุตฺต}{mtk.222}
\bookmark[dest={mtk.222},level=2]{10. ทุกฺขุทฺรยสุตฺต}
\proseitem{190}{กณฺหมคฺคญฺจ โว ภิกฺขเว เทเสสฺสามิ สุกฺกมคฺคญฺจ, ตํ สุณาถ ฯเปฯ.\csromanspacer{} กตโม จ ภิกฺขเว กโณฺห มคฺโค? ปาณาติปาโต ฯเปฯ มิจฺฉาทิฏฺฐิ.\csromanspacer{} อยํ วุจฺจติ ภิกฺขเว กโณฺห มคฺโค.}

\prose{กตโม จ ภิกฺขเว สุกฺโก มคฺโค? ปาณาติปาตา เวรมณี ฯเปฯ สมฺมาทิฏฺฐิ.\csromanspacer{} อยํ วุจฺจติ ภิกฺขเว สุกฺโก มคฺโคติ.\csromanspacer{}\csromanspacer{}ทุติยํ.}

\csromanheader{2}{3. สทฺธมฺมสุตฺต}
\csromanmatikaanchor{mtk.223}
\csromantocmarkref{subsection}{11. วิปากสุตฺต}{mtk.223}
\bookmark[dest={mtk.223},level=2]{11. วิปากสุตฺต}
\proseitem{191}{สทฺธมฺมญฺจ โว ภิกฺขเว เทเสสฺสามิ อสทฺธมฺมญฺจ, ตํ สุณาถ ฯเปฯ.\csromanspacer{} กตโม จ ภิกฺขเว อสทฺธมฺโม? ปาณาติปาโต ฯเปฯ มิจฺฉาทิฏฺฐิ.\csromanspacer{} อยํ วุจฺจติ ภิกฺขเว อสทฺธมฺโม.}

\gambhira{ทสกนิปาตปาฬิ}
\prose{\csromanfolio{486}กตโม จ ภิกฺขเว สทฺธมฺโม? ปาณาติปาตา เวรมณี ฯเปฯ สมฺมาทิฏฺฐิ.\csromanspacer{} อยํ วุจฺจติ ภิกฺขเว สทฺธมฺโมติ.\csromanspacer{}\csromanspacer{}ตติยํ.}

\csromanheader{2}{4. สปฺปุริสธมฺมสุตฺต}
\csromanmatikaanchor{mtk.224}
\csromantocmarknopage{section}{(19) 4. อริยมคฺควคฺค}{mtk.224}
\bookmark[dest={mtk.224},level=1]{(19) 4. อริยมคฺควคฺค}
\proseitem{192}{สปฺปุริสธมฺมญฺจ โว ภิกฺขเว เทเสสฺสามิ อสปฺปุริสธมฺมญฺจ, ตํ สุณาถ ฯเปฯ.\csromanspacer{} กตโม จ ภิกฺขเว อสปฺปุริสธมฺโม? ปาณาติปาโต ฯเปฯ มิจฺฉาทิฏฺฐิ.\csromanspacer{} อยํ วุจฺจติ ภิกฺขเว อสปฺปุริสธมฺโม.}

\csromanmatikaanchor{mtk.225}
\csromantocmarkref{subsection}{1. อริยมคฺคสุตฺต}{mtk.225}
\bookmark[dest={mtk.225},level=2]{1. อริยมคฺคสุตฺต}
\prose{กตโม จ ภิกฺขเว สปฺปุริสธมฺโม? ปาณาติปาตา เวรมณี ฯเปฯ สมฺมาทิฏฺฐิ.\csromanspacer{} อยํ วุจฺจติ ภิกฺขเว สปฺปุริสธมฺโมติ.\csromanspacer{}\csromanspacer{}จตุตฺถํ.}

\csromanheader{2}{5. อุปฺปาเทตพฺพธมฺมสุตฺต}
\proseitem{193}{อุปฺปาเทตพฺพญฺจ โว ภิกฺขเว ธมฺมํ เทเสสฺสามิ น อุปฺปาเทตพฺพญฺจ, ตํ สุณาถ ฯเปฯ.\csromanspacer{} กตโม จ ภิกฺขเว น อุปฺปาเทตพฺโพ ธมฺโม? ปาณาติปาโต ฯเปฯ มิจฺฉาทิฏฺฐิ.\csromanspacer{} อยํ วุจฺจติ ภิกฺขเว น อุปฺปาเทตพฺโพ ธมฺโม.}

\csromanmatikaanchor{mtk.226}
\csromantocmarkref{subsection}{2. กณฺหมคฺคสุตฺต}{mtk.226}
\bookmark[dest={mtk.226},level=2]{2. กณฺหมคฺคสุตฺต}
\prose{กตโม จ ภิกฺขเว อุปฺปาเทตพฺโพ ธมฺโม? ปาณาติปาตา เวรมณี ฯเปฯ สมฺมาทิฏฺฐิ.\csromanspacer{} อยํ วุจฺจติ ภิกฺขเว อุปฺปาเทตพฺโพ ธมฺโมติ.\csromanspacer{}\csromanspacer{}ปญฺจมํ.}

\csromanheader{2}{6. อาเสวิตพฺพธมฺมสุตฺต}
\proseitem{194}{อาเสวิตพฺพญฺจ โว ภิกฺขเว ธมฺมํ เทเสสฺสามิ นาเสวิตพฺพญฺจ, ตํ สุณาถ ฯเปฯ.\csromanspacer{} กตโม จ ภิกฺขเว นาเสวิตพฺโพ ธมฺโม? ปาณาติปาโต ฯเปฯ มิจฺฉาทิฏฺฐิ.\csromanspacer{} อยํ วุจฺจติ ภิกฺขเว นาเสวิตพฺโพ ธมฺโม.}

\csromanmatikaanchor{mtk.227}
\csromantocmarkref{subsection}{3. สทฺธมฺมสุตฺต}{mtk.227}
\bookmark[dest={mtk.227},level=2]{3. สทฺธมฺมสุตฺต}
\prose{กตโม จ ภิกฺขเว อาเสวิตพฺโพ ธมฺโม? ปาณาติปาตา เวรมณี ฯเปฯ สมฺมาทิฏฺฐิ.\csromanspacer{} อยํ วุจฺจติ ภิกฺขเว อาเสวิตพฺโพ ธมฺโมติ.\csromanspacer{}\csromanspacer{}ฉฏฺฐํ.}

\csromanheader{2}{7. ภาเวตพฺพธมฺมสุตฺต}
\proseitem{195}{ภาเวตพฺพญฺจ โว ภิกฺขเว ธมฺมํ เทเสสฺสามิ น ภาเวตพฺพญฺจ, ตํ สุณาถ ฯเปฯ.\csromanspacer{} กตโม จ ภิกฺขเว น ภาเวตพฺโพ ธมฺโม?}

\vspace{\tipitakaparskip}
\csromancenter{(19) 4.\csromanspacer{} อริยมคฺควคฺค ปาณาติปาโต ฯเปฯ มิจฺฉาทิฏฺฐิ.\csromanspacer{} อยํ วุจฺจติ ภิกฺขเว น ภาเวตพฺโพ ธมฺโม.}
\csromanmatikaanchor{mtk.228}
\csromantocmarkref{subsection}{4. สปฺปุริสธมฺมสุตฺต}{mtk.228}
\bookmark[dest={mtk.228},level=2]{4. สปฺปุริสธมฺมสุตฺต}
\prose{\csromanfolio{487}กตโม จ ภิกฺขเว ภาเวตพฺโพ ธมฺโม? ปาณาติปาตา เวรมณี ฯเปฯ สมฺมาทิฏฺฐิ.\csromanspacer{} อยํ วุจฺจติ ภิกฺขเว ภาเวตพฺโพ ธมฺโมติ.\csromanspacer{}\csromanspacer{}สตฺตมํ.}

\csromanheader{2}{8. พหุลีกาตพฺพสุตฺต}
\proseitem{196}{พหุลีกาตพฺพญฺจ โว ภิกฺขเว ธมฺมํ เทเสสฺสามิ น พหุลีกาตพฺพญฺจ, ตํ สุณาถ ฯเปฯ.\csromanspacer{} กตโม จ ภิกฺขเว น พหุลีกาตพฺโพ ธมฺโม? ปาณาติปาโต ฯเปฯ มิจฺฉาทิฏฺฐิ.\csromanspacer{} อยํ วุจฺจติ ภิกฺขเว น พหุลีกาตพฺโพ ธมฺโม.}

\csromanmatikaanchor{mtk.229}
\csromantocmarkref{subsection}{5. อุปฺปาเทตพฺพธมฺมสุตฺต}{mtk.229}
\bookmark[dest={mtk.229},level=2]{5. อุปฺปาเทตพฺพธมฺมสุตฺต}
\prose{กตโม จ ภิกฺขเว พหุลีกาตพฺโพ ธมฺโม? ปาณาติปาตา เวรมณี ฯเปฯ สมฺมาทิฏฺฐิ.\csromanspacer{} อยํ วุจฺจติ ภิกฺขเว พหุลีกาตพฺโพ ธมฺโมติ.\csromanspacer{}\csromanspacer{}อฏฺฐมํ.}

\csromanheader{2}{9. อนุสฺสริตพฺพสุตฺต}
\proseitem{197}{อนุสฺสริตพฺพญฺจ โว ภิกฺขเว ธมฺมํ เทเสสฺสามิ นานุสฺสริตพฺพญฺจ, ตํ สุณาถ ฯเปฯ.\csromanspacer{} กตโม จ ภิกฺขเว นานุสฺสริตพฺโพ ธมฺโม? ปาณาติปาโต ฯเปฯ มิจฺฉาทิฏฺฐิ.\csromanspacer{} อยํ วุจฺจติ ภิกฺขเว นานุสฺสริตพฺโพ ธมฺโม.}

\csromanmatikaanchor{mtk.230}
\csromantocmarkref{subsection}{6. อาเสวิตพฺพธมฺมสุตฺต}{mtk.230}
\bookmark[dest={mtk.230},level=2]{6. อาเสวิตพฺพธมฺมสุตฺต}
\prose{กตโม จ ภิกฺขเว อนุสฺสริตพฺโพ ธมฺโม? ปาณาติปาตา เวรมณี ฯเปฯ สมฺมาทิฏฺฐิ.\csromanspacer{} อยํ วุจฺจติ ภิกฺขเว อนุสฺสริตพฺโพ ธมฺโมติ.\csromanspacer{}\csromanspacer{}นวมํ.}

\csromanheader{2}{10. สจฺฉิกาตพฺพสุตฺต}
\proseitem{198}{สจฺฉิกาตพฺพญฺจ โว ภิกฺขเว ธมฺมํ เทเสสฺสามิ น สจฺฉิกาตพฺพญฺจ, ตํ สุณาถ ฯเปฯ.\csromanspacer{} กตโม จ ภิกฺขเว น สจฺฉิกาตพฺโพ ธมฺโม? ปาณาติปาโต ฯเปฯ มิจฺฉาทิฏฺฐิ.\csromanspacer{} อยํ วุจฺจติ ภิกฺขเว น สจฺฉิกาตพฺโพ ธมฺโม.}

\csromanmatikaanchor{mtk.231}
\csromantocmarkref{subsection}{7. ภาเวตพฺพธมฺมสุตฺต}{mtk.231}
\bookmark[dest={mtk.231},level=2]{7. ภาเวตพฺพธมฺมสุตฺต}
\gambhira{ทสกนิปาตปาฬิ}
\prose{\csromanfolio{488}กตโม จ ภิกฺขเว สจฺฉิกาตพฺโพ ธมฺโม? ปาณาติปาตา เวรมณี ฯเปฯ สมฺมาทิฏฺฐิ.\csromanspacer{} อยํ วุจฺจติ ภิกฺขเว สจฺฉิกาตพฺโพ ธมฺโมติ.\csromanspacer{}\csromanspacer{}ทสมํ.}

\csromanheader{2}{อริยมคฺควคฺโค จตุตฺโถ.}
\csromansectionrule
\csromanheader{1}{\textbf{(20) 5. อปรปุคฺคลวคฺค}}
\csromanmatikaanchor{mtk.232}
\csromantocmarkref{subsection}{8. พหุลีกาตพฺพสุตฺต}{mtk.232}
\bookmark[dest={mtk.232},level=2]{8. พหุลีกาตพฺพสุตฺต}
\csromanheader{1}{\textbf{นเสวิตพฺพาทิสุตฺตานิ}}
\proseitem{199}{ทสหิ ภิกฺขเว ธมฺเมหิ สมนฺนาคโต ปุคฺคโล น เสวิตพฺโพ.\csromanspacer{} กตเมหิ ทสหิ? ปาณาติปาตี โหติ, อทินฺนาทายี โหติ, กาเมสุมิจฺฉาจารี โหติ, มุสาวาที โหติ, ปิสุณวาโจ โหติ, ผรุสวาโจ โหติ, สมฺผปฺปลาปี โหติ, อภิชฺฌาลุ โหติ, พฺยาปนฺนจิตฺโต โหติ, มิจฺฉาทิฏฺฐิโก โหติ.\csromanspacer{} อิเมหิ โข ภิกฺขเว ทสหิ ธมฺเมหิ สมนฺนาคโต ปุคฺคโล น เสวิตพฺโพ.}

\prose{ทสหิ ภิกฺขเว ธมฺเมหิ สมนฺนาคโต ปุคฺคโล เสวิตพฺโพ.\csromanspacer{} กตเมหิ ทสหิ? ปาณาติปาตา ปฏิวิรโต โหติ, อทินฺนาทานา ปฏิวิรโต โหติ, กาเมสุมิจฺฉาจารา ปฏิวิรโต โหติ, มุสาวาทา ปฏิวิรโต โหติ, ปิสุณาย วาจาย ปฏิวิรโต โหติ, ผรุสาย วาจาย ปฏิวิรโต โหติ, สมฺผปฺปลาปา ปฏิวิรโต โหติ, อนภิชฺฌาลุ โหติ, อพฺยาปนฺนจิตฺโต โหติ, สมฺมาทิฏฺฐิโก โหติ.\csromanspacer{} อิเมหิ โข ภิกฺขเว ทสหิ ธมฺเมหิ สมนฺนาคโต ปุคฺคโล เสวิตพฺโพ. (1)}

\csromanmatikaanchor{mtk.233}
\csromantocmarkref{subsection}{9. อนุสฺสริตพฺพสุตฺต}{mtk.233}
\bookmark[dest={mtk.233},level=2]{9. อนุสฺสริตพฺพสุตฺต}
\proseitem{200-209}{ทสหิ ภิกฺขเว ธมฺเมหิ สมนฺนาคโต ปุคฺคโล น ภชิตพฺโพ ฯเปฯ.\csromanspacer{} ภชิตพฺโพ.\csromanspacer{} น ปยิรุปาสิตพฺโพ.\csromanspacer{} ปยิรุปาสิตพฺโพ.\csromanspacer{} น ปุชฺโช โหติ.\csromanspacer{} ปุชฺโช โหติ.\csromanspacer{} น ปาสํโส โหติ.\csromanspacer{} ปาสํโส โหติ.\csromanspacer{} อคารโว โหติ.\csromanspacer{} คารโว โหติ.\csromanspacer{} อปฺปติสฺโส โหติ.\csromanspacer{} สปฺปติสฺโส โหติ.\csromanspacer{} น อาราธโก โหติ.\csromanspacer{} อาราธโก โหติ.\csromanspacer{} น วิสุชฺฌติ.\csromanspacer{} วิสุชฺฌติ.\csromanspacer{} มานํ นาธิโภติ\footnote{นาภิโภติ (สี) เหฏฺฐา 458 ปิฏฺเฐปิ.}.\csromanspacer{} มานํ อธิโภติ.\csromanspacer{} ปญฺญาย น วฑฺฒติ.\csromanspacer{} ปญฺญาย วฑฺฒติ ฯเปฯ. (2-11)}

\chapterheadpageread{(20) 5. อปรปุคฺคลวคฺค}
\proseitem{210}{\csromanfolio{489}ทสหิ ภิกฺขเว ธมฺเมหิ สมนฺนาคโต ปุคฺคโล พหุํ อปุญฺญํ ปสวติ.\csromanspacer{} พหุํ ปุญฺญํ ปสวติ.\csromanspacer{} กตเมหิ ทสหิ? ปาณาติปาตา ปฏิวิรโต โหติ, อทินฺนาทานา ปฏิวิรโต โหติ, กาเมสุมิจฺฉาจารา ปฏิวิรโต โหติ, มุสาวาทา ปฏิวิรโต โหติ, ปิสุณาย วาจาย ปฏิวิรโต โหติ, ผรุสาย วาจาย ปฏิวิรโต โหติ, สมฺผปฺปลาปา ปฏิวิรโต โหติ, อนภิชฺฌาลุ โหติ, อพฺยาปนฺนจิตฺโต โหติ, สมฺมาทิฏฺฐิโก โหติ.\csromanspacer{} อิเมหิ โข ภิกฺขเว ทสหิ ธมฺเมหิ สมนฺนาคโต ปุคฺคโล พหุํ ปุญฺญํ ปสวตีติ. (12)}

\csromanmatikaanchor{mtk.234}
\csromantocmarkref{subsection}{10. สจฺฉิกาตพฺพสุตฺต}{mtk.234}
\bookmark[dest={mtk.234},level=2]{10. สจฺฉิกาตพฺพสุตฺต}
\vspace{\tipitakaparskip}
\csromancenter{อปรปุคฺคลวคฺโค ปญฺจโม.}
\csromancenter{\textbf{จตุตฺโถ ปณฺณาสโก สมตฺโต.}}
\gambhira{ทสกนิปาตปาฬิ}
\csromanheader{1}{\textbf{(21) 1. กรชกายวคฺค}}
\csromanheader{2}{1. ปฐมนิรยสคฺคสุตฺต}
\csromanmatikaanchor{mtk.235}
\csromantocmarknopage{section}{(20) 5. อปรปุคฺคลวคฺค}{mtk.235}
\bookmark[dest={mtk.235},level=1]{(20) 5. อปรปุคฺคลวคฺค}
\proseitem{211}{\csromanfolio{490}ทสหิ ภิกฺขเว ธมฺเมหิ สมนฺนาคโต ยถาภตํ นิกฺขิตฺโต เอวํ นิรเย.\csromanspacer{} กตเมหิ ทสหิ? อิธ ภิกฺขเว เอกจฺโจ ปาณาติปาตี โหติ, ลุทฺโท โลหิตปาณิ หตปหเต นิวิฏฺโฐ, อทยาปนฺโน สพฺพปาณภูเตสุ\footnote{นตฺเถตฺถ ปาฐเภโท.}.}

\csromanmatikaanchor{mtk.236}
\csromantocmarkref{section}{นเสวิตพฺพาทิสุตฺตานิ}{mtk.236}
\bookmark[dest={mtk.236},level=1]{นเสวิตพฺพาทิสุตฺตานิ}
\prose{อทินฺนาทายี โหติ, ยํ ตํ ปรสฺส ปรวิตฺตูปกรณํ คามคตํ วา อรญฺญคตํ วา, ตํ อทินฺนํ เถยฺยสงฺขาตํ อาทาตา โหติ.}

\prose{กาเมสุ มิจฺฉาจารี โหติ, ยา ตา มาตุรกฺขิตา ปิตุรกฺขิตา มาตาปิตุรกฺขิตา ภาตุรกฺขิตา ภคินิรกฺขิตา ญาติรกฺขิตา โคตฺตรกฺขิตา ธมฺมรกฺขิตา สสามิกา สปริทณฺฑา อนฺตมโส มาลาคุฬปริกฺขิตฺตาปิ, ตถารูปาสุ จาริตฺตํ อาปชฺชิตา โหติ.}

\prose{มุสาวาที โหติ, สภคฺคโต วา ปริสคฺคโต วา ญาติมชฺฌคโต วา ปูคมชฺฌคโต วา ราชกุลมชฺฌคโต วา อภินีโต สกฺขิปุฏฺโฐ “เอหมฺโภ ปุริส ยํ ชานาสิ, ตํ วเทหี”ติ, โส อชานํ วา อาห “ชานามี”ติ, ชานํ วา อาห “น ชานามี”ติ, อปสฺสํ วา อาห “ปสฺสามี”ติ, ปสฺสํ วา อาห “น ปสฺสามี”ติ, อิติ อตฺตเหตุ วา ปรเหตุ วา อามิสกิญฺจิกฺขเหตุ วา สมฺปชานมุสา ภาสิตา โหติ.}

\prose{ปิสุณวาโจ โหติ, อิโต สุตฺวา อมุตฺร อกฺขาตา อิเมสํ เภทาย, อมุตฺร วา สุตฺวา อิเมสํ อกฺขาตา อมูสํ เภทาย.\csromanspacer{} อิติ สมคฺคานํ วา เภตฺตา ภินฺนานํ วา อนุปฺปทาตา, วคฺคาราโม วคฺครโต วคฺคนนฺที วคฺคกรณิํ วาจํ ภาสิตา โหติ.}

\prose{ผรุสวาโจ โหติ, ยา สา วาจา อณฺฑกา กกฺกสา ปรกฏุกา ปราภิสชฺชนี โกธสามนฺตา อสมาธิสํวตฺตนิกา, ตถารูปิํ วาจํ ภาสิตา โหติ.}

\chapterheadpageread{(21) 1. กรชกายวคฺค}
\prose{\csromanfolio{491}สมฺผปฺปลาปี โหติ, อกาลวาที อภูตวาที อนตฺถวาที อธมฺมวาที อวินยวาที อนิธานวติํ วาจํ ภาสิตา โหติ อกาเลน อนปเทสํ อปริยนฺตวติํ อนตฺถสํหิตํ.}

\prose{อภิชฺฌาลุ โหติ, ยํ ตํ ปรสฺส ปรวิตฺตูปกรณํ, ตํ อภิชฺฌาตา โหติ “อโห วต ยํ ปรสฺส, ตํ มม อสฺสา”ติ.}

\prose{พฺยาปนฺนจิตฺโต โหติ ปทุฏฺฐมนสงฺกปฺโป “อิเม สตฺตา หญฺญนฺตุ วา พชฺฌนฺตุ วา อุจฺฉิชฺชนฺตุ วา วินสฺสนฺตุ วา มา วา อเหสุนฺ”ติ.}

\csromanmatikaanchor{mtk.237}
\csromantocmarknopage{section}{(21) 1. กรชกายวคฺค}{mtk.237}
\bookmark[dest={mtk.237},level=1]{(21) 1. กรชกายวคฺค}
\prose{มิจฺฉาทิฏฺฐิโก โหติ วิปรีตทสฺสโน “นตฺถิ ทินฺนํ, นตฺถิ ยิฏฺฐํ, นตฺถิ หุตํ, นตฺถิ สุกตทุกฺกฏานํ กมฺมานํ ผลํ วิปาโก, นตฺถิ อยํ โลโก, นตฺถิ ปโร โลโก, นตฺถิ มาตา, นตฺถิ ปิตา, นตฺถิ สตฺตา โอปปาติกา, นตฺถิ โลเก สมณพฺราหฺมณา สมฺมคฺคตา สมฺมาปฏิปนฺนา, เย อิมญฺจ โลกํ ปรญฺจ โลกํ สยํ อภิญฺญา สจฺฉิกตฺวา ปเวเทนฺตี”ติ.\csromanspacer{} อิเมหิ โข ภิกฺขเว ทสหิ ธมฺเมหิ สมนฺนาคโต ยถาภตํ นิกฺขิตฺโต เอวํ นิรเย.}

\csromanmatikaanchor{mtk.238}
\csromantocmarkref{subsection}{1. ปฐมนิรยสคฺคสุตฺต}{mtk.238}
\bookmark[dest={mtk.238},level=2]{1. ปฐมนิรยสคฺคสุตฺต}
\prose{ทสหิ ภิกฺขเว ธมฺเมหิ สมนฺนาคโต ยถาภตํ นิกฺขิตฺโต เอวํ สคฺเค.\csromanspacer{} กตเมหิ ทสหิ? อิธ ภิกฺขเว เอกจฺโจ ปาณาติปาตํ ปหาย ปาณาติปาตา ปฏิวิรโต โหติ, นิหิตทณฺโฑ นิหิตสตฺโถ ลชฺชี ทยาปนฺโน สพฺพปาณภูตหิตานุกมฺปี วิหรติ.}

\prose{อทินฺนาทานํ ปหาย อทินฺนาทานา ปฏิวิรโต โหติ, ยํ ตํ ปรสฺส ปรวิตฺตูปกรณํ คามคตํ วา อรญฺญคตํ วา, น ตํ อทินฺนํ เถยฺยสงฺขาตํ อาทาตา โหติ.}

\prose{กาเมสุมิจฺฉาจารํ ปหาย กาเมสุมิจฺฉาจารา ปฏิวิรโต โหติ, ยา ตา มาตุรกฺขิตา ฯเปฯ อนฺตมโส มาลาคุฬปริกฺขิตฺตาปิ, ตถารูปาสุ น จาริตฺตํ อาปชฺชิตา โหติ.}

\prose{มุสาวาทํ ปหาย มุสาวาทา ปฏิวิรโต โหติ, สภคฺคโต วา ปริสคฺคโต วา ญาติมชฺฌคโต วา ปูคมชฺฌคโต วา ราชกุลมชฺฌคโต วา อภินีโต สกฺขิปุฏฺโฐ “เอหมฺโภ ปุริส ยํ ชานาสิ ตํ วเทหี”ติ.\csromanspacer{} โส อชานํ วา อาห “น ชานามี”ติ, ชานํ วา อาห “ชานามี”ติ, อปสฺสํ วา อาห “น ปสฺสามี”ติ, ปสฺสํ วา อาห “ปสฺสามี”ติ, อิติ อตฺตเหตุ วา ปรเหตุ วา อามิสกิญฺจิกฺขเหตุ วา น สมฺปชานมุสา ภาสิตา โหติ.}

\gambhira{ทสกนิปาตปาฬิ}
\prose{\csromanfolio{492}ปิสุณวาจํ ปหาย ปิสุณาย วาจาย ปฏิวิรโต โหติ, น อิโต สุตฺวา อมุตฺร อกฺขาตา อิเมสํ เภทาย, อมุตฺร วา สุตฺวา อิเมสํ อกฺขาตา อมูสํ เภทาย, อิติ ภินฺนานํ วา สนฺธาตา สหิตานํ วา อนุปฺปทาตา, สมคฺคาราโม สมคฺครโต สมคฺคนนฺที สมคฺคกรณิํ วาจํ ภาสิตา โหติ.}

\prose{ผรุสวาจํ ปหาย ผรุสาย วาจาย ปฏิวิรโต โหติ, ยา สา วาจา เนลา กณฺณสุขา เปมนียา หทยงฺคมา โปรี พหุชนกนฺตา พหุชนมนาปา, ตถารูปิํ วาจํ ภาสิตา โหติ.}

\prose{สมฺผปฺปลาปํ ปหาย สมฺผปฺปลาปา ปฏิวิรโต โหติ, กาลวาที ภูตวาที อตฺถวาที ธมฺมวาที วินยวาที นิธานวติํ วาจํ ภาสิตา โหติ กาเลน สาปเทสํ ปริยนฺตวติํ อตฺถสํหิตํ.}

\prose{อนภิชฺฌาลุ โหติ, ยํ ตํ ปรสฺส ปรวิตฺตูปกรณํ, ตํ อนภิชฺฌาตา โหติ “อโห วต ยํ ปรสฺส, ตํ มม อสฺสา”ติ.}

\prose{อพฺยาปนฺนจิตฺโต โหติ อปฺปทุฏฺฐมนสงฺกปฺโป “อิเม สตฺตา อเวรา โหนฺตุ อพฺยาปชฺชา อนีฆา, สุขี อตฺตานํ ปริหรนฺตู”ติ.}

\prose{สมฺมาทิฏฺฐิโก โหติ อวิปรีตทสฺสโน “อตฺถิ ทินฺนํ, อตฺถิ ยิฏฺฐํ, อตฺถิ หุตํ, อตฺถิ สุกฏ ทุกฺกฏานํ กมฺมานํ ผลํ วิปาโก, อตฺถิ อยํ โลโก, อตฺถิ ปโร โลโก, อตฺถิ มาตา, อตฺถิ ปิตา, อตฺถิ สตฺตา, โอปปาติกา, อตฺถิ โลเก สมณพฺราหฺมณา สมฺมคฺคตา สมฺมาปฏิปนฺนา, เย อิมญฺจ โลกํ ปรญฺจ โลกํ สยํ อภิญฺญา สจฺฉิกตฺวา ปเวเทนฺตี”ติ.\csromanspacer{} อิเมหิ โข ภิกฺขเว ทสหิ ธมฺเมหิ สมนฺนาคโต ยถาภตํ นิกฺขิตฺโต เอวํ สคฺเคติ.\csromanspacer{}\csromanspacer{}ปฐมํ.}

\csromanheader{2}{2. ทุติยนิรยสคฺคสุตฺต}
\proseitem{212}{ทสหิ ภิกฺขเว ธมฺเมหิ สมนฺนาคโต ยถาภตํ นิกฺขิตฺโต เอวํ นิรเย.\csromanspacer{} กตเมหิ ทสหิ? อิธ ภิกฺขเว เอกจฺโจ ปาณาติปาตี โหติ, ลุทฺโท โลหิตปาณิ หตปหเต นิวิฏฺโฐ, อทยาปนฺโน สพฺพปาณภูเตสุ.}

\prose{อทินฺนาทายี โหติ.\csromanspacer{} กาเมสุมิจฺฉาจารี โหติ.\csromanspacer{} มุสาวาที โหติ.\csromanspacer{} ปิสุณวาโจ โหติ.\csromanspacer{} ผรุสวาโจ โหติ.\csromanspacer{} สมฺผลาปี โหติ.}

\vspace{\tipitakaparskip}
\csromancenter{(21) 1.\csromanspacer{} กรชกายวคฺค อภิชฺฌาลุ โหติ.\csromanspacer{} พฺยาปนฺนจิตฺโต โหติ.\csromanspacer{} มิจฺฉาทิฏฺฐิโก โหติ วิปรีตทสฺสโน “นตฺถิ ทินฺนํ ฯเปฯ สยํ อภิญฺญา สจฺฉิกตฺวา ปเวเทนฺตี”ติ.\csromanspacer{} อิเมหิ โข ภิกฺขเว ทสหิ ธมฺเมหิ สมนฺนาคโต ยถาภตํ นิกฺขิตฺโต เอวํ นิรเย.}
\prose{\csromanfolio{493}ทสหิ ภิกฺขเว ธมฺเมหิ สมนฺนาคโต ยถาภตํ นิกฺขิตฺโต เอวํ สคฺเค.\csromanspacer{} กตเมหิ ทสหิ? อิธ ภิกฺขเว เอกจฺโจ ปาณาติปาตํ ปหาย ปาณาติปาตา ปฏิวิรโต โหติ, นิหิตทณฺโฑ นิหิตสตฺโถ ลชฺชี ทยาปนฺโน สพฺพปาณภูตหิตานุกมฺปี วิหรติ.}

\prose{อทินฺนาทานํ ปหาย อทินฺนาทานา ปฏิวิรโต โหติ.\csromanspacer{} กาเมสุมิจฺฉาจารํ ปหาย กาเมสุมิจฺฉาจารา ปฏิวิรโต โหติ.\csromanspacer{} มุสาวาทํ ปหาย มุสาวาทา ปฏิวิรโต โหติ.\csromanspacer{} ปิสุณํ วาจํ ปหาย ปิสุณาย วาจาย ปฏิวิรโต โหติ.\csromanspacer{} ผรุสํ วาจํ ปหาย ผรุสาย วาจาย ปฏิวิรโต โหติ.\csromanspacer{} สมฺผปฺปลาปํ ปหาย สมฺผปฺปลาปา ปฏิวิรโต โหติ.\csromanspacer{} อนภิชฺฌาลุ โหติ.\csromanspacer{} อพฺยาปนฺนจิตฺโต โหติ.\csromanspacer{} สมฺมาทิฏฺฐิโก โหติ อวิปรีตทสฺสโน “อตฺถิ ทินฺนํ ฯเปฯ เย อิมญฺจ โลกํ ปรญฺจ โลกํ สยํ อภิญฺญา สจฺฉิกตฺวา ปเวเทนฺตี”ติ.\csromanspacer{} อิเมหิ โข ภิกฺขเว ทสหิ ธมฺเมหิ สมนฺนาคโต ยถาภตํ นิกฺขิตฺโต เอวํ สคฺเคติ.\csromanspacer{}\csromanspacer{}ทุติยํ.}

\csromanheader{2}{3. มาตุคามสุตฺต}
\proseitem{213}{ทสหิ ภิกฺขเว ธมฺเมหิ สมนฺนาคโต มาตุคาโม ยถาภตํ นิกฺขิตฺโต เอวํ นิรเย.\csromanspacer{} กตเมหิ ทสหิ? ปาณาติปาตี โหติ ฯเปฯ.\csromanspacer{} อทินฺนาทายี โหติ.\csromanspacer{} กาเมสุมิจฺฉาจารี โหติ.\csromanspacer{} มุสาวาที โหติ.\csromanspacer{} ปิสุณวาโจ โหติ.\csromanspacer{} ผรุสวาโจ โหติ.\csromanspacer{} สมฺผปฺปลาปี โหติ.\csromanspacer{} อภิชฺฌาลุ โหติ.\csromanspacer{} พฺยาปนฺนจิตฺโต โหติ.\csromanspacer{} มิจฺฉาทิฏฺฐิโก โหติ.\csromanspacer{} อิเมหิ โข ภิกฺขเว ทสหิ ธมฺเมหิ สมนฺนาคโต มาตุคาโม ยถาภตํ นิกฺขิตฺโต เอวํ นิรเย.}

\prose{ทสหิ ภิกฺขเว ธมฺเมหิ สมนฺนาคโต มาตุคาโม ยถาภตํ นิกฺขิตฺโต เอวํ สคฺเค.\csromanspacer{} กตเมหิ ทสหิ? ปาณาติปาตา ปฏิวิรโต โหติ ฯเปฯ.\csromanspacer{} อทินฺนาทานา ปฏิวิรโต โหติ.\csromanspacer{} กาเมสุมิจฺฉาจารา ปฏิวิรโต โหติ.\csromanspacer{} มุสาวาทา ปฏิวิรโต โหติ.\csromanspacer{} ปิสุณาย วาจาย}

\gambhira{ทสกนิปาตปาฬิ}
\prose{\csromanfolio{494}ปฏิวิรโต โหติ.\csromanspacer{} ผรุสาย วาจาย ปฏิวิรโต โหติ.\csromanspacer{} สมฺผปฺปลาปา ปฏิวิรโต โหติ.\csromanspacer{} อนภิชฺฌาลุ โหติ.\csromanspacer{} อพฺยาปนฺนจิตฺโต โหติ.\csromanspacer{} สมฺมาทิฏฺฐิโก โหติ.\csromanspacer{} อิเมหิ โข ภิกฺขเว ทสหิ ธมฺเมหิ สมนฺนาคโต มาตุคาโม ยถาภตํ นิกฺขิตฺโต เอวํ สคฺเคติ.\csromanspacer{}\csromanspacer{}ตติยํ.}

\csromanheader{2}{4. อุปาสิกาสุตฺต}
\csromanmatikaanchor{mtk.239}
\csromantocmarkref{subsection}{2. ทุติยนิรยสคฺคสุตฺต}{mtk.239}
\bookmark[dest={mtk.239},level=2]{2. ทุติยนิรยสคฺคสุตฺต}
\proseitem{214}{ทสหิ ภิกฺขเว ธมฺเมหิ สมนฺนาคตา อุปาสิกา ยถาภตํ นิกฺขิตฺตา เอวํ นิรเย.\csromanspacer{} กตเมหิ ทสหิ? ปาณาติปาตินี โหติ ฯเปฯ มิจฺฉาทิฏฺฐิกา โหติ.\csromanspacer{} อิเมหิ โข ภิกฺขเว ทสหิ ธมฺเมหิ สมนฺนาคตา อุปาสิกา ยถาภตํ นิกฺขิตฺตา เอวํ นิรเย.}

\prose{ทสหิ ภิกฺขเว ธมฺเมหิ สมนฺนาคตา อุปาสิกา ยถาภตํ นิกฺขิตฺตา เอวํสคฺเค.\csromanspacer{} กตเมหิ ทสหิ? ปาณาติปาตา ปฏิวิรตา โหติ ฯเปฯ สมฺมาทิฏฺฐิกา โหติ, อิเมหิ โข ภิกฺขเว ทสหิ ธมฺเมหิ สมนฺนาคตา อุปาสิกา ยถาภตํ นิกฺขิตฺตา เอวํ สคฺเค.\csromanspacer{}\csromanspacer{}จตุตฺถํ.}

\csromanheader{2}{5. วิสารทสุตฺต}
\proseitem{215}{ทสหิ ภิกฺขเว ธมฺเมหิ สมนฺนาคตา อุปาสิกา อวิสารทา อคารํ อชฺฌาวสติ.\csromanspacer{} กตเมหิ ทสหิ? ปาณาติปาตินี โหติ, อทินฺนาทายินีโหติ, กาเมสุมิจฺฉาจารินี โหติ, มุสาวาทินี โหติ, ปิสุณา วาจา โหติ, ผรุสวาจา โหติ, สมฺผปฺปลาปินี โหติ, อภิชฺฌาลุนี โหติ, พฺยาปนฺนจิตฺตา โหติ, มิจฺฉาทิฏฺฐิกา โหติ.\csromanspacer{} อิเมหิ โข ภิกฺขเว ทสหิ ธมฺเมหิ สมนฺนาคตา อุปาสิกา อวิสารทา อคารํ อชฺฌาวสติ.}

\prose{ทสหิ ภิกฺขเว ธมฺเมหิ สมนฺนาคตา อุปาสิกา วิสารทา อคารํ อชฺฌาวสติ.\csromanspacer{} กตเมหิ ทสหิ? ปาณาติปาตา ปฏิวิรตา โหติ, อทินฺนาทานา ปฏิวิรตา โหติ, กาเมสุมิจฺฉาจารา ปฏิวิรตา โหติ, มุสาวาทา ปฏิวิรตา โหติ, ปิสุณาย วาจาย ปฏิวิรตา โหติ, ผรุสาย วาจาย ปฏิวิรตา โหติ, สมฺผปฺปลาปา ปฏิวิรตา โหติ, อนภิชฺฌาลุนี โหติ, อพฺยาปนฺนจิตฺตา โหติ, สมฺมาทิฏฺฐิกา โหติ.\csromanspacer{} อิเมหิ โข ภิกฺขเว ทสหิ ธมฺเมหิ สมนฺนาคตา อุปาสิกา วิสารทา อคารํ อชฺฌาวสตีติ.\csromanspacer{}\csromanspacer{}ปญฺจมํ.}

\csromanheader{2}{(21) 1. กรชกายวคฺค \textbf{6. สํสปฺปนียสุตฺต}}
\csromanmatikaanchor{mtk.240}
\csromantocmarkref{subsection}{3. มาตุคามสุตฺต}{mtk.240}
\bookmark[dest={mtk.240},level=2]{3. มาตุคามสุตฺต}
\proseitem{216}{\csromanfolio{495}สํสปฺปนียปริยายํ โว ภิกฺขเว ธมฺมปริยายํ เทเสสฺสามิ, ตํ สุณาถ สาธุกํ มนสิ กโรถ, ภาสิสฺสามีติ. “เอวํ ภนฺเต”ติ โข เต ภิกฺขู ภควโต ปจฺจสฺโสสุํ.\csromanspacer{} ภควา เอตทโวจ\csromandash{}}

\prose{กตโม จ ภิกฺขเว สํสปฺปนียปริยาโย ธมฺมปริยาโย? กมฺมสฺสกา ภิกฺขเว สตฺตา กมฺมทายาทา กมฺมโยนี กมฺมพนฺธู กมฺมปฏิสรณา, ยํ กมฺมํ กโรนฺติ กลฺยาณํ วา ปาปกํ วา, ตสฺส ทายาทา ภวนฺติ.}

\prose{อิธ ภิกฺขเว เอกจฺโจ ปาณาติปาตี โหติ.\csromanspacer{} ลุทฺโท โลหิตปาณิ หตปหเต นิวิฏฺโฐ, อทยาปนฺโน สพฺพปาณภูเตสุ.\csromanspacer{} โส สํสปฺปติ กาเยน, สํสปฺปติ วาจาย, สํสปฺปติ มนสา.\csromanspacer{} ตสฺส ชิมฺหํ กายกมฺมํ โหติ, ชิมฺหํ วจีกมฺมํ, ชิมฺหํ มโนกมฺมํ, ชิมฺหา คติ, ชิมฺหุปปตฺติ.}

\prose{ชิมฺหคติกสฺส โข ปนาหํ ภิกฺขเว ชิมฺหุปปตฺติกสฺส ทฺวินฺนํ คตีนํ อญฺญตรํ คติํ วทามิ, เย วา เอกนฺตทุกฺขา นิรยา, ยา วา สํสปฺปชาติกา ติรจฺฉานโยนิ.\csromanspacer{} กตมา จ สา ภิกฺขเว สํสปฺปชาติกา ติรจฺฉานโยนิ? อหิ วิจฺฉิกา สตปที นกุลา พิฬารา มูสิกา อุลูกา, เย วา ปนญฺเญปิ เกจิ ติรจฺฉานโยนิกา สตฺตา มนุสฺเส ทิสฺวา สํสปฺปนฺติ.\csromanspacer{} อิติ โข ภิกฺขเว ภูตา ภูตสฺส อุปปตฺติ โหติ.\csromanspacer{} ยํ กโรติ เตน อุปปชฺชติ, อุปปนฺนเมนํ ผสฺสา ผุสนฺติ.\csromanspacer{} เอวมหํ ภิกฺขเว กมฺมทายาทา สตฺตาติ วทามิ.}

\prose{อิธ ปน ภิกฺขเว เอกจฺโจ อทินฺนาทายี โหติ ฯเปฯ.\csromanspacer{} กาเมสุมิจฺฉาจารี โหติ.\csromanspacer{} มุสาวาที โหติ.\csromanspacer{} ปิสุณวาโจ โหติ.\csromanspacer{} ผรุสวาโจ โหติ.\csromanspacer{} สมฺผปฺปลาปี โหติ.\csromanspacer{} อภิชฺฌาลุ โหติ.\csromanspacer{} พฺยาปนฺนจิตฺโต โหติ.\csromanspacer{} มิจฺฉาทิฏฺฐิโก โหติ วิปรีตทสฺสโน “นตฺถิ ทินฺนํ ฯเปฯ สยํ อภิญฺญา สจฺฉิกตฺวา ปเวเทนฺตี”ติ, โส สํสปฺปติ กาเยน, สํสปฺปติ วาจาย, สํสปฺปติ มนสา.\csromanspacer{} ตสฺส ชิมฺหํ กายกมฺมํ โหติ, ชิมฺหํ วจีกมฺมํ, ชิมฺหํ มโนกมฺมํ, ชิมฺหา คติ, ชิมฺหุปปตฺติ.}

\csromanmatikaanchor{mtk.241}
\csromantocmarkref{subsection}{4. อุปาสิกาสุตฺต}{mtk.241}
\bookmark[dest={mtk.241},level=2]{4. อุปาสิกาสุตฺต}
\prose{ชิมฺหคติกสฺส โข ปนาหํ ภิกฺขเว ชิมฺหุปปตฺติกสฺส ทฺวินฺนํ คตีนํ อญฺญตรํ คติํ วทามิ, เย วา เอกนฺตทุกฺขา นิรยา, ยา วา สํสปฺปชาติกา ติรจฺฉานโยนิ.\csromanspacer{} กตมา จ สา ภิกฺขเว สํสปฺปชาติกา}

\gambhira{ทสกนิปาตปาฬิ}
\prose{\csromanfolio{496}ติรจฺฉานโยนิ? อหิ วิจฺฉิกา สตปที นกุลา พิฬารา มูสิกา อุลูกา, เย วา ปนญฺเญปิ เกจิ ติรจฺฉานโยนิกา สตฺตา มนุสฺเส ทิสฺวา สํสปฺปนฺติ.\csromanspacer{} อิติ โข ภิกฺขเว ภูตา ภูตสฺส อุปปตฺติ โหติ.\csromanspacer{} ยํ กโรติ เตน อุปปชฺชติ, อุปปนฺนเมนํ ผสฺสา ผุสนฺติ.\csromanspacer{} เอวมหํ ภิกฺขเว กมฺมทายาทา สตฺตาติ วทามิ.\csromanspacer{} กมฺมสฺสกา ภิกฺขเว สตฺตา กมฺมทายาทา กมฺมโยนี กมฺมพนฺธู กมฺมปฏิสรณา, ยํ กมฺมํ กโรนฺติ กลฺยาณํ วา ปาปกํ วา, ตสฺส ทายาทา ภวนฺติ.}

\csromanmatikaanchor{mtk.242}
\csromantocmarkref{subsection}{5. วิสารทสุตฺต}{mtk.242}
\bookmark[dest={mtk.242},level=2]{5. วิสารทสุตฺต}
\prose{อิธ ภิกฺขเว เอกจฺโจ ปาณาติปาตํ ปหาย ปาณาติปาตา ปฏิวิรโต โหติ, นิหิตทณฺโฑ นิหิตสตฺโถ ลชฺชี ทยาปนฺโน สพฺพปาณภูตหิตานุกมฺปี วิหรติ.\csromanspacer{} โส น สํสปฺปติ กาเยน, น สํสปฺปติ วาจาย, น สํสปฺปติ มนสา.\csromanspacer{} ตสฺส อุชุ กายกมฺมํ โหติ, อุชุ วจีกมฺมํ, อุชุ มโนกมฺมํ, อุชุ คติ, อุชุปปตฺติ.}

\prose{อุชุคติกสฺส โข ปนาหํ ภิกฺขเว อุชุปปตฺติกสฺส ทฺวินฺนํ คตีนํ อญฺญตรํ คติํ วทามิ, เย วา เอกนฺตสุขา สคฺคา, ยานิ วา ปน ตานิ อุจฺจากุลานิ ขตฺติยมหาสาลกุลานิ วา พฺราหฺมณมหาสาลกุลานิ วา คหปติมหาสาลกุลานิ วา อฑฺฒานิ มหทฺธนานิ มหาโภคานิ ปหูตชาตรูปรชตานิ ปหูตวิตฺตูปกรณานิ ปหูตธนธญฺญานิ.\csromanspacer{} อิติ โข ภิกฺขเว ภูตา ภูตสฺส อุปปตฺติ โหติ, ยํ กโรติ เตน อุปปชฺชติ, อุปปนฺนเมนํ ผสฺสา ผุสนฺติ.\csromanspacer{} เอวมหํ ภิกฺขเว กมฺมทายาทา สตฺตาติ วทามิ.}

\prose{อิธ ปน ภิกฺขเว เอกจฺโจ อทินฺนาทานํ ปหาย อทินฺนาทานา ปฏิวิรโต โหติ ฯเปฯ.\csromanspacer{} กาเมสุมิจฺฉาจารา ปฏิวิรโต โหติ.\csromanspacer{} มุสาวาทํ ปหาย มุสาวาทา ปฏิวิรโต โหติ.\csromanspacer{} ปิสุณํ วาจํ ปหาย ปิสุณาย วาจาย ปฏิวิรโต โหติ.\csromanspacer{} ผรุสํ วาจํ ปหาย ผรุสาย วาจาย ปฏิวิรโต โหติ.\csromanspacer{} สมฺผปฺปลาปํ ปหาย สมฺผปฺปลาปา ปฏิวิรโต โหติ.\csromanspacer{} อนภิชฺฌาลุ โหติ.\csromanspacer{} อพฺยาปนฺนจิตฺโต โหติ.\csromanspacer{} สมฺมาทิฏฺฐิโก โหติ อวิปรีตทสฺสโน “อตฺถิ ทินฺนํ ฯเปฯ เย อิมญฺจ โลกํ ปรญฺจ โลกํ สยํ อภิญฺญา สจฺฉิกตฺวา ปเวเทนฺตี”ติ.\csromanspacer{} โส น สํสปฺปติ กาเยน, น สํสปฺปติ วาจาย, น สํสปฺปติ มนสา.\csromanspacer{} ตสฺส อุชุ กายกมฺมํ โหติ, อุชุ วจีกมฺมํ, อุชุ มโนกมฺมํ, อุชุ คติ, อุชุปปตฺติ.}

\csromanmatikaanchor{mtk.243}
\csromantocmarkref{subsection}{6. สํสปฺปนียสุตฺต}{mtk.243}
\bookmark[dest={mtk.243},level=2]{6. สํสปฺปนียสุตฺต}
\chapterheadpageread{(21) 1. กรชกายวคฺค}
\prose{\csromanfolio{497}อุชุคติกสฺส โข ปน อหํ ภิกฺขเว อุชุปปตฺติกสฺส ทฺวินฺนํ คตีนํ อญฺญตรํ คติํ วทามิ, เย วา เอกนฺตสุขา สคฺคา, ยานิ วา ปน ตานิ อุจฺจากุลานิ ขตฺติยมหาสาลกุลานิ วา พฺราหฺมณมหาสาลกุลานิ วา คหปติมหาสาลกุลานิ วา อฑฺฒานิ มหทฺธนานิ มหาโภคานิ ปหูตชาตรูปรชตานิ ปหูตวิตฺตูปกรณานิ ปหูตธนธญฺญานิ.\csromanspacer{} อิติ โข ภิกฺขเว ภูตา ภูตสฺส อุปปตฺติ โหติ.\csromanspacer{} ยํ กโรติ, เตน อุปปชฺชติ, อุปปนฺนเมนํ ผสฺสา ผุสนฺติ.\csromanspacer{} เอวมหํ ภิกฺขเว กมฺมทายาทา สตฺตาติ วทามิ.}

\prose{กมฺมสฺสกา ภิกฺขเว สตฺตา กมฺมทายาทา กมฺมโยนี กมฺมพนฺธู กมฺมปฏิสรณา, ยํ กมฺมํ กโรนฺติ กลฺยาณํ วา ปาปกํ วา, ตสฺส ทายาทา ภวนฺติ.\csromanspacer{} อยํ โข โส ภิกฺขเว สํสปฺปนียปริยาโย ธมฺมปริยาโยติ.\csromanspacer{}\csromanspacer{}ฉฏฺฐํ.}

\csromanheader{2}{7. ปฐมสญฺเจตนิกสุตฺต}
\proseitem{217}{นาหํ ภิกฺขเว สญฺเจตนิกานํ กมฺมานํ กตานํ อุปจิตานํ อปฺปฏิสํเวทิตฺวา\footnote{อปฺปฏิสํวิทิตฺวา (สี, สฺยา, อิ)} พฺยนฺตีภาวํ วทามิ, ตญฺจ โข ทิฏฺเฐว ธมฺเม อุปปชฺเช วา\footnote{อุปปชฺชํ วา (ก) อํ 2. 363 ปิฏฺเฐปิ ปสฺสิตพฺพํ, อุปปชฺช วา (ม 3. 257)} อปเร วา ปริยาเย.\csromanspacer{} น เตฺววาหํ ภิกฺขเว สญฺเจตนิกานํ กมฺมานํ กตานํ อุปจิตานํ อปฺปฏิสํเวทิตฺวา ทุกฺขสฺสนฺตกิริยํ วทามิ.}

\prose{ตตฺร ภิกฺขเว ติวิธา กายกมฺมนฺตสนฺโทสพฺยาปตฺติ อกุสลสญฺเจตนิกา ทุกฺขุทฺรยา ทุกฺขวิปากา\footnote{อกุสลํ สญฺเจตนิกํ ทุกฺขุทฺรยํ ทุกฺขวิปากํ (ก)} โหติ.\csromanspacer{} จตุพฺพิธา วจีกมฺมนฺตสนฺโทสพฺยาปตฺติ อกุสลสญฺเจตนิกา ทุกฺขุทฺรยา ทุกฺขวิปากา โหติ.\csromanspacer{} ติวิธา มโนกมฺมนฺตสนฺโทสพฺยาปตฺติ อกุสลสญฺเจตนิกา ทุกฺขุทฺรยา ทุกฺขวิปากา โหติ.}

\prose{กถญฺจ ภิกฺขเว ติวิธา กายกมฺมนฺตสนฺโทสพฺยาปตฺติ อกุสลสญฺเจตนิกา ทุกฺขุทฺรยา ทุกฺขวิปากา โหติ? อิธ ภิกฺขเว เอกจฺโจ ปาณาติปาตี โหติ, ลุทฺโท โลหิตปาณิ หตปหเต นิวิฏฺโฐ, อทยาปนฺโน สพฺพปาณภูเตสุ.}

\gambhira{ทสกนิปาตปาฬิ}
\prose{\csromanfolio{498}อทินฺนาทายี โหติ, ยํ ตํ ปรสฺส ปรวิตฺตูปกรณํ คามคตํ วา อรญฺญคตํ วา, ตํ อทินฺนํ เถยฺยสงฺขาตํ อาทาตา โหติ.}

\prose{กาเมสุมิจฺฉาจารี โหติ, ยา ตา มาตุรกฺขิตา ฯเปฯ อนฺตมโส มาลาคุฬปริกฺขิตฺตาปิ, ตถารูปาสุ จาริตฺตํ อาปชฺชิตา โหติ.\csromanspacer{} เอวํ โข ภิกฺขเว ติวิธา กายกมฺมนฺตสนฺโทสพฺยาปตฺติ อกุสลสญฺเจตนิกา ทุกฺขุทฺรยา ทุกฺขวิปากา โหติ.}

\prose{กถญฺจ ภิกฺขเว จตุพฺพิธา วจีกมฺมนฺตสนฺโทสพฺยาปตฺติ อกุสลสญฺเจตนิกา ทุกฺขุทฺรยา ทุกฺขวิปากา โหติ? อิธ ภิกฺขเว เอกจฺโจ มุสาวาที โหติ, สภคฺคโต วา ปริสคฺคโต วา ญาติมชฺฌคโต วา ปูคมชฺฌคโต วา ราชกุลมชฺฌคโต วา อภินีโต สกฺขิปุฏฺโฐ “เอหมฺโภ ปุริส ยํ ชานาสิ, ตํ วเทหี”ติ, โส อชานํ วา อาห “ชานามี”ติ, ชานํ วา อาห “น ชานามี”ติ, อปสฺสํ วา อาห “ปสฺสามี”ติ, ปสฺสํ วา อาห “น ปสฺสามี”ติ, อิติ อตฺตเหตุ วา ปรเหตุ วา อามิสกิญฺจิกฺขเหตุ วา สมฺปชานมุสา ภาสิตา โหติ.}

\prose{ปิสุณวาโจ โหติ, อิโต สุตฺวา อมุตฺร อกฺขาตา อิเมสํ เภทาย, อมุตฺร วา สุตฺวา อิเมสํ อกฺขาตา อมูสํ เภทาย.\csromanspacer{} อิติ สมคฺคานํ วา เภตฺตา ภินฺนานํ วา อนุปฺปทาตา วคฺคาราโม วคฺครโต วคฺคนนฺที วคฺคกรณิํ วาจํ ภาสิตา โหติ.}

\prose{ผรุสวาโจ โหติ, ยา สา วาจา อณฺฑกา กกฺกสา ปรกฏุกา ปราภิสชฺชนี โกธสามนฺตา อสมาธิสํวตฺตนิกา, ตถารูปิํ วาจํ ภาสิตา โหติ.}

\prose{สมฺผปฺปลาปี โหติ, อกาลวาที อภูตวาที อนตฺถวาที อธมฺมวาที อวินยวาที อนิธานวติํ วาจํ ภาสิตา โหติ, อกาเลน อนปเทสํ อปริยนฺตวติํ อนตฺถสํหิตํ.\csromanspacer{} เอวํ โข ภิกฺขเว จตุพฺพิธา วจีกมฺมนฺตสนฺโทสพฺยาปตฺติ อกุสลสญฺเจตนิกา ทุกฺขุทฺรยา ทุกฺขวิปากา โหติ.}

\prose{กถญฺจ ภิกฺขเว ติวิธา มโนกมฺมนฺตสนฺโทสพฺยาปตฺติ อกุสลสญฺเจตนิกา ทุกฺขุทฺรยา ทุกฺขวิปากา โหติ? อิธ ภิกฺขเว เอกจฺโจ อภิชฺฌาลุ โหติ.\csromanspacer{} ยํ ตํ ปรสฺส ปรวิตฺตุปกรณํ, ตํ อภิชฺฌาตา โหติ “อโห วต ยํ ปรสฺส, ตํ มม อสฺสา”ติ.}

\csromanmatikaanchor{mtk.244}
\csromantocmarkref{subsection}{7. ปฐมสญฺเจตนิกสุตฺต}{mtk.244}
\bookmark[dest={mtk.244},level=2]{7. ปฐมสญฺเจตนิกสุตฺต}
\chapterheadpageread{(21) 1. กรชกายวคฺค}
\prose{\csromanfolio{499}พฺยาปนฺนจิตฺโต โหติ ปทุฏฺฐมนสงฺกปฺโป “อิเม สตฺตา หญฺญนฺตุ วา พชฺฌนฺตุ วา อุจฺฉิชฺชนฺตุ วา วินสฺสนฺตุ วา มา วา อเหสุนฺ”ติ.}

\prose{มิจฺฉาทิฏฺฐิโก โหติ วิปรีตทสฺสโน “นตฺถิ ทินฺนํ ฯเปฯ เย อิมญฺจ โลกํ ปรญฺจ โลกํ สยํ อภิญฺญา สจฺฉิกตฺวา ปเวเทนฺตี”ติ.\csromanspacer{} เอวํ โข ภิกฺขเว ติวิธา มโนกมฺมนฺตสนฺโทสพฺยาปตฺติ อกุสลสญฺเจตนิกา ทุกฺขุทฺรยา ทุกฺขวิปากา โหติ.}

\prose{ติวิธ กายกมฺมนฺตสนฺโทสพฺยาปตฺติ อกุสลสญฺเจตนิกาเหตุ\footnote{...สญฺเจตนิกเหตุ (ก)} วา ภิกฺขเว สตฺตา กายสฺส เภทา ปรํ มรณา อปายํ ทุคฺคติํ วินิปาตํ นิรยํ อุปปชฺชนฺติ.\csromanspacer{} จตุพฺพิธวจีกมฺมนฺตสนฺโทสพฺยาปตฺติ อกุสลสญฺเจตนิกาเหตุ วา ภิกฺขเว สตฺตา กายสฺส เภทา ปรํ มรณา อปายํ ทุคฺคติํ วินิปาตํ นิรยํ อุปปชฺชนฺติ.\csromanspacer{} ติวิธมโนกมฺมนฺต- สนฺโทสพฺยาปตฺติ-อกุสลสญฺเจตนิกาเหตุ วา ภิกฺขเว สตฺตา กายสฺส เภทา ปรํ มรณา อปายํ ทุคฺคติํ วินิปาตํ นิรยํ อุปปชฺชนฺติ.}

\prose{เสยฺยถาปิ ภิกฺขเว อปณฺณโก มณิ อุทฺธํขิตฺโต เยน เยเนว ปติฏฺฐาติ, สุปฺปติฏฺฐิตํเยว ปติฏฺฐาติ.\csromanspacer{} เอวเมวํ โข ภิกฺขเว ติวิธกายกมฺมนฺตสนฺโทสพฺยาปตฺติ อกุสลสญฺเจตนิกาเหตุ วา สตฺตา กายสฺส เภทา ปรํ มรณา อปายํ ทุคฺคติํ วินิปาตํ นิรยํ อุปปชฺชนฺติ.\csromanspacer{} จตุพฺพิธวจีกมฺมนฺตสนฺโทสพฺยาปตฺติ- อกุสลสญฺเจตนิกาเหตุ วา สตฺตา กายสฺส เภทา ปรํ มรณา อปายํ ทุคฺคติํ วินิปาตํ นิรยํ อุปปชฺชนฺติ.\csromanspacer{} ติวิธมโนกมฺมนฺต- สนฺโทสพฺยาปตฺติ อกุสลสญฺเจตนิกาเหตุ วา สตฺตา กายสฺส เภทา ปรํ มรณา อปายํ ทุคฺคติํ วินิปาตํ นิรยํ อุปปชฺชนฺตีติ.}

\prose{นาหํ ภิกฺขเว สญฺเจตนิกานํ กมฺมานํ กตานํ อุปจิตานํ อปฺปฏิสํเวทิตฺวา พฺยนฺตีภาวํ วทามิ, ตญฺจ โข ทิฏฺเฐว ธมฺเม อุปปชฺเช วา อปเร วา ปริยาเย.\csromanspacer{} น เตฺววาหํ ภิกฺขเว สญฺเจตนิกานํ กมฺมานํ กตานํ อุปจิตานํ อปฺปฏิสํเวทิตฺวา ทุกฺขสฺสนฺตกิริยํ วทามิ.}

\prose{ตตฺร ภิกฺขเว ติวิธา กายกมฺมนฺตสมฺปตฺติ กุสลสญฺเจตนิกา สุขุทฺรยา สุขวิปากา โหติ.\csromanspacer{} จตุพฺพิธา วจีกมฺมนฺตสมฺปตฺติ กุสลสญฺเจตนิกา สุขุทฺรยา สุขวิปากา โหติ.\csromanspacer{} ติวิธา มโนกมฺมนฺตสมฺปตฺติ กุสลสญฺเจตนิกา สุขุทฺรยา สุขวิปากา โหติ.}

\gambhira{ทสกนิปาตปาฬิ}
\prose{\csromanfolio{500}กถญฺจ ภิกฺขเว ติวิธา กายกมฺมนฺตสมฺปตฺติ กุสลสญฺเจตนิกา สุขุทฺรยา สุขวิปากา โหติ? อิธ ภิกฺขเว เอกจฺโจ ปาณาติปาตา ปฏิวิรโต โหติ.\csromanspacer{} นิหิตทณฺโฑ นิหิตสตฺโถ ลชฺชี ทยาปนฺโน สพฺพปาณภูตหิตานุกมฺปี วิหรติ ฯเปฯ.}

\prose{อทินฺนาทานา ปฏิวิรโต โหติ.\csromanspacer{} ยํ ตํ ปรสฺส ปรวิตฺตูปกรณํ คามคตํ วา อรญฺญคตํ วา, น ตํ อทินฺนํ เถยฺยสงฺขาตํ อาทาตา โหติ.}

\prose{กาเมสุมิจฺฉาจารํ ปหาย กาเมสุมิจฺฉาจารา ปฏิวิรโต โหติ.\csromanspacer{} ยา ตา มาตุรกฺขิตา ฯเปฯ อนฺตมโส มาลาคุฬปริกฺขิตฺตาปิ, ตถารูปาสุ น จาริตฺตํ อาปชฺชิตา โหติ.\csromanspacer{} เอวํ โข ภิกฺขเว ติวิธา กายกมฺมนฺตสมฺปตฺติ กุสลสญฺเจตนิกา สุขุทฺรยา สุขวิปากา โหติ.}

\prose{กถญฺจ ภิกฺขเว จตุพฺพิธา วจีกมฺมนฺตสมฺปตฺติ กุสลสญฺเจตนิกา สุขุทฺรยา สุขวิปากา โหติ? อิธ ภิกฺขเว เอกจฺโจ มุสาวาทํ ปหาย มุสาวาทา ปฏิวิรโต โหติ.\csromanspacer{} สภคฺคโต วา ปริสคฺคโต วา ญาติมชฺฌคโต วา ปูคมชฺฌคโต วา ราชกุลมชฺฌคโต วา อภินีโต สกฺขิปุฏฺโฐ “เอหมฺโภ ปุริส ยํ ชานาสิ, ตํ วเทหี”ติ.\csromanspacer{} โส อชานํ วา อาห “น ชานามี”ติ, ชานํ วา อาห “ชานามี”ติ, อปสฺสํ วา อาห “น ปสฺสามี”ติ, ปสฺสํ วา อาห “ปสฺสามี”ติ, อิติ อตฺตเหตุ วา ปรเหตุ วา อามิสกิญฺจิกฺขเหตุ วา น สมฺปชานมุสา ภาสิตา โหติ.}

\prose{ปิสุณํ วาจํ ปหาย ปิสุณาย วาจาย ปฏิวิรโต โหติ, น อิโต สุตฺวา อมุตฺร อกฺขาตา อิเมสํ เภทาย, อมุตฺร วา สุตฺวา น อิเมสํ อกฺขาตา อมูสํ เภทาย.\csromanspacer{} อิติ ภินฺนานํ วา สนฺธาตา สหิตานํ วา อนุปฺปทาตา, สมคฺคาราโม สมคฺครโต สมคฺคนนฺทิํ สมคฺคกรณิํ วาจํ ภาสิตา โหติ.}

\prose{ผรุสํ วาจํ ปหาย ผรุสาย วาจาย ปฏิวิรโต โหติ.\csromanspacer{} ยา สา วาจา เนลา กณฺณสุขา เปมนียา หทยงฺคมา โปรี พหุชนกนฺตา พหุชนมนาปา, ตถารูปิํ วาจํ ภาสิตา โหติ.}

\prose{สมฺผปฺปลาปํ ปหาย สมฺผปฺปลาปา ปฏิวิรโต โหติ.\csromanspacer{} กาลวาที ภูตวาที อตฺถวาที ธมฺมวาที วินยวาที นิธานวติํ วาจํ ภาสิตา โหติ.}

\vspace{\tipitakaparskip}
\csromancenter{(21) 1.\csromanspacer{} กรชกายวคฺค กาเลน สาปเทสํ ปริยนฺตวติํ อตฺถสํหิตํ.\csromanspacer{} เอวํ โข ภิกฺขเว จตุพฺพิธา วจีกมฺมนฺตสมฺปตฺติ กุสลสญฺเจตนิกา สุขุทฺรยา สุขวิปากา โหติ.}
\prose{\csromanfolio{501}กถญฺจ ภิกฺขเว ติวิธา มโนกมฺมนฺตสมฺปตฺติ กุสลสญฺเจตนิกา สุขุทฺรยา สุขวิปากา โหติ? อิธ ภิกฺขเว เอกจฺโจ อนภิชฺฌาลุ โหติ.\csromanspacer{} ยํ ตํ ปรสฺส ปรวิตฺตูปกรณํ, ตํ อนภิชฺฌาตา โหติ “อโห วต ยํ ปรสฺส, ตํ มมสา”ติ.}

\prose{อพฺยาปนฺนจิตฺโต โหติ อปฺปทุฏฺฐมนสงฺกปฺโป “อิเม สตฺตา อเวรา โหนฺตุ อพฺยาปชฺฌา อนีฆา, สุขี อตฺตานํ ปริหรนฺตู”ติ.}

\prose{สมฺมาทิฏฺฐิโก โหติ อวิปรีตทสฺสโน “อตฺถิ ทินฺนํ, อตฺถิ ยิฏฺฐํ ฯเปฯ เย อิมญฺจ โลกํ ปรญฺจ โลกํ สยํ อภิญฺญา สจฺฉิกตฺวา ปเวเทนฺตี”ติ.\csromanspacer{} เอวํ โข ภิกฺขเว ติวิธา มโนกมฺมนฺตสมฺปตฺติ กุสลสญฺเจตนิกา สุขุทฺรยา สุขวิปากา โหติ.}

\prose{ติวิธกายกมฺมนฺตสมฺปตฺติกุสลสญฺเจตนิกาเหตุ วา ภิกฺขเว สตฺตา กายสฺส เภทา ปรํ มรณา สุคติํ สคฺคํ โลกํ อุปปชฺชนฺติ, จตุพฺพิธวจีกมฺมนฺตสมฺปตฺติกุสลสญฺเจตนิกาเหตุ วา ภิกฺขเว สตฺตา กายสฺส เภทา ปรํ มรณา สุคติํ สคฺคํ โลกํ อุปปชฺชนฺติ, ติวิธมโนกมฺมนฺตสมฺปตฺติกุสลสญฺเจตนิกาเหตุ วา ภิกฺขเว สตฺตา กายสฺส เภทา ปรํ มรณา สุคติํ สคฺคํ โลกํ อุปปชฺชนฺติ.}

\prose{เสยฺยถาปิ ภิกฺขเว อปณฺณโก มณิ อุทฺธํขิตฺโต เยน เยนว ปติฏฺฐาติ, สุปฺปติฏฺฐิตํเยว ปติฏฺฐาติ.\csromanspacer{} เอวเมวํ โข ภิกฺขเว ติวิธกายกมฺมนฺตสมฺปตฺติกุสลสญฺเจตนิกาเหตุ วา สตฺตา กายสฺส เภทา ปรํ มรณา สุคติํ สคฺคํ โลกํ อุปปชฺชนฺติ, จตุพฺพิธวจีกมฺมนฺตสมฺปตฺติกุสลสญฺเจตนิกาเหตุ วา สตฺตา กายสฺส เภทา ปรํ มรณา สุคติํ สคฺคํ โลกํ อุปปชฺชนฺติ, ติวิธมโนกมฺมนฺตสมฺปตฺติกุสลสญฺเจตนิกาเหตุ วา สตฺตา กายสฺส เภทา ปรํ มรณา สุคติํ สคฺคํ โลกํ อุปปชฺชนฺติ.\csromanspacer{} นาหํ ภิกฺขเว สญฺเจตนิกานํ กมฺมานํ กตานํ อุปจิตานํ อปฺปฏิสํเวทิตฺวา พฺยนฺตีภาวํ วทามิ, ตญฺจ โข ทิฏฺเฐว ธมฺเม อุปปชฺเช วา อปเร วา ปริยาเย.\csromanspacer{} น เตฺววาหํ ภิกฺขเว สญฺเจตนิกานํ กมฺมานํ กตานํ อุปจิตานํ อปฺปฏิสํเวทิตฺวา ทุกฺขสฺสนฺตกิริยํ วทามีติ.\csromanspacer{}\csromanspacer{}สตฺตมํ.\footnote{อฏฺฐกถายํ ปน อฏฺฐมสุตฺตมฺปิ เอตฺเถว ปริยาปนฺนํ วิย สํวณฺณนา ทิสฺสติ.}}

\gambhira{ทสกนิปาตปาฬิ}
\csromanheader{2}{8. ทุติยสญฺเจตนิกสุตฺต}
\proseitem{218}{\csromanfolio{502}นาหํ ภิกฺขเว สญฺเจตนิกานํ กมฺมานํ กตานํ อุปจิตานํ อปฺปฏิสํเวทิตฺวา พฺยนฺตีภาวํ วทามิ, ตญฺจ โข ทิฏฺเฐว ธมฺเม อุปปชฺเช วา อปเร วา ปริยาเย.\csromanspacer{} น เตฺววาหํ ภิกฺขเว สญฺเจตนิกานํ กมฺมานํ กตานํ อุปจิตานํ อปฺปฏิสํเวทิตฺวา ทุกฺขสฺสนฺตกิริยํ วทามิ.}

\prose{ตตฺร ภิกฺขเว ติวิธา กายกมฺมนฺตสนฺโทสพฺยาปตฺติ อกุสลสญฺเจตนิกา ทุกฺขุทฺรยา ทุกฺขวิปากา โหติ, จตุพฺพิธา วจีกมฺมนฺตสนฺโทสพฺยาปตฺติ อกุสลสญฺเจตนิกา ทุกฺขุทฺรยา ทุกฺขวิปากา โหติ, ติวิธา มโนกมฺมนฺตสนฺโทสพฺยาปตฺติ อกุสลสญฺเจตนิกา ทุกฺขุทฺรยา ทุกฺขวิปากา โหติ.}

\prose{กถญฺจ ภิกฺขเว ติวิธา กายกมฺมนฺตสนฺโทสพฺยาปตฺติ อกุสลสญฺเจตนิกา ทุกฺขุทฺรยา ทุกฺขวิปากา โหติ ฯเปฯ.\csromanspacer{} เอวํ โข ภิกฺขเว ติวิธา กายกมฺมนฺตสนฺโทสพฺยาปตฺติ อกุสลสญฺเจตนิกา ทุกฺขุทฺรยา ทุกฺขวิปากา โหติ.}

\prose{กถญฺจ ภิกฺขเว จตุพฺพิธา วจีกมฺมนฺตสนฺโทสพฺยาปตฺติ อกุสลสญฺเจตนิกา ทุกฺขุทฺรยา ทุกฺขวิปากา โหติ ฯเปฯ.\csromanspacer{} เอวํ โข ภิกฺขเว จตุพฺพิธา วจีกมฺมนฺตสนฺโทสพฺยาปตฺติ อกุสลสญฺเจตนิกา ทุกฺขุทฺรยา ทุกฺขวิปากา โหติ.}

\prose{กถญฺจ ภิกฺขเว ติวิธา มโนกมฺมนฺตสนฺโทสพฺยาปตฺติ อกุสลสญฺเจตนิกา ทุกฺขุทฺรยา ทุกฺขวิปากา โหติ ฯเปฯ.\csromanspacer{} เอวํ โข ภิกฺขเว ติวิธา มโนกมฺมนฺตสนฺโทสพฺยาปตฺติ อกุสลสญฺเจตนิกา ทุกฺขุทฺรยา ทุกฺขวิปากา โหติ.}

\prose{ติวิธ กายกมฺมนฺต สนฺโทส พฺยาปตฺติ อกุสล สญฺเจตนิกาเหตุ วา ภิกฺขเว สตฺตา กายสฺส เภทา ปรํ มรณา อปายํ ทุคฺคติํ วินิปาตํ นิรยํ อุปปชฺชนฺติ, จตุพฺพิธวจีกมฺมนฺต ฯเปฯ.\csromanspacer{} ติวิธมโนกมฺมนฺต- สนฺโทสพฺยาปตฺติ-อกุสลสญฺเจตนิกาเหตุ วา ภิกฺขเว สตฺตา กายสฺส เภทา ปรํ มรณา อปายํ ทุคฺคติํ วินิปาตํ นิรยํ อุปปชฺชนฺติ.}

\prose{นาหํ ภิกฺขเว สญฺเจตนิกานํ กมฺมานํ กตานํ อุปจิตานํ อปฺปฏิสํเวทิตฺวา พฺยนฺตีภาวํ วทามิ, ตญฺจ โข ทิฏฺเฐว ธมฺเม อุปปชฺเช วา อปเร วา}

\vspace{\tipitakaparskip}
\csromancenter{(21) 1.\csromanspacer{} กรชกายวคฺค ปริยาเย.\csromanspacer{} น เตฺววาหํ ภิกฺขเว สญฺเจตนิกานํ กมฺมานํ กตานํ อุปจิตานํ อปฺปฏิสํเวทิตฺวา ทุกฺขสฺสนฺตกิริยํ วทามิ.}
\prose{\csromanfolio{503}ตตฺร โข ภิกฺขเว ติวิธา กายกมฺมนฺตสมฺปตฺติ กุสลสญฺเจตนิกา สุขุทฺรยา สุขวิปากา โหติ, จตุพฺพิธา วจีกมฺมนฺตสมฺปตฺติ กุสลสญฺเจตนิกา สุขุทฺรยา สุขวิปากา โหติ, ติวิธา มโนกมฺมนฺตสมฺปตฺติ กุสลสญฺเจตนิกา สุขุทฺรยา สุขวิปากา โหติ.}

\prose{กถญฺจ ภิกฺขเว ติวิธา กายกมฺมนฺตสมฺปตฺติ กุสลสญฺเจตนิกา สุขุทฺรยา สุขวิปากา โหติ ฯเปฯ.\csromanspacer{} เอวํ โข ภิกฺขเว ติวิธา กายกมฺมนฺตสมฺปตฺติ กุสลสญฺเจตนิกา สุขุทฺรยา สุขวิปากา โหติ.}

\prose{กถญฺจ ภิกฺขเว จตุพฺพิธา วจีกมฺมนฺตสมฺปตฺติ กุสลสญฺเจตนิกา สุขุทฺรยา สุขวิปากา โหติ ฯเปฯ.\csromanspacer{} เอวํ โข ภิกฺขเว จตุพฺพิธา วจีกมฺมนฺตสมฺปตฺติ กุสลสญฺเจตนิกา สุขุทฺรยา สุขวิปากา โหติ.}

\csromanmatikaanchor{mtk.245}
\csromantocmarkref{subsection}{8. ทุติยสญฺเจตนิกสุตฺต}{mtk.245}
\bookmark[dest={mtk.245},level=2]{8. ทุติยสญฺเจตนิกสุตฺต}
\prose{กถญฺจ ภิกฺขเว ติวิธา มโนกมฺมนฺตสมฺปตฺติ กุสลสญฺเจตนิกา สุขุทฺรยา สุขวิปากา โหติ ฯเปฯ.\csromanspacer{} เอวํ โข ภิกฺขเว ติวิธา มโนกมฺมนฺตสมฺปตฺติ กุสลสญฺเจตนิกา สุขุทฺรยา สุขวิปากา โหติ.}

\prose{ติวิธกายกมฺมนฺตสมฺปตฺติกุสลสญฺเจตนิกาเหตุ วา ภิกฺขเว สตฺตา กายสฺส เภทา ปรํ มรณา สุคติํ สคฺคํ โลกํ อุปปชฺชนฺติ, จตุพฺพิธวจีกมฺมนฺตสมฺปตฺติ ฯเปฯ.\csromanspacer{} ติวิธมโนกมฺมนฺตสมฺปตฺติกุสลสญฺเจตนิกาเหตุ วา ภิกฺขเว สตฺตา กายสฺส เภทา ปรํ มรณา สุคติํ สคฺคํ โลกํ อุปปชฺชนฺติ ฯเปฯ.\footnote{อุปปชฺชนฺติ. (สฺยา, ก) ตถา สติ “นาหํ ภิกฺขเว สญฺเจตนิกาน” มิจฺจาทินา วุจฺจมานวจเนน สห เอกสุตฺตนฺติ คเหตพฺพํ. เปยฺยาเลน ปน ปุริมสุตฺเต วิย นิคมนํ ทสฺสิตํ.} .\csromanspacer{} อฏฺฐมํ.}

\csromanheader{2}{9. กรชกายสุตฺต}
\proseitem{219}{นาหํ ภิกฺขเว สญฺเจตนิกานํ กมฺมานํ กตานํ อุปจิตานํ อปฺปฏิสํเวทิตฺวา พฺยนฺตีภาวํ วทามิ, ตญฺจ โข ทิฏฺเฐว ธมฺเม อุปปชฺเช วา อปเร วา ปริยาเย.\csromanspacer{} น เตฺววาหํ ภิกฺขเว สญฺเจตนิกานํ กมฺมานํ กตานํ อุปจิตานํ อปฺปฏิสํเวทิตฺวา ทุกฺขสฺสนฺตกิริยํ วทามิ.}

\prose{ส โข โส ภิกฺขเว อริยสาวโก เอวํ วิคตาภิชฺโฌ วิคตพฺยาปาโท อสมฺมูโฬฺห สมฺปชาโน ปฏิสฺสโต เมตฺตาสหคเตน}

\gambhira{ทสกนิปาตปาฬิ}
\prose{\csromanfolio{504}เจตสา เอกํ ทิสํ ผริตฺวา วิหรติ.\csromanspacer{} ตถา ทุติยํ.\csromanspacer{} ตถา ตติยํ.\csromanspacer{} ตถา จตุตฺถํ\footnote{จตุตฺถิํ (?)}.\csromanspacer{} อิติ อุทฺธมโธ ติริยํ สพฺพธิ สพฺพตฺตตาย สพฺพาวนฺตํ โลกํ เมตฺตาสหคเตน เจตสา วิปุเลน มหคฺคเตน อปฺปมาเณน อเวเรน อพฺยาปชฺเชน ผริตฺวา วิหรติ.}

\prose{โส เอวํ ปชานาติ “ปุพฺเพ โข เม อิทํ จิตฺตํ ปริตฺตํ อโหสิ อภาวิตํ, เอตรหิ ปน เม อิทํ จิตฺตํ อปฺปมาณํ สุภาวิตํ.\csromanspacer{} ยํ โข ปน กิญฺจิ ปมาณกตํ กมฺมํ, น ตํ ตตฺราวสิสฺสติ, น ตํ ตตฺราวติฏฺฐตี”ติ.}

\prose{ตํ กิํ มญฺญถ ภิกฺขเว, ทหรตคฺเค เจ โส อยํ\footnote{เจ อยํ (สฺยา)} กุมาโร เมตฺตํ เจโตวิมุตฺติํ ภาเวยฺย, อปิ นุ โข\footnote{อปิ นุ โส (?)} ปาปกมฺมํ กเรยฺยาติ.\csromanspacer{} โน เหตํ ภนฺเต.}

\prose{อกโรนฺตํ โข ปน ปาปกมฺมํ อปิ นุ โข ทุกฺขํ ผุเสยฺยาติ.\csromanspacer{} โน เหตํ ภนฺเต, อกโรนฺตํ หิ ภนฺเต ปาปกมฺมํ กุโต ทุกฺขํ ผุสิสฺสตีติ.}

\prose{ภาเวตพฺพา โข ปนายํ ภิกฺขเว เมตฺตา เจโตวิมุตฺติ อิตฺถิยา วา ปุริเสน วา.\csromanspacer{} อิตฺถิยา วา ภิกฺขเว ปุริสสฺส วา นายํ กาโย อาทาย คมนีโย, จิตฺตนฺตโร อยํ ภิกฺขเว มจฺโจ.\csromanspacer{} โส เอวํ ปชานาติ “ยํ โข เม อิทํ กิญฺจิ ปุพฺเพ อิมินา กรชกาเยน ปาปกมฺมํ กตํ, สพฺพํ ตํ อิธ เวทนียํ, น ตํ อนุคํ ภวิสฺสตี”ติ.\csromanspacer{} เอวํ ภาวิตา โข ภิกฺขเว เมตฺตา เจโตวิมุตฺติ อนาคามิตาย สํวตฺตติ อิธปญฺญสฺส ภิกฺขุโน อุตฺตริ\footnote{อุตฺตริํ (สี, สฺยา, อิ)} วิมุตฺติํ อปฺปฏิวิชฺฌโต.}

\prose{กรุณาสหคเตน เจตสา.\csromanspacer{} มุทิตาสหคเตน เจตสา.\csromanspacer{} อุเปกฺขาสหคเตน เจตสา เอกํ ทิสํ ผริตฺวา วิหรติ.\csromanspacer{} ตถา ทุติยํ.\csromanspacer{} ตถา ตติยํ.\csromanspacer{} ตถา จตุตฺถํ.\csromanspacer{} อิติ อุทฺธมโธ ติริยํ สพฺพธิ สพฺพตฺตตาย สพฺพาวนฺตํ โลกํ อุเปกฺขาสหคเตน เจตสา วิปุเลน มหคฺคเตน อปฺปมาเณน อเวเรน อพฺยาปชฺเชน ผริตฺวา วิหรติ.}

\prose{โส เอวํ ปชานาติ “ปุพฺเพ โข เม อิทํ จิตฺตํ ปริตฺตํ อโหสิ อภาวิตํ, เอตรหิ ปน เม อิทํ จิตฺตํ อปฺปมาณํ สุภาวิตํ.\csromanspacer{} ยํ โข ปน กิญฺจิ ปมาณกตํ กมฺมํ, น ตํ ตตฺราวสิสฺสติ, น ตํ ตตฺราวติฏฺฐตี”ติ.}

\chapterheadpageread{(21) 1. กรชกายวคฺค}
\csromanmatikaanchor{mtk.246}
\csromantocmarkref{subsection}{9. กรชกายสุตฺต}{mtk.246}
\bookmark[dest={mtk.246},level=2]{9. กรชกายสุตฺต}
\prose{\csromanfolio{505}ตํ กิํ มญฺญถ ภิกฺขเว, ทหรตคฺเค เจ โส อยํ กุมาโร อุเปกฺขํ เจโตวิมุตฺติํ ภาเวยฺย, อปิ นุ โข ปาปกมฺมํ กเรยฺยาติ.\csromanspacer{} โน เหตํ ภนฺเต.}

\prose{อกโรนฺตํ โข ปน ปาปกมฺมํ อปิ นุ โข ทุกฺขํ ผุเสยฺยาติ.\csromanspacer{} โน เหตํ ภนฺเต, อกโรนฺตํ หิ ภนฺเต ปาปกมฺมํ กุโต ทุกฺขํ ผุสิสฺสตีติ.}

\prose{ภาเวตพฺพา โข ปนายํ ภิกฺขเว อุเปกฺขา เจโตวิมุตฺติ อิตฺติยา วา ปุริเสน วา.\csromanspacer{} อิตฺถิยา วา ภิกฺขเว ปุริสสฺส วา นายํ กาโย อาทาย คมนีโย, จิตฺตนฺตโร อยํ ภิกฺขเว มจฺโจ.\csromanspacer{} โส เอวํ ปชานาติ “ยํ โข เม อิทํ กิญฺจิ ปุพฺเพ อิมินา กรชกาเยน ปาปกมฺมํ กตํ, สพฺพํ ตํ อิธ เวทนียํ, น ตํ อนุคํ ภวิสฺสตี”ติ.\csromanspacer{} เอวํ ภาวิตา โข ภิกฺขเว อุเปกฺขา เจโตวิมุตฺติ อนาคามิตาย สํวตฺตติ อิธปญฺญสฺส ภิกฺขุโน อุตฺตริ วิมุตฺติํ อปฺปฏิวิชฺฌโตติ.\csromanspacer{}\csromanspacer{}นวมํ.}

\csromanheader{2}{10. อธมฺมจริยาสุตฺต}
\proseitem{220}{\csromansymbolfootnote{*}{อํ 1. 57 ปิฏฺเฐปิ.}อถ โข อญฺญตโร พฺราหฺมโณ เยน ภควา เตนุปสงฺกมิ, อุปสงฺกมิตฺวา ภควตา สทฺธิํ สมฺโมทิ, สมฺโมทนียํ กถํ สารณียํ วีติสาเรตฺวา เอกมนฺตํ นิสีทิ, เอกมนฺตํ นิสินฺโน โข โส พฺราหฺมโณ ภควนฺตํ เอตทโวจ “โก นุ โข โภ โคตม เหตุ โก ปจฺจโย, เยนมิเธกจฺเจ สตฺตา กายสฺส เภทา ปรํ มรณา อปายํ ทุคฺคติํ วินิปาตํ นิรยํ อุปปชฺชนฺตี”ติ.\csromanspacer{} อธมฺมจริยาวิสม- จริยาเหตุ โข พฺราหฺมณ เอวมิเธกจฺเจ สตฺตา กายสฺส เภทา ปรํ มรณา อปายํ ทุคฺคติํ วินิปาตํ นิรยํ อุปปชฺชนฺตีติ.}

\prose{โก ปน โภ โคตม เหตุ โก ปจฺจโย, เยนมิเธกจฺเจ สตฺตา กายสฺส เภทา ปรํ มรณา สุคติํ สคฺคํ โลกํ อุปปชฺชนฺตีติ.\csromanspacer{} ธมฺมจริยา- สมจริยาเหตุ โข พฺราหฺมณ เอวมิเธกจฺเจ สตฺตา กายสฺส เภทา ปรํ มรณา สุคติํ สคฺคํ โลกํ อุปปชฺชนฺตีติ.}

\prose{น โข อหํ อิมสฺส โภโต โคตมสฺส สํขิตฺเตน ภาสิตสฺส วิตฺถาเรน อตฺถํ อาชานามิ, สาธุ เม ภวํ โคตโม ตถา ธมฺมํ เทเสตุ, ยถาหํ อิมสฺส โภโต โคตมสฺส สํขิตฺเตน ภาสิตสฺส วิตฺถาเรน อตฺถํ อาชาเนยฺยนฺติ.\csromanspacer{} เตน หิ พฺราหฺมณ สุณาหิ สาธุกํ มนสิ}

\gambhira{ทสกนิปาตปาฬิ}
\prose{\csromanfolio{506}กโรหิ, ภาสิสฺสามีติ. “เอวํ โภ”ติ โข โส พฺราหฺมโณ ภควโต ปจฺจสฺโสสิ.\csromanspacer{} ภควา เอตทโวจ\csromandash{}}

\prose{ติวิธา โข พฺราหฺมณ กาเยน อธมฺมจริยาวิสมจริยา โหติ, จตุพฺพิธา วาจาย อธมฺมจริยาวิสมจริยา โหติ, ติวิธา มนสา อธมฺมจริยา วิสมจริยา โหติ.}

\prose{กถญฺจ พฺราหฺมณ ติวิธา กาเยน อธมฺมจริยาวิสมจริยา โหติ ฯเปฯ.\csromanspacer{} เอวํ พฺราหฺมณ ติวิธา กาเยน อธมฺมจริยา วิสมจริยา โหติ.}

\prose{กถญฺจ พฺราหฺมณ จตุพฺพิธา วาจาย อธมฺมจริยาวิสมจริยา โหติ ฯเปฯ.\csromanspacer{} เอวํ โข พฺราหฺมณ จตุพฺพิธา วาจาย อธมฺมจริยา วิสมจริยา โหติ.}

\prose{กถญฺจ พฺราหฺมณ ติวิธา มนสา อธมฺมจริยาวิสมจริยา โหติ ฯเปฯ.\csromanspacer{} เอวํ โข พฺราหฺมณ ติวิธา มนสา อธมฺมจริยาวิสมจริยา โหติ.\csromanspacer{} เอวํ อธมฺมจริยาวิสมจริยาเหตุ โข พฺราหฺมณ เอวมิเธกจฺเจ สตฺตา กายสฺส เภทา ปรํ มรณา อปายํ ทุคฺคติํ วินิปาตํ นิรยํ อุปปชฺชนฺติ.}

\prose{ติวิธา พฺราหฺมณ กาเยน ธมฺมจริยาสมจริยา โหติ, จตุพฺพิธา วาจาย ธมฺมจริยาสมจริยา โหติ, ติวิธา มนสา ธมฺมจริยาสมจริยา โหติ.}

\prose{กถญฺจ พฺราหฺมณ ติวิธา กาเยน ธมฺมจริยาสมจริยา โหติ ฯเปฯ.\csromanspacer{} เอวํ โข พฺราหฺมณ ติวิธา กาเยน ธมฺมจริยาสมจริยา โหติ.}

\csromanmatikaanchor{mtk.247}
\csromantocmarkref{subsection}{10. อธมฺมจริยาสุตฺต}{mtk.247}
\bookmark[dest={mtk.247},level=2]{10. อธมฺมจริยาสุตฺต}
\prose{กถญฺจ พฺราหฺมณ จตุพฺพิธา วาจาย ธมฺมจริยาสมจริยา โหติ ฯเปฯ.\csromanspacer{} เอวํ โข พฺราหฺมณ จตุพฺพิธา วาจาย ธมฺมจริยาสมจริยา โหติ.}

\prose{กถญฺจ พฺราหฺมณ ติวิธา มนสา ธมฺมจริยาสมจริยา โหติ ฯเปฯ.\csromanspacer{} เอวํ โข พฺราหฺมณ ติวิธา มนสา ธมฺมจริยาสมจริยา โหติ.\csromanspacer{} เอวํ ธมฺมจริยาสมจริยาเหตุ โข พฺราหฺมณ เอวมิเธกจฺเจ สตฺตา กายสฺส เภทา ปรํ มรณา สุคติํ สคฺคํ โลกํ อุปปชฺชนฺตีติ.}

\prose{อภิกฺกนฺตํ โภ โคตม, อภิกฺกนฺตํ โภ โคตม ฯเปฯ อุปาสกํ มํ ภวํ โคตโม ธาเรตุ อชฺชตคฺเค ปาณุเปตํ สรณํ คตนฺติ.\csromanspacer{}\csromanspacer{}ทสมํ.}

\vspace{\tipitakaparskip}
\csromancenter{กรชกายวคฺโค ปฐโม.}
\csromanheader{1}{\textbf{(22) 2. สามญฺญวคฺค} \textbf{(22) 2. สามญฺญวคฺค}}
\proseitem{221}{\csromanfolio{507}ทสหิ ภิกฺขเว ธมฺเมหิ สมนฺนาคโต ยถาภตํ นิกฺขิตฺโต เอวํ นิรเย.\csromanspacer{} กตเมหิ ทสหิ? ปาณาติปาตี โหติ, อทินฺนาทายี โหติ, กาเมสุมิจฺฉาจารี โหติ, มุสาวาที โหติ, ปิสุณวาโจ โหติ, ผรุสวาโจ โหติ, สมฺผปฺปลาปี โหติ, อภิชฺฌาลุ โหติ, พฺยาปนฺนจิตฺโต โหติ, มิจฺฉาทิฏฺฐิโก โหติ.\csromanspacer{} อิเมหิ โข ภิกฺขเว ทสหิ ธมฺเมหิ สมนฺนาคโต ยถาภตํ นิกฺขิตฺโต เอวํ นิรเย.}

\prose{ทสหิ ภิกฺขเว ธมฺเมหิ สมนฺนาคโต ยถาภตํ นิกฺขิตฺโต เอวํ สคฺเค.\csromanspacer{} กตเมหิ ทสหิ? ปาณาติปาตา ปฏิวิรโต โหติ, อทินฺนาทานา ปฏิวิรโต โหติ, กาเมสุมิจฺฉาจารา ปฏิวิรโต โหติ, มุสาวาทา ปฏิวิรโต โหติ, ปิสุณาย วาจาย ปฏิวิรโต โหติ, ผรุสาย วาจาย ปฏิวิรโต โหติ, สมฺผปฺปลาปา ปฏิวิรโต โหติ, อนภิชฺฌาลุ โหติ, อพฺยาปนฺนจิตฺโต โหติ, สมฺมาทิฏฺฐิโก โหติ.\csromanspacer{} อิเมหิ โข ภิกฺขเว ทสหิ ธมฺเมหิ สมนฺนาคโต ยถาภตํ นิกฺขิตฺโต เอวํ สคฺเคติ. (1)}

\proseitem{222}{วีสติยา ภิกฺขเว ธมฺเมหิ สมนฺนาคโต ยถาภตํ นิกฺขิตฺโต เอวํ นิรเย.\csromanspacer{} กตเมหิ วีสติยา? อตฺตนา จ ปาณาติปาตี โหติ, ปรญฺจ ปาณาติปาเต สมาทเปติ.\csromanspacer{} อตฺตนา จ อทินฺนาทายี โหติ, ปรญฺจ อทินฺนาทาเน สมาทเปติ.\csromanspacer{} อตฺตนา จ กาเมสุมิจฺฉาจารี โหติ, ปรญฺจ กาเมสุมิจฺฉาจาเร สมาทเปติ.\csromanspacer{} อตฺตนา จ มุสาวาที โหติ, ปรญฺจ มุสาวาเท สมาทเปติ.\csromanspacer{} อตฺตนา จ ปิสุณวาโจ โหติ, ปรญฺจ ปิสุณาย วาจาย สมาทเปติ.\csromanspacer{} อตฺตนา จ ผรุสวาโจ โหติ, ปรญฺจ ผรุสาย วาจาย สมาทเปติ.\csromanspacer{} อตฺตนา จ สมฺผปฺปลาปี โหติ, ปรญฺจ สมฺผปฺปลาเป สมาทเปติ.\csromanspacer{} อตฺตนา จ อภิชฺฌาลุ โหติ, ปรญฺจ อภิชฺฌาย สมาทเปติ.\csromanspacer{} อตฺตนา จ พฺยาปนฺนจิตฺโต โหติ, ปรญฺจ พฺยาปาเท สมาทเปติ.\csromanspacer{} อตฺตนา จ มิจฺฉาทิฏฺฐิโก โหติ, ปรญฺจ มิจฺฉาทิฏฺฐิยา สมาทเปติ.\csromanspacer{} อิเมหิ โข ภิกฺขเว วีสติยา ธมฺเมหิ สมนฺนาคโต ยถาภตํ นิกฺขิตฺโต เอวํ นิรเย.}

\prose{วีสติยา ภิกฺขเว ธมฺเมหิ สมนฺนาคโต ยถาภตํ นิกฺขิตฺโต เอวํ สคฺเค.\csromanspacer{} กตเมหิ วีสติยา? อตฺตนา จ ปาณาติปาตา ปฏิวิรโต}

\gambhira{ทสกนิปาตปาฬิ}
\prose{\csromanfolio{508}โหติ, ปรญฺจ ปาณาติปาตา เวรมณิยา สมาทเปติ.\csromanspacer{} อตฺตนา จ อทินฺนาทานา ปฏิวิรโต โหติ.\csromanspacer{} ปรญฺจ อทินฺนาทานา เวรมณิยา สมาทเปติ.\csromanspacer{} อตฺตนา จ กาเมสุมิจฺฉาจารา ปฏิวิรโต โหติ, ปรญฺจ กาเมสุมิจฺฉาจารา เวรมณิยา สมาทเปติ.\csromanspacer{} อตฺตนา จ มุสาวาทา ปฏิวิรโต โหติ, ปรญฺจ มุสาวาทา เวรมณิยา สมาทเปติ.\csromanspacer{} อตฺตนา จ ปิสุณาย วาจาย ปฏิวิรโต โหติ, ปรญฺจ ปิสุณาย วาจาย เวรมณิยา สมาทเปติ.\csromanspacer{} อตฺตนา จ ผรุสาย วาจาย ปฏิวิรโต โหติ, ปรญฺจ ผรุสาย วาจาย เวรมณิยา สมาทเปติ.\csromanspacer{} อตฺตนา จ สมฺผปฺปลาปา ปฏิวิรโต โหติ, ปรญฺจ สมฺผปฺปลาปา เวรมณิยา สมาทเปติ.\csromanspacer{} อตฺตนา จ อนภิชฺฌาลุ โหติ, ปรญฺจ อนภิชฺฌาย สมาทเปติ.\csromanspacer{} อตฺตนา จ อพฺยาปนฺนจิตฺโต โหติ, ปรญฺจ อพฺยาปาเท สมาทเปติ.\csromanspacer{} อตฺตนา จ สมฺมาทิฏฺฐิโก โหติ, ปรญฺจ สมาทิฏฺฐิยา สมาทเปติ.\csromanspacer{} อิเมหิ โข ภิกฺขเว วีสติยา ธมฺเมหิ สมนฺนาคโต ยถาภตํ นิกฺขิตฺโต เอวํ สคฺเคติ. (2)}

\proseitem{223}{ติํสาย ภิกฺขเว ธมฺเมหิ สมนฺนาคโต ยถาภตํ นิกฺขิตฺโต เอวํ นิรเย.\csromanspacer{} กตเมหิ ติํสาย? อตฺตนา จ ปาณาติปาตี โหติ, ปรญฺจ ปาณาติปาเต สมาทเปติ, ปาณาติปาเต จ สมนุญฺโญ โหติ.\csromanspacer{} อตฺตนา จ อทินฺนาทายี โหติ, ปรญฺจ อทินฺนาทาเน สมาทเปติ, อทินฺนาทาเน จ สมนุญฺโญ โหติ.\csromanspacer{} อตฺตนา จ กาเมสุมิจฺฉาจารี โหติ, ปรญฺจ กาเมสุมิจฺฉาจาเร สมาทเปติ, กาเมสุมิจฺฉาจาเร จ สมนุญฺโญ โหติ.\csromanspacer{} อตฺตนา จ มุสาวาที โหติ, ปรญฺจ มุสาวาเท สมาทเปติ, มุสาวาเท จ สมนุญฺโญ โหติ.\csromanspacer{} อตฺตนา จ ปิสุณวาโจ โหติ, ปรญฺจ ปิสุณาย วาจาย สมาทเปติ, ปิสุณาย วาจาย จ สมนุญฺโญ โหติ.\csromanspacer{} อตฺตนา จ ผรุสวาโจ โหติ, ปรญฺจ ผรุสาย วาจาย สมาทเปติ, ผรุสาย วาจาย จ สมนุญฺโญ โหติ.\csromanspacer{} อตฺตนา จ สมฺผปฺปลาปี โหติ, ปรญฺจ สมฺผปฺปลาเป สมาทเปติ, สมฺผปฺปลาเป จ สมนุญฺโญ โหติ.\csromanspacer{} อตฺตนา จ อภิชฺฌาลุ โหติ, ปรญฺจ อภิชฺฌาย สมาทเปติ, อภิชฺฌาย จ สมนุญฺโญ โหติ.\csromanspacer{} อตฺตนา จ พฺยาปนฺนจิตฺโต โหติ, ปรญฺจ พฺยาปาเท สมาทเปติ, พฺยาปาเท จ สมนุญฺโญ โหติ.\csromanspacer{} อตฺตนา จ มิจฺฉาทิฏฺฐิโก โหติ, ปรญฺจ มิจฺฉาทิฏฺฐิยา สมาทเปติ, มิจฺฉาทิฏฺฐิยา จ สมนุญฺโญ โหติ.\csromanspacer{} อิเมหิ}

\vspace{\tipitakaparskip}
\csromancenter{(22) 2.\csromanspacer{} สามญฺญวคฺค โข ภิกฺขเว ติํสาย ธมฺเมหิ สมนฺนาคโต ยถาภตํ นิกฺขิตฺโต เอวํ นิรเย.}
\prose{\csromanfolio{509}ติํสาย ภิกฺขเว ธมฺเมหิ สมนฺนาคโต ยถาภตํ นิกฺขิตฺโต เอวํ สคฺเค.\csromanspacer{} กตเมหิ ติํสาย? อตฺตนา จ ปาณาติปาตา ปฏิวิรโต โหติ, ปรญฺจ ปาณาติปาตา เวรมณิยา สมาทเปติ, ปาณาติปาตา เวรมณิยา จ สมนุญฺโญ โหติ.\csromanspacer{} อตฺตนา จ อทินฺนาทานา ปฏิวิรโต โหติ, ปรญฺจ อทินฺนาทานา เวรมณิยา สมาทเปติ, อทินฺนาทานา เวรมณิยา จ สมนุญฺโญ โหติ.\csromanspacer{} อตฺตนา จ กาเมสุมิจฺฉาจารา ปฏิวิรโต โหติ, ปรญฺจ กาเมสุมิจฺฉาจารา เวรมณิยา สมาทเปติ, กาเมสุมิจฺฉาจารา เวรมณิยา จ สมนุญฺโญ โหติ.\csromanspacer{} อตฺตนา จ มุสาวาทา ปฏิวิรโต โหติ, ปรญฺจ มุสาวาทา เวรมณิยา สมาทเปติ, มุสาวาทา เวรมณิยา จ สมนุญฺโญ โหติ.\csromanspacer{} อตฺตนา จ ปิสุณาย วาจาย ปฏิวิรโต โหติ, ปรญฺจ ปิสุณาย วาจาย เวรมณิยา สมาทเปติ, ปิสุณาย วาจาย เวรมณิยา จ สมนุญฺโญ โหติ.\csromanspacer{} อตฺตนา จ ผรุสาย วาจาย ปฏิวิรโต โหติ, ปรญฺจ ผรุสาย วาจาย เวรมณิยา สมาทเปติ, ผรุสาย วาจาย เวรมณิยา จ สมนุญฺโญ โหติ.\csromanspacer{} อตฺตนา จ สมฺผปฺปลาปา ปฏิวิรโต โหติ, ปรญฺจ สมฺผปฺปลาปา เวรมณิยา สมาทเปติ, สมฺผปฺปลาปา เวรมณิยา จ สมนุญฺโญ โหติ.\csromanspacer{} อตฺตนา จ อนภิชฺฌาลุ โหติ, ปรญฺจ อนภิชฺฌาย สมาทเปติ, อนภิชฺฌาย จ สมนุญฺโญ โหติ, อตฺตนา จ อพฺยาปนฺนจิตฺโต โหติ, ปรญฺจ อพฺยาปาเท สมาทเปติ, อพฺยาปาเท จ สมนุญฺโญ โหติ.\csromanspacer{} อตฺตนา จ สมฺมาทิฏฺฐิโก โหติ, ปรญฺจ สมฺมาทิฏฺฐิยา สมาทเปติ, สมฺมาทิฏฺฐิยา จ สมนุญฺโญ โหติ.\csromanspacer{} อิเมหิ โข ภิกฺขเว ติํสาย ธมฺเมหิ สมนฺนาคโต ยถาภตํ นิกฺขิตฺโต เอวํ สคฺเคติ. (3)}

\proseitem{224}{จตฺตารีสาย ภิกฺขเว ธมฺเมหิ สมนฺนาคโต ยถาภตํ นิกฺขิตฺโต เอวํ นิรเย.\csromanspacer{} กตเมหิ จตฺตารีสาย? อตฺตนา จ ปาณาติปาตี โหติ, ปรญฺจ ปาณาติปาเต สมาทเปติ, ปาณาติปาเต จ สมนุญฺโญ โหติ, ปาณาติปาตสฺส จ วณฺณํ ภาสติ.\csromanspacer{} อตฺตนา จ อทินฺนาทายี โหติ, ปรญฺจ อทินฺนาทาเน สมาทเปติ, อทินฺนาทาเน จ สมนุญฺโญ โหติ, อทินฺนาทานสฺส จ วณฺณํ ภาสติ.\csromanspacer{} อตฺตนา จ กาเมสุมิจฺฉาจารี}

\gambhira{ทสกนิปาตปาฬิ}
\csromanmatikaanchor{mtk.248}
\csromantocmarkref{section}{(22) 2. สามญฺญวคฺค}{mtk.248}
\bookmark[dest={mtk.248},level=1]{(22) 2. สามญฺญวคฺค}
\prose{\csromanfolio{510}โหติ, ปรญฺจ กาเมสุมิจฺฉาจาเร สมาทเปติ, กาเมสุมิจฺฉาจาเร จ สมนุญฺโญ โหติ, กาเมสุมิจฺฉาจารสฺส จ วณฺณํ ภาสติ.\csromanspacer{} อตฺตนา จ มุสาวาที โหติ, ปรญฺจ มุสาวาเท สมาทเปติ, มุสาวาเท จ สมนุญฺโญ โหติ, มุสาวาทสฺส จ วณฺณํ ภาสติ.\csromanspacer{} อตฺตนา จ ปิสุณวาโจ โหติ, ปรญฺจ ปิสุณาย วาจาย สมาทเปติ, ปิสุณาย วาจาย จ สมนุญฺโญ โหติ, ปิสุณาย วาจาย จ วณฺณํ ภาสติ.\csromanspacer{} อตฺตนา จ ผรุสวาโจ โหติ, ปรญฺจ ผรุสาย วาจาย สมาทเปติ, ผรุสาย วาจาย จ สมนุญฺโญ โหติ, ผรุสาย วาจาย จ วณฺณํ ภาสติ.\csromanspacer{} อตฺตนา จ สมฺผปฺปลาปี โหติ, ปรญฺจ สมฺผปฺปลาเป สมาทเปติ, สมฺผปฺปลาเป จ สมนุญฺโญ โหติ, สมฺผปฺปลาปาย จ วณฺณํ ภาสติ.\csromanspacer{} อตฺตนา จ อภิชฺฌาลุ โหติ, ปรญฺจ อภิชฺฌาย สมาทเปติ, อภิชฺฌาย จ สมนุญฺโญ โหติ, อภิชฺฌาย จ วณฺณํ ภาสติ.\csromanspacer{} อตฺตนา จ พฺยาปนฺนจิตฺโต โหติ, ปรญฺจ พฺยาปาเท สมาทเปติ, พฺยาปาเท จ สมนุญฺโญ โหติ, พฺยาปาทสฺส จ วณฺณํ ภาสติ.\csromanspacer{} อตฺตนา จ มิจฺฉาทิฏฺฐิโก โหติ, ปรญฺจ มิจฺฉาทิฏฺฐิยา สมาทเปติ, มิจฺฉาทิฏฺฐิยา จ สมนุญฺโญ โหติ, มิจฺฉาทิฏฺฐิยา จ วณฺณํ ภาสติ.\csromanspacer{} อิเมหิ โข ภิกฺขเว จตฺตารีสาย ธมฺเมหิ สมนฺนาคโต ยถาภตํ นิกฺขิตฺโต เอวํ นิรเย.}

\prose{จตฺตารีสาย ภิกฺขเว ธมฺเมหิ สมนฺนาคโต ยถาภตํ นิกฺขิตฺโต เอวํ สคฺเค.\csromanspacer{} กตเมหิ จตฺตารีสาย? อตฺตนา จ ปาณาติปาตา ปฏิวิรโต โหติ, ปรญฺจ ปาณาติปาตา เวรมณิยา สมาทเปติ, ปาณาติปาตา เวรมณิยา จ สมนุญฺโญ โหติ, ปาณาติปาตา เวรมณิยา จ วณฺณํ ภาสติ.\csromanspacer{} อตฺตนา จ อทินฺนาทานา ปฏิวิรโต โหติ, ปรญฺจ อทินฺนาทานา เวรมณิยา สมาทเปติ, อทินฺนาทานา เวรมณิยา จ สมนุญฺโญ โหติ, อทินฺนาทานา เวรมณิยา จ วณฺณํ ภาสติ.\csromanspacer{} อตฺตนา จ กาเมสุมิจฺฉาจารา ปฏิวิรโต โหติ, ปรญฺจ กาเมสุมิจฺฉาจารา เวรมณิยา สมาทเปติ, กาเมสุมิจฺฉาจารา เวรมณิยา จ สมนุญฺโญ โหติ, กาเมสุมิจฺฉาจารา เวรมณิยา จ วณฺณํ ภาสติ.\csromanspacer{} อตฺตนา จ มุสาวาทา ปฏิวิรโต โหติ, ปรญฺจ มุสาวาทา เวรมณิยา สมาทเปติ, มุสาวาทา เวรมณิยา จ สมนุญฺโญ โหติ, มุสาวาทา เวรมณิยา จ วณฺณํ ภาสติ.\csromanspacer{} อตฺตนา จ ปิสุณาย วาจาย ปฏิวิรโต โหติ, ปรญฺจ ปิสุณาย วาจาย}

\vspace{\tipitakaparskip}
\csromancenter{(22) 2.\csromanspacer{} สามญฺญวคฺค เวรมณิยา สมาทเปติ, ปิสุณาย วาจาย เวรมณิยา จ สมนุญฺโญ โหติ, ปิสุณาย วาจาย เวรมณิยา จ วณฺณํ ภาสติ.\csromanspacer{} อตฺตนา จ ผรุสาย วาจาย ปฏิวิรโต โหติ, ปรญฺจ ผรุสาย วาจาย เวรมณิยา จ สมาทเปติ, ผรุสาย วาจาย เวรมณิยา จ สมนุญฺโญ โหติ, ผรุสาย วาจาย เวรมณิยา จ วณฺณํ ภาสติ.\csromanspacer{} อตฺตนาจ สมฺผปฺปลาปา ปฏิวิรโต โหติ, ปรญฺจ สมฺผปฺปลาปา เวรมณิยา สมาทเปติ, สมฺผปฺปลาปา เวรมณิยา จ สมนุญฺโญ โหติ, สมฺผปฺปลาปา เวรมณิยา จ วณฺณํ ภาสติ.\csromanspacer{} อตฺตนา จ อนภิชฺฌาลุ โหติ, ปรญฺจ อนภิชฺฌาย สมาทเปติ, อนภิชฺฌาย จ สมนุญฺโญ โหติ, อนภิชฺฌาย จ วณฺณํ ภาสติ.\csromanspacer{} อตฺตนา จ อพฺยาปนฺนจิตฺโต โหติ, ปรญฺจ อพฺยาปาเท สมาทเปติ, อพฺยาปาเท จ สมนุญฺโญ โหติ, อพฺยาปาทสฺส จ วณฺณํ ภาสติ.\csromanspacer{} อตฺตนา จ สมฺมาทิฏฺฐิโก โหติ, ปรญฺจ สมฺมาทิฏฺฐิยา สมาทเปติ, สมฺมาทิฏฺฐิยา จ สมนุญฺโญ โหติ, สมฺมาทิฏฺฐิยา จ วณฺณํ ภาสติ.\csromanspacer{} อิเมหิ โข ภิกฺขเว จตฺตารีสาย ธมฺเมหิ สมนฺนาคโต ยถาภตํ นิกฺขิตฺโต เอวํ สคฺเคติ. (4)}
\proseitem{225-228}{\csromanfolio{511}ทสหิ ภิกฺขเว ธมฺเมหิ สมนฺนาคโต ขตํ อุปหตํ อตฺตานํ ปริหรติ ฯเปฯ.\csromanspacer{} อกฺขตํ อนุปหตํ อตฺตานํ ปริหรติ ฯเปฯ.\csromanspacer{} วีสติยา ภิกฺขเว ฯเปฯ.\csromanspacer{} ติํสาย ภิกฺขเว ฯเปฯ.\csromanspacer{} จตฺตารีสาย ภิกฺขเว ธมฺเมหิ สมนฺนาคโต ขตํ อุปหตํ อตฺตานํ ปริหรติ ฯเปฯ. (5-8)}

\proseitem{229-232}{ทสหิ ภิกฺขเว ธมฺเมหิ สมนฺนาคโต อิเธกจฺโจ กายสฺส เภทา ปรํ มรณา อปายํ ทุคฺคติํ วินิปาตํ นิรยํ อุปปชฺชติ ฯเปฯ อิเธกจฺโจ กายสฺส เภทา ปรํ มรณา สุคติํ สคฺคํ โลกํ อุปปชฺชติ.\csromanspacer{} วีสติยา ภิกฺขเว ฯเปฯ.\csromanspacer{} ติํสาย ภิกฺขเว ฯเปฯ.\csromanspacer{} จตฺตารีสาย ภิกฺขเว ธมฺเมหิ สมนฺนาคโต อิเธกจฺโจ กายสฺส เภทา ปรํ มรณา อปายํ ทุคฺคติํ วินิปาตํ นิรยํ อุปปชฺชติ ฯเปฯ อิเธกจฺโจ กายสฺส เภทา ปรํ มรณา สุคติํ สคฺคํ โลกํ อุปปชฺชติ. (9-12)}

\proseitem{233-236}{ทสหิ ภิกฺขเว ธมฺเมหิ สมนฺนาคโต พาโล เวทิตพฺโพ ฯเปฯ.\csromanspacer{} ปณฺฑิโต เวทิตพฺโพ ฯเปฯ.\csromanspacer{} วีสติยา ภิกฺขเว ฯเปฯ.\csromanspacer{} ติํสาย ภิกฺขเว ฯเปฯ.\csromanspacer{} จตฺตารีสาย ภิกฺขเว ธมฺเมหิ สมนฺนาคโต พาโล}

\gambhira{ทสกนิปาตปาฬิ}
\prose{\csromanfolio{512}เวทิตพฺโพ ฯเปฯ.\csromanspacer{} ปณฺฑิโต เวทิตพฺโพ ฯเปฯ.\csromanspacer{} อิเมหิ โข ภิกฺขเว จตฺตารีสาย ธมฺเมหิ สมนฺนาคโต ปณฺฑิโต เวทิตพฺโพติ. (13-16)}

\csromanheader{2}{สามญฺญวคฺโค ทุติโย.}
\csromansectionrule
\csromanheader{1}{23. ราคเปยฺยาล}
\proseitem{237}{ราคสฺส ภิกฺขเว อภิญฺญาย ทส ธมฺมา ภาเวตพฺพา.\csromanspacer{} กตเม ทส? อสุภสญฺญา มรณสญฺญา อาหาเร ปฏิกูลสญฺญา สพฺพโลเก อนภิรตสญฺญา อนิจฺจสญฺญา อนิจฺเจ ทุกฺขสญฺญา ทุกฺเข อนตฺตสญฺญา ปหานสญฺญา วิราคสญฺญา นิโรธสญฺญา.\csromanspacer{} ราคสฺส ภิกฺขเว อภิญฺญาย อิเม ทส ธมฺมา ภาเวตพฺพาติ. (1)}

\proseitem{238}{ราคสฺส ภิกฺขเว อภิญฺญาย ทส ธมฺมา ภาเวตพฺพา.\csromanspacer{} กตเม ทส? อนิจฺจสญฺญา อนตฺตสญฺญา อาหาเร ปฏิกูลสญฺญา สพฺพโลเก อนภิรตสญฺญา อฏฺฐิกสญฺญา ปุฬวกสญฺญา\footnote{ปุลวกสญฺญา (สี) ปุฬุวกสญฺญา (ก)} วินีลกสญฺญา วิปุพฺพกสญฺญา วิจฺฉิทฺทกสญฺญา อุทฺธุมาตกสญฺญา.\csromanspacer{} ราคสฺส ภิกฺขเว อภิญฺญาย อิเม ทส ธมฺมา ภาเวตพฺพาติ. (2)}

\proseitem{239}{ราคสฺส ภิกฺขเว อภิญฺญาย ทส ธมฺมา ภาเวตพฺพา.\csromanspacer{} กตเม ทส? สมฺมาทิฏฺฐิ สมฺมาสงฺกปฺโป สมฺมาวาจา สมฺมากมฺมนฺโต สมฺมา-อาชีโว สมฺมาวายาโม สมฺมาสติ สมฺมาสมาธิ สมฺมาญาณํ สมฺมาวิมุตฺติ.\csromanspacer{} ราคสฺส ภิกฺขเว อภิญฺญาย อิเม ทส ธมฺมา ภาเวตพฺพาติ. (3)}

\proseitem{240-266}{ราคสฺส ภิกฺขเว ปริญฺญาย ฯเปฯ.\csromanspacer{} ปริกฺขยาย.\csromanspacer{} ปหานาย.\csromanspacer{} ขยาย.\csromanspacer{} วยาย.\csromanspacer{} วิราคาย.\csromanspacer{} นิโรธาย. ( )\footnote{(อุปสมาย) (สี, สฺยา, อิ) อญฺเญสํ ปน นิปาตานํ ปริโยสาเน อิทํ ปทํ น ทิสฺสติ.} จาคาย.\csromanspacer{} ปฏินิสฺสคฺคาย ฯเปฯ.\csromanspacer{} อิเม ทส ธมฺมา ภาเวตพฺพา. (4-30)}

\csromanheader{2}{23. ราคเปยฺยาล}
\proseitem{267-746}{\csromanfolio{513}โทสสฺส ฯเปฯ.\csromanspacer{} โมหสฺส.\csromanspacer{} โกธสฺส.\csromanspacer{} อุปนาหสฺส.\csromanspacer{} มกฺขสฺส.\csromanspacer{} ปฬาสสฺส.\csromanspacer{} อิสฺสาย.\csromanspacer{} มจฺฉริยสฺส.\csromanspacer{} มายาย.\csromanspacer{} สาเฐยฺยสฺส.\csromanspacer{} ถมฺภสฺส.\csromanspacer{} สารมฺภสฺส.\csromanspacer{} มานสฺส.\csromanspacer{} อติมานสฺส.\csromanspacer{} มทสฺส.\csromanspacer{} ปมาทสฺส ปริญฺญาย ฯเปฯ ปริกฺขยาย.\csromanspacer{} ปหานาย.\csromanspacer{} ขยาย.\csromanspacer{} วยาย.\csromanspacer{} วิราคาย.\csromanspacer{} นิโรธาย. ( )\footnote{(อุปสมาย) (สี, สฺยา, อิ) อญฺเญสํ ปน นิปาตานํ ปริโยสาเน อิทํ ปทํ น ทิสฺสติ.} จาคาย.\csromanspacer{} ปฏินิสฺสคฺคาย ฯเปฯ.\csromanspacer{} อิเม ทส ธมฺมา ภาเวตพฺพาติ. (31-510)}

\nitthitam{ราคเปยฺยาลํ นิฏฺฐิตํ.}
\nitthitam{\textbf{ทสกนิปาตปาฬิ นิฏฺฐิตา.}}
\csromanheader{2}{\textbf{องฺคุตฺตรนิกาย}}
\gambhira{\textbf{เอกาทสกนิปาตปาฬิ}}
\namakkaram{นโม ตสฺส ภควโต อรหโต สมฺมาสมฺพุทฺธสฺส.}
\csromanmatikaanchor{mtk.249}
\csromantocmarkref{section}{(23) ราคเปยฺยาล}{mtk.249}
\bookmark[dest={mtk.249},level=1]{(23) ราคเปยฺยาล}
\csromantocmarknopage{chapter}{เอกาทสกนิปาตปาฬิ}{mtk.250}
\bookmark[dest={mtk.250},level=0]{เอกาทสกนิปาตปาฬิ}
\chapterhead{1. นิสฺสยวคฺค}
\csromanheader{2}{1. กิมตฺถิยสุตฺต}
\proseitem{1}{\csromanfolio{515}เอวํ เม สุตํ\csromandash{}เอกํ สมยํ ภควา สาวตฺถิยํ วิหรติ เชตวเน อนาถปิณฺฑิกสฺส อาราเม.\csromanspacer{} อถ โข อายสฺมา อานนฺโท เยน ภควา เตนุปสงฺกมิ, อุปสงฺกมิตฺวา ภควนฺตํ อภิวาเทตฺวา เอกมนฺตํ นิสีทิ, เอกมนฺตํ นิสินฺโน โข อายสฺมา อานนฺโท ภควนฺตํ เอตทโวจ “กิมตฺถิยานิ ภนฺเต กุสลานิ สีลานิ กิมานิสํสานี”ติ.\csromanspacer{} อวิปฺปฏิสารตฺถานิ โข อานนฺท กุสลานิ สีลานิ อวิปฺปฏิสารานิสํสานีติ.}

\prose{อวิปฺปฏิสาโร ปน ภนฺเต กิมตฺถิโย กิมานิสํโส.\csromanspacer{} อวิปฺปฏิสาโร โข อานนฺท ปาโมชฺชตฺโถ ปาโมชฺชานิสํโส.}

\prose{ปาโมชฺชํ ปน ภนฺเต กิมตฺถิยํ กิมานิสํสํ.\csromanspacer{} ปาโมชฺชํ โข อานนฺท ปีตตฺถํ ปีตานิสํสํ.}

\prose{ปีติ ปน ภนฺเต กิมตฺถิยา กิมานิสํสา.\csromanspacer{} ปีติ โข อานนฺท ปสฺสทฺธตฺถา ปสฺสทฺธานิสํสา.}

\prose{ปสฺสทฺธิ ปน ภนฺเต กิมตฺถิยา กิมานิสํสา.\csromanspacer{} ปสฺสทฺธิ โข อานนฺท สุขตฺถา สุขานิสํสา.}

\prose{สุขํ ปน ภนฺเต กิมตฺถิยํ กิมานิสํสํ.\csromanspacer{} สุขํ โข อานนฺท สมาธตฺถํ สมาธานิสํสํ.}

\gambhira{เอกาทสกนิปาตปาฬิ}
\prose{\csromanfolio{516}สมาธิ ปน ภนฺเต กิมตฺถิโย กิมานิสํโส.\csromanspacer{} สมาธิ โข อานนฺท ยถาภูตญาณทสฺสนตฺโถ ยถาภูตญาณทสฺสนานิสํโส.}

\prose{ยถาภูตญาณทสฺสนํ ปน ภนฺเต กิมตฺถิยํ กิมานิสํสํ.\csromanspacer{} ยถาภูตญาณทสฺสนํ โข อานนฺท นิพฺพิทตฺถํ นิพฺพิทานิสํสํ.}

\prose{นิพฺพิทา ปน ภนฺเต กิมตฺถิยา กิมานิสํสา.\csromanspacer{} นิพฺพิทา โข อานนฺท วิราคตฺถา วิราคานิสํสา.}

\csromanmatikaanchor{mtk.251}
\csromantocmarknopage{chapter}{1. นิสฺสยวคฺค}{mtk.251}
\bookmark[dest={mtk.251},level=0]{1. นิสฺสยวคฺค}
\prose{วิราโค ปน ภนฺเต กิมตฺถิโย กิมานิสํโส.\csromanspacer{} วิราโค โข อานนฺท วิมุตฺติญาณทสฺสนตฺโถ วิมุตฺติญาณทสฺสนานิสํโส.}

\csromanmatikaanchor{mtk.252}
\csromantocmarkref{subsection}{1. กิมตฺถิยสุตฺต}{mtk.252}
\bookmark[dest={mtk.252},level=2]{1. กิมตฺถิยสุตฺต}
\prose{อิติ โข อานนฺท กุสลานิ สีลานิ อวิปฺปฏิสารตฺถานิ อวิปฺปฏิสารานิสํสานิ.\csromanspacer{} อวิปฺปฏิสาโร ปาโมชฺชตฺโถ ปาโมชฺชานิสํโส.\csromanspacer{} ปาโมชฺชํ ปีตตฺถํ ปีตานิสํสํ.\csromanspacer{} ปีติ ปสฺสทฺธตฺถา ปสฺสทฺธานิสํสา.\csromanspacer{} ปสฺสทฺธิ สุขตฺถา สุขานิสํสา.\csromanspacer{} สุขํ สมาธตฺถํ สมาธานิสํสํ.\csromanspacer{} สมาธิ ยถาภูตญาณทสฺสนตฺโถ ยถาภูตญาณทสฺสนานิสํโส.\csromanspacer{} ยถาภูตญาณทสฺสนํ นิพฺพิทตฺถํ นิพฺพิทานิสํสํ.\csromanspacer{} นิพฺพิทา วิราคตฺถา วิราคานิสํสา.\csromanspacer{} วิราโค วิมุตฺติญาณทสฺสนตฺโถ วิมุตฺติญาณทสฺสนานิสํโส.\csromanspacer{} อิติ โข อานนฺท กุสลานิ สีลานิ อนุปุพฺเพน อคฺคาย ปเรนฺตีติ.\csromanspacer{}\csromanspacer{}ปฐมํ.}

\csromanheader{2}{2. เจตนากรณียสุตฺต}
\proseitem{2}{\csromansymbolfootnote{*}{เหฏฺฐา 258 ปิฏฺเฐปิ.}สีลวโต ภิกฺขเว สีลสมฺปนฺนสฺส น เจตนาย กรณียํ “อวิปฺปฏิสาโร เม อุปฺปชฺชตู”ติ.\csromanspacer{} ธมฺมตา เอสา ภิกฺขเว ยํ สีลวโต สีลสมฺปนฺนสฺส อวิปฺปฏิสาโร อุปฺปชฺชติ.}

\prose{อวิปฺปฏิสารสฺส ภิกฺขเว น เจตนาย กรณียํ “ปาโมชฺชํ เม อุปฺปชฺชตู”ติ.\csromanspacer{} ธมฺมตา เอสา ภิกฺขเว ยํ อวิปฺปฏิสาริสฺส ปาโมชฺชํ อุปฺปชฺชติ.}

\prose{ปมุทิตสฺส ภิกฺขเว น เจตนาย กรณียํ “ปีติ เม อุปฺปชฺชตู”ติ.\csromanspacer{} ธมฺมตา เอสา ภิกฺขเว ยํ ปมุทิตสฺส ปีติ อุปฺปชฺชติ.}

\prose{ปีติมนสฺส ภิกฺขเว น เจตนาย กรณียํ “กาโย เม ปสฺสมฺภตู”ติ.\csromanspacer{} ธมฺมตา เอสา ภิกฺขเว ยํ ปีติมนสฺส กาโย ปสฺสมฺภติ.}

\chapterheadpageread{1. นิสฺสยวคฺค}
\prose{\csromanfolio{517}ปสฺสทฺธกายสฺส ภิกฺขเว น เจตนาย กรณียํ “สุขํ เวทิยามี”ติ.\csromanspacer{} ธมฺมตา เอสา ภิกฺขเว ยํ ปสฺสทฺธกาโย สุขํ เวทิยติ.}

\prose{สุขิโน ภิกฺขเว น เจตนาย กรณียํ “จิตฺตํ เม สมาธิยตู”ติ.\csromanspacer{} ธมฺมตา เอสา ภิกฺขเว ยํ สุขิโน จิตฺตํ สมาธิยติ.}

\prose{สมาหิตสฺส ภิกฺขเว น เจตนาย กรณียํ “ยถาภูตํ ชานามิ ปสฺสามี”ติ.\csromanspacer{} ธมฺมตา เอสา ภิกฺขเว ยํ สมาหิโต ยถาภูตํ ชานาติ ปสฺสติ.}

\prose{ยถาภูตํ ภิกฺขเว ชานโต ปสฺสโต น เจตนาย กรณียํ “นิพฺพินฺทามี”ติ.\csromanspacer{} ธมฺมตา เอสา ภิกฺขเว ยํ ยถาภูตํ ชานํ ปสฺสํ นิพฺพินฺทติ.}

\prose{นิพฺพินฺนสฺส ภิกฺขเว น เจตนาย กรณียํ “วิรชฺชามี”ติ.\csromanspacer{} ธมฺมตา เอสา ภิกฺขเว ยํ นิพฺพินฺโน วิรชฺชติ.}

\prose{วิรตฺตสฺส ภิกฺขเว น เจตนาย กรณียํ “วิมุตฺติญาณทสฺสนํ สจฺฉิกโรมี”ติ.\csromanspacer{} ธมฺมตา เอสา ภิกฺขเว ยํ วิรตฺโต วิมุตฺติญาณทสฺสนํ สจฺฉิกโรติ.}

\csromanmatikaanchor{mtk.253}
\csromantocmarkref{subsection}{2. เจตนากรณียสุตฺต}{mtk.253}
\bookmark[dest={mtk.253},level=2]{2. เจตนากรณียสุตฺต}
\prose{อิติ โข ภิกฺขเว วิราโค วิมุตฺติญาณทสฺสนตฺโถ วิมุตฺติญาณทสฺสนานิสํโส, นิพฺพิทา วิราคตฺถา วิราคานิสํสา, ยถาภูตญาณทสฺสนํ นิพฺพิทตฺถํ นิพฺพิทานิสํสํ, สมาธิ ยถาภูตญาณทสฺสนตฺโถ ยถาภูตญาณทสฺสนานิสํโส, สุขํ สมาธตฺถํ สมาธานิสํสํ, ปสฺสทฺธิ สุขตฺถา สุขานิสํสา, ปีติ ปสฺสทฺธตฺถา ปสฺสทฺธานิสํสา, ปาโมชฺชํ ปีตตฺถํ ปีตานิสํสํ, อวิปฺปฏิสาโร ปาโมชฺชตฺโถ ปาโมชฺชานิสํโส, กุสลานิ สีลานิ อวิปฺปฏิสารตฺถานิ อวิปฺปฏิสารานิสํสานิ.\csromanspacer{} อิติ โข ภิกฺขเว ธมฺมา ธมฺเม อภิสนฺเทนฺติ, ธมฺมา ธมฺเม ปริปูเรนฺติ อปารา ปารํ คมนายาติ.\csromanspacer{}\csromanspacer{}ทุติยํ.}

\csromanheader{2}{3. ปฐม-อุปนิสาสุตฺต}
\proseitem{3}{\csromansymbolfootnote{*}{อํ 2. 16, เหฏฺฐา 259 ปิฏฺเฐสุปิ.}ทุสฺสีลสฺส ภิกฺขเว สีลวิปนฺนสฺส หตูปนิโส โหติ อวิปฺปฏิสาโร, อวิปฺปฏิสาเร อสติ อวิปฺปฏิสารวิปนฺนสฺส หตูปนิสํ โหติ ปาโมชฺชํ, ปาโมชฺเช อสติ ปาโมชฺชวิปนฺนสฺส หตูปนิสา โหติ ปีติ, ปีติยา อสติ ปีติวิปนฺนสฺส หตูปนิสา โหติ ปสฺสทฺธิ, ปสฺสทฺธิยา}

\gambhira{เอกาทสกนิปาตปาฬิ}
\prose{\csromanfolio{518}อสติ ปสฺสทฺธิวิปนฺนสฺส หตูปนิสํ โหติ สุขํ, สุเข อสติ สุขวิปนฺนสฺส หตูปนิโส โหติ สมฺมาสมาธิ, สมฺมาสมาธิมฺหิ อสติ สมฺมาสมาธิวิปนฺนสฺส หตูปนิสํ โหติ ยถาภูตญาณทสฺสนํ, ยถาภูตญาณทสฺสเน อสติ ยถาภูตญาณทสฺสนวิปนฺนสฺส หตูปนิสา โหติ นิพฺพิทา, นิพฺพิทาย อสติ นิพฺพิทาวิปนฺนสฺส หตูปนิโส โหติ วิราโค, วิราเค อสติ วิราควิปนฺนสฺส หตูปนิสํ โหติ วิมุตฺติญาณทสฺสนํ.}

\prose{เสยฺยถาปิ ภิกฺขเว รุกฺโข สาขาปลาสวิปนฺโน ตสฺส ปปฏิกาปิ น ปาริปูริํ คจฺฉติ, ตโจปิ.\csromanspacer{} เผคฺคุปิ.\csromanspacer{} สาโรปิ น ปาริปูริํ คจฺฉติ.\csromanspacer{} เอวเมวํ โข ภิกฺขเว ทุสฺสีลสฺส สีลวิปนฺนสฺส หตูปนิโส โหติ อวิปฺปฏิสาโร, อวิปฺปฏิสาเร อสติ อวิปฺปฏิสารวิปนฺนสฺส หตูปนิสํ โหติ ปาโมชฺชํ ฯเปฯ วิมุตฺติญาณทสฺสนํ.}

\prose{สีลวโต ภิกฺขเว สีลสมฺปนฺนสฺส อุปนิสสมฺปนฺโน โหติ อวิปฺปฏิสาโร, อวิปฺปฏิสาเร สติ อวิปฺปฏิสารสมฺปนฺนสฺส อุปนิสสมฺปนฺนํ โหติ ปาโมชฺชํ, ปาโมชฺเช สติ ปาโมชฺชสมฺปนฺนสฺส อุปนิสสมฺปนฺนา โหติ ปีติ, ปีติยา สติ ปีติสมฺปนฺนสฺส อุปนิสสมฺปนฺนา โหติ ปสฺสทฺธิ, ปสฺสทฺธิยา สติ ปสฺสทฺธิสมฺปนฺนสฺส อุปนิสสมฺปนฺนํ โหติ สุขํ, สุเข สติ สุขสมฺปนฺนสฺส อุปนิสสมฺปนฺโน โหติ สมฺมาสมาธิ, สมฺมาสมาธิมฺหิ สติ สมฺมาสมาธิสมฺปนฺนสฺส อุปนิสสมฺปนฺนํ โหติ ยถาภูตญาณทสฺสนํ, ยถาภูตญาณทสฺสเน สติ ยถาภูตญาณทสฺสนสมฺปนฺนสฺส อุปนิสสมฺปนฺนา โหติ นิพฺพิทา, นิพฺพิทาย สติ นิพฺพิทาสมฺปนฺนสฺส อุปนิสสมฺปนฺโน โหติ วิราโค, วิราเค สติ วิราคสมฺปนฺนสฺส อุปนิสสมฺปนฺนํ โหติ วิมุตฺติญาณทสฺสนํ.}

\prose{เสยฺยถาปิ ภิกฺขเว รุกฺโข สาขาปลาสสมฺปนฺโน ตสฺส ปปฏิกาปิ ปาริปูริํ คจฺฉติ, ตโจปิ.\csromanspacer{} เผคฺคุปิ.\csromanspacer{} สาโรปิ ปาริปูริํ คจฺฉติ.\csromanspacer{} เอวเมวํ โข ภิกฺขเว สีลวโต สีลสมฺปนฺนสฺส อุปนิสสมฺปนฺโน โหติ อวิปฺปฏิสาโร, อวิปฺปฏิสาเร สติ อวิปฺปฏิสารสมฺปนฺนสฺส อุปนิสสมฺปนฺนํ โหติ ฯเปฯ วิมุตฺติญาณทสฺสนนฺติ.\csromanspacer{}\csromanspacer{}ตติยํ.}

\chapterheadpageread{1. นิสฺสยวคฺค}
\csromanheader{2}{4. ทุติย-อุปนิสาสุตฺต}
\proseitem{4}{\csromanfolio{519}\csromansymbolmark{*}ตตฺร โข อายสฺมา สาริปุตฺโต ภิกฺขู อามนฺเตสิ “อาวุโส ภิกฺขเว”ติ\footnote{ภิกฺขโวติ (สี, สฺยา, อิ) เอวํ สพฺพตฺถ เหฏฺฐา 260 ปิฏฺเฐปิ.}. “อาวุโส”ติ โข เต ภิกฺขู อายสฺมโต สาริปุตฺตสฺส ปจฺจสฺโสสุํ.\csromanspacer{} อายสฺมา สาริปุตฺโต เอตทโวจ\csromandash{}}

\prose{ทุสฺสีลสฺส อาวุโส สีลวิปนฺนสฺส หตูปนิโส โหติ อวิปฺปฏิสาโร, อวิปฺปฏิสาเร อสติ อวิปฺปฏิสารวิปนฺนสฺส หตูปนิสํ โหติ ปาโมชฺชํ, ปาโมชฺเช อสติ ปาโมชฺชวิปนฺนสฺส หตูปนิสา โหติ ปีติ, ปีติยา อสติ ปีติวิปนฺนสฺส หตูปนิสา โหติ ปสฺสทฺธิ, ปสฺสทฺธิยา อสติ ปสฺสทฺธิวิปนฺนสฺส หตูปนิสํ โหติ สุขํ, สุเข อสติ สุขวิปนฺนสฺส หตูปนิโส โหติ สมฺมาสมาธิ, สมฺมาสมาธิมฺหิ อสติ สมฺมาสมาธิวิปนฺนสฺส หตูปนิสํ โหติ ยถาภูตญาณทสฺสนํ, ยถาภูตญาณทสฺสเน อสติ ยถาภูตญาณทสฺสนวิปนฺนสฺส หตูปนิสา โหติ นิพฺพิทา, นิพฺพิทาย อสติ นิพฺพิทาวิปนฺนสฺส หตูปนิโส โหติ วิราโค, วิราเค อสติ วิราควิปนฺนสฺส หตูปนิสํ โหติ วิมุตฺติญาณทสฺสนํ.}

\prose{เสยฺยถาปิ อาวุโส รุกฺโข สาขาปลาสวิปนฺโน ตสฺส ปปฏิกาปิ น ปาริปูริํ คจฺฉติ, ตโจปิ.\csromanspacer{} เผคฺคุปิ.\csromanspacer{} สาโรปิ น ปาริปูริํ คจฺฉติ.\csromanspacer{} เอวเมวํ โข อาวุโส ทุสฺสีลสฺส สีลวิปนฺนสฺส หตูปนิโส โหติ อวิปฺปฏิสาโร, อวิปฺปฏิสาเร อสติ อวิปฺปฏิสารวิปนฺนสฺส หตูปนิสํ โหติ ปาโมชฺชํ ฯเปฯ วิมุตฺติญาณทสฺสนํ.}

\csromanmatikaanchor{mtk.254}
\csromantocmarkref{subsection}{3. ปฐม-อุปนิสาสุตฺต}{mtk.254}
\bookmark[dest={mtk.254},level=2]{3. ปฐม-อุปนิสาสุตฺต}
\prose{สีลวโต อาวุโส สีลสมฺปนฺนสฺส อุปนิสสมฺปนฺโน โหติ อวิปฺปฏิสาโร, อวิปฺปฏิสาเร สติ อวิปฺปฏิสารสมฺปนฺนสฺส อุปนิสสมฺปนฺนํ โหติ ปาโมชฺชํ, ปาโมชฺเช สติ ปาโมชฺชสมฺปนฺนสฺส อุปนิสสมฺปนฺนา โหติ ปีติ, ปีติยา สติ ปีติสมฺปนฺนสฺส อุปนิสสมฺปนฺนา โหติ ปสฺสทฺธิ, ปสฺสทฺธิยา สติ ปสฺสทฺธิสมฺปนฺนสฺส อุปนิสสมฺปนฺนํ โหติ สุขํ, สุเข สติ สุขสมฺปนฺนสฺส อุปนิสสมฺปนฺโน โหติ สมฺมาสมาธิ, สมฺมาสมาธิมฺหิ สติ สมฺมาสมาธิสมฺปนฺนสฺส อุปนิสสมฺปนฺนํ โหติ ยถาภูตญาณทสฺสนํ, ยถาภูตญาณทสฺสเน สติ ยถาภูตญาณทสฺสนสมฺปนฺนสฺส อุปนิสสมฺปนฺนา โหติ นิพฺพิทา, นิพฺพิทาย สติ}

\gambhira{เอกาทสกนิปาตปาฬิ}
\prose{\csromanfolio{520}นิพฺพิทาสมฺปนฺนสฺส อุปนิสสมฺปนฺโน โหติ วิราโค, วิราเค สติ วิราคสมฺปนฺนสฺส อุปนิสสมฺปนฺนํ โหติ วิมุตฺติญาณทสฺสนํ.}

\prose{เสยฺยถาปิ อาวุโส รุกฺโข สาขาปลาสสมฺปนฺโน ตสฺส ปปฏิกาปิ ปาริปูริํ คจฺฉติ, ตโจปิ.\csromanspacer{} เผคฺคุปิ.\csromanspacer{} สาโรปิ ปาริปูริํ คจฺฉติ.\csromanspacer{} เอวเมวํ โข อาวุโส สีลวโต สีลสมฺปนฺนสฺส อุปนิสสมฺปนฺโน โหติ อวิปฺปฏิสาโร, อวิปฺปฏิสาเร สติ อวิปฺปฏิสารสมฺปนฺนสฺส อุปนิสสมฺปนฺนํ โหติ ปาโมชฺชํ ฯเปฯ วิมุตฺติญาณทสฺสนนฺติ ฯเปฯ.\csromanspacer{} จตุตฺถํ.}

\csromanheader{2}{5. ตติย-อุปนิสาสุตฺต}
\proseitem{5}{\csromansymbolfootnote{*}{เหฏฺฐา 261 ปิฏฺเฐปิ.}ตตฺร โข อายสฺมา อานนฺโท ภิกฺขู อามนฺเตสิ ฯเปฯ.\csromanspacer{} ทุสฺสีลสฺส อาวุโส สีลวิปนฺนสฺส หตูปนิโส โหติ อวิปฺปฏิสาโร อวิปฺปฏิสาเร อสติ อวิปฺปฏิสารวิปนฺนสฺส หตูปนิสํ โหติ ปาโมชฺชํ, ปาโมชฺเช อสติ ปาโมชฺชวิปนฺนสฺส หตูปนิสา โหติ ปีติ, ปีติยา อสติ ปีติวิปนฺนสฺส หตูปนิสา โหติ ปสฺสทฺธิ, ปสฺสทฺธิยา อสติ ปสฺสทฺธิวิปนฺนสฺส หตูปนิสํ โหติ สุขํ, สุเข อสติ สุขวิปนฺนสฺส หตูปนิโส โหติ สมฺมาสมาธิ, สมฺมาสมาธิมฺหิ อสติ สมฺมาสมาธิวิปนฺนสฺส หตูปนิสํ โหติ ยถาภูตญาณทสฺสนํ, ยถาภูตญาณทสฺสเน อสติ ยถาภูตญาณทสฺสนวิปนฺนสฺส หตูปนิสา โหติ นิพฺพิทา, นิพฺพิทาย อสติ นิพฺพิทาวิปนฺนสฺส หตูปนิโส โหติ วิราโค, วิราเค อสติ วิราควิปนฺนสฺส หตูปนิสํ โหติ วิมุตฺติญาณทสฺสนํ.}

\prose{เสยฺยถาปิ อาวุโส รุกฺโข สาขาปลาสวิปนฺโน ตสฺส ปปฏิกาปิ น ปาริปูริํ คจฺฉติ, ตโจปิ.\csromanspacer{} เผคฺคุปิ.\csromanspacer{} สาโรปิ น ปาริปูริํ คจฺฉติ.\csromanspacer{} เอวเมวํ โข อาวุโส ทุสฺสีลสฺส สีลวิปนฺนสฺส หตูปนิโส โหติ อวิปฺปฏิสาโร, อวิปฺปฏิสาเร อสติ อวิปฺปฏิสารวิปนฺนสฺส ปหูปนิสํ โหติ ปาโมชฺชํ ฯเปฯ วิมุตฺติญาณทสฺสนํ.}

\prose{สีลวโต อาวุโส สีลสมฺปนฺนสฺส อุปนิสสมฺปนฺโน โหติ อวิปฺปฏิสาโร, อวิปฺปฏิสาเร สติ อวิปฺปฏิสารสมฺปนฺนสฺส อุปนิสสมฺปนฺนํ โหติ ปาโมชฺชํ, ปาโมชฺเช สติ ปาโมชฺชสมฺปนฺนสฺส อุปนิสสมฺปนฺนา โหติ ปีติ, ปีติยา สติ ปีติสมฺปนฺนสฺส อุปนิสสมฺปนฺนา โหติ ปสฺสทฺธิ, ปสฺสทฺธิยา สติ ปสฺสทฺธิสมฺปนฺนสฺส อุปนิสสมฺปนฺนํ โหติ สุขํ, สุเข สติ}

\csromanmatikaanchor{mtk.255}
\csromantocmarkref{subsection}{4. ทุติย-อุปนิสาสุตฺต}{mtk.255}
\bookmark[dest={mtk.255},level=2]{4. ทุติย-อุปนิสาสุตฺต}
\chapterheadpageread{1. นิสฺสยวคฺค}
\prose{\csromanfolio{521}สุขสมฺปนฺนสฺส อุปนิสสมฺปนฺโน โหติ สมฺมาสมาธิ, สมฺมาสมาธิมฺหิ สติ สมฺมาสมาธิสมฺปนฺนสฺส อุปนิสสมฺปนฺนํ โหติ ยถาภูตญาณทสฺสนํ, ยถาภูตญาณทสฺสเน สติ ยถาภูตญาณทสฺสนสมฺปนฺนสฺส อุปนิสสมฺปนฺนา โหติ นิพฺพิทา, นิพฺพิทาย สติ นิพฺพิทาสมฺปนฺนสฺส อุปนิสสมฺปนฺโน โหติ วิราโค, วิราเค สติ วิราคสมฺปนฺนสฺส อุปนิสสมฺปนฺนํ โหติ วิมุตฺติญาณทสฺสนํ.}

\prose{เสยฺยถาปิ อาวุโส รุกฺโข สาขาปลาสสมฺปนฺโน ตสฺส ปปฏิกาปิ ปาริปูริํ คจฺฉติ, ตโจปิ.\csromanspacer{} เผคฺคุปิ.\csromanspacer{} สาโรปิ ปาริปูริํ คจฺฉติ.\csromanspacer{} เอวเมวํ โข อาวุโส สีลวโต สีลสมฺปนฺนสฺส อุปนิสสมฺปนฺโน โหติ อวิปฺปฏิสาโร, อวิปฺปฏิสาเร สติ อวิปฺปฏิสารสมฺปนฺนสฺส อุปนิสสมฺปนฺนํ โหติ ปาโมชฺชํ ฯเปฯ วิมุตฺติญาณทสฺสนนฺติ.\csromanspacer{}\csromanspacer{}ปญฺจมํ.}

\csromanheader{2}{6. พฺยสนสุตฺต}
\proseitem{6}{(โย โส ภิกฺขเว ภิกฺขุ อกฺโกสโก ปริภาสโก อริยูปวาโท สพฺรหฺมจารีนํ, ฐานเมตํ อวกาโส, ยํ โส เอกาทสนฺนํ พฺยสนานํ อญฺญตรํ พฺยสนํ นิคจฺเฉยฺย.}

\prose{กตเมสํ เอกาทสนฺนํ? อนธิคตํ นาธิคจฺฉติ, อธิคตา ปริหายติ, สทฺธมฺมสฺส น โวทายนฺติ, สทฺธมฺเมสุ วา อธิมานิโก โหติ, อนภิรโต วา พฺรหฺมจริยํ จรติ, อญฺญตรํ วา สํกิลิฏฺฐํ อาปตฺติํ อาปชฺชติ, สิกฺขํ วา ปจฺจกฺขาย หีนายาวตฺตติ, คาฬฺหํ วา โรคาตงฺกํ ผุสติ, อุมฺมาทํ วา ปาปุณาติ จิตฺตกฺเขปํ วา, สมฺมูโฬฺห กาลํ กโรติ, กายสฺส เภทา ปรํ มรณา อปายํ ทุคฺคติํ วินิปาตํ นิรยํ อุปปชฺชติ.\csromanspacer{} โย โส ภิกฺขเว ภิกฺขุ อกฺโกสโก ปริภาสโก อริยูปวาโท สพฺรหฺมจารีนํ, ฐานเมตํ อวกาโส, ยํ โส อิเมสํ เอกาทสนฺนํ พฺยสนานํ อญฺญตรํ พฺยสนํ นิคจฺเฉยฺย.)\footnote{( ) เอตฺถนฺตเร ปาโฐ สี-สฺยา-กํ-อิ-โปตฺถเกสุ น ทิสฺสติ.} โย โส ภิกฺขเว ภิกฺขุ อกฺโกสโก ปริภาสโก อริยูปวาโท สพฺรหฺมจารีนํ, อฏฺฐานเมตํ อนวกาโส, ยํ โส เอกาทสนฺนํ พฺยสนานํ อญฺญตรํ พฺยสนํ น นิคจฺเฉยฺย.}

\gambhira{เอกาทสกนิปาตปาฬิ}
\csromanmatikaanchor{mtk.256}
\csromantocmarkref{subsection}{5. ตติย-อุปนิสาสุตฺต}{mtk.256}
\bookmark[dest={mtk.256},level=2]{5. ตติย-อุปนิสาสุตฺต}
\prose{\csromanfolio{522}กตเมสํ เอกาทสนฺนํ? อนธิคตํ นาธิคจฺฉติ, อธิคตา ปริหายติ, สทฺธมฺมสฺส น โวทายนฺติ, สทฺธมฺเมสุ วา อธิมานิโก โหติ, อนภิรโต วา พฺรหฺมจริยํ จรติ, อญฺญตรํ วา สํกิลิฏฺฐํ อาปตฺติํ อาปชฺชติ, สิกฺขํ วา ปจฺจกฺขาย หีนายาวตฺตติ, คาฬฺหํ วา โรคาตงฺกํ ผุสติ, อุมฺมาทํ วา ปาปุณาติ จิตฺตกฺเขปํ วา, สมฺมูโฬฺห กาลํ กโรติ, กายสฺส เภทา ปรํ มรณา อปายํ ทุคฺคติํ วินิปาตํ นิรยํ อุปปชฺชติ.\csromanspacer{} โย โส ภิกฺขเว ภิกฺขุ อกฺโกสโก ปริภาสโก อริยูปวาโท สพฺรหฺมจารีนํ, อฏฺฐานเมตํ อนวกาโส, ยํ โส อิเมสํ เอกาทสนฺนํ พฺยสนานํ อญฺญตรํ พฺยสนํ น นิคจฺเฉยฺยาติ.\csromanspacer{}\csromanspacer{}ฉฏฺฐํ.}

\csromanheader{2}{7. สญฺญาสุตฺต}
\proseitem{7}{อถ โข อายสฺมา อานนฺโท เยน ภควา เตนุปสงฺกมิ, อุปสงฺกมิตฺวา ภควนฺตํ อภิวาเทตฺวา เอกมนฺตํ นิสีทิ, เอกมนฺตํ นิสินฺโน โข อายสฺมา อานนฺโท ภควนฺตํ เอตทโวจ\csromandash{}}

\prose{สิยา นุ โข ภนฺเต ภิกฺขุโน ตถารูโป สมาธิปฏิลาโภ, ยถา เนว ปถวิยํ ปถวิสญฺญี อสฺส, น อาปสฺมิํ อาโปสญฺญี อสฺส, น เตชสฺมิํ เตโชสญฺญี อสฺส, น วายสฺมิํ วาโยสญฺญี อสฺส, น อากาสานญฺจายตเน อากาสานญฺจายตนสญฺญี อสฺส, น วิญฺญาณญฺจายตเน วิญฺญาณญฺจายตนสญฺญี อสฺส, น อากิญฺจญฺญายตเน อากิญฺจญฺญายตนสญฺญี อสฺส, น เนวสญฺญานาสญฺญายตเน เนวสญฺญานาสญฺญายตนสญฺญี อสฺส, น อิธโลเก อิธโลกสญฺญี อสฺส, น ปรโลเก ปรโลกสญฺญี อสฺส, ยมฺปิทํ ทิฏฺฐํ สุตํ มุตํ วิญฺญาตํ ปตฺตํ ปริเยสิตํ อนุวิจริตํ มนสา, ตตฺราปิ น สญฺญี อสฺส, สญฺญี จ ปน อสฺสาติ.}

\prose{สิยา อานนฺท ภิกฺขุโน ตถารูโป สมาธิปฏิลาโภ, ยถา เนว ปถวิยํ ปถวิสญฺญี อสฺส, น อาปสฺมิํ อาโปสญฺญี อสฺส, น เตชสฺมิํ เตโชสญฺญี อสฺส, น วายสฺมิํ วาโยสญฺญี อสฺส, น อากาสานญฺจายตเน อากาสานญฺจายตนสญฺญี อสฺส, น วิญฺญาณญฺจายตเน วิญฺญาณญฺจายตนสญฺญี อสฺส, น อากิญฺจญฺญายตเน อากิญฺจญฺญายตนสญฺญี อสฺส, น เนวสญฺญานาสญฺญายตเน เนวสญฺญานาสญฺญายตนสญฺญี อสฺส, น อิธโลเก อิธโลกสญฺญี}

\chapterheadpageread{1. นิสฺสยวคฺค}
\prose{\csromanfolio{523}อสฺส, น ปรโลเก ปรโลกสญฺญี อสฺส, ยมฺปิทํ ทิฏฺฐํ สุตํ มุตํ วิญฺญาตํ ปตฺตํ ปริเยสิตํ อนุวิจริตํ มนสา, ตตฺราปิ น สญฺญี อสฺส, สญฺญี จ ปน อสฺสาติ.}

\csromanmatikaanchor{mtk.257}
\csromantocmarkref{subsection}{6. พฺยสนสุตฺต}{mtk.257}
\bookmark[dest={mtk.257},level=2]{6. พฺยสนสุตฺต}
\prose{ยถา กถํ ปน ภนฺเต สิยา ภิกฺขุโน ตถารูโป สมาธิปฏิลาโภ, ยถา เนว ปถวิยํ ปถวิสญฺญี อสฺส, น อาปสฺมิํ อาโปสญฺญี อสฺส, น เตชสฺมิํ เตโชสญฺญี อสฺส, น วายสฺมิํ วาโยสญฺญี อสฺส, น อากาสานญฺจายตเน อากาสานญฺจายตนสญฺญี อสฺส, น วิญฺญาณญฺจายตเน วิญฺญาณญฺจายตนสญฺญี อสฺส, น อากิญฺจญฺญายตเน อากิญฺจญฺญายตนสญฺญี อสฺส, น เนวสญฺญานาสญฺญายตเน เนวสญฺญานาสญฺญายตนสญฺญี อสฺส, น อิธโลเก อิธโลกสญฺญี อสฺส, น ปรโลเก ปรโลกสญฺญี อสฺส, ยมฺปิทํ ทิฏฺฐํ สุตํ มุตํ วิญฺญาตํ ปตฺตํ ปริเยสิตํ อนุวิจริตํ มนสา, ตตฺราปิ น สญฺญี อสฺส, สญฺญี จ ปน อสฺสาติ.}

\prose{อิธานนฺท ภิกฺขุ เอวํสญฺญี โหติ “เอตํ สนฺตํ เอตํ ปณีตํ, ยทิทํ สพฺพสงฺขารสมโถ สพฺพูปธิปฏินิสฺสคฺโค ตณฺหากฺขโย วิราโค นิโรโธ นิพฺพานนฺ”ติ.\csromanspacer{} เอวํ โข อานนฺท สิยา ภิกฺขุโน ตถารูโป สมาธิปฏิลาโภ, ยถา เนว ปถวิยํ ปถวิสญฺญี อสฺส, น อาปสฺมิํ อาโปสญฺญี อสฺส, น เตชสฺมิํ เตโชสญฺญี อสฺส, น วายสฺมิํ วาโยสญฺญี อสฺส, น อากาสานญฺจายตเน อากาสานญฺจายตนสญฺญี อสฺส, น วิญฺญาณญฺจายตเน วิญฺญาณญฺจายตนสญฺญี อสฺส, น อากิญฺจญฺญายตเน อากิญฺจญฺญายตนสญฺญี อสฺส, น เนวสญฺญานาสญฺญายตเน เนวสญฺญานาสญฺญายตนสญฺญี อสฺส, น อิธโลเก อิธโลกสญฺญี อสฺส, น ปรโลเก ปรโลกสญฺญี อสฺส, ยมฺปิทํ ทิฏฺฐํ สุตํ มุตํ วิญฺญาตํ ปตฺตํ ปริเยสิตํ อนุวิจริตํ มนสา, ตตฺราปิ น สญฺญี อสฺส, สญฺญี จ ปน อสฺสาติ.}

\prose{อถ โข อายสฺมา อานนฺโท ภควโต ภาสิตํ อภินนฺทิตฺวา อนุโมทิตฺวา อุฏฺฐายาสนา ภควนฺตํ อภิวาเทตฺวา ปทกฺขิณํ กตฺวา เยนายสฺมา สาริปุตฺโต เตนุปสงฺกมิ, อุปสงฺกมิตฺวา อายสฺมตา สาริปุตฺเตน สทฺธิํ สมฺโมทิ, สมฺโมทนียํ กถํ สารณียํ วีติสาเรตฺวา เอกมนฺตํ นิสีทิ, เอกมนฺตํ นิสินฺโน โข อายสฺมา อานนฺโท อายสฺมนฺตํ สาริปุตฺตํ เอตทโวจ\csromandash{}}

\gambhira{เอกาทสกนิปาตปาฬิ}
\prose{\csromanfolio{524}สิยา นุ โข อาวุโส สาริปุตฺต ภิกฺขุโน ตถารูโป สมาธิปฏิลาโภ, ยถา เนว ปถวิยํ ปถวิสญฺญี อสฺส ฯเปฯ ยมฺปิทํ ทิฏฺฐํ สุตํ มุตํ วิญฺญาตํ ปตฺตํ ปริเยสิตํ อนุวิจริตํ มนสา, ตตฺราปิ น สญฺญี อสฺส, สญฺญี ปน อสฺสาติ.\csromanspacer{} สิยา อาวุโส อานนฺท ภิกฺขุโน ตถารูโป สมาธิปฏิลาโภ, ยถา เนว ปถวิยํ ปถวิสญฺญี อสฺส ฯเปฯ ยมฺปิทํ ทิฏฺฐํ สุตํ มุตํ วิญฺญาตํ ปตฺตํ ปริเยสิตํ อนุวิจริตํ มนสา, ตตฺราปิ น สญฺญี อสฺส, สญฺญี จ ปน อสฺสาติ.}

\prose{ยถา กถํ ปนาวุโส สาริปุตฺต สิยา ภิกฺขุโน ตถารูโป สมาธิปฏิลาโภ, ยถา เนว ปถวิยํ ปถวิสญฺญี อสฺส ฯเปฯ ยมฺปิทํ ทิฏฺฐํ สุตํ มุตํ วิญฺญาตํ ปตฺตํ ปริเยสิตํ อนุวิจริตํ มนสา, ตตฺราปิ น สญฺญี อสฺส, สญฺญี จ ปน อสฺสาติ.}

\csromanmatikaanchor{mtk.258}
\csromantocmarkref{subsection}{7. สญฺญาสุตฺต}{mtk.258}
\bookmark[dest={mtk.258},level=2]{7. สญฺญาสุตฺต}
\prose{อิธ อาวุโส อานนฺท ภิกฺขุ เอวํสญฺญี โหติ “เอตํ สนฺตํ เอตํ ปณีตํ, ยทิทํ สพฺพสงฺขารสมโถ สพฺพูปธิปฏินิสฺสคฺโค ตณฺหากฺขโย วิราโค นิโรโธ นิพฺพานนฺ”ติ.\csromanspacer{} เอวํ โข อาวุโส อานนฺท สิยา ภิกฺขุโน ตถารูโป สมาธิปฏิลาโภ, ยถา เนว ปถวิยํ ปถวิสญฺญี อสฺส ฯเปฯ ยมฺปิทํ ทิฏฺฐํ สุตํ มุตํ วิญฺญาตํ ปตฺตํ ปริเยสิตํ อนุวิจริตํ มนสา, ตตฺราปิ น สญฺญี อสฺส, สญฺญี จ ปน อสฺสาติ.}

\prose{\csromansymbolfootnote{*}{อํ 2. 548 ปิฏฺเฐปิ.}อจฺฉริยํ อาวุโส, อพฺภุตํ อาวุโส, ยตฺร หิ นาม สตฺถุ เจว สาวกสฺส จ อตฺเถน อตฺโถ พฺยญฺชเนน พฺยญฺชนํ สํสนฺทิสฺสติ สเมสฺสติ น วิคฺคยฺหิสฺสติ ยทิทํ อคฺคปทสฺมิํ.\csromanspacer{} อิทานาหํ อาวุโส ภควนฺตํ อุปสงฺกมิตฺวา เอตมตฺถํ อปุจฺฉิํ, ภควาปิ เม เอเตหิ อกฺขเรหิ เอเตหิ ปเทหิ เอเตหิ พฺยญฺชเนหิ เอตมตฺถํ พฺยากาสิ เสยฺยถาปิ อายสฺมา สาริปุตฺโต.\csromanspacer{} อจฺฉริยํ อาวุโส, อพฺภุตํ อาวุโส, ยตฺร หิ นาม สตฺถุ เจว สาวกสฺส จ อตฺเถน อตฺโถ พฺยญฺชเนน พฺยญฺชนํ สํสนฺทิสฺสติ สเมสฺสติ น วิคฺคยฺหิสฺสติ ยทิทํ อคฺคปทสฺมินฺติ.\csromanspacer{}\csromanspacer{}สตฺตมํ.}

\csromanheader{2}{8. มนสิการสุตฺต}
\proseitem{8}{อถ โข อายสฺมา อานนฺโท เยน ภควา เตนุปสงฺกมิ, อุปสงฺกมิตฺวา ภควนฺตํ อภิวาเทตฺวา เอกมนฺตํ นิสีทิ, เอกมนฺตํ นิสินฺโน โข อายสฺมา อานนฺโท ภควนฺตํ เอตทโวจ\csromandash{}}

\chapterheadpageread{1. นิสฺสยวคฺค}
\prose{\csromanfolio{525}สิยา นุ โข ภนฺเต ภิกฺขุโน ตถารูโป สมาธิปฏิลาโภ, ยถา น จกฺขุํ มนสิ กเรยฺย, น รูปํ มนสิ กเรยฺย, น โสตํ มนสิ กเรยฺย, น สทฺทํ มนสิ กเรยฺย, น ฆานํ มนสิ กเรยฺย, น คนฺธํ มนสิ กเรยฺย, น ชิวฺหํ มนสิ กเรยฺย, น รสํ มนสิ กเรยฺย, น กายํ มนสิ กเรยฺย, น โผฏฺฐพฺพํ มนสิ กเรยฺย, น ปถวิํ มนสิ กเรยฺย, น อาปํ มนสิ กเรยฺย, น เตชํ มนสิ กเรยฺย, น วายํ มนสิ กเรยฺย, น อากาสานญฺจายตนํ มนสิ กเรยฺย, น วิญฺญาณญฺจายตนํ มนสิ กเรยฺย, น อากิญฺจญฺญายตนํ มนสิ กเรยฺย, น เนวสญฺญานาสญฺญายตนํ มนสิ กเรยฺย, น อิธโลกํ มนสิ กเรยฺย, น ปรโลกํ มนสิ กเรยฺย, ยมฺปิทํ ทิฏฺฐํ สุตํ มุตํ วิญฺญาตํ ปตฺตํ ปริเยสิตํ อนุวิจริตํ มนสา, ตมฺปิ น มนสิ กเรยฺย, มนสิ จ ปน กเรยฺยาติ.}

\prose{สิยา อานนฺท ภิกฺขุโน ตถารูโป สมาธิปฏิลาโภ, ยถา น จกฺขุํ มนสิ กเรยฺย, น รูปํ มนสิ กเรยฺย, น โสตํ มนสิ กเรยฺย, น สทฺทํ มนสิ กเรยฺย, น ฆานํ มนสิ กเรยฺย, น คนฺธํ มนสิ กเรยฺย, น ชิวฺหํ มนสิ กเรยฺย, น รสํ มนสิ กเรยฺย, น กายํ มนสิ กเรยฺย, น โผฏฺฐพฺพํ มนสิ กเรยฺย, น ปถวิํ มนสิ กเรยฺย, น อาปํ มนสิ กเรยฺย, น เตชํ มนสิ กเรยฺย, น วายํ มนสิ กเรยฺย, น อากาสานญฺจายตนํ มนสิ กเรยฺย, น วิญฺญาณญฺจายตนํ มนสิ กเรยฺย, น อากิญฺจญฺญายตนํ มนสิ กเรยฺย, น เนวสญฺญานาสญฺญายตนํ มนสิ กเรยฺย, น อิธโลกํ มนสิ กเรยฺย, น ปรโลกํ มนสิ กเรยฺย, ยมฺปิทํ ทิฏฺฐํ สุตํ มุตํ วิญฺญาตํ ปตฺตํ ปริเยสิตํ อนุวิจริตํ มนสา, ตมฺปิ น มนสิ กเรยฺย, มนสิ จ ปน กเรยฺยาติ.}

\prose{ยถา กถํ ปน ภนฺเต สิยา ภิกฺขุโน ตถารูโป สมาธิปฏิลาโภ, ยถา น จกฺขุํ มนสิ กเรยฺย, น รูปํ มนสิ กเรยฺย ฯเปฯ ยมฺปิทํ ทิฏฺฐํ สุตํ มุตํ วิญฺญาตํ ปตฺตํ ปริเยสิตํ อนุวิจริตํ มนสา, ตมฺปิ น มนสิ กเรยฺย, มนสิ จ ปน กเรยฺยาติ.}

\prose{อิธานนฺท ภิกฺขุ เอวํ มนสิ กโรติ “เอตํ สนฺตํ เอตํ ปณีตํ, ยทิทํ สพฺพสงฺขารสมโถ สพฺพูปธิปฏินิสฺสคฺโค ตณฺหากฺขโย วิราโค นิโรโธ นิพฺพานนฺ”ติ.\csromanspacer{} เอวํ โข อานนฺท สิยา ภิกฺขุโน ตถารูโป สมาธิปฏิลาโภ, ยถา น จกฺขุํ มนสิ กเรยฺย, น รูปํ มนสิ กเรยฺย ฯเปฯ.}

\gambhira{เอกาทสกนิปาตปาฬิ}
\prose{\csromanfolio{526}ยมฺปิทํ ทิฏฺฐํ สุตํ มุตํ วิญฺญาตํ ปตฺตํ ปริเยสิตํ อนุวิจริตํ มนสา, ตมฺปิ น มนสิ กเรยฺย, มนสิ จ ปน กเรยฺยาติ.\csromanspacer{}\csromanspacer{}อฏฺฐมํ.}

\csromanheader{2}{9. สทฺธสุตฺต}
\proseitem{9}{เอกํ สมยํ ภควา นาติเก วิหรติ คิญฺชกาวสเถ.\csromanspacer{} อถ โข อายสฺมา สทฺโธ เยน ภควา เตนุปสงฺกมิ, อุปสงฺกมิตฺวา ภควนฺตํ อภิวาเทตฺวา เอกมนฺตํ นิสีทิ, เอกมนฺตํ นิสินฺนํ โข อายสฺมนฺตํ สทฺธํ ภควา เอตทโวจ\csromandash{}}

\prose{อาชานียฌายิตํ โข สทฺธ ฌาย, มา ขฬุงฺกฌายิตํ\footnote{อาชานียชฺฌายิตํ โข สทฺธ ฌายถ, มา ขฬุงฺกชฺฌายิตํ (สี, อิ)}.\csromanspacer{} กถญฺจ ขฬุงฺกฌายิตํ โหติ? อสฺสขฬุงฺโก หิ สทฺธ โทณิยา พทฺโธ\footnote{พนฺโธ (สฺยา, ก)} “ยวสํ ยวสนฺ”ติ ฌายติ.\csromanspacer{} ตํ กิสฺส เหตุ? น หิ สทฺธ อสฺสขฬุงฺกสฺส โทณิยา พทฺธสฺส เอวํ โหติ “กิํ นุ โข มํ อชฺช อสฺสทมฺมสารถิ การณํ กาเรสฺสติ, กิมสฺสาหํ\footnote{กมฺมสฺสาหํ (ก)} ปฏิกโรมี”ติ.\csromanspacer{} โส โทณิยา พทฺโธ “ยวสํ ยวสนฺ”ติ ฌายติ.\csromanspacer{} เอวเมวํ โข สทฺธ อิเธกจฺโจ ปุริสขฬุงฺโก อรญฺญคโตปิ รุกฺขมูลคโตปิ สุญฺญาคารคโตปิ กามราคปริยุฏฺฐิเตน เจตสา วิหรติ กามราคปเรเตน, อุปฺปนฺนสฺส จ กามราคสฺส นิสฺสรณํ ยถาภูตํ นปฺปชานาติ.\csromanspacer{} โส กามราคํเยว อนฺตรํ กตฺวา ฌายติ ปชฺฌายติ นิชฺฌายติ อวชฺฌายติ.\csromanspacer{} พฺยาปาทปริยุฏฺฐิเตน เจตสา วิหรติ.\csromanspacer{} ถินมิทฺธปริยุฏฺฐิเตน เจตสา วิหรติ.\csromanspacer{} อุทฺธจฺจกุกฺกุจฺจปริยุฏฺฐิเตน เจตสา วิหรติ.\csromanspacer{} วิจิกิจฺฉาปริยุฏฺฐิเตน เจตสา วิหรติ วิจิกิจฺฉาปเรเตน, อุปฺปนฺนาย จ วิจิกิจฺฉาย นิสฺสรณํ ยถาภูตํ นปฺปชานาติ.\csromanspacer{} โส วิจิกิจฺฉํเยว อนฺตรํ กตฺวา ฌายติ ปชฺฌายติ นิชฺฌายติ อวชฺฌายติ.\csromanspacer{} โส ปถวิมฺปิ นิสฺสาย ฌายติ, อาปมฺปิ นิสฺสาย ฌายติ, เตชมฺปิ นิสฺสาย ฌายติ, วายมฺปิ นิสฺสาย ฌายติ, อากาสานญฺจายตนมฺปิ นิสฺสาย ฌายติ, วิญฺญาณญฺจายตนมฺปิ นิสฺสาย ฌายติ, อากิญฺจญฺญายตนมฺปิ นิสฺสาย ฌายติ, เนวสญฺญานาสญฺญายตนมฺปิ นิสฺสาย ฌายติ, อิธโลกมฺปิ นิสฺสาย ฌายติ, ปรโลกมฺปิ นิสฺสาย ฌายติ, ยมฺปิทํ ทิฏฺฐํ สุตํ มุตํ วิญฺญาตํ ปตฺตํ ปริเยสิตํ อนุวิจริตํ มนสา ตมฺปิ นิสฺสาย ฌายติ.\csromanspacer{} เอวํ โข สทฺธ ปุริสขฬุงฺกฌายิตํ โหติ.}

\csromanmatikaanchor{mtk.259}
\csromantocmarkref{subsection}{8. มนสิการสุตฺต}{mtk.259}
\bookmark[dest={mtk.259},level=2]{8. มนสิการสุตฺต}
\chapterheadpageread{1. นิสฺสยวคฺค}
\prose{\csromanfolio{527}กถญฺจ สทฺธ อาชานียฌายิตํ โหติ? ภโทฺร หิ สทฺธ อสฺสาชานีโย โทณิยา พทฺโธ น “ยวสํ ยวสนฺ”ติ ฌายติ.\csromanspacer{} ตํ กิสฺส เหตุ? ภทฺรสฺส หิ สทฺธ อสฺสาชานียสฺส โทณิยา พทฺธสฺส เอวํ โหติ “กิํ นุ โข มํ อชฺช อสฺสทมฺมสารถิ การณํ กาเรสฺสติ.\csromanspacer{} กิมสฺสาหํ ปฏิกโรมี”ติ.\csromanspacer{} โส โทณิยา พทฺโธ น “ยวสํ ยวสนฺ”ติ ฌายติ.\csromanspacer{} ภโทฺร หิ สทฺธ อสฺสาชานีโย ยถา อิณํ ยถา พนฺธํ ยถา ชานิํ ยถา กลิํ เอวํ ปโตทสฺส อชฺโฌหรณํ สมนุปสฺสติ.\csromanspacer{} เอวเมวํ โข สทฺธ ภโทฺร ปุริสาชานีโย อรญฺญคโตปิ รุกฺขมูลคโตปิ สุญฺญาคารคโตปิ น กามราคปริยุฏฺฐิเตน เจตสา วิหรติ น กามราคปเรเตน, อุปฺปนฺนสฺส จ กามราคสฺส นิสฺสรณํ ยถาภูตํ ปชานาติ.\csromanspacer{} น พฺยาปาทปริยุฏฺฐิเตน เจตสา วิหรติ.\csromanspacer{} น ถินมิทฺธปริยุฏฺฐิเตน เจตสา วิหรติ.\csromanspacer{} น อุทฺธจฺจกุกฺกุจฺจปริยุฏฺฐิเตน เจตสา วิหรติ.\csromanspacer{} น วิจิกิจฺฉาปริยุฏฺฐิเตน เจตสา วิหรติ น วิจิกิจฺฉาปเรเตน, อุปฺปนฺนาย จ วิจิกิจฺฉาย นิสฺสรณํ ยถาภูตํ ปชานาติ.\csromanspacer{} โส เนว ปถวิํ นิสฺสาย ฌายติ, น อาปํ นิสฺสาย ฌายติ, น เตชํ นิสฺสาย ฌายติ, น วายํ นิสฺสาย ฌายติ, น อากาสานญฺจายตนํ นิสฺสาย ฌายติ, น วิญฺญาณญฺจายตนํ นิสฺสาย ฌายติ, น อากิญฺจญฺญายตนํ นิสฺสาย ฌายติ, น เนวสญฺญานาสญฺญายตนํ นิสฺสาย ฌายติ, น อิธโลกํ นิสฺสาย ฌายติ, น ปรโลกํ นิสฺสาย ฌายติ, ยมฺปิทํ ทิฏฺฐํ สุตํ มุตํ วิญฺญาตํ ปตฺตํ ปริเยสิตํ อนุวิจริตํ มนสา ตมฺปิ นิสฺสาย น ฌายติ ฌายติ จ ปน.\csromanspacer{} เอวํ ฌายิญฺจ ปน สทฺธ ภทฺรํ ปุริสาชานียํ ส-อินฺทา เทวา สพฺรหฺมกา สปชาปติกา อารกาว นมสฺสนฺติ\csromandash{}}

\csromangathameasure{นโม เต ปุริสาชญฺญ, นโม เต ปุริสุตฺตม. \\ ยสฺส เต นาภิชานาม, ยมฺปิ นิสฺสาย ฌายสีติ.}
\csromangathagroup{\csromangathastanza{นโม เต ปุริสาชญฺญ, นโม เต ปุริสุตฺตม. \\ ยสฺส เต นาภิชานาม, ยมฺปิ นิสฺสาย ฌายสีติ.}}

\prose{เอวํ วุตฺเต อายสฺมา สทฺโธ ภควนฺตํ เอตทโวจ\csromandash{}กถํ ฌายี ปน ภนฺเต ภโทฺร ปุริสาชานีโย\footnote{ปุริสาชานีโย ฌายติ, โส (สี, สฺยา, อิ), ปุริสาชานีโย, โส (ก)} เนว ปถวิํ นิสฺสาย ฌายติ, น อาปํ นิสฺสาย ฌายติ, น เตชํ นิสฺสาย ฌายติ, น วายํ นิสฺสาย ฌายติ, น อากาสานญฺจายตนํ นิสฺสาย ฌายติ, น วิญฺญาณญฺจายตนํ นิสฺสาย ฌายติ, น อากิญฺจาญฺญายตนํ นิสฺสาย ฌายติ, น เนวสญฺญานาสญฺญายตนํ นิสฺสาย ฌายติ, น อิธโลกํ นิสฺสาย}

\gambhira{เอกาทสกนิปาตปาฬิ}
\prose{\csromanfolio{528}ฌายติ, น ปรโลกํ นิสฺสาย ฌายติ, ยมฺปิทํ ทิฏฺฐํ สุตํ มุตํ วิญฺญาตํ ปตฺตํ ปริเยสิตํ อนุวิจริตํ มนสา ตมฺปิ นิสฺสาย น ฌายติ ฌายติ จ ปน.\csromanspacer{} กถํ ฌายญฺจ ปน ภนฺเต ภทฺรํ ปุริสาชานียํ ส-อินฺทา เทวา สพฺรหฺมกา สปชาปติกา อารกาว นมสฺสนฺติ\csromandash{}}

\csromangathameasure{นโม เต ปุริสาชญฺญ, นโม เต ปุริสุตฺตม. \\ ยสฺส เต นาภิชานาม, ยมฺปิ นิสฺสาย ฌายสีติ.}
\csromangathagroup{\csromangathastanza{นโม เต ปุริสาชญฺญ, นโม เต ปุริสุตฺตม. \\ ยสฺส เต นาภิชานาม, ยมฺปิ นิสฺสาย ฌายสีติ.}}

\prose{อิธ สทฺธ ภทฺรสฺส ปุริสาชานียสฺส ปถวิยํ ปถวิสญฺญา วิภูตา โหติ, อาปสฺมิํ อาโปสญฺญา วิภูตา โหติ, เตชสฺมิํ เตโชสญฺญา วิภูตา โหติ, วายสฺมิํ วาโยสญฺญา วิภูตา โหติ, อากาสานญฺจายตเน อากาสานญฺจายตนสญฺญา วิภูตา โหติ, วิญฺญาณญฺจายตเน วิญฺญาณญฺจายตนสญฺญา วิภูตา โหติ, อากิญฺจญฺญายตเน อากิญฺจญฺญายตนสญฺญา วิภูตา โหติ, เนวสญฺญานาสญฺญายตเน เนวสญฺญานาสญฺญายตนสญฺญา วิภูตา โหติ, อิธโลเก อิธโลกสญฺญา วิภูตา โหติ, ปรโลเก ปรโลกสญฺญา วิภูตา โหติ, ยมฺปิทํ ทิฏฺฐํ สุตํ มุตํ วิญฺญาตํ ปตฺตํ ปริเยสิตํ อนุวิจริตํ มนสา ตตฺราปิ สญฺญา วิภูตา โหติ.\csromanspacer{} เอวํ ฌายี โข สทฺธ ภโทฺร ปุริสาชานีโย เนว ปถวิํ นิสฺสาย ฌายติ ฯเปฯ ยมฺปิทํ ทิฏฺฐํ สุตํ มุตํ วิญฺญาตํ ปตฺตํ ปริเยสิตํ อนุวิจริตํ มนสา ตมฺปิ นิสฺสาย น ฌายติ, ฌายติ จ ปน.\csromanspacer{} เอวํ ฌายิญฺจ ปน สทฺธ ภทฺรํ ปุริสาชานียํ ส-อินฺทา เทวา สพฺรหฺมกา สปชาปติกา อารกาว นมสฺสนฺติ\csromandash{}}

\csromangathameasure{นโม เต ปุริสาชญฺญ, นโม เต ปุริสุตฺตม. \\ ยสฺส เต นาภิชานาม, ยมฺปิ นิสฺสาย ฌายสีติ.นวมํ.}
\csromangathagroup{\csromangathastanza{นโม เต ปุริสาชญฺญ, นโม เต ปุริสุตฺตม. \\ ยสฺส เต นาภิชานาม, ยมฺปิ นิสฺสาย ฌายสีติ.\csromanspacer{}\csromanspacer{}นวมํ.}}

\csromanmatikaanchor{mtk.260}
\csromantocmarkref{subsection}{9. สทฺธสุตฺต}{mtk.260}
\bookmark[dest={mtk.260},level=2]{9. สทฺธสุตฺต}
\csromanheader{2}{10. โมรนิวาปสุตฺต}
\proseitem{10}{เอกํ สมยํ ภควา ราชคเห วิหรติ โมรนิวาเป ปริพฺพาชการาเม.\csromanspacer{} ตตฺร โข ภควา ภิกฺขู อามนฺเตสิ “ภิกฺขโว”ติ. “ภทนฺเต”ติ เต ภิกฺขู ภควโต ปจฺจสฺโสสุํ.\csromanspacer{} ภควา เอตทโวจ\csromandash{}}

\prose{ตีหิ ภิกฺขเว ธมฺเมหิ สมนฺนาคโต ภิกฺขุ อจฺจนฺตนิฏฺโฐ โหติ อจฺจนฺตโยคกฺเขมี อจฺจนฺตพฺรหฺมจารี อจฺจนฺตปริโยสาโน เสฏฺโฐ เทวมนุสฺสานํ.\csromanspacer{} กตเมหิ ตีหิ? อเสเขน สีลกฺขนฺเธน อเสเขน}

\chapterheadpageread{1. นิสฺสยวคฺค}
\prose{\csromanfolio{529}สมาธิกฺขนฺเธน อเสเขน ปญฺญากฺขนฺเธน.\csromanspacer{} อิเมหิ โข ภิกฺขเว ตีหิ ธมฺเมหิ สมนฺนาคโต ภิกฺขุ อจฺจนฺตนิฏฺโฐ โหติ อจฺจนฺตโยคกฺเขมี อจฺจนฺตพฺรหฺมจารี อจฺจนฺตปริโยสาโน เสฏฺโฐ เทวมนุสฺสานํ.}

\prose{อปเรหิปิ ภิกฺขเว ตีหิ ธมฺเมหิ สมนฺนาคโต ภิกฺขุ อจฺจนฺตนิฏฺโฐ โหติ อจฺจนฺตโยคกฺเขมี อจฺจนฺตพฺรหฺมจารี อจฺจนฺตปริโยสาโน เสฏฺโฐ เทวมนุสฺสานํ.\csromanspacer{} กตเมหิ ตีหิ? อิทฺธิปาฏิหาริเยน อาเทสนาปาฏิหาริเยน อนุสาสนีปาฏิหาริเยน.\csromanspacer{} อิเมหิ โข ภิกฺขเว ตีหิ ธมฺเมหิ สมนฺนาคโต ภิกฺขุ อจฺจนฺตนิฏฺโฐ โหติ อจฺจนฺตโยคกฺเขมี อจฺจนฺตพฺรหฺมจารี อจฺจนฺตปริโยสาโน เสฏฺโฐ เทวมนุสฺสานํ.}

\prose{อปเรหิปิ ภิกฺขเว ตีหิ ธมฺเมหิ สมนฺนาคโต ภิกฺขุ อจฺจนฺตนิฏฺโฐ โหติ อจฺจนฺตโยคกฺเขมี อจฺจนฺตพฺรหฺมจารี อจฺจนฺตปริโยสาโน เสฏฺโฐ เทวมนุสฺสานํ.\csromanspacer{} กตเมหิ ตีหิ? สมฺมาทิฏฺฐิยา สมฺมาญาเณน สมฺมาวิมุตฺติยา.\csromanspacer{} อิเมหิ โข ภิกฺขเว ตีหิ ธมฺเมหิ สมนฺนาคโต ภิกฺขุ อจฺจนฺตนิฏฺโฐ โหติ อจฺจนฺตโยคกฺเขมี อจฺจนฺตพฺรหฺมจารี อจฺจนฺตปริโยสาโน เสฏฺโฐ เทวมนุสฺสานํ.}

\prose{ทฺวีหิ ภิกฺขเว ธมฺเมหิ สมนฺนาคโต ภิกฺขุ อจฺจนฺตนิฏฺโฐ โหติ อจฺจนฺตโยคกฺเขมี อจฺจนฺตพฺรหฺมจารี อจฺจนฺตปริโยสาโน เสฏฺโฐ เทวมนุสฺสานํ.\csromanspacer{} กตเมหิ ทฺวีหิ? วิชฺชาย จรเณน.\csromanspacer{} อิเมหิ โข ภิกฺขเว ทฺวีหิ ธมฺเมหิ สมนฺนาคโต ภิกฺขุ อจฺจนฺตนิฏฺโฐ โหติ อจฺจนฺตโยคกฺเขมี อจฺจนฺตพฺรหฺมจารี อจฺจนฺตปริโยสาโน เสฏฺโฐ เทวมนุสฺสานํ.\csromanspacer{} พฺรหฺมุนาเปสา ภิกฺขเว สนงฺกุมาเรน คาถา ภาสิตา\csromandash{}}

\csromangathameasure{“ขตฺติโย เสฏฺโฐ ชเนตสฺมิํ, เย โคตฺตปฏิสาริโน.}
\csromangathagroup{\csromangathastanza{“ขตฺติโย เสฏฺโฐ ชเนตสฺมิํ, เย โคตฺตปฏิสาริโน.}}

\vspace{\tipitakaparskip}
\csromancenter{วิชฺชาจรณสมฺปนฺโน, โส เสฏฺโฐ เทวมานุเส”ติ. *}
\prose{สา โข ปเนสา ภิกฺขเว สนงฺกุมาเรน คาถา ภาสิตา สุภาสิตา, โน ทุพฺภาสิตา, อตฺถสํหิตา โน อนตฺถสํหิตา, อนุมตา มยา, อหมฺปิ ภิกฺขเว เอวํ วทามิ\csromandash{}}

\gambhira{เอกาทสกนิปาตปาฬิ}
\csromanmatikaanchor{mtk.261}
\csromanmatikaanchor{mtk.262}
\csromantocmarkref{subsection}{10. โมรนิวาปสุตฺต}{mtk.261}
\bookmark[dest={mtk.261},level=2]{10. โมรนิวาปสุตฺต}
\csromantocmarkref{section}{อุทฺทานคาถา}{mtk.262}
\bookmark[dest={mtk.262},level=1]{อุทฺทานคาถา}
\csromangathameasure{“ขตฺติโย เสฏฺโฐ ชเนตสฺมิํ, เย โคตฺตปฏิสาริโน. \\ วิชฺชาจรณสมฺปนฺโน, โส เสฏฺโฐ เทวมานุเส”ติ.ทสมํ. นิสฺสยวคฺโค ปฐโม.}
\csromangathagroup{\csromangathastanza{\csromanfolio{530}“ขตฺติโย เสฏฺโฐ ชเนตสฺมิํ, เย โคตฺตปฏิสาริโน. \\ วิชฺชาจรณสมฺปนฺโน, โส เสฏฺโฐ เทวมานุเส”ติ.\csromanspacer{}\csromanspacer{}ทสมํ.\csromanspacer{} นิสฺสยวคฺโค\footnote{นิสฺสายวคฺโค (สฺยา, กํ)} ปฐโม.}}

\titlehead{\textbf{ตสฺสุทฺทานํ}}
\csromangathameasure{กิมตฺถิยา เจตนา ตโย, อุปนิสา พฺยสเนน จ. \\ เทฺว สญฺญา มนสิกาโร, สทฺโธ โมรนิวาปกนฺติ.}
\csromangathagroup{\csromangathastanza{กิมตฺถิยา เจตนา ตโย, อุปนิสา พฺยสเนน จ. \\ เทฺว สญฺญา มนสิกาโร, สทฺโธ โมรนิวาปกนฺติ.}}

\chapterhead{2. อนุสฺสติวคฺค}
\csromanheader{2}{1. ปฐมมหานามสุตฺต}
\proseitem{11}{เอกํ สมยํ ภควา สกฺเกสุ วิหรติ กปิลวตฺถุสฺมิํ นิโคฺรธาราเม.\csromanspacer{} เตน โข ปน สมเยน สมฺพหุลา ภิกฺขู ภควโต จีวรกมฺมํ กโรนฺติ “นิฏฺฐิตจีวโร ภควา เตมาสจฺจเยน จาริกํ ปกฺกมิสฺสตี”ติ.\csromanspacer{} อสฺโสสิ โข มหานาโม สกฺโก “สมฺพหุลา กิร ภิกฺขู ภควโต จีวรกมฺมํ กโรนฺติ ‘นิฏฺฐิตจีวโร ภควา เตมาสจฺจเยน จาริกํ ปกฺกมิสฺสตี’ติ”.}

\prose{อถ โข มหานาโม สกฺโก เยน ภควา เตนุปสงฺกมิ, อุปสงฺกมิตฺวา ภควนฺตํ อภิวาเทตฺวา เอกมนฺตํ นิสีทิ, เอกมนฺตํ นิสินฺโน โข มหานาโม สกฺโก ภควนฺตํ เอตทโวจ “สุตํ เมตํ ภนฺเต สมฺพหุลา กิร ภิกฺขู ภควโต จีวรกมฺมํ กโรนฺติ, นิฏฺฐิตจีวโร ภควา เตมาสจฺจเยน จาริกํ ปกฺกมิสฺสตีติ.\csromanspacer{} เตสํ โน ภนฺเต นานาวิหาเรหิ วิหรตํ เกนสฺส\footnote{เกน (สฺยา, กํ)} วิหาเรน วิหาตพฺพนฺ”ติ.}

\prose{สาธุ สาธุ มหานาม, เอตํ โข มหานาม ตุมฺหากํ ปติรูปํ กุลปุตฺตานํ, ยํ ตุเมฺห ตถาคตํ อุปสงฺกมิตฺวา ปุจฺเฉยฺยาถ “เตสํ โน ภนฺเต นานาวิหาเรหิ วิหรตํ เกนสฺส วิหาเรน วิหาตพฺพนฺ”ติ.\csromanspacer{} สทฺโธ โข มหานาม อาราธโก โหติ, โน อสฺสทฺโธ.\csromanspacer{} อารทฺธวีริโย อาราธโก โหติ, โน กุสีโต.\csromanspacer{} อุปฏฺฐิตสฺสติ อาราธโก}

\vspace{\tipitakaparskip}
\csromancenter{อนุสฺสติวคฺค โหติ, โน มุฏฺฐสฺสติ.\csromanspacer{} สมาหิโต อาราธโก โหติ, โน อสมาหิโต.\csromanspacer{} ปญฺญวา อาราธโก โหติ, โน ทุปฺปญฺโญ.\csromanspacer{} อิเมสุ โข ตฺวํ มหานาม ปญฺจสุ ธมฺเมสุ ปติฏฺฐาย ฉ ธมฺเม อุตฺตริ\footnote{อุตฺตริํ (สี, สฺยา, กํ, อิ)} ภาเวยฺยาสิ. * อิธ ตฺวํ มหานาม ตถาคตํ อนุสฺสเรยฺยาสิ “อิติปิ โส ภควา อรหํ สมฺมาสมฺพุทฺโธ วิชฺชาจรณสมฺปนฺโน สุคโต โลกวิทู อนุตฺตโร ปุริสทมฺมสารถิ สตฺถา เทวมนุสฺสานํ พุทฺโธ ภควา”ติ.\csromanspacer{} ยสฺมิํ มหานาม สมเย อริยสาวโก ตถาคตํ อนุสฺสรติ, เนวสฺส ตสฺมิํ สมเย ราคปริยุฏฺฐิตํ จิตฺตํ โหติ, น โทสปริยุฏฺฐิตํ จิตฺตํ โหติ, น โมหปริยุฏฺฐิตํ จิตฺตํ โหติ, อุชุคตเมวสฺส ตสฺมิํสมเย จิตฺตํ โหติ ตถาคตํ อารพฺภ.\csromanspacer{} อุชุคตจิตฺโต โข ปน มหานาม อริยสาวโก ลภติ อตฺถเวทํ, ลภติ ธมฺมเวทํ, ลภติ ธมฺมูปสํหิตํ ปาโมชฺชํ, ปมุทิตสฺส ปีติ ชายติ, ปีติมนสฺส กาโย ปสฺสมฺภติ, ปสฺสทฺธกาโย สุขํ เวทิยติ, สุขิโน จิตฺตํ สมาธิยติ.\csromanspacer{} อยํ วุจฺจติ มหานาม อริยสาวโก วิสมคตาย ปชาย สมปฺปตฺโต วิหรติ, สพฺยาปชฺชาย ปชาย อพฺยาปชฺโช วิหรติ, ธมฺมโสตสมาปนฺโน พุทฺธานุสฺสติํ ภาเวติ.}
\prose{\csromanfolio{531}ปุน จปรํ ตฺวํ มหานาม ธมฺมํ อนุสฺสเรยฺยาสิ “สฺวากฺขาโต ภควตา ธมฺโม สนฺทิฏฺฐิโก อกาลิโก เอหิปสฺสิโก โอปเนยฺยิโก\footnote{โอปนยิโก (สี, สฺยา, กํ, อิ)} ปจฺจตฺตํ วิทิตพฺโพ วิญฺญูหี”ติ.\csromanspacer{} ยสฺมิํ มหานาม สมเย อริยสาวโก ธมฺมํ อนุสฺสรติ, เนวสฺส ตสฺมิํ สมเย ราคปริยุฏฺฐิตํ จิตฺตํ โหติ, น โทสปริยุฏฺฐิตํ จิตฺตํ โหติ, น โมหปริยุฏฺฐิตํ จิตฺตํ โหติ, อุชุคตเมวสฺส ตสฺมิํ สมเย จิตฺตํ โหติ ธมฺมํ อารพฺภ.\csromanspacer{} อุชุคตจิตฺโต โข ปน มหานาม อริยสาวโก ลภติ อตฺถเวทํ, ลภติ ธมฺมเวทํ, ลภติ ธมฺมูปสํหิตํ ปาโมชฺชํ, ปมุทิตสฺส ปีติ ชายติ, ปีติมนสฺส กาโย ปสฺสมฺภติ, ปสฺสทฺธกาโย สุขํ เวทิยติ, สุขิโน จิตฺตํ สมาธิยติ.\csromanspacer{} อยํ วุจฺจติ มหานาม อริยสาวโก วิสมคตาย ปชาย สมปฺปตฺโต วิหรติ, สพฺยาปชฺชาย ปชาย อพฺยาปชฺโช วิหรติ, ธมฺมโสตสมาปนฺโน ธมฺมานุสฺสติํ ภาเวติ.}

\prose{ปุน จปรํ ตฺวํ มหานาม สํฆํ อนุสฺสเรยฺยาสิ “สุปฺปฏิปนฺโน ภควโต สาวกสํโฆ อุชุปฺปฏิปนฺโน ภควโต สาวกสํโฆ ญายปฺปฏิปนฺโน}

\gambhira{เอกาทสกนิปาตปาฬิ}
\prose{\csromanfolio{532}ภควโต สาวกสํโฆ สามีจิปฺปฏิปนฺโน ภควโต สาวกสํโฆ, ยทิทํ จตฺตาริ ปุริสยุคานิ อฏฺฐ ปุริสปุคฺคลา, เอส ภควโต สาวกสํโฆ อาหุเนยฺโย ปาหุเนยฺโย ทกฺขิเณยฺโย อญฺชลิกรณีโย อนุตฺตรํ ปุญฺญกฺเขตฺตํ โลกสฺสา”ติ.\csromanspacer{} ยสฺมิํ มหานาม สมเย อริยสาวโก สํฆํ อนุสฺสรติ, เนวสฺส ตสฺมิํ สมเย ราคปริยุฏฺฐิตํ จิตฺตํ โหติ, น โทสปริยุฏฺฐิตํ จิตฺตํ โหติ, น โมหปริยุฏฺฐิตํ จิตฺตํ โหติ, อุชุคตเมวสฺส ตสฺมิํ สมเย จิตฺตํ โหติ สํฆํ อารพฺภ.\csromanspacer{} อุชุคตจิตฺโต โข ปน มหานาม อริยสาวโก ลภติ อตฺถเวทํ, ลภติ ธมฺมเวทํ, ลภติ ธมฺมูปสํหิตํ ปาโมชฺชํ, ปมุทิตสฺส ปีติ ชายติ, ปีติมนสฺส กาโย ปสฺสมฺภติ, ปสฺสทฺธกาโย สุขํ เวทิยติ, สุขิโน จิตฺตํ สมาธิยติ.\csromanspacer{} อยํ วุจฺจติ มหานาม อริยสาวโก วิสมคตาย ปชาย สมปฺปตฺโต วิหรติ, สพฺยาปชฺชาย ปชาย อพฺยาปชฺโช วิหรติ, ธมฺมโสตสมาปนฺโน สํฆานุสฺสติํ ภาเวติ.}

\prose{ปุน จปรํ ตฺวํ มหานาม อตฺตโน สีลานิ อนุสฺสเรยฺยาสิ อขณฺฑานิ อจฺฉิทฺทานิ อสพลานิ อกมฺมาสานิ ภุชิสฺสานิ วิญฺญุปฺปสตฺถานิ อปรามฏฺฐานิ สมาธิสํวตฺตนิกานิ.\csromanspacer{} ยสฺมิํ มหานาม สมเย อริยสาวโก สีลํ อนุสฺสรติ, เนวสฺส ตสฺมิํ สมเย ราคปริยุฏฺฐิตํ จิตฺตํ โหติ, น โทสปริยุฏฺฐิตํ จิตฺตํ โหติ, น โมหปริยุฏฺฐิตํ จิตฺตํ โหติ, อุชุคตเมวสฺส ตสฺมิํ สมเย จิตฺตํ โหติ สีลํ อารพฺภ.\csromanspacer{} อุชุคตจิตฺโต โข ปน มหานาม อริยสาวโก ลภติ อตฺถเวทํ, ลภติ ธมฺมเวทํ, ลภติ ธมฺมูปสํหิตํ ปาโมชฺชํ, ปมุทิตสฺส ปีติ ชายติ, ปีติมนสฺส กาโย ปสฺสมฺภติ, ปสฺสทฺธกาโย สุขํ เวทิยติ, สุขิโน จิตฺตํ สมาธิยติ.\csromanspacer{} อยํ วุจฺจติ มหานาม อริยสาวโก วิสมคตาย ปชาย สมปฺปตฺโต วิหรติ, สพฺยาปชฺชาย ปชาย อพฺยาปชฺโช วิหรติ, ธมฺมโสตสมาปนฺโน สีลานุสฺสติํ ภาเวติ.}

\prose{ปุน จปรํ ตฺวํ มหานาม อตฺตโน จาคํ อนุสฺสเรยฺยาสิ “ลาภา วต เม สุลทฺธํ วต เม, โยหํ มจฺเฉรมลปริยุฏฺฐิตาย ปชาย วิคตมลมจฺเฉเรน เจตสา อคารํ อชฺฌาวสามิ มุตฺตจาโค ปยตปาณิ โวสคฺครโต ยาจโยโค ทานสํวิภาครโต”ติ.\csromanspacer{} ยสฺมิํ}

\csromanmatikaanchor{mtk.263}
\csromantocmarknopage{chapter}{2. อนุสฺสติวคฺค}{mtk.263}
\bookmark[dest={mtk.263},level=0]{2. อนุสฺสติวคฺค}
\vspace{\tipitakaparskip}
\csromancenter{อนุสฺสติวคฺค มหานาม สมเย อริยสาวโก จาคํ อนุสฺสรติ, เนวสฺส ตสฺมิํ สมเย ราคปริยุฏฺฐิตํ จิตฺตํ โหติ, น โทสปริยุฏฺฐิตํ จิตฺตํ โหติ, น โมหปริยุฏฺฐิตํ จิตฺตํ โหติ, อุชุคตเมวสฺส ตสฺมิํ สมเย จิตฺตํ โหติ จาคํ อารพฺภ.\csromanspacer{} อุชุคตจิตฺโต โข ปน มหานาม อริยสาวโก ลภติ อตฺถเวทํ, ลภติ ธมฺมเวทํ, ลภติ ธมฺมูปสํหิตํ ปาโมชฺชํ, ปมุทิตสฺส ปีติ ชายติ, ปีติมนสฺส กาโย ปสฺสมฺภติ, ปสฺสทฺธกาโย สุขํ เวทิยติ, สุขิโน จิตฺตํ สมาธิยติ.\csromanspacer{} อยํ วุจฺจติ มหานาม อริยสาวโก วิสมคตาย ปชาย สมปฺปตฺโต วิหรติ, สพฺยาปชฺชาย ปชาย อพฺยาปชฺโช วิหรติ, ธมฺมโสตสมาปนฺโน จาคานุสฺสติํ ภาเวติ.}
\csromanmatikaanchor{mtk.264}
\csromantocmarkref{subsection}{1. ปฐมมหานามสุตฺต}{mtk.264}
\bookmark[dest={mtk.264},level=2]{1. ปฐมมหานามสุตฺต}
\prose{\csromanfolio{533}ปุน จปรํ ตฺวํ มหานาม เทวตา อนุสฺสเรยฺยาสิ “สนฺติ เทวา จาตุมหาราชิกา, สนฺติ เทวา ตาวติํสา, สนฺติ เทวา ยามา, สนฺติ เทวา ตุสิตา, สนฺติ เทวา นิมฺมานรติโน, สนฺติ เทวา ปรนิมฺมิตวสวตฺติโน, สนฺติ เทวา พฺรหฺมกายิกา, สนฺติ เทวา ตตุตฺตริ.\csromanspacer{} ยถารูปาย สทฺธาย สมนฺนาคตา ตา เทวตา อิโต จุตา ตตฺถูปปนฺนา, มยฺหมฺปิ ตถารูปา สทฺธา สํวิชฺชติ.\csromanspacer{} ยถารูเปน สีเลน สมนฺนาคตา ตา เทวตา อิโต จุตา ตตฺถูปปนฺนา, มยฺหมฺปิ ตถารูปํ สีลํ สํวิชฺชติ.\csromanspacer{} ยถารูเปน สุเตน สมนฺนาคตา ตา เทวตา อิโต จุตา ตตฺถูปปนฺนา, มยฺหมฺปิ ตถารูปํ สุตํ สํวิชฺชติ.\csromanspacer{} ยถารูเปน จาเคน สมนฺนาคตา ตา เทวตา อิโต จุตา ตตฺถูปปนฺนา, มยฺหมฺปิ ตถารูโป จาโค สํวิชฺชติ.\csromanspacer{} ยถารูปาย ปญฺญาย สมนฺนาคตา ตา เทวตา อิโต จุตา ตตฺถูปปนฺนา, มยฺหมฺปิ ตถารูปา ปญฺญา สํวิชฺชตี”ติ.\csromanspacer{} ยสฺมิํ มหานาม สมเย อริยสาวโก อตฺตโน จ ตาสญฺจ เทวตานํ สทฺธญฺจ สีลญฺจ สุตญฺจ จาคญฺจ ปญฺญญฺจ อนุสฺสรติ, เนวสฺส ตสฺมิํ สมเย ราคปริยุฏฺฐิตํ จิตฺตํ โหติ, น โทสปริยุฏฺฐิตํ จิตฺตํ โหติ, น โมหปริยุฏฺฐิตํ จิตฺตํ โหติ, อุชุคตเมวสฺส ตสฺมิํ สมเย จิตฺตํ โหติ เทวตา อารพฺภ.\csromanspacer{} อุชุคตจิตฺโต โข ปน มหานาม อริยสาวโก ลภติ อตฺถเวทํ, ลภติ ธมฺมเวทํ, ลภติ ธมฺมูปสํหิตํ ปาโมชฺชํ, ปมุทิตสฺส ปีติ ชายติ, ปีติมนสฺส กาโย ปสฺสมฺภติ, ปสฺสทฺธกาโย สุขํ เวทิยติ, สุขิโน จิตฺตํ สมาธิยติ.\csromanspacer{} อยํ วุจฺจติ มหานาม อริยสาวโก วิสมคตาย ปชาย}

\gambhira{เอกาทสกนิปาตปาฬิ}
\prose{\csromanfolio{534}สมปฺปตฺโต วิหรติ, สพฺยาปชฺชาย ปชาย อพฺยาปชฺโช วิหรติ, ธมฺมโสตสมาปนฺโน เทวตานุสฺสติํ ภาเวตีติ.\csromanspacer{}\csromanspacer{}ปฐมํ.}

\csromanheader{2}{2. ทุติยมหานามสุตฺต}
\proseitem{12}{เอกํ สมยํ ภควา สกฺเกสุ วิหรติ กปิลวตฺถุสฺมิํ นิโคฺรธาราเม.\csromanspacer{} เตน โข ปน สมเยน มหานาโม สกฺโก คิลานา วุฏฺฐิโต โหติ อจิรวุฏฺฐิโต เคลญฺญา.\csromanspacer{} เตน โข ปน สมเยน สมฺพหุลา ภิกฺขู ภควโต จีวรกมฺมํ กโรนฺติ “นิฏฺฐิตจีวโร ภควา เตมาสจฺจเยน จาริกํ ปกฺกมิสฺสตี”ติ.}

\prose{อสฺโสสิ โข มหานาโม สกฺโก “สมฺพหุลา กิร ภิกฺขู ภควโต จีวรกมฺมํ กโรนฺติ ‘นิฏฺฐิตจีวโร ภควา เตมาสจฺจเยน จาริกํ ปกฺกมิสฺสตี’ติ”.\csromanspacer{} อถ โข มหานาโม สกฺโก เยน ภควา เตนุปสงฺกมิ, อุปสงฺกมิตฺวา ภควนฺตํ อภิวาเทตฺวา เอกมนฺตํ นิสีทิ, เอกมนฺตํ นิสินฺโน โข มหานาโม สกฺโก ภควนฺตํ เอตทโวจ “สุตํ เมตํ ภนฺเต ‘สมฺพหุลา กิร ภิกฺขู ภควโต จีวรกมฺมํ กโรนฺติ นิฏฺฐิตจีวโร ภควา เตมาสจฺจเยน จาริกํ ปกฺกมิสฺสตี’ติ.\csromanspacer{} เตสํ โน ภนฺเต นานาวิหาเรหิ วิหรตํ เกนสฺส วิหาเรน วิหาตพฺพนฺ”ติ.}

\prose{สาธุ สาธุ มหานาม, เอตํ โข มหานาม ตุมฺหากํ ปติรูปํ กุลปุตฺตานํ, ยํ ตุเมฺห ตถาคตํ อุปสงฺกมิตฺวา ปุจฺเฉยฺยาถ “เตสํ โน ภนฺเต นานาวิหาเรหิ วิหรตํ เกนสฺส วิหาเรน วิหาตพฺพนฺ”ติ.\csromanspacer{} สทฺโธ โข มหานาม อาราธโก โหติ, โน อสฺสทฺโธ.\csromanspacer{} อารทฺธวีริโย อาราธโก โหติ, โน กุสีโต.\csromanspacer{} อุปฏฺฐิตสฺสติ อาราธโก โหติ, โน มุฏฺฐสฺสติ.\csromanspacer{} สมาหิโต อาราธโก โหติ, โน อสมาหิโต.\csromanspacer{} ปญฺญวา อาราธโก โหติ, โน ทุปฺปญฺโญ.\csromanspacer{} อิเมสุ โข ตฺวํ มหานาม ปญฺจสุ ธมฺเมสุ ปติฏฺฐาย ฉ ธมฺเม อุตฺตริ ภาเวยฺยาสิ.}

\prose{\csromansymbolfootnote{*}{อํ 2. 252 ปิฏฺเฐปิ.}อิธ ตฺวํ มหานาม ตถาคตํ อนุสฺสเรยฺยาสิ “อิติปิ โส ภควา ฯเปฯ สตฺถา เทวมนุสฺสานํ พุทฺโธ ภควา”ติ.\csromanspacer{} ยสฺมิํ มหานาม สมเย อริยสาวโก ตถาคตํ อนุสฺสรติ, เนวสฺส ตสฺมิํ สมเย ราคปริยุฏฺฐิตํ จิตฺตํ โหติ, น โทสปริยุฏฺฐิตํ จิตฺตํ โหติ, น โมหปริยุฏฺฐิตํ จิตฺตํ โหติ, อุชุคตเมวสฺส ตสฺมิํ สมเย จิตฺตํ โหติ}

\vspace{\tipitakaparskip}
\csromancenter{อนุสฺสติวคฺค ตถาคตํ อารพฺภ.\csromanspacer{} อุชุคตจิตฺโต โข ปน มหานาม อริยสาวโก ลภติ อตฺถเวทํ, ลภติ ธมฺมเวทํ, ลภติ ธมฺมูปสํหิตํ ปาโมชฺชํ, ปมุทิตสฺส ปีติ ชายติ, ปีติมนสฺส กาโย ปสฺสมฺภติ, ปสฺสทฺธกาโย สุขํ เวทิยติ, สุขิโน จิตฺตํ สมาธิยติ.\csromanspacer{} อิมํ โข ตฺวํ มหานาม พุทฺธานุสฺสติํ คจฺฉนฺโตปิ ภาเวยฺยาสิ, ฐิโตปิ ภาเวยฺยาสิ, นิสินฺโนปิ ภาเวยฺยาสิ, สยาโนปิ ภาเวยฺยาสิ, กมฺมนฺตํ อธิฏฺฐหนฺโตปิ ภาเวยฺยาสิ, ปุตฺตสมฺพาธสยนํ อชฺฌาวสนฺโตปิ ภาเวยฺยาสิ.}
\prose{\csromanfolio{535}ปุน จปรํ ตฺวํ มหานาม ธมฺมํ อนุสฺสเรยฺยาสิ ฯเปฯ สํฆํ อนุสฺสเรยฺยาสิ ฯเปฯ อตฺตโน สีลํ อนุสฺสเรยฺยาสิ ฯเปฯ อตฺตโน จาคํ อนุสฺสเรยฺยาสิ ฯเปฯ เทวตา อนุสฺสเรยฺยาสิ “สนฺติ เทวา จาตุมหาราชิกา ฯเปฯ สนฺติ เทวา ตตุตฺตริ.\csromanspacer{} ยถารูปาย สทฺธาย สมนฺนาคตา ตา เทวตา อิโต จุตา ตตฺถูปปนฺนา, มยฺหมฺปิ ตถารูปา สทฺธา สํวิชฺชติ.\csromanspacer{} ยถารูเปน สีเลน.\csromanspacer{} สุเตน.\csromanspacer{} จาเคน.\csromanspacer{} ปญฺญาย สมนฺนาคตา ตา เทวตา อิโต จุตา ตตฺถูปปนฺนา, มยฺหมฺปิ ตถารูปา ปญฺญา สํวิชฺชตี”ติ.\csromanspacer{} ยสฺมิํ มหานาม สมเย อริยสาวโก อตฺตโน จ ตาสญฺจ เทวตานํ สทฺธญฺจ สีลญฺจ สุตญฺจ จาคญฺจ ปญฺญญฺจ อนุสฺสรติ, เนวสฺส ตสฺมิํ สมเย ราคปริยุฏฺฐิตํ จิตฺตํ โหติ, น โทสปริยุฏฺฐิตํ จิตฺตํ โหติ, น โมหปริยุฏฺฐิตํ จิตฺตํ โหติ, อุชุคตเมวสฺส ตสฺมิํ สมเย จิตฺตํ โหติ เทวตา อารพฺภ.\csromanspacer{} อุชุคตจิตฺโต โข ปน มหานาม อริยสาวโก ลภติ อตฺถเวทํ, ลภติ ธมฺมเวทํ, ลภติ ธมฺมูปสํหิตํ ปาโมชฺชํ, ปมุทิตสฺส ปีติ ชายติ, ปีติมนสฺส กาโย ปสฺสมฺภติ, ปสฺสทฺธกาโย สุขํ เวทิยติ, สุขิโน จิตฺตํ สมาธิยติ.\csromanspacer{} อิมํ โข ตฺวํ มหานาม เทวตานุสฺสติํ คจฺฉนฺโตปิ ภาเวยฺยาสิ, ฐิโตปิ ภาเวยฺยาสิ, นิสินฺโนปิ ภาเวยฺยาสิ, สยาโนปิ ภาเวยฺยาสิ, กมฺมนฺตํ อธิฏฺฐหนฺโตปิ ภาเวยฺยาสิ, ปุตฺตสมฺพาธสยนํ อชฺฌาวสนฺโตปิ ภาเวยฺยาสีติ.\csromanspacer{}\csromanspacer{}ทุติยํ.}

\csromanheader{2}{3. นนฺทิยสุตฺต}
\proseitem{13}{เอกํ สมยํ ภควา สกฺเกสุ วิหรติ กปิลวตฺถุสฺมิํ นิโคฺรธาราเม.\csromanspacer{} เตน โข ปน สมเยน ภควา สาวตฺถิยํ วสฺสาวาสํ อุปคนฺตุกาโม โหติ\footnote{อโหสิ (ก)}.}

\gambhira{เอกาทสกนิปาตปาฬิ}
\prose{\csromanfolio{536}อสฺโสสิ โข นนฺทิโย สกฺโก “ภควา กิร สาวตฺถิยํ วสฺสาวาสํ อุปคนฺตุกาโม”ติ.\csromanspacer{} อถ โข นนฺทิยสฺส สกฺกสฺส “เอตทโหสิ ‘ยํนูนาหมฺปิ สาวตฺถิยํ วสฺสาวาสํ อุปคจฺเฉยฺยํ, ตตฺถ กมฺมนฺตญฺเจว อธิฏฺฐหิสฺสามิ, ภควนฺตญฺจ ลจฺฉามิ กาเลน กาลํ ทสฺสนายา’ติ”.}

\prose{อถ โข ภควา สาวตฺถิยํ วสฺสาวาสํ อุปคจฺฉิ\footnote{อุปคญฺฉิ (สี, อิ)}.\csromanspacer{} นนฺทิโยปิ โข สกฺโก สาวตฺถิยํ วสฺสาวาสํ อุปคจฺฉิ.\csromanspacer{} ตตฺถ กมฺมนฺตญฺเจว อธิฏฺฐาสิ\footnote{อธิฏฺฐาย (สฺยา), อธิฏฺฐาติ (ก)}, ภควนฺตญฺจ ลภิ\footnote{ลจฺฉติ (สฺยา, ก)} กาเลน กาลํ ทสฺสนาย.\csromanspacer{} เตน โข ปน สมเยน สมฺพหุลา ภิกฺขู ภควโต จีวรกมฺมํ กโรนฺติ “นิฏฺฐิตจีวโร ภควา เตมาสจฺจเยน จาริกํ ปกฺกมิสฺสตี”ติ.}

\csromanmatikaanchor{mtk.265}
\csromantocmarkref{subsection}{2. ทุติยมหานามสุตฺต}{mtk.265}
\bookmark[dest={mtk.265},level=2]{2. ทุติยมหานามสุตฺต}
\prose{อสฺโสสิ โข นนฺทิโย สกฺโก “สมฺพหุลา กิร ภิกฺขู ภควโต จีวรกมฺมํ กโรนฺติ ‘นิฏฺฐิตจีวโร ภควา เตมาสจฺจเยน จาริกํ ปกฺกมิสฺสตี’ติ”.\csromanspacer{} อถ โข นนฺทิโย สกฺโก เยน ภควา เตนุปสงฺกมิ, อุปสงฺกมิตฺวา ภควนฺตํ อภิวาเทตฺวา เอกมนฺตํ นิสีทิ, เอกมนฺตํ นิสินฺโน โข นนฺทิโย สกฺโก ภควนฺตํ เอตทโวจ “สุตํ เมตํ ภนฺเต ‘สมฺพหุลา กิร ภิกฺขู ภควโต จีวรกมฺมํ กโรนฺติ, นิฏฺฐิตจีวโร ภควา เตมาสจฺจเยน จาริกํ ปกฺกมิสฺสตี’ติ.\csromanspacer{} เตสํ โน ภนฺเต นานาวิหาเรหิ วิหรตํ เกนสฺส วิหาเรน วิหาตพฺพนฺ”ติ.}

\prose{สาธุ สาธุ นนฺทิย, เอตํ โข นนฺทิย ตุมฺหากํ ปติรูปํ กุลปุตฺตานํ, ยํ ตุเมฺห ตถาคตํ อุปสงฺกมิตฺวา ปุจฺเฉยฺยาถ “เตสํ โน ภนฺเต นานาวิหาเรหิ วิหรตํ เกนสฺส วิหาเรน วิหาตพฺพนฺ”ติ.\csromanspacer{} สทฺโธ โข นนฺทิย อาราธโก โหติ, โน อสฺสทฺโธ.\csromanspacer{} สีลวา อาราธโก โหติ, โน ทุสฺสีโล.\csromanspacer{} อารทฺธวีริโย อาราธโก โหติ, โน กุสีโต.\csromanspacer{} อุปฏฺฐิตสฺสติ อาราธโก โหติ, โน มุฏฺฐสฺสติ.\csromanspacer{} สมาหิโต อาราธโก โหติ, โน อสมาหิโต.\csromanspacer{} ปญฺญวา อาราธโก โหติ, โน ทุปฺปญฺโญ.\csromanspacer{} อิเมสุ โข เต นนฺทิย ฉสุ ธมฺเมสุ ปติฏฺฐาย ปญฺจสุ ธมฺเมสุ อชฺฌตฺตํ สติ อุปฏฺฐาเปตพฺพา.}

\prose{อิธ ตฺวํ นนฺทิย ตถาคตํ อนุสฺสเรยฺยาสิ “อิติปิ โส ภควา อรหํ สมฺมาสมฺพุทฺโธ วิชฺชาจรณสมฺปนฺโน สุคโต โลกวิทู อนุตฺตโร}

\vspace{\tipitakaparskip}
\csromancenter{อนุสฺสติวคฺค ปุริสทมฺมสารถิ สตฺถา เทวมนุสฺสานํ พุทฺโธ ภควา”ติ.\csromanspacer{} อิติ โข เต นนฺทิย ตถาคตํ อารพฺภ อชฺฌตฺตํ สติ อุปฏฺฐาเปตพฺพา.}
\prose{\csromanfolio{537}ปุน จปรํ ตฺวํ นนฺทิย ธมฺมํ อนุสฺสเรยฺยาสิ “สฺวากฺขาโต ภควตา ธมฺโม สนฺทิฏฺฐิโก อกาลิโก เอหิปสฺสิโก โอปเนยฺยิโก ปจฺจตฺตํ เวทิตพฺโพ วิญฺญูหี”ติ.\csromanspacer{} อิติ โข เต นนฺทิย ธมฺมํ อารพฺภ อชฺฌตฺตํ สติ อุปฏฺฐาเปตพฺพา.}

\prose{ปุน จปรํ ตฺวํ นนฺทิย กลฺยาณมิตฺเต อนุสฺสเรยฺยาสิ “ลาภา วต เม, สุลทฺธํ วต เม, ยสฺส เม กลฺยาณมิตฺตา อนุกมฺปกา อตฺถกามา โอวาทกา อนุสาสกา”ติ.\csromanspacer{} อิติ โข เต นนฺทิย กลฺยาณมิตฺเต อารพฺภ อชฺฌตฺตํ สติ อุปฏฺฐาเปตพฺพา.}

\prose{ปุน จปรํ ตฺวํ นนฺทิย อตฺตโน จาคํ อนุสฺสเรยฺยาสิ “ลาภา วต เม, สุลทฺธํ วต เม, โยหํ มจฺเฉรมลปริยุฏฺฐิตาย ปชาย วิคตมลมจฺเฉเรน เจตสา อคารํ อชฺฌาวสามิ มุตฺตจาโค ปยตปาณิ โวสคฺครโต ยาจโยโค ทานสํวิภาครโต”ติ.\csromanspacer{} อิติ โข เต นนฺทิย จาคํ อารพฺภ อชฺฌตฺตํ สติ อุปฏฺฐาเปตพฺพา.}

\csromanmatikaanchor{mtk.266}
\csromantocmarkref{subsection}{3. นนฺทิยสุตฺต}{mtk.266}
\bookmark[dest={mtk.266},level=2]{3. นนฺทิยสุตฺต}
\prose{ปุน จปรํ ตฺวํ นนฺทิย เทวตา อนุสฺสเรยฺยาสิ “ยา เทวตา อติกฺกมฺเมว กพฬีการาหารภกฺขานํ\footnote{กพฬิํการภกฺขานํ (สี), กพฬีการภกฺขานํ (สฺยา, กํ, อิ)} เทวตานํ สหพฺยตํ อญฺญตรํ มโนมยํ กายํ อุปปนฺนา, ตา กรณียํ อตฺตโน น สมนุปสฺสนฺติ กตสฺส วา ปติจยํ”.\csromanspacer{} เสยฺยถาปิ นนฺทิย ภิกฺขุ อสมยวิมุตฺโต กรณียํ อตฺตโน น สมนุปสฺสติ กตสฺส วา ปติจยํ.\csromanspacer{} เอวเมวํ โข นนฺทิย ยา ตา เทวตา อติกฺกมฺเมว กพฬีการาหารภกฺขานํ เทวตานํ สหพฺยตํ อญฺญตรํ มโนมยํ กายํ อุปปนฺนา, ตา กรณียํ อตฺตโน น สมนุปสฺสนฺติ กตสฺส วา ปติจยํ.\csromanspacer{} อิติ โข เต นนฺทิย เทวตา อารพฺภ อชฺฌตฺตํ สติ อุปฏฺฐาเปตพฺพา.}

\prose{อิเมหิ โข นนฺทิย เอกาทสหิ ธมฺเมหิ สมนฺนาคโต อริยสาวโก ปชหเตว ปาปเก อกุสเล ธมฺเม น อุปาทิยติ.\csromanspacer{} เสยฺยถาปิ นนฺทิย กุมฺโภ นิกฺกุชฺโช\footnote{นิกุชฺโช (ก)} วมเตว อุทกํ, โน วนฺตํ ปจฺจาวมติ\footnote{ปจฺจามสติ (สฺยา)}.\csromanspacer{} เสยฺยถาปิ วา ปน นนฺทิย สุกฺเข ติณทาเย อคฺคิ มุตฺโต ฑหญฺเญว คจฺฉติ, โน ทฑฺฒํ}

\gambhira{เอกาทสกนิปาตปาฬิ}
\prose{\csromanfolio{538}ปจฺจุทาวตฺตติ.\csromanspacer{} เอวเมวํ โข นนฺทิย อิเมหิ เอกาทสหิ ธมฺเมหิ สมนฺนาคโต อริยสาวโก ปชหเตว ปาปเก อกุสเล ธมฺเม น อุปาทิยตีติ.\csromanspacer{}\csromanspacer{}ตติยํ.}

\csromanheader{2}{4. สุภูติสุตฺต}
\proseitem{14}{อถ โข อายสฺมา สุภูติ สทฺเธน ภิกฺขุนา สทฺธิํ เยน ภควา เตนุปสงฺกมิ, อุปสงฺกมิตฺวา ภควนฺตํ อภิวาเทตฺวา เอกมนฺตํ นิสีทิ, เอกมนฺตํ นิสินฺนํ โข อายสฺมนฺตํ สุภูติํ ภควา เอตทโวจ “โก นามายํ\footnote{โก นาโม อยํ (สี, ก), โก นาม อยํ (สฺยา, กํ)} สุภูติ ภิกฺขู”ติ.\csromanspacer{} สทฺโธ นามายํ ภนฺเต ภิกฺขุ สุทตฺตสฺส\footnote{สทฺธสฺส (สี, สฺยา, กํ, อิ)} อุปาสกสฺส ปุตฺโต สทฺธา อคารสฺมา อนคาริยํ ปพฺพชิโตติ.}

\prose{กจฺจิ ปนายํ สุภูติ สทฺโธ ภิกฺขุ สุทตฺตสฺส อุปาสกสฺส ปุตฺโต สทฺธา อคารสฺมา อนคาริยํ ปพฺพชิโต สนฺทิสฺสติ สทฺธาปทาเนสูติ.\csromanspacer{} เอตสฺส ภควา กาโล เอตสฺส สุคต กาโล, ยํ ภควา สทฺธสฺส สทฺธาปทานานิ ภาเสยฺย, อิทานาหํ ชานิสฺสามิ “ยทิ วา อยํ ภิกฺขุ สนฺทิสฺสติ สทฺธาปทาเนสุ, ยทิ วา โน”ติ.}

\prose{เตน หิ สุภูติ สุณาหิ สาธุกํ มนสิ กโรหิ, ภาสิสฺสามีติ. “เอวํ ภนฺเต”ติ โข อายสฺมา สุภูติ ภควโต ปจฺจสฺโสสิ.\csromanspacer{} ภควา เอตทโวจ\csromandash{}}

\prose{อิธ สุภูติ ภิกฺขุ สีลวา โหติ ปาติโมกฺขสํวรสํวุโต วิหรติ อาจารโคจรสมฺปนฺโน, อณุมตฺเตสุ วชฺเชสุ ภยทสฺสาวี สมาทาย สิกฺขติ สิกฺขาปเทสุ.\csromanspacer{} ยมฺปิ สุภูติ ภิกฺขุ สีลวา โหติ ฯเปฯ สมาทาย สิกฺขติ สิกฺขาปเทสุ, อิทมฺปิ สุภูติ สทฺธสฺส สทฺธาปทานํ โหติ.}

\prose{ปุน จปรํ สุภูติ ภิกฺขุ พหุสฺสุโต โหติ สุตธโร สุตสนฺนิจโย, เย เต ธมฺมา อาทิกลฺยาณา มชฺเฌกลฺยาณา ปริโยสานกลฺยาณา สาตฺถํ สพฺยญฺชนํ เกวลปริปุณฺณํ ปริสุทฺธํ พฺรหฺมจริยํ อภิวทนฺติ, ตถารูปาสฺส ธมฺมา พหุสฺสุตา โหนฺติ ธาตา วจสา ปริจิตา มนสานุเปกฺขิตา ทิฏฺฐิยา สุปฺปฏิวิทฺธา, ยมฺปิ สุภูติ ภิกฺขุ}

\vspace{\tipitakaparskip}
\csromancenter{อนุสฺสติวคฺค พหุสฺสุโต โหติ ฯเปฯ ทิฏฺฐิยา สุปฺปฏิวิทฺธา, อิทมฺปิ สุภูติ สทฺธสฺส สทฺธาปทานํ โหติ.}
\prose{\csromanfolio{539}ปุน จปรํ สุภูติ ภิกฺขุ กลฺยาณมิตฺโต โหติ กลฺยาณสหาโย กลฺยาณสมฺปวงฺโก.\csromanspacer{} ยมฺปิ สุภูติ ภิกฺขุ กลฺยาณมิตฺโต โหติ กลฺยาณสหาโย กลฺยาณสมฺปวงฺโก, อิทมฺปิ สุภูติ สทฺธสฺส สทฺธาปทานํ โหติ.}

\prose{ปุน จปรํ สุภูติ ภิกฺขุ สุวโจ โหติ โสวจสฺสกรเณหิ ธมฺเมหิ สมนฺนาคโต ขโม ปทกฺขิณคฺคาหี อนุสาสนิํ.\csromanspacer{} ยมฺปิ สุภูติ ภิกฺขุ สุวโจ โหติ โสวจสฺสกรเณหิ ธมฺเมหิ สมนฺนาคโต ขโม ปทกฺขิณคฺคาหี อนุสาสนิํ, อิทมฺปิ สุภูติ สทฺธสฺส สทฺธาปทานํ โหติ.}

\prose{ปุน จปรํ สุภูติ ภิกฺขุ ยานิ ตานิ สพฺรหฺมจารีนํ อุจฺจาวจานิ กิํกรณียานิ, ตตฺร ทกฺโข โหติ อนลโส, ตตฺรุปายาย วีมํสาย สมนฺนาคโต อลํ กาตุํ อลํ สํวิธาตุํ.\csromanspacer{} ยมฺปิ สุภูติ ภิกฺขุ ยานิ ตานิ สพฺรหฺมจารีนํ อุจฺจาวจานิ กิํกรณียานิ, ตตฺร ทกฺโข โหติ อนลโส, ตตฺรุปายาย วีมํสาย สมนฺนาคโต อลํ กาตุํ อลํ สํวิธาตุํ, อิทมฺปิ สุภูติ สทฺธสฺส สทฺธาปทานํ โหติ.}

\prose{ปุน จปรํ สุภูติ ภิกฺขุ ธมฺมกาโม โหติ ปิยสมุทาหาโร อภิธมฺเม อภิวินเย อุฬารปาโมชฺโช.\csromanspacer{} ยมฺปิ สุภูติ ภิกฺขุ ธมฺมกาโม โหติ ปิยสมุหาโร อภิธมฺเม อภิวินเย อุฬารปาโมชฺโช, อิทมฺปิ สุภูติ สทฺธสฺส สทฺธาปทานํ โหติ.}

\prose{ปุน จปรํ สุภูติ ภิกฺขุ อารทฺธวีริโย วิหรติ อกุสลานํ ธมฺมานํ ปหานาย กุสลานํ ธมฺมานํ อุปสมฺปทาย, ถามวา ทฬฺหปรกฺกโม อนิกฺขิตฺตธุโร กุสเลสุ ธมฺเมสุ.\csromanspacer{} ยมฺปิ สุภูติ ภิกฺขุ อารทฺธวีริโย วิหรติ อกุสลานํ ธมฺมานํ ปหานาย กุสลานํ ธมฺมานํ อุปสมฺปทาย, ถามวา ทฬฺหปรกฺกโม อนิกฺขิตฺตธุโร กุสเลสุ ธมฺเมสุ, อิทมฺปิ สุภูติ สทฺธสฺส สทฺธาปทานํ โหติ.}

\csromanmatikaanchor{mtk.267}
\csromantocmarkref{subsection}{4. สุภูติสุตฺต}{mtk.267}
\bookmark[dest={mtk.267},level=2]{4. สุภูติสุตฺต}
\prose{ปุน จปรํ สุภูติ ภิกฺขุ จตุนฺนํ ฌานานํ อาภิเจตสิกานํ ทิฏฺฐธมฺมสุขวิหารานํ นิกามลาภี โหติ อกิจฺฉลาภี อกสิรลาภี.}

\gambhira{เอกาทสกนิปาตปาฬิ}
\prose{\csromanfolio{540}ยมฺปิ สุภูติ ภิกฺขุ จตุนฺนํ ฌานานํ อาภิเจตสิกานํ ทิฏฺฐธมฺมสุขวิหารานํ นิกามลาภี โหติ อกิจฺฉลาภี อกสิรลาภี, อิทมฺปิ สุภูติ สทฺธสฺส สทฺธาปทานํ โหติ.}

\prose{ปุน จปรํ สุภูติ ภิกฺขุ อเนกวิหิตํ ปุพฺเพนิวาสํ อนุสฺสรติ.\csromanspacer{} เสยฺยถิทํ, เอกมฺปิ ชาติํ เทฺวปิ ชาติโย ติสฺโสปิ ชาติโย จตสฺโสปิ ชาติโย ปญฺจปิ ชาติโย ทสปิ ชาติโย วีสมฺปิ ชาติโย ติํสมฺปิ ชาติโย จตฺตารีสมฺปิ ชาติโย ปญฺญาสมฺปิ ชาติโย ชาติสตมฺปิ ชาติสหสฺสมฺปิ ชาติสตสหสฺสมฺปิ อเนเกปิ สํวฏฺฏกปฺเป อเนเกปิ วิวฏฺฏกปฺเป อเนเกปิ สํวฏฺฏวิวฏฺฏกปฺเป อมุตฺราสิํ “เอวํนาโม เอวํโคตฺโต เอวํวณฺโณ เอวมาหาโร เอวํสุขทุกฺขปฺปฏิสํเวที เอวมายุปริยนฺโต.\csromanspacer{} โส ตโต จุโต อมุตฺร อุทปาทิํ, ตตฺราปาสิํ เอวํนาโม เอวํโคตฺโต เอวํวณฺโณ เอวมาหาโร เอวํสุขทุกฺขปฺปฏิสํเวที เอวมายุปริยนฺโต.\csromanspacer{} โส ตโต จุโต อิธูปปนฺโน”ติ อิติ สาการํ ส-อุทฺเทสํ อเนกวิหิตํ ปุพฺเพนิวาสํ อนุสฺสรติ.\csromanspacer{} ยมฺปิ สุภูติ ภิกฺขุ อเนกวิหิตํ ปุพฺเพนิวาสํ อนุสฺสรติ.\csromanspacer{} เสยฺยถิทํ, เอกมฺปิ ชาติํ เทฺวปิ ชาติโย ฯเปฯ อิติ สาการํ ส-อุทฺเทสํ อเนกวิหิตํ ปุพฺเพนิวาสํ อนุสฺสรติ, อิทมฺปิ สุภูติ สทฺธสฺส สทฺธาปทานํ โหติ.}

\prose{ปุน จปรํ สุภูติ ภิกฺขุ ทิพฺเพน จกฺขุนา วิสุทฺเธน อติกฺกนฺตมานุสเกน สตฺเต ปสฺสติ จวมาเน อุปปชฺชมาเน หีเน ปณีเต สุวณฺเณ ทุพฺพณฺเณ สุคเต ทุคฺคเต, ยถากมฺมูปเค สตฺเต ปชานาติ “อิเม วต โภนฺโต สตฺตา กายทุจฺจริเตน สมนฺนาคตา วจีทุจฺจริเตน สมนฺนาคตา มโนทุจฺจริเตน สมนฺนาคตา, อริยานํ อุปวาทกา มิจฺฉาทิฏฺฐิกา มิจฺฉาทิฏฺฐิกมฺมสมาทานา.\csromanspacer{} เต กายสฺส เภทา ปรํ มรณา อปายํ ทุคฺคติํ วินิปาตํ นิรยํ อุปปนฺนา.\csromanspacer{} อิเม วา ปน โภนฺโต สตฺตา กายสุจริเตน สมนฺนาคตา วจีสุจริเตน สมนฺนาคตา มโนสุจริเตน สมนฺนาคตา, อริยานํ อนุปวาทกา สมฺมาทิฏฺฐิกา สมฺมาทิฏฺฐิกมฺมสมาทานา.\csromanspacer{} เต กายสฺส เภทา ปรํ มรณา สุคติํ สคฺคํ โลกํ อุปปนฺนา”ติ.\csromanspacer{} อิติ ทิพฺเพน จกฺขุนา วิสุทฺเธน อติกฺกนฺตมานุสเกน สตฺเต ปสฺสติ จวมาเน อุปปชฺชมาเน หีเน ปณีเต สุวณฺเณ ทุพฺพณฺเณ สุคเต ทุคฺคเต, ยถากมฺมูปเค สตฺเต ปชานาติ.\csromanspacer{} ยมฺปิ สุภูติ ภิกฺขุ ทิพฺเพน จกฺขุนา วิสุทฺเธน ฯเปฯ ยถากมฺมูปเค สตฺเต ปชานาติ, อิทมฺปิ สุภูติ สทฺธสฺส สทฺธาปทานํ โหติ.}

\chapterheadpageread{2. อนุสฺสติวคฺค}
\prose{\csromanfolio{541}ปุน จปรํ สุภูติ ภิกฺขุ อาสวานํ ขยา อนาสวํ เจโตวิมุตฺติํ ปญฺญาวิมุตฺติํ ทิฏฺเฐว ธมฺเม สยํ อภิญฺญา สจฺฉิกตฺวา อุปสมฺปชฺช วิหรติ.\csromanspacer{} ยมฺปิ สุภูติ ภิกฺขุ อาสวานํ ขยา ฯเปฯ สจฺฉิกตฺวา อุปสมฺปชฺช วิหรติ, อิทมฺปิ สุภูติ สทฺธสฺส สทฺธาปทานํ โหตีติ.}

\prose{เอวํ วุตฺเต อายสฺมา สุภูติ ภควนฺตํ เอตทโวจ\csromandash{}ยานิมานิ ภนฺเต ภควตา สทฺธสฺส สทฺธาปทานานิ ภาสิตานิ, สํวิชฺชนฺติ ตานิ อิมสฺส ภิกฺขุโน, อยญฺจ ภิกฺขุ เอเตสุ สนฺทิสฺสติ.}

\prose{อยํ ภนฺเต ภิกฺขุ สีลวา โหติ, ปาติโมกฺขสํวรสํวุโต วิหรติ อาจารโคจรสมฺปนฺโน, อณุมตฺเตสุ วชฺเชสุ ภยทสฺสาวี สมาทาย สิกฺขติ สิกฺขาปเทสุ.}

\prose{อยํ ภนฺเต ภิกฺขุ พหุสฺสุโต โหติ สุตธโร สุตสนฺนิจโย, เย เต ธมฺมา อาทิกลฺยาณา มชฺเฌกลฺยาณา ปริโยสานกลฺยาณา สาตฺถํ สพฺยญฺชนํ เกวลปริปุณฺณํ ปริสุทฺธํ พฺรหฺมจริยํ อภิวทนฺติ.\csromanspacer{} ตถารูปาสฺส ธมฺมา พหุสฺสุตา โหนฺติ ธาตา วจสา ปริจิตา มนสานุเปกฺขิตา ทิฏฺฐิยา สุปฺปฏิวิทฺธา.}

\prose{อยํ ภนฺเต ภิกฺขุ กลฺยาณมิตฺโต โหติ กลฺยาณสหาโย กลฺยาณสมฺปวงฺโก.}

\prose{อยํ ภนฺเต ภิกฺขุ สุวโจ โหติ ฯเปฯ อนุสฺสาสนิํ.}

\prose{อยํ ภนฺเต ภิกฺขุ ยานิ ตานิ สพฺรหฺมจารีนํ อุจฺจาวจานิ กิํกรณียานิ, ตตฺถ ทกฺโข โหติ อนลโส.\csromanspacer{} ตตฺรุปายาย วีมํสาย สมนฺนาคโต อลํ กาตุํ อลํ สํวิธาตุํ.}

\prose{อยํ ภนฺเต ภิกฺขุ ธมฺมกาโม โหติ ปิยสมุทาหาโร อภิธมฺเม อภิวินเย อุฬารปาโมชฺโช.}

\prose{อยํ ภนฺเต ภิกฺขุ อารทฺธวีริโย วิหรติ ฯเปฯ ถามวา ทฬฺหปรกฺกโม อนิกฺขิตฺตธุโร กุสเลสุ ธมฺเมสุ.}

\prose{อยํ ภนฺเต ภิกฺขุ จตุนฺนํ ฌานานํ อาภิเจตสิกานํ ทิฏฺฐธมฺมสุขวิหารานํ นิกามลาภี โหติ อกิจฺฉลาภี อกสิรลาภี.}

\gambhira{เอกาทสกนิปาตปาฬิ}
\prose{\csromanfolio{542}อยํ ภนฺเต ภิกฺขุ อเนกวิหิตํ ปุพฺเพนิวาสํ อนุสฺสรติ.\csromanspacer{} เสยฺยถิทํ, เอกมฺปิ ชาติํ เทฺวปิ ชาติโย ฯเปฯ อิติ สาการํ ส-อุทฺเทสํ อเนกวิหิตํ ปุพฺเพนิวาสํ อนุสฺสรติ.}

\prose{อยํ ภนฺเต ภิกฺขุ ทิพฺเพน จกฺขุนา วิสุทฺเธน อติกฺกนฺตมานุสเกน ฯเปฯ ยถากมฺมูปเค สตฺเต ปชานาติ.}

\prose{อยํ ภนฺเต ภิกฺขุ อาสวานํ ขยา ฯเปฯ สจฺฉิกตฺวา อุปสมฺปชฺช วิหรติ, ยานิมานิ ภนฺเต ภควตา สทฺธสฺส สทฺธาปทานานิ ภาสิตานิ, สํวิชฺชนฺติ ตานิ อิมสฺส ภิกฺขุโน, อยญฺจ ภิกฺขุ เอเตสุ สนฺทิสฺสตีติ.}

\prose{สาธุ สาธุ สุภูติ, เตน หิ ตฺวํ สุภูติ อิมินา จ สทฺเธน ภิกฺขุนา สทฺธิํ วิหเรยฺยาสิ.\csromanspacer{} ยทา จ ตฺวํ สุภูติ อากงฺเขยฺยาสิ ตถาคตํ ทสฺสนาย.\csromanspacer{} อิมินา สทฺเธน ภิกฺขุนา สทฺธิํ อุปสงฺกเมยฺยาสิ ตถาคตํ ทสฺสนายาติ.\csromanspacer{}\csromanspacer{}จตุตฺถํ.}

\csromanheader{2}{5. เมตฺตาสุตฺต}
\proseitem{15}{\csromansymbolfootnote{*}{ขุ 9. 314; ขุ 11. 195 ปิฏฺเฐสุ.}เมตฺตาย ภิกฺขเว เจโตวิมุตฺติยา อาเสวิตาย ภาวิตาย พหุลีกตาย ยานีกตาย วตฺถุกตาย อนุฏฺฐิตาย ปริจิตาย สุสมารทฺธาย เอกาทสานิสํสา ปาฏิกงฺขา\csromandash{}}

\prose{กตเม เอกาทส? สุขํ สุปติ, สุขํ ปฏิพุชฺฌติ, น ปาปกํ สุปินํ ปสฺสติ, มนุสฺสานํ ปิโย โหติ, อมนุสฺสานํ ปิโย โหติ, เทวตา รกฺขนฺติ, นาสฺส อคฺคิ วา วิสํ วา สตฺถํ วา กมติ, ตุวฏํ จิตฺตํ สมาธิยติ, มุขวณฺโณ วิปฺปสีทติ, อสมฺมูโฬฺห กาลํ กโรติ, อุตฺตริ อปฺปฏิวิชฺฌนฺโต พฺรหฺมโลกูปโค โหติ.\csromanspacer{} เมตฺตาย ภิกฺขเว เจโตวิมุตฺติยา อาเสวิตาย ภาวิตาย พหุลีกตาย ยานีกตาย วตฺถุกตาย อนุฏฺฐิตาย ปริจิตาย สุสมารทฺธาย อิเม เอกาทสานิสํสา ปาฏิกงฺขาติ.\csromanspacer{}\csromanspacer{}ปญฺจมํ.}

\csromanheader{2}{6. อฏฺฐกนาครสุตฺต}
\proseitem{16}{เอกํ สมยํ อายสฺมา อานนฺโท เวสาลิยํ วิหรติ เพลุวคามเก\footnote{เวฬุวคามเก (สฺยา, กํ, ก)}, เตน โข ปน สมเยน ทสโม คหปติ อฏฺฐกนาคโร ปาฏลิปุตฺตํ อนุปฺปตฺโต โหติ เกนจิเทว กรณีเยน.}

\chapterheadpageread{2. อนุสฺสติวคฺค}
\prose{\csromanfolio{543}อถ โข ทสโม คหปติ อฏฺฐกนาคโร เยน กุกฺกุฏาราโม, เยน อญฺญตโร ภิกฺขุ เตนุปสงฺกมิ, อุปสงฺกมิตฺวา ตํ ภิกฺขุํ เอตทโวจ “กหํ นุ โข ภนฺเต อายสฺมา อานนฺโท เอตรหิ วิหรติ, ทสฺสนกามา หิ มยํ ภนฺเต อายสฺมนฺตํ อานนฺทนฺ”ติ.\csromanspacer{} เอโส คหปติ อายสฺมา อานนฺโท เวสาลิยํ วิหรติ เพลุวคามเกติ.}

\prose{อถ โข ทสโม คหปติ อฏฺฐกนาคโร ปาฏลิปุตฺเต ตํ กรณียํ ตีเรตฺวา เยน เวสาลี เพลุวคามโก, เยนายสฺมา อานนฺโท เตนุปสงฺกมิ, อุปสงฺกมิตฺวา อายสฺมนฺตํ อานนฺทํ อภิวาเทตฺวา เอกมนฺตํ นิสีทิ, เอกมนฺตํ นิสินฺโน โข ทสโม คหปติ อฏฺฐกนาคโร อายสฺมนฺตํ อานนฺทํ เอตทโวจ “อตฺถิ นุ โข ภนฺเต อานนฺท เตน ภควตา ชานตา ปสฺสตา อรหตา สมฺมาสมฺพุทฺเธน เอกธมฺโม สมฺมทกฺขาโต, ยตฺถ ภิกฺขุโน อปฺปมตฺตสฺส อาตาปิโน ปหิตตฺตสฺส วิหรโต อวิมุตฺตํ วา จิตฺตํ วิมุจฺจติ, อปริกฺขีณา วา อาสวา ปริกฺขยํ คจฺฉนฺติ, อนนุปฺปตฺตํ วา อนุตฺตรํ โยคกฺเขมํ อนุปาปุณาตี”ติ.\csromanspacer{} อตฺถิ โข คหปติ เตน ภควตา ชานตา ปสฺสตา อรหตา สมฺมาสมฺพุทฺเธน เอกธมฺโม สมฺมทกฺขาโต, ยตฺถ ภิกฺขุโน อปฺปมตฺตสฺส อาตาปิโน ปหิตตฺตสฺส วิหรโต อวิมุตฺตํ วา จิตฺตํ วิมุจฺจติ, อปริกฺขีณา วา อาสวา ปริกฺขยํ คจฺฉนฺติ, อนนุปฺปตฺตํ วา อนุตฺตรํ โยคกฺเขมํ อนุปาปุณาตีติ.}

\prose{กตโม ปน ภนฺเต อานนฺท เตน ภควตา ชานตา ปสฺสตา อรหตา สมฺมาสมฺพุทฺเธน เอกธมฺโม สมฺมทกฺขาโต, ยตฺถ ภิกฺขุโน อปฺปมตฺตสฺส อาตาปิโน ปหิตตฺตสฺส วิหรโต อวิมุตฺตํ วา จิตฺตํ วิมุจฺจติ, อปริกฺขีณา วา อาสวา ปริกฺขยํ คจฺฉนฺติ, อนนุปฺปตฺตํ วา อนุตฺตรํ โยคกฺเขมํ อนุปาปุณาตีติ.\csromanspacer{} อิธ คหปติ ภิกฺขุ วิวิจฺเจว กาเมหิ วิวิจฺจ อกุสเลหิ ธมฺเมหิ สวิตกฺกํ สวิจารํ วิเวกชํ ปีติสุขํ ปฐมํ ฌานํ อุปสมฺปชฺช วิหรติ.\csromanspacer{} โส อิติ ปฏิสญฺจิกฺขติ “อิทมฺปิ โข ปฐมํ ฌานํ อภิสงฺขตํ อภิสญฺเจตยิตํ, ยํ โข ปน กิญฺจิ อภิสงฺขตํ อภิสญฺเจตยิตํ, ตทนิจฺจํ นิโรธธมฺมนฺ”ติ ปชานาติ, โส ตตฺถ ฐิโต อาสวานํ ขยํ ปาปุณาติ, โน เจ อาสวานํ ขยํ ปาปุณาติ, เตเนว ธมฺมราเคน ตาย ธมฺมนนฺทิยา ปญฺจนฺนํ โอรมฺภาคิยานํ สํโยชนานํ ปริกฺขยา โอปปาติโก โหติ, ตตฺถ ปรินิพฺพายี อนาวตฺติธมฺโม ตสฺมา โลกา.\csromanspacer{} อยมฺปิ โข คหปติ เตน ภควตา ชานตา ปสฺสตา อรหตา}

\gambhira{เอกาทสกนิปาตปาฬิ}
\prose{\csromanfolio{544}สมฺมาสมฺพุทฺเธน เอกธมฺโม สมฺมทกฺขาโต, ยตฺถ ภิกฺขุโน อปฺปมตฺตสฺส อาตาปิโน ปหิตตฺตสฺส วิหรโต อวิมุตฺตํ วา จิตฺตํ วิมุจฺจติ, อปริกฺขีณา วา อาสวา ปริกฺขยํ คจฺฉนฺติ, อนนุปฺปตฺตํ วา อนุตฺตรํ โยคกฺเขมํ อนุปาปุณาติ.}

\prose{ปุน จปรํ คหปติ ภิกฺขุ วิตกฺกวิจารานํ วูปสมา อชฺฌตฺตํ สมฺปสาทนํ เจตโส เอโกทิภาวํ อวิตกฺกํ อวิจารํ สมาธิชํ ปีติสุขํ ทุติยํ ฌานํ ฯเปฯ ตติยํ ฌานํ ฯเปฯ จตุตฺถํ ฌานํ อุปสมฺปชฺช วิหรติ, โส อิติ ปฏิสญฺจิกฺขติ “อิทมฺปิ โข จตุตฺถํ ฌานํ อภิสงฺขตํ อภิสญฺเจตยิตํ, ยํ โข ปน กิญฺจิ อภิสงฺขตํ อภิสญฺเจตยิตํ, ตทนิจฺจํ นิโรธธมฺมนฺ”ติ ปชานาติ.\csromanspacer{} โส ตตฺถ ฐิโต อาสวานํ ขยํ ปาปุณาติ, โน เจ อาสวานํ ขยํ ปาปุณาติ, เตเนว ธมฺมราเคน ตาย ธมฺมนนฺทิยา ปญฺจนฺนํ โอรมฺภาคิยานํ สํโยชนานํ ปริกฺขยา โอปปาติโก โหติ, ตตฺถ ปรินิพฺพายี อนาวตฺติธมฺโม ตสฺมา โลกา.\csromanspacer{} อยมฺปิ โข คหปติ เตน ภควตา ชานตา ปสฺสตา อรหตา สมฺมาสมฺพุทฺเธน เอกธมฺโม สมฺมทกฺขาโต, ยตฺถ ภิกฺขุโน อปฺปมตฺตสฺส อาตาปิโน ปหิตตฺตสฺส วิหรโต อวิมุตฺตํ วา จิตฺตํ วิมุจฺจติ, อปริกฺขีณา วา อาสวา ปริกฺขยํ คจฺฉนฺติ, อนนุปฺปตฺตํ วา อนุตฺตรํ โยคกฺเขมํ อนุปาปุณาติ.}

\csromanmatikaanchor{mtk.268}
\csromantocmarkref{subsection}{5. เมตฺตาสุตฺต}{mtk.268}
\bookmark[dest={mtk.268},level=2]{5. เมตฺตาสุตฺต}
\prose{ปุน จปรํ คหปติ ภิกฺขุ เมตฺตาสหคเตน เจตสา เอกํ ทิสํ ผริตฺวา วิหรติ.\csromanspacer{} ตถา ทุติยํ.\csromanspacer{} ตถา ตติยํ.\csromanspacer{} ตถา จตุตฺถํ.\csromanspacer{} อิติ อุทฺธมโธ ติริยํ สพฺพธิ สพฺพตฺตตาย สพฺพาวนฺตํ โลกํ เมตฺตาสหคเตน เจตสา เอกํ ทิสํ ผริตฺวา วิหรติ วิปุเลน มหคฺคเตน อปฺปมาเณน อเวเรน อพฺยาปชฺเชน ผริตฺวา วิหรติ, โส อิติ ปฏิสญฺจิกฺขติ “อยมฺปิ โข เมตฺตา เจโตวิมุตฺติ อภิสงฺขตา อภิสญฺเจตยิตา, ยํ โข ปน กิญฺจิ อภิสงฺขตํ อภิสญฺเจตยิตํ, ตทนิจฺจํ นิโรธธมฺมนฺ”ติ ปชานาติ, โส ตตฺถ ฐิโต อาสวานํ ขยํ ปาปุณาติ, โน เจ อาสวานํ ขยํ ปาปุณาติ.\csromanspacer{} เตเนว ธมฺมราเคน ตาย ธมฺมนนฺทิยา ปญฺจนฺนํ โอรมฺภาคิยานํ สํโยชนานํ ปริกฺขยา โอปปาติโก โหติ.\csromanspacer{} ตตฺถ ปรินิพฺพายี อนาวตฺติธมฺโม ตสฺมา โลกา.\csromanspacer{} อยมฺปิ โข คหปติ เตน ภควตา ชานตา ฯเปฯ อนนุปฺปตฺตํ วา อนุตฺตรํ โยคกฺเขมํ อนุปาปุณาติ.}

\prose{ปุน จปรํ คหปติ ภิกฺขุ กรุณาสหคเตน เจตสา ฯเปฯ มุทิตา สหคเตน เจตสา ฯเปฯ อุเปกฺขาสหคเตน เจตสา เอกํ ทิสํ ผริตฺวา}

\vspace{\tipitakaparskip}
\csromancenter{อนุสฺสติวคฺค วิหรติ.\csromanspacer{} ตถา ทุติยํ.\csromanspacer{} ตถา ตติยํ.\csromanspacer{} ตถา จตุตฺถํ.\csromanspacer{} อิติ อุทฺธมโธ ติริยํ สพฺพธิ สพฺพตฺตตาย สพฺพาวนฺตํ โลกํ อุเปกฺขาสหคเตน เจตสา วิปุเลน มหคฺคเตน อปฺปมาเณน อเวเรน อพฺยาปชฺเชน ผริตฺวา วิหรติ.\csromanspacer{} โส อิติ ปฏิสญฺจิกฺขติ “อยมฺปิ โข อุเปกฺขาเจโตวิมุตฺติ อภิสงฺขตา อภิสญฺเจตยิตา, ยํ โข ปน กิญฺจิ อภิสงฺขตํ อภิสญฺเจตยิตํ, ตทนิจฺจํ นิโรธธมฺมนฺ”ติ ปชานาติ.\csromanspacer{} โส ตตฺถ ฐิโต อาสวานํ ขยํ ปาปุณาติ.\csromanspacer{} โน เจ อาสวานํ ขยํ ปาปุณาติ, เตเนว ธมฺมราเคน ตาย ธมฺมนนฺทิยา ปญฺจนฺนํ โอรมฺภาคิยานํ สํโยชนานํ ปริกฺขยา โอปปาติโก โหติ ตตฺถ ปรินิพฺพายี อนาวตฺติธมฺโม ตสฺมา โลกา.\csromanspacer{} อยมฺปิ โข คหปติ เตน ภควตา ชานตา ฯเปฯ อนนุปฺปตฺตํ วา อนุตฺตรํ โยคกฺเขมํ อนุปาปุณาติ.}
\csromanmatikaanchor{mtk.269}
\csromantocmarkref{subsection}{6. อฏฺฐกนาครสุตฺต}{mtk.269}
\bookmark[dest={mtk.269},level=2]{6. อฏฺฐกนาครสุตฺต}
\prose{\csromanfolio{545}ปุน จปรํ คหปติ ภิกฺขุ สพฺพโส รูปสญฺญานํ สมติกฺกมา ปฏิฆสญฺญานํ อตฺถงฺคมา นานตฺตสญฺญานํ อมนสิการา “อนนฺโต อากาโส”ติ อากาสานญฺจายตนํ อุปสมฺปชฺช วิหรติ.\csromanspacer{} โส อิติ ปฏิสญฺจิกฺขติ “อยมฺปิ โข อากาสานญฺจายตนสมาปตฺติ อภิสงฺขตา อภิสญฺเจตยิตา, ยํ โข ปน กิญฺจิ อภิสงฺขตํ อภิสญฺเจตยิตํ, ตทนิจฺจํ นิโรธธมฺมนฺ”ติ ปชานาติ.\csromanspacer{} โส ตตฺถ ฐิโต อาสวานํ ขยํ ปาปุณาติ.\csromanspacer{} โน เจ อาสวานํ ขยํ ปาปุณาติ, เตเนว ธมฺมราเคน ตาย ธมฺมนนฺทิยา ปญฺจนฺนํ โอรมฺภาคิยานํ สํโยชนานํ ปริกฺขยา โอปปาติโก โหติ ตตฺถ ปรินิพฺพายี อนาวตฺติธมฺโม ตสฺมา โลกา.\csromanspacer{} อยมฺปิ โข คหปติ เตน ภควตา ชานตาฯเปฯ อนนุปฺปตฺตํ วา อนุตฺตรํ โยคกฺเขมํ อนุปาปุณาติ.}

\prose{ปุน จปรํ คหปติ ภิกฺขุ สพฺพโส อากาสานญฺจายตนํ สมติกฺกมฺม “อนนฺตํ วิญฺญาณนฺ”ติ วิญฺญาณญฺจายตนํ อุปสมฺปชฺช วิหรติ ฯเปฯ สพฺพโส วิญฺญาณญฺจายตนํ สมติกฺกมฺม “นตฺถิ กิญฺจี”ติ อากิญฺจญฺญายตนํ อุปสมฺปชฺชวิหรติ.\csromanspacer{} โส อิติ ปฏิสญฺจิกฺขติ “อยมฺปิ โข อากิญฺจญฺญายตนสมาปตฺติ อภิสงฺขตา อภิสญฺเจตยิตา, ยํ โข ปน กิญฺจิ อภิสงฺขตํ อภิสญฺเจตยิตํ, ตทนิจฺจํ นิโรธธมฺมนฺ”ติ ปชานาติ.\csromanspacer{} โส ตตฺถ ฐิโต อาสวานํ ขยํ ปาปุณาติ.\csromanspacer{} โน เจ อาสวานํ ขยํ ปาปุณาติ, เตเนว ธมฺมราเคน ตาย ธมฺมนนฺทิยา ปญฺจนฺนํ โอรมฺภาคิยานํ สํโยชนานํ ปริกฺขยา โอปปาติโก โหติ ตตฺถ ปรินิพฺพายี อนาวตฺติธมฺโม ตสฺมา โลกา.}

\gambhira{เอกาทสกนิปาตปาฬิ}
\prose{\csromanfolio{546}อยมฺปิ โข คหปติ เตน ภควตา ชานตา ฯเปฯ อนนุปฺปตฺตํ วา อนุตฺตรํ โยคกฺเขมํ อนุปาปุณาตีติ.}

\prose{เอวํ วุตฺเต ทสโม คหปติ อฏฺฐกนาคโร อายสฺมนฺตํ อานนฺทํ เอตทโวจ “เสยฺยถาปิ ภนฺเต อานนฺท ปุริโส เอกํ นิธิมุขํ คเวสนฺโต สกิเทว\footnote{สพฺพตฺถปิ เอวเมว ทิสฺสติ.} เอกาทส นิธิมุขานิ อธิคจฺเฉยฺย.\csromanspacer{} เอวเมวํ โข อหํ ภนฺเต เอกํ อมตทฺวารํ คเวสนฺโต สกิเทว เอกาทส อมตทฺวารานิ\footnote{เอกาทสนฺนํ อมตทฺวารานํ (สพฺพตฺถ) ม 2. 15 ปิฏฺเฐ ปสฺสิตพฺพํ.} อลตฺถํ เสวนาย\footnote{สวนาย (สฺยา) สี-อิ-มชฺฌิมปณฺณาสกทุติยสุตฺเตปิ, ภาวนาย (ม 2. 15 ปิฏฺเฐ.)}.\csromanspacer{} เสยฺยถาปิ ภนฺเต ปุริสสฺส อคารํ เอกาทส ทฺวารํ.\csromanspacer{} โส ตสฺมิํ อคาเร อาทิตฺเต เอกเมเกนปิ ทฺวาเรน สกฺกุเณยฺย อตฺตานํ โสตฺถิํ กาตุํ.\csromanspacer{} เอวเมวํ โข อหํ ภนฺเต อิเมสํ เอกาทสนฺนํ อมตทฺวารานํ เอกเมเกนปิ อมตทฺวาเรน สกฺกุณิสฺสามิ อตฺตานํ โสตฺถิํ กาตุํ.\csromanspacer{} อิเม หิ นาม ภนฺเต อญฺญติตฺถิยา อาจริยสฺส อาจริยธนํ ปริเยสิสฺสนฺติ, กิํ\footnote{กิมงฺคํ(ม 2. 15)} ปนาหํ อายสฺมโต อานนฺทสฺส ปูชํ น กริสฺสามี”ติ.}

\prose{อถ โข ทสโม คหปติ อฏฺฐกนาคโร เวสาลิกญฺจ ปาฏลิปุตฺตกญฺจ ภิกฺขุสํฆํ สนฺนิปาตาเปตฺวา ปณีเตน ขาทนีเยน โภชนีเยน สหตฺถา สนฺตปฺเปสิ สมฺปวาเรสิ.\csromanspacer{} เอกเมกญฺจ ภิกฺขุํ ปจฺเจกํ ทุสฺสยุเคน อจฺฉาเทสิ อายสฺมนฺตญฺจ อานนฺทํ ติจีวเรน, อายสฺมโต อานนฺทสฺส ปญฺจสตํ วิหารํ การาเปสีติ.\csromanspacer{}\csromanspacer{}ฉฏฺฐํ.}

\csromanheader{2}{7. โคปาลสุตฺต}
\proseitem{17}{เอกาทสหิ ภิกฺขเว องฺเคหิ สมนฺนาคโต โคปาลโก อภพฺโพ โคคณํ ปริหริตุํ ผาติํ กาตุํ\footnote{ผาติกตฺตุํ (สี), ผาติกาตุํ (สฺยา, อิ)}.\csromanspacer{} กตเมหิ เอกาทสหิ? อิธ ภิกฺขเว โคปาลโก น รูปญฺญู โหติ, น ลกฺขณกุสโล โหติ, น อาสาฏิกํ หาเรตา\footnote{สาเฏตา (สี, สฺยา, อิ)} โหติ, น วณํ ปฏิจฺฉาเทตา โหติ, น ธูมํ กตฺตา โหติ, น ติตฺถํ ชานาติ, น ปีตํ ชานาติ, น วีถิํ ชานาติ, น โคจรกุสโล โหติ, อนวเสสโทหี จ โหติ, เย เต อุสภา โคปิตโร โคปริณายกา, เต น อติเรกปูชาย ปูเชตา}

\vspace{\tipitakaparskip}
\csromancenter{อนุสฺสติวคฺค โหติ.\csromanspacer{} อิเมหิ โข ภิกฺขเว เอกาทสหิ องฺเคหิ สมนฺนาคโต โคปาลโก อภพฺโพ โคคณํ ปริหริตุํ ผาติํ กาตุํ.}
\prose{\csromanfolio{547}เอวเมวํ โข ภิกฺขเว เอกาทสหิ ธมฺเมหิ สมนฺนาคโต ภิกฺขุ อภพฺโพ อิมสฺมิํ ธมฺมวินเย วุทฺธิํ วิรูฬฺหิํ เวปุลฺลํ อาปชฺชิตุํ.\csromanspacer{} กตเมหิ เอกาทสหิ? อิธ ภิกฺขเว ภิกฺขุ น รูปญฺญู โหติ, น ลกฺขณกุสโล โหติ, น อาสาฏิกํ หาเรตา โหติ, น วณํ ปฏิจฺฉาเทตา โหติ, น ธูมํ กตฺตา โหติ, น ติตฺถํ ชานาติ, น ปีตํ ชานาติ, น วีถิํ ชานาติ, น โคจรกุสโล โหติ, อนวเสสโทหี จ โหติ, เย เต ภิกฺขู เถรา รตฺตญฺญู จิรปพฺพชิตา สํฆปิตโร สํฆปริณายกา, เต น อติเรกปูชาย ปูเชตา โหติ.}

\prose{กถญฺจ ภิกฺขเว ภิกฺขุ น รูปญฺญู โหติ? อิธ ภิกฺขเว ภิกฺขุ “ยํ กิญฺจิ รูปํ ( )\footnote{(สพฺพํ รูปํ) ม 1. 282 ปิฏฺเฐ. ( ) กตฺถจิ ทิสฺสติ.} จตฺตาริ มหาภูตานิ จตุนฺนญฺจ มหาภูตานํ อุปาทายรูปนฺ”ติ ยถาภูตํ นปฺปชานาติ.\csromanspacer{} เอวํ โข ภิกฺขเว ภิกฺขุ น รูปญฺญู โหติ. (1)}

\prose{กถญฺจ ภิกฺขเว ภิกฺขุ น ลกฺขณกุสโล โหติ? อิธ ภิกฺขเว ภิกฺขุ “กมฺมลกฺขโณ พาโล กมฺมลกฺขโณ ปณฺฑิโต”ติ ยถาภูตํ นปฺปชานาติ.\csromanspacer{} เอวํ โข ภิกฺขเว ภิกฺขุ น ลกฺขณกุสโล โหติ. (2)}

\prose{กถญฺจ ภิกฺขเว ภิกฺขุ น อาสาฏิกํ หาเรตา โหติ? อิธ ภิกฺขเว ภิกฺขุ อุปฺปนฺนํ กามวิตกฺกํ อธิวาเสติ นปฺปชหติ น วิโนเทติ, น พฺยนฺตีกโรติ น อนภาวํ คเมติ.\csromanspacer{} อุปฺปนฺนํ พฺยาปาทวิตกฺกํ.\csromanspacer{} อุปฺปนฺนํ วิหิํสาวิตกฺกํ.\csromanspacer{} อุปฺปนฺนุปฺปนฺเน ปาปเก อกุสเล ธมฺเม อธิวาเสติ นปฺปชหติ น วิโนเทติ น พฺยนฺตีกโรติ น อนภาวํ คเมติ.\csromanspacer{} เอวํ โข ภิกฺขเว ภิกฺขุ น อาสาฏิกํ หาเรตา โหติ. (3)}

\prose{กถญฺจ ภิกฺขเว ภิกฺขุ น วณํ ปฏิจฺฉาเทตา โหติ? อิธ ภิกฺขเว ภิกฺขุ จกฺขุนา รูปํ ทิสฺวา นิมิตฺตคฺคาหี โหติ อนุพฺยญฺชนคฺคาหี, ยตฺวาธิกรณเมนํ จกฺขุนฺทฺริยํ อสํวุตํ วิหรนฺตํ อภิชฺฌาโทมนสฺสา ปาปกา อกุสลา ธมฺมา อนฺวาสฺสเวยฺยุํ, ตสฺส สํวราย น ปฏิปชฺชติ, น รกฺขติ จกฺกุนฺทฺริยํ, จกฺขุนฺทฺริเย สํวรํ นาปชฺชติ.\csromanspacer{} โสเตน สทฺทํ สุตฺวา.\csromanspacer{} ฆาเนน คนฺธํ ฆายิตฺวา.\csromanspacer{} ชิวฺหาย รสํ สายิตฺวา.\csromanspacer{} กาเยน โผฏฺฐพฺพํ ผุสิตฺวา.}

\gambhira{เอกาทสกนิปาตปาฬิ}
\prose{\csromanfolio{548}มนสา ธมฺมํ วิญฺญาย นิมิตฺตคฺคาหี โหติ อนุพฺยญฺชนคฺคาหี, ยตฺวาธิกรณเมนํ มนินฺทฺริยํ อสํวุตํ วิหรนฺตํ อภิชฺฌาโทมนสฺสา ปาปกา อกุสลา ธมฺมา อนฺวาสฺสเวยฺยุํ, ตสฺส สํวราย น ปฏิปชฺชติ, น รกฺขติ มนินฺทฺริยํ, มนินฺทฺริเย สํวรํ นาปชฺชติ.\csromanspacer{} เอวํ โข ภิกฺขเว ภิกฺขุ น วณํ ปฏิจฺฉาเทตา โหติ. (4)}

\prose{กถญฺจ ภิกฺขเว ภิกฺขุ น ธูมํ กตฺตา โหติ? อิธ ภิกฺขเว ภิกฺขุ น\footnote{ม 1. 283 ปิฏฺเฐ ปน อยํ นกาโร ธมฺมนฺติปทสฺส อนนฺตรํ ทิสฺสติ.} ยถาสุตํ ยถาปริยตฺตํ ธมฺมํ วิตฺถาเรน ปเรสํ เทเสตา โหติ.\csromanspacer{} เอวํ โข ภิกฺขเว ภิกฺขุ น ธูมํ กตฺตา โหติ. (5)}

\prose{กถญฺจ ภิกฺขเว ภิกฺขุ น ติตฺถํ ชานาติ? อิธ ภิกฺขเว ภิกฺขุ เย เต ภิกฺขู พหุสฺสุตา อาคตาคมา ธมฺมธรา วินยธรา มาติกาธรา, เต กาเลน กาลํ อุปสงฺกมิตฺวา น ปริปุจฺฉาติ น ปริปญฺหติ “อิทํ ภนฺเต กถํ อิมสฺส โก อตฺโถ”ติ.\csromanspacer{} ตสฺส เต อายสฺมนฺโต อวิวฏญฺเจว น วิวรนฺติ, อนุตฺตานีกตญฺจ น อุตฺตานีกโรนฺติ, อเนกวิหิเตสุ จ กงฺขาฐานิเยสุ ธมฺเมสุ กงฺขํ น ปฏิวิโนเทนฺติ.\csromanspacer{} เอวํ โข ภิกฺขเว ภิกฺขุ น ติตฺถํ ชานาติ. (6)}

\csromanmatikaanchor{mtk.270}
\csromantocmarkref{subsection}{7. โคปาลสุตฺต}{mtk.270}
\bookmark[dest={mtk.270},level=2]{7. โคปาลสุตฺต}
\prose{กถญฺจ ภิกฺขเว ภิกฺขุ น ปีตํ ชานาติ? อิธ ภิกฺขเว ภิกฺขุ ตถาคตปฺปเวทิเต ธมฺมวินเย เทสิยมาเน น ลภติ อตฺถเวทํ, น ลภติ ธมฺมเวทํ, น ลภติ ธมฺมูปสํหิตํ ปาโมชฺชํ.\csromanspacer{} เอวํ โข ภิกฺขเว ภิกฺขุ น ปีตํ ชานาติ. (7)}

\prose{กถญฺจ ภิกฺขเว ภิกฺขุ น วีถิํ ชานาติ? อิธ ภิกฺขเว ภิกฺขุ อริยํ อฏฺฐงฺคิกํ มคฺคํ ยถาภูตํ นปฺปชานาติ.\csromanspacer{} เอวํ โข ภิกฺขเว ภิกฺขุ น วีถิํ ชานาติ. (8)}

\prose{กถญฺจ ภิกฺขเว ภิกฺขุ น โคจรกุสโล โหติ? อิธ ภิกฺขเว ภิกฺขุ จตฺตาโร สติปฏฺฐาเน ยถาภูตํ นปฺปชานาติ.\csromanspacer{} เอวํ โข ภิกฺขเว ภิกฺขุ น โคจรกุสโล โหติ. (9)}

\prose{กถญฺจ ภิกฺขเว ภิกฺขุ อนวเสสโทหี โหติ? อิธ ภิกฺขเว ภิกฺขุํ สทฺธา คหปติกา อภิหฏฺฐุํ ปวาเรนฺติ จีวรปิณฺฑปาต- เสนาสนคิลานปจฺจยเภสชฺชปริกฺขาเรหิ.\csromanspacer{} ตตฺร ภิกฺขุ มตฺตํ น ชานาติ ปฏิคฺคหณาย.\csromanspacer{} เอวํ โข ภิกฺขเว ภิกฺขุ อนวเสสโทหี โหติ. (10)}

\chapterheadpageread{2. อนุสฺสติวคฺค}
\prose{\csromanfolio{549}กถญฺจ ภิกฺขเว ภิกฺขุ เย เต ภิกฺขู เถรา รตฺตญฺญู จิรปพฺพชิตา สํฆปิตโร สํฆปริณายกา, เต น อติเรกปูชาย ปูเชตา โหติ? อิธ ภิกฺขเว ภิกฺขุ เย เต ภิกฺขู เถรา รตฺตญฺญู จิรปพฺพชิตา สํฆปิตโร สํฆปริณายกา, เตสุ น เมตฺตํ กายกมฺมํ ปจฺจุปฏฺฐาเปติ อาวิ เจว รโห จ.\csromanspacer{} น เมตฺตํ วจีกมฺมํ.\csromanspacer{} น เมตฺตํ มโนกมฺมํ ปจฺจุปฏฺฐาเปติ อาวิ เจว รโห จ.\csromanspacer{} เอวํ โข ภิกฺขเว ภิกฺขุ เย เต ภิกฺขู เถรา รตฺตญฺญู จิรปพฺพชิตา สํฆปิตโร สํฆปริณายกา, น เต อติเรกปูชาย ปูเชตา โหติ. (11)}

\prose{อิเมหิ โข ภิกฺขเว เอกาทสหิ ธมฺเมหิ สมนฺนาคโต ภิกฺขุ อภพฺโพ อิมสฺมิํ ธมฺมวินเย วุทฺธิํ วิรูฬฺหิํ เวปุลฺลํ อาปชฺชิตุํ.}

\prose{เอกกาทสหิ ภิกฺขเว องฺเคหิ สมนฺนาคโต โคปาลโก ภพฺโพ โคคณํ ปริหริตุํ ผาติํ กาตุํ.\csromanspacer{} กตเมหิ เอกาทสหิ? อิธ ภิกฺขเว โคปาลโก รูปญฺญู โหติ, ลกฺขณกุสโล โหติ, อาสาฏิกํ หาเรตา โหติ, วณํ ปฏิจฺฉาเทตา โหติ, ธูมํ กตฺตา โหติ, ติตฺถํ ชานาติ, ปีตํ ชานาติ, วีถิํ ชานาติ, โคจรกุสโล โหติ, สาวเสสโทหี จ โหติ, เย เต อุสภา โคปิตโร โคปริณายกา, เต อติเรกปูชาย ปูเชตา โหติ.\csromanspacer{} อิเมหิ โข ภิกฺขเว เอกาทสหิ องฺเคหิ สมนฺนาคโต โคปาลโก ภพฺโพ โคคณํ ปริหริตุํ ผาติํ กาตุํ.}

\prose{เอวเมวํ โข ภิกฺขเว เอกาทสหิ ธมฺเมหิ สมนฺนาคโต ภิกฺขุ ภพฺโพ อิมสฺมิํ ธมฺมวินเย วุทฺธิํ วิรูฬฺหิํ เวปุลฺลํ อาปชฺชิตุํ.\csromanspacer{} กตเมหิ เอกาทสหิ? อิธ ภิกฺขเว ภิกฺขุ รูปญฺญู โหติ, ลกฺขณกุสโล โหติ, อาสาฏิกํ หาเรตา โหติ, วณํ ปฏิจฺฉาเทตา โหติ, ธูมํ กตฺตา โหติ, ติตฺถํ ชานาติ, ปีตํ ชานาติ, วีถิํชานาติ, โคจรกุสโล โหติ, สาวเสสโทหี จ โหติ, เย เต ภิกฺขู เถรา รตฺตญฺญู จิรปพฺพชิตา สํฆปิตโร สํฆปริณายกา, เต อติเรกปูชาย ปูเชตา โหติ.}

\prose{กถญฺจ ภิกฺขเว ภิกฺขุ รูปญฺญู โหติ? อิธ ภิกฺขเว ภิกฺขุ ยํ กิญฺจิ รูปํ จตฺตาริ มหาภูตานิ จตุนฺนญฺจ มหาภูตานํ อุปทายรูปนฺติ ยถาภูตํ ปชานาติ.\csromanspacer{} เอวํ โข ภิกฺขเว ภิกฺขุ รูปญฺญู โหติ. (1)}

\gambhira{เอกาทสกนิปาตปาฬิ}
\prose{\csromanfolio{550}กถญฺจ ภิกฺขเว ภิกฺขุ ลกฺขณกุสโล โหติ? อิธ ภิกฺขเว ภิกฺขุ “กมฺมลกฺขโณ พาโล กมฺมลกฺขโณ ปณฺฑิโต”ติ ยถาภูตํ ปชานาติ.\csromanspacer{} เอวํ โข ภิกฺขเว ภิกฺขุ ลกฺขณกุสโล โหติ. (2)}

\prose{กถญฺจ ภิกฺขเว ภิกฺขุ อาสาฏิกํ หาเรตา โหติ? อิธ ภิกฺขเว ภิกฺขุ อุปฺปนฺนํ กามวิตกฺกํ นาธิวาเสติ ปชหติ วิโนเทติ พฺยนฺตีกโรติ อนภาวํ คเมติ.\csromanspacer{} อุปฺปนฺนํ พฺยาปาทวิตกฺกํ.\csromanspacer{} อุปฺปนฺนํ วิหิํสาวิตกฺกํ.\csromanspacer{} อุปฺปนฺนุปฺปนฺเน ปาปเก อกุสเล ธมฺเม นาธิวาเสติ ปชหติ วิโนเทติ พฺยนฺตีกโรติ อนภาวํ คเมติ.\csromanspacer{} เอวํ โข ภิกฺขเว ภิกฺขุ อาสาฏิกํ หาเรตา โหติ. (3)}

\prose{กถญฺจ ภิกฺขเว ภิกฺขุ วณํ ปฏิจฺฉาเทตา โหติ? อิธ ภิกฺขเว ภิกฺขุ จกฺขุนา รูปํ ทิสฺวา น นิมิตฺตคฺคาหี โหติ นานุพฺยญฺชนคฺคาหี, ยตฺวาธิกรณเมนํ จกฺขุนฺทฺริยํ อสํวุตํ วิหรนฺตํ อภิชฺฌาโทมนสฺสา ปาปกา อกุสลา ธมฺมา อนฺวาสฺสเวยฺยุํ, ตสฺส สํวราย ปฏิปชฺชติ, รกฺขติ จกฺขุนฺทฺริยํ, จกฺขุนฺทฺริเย สํวรํ อาปชฺชติ.\csromanspacer{} โสเตน สทฺทํ สุตฺวา.\csromanspacer{} ฆาเนน คนฺธํ ฆายิตฺวา.\csromanspacer{} ชิวฺหาย รสํ สายิตฺวา.\csromanspacer{} กาเยน โผฏฺฐพฺพํ ผุสิตฺวา.\csromanspacer{} มนสา ธมฺมํ วิญฺญาย น นิมิตฺตคฺคาหี โหติ นานุพฺยญฺชนคฺคาหี, ยตฺวาธิกรณเมนํ มนินฺทฺริยํ อสํวุตํ วิหรนฺตํ อภิชฺฌาโทมนสฺสา ปาปกา อกุสลา ธมฺมา อนฺวาสฺสเวยฺยุํ, ตสฺส สํวราย ปฏิปชฺชติ, รกฺขติ มนินฺทฺริยํ, มนินฺทฺริเย สํวรํ อาปชฺชติ.\csromanspacer{} เอวํ โข ภิกฺขเว ภิกฺขุ วณํ ปฏิจฺฉาเทตา โหติ. (4)}

\prose{กถญฺจ ภิกฺขเว ภิกฺขุ ธูมํ กตฺตา โหติ? อิธ ภิกฺขเว ภิกฺขุ ยถาสุตํ ยถาปริยตฺตํ ธมฺมํ วิตฺถาเรน ปเรสํ เทเสตา โหติ.\csromanspacer{} เอวํ โข ภิกฺขเว ภิกฺขุ ธูมํ กตฺตา โหติ. (5)}

\prose{กถญฺจ ภิกฺขเว ภิกฺขุ ติตฺถํ ชานติ? อิธ ภิกฺขเว ภิกฺขุ เย เต ภิกฺขู พหุสฺสุตา อาคตาคมา ธมฺมธรา วินยธรา มาติกาธรา, เต กาเลน กาลํ อุปสงฺกมิตฺวา ปริปุจฺฉติ ปริปญฺหติ “อิทํ ภนฺเต กถํ อิมสฺส โก อตฺโถ”ติ.\csromanspacer{} ตสฺส เต อายสฺมนฺโต อวิวฏญฺเจว วิวรนฺติ, อนุตฺตานีกตญฺจ อุตฺตานีกโรนฺติ, อเนกวิหิเตสุ จ กงฺขาฐานิเยสุ ธมฺเมสุ กงฺขํ ปฏิวิโนเทนฺติ.\csromanspacer{} เอวํ โข ภิกฺขเว ภิกฺขุ ติตฺถํ ชานาติ. (6)}

\prose{กถญฺจ ภิกฺขเว ภิกฺขุ ปีตํ ชานาติ? อิธ ภิกฺขเว ภิกฺขุ ตถาคตปฺปเวทิเต ธมฺมวินเย เทสิยมาเน ลภติ อตฺถเวทํ, ลภติ ธมฺมเวทํ, ลภติ ธมฺมูปสํหิตํ ปาโมชฺชํ.\csromanspacer{} เอวํ โข ภิกฺขเว ภิกฺขุ ปีตํ ชานาติ. (7)}

\chapterheadpageread{2. อนุสฺสติวคฺค}
\prose{\csromanfolio{551}กถญฺจ ภิกฺขเว ภิกฺขุ วีถิํ ชานาติ? อิธ ภิกฺขเว ภิกฺขุ อริยํ อฏฺฐงฺคิกํ มคฺคํ ยถาภูตํ ปชานาติ.\csromanspacer{} เอวํ โข ภิกฺขเว ภิกฺขุ วีถิํ ชานาติ. (8)}

\prose{กถญฺจ ภิกฺขเว ภิกฺขุ โคจรกุสโล โหติ? อิธ ภิกฺขเว ภิกฺขุ จตฺตาโร สติปฏฺฐาเน ยถาภูตํ ปชานาติ.\csromanspacer{} เอวํ โข ภิกฺขเว ภิกฺขุ โคจรกุสโล โหติ. (9)}

\prose{กถญฺจ ภิกฺขเว ภิกฺขุ สาวเสสโทหี โหติ? อิธ ภิกฺขเว ภิกฺขุ สทฺธา คหปติกา อภิหฏฺฐุํ ปวาเรนฺติ จีวรปิณฺฑปาตเสนาสนคิลานปจฺจยเภสชฺชปริกฺขาเรหิ.\csromanspacer{} ตตฺร ภิกฺขุ มตฺตํ ชานาติ ปฏิคฺคหณาย.\csromanspacer{} เอวํ โข ภิกฺขเว ภิกฺขุ สาวเสสโทหี โหติ. (10)}

\prose{กถญฺจ ภิกฺขเว ภิกฺขุ เย เต ภิกฺขู เถรา รตฺตญฺญู จิรปพฺพชิตา สํฆปิตโร สํฆปริณายกา, เต อติเรกปูชาย ปูเชตา โหติ, อิธ ภิกฺขเว ภิกฺขุ เย เต เถรา รตฺตญฺญู จิรปพฺพชิตา สํฆปิตโร สํฆปริณายกา, เตสุ เมตฺตํ กายกมฺมํ ปจฺจุปฏฺฐาเปติ อาวิ เจว รโห จ.\csromanspacer{} เมตฺตํ วจีกมฺมํ.\csromanspacer{} เมตฺตํ มโนกมฺมํ ปจฺจุปฏฺฐาเปติ อาวิ เจว รโห จ.\csromanspacer{} เอวํ โข ภิกฺขเว ภิกฺขุ เย เต ภิกฺขู เถรา รตฺตญฺญู จิรปพฺพชิตา สํฆปิตโร สํฆปริณายกา, เต อติเรกปูชาย ปูเชตา โหติ. (11)}

\prose{อิเมหิ โข ภิกฺขเว เอกาทสหิ ธมฺเมหิ สมนฺนาคโต ภิกฺขุ ภพฺโพ อิมสฺมิํ ธมฺมวินเย วุทฺธิํ วิรูฬฺหิํ เวปุลฺลํ อาปชฺชิตุนฺติ.\csromanspacer{}\csromanspacer{}สตฺตมํ.}

\csromanheader{2}{8. ปฐมสมาธิสุตฺต}
\proseitem{18}{\csromansymbolfootnote{*}{เหฏฺฐา 262 ปิฏฺเฐปิ.}อถ โข สมฺพหุลา ภิกฺขู เยน ภควา เตนุปสงฺกมิํสุ, อุปสงฺกมิตฺวา ภควนฺตํ อภิวาเทตฺวา เอกมนฺตํ นิสีทิํสุ, เอกมนฺตํ นิสินฺนา โข เต ภิกฺขู ภควนฺตํ เอตทโวจุํ\csromandash{}}

\prose{สิยา นุ โข ภนฺเต ภิกฺขุโน ตถารูโป สมาธิปฏิลาโภ, ยถา เนว ปถวิยํ ปถวิสญฺญี อสฺส, น อาปสฺมิํ อาโปสญฺญี อสฺส, น เตชสฺมิํ เตโชสญฺญี อสฺส, น วายสฺมิํ วาโยสญฺญี อสฺส, น}

\gambhira{เอกาทสกนิปาตปาฬิ}
\prose{\csromanfolio{552}อากาสานญฺจายตเน อากาสานญฺจายตนสญฺญี อสฺส, น วิญฺญาณญฺจายตเน วิญฺญาณญฺจายตนสญฺญี อสฺส, น อากิญฺจญฺญายตเน อากิญฺจญฺญายตนสญฺญี อสฺส, น เนวสญฺญานาสญฺญายตเน เนวสญฺญานาสญฺญายตนสญฺญี อสฺส, น อิธโลเก อิธโลกสญฺญี อสฺส, น ปรโลเก ปรโลกสญฺญี อสฺส, ยมฺปิทํ ทิฏฺฐํ สุตํ มุตํ วิญฺญาตํ ปตฺตํ ปริเยสิตํ อนุวิจริตํ มนสา, ตตฺราปิ น สญฺญี อสฺส, สญฺญี จ ปน อสฺสาติ.}

\prose{สิยา ภิกฺขเว ภิกฺขุโน ตถารูโป สมาธิปฏิลาโภ, ยถา เนว ปถวิยํ ปถวิสญฺญี อสฺส ฯเปฯ ยมฺปิทํ ทิฏฺฐํ สุตํ มุตํ วิญฺญาตํ ปตฺตํ ปริเยสิตํ อนุวิจริตํ มนสา, ตตฺราปิ น สญฺญี อสฺส, สญฺญี จ ปน อสฺสาติ.}

\prose{ยถา กถํ ปน ภนฺเต สิยา ภิกฺขุโน ตถารูโป สมาธิปฏิลาโภ, ยถา เนว ปถวิยํ ปถวิสญฺญี อสฺส ฯเปฯ ยมฺปิทํ ทิฏฺฐํ สุตํ มุตํ วิญฺญาตํ ปตฺตํ ปริเยสิตํ อนุวิจริตํ มนสา, ตตฺราปิ น สญฺญี อสฺส, สญฺญี จ ปน อสฺสาติ.}

\prose{อิธ ภิกฺขเว ภิกฺขุ เอวํสญฺญี โหติ “เอตํ สนฺตํ เอตํ ปณีตํ, ยทิทํ สพฺพสงฺขารสมโถ สพฺพูปธิปฏินิสฺสคฺโค ตณฺหากฺขโย วิราโค นิโรโธ นิพฺพานนฺ”ติ.\csromanspacer{} เอวํ โข ภิกฺขเว สิยา ภิกฺขุโน ตถารูโป สมาธิปฏิลาโภ, ยถา เนว ปถวิยํ ปถวิสญฺญี อสฺส, น อาปสฺมิํ อาโปสญฺญี อสฺส, น เตชสฺมิํ เตโชสญฺญี อสฺส, น วายสฺมิํ วาโยสญฺญี อสฺส, น อากาสานญฺจายตเน อากาสานญฺจายตนสญฺญี อสฺส, น วิญฺญาณญฺจายตเน วิญฺญาณญฺจายตนสญฺญี อสฺส, น อากิญฺจญฺญายตเน อากิญฺจญฺญายตนสญฺญี อสฺส, น เนวสญฺญานาสญฺญายตเน เนวสญฺญานาสญฺญายตนสญฺญี อสฺส, น อิธโลเก อิธโลกสญฺญี อสฺส, น ปรโลเก ปรโลกสญฺญี อสฺส, ยมฺปิทํ ทิฏฺฐํ สุตํ มุตํ วิญฺญาตํ ปตฺตํ ปริเยสิตํ อนุวิจริตํ มนสา, ตตฺราปิ น สญฺญี อสฺส, สญฺญี จ ปน อสฺสาติ.\csromanspacer{}\csromanspacer{}อฏฺฐมํ.}

\csromanheader{2}{9. ทุติยสมาธิสุตฺต}
\proseitem{19}{ตตฺร โข ภควา ภิกฺขู อามนฺเตสิ “ภิกฺขโว”ติ. “ภทนฺเต”ติ เต ภิกฺขู ภควโต ปจฺจสฺโสสุํ.\csromanspacer{} ภควา เอตทโวจ\csromandash{}}

\chapterheadpageread{2. อนุสฺสติวคฺค}
\prose{\csromanfolio{553}สิยา นุ โข ภิกฺขเว ภิกฺขุโน ตถารูโป สมาธิปฏิลาโภ, ยถา เนว ปถวิยํ ปถวิสญฺญี อสฺส, น อาปสฺมิํ อาโปสญฺญี อสฺส ฯเปฯ น อากิญฺจญฺญายตเน อากิญฺจญฺญายตนสญฺญี อสฺส, น เนวสญฺญานาสญฺญายตเน เนวสญฺญานาสญฺญายตนสญฺญี อสฺส, น อิธโลเก อิธโลกสญฺญี อสฺส, น ปรโลเก ปรโลกสญฺญี อสฺส, ยมฺปิทํ ทิฏฺฐํ สุตํ มุตํ วิญฺญาตํ ปตฺตํ ปริเยสิตํ อนุวิจริตํ มนสา, ตตฺราปิ น สญฺญี อสฺส, สญฺญี จ ปน อสฺสาติ.\csromanspacer{} ภควํมูลกา โน ภนฺเต ธมฺมา ภควํเนตฺติกา ภควํปฏิสรณา.\csromanspacer{} สาธุ วต ภนฺเต ภควนฺตํเยว ปฏิภาตุ เอตสฺส ภาสิตสฺส อตฺโถ, ภควโต สุตฺวา ภิกฺขู ธาเรสฺสนฺตีติ.}

\csromanmatikaanchor{mtk.271}
\csromantocmarkref{subsection}{8. ปฐมสมาธิสุตฺต}{mtk.271}
\bookmark[dest={mtk.271},level=2]{8. ปฐมสมาธิสุตฺต}
\prose{เตน หิ ภิกฺขเว สุณาถ สาธุกํ มนสิ กโรถ, ภาสิสฺสามีติ. “เอวํ ภนฺเต”ติ โข เต ภิกฺขู ภควโต ปจฺจสฺโสสุํ.\csromanspacer{} ภควา เอตทโวจ\csromandash{}}

\prose{สิยา ภิกฺขเว ภิกฺขุโน ตถารูโป สมาธิปฏิลาโภ, ยถา เนว ปถวิยํ ปถวิสญฺญี อสฺส ฯเปฯ ยมฺปิทํ ทิฏฺฐํ สุตํ มุตํ วิญฺญาตํ ปตฺตํ ปริเยสิตํ อนุวิจริตํ มนสา, ตตฺราปิ น สญฺญี อสฺส, สญฺญี จ ปน อสฺสาติ.}

\prose{ยถา กถํ ปน ภนฺเต สิยา ภิกฺขุโน ตถารูโป สมาธิปฏิลาโภ, ยถา เนว ปถวิยํ ปถวิสญฺญี อสฺส ฯเปฯ ยมฺปิทํ ทิฏฺฐํ สุตํ มุตํ วิญฺญาตํ ปตฺตํ ปริเยสิตํ อนุวิจริตํ มนสา, ตตฺราปิ น สญฺญี อสฺส, สญฺญี จ ปน อสฺสาติ.}

\prose{อิธ ภิกฺขเว ภิกฺขุ เอวํสญฺญี โหติ “เอตํ สนฺตํ เอตํ ปณีตํ, ยทิทํ สพฺพสงฺขารสมโถ สพฺพูปธิปฏินิสฺสคฺโค ตณฺหากฺขโย วิราโค นิโรโธ นิพฺพานนฺ”ติ.\csromanspacer{} เอวํ โข ภิกฺขเว สิยา ภิกฺขุโน ตถารูโป สมาธิปฏิลาโภ, ยถา เนว ปถวิยํ ปถวิสญฺญี อสฺส ฯเปฯ ยมฺปิทํ ทิฏฺฐํ สุตํ มุตํ วิญฺญาตํ ปตฺตํ ปริเยสิตํ อนุวิจริตํ มนสา, ตตฺราปิ น สญฺญี อสฺส, สญฺญี จ ปน อสฺสาติ.\csromanspacer{}\csromanspacer{}นวมํ.}

\csromanheader{2}{10. ตติยสมาธิสุตฺต}
\proseitem{20}{\csromansymbolfootnote{*}{เหฏฺฐา 263 ปิฏฺเฐปิ.}อถ โข สมฺพหุลา ภิกฺขู เยนายสฺมา สาริปุตฺโต เตนุปสงฺกมิํสุ, อุปสงฺกมิตฺวา อายสฺมตา สาริปุตฺเตน สทฺธิํ สมฺโมทิํสุ,}

\gambhira{เอกาทสกนิปาตปาฬิ}
\prose{\csromanfolio{554}สมฺโมทนียํ กถํ สารณียํ วีติสาเรตฺวา เอกมนฺตํ นิสีทิํสุ, เอกมนฺตํ นิสินฺนา โข เต ภิกฺขู อายสฺมนฺตํ สาริปุตฺตํ เอตทโวจุํ\csromandash{}}

\csromanmatikaanchor{mtk.272}
\csromantocmarkref{subsection}{9. ทุติยสมาธิสุตฺต}{mtk.272}
\bookmark[dest={mtk.272},level=2]{9. ทุติยสมาธิสุตฺต}
\prose{สิยา นุ โข อาวุโส สาริปุตฺต ภิกฺขุโน ตถารูโป สมาธิปฏิลาโภ, ยถา เนว ปถวิยํ ปถวิสญฺญี อสฺส ฯเปฯ ยมฺปิทํ ทิฏฺฐํ สุตํ มุตํ วิญฺญาตํ ปตฺตํ ปริเยสิตํ อนุวิจริตํ มนสา, ตตฺราปิ น สญฺญี อสฺส, สญฺญี จ ปน อสฺสาติ.\csromanspacer{} สิยา อาวุโส ภิกฺขุโน ตถารูโป สมาธิปฏิลาโภ, ยถา เนว ปถวิยํ ปถวิสญฺญี อสฺส ฯเปฯ ยมฺปิทํ ทิฏฺฐํ สุตํ มุตํ วิญฺญาตํ ปตฺตํ ปริเยสิตํ อนุวิจริตํ มนสา, ตตฺราปิ น สญฺญี อสฺส, สญฺญี จ ปน อสฺสาติ.}

\prose{ยถา กถํ ปน อาวุโส สาริปุตฺต สิยา ภิกฺขุโน ตถารูโป สมาธิปฏิลาโภ, ยถา เนว ปถวิยํ ปถวิสญฺญี อสฺส ฯเปฯ ยมฺปิทํ ทิฏฺฐํ สุตํ มุตํ วิญฺญาตํ ปตฺตํ ปริเยสิตํ อนุวิจริตํ มนสา, ตตฺราปิ น สญฺญี อสฺส, สญฺญี จ ปน อสฺสาติ.}

\prose{อิธ อาวุโส ภิกฺขุ เอวํสญฺญี โหติ “เอตํ สนฺตํ เอตํ ปณีตํ, ยทิทํ สพฺพสงฺขารสมโถ สพฺพูปธิปฏินิสฺสคฺโค ตณฺหากฺขโย วิราโค นิโรโธ นิพฺพานนฺ”ติ.\csromanspacer{} เอวํ โข อาวุโส สิยา ภิกฺขุโน ตถารูโป สมาธิปฏิลาโภ, ยถา เนว ปถวิยํ ปถวิสญฺญี อสฺส ฯเปฯ ยมฺปิทํ ทิฏฺฐํ สุตํ มุตํ วิญฺญาตํ ปตฺตํ ปริเยสิตํ อนุวิจริตํ มนสา, ตตฺราปิ น สญฺญี อสฺส, สญฺญี จ ปน อสฺสาติ.\csromanspacer{}\csromanspacer{}ทสมํ.}

\csromanheader{2}{11. จตุตฺถสมาธิสุตฺต}
\proseitem{21}{ตตฺร โข อายสฺมา สาริปุตฺโต ภิกฺขู อามนฺเตสิ\csromandash{}สิยา นุ โข อาวุโส ภิกฺขุโน ตถารูโป สมาธิปฏิลาโภ, ยถา เนว ปถวิยํ ปถวิสญฺญี อสฺส, น อาปสฺมิํ อาโปสญฺญี อสฺส, น เตชสฺมิํ เตโชสญฺญี อสฺส, น วายสฺมิํ วาโยสญฺญี อสฺส, น อากาสานญฺจายตเน อากาสานญฺจายตนสญฺญี อสฺส, น วิญฺญาณญฺจายตเน วิญฺญาณญฺจายตนสญฺญี อสฺส, น อากิญฺจญฺญายตเน อากิญฺจญฺญายตนสญฺญี อสฺส, น เนวสญฺญานาสญฺญายตเน เนวสญฺญานาสญฺญา- ยตนสญฺญี อสฺส, น อิธโลเก อิธโลกสญฺญี อสฺส, น ปรโลเก ปรโลกสญฺญี อสฺส, ยมฺปิทํ ทิฏฺฐํ สุตํ มุตํ วิญฺญาตํ ปตฺตํ ปริเยสิตํ อนุวิจริตํ มนสา, ตตฺราปิ น สญฺญี อสฺส, สญฺญี จ ปน อสฺสาติ.}

\chapterheadpageread{2. อนุสฺสติวคฺค}
\prose{\csromanfolio{555}ทูรโตปิ โข มยํ อาวุโส อาคจฺเฉยฺยาม อายสฺมโต สาริปุตฺตสฺส สนฺติเก เอตสฺส ภาสิตสฺส อตฺถมญฺญาตุํ, สาธุ วตายสฺมนฺตํเยว สาริปุตฺตํ ปฏิภาตุ เอตสฺส ภาสิตสฺส อตฺโถ, อายสฺมโต สาริปุตฺตสฺส สุตฺวา ภิกฺขู ธาเรสฺสนฺตีติ.}

\prose{เตนหาวุโส สุณาถ สาธุกํ มนสิ กโรถ, ภาสิสฺสามีติ. “เอวมาวุโส”ติ โข เต ภิกฺขู อายสฺมโต สาริปุตฺตสฺส ปจฺจสฺโสสุํ.\csromanspacer{} อายสฺมา สาริปุตฺโต เอตทโวจ\csromandash{}}

\csromanmatikaanchor{mtk.273}
\csromantocmarkref{subsection}{10. ตติยสมาธิสุตฺต}{mtk.273}
\bookmark[dest={mtk.273},level=2]{10. ตติยสมาธิสุตฺต}
\prose{สิยา อาวุโส ภิกฺขุโน ตถารูโป สมาธิปฏิลาโภ, ยถา เนว ปถวิยํ ปถวิสญฺญี อสฺส ฯเปฯ ยมฺปิทํ ทิฏฺฐํ สุตํ มุตํ วิญฺญาตํ ปตฺตํ ปริเยสิตํ อนุวิจริตํ มนสา, ตตฺราปิ น สญฺญี อสฺส, สญฺญี จ ปน อสฺสาติ.}

\prose{ยถา กถํ ปนาวุโส สิยา ภิกฺขุโน ตถารูโป สมาธิปฏิลาโภ, ยถา เนว ปถวิยํ ปถวิสญฺญี อสฺส ฯเปฯ ยมฺปิทํ ทิฏฺฐํ สุตํ มุตํ วิญฺญาตํ ปตฺตํ ปริเยสิตํ อนุวิจริตํ มนสา, ตตฺราปิ น สญฺญี อสฺส, สญฺญี จ ปน อสฺสาติ.}

\prose{อิธ อาวุโส ภิกฺขุ เอวํสญฺญี โหติ “เอตํ สนฺตํ เอตํ ปณีตํ, ยทิทํ สพฺพสงฺขารสมโถ สพฺพูปธิปฏินิสฺสคฺโค ตณฺหากฺขโย วิราโค นิโรโธ นิพฺพานนฺ”ติ.\csromanspacer{} เอวํ โข อาวุโส สิยา ภิกฺขุโน ตถารูโป สมาธิปฏิลาโภ, ยถา เนว ปถวิยํ ปถวิสญฺญี อสฺส, น อาปสฺมิํ อาโปสญฺญี อสฺส, น เตชสฺมิํ เตโชสญฺญี อสฺส, น วายสฺมิํ วาโยสญฺญี อสฺส, น อากาสานญฺจายตเน อากาสานญฺจายตนสญฺญี อสฺส, น วิญฺญาณญฺจายตเน วิญฺญาณญฺจายตนสญฺญี อสฺส, น อากิญฺจญฺญายตเน อากิญฺจญฺญายตนสญฺญี อสฺส, น เนวสญฺญานาสญฺญายตเน เนวสญฺญานาสญฺญายตนสญฺญี อสฺส, น อิธโลเก อิธโลกสญฺญี อสฺส, น ปรโลเก ปรโลกสญฺญี อสฺส, ยมฺปิทํ ทิฏฺฐํ สุตํ มุตํ วิญฺญาตํ ปตฺตํ ปริเยสิตํ อนุวิจริตํ มนสา, ตตฺราปิ น สญฺญี อสฺส, สญฺญี จ ปน อสฺสาติ.\csromanspacer{}\csromanspacer{}เอกาทสมํ.}

\vspace{\tipitakaparskip}
\csromancenter{อนุสฺสติวคฺโค ทุติโย.}
\gambhira{เอกาทสกนิปาตปาฬิ}
\titlehead{\textbf{ตสฺสุทฺทานํ}}
\csromanmatikaanchor{mtk.274}
\csromanmatikaanchor{mtk.275}
\csromantocmarkref{subsection}{11. จตุตฺถสมาธิสุตฺต}{mtk.274}
\bookmark[dest={mtk.274},level=2]{11. จตุตฺถสมาธิสุตฺต}
\csromantocmarkref{section}{อุทฺทานคาถา}{mtk.275}
\bookmark[dest={mtk.275},level=1]{อุทฺทานคาถา}
\csromangathameasure{เทฺว วุตฺตา มหานาเมน, นนฺทิเยน สุภูตินา. \\ เมตฺตา อฏฺฐโก โคปาโล, จตฺตาโร จ สมาธินาติ.}
\csromangathagroup{\csromangathastanza{\csromanfolio{556}เทฺว วุตฺตา มหานาเมน, นนฺทิเยน สุภูตินา. \\ เมตฺตา อฏฺฐโก โคปาโล, จตฺตาโร จ สมาธินาติ.}}

\chapterhead{3. สามญฺญวคฺค}
\proseitem{22-29}{เอกาทสหิ ภิกฺขเว องฺเคหิ สมนฺนาคโต โคปาลโก อภพฺโพ โคคณํ ปริหริตุํ ผาติํ กาตุํ.\csromanspacer{} กตเมหิ เอกาทสหิ? อิธ ภิกฺขเว โคปาลโก น รูปญฺญู โหติ, น ลกฺขณกุสโล โหติ, น อาสาฏิกํ หาเรตา โหติ, น วณํ ปฏิจฺฉาเทตา โหติ, น ธูมํ กตฺตา โหติ, น ติตฺถํ ชานาติ, น ปีตํ ชานาติ, น วีถิํ ชานาติ, น โคจรกุสโล โหติ, อนวเสสโทหี จ โหติ, เย เต อุสภา โคปิตโร โคปริณายกา, เต น อติเรกปูชาย ปูเชตา โหติ.\csromanspacer{} อิเมหิ โข ภิกฺขเว เอกาทสหิ องฺเคหิ สมนฺนาคโต โคปาลโก อภพฺโพ โคคณํ ปริหริตุํ ผาติํ กาตุํ.}

\prose{เอวเมวํ โข ภิกฺขเว เอกาทสหิ ธมฺเมหิ สมนฺนาคโต ภิกฺขุ อภพฺโพ จกฺขุสฺมิํ อนิจฺจานุปสฺสี วิหริตุํ ฯเปฯ อภพฺโพ จกฺขุสฺมิํ ทุกฺขานุปสฺสี วิหริตุํ.\csromanspacer{} อภพฺโพ จกฺขุสฺมิํ อนตฺตานุปสฺสี วิหริตุํ.\csromanspacer{} อภพฺโพ จกฺขุสฺมิํ ขยานุปสฺสี วิหริตุํ.\csromanspacer{} อภพฺโพ จกฺขุสฺมิํ วยานุปสฺสี วิหริตุํ.\csromanspacer{} อภพฺโพ จกฺขุสฺมิํ วิราคานุปสฺสี วิหริตุํ.\csromanspacer{} อภพฺโพ จกฺขุสฺมิํ นิโรธานุปสฺสี วิหริตุํ.\csromanspacer{} อภพฺโพ จกฺขุสฺมิํ ปฏินิสฺสคฺคานุปสฺสี วิหริตุํ. (1-8)}

\csromanheader{2}{30-69. โสตสฺมิํ. ฆานสฺมิํ. ชิวฺหาย. กายสฺมิํ. มนสฺมิํ. (9-48)}
\proseitem{70-117}{รูเปสุ.\csromanspacer{} สทฺเทสุ.\csromanspacer{} คนฺเธสุ.\csromanspacer{} รเสสุ.\csromanspacer{} โผฏฺฐพฺเพสุ.\csromanspacer{} ธมฺเมสุ. (49-96)}

\proseitem{118-165}{จกฺขุวิญฺญาเณ.\csromanspacer{} โสตวิญฺญาเณ.\csromanspacer{} ฆานวิญฺญาเณ.\csromanspacer{} ชิวฺหาวิญฺญาเณ.\csromanspacer{} กายวิญฺญาเณ.\csromanspacer{} มโนวิญฺญาเณ. (97-144)}

\proseitem{116-213}{จกฺขุสมฺผสฺเส.\csromanspacer{} โสตสมฺผสฺเส.\csromanspacer{} ฆานสมฺผสฺเส.\csromanspacer{} ชิวฺหาสมฺผสฺเส.\csromanspacer{} กายสมฺผสฺเส.\csromanspacer{} มโนสมฺผสฺเส. (145-192)}

\chapterheadpageread{4. ราคเปยฺยาล}
\proseitem{214-261}{\csromanfolio{557}จกฺขุสมฺผสฺสชาย เวทนาย.\csromanspacer{} โสตสมฺผสฺสชาย เวทนาย.\csromanspacer{} ฆานสมฺผสฺสชาย เวทนาย.\csromanspacer{} ชิวฺหาสมฺผสฺสชาย เวทนาย.\csromanspacer{} กายสมฺผสฺสชาย เวทนาย.\csromanspacer{} มโนสมฺผสฺสชาย เวทนาย. (193-240)}

\proseitem{262-309}{รูปสญฺญาย.\csromanspacer{} สทฺทสญฺญาย.\csromanspacer{} คนฺธสญฺญาย.\csromanspacer{} รสสญฺญาย.\csromanspacer{} โผฏฺฐพฺพสญฺญาย.\csromanspacer{} ธมฺมสญฺญาย. (241-288)}

\proseitem{310-357}{รูปสญฺเจตนาย.\csromanspacer{} สทฺทสญฺเจตนาย.\csromanspacer{} คนฺธสญฺเจตนาย.\csromanspacer{} รสสญฺเจตนาย.\csromanspacer{} โผฏฺฐพฺพสญฺเจตนาย.\csromanspacer{} ธมฺมสญฺเจตนาย. (289-336)}

\proseitem{358-405}{รูปตณฺหาย.\csromanspacer{} สทฺทตณฺหาย.\csromanspacer{} คนฺธตณฺหาย.\csromanspacer{} รสตณฺหาย.\csromanspacer{} โผฏฺฐพฺพตณฺหาย.\csromanspacer{} ธมฺมตณฺหาย. (337-384)}

\csromanmatikaanchor{mtk.276}
\csromantocmarkref{chapter}{3. สามญฺญวคฺค}{mtk.276}
\bookmark[dest={mtk.276},level=0]{3. สามญฺญวคฺค}
\proseitem{406-453}{รูปวิตกฺเก.\csromanspacer{} สทฺทวิตกฺเก.\csromanspacer{} คนฺธวิตกฺเก.\csromanspacer{} รสวิตกฺเก.\csromanspacer{} โผฏฺฐพฺพวิตกฺเก.\csromanspacer{} ธมฺมวิตกฺเก. (385-432)}

\proseitem{454-501}{รูปวิจาเร.\csromanspacer{} สทฺทวิจาเร.\csromanspacer{} คนฺธวิจาเร.\csromanspacer{} รสวิจาเร.\csromanspacer{} โผฏฺฐพฺพวิจาเร.\csromanspacer{} ธมฺมวิจาเร อนิจฺจานุปสฺสี วิหริตุํ.\csromanspacer{} ทุกฺขานุปสฺสี วิหริตุํ.\csromanspacer{} อนตฺตานุปสฺสี วิหริตุํ.\csromanspacer{} ขยานุปสฺสี วิหริตุํ.\csromanspacer{} วยานุปสฺสี วิหริตุํ.\csromanspacer{} วิราคามุปสฺสี วิหริตุํ.\csromanspacer{} นิโรธานุปสฺสี วิหริตุํ.\csromanspacer{} ปฏินิสฺสคฺคานุปสฺสี วิหริตุํ ฯเปฯ. (433-480)}
\csromansectionrule

\csromanheader{2}{4. ราคเปยฺยาล}
\proseitem{502}{ราคสฺส ภิกฺขเว อภิญฺญาย เอกาทส ธมฺมา ภาเวตพฺพา.\csromanspacer{} กตเม เอกาทส? ปฐมํ ฌานํ ทุติยํ ฌานํ ตติยํ ฌานํ จตุตฺถํ ฌานํ เมตฺตาเจโตวิมุตฺติ กรุณาเจโตวิมุตฺติ มุทิตาเจโตวิมุตฺติ อุเปกฺขาเจโตวิมุตฺติ อากาสานญฺจายตนํ วิญฺญาณญฺจายตนํ อากิญฺจญฺญายตนํ.\csromanspacer{} ราคสฺส ภิกฺขเว อภิญฺญาย อิเม เอกาทส ธมฺมา ภาเวตพฺพา. (1)}

\proseitem{503-511}{ราคสฺส ภิกฺขเว ปริญฺญาย.\csromanspacer{} ปริกฺขยาย.\csromanspacer{} ปหานาย.\csromanspacer{} ขยาย.\csromanspacer{} วยาย.\csromanspacer{} วิราคาย.\csromanspacer{} นิโรธาย.\csromanspacer{} จาคาย.\csromanspacer{} ปฏินิสฺสคฺคาย อิเม เอกาทส ธมฺมา ภาเวตพฺพา. (2-10)}

\gambhira{เอกาทสกนิปาตปาฬิ}
\proseitem{512-671}{\csromanfolio{558}โทสสฺส ฯเปฯ.\csromanspacer{} โมหสฺส.\csromanspacer{} โกธสฺส.\csromanspacer{} อุปนาหสฺส.\csromanspacer{} มกฺขสฺส.\csromanspacer{} ปฬาสสฺส.\csromanspacer{} อิสฺสาย.\csromanspacer{} มจฺฉริยสฺส.\csromanspacer{} มายาย.\csromanspacer{} สาเฐยฺยสฺส.\csromanspacer{} ถมฺภสฺส.\csromanspacer{} สารมฺภสฺส.\csromanspacer{} มานสฺส.\csromanspacer{} อติมานสฺส.\csromanspacer{} มทสฺส.\csromanspacer{} ปมาทสฺส อภิญฺญาย ฯเปฯ ปริญฺญาย.\csromanspacer{} ปริกฺขยาย.\csromanspacer{} ปหานาย.\csromanspacer{} ขยาย.\csromanspacer{} วยาย.\csromanspacer{} วิราคาย.\csromanspacer{} นิโรธาย.\csromanspacer{} จาคาย.\csromanspacer{} ปฏินิสฺสคฺคาย ฯเปฯ.\csromanspacer{} อิเม เอกาทส ธมฺมา ภาเวตพฺพาติ. (11. 170)}

\csromanmatikaanchor{mtk.277}
\csromantocmarkref{chapter}{4. ราคเปยฺยาล}{mtk.277}
\bookmark[dest={mtk.277},level=0]{4. ราคเปยฺยาล}
\prose{อิทมโวจ ภควา, อตฺตมนา เต ภิกฺขู ภควโต ภาสิตํ อภินนฺทุนฺติ.}

\nitthitam{\textbf{ราคเปยฺยาลํ นิฏฺฐิตํ.}}
\csromangathameasure{นว สุตฺตสหสฺสานิ, ภิยฺโย ปญฺจสตานิ จ. \\ สตฺตปญฺญาส สุตฺตนฺตา, องฺคุตฺตรสมายุตา3ติ.}
\csromangathagroup{\csromangathastanza{นว สุตฺตสหสฺสานิ, ภิยฺโย ปญฺจสตานิ จ\footnote{ปญฺจ สุตฺตสตานิ จ (ฏฺฐ)}. \\ สตฺตปญฺญาส สุตฺตนฺตา\footnote{ปญฺจ สุตฺตสตานิ จ (ฏฺฐ)}, องฺคุตฺตรสมายุตา3ติ.}}

\nitthitam{\textbf{เอกาทสกนิปาตปาฬิ นิฏฺฐิตา.}}
\vspace{\tipitakaparskip}
\csromancenter{\textbf{องฺคุตฺตรนิกาโย สมตฺโต.}}
\orphannote{วิ 1. 1 ปิฏฺฐาทีสุปิ.}

\orphannote{อํ 1. 341; ที 3. 193 ปิฏฺเฐสุปิ.}

\orphannote{อุปริ 397; ขุ 9. 355 ปิฏฺเฐปิ.}

\orphannote{ขุ 8. 265 ปิฏฺเฐปิ.}

\orphannote{เหฏฺฐา 107, 111 ปิฏฺเฐสุปิ.}

\orphannote{เอตฺถ “อถ โข”ติ จ, “อุปาสิกา”ติ จ อิทํ อฏฺฐกถายเมว ทิสฺสติ, น ปาฬิโปตฺถเกสุ.}

\orphannote{เหฏฺฐา 133, อุปริ 264 ปิฏฺเฐสุปิ.}

\orphannote{อจิรีติ (สฺยา), เอวํ จริติ, จริสฺสติ” ปเทสุปิ.}

\orphannote{อภิ 2. 364 ปิฏฺเฐปิ.}

\orphannote{ที 1. 7; ม 2. 181; สํ 3. 368; วิ 2. 213 ปิฏฺเฐสุ.}

\orphannote{อภิ 4. 115 ปิฏฺเฐปิ.}

\orphannote{เหฏฺฐา 58; ขุ 9. 355 ปิฏฺเฐสุปิ.}

\orphannote{โส อชานํ วา อหํ ชานามีติ, ชานํ วา อหํ น ชานามีติ, อปสฺสํ วา อหํ ปสฺสามีติ, ปสฺสํ วา อหํ น ปสฺสามีติ (อิ, ก) เอวมุปริปิ.}

\orphannote{เหฏฺฐา 257 ปิฏฺเฐปิ.}

\orphannote{ที 1. 92; สํ 1. 155; 471 ปิฏฺเฐสุปิ.}

\orphannote{อํ 2. 252 ปิฏฺเฐปิ.}

\orphannote{โหนฺติ องฺคุตฺตราคเม (ฏฺฐ)}

